\documentclass[11pt,letterpaper,reqno]{amsart}

%% Solutions to I. Karatzas and S. E. Shreve, Brownian Motion and Stochastic Calculus
%% (Springer, second edition, 1991).  This document contains the solutions only: the
%% statements are not reproduced, and each item is identified by its number in the book.
%% Self-contained: the preamble below is kssol.sty inlined verbatim, so that the
%% document compiles with pdflatex alone.

\makeatletter
\usepackage{amsmath}
\usepackage{amssymb}
\usepackage{amsthm}
\usepackage{mathtools}
\usepackage{mathrsfs}
\usepackage[margin=1in]{geometry}
\usepackage{enumitem}
\usepackage{booktabs}
\usepackage{setspace}
\usepackage[expansion=false]{microtype}  % expansion off: it can make pdfTeX misplace text
\usepackage{xcolor}
\usepackage{graphicx}
\usepackage{fancyvrb}
\usepackage{bm}

\definecolor{NavyBlue}{RGB}{0,51,153}
\definecolor{RuleGrey}{RGB}{120,120,120}

\usepackage[colorlinks=true,linkcolor=NavyBlue,urlcolor=NavyBlue,
                citecolor=NavyBlue,linktoc=all]{hyperref}

%% ---------------------------------------------------------------------------
%% Layout
%% ---------------------------------------------------------------------------
\onehalfspacing
\allowdisplaybreaks[2]
\setlength{\parindent}{0pt}
\setlength{\parskip}{0.6ex plus 0.2ex minus 0.1ex}
\setlength{\emergencystretch}{2em}
\setlist{nosep,topsep=0.5ex}
\raggedbottom

%% ---------------------------------------------------------------------------
%% Notation.  The book writes E and P in plain italic (E^x, P^x, P^\mu), script
%% letters for sigma-fields and spaces of processes, and <M>, <M,N> for the
%% quadratic variation and the cross-variation; the macros follow it.
%% ---------------------------------------------------------------------------
\newcommand{\R}{\mathbb{R}}
\newcommand{\N}{\mathbb{N}}
\newcommand{\Z}{\mathbb{Z}}
\newcommand{\Q}{\mathbb{Q}}
\newcommand{\F}{\mathscr{F}}                 % \F_t, \F_\infty, \F_T, \F^X_t, \F_{t+}
\newcommand{\G}{\mathscr{G}}
\newcommand{\B}{\mathscr{B}}                 % Borel sigma-field: \B(\R^d)
\newcommand{\cH}{\mathscr{H}}
\newcommand{\cM}{\mathscr{M}}                % \cM_2, \cM_2^c, \cM^{c,\mathrm{loc}}
\newcommand{\cL}{\mathscr{L}}                % \cL, \cL^*, \cL_0
\newcommand{\cP}{\mathscr{P}}
\newcommand{\cN}{\mathscr{N}}                % the null sets
\newcommand{\cS}{\mathscr{S}}
\newcommand{\cA}{\mathscr{A}}
\newcommand{\cC}{\mathscr{C}}
\newcommand{\cD}{\mathscr{D}}
\newcommand{\loc}{\mathrm{loc}}
\newcommand{\ind}{\mathbf{1}}                % indicator: \ind_A, \ind_{\{T\le t\}}
\newcommand{\qv}[1]{\langle #1\rangle}       % \qv{M}_t, \qv{M,N}_t
\newcommand{\eps}{\varepsilon}
\newcommand{\defeq}{\triangleq}
\newcommand{\wh}[1]{\widehat{#1}}
\newcommand{\wt}[1]{\widetilde{#1}}
\newcommand{\ol}[1]{\overline{#1}}
\newcommand{\ul}[1]{\underline{#1}}
\newcommand{\bs}{\boldsymbol}
\newcommand{\Var}{\mathrm{Var}}
\newcommand{\Cov}{\mathrm{Cov}}
\newcommand{\as}{\text{a.s.}}
\DeclarePairedDelimiter{\abs}{\lvert}{\rvert}
\DeclarePairedDelimiter{\norm}{\lVert}{\rVert}
\DeclareMathOperator{\sgn}{sgn}
\DeclareMathOperator{\supp}{supp}
\DeclareMathOperator{\tr}{tr}
\DeclareMathOperator{\meas}{meas}
\DeclareMathOperator{\erf}{erf}
\DeclareMathOperator{\erfc}{erfc}
\DeclareMathOperator*{\esssup}{ess\,sup}
\DeclareMathOperator*{\essinf}{ess\,inf}
\DeclareMathOperator*{\argmin}{arg\,min}
\DeclareMathOperator*{\argmax}{arg\,max}
\DeclareMathOperator*{\plim}{\mathnormal{P}\text{-}lim}

%% Unnumbered statements inside a solution.
\theoremstyle{plain}
\newtheorem*{lem}{Lemma}
\newtheorem*{claim}{Claim}
\theoremstyle{remark}
\newtheorem*{rem}{Remark}

%% ---------------------------------------------------------------------------
%% Problem / Exercise environment
%%
%%   \begin{kex}{Problem}{2.2}{Short title}  ... \solution ... \end{kex}
%%
%% The first argument is the book's own word for the item (Problem or Exercise),
%% the second its number within the chapter.  The statement is set in italics,
%% the solution upright.  The short title may be empty.
%% ---------------------------------------------------------------------------
\newcommand{\ks@head}[3]{%
  \par\addvspace{1.6ex}\penalty-200\relax
  \phantomsection
  % The bookmark carries the kind and the number only: the short titles contain
  % mathematics, which hyperref cannot turn into a PDF string.
  \addcontentsline{toc}{subsection}{#1 #2\texorpdfstring{\ifx\relax#3\relax\else\quad #3\fi}{}}%
  {\color{RuleGrey}\hrule height 0.4pt}%
  \nobreak\vspace{0.8ex}\nobreak
  \noindent{\bfseries #1 #2}%
  \ifx\relax#3\relax\else\quad{\bfseries\itshape #3}\fi
  \par\nobreak\vspace{0.4ex}\nobreak
  \ks@beginstatement
}

% The statement is set in italics, in a group that \solution closes.  The public
% document carries no statements (tools/build_main.py strips them), and redefines
% this to \ignorespaces so that no group is opened and none is left open.
\newcommand{\ks@beginstatement}{\begingroup\itshape\ignorespaces}

\newenvironment{kex}[3]{\ks@head{#1}{#2}{#3}}{\par\addvspace{1.2ex}}

\newcommand{\solution}{%
  \endgroup
  \par\vspace{0.6ex}%
  \noindent\textbf{Solution.}\hspace{0.6em}\ignorespaces
}

% A short aside: a caveat, a check, or a note on the text of the book.
\newcommand{\rmk}[1]{%
  \par\vspace{0.6ex}%
  {\small\noindent\textit{Remark.}\hspace{0.5em}#1\par}%
  \vspace{0.4ex}%
}

% Reference to a numbered result, equation or section of the book (upright small
% capitals, so that it reads correctly inside the italic statements).
\newcommand{\bk}[1]{\textup{\textsc{#1}}}

\DefineVerbatimEnvironment{codeblock}{Verbatim}%
  {fontsize=\small,xleftmargin=1.2em,samepage=false}

%% ---------------------------------------------------------------------------
%% Chapter and section headings
%% ---------------------------------------------------------------------------
\newcounter{solchap}
\renewcommand{\theHequation}{\arabic{solchap}.\arabic{equation}}

\newcommand{\solchapter}[2]{%
  \clearpage
  \setcounter{solchap}{#1}\setcounter{equation}{0}%
  \phantomsection\label{ch:#1}%
  \addcontentsline{toc}{section}{Chapter #1.\quad #2}%
  \markboth{Solutions to Karatzas and Shreve, \emph{Brownian Motion and Stochastic Calculus}}{Chapter #1: #2}%
  \thispagestyle{plain}%
  \begin{center}
    {\large\bfseries Chapter #1\par}\vspace{0.8ex}
    {\LARGE\bfseries #2\par}
  \end{center}
  \vspace{1.2ex}%
}

% \solsection{1.2}{Stopping Times}
\newcommand{\solsection}[2]{%
  \par\addvspace{2.2ex}\penalty-300\relax
  \phantomsection
  \addcontentsline{toc}{subsection}{Section #1\quad #2}%
  \noindent{\large\bfseries Section #1\quad #2}%
  \par\nobreak\vspace{0.4ex}\nobreak
}

\hypersetup{
  pdftitle={Solutions to Karatzas and Shreve, Brownian Motion and Stochastic Calculus},
  pdfauthor={Zijun Meng},
  bookmarksdepth=2, bookmarksopen=true, bookmarksopenlevel=0}
%% This document carries the solutions alone, so a heading is followed directly by
%% the solution: no italic statement group to open, and no \solution to close one.
\renewcommand{\ks@beginstatement}{\ignorespaces}
\makeatother

\hypersetup{
  pdftitle={Solutions to Karatzas and Shreve, Brownian Motion and Stochastic Calculus},
  pdfauthor={Zijun Meng and Claude Opus 5.5},
  pdfsubject={Solutions to the problems and exercises of I. Karatzas and S. E. Shreve, Brownian Motion and Stochastic Calculus (Springer, second edition, 1991) that the book leaves unsolved; statements not reproduced},
  bookmarksdepth=2, bookmarksopen=true, bookmarksopenlevel=0}

\begin{document}

%% ---- title page -----------------------------------------------------------
\begin{titlepage}
\centering
\vspace*{3.5cm}
{\Huge\bfseries Solutions to\par}
\vspace{1.2cm}
{\LARGE Ioannis Karatzas and Steven E. Shreve\par}
\vspace{0.6cm}
{\LARGE\itshape Brownian Motion and Stochastic Calculus\par}
\vspace{0.4cm}
{\large Springer, Graduate Texts in Mathematics 113, second edition, 1991\par}
\vspace{2.2cm}
{\Large Zijun Meng and Claude Opus 5.5\par}
\vfill
{\large Chapters 1--6\par}
\vspace{0.6cm}
{\small The 208 problems and exercises that the book's own\par
 ``Solutions to Selected Problems'' leave unsolved: 193 for which it prints nothing,\par
 and 15 for which it prints only a hint or a part of the solution.\par
 Each chapter opens with the list of the items it covers\par
 and of those the book already solves.\par}
\vspace{0.8cm}
{\small The statements are not reproduced here: read them in Karatzas and Shreve's
\emph{Brownian Motion and Stochastic Calculus}.\par
 Each solution is headed by the book's own word for the item (Problem or Exercise) and its\par
 number; as in the book, a result is cited by two numbers inside its own chapter\par
 (Theorem 3.8) and by three from another chapter (Theorem 1.3.8).\par}
\vspace{2cm}
\end{titlepage}

%% ---- table of contents ------------------------------------------------------
\begingroup
\setcounter{tocdepth}{2}
\tableofcontents
\endgroup
\clearpage
\solchapter{1}{Martingales, Stopping Times, and Filtrations}

\begin{center}
\begin{minipage}{0.86\textwidth}
\small
The 29 problems and exercises of Chapter 1 for which the book prints no solution: 1.5, 1.7, 1.10, 1.16, 1.17, 2.2, 2.10, 2.13, 2.14, 2.19, 2.21, 2.24, 3.4, 3.9, 3.16, 3.18, 3.21, 3.23, 3.24, 3.25, 3.29, 3.30, 4.13, 5.7, 5.14, 5.19, 5.20, 5.21, 5.26.  For 2 more, 2.6, 2.22, the book prints only a hint or solves only some of the parts; what it leaves out is supplied here.  The book's own ``Solutions to Selected Problems'' fully covers 1.8, 2.7, 2.17, 2.23, 3.2, 3.7, 3.11, 3.19, 3.20, 3.26, 3.27, 3.28, 4.9, 4.11, 4.15, 5.11, 5.12, 5.17, 5.24, 5.25, which are therefore omitted.
\end{minipage}
\end{center}

\bigskip
\hrule
\bigskip

\solsection{1.1}{Stochastic Processes and $\sigma$-Fields}

\begin{kex}{Problem}{1.5}{}
Both processes are defined on the same probability space $(\Omega,\F,P)$ and take values in $\R^d$. The hypothesis on the paths means that there are events $\Omega_X,\Omega_Y\in\F$ with $P(\Omega_X)=P(\Omega_Y)=1$ such that $t\mapsto X_t(\omega)$ is right-continuous on $[0,\infty)$ for every $\omega\in\Omega_X$, and $t\mapsto Y_t(\omega)$ is right-continuous on $[0,\infty)$ for every $\omega\in\Omega_Y$.

Let $\Q_+\defeq\Q\cap[0,\infty)$. For each $q\in\Q_+$ the set $N_q\defeq\{X_q\neq Y_q\}$ belongs to $\F$, because $X_q-Y_q$ is an $\R^d$-valued random variable and $\R^d\setminus\{0\}$ is a Borel set, and $P(N_q)=0$ by \bk{Definition 1.1}. We set
\[
  \Omega^*\defeq\Omega_X\cap\Omega_Y\cap\bigcap_{q\in\Q_+}N_q^c\in\F .
\]
Its complement is the union of the countably many $P$-null events $\Omega_X^c$, $\Omega_Y^c$ and $N_q$, $q\in\Q_+$, so countable subadditivity gives $P(\Omega^*)=1$.

Now fix $\omega\in\Omega^*$ and $t\in[0,\infty)$, and choose a sequence $\{q_n\}_{n=1}^\infty\subseteq\Q_+$ with $q_n>t$ for every $n$ and $q_n\to t$ as $n\to\infty$; for instance, any rational $q_n\in(t,t+\frac1n)$ will do. Since $\omega\notin N_{q_n}$, we have $X_{q_n}(\omega)=Y_{q_n}(\omega)$ for every $n\ge1$. Since $\omega\in\Omega_X\cap\Omega_Y$, both paths are right-continuous at $t$, so each of them converges along any sequence of times that approaches $t$ from the right; therefore
\[
  X_t(\omega)=\lim_{n\to\infty}X_{q_n}(\omega)=\lim_{n\to\infty}Y_{q_n}(\omega)=Y_t(\omega).
\]
As $t\ge0$ was arbitrary, we have shown
\begin{equation}\label{eq:1-1-05:a}
  \Omega^*\subseteq\{\omega\in\Omega;\ X_t(\omega)=Y_t(\omega),\ \forall\,0\le t<\infty\},\qquad P(\Omega^*)=1 .
\end{equation}
Thus almost all sample paths of $X$ and $Y$ agree, i.e., $X$ and $Y$ are indistinguishable in the sense of \bk{Definition 1.3}.

\rmk{The set $A\defeq\{X_t=Y_t,\ \forall\,0\le t<\infty\}$ is an uncountable intersection of events, and \bk{Definition 1.3} tacitly presupposes that its probability makes sense. By \eqref{eq:1-1-05:a} we have $\Omega^*\subseteq A$ and $A\setminus\Omega^*\subseteq(\Omega^*)^c$, a $P$-null event. Hence $A$ differs from the event $\Omega^*$ by a subset of a null event: it belongs to the $P$-completion of $\F$ and has probability one there, and it belongs to $\F$ itself whenever $(\Omega,\F,P)$ is complete or \emph{every} path of $X$ and $Y$ is right-continuous (then $\Omega_X=\Omega_Y=\Omega$ and $A=\bigcap_{q\in\Q_+}N_q^c$). The same remark applies to the phrase ``a.s.\ right-continuous sample paths'', which we have read as the existence of the events $\Omega_X$, $\Omega_Y$. The argument uses the modification property only at the countably many rational times; \bk{Example 1.4} shows that the conclusion fails without a path regularity assumption.}
\end{kex}

\begin{kex}{Exercise}{1.7}{}
The process $X$ takes values in $\R^d$, and $\abs{\cdot}$ denotes the Euclidean norm. If $t_0=0$ the interval $[0,t_0)$ is empty and $A=\Omega\in\F^X_0$, so we assume $t_0>0$. We write $X_{t-}(\omega)\defeq\lim_{s\uparrow t}X_s(\omega)$ for $t>0$; this limit exists in $\R^d$ for every $\omega$ by hypothesis. Put $s_m\defeq t_0\,m/(m+1)$ for $m\ge1$, so that $0<s_m<t_0$ and $s_m\uparrow t_0$, and define, for integers $m,k,n\ge1$,
\[
  B_{m,k,n}\defeq\bigcap_{\substack{q,r\in\Q\cap[0,s_m]\\ \abs{q-r}<1/n}}
  \Bigl\{\abs{X_q-X_r}\le\frac1k\Bigr\},
  \qquad
  B\defeq\bigcap_{m\ge1}\,\bigcap_{k\ge1}\,\bigcup_{n\ge1}B_{m,k,n}.
\]
For $q,r\in[0,s_m]\subseteq[0,t_0]$ the random vectors $X_q$ and $X_r$ are $\F^X_{t_0}$-measurable by the definition of $\F^X_{t_0}=\sigma(X_s;\,0\le s\le t_0)$, and $(x,y)\mapsto\abs{x-y}$ is continuous on $\R^d\times\R^d$; hence $\{\abs{X_q-X_r}\le1/k\}\in\F^X_{t_0}$. Since $B$ is obtained from such events by countable intersections and unions, $B\in\F^X_{t_0}$. It therefore suffices to prove that $A=B$.

\emph{$A\subseteq B$.} Let $\omega\in A$ and fix $m,k\ge1$. The function $t\mapsto X_t(\omega)$ is continuous on the compact interval $[0,s_m]\subset[0,t_0)$, hence uniformly continuous there: there is an integer $n\ge1$ such that $\abs{X_u(\omega)-X_v(\omega)}\le1/k$ whenever $u,v\in[0,s_m]$ and $\abs{u-v}<1/n$. In particular $\omega\in B_{m,k,n}$. As $m$ and $k$ were arbitrary, $\omega\in B$.

\emph{$B\subseteq A$.} Let $\omega\in B$. The path $t\mapsto X_t(\omega)$ is right-continuous at every point of $[0,\infty)$; in particular it is continuous at $t=0$ as a function on $[0,t_0)$, and it is continuous at a point $t\in(0,t_0)$ as soon as $X_{t-}(\omega)=X_t(\omega)$. So fix $t\in(0,t_0)$, choose $m$ with $t<s_m$, and let $k\ge1$ be arbitrary. Since $\omega\in B$, there is $n\ge1$ with $\omega\in B_{m,k,n}$. Choose rationals $q_j\in[0,t)$ with $q_j\uparrow t$ and rationals $r_j\in(t,s_m]$ with $r_j\downarrow t$; this is possible because $0<t<s_m$. For all sufficiently large $j$ we have $\abs{q_j-r_j}<1/n$, and then, because $q_j,r_j\in\Q\cap[0,s_m]$,
\[
  \abs{X_{q_j}(\omega)-X_{r_j}(\omega)}\le\frac1k .
\]
Letting $j\to\infty$, we have $X_{q_j}(\omega)\to X_{t-}(\omega)$ by the existence of the left-hand limit at $t$ (the $q_j$ are strictly less than $t$), and $X_{r_j}(\omega)\to X_t(\omega)$ by right-continuity at $t$. Hence $\abs{X_{t-}(\omega)-X_t(\omega)}\le1/k$. Since $k$ was arbitrary, $X_{t-}(\omega)=X_t(\omega)$, and the path is continuous at $t$. This holds for every $t\in(0,t_0)$, so $\omega\in A$.

Thus $A=B\in\F^X_{t_0}$.

\rmk{The argument uses the hypothesis that \emph{every} path is RCLL in the inclusion $B\subseteq A$: on $B$ we control the path only along rational times, and it is the right-continuity and the left-hand limits of the path of the particular $\omega$ under consideration that transfer this control to all times. Moreover, only the variables $X_q$ with $q\in\Q\cap[0,t_0)$ enter the definition of $B$, so that in fact $A\in\sigma(X_s;\,0\le s<t_0)$.}
\end{kex}

\begin{kex}{Exercise}{1.10}{}
The process $X$ takes values in $\R^d$, and $\abs{\cdot}$ denotes the Euclidean norm. We read ``$X$ is continuous on $[0,t_0]$'' as: the path $t\mapsto X_t(\omega)$, a function on $[0,\infty)$, is continuous at every point of $[0,t_0]$; in particular it is right-continuous at $t_0$. (The other possible reading is treated in the remark at the end.) For a function $x\colon[0,\infty)\to\R^d$ which is LCRL we write $x(t+)\defeq\lim_{s\downarrow t}x(s)$, $t\ge 0$. Since $x$ is left-continuous at every $t>0$, such an $x$ is continuous at a point $t\ge 0$ if and only if $x(t+)=x(t)$.

For integers $m,n\ge 1$ let
\[
  B_{m,n}\defeq\bigcap_{\substack{q,r\in\Q\cap[0,\,t_0+1/n)\\ \abs{q-r}<1/n}}
  \Bigl\{\omega\in\Omega:\ \abs{X_q(\omega)-X_r(\omega)}\le\tfrac1m\Bigr\}.
\]
This is a countable intersection, and for $q,r<t_0+\frac1n$ the random variables $X_q$, $X_r$ are $\F_{t_0+1/n}$-measurable because $X$ is adapted; hence $B_{m,n}\in\F_{t_0+1/n}$. Moreover $B_{m,n}\subseteq B_{m,n+1}$, since passing from $n$ to $n+1$ shrinks the set of pairs $(q,r)$ over which the intersection is taken.

\begin{claim}
$A=\bigcap_{m=1}^\infty\bigcup_{n=1}^\infty B_{m,n}$.
\end{claim}

\begin{proof}
Fix $\omega$ and write $x(t)=X_t(\omega)$; this is an LCRL function.

\emph{The inclusion $\subseteq$.} Let $x$ be continuous at every point of $[0,t_0]$, and let $m\ge1$. The restriction of $x$ to the compact interval $[0,t_0]$ is continuous, hence uniformly continuous: there is $\delta_1>0$ such that $\abs{x(s)-x(t)}\le\frac1{3m}$ whenever $s,t\in[0,t_0]$ and $\abs{s-t}<\delta_1$. By right-continuity at $t_0$ there is $\delta_2>0$ such that $\abs{x(s)-x(t_0)}\le\frac1{3m}$ for $t_0\le s<t_0+\delta_2$. Choose $n$ with $\frac1n<\min(\delta_1,\delta_2)$, and let $q\le r$ be rationals in $[0,t_0+\frac1n)$ with $r-q<\frac1n$. If $r\le t_0$, then $\abs{x(q)-x(r)}\le\frac1{3m}$. If $q\ge t_0$, then $\abs{x(q)-x(r)}\le\abs{x(q)-x(t_0)}+\abs{x(t_0)-x(r)}\le\frac2{3m}$. If $q<t_0<r$, then $t_0-q<\frac1n<\delta_1$ and $r-t_0<\frac1n<\delta_2$, so that again $\abs{x(q)-x(r)}\le\abs{x(q)-x(t_0)}+\abs{x(t_0)-x(r)}\le\frac2{3m}$. In all cases $\abs{x(q)-x(r)}\le\frac1m$, that is, $\omega\in B_{m,n}$.

\emph{The inclusion $\supseteq$.} Let $\omega\in\bigcap_m\bigcup_n B_{m,n}$ and fix $t\in[0,t_0]$; we have to show $x(t+)=x(t)$. Let $m\ge1$ and choose $n$ with $\omega\in B_{m,n}$. Choose rationals $r_j\downarrow t$ with $t<r_j<t+\frac1{2n}$; then $r_j<t_0+\frac1n$. If $t>0$, choose also rationals $q_j\uparrow t$ with $\max(0,t-\frac1{2n})<q_j<t$; if $t=0$, put $q_j=0$ for all $j$. In both cases $q_j,r_j\in\Q\cap[0,t_0+\frac1n)$ and $\abs{q_j-r_j}<\frac1n$, so $\abs{x(q_j)-x(r_j)}\le\frac1m$ for every $j$. As $j\to\infty$ we have $x(r_j)\to x(t+)$ by the existence of the right-hand limit, and $x(q_j)\to x(t)$ (by left-continuity at $t$ if $t>0$, trivially if $t=0$). Therefore $\abs{x(t)-x(t+)}\le\frac1m$. Since $m$ is arbitrary, $x(t+)=x(t)$, and $x$ is continuous at $t$.
\end{proof}

Now fix an integer $k\ge1$. Because $B_{m,n}$ increases with $n$, we have $\bigcup_{n\ge1}B_{m,n}=\bigcup_{n\ge k}B_{m,n}$, and for $n\ge k$ we know $B_{m,n}\in\F_{t_0+1/n}\subseteq\F_{t_0+1/k}$. By the Claim,
\[
  A=\bigcap_{m=1}^\infty\ \bigcup_{n=k}^\infty B_{m,n}\in\F_{t_0+1/k}\qquad\text{for every }k\ge1 .
\]
Since the $\sigma$-fields $\F_{t_0+\eps}$ decrease as $\eps\downarrow0$, this gives $A\in\bigcap_{k\ge1}\F_{t_0+1/k}=\bigcap_{\eps>0}\F_{t_0+\eps}=\F_{t_0+}$, and $\F_{t_0+}=\F_{t_0}$ by the right-continuity of the filtration. Hence $A\in\F_{t_0}$.

\rmk{\emph{(a) The other reading.} If ``continuous on $[0,t_0]$'' is understood as continuity of the restriction of the path to $[0,t_0]$, no condition is imposed to the right of $t_0$, and the corresponding event $A'$ is $\{X_{t+}=X_t \text{ for all } t\in[0,t_0)\}$. The same argument, with $\Q\cap[0,t_0+\frac1n)$ replaced by $\Q\cap[0,t_0]$ in the definition of $B_{m,n}$ (in the second inclusion one takes, for $t<t_0$, the $r_j$ in $(t,t_0)$ as well; for $t_0=0$ there is nothing to prove), shows that $A'=\bigcap_m\bigcup_n B'_{m,n}$ with $B'_{m,n}\in\F^X_{t_0}$. Thus $A'\in\F^X_{t_0}\subseteq\F_{t_0}$ for every filtration to which $X$ is adapted, right-continuous or not.

\emph{(b) Right-continuity of the filtration cannot be omitted for the event $A$.} Let $\Omega=\{0,1\}$, with $\F$ the collection of all subsets of $\Omega$, and $X_t(\omega)=\omega\,\ind_{(t_0,\infty)}(t)$. Every path is LCRL, and $A=\{0\}$. But $X_t\equiv0$ for $t\le t_0$, so $\F^X_{t_0}=\{\emptyset,\Omega\}$ and $A\notin\F^X_{t_0}$, although $A\in\F^X_{t_0+}$.}
\end{kex}

\begin{kex}{Problem}{1.16}{}
Let $X$ have state space $(\R^d,\B(\R^d))$ and be defined on $(\Omega,\F)$. ``Finite'' means $T(\omega)<\infty$ for every $\omega\in\Omega$, so that, by \bk{Definition 1.15}, $X_T(\omega)=X_{T(\omega)}(\omega)$ is defined for every $\omega\in\Omega$; we have to show that $X_T\colon\Omega\to\R^d$ is $\F/\B(\R^d)$-measurable. We write $X_T$ as a composition of two measurable maps.

\emph{Step 1.} Define
\[
\varphi\colon\Omega\to[0,\infty)\times\Omega,\qquad \varphi(\omega)\defeq\bigl(T(\omega),\omega\bigr).
\]
We claim that $\varphi$ is $\F/\bigl(\B([0,\infty))\otimes\F\bigr)$-measurable. The collection
\[
\cD\defeq\bigl\{C\subseteq[0,\infty)\times\Omega:\ \varphi^{-1}(C)\in\F\bigr\}
\]
is a $\sigma$-field on $[0,\infty)\times\Omega$, because taking preimages commutes with complements and countable unions and $\F$ is a $\sigma$-field. For a measurable rectangle $B\times F$ with $B\in\B([0,\infty))$ and $F\in\F$ we have
\[
\varphi^{-1}(B\times F)=\{\omega\in\Omega:\ T(\omega)\in B\}\cap F .
\]
Now $T$ is an $\F$-measurable random variable with values in $[0,\infty]$, and $B$, being a Borel subset of $[0,\infty)$, is also a Borel subset of $[0,\infty]$; hence $\{T\in B\}\in\F$ and $\varphi^{-1}(B\times F)\in\F$. Thus $\cD$ contains all measurable rectangles, and since these generate $\B([0,\infty))\otimes\F$, we conclude $\B([0,\infty))\otimes\F\subseteq\cD$, which is the claim.

\emph{Step 2.} By \bk{Definition 1.6}, the measurability of the process $X$ means that the map
\[
\Phi\colon[0,\infty)\times\Omega\to\R^d,\qquad \Phi(t,\omega)\defeq X_t(\omega),
\]
is $\bigl(\B([0,\infty))\otimes\F\bigr)/\B(\R^d)$-measurable. For every $\omega\in\Omega$,
\[
(\Phi\circ\varphi)(\omega)=\Phi\bigl(T(\omega),\omega\bigr)=X_{T(\omega)}(\omega)=X_T(\omega),
\]
so that for $A\in\B(\R^d)$,
\[
\{X_T\in A\}=\varphi^{-1}\bigl(\Phi^{-1}(A)\bigr)\in\F
\]
by Step 1, because $\Phi^{-1}(A)\in\B([0,\infty))\otimes\F$. Hence $X_T$ is $\F/\B(\R^d)$-measurable, i.e., a random variable on $(\Omega,\F)$.

\rmk{The same argument handles a random time $T$ which may take the value $\infty$: with $\Omega_0\defeq\{T<\infty\}\in\F$ and the trace $\sigma$-field $\F_0\defeq\{F\cap\Omega_0:\ F\in\F\}\subseteq\F$, the map $\varphi$ restricted to $\Omega_0$ is $\F_0/\bigl(\B([0,\infty))\otimes\F\bigr)$-measurable by Step 1, and therefore $\{X_T\in A\}=\{\omega\in\Omega_0:\ X_{T(\omega)}(\omega)\in A\}$ belongs to $\F$ for every $A\in\B(\R^d)$. In particular, if $T$ is only almost surely finite, then $X_T$, defined on $\{T<\infty\}$ and set equal to an arbitrary constant on the null event $\{T=\infty\}$, is a random variable. This is the measurability needed in \bk{Problem 1.17}.}
\end{kex}

\begin{kex}{Problem}{1.17}{}
The processes of this section take values in $\R^d$, so we let $A$ range over $\B(\R^d)$; for $d=1$ this is the statement as printed. No variable $X_\infty$ is given, so by \bk{Definition 1.15} the function $X_T$ is defined on the event $\Omega_0\defeq\{T<\infty\}$ only, and
\[
  E_A\defeq\{X_T\in A\}=\{\omega\in\Omega_0:\ X_{T(\omega)}(\omega)\in A\},\qquad A\in\B(\R^d).
\]
We write $N\defeq\{T=\infty\}=\Omega\setminus\Omega_0$; both $\Omega_0$ and $N$ belong to $\F$ because $T$ is an $\F$-measurable map into $[0,\infty]$. The collection in question is
\[
  \G\defeq\{E_A:\ A\in\B(\R^d)\}\cup\{E_A\cup N:\ A\in\B(\R^d)\}.
\]

\emph{Step 1: $\G\subseteq\F$.} Let $S\defeq T\,\ind_{\Omega_0}$, that is, $S=T$ on $\Omega_0$ and $S=0$ on $N$. For $b\ge 0$ we have $\{S\le b\}=\{T\le b\}\cup N\in\F$, so $S$ is a finite random time. The map $\omega\mapsto(S(\omega),\omega)$ from $(\Omega,\F)$ into $([0,\infty)\times\Omega,\ \B([0,\infty))\otimes\F)$ is measurable, since the preimage of a rectangle $B\times F$ with $B\in\B([0,\infty))$, $F\in\F$, is $\{S\in B\}\cap F\in\F$, and such rectangles generate the product $\sigma$-field. Composing it with the map $(s,\omega)\mapsto X_s(\omega)$, which is $\B([0,\infty))\otimes\F$-measurable because $X$ is a measurable process (\bk{Definition 1.6}), we see that $X_S:\omega\mapsto X_{S(\omega)}(\omega)$ is an $\F$-measurable random variable with values in $\R^d$; this is \bk{Problem 1.16}. Since $X_S=X_T$ on $\Omega_0$,
\begin{equation}\label{eq:1-1-17:a}
  E_A=\{X_S\in A\}\cap\Omega_0\in\F,\qquad A\in\B(\R^d),
\end{equation}
and consequently $E_A\cup N\in\F$ as well.

\emph{Step 2: $\G$ is a $\sigma$-field.} From \eqref{eq:1-1-17:a} and the fact that preimages under $X_S$ commute with set operations, we obtain for $A,A_1,A_2,\dots\in\B(\R^d)$:
\begin{equation}\label{eq:1-1-17:b}
  E_{\R^d}=\Omega_0,\qquad
  \Omega_0\setminus E_A=E_{A^c},\qquad
  \bigcup_{n=1}^\infty E_{A_n}=E_{\bigcup_{n}A_n},
\end{equation}
where $A^c=\R^d\setminus A$. Moreover $E_A\subseteq\Omega_0$, so that $E_A$ and $N$ are disjoint.

\textbf{(i)} $\Omega=\Omega_0\cup N=E_{\R^d}\cup N\in\G$.

\textbf{(ii)} $\G$ is closed under complementation. Indeed, by \eqref{eq:1-1-17:b} and the disjointness of $E_A$ and $N$,
\[
  \Omega\setminus E_A=(\Omega_0\setminus E_A)\cup N=E_{A^c}\cup N\in\G,
  \qquad
  \Omega\setminus(E_A\cup N)=\Omega_0\setminus E_A=E_{A^c}\in\G .
\]

\textbf{(iii)} $\G$ is closed under countable unions. Let $G_1,G_2,\dots\in\G$; for each $n$ there is $A_n\in\B(\R^d)$ such that $G_n=E_{A_n}$ or $G_n=E_{A_n}\cup N$. With $A\defeq\bigcup_n A_n\in\B(\R^d)$, \eqref{eq:1-1-17:b} gives
\[
  \bigcup_{n=1}^\infty G_n=
  \begin{cases}
    E_A, & \text{if } G_n=E_{A_n} \text{ for every } n,\\[2pt]
    E_A\cup N, & \text{if } G_n=E_{A_n}\cup N \text{ for at least one } n,
  \end{cases}
\]
and in either case the union belongs to $\G$.

By (i)--(iii), $\G$ is a $\sigma$-field, and by Step 1 it is contained in $\F$.

\rmk{The book prints $A\in\B(\R)$ although the state space of the section is $\R^d$; the argument is the same for every $d$. If $T$ is finite, then $N=\emptyset$ and $\G=\{\{X_T\in A\}:A\in\B(\R^d)\}$ is the usual $\sigma$-field $\sigma(X_T)$ generated by the random variable $X_T$. In general, $\G$ is the $\sigma$-field generated by the map which equals $X_T$ on $\{T<\infty\}$ and takes an extra value $\Delta\notin\R^d$ on $\{T=\infty\}$, the space $\R^d\cup\{\Delta\}$ carrying the $\sigma$-field $\{A,\,A\cup\{\Delta\}:A\in\B(\R^d)\}$. Note that $N=E_\emptyset\cup N$ and $\Omega_0=E_{\R^d}$ always belong to $\G$, also when one of them is empty.}
\end{kex}

\solsection{1.2}{Stopping Times}

\begin{kex}{Problem}{2.2}{}
Throughout, $\omega$ and $\omega'$ are the two fixed points of the statement, $X$ takes values in $(\R^d,\B(\R^d))$, and $\F_t^X=\sigma(X_s;\,0\le s\le t)$ as in \bk{Section 1.1}. The argument rests on the fact that an event of $\F_t^X$ cannot separate two points of $\Omega$ at which the path of $X$ agrees up to time $t$.

\begin{lem}
Let $0\le t<\infty$ and suppose that $X_s(\omega)=X_s(\omega')$ for every $s\in[0,t]$. Then $\ind_A(\omega)=\ind_A(\omega')$ for every $A\in\F_t^X$.
\end{lem}

\begin{proof}
Let $\G$ be the collection of all subsets $A$ of $\Omega$ such that $\ind_A(\omega)=\ind_A(\omega')$, that is, such that $A$ contains either both of the points $\omega$, $\omega'$ or neither of them. Then $\G$ is a $\sigma$-field: it contains $\Omega$; if $A\in\G$, then $\ind_{A^c}=1-\ind_A$ takes the same value at $\omega$ and at $\omega'$, so $A^c\in\G$; and if $A_1,A_2,\ldots\in\G$ and $A=\bigcup_{n}A_n$, then $\ind_A=\sup_n\ind_{A_n}$ takes the same value at $\omega$ and at $\omega'$, so $A\in\G$. Moreover, for $s\in[0,t]$ and $\Gamma\in\B(\R^d)$ the hypothesis $X_s(\omega)=X_s(\omega')$ gives
\[
\ind_{\{X_s\in\Gamma\}}(\omega)=\ind_\Gamma\bigl(X_s(\omega)\bigr)=\ind_\Gamma\bigl(X_s(\omega')\bigr)=\ind_{\{X_s\in\Gamma\}}(\omega'),
\]
so $\{X_s\in\Gamma\}\in\G$. Thus every $X_s$, $0\le s\le t$, is measurable with respect to the $\sigma$-field $\G$, and since $\F_t^X$ is the smallest $\sigma$-field with this property, $\F_t^X\subseteq\G$.
\end{proof}

We now distinguish two cases, according to whether $T(\omega)$ is finite.

\emph{Case 1: $t\defeq T(\omega)<\infty$.} The hypothesis says that $X_s(\omega)=X_s(\omega')$ for all $s\in[0,t]$, so the Lemma applies with this $t$. Since $T$ is a stopping time of $\{\F_t^X\}$, \bk{Definition 2.1} gives $\{T\le t\}\in\F_t^X$ and, for every integer $n\ge1$ with $t-\frac1n\ge0$, $\{T\le t-\frac1n\}\in\F^X_{t-(1/n)}\subseteq\F_t^X$. Because a random time takes values in $[0,\infty]$ (\bk{Section 1.1}), so that $T\ge0$,
\[
\{T=t\}=\{T\le t\}\setminus\bigcup_{n\ge1,\ t-(1/n)\ge0}\Bigl\{T\le t-\frac1n\Bigr\}
\]
(the union being empty if $t=0$), and therefore $\{T=t\}\in\F_t^X$. The Lemma yields $\ind_{\{T=t\}}(\omega')=\ind_{\{T=t\}}(\omega)=1$, that is, $T(\omega')=t=T(\omega)$.

\emph{Case 2: $T(\omega)=\infty$.} Now $[0,T(\omega)]\cap[0,\infty)=[0,\infty)$, so $X_s(\omega)=X_s(\omega')$ for all $s\in[0,\infty)$, and the Lemma applies with every $t\in[0,\infty)$. For each such $t$ we have $\{T\le t\}\in\F_t^X$, again by \bk{Definition 2.1}, and $\omega\notin\{T\le t\}$; hence $\ind_{\{T\le t\}}(\omega')=\ind_{\{T\le t\}}(\omega)=0$, i.e.\ $T(\omega')>t$. As $t\in[0,\infty)$ is arbitrary, $T(\omega')=\infty=T(\omega)$.

In both cases $T(\omega)=T(\omega')$.
\end{kex}

\begin{kex}{Problem}{2.6}{}
The book's Section 1.6 only suggests the identity \eqref{eq:1-2-06:id} below and notes that one inclusion is obvious; we prove both inclusions and draw the conclusion. Recall the setting of \bk{Example 2.5}: $X$ is an $\R^d$-valued process, adapted to the filtration $\{\F_t\}$, every path $t\mapsto X_t(\omega)$ is right-continuous on $[0,\infty)$, and $H_\Gamma(\omega)=\inf\{t\ge 0;\ X_t(\omega)\in\Gamma\}$ with $\inf\emptyset=\infty$. Since the definition of an optional time (\bk{Definition 2.1}) presupposes a random time, i.e.\ an $\F$-measurable map into $[0,\infty]$, and since $\{H_\Gamma<\infty\}=\bigcup_{n\ge 1}\{H_\Gamma<n\}$ and $\F_t\subseteq\F$, everything follows once we show that $\{H_\Gamma<t\}\in\F_t$ for every $t\ge 0$.

For $t=0$ we have $\{H_\Gamma<0\}=\emptyset\in\F_0$. Fix $t>0$. We claim that
\begin{equation}\label{eq:1-2-06:id}
  \{H_\Gamma<t\}\;=\;\bigcup_{s\in\Q\cap[0,t)}\{X_s\in\Gamma\}.
\end{equation}

\emph{The inclusion $\supseteq$.} If $X_s(\omega)\in\Gamma$ for some $s\in[0,t)$, then $s$ belongs to the set whose infimum is $H_\Gamma(\omega)$, so $H_\Gamma(\omega)\le s<t$. (This uses neither the openness of $\Gamma$ nor the rationality of $s$.)

\emph{The inclusion $\subseteq$.} Let $H_\Gamma(\omega)<t$. Then the set $\{u\ge 0;\ X_u(\omega)\in\Gamma\}$ is nonempty and its infimum is strictly smaller than $t$; hence $t$ is not a lower bound of this set, and there exists $u\in[0,t)$ with $X_u(\omega)\in\Gamma$. Because $\Gamma$ is open, there is $\eps>0$ such that every $y\in\R^d$ with $\norm{y-X_u(\omega)}<\eps$ lies in $\Gamma$. Because the path $s\mapsto X_s(\omega)$ is right-continuous at $u$, there is $\delta>0$, which we may take so small that $u+\delta\le t$, such that $\norm{X_s(\omega)-X_u(\omega)}<\eps$ for all $s\in[u,u+\delta)$. The interval $[u,u+\delta)$ has nonempty interior and therefore contains a rational number $s$; for it, $0\le s<t$ and $X_s(\omega)\in\Gamma$. Thus $\omega$ belongs to the right-hand side of \eqref{eq:1-2-06:id}.

\emph{Conclusion.} An open set is a Borel set, and $X$ is adapted, so for every $s\in[0,t)$ the event $\{X_s\in\Gamma\}=X_s^{-1}(\Gamma)$ belongs to $\F_s\subseteq\F_t$. The union in \eqref{eq:1-2-06:id} is countable, whence $\{H_\Gamma<t\}\in\F_t$. As $t\ge 0$ was arbitrary, $H_\Gamma$ is an optional time of $\{\F_t\}$.

\rmk{``Optional'' cannot be improved to ``stopping'' without right-continuity of the filtration (\bk{Proposition 2.3}). Take $\Omega=\{a,b\}$ with $\F$ the collection of all its subsets, $X_t(a)=(t-1)^+$, $X_t(b)=0$, the filtration $\F_t=\F^X_t=\sigma(X_s;\ 0\le s\le t)$, and $\Gamma=(0,\infty)$. The paths are continuous, $H_\Gamma(a)=1$ and $H_\Gamma(b)=\infty$; but $X_s(a)=X_s(b)=0$ for $s\le 1$, so $\F^X_1=\{\emptyset,\Omega\}$ does not contain $\{H_\Gamma\le 1\}=\{a\}$.}
\end{kex}

\begin{kex}{Problem}{2.10}{}
Random times take values in $[0,\infty]$, and the conditions $T>0$, $S>0$ are understood to hold at every $\omega\in\Omega$. We write $\Q$ for the set of rational numbers and record first a fact used repeatedly.

\begin{claim}
If $R$ is an optional time and $0\le u<t$, then $\{R\le u\}\in\F_t$ and $\{R>u\}\in\F_t$.
\end{claim}
\begin{proof}
Choose an integer $m\ge 1$ with $u+\frac1m\le t$. Then
$\{R\le u\}=\bigcap_{n\ge m}\{R<u+\frac1n\}$, and each set of this intersection belongs to
$\F_{u+(1/n)}\subseteq\F_{u+(1/m)}\subseteq\F_t$. The second assertion follows by taking complements.
\end{proof}

\emph{$T+S$ is optional.} For $t=0$ we have $\{T+S<0\}=\emptyset\in\F_0$. For $t>0$ we show that
\begin{equation}\label{eq:1-2-10:a}
\{T+S<t\}=\bigcup_{\substack{r\in\Q\\ 0<r<t}}\{T<r,\ S<t-r\}.
\end{equation}
If $T(\omega)<r$ and $S(\omega)<t-r$, then $T(\omega)+S(\omega)<t$. Conversely, if $T(\omega)+S(\omega)<t$, then both
values are finite and $0\le T(\omega)<t-S(\omega)\le t$, so there is a rational $r$ with $T(\omega)<r<t-S(\omega)$; it
satisfies $0<r<t$, $T(\omega)<r$ and $S(\omega)<t-r$. This proves \eqref{eq:1-2-10:a}. For each $r$ in the union,
$\{T<r\}\in\F_r\subseteq\F_t$ and $\{S<t-r\}\in\F_{t-r}\subseteq\F_t$ because $T$ and $S$ are optional; the union being
countable, $\{T+S<t\}\in\F_t$. Hence $T+S$ is optional.

\emph{A decomposition.} Assume now that $T>0$, as in both (i) and (ii). Then $T+S>0$, so
$\{T+S\le 0\}=\emptyset\in\F_0$, and it remains to show $\{T+S>t\}\in\F_t$ for every fixed $t>0$. We claim that
\begin{equation}\label{eq:1-2-10:b}
\{T<t,\ T+S>t\}=\bigcup_{\substack{r\in\Q\\ 0<r<t}}\{r<T<t,\ S>t-r\}\in\F_t .
\end{equation}
If $r<T(\omega)<t$ and $S(\omega)>t-r$, then $T(\omega)+S(\omega)>r+(t-r)=t$. Conversely, let $T(\omega)<t$ and
$T(\omega)+S(\omega)>t$. Then $T(\omega)$ is finite and $T(\omega)>t-S(\omega)$ (the right-hand side being $-\infty$ if
$S(\omega)=\infty$); since also $T(\omega)>0$, there is a rational $r$ with
$\max\{0,\,t-S(\omega)\}<r<T(\omega)$, and then $0<r<t$, $r<T(\omega)<t$ and $S(\omega)>t-r$. This proves the equality
in \eqref{eq:1-2-10:b}. As for the measurability: $\{T<t\}\in\F_t$ by optionality, while $\{T>r\}\in\F_t$ and
$\{S>t-r\}\in\F_t$ by the Claim, applied with $u=r<t$ and with $u=t-r<t$ respectively; the latter inequality is where
$r>0$, and thus the hypothesis $T>0$, is needed.

\textbf{(i)} Let $T>0$ and $S>0$. On $\{T\ge t\}$ we have $T+S>t$ because $S>0$; therefore
\[
\{T+S>t\}=\{T\ge t\}\cup\{T<t,\ T+S>t\}.
\]
The first set is the complement of $\{T<t\}\in\F_t$, and the second is in $\F_t$ by \eqref{eq:1-2-10:b}. Thus
$\{T+S\le t\}\in\F_t$ for every $t\ge 0$, and $T+S$ is a stopping time.

\textbf{(ii)} Let $T>0$ be a stopping time and $S$ an optional time. Splitting according to $T>t$, $T=t$, $T<t$, we
obtain
\[
\{T+S>t\}=\{T>t\}\cup\{T=t,\ S>0\}\cup\{T<t,\ T+S>t\};
\]
indeed $T+S>t$ holds everywhere on $\{T>t\}$, and on $\{T=t\}$ it holds exactly when $S>0$. Now
$\{T>t\}\in\F_t$ and $\{T=t\}=\{T\le t\}\setminus\{T<t\}\in\F_t$, because $T$ is a stopping time and hence optional
(\bk{Proposition 2.3}); $\{S>0\}\in\F_t$ by the Claim with $u=0<t$; and the third set is in $\F_t$ by
\eqref{eq:1-2-10:b}. Hence $\{T+S\le t\}\in\F_t$ for every $t\ge 0$, and $T+S$ is a stopping time.

\rmk{The extra hypotheses cannot simply be dropped. Suppose that the filtration admits an optional time $R$ with
$\{R\le t_0\}\notin\F_{t_0}$ for some $t_0>0$ (by \bk{Proposition 2.3} this requires a filtration which is not
right-continuous). The random time $T\defeq R\vee\frac{t_0}{2}$ is positive and optional, since
$\{T<t\}=\{R<t\}\cap\{\frac{t_0}{2}<t\}\in\F_t$, but $\{T\le t_0\}=\{R\le t_0\}\notin\F_{t_0}$. With $S\equiv 0$ the
sum $T+S=T$ is not a stopping time: thus $S>0$ cannot be omitted in (i), nor can ``$T$ is a stopping time'' in (ii).
With $T\equiv 0$ (a stopping time) and $S=R$, the sum $T+S=R$ is not a stopping time: thus $T>0$ cannot be omitted in
(ii).}
\end{kex}

\begin{kex}{Problem}{2.13}{}
Throughout, $T$ is a stopping time of the filtration $\{\F_t\}$ in the sense of \bk{Definition 2.1}: $T$ takes values in $[0,\infty]$ and $\{T\le s\}\in\F_s$ for every $s\ge 0$. By \bk{Definition 2.12},
\[
  \F_T=\bigl\{A\in\F:\ A\cap\{T\le s\}\in\F_s \ \text{for every } s\ge 0\bigr\}.
\]

\emph{Step 1: $\F_T$ is a $\sigma$-field.} Fix $s\ge 0$.

\textbf{(i)} $\Omega\in\F$ and $\Omega\cap\{T\le s\}=\{T\le s\}\in\F_s$, because $T$ is a stopping time; hence $\Omega\in\F_T$.

\textbf{(ii)} Let $A\in\F_T$. Then $A^c\in\F$, and
\[
  A^c\cap\{T\le s\}=\{T\le s\}\cap\bigl(A\cap\{T\le s\}\bigr)^c
\]
is the intersection of two sets of $\F_s$, the first because $T$ is a stopping time and the second because $A\in\F_T$ and $\F_s$ is closed under complementation. Hence $A^c\in\F_T$.

\textbf{(iii)} Let $A_1,A_2,\dots\in\F_T$. Then $\bigcup_{n\ge1}A_n\in\F$, and
\[
  \Bigl(\bigcup_{n=1}^\infty A_n\Bigr)\cap\{T\le s\}=\bigcup_{n=1}^\infty\bigl(A_n\cap\{T\le s\}\bigr)\in\F_s
\]
as a countable union of sets of $\F_s$. Hence $\bigcup_{n\ge1}A_n\in\F_T$.

Since $s\ge0$ was arbitrary, (i)--(iii) show that $\F_T$ is a $\sigma$-field (a sub-$\sigma$-field of $\F$).

\emph{Step 2: $T$ is $\F_T$-measurable.} The random time $T$ takes values in $[0,\infty]$, and the Borel $\sigma$-field $\B([0,\infty])$ is generated by the intervals $[0,r]$, $0\le r<\infty$; indeed, the $\sigma$-field they generate contains $(a,b]=[0,b]\setminus[0,a]$ for $0\le a<b<\infty$, the set $\{0\}=[0,0]$ and the set $\{\infty\}=[0,\infty]\setminus\bigcup_{n\ge1}[0,n]$, hence every open subset of $[0,\infty]$, each of which is a countable union of sets of the forms $[0,b)=\bigcup_{n>1/b}[0,b-\tfrac1n]$, $(a,b)=\bigcup_{n>1/(b-a)}(a,b-\tfrac1n]$ and $(a,\infty]=\bigcup_{n}(a,a+n]\cup\{\infty\}$. Because $\{B\in\B([0,\infty]):\ T^{-1}(B)\in\F_T\}$ is a $\sigma$-field (by Step~1), it therefore suffices to show that
\[
  \{T\le r\}\in\F_T\qquad\text{for every } r\in[0,\infty).
\]
Fix such an $r$. First, $\{T\le r\}\in\F_r\subseteq\F$. Next, for every $s\ge0$,
\[
  \{T\le r\}\cap\{T\le s\}=\{T\le r\wedge s\}\in\F_{r\wedge s}\subseteq\F_s ,
\]
because $T$ is a stopping time and the filtration is nondecreasing. By \bk{Definition 2.12}, $\{T\le r\}\in\F_T$.

\emph{Step 3: the case of a constant time.} Let $T(\omega)=t$ for every $\omega\in\Omega$, where $t\in[0,\infty)$ is a constant. Then $T$ is a stopping time, since $\{T\le s\}$ is $\emptyset$ for $s<t$ and $\Omega$ for $s\ge t$. For $A\in\F$ this gives
\[
  A\cap\{T\le s\}=
  \begin{cases}
    \emptyset, & 0\le s<t,\\
    A, & s\ge t.
  \end{cases}
\]
If $A\in\F_T$, then taking $s=t$ in \bk{Definition 2.12} yields $A=A\cap\{T\le t\}\in\F_t$; thus $\F_T\subseteq\F_t$. Conversely, if $A\in\F_t$, then $A\in\F$, and $A\cap\{T\le s\}$ equals $\emptyset\in\F_s$ for $s<t$ and equals $A\in\F_t\subseteq\F_s$ for $s\ge t$; thus $A\in\F_T$ and $\F_t\subseteq\F_T$. Therefore $\F_T=\F_t$, so that the notation $\F_T$ is consistent with the notation for the filtration.
\end{kex}

\begin{kex}{Exercise}{2.14}{}
Fix $t\ge 0$; by \bk{Definition 2.1} we have to show that $\{S\le t\}\in\F_t$. The random time $S$ takes values in $[0,\infty]$, and $[0,t]$ is a Borel subset of $[0,\infty]$, so the $\F_T$-measurability of $S$ gives
\[
A\defeq\{S\le t\}=S^{-1}([0,t])\in\F_T .
\]
By \bk{Definition 2.12}, applied to this event $A$ and to this same $t$, we obtain $A\cap\{T\le t\}\in\F_t$. On the other hand, $T(\omega)\le S(\omega)$ for every $\omega\in\Omega$, so $S(\omega)\le t$ implies $T(\omega)\le t$; that is, $\{S\le t\}\subseteq\{T\le t\}$. Consequently
\[
\{S\le t\}=\{S\le t\}\cap\{T\le t\}=A\cap\{T\le t\}\in\F_t .
\]
Since $t\ge 0$ was arbitrary, $S$ is a stopping time of $\{\F_t\}$.

\rmk{The hypothesis $S\ge T$ is used on all of $\Omega$, and no completeness or right-continuity of the filtration is needed. It cannot be dropped: if $T\equiv 1$, then $\F_T=\F_1$ by \bk{Problem 2.13}, and $S=\tfrac12\ind_A$ with $A\in\F_1\setminus\F_{1/2}$ (whenever such an event exists) is $\F_T$-measurable, but $\{S\le 0\}=A^c\notin\F_0$ because $\F_0\subseteq\F_{1/2}$; so $S$ is not a stopping time.}
\end{kex}

\begin{kex}{Problem}{2.19}{}
The assumptions are that $X=\{X_t,\F_t;\,0\le t<\infty\}$ is an $\R^d$-valued progressively measurable process and that $T$ is a stopping time of $\{\F_t\}$. As in \bk{Proposition 2.18}, $Y_T$ is the function of \bk{Definition 1.15}, defined on $\{T<\infty\}$. We fix $t\ge 0$ throughout Steps 1--3 and write $K\defeq\sup\abs{f}<\infty$.

\emph{Step 1: the integrand.} We claim that
\[
g(s,\omega)\defeq f\bigl(s,X_s(\omega)\bigr),\qquad (s,\omega)\in[0,t]\times\Omega,
\]
is $\B([0,t])\otimes\F_t$-measurable. Consider the map $\Phi(s,\omega)\defeq\bigl(s,X_s(\omega)\bigr)$ from $[0,t]\times\Omega$ into $[0,\infty)\times\R^d$. Its first component $(s,\omega)\mapsto s$ is $\B([0,t])\otimes\F_t/\B([0,\infty))$-measurable, since the preimage of $A\in\B([0,\infty))$ is $(A\cap[0,t])\times\Omega$; its second component $(s,\omega)\mapsto X_s(\omega)$ is $\B([0,t])\otimes\F_t/\B(\R^d)$-measurable by progressive measurability of $X$ (\bk{Definition 1.11}). Hence $\Phi^{-1}(A\times B)\in\B([0,t])\otimes\F_t$ for all $A\in\B([0,\infty))$, $B\in\B(\R^d)$, and since these rectangles generate $\B([0,\infty))\otimes\B(\R^d)$, the map $\Phi$ is $\B([0,t])\otimes\F_t/\B([0,\infty))\otimes\B(\R^d)$-measurable. Therefore $g=f\circ\Phi$ is $\B([0,t])\otimes\F_t$-measurable, and $\abs{g}\le K$.

In particular, for every $\omega\in\Omega$ the section $s\mapsto f(s,X_s(\omega))$ is a Borel-measurable function on $[0,t]$ bounded by $K$, so the Lebesgue integral $Y_u(\omega)=\int_0^u f(s,X_s(\omega))\,ds$ exists and is finite for every $u\in[0,t]$ and every $\omega$; as $t$ is arbitrary, $Y$ is well defined on $[0,\infty)\times\Omega$, with $Y_0=0$.

\emph{Step 2: a measurability lemma.} We use the measurability part of the Fubini--Tonelli theorem in the following form, which involves no measure on the second factor.

\begin{lem}
Let $(E,\mathscr{E})$ be a measurable space and $h\colon[0,t]\times E\to\R$ a bounded, $\B([0,t])\otimes\mathscr{E}$-measurable function. Then $e\mapsto\int_0^t h(s,e)\,ds$ is $\mathscr{E}$-measurable.
\end{lem}
\begin{proof}
Let $\cH$ be the class of bounded, $\B([0,t])\otimes\mathscr{E}$-measurable $h$ for which the conclusion holds; each section $h(\cdot,e)$ of such an $h$ is Borel measurable and bounded, so the integral makes sense. The class $\cH$ is a vector space, it contains the constants, and it is closed under bounded monotone limits by the bounded convergence theorem on $[0,t]$ (a pointwise limit of $\mathscr{E}$-measurable functions being $\mathscr{E}$-measurable). It contains the indicators of the rectangles $A\times C$, $A\in\B([0,t])$, $C\in\mathscr{E}$, because $\int_0^t\ind_{A\times C}(s,e)\,ds=\meas(A)\,\ind_C(e)$. These rectangles form a $\pi$-system generating $\B([0,t])\otimes\mathscr{E}$, so the monotone class theorem shows that $\cH$ contains every bounded $\B([0,t])\otimes\mathscr{E}$-measurable function.
\end{proof}

\emph{Step 3: joint measurability of $Y$ on $[0,t]\times\Omega$.} Take $(E,\mathscr{E})\defeq\bigl([0,t]\times\Omega,\;\B([0,t])\otimes\F_t\bigr)$ and
\[
h\bigl(s,(u,\omega)\bigr)\defeq\ind_{\{s\le u\}}\,g(s,\omega),\qquad s\in[0,t],\ (u,\omega)\in[0,t]\times\Omega .
\]
The function $(s,u,\omega)\mapsto g(s,\omega)$ is the composition of $g$ with the projection $(s,u,\omega)\mapsto(s,\omega)$, which is $\B([0,t])\otimes\B([0,t])\otimes\F_t/\B([0,t])\otimes\F_t$-measurable; the set $\{(s,u,\omega):s\le u\}=\{(s,u)\in[0,t]^2:s\le u\}\times\Omega$ belongs to $\B([0,t])\otimes\B([0,t])\otimes\F_t$, because $\{s\le u\}$ is closed in $[0,t]^2$ and $\B([0,t]^2)=\B([0,t])\otimes\B([0,t])$. Since $\B([0,t])\otimes\B([0,t])\otimes\F_t=\B([0,t])\otimes\mathscr{E}$, the function $h$ is bounded and $\B([0,t])\otimes\mathscr{E}$-measurable. The lemma now shows that
\[
(u,\omega)\longmapsto\int_0^t h\bigl(s,(u,\omega)\bigr)\,ds=\int_0^u f\bigl(s,X_s(\omega)\bigr)\,ds=Y_u(\omega)
\]
is $\B([0,t])\otimes\F_t$-measurable on $[0,t]\times\Omega$. As $t\ge0$ was arbitrary, this is precisely the statement that $Y$ is progressively measurable with respect to $\{\F_t\}$ (\bk{Definition 1.11}); in particular $Y$ is adapted, each $Y_t$ being the section at $u=t$ of this map.

\emph{Step 4: $Y_T$.} The real-valued process $Y=\{Y_t,\F_t;\,0\le t<\infty\}$ is progressively measurable by Step 3, and $T$ is a stopping time of $\{\F_t\}$. \bk{Proposition 2.18}, applied to $Y$ in place of $X$ (with $d=1$), yields that the random variable $Y_T=\int_0^T f(s,X_s)\,ds$, defined on $\{T<\infty\}\in\F_T$, is $\F_T$-measurable. Explicitly, the proposition gives $\{Y_T\in B\}\cap\{T\le t\}\in\F_t$ for every $B\in\B(\R)$ and $t\ge0$. Consequently the event
\[
\{Y_T\in B\}\cap\{T<\infty\}=\bigcup_{n\ge1}\bigl(\{Y_T\in B\}\cap\{T\le n\}\bigr)
\]
lies in $\F_\infty\subseteq\F$, each set of the union being in $\F_n$; and its intersection with $\{T\le t\}$ is $\{Y_T\in B\}\cap\{T\le t\}\in\F_t$ for every $t\ge0$. Hence $\{Y_T\in B\}\cap\{T<\infty\}\in\F_T$ by \bk{Definition 2.12}. The same proposition shows moreover that the stopped process $\{Y_{T\wedge t},\F_t;\,0\le t<\infty\}$ is progressively measurable.

\rmk{Since $\abs{Y_u-Y_v}\le K\abs{u-v}$, the process $Y$ has continuous paths; once adaptedness is known (the lemma with $E=\Omega$, $\mathscr{E}=\F_t$, $h=g$), progressive measurability also follows from \bk{Proposition 1.13}. On $\{T=\infty\}$ the variable $Y_T$ is in general undefined, because $\int_0^\infty f(s,X_s)\,ds$ need not exist for a merely bounded $f$ (take $f\equiv1$); if $Y_\infty(\omega)=\lim_{u\to\infty}Y_u(\omega)$ exists in $\R$ for every $\omega$ (for instance if $\abs{f(s,x)}\le\varphi(s)$ for all $(s,x)$, with $\varphi\ge0$ Borel measurable and $\int_0^\infty\varphi(s)\,ds<\infty$, in which case $s\mapsto f(s,X_s(\omega))$ is Lebesgue integrable on $[0,\infty)$ and $Y_u(\omega)\to\int_0^\infty f(s,X_s(\omega))\,ds$ by dominated convergence), then $Y_\infty$ is $\F_\infty$-measurable as a limit of the $\F_\infty$-measurable variables $Y_n$, the extension $Y_T\defeq Y_\infty$ on $\{T=\infty\}$ of \bk{Definition 1.15} is available, and $Y_T$ remains $\F_T$-measurable on all of $\Omega$: for $B\in\B(\R)$ the event $\{Y_T\in B\}=\bigl(\{Y_T\in B\}\cap\{T<\infty\}\bigr)\cup\bigl(\{Y_\infty\in B\}\cap\{T=\infty\}\bigr)$ lies in $\F_\infty\subseteq\F$, and its intersection with $\{T\le t\}$ is $\{Y_T\in B\}\cap\{T\le t\}\in\F_t$.}
\end{kex}

\begin{kex}{Problem}{2.21}{}
Throughout, $T$ is an optional time of $\{\F_t\}$ with values in $[0,\infty]$, and $\F_{t+}=\bigcap_{\eps>0}\F_{t+\eps}$. Since the filtration is nondecreasing, we have
\begin{equation}\label{eq:1-2-21:a}
\F_{t+}=\bigcap_{n\ge m}\F_{t+(1/n)}\quad\text{for every integer } m\ge 1,
\qquad\text{and}\qquad \F_s\subseteq\F_{s+}\subseteq\F_t\quad\text{for } 0\le s<t .
\end{equation}
By \bk{Corollary 2.4}, $T$ is a stopping time of $\{\F_{t+}\}$, that is,
\begin{equation}\label{eq:1-2-21:b}
\{T\le t\}\in\F_{t+}\qquad\text{for every } t\ge 0 .
\end{equation}

\emph{Step 1: $\F_{T+}$ is a $\sigma$-field.} Fix $t\ge0$. We have $\Omega\in\F$ and $\Omega\cap\{T\le t\}=\{T\le t\}\in\F_{t+}$ by \eqref{eq:1-2-21:b}, so $\Omega\in\F_{T+}$. If $A\in\F_{T+}$, then $A^c\in\F$ and
\[
A^c\cap\{T\le t\}=\{T\le t\}\setminus\bigl(A\cap\{T\le t\}\bigr)\in\F_{t+},
\]
as a difference of two members of the $\sigma$-field $\F_{t+}$; hence $A^c\in\F_{T+}$. If $A_1,A_2,\dots\in\F_{T+}$, then $\bigcup_k A_k\in\F$ and
\[
\Bigl(\bigcup_{k=1}^\infty A_k\Bigr)\cap\{T\le t\}=\bigcup_{k=1}^\infty\bigl(A_k\cap\{T\le t\}\bigr)\in\F_{t+},
\]
so $\bigcup_k A_k\in\F_{T+}$. Thus $\F_{T+}$ is a $\sigma$-field.

\emph{Step 2: $T$ is $\F_{T+}$-measurable.} The intervals $[0,s]$, $s\ge0$, generate the Borel $\sigma$-field of $[0,\infty]$ (indeed $\{\infty\}=[0,\infty]\setminus\bigcup_{n}[0,n]$ and $(a,b]=[0,b]\setminus[0,a]$), so by Step 1 it suffices to show $\{T\le s\}\in\F_{T+}$ for every $s\ge0$. Fix such an $s$. First, $\{T\le s\}\in\F_{s+}\subseteq\F$ by \eqref{eq:1-2-21:b}. Next, for every $t\ge0$,
\[
\{T\le s\}\cap\{T\le t\}=\{T\le s\wedge t\}\in\F_{(s\wedge t)+}\subseteq\F_{t+},
\]
again by \eqref{eq:1-2-21:b} and because $u\mapsto\F_{u+}$ is nondecreasing. Hence $\{T\le s\}\in\F_{T+}$.

\emph{Step 3: the alternative description.} Let $\G\defeq\{A\in\F;\ A\cap\{T<t\}\in\F_t,\ \forall\,t\ge0\}$.

Let $A\in\F_{T+}$ and $t\ge0$. If $t=0$, then $A\cap\{T<0\}=\emptyset\in\F_0$. If $t>0$, choose an integer $m$ with $1/m<t$; then
\[
A\cap\{T<t\}=\bigcup_{n\ge m}\Bigl(A\cap\Bigl\{T\le t-\frac1n\Bigr\}\Bigr),
\]
and each set of the union belongs to $\F_{(t-(1/n))+}\subseteq\F_t$ by the definition of $\F_{T+}$ and \eqref{eq:1-2-21:a}. Hence $A\cap\{T<t\}\in\F_t$, and $A\in\G$.

Conversely, let $A\in\G$ and $t\ge0$. For every integer $m\ge1$,
\[
A\cap\{T\le t\}=\bigcap_{n\ge m}\Bigl(A\cap\Bigl\{T<t+\frac1n\Bigr\}\Bigr),
\]
and for $n\ge m$ the $n$-th set of the intersection belongs to $\F_{t+(1/n)}\subseteq\F_{t+(1/m)}$. Therefore $A\cap\{T\le t\}\in\F_{t+(1/m)}$ for every $m\ge1$, that is, $A\cap\{T\le t\}\in\F_{t+}$ by \eqref{eq:1-2-21:a}. Thus $A\in\F_{T+}$, and we have shown $\F_{T+}=\G$.

\emph{Step 4: $\F_T\subseteq\F_{T+}$ for a stopping time $T$.} A stopping time is optional by \bk{Proposition 2.3}, so $\F_{T+}$ is defined. If $A\in\F_T$, then $A\in\F$ and $A\cap\{T\le t\}\in\F_t\subseteq\F_{t+}$ for every $t\ge0$, by \bk{Definition 2.12} and \eqref{eq:1-2-21:a}; this is precisely the condition of \bk{Definition 2.20}, so $A\in\F_{T+}$.

\rmk{By \eqref{eq:1-2-21:b}, $\F_{T+}$ is nothing but the $\sigma$-field of \bk{Definition 2.12} for $T$ regarded as a stopping time of the filtration $\{\F_{t+}\}$; Steps 1 and 2 are therefore also a consequence of \bk{Problem 2.13} applied to that filtration.}
\end{kex}

\begin{kex}{Problem}{2.22}{}
The book's Section 1.6 treats only the second claim, and leaves two verifications in it to the reader; here we prove the first claim, and give the second one with all details.

\emph{First claim.} The statements to be proved are the following. Let $T$ and $S$ be optional times of $\{\F_t\}$.
\begin{enumerate}
\item[\textbf{(a)}] For every $A\in\F_{S+}$ we have $A\cap\{S\le T\}\in\F_{T+}$. In particular, if $S\le T$ on $\Omega$, then $\F_{S+}\subseteq\F_{T+}$.
\item[\textbf{(b)}] $T\wedge S$ is an optional time, $\F_{(T\wedge S)+}=\F_{T+}\cap\F_{S+}$, and each of the events
\[
\{T<S\},\quad \{S<T\},\quad \{T\le S\},\quad \{S\le T\},\quad \{T=S\}
\]
belongs to $\F_{T+}\cap\F_{S+}$.
\end{enumerate}
We set $\G_t\defeq\F_{t+}=\bigcap_{\eps>0}\F_{t+\eps}$ for $t\ge0$. Each $\G_t$ is a sub-$\sigma$-field of $\F$, and for $0\le s<t$ we have $\G_s\subseteq\F_t\subseteq\G_t$; thus $\{\G_t\}$ is a filtration on $(\Omega,\F)$. \bk{Definition 2.1}, \bk{Lemma 2.9}, \bk{Definition 2.12}, \bk{Lemma 2.15} and \bk{Lemma 2.16} are formulated for an arbitrary filtration, so we may apply them to $\{\G_t\}$.

By \bk{Corollary 2.4}, a random time $R$ is an optional time of $\{\F_t\}$ if and only if it is a stopping time of $\{\G_t\}$. For such an $R$, the $\sigma$-field of \bk{Definition 2.12} relative to the filtration $\{\G_t\}$, which we denote by $\G_R$, consists of the events $A\in\F$ with $A\cap\{R\le t\}\in\G_t=\F_{t+}$ for every $t\ge0$; comparing with \bk{Definition 2.20}, we see that
\begin{equation}\label{eq:1-2-22:a}
\G_R=\F_{R+}\qquad\text{for every optional time $R$ of $\{\F_t\}$.}
\end{equation}

Now $T$ and $S$ are stopping times of $\{\G_t\}$. \bk{Lemma 2.15} for the filtration $\{\G_t\}$ states that $A\cap\{S\le T\}\in\G_T$ for every $A\in\G_S$, and that $\G_S\subseteq\G_T$ if $S\le T$ on $\Omega$; in view of \eqref{eq:1-2-22:a} with $R=S$ and $R=T$, this is (a).

By \bk{Lemma 2.9} for the filtration $\{\G_t\}$, the random time $T\wedge S$ is a stopping time of $\{\G_t\}$, hence an optional time of $\{\F_t\}$ by \bk{Corollary 2.4}; in particular $\F_{(T\wedge S)+}$ is defined, and \eqref{eq:1-2-22:a} holds also for $R=T\wedge S$. \bk{Lemma 2.16} for the filtration $\{\G_t\}$ states that $\G_{T\wedge S}=\G_T\cap\G_S$ and that each of the five events displayed in (b) belongs to $\G_T\cap\G_S$. By \eqref{eq:1-2-22:a} with $R=T\wedge S$, $R=T$ and $R=S$, this is (b).

\emph{Second claim.} Let $S$ be an optional time and $T$ a stopping time of $\{\F_t\}$ with $S\le T$ on $\Omega$ and $S<T$ on $\{S<\infty\}$, and let $A\in\F_{S+}$. By \bk{Problem 2.21},
\begin{equation}\label{eq:1-2-22:b}
A\in\F\qquad\text{and}\qquad A\cap\{S<r\}\in\F_r\quad\text{for every } r\ge0 .
\end{equation}
We fix $t\ge0$ and show that
\begin{equation}\label{eq:1-2-22:c}
A\cap\{T\le t\}=\bigcup_{r\in\Q,\ 0\le r<t}\Bigl(\bigl[A\cap\{S<r\}\bigr]\cap\{r<T\le t\}\Bigr),
\end{equation}
where the union is empty if $t=0$. The right-hand side is clearly contained in the left-hand side. Conversely, let $\omega\in A$ with $T(\omega)\le t$. Then $S(\omega)\le T(\omega)\le t<\infty$, so that $\omega\in\{S<\infty\}$ and therefore $S(\omega)<T(\omega)$. We choose a rational number $r$ with $S(\omega)<r<T(\omega)$; then $0\le S(\omega)<r<T(\omega)\le t$, and $\omega$ belongs to the set with index $r$ on the right-hand side of \eqref{eq:1-2-22:c}.

For each rational $r\in[0,t)$ we have $A\cap\{S<r\}\in\F_r\subseteq\F_t$ by \eqref{eq:1-2-22:b}, and
\[
\{r<T\le t\}=\{T\le t\}\setminus\{T\le r\}\in\F_t,
\]
because $T$ is a stopping time and $\F_r\subseteq\F_t$. The union in \eqref{eq:1-2-22:c} is countable, hence $A\cap\{T\le t\}\in\F_t$. Since $t\ge0$ was arbitrary and $A\in\F$, we conclude from \bk{Definition 2.12} that $A\in\F_T$. This proves $\F_{S+}\subseteq\F_T$.

\rmk{The positivity of $T$ was not used, and indeed it follows from the other hypotheses: on $\{S<\infty\}$ we have $T>S\ge0$, and on $\{S=\infty\}$ we have $T\ge S=\infty$. The second claim sharpens (a), which under $S\le T$ gives only $\F_{S+}\subseteq\F_{T+}$; by \bk{Problem 2.21}, $\F_T\subseteq\F_{T+}$. The strict inequality cannot be dropped: for $S=T$ a stopping time, $\F_{S+}\subseteq\F_T$ would mean $\F_{T+}=\F_T$, which fails for the constant time $T\equiv t$ whenever $\F_{t+}\ne\F_t$, since then $\F_T=\F_t$ by \bk{Problem 2.13} and $\F_{T+}=\F_{t+}$ by \bk{Definition 2.20}.}
\end{kex}

\begin{kex}{Problem}{2.24}{}
For fixed $n\ge1$ the sets $\{T=\infty\}$ and $\{(k-1)/2^n\le T<k/2^n\}$, $k\ge1$, form a partition of $\Omega$, so $T_n$ is well defined, with values in $\{k/2^n;\ k\ge1\}\cup\{\infty\}$. Since $T$ is optional, $\{T<t\}\in\F_t$ for every $t\ge0$ (\bk{Definition 2.1}; for $t=0$ this set is empty), and in particular $T$ and $T_n$ are $\F$-measurable. By the definition of $T_n$,
\begin{equation}\label{eq:1-2-24:a}
\Bigl\{T_n=\frac{k}{2^n}\Bigr\}=\Bigl\{\frac{k-1}{2^n}\le T<\frac{k}{2^n}\Bigr\}
=\Bigl\{T<\frac{k}{2^n}\Bigr\}\setminus\Bigl\{T<\frac{k-1}{2^n}\Bigr\},
\qquad n,k\ge1 .
\end{equation}

\emph{Step 1: each $T_n$ is a stopping time.} Fix $n\ge1$ and $t\ge0$, and let $j$ be the integer part of $2^nt$, so that $j/2^n\le t<(j+1)/2^n$. As $T_n$ takes finite values only in $\{k/2^n;\ k\ge1\}$, we obtain from \eqref{eq:1-2-24:a}
\[
\{T_n\le t\}=\bigcup_{k=1}^{j}\Bigl\{T_n=\frac{k}{2^n}\Bigr\}
=\bigcup_{k=1}^{j}\Bigl\{\frac{k-1}{2^n}\le T<\frac{k}{2^n}\Bigr\}
=\Bigl\{T<\frac{j}{2^n}\Bigr\}\in\F_{j/2^n}\subseteq\F_t ,
\]
the union being empty when $j=0$. Hence $T_n$ is a stopping time of $\{\F_t\}$.

\emph{Step 2: monotonicity and convergence.} On $\{T=\infty\}$ we have $T_n=T=\infty$ for every $n$. Let $\omega\in\{T<\infty\}$ and $n\ge1$, and let $k\ge1$ be the integer with $(k-1)/2^n\le T(\omega)<k/2^n$. Then $T_n(\omega)=k/2^n$, and therefore
\begin{equation}\label{eq:1-2-24:b}
T(\omega)<T_n(\omega)\le T(\omega)+\frac{1}{2^n}.
\end{equation}
Moreover $T(\omega)<2k/2^{n+1}$, so the integer $k'\ge1$ with $(k'-1)/2^{n+1}\le T(\omega)<k'/2^{n+1}$ satisfies $k'\le 2k$, that is, $T_{n+1}(\omega)=k'/2^{n+1}\le k/2^n=T_n(\omega)$. Thus $T_n\ge T_{n+1}\ge T$ on $\Omega$, and letting $n\to\infty$ in \eqref{eq:1-2-24:b} gives $\lim_{n\to\infty}T_n=T$ on $\{T<\infty\}$; on $\{T=\infty\}$ the convergence is trivial. Consequently $T_n\downarrow T$ everywhere on $\Omega$.

\emph{Step 3: the last assertion.} Let $A\in\F_{T+}$. By \bk{Problem 2.21} we have $A\cap\{T<t\}\in\F_t$ for every $t\ge0$. For $n,k\ge1$, \eqref{eq:1-2-24:a} yields
\[
A\cap\Bigl\{T_n=\frac{k}{2^n}\Bigr\}
=\Bigl(A\cap\Bigl\{T<\frac{k}{2^n}\Bigr\}\Bigr)\setminus\Bigl(A\cap\Bigl\{T<\frac{k-1}{2^n}\Bigr\}\Bigr),
\]
where the first set belongs to $\F_{k/2^n}$ and the second to $\F_{(k-1)/2^n}\subseteq\F_{k/2^n}$. Hence $A\cap\{T_n=k/2^n\}\in\F_{k/2^n}$.

\rmk{By \eqref{eq:1-2-24:b} the inequality $T_n>T$ is strict on $\{T<\infty\}$, and each $T_n$ is a positive stopping time with countably many values. Step 3 says moreover that $\F_{T+}\subseteq\F_{T_n}$: for $A\in\F_{T+}$ and $t\ge0$, with $j$ the integer part of $2^nt$, the set $A\cap\{T_n\le t\}=\bigcup_{k=1}^{j}\bigl(A\cap\{T_n=k/2^n\}\bigr)$ belongs to $\F_{j/2^n}\subseteq\F_t$.}
\end{kex}

\solsection{1.3}{Continuous-Time Martingales}

\begin{kex}{Problem}{3.4}{}
Let $N=\{N_t,\F_t;\ 0\le t<\infty\}$ be a Poisson process with intensity $\lambda>0$ in the sense of \bk{Definition 3.3}, and $M_t=N_t-\lambda t$. According to \bk{Definition 3.1} and the standing conventions of the book's Section 1.3, we have to check that $M$ is adapted to $\{\F_t\}$, that $E\abs{M_t}<\infty$ for every $t\ge 0$, and that $E(M_t\,|\,\F_s)=M_s$ a.s.\ for $0\le s<t<\infty$.

\emph{Adaptedness.} For each $t\ge 0$ the random variable $N_t$ is $\F_t$-measurable by \bk{Definition 3.3}, and $\lambda t$ is a constant; hence $M_t$ is $\F_t$-measurable.

\emph{Integrability.} We have $N_0=0$ a.s., so $M_0=0$ a.s.\ is integrable. For $t>0$, \bk{Definition 3.3} with $s=0$ shows that $N_t-N_0$ is Poisson distributed with mean $\lambda t$; in particular it is nonnegative and
\[
  E\abs{N_t-N_0}\;=\;\sum_{k=0}^\infty k\,e^{-\lambda t}\frac{(\lambda t)^k}{k!}
  \;=\;\lambda t\,e^{-\lambda t}\sum_{k=1}^\infty\frac{(\lambda t)^{k-1}}{(k-1)!}\;=\;\lambda t\;<\;\infty .
\]
Since $N_t=N_t-N_0$ a.s., we obtain $E\abs{N_t}=\lambda t$ and $E\abs{M_t}\le E\abs{N_t}+\lambda t=2\lambda t<\infty$.

\emph{The martingale property.} Fix $0\le s<t<\infty$. By \bk{Definition 3.3} the increment $N_t-N_s$ is independent of $\F_s$ and is Poisson distributed with mean $\lambda(t-s)$; by the computation just made (with $\lambda(t-s)$ in place of $\lambda t$) it is integrable with $E(N_t-N_s)=\lambda(t-s)$. The conditional expectation of an integrable random variable given a $\sigma$-field of which it is independent is a.s.\ equal to its expectation: indeed, for every $A\in\F_s$,
\[
  \int_A (N_t-N_s)\,dP\;=\;E\big[\ind_A\,(N_t-N_s)\big]\;=\;P(A)\,E(N_t-N_s)\;=\;\int_A \lambda(t-s)\,dP ,
\]
and the constant $\lambda(t-s)$ is $\F_s$-measurable. Therefore
\[
  E(N_t-N_s\,|\,\F_s)\;=\;\lambda(t-s)\quad\text{a.s.}
\]
Since $N_s$ is integrable and $\F_s$-measurable, $E(N_s\,|\,\F_s)=N_s$ a.s., and linearity of conditional expectation gives, a.s.,
\[
  E(M_t\,|\,\F_s)\;=\;E(N_t-N_s\,|\,\F_s)+N_s-\lambda t\;=\;\lambda(t-s)+N_s-\lambda t\;=\;N_s-\lambda s\;=\;M_s .
\]
Thus $\{M_t,\F_t;\ 0\le t<\infty\}$ is both a submartingale and a supermartingale, i.e., a martingale.

\rmk{The argument uses only that $N$ is adapted, that $N_0$ is integrable, and that the increments $N_t-N_s$ are integrable, independent of $\F_s$, and have mean $\lambda(t-s)$; neither the precise form of the Poisson distribution nor the RCLL property of the paths enters. In particular the conclusion holds for the process $\{N_t,\F^N_t\}$ of \bk{(3.1)}, which is Poisson by \bk{Problem 3.2}. Since $N_t-N_s\ge 0$ a.s.\ and hence $E(N_t\,|\,\F_s)=N_s+\lambda(t-s)\ge N_s$, the same computation shows that $N$ itself is a submartingale, as asserted in \bk{Remark 3.5}.}
\end{kex}

\begin{kex}{Problem}{3.9}{}
Let $N=\{N_t,\F_t;\ 0\le t<\infty\}$ be as in \bk{Definition 3.3}, and let $M_t\defeq N_t-\lambda t$ be the compensated Poisson process. By \bk{Problem 3.4}, $\{M_t,\F_t;\ 0\le t<\infty\}$ is a martingale, and every path of $M$ is right-continuous, because $N$ is an RCLL process. Thus \bk{Theorem 3.8} applies to $M$, and to every submartingale of the form $\varphi(M)$ with $\varphi$ continuous. Since $N_0=0$ a.s.\ and $N_t-N_0$ is Poisson distributed with mean $\lambda t$, so is $N_t$ for every $t>0$; in particular
\begin{equation}\label{eq:1-3-09:mom}
E(M_0)=0,\qquad E(M_t^2)=\Var(N_t)=\lambda t,\qquad 0\le t<\infty .
\end{equation}
The suprema and infima below are random variables: by right-continuity they are not changed if $s$ is restricted to the rational points of the interval together with its right endpoint.

\emph{Step 1: the hint.} Let $Z$ be a Poisson random variable with mean $\mu>0$, and let $m\defeq\lfloor\mu\rfloor$ be the integer part of $\mu$. We claim that
\begin{equation}\label{eq:1-3-09:exact}
E(Z-\mu)^+=e^{-\mu}\,\frac{\mu^{m+1}}{m!}.
\end{equation}
Indeed, an integer $k$ satisfies $k>\mu$ if and only if $k\ge m+1$, and for $k\ge 1$ we have
\[
(k-\mu)\frac{\mu^k}{k!}=\mu\Bigl(\frac{\mu^{k-1}}{(k-1)!}-\frac{\mu^k}{k!}\Bigr).
\]
Therefore the series for $E(Z-\mu)^+$ telescopes, and since $\mu^K/K!\to0$ as $K\to\infty$,
\[
E(Z-\mu)^+=e^{-\mu}\sum_{k=m+1}^\infty (k-\mu)\frac{\mu^k}{k!}
=\mu e^{-\mu}\lim_{K\to\infty}\Bigl(\frac{\mu^{m}}{m!}-\frac{\mu^{K}}{K!}\Bigr)
=e^{-\mu}\,\frac{\mu^{m+1}}{m!}.
\]
Now let $\mu\ge 1$, so that $m\ge1$, and write $\mu=m+\theta$ with $0\le\theta<1$. From \eqref{eq:1-3-09:exact},
\begin{equation}\label{eq:1-3-09:fact}
\frac{E(Z-\mu)^+}{\sqrt{\mu}}=e^{-\mu}\frac{\mu^{m+\frac12}}{m!}
=\frac{1}{\sqrt{2\pi}}\cdot\frac{\sqrt{2\pi m}\,m^me^{-m}}{m!}\cdot\Bigl(1+\frac{\theta}{m}\Bigr)^{m+\frac12}e^{-\theta}.
\end{equation}
By Stirling's approximation, the second factor on the right converges to $1$ as $m\to\infty$. For the third factor we use $x-\tfrac12x^2\le\log(1+x)\le x$, valid for $x\ge0$: its logarithm
$\ell\defeq\bigl(m+\tfrac12\bigr)\log\bigl(1+\tfrac{\theta}{m}\bigr)-\theta$ satisfies
\[
\ell\le\Bigl(m+\frac12\Bigr)\frac{\theta}{m}-\theta=\frac{\theta}{2m}\le\frac{1}{2m},
\qquad
\ell\ge\frac{\theta}{2m}-\Bigl(m+\frac12\Bigr)\frac{\theta^2}{2m^2}\ge-\frac{2m+1}{4m^2}\ge-\frac{1}{m},
\]
so that $\abs{\ell}\le 1/m$ whatever the value of $\theta\in[0,1)$. As $\mu\to\infty$ we have $m=\lfloor\mu\rfloor\to\infty$, and \eqref{eq:1-3-09:fact} yields $\mu^{-1/2}E(Z-\mu)^+\to 1/\sqrt{2\pi}$. Applying this with $Z=N_t$, $\mu=\lambda t$, we obtain
\begin{equation}\label{eq:1-3-09:lim}
\lim_{t\to\infty}\frac{E(M_t^+)}{\sqrt{\lambda t}}=\lim_{t\to\infty}\frac{E(N_t-\lambda t)^+}{\sqrt{\lambda t}}=\frac{1}{\sqrt{2\pi}} .
\end{equation}

\textbf{(a)} Fix $c>0$ and $t>0$. The martingale $M$ is in particular a submartingale with right-continuous paths, so the first submartingale inequality, \bk{Theorem 3.8}\,(i), on the interval $[0,t]$ and at the level $c\sqrt{\lambda t}>0$ gives
\[
P\biggl[\sup_{0\le s\le t}(N_s-\lambda s)\ge c\sqrt{\lambda t}\biggr]\le\frac{E(M_t^+)}{c\sqrt{\lambda t}} .
\]
Letting $t\to\infty$ and using \eqref{eq:1-3-09:lim}, we see that the upper limit of the left-hand side is at most $1/(c\sqrt{2\pi})$.

\textbf{(b)} Likewise, the second submartingale inequality, \bk{Theorem 3.8}\,(ii), on $[0,t]$ at the level $c\sqrt{\lambda t}$, together with $E(M_0)=0$ from \eqref{eq:1-3-09:mom}, gives
\[
P\biggl[\inf_{0\le s\le t}(N_s-\lambda s)\le -c\sqrt{\lambda t}\biggr]\le\frac{E(M_t^+)-E(M_0)}{c\sqrt{\lambda t}}=\frac{E(M_t^+)}{c\sqrt{\lambda t}},
\]
and the claim follows again from \eqref{eq:1-3-09:lim}.

\textbf{(c)} Since $E\abs{M_t}<\infty$ for every $t\ge0$ and $x\mapsto\abs{x}$ is convex, \bk{Proposition 3.6} shows that $\{\abs{M_t},\F_t;\ 0\le t<\infty\}$ is a submartingale; it is nonnegative, its paths are right-continuous, and $E(\abs{M_\tau}^2)=\lambda\tau<\infty$ by \eqref{eq:1-3-09:mom}. Doob's maximal inequality, \bk{Theorem 3.8}\,(iv), with $p=2$ on the interval $[\sigma,\tau]$ therefore gives
\[
E\Bigl(\sup_{\sigma\le t\le\tau}\abs{M_t}\Bigr)^2\le 4\,E(M_\tau^2)=4\lambda\tau .
\]
For $\sigma\le t\le\tau$ we have $\bigl(\frac{N_t}{t}-\lambda\bigr)^2=M_t^2/t^2\le M_t^2/\sigma^2$, because $t\ge\sigma>0$; hence
\[
\sup_{\sigma\le t\le\tau}\Bigl(\frac{N_t}{t}-\lambda\Bigr)^2\le\frac{1}{\sigma^2}\sup_{\sigma\le t\le\tau}M_t^2=\frac{1}{\sigma^2}\Bigl(\sup_{\sigma\le t\le\tau}\abs{M_t}\Bigr)^2 ,
\]
and taking expectations we conclude that
\[
E\biggl[\sup_{\sigma\le t\le\tau}\Bigl(\frac{N_t}{t}-\lambda\Bigr)^2\biggr]\le\frac{4\tau\lambda}{\sigma^2} .
\]

\rmk{The limit \eqref{eq:1-3-09:lim} is the one suggested by the central limit theorem: $(N_t-\lambda t)/\sqrt{\lambda t}$ converges in distribution to a standard normal variable $G$, and $E(G^+)=1/\sqrt{2\pi}$. The exact formula \eqref{eq:1-3-09:exact} makes this argument unnecessary.}
\end{kex}

\begin{kex}{Problem}{3.16}{}
According to the convention stated before \bk{Proposition 3.6} (in its supermartingale form), the second assertion means that $X_\infty$ is an integrable, $\F_\infty$-measurable random variable, where $\F_\infty=\sigma\bigl(\bigcup_{t\ge0}\F_t\bigr)$, and that
\begin{equation}\label{eq:1-3-16:a}
E(X_\infty\,|\,\F_t)\le X_t\quad\text{a.s.\ $P$, for every } 0\le t<\infty .
\end{equation}
Since the limit in the statement is only asserted to exist almost surely, we must also say how $X_\infty$ is defined on the exceptional set; we do this in Step~2 in such a way that $X_\infty$ is $\F_\infty$-measurable without any completeness assumption on the filtration.

\emph{Step 1: almost sure convergence.} The process $Y\defeq-X$ is adapted to $\{\F_t\}$, has right-continuous paths, satisfies $E\abs{Y_t}=E X_t<\infty$, and $E(Y_t\,|\,\F_s)=-E(X_t\,|\,\F_s)\ge -X_s=Y_s$ a.s.\ for $0\le s<t<\infty$; thus $\{Y_t,\F_t;\ 0\le t<\infty\}$ is a right-continuous submartingale. Because $X_t\ge0$ we have $Y_t^+=0$, hence $\sup_{t\ge0}E(Y_t^+)=0<\infty$. \bk{Theorem 3.15} now provides an event $\Omega_0$ with $P(\Omega_0)=1$ such that, for every $\omega\in\Omega_0$, the limit $\lim_{t\to\infty}Y_t(\omega)$ exists, and this limit is an integrable random variable, in particular finite a.s.; discarding a further null set, we may assume that it is finite for every $\omega\in\Omega_0$. Consequently
\begin{equation}\label{eq:1-3-16:b}
\lim_{t\to\infty}X_t(\omega)\ \text{exists in } [0,\infty)\quad\text{for every }\omega\in\Omega_0,\qquad P(\Omega_0)=1 .
\end{equation}

\emph{Step 2: definition and measurability of $X_\infty$.} Let $L\defeq\liminf_{n\to\infty}X_n$, the lower limit being taken along the integers $n\in\N$. Each $X_n$ is $\F_n$-measurable, hence $\F_\infty$-measurable, so $L$ is an $\F_\infty$-measurable random variable with values in $[0,\infty]$, and we set
\[
X_\infty\defeq L\,\ind_{\{L<\infty\}} .
\]
This is an $\F_\infty$-measurable random variable with values in $[0,\infty)$. For $\omega\in\Omega_0$ the limit in \eqref{eq:1-3-16:b} coincides with the limit along the integers, so that $L(\omega)<\infty$ and
\[
X_\infty(\omega)=\lim_{t\to\infty}X_t(\omega),\qquad\omega\in\Omega_0 .
\]
This is the first assertion of the problem.

\emph{Step 3: the supermartingale property at $t=\infty$.} Fix $0\le t<\infty$ and $A\in\F_t$. For every integer $n\ge t$ the supermartingale property gives $E(X_n\,|\,\F_t)\le X_t$ a.s.\ (with equality if $n=t$), and integrating over $A\in\F_t$ we get
\[
\int_A X_n\,dP\le\int_A X_t\,dP,\qquad n\in\N,\ n\ge t .
\]
The random variables $X_n\ind_A$ are nonnegative and converge to $X_\infty\ind_A$ on $\Omega_0$, that is, almost surely. Fatou's lemma therefore yields
\begin{equation}\label{eq:1-3-16:c}
\int_A X_\infty\,dP\le\liminf_{n\to\infty}\int_A X_n\,dP\le\int_A X_t\,dP,\qquad A\in\F_t,\ 0\le t<\infty .
\end{equation}
With $t=0$ and $A=\Omega$ this reads $E X_\infty\le E X_0<\infty$; as $X_\infty\ge0$, the random variable $X_\infty$ is integrable, and the conditional expectation $E(X_\infty\,|\,\F_t)$ is well defined. Now fix $t$ and take in \eqref{eq:1-3-16:c} the event $A\defeq\{E(X_\infty\,|\,\F_t)>X_t\}$, which belongs to $\F_t$. By the definition of conditional expectation, $\int_A X_\infty\,dP=\int_A E(X_\infty\,|\,\F_t)\,dP$, and \eqref{eq:1-3-16:c} becomes
\[
\int_A\bigl(E(X_\infty\,|\,\F_t)-X_t\bigr)\,dP\le0 ,
\]
both terms of the integrand being integrable. The integrand is strictly positive on $A$, so $P(A)=0$, which is \eqref{eq:1-3-16:a}.

We have shown that $X_\infty$ is integrable and $\F_\infty$-measurable and that \eqref{eq:1-3-16:a} holds. Together with the supermartingale property of $\{X_t,\F_t;\ 0\le t<\infty\}$, this says that $\{X_t,\F_t;\ 0\le t\le\infty\}$ is a supermartingale with last element $X_\infty$.

\rmk{Right-continuity of the paths was used only through \bk{Theorem 3.15}, to pass from convergence along a countable set of times to convergence as $t\to\infty$ through all of $[0,\infty)$. No uniform integrability is needed, and none holds in general: the inequality in \eqref{eq:1-3-16:a} may be strict even when $X$ is a martingale, since a nonnegative martingale may converge to zero almost surely while $E X_t$ stays constant and positive.}
\end{kex}

\begin{kex}{Exercise}{3.18}{}
\emph{Step 1: the limit $X_\infty$.} A uniformly integrable family is bounded in $L^1$, so $\sup_{t\ge0}E\abs{X_t}<\infty$. The process $-X$ is a right-continuous submartingale with $\sup_{t\ge0}E\bigl((-X_t)^+\bigr)\le\sup_{t\ge0}E\abs{X_t}<\infty$, and \bk{Theorem 3.15} provides an event $\Omega_0\in\F$ with $P(\Omega_0)=1$ on which $\lim_{t\to\infty}X_t(\omega)$ exists in $\R$. We set $X_\infty\defeq\lim_{n\to\infty}X_n$ on $\Omega_0$ and $X_\infty\defeq0$ on $\Omega\setminus\Omega_0$; this is an $\F$-measurable random variable, $X_t\to X_\infty$ a.s.\ as $t\to\infty$, and $E\abs{X_\infty}<\infty$ by \bk{Theorem 3.15}. For every sequence $t_n\uparrow\infty$ the sequence $\{X_{t_n}\}_{n=1}^\infty$ is uniformly integrable and converges a.s.\ to $X_\infty$, hence converges in $L^1$ (Vitali's convergence theorem). Therefore
\begin{equation}\label{eq:1-3-18:a}
\lim_{t\to\infty}E\abs{X_t-X_\infty}=0 .
\end{equation}
Let $0\le s<\infty$ and $A\in\F_s$. The supermartingale property gives $\int_A X_s\,dP\ge\int_A X_t\,dP$ for all $t\ge s$, and letting $t\to\infty$ with the help of \eqref{eq:1-3-18:a} we obtain $\int_A X_s\,dP\ge\int_A X_\infty\,dP=\int_A E[X_\infty\mid\F_s]\,dP$. Since $X_s$ and $E[X_\infty\mid\F_s]$ are both $\F_s$-measurable and $A\in\F_s$ is arbitrary, this means
\begin{equation}\label{eq:1-3-18:b}
X_s\ \ge\ E[X_\infty\mid\F_s]\qquad\text{a.s. } P,\ \text{ for every } 0\le s<\infty .
\end{equation}

\emph{Step 2: the martingale part.} For each $t\ge0$ let $N_t$ be a version of $E[X_\infty\mid\F_t]$. By the tower property $\{N_t,\F_t;\ 0\le t<\infty\}$ is a martingale, in particular a submartingale, and $t\mapsto EN_t=EX_\infty$ is constant, hence right-continuous. As $\{\F_t\}$ satisfies the usual conditions, \bk{Theorem 3.13} yields a right-continuous modification $\wt M$ of $N$ which is adapted to $\{\F_t\}$. For $0\le s<t$ we have $\wt M_t=N_t$ and $\wt M_s=N_s$ a.s., so $E[\wt M_t\mid\F_s]=E[N_t\mid\F_s]=N_s=\wt M_s$ a.s.; thus $\{\wt M_t,\F_t\}$ is a right-continuous martingale with
\begin{equation}\label{eq:1-3-18:c}
\wt M_t=E[X_\infty\mid\F_t]\qquad\text{a.s. } P,\ \text{ for every } 0\le t<\infty .
\end{equation}

The process $X-\wt M$ is right-continuous, and by \eqref{eq:1-3-18:b}, \eqref{eq:1-3-18:c} the event
\[
\Lambda\defeq\bigcup_{q\in\Q\cap[0,\infty)}\{X_q-\wt M_q<0\}
\]
has probability zero. If $\omega\notin\Lambda$ and $t\ge0$, we choose rationals $q_n\downarrow t$ and get $X_t(\omega)-\wt M_t(\omega)=\lim_n\bigl(X_{q_n}(\omega)-\wt M_{q_n}(\omega)\bigr)\ge0$. Because $\F_0$ contains all $P$-negligible events, $\Lambda\in\F_0\subseteq\F_t$ for every $t$. We define
\[
M_t\defeq\wt M_t\,\ind_{\Omega\setminus\Lambda}+X_t\,\ind_{\Lambda},\qquad
Z_t\defeq X_t-M_t=(X_t-\wt M_t)\,\ind_{\Omega\setminus\Lambda},\qquad 0\le t<\infty .
\]
Then $M$ is adapted to $\{\F_t\}$, every path of $M$ is right-continuous, and $M_t=\wt M_t$ for all $t\ge0$ on the event $\Omega\setminus\Lambda$ of probability one. Consequently $M_t$ is integrable, $E[M_t\mid\F_s]=E[\wt M_t\mid\F_s]=\wt M_s=M_s$ a.s.\ for $s<t$, and \eqref{eq:1-3-18:c} holds with $M$ in place of $\wt M$: $\{M_t,\F_t;\ 0\le t<\infty\}$ is a right-continuous martingale.

It is uniformly integrable. Indeed, for $t\ge0$ and $\lambda>0$ the event $\{\abs{M_t}>\lambda\}$ belongs to $\F_t$ and $\abs{M_t}\le E[\abs{X_\infty}\mid\F_t]$ a.s.\ (Jensen's inequality), so that
\[
\int_{\{\abs{M_t}>\lambda\}}\abs{M_t}\,dP\le\int_{\{\abs{M_t}>\lambda\}}\abs{X_\infty}\,dP,
\qquad
P[\abs{M_t}>\lambda]\le\frac{E\abs{M_t}}{\lambda}\le\frac{E\abs{X_\infty}}{\lambda}.
\]
Given $\eps>0$, the integrability of $X_\infty$ provides $\delta>0$ such that $\int_B\abs{X_\infty}\,dP<\eps$ whenever $P(B)<\delta$; for $\lambda>E\abs{X_\infty}/\delta$ we then have $\sup_{t\ge0}\int_{\{\abs{M_t}>\lambda\}}\abs{M_t}\,dP\le\eps$.

\emph{Step 3: the potential.} By construction $X_t(\omega)=M_t(\omega)+Z_t(\omega)$ for \emph{all} $t\ge0$ and $\omega\in\Omega$. The process $Z$ is adapted, has right-continuous paths (as the difference of two such processes), and is nonnegative: $Z_t(\omega)=0$ for $\omega\in\Lambda$, while $Z_t(\omega)=X_t(\omega)-\wt M_t(\omega)\ge0$ for $\omega\notin\Lambda$, as shown in Step 2. Each $Z_t$ is integrable, and for $0\le s<t$,
\[
E[Z_t\mid\F_s]=E[X_t\mid\F_s]-E[M_t\mid\F_s]\le X_s-M_s=Z_s\qquad\text{a.s. } P,
\]
so $\{Z_t,\F_t;\ 0\le t<\infty\}$ is a right-continuous, nonnegative supermartingale. Finally, $EM_t=EX_\infty$ by \eqref{eq:1-3-18:c}, whence by \eqref{eq:1-3-18:a}
\[
0\le EZ_t=EX_t-EX_\infty\le E\abs{X_t-X_\infty}\xrightarrow[t\to\infty]{}0 .
\]
Thus $Z$ is a potential in the sense of \bk{Definition 3.17}, and $X=M+Z$ is the desired decomposition.

\rmk{The decomposition is unique up to indistinguishability. Suppose $X_t=M_t+Z_t=M_t'+Z_t'$ a.s.\ for every $t$, with $M,M'$ right-continuous martingales and $Z,Z'$ potentials. Then $D\defeq M-M'$ is a right-continuous martingale with $D_t=Z_t'-Z_t$ a.s., so $E\abs{D_t}\le EZ_t'+EZ_t\to0$ as $t\to\infty$. Since $\abs{D}$ is a submartingale (Jensen's inequality), $E\abs{D_s}\le E\abs{D_t}$ for $s\le t$; letting $t\to\infty$ gives $E\abs{D_s}=0$, i.e.\ $M_s=M_s'$ a.s.\ for every $s$. By right-continuity, $M$ and $M'$ agree for all $t$ outside the null event $\bigcup_{q\in\Q\cap[0,\infty)}\{M_q\ne M_q'\}$, and then so do $Z$ and $Z'$ outside a null event. In particular the uniform integrability of $M$ is automatic.}
\end{kex}

\begin{kex}{Problem}{3.21}{}
Here $N$ is a Poisson process with intensity $\lambda$ in the sense of \bk{Definition 3.3}: it is adapted to $\{\F_t\}$, integer-valued, with RCLL paths, $N_0=0$ a.s., and for $0\le s<t$ the increment $N_t-N_s$ is independent of $\F_s$ and Poisson distributed with mean $\lambda(t-s)$. We write $\theta\defeq e^{iu}\in\mathbb{C}\setminus\{0\}$, so that, $N_t$ being an integer,
\[
X_t=\theta^{N_t}\,e^{-\lambda t(\theta-1)},\qquad 0\le t<\infty .
\]
We shall use the following elementary computation. If $K$ is a Poisson random variable with mean $\mu\ge 0$ and $z\in\mathbb{C}$, then $E\abs{z}^{K}=\sum_{k\ge0}e^{-\mu}(\mu\abs{z})^k/k!=e^{\mu(\abs{z}-1)}<\infty$, so $z^K$ is integrable and, the series being absolutely convergent,
\begin{equation}\label{eq:1-3-21:pgf}
E\bigl[z^{K}\bigr]=\sum_{k=0}^\infty e^{-\mu}\frac{(\mu z)^k}{k!}=e^{\mu(z-1)} .
\end{equation}

\textbf{(i)} \emph{Adaptedness and integrability.} For each $t$, $X_t$ is a continuous function of the $\F_t$-measurable random variable $N_t$, so $\mathrm{Re}(X_t)$ and $\mathrm{Im}(X_t)$ are $\F_t$-measurable. Since $N_0=0$ a.s., $N_t=N_t-N_0$ a.s.\ is Poisson with mean $\lambda t$ (for $t=0$ this reads $N_0=0$ a.s.), and therefore
\[
E\abs{X_t}=e^{-\lambda t\,\mathrm{Re}(\theta-1)}\,E\abs{\theta}^{N_t}
=e^{-\lambda t\,\mathrm{Re}(\theta-1)}\,e^{\lambda t(\abs{\theta}-1)}<\infty .
\]
Hence $\mathrm{Re}(X_t)$ and $\mathrm{Im}(X_t)$, which are bounded in absolute value by $\abs{X_t}$, are integrable.

\emph{The martingale property.} Fix $0\le s<t$ and $A\in\F_s$, and set
\[
Z\defeq\theta^{N_t-N_s},\qquad V\defeq \ind_A X_s ,
\]
so that $\ind_AX_t=e^{-\lambda(t-s)(\theta-1)}\,V Z$. The complex random variable $V$ is $\F_s$-measurable and integrable; $Z$ is a function of $N_t-N_s$, hence independent of $\F_s$, and it is integrable with $E[Z]=e^{\lambda(t-s)(\theta-1)}$ by \eqref{eq:1-3-21:pgf}. Writing $V=V_1+iV_2$ and $Z=Z_1+iZ_2$ with real $V_j$, $Z_k$, each pair $(V_j,Z_k)$ consists of independent integrable real random variables, so $V_jZ_k$ is integrable and $E[V_jZ_k]=E[V_j]\,E[Z_k]$; expanding the product $VZ$ and using the linearity of the expectation we obtain $E[VZ]=E[V]\,E[Z]$. Consequently
\begin{equation}\label{eq:1-3-21:mart}
E[\ind_AX_t]=e^{-\lambda(t-s)(\theta-1)}\,E[V]\,E[Z]=E[\ind_AX_s],\qquad A\in\F_s .
\end{equation}
Taking real and imaginary parts in \eqref{eq:1-3-21:mart} (the expectation of a complex integrable random variable being defined through its real and imaginary parts), we get
\[
E\bigl[\ind_A\,\mathrm{Re}(X_t)\bigr]=E\bigl[\ind_A\,\mathrm{Re}(X_s)\bigr],\qquad
E\bigl[\ind_A\,\mathrm{Im}(X_t)\bigr]=E\bigl[\ind_A\,\mathrm{Im}(X_s)\bigr]
\]
for every $A\in\F_s$. As $\mathrm{Re}(X_s)$ and $\mathrm{Im}(X_s)$ are $\F_s$-measurable and integrable, this says precisely that $E[\mathrm{Re}(X_t)\mid\F_s]=\mathrm{Re}(X_s)$ and $E[\mathrm{Im}(X_t)\mid\F_s]=\mathrm{Im}(X_s)$ a.s. Thus $\{\mathrm{Re}(X_t),\F_t\}$ and $\{\mathrm{Im}(X_t),\F_t\}$ are martingales in the sense of \bk{Definition 3.1}.

\textbf{(ii)} The answer is no. For $u=-i$ we have $iu=1$ and $\theta=e$, so
\[
X_t=\exp\bigl[N_t-\lambda t(e-1)\bigr],\qquad 0\le t<\infty,
\]
is real and strictly positive; by part (i) it is a martingale, with $E X_t=E X_0=1$ for all $t\ge0$ because $N_0=0$ a.s. Its paths are right-continuous, since those of $N$ are and $t\mapsto\lambda t(e-1)$ is continuous. So \bk{Problem 3.20} applies to $X$, and it suffices to show that one of its four equivalent conditions fails; we show that (c) fails, and then indicate directly why (a) and (b) fail as well.

\emph{Step 1: $E\sqrt{X_t}\to0$ as $t\to\infty$.} By \eqref{eq:1-3-21:pgf} with $z=e^{1/2}$ and $\mu=\lambda t$,
\[
E\sqrt{X_t}=e^{-\lambda t(e-1)/2}\,E\bigl[e^{N_t/2}\bigr]
=\exp\Bigl[\lambda t\Bigl(\sqrt e-1-\frac{e-1}{2}\Bigr)\Bigr]
=\exp\Bigl[-\frac{\lambda t}{2}\bigl(\sqrt e-1\bigr)^2\Bigr]\xrightarrow[t\to\infty]{}0 .
\]

\emph{Step 2: $\lim_{t\to\infty}X_t=0$ a.s.} Being a right-continuous, nonnegative martingale, $X$ is in particular a right-continuous, nonnegative supermartingale, so by \bk{Problem 3.16} the limit $X_\infty\defeq\lim_{t\to\infty}X_t$ exists a.s.\ (and is nonnegative). By Fatou's lemma along the integers and Step 1,
\[
E\sqrt{X_\infty}=E\Bigl[\lim_{n\to\infty}\sqrt{X_n}\Bigr]\le\liminf_{n\to\infty}E\sqrt{X_n}=0 ,
\]
whence $X_\infty=0$ a.s.

\emph{Step 3: the conditions of \bk{Problem 3.20} fail.} Suppose (c) held: $X_t$ converges a.s.\ to an integrable random variable $X'_\infty$ such that $\{X_t,\F_t;\ 0\le t\le\infty\}$ (with $X'_\infty$ as last element) is a martingale. By Step 2, $X'_\infty=0$ a.s., and the martingale property at $t=\infty$ would give $X_0=E(X'_\infty\mid\F_0)=0$ a.s., contradicting $X_0=1$ a.s. Hence (c) fails, and by the equivalence asserted in \bk{Problem 3.20}, so do (a), (b) and (d): the martingale $X$ is not uniformly integrable, it does not converge in $L^1$, and there is no integrable $Y$ with $X_t=E(Y\mid\F_t)$ a.s.\ for every $t\ge0$.

The failure of (a) and (b) can also be seen without appealing to that equivalence. If $X_t$ converged in $L^1$ to some $\xi$ as $t\to\infty$, then a subsequence of $\{X_n\}_{n\ge1}$ would converge a.s.\ to $\xi$, so $\xi=X_\infty=0$ a.s.\ by Step 2, and we would have $E X_t=E\abs{X_t-\xi}\to0$, contradicting $E X_t=1$; thus (b) fails. If the family $\{X_t\}_{t\ge0}$ were uniformly integrable, then the a.s.\ convergence $X_n\to0$ would imply $E X_n\to0$ (Vitali's convergence theorem), again contradicting $E X_n=1$; thus (a) fails.

\rmk{The same argument shows what happens for a general real exponent: for $u=-i\alpha$ with $\alpha\in\R\setminus\{0\}$, the positive martingale $X_t=\exp[\alpha N_t-\lambda t(e^{\alpha}-1)]$ satisfies $E\sqrt{X_t}=\exp[-\tfrac{\lambda t}{2}(e^{\alpha/2}-1)^2]\to0$, hence $X_t\to0$ a.s.\ while $E X_t=1$, and none of the conditions of \bk{Problem 3.20} holds. For real $u$, on the other hand, $\abs{X_t}=e^{\lambda t(1-\cos u)}$ is deterministic, and is bounded in $t$ only in the trivial case $e^{iu}=1$, $X\equiv1$.}
\end{kex}

\begin{kex}{Problem}{3.23}{}
The conclusion of \bk{Theorem 3.22} to be established is: $X_S$ and $X_T$ are integrable, and
\begin{equation}\label{eq:1-3-23:os}
E(X_T\mid\F_{S+})\ge X_S\quad\text{a.s. } P,
\end{equation}
where $\F_{S+}$ may be replaced by $\F_S$ if $S$ is a stopping time; in particular $E X_T\ge E X_0$. In case (i) the random variables $X_S$, $X_T$ are defined on all of $\Omega$, because $S\le T\le a<\infty$. In case (ii) the times $S$, $T$ may take the value $+\infty$, and part of the task is to say what $X_\infty$ is; we shall show that $X_\infty\defeq\lim_{t\to\infty}X_t$ exists a.s.\ and that \eqref{eq:1-3-23:os} holds with this $X_\infty$. In both cases we produce a right-continuous submartingale with a last element, in the sense of the convention preceding \bk{Proposition 3.6}, and apply \bk{Theorem 3.22} to it.

\textbf{(i)} Let $T\le a$, and define
\[
Z_t\defeq X_{t\wedge a}\quad(0\le t<\infty),\qquad Z_\infty\defeq X_a .
\]
Every path of $Z$ is right-continuous on $[0,\infty)$, each $Z_t$ is integrable, and $Z_t$ is measurable with respect to $\F_{t\wedge a}\subseteq\F_t$; moreover $Z_\infty$ is integrable and $\F_a$-measurable, hence $\F_\infty$-measurable. Let $0\le s<t\le\infty$, with the understanding $t\wedge a=a$ for $t=\infty$. If $s\ge a$, then $Z_t=X_a=Z_s$ and $Z_s$ is $\F_s$-measurable, so $E(Z_t\mid\F_s)=Z_s$ a.s. If $s<a$, then $s<t\wedge a$ and the submartingale property of $X$ gives
\[
E(Z_t\mid\F_s)=E(X_{t\wedge a}\mid\F_s)\ge X_s=Z_s\quad\text{a.s.}
\]
Thus $\{Z_t,\F_t;\ 0\le t\le\infty\}$ is a right-continuous submartingale with last element $Z_\infty$. Since $S\le T$ are optional times of $\{\F_t\}$, \bk{Theorem 3.22} shows that $Z_S$, $Z_T$ are integrable and that $E(Z_T\mid\F_{S+})\ge Z_S$ a.s., with $\F_S$ in place of $\F_{S+}$ if $S$ is a stopping time. But $S(\omega)\le T(\omega)\le a$ for every $\omega$, so that $Z_T=X_{T\wedge a}=X_T$ and $Z_S=X_{S\wedge a}=X_S$ identically, and \eqref{eq:1-3-23:os} is proved. Taking $S\equiv0$ and expectations gives $E X_T\ge E X_0$. If $X$ is a martingale, then $-X$ is also a right-continuous submartingale and $T\le a$, so \eqref{eq:1-3-23:os} holds for $X$ and for $-X$; hence $E(X_T\mid\F_{S+})=X_S$ a.s., and $E X_T=E X_0$.

\textbf{(ii)} \emph{Step 1: the limit $X_\infty$.} For every $t\ge0$ we have, a.s., $X_t\le E(Y\mid\F_t)\le E(Y^+\mid\F_t)$ and $0\le E(Y^+\mid\F_t)$, so that
\begin{equation}\label{eq:1-3-23:plus}
0\le X_t^+\le E(Y^+\mid\F_t)\quad\text{a.s.},\qquad\text{and}\qquad \sup_{t\ge0}E(X_t^+)\le E(Y^+)<\infty .
\end{equation}
By \bk{Theorem 3.15}, there is an event $\Omega^*\in\F$ with $P(\Omega^*)=1$ on which $\lim_{t\to\infty}X_t$ exists in $\R$. In order to have a random variable defined everywhere and measurable with respect to $\F_\infty$, we set
\[
X_\infty(\omega)\defeq\begin{cases}\lim_{n\to\infty}X_n(\omega), & \text{if this limit along the integers exists in }\R,\\[0.3ex] 0, & \text{otherwise.}\end{cases}
\]
The set where the sequence $\{X_n\}_{n=1}^\infty$ converges in $\R$ belongs to $\F_\infty$, since every $X_n$ is $\F_\infty$-measurable; hence $X_\infty$ is $\F_\infty$-measurable, and $X_\infty=\lim_{t\to\infty}X_t$ on $\Omega^*$. Fatou's lemma, the identity $\abs{x}=2x^+-x$, the inequality $E X_n\ge E X_0$ and \eqref{eq:1-3-23:plus} give
\[
E\abs{X_\infty}\le\liminf_{n\to\infty}E\abs{X_n}=\liminf_{n\to\infty}\bigl(2E(X_n^+)-E X_n\bigr)\le 2E(Y^+)-E X_0<\infty .
\]

\emph{Step 2: $X_\infty$ is a last element.} We have to show that $E(X_\infty\mid\F_t)\ge X_t$ a.s.\ for every $t\ge0$, that is,
\begin{equation}\label{eq:1-3-23:last}
\int_A X_t\,dP\le\int_A X_\infty\,dP\qquad\text{for every } t\ge0,\ A\in\F_t .
\end{equation}
The family $\{E(Y^+\mid\F_t);\ t\ge0\}$ of conditional expectations of one integrable random variable is uniformly integrable; by \eqref{eq:1-3-23:plus}, so is $\{X_n^+\}_{n=1}^\infty$. Since $X_n^+\to X_\infty^+$ a.s., the convergence takes place in $L^1$ as well, and therefore
\[
\lim_{n\to\infty}\int_A X_n^+\,dP=\int_A X_\infty^+\,dP\qquad\text{for every } A\in\F .
\]
On the other hand $X_n^-\to X_\infty^-$ a.s., and Fatou's lemma yields
\[
\int_A X_\infty^-\,dP\le\liminf_{n\to\infty}\int_A X_n^-\,dP\qquad\text{for every } A\in\F .
\]
Now fix $t\ge0$ and $A\in\F_t$. For every integer $n>t$ the submartingale property gives $\int_A X_t\,dP\le\int_A X_n\,dP=\int_A X_n^+\,dP-\int_A X_n^-\,dP$, and consequently
\[
\int_A X_t\,dP\le\lim_{n\to\infty}\int_A X_n^+\,dP-\liminf_{n\to\infty}\int_A X_n^-\,dP
\le\int_A X_\infty^+\,dP-\int_A X_\infty^-\,dP=\int_A X_\infty\,dP,
\]
which is \eqref{eq:1-3-23:last}. Hence $\{X_t,\F_t;\ 0\le t\le\infty\}$ is a right-continuous submartingale with last element $X_\infty$.

\emph{Step 3: conclusion.} \bk{Theorem 3.22}, applied to $\{X_t,\F_t;\ 0\le t\le\infty\}$ and the optional times $S\le T$, shows that $X_S$ and $X_T$ (equal to $X_\infty$ on $\{S=\infty\}$, respectively $\{T=\infty\}$) are integrable, that \eqref{eq:1-3-23:os} holds, that $\F_{S+}$ may be replaced by $\F_S$ when $S$ is a stopping time, and that $E X_T\ge E X_0$. No boundedness of $S$ or $T$ is needed.

\rmk{(a) In (ii), the random variable $Y'\defeq E(Y\mid\F_\infty)$ is integrable, $\F_\infty$-measurable and satisfies $E(Y'\mid\F_t)=E(Y\mid\F_t)\ge X_t$ a.s., so it is also a last element, and \bk{Theorem 3.22} applies at once with $X_\infty$ replaced by $Y'$. The two versions coincide when $T<\infty$ a.s. In general the version proved above uses the smaller, and more natural, value on $\{T=\infty\}$: by hypothesis (ii) and L\'evy's theorem on the a.s.\ convergence of $E(Y\mid\F_n)$ to $Y'$, we have $X_\infty=\lim_n X_n\le Y'$ a.s.

(b) If $X$ is a martingale satisfying (ii), then only the inequality $E(X_T\mid\F_{S+})\ge X_S$ is obtained, because $-X$ need not satisfy (ii); equality holds if in addition $X_t\ge E(Y''\mid\F_t)$ a.s.\ for some integrable $Y''$ and every $t\ge 0$, since Steps 1--3 then apply to $-X$, with the same limit up to sign. Equality can fail otherwise: let $\xi_1,\xi_2,\dots$ be independent with $P[\xi_k=2]=P[\xi_k=0]=\tfrac12$, let $\F_t\defeq\sigma(\xi_k;\ k\le t)$ and $X_t\defeq-\prod_{k\le t}\xi_k$ (with $X_t=-1$ for $t<1$). This is a right-continuous martingale with $X_t\le 0=E(0\mid\F_t)$, and $X_\infty=0$ a.s.\ because a.s.\ some $\xi_k$ vanishes; for $T\equiv\infty$ we get $E X_T=0>-1=E X_0$.}
\end{kex}

\begin{kex}{Problem}{3.24}{}
Every sample path of $X$ is right-continuous and $X$ is adapted, so $X$ is progressively measurable with respect to
$\{\F_t\}$ by \bk{Proposition 1.13}. The stopping times $S$, $T$ may take the value $+\infty$; no last element of
$X$ is needed, because only the finite times $T\wedge t$ and $S\wedge t$ enter the statement. Both parts follow from one
inequality.

\begin{claim}
Let $V$ be a bounded stopping time of $\{\F_t\}$ and $R$ an arbitrary stopping time of $\{\F_t\}$. Then $X_V$ and
$X_{R\wedge V}$ are integrable, $X_V$ is $\F_V$-measurable, and
\begin{equation}\label{eq:1-3-24:a}
E[X_V \mid \F_R] \ge X_{R\wedge V} \quad\text{a.s. } P .
\end{equation}
\end{claim}
\begin{proof}
Let $a>0$ be a constant with $V\le a$. By \bk{Lemma 2.9}, $R\wedge V$ is a stopping time; it satisfies
$R\wedge V\le V\le a$. Every stopping time is optional, so \bk{Problem 3.23}\,(i), applied to the right-continuous
submartingale $X$ and to the pair $R\wedge V\le V$ with $V$ bounded, yields that $X_{R\wedge V}$ and $X_V$ are
integrable and, $R\wedge V$ being a stopping time, that
\begin{equation}\label{eq:1-3-24:b}
E[X_V \mid \F_{R\wedge V}] \ge X_{R\wedge V} \quad\text{a.s. } P .
\end{equation}
Since $V$ is finite, $X_V$ is defined on all of $\Omega$, and it is $\F_V$-measurable by \bk{Proposition 2.18}; being
integrable, it satisfies $E[X_V\mid\F_V]=X_V$ a.s. We now apply \bk{Problem 2.17}\,(ii) to the stopping times $V$, $R$
and the integrable random variable $Z=X_V$:
\[
E[X_V\mid\F_R]=E\bigl[\,E[X_V\mid\F_V]\bigm|\F_R\bigr]=E[X_V\mid\F_{R\wedge V}] \quad\text{a.s. } P .
\]
Together with \eqref{eq:1-3-24:b} this is \eqref{eq:1-3-24:a}.
\end{proof}

\textbf{(i)} Fix $t\ge 0$. By \bk{Lemma 2.9} (the constant $t$ is a stopping time), $T\wedge t$ is a stopping time,
and it is bounded by $t$. By the Claim, $X_{T\wedge t}$ is integrable and $\F_{T\wedge t}$-measurable; since
$T\wedge t\le t$, \bk{Lemma 2.15} and \bk{Problem 2.13} give $\F_{T\wedge t}\subseteq\F_t$, so that $X_{T\wedge t}$ is
$\F_t$-measurable. Thus the process $\{X_{T\wedge t};\,0\le t<\infty\}$ is adapted to $\{\F_t\}$ and integrable.

Now let $0\le s<t$. We apply the Claim with $V=T\wedge t$ and with the constant stopping time $R=s$, for which
$\F_R=\F_s$ by \bk{Problem 2.13} and $R\wedge V=s\wedge T\wedge t=T\wedge s$. Inequality \eqref{eq:1-3-24:a} reads
\[
E[X_{T\wedge t}\mid\F_s]\ge X_{T\wedge s} \quad\text{a.s. } P ,
\]
which is the submartingale property of $\{X_{T\wedge t},\F_t;\,0\le t<\infty\}$.

\textbf{(ii)} Fix $t\ge 0$ and apply the Claim with $V=T\wedge t$ and $R=S$. Because $S\le T$, we have
$R\wedge V=S\wedge T\wedge t=S\wedge t$, and \eqref{eq:1-3-24:a} becomes
\[
E[X_{T\wedge t}\mid\F_S]\ge X_{S\wedge t} \quad\text{a.s. } P .
\]

\rmk{The stopped process in (i) is again right-continuous, since $t\mapsto T(\omega)\wedge t$ is continuous and
nondecreasing. The hypothesis $S\le T$ is used in (ii) only to identify $S\wedge T\wedge t$ with $S\wedge t$: for two
arbitrary stopping times the same argument gives $E[X_{T\wedge t}\mid\F_S]\ge X_{S\wedge T\wedge t}$ a.s.}
\end{kex}

\begin{kex}{Problem}{3.25}{}
Let $X=\{X_t,\F_t;\ 0\le t<\infty\}$ be the submartingale in question; by the standing conventions of the book's Section 1.3, $X$ is adapted to the given filtration $\{\F_t\}$ and $E\abs{X_t}<\infty$ for every $t\ge 0$, so that all expectations and conditional expectations below are well defined and finite. No regularity of the paths is needed. By \bk{Definition 3.1} we have to show that $E(X_t\,|\,\F_s)\le X_s$ a.s.\ for all $0\le s<t<\infty$; together with the submartingale inequality this gives $E(X_t\,|\,\F_s)=X_s$ a.s.

Fix $0\le s<t<\infty$ and set
\[
  Y\;\defeq\;E(X_t\,|\,\F_s)-X_s ,
\]
where we choose a version of the conditional expectation that is finite everywhere. The random variable $Y$ is $\F_s$-measurable and integrable, and $Y\ge 0$ a.s.\ by the submartingale property. By the defining property of conditional expectation (applied to the set $\Omega\in\F_s$) and the hypothesis of constant expectation,
\[
  E(Y)\;=\;E\big[E(X_t\,|\,\F_s)\big]-E(X_s)\;=\;E(X_t)-E(X_s)\;=\;E(X_0)-E(X_0)\;=\;0 .
\]
A random variable which is a.s.\ nonnegative and has expectation zero vanishes a.s.: for every $n\ge 1$,
$P[Y\ge 1/n]\le n\,E\big(Y\ind_{\{Y\ge 0\}}\big)=n\,E(Y)=0$ by the \v{C}eby\v{s}ev inequality, and so $P[Y>0]=\lim_{n\to\infty}P[Y\ge 1/n]=0$. Hence $Y=0$ a.s., that is,
\[
  E(X_t\,|\,\F_s)\;=\;X_s\quad\text{a.s.}
\]
Since $0\le s<t<\infty$ were arbitrary, $X$ is both a submartingale and a supermartingale, i.e., a martingale with respect to $\{\F_t\}$.

\rmk{The argument shows slightly more: a submartingale $X$ satisfies $E(X_t\,|\,\F_s)=X_s$ a.s.\ for a given pair $s<t$ if and only if $E(X_s)=E(X_t)$. Applying the result to $-X$, a supermartingale of constant expectation is likewise a martingale.}
\end{kex}

\begin{kex}{Problem}{3.29}{}
No condition on the filtration $\{\F_t\}$ is needed. We recall the convention $\inf\emptyset=\infty$.

\emph{Step 1: $T$ is a stopping time, and $X_T=0$ on $\{T<\infty\}$.}
The process $X$ is adapted with continuous paths and $T$ is the hitting time $H_\Gamma$ of the closed set $\Gamma=\{0\}$, so $T$ is a stopping time of $\{\F_t\}$ by \bk{Problem 2.7}. For every $\omega$ the set $\{t\ge 0;\ X_t(\omega)=0\}$ is closed in $[0,\infty)$, because $t\mapsto X_t(\omega)$ is continuous; if it is nonempty, it therefore contains its infimum. Hence
\begin{equation}\label{eq:1-3-29:a}
X_{T(\omega)}(\omega)=0\qquad\text{for every }\omega\in\{T<\infty\}.
\end{equation}

\emph{Step 2: a fixed $t$.}
Fix $t\ge 0$ and an integer $n\ge 1$, and set
\[
S\defeq T\wedge n,\qquad R\defeq (T\wedge n)+t .
\]
By \bk{Lemma 2.9}, $S$ and $R$ are stopping times (a constant is a stopping time); they satisfy $S\le R\le n+t$. The process $-X$ is a right-continuous submartingale, so \bk{Problem 3.23}(i), applied to $-X$ and to the bounded stopping times $S\le R$, shows that $X_S$ and $X_R$ are integrable and
\begin{equation}\label{eq:1-3-29:b}
E\big(X_R\,\big|\,\F_S\big)\le X_S\qquad\text{a.s. } P.
\end{equation}
The event $A_n\defeq\{T\le n\}$ belongs to $\F_S$: for $s\ge n$ we have $A_n\cap\{S\le s\}=\{T\le n\}\in\F_n\subseteq\F_s$, and for $0\le s<n$ we have $A_n\cap\{S\le s\}=\{T\le s\}\in\F_s$. On $A_n$ we have $S=T$ and $R=T+t$, so that $X_S=X_T=0$ there by \eqref{eq:1-3-29:a}, and $X_R=X_{T+t}$. Integrating \eqref{eq:1-3-29:b} over $A_n$ we obtain
\[
E\big[X_{T+t}\,\ind_{\{T\le n\}}\big]=E\big[X_R\,\ind_{A_n}\big]\le E\big[X_S\,\ind_{A_n}\big]=0 .
\]
Since $X_R\ind_{A_n}\ge 0$, this forces $X_{T+t}\,\ind_{\{T\le n\}}=0$ a.s. Thus there is a $P$-null event $N_{t,n}$ with $X_{T+t}(\omega)=0$ for every $\omega\in\{T\le n\}\setminus N_{t,n}$.

\emph{Step 3: all $t$ simultaneously.}
Let $N\defeq\bigcup_{n\ge 1}\bigcup_{q\in\Q\cap[0,\infty)}N_{q,n}$, a $P$-null event. If $\omega\in\{T<\infty\}\setminus N$, then $\omega\in\{T\le n\}$ for some integer $n$, and therefore $X_{T(\omega)+q}(\omega)=0$ for every rational $q\ge 0$. The function $t\mapsto X_{T(\omega)+t}(\omega)$ is continuous on $[0,\infty)$ and vanishes on the dense set $\Q\cap[0,\infty)$, hence it vanishes identically. Consequently
\[
X_{T+t}=0\quad\text{for all }0\le t<\infty,\qquad\text{on }\{T<\infty\}\setminus N,
\]
which is the assertion.

\rmk{Continuity of the paths is essential only at one point, namely to obtain from \bk{Problem 2.7} that $T$ is a stopping time. The other two uses need right-continuity alone: if $T(\omega)<\infty$, there are $t_k\ge T(\omega)$ with $t_k\to T(\omega)$ and $X_{t_k}(\omega)=0$, so that $X_{T(\omega)}(\omega)=0$; and in Step 3 a right-continuous function vanishing on $\Q\cap[0,\infty)$ vanishes identically. If ``nonnegative'' is understood only as $X_t\ge 0$ a.s.\ for each $t$, continuity gives $X_t\ge 0$ for all $t$ outside a single null event, and the argument is unchanged.}
\end{kex}

\begin{kex}{Exercise}{3.30}{}
We read the hypothesis ``increasing'' in its weakest form: for every $t\ge 0$ and $n\ge 1$ we have $X^{(n)}_t\le X^{(n+1)}_t$ a.s. For fixed $t$ the limit $\lim_n X^{(n)}_t$ then exists in $(-\infty,\infty]$ outside a $P$-null set, and $\xi_t$ is a real-valued random variable which agrees with this limit a.s.; in particular
\begin{equation}\label{eq:1-3-30:a}
  X^{(n)}_t\;\le\;\xi_t\quad\text{a.s., for every } n\ge 1,\ t\ge 0 .
\end{equation}
We show first that $\xi$ is a supermartingale, then that $t\mapsto E\xi_t$ is right-continuous, and finally we apply \bk{Theorem 3.13}.

\emph{Step 1: $\xi$ is a supermartingale.} Fix $t\ge 0$. Each $X^{(n)}_t$ is $\F_t$-measurable, so $\limsup_n X^{(n)}_t$ is an $\F_t$-measurable extended random variable; $\xi_t$ agrees with it a.s., and $\F_t\supseteq\F_0$ contains all $P$-null sets of $\F$ by the usual conditions (\bk{Definition 2.25}), so $\xi_t$ is $\F_t$-measurable. It is integrable by assumption. Next, the random variables $X^{(n)}_t-X^{(1)}_t$, $n\ge 1$, are a.s.\ nonnegative and a.s.\ increase to $\xi_t-X^{(1)}_t$. For every $A\in\F$ the monotone convergence theorem gives $\int_A (X^{(n)}_t-X^{(1)}_t)\,dP\uparrow\int_A(\xi_t-X^{(1)}_t)\,dP$, and since $X^{(1)}_t$ is integrable we may add $\int_A X^{(1)}_t\,dP$ to both sides:
\begin{equation}\label{eq:1-3-30:b}
  \int_A X^{(n)}_t\,dP\;\uparrow\;\int_A \xi_t\,dP\qquad (n\to\infty),\quad A\in\F,\ t\ge 0 .
\end{equation}
Now let $0\le s<t<\infty$ and $A\in\F_s$. The supermartingale property of $X^{(n)}$ reads $\int_A X^{(n)}_t\,dP\le\int_A X^{(n)}_s\,dP$; letting $n\to\infty$ and using \eqref{eq:1-3-30:b} at the times $t$ and $s$ we obtain $\int_A\xi_t\,dP\le\int_A\xi_s\,dP$. As $\xi_s$ is $\F_s$-measurable and $A\in\F_s$ is arbitrary, this means $E(\xi_t\,|\,\F_s)\le\xi_s$ a.s. Hence $\xi=\{\xi_t,\F_t;\ 0\le t<\infty\}$ is a supermartingale, and in particular the function $t\mapsto E\xi_t$ is nonincreasing.

\emph{Step 2: $t\mapsto E\xi_t$ is right-continuous.} Fix $n\ge 1$. The process $-X^{(n)}$ is a submartingale with respect to $\{\F_t\}$, which satisfies the usual conditions, and it has a right-continuous modification, namely itself. By the ``only if'' part of \bk{Theorem 3.13}, the function $t\mapsto E X^{(n)}_t$ is right-continuous on $[0,\infty)$.

Fix now $t\ge 0$. Because $u\mapsto E\xi_u$ is nonincreasing, the limit $\ell\defeq\lim_{s\downarrow t}E\xi_s=\sup_{s>t}E\xi_s$ exists and satisfies $\ell\le E\xi_t$. On the other hand, for every $n\ge 1$ and $s>t$ we have $EX^{(n)}_s\le E\xi_s\le \ell$ by \eqref{eq:1-3-30:a}, and letting $s\downarrow t$ we get
\[
  E X^{(n)}_t\;=\;\lim_{s\downarrow t}E X^{(n)}_s\;\le\;\ell .
\]
Letting $n\to\infty$ and using \eqref{eq:1-3-30:b} with $A=\Omega$, we conclude that $E\xi_t\le\ell$. Therefore $\lim_{s\downarrow t}E\xi_s=E\xi_t$.

\emph{Step 3: the modification.} The process $-\xi=\{-\xi_t,\F_t;\ 0\le t<\infty\}$ is a submartingale by Step 1, the filtration satisfies the usual conditions, and $t\mapsto E(-\xi_t)$ is right-continuous by Step 2. According to \bk{Theorem 3.13}, $-\xi$ has a modification $Y$ which is RCLL and adapted to $\{\F_t\}$. Put $X\defeq -Y$. Then $X$ has RCLL paths, is adapted to $\{\F_t\}$, and $P[X_t=\xi_t]=1$ for every $t\ge 0$, i.e., $X$ is a modification of $\xi$. Consequently $E\abs{X_t}=E\abs{\xi_t}<\infty$, and for $0\le s<t<\infty$,
\[
  E(X_t\,|\,\F_s)\;=\;E(\xi_t\,|\,\F_s)\;\le\;\xi_s\;=\;X_s\quad\text{a.s.,}
\]
so that $X=\{X_t,\F_t;\ 0\le t<\infty\}$ is an RCLL supermartingale.

\rmk{The nonnegativity of $\xi_t$ has not been used; integrability of $\xi_t$ for every $t$ suffices, since the sequence is bounded from below by the integrable random variable $X^{(1)}_t$. The right-continuity of the $X^{(n)}$ entered only through the right-continuity of the functions $t\mapsto EX^{(n)}_t$: the essential point of Step 2 is that an increasing limit of right-continuous functions is lower semicontinuous from the right, while the supermartingale property of the limit makes $t\mapsto E\xi_t$ nonincreasing, and a nonincreasing function with this property is right-continuous.}
\end{kex}

\solsection{1.4}{The Doob--Meyer Decomposition}

\begin{kex}{Problem}{4.13}{}
Let $X=\{X_t,\F_t;\ 0\le t<\infty\}$ be a submartingale all of whose sample paths are continuous and which satisfies $X_t\ge 0$ a.s.\ for every $t\ge 0$. Fix $a>0$ and a nondecreasing sequence $\{T_n\}_{n=1}^\infty\subseteq\cS_a$ of stopping times of $\{\F_t\}$, and set $T\defeq\lim_{n\to\infty}T_n=\sup_{n\ge 1}T_n$. According to \bk{Definition 4.12} we have to show that $\lim_{n\to\infty}E(X_{T_n})=E(X_T)$.

\emph{Step 1: $T\in\cS_a$, and the random variables $X_{T_n}$, $X_T$ are well defined.} For every $t\ge 0$ we have $\{T\le t\}=\bigcap_{n\ge 1}\{T_n\le t\}\in\F_t$, so $T$ is a stopping time of $\{\F_t\}$. By the definition of $\cS_a$ (\bk{Definition 4.8}), $P(T_n\le a)=1$ for every $n$; hence the event
\[
  \Omega_0\;\defeq\;\bigcap_{n\ge 1}\{T_n\le a\}\;=\;\{T\le a\}
\]
has probability one, and $T\in\cS_a$. Being adapted with continuous paths, $X$ is progressively measurable (\bk{Proposition 1.13}), so for every stopping time $S$ the function $X_S$, defined on $\{S<\infty\}$, is measurable (\bk{Proposition 2.18}). The random variables $X_{T_n}$ and $X_T$ are thus defined on $\Omega_0$, that is, almost surely, which is all that is needed to speak of their expectations and of their uniform integrability.

\emph{Step 2: almost sure convergence.} Let $\omega\in\Omega_0$. Then $T_n(\omega)\uparrow T(\omega)$ and all these numbers lie in $[0,a]$; in particular the limit $T(\omega)$ is finite. Since the path $t\mapsto X_t(\omega)$ is continuous on $[0,\infty)$, and in particular at the point $T(\omega)$, we obtain
\begin{equation}\label{eq:1-4-13:a}
  \lim_{n\to\infty}X_{T_n(\omega)}(\omega)\;=\;X_{T(\omega)}(\omega)\qquad\text{for every }\omega\in\Omega_0 ,
\end{equation}
i.e., $X_{T_n}\to X_T$ almost surely.

\emph{Step 3: uniform integrability and conclusion.} The process $X$ is a right-continuous submartingale with $X_t\ge 0$ a.s.\ for every $t\ge 0$; by \bk{Problem 4.9}(a) it is of class $DL$, which means (\bk{Definition 4.8}) that the family $\{X_S\}_{S\in\cS_a}$ is uniformly integrable. Since $T_n\in\cS_a$ for every $n$ and $T\in\cS_a$ by Step 1, the sequence $\{X_{T_n}\}_{n=1}^\infty$ is uniformly integrable and $X_T$ is integrable. A uniformly integrable sequence which converges almost surely to an integrable random variable converges to it in $L^1$ (Vitali's convergence theorem); from \eqref{eq:1-4-13:a} we conclude that $E\abs{X_{T_n}-X_T}\to 0$, and therefore
\[
  \abs[\big]{E(X_{T_n})-E(X_T)}\;\le\;E\abs{X_{T_n}-X_T}\;\xrightarrow[n\to\infty]{}\;0 .
\]
As $a>0$ and the sequence $\{T_n\}$ were arbitrary, $X$ is regular.

\rmk{The uniform integrability of Step 3 can also be seen directly, which makes clear where nonnegativity enters. For $S\in\cS_a$, the optional sampling theorem for bounded stopping times (\bk{Problem 3.23}(i), applied to the stopping times $S\wedge a\le a$; note that $X_S=X_{S\wedge a}$ a.s.) gives $X_S\le E(X_a\,|\,\F_{S\wedge a})$ a.s. Moreover $X_S\ge 0$ a.s., because by right-continuity the null set $\bigcup_{q\in\Q\cap[0,\infty)}\{X_q<0\}$ contains every $\omega$ for which the path $X_\cdot(\omega)$ takes a negative value. Hence $0\le X_S\le E(X_a\,|\,\F_{S\wedge a})$ a.s., and the family of conditional expectations of the single integrable random variable $X_a$ is uniformly integrable. Without nonnegativity only the upper bound survives, and it does not control the negative parts.}

\rmk{Nonnegativity was used only to obtain membership in the class $DL$, and continuity only through \eqref{eq:1-4-13:a}. The argument therefore shows more generally that every submartingale of class $DL$ with continuous paths is regular; for instance, by \bk{Problem 4.9}(b), every continuous submartingale of the form \bk{(4.6)}, and in particular every continuous martingale.}
\end{kex}

\solsection{1.5}{Continuous, Square-Integrable Martingales}

\begin{kex}{Problem}{5.7}{}
Throughout, $\{\F_t\}$ satisfies the usual conditions, as in the book's Section 1.5. Quadratic variations are determined only up to indistinguishability (\bk{Definition 5.3}), so (i) and (ii) are understood as identities between indistinguishable processes, and (iii) as an inequality valid for all $t\ge 0$ simultaneously, for $P$-a.e.\ $\omega$. We call a process $A$ \emph{of class $\mathscr{V}$} if $A=A^{(1)}-A^{(2)}$, where $A^{(1)}$, $A^{(2)}$ are natural increasing processes (\bk{Definitions 4.4, 4.5}).

\emph{Step 0. Preliminaries.} \textbf{(a)} If $A$, $A'$ are natural increasing processes and $c,c'\ge 0$, then $cA+c'A'$ is a natural increasing process. Indeed, it is adapted, a.s.\ null at $t=0$ with nondecreasing, right-continuous paths, and $E(cA_t+c'A'_t)<\infty$; and for a.e.\ $\omega$ the Lebesgue--Stieltjes measure of $cA(\omega)+c'A'(\omega)$ is $c\,dA(\omega)+c'\,dA'(\omega)$, so both sides of \bk{(4.4)} for $cA+c'A'$ are the corresponding combinations of those for $A$ and $A'$ (all integrals being finite, because the martingale in \bk{(4.4)} is bounded and $E(A_t)$, $E(A'_t)$ are finite). Consequently, a linear combination $\sum_{i=1}^n c_iA^{(i)}$ of natural increasing processes with real coefficients is of class $\mathscr{V}$: write it as $\sum_{i:\,c_i\ge 0}c_iA^{(i)}-\sum_{i:\,c_i<0}\abs{c_i}A^{(i)}$. In particular, linear combinations of processes of class $\mathscr{V}$ are of class $\mathscr{V}$, and so is $\qv{X,Y}=\frac14\qv{X+Y}-\frac14\qv{X-Y}$ for $X,Y\in\cM_2$.

\textbf{(b)} $\cM_2$ is a vector space: a linear combination of right-continuous, square-integrable martingales which vanish a.s.\ at $t=0$ has the same properties. For $X,Y\in\cM_2$ the product $X_tY_t$ is integrable by the Cauchy--Schwarz inequality.

\begin{lem}
Let $B$ be a process of class $\mathscr{V}$ which is a martingale. Then $B$ is indistinguishable from the zero process.
\end{lem}
\begin{proof}
This is the uniqueness argument in the proof of \bk{Theorem 4.10}, which uses only the hypotheses of the lemma; we recall it. Write $B=A'-A''$ with $A'$, $A''$ natural increasing, fix $t>0$, and let $\{\eta_u,\F_u;\ 0\le u<\infty\}$ be a bounded, right-continuous martingale, $\abs{\eta}\le C$. By \bk{(4.4)$'$} of \bk{Lemma 4.7}, applied to $A'$ and to $A''$,
\[
E\bigl[\eta_t(A'_t-A''_t)\bigr]=E\int_{(0,t]}\eta_{u-}\,dA'_u-E\int_{(0,t]}\eta_{u-}\,dA''_u .
\]
Let $\Pi_n=\{t^{(n)}_0,\dots,t^{(n)}_{m_n}\}$ be partitions of $[0,t]$ with mesh tending to zero, and put $\eta^{(n)}_u\defeq\sum_{j=1}^{m_n}\eta_{t^{(n)}_{j-1}}\ind_{(t^{(n)}_{j-1},t^{(n)}_j]}(u)$. Since every path of the martingale $\eta$ is right-continuous, \bk{Theorem 3.8(v)} applies: for a.e.\ $\omega$ the path of $\eta$ has left-hand limits everywhere on $(0,\infty)$, so $\eta^{(n)}_u\to\eta_{u-}$ for every $u\in(0,t]$; since $\abs{\eta^{(n)}}\le C$ and $E(A'_t+A''_t)<\infty$, dominated convergence (first for the Lebesgue--Stieltjes integrals, then for the expectations) gives
\[
E\bigl[\eta_t(A'_t-A''_t)\bigr]
=\lim_{n\to\infty}E\sum_{j=1}^{m_n}\eta_{t^{(n)}_{j-1}}\bigl[B_{t^{(n)}_j}-B_{t^{(n)}_{j-1}}\bigr]=0,
\]
because $\eta_{t^{(n)}_{j-1}}$ is bounded and $\F_{t^{(n)}_{j-1}}$-measurable and $B$ is a martingale. Now let $\zeta$ be a random variable with $\abs{\zeta}\le C$, and choose for every $u\ge 0$ a version $\zeta_u$ of $E[\zeta\,|\,\F_u]$; then $\{\zeta_u,\F_u\}$ is a martingale and $u\mapsto E\zeta_u=E\zeta$ is constant, in particular right-continuous. As $\{\F_u\}$ satisfies the usual conditions, \bk{Theorem 3.13} provides a modification $\tilde\eta$ of $\{\zeta_u\}$ which is right-continuous, adapted to $\{\F_u\}$ and a martingale. For each $u$ we have $\abs{\tilde\eta_u}=\abs{E[\zeta\,|\,\F_u]}\le C$ a.s., so the event
\[
N_0\defeq\bigcup_{u\in\Q\cap[0,\infty)}\bigl\{\abs{\tilde\eta_u}>C\bigr\}
\]
has probability zero, and by right-continuity $\abs{\tilde\eta_u(\omega)}\le C$ for \emph{all} $u\ge 0$ when $\omega\notin N_0$. The usual conditions give $N_0\in\F_0$, so $\eta_u\defeq\tilde\eta_u\ind_{\Omega\setminus N_0}$ defines an adapted process, all of whose paths are right-continuous and bounded by $C$; it is a modification of $\tilde\eta$, hence a martingale with $\eta_t=E[\zeta\,|\,\F_t]$ a.s. Thus $\eta$ is a bounded, right-continuous martingale, to which the first part of the proof applies, and we obtain $E[\zeta B_t]=E[\eta_tB_t]=0$, since $B_t$ is $\F_t$-measurable and integrable. With $\zeta=\sgn(B_t)$ this yields $E\abs{B_t}=0$, i.e.\ $B_t=0$ a.s. Thus, almost surely, $B_t=0$ for all rational $t$, hence for all $t\ge 0$ by the right-continuity of $A'$ and $A''$.
\end{proof}

\emph{Step 1. Characterization of the cross-variation.} Let $X,Y\in\cM_2$. We claim: \emph{if $A$ is of class $\mathscr{V}$ and $XY-A$ is a martingale, then $A$ is indistinguishable from $\qv{X,Y}$.} Indeed, $XY-\qv{X,Y}$ is a martingale, being one quarter of the difference of the martingales $(X+Y)^2-\qv{X+Y}$ and $(X-Y)^2-\qv{X-Y}$ (recall that $X\pm Y\in\cM_2$). Hence $A-\qv{X,Y}=(XY-\qv{X,Y})-(XY-A)$ is a martingale; it is of class $\mathscr{V}$ by Step 0(a), and the lemma applies. This is the assertion made in the book after \bk{Definition 5.5}.

Two consequences will be used. First, for $c\in\R$ and $X\in\cM_2$,
\begin{equation}\label{eq:1-5-07:a}
\qv{cX}=c^2\qv{X},
\end{equation}
because $c^2\qv{X}$ is a natural increasing process and $(cX)^2-c^2\qv{X}=c^2(X^2-\qv{X})$ is a martingale, so the uniqueness in \bk{Definition 5.3} applies. Secondly,
\begin{equation}\label{eq:1-5-07:b}
\qv{X,X}=\tfrac14\qv{2X}-\tfrac14\qv{0}=\qv{X},
\end{equation}
by \eqref{eq:1-5-07:a} with $c=2$ and $c=0$.

\textbf{(ii)} Since $Y-X=-(X-Y)$, \eqref{eq:1-5-07:a} with $c=-1$ gives $\qv{Y-X}=\qv{X-Y}$, and trivially $\qv{Y+X}=\qv{X+Y}$. Hence $\qv{Y,X}=\frac14[\qv{Y+X}-\qv{Y-X}]=\qv{X,Y}$.

\textbf{(i)} Put $A\defeq\alpha\qv{X,Z}+\beta\qv{Y,Z}$. By Step 0(a), $A$ is of class $\mathscr{V}$. By Step 0(b), $\alpha X+\beta Y\in\cM_2$, and
\[
(\alpha X+\beta Y)Z-A=\alpha\bigl(XZ-\qv{X,Z}\bigr)+\beta\bigl(YZ-\qv{Y,Z}\bigr)
\]
is a martingale, as a linear combination of two martingales. Step 1, applied to the pair $\alpha X+\beta Y$, $Z$, shows that $A$ is indistinguishable from $\qv{\alpha X+\beta Y,Z}$.

Together with (ii), this gives linearity in the second argument as well, so $\qv{\cdot,\cdot}$ is a symmetric bilinear form on $\cM_2$ (with values in the processes of class $\mathscr{V}$, modulo indistinguishability).

\textbf{(iii) and (iv).} \emph{Step 2. An inequality for increments.} Write $\xi=\qv{X,Y}$. For every $\lambda\in\R$, bilinearity, symmetry and \eqref{eq:1-5-07:b} give
\[
\qv{X+\lambda Y}=\qv{X+\lambda Y,X+\lambda Y}=\qv{X}+2\lambda\,\xi+\lambda^2\qv{Y}
\]
up to indistinguishability, and $\qv{X+\lambda Y}$ has a.s.\ nondecreasing paths. Hence there is an event $N$ with $P(N)=0$ (a countable union of null events, one for each rational $\lambda$) such that for every $\omega\notin N$, every \emph{rational} $\lambda$ and all $0\le s<t<\infty$,
\begin{equation}\label{eq:1-5-07:c}
a+2\lambda b+\lambda^2c\ \ge\ 0,
\end{equation}
where
\[
a\defeq\qv{X}_t(\omega)-\qv{X}_s(\omega),\qquad
b\defeq\xi_t(\omega)-\xi_s(\omega),\qquad
c\defeq\qv{Y}_t(\omega)-\qv{Y}_s(\omega);
\]
enlarging $N$ by a null event, we may also assume that $\qv{X}_0(\omega)=\qv{Y}_0(\omega)=\xi_0(\omega)=0$ and that $\qv{X}(\omega)$, $\qv{Y}(\omega)$ are nondecreasing for $\omega\notin N$, so that $a,c\ge 0$. The left-hand side of \eqref{eq:1-5-07:c} is continuous in $\lambda$, so \eqref{eq:1-5-07:c} holds for all real $\lambda$. If $c>0$, the choice $\lambda=-b/c$ gives $a-b^2/c\ge 0$; if $c=0$, then $a+2\lambda b\ge 0$ for all $\lambda\in\R$ forces $b=0$. In both cases $b^2\le ac$, that is: for every $\omega\notin N$,
\begin{equation}\label{eq:1-5-07:d}
\abs{\xi_t-\xi_s}^2\le\bigl(\qv{X}_t-\qv{X}_s\bigr)\bigl(\qv{Y}_t-\qv{Y}_s\bigr);\qquad 0\le s<t<\infty .
\end{equation}

\emph{Step 3. Proof of} (iii). Taking $s=0$ in \eqref{eq:1-5-07:d} we obtain $\abs{\qv{X,Y}_t}^2\le\qv{X}_t\qv{Y}_t$ for all $t\ge 0$ and all $\omega\notin N$.

\emph{Step 4. Proof of} (iv). Fix $\omega\notin N$ and $0\le s<t<\infty$. For $0\le u<v$, \eqref{eq:1-5-07:d} and the inequality $\sqrt{pq}\le\frac12(p+q)$ for $p,q\ge 0$ give
\[
\abs{\xi_v-\xi_u}\le\tfrac12\bigl[\qv{X}_v-\qv{X}_u+\qv{Y}_v-\qv{Y}_u\bigr].
\]
Summing over the intervals of an arbitrary partition $s=t_0<t_1<\dots<t_m=t$ of $[s,t]$, the right-hand sides telescope:
\[
\sum_{k=1}^m\abs{\xi_{t_k}-\xi_{t_{k-1}}}\le\tfrac12\bigl[\qv{X}_t-\qv{X}_s+\qv{Y}_t-\qv{Y}_s\bigr].
\]
Taking the supremum over all such partitions, the total variation $V_{[s,t]}$ of $\xi(\omega)$ on $[s,t]$ satisfies
\[
V_{[s,t]}\le\tfrac12\bigl[\qv{X}_t-\qv{X}_s+\qv{Y}_t-\qv{Y}_s\bigr]<\infty .
\]
In particular $\check{\xi}_t=V_{[0,t]}<\infty$ for every $t$, and since the total variation of a function is additive over adjacent intervals, $\check{\xi}_t=\check{\xi}_s+V_{[s,t]}$. Therefore
\[
\check{\xi}_t(\omega)-\check{\xi}_s(\omega)\le\tfrac12\bigl[\qv{X}_t(\omega)-\qv{X}_s(\omega)+\qv{Y}_t(\omega)-\qv{Y}_s(\omega)\bigr];\qquad 0\le s<t<\infty,
\]
for every $\omega\notin N$, which is (iv).

\rmk{Applying the Cauchy--Schwarz inequality to the sum $\sum_k\abs{\xi_{t_k}-\xi_{t_{k-1}}}$, with each term bounded by means of \eqref{eq:1-5-07:d}, yields the sharper bound $\check{\xi}_t-\check{\xi}_s\le\bigl(\qv{X}_t-\qv{X}_s\bigr)^{1/2}\bigl(\qv{Y}_t-\qv{Y}_s\bigr)^{1/2}$, valid for all $0\le s<t<\infty$ outside the same null event $N$.}
\end{kex}

\begin{kex}{Problem}{5.14}{}
We reduce the claim to \bk{Theorem 5.8} by polarization. As in that theorem, $t>0$ is fixed,
partitions are $0=t_0\le t_1\le\dots\le t_m=t$, and the assertion to be proved is: for every
$\eps>0$, $\eta>0$ there exists $\delta>0$ such that $\norm{\Pi}<\delta$ implies
$P\bigl[\abs{C_t(\Pi)-\qv{X,Y}_t}>\eps\bigr]<\eta$, where
\[
C_t(\Pi)\defeq\sum_{k=1}^{m}(X_{t_k}-X_{t_{k-1}})(Y_{t_k}-Y_{t_{k-1}}).
\]

\emph{Step 1: $X+Y$ and $X-Y$ belong to $\cM_2^c$.} Both processes are adapted to $\{\F_t\}$,
continuous, and vanish at time $0$ a.s.; they are martingales as linear combinations of
martingales with respect to the same filtration; and
$E(X_s\pm Y_s)^2\le 2EX_s^2+2EY_s^2<\infty$ for every $s\ge 0$. Hence \bk{Theorem 5.8} applies
to each of them. Writing $V_t^{(2)}(\Pi;Z)\defeq\sum_{k=1}^m\abs{Z_{t_k}-Z_{t_{k-1}}}^2$ for the
quadratic variation of a process $Z$ over $\Pi$, we thus have
\begin{equation}\label{eq:1-5-14:a}
\lim_{\norm{\Pi}\to 0}V_t^{(2)}(\Pi;X+Y)=\qv{X+Y}_t,\qquad
\lim_{\norm{\Pi}\to 0}V_t^{(2)}(\Pi;X-Y)=\qv{X-Y}_t\quad\text{(in probability)}.
\end{equation}

\emph{Step 2: polarization.} For real numbers $a$, $b$ we have $4ab=(a+b)^2-(a-b)^2$. Applying
this with $a=X_{t_k}-X_{t_{k-1}}$, $b=Y_{t_k}-Y_{t_{k-1}}$ and summing over $k$, we obtain, for
every partition $\Pi$ and every $\omega$,
\[
C_t(\Pi)=\tfrac14\bigl[V_t^{(2)}(\Pi;X+Y)-V_t^{(2)}(\Pi;X-Y)\bigr],
\]
whereas $\qv{X,Y}_t=\tfrac14\bigl[\qv{X+Y}_t-\qv{X-Y}_t\bigr]$ by \bk{Definition 5.5}.
Consequently
\begin{equation}\label{eq:1-5-14:b}
\abs{C_t(\Pi)-\qv{X,Y}_t}\le
\tfrac14\abs[\big]{V_t^{(2)}(\Pi;X+Y)-\qv{X+Y}_t}+\tfrac14\abs[\big]{V_t^{(2)}(\Pi;X-Y)-\qv{X-Y}_t}.
\end{equation}

\emph{Step 3: conclusion.} Let $\eps>0$ and $\eta>0$ be given. By \eqref{eq:1-5-14:a} there
are $\delta_+>0$ and $\delta_->0$ such that
\[
\norm{\Pi}<\delta_\pm\quad\Longrightarrow\quad
P\bigl[\abs{V_t^{(2)}(\Pi;X\pm Y)-\qv{X\pm Y}_t}>2\eps\bigr]<\frac{\eta}{2}.
\]
Set $\delta\defeq\delta_+\wedge\delta_-$ and let $\norm{\Pi}<\delta$. If both differences on the
right-hand side of \eqref{eq:1-5-14:b} are at most $2\eps$ in absolute value, then the left-hand
side is at most $\eps$; therefore
\[
P\bigl[\abs{C_t(\Pi)-\qv{X,Y}_t}>\eps\bigr]
\le\sum_{\pm}P\bigl[\abs{V_t^{(2)}(\Pi;X\pm Y)-\qv{X\pm Y}_t}>2\eps\bigr]<\eta .
\]
This is the asserted convergence in probability.
\end{kex}

\begin{kex}{Problem}{5.19}{}
Throughout, $X=\{X_t,\F_t;\ 0\le t<\infty\}$ is a local martingale in the sense of \bk{Definition 5.15}: $X$ is adapted, and there is a nondecreasing sequence $\{T_n\}_{n=1}^\infty$ of stopping times of $\{\F_t\}$ such that each $X^{(n)}=\{X_{t\wedge T_n},\F_t;\ 0\le t<\infty\}$ is a martingale and the event $\Omega_0\defeq\{\lim_{n\to\infty}T_n=\infty\}$ has probability one. We shall use two consequences of this definition. First, $X_0=X^{(n)}_0$ is integrable. Secondly, for every $\omega\in\Omega_0$ and every $t\ge 0$ there is an integer $n_0(\omega,t)$ with $T_n(\omega)\ge t$ for $n\ge n_0(\omega,t)$, so that
\begin{equation}\label{eq:1-5-19:a}
  X_{t\wedge T_n}=X_t\ \text{ for all } n\ge n_0(\cdot,t),\qquad\text{and in particular}\qquad
  \lim_{n\to\infty}X_{t\wedge T_n}=X_t\quad\text{on }\Omega_0 .
\end{equation}
Finally, the martingale property of $X^{(n)}$ reads
\begin{equation}\label{eq:1-5-19:b}
  E\big(X_{t\wedge T_n}\,\big|\,\F_s\big)=X_{s\wedge T_n}\quad\text{a.s.,}\qquad 0\le s<t<\infty,\ n\ge 1 .
\end{equation}

\textbf{(i)} Let $X$ be of class $DL$; according to \bk{Definition 4.8} this means that $X$ is right-continuous and that the family $\{X_T\}_{T\in\cS_a}$ is uniformly integrable for every $0<a<\infty$, where $\cS_a$ is the class of stopping times $T$ of $\{\F_t\}$ with $P[T\le a]=1$. Fix $0\le s<t<\infty$. For every $n\ge 1$ the random time $t\wedge T_n$ is a stopping time bounded by $t$, i.e., a member of $\cS_t$; hence the sequence $\{X_{t\wedge T_n}\}_{n=1}^\infty$ is uniformly integrable. By \eqref{eq:1-5-19:a} it converges a.s.\ to $X_t$. A uniformly integrable sequence which converges a.s.\ has an integrable limit and converges to it in $L^1$; thus $E\abs{X_t}<\infty$ and $E\abs{X_{t\wedge T_n}-X_t}\to 0$ as $n\to\infty$. Since conditional expectation is a contraction on $L^1$,
\[
  E\,\Big|E\big(X_{t\wedge T_n}\,\big|\,\F_s\big)-E\big(X_t\,\big|\,\F_s\big)\Big|
  \;\le\; E\abs{X_{t\wedge T_n}-X_t}\;\xrightarrow[n\to\infty]{}\;0 ,
\]
and along a subsequence $\{n_k\}$ we have $E(X_{t\wedge T_{n_k}}\,|\,\F_s)\to E(X_t\,|\,\F_s)$ a.s. On the other hand, the right-hand side of \eqref{eq:1-5-19:b} converges a.s.\ to $X_s$, by \eqref{eq:1-5-19:a} with $t$ replaced by $s$. Letting $n\to\infty$ along $\{n_k\}$ in \eqref{eq:1-5-19:b}, we obtain
\[
  E\big(X_t\,\big|\,\F_s\big)=X_s\quad\text{a.s.}
\]
Together with the adaptedness of $X$ and the integrability of every $X_t$ (the argument for the latter applies to each $t>0$, and $X_0$ is integrable as noted above), this shows that $X$ is a martingale.

\textbf{(ii)} Let $X$ be nonnegative, in the sense that $P[X_t\ge 0;\ \forall\,0\le t<\infty]=1$. Then for every $t\ge 0$ and $n\ge 1$ we have $X_{t\wedge T_n}\ge 0$ a.s. Fix $0\le s<t<\infty$. From \eqref{eq:1-5-19:a}, Fatou's lemma and the martingale property of $X^{(n)}$,
\[
  E(X_t)\;=\;E\Big(\lim_{n\to\infty}X_{t\wedge T_n}\Big)\;\le\;\varliminf_{n\to\infty}E\big(X_{t\wedge T_n}\big)
  \;=\;\varliminf_{n\to\infty}E\big(X_{0\wedge T_n}\big)\;=\;E(X_0)\;<\;\infty ,
\]
so that every $X_t$ is integrable. Next, Fatou's lemma for conditional expectations (valid for a.s.\ nonnegative random variables), \eqref{eq:1-5-19:b} and \eqref{eq:1-5-19:a} give
\[
  E\big(X_t\,\big|\,\F_s\big)\;=\;E\Big(\varliminf_{n\to\infty}X_{t\wedge T_n}\,\Big|\,\F_s\Big)
  \;\le\;\varliminf_{n\to\infty}E\big(X_{t\wedge T_n}\,\big|\,\F_s\big)
  \;=\;\varliminf_{n\to\infty}X_{s\wedge T_n}\;=\;X_s\quad\text{a.s.}
\]
Since $X$ is adapted, it is a supermartingale.

\rmk{If $X$ is right-continuous, the nonnegativity used in (ii) follows from the apparently weaker condition $P[X_t\ge 0]=1$ for every $t\ge 0$: then a.s.\ $X_q\ge 0$ for all rational $q\ge 0$, and right-continuity extends this to all $t\ge 0$ simultaneously. Without any regularity of the paths, nonnegativity at fixed times gives no control over $X_{t\wedge T_n}$, and the pathwise formulation is the appropriate one.}

\textbf{(iii)} Let $M\in\cM^{c,\loc}$. By \bk{Problem 5.17} and \bk{Definition 5.18}, $\qv{M}$ is an adapted, continuous process with $\qv{M}_0=0$ a.s., whose paths are a.s.\ nondecreasing, and
\[
  N\;\defeq\;M^2-\qv{M}\ \in\ \cM^{c,\loc} .
\]
Consequently $\qv{M}_\infty\defeq\lim_{t\to\infty}\qv{M}_t$ exists in $[0,\infty]$ a.s., and this is the meaning of $\qv{M}_S$ on $\{S=\infty\}$. The processes $M$ and $\qv{M}$, being continuous and adapted, are progressively measurable, so that $M_S\ind_{\{S<\infty\}}$ and $\qv{M}_S\ind_{\{S<\infty\}}$ are random variables; with the conventions just made, $M_S^2$ and $\qv{M}_S$ are random variables with values in $[0,\infty]$ (the latter outside a null set), and the two expectations in the assertion are well defined in $[0,\infty]$.

Let $\{T_n\}_{n=1}^\infty$ be a nondecreasing sequence of stopping times with $P[\lim_{n\to\infty}T_n=\infty]=1$ such that $N^{(n)}=\{N_{t\wedge T_n},\F_t;\ 0\le t<\infty\}$ is a martingale for every $n\ge 1$; it is continuous, and $N^{(n)}_0=M_0^2-\qv{M}_0=0$ a.s. Set
\[
  \sigma_n\;\defeq\;S\wedge n\wedge T_n ,\qquad n\ge 1 .
\]
The random time $S\wedge n$ is a stopping time bounded by $n$. The optional sampling theorem for bounded stopping times, \bk{Problem 3.23}(i), applied to the right-continuous martingale $N^{(n)}$ (i.e., to the submartingales $N^{(n)}$ and $-N^{(n)}$) and to the stopping times $0\le S\wedge n$, shows that $N^{(n)}_{S\wedge n}=N_{\sigma_n}$ is integrable and
\[
  E\big(N_{\sigma_n}\big)\;=\;E\big(N^{(n)}_{S\wedge n}\big)\;=\;E\big(N^{(n)}_0\big)\;=\;0 .
\]
Now $M^2_{\sigma_n}=N_{\sigma_n}+\qv{M}_{\sigma_n}$, where $N_{\sigma_n}$ is integrable and $\qv{M}_{\sigma_n}\ge 0$ a.s.; hence, as an identity in $[0,\infty]$,
\begin{equation}\label{eq:1-5-19:c}
  E\big(M^2_{\sigma_n}\big)\;=\;E\big(N_{\sigma_n}\big)+E\qv{M}_{\sigma_n}\;=\;E\qv{M}_{\sigma_n}\;\le\;E\qv{M}_S ,
\end{equation}
the inequality because $\sigma_n\le S$ and the paths of $\qv{M}$ are a.s.\ nondecreasing (on $\{S=\infty\}$ we have $\qv{M}_{\sigma_n}\le\sup_{t\ge 0}\qv{M}_t=\qv{M}_\infty$).

We claim that
\begin{equation}\label{eq:1-5-19:d}
  M_S^2\;\le\;\varliminf_{n\to\infty}M^2_{\sigma_n}\quad\text{on }\ \Omega_0=\Big\{\lim_{n\to\infty}T_n=\infty\Big\} .
\end{equation}
Indeed, let $\omega\in\Omega_0$. If $S(\omega)<\infty$, then $n\wedge T_n(\omega)\ge S(\omega)$ for all sufficiently large $n$, so $\sigma_n(\omega)=S(\omega)$ and $M^2_{\sigma_n}(\omega)=M^2_S(\omega)$ for all such $n$. If $S(\omega)=\infty$, then $\sigma_n(\omega)=n\wedge T_n(\omega)$ is finite and tends to infinity as $n\to\infty$; therefore, for every $u\ge 0$, we have $M^2_{\sigma_n}(\omega)\ge\inf_{t\ge u}M_t^2(\omega)$ for all sufficiently large $n$, whence
$\varliminf_{n\to\infty}M^2_{\sigma_n}(\omega)\ge\sup_{u\ge 0}\inf_{t\ge u}M^2_t(\omega)=\varliminf_{t\to\infty}M_t^2(\omega)=M_\infty^2(\omega)=M_S^2(\omega)$. This proves \eqref{eq:1-5-19:d}.

Since $P(\Omega_0)=1$, Fatou's lemma together with \eqref{eq:1-5-19:d} and \eqref{eq:1-5-19:c} yields
\[
  E\big(M_S^2\big)\;\le\;E\Big(\varliminf_{n\to\infty}M^2_{\sigma_n}\Big)\;\le\;\varliminf_{n\to\infty}E\big(M^2_{\sigma_n}\big)\;\le\;E\qv{M}_S ,
\]
which is the assertion. (If $E\qv{M}_S=\infty$ there is nothing to prove, but the argument covers this case as well.)

\rmk{The choice of the lower limit in the definition of $M_\infty^2$ is what makes Fatou's lemma applicable on $\{S=\infty\}$, where $\lim_{t\to\infty}M_t$ need not exist.}
\end{kex}

\begin{kex}{Exercise}{5.20}{}
We write $\F^X_t\defeq\sigma(X_u;\ 0\le u\le t)$ for the filtration generated by $X$, and $\sigma^2\defeq EX_1^2<\infty$. The hypotheses are understood as follows: $X$ has right-continuous paths, $X_0=0$ a.s., $EX_t^2<\infty$ for every $t\ge 0$, $X$ is a martingale with respect to $\{\F^X_t\}$, and for all $0\le s<t<\infty$
\begin{equation}\label{eq:1-5-20:a}
X_t-X_s\ \text{is independent of }\F^X_s,\qquad X_t-X_s\ \text{has the same distribution as }X_{t-s}-X_0 .
\end{equation}

\rmk{Two points of interpretation. (a) If ``independent increments'' is taken to mean that $X_{t_1}-X_{t_0},\dots,X_{t_n}-X_{t_{n-1}}$ are independent for all $0=t_0<t_1<\dots<t_n$, the first part of \eqref{eq:1-5-20:a} follows. Indeed, given $0\le s<t$ and $0<s_1<\dots<s_n=s$, the vector $(X_{s_1}-X_0,\dots,X_{s_n}-X_{s_{n-1}})$ is independent of $X_t-X_s$; since $X_0=0$ a.s., the vector $(X_0,X_{s_1},\dots,X_{s_n})$ is a.s.\ equal to a Borel function of the former, and is therefore independent of $X_t-X_s$. The events $\{X_0\in B_0,X_{s_1}\in B_1,\dots,X_{s_n}\in B_n\}$, with $B_i\in\B(\R)$ and $0<s_1<\dots<s_n=s$, form a $\pi$-system generating $\F^X_s$, and Dynkin's $\pi$-$\lambda$ theorem gives the independence of $X_t-X_s$ and $\F^X_s$. (For $s=0$, $\F^X_0=\sigma(X_0)$ consists of events of probability $0$ or $1$.) (b) The filtration $\{\F^X_t\}$ need not satisfy the usual conditions, which are in force throughout \bk{Section 1.5} and under which $\qv{X}$ is defined (\bk{Definition 5.3}). We therefore first pass to the smallest filtration containing $\{\F^X_t\}$ that satisfies the usual conditions, show in Step~1 that nothing is lost, and compute $\qv{X}$ there. The argument of Steps 2--4 shows more generally that $\qv{X}_t=t\,EX_1^2$ whenever $X\in\cM_2$ relative to a filtration $\{\F_t\}$ satisfying the usual conditions, $X_t-X_s$ is independent of $\F_s$ and distributed as $X_{t-s}$ for $0\le s<t$.}

Let $\cN\defeq\{N\in\F;\ P(N)=0\}$ and
\[
\F_t\defeq\bigcap_{u>t}\sigma\big(\F^X_u\cup\cN\big),\qquad 0\le t<\infty .
\]
Then $\{\F_t\}$ is a filtration, $\F^X_t\subseteq\F_t$, $\cN\subseteq\F_0$, and $\bigcap_{v>t}\F_v=\bigcap_{v>t}\bigcap_{u>v}\sigma(\F^X_u\cup\cN)=\F_t$; thus $\{\F_t\}$ satisfies the usual conditions (\bk{Definition 2.25}).

\emph{Step 1: for $0\le s<t$, the increment $X_t-X_s$ is independent of $\F_s$; consequently $X\in\cM_2$ relative to $\{\F_t\}$.}
Fix $0\le s<t$, $A\in\F_s$ and a bounded continuous function $g:\R\to\R$. Let $u_n\defeq s+\frac1n$, for $n$ so large that $u_n<t$. For each $u$, the class of events $A'\in\F$ for which there is an $A''\in\F^X_u$ with $P(A'\,\triangle\,A'')=0$ is a $\sigma$-field containing $\F^X_u$ and $\cN$, hence containing $\sigma(\F^X_u\cup\cN)$. Since $A\in\sigma(\F^X_{u_n}\cup\cN)$, there is an $A_n\in\F^X_{u_n}$ with $P(A\,\triangle\,A_n)=0$, and \eqref{eq:1-5-20:a} gives
\[
E\big[\ind_A\,g(X_t-X_{u_n})\big]=E\big[\ind_{A_n}\,g(X_t-X_{u_n})\big]=P(A_n)\,E\,g(X_t-X_{u_n})=P(A)\,E\,g(X_t-X_{u_n}).
\]
As $n\to\infty$ we have $X_{u_n}\to X_s$ by right-continuity of the paths, so $g(X_t-X_{u_n})\to g(X_t-X_s)$ boundedly, and the bounded convergence theorem yields
\[
E\big[\ind_A\,g(X_t-X_s)\big]=P(A)\,E\,g(X_t-X_s).
\]
The finite Borel measures $B\mapsto P(A\cap\{X_t-X_s\in B\})$ and $B\mapsto P(A)\,P[X_t-X_s\in B]$ on $\R$ thus assign the same integral to every bounded continuous function, so they coincide. This is the asserted independence.

Now $X$ is adapted to $\{\F_t\}$, and $E(X_t)=E(X_s)$ because $X$ is a martingale for $\{\F^X_t\}$. By independence, $E[X_t-X_s\,|\,\F_s]=E(X_t-X_s)=0$ a.s. Hence $X$ is a right-continuous, square-integrable martingale relative to $\{\F_t\}$ with $X_0=0$ a.s., i.e., $X\in\cM_2$ relative to $\{\F_t\}$, and $\qv{X}$ is well defined by \bk{Definition 5.3}.

\emph{Step 2: $EX_t^2=\sigma^2t$ for every $t\ge 0$.}
Let $f(t)\defeq EX_t^2$; then $f(0)=0$. For $0\le s<t$ the product $X_s(X_t-X_s)$ is integrable by the Cauchy--Schwarz inequality, $X_s$ is $\F_s$-measurable, and so $E[X_s(X_t-X_s)]=E\big[X_s\,E[X_t-X_s\,|\,\F_s]\big]=0$. Therefore
\[
f(t)=E\big[X_s+(X_t-X_s)\big]^2=f(s)+E(X_t-X_s)^2=f(s)+f(t-s),
\]
where the last equality uses the second part of \eqref{eq:1-5-20:a} and $X_0=0$ a.s. Thus $f$ is nonnegative, additive on $[0,\infty)$, and hence nondecreasing. For integers $m,n\ge 1$, additivity gives $n f(m/n)=f(m)=m f(1)$, i.e., $f(q)=q\sigma^2$ for every rational $q\ge 0$. For real $t\ge 0$ and rationals $0\le p\le t\le q$, monotonicity gives $p\sigma^2\le f(t)\le q\sigma^2$; letting $p\uparrow t$ and $q\downarrow t$ we obtain $f(t)=\sigma^2 t$. In particular
\begin{equation}\label{eq:1-5-20:b}
E(X_t-X_s)^2=f(t-s)=\sigma^2(t-s),\qquad 0\le s<t .
\end{equation}

\emph{Step 3: $\{X_t^2-\sigma^2t,\ \F_t;\ 0\le t<\infty\}$ is a martingale.}
This process is adapted and integrable. For $0\le s<t$ we write $X_t^2-X_s^2=(X_t-X_s)^2+2X_s(X_t-X_s)$; all three terms are integrable, and
\[
E\big[X_s(X_t-X_s)\,\big|\,\F_s\big]=X_s\,E\big[X_t-X_s\,\big|\,\F_s\big]=0\quad\text{a.s.}
\]
The random variable $(X_t-X_s)^2$ is integrable and, by Step~1, independent of $\F_s$, so that with \eqref{eq:1-5-20:b}
\[
E\big[X_t^2-X_s^2\,\big|\,\F_s\big]=E\big[(X_t-X_s)^2\,\big|\,\F_s\big]=E(X_t-X_s)^2=\sigma^2(t-s)\quad\text{a.s.,}
\]
which is the martingale property of $X_t^2-\sigma^2t$.

\emph{Step 4: identification of $\qv{X}$.}
The process $A_t\defeq\sigma^2t$ is deterministic, hence adapted; it satisfies $A_0=0$, its paths are nondecreasing and continuous, and $EA_t<\infty$. It is therefore an increasing process in the sense of \bk{Definition 4.4}, and it is natural by \bk{Remark 4.6}(ii), being continuous. By Step~3, $X^2-A$ is a martingale. According to \bk{Definition 5.3}, $\qv{X}$ is the unique (up to indistinguishability) adapted, natural increasing process with $\qv{X}_0=0$ a.s.\ for which $X^2-\qv{X}$ is a martingale. Consequently, almost surely,
\[
\qv{X}_t=t\,(EX_1^2)\qquad\text{for all }0\le t<\infty .
\]

\rmk{Since $\F^X_s\subseteq\F_s$ and $X_t^2-\sigma^2t$ is $\F^X_t$-measurable, the tower property shows that $X_t^2-t\,EX_1^2$ is also a martingale relative to the generated filtration $\{\F^X_t\}$ itself. The result is consistent with \bk{Example 5.4}: for the compensated Poisson process $M_t=N_t-\lambda t$ we have $EM_1^2=\Var(N_1)=\lambda$ and $\qv{M}_t=\lambda t$.}
\end{kex}

\begin{kex}{Exercise}{5.21}{}
Throughout, $\{\F_t\}$ satisfies the usual conditions (the standing assumption of the book's Section 1.5), $X\in\cM^{c,\loc}$, and $\qv{X}$ is the quadratic variation of \bk{Problem 5.17} and \bk{Definition 5.18}: an adapted, continuous, nondecreasing process with $\qv{X}_0=0$ a.s., such that $N\defeq X^2-\qv{X}\in\cM^{c,\loc}$. Since $\qv{X}$ is determined only up to indistinguishability, we fix an event $\Omega_0$ with $P(\Omega_0)=1$ on which $X_0=0$, $\qv{X}_0=0$ and $t\mapsto\qv{X}_t$ is nondecreasing; all paths of $X$ and of $\qv{X}$ are continuous. On $\Omega_0$ we set $\qv{X}_\infty\defeq\lim_{t\to\infty}\qv{X}_t\in[0,\infty]$, and for a stopping time $T$ the random variable $\qv{X}_T$ is understood with this convention on $\{T=\infty\}$, as in \bk{Problem 5.12}. For a process $Y$ and a partition $\Pi=\{t_0,\dots,t_m\}$ of $[0,t]$ we write $V^{(2)}_t(\Pi;Y)=\sum_{k=1}^m\abs{Y_{t_k}-Y_{t_{k-1}}}^2$. The three assertions to be proved are the following.
\begin{enumerate}[label=\textup{(\alph*)}]
\item For every $t>0$, $\lim_{\norm{\Pi}\to0}V^{(2)}_t(\Pi;X)=\qv{X}_t$ in probability: for every $\eps>0$, $\eta>0$ there is $\delta>0$ such that $P\big[\abs{V^{(2)}_t(\Pi;X)-\qv{X}_t}>\eps\big]<\eta$ for every partition $\Pi$ of $[0,t]$ with $\norm{\Pi}<\delta$.
\item If $T$ is a stopping time of $\{\F_t\}$ and $\qv{X}_T=0$ a.s., then $P[X_{T\wedge t}=0;\ \forall\,0\le t<\infty]=1$.
\item If, for $P$-a.e.\ $\omega$, the path $s\mapsto X_s(\omega)$ has finite first variation on every interval $[0,t]$, then $P[X_t=0;\ \forall\,0\le t<\infty]=1$.
\end{enumerate}

\emph{Step 1: a localizing sequence.} We construct stopping times $R_1\le R_2\le\cdots$ of $\{\F_t\}$ with $R_n\uparrow\infty$ a.s., such that for every $n\ge1$ the process $X^{(n)}\defeq\{X_{t\wedge R_n},\F_t;\ 0\le t<\infty\}$ belongs to $\cM_2^c$ and
\begin{equation}\label{eq:1-5-21:a}
P\big[\qv{X^{(n)}}_t=\qv{X}_{t\wedge R_n};\ \forall\,0\le t<\infty\big]=1 .
\end{equation}
By \bk{Definition 5.15} there are nondecreasing sequences $\{S_n\}$, $\{U_n\}$ of stopping times, tending to infinity a.s., such that $\{X_{t\wedge S_n},\F_t\}$ and $\{N_{t\wedge U_n},\F_t\}$ are martingales for every $n$. The $\R^2$-valued process $(X,\qv{X})$ is adapted with continuous paths, and $\Gamma_n\defeq\{(x,a)\in\R^2;\ \abs{x}\ge n\text{ or }a\ge n\}$ is closed, so by \bk{Problem 2.7}
\[
H_n\defeq\inf\{t\ge0;\ \abs{X_t}\ge n\ \text{or}\ \qv{X}_t\ge n\}
\]
is a stopping time; hence so is $R_n\defeq S_n\wedge U_n\wedge H_n$ (\bk{Lemma 2.9}). Since $\Gamma_{n+1}\subseteq\Gamma_n$, the sequence $\{H_n\}$ is nondecreasing, and so is $\{R_n\}$. For every $\omega$ and every $\tau<\infty$ the continuous functions $X_\cdot(\omega)$, $\qv{X}_\cdot(\omega)$ are bounded on $[0,\tau]$, whence $H_n(\omega)>\tau$ for all large $n$; thus $H_n\uparrow\infty$ everywhere and $R_n\uparrow\infty$ a.s.

Fix $n$. We have $X_{t\wedge R_n}=Y_{t\wedge R_n}$ with $Y_t\defeq X_{t\wedge S_n}$, a continuous martingale; applying \bk{Problem 3.24}(i) to the submartingales $Y$ and $-Y$ we see that $X^{(n)}$ is a martingale. In the same way $N^{(n)}\defeq\{N_{t\wedge R_n},\F_t;\ 0\le t<\infty\}$ is a martingale. By continuity of paths and the definition of $H_n$ we have $\abs{X_s}\le n$ and $\qv{X}_s\le n$ for $0\le s\le H_n$, provided $H_n>0$; and on $\Omega_0\cap\{H_n=0\}$ we have $X_{s\wedge R_n}=X_0=0$, $\qv{X}_{s\wedge R_n}=\qv{X}_0=0$. Consequently
\begin{equation}\label{eq:1-5-21:b}
\abs{X_{t\wedge R_n}}\le n,\qquad 0\le\qv{X}_{t\wedge R_n}\le n\qquad\text{for all }t\ge0,\ \text{on }\Omega_0 .
\end{equation}
Therefore $X^{(n)}$ is a continuous martingale, a.s.\ bounded and with $X^{(n)}_0=0$ a.s.; that is, $X^{(n)}\in\cM_2^c$. The process $A_t\defeq\qv{X}_{t\wedge R_n}=(X^{(n)}_t)^2-N^{(n)}_t$ is adapted (as a difference of adapted processes), continuous, nondecreasing on $\Omega_0$ (hence of bounded variation), with $A_0=0$ a.s., and $(X^{(n)})^2-A=N^{(n)}$ is a martingale. The uniqueness assertion of \bk{Theorem 5.13}, applied with both martingales equal to $X^{(n)}$ and recalling $\qv{X^{(n)},X^{(n)}}=\qv{X^{(n)}}$, shows that $A$ and $\qv{X^{(n)}}$ are indistinguishable, which is \eqref{eq:1-5-21:a}. (This is relation \bk{(5.4)} in the present setting.)

\emph{Step 2: proof of} (a). Fix $t>0$, $\eps>0$, $\eta>0$. The events $\{R_n<t\}$ decrease with $n$, and their intersection is contained in the null event $\{\lim_n R_n\le t\}$; hence we can fix $n$ with $P[R_n<t]<\eta/2$. \bk{Theorem 5.8}, applied to $X^{(n)}\in\cM_2^c$, provides $\delta>0$ such that
\[
P\big[\,\abs{V^{(2)}_t(\Pi;X^{(n)})-\qv{X^{(n)}}_t}>\eps\,\big]<\frac{\eta}{2}\qquad\text{for every partition $\Pi$ of $[0,t]$ with }\norm{\Pi}<\delta .
\]
On $\{R_n\ge t\}$ we have $X^{(n)}_s=X_s$ for all $0\le s\le t$, so that $V^{(2)}_t(\Pi;X^{(n)})=V^{(2)}_t(\Pi;X)$ for \emph{every} partition $\Pi$ of $[0,t]$; and, by \eqref{eq:1-5-21:a}, $\qv{X^{(n)}}_t=\qv{X}_{t\wedge R_n}=\qv{X}_t$ a.s.\ on $\{R_n\ge t\}$. Thus, for $\norm{\Pi}<\delta$,
\[
P\big[\,\abs{V^{(2)}_t(\Pi;X)-\qv{X}_t}>\eps\,\big]
\le P[R_n<t]+P\big[\,\abs{V^{(2)}_t(\Pi;X^{(n)})-\qv{X^{(n)}}_t}>\eps\,\big]<\eta ,
\]
which is (a).

\emph{Step 3: proof of} (b). Let $T$ be a stopping time with $\qv{X}_T=0$ a.s. Fix $n$. On $\Omega_0$ the path of $\qv{X}$ is nondecreasing, so $\qv{X}_{t\wedge R_n}\le\qv{X}_t$ for all $t$ and, letting $t\to\infty$, $\lim_{t\to\infty}\qv{X}_{t\wedge R_n}\le\qv{X}_\infty$. Together with \eqref{eq:1-5-21:a} this gives, a.s.,
\[
0\le\qv{X^{(n)}}_T=\qv{X}_{T\wedge R_n}\le\qv{X}_T=0
\]
(on $\{T=\infty\}$ the middle term is to be read as $\lim_{t\to\infty}\qv{X}_{t\wedge R_n}$). Hence $\qv{X^{(n)}}_T=0$ a.s., and \bk{Problem 5.12}, applied to $X^{(n)}\in\cM_2^c$ and $T$, yields an event $\Omega_n$ of probability one on which
\[
X_{T\wedge R_n\wedge t}=X^{(n)}_{T\wedge t}=0\qquad\text{for all }0\le t<\infty .
\]
Let $\Omega^*\defeq\{\lim_n R_n=\infty\}\cap\bigcap_{n\ge1}\Omega_n$, so that $P(\Omega^*)=1$. For $\omega\in\Omega^*$ and $t\ge0$ we choose $n$ with $R_n(\omega)\ge t$ and obtain $X_{T\wedge t}(\omega)=X_{T\wedge R_n\wedge t}(\omega)=0$. This proves (b).

\emph{Step 4: proof of} (c). Let $\Omega_1$ be an event of probability one such that for $\omega\in\Omega_1$ and every integer $j\ge1$ the first variation
\[
\check X_j(\omega)\defeq\sup\Big\{\sum_{k=1}^m\abs{X_{t_k}(\omega)-X_{t_{k-1}}(\omega)};\ 0=t_0\le t_1\le\dots\le t_m=j\Big\}
\]
is finite. Fix an integer $j\ge1$ and a sequence $\{\Pi_i\}_{i=1}^\infty$ of partitions of $[0,j]$ with $\norm{\Pi_i}\to0$. With the modulus of continuity $m_j(X;\delta)$ of \bk{(5.3)} we have, for every $\omega\in\Omega_1$,
\[
V^{(2)}_j(\Pi_i;X)\le m_j\big(X;\norm{\Pi_i}\big)\sum_{k}\abs{X_{t_k}-X_{t_{k-1}}}
\le m_j\big(X;\norm{\Pi_i}\big)\,\check X_j\xrightarrow[i\to\infty]{}0,
\]
because the continuous path $X_\cdot(\omega)$ is uniformly continuous on $[0,j]$. Thus $V^{(2)}_j(\Pi_i;X)\to0$ a.s., hence in probability; by (a), $V^{(2)}_j(\Pi_i;X)\to\qv{X}_j$ in probability. Limits in probability are a.s.\ unique, so $\qv{X}_j=0$ a.s. This holds for every integer $j\ge 1$, and $\qv{X}$ is nondecreasing on $\Omega_0$; therefore $\qv{X}_\infty=\lim_{j\to\infty}\qv{X}_j=0$ a.s. Now (b), applied with the stopping time $T\equiv\infty$, gives $P[X_t=0;\ \forall\,0\le t<\infty]=1$.

\rmk{The convention for $\qv{X}_T$ on $\{T=\infty\}$ cannot be dispensed with in (b): with $T\equiv\infty$ a hypothesis bearing only on $\{T<\infty\}$ would be void, whereas the conclusion asserts $X\equiv0$. Since \bk{Problem 5.11} concerns arbitrary continuous processes, (a) also shows that for $X\in\cM^{c,\loc}$ the variations of order $q>2$ over $[0,t]$ tend to zero in probability, and those of order $q<2$ tend to infinity in probability on $\{\qv{X}_t>0\}$.}
\end{kex}

\begin{kex}{Problem}{5.26}{}
We use \bk{Definition 5.15} of a local martingale; no property of the filtrations $\{\F_t\}$,
$\{\G_t\}$ beyond the independence of $\F_\infty=\sigma\bigl(\bigcup_{t\ge0}\F_t\bigr)$ and
$\G_\infty=\sigma\bigl(\bigcup_{t\ge0}\G_t\bigr)$ will be needed. Let $\{S_n\}_{n=1}^\infty$ be
a nondecreasing sequence of stopping times of $\{\F_t\}$ with $P[\lim_n S_n=\infty]=1$ such that
$X^{(n)}\defeq\{M_{t\wedge S_n},\F_t;\ 0\le t<\infty\}$ is a martingale for every $n$, and let
$\{T_n\}_{n=1}^\infty$ be a nondecreasing sequence of stopping times of $\{\G_t\}$ with
$P[\lim_n T_n=\infty]=1$ such that $Y^{(n)}\defeq\{N_{t\wedge T_n},\G_t;\ 0\le t<\infty\}$ is a
martingale for every $n$.

\emph{Step 1: conditioning on $\cH_s$.}

\begin{lem}
Let $s\ge0$, let $X$ be an integrable, $\F_\infty$-measurable random variable and $Y$ an
integrable, $\G_\infty$-measurable random variable. Then $XY$ is integrable and
\begin{equation}\label{eq:1-5-26:a}
E[XY\,|\,\cH_s]=E[X\,|\,\F_s]\,E[Y\,|\,\G_s]\qquad\text{a.s.}
\end{equation}
In particular (take $Y\equiv1$, resp.\ $X\equiv 1$), $E[X\,|\,\cH_s]=E[X\,|\,\F_s]$ and
$E[Y\,|\,\cH_s]=E[Y\,|\,\G_s]$ a.s.
\end{lem}

\begin{proof}
Since $X$ and $Y$ are independent, $E\abs{XY}=E\abs{X}\,E\abs{Y}<\infty$. Put
$X'\defeq E[X\,|\,\F_s]$ and $Y'\defeq E[Y\,|\,\G_s]$; these are integrable, measurable with
respect to $\F_s\subseteq\F_\infty$ and $\G_s\subseteq\G_\infty$ respectively, hence
independent, so that $X'Y'$ is integrable, and it is $\cH_s$-measurable. Let $A\in\F_s$ and
$B\in\G_s$. The random variables $X\ind_A$ and $Y\ind_B$ are independent and integrable, and so
are $X'\ind_A$ and $Y'\ind_B$; therefore
\[
E[XY\ind_{A\cap B}]=E[X\ind_A]\,E[Y\ind_B]=E[X'\ind_A]\,E[Y'\ind_B]=E[X'Y'\ind_{A\cap B}].
\]
The class $\mathscr{D}$ of sets $C\in\cH_s$ with $E[XY\ind_C]=E[X'Y'\ind_C]$ contains $\Omega$,
is closed under proper differences (by linearity) and under increasing countable unions (by
dominated convergence); it is a Dynkin system. It contains the class
$\{A\cap B;\ A\in\F_s,\ B\in\G_s\}$, which is closed under finite intersections and generates
$\cH_s$. By Dynkin's $\pi$-$\lambda$ theorem, $\mathscr{D}=\cH_s$, and \eqref{eq:1-5-26:a}
follows.
\end{proof}

\emph{Step 2: $M$, $N$ and $MN$ are local martingales relative to $\{\cH_t\}$.} Because
$\F_t\subseteq\cH_t$ and $\G_t\subseteq\cH_t$, the processes $M$, $N$, $MN$ are adapted to
$\{\cH_t\}$, and every stopping time of $\{\F_t\}$ or of $\{\G_t\}$ is a stopping time of
$\{\cH_t\}$. Fix $n\ge1$ and $0\le s<t<\infty$. The random variable $X^{(n)}_t=M_{t\wedge S_n}$
is integrable and $\F_t$-measurable, hence $\F_\infty$-measurable, and the lemma gives
\[
E\bigl[X^{(n)}_t\,\big|\,\cH_s\bigr]=E\bigl[X^{(n)}_t\,\big|\,\F_s\bigr]=X^{(n)}_s\qquad\text{a.s.}
\]
Thus $\{M_{t\wedge S_n},\cH_t;\ 0\le t<\infty\}$ is a martingale for every $n$, and
$\{M_t,\cH_t\}$ is a continuous local martingale, with the localizing sequence $\{S_n\}$. In the same
way, $\{N_{t\wedge T_n},\cH_t;\ 0\le t<\infty\}$ is a martingale for every $n$ and
$\{N_t,\cH_t\}$ is a continuous local martingale.

For the product, consider $Z^{(n)}_t\defeq X^{(n)}_tY^{(n)}_t=M_{t\wedge S_n}N_{t\wedge T_n}$.
This process is continuous and adapted to $\{\cH_t\}$, and by the lemma, applied to
$X=X^{(n)}_t$ and $Y=Y^{(n)}_t$, each $Z^{(n)}_t$ is integrable and, for $0\le s<t$,
\[
E\bigl[Z^{(n)}_t\,\big|\,\cH_s\bigr]
=E\bigl[X^{(n)}_t\,\big|\,\F_s\bigr]\,E\bigl[Y^{(n)}_t\,\big|\,\G_s\bigr]
=X^{(n)}_sY^{(n)}_s=Z^{(n)}_s\qquad\text{a.s.}
\]
Hence $\{Z^{(n)}_t,\cH_t;\ 0\le t<\infty\}$ is a continuous martingale. Now
$R_n\defeq S_n\wedge T_n$ is a stopping time of $\{\cH_t\}$ (\bk{Lemma 2.9}), the sequence
$\{R_n\}$ is nondecreasing, and $\lim_nR_n=\infty$ on the event
$\{\lim_nS_n=\infty\}\cap\{\lim_nT_n=\infty\}$ of probability one. By \bk{Problem 3.24}(i),
applied to the right-continuous submartingales $Z^{(n)}$ and $-Z^{(n)}$ and the stopping time
$R_n$ (that problem imposes no condition on the filtration), the process
$\{Z^{(n)}_{t\wedge R_n},\cH_t;\ 0\le t<\infty\}$ is a martingale. Since $t\wedge R_n\le S_n$
and $t\wedge R_n\le T_n$,
\[
Z^{(n)}_{t\wedge R_n}=M_{t\wedge R_n\wedge S_n}N_{t\wedge R_n\wedge T_n}
=M_{t\wedge R_n}N_{t\wedge R_n}.
\]
Therefore $\{M_{t\wedge R_n}N_{t\wedge R_n},\cH_t;\ 0\le t<\infty\}$ is a martingale for every
$n$, i.e., $\{M_tN_t,\cH_t\}$ is a continuous local martingale with the localizing sequence
$\{R_n\}$.

\emph{Step 3: $\{\wt\cH_t\}$ satisfies the usual conditions.} Write
$\cH^{\cN}_s\defeq\sigma(\cH_s\cup\cN)$, so that $\wt\cH_t=\bigcap_{s>t}\cH^{\cN}_s$. Each
$\wt\cH_t$ is a sub-$\sigma$-field of $\F$, and for $t<t'$ we have
$\wt\cH_t\subseteq\cH^{\cN}_{u}\subseteq\cH^{\cN}_s$ whenever $t<u\le t'<s$ (the family
$\{\cH^{\cN}_s\}$ being nondecreasing), whence $\wt\cH_t\subseteq\wt\cH_{t'}$; so
$\{\wt\cH_t\}$ is a filtration, and
\begin{equation}\label{eq:1-5-26:b}
\cH_t\subseteq\cH^{\cN}_t\subseteq\wt\cH_t\subseteq\cH^{\cN}_s\qquad\text{for all }0\le t<s.
\end{equation}
It is right-continuous: the inclusion $\wt\cH_t\subseteq\bigcap_{u>t}\wt\cH_u$ is clear, and if
$A\in\bigcap_{u>t}\wt\cH_u$ and $s>t$, then choosing $u\in(t,s)$ we get
$A\in\wt\cH_u\subseteq\cH^{\cN}_s$ by \eqref{eq:1-5-26:b}; as $s>t$ is arbitrary,
$A\in\wt\cH_t$. Finally $\cN\subseteq\cH^{\cN}_s$ for every $s>0$, so $\cN\subseteq\wt\cH_0$.
By \bk{Definition 2.25}, $\{\wt\cH_t\}$ satisfies the usual conditions.

\emph{Step 4: passage from $\{\cH_t\}$ to $\{\wt\cH_t\}$.}

\begin{claim}
If $\{Z_t,\cH_t;\ 0\le t<\infty\}$ is a right-continuous martingale, then so is
$\{Z_t,\wt\cH_t;\ 0\le t<\infty\}$.
\end{claim}

\begin{proof}
By \eqref{eq:1-5-26:b}, $Z$ is adapted to $\{\wt\cH_t\}$. Fix $u\ge0$. The class
\[
\mathscr{E}_u\defeq\bigl\{C\in\F;\ P(C\,\triangle\,D)=0\text{ for some }D\in\cH_u\bigr\}
\]
contains $\cH_u$ and $\cN$ (take $D=\emptyset$), and it is a $\sigma$-field: it is closed
under complementation because $C^c\triangle D^c=C\triangle D$, and under countable unions
because $\bigl(\bigcup_kC_k\bigr)\triangle\bigl(\bigcup_kD_k\bigr)\subseteq
\bigcup_k(C_k\triangle D_k)$. Hence $\cH^{\cN}_u\subseteq\mathscr{E}_u$. Consequently, if
$u<t$ and $C\in\cH^{\cN}_u$, then with $D\in\cH_u$ such that $P(C\triangle D)=0$,
\begin{equation}\label{eq:1-5-26:c}
E[Z_t\ind_C]=E[Z_t\ind_D]=E[Z_u\ind_D]=E[Z_u\ind_C].
\end{equation}
Now let $0\le s<t$ and $C\in\wt\cH_s$. For every integer $k$ with $u_k\defeq s+1/k<t$ we have
$C\in\cH^{\cN}_{u_k}$ by \eqref{eq:1-5-26:b}, so \eqref{eq:1-5-26:c} gives
$E[Z_t\ind_C]=E[Z_{u_k}\ind_C]$. As $k\to\infty$, $Z_{u_k}\to Z_s$ pointwise by
right-continuity; moreover the random variables $Z_{u_k}=E[Z_t\,|\,\cH_{u_k}]$ are conditional
expectations of the single integrable random variable $Z_t$, hence form a uniformly integrable
family. Therefore $Z_{u_k}\to Z_s$ in $L^1$, and we obtain $E[Z_t\ind_C]=E[Z_s\ind_C]$ for
every $C\in\wt\cH_s$. Since $Z_s$ is $\wt\cH_s$-measurable and integrable, this is
$E[Z_t\,|\,\wt\cH_s]=Z_s$ a.s.
\end{proof}

By \eqref{eq:1-5-26:b}, every stopping time of $\{\cH_t\}$ is a stopping time of
$\{\wt\cH_t\}$; in particular so are $S_n$, $T_n$ and $R_n$. By Step 2 the continuous
processes $\{M_{t\wedge S_n}\}$, $\{N_{t\wedge T_n}\}$ and $\{M_{t\wedge R_n}N_{t\wedge R_n}\}$
are martingales relative to $\{\cH_t\}$, hence, by the claim, relative to $\{\wt\cH_t\}$.
Thus $M$, $N$ and $MN$ are continuous local martingales relative to $\{\wt\cH_t\}$, with the
localizing sequences $\{S_n\}$, $\{T_n\}$ and $\{R_n\}$, respectively.

\emph{Step 5: $\qv{M,N}\equiv0$.} The cross-variation is defined in \bk{Problem 5.17} and
\bk{Definition 5.18} for two elements of $\cM^{c,\loc}$ relative to one and the same filtration
satisfying the usual conditions; here this filtration is $\{\wt\cH_t\}$. Assume first that
$M_0=N_0=0$ a.s. By Step 4, $M$, $N$ and $MN$ belong to $\cM^{c,\loc}$ relative to
$\{\wt\cH_t\}$. The process $A\equiv0$ is adapted, continuous, of bounded variation, vanishes
at $t=0$, and $MN-A=MN\in\cM^{c,\loc}$. By the uniqueness assertion of \bk{Problem 5.17},
$\qv{M,N}$ is indistinguishable from $A$, i.e.,
\[
P\bigl[\qv{M,N}_t=0\ \text{ for all }0\le t<\infty\bigr]=1 .
\]

\rmk{The statement does not assume $M_0=N_0=0$, whereas $\qv{M,N}$ is defined in the book only
for processes in $\cM^{c,\loc}$. For general initial values one sets, as usual,
$\qv{M,N}\defeq\qv{M-M_0,N-N_0}$, and the conclusion persists. Indeed, $M_0=M_{0\wedge S_1}$ is
integrable and $\F_0$-measurable, so $\{M_{t\wedge S_n}-M_0,\F_t\}$ is a martingale for every
$n$; thus $\{M_t-M_0,\F_t\}$ is a continuous local martingale vanishing at $t=0$, and likewise
$\{N_t-N_0,\G_t\}$. Applying Steps 2, 4 and 5 to this pair we get $\qv{M-M_0,N-N_0}\equiv0$.}
\end{kex}


%% ========================================================================
\solchapter{2}{Brownian Motion}

\begin{center}
\begin{minipage}{0.86\textwidth}
\small
The 35 problems and exercises of Chapter 2 for which the book prints no solution: 2.7, 2.10, 3.3, 4.1, 4.5, 4.8, 4.12, 4.14, 5.4, 5.5, 5.17, 5.18, 5.21, 5.22, 6.1, 6.11, 6.21, 7.11, 7.12, 7.13, 7.14, 8.2, 8.4, 8.6, 8.7, 8.9, 8.11, 8.13, 8.14, 8.17, 8.18, 8.19, 8.20, 9.17, 9.21.  For 5 more, 5.9, 6.9, 7.6, 9.5, 9.22, the book prints only a hint or solves only some of the parts; what it leaves out is supplied here.  The book's own ``Solutions to Selected Problems'' fully covers 1.4, 2.3, 2.4, 2.5, 2.9, 4.2, 4.11, 4.16, 5.2, 5.7, 7.1, 7.3, 7.4, 7.5, 7.18, 7.19, 8.8, 8.12, 9.3, 9.8, 9.11, 9.26, which are therefore omitted.
\end{minipage}
\end{center}

\bigskip
\hrule
\bigskip

\solsection{2.2}{First Construction of Brownian Motion}

\begin{kex}{Exercise}{2.7}{}
We use the hint in the following form, which is all that the argument needs. Let us say that a set $E\subseteq\R^{[0,\infty)}$ is \emph{determined by} a set $S\subseteq[0,\infty)$ if
\begin{equation}\label{eq:2-2-07:a}
\omega\in E,\quad \omega'\in\R^{[0,\infty)},\quad \omega'(t)=\omega(t)\ \text{ for every } t\in S
\qquad\Longrightarrow\qquad \omega'\in E ,
\end{equation}
and let $\G$ be the collection of all sets $E\subseteq\R^{[0,\infty)}$ which are determined by some countable (finite, possibly empty, or countably infinite) set $S\subseteq[0,\infty)$.

\begin{lem}
$\G$ is a $\sigma$-field containing the field $\cC$ of cylinder sets; consequently $\B(\R^{[0,\infty)})\subseteq\G$.
\end{lem}

\begin{proof}
\emph{Cylinder sets.} A cylinder set $C=\{\omega;\ (\omega(t_1),\dots,\omega(t_n))\in A\}$ as in \bk{(2.1)} is determined by the finite set $S=\{t_1,\dots,t_n\}$: if $\omega\in C$ and $\omega'=\omega$ on $S$, then $(\omega'(t_1),\dots,\omega'(t_n))=(\omega(t_1),\dots,\omega(t_n))\in A$. Thus $\cC\subseteq\G$; in particular $\R^{[0,\infty)}\in\G$.

\emph{Complements.} Let $E\in\G$ be determined by the countable set $S$; we claim that $E^c$ is determined by the same $S$. Let $\omega\in E^c$ and let $\omega'$ agree with $\omega$ on $S$. If $\omega'$ were in $E$, then \eqref{eq:2-2-07:a}, applied with the roles of $\omega$ and $\omega'$ interchanged (agreement on $S$ is a symmetric relation), would give $\omega\in E$, which is false. Hence $\omega'\in E^c$.

\emph{Countable unions.} Let $E_n\in\G$ be determined by the countable set $S_n$, $n\ge1$, and put $S\defeq\bigcup_{n\ge1}S_n$, a countable set. If $\omega\in\bigcup_n E_n$ and $\omega'=\omega$ on $S$, choose $n$ with $\omega\in E_n$; since $\omega'=\omega$ on $S_n\subseteq S$, we get $\omega'\in E_n\subseteq\bigcup_n E_n$. Hence $\bigcup_n E_n\in\G$.

So $\G$ is a $\sigma$-field containing $\cC$, and $\B(\R^{[0,\infty)})$, being the smallest such $\sigma$-field, is contained in $\G$.
\end{proof}

Now let $E\in\B(\R^{[0,\infty)})$ with $E\subseteq C[0,\infty)$, and suppose that $E$ is not empty; fix $\omega\in E$. By the Lemma, $E$ is determined by some countable set $S\subseteq[0,\infty)$. Since $[0,\infty)$ is uncountable, there exists $t_0\in[0,\infty)\setminus S$. Define $\omega'\in\R^{[0,\infty)}$ by
\[
\omega'(t)\defeq
\begin{cases}
\omega(t), & t\neq t_0,\\
\omega(t_0)+1, & t=t_0 .
\end{cases}
\]
Then $\omega'=\omega$ on $S$, so $\omega'\in E$ by \eqref{eq:2-2-07:a}. But $\omega'$ is not continuous at $t_0$: the point $t_0$ is not isolated in $[0,\infty)$, and for any sequence $s_k\to t_0$ with $s_k\neq t_0$ the continuity of $\omega$ (recall $\omega\in E\subseteq C[0,\infty)$) gives
\[
\lim_{k\to\infty}\omega'(s_k)=\lim_{k\to\infty}\omega(s_k)=\omega(t_0)\neq\omega(t_0)+1=\omega'(t_0).
\]
Thus $\omega'\in E\setminus C[0,\infty)$, contradicting $E\subseteq C[0,\infty)$. Therefore $E=\emptyset$.

\rmk{Since $C[0,\infty)$ is a nonempty subset of itself, the result shows in particular that $C[0,\infty)\notin\B(\R^{[0,\infty)})$, as announced in the text preceding the exercise. The representation in the hint is a sharper form of the Lemma: each $E\in\B(\R^{[0,\infty)})$ can indeed be written as $\{\omega;\ (\omega(t_1),\omega(t_2),\ldots)\in A\}$ with $A\in\B(\R\times\R\times\cdots)$, and every set of this form is determined by the countable set $\{t_1,t_2,\ldots\}$. Only this last property was needed, so we have proved it directly rather than the full representation.}
\end{kex}

\begin{kex}{Problem}{2.10}{}
We show that the identity holds with
\begin{equation}\label{eq:2-2-10:C}
C_n \;=\; (2n-1)(2n-3)\cdots 3\cdot 1 \;=\; \frac{(2n)!}{2^n\,n!},
\end{equation}
which is positive and depends neither on $s$ nor on $t$.

\emph{Step 1: reduction to the standard normal law.}
Fix $0\le s<t$ and put $Z\defeq (B_t-B_s)/\sqrt{t-s}$. Since $B_t-B_s$ has the density
$x\mapsto (2\pi(t-s))^{-1/2}e^{-x^2/2(t-s)}$, the change of variable $x=\sqrt{t-s}\,z$ shows that $Z$ has
the standard normal density $\varphi(z)\defeq (2\pi)^{-1/2}e^{-z^2/2}$, $z\in\R$. Consequently, for every
positive integer $n$,
\begin{equation}\label{eq:2-2-10:scale}
E\abs{B_t-B_s}^{2n} \;=\; (t-s)^n\,E\abs{Z}^{2n} \;=\; (t-s)^n\, m_n,
\qquad m_n\defeq\int_{-\infty}^{\infty} z^{2n}\varphi(z)\,dz\in[0,\infty],
\end{equation}
and $m_n$ depends only on $n$. It remains to prove that $m_n$ equals the number $C_n$
of \eqref{eq:2-2-10:C}; we set also $m_0\defeq\int_{-\infty}^\infty\varphi(z)\,dz=1$.

\emph{Step 2: the recursion $m_n=(2n-1)\,m_{n-1}$ for $n\ge 1$.}
We have $\varphi'(z)=-z\varphi(z)$. Fix $n\ge 1$ and $a>0$. Integrating by parts on $[-a,a]$, with the
continuously differentiable functions $z\mapsto z^{2n-1}$ and $z\mapsto -\varphi(z)$, we obtain
\begin{equation}\label{eq:2-2-10:parts}
\int_{-a}^{a} z^{2n}\varphi(z)\,dz
=\int_{-a}^{a} z^{2n-1}\bigl(-\varphi'(z)\bigr)\,dz
=-2a^{2n-1}\varphi(a)+(2n-1)\int_{-a}^{a} z^{2n-2}\varphi(z)\,dz ,
\end{equation}
where we have used that $z^{2n-1}\varphi(z)$ is an odd function of $z$. As $a\to\infty$, the two integrals in
\eqref{eq:2-2-10:parts} have nonnegative integrands and therefore increase to $m_n$ and $m_{n-1}$,
respectively, by monotone convergence; moreover $a^{2n-1}\varphi(a)\to 0$, because
$e^{a^2/2}\ge (a^2/2)^{n}/n!$ gives $a^{2n-1}\varphi(a)\le 2^n n!\,(2\pi)^{-1/2}a^{-1}$. Hence
\[
m_n=(2n-1)\,m_{n-1}\qquad\text{in }[0,\infty],\quad n\ge 1 .
\]

\emph{Step 3: conclusion.}
Since $m_0=1$, induction on $n$ in the recursion of Step 2 shows that $m_n$ is finite for every $n$ and that
\[
m_n=(2n-1)(2n-3)\cdots 3\cdot 1
=\frac{(2n)(2n-1)(2n-2)\cdots 2\cdot 1}{(2n)(2n-2)\cdots 2}
=\frac{(2n)!}{2^n\,n!}=C_n>0 .
\]
Substituting into \eqref{eq:2-2-10:scale}, and recalling $t-s=\abs{t-s}$, we conclude
\[
E\abs{B_t-B_s}^{2n}=C_n\abs{t-s}^n,\qquad 0\le s<t,\ n=1,2,\dots
\]
\rmk{For $n=1$ we recover $C_1=1$, the variance, and $C_2=3$ is the familiar fourth moment of the standard
normal law. With $n\ge 2$ the identity is condition \bk{(2.7)} of \bk{Theorem 2.8}, with $\alpha=2n$,
$\beta=n-1>0$ and $C=C_n$; this is how the problem is used in the proof of \bk{Corollary 2.11} and in
\bk{Remark 2.12}, where it yields H\"older continuity of every exponent
$\gamma<\beta/\alpha=(n-1)/2n$, hence, $n$ being arbitrary, of every exponent $\gamma<1/2$.}
\end{kex}

\solsection{2.3}{Second Construction of Brownian Motion}

\begin{kex}{Exercise}{3.3}{}
Throughout, $\{\xi^{(n)}_k;\ k\in I(n),\ n\ge 0\}$ are the independent standard normal random variables of \bk{Section 3}, and $B^{(n)}$ is defined by the recursion preceding \bk{(3.2)}: $B^{(n)}_0=0$, $B^{(0)}_1=\xi^{(0)}_1$,
\begin{equation}\label{eq:2-3-03:rec}
\begin{aligned}
B^{(n)}_{j/2^{n-1}}&=B^{(n-1)}_{j/2^{n-1}}, &&0\le j\le 2^{n-1},\\
B^{(n)}_{k/2^n}&=\tfrac12\Bigl(B^{(n-1)}_{(k-1)/2^n}+B^{(n-1)}_{(k+1)/2^n}\Bigr)+2^{-(n+1)/2}\,\xi^{(n)}_k, &&k\in I(n),
\end{aligned}
\end{equation}
for $n\ge1$, with linear interpolation between the points $k/2^n$. \bk{Lemma 3.1} provides an event $\wt\Omega\in\F$ with $P(\wt\Omega)=1$ on which $B^{(n)}$ converges uniformly on $[0,1]$ to a continuous function $B$. We set $B_t(\omega)\defeq 0$ for all $t\in[0,1]$ when $\omega\notin\wt\Omega$; then $B_t=\lim_n B^{(n)}_t\ind_{\wt\Omega}$ is a random variable for every $t$, \emph{every} path of $B$ is continuous, and $B_0=0$ identically. From the first relation in \eqref{eq:2-3-03:rec}, $B^{(m)}_{k/2^n}=B^{(n)}_{k/2^n}$ for all $m\ge n$ and $0\le k\le 2^n$, and therefore
\begin{equation}\label{eq:2-3-03:dy}
B_{k/2^n}=B^{(n)}_{k/2^n}\quad\text{on }\wt\Omega,\qquad 0\le k\le 2^n,\ n\ge 0 .
\end{equation}
We shall use two standard facts about Gaussian vectors: a linear image of a random vector with independent normal components is a Gaussian vector, and the components of a Gaussian vector are independent as soon as they are pairwise uncorrelated.

\textbf{(a)} We argue by induction on $n$, proving in addition that the increments of $B^{(n)}$ in question are measurable with respect to $\G_n\defeq\sigma\bigl(\xi^{(m)}_k;\ k\in I(m),\ 0\le m\le n\bigr)$. For $n=0$ the only increment is $B^{(0)}_1-B^{(0)}_0=\xi^{(0)}_1$, and the claim holds. Let $n\ge1$ and suppose the claim holds for $n-1$: the random variables
\[
D_j\defeq B^{(n-1)}_{j/2^{n-1}}-B^{(n-1)}_{(j-1)/2^{n-1}},\qquad 1\le j\le 2^{n-1},
\]
are independent, normal with mean zero and variance $2^{-(n-1)}$, and $\G_{n-1}$-measurable. Every $k\in I(n)$ is of the form $k=2j-1$ with $1\le j\le 2^{n-1}$, and then $(k-1)/2^n=(j-1)/2^{n-1}$, $(k+1)/2^n=j/2^{n-1}$. With $\sigma\defeq 2^{-(n+1)/2}$, the relations \eqref{eq:2-3-03:rec} give, for the two increments of $B^{(n)}$ over the halves of $[(j-1)/2^{n-1},j/2^{n-1}]$,
\begin{align*}
X_{2j-1}&\defeq B^{(n)}_{(2j-1)/2^n}-B^{(n)}_{(2j-2)/2^n}\\
&=\tfrac12\Bigl(B^{(n-1)}_{(j-1)/2^{n-1}}+B^{(n-1)}_{j/2^{n-1}}\Bigr)+\sigma\xi^{(n)}_{2j-1}-B^{(n-1)}_{(j-1)/2^{n-1}}
=\tfrac12D_j+\sigma\xi^{(n)}_{2j-1},\\
X_{2j}&\defeq B^{(n)}_{2j/2^n}-B^{(n)}_{(2j-1)/2^n}\\
&=B^{(n-1)}_{j/2^{n-1}}-\tfrac12\Bigl(B^{(n-1)}_{(j-1)/2^{n-1}}+B^{(n-1)}_{j/2^{n-1}}\Bigr)-\sigma\xi^{(n)}_{2j-1}
=\tfrac12D_j-\sigma\xi^{(n)}_{2j-1}.
\end{align*}
These are $\G_n$-measurable. The family $\{\xi^{(n)}_k\}_{k\in I(n)}$ is independent of $\G_{n-1}$, and its members are independent of one another; together with the induction hypothesis this shows that the $2^n$ random variables
\[
D_1,\dots,D_{2^{n-1}},\ \xi^{(n)}_1,\xi^{(n)}_3,\dots,\xi^{(n)}_{2^n-1}
\]
are independent and normal with mean zero. Consequently $(X_1,\dots,X_{2^n})$, a linear image of this vector, is a Gaussian vector with mean zero. Its covariances are computed from $E[D_j^2]=2^{-(n-1)}$, $E[(\xi^{(n)}_k)^2]=1$ and the vanishing of all mixed second moments of the above $2^n$ variables:
\begin{align*}
E\bigl[X_{2j-1}^2\bigr]=E\bigl[X_{2j}^2\bigr]&=\tfrac14\,2^{-(n-1)}+\sigma^2=2^{-(n+1)}+2^{-(n+1)}=2^{-n},\\
E\bigl[X_{2j-1}X_{2j}\bigr]&=\tfrac14\,2^{-(n-1)}-\sigma^2=0,
\end{align*}
and $E[X_kX_l]=0$ when $X_k$, $X_l$ belong to different indices $j\ne j'$, since then they are linear combinations of disjoint sets of the independent variables. Thus $X_1,\dots,X_{2^n}$ are uncorrelated components of a Gaussian vector, hence independent, each normal with mean zero and variance $2^{-n}$. This completes the induction.

\textbf{(b)} Let $0=t_0<t_1<\dots<t_n\le1$ be dyadic rationals. Choose $N\ge0$ so large that $t_j=k_j/2^N$ with integers $0=k_0<k_1<\dots<k_n\le 2^N$. By \eqref{eq:2-3-03:dy}, on $\wt\Omega$,
\[
B_{t_j}-B_{t_{j-1}}=B^{(N)}_{k_j/2^N}-B^{(N)}_{k_{j-1}/2^N}
=\sum_{k=k_{j-1}+1}^{k_j}\Bigl(B^{(N)}_{k/2^N}-B^{(N)}_{(k-1)/2^N}\Bigr),\qquad 1\le j\le n .
\]
The $n$ sums on the right involve pairwise disjoint blocks of the increments of part (a) (with $N$ in place of $n$), which are independent $N(0,2^{-N})$ random variables. Hence these sums are independent, and the $j$-th is normal with mean zero and variance $(k_j-k_{j-1})2^{-N}=t_j-t_{j-1}$. Since $P(\wt\Omega)=1$, the vector $(B_{t_j}-B_{t_{j-1}})_{j=1}^n$ agrees almost surely with the vector of these sums, so it has the same distribution; this is the assertion.

\textbf{(c)} Let $0=t_0<t_1<\dots<t_n\le1$ be arbitrary. For $m\ge1$ set $t^{(m)}_0\defeq0$ and
\[
t^{(m)}_j\defeq 2^{-m}\lceil 2^mt_j\rceil,\qquad 1\le j\le n,
\]
so that $t^{(m)}_j$ is a dyadic rational with $t_j\le t^{(m)}_j<t_j+2^{-m}$, $t^{(m)}_j\le1$ (because $2^mt_j\le 2^m$), and $t^{(m)}_1>0$. If $m$ is so large that $2^{-m}<\min_{1\le j\le n}(t_j-t_{j-1})$, then $t^{(m)}_{j-1}<t_{j-1}+2^{-m}<t_j\le t^{(m)}_j$, and so $0=t^{(m)}_0<t^{(m)}_1<\dots<t^{(m)}_n\le1$. For such $m$, part (b) and the form of the characteristic function of independent centred normal random variables give, for all $\lambda_1,\dots,\lambda_n\in\R$,
\begin{equation}\label{eq:2-3-03:cf}
E\biggl[\exp\Bigl\{i\sum_{j=1}^n\lambda_j\bigl(B_{t^{(m)}_j}-B_{t^{(m)}_{j-1}}\bigr)\Bigr\}\biggr]
=\prod_{j=1}^n\exp\Bigl\{-\tfrac12\lambda_j^2\bigl(t^{(m)}_j-t^{(m)}_{j-1}\bigr)\Bigr\}.
\end{equation}
Let $m\to\infty$. Since $t^{(m)}_j\to t_j$ and every path of $B$ is continuous, $B_{t^{(m)}_j}\to B_{t_j}$ pointwise on $\Omega$ for each $j$; the integrand on the left of \eqref{eq:2-3-03:cf} has modulus one, so by the bounded convergence theorem the left-hand side converges to the corresponding expression with $t_j$ in place of $t^{(m)}_j$. The right-hand side converges trivially. Hence
\[
E\biggl[\exp\Bigl\{i\sum_{j=1}^n\lambda_j\bigl(B_{t_j}-B_{t_{j-1}}\bigr)\Bigr\}\biggr]
=\prod_{j=1}^n\exp\Bigl\{-\tfrac12\lambda_j^2(t_j-t_{j-1})\Bigr\},\qquad(\lambda_1,\dots,\lambda_n)\in\R^n,
\]
which is \bk{(3.5)}. The right-hand side is the characteristic function of a vector of $n$ independent normal random variables with means zero and variances $t_j-t_{j-1}$; by the uniqueness theorem for characteristic functions on $\R^n$, the vector $(B_{t_j}-B_{t_{j-1}})_{j=1}^n$ has this distribution. This proves (c).

\textbf{Proof of \bk{Theorem 3.2}.} We verify \bk{Definition 1.1} on $[0,1]$ for $\{B_t,\F^B_t;\ 0\le t\le1\}$. The process $B$ is continuous, adapted to its own filtration, and $B_0=0$. Let $0\le s<t\le1$. Applying (c) to the points $0<t$ (if $s=0$) or $0<s<t$ (if $s>0$), we see that $B_t-B_s$ is normal with mean zero and variance $t-s$. It remains to show that $B_t-B_s$ is independent of $\F^B_s=\sigma(B_u;\ 0\le u\le s)$. If $s=0$ this is clear, because $B_0\equiv0$ makes $\F^B_0$ trivial. Let $s>0$ and consider the collection $\cC$ of all events of the form
\[
\bigl\{(B_{s_1},\dots,B_{s_m})\in\Gamma\bigr\},\qquad m\ge1,\quad 0<s_1<\dots<s_m=s,\quad\Gamma\in\B(\R^m).
\]
Each such event is independent of $B_t-B_s$: by (c), applied to $0<s_1<\dots<s_m<t$, the random variable $B_t-B_{s_m}$ is independent of the vector $(B_{s_1}-B_0,\,B_{s_2}-B_{s_1},\dots,B_{s_m}-B_{s_{m-1}})$, and $(B_{s_1},\dots,B_{s_m})$ is a linear, hence Borel, function of that vector (recall $B_0=0$). The collection $\cC$ is closed under finite intersections, since two events of this form can both be written in terms of the union of their two sets of time points; and $\cC$ generates $\F^B_s$, because it contains $\{B_u\in A\}$ for every $u\in(0,s]$, $A\in\B(\R)$ (take the time points $u,s$, or $s$ alone if $u=s$), while $\sigma(B_0)$ is trivial. For fixed $A\in\B(\R)$, the events $F\in\F$ with $P[F\cap\{B_t-B_s\in A\}]=P[F]\,P[B_t-B_s\in A]$ form a Dynkin system containing $\cC$; by Dynkin's $\pi$-$\lambda$ theorem it contains $\F^B_s$. Therefore $B_t-B_s$ is independent of $\F^B_s$, and $\{B_t,\F^B_t;\ 0\le t\le1\}$ is a Brownian motion on $[0,1]$. Neither \bk{(3.4)} nor the completeness of the Haar system has been used.

\rmk{\bk{Lemma 3.1} defines $B$ only for almost every $\omega$, whereas \bk{Definition 1.1} requires a continuous process; this is why we put $B\equiv0$ off $\wt\Omega$. The set $\wt\Omega$ of the proof of \bk{Lemma 3.1} is the event that $\max_{k\in I(n)}\abs{\xi^{(n)}_k}\le n$ for all but finitely many $n$, and so belongs to $\F$.}

\rmk{The identity \bk{(3.4)}, and with it the completeness of the Haar system, now follows; this is the sense of the reference to this exercise made before \bk{(3.4)}. Fix $s,t\in[0,1]$. By \bk{(3.2)} and the orthonormality of the $\xi^{(m)}_k$ in $L^2(P)$,
\[
E\bigl[B^{(M)}_sB^{(M)}_t\bigr]=\sum_{m=0}^M\sum_{k\in I(m)}S^{(m)}_k(s)S^{(m)}_k(t).
\]
For each $M$, $B^{(M)}_t$ is a convex combination of two of the variables $B^{(M)}_{k/2^M}$, which by (a) have variances $k/2^M\le1$; hence $\sum_{m\le M}\sum_{k\in I(m)}S^{(m)}_k(t)^2=E[(B^{(M)}_t)^2]\le1$ for every $M$. Thus the orthogonal series $\sum_m\sum_k\xi^{(m)}_kS^{(m)}_k(t)$ converges in $L^2(P)$, and its $L^2$-limit coincides a.s.\ with its a.s.\ limit $B_t$. Passing to the limit in the display, and using $E[B_sB_t]=E[B_{s\wedge t}^2]=s\wedge t$ from (c), we obtain \bk{(3.4)}. With $s=t$ it reads $\sum_{n}\sum_{k}\langle\ind_{[0,t]},H^{(n)}_k\rangle^2=\norm{\ind_{[0,t]}}^2$, i.e.\ equality holds in Bessel's inequality for $\ind_{[0,t]}$, so that $\ind_{[0,t]}$ lies in the closed linear span of the (orthonormal) Haar functions; since the linear combinations of the functions $\ind_{[0,t]}$, $0\le t\le1$, are dense in $L^2[0,1]$, the Haar system is complete.}
\end{kex}

\solsection{2.4}{The Space $C[0,\infty)$, Weak Convergence, and the Wiener Measure}

\begin{kex}{Problem}{4.1}{}
For $\omega_1,\omega_2\in C[0,\infty)$ and $n\ge 1$ we write
\[
  m_n(\omega_1,\omega_2)\defeq\max_{0\le t\le n}\abs{\omega_1(t)-\omega_2(t)},
  \qquad
  d_n(\omega_1,\omega_2)\defeq\max_{0\le t\le n}\bigl(\abs{\omega_1(t)-\omega_2(t)}\wedge 1\bigr)
  = m_n(\omega_1,\omega_2)\wedge 1 .
\]
Both maxima are attained, because $t\mapsto\abs{\omega_1(t)-\omega_2(t)}$ is continuous on the compact interval
$[0,n]$, and the last equality holds because $a\mapsto a\wedge 1$ is nondecreasing. Thus
$\rho=\sum_{n\ge1}2^{-n}d_n$ with $0\le d_n\le 1$, so the series converges and $0\le\rho\le 1$. Since
$d_n\le d_{n+1}$, we have for every $N\ge 1$
\begin{equation}\label{eq:2-4-01:a}
  2^{-N}d_N(\omega_1,\omega_2)\;\le\;\rho(\omega_1,\omega_2)\;\le\;
  \sum_{n=1}^{N}2^{-n}d_N(\omega_1,\omega_2)+\sum_{n>N}2^{-n}
  \;\le\; d_N(\omega_1,\omega_2)+2^{-N}.
\end{equation}

\emph{Step 1: $\rho$ is a metric.} Symmetry is evident, and $\rho(\omega,\omega)=0$. If
$\rho(\omega_1,\omega_2)=0$, then $d_n(\omega_1,\omega_2)=0$ for every $n$, i.e.\ $\omega_1=\omega_2$ on $[0,n]$
for every $n$, so $\omega_1=\omega_2$. For the triangle inequality we note that $\varphi(a)\defeq a\wedge 1$,
$a\ge 0$, is nondecreasing and subadditive: if $a+b\le 1$ then $\varphi(a+b)=a+b=\varphi(a)+\varphi(b)$, and if
$a+b>1$ then either one of $a,b$ is at least $1$, or both are less than $1$ and $\varphi(a)+\varphi(b)=a+b>1$;
in both cases $\varphi(a)+\varphi(b)\ge 1=\varphi(a+b)$. Hence, for $\omega_1,\omega_2,\omega_3\in C[0,\infty)$
and $0\le t\le n$,
\[
  \abs{\omega_1(t)-\omega_3(t)}\wedge 1
  \le \varphi\bigl(\abs{\omega_1(t)-\omega_2(t)}+\abs{\omega_2(t)-\omega_3(t)}\bigr)
  \le d_n(\omega_1,\omega_2)+d_n(\omega_2,\omega_3).
\]
Taking the maximum over $t\in[0,n]$ gives $d_n(\omega_1,\omega_3)\le d_n(\omega_1,\omega_2)+d_n(\omega_2,\omega_3)$;
multiplying by $2^{-n}$ and summing over $n$ yields
$\rho(\omega_1,\omega_3)\le\rho(\omega_1,\omega_2)+\rho(\omega_2,\omega_3)$.

\emph{Step 2: completeness.} Let $\{\omega_k\}_{k\ge1}$ be a Cauchy sequence for $\rho$. Fix $n\ge 1$ and
$0<\eps<1$. There is $K$ such that $\rho(\omega_k,\omega_l)<2^{-n}\eps$ for $k,l\ge K$, and then
\eqref{eq:2-4-01:a} gives $d_n(\omega_k,\omega_l)<\eps<1$, that is,
$m_n(\omega_k,\omega_l)<\eps$. So $\{\omega_k\}$ is uniformly Cauchy on $[0,n]$. In particular
$\{\omega_k(t)\}_{k\ge 1}$ is a Cauchy sequence of real numbers for every $t\ge 0$, and we may define
$\omega(t)\defeq\lim_{k\to\infty}\omega_k(t)$. Letting $l\to\infty$ in
$\abs{\omega_k(t)-\omega_l(t)}<\eps$ we obtain $\abs{\omega_k(t)-\omega(t)}\le\eps$ for all $t\in[0,n]$ and
$k\ge K$; thus $\omega_k\to\omega$ uniformly on $[0,n]$. As a uniform limit of continuous functions, $\omega$ is
continuous on $[0,n]$; $n$ being arbitrary, $\omega\in C[0,\infty)$, and $m_n(\omega_k,\omega)\to 0$ as
$k\to\infty$, for every $n$. Finally, given $\eps>0$, choose $N$ with $2^{-N}<\eps/2$ and then $K$ with
$m_N(\omega_k,\omega)<\eps/2$ for $k\ge K$; by \eqref{eq:2-4-01:a},
$\rho(\omega_k,\omega)\le d_N(\omega_k,\omega)+2^{-N}<\eps$ for $k\ge K$. Hence $\rho(\omega_k,\omega)\to 0$, and
$\bigl(C[0,\infty),\rho\bigr)$ is complete.

\emph{Step 3: separability.} For integers $N,m\ge 1$ and $q=(q_0,\dots,q_{Nm})\in\Q^{Nm+1}$, let
$g_{N,m,q}$ be the function on $[0,\infty)$ which takes the value $q_j$ at $j/m$, $0\le j\le Nm$, is linear on
each interval $[j/m,(j+1)/m]$, $0\le j<Nm$, and equals $q_{Nm}$ on $[N,\infty)$. It is continuous, and the
collection $D$ of all such functions is countable, being indexed by the countable set
$\bigcup_{N,m\ge1}\{(N,m)\}\times\Q^{Nm+1}$. We show that $D$ is dense.

Let $\omega\in C[0,\infty)$ and $\eps>0$. Choose $N$ with $2^{-N}<\eps/2$. Since $\omega$ is uniformly
continuous on the compact interval $[0,N]$, there is $m\ge1$ such that $\abs{\omega(s)-\omega(t)}<\eps/6$
whenever $s,t\in[0,N]$ and $\abs{s-t}\le 1/m$. Choose rationals $q_j$ with
$\abs{q_j-\omega(j/m)}<\eps/6$, $0\le j\le Nm$, and put $g\defeq g_{N,m,q}\in D$. Let $t\in[0,N]$ and pick $j<Nm$
with $j/m\le t\le (j+1)/m$. For $i\in\{j,j+1\}$,
\[
  \abs{q_i-\omega(t)}\le\abs{q_i-\omega(i/m)}+\abs{\omega(i/m)-\omega(t)}<\frac{\eps}{6}+\frac{\eps}{6}=\frac{\eps}{3},
\]
and $g(t)=\lambda q_j+(1-\lambda)q_{j+1}$ for some $\lambda\in[0,1]$, so that
$\abs{g(t)-\omega(t)}\le\lambda\abs{q_j-\omega(t)}+(1-\lambda)\abs{q_{j+1}-\omega(t)}<\eps/3$. Therefore
$d_N(g,\omega)\le m_N(g,\omega)<\eps/2$, and \eqref{eq:2-4-01:a} gives
$\rho(g,\omega)\le d_N(g,\omega)+2^{-N}<\eps$. Thus $D$ is a countable dense subset of
$\bigl(C[0,\infty),\rho\bigr)$, which is therefore separable.

\rmk{The inequalities \eqref{eq:2-4-01:a} show, as in Step 2, that $\rho(\omega_k,\omega)\to0$ if and only if
$\omega_k\to\omega$ uniformly on every compact subset of $[0,\infty)$: $\rho$ metrizes the topology of
uniform convergence on compacts. The same proof, with $\abs{\,\cdot\,}$ the Euclidean norm and $\Q^d$-valued
nodes, applies to $C[0,\infty)^d$.}
\end{kex}

\begin{kex}{Problem}{4.5}{}
We use the notation of \bk{Definition 4.4}: $X_n$ is defined on a probability space $(\Omega_n,\F_n,P_n)$ and is $\F_n/\B(S_1)$-measurable, $X$ is defined on $(\Omega,\F,P)$ and is $\F/\B(S_1)$-measurable, and $E_n$, $E$ denote expectations with respect to $P_n$, $P$.

\emph{Step 1: $Y_n$ and $Y$ are random variables with values in $(S_2,\B(S_2))$.}
The collection $\mathscr{D}\defeq\{B\subseteq S_2;\ \varphi^{-1}(B)\in\B(S_1)\}$ is a $\sigma$-field on $S_2$, because taking preimages commutes with complements and countable unions and $\varphi^{-1}(S_2)=S_1$. Since $\varphi$ is continuous, $\varphi^{-1}(G)$ is open in $S_1$, hence belongs to $\B(S_1)$, for every open $G\subseteq S_2$; thus $\mathscr{D}$ contains the open sets of $S_2$ and therefore $\B(S_2)\subseteq\mathscr{D}$. In other words, $\varphi$ is $\B(S_1)/\B(S_2)$-measurable. Consequently, for every $B\in\B(S_2)$,
\[
  \{Y_n\in B\}=\{X_n\in\varphi^{-1}(B)\}\in\F_n ,\qquad \{Y\in B\}=\{X\in\varphi^{-1}(B)\}\in\F ,
\]
so $Y_n$ is $\F_n/\B(S_2)$-measurable and $Y$ is $\F/\B(S_2)$-measurable. Moreover these identities show that the induced measures of \bk{(4.2)} are related by
\begin{equation}\label{eq:2-4-05:a}
  P_nY_n^{-1}(B)=P_nX_n^{-1}\bigl(\varphi^{-1}(B)\bigr),\qquad PY^{-1}(B)=PX^{-1}\bigl(\varphi^{-1}(B)\bigr),\qquad B\in\B(S_2);
\end{equation}
that is, $P_nY_n^{-1}$ and $PY^{-1}$ are the images under $\varphi$ of $P_nX_n^{-1}$ and $PX^{-1}$, respectively.

\emph{Step 2: convergence.}
Let $f\colon S_2\to\R$ be bounded and continuous, and set $g\defeq f\circ\varphi\colon S_1\to\R$. Then $g$ is continuous, as a composition of continuous functions, and bounded, since $\sup_{S_1}\abs{g}\le\sup_{S_2}\abs{f}<\infty$. By \eqref{eq:2-4-05:a} and the change-of-variables formula for image measures (valid for the bounded Borel function $f$),
\[
  \int_{S_2} f\,d\bigl(P_nY_n^{-1}\bigr)=\int_{S_1} f\circ\varphi\;d\bigl(P_nX_n^{-1}\bigr)=\int_{S_1} g\,d\bigl(P_nX_n^{-1}\bigr),
\]
and in the same way $\int_{S_2} f\,d(PY^{-1})=\int_{S_1} g\,d(PX^{-1})$. By hypothesis $X_n\xrightarrow{\mathscr{D}}X$, which by \bk{Definition 4.4} means that $P_nX_n^{-1}$ converges weakly to $PX^{-1}$ on $(S_1,\B(S_1))$; applying \bk{Definition 4.3} to the bounded continuous function $g$ we get
\[
  \lim_{n\to\infty}\int_{S_2} f\,d\bigl(P_nY_n^{-1}\bigr)
  =\lim_{n\to\infty}\int_{S_1} g\,d\bigl(P_nX_n^{-1}\bigr)
  =\int_{S_1} g\,d\bigl(PX^{-1}\bigr)
  =\int_{S_2} f\,d\bigl(PY^{-1}\bigr).
\]
Since $f$ was an arbitrary bounded, continuous, real-valued function on $S_2$, the sequence $\{P_nY_n^{-1}\}_{n=1}^\infty$ converges weakly to $PY^{-1}$ in the sense of \bk{Definition 4.3}, i.e., $Y_n\xrightarrow{\mathscr{D}}Y$ in the sense of \bk{Definition 4.4}.

\rmk{In the equivalent formulation of \bk{Definition 4.4} the argument reads $E_nf(Y_n)=E_n\,g(X_n)\to E\,g(X)=Ef(Y)$. Only two properties of $\varphi$ were used: it is Borel measurable, and $f\circ\varphi$ is bounded and continuous on $S_1$ whenever $f$ is bounded and continuous on $S_2$.}
\end{kex}

\begin{kex}{Problem}{4.8}{}
Fix $T>0$ and, for $\delta>0$, write
\[
D_\delta\defeq\bigl\{(s,t)\in[0,T]^2;\ \abs{s-t}\le\delta\bigr\},
\qquad\text{so that}\qquad
m^T(\omega,\delta)=\max_{(s,t)\in D_\delta}\abs{\omega(s)-\omega(t)}
\]
by \bk{(4.3)}. The set $D_\delta$ is a nonempty, closed and bounded subset of $\R^2$, hence compact, and
$(s,t)\mapsto\abs{\omega(s)-\omega(t)}$ is continuous on it for every $\omega\in C[0,\infty)$; thus the
maximum in \bk{(4.3)} is attained, and $m^T(\omega,\delta)$ is a well-defined number in $[0,\infty)$.

\emph{Monotonicity in $\delta$.} If $0<\delta_1\le\delta_2$, then $D_{\delta_1}\subseteq D_{\delta_2}$, and a
maximum over a larger set is at least as large: $m^T(\omega,\delta_1)\le m^T(\omega,\delta_2)$.

\emph{Continuity in $\omega$.} Fix $\delta>0$ and an integer $n\ge T$. For $\omega_1,\omega_2\in C[0,\infty)$ put
\[
d_n(\omega_1,\omega_2)\defeq\max_{0\le t\le n}\abs{\omega_1(t)-\omega_2(t)}<\infty .
\]
For every $(s,t)\in D_\delta$ we have $s,t\in[0,T]\subseteq[0,n]$, so the triangle inequality gives
\[
\abs{\omega_1(s)-\omega_1(t)}
\le\abs{\omega_1(s)-\omega_2(s)}+\abs{\omega_2(s)-\omega_2(t)}+\abs{\omega_2(t)-\omega_1(t)}
\le m^T(\omega_2,\delta)+2\,d_n(\omega_1,\omega_2).
\]
Taking the maximum over $(s,t)\in D_\delta$ and then interchanging the roles of $\omega_1$ and $\omega_2$, we obtain
\begin{equation}\label{eq:2-4-08:a}
\abs{m^T(\omega_1,\delta)-m^T(\omega_2,\delta)}\le 2\,d_n(\omega_1,\omega_2).
\end{equation}
It remains to compare $d_n$ with $\rho$. Since $x\mapsto x\wedge1$ is nondecreasing and continuous on $[0,\infty)$,
\[
\max_{0\le t\le n}\bigl(\abs{\omega_1(t)-\omega_2(t)}\wedge1\bigr)=d_n(\omega_1,\omega_2)\wedge1 ,
\]
and all the terms of the series \bk{(4.1)} are nonnegative; keeping only the $n$-th one,
\begin{equation}\label{eq:2-4-08:b}
d_n(\omega_1,\omega_2)\wedge1\le 2^n\rho(\omega_1,\omega_2).
\end{equation}
Consequently, if $\rho(\omega_1,\omega_2)<2^{-n}$, then $d_n(\omega_1,\omega_2)\wedge1<1$, which forces
$d_n(\omega_1,\omega_2)<1$ and therefore $d_n(\omega_1,\omega_2)=d_n(\omega_1,\omega_2)\wedge1\le2^n\rho(\omega_1,\omega_2)$.
Together with \eqref{eq:2-4-08:a} this yields
\begin{equation}\label{eq:2-4-08:c}
\abs{m^T(\omega_1,\delta)-m^T(\omega_2,\delta)}\le 2^{n+1}\rho(\omega_1,\omega_2)
\qquad\text{whenever }\rho(\omega_1,\omega_2)<2^{-n}.
\end{equation}
Given $\eps>0$, the right-hand side of \eqref{eq:2-4-08:c} is less than $\eps$ as soon as
$\rho(\omega_1,\omega_2)<2^{-n}\min\{1,\eps/2\}$. Hence $\omega\mapsto m^T(\omega,\delta)$ is continuous on
$(C[0,\infty),\rho)$; in fact it is uniformly continuous, with a modulus that depends on $T$ (through $n$) but
neither on $\delta$ nor on the point $\omega$.

\emph{The limit as $\delta\downarrow0$.} Fix $\omega\in C[0,\infty)$. Its restriction to the compact interval
$[0,T]$ is continuous, hence uniformly continuous (Heine--Cantor): for every $\eps>0$ there is a
$\delta_0=\delta_0(\eps,\omega,T)>0$ such that $\abs{\omega(s)-\omega(t)}\le\eps$ for all $s,t\in[0,T]$ with
$\abs{s-t}\le\delta_0$, that is, $m^T(\omega,\delta_0)\le\eps$. By the monotonicity already proved,
$0\le m^T(\omega,\delta)\le m^T(\omega,\delta_0)\le\eps$ for all $0<\delta\le\delta_0$. Since $\eps>0$ is arbitrary,
$\lim_{\delta\downarrow0}m^T(\omega,\delta)=0$.

\rmk{These are exactly the two facts invoked in the proof of \bk{Theorem 4.9}. For $\eps>0$ the set
$K_\delta=\{\omega\in\ol{A};\ m^T(\omega,\delta)\ge\eps\}$ is the intersection of the closed set $\ol{A}$ with the
preimage of the closed set $[\eps,\infty)$ under the continuous map $m^T(\,\cdot\,,\delta)$, hence closed; and
$\bigcap_{\delta>0}K_\delta=\emptyset$, because every $\omega$ satisfies $m^T(\omega,\delta)<\eps$ for all
sufficiently small $\delta$. Monotonicity in $\delta$ makes the family $\{K_\delta\}_{\delta>0}$ decrease as
$\delta\downarrow0$, which is what allows the finite intersection property of compact sets to produce a single
$\delta(\eps)$ with $K_{\delta(\eps)}=\emptyset$.}
\end{kex}

\begin{kex}{Problem}{4.12}{}
We write $S \defeq C[0,\infty)$, equipped with the metric $\rho$ of \bk{(4.1)}, and let $M<\infty$ be a constant with $\abs{f_n(\omega)}\le M$ for all $n\ge 1$ and $\omega\in S$. The functions $f_n$ and $f$ are continuous, hence $\B(S)$-measurable. Since every singleton $\{\omega\}$ is a compact subset of $S$, the hypothesis gives $f_n(\omega)\to f(\omega)$ for every $\omega\in S$, and therefore $\abs{f(\omega)}\le M$ for every $\omega \in S$. Thus $f$ is bounded and continuous, all the integrals in \bk{(4.8)} are well defined, and
\begin{equation}\label{eq:2-4-12:bd}
\abs{f_n(\omega)-f(\omega)}\le 2M ,\qquad n\ge 1,\ \omega\in S .
\end{equation}

\emph{Step 1: the family $\Pi\defeq\{P_n;\ n\ge 1\}$ is tight.} By \bk{Problem 4.1}, $S$ is a complete, separable metric space, so by Prohorov's \bk{Theorem 4.7} it suffices to show that $\Pi$ is relatively compact in the sense of \bk{Definition 4.6}. Let $\{Q_k\}_{k=1}^\infty$ be a sequence of elements of $\Pi$, and choose indices $n_k\ge 1$ with $Q_k=P_{n_k}$. If the sequence $\{n_k\}_{k=1}^\infty$ is bounded, some integer $m$ satisfies $n_k=m$ for infinitely many $k$; along these $k$ the sequence $\{Q_k\}$ is constant, equal to $P_m$, and thus converges weakly to $P_m$. If $\{n_k\}_{k=1}^\infty$ is unbounded, we can select $k_1<k_2<\cdots$ with $n_{k_1}<n_{k_2}<\cdots$; then $\{Q_{k_j}\}_{j=1}^\infty=\{P_{n_{k_j}}\}_{j=1}^\infty$ is a subsequence of $\{P_n\}_{n=1}^\infty$, and for every bounded, continuous $g:S\to\R$ the numbers $\int_S g\,dQ_{k_j}$ form a subsequence of the convergent sequence $\int_S g\,dP_n$, so that $Q_{k_j}$ converges weakly to $P$ as $j\to\infty$ (\bk{Definition 4.3}). In either case $\{Q_k\}$ contains a weakly convergent subsequence. Hence $\Pi$ is relatively compact, and consequently tight.

\emph{Step 2: proof of \bk{(4.8)}.} Fix $\eps>0$. By Step 1 there is a compact set $K\subseteq S$ such that
\[
P_n(S\setminus K)\le \eps \qquad\text{for every } n\ge 1 ;
\]
$K$ is closed, hence a Borel set. For every $n\ge 1$ we have
\begin{equation}\label{eq:2-4-12:split}
\abs[\Big]{\int_S f_n\,dP_n-\int_S f\,dP}
\le \int_S \abs{f_n-f}\,dP_n+\abs[\Big]{\int_S f\,dP_n-\int_S f\,dP} .
\end{equation}
Splitting the first integral on the right over $K$ and $S\setminus K$ and using \eqref{eq:2-4-12:bd} on $S\setminus K$, we obtain
\[
\int_S \abs{f_n-f}\,dP_n
\le \sup_{\omega\in K}\abs{f_n(\omega)-f(\omega)}\cdot P_n(K)+2M\,P_n(S\setminus K)
\le \sup_{\omega\in K}\abs{f_n(\omega)-f(\omega)}+2M\eps .
\]
Since $f_n\to f$ uniformly on the compact set $K$, the supremum on the right tends to zero as $n\to\infty$. The second term on the right-hand side of \eqref{eq:2-4-12:split} tends to zero as $n\to\infty$, because $f$ is bounded and continuous and $P_n$ converges weakly to $P$ (\bk{Definition 4.3}). Therefore
\[
\limsup_{n\to\infty}\abs[\Big]{\int_S f_n\,dP_n-\int_S f\,dP}\le 2M\eps ,
\]
and since $\eps>0$ is arbitrary, \bk{(4.8)} follows.

\rmk{The only properties of $C[0,\infty)$ used are that it is a complete, separable metric space; the same proof gives the extension to $C[0,\infty)^d$ mentioned in \bk{Remark 4.13}, and indeed to any complete, separable metric space.}
\end{kex}

\begin{kex}{Exercise}{4.14}{}
For $n\ge 1$ let $x_n\colon[0,\infty)\to\R$ be the function $x_n(t)\defeq nt\,\ind_{[0,1/2n]}(t)+(1-nt)\,\ind_{(1/2n,1/n]}(t)$. It is piecewise affine, with $x_n(0)=0$, $x_n(1/2n)=\tfrac12=\lim_{t\downarrow 1/2n}(1-nt)$ and $x_n(1/n)=0=x_n(t)$ for $t>1/n$; hence $x_n\in C[0,\infty)$, and
\begin{equation}\label{eq:2-4-14:a}
0\le x_n(t)\le \tfrac12 \ \ (t\ge 0),\qquad x_n\bigl(\tfrac{1}{2n}\bigr)=\tfrac12,\qquad x_n(t)=0\ \text{ for } t=0 \text{ and for } t\ge \tfrac1n .
\end{equation}
The process $X^{(n)}$ is the constant mapping $\omega\mapsto x_n$ on some probability space $(\Omega_n,\F_n,P_n)$, and $X$ is the constant mapping $\omega\mapsto 0$ on some $(\Omega,\F,P)$; constant mappings are measurable, so these are random variables with values in $\bigl(C[0,\infty),\B(C[0,\infty))\bigr)$, and their laws are the unit point masses at $x_n$ and at the zero function, respectively. Consequently, for every bounded measurable $g$ on $C[0,\infty)$ and every bounded measurable $f$ on $\R^d$,
\begin{equation}\label{eq:2-4-14:b}
E_n\,g(X^{(n)})=g(x_n),\quad E\,g(X)=g(0),\qquad
E_n f\bigl(X^{(n)}_{t_1},\dots,X^{(n)}_{t_d}\bigr)=f\bigl(x_n(t_1),\dots,x_n(t_d)\bigr).
\end{equation}

\emph{Convergence of the finite-dimensional distributions.} Let $d\ge1$ and $t_1,\dots,t_d\in[0,\infty)$, and let $f\colon\R^d\to\R$ be bounded and continuous. Choose an integer $n_0$ with $n_0\ge 1/t_j$ for every $j$ such that $t_j>0$. For $n\ge n_0$ and each $j$ we have either $t_j=0$ or $t_j\ge 1/n$, so $x_n(t_j)=0$ by \eqref{eq:2-4-14:a}. By \eqref{eq:2-4-14:b},
\[
E_n f\bigl(X^{(n)}_{t_1},\dots,X^{(n)}_{t_d}\bigr)=f(0,\dots,0)=E f\bigl(X_{t_1},\dots,X_{t_d}\bigr)\qquad\text{for all } n\ge n_0 .
\]
The left-hand side is thus eventually equal to the right-hand side, and in particular converges to it. By \bk{Definition 4.4} (with $S=\R^d$), $\bigl(X^{(n)}_{t_1},\dots,X^{(n)}_{t_d}\bigr)\xrightarrow{\ \mathscr{D}\ }\bigl(X_{t_1},\dots,X_{t_d}\bigr)$, i.e., the finite-dimensional distributions of $X^{(n)}$ converge weakly to those of $X$.

\emph{Failure of convergence in distribution.} Define $g\colon C[0,\infty)\to\R$ by
\[
g(\omega)\defeq \Bigl(\max_{0\le t\le 1}\abs{\omega(t)}\Bigr)\wedge 1 .
\]
Then $0\le g\le 1$, and $g$ is continuous for the metric $\rho$ of \bk{(4.1)}: for $\omega_1,\omega_2\in C[0,\infty)$, the triangle inequality for the supremum norm on $[0,1]$ and the elementary inequality $\abs{a\wedge1-b\wedge1}\le\abs{a-b}\wedge1$ $(a,b\ge0)$ give
\[
\abs{g(\omega_1)-g(\omega_2)}
\le \Bigl(\max_{0\le t\le 1}\abs{\omega_1(t)-\omega_2(t)}\Bigr)\wedge 1
=\max_{0\le t\le 1}\bigl(\abs{\omega_1(t)-\omega_2(t)}\wedge 1\bigr)
\le 2\rho(\omega_1,\omega_2),
\]
the last step because the middle quantity, multiplied by $\tfrac12$, is the term $n=1$ of the series \bk{(4.1)}, all of whose terms are nonnegative. So $g$ is a bounded, continuous function on $C[0,\infty)$. Since $1/2n\in[0,1]$ for every $n\ge1$, \eqref{eq:2-4-14:a} yields $\max_{0\le t\le1}\abs{x_n(t)}=\tfrac12$, and therefore, by \eqref{eq:2-4-14:b},
\[
E_n\,g(X^{(n)})=g(x_n)=\tfrac12\quad\text{for every } n\ge1,\qquad\text{whereas}\qquad E\,g(X)=g(0)=0 .
\]
Thus $\lim_{n\to\infty}E_n\,g(X^{(n)})=\tfrac12\ne E\,g(X)$ for a bounded continuous $g$, and by \bk{Definition 4.4} the sequence $\{X^{(n)}\}_{n=1}^\infty$ does not converge in distribution to $X$.

\rmk{The same function $g$ shows that no subsequence of $\{X^{(n)}\}$ converges in distribution to $X$. In fact no subsequence converges in distribution at all: by the discussion preceding the exercise, a limit would have the finite-dimensional distributions of $X$, hence would have the law of $X$, since the law of a continuous process is determined by its finite-dimensional distributions (\bk{Problem 4.2}: the cylinder sets form a $\pi$-system generating $\B(C[0,\infty))$). By the Prohorov theorem \bk{4.7}, the sequence $\{X^{(n)}\}_{n=1}^\infty$ is therefore not tight, which is why \bk{Theorem 4.15} does not apply to it.}
\end{kex}

\solsection{2.5}{The Markov Property}

\begin{kex}{Problem}{5.4}{}
The assertions to be proved are: (I) the process $B$ built on the product space
$\R^d\times\Omega^{(1)}\times\dots\times\Omega^{(d)}$ is a $d$-dimensional Brownian motion with initial
distribution $\mu$; (II) under $P^0$ the coordinate mapping process on $C[0,\infty)^d$, with the filtration
$\{\F^B_t\}$, is a $d$-dimensional Brownian motion starting at the origin; (III) under $P^x$ of \bk{(5.1)} it is
a $d$-dimensional Brownian motion starting at $x$; (IV) \bk{Proposition 5.3}. We reduce all four to one
criterion on the joint law of the initial position and the increments.

\emph{Notation.} For $h>0$ we write $\gamma_h$ for the normal law on $\R$ with mean $0$ and variance $h$,
and $\nu_h\defeq\gamma_h\times\dots\times\gamma_h$ ($d$ factors) for the normal law on $\R^d$ with mean zero
and covariance matrix $hI_d$; $\delta_x$ is the unit mass at $x$. For a probability measure $\mu$ on
$(\R^d,\B(\R^d))$ and an $\R^d$-valued process $B$ on a probability space $(\Omega,\F,P)$ we consider the
following property:
\begin{equation}\label{eq:2-5-04:star}
\begin{gathered}
P\bigl[B_0\in\Gamma_0,\ B_{t_1}-B_{t_0}\in\Gamma_1,\ \dots,\ B_{t_n}-B_{t_{n-1}}\in\Gamma_n\bigr]
 =\mu(\Gamma_0)\prod_{k=1}^n\nu_{t_k-t_{k-1}}(\Gamma_k)\\
\text{for all } n\ge1,\quad 0=t_0<t_1<\dots<t_n<\infty,\quad \Gamma_0,\dots,\Gamma_n\in\B(\R^d).
\end{gathered}
\end{equation}

\begin{lem}[A]
\textup{(a)} If $B=\{B_t,\F_t;\,t\ge0\}$ is a $d$-dimensional Brownian motion with initial distribution
$\mu$, relative to any filtration $\{\F_t\}$, then \eqref{eq:2-5-04:star} holds.

\textup{(b)} If $B$ is a continuous, $\R^d$-valued process satisfying \eqref{eq:2-5-04:star}, then
$\{B_t,\F^B_t;\,t\ge0\}$ is a $d$-dimensional Brownian motion with initial distribution $\mu$.
\end{lem}

\begin{proof}
(a) We use induction on $n$. Let $D_m\defeq\{B_0\in\Gamma_0,\,B_{t_k}-B_{t_{k-1}}\in\Gamma_k,\ 1\le k\le m\}$
for $0\le m\le n$. By adaptedness $D_m\in\F_{t_m}$, and $P(D_0)=\mu(\Gamma_0)$ by \bk{Definition 5.1}(i).
By \bk{Definition 5.1}(ii) the increment $B_{t_m}-B_{t_{m-1}}$ is independent of $\F_{t_{m-1}}\ni D_{m-1}$ and
has law $\nu_{t_m-t_{m-1}}$, so $P(D_m)=P(D_{m-1})\,\nu_{t_m-t_{m-1}}(\Gamma_m)$ for $1\le m\le n$, and
\eqref{eq:2-5-04:star} follows.

(b) The process $B$ is continuous by assumption and is adapted to $\{\F^B_t\}$. Taking $n=1$ and
$\Gamma_1=\R^d$ in \eqref{eq:2-5-04:star} gives $P[B_0\in\Gamma_0]=\mu(\Gamma_0)$, which is (i) of
\bk{Definition 5.1}. Fix $0\le s<t$ and $\Gamma\in\B(\R^d)$. Taking in \eqref{eq:2-5-04:star} the partition
$\{0,s,t\}$ (or $\{0,t\}$ if $s=0$) and all sets but the last equal to $\R^d$, we see that $B_t-B_s$ has law
$\nu_{t-s}$. It remains to show that
\begin{equation}\label{eq:2-5-04:ind}
P\bigl[D\cap\{B_t-B_s\in\Gamma\}\bigr]=P(D)\,P[B_t-B_s\in\Gamma]
\end{equation}
for every $D\in\F^B_s$.

Let $\pi=\{0=s_0<s_1<\dots<s_m=s\}$ be a partition of $[0,s]$ (with $m=0$ if $s=0$) and put
$\xi_\pi\defeq(B_0,B_{s_1}-B_{s_0},\dots,B_{s_m}-B_{s_{m-1}})$, a random vector with values in
$(\R^d)^{m+1}$. Applying \eqref{eq:2-5-04:star} to the partition $s_0<\dots<s_m<t$ of $[0,t]$, once with
the sets $\Gamma_0,\dots,\Gamma_m,\Gamma$ and once with $\Gamma_0,\dots,\Gamma_m,\R^d$, we obtain
\[
P\bigl[\xi_\pi\in G,\ B_t-B_s\in\Gamma\bigr]=P[\xi_\pi\in G]\;\nu_{t-s}(\Gamma)
=P[\xi_\pi\in G]\;P[B_t-B_s\in\Gamma]
\]
for every rectangle $G=\Gamma_0\times\dots\times\Gamma_m$ with $\Gamma_j\in\B(\R^d)$. These rectangles are
closed under pairwise intersection and generate $\B((\R^d)^{m+1})$, and the collection of all
$G\in\B((\R^d)^{m+1})$ for which the equality between the first and the last member holds is a Dynkin system
(it contains the whole space and is closed under proper differences and increasing unions, by additivity and
continuity from below of $P$). By the Dynkin System \bk{Theorem 1.3} the equality holds for every
$G\in\B((\R^d)^{m+1})$. Now $(B_{s_0},\dots,B_{s_m})=L\xi_\pi$, where
$L(y_0,\dots,y_m)=(y_0,y_0+y_1,\dots,y_0+\dots+y_m)$ is a linear bijection of $(\R^d)^{m+1}$; hence
$\{(B_{s_0},\dots,B_{s_m})\in A\}=\{\xi_\pi\in L^{-1}A\}$ for $A\in\B((\R^d)^{m+1})$, and
\eqref{eq:2-5-04:ind} holds for every $D$ in the collection
\[
\cC\defeq\bigl\{\{(B_{s_0},\dots,B_{s_m})\in A\};\ \ m\ge0,\ \ 0=s_0<\dots<s_m=s,\ \ A\in\B((\R^d)^{m+1})\bigr\}.
\]
The collection $\cC$ is closed under pairwise intersection: an event of $\cC$ written with a partition
$\pi$ can also be written with any finer partition $\pi'$ of $[0,s]$ (replace $A$ by its inverse image under
the projection which keeps the coordinates indexed by $\pi$), so two events of $\cC$ can be written with the
same partition, and then their intersection corresponds to the intersection of the two Borel sets. Moreover
$\cC\subseteq\F^B_s$, and $\cC$ contains every event $\{B_u\in\Gamma'\}$ with $0\le u\le s$,
$\Gamma'\in\B(\R^d)$ (use the partition with points $0,u,s$); thus $\sigma(\cC)=\sigma(B_u;\,0\le u\le s)=\F^B_s$.
Finally, the collection of all $D\in\F$ satisfying \eqref{eq:2-5-04:ind} is a Dynkin system, for the same
reason as before. By \bk{Theorem 1.3} it contains $\F^B_s$. Hence $B_t-B_s$ is independent of $\F^B_s$, and
(ii) of \bk{Definition 5.1} holds.
\end{proof}

\begin{lem}[B]
On a probability space $(\Omega,\F,P)$ let $X$ be an $\R^d$-valued random variable with distribution $\mu$,
and for $i=1,\dots,d$ let $\wt B^{(i)}=\{\wt B^{(i)}_t,\wt\F^{(i)}_t;\,t\ge0\}$ be a standard, one-dimensional
Brownian motion. Assume that the $\sigma$-fields
$\G_0\defeq\sigma(X)$ and $\G_i\defeq\sigma(\wt B^{(i)}_t;\,t\ge0)$, $1\le i\le d$, are independent. Then
\[
B_t\defeq X+\bigl(\wt B^{(1)}_t,\dots,\wt B^{(d)}_t\bigr),\qquad 0\le t<\infty,
\]
together with $\{\F^B_t\}$, is a $d$-dimensional Brownian motion with initial distribution $\mu$.
\end{lem}

\begin{proof}
$B$ is a continuous, $\R^d$-valued process, because each $\wt B^{(i)}$ is continuous. Fix
$0=t_0<t_1<\dots<t_n$, write $h_k\defeq t_k-t_{k-1}$ and
$\eta^i_k\defeq\wt B^{(i)}_{t_k}-\wt B^{(i)}_{t_{k-1}}$, so that
$B_{t_k}-B_{t_{k-1}}=(\eta^1_k,\dots,\eta^d_k)$ for $1\le k\le n$. A standard, one-dimensional Brownian
motion (\bk{Definition 1.1}) is a one-dimensional Brownian motion with initial distribution $\delta_0$ in
the sense of \bk{Definition 5.1}, so Lemma~A(a) with $d=1$ gives, for each $i$ and all
$\Gamma^i_1,\dots,\Gamma^i_n\in\B(\R)$,
\[
P\bigl[\wt B^{(i)}_0=0\bigr]=1,\qquad
P\bigl[\eta^i_k\in\Gamma^i_k,\ 1\le k\le n\bigr]=\prod_{k=1}^n\gamma_{h_k}(\Gamma^i_k).
\]
In particular $B_0=X$ a.s. Let $\Gamma_0\in\B(\R^d)$ and
$\Gamma_k=\Gamma^1_k\times\dots\times\Gamma^d_k$ with $\Gamma^i_k\in\B(\R)$, $1\le k\le n$. The events
$\{X\in\Gamma_0\}\in\G_0$ and $\{\eta^i_k\in\Gamma^i_k,\ 1\le k\le n\}\in\G_i$, $1\le i\le d$, are independent,
whence
\begin{align*}
P\bigl[B_0\in\Gamma_0,\ B_{t_k}-B_{t_{k-1}}\in\Gamma_k,\ 1\le k\le n\bigr]
&=P[X\in\Gamma_0]\prod_{i=1}^dP\bigl[\eta^i_k\in\Gamma^i_k,\ 1\le k\le n\bigr]\\
&=\mu(\Gamma_0)\prod_{i=1}^d\prod_{k=1}^n\gamma_{h_k}(\Gamma^i_k)
 =\mu(\Gamma_0)\prod_{k=1}^n\nu_{h_k}(\Gamma_k).
\end{align*}
Thus the law of $(B_0,B_{t_1}-B_{t_0},\dots,B_{t_n}-B_{t_{n-1}})$ on $(\R^d)^{n+1}$ and the product measure
$\mu\times\nu_{h_1}\times\dots\times\nu_{h_n}$ agree on all sets
$\Gamma_0\times\Gamma_1\times\dots\times\Gamma_n$ of the above type. These sets are closed under pairwise
intersection and generate $\B((\R^d)^{n+1})$, and the sets on which two probability measures agree form a
Dynkin system; by \bk{Theorem 1.3} the two measures coincide. This is \eqref{eq:2-5-04:star}, and Lemma~A(b)
completes the proof.
\end{proof}

\textbf{(I) The product-space construction.} Let
$\Omega=\R^d\times\Omega^{(1)}\times\dots\times\Omega^{(d)}$,
$\F=\B(\R^d)\otimes\F^{(1)}\otimes\dots\otimes\F^{(d)}$ and
$P=\mu\times P^{(1)}\times\dots\times P^{(d)}$, and let $p_0,p_1,\dots,p_d$ be the coordinate projections
of $\Omega$ onto its factors; they are measurable, and by the definition of the product measure
\begin{equation}\label{eq:2-5-04:prod}
P\bigl[p_0^{-1}(A_0)\cap p_1^{-1}(A_1)\cap\dots\cap p_d^{-1}(A_d)\bigr]=\mu(A_0)\prod_{i=1}^dP^{(i)}(A_i),
\qquad A_0\in\B(\R^d),\ A_i\in\F^{(i)}.
\end{equation}
Taking all sets but one equal to the whole factor, we see that $Pp_0^{-1}=\mu$ and $Pp_i^{-1}=P^{(i)}$; then
\eqref{eq:2-5-04:prod} says that the $\sigma$-fields $p_0^{-1}(\B(\R^d))$, $p_1^{-1}(\F^{(1)})$, \dots,
$p_d^{-1}(\F^{(d)})$ are independent under $P$. The random variable $X\circ p_0=p_0$ has distribution $\mu$
and generates $p_0^{-1}(\B(\R^d))$. For each $i$, the process $\wt B^{(i)}_t\circ p_i$, $t\ge0$, is continuous
and adapted to the filtration $\{p_i^{-1}(\wt\F^{(i)}_t)\}$ of $\F$; because $Pp_i^{-1}=P^{(i)}$, we have for
$0\le s<t$, $A\in\wt\F^{(i)}_s$ and $\Gamma\in\B(\R)$
\begin{align*}
P\bigl[p_i^{-1}(A)\cap\{\wt B^{(i)}_t\circ p_i-\wt B^{(i)}_s\circ p_i\in\Gamma\}\bigr]
&=P^{(i)}\bigl[A\cap\{\wt B^{(i)}_t-\wt B^{(i)}_s\in\Gamma\}\bigr]\\
&=P^{(i)}(A)\,\gamma_{t-s}(\Gamma)=P\bigl[p_i^{-1}(A)\bigr]\gamma_{t-s}(\Gamma),
\end{align*}
and likewise $P[\wt B^{(i)}_0\circ p_i=0]=1$. Hence $\{\wt B^{(i)}_t\circ p_i,\ p_i^{-1}(\wt\F^{(i)}_t);\,t\ge0\}$
is a standard, one-dimensional Brownian motion on $(\Omega,\F,P)$, and
$\sigma(\wt B^{(i)}_t\circ p_i;\,t\ge0)\subseteq p_i^{-1}(\F^{(i)})$. All hypotheses of Lemma~B are satisfied,
and its conclusion is that
$B_t(\omega)=X(\omega_0)+(\wt B^{(1)}_t(\omega_1),\dots,\wt B^{(d)}_t(\omega_d))$ with the filtration
$\F_t=\F^B_t$ is a $d$-dimensional Brownian motion with initial distribution $\mu$.

\textbf{(II) The measure $P^0$.} We regard $C[0,\infty)^d$ as the $d$-fold product of the metric space
$C[0,\infty)$ of \bk{(4.1)}, with the product topology; its elements $\omega=(\omega_1,\dots,\omega_d)$ are
the continuous functions from $[0,\infty)$ into $\R^d$, and $B_t(\omega)=\omega(t)
=(\omega_1(t),\dots,\omega_d(t))$. Let $p_i\colon C[0,\infty)^d\to C[0,\infty)$ be the $i$-th coordinate
projection and $W_t(w)=w(t)$ the coordinate mapping process on $C[0,\infty)$. We first check that
\begin{equation}\label{eq:2-5-04:sig}
\B\bigl(C[0,\infty)^d\bigr)=\B(C[0,\infty))\otimes\dots\otimes\B(C[0,\infty))=\sigma(B_t;\,0\le t<\infty),
\end{equation}
so that the product measure $P^0=P^{(1)}\times\dots\times P^{(d)}$ is indeed defined on
$\B(C[0,\infty)^d)$. The projections $p_i$ are continuous, hence Borel measurable, and this gives
``$\supseteq$'' in the first equality. Conversely, $C[0,\infty)$ is a separable metric space
(\bk{Problem 4.1}) and so has a countable base $\mathscr{U}$ for its topology; the sets
$U_1\times\dots\times U_d$ with $U_i\in\mathscr{U}$ form a countable base for the product topology, so every
open subset of $C[0,\infty)^d$ is a countable union of such sets and belongs to the product
$\sigma$-field. For the second equality, $\B(C[0,\infty))=\sigma(W_t;\,t\ge0)$ by \bk{Problem 4.2}, so the
product $\sigma$-field, which is generated by $p_1,\dots,p_d$, is generated by the mappings
$W_t\circ p_i\colon\omega\mapsto\omega_i(t)$, $t\ge0$, $1\le i\le d$; these are the components of the
mappings $B_t$, $t\ge0$.

Now we argue as in (I), without the factor $\R^d$. Under $P^0$ the $\sigma$-fields
$p_i^{-1}(\B(C[0,\infty)))$, $1\le i\le d$, are independent and $P^0p_i^{-1}=P^{(i)}$ is Wiener measure.
By \bk{Definition 4.21}, $W$ is a standard, one-dimensional Brownian motion on
$(C[0,\infty),\B(C[0,\infty)),P^{(i)})$ relative to some filtration $\{\wt\F_t\}$, and the computation
displayed in (I) shows that $\{W_t\circ p_i,\ p_i^{-1}(\wt\F_t);\,t\ge0\}$ is a standard, one-dimensional
Brownian motion on $(C[0,\infty)^d,\B(C[0,\infty)^d),P^0)$, with
$\sigma(W_t\circ p_i;\,t\ge0)=p_i^{-1}(\B(C[0,\infty)))$. Lemma~B, applied with $X\equiv0$ (so that
$\G_0$ is trivial and $\mu=\delta_0$) and $\wt B^{(i)}_t=W_t\circ p_i$, shows that
$B_t=(W_t\circ p_1,\dots,W_t\circ p_d)$ with the filtration $\{\F^B_t\}$ is, under $P^0$, a $d$-dimensional
Brownian motion with initial distribution $\delta_0$, that is, starting at the origin.

\textbf{(III) The measures $P^x$.} Fix $x\in\R^d$ and let
$T_x\colon C[0,\infty)^d\to C[0,\infty)^d$, $T_x\omega\defeq\omega(\cdot)+x$. For every $t\ge0$ the mapping
$B_t\circ T_x=B_t+x$ is $\B(C[0,\infty)^d)/\B(\R^d)$-measurable, so by \eqref{eq:2-5-04:sig} $T_x$ is
$\B(C[0,\infty)^d)/\B(C[0,\infty)^d)$-measurable. Since $F-x=T_x^{-1}(F)$, the set $F-x$ is a Borel set
whenever $F$ is, and \bk{(5.1)} reads $P^x=P^0T_x^{-1}$; in particular $P^x$ is a probability measure on
$(C[0,\infty)^d,\B(C[0,\infty)^d))$. For $0=t_0<t_1<\dots<t_n$ and
$\Gamma_0,\dots,\Gamma_n\in\B(\R^d)$, the increments of $B\circ T_x=B+x$ are those of $B$, and therefore,
by (II) and Lemma~A(a) (the index $k$ runs over $1,\dots,n$, and $\Gamma_0-x\defeq\{y-x;\,y\in\Gamma_0\}$),
\begin{align}
P^x\bigl[B_0\in\Gamma_0,\ B_{t_k}-B_{t_{k-1}}\in\Gamma_k,\ k\le n\bigr]
&=P^0\bigl[B_0+x\in\Gamma_0,\ B_{t_k}-B_{t_{k-1}}\in\Gamma_k,\ k\le n\bigr]\notag\\
&=\delta_0(\Gamma_0-x)\prod_{k=1}^n\nu_{t_k-t_{k-1}}(\Gamma_k)
 =\delta_x(\Gamma_0)\prod_{k=1}^n\nu_{t_k-t_{k-1}}(\Gamma_k).\label{eq:2-5-04:px}
\end{align}
This is \eqref{eq:2-5-04:star} for the continuous process $B$ on
$(C[0,\infty)^d,\B(C[0,\infty)^d),P^x)$, with $\mu=\delta_x$. By Lemma~A(b), under $P^x$ the process
$\{B_t,\F^B_t;\,t\ge0\}$ is a $d$-dimensional Brownian motion with initial distribution $\delta_x$, i.e.,
starting at $x$.

\textbf{(IV) Proof of Proposition 5.3.} By \bk{Problem 5.2}, for each $F\in\B(C[0,\infty)^d)$ the function
$x\mapsto P^x(F)$ is Borel measurable with values in $[0,1]$, so the integral in \bk{(5.2)} is well defined
and lies in $[0,1]$. The set function $P^\mu$ is a probability measure: $P^\mu(C[0,\infty)^d)=\int 1\,\mu(dx)=1$,
and if $F_1,F_2,\dots\in\B(C[0,\infty)^d)$ are pairwise disjoint, then
$P^x(\bigcup_jF_j)=\sum_jP^x(F_j)$ for every $x$, and the monotone convergence theorem yields
$P^\mu(\bigcup_jF_j)=\sum_jP^\mu(F_j)$. Let now $0=t_0<t_1<\dots<t_n$ and
$\Gamma_0,\dots,\Gamma_n\in\B(\R^d)$. The event
$F=\{B_0\in\Gamma_0,\ B_{t_k}-B_{t_{k-1}}\in\Gamma_k,\ 1\le k\le n\}$ belongs to $\B(C[0,\infty)^d)$ by
\eqref{eq:2-5-04:sig}, and integrating \eqref{eq:2-5-04:px} with respect to $\mu(dx)$ we find
\[
P^\mu(F)=\int_{\R^d}\delta_x(\Gamma_0)\,\mu(dx)\,\prod_{k=1}^n\nu_{t_k-t_{k-1}}(\Gamma_k)
=\mu(\Gamma_0)\prod_{k=1}^n\nu_{t_k-t_{k-1}}(\Gamma_k),
\]
because $\delta_x(\Gamma_0)=\ind_{\Gamma_0}(x)$. Thus the coordinate mapping process, which is continuous,
satisfies \eqref{eq:2-5-04:star} on $(C[0,\infty)^d,\B(C[0,\infty)^d),P^\mu)$, and Lemma~A(b) shows that
$B=\{B_t,\F^B_t;\,t\ge0\}$ is, under $P^\mu$, a $d$-dimensional Brownian motion with initial distribution
$\mu$. This proves \bk{Proposition 5.3}.

\rmk{The book does not say which topology $C[0,\infty)^d$ carries; we have taken the product topology,
which is the topology of uniform convergence on compact sets of $\R^d$-valued functions and is induced,
for instance, by the metric $\sum_{i=1}^d\rho(\omega_i,\omega_i')$ with $\rho$ as in \bk{(4.1)}. The identity
\eqref{eq:2-5-04:sig} is what makes the product of $d$ Wiener measures a measure on the Borel
$\sigma$-field of this space. Lemma~A(b) is the $d$-dimensional form of \bk{Problem 1.4}.}
\end{kex}

\begin{kex}{Problem}{5.5}{}
Throughout, $B$ is a $d$-dimensional Brownian motion with some initial distribution $\mu$ on a probability space $(\Omega,\F,P)$, in the sense of \bk{Definition 5.1}; no moment assumption is made on $\mu$, and none is needed, because only the increments of $B$ enter. For $0\le s<t$ we write $\Delta_{s,t}\defeq B_t-B_s=M_t-M_s$, with components $\Delta^{(i)}_{s,t}$. By \bk{Definition 5.1}(ii), $\Delta_{s,t}$ is independent of $\F_s$ and normal with mean zero and covariance matrix $(t-s)I_d$; in particular
\begin{equation}\label{eq:2-5-05:mom}
  E\big[\Delta^{(i)}_{s,t}\big]=0,\qquad
  E\big[\Delta^{(i)}_{s,t}\Delta^{(j)}_{s,t}\big]=(t-s)\delta_{ij},\qquad 1\le i,j\le d .
\end{equation}

\emph{Step 1: each $M^{(i)}$ is a continuous, square-integrable martingale.} Every path of $B$ is continuous, hence so is every path of $M^{(i)}$, and $M^{(i)}_0=0$. For each $t\ge 0$ the random variable $M^{(i)}_t=B^{(i)}_t-B^{(i)}_0$ is $\F_t$-measurable, because $B$ is adapted and $\F_0\subseteq\F_t$. For $t>0$ we have $M^{(i)}_t=\Delta^{(i)}_{0,t}$, so \eqref{eq:2-5-05:mom} gives $E\big[(M^{(i)}_t)^2\big]=t<\infty$; thus $M^{(i)}_t$ is square-integrable, and a fortiori integrable. Finally, for $0\le s<t$, the integrable random variable $M^{(i)}_t-M^{(i)}_s=\Delta^{(i)}_{s,t}$ is independent of $\F_s$, so that
\[
  E\big[M^{(i)}_t-M^{(i)}_s\,\big|\,\F_s\big]=E\big[\Delta^{(i)}_{s,t}\big]=0,\quad\text{a.s. } P .
\]
Hence $\{M^{(i)}_t,\F_t;\ 0\le t<\infty\}$ is a continuous martingale with $M^{(i)}_0=0$ and $E\big[(M^{(i)}_t)^2\big]<\infty$ for every $t$, i.e., a continuous, square-integrable martingale in the sense of \bk{Definition 1.5.1}.

\emph{Step 2: $M^{(i)}_tM^{(j)}_t-t\delta_{ij}$ is a martingale.} Fix $1\le i,j\le d$ and set $Z_t\defeq M^{(i)}_tM^{(j)}_t-t\delta_{ij}$. This process is adapted, continuous, and integrable by the Cauchy--Schwarz inequality and Step~1. For $0\le s<t$ we have the algebraic identity
\[
  M^{(i)}_tM^{(j)}_t-M^{(i)}_sM^{(j)}_s
  =\Delta^{(i)}_{s,t}\Delta^{(j)}_{s,t}+M^{(i)}_s\Delta^{(j)}_{s,t}+M^{(j)}_s\Delta^{(i)}_{s,t},
\]
in which each of the three terms on the right is integrable (again by the Cauchy--Schwarz inequality). The product $\Delta^{(i)}_{s,t}\Delta^{(j)}_{s,t}$ is a Borel function of $\Delta_{s,t}$, hence independent of $\F_s$, and \eqref{eq:2-5-05:mom} yields $E\big[\Delta^{(i)}_{s,t}\Delta^{(j)}_{s,t}\,\big|\,\F_s\big]=(t-s)\delta_{ij}$ a.s. Since $M^{(i)}_s$ is $\F_s$-measurable and both $\Delta^{(j)}_{s,t}$ and $M^{(i)}_s\Delta^{(j)}_{s,t}$ are integrable,
\[
  E\big[M^{(i)}_s\Delta^{(j)}_{s,t}\,\big|\,\F_s\big]=M^{(i)}_s\,E\big[\Delta^{(j)}_{s,t}\,\big|\,\F_s\big]
  =M^{(i)}_s\,E\big[\Delta^{(j)}_{s,t}\big]=0,\quad\text{a.s.},
\]
and likewise with $i$ and $j$ interchanged. Adding up, we obtain
\begin{equation}\label{eq:2-5-05:mart}
  E\big[M^{(i)}_tM^{(j)}_t-M^{(i)}_sM^{(j)}_s\,\big|\,\F_s\big]=(t-s)\delta_{ij},\quad\text{a.s. } P;\qquad 0\le s<t<\infty,
\end{equation}
which is the martingale property of $\{Z_t,\F_t;\ 0\le t<\infty\}$.

\emph{Step 3: the cross-variations.} The quadratic variation and the cross-variation of \bk{Definitions 1.5.3} and \bk{1.5.5} are defined in \bk{Section 1.5} under the standing assumption that the filtration satisfies the usual conditions, whereas \bk{Definition 5.1} imposes no such condition on $\{\F_t\}$. We show first that this causes no difficulty.

Let $\ol{\F}$ be the completion of $\F$ with respect to $P$ (we denote the extended measure again by $P$), let $\cN$ be the collection of the $P$-null sets of $\ol{\F}$, and set
\[
  \wt{\F}_t\defeq\bigcap_{u>t}\sigma(\F_u\cup\cN),\qquad 0\le t<\infty .
\]
Then $\{\wt{\F}_t\}$ is a filtration on $(\Omega,\ol{\F},P)$ with $\F_t\subseteq\wt{\F}_t$; it is right-continuous, since $\bigcap_{v>t}\wt{\F}_v=\bigcap_{v>t}\bigcap_{u>v}\sigma(\F_u\cup\cN)=\bigcap_{u>t}\sigma(\F_u\cup\cN)=\wt{\F}_t$, and $\wt{\F}_0$ contains $\cN$; that is, $\{\wt{\F}_t\}$ satisfies the usual conditions. (If $\{\F_t\}$ itself satisfies the usual conditions on a complete space, then $\wt{\F}_t=\F_t$.)

\begin{claim}
$\{B_t,\wt{\F}_t;\ 0\le t<\infty\}$ is again a $d$-dimensional Brownian motion with initial distribution $\mu$.
\end{claim}
\begin{proof}
The process $B$ is continuous and adapted to $\{\wt{\F}_t\}\supseteq\{\F_t\}$, and conditions (i) and the distributional part of (ii) in \bk{Definition 5.1} do not involve the filtration. It remains to show that $\Delta_{s,t}$ is independent of $\wt{\F}_s$ for $0\le s<t$.

For $u\ge 0$, the collection of all $A\in\ol{\F}$ for which there exists $A'\in\F_u$ with $P(A\,\triangle\,A')=0$ is a $\sigma$-field (it is closed under complements and countable unions, because $A^c\triangle(A')^c=A\triangle A'$ and $\big(\bigcup_nA_n\big)\triangle\big(\bigcup_nA_n'\big)\subseteq\bigcup_n(A_n\triangle A_n')$), and it contains $\F_u$ and $\cN$; therefore it contains $\sigma(\F_u\cup\cN)$. Now let $A\in\wt{\F}_s$, let $f:\R^d\to\R$ be bounded and continuous, and let $s<u<t$. Since $A\in\sigma(\F_u\cup\cN)$, there is $A'\in\F_u$ with $P(A\,\triangle\,A')=0$, and the independence of $B_t-B_u$ from $\F_u$ gives
\[
  E\big[f(B_t-B_u)\ind_A\big]=E\big[f(B_t-B_u)\ind_{A'}\big]=E\big[f(B_t-B_u)\big]P(A')=E\big[f(B_t-B_u)\big]P(A).
\]
As $u\downarrow s$ we have $f(B_t-B_u)\to f(B_t-B_s)$ pointwise, by the continuity of the paths of $B$ and of $f$, and the bounded convergence theorem yields
\begin{equation}\label{eq:2-5-05:ind}
  E\big[f(\Delta_{s,t})\ind_A\big]=E\big[f(\Delta_{s,t})\big]P(A).
\end{equation}
For a closed set $F\subseteq\R^d$, the bounded continuous functions $f_n(x)\defeq\big(1-n\,\mathrm{dist}(x,F)\big)^+$ decrease to $\ind_F$ as $n\to\infty$, so \eqref{eq:2-5-05:ind} and the bounded convergence theorem give $P[\Delta_{s,t}\in F,\ A]=P[\Delta_{s,t}\in F]\,P(A)$. The collection of all $\Gamma\in\B(\R^d)$ with $P[\Delta_{s,t}\in\Gamma,\ A]=P[\Delta_{s,t}\in\Gamma]\,P(A)$ is a Dynkin system, and it contains the closed sets, which are closed under finite intersections and generate $\B(\R^d)$. By the Dynkin system theorem (\bk{Theorem 1.3}) it is all of $\B(\R^d)$. Since $A\in\wt{\F}_s$ was arbitrary, $\Delta_{s,t}$ is independent of $\wt{\F}_s$.
\end{proof}

Steps 1 and 2 used nothing about $\{\F_t\}$ beyond \bk{Definition 5.1}, so by the Claim they apply verbatim with $\{\wt{\F}_t\}$ in place of $\{\F_t\}$: on $(\Omega,\ol{\F},P)$, with the filtration $\{\wt{\F}_t\}$ satisfying the usual conditions, each $M^{(i)}$ belongs to $\cM_2^c$, and $\{M^{(i)}_tM^{(j)}_t-t\delta_{ij},\wt{\F}_t;\ 0\le t<\infty\}$ is a martingale. The deterministic process $A_t\defeq t\delta_{ij}$ is adapted, continuous, of bounded variation on every bounded interval (it is either identically zero or increasing), and $A_0=0$. By \bk{Theorem 1.5.13}, the cross-variation $\qv{M^{(i)},M^{(j)}}$ of \bk{Definition 1.5.5} is, up to indistinguishability, the unique process with these properties for which $M^{(i)}M^{(j)}-A$ is a martingale. Therefore
\[
  \qv{M^{(i)},M^{(j)}}_t=t\,\delta_{ij};\qquad 0\le t<\infty,\ 1\le i,j\le d,\quad\text{a.s. } P,
\]
and in particular $\qv{M^{(i)}}_t=\qv{M^{(i)},M^{(i)}}_t=t$. When $\{\F_t\}$ itself satisfies the usual conditions this is the assertion of the problem as it stands; in general, \eqref{eq:2-5-05:mart} shows that $M^{(i)}M^{(j)}-\qv{M^{(i)},M^{(j)}}$ is a martingale also with respect to the original filtration $\{\F_t\}$, the cross-variation $t\delta_{ij}$ being deterministic and hence adapted to any filtration.

\emph{Step 4: $M$ is independent of $\F_0$.} The assertion means that the $\sigma$-field $\sigma(M)\defeq\sigma(M_t;\ 0\le t<\infty)$ is independent of $\F_0$. (Regarding $M$ as a random variable with values in $\big(C[0,\infty)^d,\B(C[0,\infty)^d)\big)$, we have $\sigma(M)=\{\{M\in G\};\ G\in\B(C[0,\infty)^d)\}$, because the Borel $\sigma$-field of $C[0,\infty)^d$ is generated by the finite-dimensional cylinder sets; this is \bk{Problem 4.2}, whose proof carries over to $d$ dimensions.)

Fix $n\ge 1$ and $0=t_0<t_1<\cdots<t_n$, and write $\Delta_k\defeq\Delta_{t_{k-1},t_k}$, $1\le k\le n$. We show by induction on $m\in\{0,1,\ldots,n\}$ that for all $A\in\F_0$ and $\Gamma_1,\ldots,\Gamma_n\in\B(\R^d)$,
\begin{equation}\label{eq:2-5-05:prod}
  P\Big[A\cap\bigcap_{k=1}^m\{\Delta_k\in\Gamma_k\}\Big]=P(A)\prod_{k=1}^mP[\Delta_k\in\Gamma_k].
\end{equation}
For $m=0$ this is trivial. If \eqref{eq:2-5-05:prod} holds for $m-1$, then, since the event $A\cap\bigcap_{k=1}^{m-1}\{\Delta_k\in\Gamma_k\}$ belongs to $\F_{t_{m-1}}$ ($B$ is adapted) and $\Delta_m$ is independent of $\F_{t_{m-1}}$,
\[
  P\Big[A\cap\bigcap_{k=1}^m\{\Delta_k\in\Gamma_k\}\Big]
  =P\Big[A\cap\bigcap_{k=1}^{m-1}\{\Delta_k\in\Gamma_k\}\Big]\,P[\Delta_m\in\Gamma_m]
  =P(A)\prod_{k=1}^mP[\Delta_k\in\Gamma_k].
\]
Taking $A=\Omega$ in \eqref{eq:2-5-05:prod} with $m=n$ shows $P\big[\bigcap_{k=1}^n\{\Delta_k\in\Gamma_k\}\big]=\prod_{k=1}^nP[\Delta_k\in\Gamma_k]$, and so \eqref{eq:2-5-05:prod} reads
\[
  P\big[A\cap\{(\Delta_1,\ldots,\Delta_n)\in\Gamma_1\times\cdots\times\Gamma_n\}\big]
  =P(A)\,P\big[(\Delta_1,\ldots,\Delta_n)\in\Gamma_1\times\cdots\times\Gamma_n\big].
\]
For fixed $A\in\F_0$, the collection of all $D\in\B((\R^d)^n)$ such that $P[A\cap\{(\Delta_1,\ldots,\Delta_n)\in D\}]=P(A)\,P[(\Delta_1,\ldots,\Delta_n)\in D]$ is a Dynkin system containing the measurable rectangles, which are closed under finite intersections and generate $\B((\R^d)^n)$; by the Dynkin system theorem it equals $\B((\R^d)^n)$. Because $M_{t_k}=\Delta_1+\cdots+\Delta_k$, the vector $(M_{t_1},\ldots,M_{t_n})$ is the image of $(\Delta_1,\ldots,\Delta_n)$ under a continuous (linear) map of $(\R^d)^n$ into itself, and we conclude that
\begin{equation}\label{eq:2-5-05:cyl}
  P\big[A\cap\{(M_{t_1},\ldots,M_{t_n})\in D\}\big]=P(A)\,P\big[(M_{t_1},\ldots,M_{t_n})\in D\big];
  \qquad A\in\F_0,\ D\in\B((\R^d)^n).
\end{equation}

Let $\cC$ be the collection of all events $\{(M_{t_1},\ldots,M_{t_n})\in D\}$ with $n\ge1$, $0<t_1<\cdots<t_n$ and $D\in\B((\R^d)^n)$. Since $M_0=0$, these events generate $\sigma(M)$. The collection $\cC$ is closed under finite intersections: two such events, based on the time sets $S_1$ and $S_2$, can both be written as events of the same form based on the time set $S_1\cup S_2$, and then their intersection corresponds to the intersection of the two Borel sets. For fixed $A\in\F_0$, the collection of all $C\in\sigma(M)$ with $P(A\cap C)=P(A)P(C)$ is a Dynkin system, and it contains $\cC$ by \eqref{eq:2-5-05:cyl}; by the Dynkin system theorem it contains $\sigma(\cC)=\sigma(M)$. Thus $P(A\cap C)=P(A)P(C)$ for all $A\in\F_0$ and $C\in\sigma(M)$, i.e., $M$ is independent of $\F_0$.

\rmk{Step~4, applied with $\{\wt{\F}_t\}$ in place of $\{\F_t\}$ (which is legitimate by the Claim), shows that $M$ is even independent of $\wt{\F}_0\supseteq\F_{0+}$. In particular $M=B-B_0$ is independent of $B_0$, and, since $M_0=0$ and $M_t-M_s=B_t-B_s$, the process $\{M_t,\F_t;\ 0\le t<\infty\}$ is a $d$-dimensional Brownian motion starting at the origin.}
\end{kex}

\begin{kex}{Problem}{5.9}{}
Equation \bk{(5.5)} is proved in full in the book's Section 2.10; the book only asserts that \bk{(5.6)} can be proved similarly, and we prove \bk{(5.6)} here.

\emph{The meaning of \bk{(5.6)}.} The random variable $P[X+Y\in\Gamma\mid Y]$ is, by definition, any $\sigma(Y)$-measurable version of the conditional expectation of $\ind_{\{X+Y\in\Gamma\}}$ given $\sigma(Y)$. Every $\sigma(Y)$-measurable real random variable has the form $h(Y)$ for some Borel-measurable $h:\R^d\to\R$ (the Doob--Dynkin lemma), and one writes $P[X+Y\in\Gamma\mid Y=y]\defeq h(y)$ for any Borel-measurable $h$ such that
\begin{equation}\label{eq:2-5-09:def}
  h(Y)=P[X+Y\in\Gamma\mid Y],\quad\text{a.s. } P .
\end{equation}
This function is determined only up to a $PY^{-1}$-null set: if $h_1,h_2$ are Borel-measurable and both satisfy \eqref{eq:2-5-09:def}, then $h_1(Y)=h_2(Y)$ a.s., so $PY^{-1}\{y;\,h_1(y)\neq h_2(y)\}=P[h_1(Y)\neq h_2(Y)]=0$ by \bk{(4.2)}. Consequently, \bk{(5.6)} is the assertion that the function
\[
  g(y)\defeq P[X+y\in\Gamma],\qquad y\in\R^d,
\]
is Borel-measurable and satisfies \eqref{eq:2-5-09:def}; every other admissible $h$ then agrees with $g$ for $PY^{-1}$-a.e.\ $y$.

\emph{Step 1: an identity for sets in the product space.} For $D\in\B(\R^d\times\R^d)$ and $y\in\R^d$, the section $D^y\defeq\{x\in\R^d;\ (x,y)\in D\}$ is the inverse image of $D$ under the continuous map $x\mapsto(x,y)$, hence $D^y\in\B(\R^d)$, and we may define
\[
  g_D(y)\defeq P[(X,y)\in D]=P[X\in D^y]\in[0,1],\qquad y\in\R^d .
\]
Let $\cD$ be the collection of all $D\in\B(\R^d\times\R^d)$ such that $g_D$ is Borel-measurable and
\begin{equation}\label{eq:2-5-09:id}
  P[(X,Y)\in D,\ Y\in C]=\int_C g_D(y)\,PY^{-1}(dy),\qquad\forall\, C\in\B(\R^d).
\end{equation}
We claim that $\cD=\B(\R^d\times\R^d)$.

\textbf{(a)} $\cD$ contains every measurable rectangle $D=A\times B$ with $A,B\in\B(\R^d)$. Indeed, $D^y=A$ if $y\in B$ and $D^y=\emptyset$ otherwise, so $g_D(y)=P[X\in A]\,\ind_B(y)$ is Borel-measurable, and the right-hand side of \eqref{eq:2-5-09:id} equals $P[X\in A]\,PY^{-1}(B\cap C)=P[X\in A]\,P[Y\in B\cap C]$. On the other hand, $\{Y\in B\cap C\}\in\G$ because $Y$ is $\G$-measurable, and $X$ is independent of $\G$; therefore the left-hand side is
\[
  P[X\in A,\ Y\in B\cap C]=P[X\in A]\,P[Y\in B\cap C].
\]
In particular $\R^d\times\R^d\in\cD$.

\textbf{(b)} $\cD$ is closed under proper differences. Let $D_1,D_2\in\cD$ with $D_1\subseteq D_2$ and $D=D_2\setminus D_1$. Then $D^y=D_2^y\setminus D_1^y$ with $D_1^y\subseteq D_2^y$, so $g_D=g_{D_2}-g_{D_1}$ is Borel-measurable; and subtracting \eqref{eq:2-5-09:id} for $D_1$ from \eqref{eq:2-5-09:id} for $D_2$ (all quantities are finite) gives \eqref{eq:2-5-09:id} for $D$.

\textbf{(c)} $\cD$ is closed under increasing countable unions. Let $D_1\subseteq D_2\subseteq\cdots$ be in $\cD$ and $D=\bigcup_n D_n$. Then $D_n^y\uparrow D^y$ for every $y$, so $g_{D_n}(y)\uparrow g_D(y)$ by continuity from below of $P$; thus $g_D$ is Borel-measurable as a pointwise limit of Borel-measurable functions. The events $\{(X,Y)\in D_n,\ Y\in C\}$ increase to $\{(X,Y)\in D,\ Y\in C\}$, so the left-hand sides of \eqref{eq:2-5-09:id} for $D_n$ converge to that for $D$, and the right-hand sides converge by the monotone convergence theorem.

By (a)--(c), $\cD$ is a Dynkin system containing the measurable rectangles, which form a collection closed under finite intersections and generating $\B(\R^d)\otimes\B(\R^d)=\B(\R^d\times\R^d)$ (the equality of the product and the Borel $\sigma$-fields holds because $\R^d$ is separable). The Dynkin system theorem (\bk{Theorem 1.3}) yields $\cD=\B(\R^d\times\R^d)$.

\emph{Step 2: proof of \bk{(5.6)}.} Fix $\Gamma\in\B(\R^d)$ and let $D\defeq\{(x,y)\in\R^d\times\R^d;\ x+y\in\Gamma\}$. This is the inverse image of $\Gamma$ under the continuous map $(x,y)\mapsto x+y$, so $D\in\B(\R^d\times\R^d)=\cD$. For this $D$ we have $\{(X,Y)\in D\}=\{X+Y\in\Gamma\}$ and $g_D(y)=P[X+y\in\Gamma]=g(y)$. Step~1 therefore shows that $g$ is Borel-measurable and that, for every $C\in\B(\R^d)$,
\[
  E\big[\ind_{\{X+Y\in\Gamma\}}\ind_{\{Y\in C\}}\big]
  =\int_C g(y)\,PY^{-1}(dy)
  =E\big[g(Y)\,\ind_{\{Y\in C\}}\big],
\]
the second equality being the change-of-variables formula for the induced measure $PY^{-1}$ of \bk{(4.2)}, applied to the bounded Borel function $g\,\ind_C$. Since $\sigma(Y)=\{\{Y\in C\};\ C\in\B(\R^d)\}$, this says that
\[
  E\big[\ind_{\{X+Y\in\Gamma\}}\ind_F\big]=E\big[g(Y)\,\ind_F\big],\qquad\forall\,F\in\sigma(Y).
\]
Moreover $g(Y)$ is bounded and $\sigma(Y)$-measurable, as the composition of the $\sigma(Y)$-measurable map $Y$ with the Borel function $g$. Hence $g(Y)$ is a version of $P[X+Y\in\Gamma\mid Y]$; that is, $g$ satisfies \eqref{eq:2-5-09:def}. By the discussion preceding Step~1, this is \bk{(5.6)}: every version $h$ of $y\mapsto P[X+Y\in\Gamma\mid Y=y]$ satisfies $h(y)=P[X+y\in\Gamma]$ for $PY^{-1}$-a.e.\ $y\in\R^d$.

\rmk{Combining \bk{(5.5)} and \bk{(5.6)}, we obtain $P[X+Y\in\Gamma\mid\G]=g(Y)$ a.s., with $g(y)=P[X+y\in\Gamma]$. With $X=B_t-B_s$, $Y=B_s$ and $\G=\F_s$ for a Brownian motion $B$ with initial distribution $\mu$, these are precisely \bk{(5.3)} and \bk{(5.4)}, because $P^\mu[(B_t-B_s)+y\in\Gamma]=P^y[B_{t-s}\in\Gamma]$.}
\end{kex}

\begin{kex}{Exercise}{5.17}{}
Throughout we fix $x\in\R^d$, $s\ge 0$ and $F\in\F^X_\infty$, and we write $A\defeq\theta_s^{-1}F$. Conditional probabilities given $X_s$ are conditional probabilities given the $\sigma$-field $\sigma(X_s)$, and all conditional expectations below are of bounded random variables.

\emph{Preliminaries.} Since $X$ is adapted to $\{\F_t\}$ and $\F_t\subseteq\F$, every $X_t$ is $\F$-measurable, so $\F^X_\infty=\sigma(X_t;\,t\ge 0)\subseteq\F$; as $\theta_s$ is $\F/\F$-measurable, we get $A\in\F$, and all the conditional probabilities in (c$''$) and (c$'''$) are well defined. Adaptedness also gives
\begin{equation}\label{eq:2-5-17:a}
  \sigma(X_s)\subseteq\F_s .
\end{equation}
Finally, (c$''$) does hold for a Markov family with shift-operators: it is \bk{Proposition 5.15}\,(c$'$) rewritten by means of the identity $\{X_{s+\cdot}\in \Gamma\}=\theta_s^{-1}\{X_\cdot\in \Gamma\}$, $\Gamma\in\B\bigl((\R^d)^{[0,\infty)}\bigr)$, which follows from \bk{(5.14)}, as explained in the book just before \bk{Theorem 5.16}.

\textbf{(i) (c$''$) implies (c$'''$).} Let $G\in\F_s$, and let $Z$ be a fixed version of $P^x[A\mid X_s]$; thus $Z$ is $\sigma(X_s)$-measurable and we may take $0\le Z\le 1$. By (c$''$), $Z$ is also a version of $P^x[A\mid\F_s]$. Using successively the tower property (justified by \eqref{eq:2-5-17:a}), the $\F_s$-measurability of $\ind_G$, relation (c$''$), and the $\sigma(X_s)$-measurability of the bounded random variable $Z$, we obtain, $P^x$-a.s.,
\begin{align*}
  P^x[G\cap A\mid X_s]
  &=E^x\bigl[\,E^x[\ind_G\ind_A\mid\F_s]\bigm| X_s\bigr]
   =E^x\bigl[\ind_G\,P^x[A\mid\F_s]\bigm| X_s\bigr]\\
  &=E^x[\ind_G\,Z\mid X_s]
   =Z\,E^x[\ind_G\mid X_s]
   =P^x[G\mid X_s]\,P^x[A\mid X_s].
\end{align*}
This is (c$'''$).

\textbf{(ii) (c$'''$) implies (c$''$).} For this implication we use only that $X$ is adapted and that $\theta_s$ is $\F/\F$-measurable (the preliminaries above, without the last sentence), and we assume that (c$'''$) holds, for the fixed $x$, $s$, $F$, for every $G\in\F_s$. Let again $Z$ be a version of $P^x[A\mid X_s]$ with $0\le Z\le 1$. We have to show that $Z$ is a version of $P^x[A\mid\F_s]$. First, $Z$ is $\sigma(X_s)$-measurable, hence $\F_s$-measurable by \eqref{eq:2-5-17:a}. Second, for every $G\in\F_s$, since $Z$ is bounded and $\sigma(X_s)$-measurable,
\begin{align*}
  E^x[\ind_G\,Z]
  &=E^x\bigl[E^x[\ind_G Z\mid X_s]\bigr]
  =E^x\bigl[Z\,P^x[G\mid X_s]\bigr]\\
  &=E^x\bigl[P^x[G\cap A\mid X_s]\bigr]
  =P^x(G\cap A),
\end{align*}
where the third equality is (c$'''$). Thus $\int_G Z\,dP^x=P^x(G\cap A)$ for every $G\in\F_s$, and these two properties characterize $P^x[A\mid\F_s]$ up to $P^x$-null sets. Therefore
\[
  P^x[\theta_s^{-1}F\mid\F_s]=Z=P^x[\theta_s^{-1}F\mid X_s],\qquad P^x\text{-a.s.},
\]
which is (c$''$).

\rmk{The argument is the general fact that, for $\sigma$-fields $\cH\subseteq\G$ and an event $A$ on a probability space, $P[A\mid\G]=P[A\mid\cH]$ a.s.\ if and only if $P[G\cap A\mid\cH]=P[G\mid\cH]\,P[A\mid\cH]$ a.s.\ for every $G\in\G$; here $\cH=\sigma(X_s)$ and $\G=\F_s$. In particular, in (ii) the equivalence holds separately for each triple $(x,s,F)$.}
\end{kex}

\begin{kex}{Problem}{5.18}{}
We read the hypothesis in the sense of \bk{Definition 5.10}: $X$ is an $\{\F_t\}$-adapted, $d$-dimensional process
on $(\Omega,\F,P)$, its initial distribution is $\mu\defeq PX_0^{-1}$ (so that condition (i) of
\bk{Definition 5.10} holds by the choice of $\mu$), and for all $s,t\ge0$ and $\Gamma\in\B(\R^d)$,
\begin{equation}\label{eq:2-5-18:a}
P[X_{t+s}\in\Gamma\mid\F_s]=P[X_{t+s}\in\Gamma\mid X_s],\qquad P\text{-a.s.},
\end{equation}
where conditioning on $X_s$ means conditioning on the $\sigma$-field
$\sigma(X_s)=\{\{X_s\in\Lambda\};\ \Lambda\in\B(\R^d)\}\subseteq\F_s$. We have to show that $Y$ is adapted to
$\{\F_t\}$ and satisfies (i) and (ii) of \bk{Definition 5.10} for the same measure $P$ and a suitable initial
distribution $\nu$.

For $u\ge0$ let $A_u\colon\R^d\to\R^d$ be the affine map $A_u(x)\defeq\varphi(u)+\Psi(u)x$, so that
$Y_u=A_u(X_u)$. Because $\Psi(u)$ is nonsingular, $A_u$ is a bijection of $\R^d$ with inverse
$A_u^{-1}(y)=\Psi(u)^{-1}\bigl(y-\varphi(u)\bigr)$. Both $A_u$ and $A_u^{-1}$ are continuous, hence
$\B(\R^d)/\B(\R^d)$-measurable. Consequently, for every $\Gamma\in\B(\R^d)$,
\begin{equation}\label{eq:2-5-18:b}
\{Y_u\in\Gamma\}=\{X_u\in A_u^{-1}(\Gamma)\}
\qquad\text{and}\qquad
\{X_u\in\Gamma\}=\{Y_u\in A_u(\Gamma)\},
\end{equation}
where $A_u^{-1}(\Gamma)$ is the preimage of $\Gamma$ under the Borel map $A_u$, and
$A_u(\Gamma)=(A_u^{-1})^{-1}(\Gamma)$ is the preimage of $\Gamma$ under the Borel map $A_u^{-1}$; both sets belong
to $\B(\R^d)$.

\emph{Adaptedness.} By the first identity in \eqref{eq:2-5-18:b} and the adaptedness of $X$, we have
$\{Y_u\in\Gamma\}\in\F_u$ for every $u\ge0$ and $\Gamma\in\B(\R^d)$; thus $Y$ is a $d$-dimensional process adapted
to $\{\F_t\}$.

\emph{Initial distribution.} Define $\nu\defeq PY_0^{-1}$, a probability measure on $(\R^d,\B(\R^d))$; then (i) of
\bk{Definition 5.10} holds for $Y$ and $\nu$. Explicitly, since $\varphi(0)=0$ we have $Y_0=\Psi(0)X_0$ and
$\nu(\Gamma)=\mu\bigl(\Psi(0)^{-1}\Gamma\bigr)$ for $\Gamma\in\B(\R^d)$.

\emph{The Markov property.} The two identities in \eqref{eq:2-5-18:b} show that every set of $\sigma(Y_u)$
belongs to $\sigma(X_u)$ and conversely, that is,
\begin{equation}\label{eq:2-5-18:c}
\sigma(Y_u)=\sigma(X_u),\qquad u\ge0 .
\end{equation}
Fix now $s,t\ge0$ and $\Gamma\in\B(\R^d)$, and put $\Gamma'\defeq A_{t+s}^{-1}(\Gamma)\in\B(\R^d)$. Using, in this
order, \eqref{eq:2-5-18:b} with $u=t+s$, the Markov property \eqref{eq:2-5-18:a} of $X$ applied to the Borel set
$\Gamma'$, the equality \eqref{eq:2-5-18:c} with $u=s$, and \eqref{eq:2-5-18:b} once more, we obtain $P$-a.s.
\begin{align*}
P[Y_{t+s}\in\Gamma\mid\F_s]
&=P[X_{t+s}\in\Gamma'\mid\F_s]
=P[X_{t+s}\in\Gamma'\mid\sigma(X_s)]\\
&=P[X_{t+s}\in\Gamma'\mid\sigma(Y_s)]
=P[Y_{t+s}\in\Gamma\mid Y_s].
\end{align*}
(The first and the last equalities are between conditional probabilities of one and the same event, and the
third is between conditional probabilities given one and the same $\sigma$-field.) This is condition (ii) of
\bk{Definition 5.10} for $Y$. Hence $Y=\{Y_t,\F_t;\ t\ge0\}$ is a Markov process on $(\Omega,\F,P)$ with initial
distribution $\nu$.

\rmk{(a) The hypothesis $\varphi(0)=0$ was used only to write $\nu$ in terms of $\mu$; without it the proof goes
through verbatim, with $\nu(\Gamma)=\mu\bigl(\Psi(0)^{-1}(\Gamma-\varphi(0))\bigr)$. No regularity of
$\varphi$ and $\Psi$ in $t$ is needed either. The normalization $\varphi(0)=0$ matters for
\bk{Definition 5.19}, where it makes $Y_0=\sigma B_0$, so that $Y$ starts at $x$ under $P^{\sigma^{-1}x}$.

(b) Nonsingularity of $\Psi(t)$ for \emph{every} $t$ cannot be dropped: what the proof needs is
$\sigma(Y_s)\supseteq\sigma(X_s)$. Let $d=1$, let $X=B$ be a standard one-dimensional Brownian motion with
$B_0=0$ and the filtration $\{\F^B_t\}$ (a Markov process by \bk{Theorem 5.12}), let $\varphi\equiv0$, and let
$\Psi(t)=1$ for $t\neq1$, $\Psi(1)=0$. Then $Y_1=0$, so $\sigma(Y_1)$ is trivial and
$P[Y_2>0\mid Y_1]=P[B_2>0]=\tfrac12$ a.s. On the other hand $B_2-B_1$ is independent of $\F^B_1$ and normal with
mean zero and variance one, so \bk{Problem 5.9} gives
$P[Y_2>0\mid\F^B_1]=P[B_2>0\mid\F^B_1]=\Phi(B_1)$ a.s., where
$\Phi(y)=P[B_2-B_1>-y]$ is the standard normal distribution function. Since $\Phi(B_1)\neq\tfrac12$ on
$\{B_1\neq0\}$, an event of probability one, condition (ii) of \bk{Definition 5.10} fails for $Y$ with
$s=t=1$, $\Gamma=(0,\infty)$.}
\end{kex}

\begin{kex}{Exercise}{5.21}{}
Let $N=\{N_t,\F_t;\ t\ge0\}$, $(\Omega,\F)$, $\{P^x\}_{x\in\R}$ be a Poisson family with intensity $\lambda>0$, and $\tilde N_t\defeq N_t-N_0$. We have to verify conditions (a)--(d) of \bk{Definition 5.11} with $d=1$ and $X=N$. The process $N$ is adapted by assumption, condition (a) is \bk{Definition 5.20}(i), and condition (b) is \bk{Definition 5.20}(ii). For $t\ge0$ and $k=0,1,2,\dots$ we write
\[
p_k(t)\defeq e^{-\lambda t}\,\frac{(\lambda t)^k}{k!}\quad(t>0),\qquad p_k(0)\defeq \ind_{\{k=0\}} ,
\]
and, for $\Gamma\in\B(\R)$,
\begin{equation}\label{eq:2-5-21:h}
h_{t,\Gamma}(y)\defeq\sum_{k=0}^\infty p_k(t)\,\ind_\Gamma(y+k),\qquad y\in\R .
\end{equation}
Each $y\mapsto\ind_\Gamma(y+k)$ is Borel-measurable, so $h_{t,\Gamma}:\R\to[0,1]$ is Borel-measurable.

\emph{Step 1: the increments.} Fix $x\in\R$ and $s,t\ge0$, and put
\[
X\defeq N_{t+s}-N_s=\tilde N_{t+s}-\tilde N_s,\qquad Y\defeq N_s .
\]
By \bk{Definition 5.20}(iii) and \bk{Definition 1.3.3}, under $P^x$ the process $\{\tilde N_u,\F_u;\ u\ge0\}$ is a Poisson process with intensity $\lambda$; hence, if $t>0$, the random variable $X$ is independent of $\F_s$ under $P^x$ and $P^x[X=k]=p_k(t)$ for $k=0,1,2,\dots$. If $t=0$, then $X\equiv0$, which is again independent of $\F_s$ and satisfies $P^x[X=k]=p_k(0)$. In both cases $\sum_{k\ge0}P^x[X=k]=1$. The same argument with $s=0$, applied under $P^y$ and using $\tilde N_0\equiv0$, gives
\begin{equation}\label{eq:2-5-21:law}
P^y[\tilde N_t=k]=p_k(t),\qquad k=0,1,2,\dots,\quad y\in\R,\ t\ge0 .
\end{equation}

\emph{Step 2: condition \textup{(c)}.} The random variable $Y=N_s$ is $\F_s$-measurable, and $X$ is independent of $\F_s$ under $P^x$. \bk{Problem 5.9}, equation \bk{(5.5)}, applied on $(\Omega,\F,P^x)$ with $\G=\F_s$, yields for every $\Gamma\in\B(\R)$
\[
P^x[N_{t+s}\in\Gamma\mid\F_s]=P^x[X+Y\in\Gamma\mid\F_s]=P^x[X+Y\in\Gamma\mid Y]=P^x[N_{t+s}\in\Gamma\mid N_s],\qquad P^x\text{-a.s.},
\]
which is (c).

\emph{Step 3: condition \textup{(d)}.} By \bk{Problem 5.9}, equation \bk{(5.6)}, in the notation of \bk{(4.2)},
\begin{equation}\label{eq:2-5-21:d1}
P^x[N_{t+s}\in\Gamma\mid N_s=y]=P^x[X+y\in\Gamma]\qquad\text{for } P^xN_s^{-1}\text{-a.e. } y\in\R .
\end{equation}
For every $y\in\R$, since $P^x[X\in\{0,1,2,\dots\}]=1$, Step 1 and \eqref{eq:2-5-21:h} give
\[
P^x[X+y\in\Gamma]=\sum_{k=0}^\infty P^x[X=k]\,\ind_\Gamma(y+k)=h_{t,\Gamma}(y).
\]
On the other hand, $P^y[N_0=y]=1$ by \bk{Definition 5.20}(ii), so that $N_t=y+\tilde N_t$ holds $P^y$-a.s., and \eqref{eq:2-5-21:law} gives, for every $y\in\R$,
\[
P^y[N_t\in\Gamma]=P^y[y+\tilde N_t\in\Gamma]=\sum_{k=0}^\infty P^y[\tilde N_t=k]\,\ind_\Gamma(y+k)=h_{t,\Gamma}(y).
\]
Thus $P^x[X+y\in\Gamma]=P^y[N_t\in\Gamma]$ for \emph{every} $y\in\R$, and \eqref{eq:2-5-21:d1} becomes
\[
P^x[N_{t+s}\in\Gamma\mid N_s=y]=P^y[N_t\in\Gamma]\qquad\text{for } P^xN_s^{-1}\text{-a.e. } y\in\R,
\]
which is (d). Hence the Poisson family is a Markov family.

\rmk{Steps 2 and 3 can also be checked by hand, without \bk{Problem 5.9}: for $A\in\F_s$, the independence of $X$ and $\F_s$ gives
$P^x[\{X+Y\in\Gamma\}\cap A]=\sum_{k\ge0}P^x[X=k]\,P^x[\{Y+k\in\Gamma\}\cap A]=E^x\bigl[\ind_A\,h_{t,\Gamma}(N_s)\bigr]$,
the interchange of sum and expectation being justified by monotone convergence. Since $h_{t,\Gamma}(N_s)$ is $\sigma(N_s)$-measurable, it is a version of both conditional probabilities in \textup{(c)}, and $h_{t,\Gamma}$ is a version of $y\mapsto P^x[N_{t+s}\in\Gamma\mid N_s=y]$. In particular, in the notation of \bk{(5.7)}, $U_t\ind_\Gamma=h_{t,\Gamma}$ is Borel-measurable and does not depend on $x$.}

\emph{Step 4: independence of $\F^{\tilde N}_\infty$ and $\F_0$.} Fix $x\in\R$; all probabilities below are taken under $P^x$. Recall that $\F^{\tilde N}_\infty=\sigma(\tilde N_t;\ 0\le t<\infty)$. Since $N$ is adapted and $\F_0\subseteq\F_t$, each $\tilde N_t=N_t-N_0$ is $\F_t$-measurable.

Let $0=t_0<t_1<\dots<t_n$ and $D_i\defeq\tilde N_{t_i}-\tilde N_{t_{i-1}}$, $1\le i\le n$. We claim that for all $A\in\F_0$ and $\Gamma_1,\dots,\Gamma_n\in\B(\R)$,
\begin{equation}\label{eq:2-5-21:prod}
P^x\bigl[A\cap\{D_1\in\Gamma_1,\dots,D_n\in\Gamma_n\}\bigr]=P^x(A)\prod_{i=1}^nP^x[D_i\in\Gamma_i].
\end{equation}
We argue by induction on $n$, the case $n=0$ being trivial. For $n\ge1$, the event $A\cap\{D_1\in\Gamma_1,\dots,D_{n-1}\in\Gamma_{n-1}\}$ belongs to $\F_{t_{n-1}}$, because $A\in\F_0\subseteq\F_{t_{n-1}}$ and $D_1,\dots,D_{n-1}$ are $\F_{t_{n-1}}$-measurable; and $D_n$ is independent of $\F_{t_{n-1}}$, since $\tilde N$ is a Poisson process with respect to $\{\F_t\}$ under $P^x$ (\bk{Definition 1.3.3}). Therefore
\[
P^x\bigl[A\cap\{D_1\in\Gamma_1,\dots,D_n\in\Gamma_n\}\bigr]
=P^x\bigl[A\cap\{D_1\in\Gamma_1,\dots,D_{n-1}\in\Gamma_{n-1}\}\bigr]\,P^x[D_n\in\Gamma_n],
\]
and the induction hypothesis gives \eqref{eq:2-5-21:prod}. Taking $A=\Omega$ in \eqref{eq:2-5-21:prod} and substituting back, we obtain
\begin{equation}\label{eq:2-5-21:ind}
P^x(A\cap C)=P^x(A)\,P^x(C)
\end{equation}
for all $A\in\F_0$ and all $C$ in the class $\cP_{t_1,\dots,t_n}$ of events $\{D_1\in\Gamma_1,\dots,D_n\in\Gamma_n\}$, $\Gamma_i\in\B(\R)$. This class is closed under finite intersections and generates
\[
\G_{t_1,\dots,t_n}\defeq\sigma(D_1,\dots,D_n)=\sigma(\tilde N_{t_1},\dots,\tilde N_{t_n}),
\]
the last equality because $\tilde N_{t_0}=\tilde N_0\equiv0$, so that $\tilde N_{t_k}=D_1+\dots+D_k$ and $D_k=\tilde N_{t_k}-\tilde N_{t_{k-1}}$. For fixed $A\in\F_0$, the collection of all $C\in\F$ satisfying \eqref{eq:2-5-21:ind} is a Dynkin system (it contains $\Omega$ and is closed under proper differences and increasing countable unions), so by the Dynkin System Theorem (\bk{Theorem 1.3}) \eqref{eq:2-5-21:ind} holds for all $A\in\F_0$ and $C\in\G_{t_1,\dots,t_n}$.

Finally, let $\cP$ be the union of the $\sigma$-fields $\G_{t_1,\dots,t_n}$ over all $n\ge1$ and all $0<t_1<\dots<t_n$. If $C\in\G_{t_1,\dots,t_n}$ and $C'\in\G_{t'_1,\dots,t'_m}$, then both events belong to the $\sigma$-field $\G_{u_1,\dots,u_r}$ associated with the ordered union $\{u_1<\dots<u_r\}$ of the two sets of times, and so does $C\cap C'$; hence $\cP$ is closed under finite intersections. Moreover $\sigma(\cP)=\F^{\tilde N}_\infty$: every $\tilde N_t$ with $t>0$ is $\G_t$-measurable, $\tilde N_0\equiv0$, and $\cP\subseteq\F^{\tilde N}_\infty$. Since \eqref{eq:2-5-21:ind} holds for all $A\in\F_0$ and $C\in\cP$, a second application of the Dynkin System Theorem shows that it holds for all $A\in\F_0$ and $C\in\F^{\tilde N}_\infty$. Thus $\F^{\tilde N}_\infty$ and $\F_0$ are independent under $P^x$, for every $x\in\R$.
\end{kex}

\begin{kex}{Exercise}{5.22}{}
Let $W=\{W_t,\F_t;\,0\le t<\infty\}$ be a standard, one-dimensional Brownian motion on a probability space $(\Omega,\F,P)$ (\bk{Definition 1.1}), and set $A\defeq\{W_1\le 0\}\in\F_1$. We define
\[
  X_t\defeq W_{t\wedge 1}+\ind_{A^c}\,\bigl(W_t-W_{t\wedge 1}\bigr),\qquad 0\le t<\infty .
\]
Thus $X$ follows $W$ up to time $1$; after time $1$ it keeps following $W$ if $W_1>0$, and it is frozen at the value $W_1$ if $W_1\le 0$. The process $X$ has continuous paths and $X_0=0$. We show that $X$ is a martingale with respect to $\{\F_t\}$ (hence also with respect to $\{\F^X_t\}$), but that $X$ is not a Markov process in the sense of \bk{Definition 5.10}, with respect to any filtration to which it is adapted.

\emph{Step 1: $X$ is an $\{\F_t\}$-martingale.} For $t\le 1$ we have $X_t=W_t$, and for $t>1$ the random variable $X_t=W_1+\ind_{A^c}(W_t-W_1)$ is $\F_t$-measurable because $A\in\F_1\subseteq\F_t$; so $X$ is adapted. Moreover $\abs{X_t}\le\abs{W_{t\wedge1}}+\abs{W_t-W_{t\wedge 1}}$, so $E\abs{X_t}<\infty$. Let $0\le s<t$. If $t\le 1$, then $E[X_t\mid\F_s]=E[W_t\mid\F_s]=W_s=X_s$ a.s., because $W_t-W_s$ is independent of $\F_s$ and has mean zero. If $1\le s<t$, then $X_t-X_s=\ind_{A^c}(W_t-W_s)$ with $\ind_{A^c}$ bounded and $\F_s$-measurable, whence
\[
  E[X_t-X_s\mid\F_s]=\ind_{A^c}\,E[W_t-W_s\mid\F_s]=\ind_{A^c}\,E[W_t-W_s]=0\quad\text{a.s.}
\]
Finally, if $s<1<t$, the two cases already treated and the tower property give
$E[X_t\mid\F_s]=E\bigl[E[X_t\mid\F_1]\bigm|\F_s\bigr]=E[X_1\mid\F_s]=X_s$ a.s. Hence $\{X_t,\F_t\}$ is a martingale. Since $X$ is $\{\F_t\}$-adapted, $\F^X_s\subseteq\F_s$, and conditioning the identity $E[X_t\mid\F_s]=X_s$ on $\F^X_s$ shows that $\{X_t,\F^X_t\}$ is a martingale as well.

\emph{Step 2: a conditional probability.} We claim that
\begin{equation}\label{eq:2-5-22:a}
  P[X_3>0\mid\F^X_2]=Y\defeq\ind_{\{X_1>0\}}\,\Phi(X_2)\quad\text{a.s.},
  \qquad\text{where }\ \Phi(y)\defeq\int_{-\infty}^{y}\frac{e^{-z^2/2}}{\sqrt{2\pi}}\,dz .
\end{equation}
On $A$ we have $X_3=W_1\le 0$, and on $A^c$ we have $X_3=W_3$; therefore $\ind_{\{X_3>0\}}=\ind_{A^c}\ind_{\{W_3>0\}}$. We apply \bk{Problem 5.9} with the $\sigma$-field $\F_2$, the $\F_2$-measurable random variable $W_2$ and the random variable $W_3-W_2$, which is independent of $\F_2$ and normal with mean zero and variance one: by \bk{(5.5)} and \bk{(5.6)},
\[
  P[W_3>0\mid\F_2]=h(W_2)\ \ \text{a.s.},\qquad h(y)=P[W_3-W_2+y>0]=1-\Phi(-y)=\Phi(y).
\]
As $A^c\in\F_2$, it follows that $P[X_3>0\mid\F_2]=\ind_{A^c}\Phi(W_2)$ a.s. Now $A^c=\{X_1>0\}$, and $W_2=X_2$ on $A^c$, so $\ind_{A^c}\Phi(W_2)=Y$. The random variable $Y$ is $\F^X_2$-measurable, and $\F^X_2\subseteq\F_2$; hence
$P[X_3>0\mid\F^X_2]=E\bigl[P[X_3>0\mid\F_2]\bigm|\F^X_2\bigr]=E[Y\mid\F^X_2]=Y$ a.s., which is \eqref{eq:2-5-22:a}.

\emph{Step 3: $\{X_t,\F^X_t\}$ is not a Markov process.} Suppose that condition (ii) of \bk{Definition 5.10} held for $s=2$, $t=1$, $\Gamma=(0,\infty)$ and the filtration $\{\F^X_t\}$. By \eqref{eq:2-5-22:a}, $Y$ would then be a version of $P[X_3>0\mid X_2]$, and so there would exist a Borel-measurable function $g:\R\to[0,1]$ with
\begin{equation}\label{eq:2-5-22:b}
  \ind_{\{X_1>0\}}\,\Phi(X_2)=g(X_2)\quad\text{a.s.}
\end{equation}
(every $\sigma(X_2)$-measurable random variable is a Borel function of $X_2$). Let $N\defeq\{y\le 0:\ g(y)\ne 0\}\in\B(\R)$.

On the event $A=\{X_1\le 0\}$ we have $X_2=W_1$ and the left-hand side of \eqref{eq:2-5-22:b} vanishes, so $g(W_1)=0$ a.s.\ on $\{W_1\le0\}$; that is, $P[W_1\in N]=0$. Since $W_1$ has a strictly positive density with respect to Lebesgue measure, $N$ has Lebesgue measure zero. Consequently
\begin{equation}\label{eq:2-5-22:c}
  P[W_2\in N]=0,
\end{equation}
because $W_2$, being normal with variance $2$, has a density.

On the event $D\defeq\{W_1>0,\ W_2\le 0\}\subseteq A^c$ we have $X_1>0$ and $X_2=W_2\le0$, so the left-hand side of \eqref{eq:2-5-22:b} equals $\Phi(W_2)>0$; hence $g(W_2)>0$, i.e.\ $W_2\in N$, a.s.\ on $D$. Together with \eqref{eq:2-5-22:c} this yields $P(D)\le P[W_2\in N]=0$. On the other hand, $W_1$ and $W_2-W_1$ are independent standard normal random variables, so
\[
  P(D)\ \ge\ P[\,0<W_1<1,\ W_2-W_1\le -1\,]=P[0<W_1<1]\cdot P[W_2-W_1\le-1]>0 .
\]
This contradiction shows that
\[
  P[X_3\in(0,\infty)\mid\F^X_2]\ \ne\ P[X_3\in(0,\infty)\mid X_2]
\]
on a set of positive probability, so $\{X_t,\F^X_t;\,t\ge 0\}$ is not a Markov process.

\emph{Step 4: no other filtration helps.} Let $\{\G_t\}$ be any filtration to which $X$ is adapted, and suppose that $P[X_3>0\mid\G_2]=P[X_3>0\mid X_2]$ a.s. Since $\sigma(X_2)\subseteq\F^X_2\subseteq\G_2$, conditioning both sides on $\F^X_2$ gives $P[X_3>0\mid\F^X_2]=P[X_3>0\mid X_2]$ a.s., which was excluded in Step~3. In particular, $\{X_t,\F_t;\,t\ge 0\}$ is a continuous martingale which is not a Markov process.

\rmk{The mechanism is that the present position $X_2$ does not reveal whether the path has been frozen at time $1$, whereas the past does: a negative value of $X_2$ may come from a frozen path, which stays put, or from a path with $W_1>0$, which keeps moving. The martingale property only constrains the conditional \emph{mean} of the future given the past, and both behaviours have the same conditional mean $X_2$.}
\end{kex}

\solsection{2.6}{The Strong Markov Property and the Reflection Principle}

\begin{kex}{Exercise}{6.1}{}
We take
\[
  S\defeq \ind_{\{B_1<0\}},
\]
that is, $S=0$ on $\{B_1\ge 0\}$ and $S=1$ on $\{B_1<0\}$. This is an $\F$-measurable random variable with
values in $\{0,1\}$, hence a random time with $P[0\le S<\infty]=1$. For every $t\ge 0$,
\[
  W_t=B_{S+t}-B_S
     = B_t\,\ind_{\{B_1\ge 0\}}+\bigl(B_{1+t}-B_1\bigr)\ind_{\{B_1<0\}} ,
\]
on the event $\{B_0=0\}$ of probability one (all identities between random variables and between events
below are understood up to this null event); in particular each $W_t$ is a random variable, and $W$ is a
continuous process, adapted to its own filtration $\{\F^W_t\}$, with $W_0=0$ identically.

If $W$ were a standard Brownian motion, then by \bk{Definition 1.1} the random variable $W_1=W_1-W_0$ would
be normally distributed with mean zero and variance one, and so $P[W_1<0]=\tfrac12$. But the formula above
with $t=1$ shows that $W_1=B_1\ge 0$ on $\{B_1\ge 0\}$ and $W_1=B_2-B_1$ on $\{B_1<0\}$, whence
\[
  \{W_1<0\}=\{B_1<0\}\cap\{B_2-B_1<0\}.
\]
The event $\{B_1<0\}$ belongs to $\F_1$, the increment $B_2-B_1$ is independent of $\F_1$, and both $B_1$
and $B_2-B_1$ are normal with mean zero and variance one (\bk{Definition 1.1}). Therefore
\[
  P[W_1<0]=P[B_1<0]\cdot P[B_2-B_1<0]=\tfrac12\cdot\tfrac12=\tfrac14\neq\tfrac12 .
\]
Thus $W_1$ is not normally distributed with mean zero and variance one, and $W$ is not a Brownian motion,
whatever filtration it is paired with. (Nor is $W$ a Brownian motion with a different variance parameter or a
different initial distribution: $W_0=0$ a.s., and $P[W_1<0]=\tfrac12$ for every centered nondegenerate normal
law.)

\rmk{The random time $S$ is not an optional time of $\{\F_t\}$, in agreement with the discussion preceding
the exercise. Indeed, if it were, we would have $\{B_1\ge 0\}=\{S<\tfrac12\}\in\F_{1/2}$, and this event would
be independent of the increment $Y\defeq B_1-B_{1/2}$, so that $P[B_1\ge 0,\,Y>0]=\tfrac14$. However, with
$X\defeq B_{1/2}$, the random variables $X$ and $Y$ are independent and normal with mean zero and variance
$\tfrac12$, so the law of $(X,Y)$ is invariant under rotations of the plane and gives each sector of opening
angle $\theta$ with vertex at the origin the mass $\theta/2\pi$. Hence
\[
  P[B_1\ge 0,\,Y>0]=P[Y>0]-P[Y>0,\ X<-Y]=\tfrac12-\tfrac18=\tfrac38\neq\tfrac14 ,
\]
the set $\{(x,y):\,y>0,\ x<-y\}$ being a sector of opening angle $\pi/4$. The example exploits exactly this:
$S$ looks into the future of $B$ in order to decide where to restart, and thereby biases the sign of the
path that follows.}
\end{kex}

\begin{kex}{Problem}{6.9}{}
The book's Section 2.10 derives (ii) from (i) and \bk{Problem 1.2.17}(i), and proves (iii) for condition (e), saying only that (e$'$) is handled in the same way; part (i) is used there but not proved. We prove (i), record the two points on which the printed argument for (ii) silently relies, and write out (iii) in a form that covers (e) and (e$'$) at once.

\textbf{(i)} Let $Y_i$ be a version of $E[Z_i\mid\F_{S+}]$, $i=1,2$, and put $D\defeq\ind_A(Y_1-Y_2)$. Since $A\in\F_{S+}$, the random variable $D$ is $\F_{S+}$-measurable, and it is integrable. For every $G\in\F_{S+}$ we have $A\cap G\in\F_{S+}$, so the defining property of conditional expectation gives
\[
  E[D\,\ind_G]=E[Y_1\ind_{A\cap G}]-E[Y_2\ind_{A\cap G}]
  =E[Z_1\ind_{A\cap G}]-E[Z_2\ind_{A\cap G}]=E[(Z_1-Z_2)\ind_A\ind_G]=0,
\]
because $(Z_1-Z_2)\ind_A=0$ identically. Taking $G=\{D>0\}$ and then $G=\{D<0\}$, both of which belong to $\F_{S+}$, we obtain $E[D^+]=E[D^-]=0$, hence $D=0$ a.s., i.e.\ $Y_1=Y_2$ a.s.\ on $A$.

\rmk{The argument uses nothing about $\F_{S+}$ except that it is a sub-$\sigma$-field of $\F$ containing $A$; and it suffices that $Z_1=Z_2$ hold a.s.\ on $A$.}

\textbf{(ii)} The printed solution applies \bk{Problem 1.2.17}(i), which is a statement about stopping times $T$ and the $\sigma$-fields $\F_T$ of \bk{Definition 1.2.12}, to the optional time $S$ and to $\F_{S+}$. This is legitimate for the following reason. Consider the filtration $\{\G_t\}$ with $\G_t\defeq\F_{t+}$. By \bk{Corollary 1.2.4}, $S$ is a stopping time of $\{\G_t\}$, and so are the constant $s$ and $S\wedge s$. The $\sigma$-field of events determined prior to the $\{\G_t\}$-stopping time $S$ is, by \bk{Definition 1.2.12}, the class of all $B\in\F$ with $B\cap\{S\le t\}\in\G_t=\F_{t+}$ for every $t\ge0$; by \bk{Definition 1.2.20} this is exactly $\F_{S+}$. In the same way $\G_{S\wedge s}=\F_{(S\wedge s)+}$; here $S\wedge s$ is an optional time of $\{\F_t\}$, since $\{S\wedge s<t\}$ equals $\{S<t\}\in\F_t$ for $t\le s$ and $\Omega$ for $t>s$. Thus \bk{Problem 1.2.17}(i), applied in the filtration $\{\G_t\}$ to the stopping times $S$ and $s$ (in the roles of $T$ and $S$ there) and to the integrable random variable $Z_2$, reads
\[
  E[Z_2\mid\F_{S+}]=E[Z_2\mid\F_{(S\wedge s)+}],\quad\text{a.s. on }\{S\le s\},
\]
and combining this with (i) on $A$ gives the assertion on $\{S\le s\}\cap A$, as in the book.

\textbf{(iii)} We are in the setting of \bk{Proposition 6.7}: $X=\{X_t,\F_t;\ t\ge0\}$ is progressively measurable on $(\Omega,\F)$ and $\{P^x\}_{x\in\R^d}$ satisfies (a) and (b) of \bk{Definition 6.3}. Fix $x\in\R^d$, $t\ge0$, and a bounded function $g:\R^d\to\R$ which is either $\ind_\Gamma$ with $\Gamma\in\B(\R^d)$ (case (e)) or continuous (case (e$'$)); then $U_tg$ is the function of \bk{(5.7)}, and with the convention $E^x[\ind_\Gamma(X_{S+t})\mid\F_{S+}]=P^x[X_{S+t}\in\Gamma\mid\F_{S+}]$ both conditions take the form
\begin{equation}\label{eq:2-6-09:e}
  E^x[g(X_{S+t})\mid\F_{S+}]=(U_tg)(X_S),\quad P^x\text{-a.s. on }\{S<\infty\}.
\end{equation}
For an optional time $S$ that may take the value $+\infty$, the random variables $X_S$ and $X_{S+t}$ are defined on $\{S<\infty\}$ only, and the left-hand side of \eqref{eq:2-6-09:e} is understood as $E^x[Z_1\mid\F_{S+}]$ with
\[
  Z_1\defeq g(X_{S+t})\,\ind_{\{S<\infty\}}.
\]
In case (e), $g=\ind_\Gamma$, this is not a choice of ours: $Z_1$ is the indicator of the event $\{X_{S+t}\in\Gamma\}\defeq\{S<\infty,\ X_{S+t}\in\Gamma\}$, which is the convention laid down in \bk{Remark 6.4} for \bk{Definitions 6.2} and \bk{6.3} and hence for the identical expression in \bk{Proposition 6.7}(e). In case (e$'$) the book does not say how $f(X_{S+t})$ is to be read on $\{S=\infty\}$, but the reading is immaterial: $\{S<\infty\}\in\F_{S+}$ because $S$ is $\F_{S+}$-measurable (\bk{Problem 1.2.21}), so by (i) any other integrable random variable which agrees with $g(X_{S+t})$ on $\{S<\infty\}$ has the same conditional expectation, $P^x$-a.s.\ on $\{S<\infty\}$. That $Z_1$ is a random variable follows from \bk{Proposition 1.2.18}, applied in the filtration $\{\F_{u+}\}$, with respect to which $X$ is still progressively measurable and $S+t$ is a stopping time (\bk{Corollary 1.2.4}); it is bounded, hence $P^x$-integrable.

Suppose now that \eqref{eq:2-6-09:e} holds, for our fixed $x$, $t$, $g$, whenever $S$ is a bounded optional time of $\{\F_t\}$, and let $S$ be an arbitrary optional time of $\{\F_t\}$. For each positive integer $n$, $S\wedge n$ is a bounded optional time (as noted in (ii)), so
\begin{equation}\label{eq:2-6-09:bdd}
  E^x[Z_2^{(n)}\mid\F_{(S\wedge n)+}]=(U_tg)(X_{S\wedge n}),\quad P^x\text{-a.s.},
  \qquad\text{where } Z_2^{(n)}\defeq g(X_{(S\wedge n)+t}).
\end{equation}
On $\{S\le n\}$ we have $S\wedge n=S<\infty$, hence $Z_1=Z_2^{(n)}$ and $X_{S\wedge n}=X_S$ there. Moreover $\{S\le n\}\in\F_{S+}$. We may therefore apply (ii) on the probability space $(\Omega,\F,P^x)$ with $Z_2=Z_2^{(n)}$, $s=n$ and $A=\{S\le n\}$, and then \eqref{eq:2-6-09:bdd}:
\[
  E^x[Z_1\mid\F_{S+}]=E^x[Z_2^{(n)}\mid\F_{(S\wedge n)+}]=(U_tg)(X_{S\wedge n})=(U_tg)(X_S),
  \quad P^x\text{-a.s. on }\{S\le n\}.
\]
Thus there is a $P^x$-null event $N_n$ such that $E^x[Z_1\mid\F_{S+}]=(U_tg)(X_S)$ on $\{S\le n\}\setminus N_n$. The event $N\defeq\bigcup_{n\ge1}N_n$ is $P^x$-null, and since $\{S<\infty\}=\bigcup_{n\ge1}\{S\le n\}$, the equality holds on $\{S<\infty\}\setminus N$. This is \eqref{eq:2-6-09:e} for $S$. As $x$, $t$ and $g$ were arbitrary within the class considered (indicators of Borel sets for (e), bounded continuous functions for (e$'$)), condition (e), respectively (e$'$), holds for every optional time as soon as it holds for every bounded one.
\end{kex}

\begin{kex}{Problem}{6.11}{}
The setting is that of \bk{Theorem 6.10}: $X=\{X_t,\F_t;\,t\ge 0\}$ is a progressively measurable, $\R^d$-valued process on $(\Omega,\F)$, $\{P^x\}_{x\in\R^d}$ is a family of probability measures on $(\Omega,\F)$, and $\{\theta_s\}_{s\ge 0}$ are $\F/\F$-measurable shift operators satisfying \bk{(5.14)}. Since $X$ is adapted, $\F^X_\infty=\sigma(X_t;\,t\ge 0)\subseteq\F$, so that $P^y(F)$ and $E^y(Y)=\int_\Omega Y\,dP^y$ are defined for every $y\in\R^d$, every $F\in\F^X_\infty$ and every bounded, $\F^X_\infty$-measurable $Y$.

\emph{Step 0: the meaning of the two conditions.} Let $S$ be an optional time of $\{\F_t\}$. Then $\{S<\infty\}=\bigcup_{n}\{S<n\}\in\F$, and $S$ is $\F$-measurable. The random shift $\theta_S$ is defined on $\{S<\infty\}$ only, and \bk{(5.14)} gives
\begin{equation}\label{eq:2-6-11:a}
  X_t\bigl(\theta_S(\omega)\bigr)=X_{S(\omega)+t}(\omega),\qquad \omega\in\{S<\infty\},\ t\ge 0 .
\end{equation}
We claim that
\begin{equation}\label{eq:2-6-11:b}
  \theta_S^{-1}F\defeq\{\omega\in\{S<\infty\};\ \theta_S(\omega)\in F\}\in\F\qquad\text{for every }F\in\F^X_\infty .
\end{equation}
Indeed, progressive measurability implies that $(s,\omega)\mapsto X_s(\omega)$ is $\B([0,\infty))\otimes\F$-measurable, because for every Borel set $\Gamma$ the set $\{(s,\omega);\,X_s(\omega)\in\Gamma\}$ is the union over $n\ge 1$ of the sets $\{(s,\omega)\in[0,n]\times\Omega;\,X_s(\omega)\in\Gamma\}\in\B([0,n])\otimes\F_n$. For fixed $t\ge 0$ the mapping $\omega\mapsto(S(\omega)+t,\omega)$ from $\{S<\infty\}$ into $[0,\infty)\times\Omega$ is measurable with respect to the trace of $\F$ on $\{S<\infty\}$ and $\B([0,\infty))\otimes\F$; composing, and using \eqref{eq:2-6-11:a}, we get $\theta_S^{-1}\{X_t\in\Gamma\}=\{S<\infty,\ X_{S+t}\in\Gamma\}\in\F$ for $\Gamma\in\B(\R^d)$. The collection of sets $F\in\F^X_\infty$ with $\theta_S^{-1}F\in\F$ is a $\sigma$-field: it contains $\Omega$ since $\theta_S^{-1}\Omega=\{S<\infty\}$, it is closed under complementation since $\theta_S^{-1}(F^c)=\{S<\infty\}\setminus\theta_S^{-1}F$, and it is closed under countable unions. As it contains the generators $\{X_t\in\Gamma\}$ of $\F^X_\infty$, it is all of $\F^X_\infty$, which proves \eqref{eq:2-6-11:b}.

Consequently, for a bounded, $\F^X_\infty$-measurable $Y$, the function $Y\circ\theta_S$ is a bounded function on $\{S<\infty\}$, measurable with respect to the trace of $\F$ on that set. In accordance with \bk{Remark 6.4}, where $\{X_{S+t}\in\Gamma\}$ stands for $\{S<\infty,\,X_{S+t}\in\Gamma\}$, the symbol $E^x[Y\circ\theta_S\mid\F_{S+}]$ denotes the conditional expectation of the bounded, $\F$-measurable random variable
\[
  \wt Y_S\defeq (Y\circ\theta_S)\,\ind_{\{S<\infty\}}
\]
(equal to $0$ on $\{S=\infty\}$); for $Y=\ind_F$ this is $\ind_{\theta_S^{-1}F}$, in agreement with the left-hand side of (e$''$).

Finally, for a function $u\colon\R^d\to\R$ and an integrable random variable $Z$, the statement
\[
  E^x[Z\mid\F_{S+}]=u(X_S),\qquad P^x\text{-a.s. on }\{S<\infty\},
\]
means: for one (equivalently, every) version $\zeta$ of $E^x[Z\mid\F_{S+}]$ there is a set $N\in\F$ with $P^x(N)=0$ such that $\zeta(\omega)=u\bigl(X_{S(\omega)}(\omega)\bigr)$ for all $\omega\in\{S<\infty\}\setminus N$. This is how (e$''$), with $u(y)=P^y(F)$, and (e$'''$), with $u(y)=E^y(Y)$, are to be read; the \emph{Note} in the statement says exactly this, since $E^y(Y)=\int_\Omega Y(\omega')\,P^y(d\omega')$ evaluated at $y=X_{S(\omega)}(\omega)$ is the right-hand side displayed there. Two versions of a conditional expectation differ on a $P^x$-null set of $\F$, so the choice of version is immaterial. This formulation does not require any measurability of $u$; see the remark at the end.

\emph{Step 1: \textup{(e$'''$)} implies \textup{(e$''$)}.} Let $x\in\R^d$, $F\in\F^X_\infty$ and an optional time $S$ of $\{\F_t\}$ be given. The random variable $Y\defeq\ind_F$ is bounded and $\F^X_\infty$-measurable, $\wt Y_S=\ind_{\theta_S^{-1}F}$, and $E^y(Y)=P^y(F)$ for every $y\in\R^d$. Hence (e$'''$) for this $Y$ reads
\[
  P^x[\theta_S^{-1}F\mid\F_{S+}]=E^x\bigl[\ind_{\theta_S^{-1}F}\bigm|\F_{S+}\bigr]=P^{X_S}(F),\qquad P^x\text{-a.s. on }\{S<\infty\},
\]
which is (e$''$).

\emph{Step 2: \textup{(e$''$)} implies \textup{(e$'''$)}.} Assume (e$''$), and fix $x\in\R^d$ and an optional time $S$ of $\{\F_t\}$.

Suppose first that $Y$ is simple: $Y=\sum_{i=1}^m a_i\ind_{F_i}$ with $a_i\in\R$ and $F_i\in\F^X_\infty$. Then
\[
  \wt Y_S=\sum_{i=1}^m a_i\,\ind_{\theta_S^{-1}F_i}
  \qquad\text{and}\qquad
  E^y(Y)=\sum_{i=1}^m a_i\,P^y(F_i)\quad\text{for every }y\in\R^d .
\]
Let $\zeta_i$ be a version of $P^x[\theta_S^{-1}F_i\mid\F_{S+}]$. By (e$''$) there are sets $N_i\in\F$ with $P^x(N_i)=0$ and $\zeta_i(\omega)=P^{X_{S(\omega)}(\omega)}(F_i)$ for $\omega\in\{S<\infty\}\setminus N_i$. By the linearity of conditional expectation, $\zeta\defeq\sum_{i=1}^m a_i\zeta_i$ is a version of $E^x[\wt Y_S\mid\F_{S+}]$, and for $\omega\in\{S<\infty\}\setminus\bigcup_{i=1}^m N_i$ we have
\[
  \zeta(\omega)=\sum_{i=1}^m a_i\,P^{X_{S(\omega)}(\omega)}(F_i)=E^{X_{S(\omega)}(\omega)}(Y).
\]
Thus (e$'''$) holds for simple $Y$.

Now let $Y$ be bounded and $\F^X_\infty$-measurable, say $\abs{Y}\le K$, and put $Y_n\defeq 2^{-n}\lfloor 2^nY\rfloor$, $n\ge 1$. Each $Y_n$ is $\F^X_\infty$-measurable and takes only finitely many values (multiples of $2^{-n}$ in $[-K-1,K]$), so it is simple, and
\begin{equation}\label{eq:2-6-11:c}
  0\le Y(\omega')-Y_n(\omega')\le 2^{-n}\qquad\text{for all }\omega'\in\Omega,\ n\ge 1 .
\end{equation}
We draw two consequences from \eqref{eq:2-6-11:c}. First, integrating with respect to $P^y$,
\begin{equation}\label{eq:2-6-11:d}
  \abs[\big]{E^y(Y)-E^y(Y_n)}\le 2^{-n}\qquad\text{for all }y\in\R^d,\ n\ge 1 .
\end{equation}
Second, applying \eqref{eq:2-6-11:c} at $\omega'=\theta_S(\omega)$, $\omega\in\{S<\infty\}$, we find $0\le\wt Y_S-(\wt{Y_n})_S\le 2^{-n}$ everywhere on $\Omega$, where $(\wt{Y_n})_S\defeq(Y_n\circ\theta_S)\ind_{\{S<\infty\}}$. Let $\zeta$ be a version of $E^x[\wt Y_S\mid\F_{S+}]$ and $\zeta_n$ a version of $E^x[(\wt{Y_n})_S\mid\F_{S+}]$. By the monotonicity and linearity of conditional expectation, there are sets $M_n\in\F$ with $P^x(M_n)=0$ and
\begin{equation}\label{eq:2-6-11:e}
  \abs{\zeta(\omega)-\zeta_n(\omega)}\le 2^{-n}\qquad\text{for }\omega\notin M_n .
\end{equation}
By the case of simple random variables already treated, there are sets $N_n\in\F$ with $P^x(N_n)=0$ and
\begin{equation}\label{eq:2-6-11:f}
  \zeta_n(\omega)=E^{X_{S(\omega)}(\omega)}(Y_n)\qquad\text{for }\omega\in\{S<\infty\}\setminus N_n .
\end{equation}
Set $N\defeq\bigcup_{n\ge 1}(M_n\cup N_n)\in\F$; then $P^x(N)=0$. For $\omega\in\{S<\infty\}\setminus N$ and every $n\ge 1$, we obtain from \eqref{eq:2-6-11:e}, \eqref{eq:2-6-11:f} and \eqref{eq:2-6-11:d} (with $y=X_{S(\omega)}(\omega)$)
\[
  \abs[\big]{\zeta(\omega)-E^{X_{S(\omega)}(\omega)}(Y)}
  \le\abs{\zeta(\omega)-\zeta_n(\omega)}
   +\abs[\big]{E^{X_{S(\omega)}(\omega)}(Y_n)-E^{X_{S(\omega)}(\omega)}(Y)}
  \le 2\cdot 2^{-n}.
\]
Letting $n\to\infty$ yields $\zeta(\omega)=E^{X_{S(\omega)}(\omega)}(Y)$ for all $\omega\in\{S<\infty\}\setminus N$, i.e.,
\[
  E^x[Y\circ\theta_S\mid\F_{S+}]=E^{X_S}(Y),\qquad P^x\text{-a.s. on }\{S<\infty\}.
\]
Since $x$, $S$ and $Y$ were arbitrary, (e$'''$) holds.

\rmk{The equivalence holds separately for each pair $(x,S)$, and neither (a) nor (b) of \bk{Definition 6.3} was used. If (a) holds, then the functions $u(y)=P^y(F)$ and, by \eqref{eq:2-6-11:d} and the fact that a uniform limit of linear combinations of universally measurable functions is universally measurable, $u(y)=E^y(Y)$ are universally measurable. By \bk{Problem 5.7}, applied to the finite measure $\mu(\Gamma)=P^x[S<\infty,\,X_S\in\Gamma]$ on $\B(\R^d)$, there is then a Borel-measurable $g$ with $u(X_S)=g(X_S)$ on $\{S<\infty\}$ outside a $P^x$-null set of $\F$, exactly as in \bk{(5.8)}, \bk{(5.9)}; and $g(X_S)\ind_{\{S<\infty\}}$ is $\F_{S+}$-measurable by \bk{Remark 6.5}. So under (a) the right-hand sides of (e$''$) and (e$'''$) agree $P^x$-a.s.\ on $\{S<\infty\}$ with genuine $\F_{S+}$-measurable random variables, and our reading of the two conditions is the one used in the proof of \bk{Proposition 5.13}.}
\end{kex}

\begin{kex}{Problem}{6.21}{}
Both assertions follow from one lemma, which we prove by the discretization of \bk{Problem 1.2.24}. As in \bk{Definition 1.3.3} and in \bk{Section 1.1}, an ``integer-valued RCLL process'' is a process \emph{every} path of which is integer-valued, right-continuous and has left limits. For $t\ge 0$ and a bounded Borel function $g:\R\to\R$ we set
\begin{equation}\label{eq:2-6-21:Q}
  (Q_tg)(y)\defeq\sum_{k=0}^{\infty}p_k(t)\,g(y+k),\qquad
  p_k(t)\defeq e^{-\lambda t}\frac{(\lambda t)^k}{k!}\quad(t>0),\qquad p_k(0)\defeq\ind_{\{k=0\}} .
\end{equation}
Since $\sum_kp_k(t)=1$, the series converges absolutely, $\abs{Q_tg}\le\sup\abs{g}$, and $Q_tg$ is Borel measurable as a pointwise limit of Borel functions.

\begin{lem}
Let $M=\{M_t,\F_t;\,0\le t<\infty\}$ be a Poisson process with intensity $\lambda>0$ on a probability space $(\Omega,\F,P)$, let $\xi$ be a real-valued, $\F_0$-measurable random variable, and set $X_t\defeq\xi+M_t$. Then $X=\{X_t,\F_t;\,t\ge0\}$ is progressively measurable, and for every optional time $S$ of $\{\F_t\}$, every $t\ge0$ and every bounded Borel function $g:\R\to\R$,
\begin{equation}\label{eq:2-6-21:lem}
  E\bigl[g(X_{S+t})\ind_{\{S<\infty\}}\bigm|\F_{S+}\bigr]=(Q_tg)(X_S)\,\ind_{\{S<\infty\}},\qquad P\text{-a.s.}
\end{equation}
\end{lem}

\begin{proof}
\emph{Step 1: measurability.} The process $X$ is adapted to $\{\F_t\}$ and all its paths are right-continuous, so it is progressively measurable with respect to $\{\F_t\}$ by \bk{Proposition 1.1.13}, hence also with respect to the larger filtration $\{\F_{t+}\}$. The optional time $S$ is a stopping time of $\{\F_{t+}\}$ (\bk{Corollary 1.2.4}), and the $\sigma$-field attached to it by this filtration is $\F_{S+}$ (\bk{Definition 1.2.20}). \bk{Proposition 1.2.18}, applied with the filtration $\{\F_{t+}\}$, therefore shows that $\{S<\infty\}\in\F_{S+}$ and that $X_S$, defined on $\{S<\infty\}$, is $\F_{S+}$-measurable. Consequently the right-hand side of \eqref{eq:2-6-21:lem} is a bounded, $\F_{S+}$-measurable random variable, and \eqref{eq:2-6-21:lem} is equivalent to
\begin{equation}\label{eq:2-6-21:int}
  E\bigl[g(X_{S+t})\ind_{A\cap\{S<\infty\}}\bigr]=E\bigl[(Q_tg)(X_S)\ind_{A\cap\{S<\infty\}}\bigr]\qquad\text{for every }A\in\F_{S+}.
\end{equation}
For $t=0$ we have $Q_0g=g$ and \eqref{eq:2-6-21:int} is trivial; so let $t>0$.

\emph{Step 2: deterministic times.} Let $s\ge0$ and $B\in\F_s$. We claim that
\begin{equation}\label{eq:2-6-21:det}
  E\bigl[g(X_{s+t})\ind_B\bigr]=E\bigl[(Q_tg)(X_s)\ind_B\bigr].
\end{equation}
By \bk{Definition 1.3.3} the increment $D\defeq M_{s+t}-M_s$ is independent of $\F_s$ and $P[D=k]=p_k(t)$ for $k=0,1,2,\dots$; in particular $P[D\in\{0,1,2,\dots\}]=1$. Since $X_{s+t}=X_s+D$, we have, $P$-a.s.,
\[
  g(X_{s+t})\ind_B=\sum_{k=0}^{\infty}\ind_{\{D=k\}}\,g(X_s+k)\ind_B .
\]
The partial sums are bounded in absolute value by $\sup\abs g$, and for each $k$ the random variable $g(X_s+k)\ind_B$ is bounded and $\F_s$-measurable, hence independent of $\ind_{\{D=k\}}$. By bounded convergence, used twice,
\[
  E\bigl[g(X_{s+t})\ind_B\bigr]=\sum_{k=0}^{\infty}p_k(t)\,E\bigl[g(X_s+k)\ind_B\bigr]=E\bigl[(Q_tg)(X_s)\ind_B\bigr].
\]

\emph{Step 3: discretization.} For $n\ge1$ let $S_n$ be the random time of \bk{Problem 1.2.24}: $S_n=\infty$ on $\{S=\infty\}$ and $S_n=k2^{-n}$ on $\{(k-1)2^{-n}\le S<k2^{-n}\}$, $k\ge1$. By that problem, $S_n\downarrow S$, and $A\cap\{S_n=k2^{-n}\}\in\F_{k2^{-n}}$ for every $A\in\F_{S+}$ and $k\ge1$. Fix $A\in\F_{S+}$. The event $\{S<\infty\}$ is the disjoint union of the events $\{S_n=k2^{-n}\}$, $k\ge1$; hence, by bounded convergence and by \eqref{eq:2-6-21:det} with $s=k2^{-n}$ and $B=A\cap\{S_n=k2^{-n}\}$,
\begin{align}
  E\bigl[g(X_{S_n+t})\ind_{A\cap\{S<\infty\}}\bigr]
  &=\sum_{k=1}^{\infty}E\bigl[g(X_{k2^{-n}+t})\ind_{A\cap\{S_n=k2^{-n}\}}\bigr]
   =\sum_{k=1}^{\infty}E\bigl[(Q_tg)(X_{k2^{-n}})\ind_{A\cap\{S_n=k2^{-n}\}}\bigr]\notag\\
  &=E\bigl[(Q_tg)(X_{S_n})\ind_{A\cap\{S<\infty\}}\bigr].\label{eq:2-6-21:n}
\end{align}
(The random variables $X_{S_n+t}$ and $X_{S_n}$ on $\{S<\infty\}$ are measurable, because on $\{S_n=k2^{-n}\}$ they coincide with $X_{k2^{-n}+t}$ and $X_{k2^{-n}}$.)

\emph{Step 4: passage to the limit.} Fix $\omega\in\{S<\infty\}$. The numbers $S_n(\omega)$ decrease to $S(\omega)$, so $S_n(\omega)+t\downarrow S(\omega)+t$, and by right-continuity of the path $M_\cdot(\omega)$ we have $M_{S_n(\omega)+t}(\omega)\to M_{S(\omega)+t}(\omega)$ and $M_{S_n(\omega)}(\omega)\to M_{S(\omega)}(\omega)$. A convergent sequence of integers is eventually constant; thus there is $n_0(\omega)$ such that for all $n\ge n_0(\omega)$
\[
  X_{S_n+t}(\omega)=\xi(\omega)+M_{S_n(\omega)+t}(\omega)=X_{S+t}(\omega)\qquad\text{and}\qquad X_{S_n}(\omega)=X_S(\omega).
\]
Therefore $g(X_{S_n+t})\to g(X_{S+t})$ and $(Q_tg)(X_{S_n})\to(Q_tg)(X_S)$ pointwise on $\{S<\infty\}$, although neither $g$ nor $Q_tg$ need be continuous. Both sequences are bounded by $\sup\abs g$, so the bounded convergence theorem allows us to let $n\to\infty$ in \eqref{eq:2-6-21:n}; this gives \eqref{eq:2-6-21:int}.
\end{proof}

\textbf{The Poisson family is a strong Markov family.} Let $N=\{N_t,\F_t;\,t\ge0\}$, $(\Omega,\F)$, $\{P^x\}_{x\in\R}$ be a Poisson family with intensity $\lambda$, and $\tilde N_t=N_t-N_0$. We verify the requirements of \bk{Definition 6.3} with $d=1$. Fix $x\in\R$. On $(\Omega,\F,P^x)$ the process $\{\tilde N_t,\F_t\}$ is a Poisson process with intensity $\lambda$ by \bk{Definition 5.20}(iii), and $\xi\defeq N_0$ is $\F_0$-measurable because $N$ is adapted. The Lemma applies with $M=\tilde N$ and $X=\xi+\tilde N=N$: the process $N$ is progressively measurable with respect to $\{\F_t\}$ (a property which does not involve the measure), and for every optional time $S$ of $\{\F_t\}$, $t\ge0$ and $\Gamma\in\B(\R)$,
\begin{equation}\label{eq:2-6-21:fam}
  P^x\bigl[N_{S+t}\in\Gamma,\ S<\infty\bigm|\F_{S+}\bigr]=(Q_t\ind_\Gamma)(N_S)\,\ind_{\{S<\infty\}},\qquad P^x\text{-a.s.}
\end{equation}
Conditions (a) and (b) of \bk{Definition 6.3} are (i) and (ii) of \bk{Definition 5.20}. It remains to identify $Q_t$ with the operator $U_t$ of \bk{(5.7)}. Let $y\in\R$. Under $P^y$ we have $N_0=y$ a.s., and $\tilde N_t=\tilde N_t-\tilde N_0$ has the Poisson distribution with mean $\lambda t$ if $t>0$ (\bk{Definition 1.3.3}), while $\tilde N_0=0$; hence $P^y[\tilde N_t=k]=p_k(t)$ for all $t\ge0$, $k\ge0$, and
\[
  (U_t\ind_\Gamma)(y)=P^y[N_t\in\Gamma]=P^y[\,y+\tilde N_t\in\Gamma\,]=\sum_{k=0}^{\infty}p_k(t)\ind_\Gamma(y+k)=(Q_t\ind_\Gamma)(y).
\]
Thus \eqref{eq:2-6-21:fam} is precisely condition (e) of \bk{Proposition 6.7}, in the sense of \bk{Remark 6.4}, for every optional time $S$, every $t\ge0$ and every $x\in\R$. By \bk{Proposition 6.7}, $N$, $(\Omega,\F)$, $\{P^x\}_{x\in\R}$ is a strong Markov family.

\textbf{The Poisson process is a strong Markov process.} Let $N=\{N_t,\F_t;\,t\ge0\}$ be a Poisson process with intensity $\lambda$ on $(\Omega,\F,P)$. We verify \bk{Definition 6.2} with $d=1$ and $\mu=\delta_0$, the unit mass at the origin. By the Lemma with $M=N$ and $\xi\equiv0$, $N$ is progressively measurable, and (i) holds because $N_0=0$ a.s. For (ii), let $S$ be an optional time of $\{\F_t\}$, $t\ge0$, $\Gamma\in\B(\R)$, and write $h\defeq Q_t\ind_\Gamma$, a Borel function. By \eqref{eq:2-6-21:lem},
\begin{equation}\label{eq:2-6-21:proc}
  P\bigl[N_{S+t}\in\Gamma,\ S<\infty\bigm|\F_{S+}\bigr]=Y\defeq h(N_S)\ind_{\{S<\infty\}},\qquad P\text{-a.s.}
\end{equation}
Let $\G$ be the $\sigma$-field generated by $N_S$ in the sense of \bk{Problem 1.1.17}, i.e.\ the collection of the sets $\{N_S\in A\}$ and $\{N_S\in A\}\cup\{S=\infty\}$, $A\in\B(\R)$, where $\{N_S\in A\}=\{S<\infty,\ N_S\in A\}$. By Step~1 of the proof of the Lemma, $\{S=\infty\}\in\F_{S+}$ and $\{N_S\in A\}\in\F_{S+}$, so $\G\subseteq\F_{S+}$. Moreover $Y$ is $\G$-measurable: for $C\in\B(\R)$ the event $\{Y\in C\}$ equals $\{N_S\in h^{-1}(C)\}$ if $0\notin C$ and $\{N_S\in h^{-1}(C)\}\cup\{S=\infty\}$ if $0\in C$. Conditioning \eqref{eq:2-6-21:proc} on $\G$ and using the tower property, we obtain
\[
  P\bigl[N_{S+t}\in\Gamma,\ S<\infty\bigm|N_S\bigr]=E[\,Y\mid\G\,]=Y=P\bigl[N_{S+t}\in\Gamma,\ S<\infty\bigm|\F_{S+}\bigr],\qquad P\text{-a.s.},
\]
which, with the convention of \bk{Remark 6.4}, is condition (ii) of \bk{Definition 6.2}. Hence $N$ is a strong Markov process with initial distribution $\delta_0$.

\rmk{The second assertion can also be deduced from the first by \bk{Remark 6.8}, but only for a Poisson process which is part of a Poisson family; the argument above covers every Poisson process in the sense of \bk{Definition 1.3.3}.}
\end{kex}

\solsection{2.7}{Brownian Filtrations}

\begin{kex}{Problem}{7.6}{}
The book's Section 2.10 gives only a one-line indication (repeat the argument of the solution of \bk{Problem 7.5}, using \bk{Problem 7.1}); we write out the whole argument.

Recall from Section 2.7 that, for a filtration $\{\F_t\}$, $\F_{t-}=\sigma\bigl(\bigcup_{s<t}\F_s\bigr)$ for $t>0$ and $\F_{0-}=\F_0$, and that left-continuity means $\F_{t-}=\F_t$ for every $0\le t<\infty$. Recall also from \bk{Definition 7.2} that $\F^\mu_t=\sigma(\F^X_t\cup\cN^\mu)$, where $\cN^\mu$ is the collection of $P^\mu$-null sets, which does not depend on $t$. Since $\F^X_s\subseteq\F^X_t$ for $s\le t$, we have $\F^\mu_s\subseteq\F^\mu_t$ for $s\le t$, so $\{\F^\mu_t\}$ is indeed a filtration, and
\[
  \F^\mu_{t-}\defeq\sigma\Bigl(\bigcup_{0\le s<t}\F^\mu_s\Bigr),\quad t>0;\qquad \F^\mu_{0-}\defeq\F^\mu_0 .
\]
For $t=0$ there is nothing to prove. We fix $t>0$ and prove the two inclusions.

\emph{The inclusion $\F^\mu_{t-}\subseteq\F^\mu_t$.} For every $0\le s<t$ we have $\F^\mu_s\subseteq\F^\mu_t$, hence $\bigcup_{s<t}\F^\mu_s\subseteq\F^\mu_t$; as $\F^\mu_t$ is a $\sigma$-field, it contains the $\sigma$-field generated by this union, which is $\F^\mu_{t-}$. (This inclusion holds for every filtration and does not use the left-continuity of the paths.)

\emph{The inclusion $\F^\mu_t\subseteq\F^\mu_{t-}$.} It suffices to show that the $\sigma$-field $\F^\mu_{t-}$ contains both $\cN^\mu$ and $\F^X_t$, for then it contains $\sigma(\F^X_t\cup\cN^\mu)=\F^\mu_t$.

First, $\cN^\mu\subseteq\F^\mu_0$ by the definition of $\F^\mu_0$, and $\F^\mu_0\subseteq\F^\mu_{t-}$ because $0<t$. Thus $\cN^\mu\subseteq\F^\mu_{t-}$.

Secondly, for every $0\le s<t$ we have $\F^X_s\subseteq\F^\mu_s\subseteq\F^\mu_{t-}$, so that $\bigcup_{s<t}\F^X_s\subseteq\F^\mu_{t-}$ and therefore
\[
  \F^X_{t-}=\sigma\Bigl(\bigcup_{0\le s<t}\F^X_s\Bigr)\subseteq\F^\mu_{t-}.
\]
Because $X$ has left-continuous paths, the filtration $\{\F^X_t\}$ is left-continuous by \bk{Problem 7.1}\,(ii), that is, $\F^X_t=\F^X_{t-}$. Consequently $\F^X_t\subseteq\F^\mu_{t-}$.

The two inclusions give $\F^\mu_{t-}=\F^\mu_t$ for every $t>0$, and the filtration $\{\F^\mu_t\}$ is left-continuous.

\rmk{The book's indication asks to replace, in the solution of \bk{Problem 7.5}, the $\sigma$-field $\F^\mu_\infty$ by ``$\F^X_{t-}$''; the $\sigma$-field that plays the role of $\F^\mu_\infty$ is $\F^\mu_{t-}$. With this reading the argument of that solution runs as follows: by \bk{Problem 7.3} applied to $\G=\F^X_t$, every $F\in\F^\mu_t$ can be written as $F=G\,\triangle\,N$ with $G\in\F^X_t=\F^X_{t-}\subseteq\F^\mu_{t-}$ and $N\in\cN^\mu\subseteq\F^\mu_{t-}$, whence $F\in\F^\mu_{t-}$. The argument above shows that the representation of \bk{Problem 7.3} is not needed, since $\F^\mu_t$ is by definition the smallest $\sigma$-field containing $\F^X_t$ and $\cN^\mu$.}
\end{kex}

\begin{kex}{Exercise}{7.11}{}
Throughout, as stipulated before the exercise, $X=\{X_t,\F^X_t;\ 0\le t<\infty\}$ is a strong Markov
process with initial distribution $\mu$ on $(\Omega,\F^X_\infty,P^\mu)$; the filtration $\{\F^\mu_t\}$
is regarded as a filtration on the measurable space $(\Omega,\F^\mu)$, so that the $\sigma$-fields
$\F^\mu_S$ and $\F^\mu_{S+}$ of \bk{Definitions 1.2.12} and \bk{1.2.20} consist of sets
$A\in\F^\mu$. We shall use two elementary properties of the collection
$\cN^\mu=\{F\subseteq\Omega;\ \exists\, G\in\F^X_\infty \text{ with } F\subseteq G,\ P^\mu(G)=0\}$:
a subset of a set in $\cN^\mu$ belongs to $\cN^\mu$, and a countable union of sets in $\cN^\mu$
belongs to $\cN^\mu$ (take the union of the corresponding sets $G$). Moreover,
$\cN^\mu\subseteq\F^\mu_t$ for every $0\le t<\infty$, by \bk{Definition 7.2}.

\emph{Step 1: $S$ is a stopping time of $\{\F^\mu_t\}$.} Since $X$ is strongly Markov, the augmented
filtration $\{\F^\mu_t\}$ is right-continuous by \bk{Proposition 7.7}. For a right-continuous
filtration the notions of optional time and stopping time coincide (\bk{Proposition 1.2.3}); hence
the optional time $S$ of $\{\F^\mu_t\}$ is a stopping time of $\{\F^\mu_t\}$.

\emph{Step 2: construction of $T$.} Let $\Q_+$ denote the set of strictly positive rational numbers.
For every $q\in\Q_+$ we have $\{S<q\}\in\F^\mu_q=\sigma(\F^X_q\cup\cN^\mu)$, and \bk{Problem 7.3},
applied to the sub-$\sigma$-field $\G=\F^X_q$ of $\F^X_\infty$, provides a set $G_q\in\F^X_q$ with
\[
  N_q\defeq \{S<q\}\,\triangle\, G_q\in\cN^\mu .
\]
We define
\[
  T(\omega)\defeq \inf\{q\in\Q_+;\ \omega\in G_q\},\qquad \omega\in\Omega,
\]
with the convention $\inf\emptyset=\infty$. For every $0\le t<\infty$, the infimum of a set of
real numbers is strictly less than $t$ if and only if the set contains an element strictly less
than $t$, so
\[
  \{T<t\}=\bigcup_{q\in\Q_+,\ q<t} G_q\ \in\F^X_t ,
\]
because $G_q\in\F^X_q\subseteq\F^X_t$ for $q<t$ and the union is countable (it is empty for $t=0$).
In particular $T$ is an $\F^X_\infty$-measurable random variable with values in $[0,\infty]$, i.e., a
random time, and by \bk{Definition 1.2.1} it is an optional time of $\{\F^X_t\}$.

Next, for every $\omega\in\Omega$ we have
\begin{equation}\label{eq:2-7-11:a}
  S(\omega)=\inf\{q\in\Q_+;\ S(\omega)<q\}.
\end{equation}
Indeed, if $S(\omega)=\infty$ the set on the right is empty and its infimum is $\infty$; if
$S(\omega)<\infty$, this set consists of all rationals in $(S(\omega),\infty)$, whose infimum is
$S(\omega)$ by the density of the rationals. Now set
\[
  N\defeq\bigcup_{q\in\Q_+}N_q\in\cN^\mu .
\]
For $\omega\notin N$ and every $q\in\Q_+$ we have $\omega\notin\{S<q\}\triangle G_q$, that is,
$S(\omega)<q$ if and only if $\omega\in G_q$. Thus the sets $\{q\in\Q_+;\ S(\omega)<q\}$ and
$\{q\in\Q_+;\ \omega\in G_q\}$ coincide, and \eqref{eq:2-7-11:a} together with the definition of
$T$ gives $S(\omega)=T(\omega)$. Consequently $\{S\ne T\}\subseteq N$, and therefore
$\{S\ne T\}\in\cN^\mu$.

\emph{Step 3: $\F^\mu_{S+}=\F^\mu_S$.} By Step 1, $S$ is a stopping time of $\{\F^\mu_t\}$, so both
$\sigma$-fields are defined. By \bk{Definition 1.2.20}, $\F^\mu_{S+}$ consists of the sets
$A\in\F^\mu$ with $A\cap\{S\le t\}\in\F^\mu_{t+}$ for every $t\ge 0$, and by
\bk{Definition 1.2.12}, $\F^\mu_S$ consists of the sets $A\in\F^\mu$ with
$A\cap\{S\le t\}\in\F^\mu_t$ for every $t\ge 0$. Since $\F^\mu_{t+}=\F^\mu_t$ for every $t\ge0$
(\bk{Proposition 7.7}), the two descriptions are identical, and $\F^\mu_{S+}=\F^\mu_S$.

\emph{Step 4: $\F^\mu_S=\F^\mu_T$.} Let $A\in\F^\mu$ and $0\le t<\infty$. The sets
$A\cap\{S\le t\}$ and $A\cap\{T\le t\}$ coincide on $\{S=T\}$, so the two sets
\[
  D_1\defeq \bigl(A\cap\{S\le t\}\bigr)\setminus\bigl(A\cap\{T\le t\}\bigr),
  \qquad
  D_2\defeq \bigl(A\cap\{T\le t\}\bigr)\setminus\bigl(A\cap\{S\le t\}\bigr)
\]
are subsets of $\{S\ne T\}\in\cN^\mu$; hence $D_1,D_2\in\cN^\mu\subseteq\F^\mu_t$. From the
identities
\[
  A\cap\{T\le t\}=\Bigl(\bigl(A\cap\{S\le t\}\bigr)\setminus D_1\Bigr)\cup D_2,
  \qquad
  A\cap\{S\le t\}=\Bigl(\bigl(A\cap\{T\le t\}\bigr)\setminus D_2\Bigr)\cup D_1
\]
we see that $A\cap\{S\le t\}\in\F^\mu_t$ if and only if $A\cap\{T\le t\}\in\F^\mu_t$. As this
holds for every $t\ge 0$, a set $A\in\F^\mu$ belongs to $\F^\mu_S$ if and only if it belongs to
the collection $\F^\mu_T$ defined in the statement. Together with Step 3,
\[
  \F^\mu_{S+}=\F^\mu_S=\F^\mu_T .
\]

\rmk{The optional time $T$ of $\{\F^X_t\}$ need not be a stopping time of $\{\F^X_t\}$, which is
why $\F^\mu_T$ has to be defined explicitly in the statement. It is, however, a stopping time of
$\{\F^\mu_t\}$: taking $A=\Omega$ in Step 4 gives $\{T\le t\}\in\F^\mu_t$ for every $t\ge0$,
because $\{S\le t\}\in\F^\mu_t$ by Step 1. Hence the collection $\F^\mu_T$ of the statement is
exactly the $\sigma$-field of \bk{Definition 1.2.12} for the stopping time $T$ of the filtration
$\{\F^\mu_t\}$ on $(\Omega,\F^\mu)$.}
\end{kex}

\begin{kex}{Exercise}{7.12}{}
In the definition of $T_n$ the index $k$ runs through the positive integers, as in \bk{Problem 1.2.24}; the events $\{T=\infty\}$ and $\{(k-1)2^{-n}\le T<k2^{-n}\}$, $k\ge 1$, form a partition of $\Omega$, so $T_n$ is well defined with values in $\{k2^{-n};\,k\ge 1\}\cup\{\infty\}$. The underlying measurable space is $(\Omega,\F^X_\infty)$, so that
\begin{align*}
  \F^X_{T_n}&=\bigl\{G\in\F^X_\infty;\ G\cap\{T_n\le t\}\in\F^X_t,\ \forall\, t\ge 0\bigr\},\\
  \F^X_{T+}&=\bigl\{G\in\F^X_\infty;\ G\cap\{T\le t\}\in\F^X_{t+},\ \forall\, t\ge 0\bigr\}
\end{align*}
(\bk{Definitions 1.2.12} and \bk{1.2.20}), and, as in \bk{Exercise 7.11}, $\F^\mu_T$ is the collection of sets $A\in\F^\mu=\sigma(\F^X_\infty\cup\cN^\mu)$ with $A\cap\{T\le t\}\in\F^\mu_t$ for every $0\le t<\infty$. We recall that $\cN^\mu$ consists of the subsets of the $P^\mu$-null sets of $\F^X_\infty$; consequently every subset of a member of $\cN^\mu$ is in $\cN^\mu$, and so is every countable union of members of $\cN^\mu$. These two facts are used repeatedly below.

\emph{Step 1. $T_n$ is a stopping time of $\{\F^X_t\}$.} Fix $t\ge 0$ and let $j\ge 0$ be the integer with $j2^{-n}\le t<(j+1)2^{-n}$. Since $T_n$ takes the value $k2^{-n}$ exactly on $\{(k-1)2^{-n}\le T<k2^{-n}\}$, we have
\[
  \{T_n\le t\}=\bigcup_{k=1}^{j}\Bigl\{\frac{k-1}{2^n}\le T<\frac{k}{2^n}\Bigr\}=\Bigl\{T<\frac{j}{2^n}\Bigr\},
\]
which is empty if $j=0$ and belongs to $\F^X_{j2^{-n}}\subseteq\F^X_t$ if $j\ge 1$, because $T$ is optional. Hence $T_n$ is a stopping time of $\{\F^X_t\}$ (this is part of \bk{Problem 1.2.24}). In particular, for every $k\ge 1$,
\begin{equation}\label{eq:2-7-12:a}
  \Bigl\{T_n=\frac{k}{2^n}\Bigr\}=\Bigl\{T_n\le\frac{k}{2^n}\Bigr\}\setminus\Bigl\{T_n\le\frac{k-1}{2^n}\Bigr\}\in\F^X_{k2^{-n}},
  \qquad \{T_n=\infty\}=\{T=\infty\}\in\F^X_\infty .
\end{equation}

\emph{Step 2. $\F^\mu_T\subseteq\sigma(\F^X_{T_n}\cup\cN^\mu)$.} Let $A\in\F^\mu_T$, and put
\[
  A_k\defeq A\cap\Bigl\{T_n=\frac{k}{2^n}\Bigr\}\quad(k\ge 1),\qquad A_\infty\defeq A\cap\{T=\infty\},
\]
so that $A$ is the disjoint union of $A_\infty$ and the sets $A_k$, $k\ge 1$.

For every $t>0$ we have
\[
  A\cap\{T<t\}=\bigcup_{m>1/t}\Bigl(A\cap\Bigl\{T\le t-\frac1m\Bigr\}\Bigr)\in\F^\mu_t ,
\]
because the $m$-th set of the union is in $\F^\mu_{t-(1/m)}\subseteq\F^\mu_t$; for $t=0$ the set $A\cap\{T<0\}$ is empty. Therefore
\[
  A_k=\Bigl(A\cap\Bigl\{T<\frac{k}{2^n}\Bigr\}\Bigr)\setminus\Bigl(A\cap\Bigl\{T<\frac{k-1}{2^n}\Bigr\}\Bigr)\in\F^\mu_{k2^{-n}},\qquad k\ge 1 .
\]
Moreover $A_\infty=A\setminus\bigcup_{m\ge 1}\bigl(A\cap\{T\le m\}\bigr)\in\F^\mu$, since $A\in\F^\mu$ and $A\cap\{T\le m\}\in\F^\mu_m\subseteq\F^\mu$.

We now apply \bk{Problem 7.3} to the sub-$\sigma$-fields $\F^X_{k2^{-n}}$ and $\F^X_\infty$ of $\F^X_\infty$, whose augmentations $\sigma(\,\cdot\,\cup\cN^\mu)$ are $\F^\mu_{k2^{-n}}$ and $\F^\mu$: there exist sets
\[
  H_k\in\F^X_{k2^{-n}}\ \ (k\ge 1),\quad H_\infty\in\F^X_\infty,\qquad\text{with}\quad A_k\,\triangle\,H_k\in\cN^\mu,\quad A_\infty\,\triangle\,H_\infty\in\cN^\mu .
\]
We set $G_k\defeq H_k\cap\{T_n=k2^{-n}\}$ and $G_\infty\defeq H_\infty\cap\{T_n=\infty\}$. By \eqref{eq:2-7-12:a}, $G_k\in\F^X_{k2^{-n}}$ and $G_\infty\in\F^X_\infty$. Since $A_k\subseteq\{T_n=k2^{-n}\}$, we have $A_k\,\triangle\,G_k=(A_k\,\triangle\,H_k)\cap\{T_n=k2^{-n}\}\subseteq A_k\,\triangle\,H_k$, and likewise $A_\infty\,\triangle\,G_\infty\subseteq A_\infty\,\triangle\,H_\infty$; hence these symmetric differences are in $\cN^\mu$. Let
\[
  G\defeq G_\infty\cup\bigcup_{k=1}^\infty G_k .
\]
Then $G\in\F^X_\infty$, and for $t\ge 0$, because $G_k\subseteq\{T_n=k2^{-n}\}$ and $G_\infty\subseteq\{T_n=\infty\}$,
\[
  G\cap\{T_n\le t\}=\bigcup_{k\ge 1,\ k2^{-n}\le t}G_k\in\F^X_t ,
\]
a finite union of sets of $\F^X_{k2^{-n}}\subseteq\F^X_t$. Thus $G\in\F^X_{T_n}$. Furthermore, a point which lies in exactly one of the sets $A=A_\infty\cup\bigcup_kA_k$ and $G=G_\infty\cup\bigcup_kG_k$ lies, for some index $i\in\{1,2,\dots\}\cup\{\infty\}$, in exactly one of $A_i$, $G_i$; that is,
\[
  A\,\triangle\,G\subseteq(A_\infty\,\triangle\,G_\infty)\cup\bigcup_{k=1}^\infty(A_k\,\triangle\,G_k)\in\cN^\mu .
\]
Consequently $N\defeq A\,\triangle\,G\in\cN^\mu$ and $A=G\,\triangle\,N\in\sigma(\F^X_{T_n}\cup\cN^\mu)$. We have in fact shown more:
\begin{equation}\label{eq:2-7-12:b}
  \text{for every }A\in\F^\mu_T\text{ and }n\ge 1\text{ there is }G^{(n)}\in\F^X_{T_n}\text{ with }A\,\triangle\,G^{(n)}\in\cN^\mu .
\end{equation}

\emph{Step 3. $\F^\mu_T\subseteq\sigma(\F^X_{T+}\cup\cN^\mu)$.} We first check the hypotheses of the second part of \bk{Problem 1.2.23}. Each $T_n$ is a stopping time of $\{\F^X_t\}$ by Step 1, and it is positive since $T_n\ge 2^{-n}$. On $\{T<\infty\}$ we have $T<T_n\le T+2^{-n}$, and on $\{T=\infty\}$ we have $T_n=\infty$; hence $T<T_n$ on $\{T<\infty\}$ and $T=\inf_{n\ge1}T_n$ on $\Omega$. \bk{Problem 1.2.23} gives
\begin{equation}\label{eq:2-7-12:c}
  \F^X_{T+}=\bigcap_{n=1}^\infty\F^X_{T_n}.
\end{equation}
Moreover $T_{n+1}\le T_n$: this is clear on $\{T=\infty\}$, and if $(k-1)2^{-n}\le T<k2^{-n}$, then $T<2k\cdot2^{-(n+1)}$, so $T_{n+1}\le 2k\cdot2^{-(n+1)}=T_n$. By \bk{Lemma 1.2.15}, the $\sigma$-fields $\F^X_{T_n}$ decrease with $n$:
\begin{equation}\label{eq:2-7-12:d}
  \F^X_{T_{n}}\subseteq\F^X_{T_m},\qquad n\ge m\ge 1 .
\end{equation}

Let now $A\in\F^\mu_T$, choose the sets $G^{(n)}$, $n\ge1$, as in \eqref{eq:2-7-12:b}, and define, as in the proof of \bk{Proposition 7.7},
\[
  G\defeq\bigcap_{m=1}^\infty\bigcup_{n=m}^\infty G^{(n)} .
\]
The sets $\bigcup_{n\ge m}G^{(n)}$ decrease with $m$, so $G=\bigcap_{m\ge M}\bigcup_{n\ge m}G^{(n)}$ for every $M\ge1$; by \eqref{eq:2-7-12:d}, $\bigcup_{n\ge m}G^{(n)}\in\F^X_{T_m}\subseteq\F^X_{T_M}$ for $m\ge M$, whence $G\in\F^X_{T_M}$. As $M$ is arbitrary, \eqref{eq:2-7-12:c} yields $G\in\F^X_{T+}$. Finally, if $\omega\in G\setminus A$, then $\omega\in G^{(n)}\setminus A$ for some $n$; and if $\omega\in A\setminus G$, then there is $m$ with $\omega\notin\bigcup_{n\ge m}G^{(n)}$, so $\omega\in A\setminus G^{(m)}$. Therefore
\[
  A\,\triangle\,G\subseteq\bigcup_{n=1}^\infty\bigl(A\,\triangle\,G^{(n)}\bigr)\in\cN^\mu ,
\]
and $A=G\,\triangle\,(A\,\triangle\,G)\in\sigma(\F^X_{T+}\cup\cN^\mu)$. This proves $\F^\mu_T\subseteq\sigma(\F^X_{T+}\cup\cN^\mu)$.

\rmk{The strong Markov property of $X$ was not used: the argument is valid for any process $X$ and any optional time $T$ of $\{\F^X_t\}$, with $\F^\mu_T$ defined by the formula of \bk{Exercise 7.11}. The strong Markov property enters in \bk{Exercise 7.11} (through \bk{Proposition 7.7}), which shows that this formula defines the usual $\sigma$-field of the $\{\F^\mu_t\}$-stopping time $T$.}
\end{kex}

\begin{kex}{Exercise}{7.13}{}
\emph{Conventions.} As stipulated before \bk{Exercise 7.11}, $X=\{X_t,\F^X_t;\ 0\le t<\infty\}$ is a
strong Markov process with initial distribution $\mu$ on $(\Omega,\F^X_\infty,P^\mu)$; in particular
$X$ is progressively measurable with respect to $\{\F^X_t\}$ (\bk{Definition 6.2}), hence the map
$(s,\omega)\mapsto X_s(\omega)$ is $\B([0,\infty))\otimes\F^X_\infty$-measurable. The measure
$P^\mu$ is extended to $\F^\mu=\sigma(\F^X_\infty\cup\cN^\mu)$ as in \bk{Problem 7.3} (with
$\G=\F^X_\infty$); the extension agrees with $P^\mu$ on $\F^X_\infty$ and every set of $\cN^\mu$
has measure zero. The conditional probabilities in \bk{(7.1)} are computed on
$(\Omega,\F^\mu,P^\mu)$, those in \bk{(7.2)} on $(\Omega,\F^X_\infty,P^\mu)$; since the two measures
agree on $\F^X_\infty$, a conditional probability of a set of $\F^X_\infty$ given a
sub-$\sigma$-field of $\F^X_\infty$ is the same object on either space. Following \bk{Remark 6.4},
for a random time $R$ we write $\{X_{R+t}\in\Gamma\}\defeq\{R<\infty,\ X_{R+t}\in\Gamma\}$, and
$\sigma(X_R)$ denotes the $\sigma$-field generated by $X_R$ in the sense of \bk{Problem 1.1.17}: the
collection of all sets
\[
  \{R<\infty,\ X_R\in B\}\qquad\text{and}\qquad \{R<\infty,\ X_R\in B\}\cup\{R=\infty\},
  \qquad B\in\B(\R^d).
\]
If $R$ is $\F^X_\infty$-measurable (respectively, $\F^\mu$-measurable) and $u\ge0$, then
$\omega\mapsto(R(\omega)+u,\omega)$ is a measurable map from $\{R<\infty\}$, with the trace of
$\F^X_\infty$ (resp.\ $\F^\mu$), into $[0,\infty)\times\Omega$ with
$\B([0,\infty))\otimes\F^X_\infty$; composing with $X$ we see that
$\{R<\infty,\ X_{R+u}\in B\}$ belongs to $\F^X_\infty$ (resp.\ $\F^\mu$) for every
$B\in\B(\R^d)$. Thus $\sigma(X_R)$ is a sub-$\sigma$-field of $\F^X_\infty$ (resp.\ of $\F^\mu$).

We begin with a general fact about conditioning on $\sigma$-fields which differ only by null sets.

\begin{lem}
Let $\G_1,\G_2$ be sub-$\sigma$-fields of $\F^\mu$ such that for every $C\in\G_1$ there is
$C'\in\G_2$ with $C\triangle C'\in\cN^\mu$, and for every $C'\in\G_2$ there is $C\in\G_1$ with
$C\triangle C'\in\cN^\mu$. Let $A_1,A_2\in\F^\mu$ satisfy $A_1\triangle A_2\in\cN^\mu$. If $W_i$ is a
version of $P^\mu[A_i\mid\G_i]$, $i=1,2$, then $W_1=W_2$, $P^\mu$-a.s.
\end{lem}

\begin{proof}
For $i=1,2$ let
$\cH_i\defeq\{F\subseteq\Omega;\ \exists\, G\in\G_i \text{ with } F\triangle G\in\cN^\mu\}$. This
is a $\sigma$-field: $\Omega\in\cH_i$; if $F\triangle G\in\cN^\mu$ then
$F^c\triangle G^c=F\triangle G\in\cN^\mu$; and if $F_n\triangle G_n\in\cN^\mu$ for all $n\ge1$, then
$\bigl(\bigcup_n F_n\bigr)\triangle\bigl(\bigcup_n G_n\bigr)\subseteq\bigcup_n(F_n\triangle G_n)
\in\cN^\mu$, because $\cN^\mu$ is closed under countable unions and under passage to subsets. It
contains $\G_i$ (take $G=F$) and $\cN^\mu$ (take $G=\emptyset$), so
$\sigma(\G_i\cup\cN^\mu)\subseteq\cH_i$; conversely, each $F\in\cH_i$ can be written
$F=G\triangle(F\triangle G)$ with $G\in\G_i$ and $F\triangle G\in\cN^\mu$, so
$\cH_i\subseteq\sigma(\G_i\cup\cN^\mu)$. Hence $\cH_i=\sigma(\G_i\cup\cN^\mu)\subseteq\F^\mu$. The
hypothesis says $\G_1\subseteq\cH_2$ and $\G_2\subseteq\cH_1$; as $\cN^\mu\subseteq\cH_1\cap\cH_2$,
we obtain $\cH_1\subseteq\cH_2$ and $\cH_2\subseteq\cH_1$. Let $\mathscr{K}\defeq\cH_1=\cH_2$.

Fix $i\in\{1,2\}$. The random variable $W_i$ is $\G_i$-measurable, hence
$\mathscr{K}$-measurable. Let $F\in\mathscr{K}=\cH_i$ and choose $G\in\G_i$ with
$F\triangle G\in\cN^\mu$. Since $\ind_F=\ind_G$ and $\ind_{A_i}=\ind_{A_1}$ hold $P^\mu$-a.s., and
$0\le W_i\le1$ a.s.,
\[
  E^\mu[W_i\ind_F]=E^\mu[W_i\ind_G]=P^\mu(A_i\cap G)=P^\mu(A_1\cap F).
\]
Therefore both $W_1$ and $W_2$ are versions of $P^\mu[A_1\mid\mathscr{K}]$, and the a.s.\ uniqueness
of conditional expectation gives $W_1=W_2$, $P^\mu$-a.s.
\end{proof}

Now let $S$ be an optional time of $\{\F^\mu_t\}$, and fix $t\ge0$ and $\Gamma\in\B(\R^d)$. By
\bk{Exercise 7.11} there is an optional time $T$ of $\{\F^X_t\}$ with
\[
  N\defeq\{S\ne T\}\in\cN^\mu,\qquad\text{and}\qquad \F^\mu_{S+}=\F^\mu_S=\F^\mu_T ,
\]
where $\F^\mu_T=\{A\in\F^\mu;\ A\cap\{T\le u\}\in\F^\mu_u,\ \forall\, 0\le u<\infty\}$. We set
\[
  A_S\defeq\{S<\infty,\ X_{S+t}\in\Gamma\}\in\F^\mu,\qquad
  A_T\defeq\{T<\infty,\ X_{T+t}\in\Gamma\}\in\F^X_\infty .
\]
These two sets coincide on $\{S=T\}$, so $A_S\triangle A_T\subseteq N$ and
$A_S\triangle A_T\in\cN^\mu$.

\emph{Step 1: the hypothesis \bk{(7.2)} holds a.s.\ on all of $\Omega$.} The set $D\defeq\{T=\infty\}$
belongs to $\sigma(X_T)$ (take $B=\emptyset$ above), and also to $\F^X_{T+}$, since
$D\in\F^X_\infty$ and $D\cap\{T\le u\}=\emptyset\in\F^X_{u+}$ for all $u\ge0$
(\bk{Definition 1.2.20}). If $\G$ is either of these two $\sigma$-fields and $W$ a version of
$P^\mu[A_T\mid\G]$, then $W\ge0$ a.s.\ and $E^\mu[W\ind_D]=P^\mu(A_T\cap D)=0$, whence $W=0$ a.s.\
on $D$. Together with \bk{(7.2)} this shows: if $W'$ is a version of $P^\mu[A_T\mid\F^X_{T+}]$ and
$W''$ a version of $P^\mu[A_T\mid\sigma(X_T)]$, then
\begin{equation}\label{eq:2-7-13:a}
  W'=W'',\qquad P^\mu\text{-a.s.}
\end{equation}

\emph{Step 2: $\F^\mu_{S+}$ and $\F^X_{T+}$ differ only by null sets.} We verify the hypotheses of
the Lemma for $\G_1=\F^X_{T+}$ and $\G_2=\F^\mu_{S+}=\F^\mu_T$. First, $\F^X_{T+}\subseteq\F^\mu_T$:
if $C\in\F^X_{T+}$, then $C\in\F^X_\infty\subseteq\F^\mu$ and, for every $u\ge0$,
\[
  C\cap\{T\le u\}\in\F^X_{u+}=\bigcap_{\eps>0}\F^X_{u+\eps}
  \subseteq\bigcap_{\eps>0}\F^\mu_{u+\eps}=\F^\mu_{u+}=\F^\mu_u ,
\]
the last equality being the right-continuity of the augmented filtration
(\bk{Proposition 7.7}). So for $C\in\G_1$ we may take $C'=C$. Secondly, \bk{Exercise 7.12} asserts
$\F^\mu_T\subseteq\sigma(\F^X_{T+}\cup\cN^\mu)$, and by \bk{Problem 7.3} (applied to the
sub-$\sigma$-field $\F^X_{T+}$ of $\F^X_\infty$) every set $C'$ of the latter $\sigma$-field
admits $C\in\F^X_{T+}$ with $C\triangle C'\in\cN^\mu$. The Lemma, with $A_1=A_T$ and $A_2=A_S$,
now yields: if $W'$ is a version of $P^\mu[A_T\mid\F^X_{T+}]$ and $V'$ a version of
$P^\mu[A_S\mid\F^\mu_{S+}]$, then
\begin{equation}\label{eq:2-7-13:b}
  V'=W',\qquad P^\mu\text{-a.s.}
\end{equation}

\emph{Step 3: $\sigma(X_S)$ and $\sigma(X_T)$ differ only by null sets.} For $B\in\B(\R^d)$ the sets
$\{S<\infty,\ X_S\in B\}$ and $\{T<\infty,\ X_T\in B\}$ coincide on $\{S=T\}$, and so do
$\{S=\infty\}$ and $\{T=\infty\}$. Consequently, the correspondence
\begin{align*}
  \{T<\infty,\ X_T\in B\}&\leftrightarrow\{S<\infty,\ X_S\in B\},\\
  \{T<\infty,\ X_T\in B\}\cup\{T=\infty\}&\leftrightarrow\{S<\infty,\ X_S\in B\}\cup\{S=\infty\}
\end{align*}
associates with every set $C\in\sigma(X_T)$ a set $C'\in\sigma(X_S)$ with
$C\triangle C'\subseteq N$, hence $C\triangle C'\in\cN^\mu$, and every set of $\sigma(X_S)$ arises
in this way. The hypotheses of the Lemma hold for $\G_1=\sigma(X_T)$, $\G_2=\sigma(X_S)$,
$A_1=A_T$, $A_2=A_S$: if $W''$ is a version of $P^\mu[A_T\mid\sigma(X_T)]$ and $V''$ a version of
$P^\mu[A_S\mid\sigma(X_S)]$, then
\begin{equation}\label{eq:2-7-13:c}
  V''=W'',\qquad P^\mu\text{-a.s.}
\end{equation}

\emph{Conclusion.} Combining \eqref{eq:2-7-13:b}, \eqref{eq:2-7-13:a} and \eqref{eq:2-7-13:c}, we
obtain $V'=W'=W''=V''$, $P^\mu$-a.s., that is,
\[
  P^\mu[X_{S+t}\in\Gamma\mid\F^\mu_{S+}]=P^\mu[X_{S+t}\in\Gamma\mid X_S],\qquad P^\mu\text{-a.s.}
\]
on $\Omega$, and in particular on $\{S<\infty\}$. This is \bk{(7.1)}. Since $X$ is also
progressively measurable with respect to the larger filtration $\{\F^\mu_t\}$ and
$P^\mu[X_0\in\Gamma]=\mu(\Gamma)$, the process $\{X_t,\F^\mu_t;\ t\ge0\}$ on
$(\Omega,\F^\mu,P^\mu)$ is a strong Markov process with initial distribution $\mu$ in the sense of
\bk{Definition 6.2}.

\rmk{The hypothesis \bk{(7.2)} is exactly condition (ii) of \bk{Definition 6.2} for
$\{X_t,\F^X_t\}$, which is in force throughout \bk{Exercises 7.11--7.13}; it entered above both
directly, in Step 1, and through \bk{Proposition 7.7}, on which Step 2 and \bk{Exercise 7.11}
rest.}
\end{kex}

\begin{kex}{Problem}{7.14}{}
We write $\Omega=C[0,\infty)$, so that $B_t(\omega)=\omega(t)$. Each coordinate mapping $B_t$ is
continuous on $\Omega$, hence $\B(\Omega)$-measurable, and therefore
$\F^B_t=\sigma(B_s;\ 0\le s\le t)\subseteq\B(\Omega)$ for every $t\ge0$. Under $P^0$ the process
$\{B_t,\F^B_t;\ t\ge0\}$ is a standard Brownian motion (\bk{Definition 4.21}, where Wiener
measure is written $P_*$); in particular, by \bk{Definition 1.1}, $P^0[B_0=0]=1$ and $B_t-B_s$ is
normally distributed with mean zero and variance $t-s$ for $0\le s<t$. In the notation of \bk{Definition 7.2}, with $X=B$ and $P^\mu=P^0$ (that is,
$\mu=\delta_0$), we have
\[
  \ol{\F}_0=\sigma\bigl(\F^B_0\cup\cN_0\bigr),\qquad
  \cN_0=\{A\subseteq\Omega;\ \exists\, G\in\F^B_0 \text{ with } A\subseteq G,\ P^0(G)=0\}.
\]

For every integer $n\ge1$ we set
\[
  F_n\defeq\{\omega\in\Omega;\ \omega(t)=\omega(0) \text{ for all } t\in[0,1/n]\}.
\]
Since every $\omega\in\Omega$ is continuous and the rationals of $[0,1/n]$ are dense in this
interval, $\omega$ is constant on $[0,1/n]$ if and only if $\omega(q)=\omega(0)$ for every rational
$q\in[0,1/n]$. Thus
\begin{equation}\label{eq:2-7-14:a}
  F_n=\bigcap_{q\in\Q\cap[0,1/n]}\{B_q-B_0=0\}\ \in\F^B_{1/n},
\end{equation}
a countable intersection of sets of $\F^B_{1/n}$. Clearly $F_n\subseteq F_{n+1}\subseteq F$ for
every $n\ge1$. Conversely, if $\omega\in F$ is constant on $[0,\eps]$ with $\eps>0$, then
$\omega\in F_n$ for every integer $n\ge1/\eps$. Consequently
\begin{equation}\label{eq:2-7-14:b}
  F=\bigcup_{n\ge m}F_n\qquad\text{for every integer } m\ge1 .
\end{equation}

\textbf{(i)} By \eqref{eq:2-7-14:a} and \eqref{eq:2-7-14:b} with $m=1$, the set $F$ belongs to
$\F^B_1\subseteq\B(\Omega)$, so $P^0(F)$ is defined. For every $n\ge1$ we have
$F_n\subseteq\{B_{1/n}-B_0=0\}$, and under $P^0$ the increment $B_{1/n}-B_0$ is, as recalled
above, normally distributed with mean zero and variance $1/n>0$; its distribution has a density,
whence
$P^0[B_{1/n}-B_0=0]=0$ and $P^0(F_n)=0$. Countable subadditivity in \eqref{eq:2-7-14:b} gives
$P^0(F)\le\sum_{n\ge1}P^0(F_n)=0$.

\textbf{(ii)} Fix $\eps>0$ and choose an integer $m\ge1/\eps$. For $n\ge m$ we have
$F_n\in\F^B_{1/n}\subseteq\F^B_{1/m}\subseteq\F^B_\eps$ by \eqref{eq:2-7-14:a}, and
\eqref{eq:2-7-14:b} exhibits $F$ as a countable union of such sets; hence $F\in\F^B_\eps$. As
$\eps>0$ is arbitrary, $F\in\bigcap_{\eps>0}\F^B_\eps=\F^B_{0+}$.

\textbf{(iii)} Let $\Omega_0\defeq\{\omega\in\Omega;\ \omega(0)=0\}$ and consider the collection
\[
  \cD\defeq\{A\subseteq\Omega;\ A\cap\Omega_0=\emptyset \ \text{ or }\ A\cap\Omega_0=\Omega_0\}.
\]
This is a $\sigma$-field of subsets of $\Omega$: it contains $\Omega$; if $A\in\cD$ then
$A^c\cap\Omega_0=\Omega_0\setminus(A\cap\Omega_0)$ is $\Omega_0$ or $\emptyset$, so $A^c\in\cD$;
and if $A_1,A_2,\dots\in\cD$, then $\bigl(\bigcup_kA_k\bigr)\cap\Omega_0=\bigcup_k(A_k\cap\Omega_0)$
is $\Omega_0$ if $A_k\cap\Omega_0=\Omega_0$ for some $k$, and is $\emptyset$ otherwise.

We show that $\F^B_0\cup\cN_0\subseteq\cD$. First,
$\F^B_0=\sigma(B_0)=\{B_0^{-1}(\Gamma);\ \Gamma\in\B(\R)\}$, and for $\Gamma\in\B(\R)$,
\[
  B_0^{-1}(\Gamma)\cap\Omega_0=\{\omega;\ \omega(0)\in\Gamma,\ \omega(0)=0\}=
  \begin{cases}\Omega_0, & 0\in\Gamma,\\ \emptyset, & 0\notin\Gamma,\end{cases}
\]
so $\F^B_0\subseteq\cD$. Secondly, let $A\in\cN_0$, and let $G=B_0^{-1}(\Gamma)\in\F^B_0$,
$\Gamma\in\B(\R)$, satisfy $A\subseteq G$ and $P^0(G)=0$. Because $P^0[B_0=0]=1$, we have
$P^0(G)=P^0[B_0\in\Gamma]=\ind_\Gamma(0)$, so $0\notin\Gamma$ and $G\cap\Omega_0=\emptyset$ by the
preceding display; a fortiori $A\cap\Omega_0=\emptyset$ and $A\in\cD$.

Since $\cD$ is a $\sigma$-field containing $\F^B_0\cup\cN_0$, it contains
$\ol{\F}_0=\sigma(\F^B_0\cup\cN_0)$. It remains to observe that $F\notin\cD$. The function
$\omega_1\equiv0$ belongs to $\Omega_0$ and is constant on $[0,1]$, so
$\omega_1\in F\cap\Omega_0$ and $F\cap\Omega_0\ne\emptyset$. The function $\omega_2(t)=t$,
$t\ge0$, belongs to $\Omega_0$, but for every $\eps>0$ it takes the distinct values $0$ and $\eps$
on $[0,\eps]$, so $\omega_2\notin F$ and $F\cap\Omega_0\ne\Omega_0$. Hence $F\notin\cD$, and
therefore $F\notin\ol{\F}_0$.

\rmk{By (ii) and the inclusion $\F^B_\eps\subseteq\ol{\F}_\eps$, we have
$F\in\bigcap_{\eps>0}\ol{\F}_\eps=\ol{\F}_{0+}$, whereas $F\notin\ol{\F}_0$ by (iii): the
filtration $\{\ol{\F}_t\}$ obtained by completing each $\F^B_t$ separately is not
right-continuous at $t=0$. By (i), on the other hand, $F$ is a $P^0$-null set of $\F^B_\infty$,
so it belongs to $\cN^\mu$ and thus to the augmentation $\F^\mu_0$ of
\bk{Definition 7.2}\,(ii).}
\end{kex}

\solsection{2.8}{Computations Based on Passage Times}

\begin{kex}{Problem}{8.2}{}
We fix $t>0$ and write
\[
  p(x)\defeq\frac{1}{\sqrt{2\pi t}}\,e^{-x^2/2t},\quad x\in\R,
  \qquad\qquad
  g(b)\defeq 2p(b)=\sqrt{\frac{2}{\pi t}}\,e^{-b^2/2t},\quad b>0 ,
\]
so that $\int_0^\infty g(b)\,db=\int_{-\infty}^{\infty}p(x)\,dx=1$. Throughout Section 2.8,
$\{W_t,\F_t\}$, $(\Omega,\F)$, $\{P^x\}_{x\in\R}$ is a one-dimensional Brownian family; by
\bk{Definition 5.8}(iii), under $P^0$ this process is a one-dimensional Brownian motion starting at
the origin, that is, a Brownian motion with initial distribution $\mu=\delta_0$ in the sense of
\bk{Definition 5.1} with $d=1$: it is continuous and adapted, $P^0[W_0=0]=1$, and for $0\le s<u$ the
increment $W_u-W_s$ is independent of $\F_s$ and normal with mean zero and variance $u-s$. These are
precisely the requirements of \bk{Definition 1.1}, so under $P^0$ the process $\{W_t,\F_t\}$ is a
standard, one-dimensional Brownian motion. In particular $W_0=0$, $P^0$-a.s., and $W_t=W_t-W_0$ has under $P^0$ the normal
distribution with mean zero and variance $t$, that is, the density $p$.
The assertion \bk{(8.3)} means that each of the three random variables $M_t$, $\abs{W_t}$, $M_t-W_t$ has,
under $P^0$, the density $g$ on $(0,\infty)$: its law gives the mass $\int_B g(b)\,db$ to every
$B\in\B((0,\infty))$, and consequently no mass to $(-\infty,0]$. We note that $M_t=\max_{0\le s\le t}W_s
=\sup_{s\in[0,t]\cap\Q}W_s$ is a random variable, by the continuity of the paths, and that
\begin{equation}\label{eq:2-8-02:a}
  M_t\ge W_t \ \text{ everywhere,}\qquad M_t\ge W_0=0 \quad P^0\text{-a.s.}
\end{equation}

\emph{Step 1: the law of $\abs{W_t}$.} For $B\in\B((0,\infty))$ the sets $B$ and $-B$ are disjoint, and
$p$ is even; hence
\[
  P^0[\abs{W_t}\in B]=P^0[W_t\in B]+P^0[W_t\in -B]=\int_B p(x)\,dx+\int_{-B}p(x)\,dx=\int_B g(b)\,db .
\]

\emph{Step 2: the law of $M_t$.} Let $D\defeq\{(a,b)\in\R^2:\ b\ge 0,\ a\le b\}$ and
\[
  f(a,b)\defeq\frac{2(2b-a)}{\sqrt{2\pi t^3}}\exp\Bigl\{-\frac{(2b-a)^2}{2t}\Bigr\},\qquad (a,b)\in D ,
\]
a nonnegative function on $D$. \bk{Proposition 8.1} states that
\begin{equation}\label{eq:2-8-02:b}
  P^0[(W_t,M_t)\in A]=\iint_A f(a,b)\,da\,db\qquad\text{for every Borel set } A\subseteq D .
\end{equation}
For fixed $b\ge0$ we have
$\frac{\partial}{\partial a}\bigl[2p(2b-a)\bigr]=\frac{2(2b-a)}{t}\,p(2b-a)=f(a,b)$ and
$p(2b-a)\to0$ as $a\to-\infty$, so that
\begin{equation}\label{eq:2-8-02:c}
  \int_{-\infty}^{c}f(a,b)\,da=2p(2b-c),\qquad b\ge 0,\ c\le b .
\end{equation}
Let $B\in\B((0,\infty))$. By \eqref{eq:2-8-02:a} the pair $(W_t,M_t)$ lies in $D$, $P^0$-a.s., so the event
$\{M_t\in B\}$ differs from $\{(W_t,M_t)\in A\}$, where $A\defeq\{(a,b)\in D:\ b\in B\}$, by a $P^0$-null
set. From \eqref{eq:2-8-02:b}, Tonelli's theorem and \eqref{eq:2-8-02:c} with $c=b$ we obtain
\[
  P^0[M_t\in B]=\int_B\Bigl(\int_{-\infty}^{b}f(a,b)\,da\Bigr)db=\int_B 2p(b)\,db=\int_B g(b)\,db .
\]
In particular, for every $b>0$,
\begin{equation}\label{eq:2-8-02:d}
  P^0[M_t\ge b]=\int_b^\infty 2p(x)\,dx=2P^0[W_t\ge b].
\end{equation}

\emph{Step 3: the law of $M_t-W_t$.} Let $c>0$. The set $A_c\defeq\{(a,b):\ b\ge 0,\ a\le b-c\}$ is a
Borel subset of $D$, and by \eqref{eq:2-8-02:a} the events $\{M_t-W_t\ge c\}$ and $\{(W_t,M_t)\in A_c\}$
differ by a $P^0$-null set. Thus \eqref{eq:2-8-02:b}, Tonelli's theorem, \eqref{eq:2-8-02:c} with $b-c$ in
place of $c$, and the substitution $x=b+c$ give
\[
  P^0[M_t-W_t\ge c]=\int_0^\infty\Bigl(\int_{-\infty}^{b-c}f(a,b)\,da\Bigr)db
  =\int_0^\infty 2p(b+c)\,db=\int_c^\infty g(x)\,dx .
\]
Let $\mu$ be the law of $M_t-W_t$ under $P^0$ and $\nu$ the measure with density $g$ on $(0,\infty)$. We
have just shown that $\mu([c,\infty))=\nu([c,\infty))$ for every $c>0$; letting $c\downarrow0$ we get
$\mu((0,\infty))=\int_0^\infty g(x)\,dx=1$. The intervals $[c,\infty)$, $c>0$, form a $\pi$-system which
generates $\B((0,\infty))$, and the restrictions of $\mu$ and $\nu$ to $(0,\infty)$ are finite measures of
the same total mass which agree on it; by Dynkin's $\pi$-$\lambda$ theorem, $\mu(B)=\nu(B)=\int_B g(b)\,db$
for every $B\in\B((0,\infty))$. Together with Steps 1 and 2, this proves \bk{(8.3)}.

\emph{Step 4: the first inequality of \bk{(8.3)$'$}.} Let $b>0$ and put
$M_t^-\defeq\max_{0\le u\le t}(-W_u)$. Since $\abs{W_u}=\max\{W_u,-W_u\}$,
\begin{equation}\label{eq:2-8-02:e}
  P^0\Bigl[\max_{0\le u\le t}\abs{W_u}\ge b\Bigr]=P^0\bigl[\{M_t\ge b\}\cup\{M_t^-\ge b\}\bigr]
  \le P^0[M_t\ge b]+P^0[M_t^-\ge b].
\end{equation}
We claim that $M_t^-$ and $M_t$ have the same law under $P^0$. As recalled at the outset, under $P^0$ the
process $W$ is a standard Brownian motion, and then so is $-W$ (\bk{Definition 1.1}: it is continuous and adapted, starts at zero
a.s., and its increments $-(W_v-W_u)$ are independent of $\F_u$ and normal with mean zero and variance
$v-u$). Hence $W$ and $-W$ have the same finite-dimensional distributions under $P^0$. Let
$\{q_1,q_2,\dots\}$ be an enumeration of $[0,t]\cap\Q$. For every $n\ge1$ and $x\in\R$,
\[
  P^0\Bigl[\max_{1\le k\le n}W_{q_k}\le x\Bigr]=P^0\Bigl[\max_{1\le k\le n}(-W_{q_k})\le x\Bigr] ,
\]
and as $n\to\infty$ the events on the left decrease to $\{M_t\le x\}$ and those on the right to
$\{M_t^-\le x\}$, because of the continuity of the paths. Therefore $P^0[M_t\le x]=P^0[M_t^-\le x]$ for all
$x\in\R$, which proves the claim. Now \eqref{eq:2-8-02:e} and \eqref{eq:2-8-02:d} yield
\[
  P^0\Bigl[\max_{0\le u\le t}\abs{W_u}\ge b\Bigr]\le 2P^0[M_t\ge b]=4P^0[W_t\ge b].
\]

\emph{Step 5: the second inequality of \bk{(8.3)$'$}.} For $x\ge b>0$ we have $1\le x/b$, whence
\[
  P^0[W_t\ge b]=\int_b^\infty p(x)\,dx\le\frac1b\int_b^\infty x\,p(x)\,dx
  =\frac{1}{b\sqrt{2\pi t}}\Bigl[-t\,e^{-x^2/2t}\Bigr]_{x=b}^{\infty}
  =\sqrt{\frac{t}{2\pi}}\,\frac1b\,e^{-b^2/2t}.
\]
Multiplying by $4$ we obtain the second inequality of \bk{(8.3)$'$}.

\rmk{The identity \eqref{eq:2-8-02:d} is, after the substitution $x\mapsto x\sqrt t$,
the second equality in \bk{(8.4)} of \bk{Remark 8.3}; the first one holds because
$\{T_b\le t\}=\{M_t\ge b\}$ on $\{W_0=0\}$, by the continuity of the paths and the intermediate value
theorem. The equality of the law of $M_t$ with the other two in \bk{(8.3)} holds for each fixed $t$ only: the
process $M$ has nondecreasing paths, whereas $P^0[\abs{W_2}<\abs{W_1}]>0$, so $M$ cannot have the law of
the process $\abs{W}$. The processes $M-W$ and $\abs{W}$, on the other hand, do have the same law under
$P^0$; this is part of L\'evy's theorem, \bk{Theorem 3.6.17}, and consequently $M$ does not have the law of
the process $M-W$ either.}
\end{kex}

\begin{kex}{Exercise}{8.4}{}
We fix $b>0$, $\alpha>0$, put $\lambda=\sqrt{2\alpha}$ and work under $P^0$, under which $W=\{W_t,\F_t;\,0\le t<\infty\}$ is a standard Brownian motion: it has continuous paths, $P^0[W_0=0]=1$, and for $0\le s<t$ the increment $W_t-W_s$ is independent of $\F_s$ and normal with mean zero and variance $t-s$. On $\{T_b=\infty\}$ we read $e^{-\alpha T_b}$ as $0$, so that \bk{(8.6)} is the assertion
\begin{equation}\label{eq:2-8-04:goal}
  E^0\bigl[e^{-\alpha T_b}\,\ind_{\{T_b<\infty\}}\bigr]=e^{-b\sqrt{2\alpha}} .
\end{equation}

\emph{Step 1. $X$ is a continuous martingale.} The process $X$ is adapted, continuous and strictly positive. Since $W_t$ is normal with mean zero and variance $t$ under $P^0$, we have $E^0X_t=e^{-\lambda^2t/2}E^0e^{\lambda W_t}=1<\infty$. For $0\le s<t$, the random variable $X_s$ is $\F_s$-measurable and $W_t-W_s$ is independent of $\F_s$, so
\[
  E^0[X_t\mid\F_s]
  =X_s\,e^{-\lambda^2(t-s)/2}\,E^0\bigl[e^{\lambda(W_t-W_s)}\bigr]
  =X_s\,e^{-\lambda^2(t-s)/2}\,e^{\lambda^2(t-s)/2}=X_s\qquad\text{a.s. }P^0 .
\]

\emph{Step 2. Optional sampling at $T_b\wedge n$.} The passage time $T_b=\inf\{t\ge0;\,W_t=b\}$ of \bk{(6.1)} is a stopping time of $\{\F_t\}$, because $W$ is continuous and adapted and $\{b\}$ is closed (\bk{Problem 1.2.7}); hence so is the bounded time $T_b\wedge n$ for every integer $n\ge1$. Both $X$ and $-X$ are right-continuous submartingales, and \bk{Problem 1.3.23}\,(i), applied to each of them with the bounded optional times $S=0\le T=T_b\wedge n$, gives $E^0X_{T_b\wedge n}\ge E^0X_0$ and $E^0X_{T_b\wedge n}\le E^0X_0$. As $X_0=e^{\lambda W_0}=1$ a.s.\ $P^0$,
\begin{equation}\label{eq:2-8-04:os}
  E^0X_{T_b\wedge n}=1\qquad\text{for every } n\ge1 .
\end{equation}

\emph{Step 3. Passage to the limit.} Let $\Omega_0=\{W_0=0\}$, an event of $P^0$-probability one, and fix $\omega\in\Omega_0$. If we had $W_t(\omega)>b$ for some $t<T_b(\omega)$, then, as $W_0(\omega)=0<b$, the intermediate value theorem would produce $s\in(0,t)$ with $W_s(\omega)=b$, contradicting $s<T_b(\omega)$. Hence $W_t(\omega)\le b$ for $0\le t<T_b(\omega)$; moreover, if $T_b(\omega)<\infty$, then $W_{T_b(\omega)}(\omega)=b$, because the set $\{t\ge0;\,W_t(\omega)=b\}$ is closed and nonempty and so contains its infimum. Consequently $W_{T_b\wedge n}\le b$ on $\Omega_0$, and since $\lambda>0$,
\begin{equation}\label{eq:2-8-04:bd}
  0<X_{T_b\wedge n}
  =\exp\bigl\{\lambda W_{T_b\wedge n}-\tfrac12\lambda^2(T_b\wedge n)\bigr\}
  \le e^{\lambda b}\qquad\text{on }\Omega_0,\ n\ge1 .
\end{equation}
We identify the limit of $X_{T_b\wedge n}$ as $n\to\infty$ on $\Omega_0$. On $\Omega_0\cap\{T_b<\infty\}$ we have $T_b\wedge n=T_b$ for all $n\ge T_b$, so
\[
  \lim_{n\to\infty}X_{T_b\wedge n}=X_{T_b}=\exp\bigl\{\lambda b-\tfrac12\lambda^2T_b\bigr\}=e^{\lambda b}e^{-\alpha T_b},
\]
because $\tfrac12\lambda^2=\alpha$. On $\Omega_0\cap\{T_b=\infty\}$ we have $T_b\wedge n=n$ and $W_n\le b$, so
\[
  0<X_{T_b\wedge n}=X_n\le\exp\bigl\{\lambda b-\tfrac12\lambda^2n\bigr\}\xrightarrow[n\to\infty]{}0 .
\]
Thus $X_{T_b\wedge n}\to e^{\lambda b}e^{-\alpha T_b}\ind_{\{T_b<\infty\}}$ a.s.\ $P^0$, and by \eqref{eq:2-8-04:bd} the convergence is dominated by the constant $e^{\lambda b}$. The bounded convergence theorem and \eqref{eq:2-8-04:os} yield
\[
  1=\lim_{n\to\infty}E^0X_{T_b\wedge n}=e^{\lambda b}\,E^0\bigl[e^{-\alpha T_b}\ind_{\{T_b<\infty\}}\bigr],
\]
which is \eqref{eq:2-8-04:goal}, i.e.\ \bk{(8.6)}, since $\lambda b=b\sqrt{2\alpha}$.

Letting $\alpha\downarrow0$ in \eqref{eq:2-8-04:goal}, the integrands $e^{-\alpha T_b}\ind_{\{T_b<\infty\}}$ increase to $\ind_{\{T_b<\infty\}}$, and the monotone convergence theorem gives
\begin{equation}\label{eq:2-8-04:finite}
  P^0[T_b<\infty]=\lim_{\alpha\downarrow0}e^{-b\sqrt{2\alpha}}=1 .
\end{equation}

\emph{Step 4. The Laplace transform of the density in \bk{(8.5)}.} Let
\[
  f(t)\defeq\frac{b}{\sqrt{2\pi t^3}}\,e^{-b^2/2t},\qquad t>0 .
\]
We show that
\begin{equation}\label{eq:2-8-04:lt}
  \int_0^\infty e^{-\alpha t}f(t)\,dt=e^{-b\sqrt{2\alpha}}\qquad\text{for every }\alpha>0 .
\end{equation}
With $\lambda=\sqrt{2\alpha}$ we have $\frac{b^2}{2t}+\alpha t=\frac{(b-\lambda t)^2}{2t}+\lambda b$, so the left-hand side of \eqref{eq:2-8-04:lt} equals $e^{-\lambda b}A$, where
\[
  A\defeq\int_0^\infty\frac{b}{\sqrt{2\pi}}\,t^{-3/2}\exp\Bigl\{-\frac{(b-\lambda t)^2}{2t}\Bigr\}dt,
  \qquad
  B\defeq\int_0^\infty\frac{\lambda}{\sqrt{2\pi}}\,t^{-1/2}\exp\Bigl\{-\frac{(b-\lambda t)^2}{2t}\Bigr\}dt ;
\]
both integrands are nonnegative, so $A,B\in[0,\infty]$ are well defined, and we have to prove $A=1$.

First, $A=B$. The map $t\mapsto s=b^2/(\lambda^2t)$ is a decreasing $C^1$-bijection of $(0,\infty)$ onto itself, with $t=b^2/(\lambda^2s)$, $dt=-b^2\lambda^{-2}s^{-2}\,ds$ and $t^{-3/2}=\lambda^3b^{-3}s^{3/2}$, so that $b\,t^{-3/2}\abs{dt}=\lambda\,s^{-1/2}\,ds$; furthermore
\[
  \frac{(b-\lambda t)^2}{t}=\frac{b^2}{t}-2\lambda b+\lambda^2t=\lambda^2s-2\lambda b+\frac{b^2}{s}=\frac{(b-\lambda s)^2}{s} .
\]
The change of variables formula therefore transforms $A$ into $B$.

Secondly, $A+B=2$. The function $u(t)=b\,t^{-1/2}-\lambda\,t^{1/2}$ has derivative $u'(t)=-\tfrac12\bigl(b\,t^{-3/2}+\lambda\,t^{-1/2}\bigr)<0$ and satisfies $u(0+)=+\infty$, $u(\infty)=-\infty$; it is thus a decreasing $C^1$-bijection of $(0,\infty)$ onto $\R$, and $u(t)^2=(b-\lambda t)^2/t$. Hence
\[
  \frac{A+B}{2}=\int_0^\infty\frac{1}{\sqrt{2\pi}}\,e^{-u(t)^2/2}\,\abs{u'(t)}\,dt=\int_{-\infty}^\infty\frac{1}{\sqrt{2\pi}}\,e^{-u^2/2}\,du=1 .
\]
The two facts give $A=B=1$, and \eqref{eq:2-8-04:lt} is proved.

\emph{Step 5. Derivation of \bk{(8.5)}.} Let $\mu(\Gamma)\defeq P^0[T_b\in\Gamma]$ and $\nu(\Gamma)\defeq\int_\Gamma f(t)\,dt$ for $\Gamma\in\B([0,\infty))$. These are finite measures on $[0,\infty)$: $\mu$ has total mass $P^0[T_b<\infty]\le1$, and $\nu$ has total mass $\lim_{\alpha\downarrow0}\int_0^\infty e^{-\alpha t}f(t)\,dt=1$ by monotone convergence and \eqref{eq:2-8-04:lt}. By \eqref{eq:2-8-04:goal} and \eqref{eq:2-8-04:lt},
\[
  \int_{[0,\infty)}e^{-\alpha t}\,\mu(dt)=e^{-b\sqrt{2\alpha}}=\int_{[0,\infty)}e^{-\alpha t}\,\nu(dt)\qquad\text{for every }\alpha>0 .
\]
By the uniqueness theorem for Laplace transforms (two finite measures on $[0,\infty)$ whose Laplace transforms agree on $(0,\infty)$ coincide), $\mu=\nu$. Together with \eqref{eq:2-8-04:finite} this says that $T_b$ is a.s.\ finite under $P^0$ and has the density $f$:
\[
  P^0[T_b\in dt]=\frac{b}{\sqrt{2\pi t^3}}\,e^{-b^2/2t}\,dt;\qquad b>0,\ t>0,
\]
which is \bk{(8.5)}.

\rmk{The derivation uses neither the reflection principle nor \bk{(8.3)}, so it is independent of the route by which the book arrives at \bk{(8.5)}. The sign of $\lambda$ matters in Step~3: with $\lambda=-\sqrt{2\alpha}<0$ the process \bk{(8.7)} is still a martingale and \eqref{eq:2-8-04:os} still holds, but $X_{T_b\wedge n}$ is no longer bounded and the limit cannot be taken under the expectation; indeed $X_{T_b\wedge n}\to e^{-\abs{\lambda}b}e^{-\alpha T_b}$ a.s.\ $P^0$, whose expectation is $e^{-2\abs{\lambda}b}<1$ by \bk{(8.6)}.}
\end{kex}

\begin{kex}{Problem}{8.6}{}
Here $\{W_t,\F_t;\ 0\le t<\infty\}$, $(\Omega,\F)$, $\{P^x\}_{x\in\R}$ is the one-dimensional Brownian
family of this section, $p(t;x,y)=(2\pi t)^{-1/2}e^{-(x-y)^2/2t}$ is the Gaussian kernel of
\bk{(2.6)}, and $T_0=\inf\{t\ge 0;\ W_t=0\}$. The assertion \bk{(8.9)} means that
\begin{equation}\label{eq:2-8-06:a}
  P^x[W_t\in\Gamma,\ T_0>t]=\int_\Gamma p_-(t;x,y)\,dy;\qquad t>0,\ x>0,\ \Gamma\in\B((0,\infty)).
\end{equation}
We prove \eqref{eq:2-8-06:a} in Steps 1 and 2, and in Step 3 we show that it does furnish the
transition density of the absorbed process.

We shall use the following facts. The paths of $W$ are continuous and $P^x[W_0=x]=1$
(\bk{Definition 5.8}). The set $\{0\}$ is closed, so $T_0$ is a stopping time of $\{\F_t\}$
(\bk{Problem 1.2.7}), hence also an optional time; since the set $\{u\ge 0;\ W_u(\omega)=0\}$ is
closed, the infimum defining $T_0(\omega)$ is attained when it is finite, and therefore
\begin{equation}\label{eq:2-8-06:b}
  W_{T_0}=0 \ \text{ on } \{T_0<\infty\},\qquad
  \{T_0>t\}=\{W_u\neq 0 \text{ for all } u\in[0,t]\};\quad 0\le t<\infty .
\end{equation}

\emph{Step 1: reflection at $T_0$.} Fix $t>0$, $x>0$ and $a>0$. We claim that
\begin{equation}\label{eq:2-8-06:c}
  P^x[W_t>a,\ T_0\le t]=P^x[W_t<-a,\ T_0\le t]=P^x[W_t<-a].
\end{equation}
The Brownian family is a strong Markov family (\bk{Theorem 6.15}) and $W$ is continuous, so
\bk{Corollary 6.18} is applicable with the optional time $S=T_0$ and the random time
\[
  T\defeq t \ \text{ on } \{T_0\le t\},\qquad T\defeq\infty \ \text{ on } \{T_0>t\}.
\]
Indeed, $T\ge S$ everywhere, and $T$ is $\F_{S+}$-measurable because it takes only the values $t$
and $\infty$ and, for every $u\ge 0$,
$\{T=t\}\cap\{T_0\le u\}=\{T_0\le t\wedge u\}\in\F_{t\wedge u}\subseteq\F_{u+}$, i.e.,
$\{T_0\le t\}\in\F_{T_0+}$. On $\{T<\infty\}=\{T_0\le t\}$ we have $W_T=W_t$ and, by
\eqref{eq:2-8-06:b}, $W_S=0$. Thus \bk{Corollary 6.18} yields, for every $\Gamma\in\B(\R)$,
\begin{equation}\label{eq:2-8-06:d}
  P^x[W_t\in\Gamma\,|\,\F_{T_0+}](\omega)=\bigl(U_{t-T_0(\omega)}\ind_\Gamma\bigr)(0),
  \qquad P^x\text{-a.e. } \omega\in\{T_0\le t\},
\end{equation}
where $(U_r\ind_\Gamma)(0)=P^0[W_r\in\Gamma]$ by \bk{(5.7)}. For $r>0$ the random variable $W_r$ is,
under $P^0$, normal with mean zero and variance $r$, hence symmetric, and
$P^0[W_r>a]=P^0[W_r<-a]$; for $r=0$ both probabilities vanish, because $P^0[W_0=0]=1$ and $a>0$.
Consequently
\[
  \bigl(U_{r}\ind_{(a,\infty)}\bigr)(0)=\bigl(U_{r}\ind_{(-\infty,-a)}\bigr)(0);\qquad 0\le r\le t,
\]
and \eqref{eq:2-8-06:d}, used with $\Gamma=(a,\infty)$ and with $\Gamma=(-\infty,-a)$, gives
\[
  P^x[W_t>a\,|\,\F_{T_0+}]=P^x[W_t<-a\,|\,\F_{T_0+}],\qquad P^x\text{-a.s. on } \{T_0\le t\}.
\]
Integrating over the set $\{T_0\le t\}\in\F_{T_0+}$ with respect to $P^x$ we obtain the first
equality in \eqref{eq:2-8-06:c}. For the second, let $\omega$ be such that $W_0(\omega)=x>0$ and
$W_t(\omega)<-a<0$; by the intermediate value theorem the continuous path $u\mapsto W_u(\omega)$
vanishes somewhere in $(0,t)$, so $T_0(\omega)\le t$. Hence
$\{W_t<-a\}\cap\{W_0=x\}\subseteq\{T_0\le t\}$, and since $P^x[W_0=x]=1$ the second equality in
\eqref{eq:2-8-06:c} follows.

\emph{Step 2: proof of \eqref{eq:2-8-06:a}.} Under $P^x$ the random variable $W_t$ is normal with
mean $x$ and variance $t$, i.e., has density $p(t;x,\cdot)$. From \eqref{eq:2-8-06:c} and the
change of variable $z=-y$ we get, for every $a>0$,
\begin{align}
  P^x[W_t>a,\ T_0>t]
  &=P^x[W_t>a]-P^x[W_t>a,\ T_0\le t]
   =P^x[W_t>a]-P^x[W_t<-a] \notag\\
  &=\int_a^\infty p(t;x,y)\,dy-\int_{-\infty}^{-a}p(t;x,z)\,dz
   =\int_a^\infty p_-(t;x,y)\,dy. \label{eq:2-8-06:e}
\end{align}
The function $p_-(t;x,\cdot)$ is continuous, and it is strictly positive on $(0,\infty)$, since
for $x,y>0$
\begin{equation}\label{eq:2-8-06:f}
  p_-(t;x,y)=\frac{1}{\sqrt{2\pi t}}\,e^{-(x-y)^2/2t}\bigl(1-e^{-2xy/t}\bigr)>0 .
\end{equation}
Therefore
\[
  \mu(\Gamma)\defeq P^x[W_t\in\Gamma,\ T_0>t],\qquad
  \nu(\Gamma)\defeq\int_\Gamma p_-(t;x,y)\,dy;\qquad \Gamma\in\B((0,\infty)),
\]
are finite measures on $((0,\infty),\B((0,\infty)))$ (we have $\nu((0,\infty))\le\int_\R p(t;x,y)\,dy=1$).
By \eqref{eq:2-8-06:e} they agree on the intervals $(a,\infty)$, $a>0$, and, letting $a\downarrow 0$
and using continuity from below, also on $(0,\infty)$. These sets form a class closed under finite
intersections which generates $\B((0,\infty))$, and the collection of $\Gamma\in\B((0,\infty))$ with
$\mu(\Gamma)=\nu(\Gamma)$ is a Dynkin system because $\mu((0,\infty))=\nu((0,\infty))<\infty$. By
Dynkin's $\pi$-$\lambda$ theorem, $\mu=\nu$ on $\B((0,\infty))$, which is \eqref{eq:2-8-06:a}, i.e.\ \bk{(8.9)}.

\rmk{The measure $\Gamma\mapsto P^x[W_t\in\Gamma,\ T_0>t]$ on $\B(\R)$ charges only $(0,\infty)$: if
$W_0(\omega)=x>0$ and $T_0(\omega)>t$, then $W_t(\omega)>0$ by \eqref{eq:2-8-06:b} and the intermediate
value theorem. Its total mass is, by \eqref{eq:2-8-06:a},
\[
  P^x[T_0>t]=\int_0^\infty p(t;x,y)\,dy-\int_{-\infty}^0p(t;x,z)\,dz
  =P^0[W_t>-x]-P^0[W_t<-x]=1-2P^0[W_t> x],
\]
so that $P^x[T_0\le t]=2P^0[W_t\ge x]=P^0[T_x\le t]$, in agreement with \bk{(8.4)} and with the
symmetry and spatial homogeneity of Brownian motion.}

\emph{Step 3: the absorbed process.} Set $X_t\defeq W_{t\wedge T_0}$, $0\le t<\infty$; this process
is continuous and adapted to $\{\F_t\}$, because $W$ is progressively measurable and $t\wedge T_0$
is a stopping time (\bk{Proposition 1.2.18}). By \eqref{eq:2-8-06:b},
\begin{equation}\label{eq:2-8-06:g}
  X_t=W_t\neq 0 \ \text{ on } \{T_0>t\},\qquad X_t=W_{T_0}=0 \ \text{ on } \{T_0\le t\}.
\end{equation}
We show that for $x>0$, $s\ge 0$, $t>0$ and $\Gamma\in\B((0,\infty))$,
\begin{equation}\label{eq:2-8-06:h}
  P^x[X_{s+t}\in\Gamma\,|\,\F_s]=\ind_{\{X_s>0\}}\int_\Gamma p_-(t;X_s,y)\,dy,\qquad P^x\text{-a.s.}
\end{equation}
(with the convention that the right-hand side is zero on $\{X_s\le 0\}$). Taking $s=0$ and
expectations, or directly from \eqref{eq:2-8-06:a} and \eqref{eq:2-8-06:g}, we also have
$P^x[X_t\in\Gamma]=\int_\Gamma p_-(t;x,y)\,dy$.

Let $D$ be the set of rational numbers in $[0,t]$ and consider the subset
\[
  F\defeq\bigcup_{n=1}^\infty\ \bigcap_{q\in D}\bigl\{\omega\in\R^{[0,\infty)};\ \omega(q)\ge 1/n\bigr\}
  \cap\bigl\{\omega\in\R^{[0,\infty)};\ \omega(t)\in\Gamma\bigr\}
\]
of $\R^{[0,\infty)}$, which belongs to $\B(\R^{[0,\infty)})$ as a countable combination of
one-dimensional cylinder sets. A \emph{continuous} function $\omega$ belongs to $F$ if and only if
$\omega(t)\in\Gamma$ and $\omega(u)>0$ for all $u\in[0,t]$: if $\omega(q)\ge 1/n$ for all $q\in D$,
then $\omega\ge 1/n$ on $[0,t]$ by density and continuity; conversely, a continuous function which is
strictly positive on the compact interval $[0,t]$ attains there a minimum $m>0$, and
$\omega(q)\ge 1/n$ on $D$ as soon as $1/n\le m$. We claim that
\begin{equation}\label{eq:2-8-06:i}
  \{X_{s+t}\in\Gamma\}\cap\{W_0=x\}=\{T_0>s\}\cap\{W_{s+\cdot}\in F\}\cap\{W_0=x\}.
\end{equation}
Let $\omega$ satisfy $W_0(\omega)=x$. If $X_{s+t}(\omega)\in\Gamma\subseteq(0,\infty)$, then by
\eqref{eq:2-8-06:g} we have $T_0(\omega)>s+t$ and $W_{s+t}(\omega)\in\Gamma$; by
\eqref{eq:2-8-06:b} the path of $\omega$ has no zero in $[0,s+t]$ and starts at $x>0$, so it is
strictly positive on $[0,s+t]$; thus $T_0(\omega)>s$ and the continuous function $W_{s+\cdot}(\omega)$
lies in $F$. Conversely, if $T_0(\omega)>s$ and $W_{s+\cdot}(\omega)\in F$, then the path has no zero
in $[0,s]$ and is strictly positive on $[s,s+t]$, so $T_0(\omega)>s+t$ by \eqref{eq:2-8-06:b}, and
$X_{s+t}(\omega)=W_{s+t}(\omega)\in\Gamma$. This proves \eqref{eq:2-8-06:i}.

The Brownian family is a Markov family (\bk{Theorem 5.12}), so by \bk{Proposition 5.15}\,(c$'$),\,(d$'$)
\[
  P^x[W_{s+\cdot}\in F\,|\,\F_s]=h(W_s),\quad P^x\text{-a.s.},\qquad\text{where}\quad
  h(y)\defeq P^y[W_\cdot\in F];\ y\in\R .
\]
For $y>0$ the same reasoning as for \eqref{eq:2-8-06:i} (with $s=0$) shows that
$\{W_\cdot\in F\}$ and $\{W_t\in\Gamma,\ T_0>t\}$ differ by a subset of the $P^y$-null set
$\{W_0\ne y\}$, whence $h(y)=\int_\Gamma p_-(t;y,z)\,dz$ by \eqref{eq:2-8-06:a}. Since
$\{T_0>s\}\in\F_s$ and $P^x[W_0=x]=1$, \eqref{eq:2-8-06:i} gives
\[
  P^x[X_{s+t}\in\Gamma\,|\,\F_s]=\ind_{\{T_0>s\}}\,P^x[W_{s+\cdot}\in F\,|\,\F_s]
  =\ind_{\{T_0>s\}}\,h(W_s),\qquad P^x\text{-a.s.}
\]
Finally, on $\{T_0>s\}\cap\{W_0=x\}$ we have $X_s=W_s>0$, and on $\{T_0\le s\}$ we have $X_s=0$, by
\eqref{eq:2-8-06:g}; hence $\ind_{\{T_0>s\}}h(W_s)=\ind_{\{X_s>0\}}h(X_s)$, $P^x$-a.s., and
\eqref{eq:2-8-06:h} is proved.

The right-hand side of \eqref{eq:2-8-06:h} is a Borel-measurable function of $X_s$ (the map
$y\mapsto\int_\Gamma p_-(t;y,z)\,dz$ is continuous on $(0,\infty)$ by dominated convergence, locally
in $y$). As $X_{s+t}$ takes values in $[0,\infty)$ $P^x$-a.s., the conditional probability of
$\{X_{s+t}=0\}$ is one minus that of $\{X_{s+t}\in(0,\infty)\}$, so
$P^x[X_{s+t}\in\Gamma\,|\,\F_s]$ is a function of $X_s$ for every $\Gamma\in\B([0,\infty))$.
Thus, under each $P^x$, $x>0$, the absorbed Brownian motion $\{W_{t\wedge T_0},\F_t;\ 0\le t<\infty\}$ is a
Markov process with state space $[0,\infty)$, whose transition probabilities
\[
  P_t(y,\Gamma)=\int_{\Gamma\cap(0,\infty)}p_-(t;y,z)\,dz+\ind_\Gamma(0)\,\Bigl(1-\int_0^\infty p_-(t;y,z)\,dz\Bigr)
  \ \ (y>0),\qquad P_t(0,\Gamma)=\ind_\Gamma(0),
\]
have on $(0,\infty)$ the density $p_-(t;y,z)$ of \bk{(8.9)}; the remaining mass
$P^y[T_0\le t]=2P^0[W_t\ge y]$ sits at the absorbing state $0$.
\end{kex}

\begin{kex}{Problem}{8.7}{}
Here $\{W_t,\F_t;\ 0\le t<\infty\}$, $(\Omega,\F)$, $\{P^x\}_{x\in\R}$ is the one-dimensional Brownian family of Section~2.8, and $p$ is the Gaussian kernel of \bk{(2.6)}. We have to verify the two conditions of \bk{Definition 5.10} for the real-valued process $\abs{W}$ on $(\Omega,\F,P^0)$ with initial distribution $\mu=\delta_0$, the unit mass at the origin (in that definition the letters $s$ and $t$ play the roles which $t$ and $s$ play in \bk{(8.10)}), and to identify the conditional law of $\abs{W_{t+s}}$ given $\abs{W_t}$. The precise meaning of \bk{(8.10)}, in the notation of \bk{(4.2)}, is: for all $s>0$, $t\ge0$ and $\Gamma\in\B(\R)$,
\begin{equation}\label{eq:2-8-07:meaning}
  P^0\big[\abs{W_{t+s}}\in\Gamma \,\big|\, \abs{W_t}=x\big]=\int_{\Gamma\cap[0,\infty)}p_+(s;x,y)\,dy
  \quad\text{for } P^0\abs{W_t}^{-1}\text{-a.e.\ } x\in[0,\infty).
\end{equation}

\emph{Step 1: adaptedness and initial distribution.} For each $t\ge0$ the random variable $\abs{W_t}$ is $\F_t$-measurable, being a continuous function of the $\F_t$-measurable random variable $W_t$; thus $\abs{W}$ is adapted to $\{\F_t\}$. Since $P^0[W_0=0]=1$, we have $P^0[\abs{W_0}\in\Gamma]=\ind_\Gamma(0)=\delta_0(\Gamma)$ for every $\Gamma\in\B(\R)$, which is condition (i) of \bk{Definition 5.10}.

\emph{Step 2: the conditional law given $\F_t$.} Fix $s>0$, $t\ge0$ and $\Gamma\in\B(\R)$, and put
\[
  \Gamma_+\defeq\Gamma\cap[0,\infty),\qquad \Gamma'\defeq\{z\in\R;\ \abs{z}\in\Gamma\}.
\]
The set $\Gamma'$ is the inverse image of $\Gamma$ under the continuous map $z\mapsto\abs{z}$, so $\Gamma'\in\B(\R)$, and $\{\abs{W_{t+s}}\in\Gamma\}=\{W_{t+s}\in\Gamma'\}$. A Brownian family is a Markov family (\bk{Theorem 5.12}), so \bk{Proposition 5.13}\,(e) gives
\begin{equation}\label{eq:2-8-07:markov}
  P^0\big[\abs{W_{t+s}}\in\Gamma \,\big|\, \F_t\big]=P^0\big[W_{t+s}\in\Gamma' \,\big|\, \F_t\big]=h(W_t),\quad\text{a.s. } P^0,
\end{equation}
where $h(z)\defeq(U_s\ind_{\Gamma'})(z)=P^z[W_s\in\Gamma']$, $z\in\R$, in the notation of \bk{(5.7)}. Under $P^z$ the process $W$ is a Brownian motion with $P^z[W_0=z]=1$, so $W_s=z+(W_s-W_0)$ is normally distributed with mean $z$ and variance $s>0$; that is, $W_s$ has density $p(s;z,\cdot)$ under $P^z$, and
\begin{equation}\label{eq:2-8-07:h}
  h(z)=\int_{\Gamma'}p(s;z,y)\,dy,\qquad z\in\R.
\end{equation}

We now express $h$ through $p_+$. The set $\Gamma'$ is the disjoint union of $\Gamma_+$ and $-(\Gamma\cap(0,\infty))\defeq\{-y;\ y\in\Gamma\cap(0,\infty)\}$. The substitution $y\mapsto-y$, which preserves Lebesgue measure, and the fact that the single point $\{0\}$ is Lebesgue-null yield
\[
  \int_{-(\Gamma\cap(0,\infty))}p(s;z,y)\,dy=\int_{\Gamma\cap(0,\infty)}p(s;z,-y)\,dy=\int_{\Gamma_+}p(s;z,-y)\,dy,
\]
and consequently, from \eqref{eq:2-8-07:h},
\begin{equation}\label{eq:2-8-07:h2}
  h(z)=\int_{\Gamma_+}\big[p(s;z,y)+p(s;z,-y)\big]\,dy,\qquad z\in\R.
\end{equation}
From \bk{(2.6)} we read off $p(s;-z,y)=p(s;z,-y)$ and $p(s;-z,-y)=p(s;z,y)$, so the integrand in \eqref{eq:2-8-07:h2} is unchanged when $z$ is replaced by $-z$. Hence $h(-z)=h(z)$, i.e.,
\begin{equation}\label{eq:2-8-07:q}
  h(z)=q(\abs{z}),\quad z\in\R,\qquad\text{where}\quad q(x)\defeq\int_{\Gamma_+}p_+(s;x,y)\,dy,\quad x\ge0.
\end{equation}
The function $q:[0,\infty)\to[0,1]$ is Borel-measurable: the map $(x,y)\mapsto p_+(s;x,y)\ind_{\Gamma_+}(y)$ is nonnegative and Borel-measurable on $[0,\infty)\times\R$ (a continuous function times the indicator of a Borel set), and the Tonelli theorem asserts the measurability of its partial integral with respect to $y$. Substituting \eqref{eq:2-8-07:q} into \eqref{eq:2-8-07:markov} we obtain
\begin{equation}\label{eq:2-8-07:key}
  P^0\big[\abs{W_{t+s}}\in\Gamma \,\big|\, \F_t\big]=q(\abs{W_t}),\quad\text{a.s. } P^0.
\end{equation}

\emph{Step 3: the Markov property.} The random variable $q(\abs{W_t})$ is $\sigma(\abs{W_t})$-measurable, and $\sigma(\abs{W_t})\subseteq\F_t$. Conditioning both sides of \eqref{eq:2-8-07:key} on $\sigma(\abs{W_t})$ and using the tower property, we get
\begin{equation}\label{eq:2-8-07:given}
  P^0\big[\abs{W_{t+s}}\in\Gamma \,\big|\, \abs{W_t}\big]
  =E^0\big[q(\abs{W_t}) \,\big|\, \abs{W_t}\big]=q(\abs{W_t}),\quad\text{a.s. } P^0,
\end{equation}
and a comparison of \eqref{eq:2-8-07:key} with \eqref{eq:2-8-07:given} shows
\begin{equation}\label{eq:2-8-07:mp}
  P^0\big[\abs{W_{t+s}}\in\Gamma \,\big|\, \F_t\big]=P^0\big[\abs{W_{t+s}}\in\Gamma \,\big|\, \abs{W_t}\big],\quad\text{a.s. } P^0,
\end{equation}
for $s>0$, $t\ge0$, $\Gamma\in\B(\R)$. For $s=0$ the identity \eqref{eq:2-8-07:mp} holds trivially, because $\ind_{\{\abs{W_t}\in\Gamma\}}$ is measurable with respect to both $\F_t$ and $\sigma(\abs{W_t})$, so both sides are equal to it a.s. This is condition (ii) of \bk{Definition 5.10}; together with Step~1, it shows that $\abs{W}$ is a Markov process under $P^0$ (with initial distribution $\delta_0$).

\emph{Step 4: the transition density.} By definition, $x\mapsto P^0[\abs{W_{t+s}}\in\Gamma\mid\abs{W_t}=x]$ is any Borel-measurable function $g$ with $g(\abs{W_t})=P^0[\abs{W_{t+s}}\in\Gamma\mid\abs{W_t}]$, a.s.\ $P^0$, and any two such functions agree $P^0\abs{W_t}^{-1}$-a.e. According to \eqref{eq:2-8-07:given}, the Borel function $q$ of \eqref{eq:2-8-07:q} (extended by zero to $(-\infty,0)$, a $P^0\abs{W_t}^{-1}$-null set) is one of them; this is exactly \eqref{eq:2-8-07:meaning}, i.e., \bk{(8.10)}.

Finally, $p_+(s;x,\cdot)$ is indeed a probability density on $[0,\infty)$ for every $s>0$, $x\ge0$: it is nonnegative, and \eqref{eq:2-8-07:h2} with $\Gamma=\R$ (so that $\Gamma'=\R$, $\Gamma_+=[0,\infty)$) gives
\[
  \int_0^\infty p_+(s;x,y)\,dy=\int_{-\infty}^\infty p(s;x,y)\,dy=1.
\]

\rmk{For $t>0$ the law of $\abs{W_t}$ under $P^0$ has the strictly positive density $p_+(t;0,\cdot)=2p(t;0,\cdot)$ on $[0,\infty)$, so in \eqref{eq:2-8-07:meaning} ``$P^0\abs{W_t}^{-1}$-a.e.\ $x$'' means ``Lebesgue-a.e.\ $x\ge0$''; for $t=0$ this law is $\delta_0$, and \eqref{eq:2-8-07:meaning} is a statement about $x=0$ only, namely $P^0[\abs{W_s}\in\Gamma]=\int_{\Gamma_+}2p(s;0,y)\,dy$. The argument of Steps~2--4 uses nothing about the initial point: with $P^0$ replaced by $P^a$, $a\in\R$, it shows that \eqref{eq:2-8-07:key}--\eqref{eq:2-8-07:mp} and \eqref{eq:2-8-07:meaning} hold under every $P^a$, with the same function $q$.}
\end{kex}

\begin{kex}{Exercise}{8.9}{}
Here $a\in(0,\infty)$ and $\varphi:\R\to[0,a]$ is the function introduced just before the exercise: $\varphi(2na)=0$, $\varphi((2n+1)a)=a$ for $n\in\Z$, and $\varphi$ is linear between these points. Explicitly, for every $n\in\Z$,
\begin{equation}\label{eq:2-8-09:phi}
\varphi(z)=z-2na \quad\text{for } 2na\le z\le(2n+1)a,
\qquad
\varphi(z)=2na-z \quad\text{for } (2n-1)a\le z\le 2na .
\end{equation}
Recall also $p(t;x,y)=(2\pi t)^{-1/2}e^{-(x-y)^2/2t}$ from \bk{(2.6)} and $p_+(t;x,y)=p(t;x,y)+p(t;x,-y)$ from \bk{(8.10)}. The assertion to be proved is that, for $t>0$ and $x\in(0,a)$,
\begin{equation}\label{eq:2-8-09:claim}
P^x[\varphi(W_t)\in A]=\int_A\ \sum_{n=-\infty}^{\infty}p_+(t;x,y+2na)\,dy
\qquad\text{for every } A\in\B((0,a)).
\end{equation}
We prove \eqref{eq:2-8-09:claim} for every $x\in\R$; the restriction $0<x<a$ plays no role.

\emph{Step 1: the preimage of a Borel set.} Fix $A\in\B((0,a))$ and put, for $n\in\Z$,
\[
\begin{aligned}
A_n^+&\defeq A+2na=\{y+2na;\ y\in A\}\subseteq\bigl(2na,(2n+1)a\bigr),\\
A_n^-&\defeq 2na-A=\{2na-y;\ y\in A\}\subseteq\bigl((2n-1)a,2na\bigr).
\end{aligned}
\]
These are Borel sets, and they are pairwise disjoint, because the open intervals $(ka,(k+1)a)$, $k\in\Z$, that contain them are pairwise disjoint ($A_n^+$ lies in the interval with $k=2n$, $A_n^-$ in the one with $k=2n-1$). We claim that
\begin{equation}\label{eq:2-8-09:pre}
\varphi^{-1}(A)=\bigcup_{n\in\Z}\bigl(A_n^+\cup A_n^-\bigr).
\end{equation}
Indeed, let $z\in\R$. If $z\in a\Z$, then $\varphi(z)\in\{0,a\}$, so $z\notin\varphi^{-1}(A)$, and $z$ belongs to none of the sets $A_n^\pm$. Otherwise $z$ lies in exactly one interval $(ka,(k+1)a)$. If $k=2n$ is even, then $\varphi(z)=z-2na$ by \eqref{eq:2-8-09:phi}, so $\varphi(z)\in A$ if and only if $z\in A_n^+$; moreover $z$ lies in no other set of the family. If $k=2n-1$ is odd, then $\varphi(z)=2na-z$, so $\varphi(z)\in A$ if and only if $z\in A_n^-$, and again $z$ lies in no other set of the family. This proves \eqref{eq:2-8-09:pre}.

\emph{Step 2: the distribution of $\varphi(W_t)$.} Under $P^x$ the random variable $W_t$ is normal with mean $x$ and variance $t>0$, i.e.\ $P^x[W_t\in B]=\int_B p(t;x,z)\,dz$ for $B\in\B(\R)$. From \eqref{eq:2-8-09:pre}, the disjointness of the sets $A_n^\pm$ and countable additivity,
\[
P^x[\varphi(W_t)\in A]=P^x\bigl[W_t\in\varphi^{-1}(A)\bigr]
=\sum_{n\in\Z}\biggl(\int_{A_n^+}p(t;x,z)\,dz+\int_{A_n^-}p(t;x,z)\,dz\biggr).
\]
The substitutions $z=y+2na$ in the first integral and $z=2na-y$ in the second (translations and reflections preserve Lebesgue measure) turn this into
\[
\begin{aligned}
P^x[\varphi(W_t)\in A]
&=\sum_{n\in\Z}\int_A\bigl[p(t;x,y+2na)+p(t;x,2na-y)\bigr]\,dy\\
&=\int_A\ \sum_{n\in\Z}\bigl[p(t;x,y+2na)+p(t;x,2na-y)\bigr]\,dy ,
\end{aligned}
\]
where the interchange of summation and integration is justified by Tonelli's theorem, all terms being nonnegative.

\emph{Step 3: identification of the density.} Taking $A=(0,a)$ in Step 2 shows that the nonnegative function $y\mapsto\sum_{n\in\Z}[p(t;x,y+2na)+p(t;x,2na-y)]$ has integral $P^x[W_t\notin a\Z]=1$ over $(0,a)$; in fact the series converges for \emph{every} $y$, since $p(t;x,y\pm 2na)\le(2\pi t)^{-1/2}\exp\{-(2\abs{n}a-\abs{x}-\abs{y})^2/2t\}$ as soon as $2\abs{n}a\ge\abs{x}+\abs{y}$. A convergent series of nonnegative terms may be rearranged and regrouped at will; replacing $n$ by $-n$ in the terms $p(t;x,2na-y)$, we obtain
\[
\begin{aligned}
\sum_{n\in\Z}\bigl[p(t;x,y+2na)+p(t;x,2na-y)\bigr]
&=\sum_{n\in\Z}\bigl[p(t;x,y+2na)+p(t;x,-(y+2na))\bigr]\\
&=\sum_{n\in\Z}p_+(t;x,y+2na).
\end{aligned}
\]
Together with Step 2 this is \eqref{eq:2-8-09:claim}. Finally, $P^x[\varphi(W_t)\in\{0,a\}]=P^x[W_t\in a\Z]=0$, so the distribution of $\varphi(W_t)$ under $P^x$ is carried by $(0,a)$ and
\[
q(t;x,y)\defeq\sum_{n=-\infty}^{\infty}p_+(t;x,y+2na),\qquad 0<y<a,
\]
is its density there: $P^x[\varphi(W_t)\in dy]=q(t;x,y)\,dy$. \qed

\rmk{The function $x\mapsto q(t;x,y)$ is even and $2a$-periodic on $\R$: since $p(t;-x,z)=p(t;x,-z)$ and $p(t;x+2a,z)=p(t;x,z-2a)$, replacing $x$ by $-x$ merely interchanges the two families $\{p(t;x,y+2na)\}_n$ and $\{p(t;x,-y-2na)\}_n$, and replacing $x$ by $x+2a$ shifts the summation index by one. As $\varphi$ is itself even and $2a$-periodic with $\varphi(z)=z$ on $[0,a]$, every $x\in\R$ is of the form $\pm\varphi(x)+2ka$, and therefore $q(t;x,y)=q(t;\varphi(x),y)$ for all $x\in\R$. Thus the law of $\varphi(W_t)$ under $P^x$ depends on $x$ only through $\varphi(x)$, and the formula of the exercise, which we proved for all $x\in\R$, is the transition density of the doubly reflected Brownian motion: by the Markov property of the Brownian family $\{W_t,\F_t\}$, $\{P^x\}$ (\bk{Section 2.5}), $P^x[\varphi(W_{t+s})\in A\mid\F_s]=\int_A q(t;\varphi(W_s),y)\,dy$, $P^x$-a.s.}
\end{kex}

\begin{kex}{Exercise}{8.11}{}
We write $T\defeq T_0\wedge T_a$, and for $t>0$, $c\in\R$,
\[
  \psi(t,c)\defeq\frac{c}{\sqrt{2\pi t^3}}\,\exp\Bigl\{-\frac{c^2}{2t}\Bigr\},
  \qquad
  S(t,b)\defeq\sum_{n=-\infty}^{\infty}\psi(t,2na+b),\quad b\in\R ,
\]
so that the right-hand side of \bk{(8.24)} is $f(t)\,dt$ with $f(t)\defeq S(t,x)+S(t,a-x)$. The assertion to be proved is that
\begin{equation}\label{eq:2-8-11:claim}
  P^x[T\in B]=\int_B f(t)\,dt\qquad\text{for every Borel set } B\subseteq(0,\infty).
\end{equation}

\emph{Step 1. The series $S$.} Fix $b\in\R$ and $0<t_1<t_2<\infty$. For $t\in[t_1,t_2]$ and every integer $n$ with $\abs{n}\ge\abs{b}/a$ we have $\abs{n}a\le\abs{2na+b}\le 3\abs{n}a$, hence
\begin{equation}\label{eq:2-8-11:bound}
  \sup_{t_1\le t\le t_2}\abs{\psi(t,2na+b)}\le\frac{3\abs{n}a}{\sqrt{2\pi t_1^3}}\,\exp\Bigl\{-\frac{n^2a^2}{2t_2}\Bigr\},
  \qquad\text{and so}\qquad
  \sum_{n=-\infty}^{\infty}\,\sup_{t_1\le t\le t_2}\abs{\psi(t,2na+b)}<\infty .
\end{equation}
Thus the series defining $S(t,b)$ converges absolutely for every $t>0$ and uniformly on $[t_1,t_2]$; since each term is continuous in $t$, the function $t\mapsto S(t,b)$ is continuous on $(0,\infty)$, and so is $f$. Absolute convergence permits us to re-index the series: replacing $n$ by $n+1$, and replacing $n$ by $-n$ and using $\psi(t,-c)=-\psi(t,c)$, we obtain
\begin{equation}\label{eq:2-8-11:sym}
  S(t,b+2a)=S(t,b),\qquad S(t,-b)=-S(t,b);\qquad t>0,\ b\in\R .
\end{equation}

\emph{Step 2. The cases $x=0$ and $x=a$.} If $x\in\{0,a\}$, then $P^x[T=0]=1$, so the left-hand side of \eqref{eq:2-8-11:claim} is zero for every Borel $B\subseteq(0,\infty)$. On the other hand, \eqref{eq:2-8-11:sym} gives $S(t,0)=-S(t,0)$ and $S(t,a)=-S(t,-a)=-S(t,a)$, so $S(t,0)=S(t,a)=0$ and $f(t)=S(t,0)+S(t,a)=0$ for every $t>0$, both when $x=0$ and when $x=a$. Hence \eqref{eq:2-8-11:claim} holds in these two cases.

\emph{Step 3. The survival probability for $0<x<a$.} From now on $0<x<a$. Let $\Phi$ be the standard normal distribution function and recall $p(t;x,y)=(2\pi t)^{-1/2}\exp\{-(y-x)^2/2t\}$ and $p_-(t;x,y)=p(t;x,y)-p(t;x,-y)$ from \bk{(8.9)}. Fix $t>0$. On $\{T>t\}$ the continuous path $s\mapsto W_s$ starts at $x\in(0,a)$ and takes neither of the values $0$, $a$ on $[0,t]$, so by the intermediate value theorem $W_t\in(0,a)$. Therefore \bk{Proposition 8.10} gives
\[
  P^x[T>t]=P^x[W_t\in(0,a),\,T>t]=\int_0^a\sum_{n=-\infty}^{\infty}\bigl[p(t;x,y+2na)-p(t;x,-y-2na)\bigr]\,dy .
\]
The nonnegative functions $y\mapsto p(t;x,y+2na)$, $n\in\Z$, satisfy $\sum_n\int_0^a p(t;x,y+2na)\,dy=\sum_n\int_{2na}^{2na+a}p(t;x,z)\,dz\le\int_{\R}p(t;x,z)\,dz=1$, and likewise $\sum_n\int_0^a p(t;x,-y-2na)\,dy\le 1$. By Fubini--Tonelli we may thus integrate term by term, and treat the two families separately:
\begin{equation}\label{eq:2-8-11:surv}
  P^x[T>t]=\sum_{n=-\infty}^{\infty}A_n(t)-\sum_{n=-\infty}^{\infty}B_n(t),
\end{equation}
where both series have nonnegative terms and sums at most $1$, and
\begin{align*}
  A_n(t)&\defeq\int_{2na}^{2na+a}p(t;x,z)\,dz=\Phi\Bigl(\frac{2na+a-x}{\sqrt t}\Bigr)-\Phi\Bigl(\frac{2na-x}{\sqrt t}\Bigr),\\
  B_n(t)&\defeq\int_{-2na-a}^{-2na}p(t;x,z)\,dz=\Phi\Bigl(\frac{-2na-x}{\sqrt t}\Bigr)-\Phi\Bigl(\frac{-2na-a-x}{\sqrt t}\Bigr)\\
  &=\Phi\Bigl(\frac{2na+a+x}{\sqrt t}\Bigr)-\Phi\Bigl(\frac{2na+x}{\sqrt t}\Bigr);
\end{align*}
the last equality uses $\Phi(-u)=1-\Phi(u)$. (In $B_n$ we substituted $z=-y-2na$.)

\emph{Step 4. Integration in time.} For fixed $c\in\R$ the function $t\mapsto\Phi(c/\sqrt t)$ is continuously differentiable on $(0,\infty)$ with
\[
  \frac{d}{dt}\,\Phi\Bigl(\frac{c}{\sqrt t}\Bigr)=-\frac{c}{2t^{3/2}}\cdot\frac{1}{\sqrt{2\pi}}\exp\Bigl\{-\frac{c^2}{2t}\Bigr\}=-\frac12\,\psi(t,c).
\]
Hence, for $0<t_1<t_2<\infty$ and every $n\in\Z$,
\begin{align*}
  A_n(t_1)-A_n(t_2)&=\frac12\int_{t_1}^{t_2}\bigl[\psi(s,2na+a-x)-\psi(s,2na-x)\bigr]\,ds,\\
  B_n(t_1)-B_n(t_2)&=\frac12\int_{t_1}^{t_2}\bigl[\psi(s,2na+a+x)-\psi(s,2na+x)\bigr]\,ds .
\end{align*}
We subtract \eqref{eq:2-8-11:surv} at $t_2$ from \eqref{eq:2-8-11:surv} at $t_1$ (all four series converge) and sum these identities over $n$. By \eqref{eq:2-8-11:bound}, applied with $b=a-x,\,-x,\,a+x,\,x$, we have $\sum_n\int_{t_1}^{t_2}\abs{\psi(s,2na+b)}\,ds<\infty$ for each of these four values of $b$, so Fubini's theorem allows us to interchange summation and integration and to sum each of the four families separately. This yields
\[
  P^x[t_1<T\le t_2]=P^x[T>t_1]-P^x[T>t_2]
  =\frac12\int_{t_1}^{t_2}\bigl[S(s,a-x)-S(s,-x)-S(s,a+x)+S(s,x)\bigr]\,ds .
\]
By \eqref{eq:2-8-11:sym}, $-S(s,-x)=S(s,x)$ and $-S(s,a+x)=S(s,-a-x)=S(s,a-x)$. The integrand is therefore $2S(s,x)+2S(s,a-x)=2f(s)$, and we have proved
\begin{equation}\label{eq:2-8-11:int}
  P^x[t_1<T\le t_2]=\int_{t_1}^{t_2}f(s)\,ds;\qquad 0<t_1<t_2<\infty .
\end{equation}

\emph{Step 5. Conclusion.} The left-hand side of \eqref{eq:2-8-11:int} is nonnegative for all $0<t_1<t_2$, and $f$ is continuous on $(0,\infty)$ by Step 1; consequently $f\ge0$ on $(0,\infty)$ (if $f(t_0)<0$, then $f<0$ on a neighbourhood of $t_0$, and the integral of $f$ over it would be negative). Thus $B\mapsto\int_Bf(t)\,dt$ and $B\mapsto P^x[T\in B]$ are two measures on $\bigl((0,\infty),\B((0,\infty))\bigr)$ which agree, by \eqref{eq:2-8-11:int}, on the $\pi$-system $\{(t_1,t_2];\,0<t_1<t_2<\infty\}\cup\{\emptyset\}$; this $\pi$-system generates $\B((0,\infty))$, and both measures are finite on each of the sets $(1/k,k]$, $k\ge2$, which belong to it and increase to $(0,\infty)$. By Dynkin's $\pi$-$\lambda$ theorem (uniqueness of $\sigma$-finite measures agreeing on a generating $\pi$-system), the two measures coincide. This is \eqref{eq:2-8-11:claim}, i.e.\ \bk{(8.24)}, for $0<x<a$; together with Step 2 the exercise is solved.

\rmk{For $0<x<a$ the function $f$ is a genuine probability density on $(0,\infty)$: $P^x[T>0]=1$ by continuity of the paths, and $P^x[T>t]\le P^x[0<W_t<a]\le a/\sqrt{2\pi t}\to0$ as $t\to\infty$, so that $\int_0^\infty f(t)\,dt=P^x[0<T<\infty]=1$. For $x\in\{0,a\}$ the law of $T$ is the unit mass at $t=0$, and \bk{(8.24)} holds in the sense that this law has no mass on $(0,\infty)$, where the series vanishes identically. A numerical comparison of the series in \bk{(8.24)} with the eigenfunction expansion $\sum_{k\ \mathrm{odd}}(2k\pi/a^2)\sin(k\pi x/a)\exp\{-k^2\pi^2t/2a^2\}$ of the same density, and of its integral with $1$, shows agreement to machine precision.}
\end{kex}

\begin{kex}{Exercise}{8.13}{}
Throughout, $\{W_t,\F_t;\ 0\le t<\infty\}$, $(\Omega,\F)$, $\{P^x\}_{x\in\R}$ is the one-dimensional Brownian family of Section 2.8, $T_b=\inf\{t\ge 0;\ W_t=b\}$ as in \bk{(6.1)}, and we fix $a>0$, $x\in[0,a]$ and put
\[
\tau\defeq T_0\wedge T_a=\inf\{t\ge 0;\ W_t\in\{0,a\}\}.
\]
The paths of $W$ are continuous and $\{0\}$, $\{a\}$ are closed, so $T_0$, $T_a$ are stopping times of $\{\F_t\}$ by \bk{Problem 1.2.7}, and so is $\tau$ by \bk{Lemma 1.2.9}. Moreover, for $b\in\R$ the set $\{t\ge 0;\ W_t(\omega)=b\}$ is closed, hence contains its infimum when it is nonempty:
\begin{equation}\label{eq:2-8-13:hit}
W_{T_b}=b\quad\text{on }\{T_b<\infty\}.
\end{equation}

\emph{Step 1: $P^x[\tau<\infty]=1$.} Let $t>0$ and $\omega\in\{W_0=x\}\cap\{\tau>t\}$. Then $W_s(\omega)\notin\{0,a\}$ for all $s\in[0,t]$; in particular $x=W_0(\omega)\in(0,a)$, and by the intermediate value theorem the continuous function $s\mapsto W_s(\omega)$ cannot leave $(0,a)$ on $[0,t]$ without taking one of the values $0$, $a$. Hence $0<W_t(\omega)<a$. Since $P^x[W_0=x]=1$ and $W_t$ is normal with mean $x$ and variance $t$ under $P^x$,
\[
P^x[\tau>t]\le P^x[0<W_t<a]=\int_0^a\frac{1}{\sqrt{2\pi t}}\,e^{-(y-x)^2/2t}\,dy\le\frac{a}{\sqrt{2\pi t}} .
\]
Letting $t\to\infty$ we obtain $P^x[\tau=\infty]=\lim_{t\to\infty}P^x[\tau>t]=0$.

\emph{Step 2: the two events partition $\{\tau<\infty\}$.} On $\{\tau<\infty\}$ at least one of $T_0$, $T_a$ is finite, and they cannot be equal: if $T_0=T_a<\infty$, then \eqref{eq:2-8-13:hit} would give $0=W_{T_0}=W_{T_a}=a$. Conversely, each of the events $\{T_0<T_a\}$, $\{T_a<T_0\}$ is contained in $\{\tau<\infty\}$. Thus $\{\tau<\infty\}$ is the disjoint union of $\{T_0<T_a\}$ and $\{T_a<T_0\}$, and by Step 1
\begin{equation}\label{eq:2-8-13:sum}
P^x[T_0<T_a]+P^x[T_a<T_0]=1 .
\end{equation}
Furthermore, on $\{T_a<T_0\}$ we have $\tau=T_a<\infty$, and on $\{T_0<T_a\}$ we have $\tau=T_0<\infty$, so \eqref{eq:2-8-13:hit} yields
\begin{equation}\label{eq:2-8-13:Wtau}
W_\tau=a\,\ind_{\{T_a<T_0\}}\quad\text{on }\{\tau<\infty\}.
\end{equation}

\emph{Step 3: optional sampling.} Under $P^x$ the process $W$ is a Brownian motion starting at $x$ (\bk{Definition 5.8}), so by \bk{Problem 5.5} the process $\{W_t-W_0,\F_t;\ 0\le t<\infty\}$ is a continuous martingale on $(\Omega,\F,P^x)$; as $W_0=x$ holds $P^x$-a.s., the process $\{W_t,\F_t;\ 0\le t<\infty\}$ is itself a continuous martingale under $P^x$, with $E^xW_0=x$. Applying \bk{Problem 1.3.24}\,(i) to the submartingales $W$ and $-W$ and to the stopping time $\tau$, we see that $\{W_{\tau\wedge t},\F_t;\ 0\le t<\infty\}$ is a martingale under $P^x$; in particular
\begin{equation}\label{eq:2-8-13:os}
E^x\bigl[W_{\tau\wedge t}\bigr]=E^x\bigl[W_{\tau\wedge 0}\bigr]=E^x[W_0]=x;\qquad 0\le t<\infty .
\end{equation}
We now let $t\to\infty$ along the integers. On $\{\tau<\infty\}$ we have $W_{\tau\wedge t}=W_\tau$ as soon as $t\ge\tau$, so $W_{\tau\wedge t}\to W_\tau$, $P^x$-a.s.\ by Step 1. Moreover $\abs{W_{\tau\wedge t}}\le a$ on $\{W_0=x\}$: this is clear if $\tau=0$, since then $W_{\tau\wedge t}=W_0=x\in[0,a]$; and if $\tau>0$, then $x\in(0,a)$ and, as in Step 1, $W_s\in(0,a)$ for all $s<\tau$, whence $W_{\tau\wedge t}\in[0,a]$ by continuity. The bounded convergence theorem, \eqref{eq:2-8-13:os} and \eqref{eq:2-8-13:Wtau} give
\[
x=\lim_{t\to\infty}E^x\bigl[W_{\tau\wedge t}\bigr]=E^x\bigl[W_\tau\ind_{\{\tau<\infty\}}\bigr]=a\,P^x[T_a<T_0] .
\]
Therefore $P^x[T_a<T_0]=x/a$, and \eqref{eq:2-8-13:sum} gives $P^x[T_0<T_a]=1-x/a=(a-x)/a$.

\rmk{The result agrees with the trivial boundary cases ($P^0[T_0=0]=1$ and $T_a>0$ when $W_0=0$, so $P^0[T_0<T_a]=1$; similarly $P^a[T_a<T_0]=1$), and with the transforms \bk{(8.27)}, \bk{(8.28)}: letting $\alpha\downarrow 0$ there, monotone convergence turns the left-hand sides into $P^x[T_0<T_a]$ and $P^x[T_a<T_0]$, while $\sinh((a-x)\sqrt{2\alpha})/\sinh(a\sqrt{2\alpha})\to(a-x)/a$ and $\sinh(x\sqrt{2\alpha})/\sinh(a\sqrt{2\alpha})\to x/a$. The argument above does not use these formulas.}
\end{kex}

\begin{kex}{Problem}{8.14}{}
As in \bk{Exercise 8.13}, $a>0$ is fixed, $\{W_t,\F_t;\ 0\le t<\infty\}$, $(\Omega,\F)$, $\{P^x\}_{x\in\R}$ is a one-dimensional Brownian family (\bk{Definition 5.8}), and $T_b=\inf\{t\ge0;\ W_t=b\}$ is the passage time of \bk{(6.1)}, with $\inf\emptyset=\infty$. We fix $x\in[0,a]$ and write
\[
T\defeq T_0\wedge T_a=\inf\{t\ge 0;\ W_t\in\{0,a\}\}.
\]
Since $W$ is continuous and adapted and $\{0,a\}$ is closed, $T$ is a stopping time of $\{\F_t\}$ (\bk{Problem 1.2.7}). We do not assume $P^x[T<\infty]=1$; this will come out of the argument.

\emph{Step 1: two martingales.} Under $P^x$, the process $W$ is a Brownian motion starting at $x$: we have $P^x[W_0=x]=1$, and for $0\le s<t$ the increment $W_t-W_s$ is independent of $\F_s$ and normal with mean zero and variance $t-s$. Put
\[
M_t\defeq W_t-x,\qquad N_t\defeq (W_t-x)^2-t;\qquad 0\le t<\infty .
\]
Both processes are continuous and adapted, and $M_t=(W_t-W_0)+(W_0-x)$ is, $P^x$-a.s., equal to $W_t-W_0$, which is normal with mean zero and variance $t$; thus $M_t$ and $N_t$ are $P^x$-integrable. For $0\le s<t$, the random variable $M_s$ is $\F_s$-measurable and square-integrable, $W_t-W_s$ is independent of $\F_s$, and therefore, $P^x$-a.s.,
\begin{align*}
E^x[M_t-M_s\mid\F_s]&=E^x[W_t-W_s]=0,\\
E^x[N_t-N_s\mid\F_s]&=E^x\bigl[(W_t-W_s)^2\bigm|\F_s\bigr]+2M_s\,E^x[W_t-W_s\mid\F_s]-(t-s)\\
&=(t-s)+0-(t-s)=0,
\end{align*}
where we used $M_t^2-M_s^2=(M_t-M_s)^2+2M_s(M_t-M_s)$ and the integrability of $M_s(W_t-W_s)$ (a product of independent integrable random variables). Hence $M$ and $N$ are continuous martingales with respect to $\{\F_t\}$ under $P^x$, and $M_0=N_0=0$, $P^x$-a.s. (This is also the content of \bk{Problem 5.5}.)

\emph{Step 2: optional sampling at $T\wedge t$.} Fix $t\in[0,\infty)$. The stopping time $T\wedge t$ is bounded by $t$, so the optional sampling theorem for bounded stopping times (\bk{Problem 1.3.23}(i), applied to the continuous martingales $\pm M$, $\pm N$ and the stopping times $0\le T\wedge t$) gives $E^xM_{T\wedge t}=E^xM_0=0$ and $E^xN_{T\wedge t}=E^xN_0=0$, that is,
\begin{equation}\label{eq:2-8-14:os}
E^x\bigl[W_{T\wedge t}\bigr]=x,\qquad
E^x\bigl[T\wedge t\bigr]=E^x\bigl[(W_{T\wedge t}-x)^2\bigr];\qquad 0\le t<\infty .
\end{equation}

\emph{Step 3: the stopped process stays in $[0,a]$.} Let $\Omega_0\defeq\{W_0=x\}$, so that $P^x(\Omega_0)=1$. We claim that on $\Omega_0$
\begin{equation}\label{eq:2-8-14:bd}
0\le W_s\le a\quad\text{for every } s\in[0,T]\cap[0,\infty).
\end{equation}
Indeed, let $\omega\in\Omega_0$ and suppose that $W_s(\omega)>a$ for some finite $s\le T(\omega)$. Since $W_0(\omega)=x\le a$, the intermediate value theorem applied to the continuous function $u\mapsto W_u(\omega)$ on $[0,s]$ yields $u\in[0,s)$ with $W_u(\omega)=a$, whence $T(\omega)\le u<s$, a contradiction. The case $W_s(\omega)<0$ is excluded in the same way. In particular,
\begin{equation}\label{eq:2-8-14:bd2}
0\le W_{T\wedge t}\le a\quad\text{and}\quad (W_{T\wedge t}-x)^2\le a^2,\qquad P^x\text{-a.s., for every } t\ge0 .
\end{equation}

\emph{Step 4: $T$ is integrable.} From \eqref{eq:2-8-14:os} and \eqref{eq:2-8-14:bd2}, $E^x[T\wedge t]\le a^2$ for every $t\ge 0$. Letting $t\uparrow\infty$, the monotone convergence theorem gives
\[
E^x[T]\le a^2<\infty ,\qquad\text{and in particular}\qquad P^x[T<\infty]=1 .
\]
On $\{T<\infty\}$ the infimum defining $T$ is attained, because $\{u\ge 0;\ W_u\in\{0,a\}\}$ is a closed subset of $[0,\infty)$ by the continuity of the paths; hence $W_T\in\{0,a\}$ on $\{T<\infty\}$. Moreover, on $\{T<\infty\}$ we have $W_{T\wedge t}=W_T$ for all $t\ge T$, so $W_{T\wedge t}\to W_T$ as $t\to\infty$, $P^x$-a.s.

\emph{Step 5: passage to the limit.} Let $p\defeq P^x[W_T=a,\ T<\infty]$, so that $P^x[W_T=0,\ T<\infty]=1-p$ by Step~4. Because of the bounds \eqref{eq:2-8-14:bd2}, the bounded convergence theorem applies to both right-hand quantities in \eqref{eq:2-8-14:os}, and monotone convergence applies to $E^x[T\wedge t]$. Letting $t\to\infty$ in the first relation of \eqref{eq:2-8-14:os} we obtain
\[
x=E^x[W_T]=a\,p,\qquad\text{i.e.}\qquad p=\frac{x}{a},\quad 1-p=\frac{a-x}{a},
\]
which are the formulas of \bk{Exercise 8.13} (note that $\{W_T=a\}=\{T_a<T_0\}$ and $\{W_T=0\}=\{T_0<T_a\}$ on $\{T<\infty\}$, since $T_0\neq T_a$ there). Letting $t\to\infty$ in the second relation of \eqref{eq:2-8-14:os},
\begin{align*}
E^x[T]&=E^x\bigl[(W_T-x)^2\bigr]=(a-x)^2\,p+x^2\,(1-p)\\
&=\frac{x(a-x)^2+x^2(a-x)}{a}=\frac{x(a-x)\bigl[(a-x)+x\bigr]}{a}=x(a-x).
\end{align*}
This proves $E^x(T_0\wedge T_a)=x(a-x)$ for $0\le x\le a$. \qed

\rmk{The result agrees with \bk{(8.29)}. For $\alpha>0$ we have $0\le(1-e^{-\alpha T})/\alpha\uparrow T$ as $\alpha\downarrow0$, so by monotone convergence $E^x[T]=\lim_{\alpha\downarrow0}\alpha^{-1}\bigl(1-E^x[e^{-\alpha T}]\bigr)$. With $\cosh z=1+z^2/2+O(z^4)$, the right-hand side of \bk{(8.29)} equals
\[
\frac{1+\alpha\,(x-\tfrac a2)^2+O(\alpha^2)}{1+\alpha\,\tfrac{a^2}{4}+O(\alpha^2)}
=1-\alpha\Bigl[\frac{a^2}{4}-\Bigl(x-\frac a2\Bigr)^2\Bigr]+O(\alpha^2)
=1-\alpha\,x(a-x)+O(\alpha^2),
\]
which again gives $E^x[T]=x(a-x)$. In the boundary cases $x=0$ and $x=a$ the formula reads $E^x[T]=0$, as it must, because then $T=0$ on $\{W_0=x\}$.}
\end{kex}

\begin{kex}{Problem}{8.17}{}
Fix $t>0$. Here $M_t=\max_{0\le u\le t}W_u$ as in \bk{(8.1)} and $\theta_t=\sup\{0\le s\le t;\ W_s=M_t\}$ as in \bk{(8.30)}. Put
\[
D\defeq\{(a,b,s)\in\R^3;\ b\ge a^+,\ 0<s<t\}
\]
and
\[
f(a,b,s)\defeq\frac{b(b-a)}{\pi\sqrt{s^3(t-s)^3}}\exp\Bigl\{-\frac{b^2}{2s}-\frac{(b-a)^2}{2(t-s)}\Bigr\},\quad (a,b,s)\in D .
\]
\bk{Proposition 8.15} asserts that $f$ is the density of $(W_t,M_t,\theta_t)$ on $D$, i.e.
\begin{equation}\label{eq:2-8-17:tri}
P^0\bigl[(W_t,M_t,\theta_t)\in\Gamma\bigr]=\int_\Gamma f(a,b,s)\,da\,db\,ds\qquad\text{for every Borel set }\Gamma\subseteq D .
\end{equation}
All integrands below are nonnegative and Borel measurable, so every iterated integral may be computed in any order (Tonelli's theorem). We set
\[
g(b,s)\defeq\frac{b}{\pi\sqrt{s^3(t-s)}}\,e^{-b^2/2s},\qquad
h(s)\defeq\frac{1}{\pi\sqrt{s(t-s)}};\qquad b\ge0,\ 0<s<t.
\]

\emph{Step 1: the joint density of $(M_t,\theta_t)$.} Let $B\in\B([0,\infty))$ and $C\in\B((0,t))$. Since $W_t\le M_t$ everywhere, and since ``$b\ge0$ and $a\le b$'' is the same as $b\ge a^+$, we have
\[
\{M_t\in B,\ \theta_t\in C\}=\{(W_t,M_t,\theta_t)\in\Gamma\},\qquad
\Gamma\defeq\{(a,b,s);\ b\in B,\ s\in C,\ a\le b\}\subseteq D,
\]
and $\Gamma$ is a Borel set. For fixed $b\ge0$ and $s\in(0,t)$, the substitution $x=b-a$ gives
\[
\int_{-\infty}^{b}(b-a)\exp\Bigl\{-\frac{(b-a)^2}{2(t-s)}\Bigr\}\,da
=\int_0^\infty x\,e^{-x^2/2(t-s)}\,dx
=\Bigl[-(t-s)\,e^{-x^2/2(t-s)}\Bigr]_{x=0}^{x=\infty}=t-s ,
\]
so that $\int_{-\infty}^b f(a,b,s)\,da=g(b,s)$. Therefore \eqref{eq:2-8-17:tri} yields
\begin{equation}\label{eq:2-8-17:joint}
P^0[M_t\in B,\ \theta_t\in C]=\int_C\int_B\int_{-\infty}^b f(a,b,s)\,da\,db\,ds=\int_C\int_B g(b,s)\,db\,ds
\end{equation}
for all $B\in\B([0,\infty))$, $C\in\B((0,t))$. This is the first formula: $P^0[M_t\in db,\ \theta_t\in ds]=g(b,s)\,db\,ds$ for $b\ge0$, $0<s<t$.

\emph{Step 2: the density of $\theta_t$.} For $0<s<t$,
\begin{equation}\label{eq:2-8-17:ray}
\int_0^\infty b\,e^{-b^2/2s}\,db=\Bigl[-s\,e^{-b^2/2s}\Bigr]_{b=0}^{b=\infty}=s,
\qquad\text{hence}\qquad
\int_0^\infty g(b,s)\,db=\frac{s}{\pi\sqrt{s^3(t-s)}}=h(s).
\end{equation}
Because $M_t\ge W_0=0$, the choice $B=[0,\infty)$ in \eqref{eq:2-8-17:joint} gives
\begin{equation}\label{eq:2-8-17:marg}
P^0[\theta_t\in C]=\int_C h(s)\,ds\qquad\text{for every } C\in\B((0,t)).
\end{equation}

\emph{Step 3: the distribution function.} Let $0<s\le t$ and $\varphi_s\defeq\arcsin\sqrt{s/t}\in(0,\pi/2]$. The map $\varphi\mapsto u=t\sin^2\varphi$ is a continuously differentiable, strictly increasing bijection of $(0,\varphi_s)$ onto $(0,s)$, with $du=2t\sin\varphi\cos\varphi\,d\varphi$ and $\sqrt{u(t-u)}=t\sin\varphi\cos\varphi$ (both $\sin\varphi$ and $\cos\varphi$ are positive on $(0,\pi/2)$). Hence
\begin{equation}\label{eq:2-8-17:int}
\int_0^s h(u)\,du=\int_0^{\varphi_s}\frac{2t\sin\varphi\cos\varphi}{\pi\,t\sin\varphi\cos\varphi}\,d\varphi
=\frac{2}{\pi}\arcsin\sqrt{\frac{s}{t}},\qquad 0<s\le t .
\end{equation}
For $s=t$ the right-hand side is $1$; by \eqref{eq:2-8-17:marg} with $C=(0,t)$ this means $P^0[0<\theta_t<t]=1$, and since $0\le\theta_t\le t$ by definition,
\begin{equation}\label{eq:2-8-17:ends}
P^0[\theta_t=0]=P^0[\theta_t=t]=0 .
\end{equation}
(In particular $P^0[(W_t,M_t,\theta_t)\in D]=1$, so \eqref{eq:2-8-17:tri} describes the whole distribution of the triple; and by \eqref{eq:2-8-17:marg}, \eqref{eq:2-8-17:ends} the function $h$, extended by zero outside $(0,t)$, is the density of $\theta_t$ on $\R$.) Now for $0<s<t$ we obtain from \eqref{eq:2-8-17:ends}, \eqref{eq:2-8-17:marg} with $C=(0,s]$, and \eqref{eq:2-8-17:int}:
\[
P^0[\theta_t\le s]=P^0[0<\theta_t\le s]=\int_0^s h(u)\,du=\frac{2}{\pi}\arcsin\sqrt{\frac{s}{t}} .
\]
For $s=0$ both sides vanish by \eqref{eq:2-8-17:ends}, and for $s=t$ both sides equal $1$. This proves the arc-sine law for all $0\le s\le t$.

\emph{Step 4: the conditional distribution of $M_t$ given $\theta_t$.} For $s>0$ and $B\in\B([0,\infty))$ let
\[
k(b;s)\defeq\frac{b}{s}\,e^{-b^2/2s},\quad b\ge0;\qquad Q(s,B)\defeq\int_B k(b;s)\,db .
\]
By \eqref{eq:2-8-17:ray}, $Q(s,\cdot)$ is a probability measure on $[0,\infty)$ for every $s>0$, and $s\mapsto Q(s,B)$ is Borel measurable for every $B$ by Tonelli's theorem, $k$ being jointly continuous. Moreover $g(b,s)=k(b;s)\,h(s)$ for $b\ge0$, $0<s<t$, so \eqref{eq:2-8-17:joint} and \eqref{eq:2-8-17:marg} give, for $C\in\B((0,t))$,
\[
P^0[M_t\in B,\ \theta_t\in C]=\int_C Q(s,B)\,h(s)\,ds=E^0\bigl[Q(\theta_t,B)\,\ind_{\{\theta_t\in C\}}\bigr].
\]
By \eqref{eq:2-8-17:ends} the same identity holds for every $C\in\B([0,t])$ (define $Q(0,\cdot)$ to be any fixed probability measure; the points $0$ and $t$ carry no mass). Since the events $\{\theta_t\in C\}$, $C\in\B([0,t])$, constitute $\sigma(\theta_t)$ and $Q(\theta_t,B)$ is $\sigma(\theta_t)$-measurable, this says precisely that
\[
P^0[M_t\in B\mid\theta_t]=Q(\theta_t,B)\quad\text{a.s.}\ P^0,\qquad B\in\B([0,\infty)),
\]
i.e.\ $Q$ is a regular conditional distribution of $M_t$ given $\theta_t$:
\[
P^0[M_t\in db\mid\theta_t=s]=\frac{b}{s}\,e^{-b^2/2s}\,db,\qquad b\ge0,\ 0<s<t .
\]
The kernel $Q$ was defined without reference to $t$; thus the conditional density of $M_t$ given $\theta_t=s$ is the same for all $t>s$. \qed

\rmk{As a check, the conditional law found in Step 4 has second moment $\int_0^\infty (b^3/s)e^{-b^2/2s}\,db=2s$, so that $E^0[M_t^2]=2E^0[\theta_t]=2\int_0^t s\,h(s)\,ds=t$ (the arc-sine law is symmetric about $t/2$), in agreement with the fact that $M_t$ has the same distribution as $\abs{W_t}$ under $P^0$, by \bk{(8.3)}.}
\end{kex}

\begin{kex}{Problem}{8.18}{}
Fix $t>0$. For a continuous function $x\colon[0,\infty)\to\R$ we set
\[
g(x)\defeq\sup\{0\le u\le t;\ x(u)=0\},\qquad\text{with the convention }\sup\emptyset\defeq0 ,
\]
so that $\gamma_t=g(W)$. (The convention matters only on the $P^0$-null event $\{W_0\neq0\}$.) The set $\{0\le u\le t;\ x(u)=0\}$ is compact; if it is nonempty, it contains its supremum $g(x)$. Consequently, for $0<s\le t$,
\begin{equation}\label{eq:2-8-18:a}
g(x)<s\quad\Longleftrightarrow\quad x(u)\neq0\ \text{ for every } u\in[s,t].
\end{equation}
Indeed, if $x$ has a zero in $[s,t]$ then $g(x)\ge s$; if it has none, then either $x$ has no zero in $[0,t]$ and $g(x)=0<s$, or $g(x)$ is a zero of $x$ in $[0,t]$, hence lies in $[0,s)$.

\emph{Step 1: $\theta_t$ is the last zero of $Y=M-W$ before $t$.} Let $M_u=\max_{0\le r\le u}W_r$, $Y_u=M_u-W_u\ge0$ and $\theta_t=\sup\{0\le u\le t;\ W_u=M_t\}$, as in \bk{(8.1)}, \bk{Problem 8.8} and \bk{(8.30)}. We claim that, for every $\omega$,
\begin{equation}\label{eq:2-8-18:b}
\theta_t=g(Y).
\end{equation}
The continuous function $W$ attains its maximum over $[0,t]$, so the compact set $\{0\le u\le t;\ W_u=M_t\}$ is nonempty and contains $\theta_t$: $W_{\theta_t}=M_t$. Since $M_{\theta_t}\le M_t=W_{\theta_t}\le M_{\theta_t}$, we have $Y_{\theta_t}=0$, whence $\theta_t\le g(Y)$. Conversely, put $\rho\defeq g(Y)$; then $Y_\rho=0$, i.e.\ $W_\rho=M_\rho$, and $W_u<M_u$ for $\rho<u\le t$. If we had $M_t>M_\rho$, the maximum of $W$ over $[0,t]$ would be attained at some $u\in(\rho,t]$, and then $W_u=M_t\ge M_u$, i.e.\ $Y_u=0$, which is impossible. Thus $W_\rho=M_\rho=M_t$, and so $\rho\le\theta_t$. This proves \eqref{eq:2-8-18:b}.

\emph{Step 2: $g(Y)$ and $g(\abs{W})$ have the same distribution under $P^0$.} Fix $0<s\le t$, and let $(u_m)_{m\ge1}$ be a sequence, repetitions being allowed, whose range is the countable dense subset $(\Q\cap[s,t])\cup\{s,t\}$ of $[s,t]$ (for $s=t$ this is the constant sequence $u_m=t$). Let $X$ be a process with continuous, nonnegative paths (we shall take $X=Y$ and $X=\abs{W}$). A continuous function has no zero on the compact interval $[s,t]$ and is nonnegative there if and only if it is bounded below on $[s,t]$ by $1/n$ for some $n\ge1$, and by continuity it suffices to test this bound at the points $u_m$. Hence \eqref{eq:2-8-18:a} gives
\[
\{g(X)<s\}=\bigcup_{n=1}^\infty\,\bigcap_{m=1}^\infty A_{n,m}(X),\qquad
A_{n,m}(X)\defeq\Bigl\{X_{u_1}\ge\frac1n,\dots,X_{u_m}\ge\frac1n\Bigr\}.
\]
In particular $\{g(X)<s\}$ is an event for $0<s\le t$; as $\{g(X)<s\}$ is empty for $s\le0$ and is the whole space for $s>t$, $g(X)$ is a random variable. The events $A_{n,m}(X)$ decrease in $m$, and $\bigcap_m A_{n,m}(X)$ increases in $n$; by the continuity of probability measures along monotone sequences,
\begin{equation}\label{eq:2-8-18:c}
P^0[g(X)<s]=\lim_{n\to\infty}\,\lim_{m\to\infty}P^0\bigl[A_{n,m}(X)\bigr].
\end{equation}
Now $P^0[A_{n,m}(X)]$ is determined by the distribution of $(X_{u_1},\dots,X_{u_m})$, and by \bk{Problem 8.8} the processes $Y$ and $\abs{W}$ have the same finite-dimensional distributions under $P^0$. Therefore the right-hand side of \eqref{eq:2-8-18:c} is the same for $X=Y$ and for $X=\abs{W}$. Since $g(\abs{W})=g(W)=\gamma_t$, we obtain with \eqref{eq:2-8-18:b}
\begin{equation}\label{eq:2-8-18:d}
P^0[\gamma_t<s]=P^0[\theta_t<s];\qquad 0<s\le t.
\end{equation}

\emph{Step 3: the arc-sine law.} Let $F(s)\defeq\frac2\pi\arcsin\sqrt{s/t}$, $0\le s\le t$. By \bk{Problem 8.17}, $P^0[\theta_t\le s]=F(s)$ for $0\le s\le t$. The function $F$ is continuous on $[0,t]$, so $P^0[\theta_t<s]=\lim_{r\uparrow s}F(r)=F(s)$ for $0<s\le t$, and \eqref{eq:2-8-18:d} yields $P^0[\gamma_t<s]=F(s)$ for $0<s\le t$. Letting $r\downarrow s$ in $P^0[\gamma_t<r]=F(r)$ we get
\[
P^0[\gamma_t\le s]=\frac2\pi\arcsin\sqrt{\frac st};\qquad 0\le s<t,
\]
and this holds for $s=t$ as well, because $\gamma_t\le t$ and $F(t)=1$. In particular $P^0[\gamma_t=0]=F(0)=0$ and $P^0[\gamma_t=t]=1-\lim_{s\uparrow t}F(s)=0$. Finally, $F$ is continuous on $[0,t]$ and continuously differentiable on $(0,t)$, with
\[
F'(s)=\frac2\pi\cdot\frac{1}{\sqrt{1-s/t}}\cdot\frac{1}{2\sqrt{st}}=\frac{1}{\pi\sqrt{s(t-s)}}\,;\qquad 0<s<t,
\]
so that $F(s)=\int_0^sF'(u)\,du$ for $0\le s\le t$ (the fundamental theorem of calculus on $[\eps,s]$, and monotone convergence as $\eps\downarrow0$ and, when $s=t$, as the upper limit increases to $t$). The distribution function of $\gamma_t$ is thus that of the measure with density $F'$ on $(0,t)$, which means
\[
P^0[\gamma_t\in C]=\int_C\frac{ds}{\pi\sqrt{s(t-s)}}\quad\text{for every } C\in\B((0,t)),
\]
which is the assertion $P^0[\gamma_t\in ds]=ds/\pi\sqrt{s(t-s)}$; $0<s<t$. \qed

\rmk{The result can be checked without \bk{Problem 8.8}. By \eqref{eq:2-8-18:a} and the Markov property, $P^0[\gamma_t<s]=E^0\bigl[\varphi(W_s)\bigr]$ with $\varphi(x)=P^x[T_0>t-s]$, and $P^x[T_0>r]=P^0[\abs{W_r}<\abs{x}]$ by \bk{(8.4)} and symmetry. Hence, with $Z_1,Z_2$ independent standard normal random variables, $P^0[\gamma_t<s]=P\bigl[\sqrt{t-s}\,\abs{Z_2}<\sqrt{s}\,\abs{Z_1}\bigr]$. The law of $(Z_1,Z_2)$ is invariant under rotations, so this probability is the fraction of the directions $(\cos\alpha,\sin\alpha)$ with $\abs{\tan\alpha}<\sqrt{s/(t-s)}$, namely $\frac2\pi\arctan\sqrt{s/(t-s)}=\frac2\pi\arcsin\sqrt{s/t}$, in agreement with Step 3.}
\end{kex}

\begin{kex}{Exercise}{8.19}{}
Fix $t>0$. Recall $M_t=\max_{0\le v\le t}W_v$ from \bk{(8.1)}, $\theta_t=\sup\{0\le v\le t;\ W_v=M_t\}$ from \bk{(8.30)} and $\gamma_t=\sup\{0\le v\le t;\ W_v=0\}$ from \bk{(8.35)}. Let
\[
D\defeq\{(a,b,s,u)\in\R^4;\ a<0<b,\ 0<u<s<t\},
\]
and let $h(a,b,s,u)$ denote the function on $D$ displayed in the statement. The assertion to be proved is
\begin{equation}\label{eq:2-8-19:goal}
P^0\bigl[(W_t,M_t,\gamma_t,\theta_t)\in\Gamma\bigr]=\int_\Gamma h(a,b,s,u)\,da\,db\,ds\,du\qquad\text{for every Borel set }\Gamma\subseteq D.
\end{equation}
(The random variables $M_t$, $\theta_t$ are $\F_t$-measurable, because $\{\theta_t<r\}=\{\max_{r\le v\le t}W_v<M_t\}$ for $0<r\le t$; the measurability of $\gamma_t$ is checked in Step 1.)

\emph{Notation.} For $v>0$ and $x\in\R$ we write
\[
p_v(x)\defeq\frac{1}{\sqrt{2\pi v}}\,e^{-x^2/2v},\qquad
k_v(x)\defeq\frac{x}{\sqrt{2\pi v^3}}\,e^{-x^2/2v}=\frac{x}{v}\,p_v(x)=-\frac{\partial}{\partial x}p_v(x),
\]
so that $p(v;x,y)=p_v(y-x)$ is the Gaussian kernel and, for $x>0$, $k_v(x)$ is the passage-time density of \bk{(8.5)}. For $\tau>0$ and $y,a\in\R$ we put
\[
q_\tau(y,a)\defeq p_\tau(y-a)-p_\tau(y+a).
\]
Since $\frac1u+\frac1{s-u}=\frac{s}{u(s-u)}$, the function $h$ factorizes as
\begin{equation}\label{eq:2-8-19:fact}
h(a,b,s,u)=2\,k_u(b)\,k_{s-u}(b)\,k_{t-s}(-a);\qquad (a,b,s,u)\in D.
\end{equation}

\emph{Step 1: reduction to an event determined at a fixed time.} Let $\Omega_0\defeq\{W_0=0\}$, so that $P^0(\Omega_0)=1$, and fix $0<r<t$. On $\Omega_0$ the set $Z\defeq\{v\in[0,t];\ W_v=0\}$ is closed and contains $0$, so $\gamma_t=\max Z$ and
\[
\{\gamma_t<r\}\cap\Omega_0=\{W_v\neq0\ \text{for all } v\in[r,t]\}\cap\Omega_0
=\Bigl\{\min_{r\le v\le t}\abs{W_v}>0\Bigr\}\cap\Omega_0 .
\]
In particular $\gamma_t$ agrees on $\Omega_0$ with an $\F_t$-measurable random variable. By the continuity of the paths (intermediate value theorem),
\begin{equation}\label{eq:2-8-19:event}
\{W_t<0,\ \gamma_t<r\}\cap\Omega_0=E_r\cap\Omega_0,\qquad E_r\defeq\{W_v<0\ \text{for all } v\in[r,t]\}.
\end{equation}
On $E_r\cap\Omega_0$ we have $M_r\ge W_0=0>W_v$ for every $v\in[r,t]$; consequently $M_t=M_r$, the sets $\{v\in[0,t];\ W_v=M_t\}$ and $\{v\in[0,r];\ W_v=M_r\}$ coincide and do not contain $r$, and therefore
\begin{equation}\label{eq:2-8-19:same}
M_t=M_r,\qquad \theta_t=\theta_r<r\qquad\text{on } E_r\cap\Omega_0 .
\end{equation}

\emph{Step 2: the Markov property.} Fix $0<r<t$ and Borel sets $A\subseteq(-\infty,0)$, $B\subseteq(0,\infty)$, $C\subseteq(0,t)$, and put $\tau\defeq t-r$. Let
\[
F\defeq\Bigl\{\omega\in\R^{[0,\infty)};\ \sup_{v\in[0,\tau]\cap\Q}\omega(v)<0,\ \omega(\tau)\in A\Bigr\}\in\B\bigl(\R^{[0,\infty)}\bigr);
\]
this set is determined by countably many coordinates, and it is the strict inequality for the supremum (rather than for each rational coordinate separately) that is needed below. If $\omega$ is continuous, then $\sup_{v\in[0,\tau]\cap\Q}\omega(v)=\max_{0\le v\le\tau}\omega(v)$, because $[0,\tau]\cap\Q$ is dense in $[0,\tau]$ and the maximum is attained; hence a continuous $\omega$ belongs to $F$ if and only if $\omega(v)<0$ for every $v\in[0,\tau]$ and $\omega(\tau)\in A$. Since the paths of $W$ are continuous, we have exactly $\{W_{r+\cdot}\in F\}=E_r\cap\{W_t\in A\}$, and \eqref{eq:2-8-19:event}, \eqref{eq:2-8-19:same} give
\begin{equation}\label{eq:2-8-19:m1}
P^0[W_t\in A,\ M_t\in B,\ \theta_t\in C,\ \gamma_t<r]
=P^0\bigl[M_r\in B,\ \theta_r\in C\cap(0,r),\ W_r<0,\ W_{r+\cdot}\in F\bigr].
\end{equation}
The Brownian family is a Markov family (\bk{Theorem 5.12}), so by \bk{Proposition 5.15}, (c$'$) and (d$'$),
\[
P^0[W_{r+\cdot}\in F\mid\F_r]=g(W_r)\quad P^0\text{-a.s.},\qquad g(y)\defeq P^y[W_\cdot\in F].
\]
Let $y<0$, so that $P^y[W_0=y]=1$. On $\{W_0=y\}$ the continuous path $W_\cdot$ is negative throughout $[0,\tau]$ if and only if it has no zero in $[0,\tau]$ (intermediate value theorem), that is, if and only if $T_0>\tau$; by the description of $F$ just given, $\{W_\cdot\in F\}\cap\{W_0=y\}=\{T_0>\tau,\ W_\tau\in A\}\cap\{W_0=y\}$, with no exceptional set. Under $P^y$ the process $-W$ is a Brownian motion starting at $-y>0$, so it has the same law as $W$ under $P^{-y}$, and $T_0$ is the same for $W$ and $-W$; thus \bk{(8.9)} of \bk{Problem 8.6} yields
\begin{equation}\label{eq:2-8-19:g}
g(y)=P^{-y}[W_\tau\in -A,\ T_0>\tau]=\int_{-A}\bigl[p_\tau(z+y)-p_\tau(z-y)\bigr]\,dz=\int_A q_\tau(y,a)\,da;\qquad y<0 .
\end{equation}
In particular $q_\tau(y,a)\ge0$ for $y,a<0$ (this is also clear from $\abs{y-a}\le\abs{y+a}$), and $g$ is Borel measurable on $(-\infty,0)$. The random variables $M_r$, $\theta_r$, $W_r$ are $\F_r$-measurable, so conditioning on $\F_r$ in \eqref{eq:2-8-19:m1} gives
\begin{equation}\label{eq:2-8-19:m2}
P^0[W_t\in A,\ M_t\in B,\ \theta_t\in C,\ \gamma_t<r]
=E^0\bigl[\ind_{\{M_r\in B,\ \theta_r\in C\cap(0,r),\ W_r<0\}}\;g(W_r)\bigr].
\end{equation}
By \bk{Proposition 8.15}, applied with $r$ in place of $t$, the triple $(W_r,M_r,\theta_r)$ has on the set $\{(y,b,u);\ b\ge y^+,\ 0<u<r\}$, which contains $(-\infty,0)\times B\times(C\cap(0,r))$, the density
\[
\frac{b(b-y)}{\pi\sqrt{u^3(r-u)^3}}\exp\Bigl\{-\frac{b^2}{2u}-\frac{(b-y)^2}{2(r-u)}\Bigr\}=2\,k_u(b)\,k_{r-u}(b-y).
\]
Inserting this and \eqref{eq:2-8-19:g} into \eqref{eq:2-8-19:m2} and using Tonelli's theorem (all integrands are nonnegative and Borel measurable), we obtain
\begin{equation}\label{eq:2-8-19:m3}
\begin{split}
P^0[&W_t\in A,\ M_t\in B,\ \theta_t\in C,\ \gamma_t<r]\\
&=\int_A\int_B\int_{C\cap(0,r)}2\,k_u(b)\biggl(\int_{-\infty}^0k_{r-u}(b-y)\,q_{t-r}(y,a)\,dy\biggr)du\,db\,da .
\end{split}
\end{equation}

\emph{Step 3: evaluation of the inner integral.}
\begin{lem}
Let $b>0$, $a<0$ and $T>0$. Then for every $\sigma\in(0,T)$,
\begin{equation}\label{eq:2-8-19:lem}
L(\sigma)\defeq\int_{-\infty}^0k_\sigma(b-y)\,q_{T-\sigma}(y,a)\,dy=\int_0^\sigma k_v(b)\,k_{T-v}(-a)\,dv .
\end{equation}
\end{lem}
\begin{proof}
On $(0,T)\times\R$ define $\varphi(\sigma,y)\defeq k_\sigma(b-y)$ and $\psi(\sigma,y)\defeq q_{T-\sigma}(y,a)$. The function $(v,x)\mapsto p_v(x)$ satisfies the heat equation $\partial_vp=\frac12\partial_{xx}p$ on $(0,\infty)\times\R$, and so does its derivative $k=-\partial_xp$; hence
\begin{equation}\label{eq:2-8-19:heat}
\partial_\sigma\varphi=\tfrac12\,\partial_{yy}\varphi,\qquad \partial_\sigma\psi=-\tfrac12\,\partial_{yy}\psi\qquad\text{on }(0,T)\times\R .
\end{equation}
Moreover, since $p_\tau$ is even and $\partial_xp_\tau=-k_\tau$ is odd,
\begin{equation}\label{eq:2-8-19:bdry}
\psi(\sigma,0)=0,\qquad \partial_y\psi(\sigma,0)=-k_{T-\sigma}(-a)+k_{T-\sigma}(a)=-2\,k_{T-\sigma}(-a).
\end{equation}
(Indeed $\partial_y\psi(\sigma,y)=-k_{T-\sigma}(y-a)+k_{T-\sigma}(y+a)$, and $k_{T-\sigma}(a)=-k_{T-\sigma}(-a)$.)

Fix $0<\sigma_0<\sigma_1<T$. Each of $\varphi$, $\partial_y\varphi$, $\partial_{yy}\varphi$ is the product of $p_\sigma(b-y)$ with a polynomial in $b-y$ of degree at most $3$ whose coefficients are continuous functions of $\sigma\in(0,T)$; since $(b-y)^2\ge y^2$ for $y\le0$, there is a constant $K$ such that
\[
\abs{\varphi}+\abs{\partial_y\varphi}+\abs{\partial_{yy}\varphi}\le K(1+\abs{y})^3e^{-y^2/2\sigma_1}\qquad\text{for }\sigma\in[\sigma_0,\sigma_1],\ y\le0 .
\]
Likewise each of $\psi$, $\partial_y\psi$, $\partial_{yy}\psi$ is a sum of two terms of the form (polynomial in $x$ of degree at most $2$, with coefficients continuous in $\tau$) times $p_\tau(x)$, with $\tau=T-\sigma\in[T-\sigma_1,T-\sigma_0]$ and $x=y\mp a$, and is therefore bounded by a constant $K'$ on $[\sigma_0,\sigma_1]\times\R$. By \eqref{eq:2-8-19:heat}, $\abs{\partial_\sigma(\varphi\psi)}\le KK'(1+\abs{y})^3e^{-y^2/2\sigma_1}$ on $[\sigma_0,\sigma_1]\times(-\infty,0]$, which is integrable in $y$. Hence we may differentiate under the integral sign: $L$ is continuously differentiable on $(\sigma_0,\sigma_1)$ and, integrating by parts twice,
\[
L'(\sigma)=\frac12\int_{-\infty}^0\bigl(\partial_{yy}\varphi\cdot\psi-\varphi\cdot\partial_{yy}\psi\bigr)\,dy
=\frac12\Bigl[\partial_y\varphi\cdot\psi-\varphi\cdot\partial_y\psi\Bigr]_{y=-\infty}^{y=0}
=-\frac12\,\varphi(\sigma,0)\,\partial_y\psi(\sigma,0),
\]
where the boundary terms at $-\infty$ vanish by the bounds just established and we used $\psi(\sigma,0)=0$. By \eqref{eq:2-8-19:bdry},
\[
L'(\sigma)=k_\sigma(b)\,k_{T-\sigma}(-a);\qquad 0<\sigma<T,
\]
because $\sigma_0,\sigma_1$ were arbitrary. The right-hand side is continuous on $(0,T)$ and tends to $0$ as $\sigma\downarrow0$ (as $k_\sigma(b)\to0$ for $b>0$), so it is integrable on $(0,\sigma]$ and
\[
L(\sigma)-L(\sigma_0)=\int_{\sigma_0}^\sigma k_v(b)\,k_{T-v}(-a)\,dv;\qquad 0<\sigma_0<\sigma<T .
\]
It remains to show that $L(\sigma_0)\to0$ as $\sigma_0\downarrow0$. For $y,a<0$ we have $0\le q_\tau(y,a)\le p_\tau(y-a)$, so $\int_{-\infty}^0q_\tau(y,a)\,dy\le1$. The function $x\mapsto x\,e^{-x^2/2\sigma}$ is decreasing on $[\sqrt\sigma,\infty)$, so for $0<\sigma<b^2$ and $y\le0$ we have $0\le k_\sigma(b-y)\le k_\sigma(b)$. Therefore $0\le L(\sigma)\le k_\sigma(b)\to0$ as $\sigma\downarrow0$, and \eqref{eq:2-8-19:lem} follows.
\end{proof}

We apply the Lemma in \eqref{eq:2-8-19:m3} with $T=t-u$ and $\sigma=r-u\in(0,T)$, for $u\in C\cap(0,r)$, and substitute $v=s-u$:
\[
\int_{-\infty}^0k_{r-u}(b-y)\,q_{t-r}(y,a)\,dy=\int_u^rk_{s-u}(b)\,k_{t-s}(-a)\,ds .
\]
In view of \eqref{eq:2-8-19:fact}, the identity \eqref{eq:2-8-19:m3} becomes
\begin{equation}\label{eq:2-8-19:rect}
P^0[W_t\in A,\ M_t\in B,\ \theta_t\in C,\ \gamma_t<r]
=\int_A\int_B\int_C\int_0^r\ind_{\{u<s\}}\,h(a,b,s,u)\,ds\,du\,db\,da ,
\end{equation}
valid for all $0<r<t$ and all Borel sets $A\subseteq(-\infty,0)$, $B\subseteq(0,\infty)$, $C\subseteq(0,t)$.

\emph{Step 4: conclusion.} On the Borel subsets of $D$ consider the measures
\[
\mu(\Gamma)\defeq P^0\bigl[(W_t,M_t,\gamma_t,\theta_t)\in\Gamma\bigr],\qquad \nu(\Gamma)\defeq\int_\Gamma h(a,b,s,u)\,da\,db\,ds\,du .
\]
The sets
\[
\Gamma_{A,B,C,r}\defeq\{(a,b,s,u)\in D;\ a\in A,\ b\in B,\ u\in C,\ s<r\},
\]
with $A,B,C$ as above and $0<r<t$, form a class $\mathscr{C}$ of Borel subsets of $D$ which is closed under finite intersections, since $\Gamma_{A,B,C,r}\cap\Gamma_{A',B',C',r'}=\Gamma_{A\cap A',\,B\cap B',\,C\cap C',\,r\wedge r'}$.

We claim that $\mu=\nu$ on $\mathscr{C}$. The right-hand side of \eqref{eq:2-8-19:rect} is $\nu(\Gamma_{A,B,C,r})$, because for $a\in A$, $b\in B$, $u\in C$ and $0<s<r$ the point $(a,b,s,u)$ belongs to $D$ exactly when $u<s$. As for the left-hand side, let $\omega\in\Omega_0$ satisfy $W_t(\omega)\in A$ and $M_t(\omega)\in B$, i.e.\ $W_t<0<M_t$. Then $\gamma_t<t$, because $W_{\gamma_t}=0\ne W_t$; next, $W_v<0$ for $\gamma_t<v\le t$ by the intermediate value theorem, whereas $W_{\theta_t}=M_t>0$, so that $\theta_t\le\gamma_t$, and in fact $\theta_t<\gamma_t$ since $W_{\gamma_t}=0\neq M_t$; finally $\theta_t>0$ since $W_0=0\ne M_t$. Thus $0<\theta_t<\gamma_t<t$ at $\omega$, and consequently
\[
\bigl\{(W_t,M_t,\gamma_t,\theta_t)\in\Gamma_{A,B,C,r}\bigr\}\cap\Omega_0=\{W_t\in A,\ M_t\in B,\ \theta_t\in C,\ \gamma_t<r\}\cap\Omega_0 .
\]
Hence the left-hand side of \eqref{eq:2-8-19:rect} is $\mu(\Gamma_{A,B,C,r})$, and the claim is proved.

The class $\mathscr{C}$ generates $\B(D)$. Indeed, choose $r_n\uparrow t$ and put $D_n\defeq\Gamma_{(-\infty,0),(0,\infty),(0,t),r_n}=D\cap\{s<r_n\}\in\mathscr{C}$, so that $D_n\uparrow D$. For open intervals $A\subseteq(-\infty,0)$, $B\subseteq(0,\infty)$, $C\subseteq(0,t)$ and $0\le r'<r<t$, the set $\{(a,b,s,u)\in D;\ a\in A,\ b\in B,\ u\in C,\ r'<s<r\}$ belongs to $\sigma(\mathscr{C})$, being $\Gamma_{A,B,C,r}\setminus\bigcap_{m>1/(t-r')}\Gamma_{A,B,C,r'+1/m}$ (or $\Gamma_{A,B,C,r}$ itself if $r'=0$); every open subset of $D$ is a countable union of such sets (those with rational endpoints suffice), and the open subsets of $D$ generate $\B(D)$.

Now fix $n$. The set functions $\Gamma\mapsto\mu(\Gamma\cap D_n)$ and $\Gamma\mapsto\nu(\Gamma\cap D_n)$ are measures on $\B(D)$ which agree on $\mathscr{C}$ (as $\Gamma\cap D_n\in\mathscr{C}$ for $\Gamma\in\mathscr{C}$) and have the same finite total mass $\mu(D_n)=\nu(D_n)\le1$; the class $\mathscr{C}\cup\{D\}$ is closed under finite intersections and generates $\B(D)$, so by Dynkin's $\pi$-$\lambda$ theorem these two measures coincide on $\B(D)$. Letting $n\to\infty$ and using the continuity from below of $\mu$ and $\nu$, we obtain $\mu(\Gamma)=\nu(\Gamma)$ for every $\Gamma\in\B(D)$, which is \eqref{eq:2-8-19:goal}. \qed

\rmk{\emph{Check.} The argument just given shows that $\{W_t<0<M_t\}$ and $\{(W_t,M_t,\gamma_t,\theta_t)\in D\}$ differ by a $P^0$-null set, and $P^0[M_t=0]=0$ by \bk{(8.2)}; so the total mass of the density must be $P^0[W_t<0]=\frac12$, not $1$. This is confirmed by integrating \eqref{eq:2-8-19:fact}: we have $\int_{-\infty}^0k_{t-s}(-a)\,da=p_{t-s}(0)=(2\pi(t-s))^{-1/2}$ and, with $c\defeq u(s-u)/s$ and $\int_0^\infty b^2e^{-b^2/2c}\,db=\frac12\sqrt{2\pi}\,c^{3/2}$,
\[
\int_0^\infty2\,k_u(b)\,k_{s-u}(b)\,db=\frac{1}{\pi\,(u(s-u))^{3/2}}\cdot\frac{\sqrt{2\pi}}{2}\Bigl(\frac{u(s-u)}{s}\Bigr)^{3/2}=\frac{1}{\sqrt{2\pi s^3}} ,
\]
which does not depend on $u$. Hence
\[
P^0[W_t<0,\ \gamma_t\in ds,\ \theta_t\in du]=\frac{du\,ds}{2\pi s\sqrt{s(t-s)}}\,;\qquad 0<u<s<t,
\]
and integrating over $u\in(0,s)$ gives $P^0[W_t<0,\ \gamma_t\in ds]=ds/\bigl(2\pi\sqrt{s(t-s)}\bigr)$, one half of the arc-sine density of \bk{Problem 8.18}, as it must be by the symmetry of Brownian motion; its integral over $(0,t)$ is $\frac12$.}
\end{kex}

\begin{kex}{Exercise}{8.20}{}
We fix $t>0$ and write $U\defeq M_t(M_t-W_t)$,
\[
  p(a)\defeq\frac{1}{\sqrt{2\pi t}}\,e^{-a^2/2t},\quad a\in\R,
  \qquad\qquad
  q(u)\defeq\frac{2}{t}\,e^{-2u/t},\quad u>0 ,
\]
so that $p$ is the density of $W_t$ under $P^0$. Let $D\defeq\{(a,b)\in\R^2:\ b\ge 0,\ a\le b\}$ and
\[
  f(a,b)\defeq\frac{2(2b-a)}{\sqrt{2\pi t^3}}\exp\Bigl\{-\frac{(2b-a)^2}{2t}\Bigr\},\qquad (a,b)\in D ;
\]
$f\ge 0$ on $D$ because $2b-a=b+(b-a)\ge 0$ there. By the continuity of the paths,
$M_t=\sup_{s\in[0,t]\cap\Q}W_s$ is a random variable, $M_t\ge W_t$ everywhere and
$M_t\ge W_0=0$, $P^0$-a.s.; thus $(W_t,M_t)\in D$, $P^0$-a.s., and \bk{Proposition 8.1} states that
\begin{equation}\label{eq:2-8-20:a}
  P^0[(W_t,M_t)\in A]=\iint_{A\cap D} f(a,b)\,da\,db\qquad\text{for every } A\in\B(\R^2) .
\end{equation}
Consequently $E^0\bigl[\Phi(W_t,M_t)\bigr]=\iint_D\Phi(a,b)f(a,b)\,da\,db$ for every Borel function
$\Phi:\R^2\to[0,\infty)$ (by \eqref{eq:2-8-20:a} for indicators, by linearity for simple functions, and by
monotone convergence in general).

\emph{Step 1: a change of variable.} Fix $a\in\R$, put $a^+\defeq\max(a,0)$, and consider
\[
  \varphi_a(b)\defeq b(b-a),\qquad b\in(a^+,\infty).
\]
For $b>a^+$ we have $b>0$ and $b-a>0$, hence $\varphi_a'(b)=2b-a=b+(b-a)>0$. Moreover
$\varphi_a(b)\to a^+(a^+-a)=0$ as $b\downarrow a^+$ (one of the two factors vanishes, according as $a\ge0$
or $a<0$), and $\varphi_a(b)\to\infty$ as $b\to\infty$. Thus $\varphi_a$ is a continuously differentiable,
strictly increasing bijection of $(a^+,\infty)$ onto $(0,\infty)$. Furthermore
\begin{equation}\label{eq:2-8-20:b}
  (2b-a)^2=4b(b-a)+a^2=4\varphi_a(b)+a^2 ,
\end{equation}
so that, for $b>a^+$,
\begin{equation}\label{eq:2-8-20:c}
  f(a,b)=\frac{2}{\sqrt{2\pi t^3}}\exp\Bigl\{-\frac{a^2}{2t}-\frac{2\varphi_a(b)}{t}\Bigr\}\varphi_a'(b)
  =p(a)\,q\bigl(\varphi_a(b)\bigr)\,\varphi_a'(b).
\end{equation}
Hence, for every Borel function $h:[0,\infty)\to[0,\infty)$, the change of variable $u=\varphi_a(b)$ gives
\begin{equation}\label{eq:2-8-20:d}
  \int_{a^+}^\infty h\bigl(b(b-a)\bigr)f(a,b)\,db=p(a)\int_0^\infty h(u)\,q(u)\,du ,\qquad a\in\R.
\end{equation}

\emph{Step 2: the joint law of $(W_t,U)$.} Let $g:\R\to[0,\infty)$ and $h:[0,\infty)\to[0,\infty)$ be Borel
functions. On the event $\{(W_t,M_t)\in D\}$, of $P^0$-probability one, we have $U\ge 0$ and
$g(W_t)h(U)=\Phi(W_t,M_t)$ with the nonnegative Borel function $\Phi(a,b)\defeq g(a)\,h\bigl(b(b-a)\bigr)\ind_D(a,b)$.
The section of $D$ at $a\in\R$ is $\{b:\ (a,b)\in D\}=[a^+,\infty)$, which differs from $(a^+,\infty)$ by one
point. By the formula following \eqref{eq:2-8-20:a}, Tonelli's theorem and \eqref{eq:2-8-20:d},
\begin{equation}\label{eq:2-8-20:e}
\begin{split}
  E^0\bigl[g(W_t)h(U)\bigr]
  &=\int_{-\infty}^\infty g(a)\Bigl(\int_{a^+}^\infty h\bigl(b(b-a)\bigr)f(a,b)\,db\Bigr)da\\
  &=\Bigl(\int_{-\infty}^\infty g(a)\,p(a)\,da\Bigr)\Bigl(\int_0^\infty h(u)\,q(u)\,du\Bigr).
\end{split}
\end{equation}

\emph{Step 3: conclusions.} Taking $g\equiv 1$ and $h=\ind_B$, $B\in\B([0,\infty))$, in
\eqref{eq:2-8-20:e}, and using $\int_{-\infty}^\infty p(a)\,da=1$, we find
\[
  P^0[U\in B]=\int_{B\cap(0,\infty)}\frac{2}{t}\,e^{-2u/t}\,du ;
\]
that is, $U$ has under $P^0$ the exponential distribution with parameter $2/t$. In particular, for $x\ge0$,
\[
  P^0\bigl[M_t(M_t-W_t)\ge tx/2\bigr]=\int_{tx/2}^\infty\frac{2}{t}\,e^{-2u/t}\,du=e^{-x},
\]
which says that $(2/t)M_t(M_t-W_t)$ has the standard exponential distribution. Taking $h\equiv1$ and
$g=\ind_A$, $A\in\B(\R)$, we recover $P^0[W_t\in A]=\int_A p(a)\,da$. Finally, with $g=\ind_A$ and
$h=\ind_B$, \eqref{eq:2-8-20:e} and the two identities just obtained give
\[
  P^0[W_t\in A,\ U\in B]=P^0[W_t\in A]\,P^0[U\in B],\qquad A\in\B(\R),\ B\in\B([0,\infty)).
\]
Since $U\ge0$, $P^0$-a.s., the same product formula holds for every $B\in\B(\R)$ (replace $B$ by
$B\cap[0,\infty)$). Therefore $W_t$ and $M_t(M_t-W_t)$ are independent under $P^0$.

\rmk{The identity \eqref{eq:2-8-20:b} explains the result: in terms of the reflected endpoint $2M_t-W_t$,
we have $(2M_t-W_t)^2=W_t^2+4M_t(M_t-W_t)$, and the density \bk{(8.2)} depends on $(a,b)$ only through
$2b-a$. As a check, \eqref{eq:2-8-20:e} with $g\equiv h\equiv1$ gives $\iint_D f=1$, in agreement with
$P^0[(W_t,M_t)\in D]=1$; and $E^0[U]=t/2$ agrees with
$E^0[M_t^2]-E^0[M_tW_t]=t-t/2$, where $E^0[M_t^2]=E^0[W_t^2]=t$ by \bk{(8.3)} of \bk{Problem 8.2} and
$2E^0[M_tW_t]=E^0[M_t^2]+E^0[W_t^2]-E^0[(M_t-W_t)^2]=t$, again by \bk{(8.3)}.}
\end{kex}

\solsection{2.9}{The Brownian Sample Paths}

\begin{kex}{Problem}{9.5}{}
The book's Section 2.10 gives only the hint ``use \bk{Problem 7.18} and \bk{Lemma 9.4 (ii)}''; we
write out the whole argument. Throughout, $W=\{W_t,\F_t;\ 0\le t<\infty\}$ is the standard,
one-dimensional Brownian motion of this section. Its paths are continuous in the sense of
\bk{Definition 1.1} as it is read in this section (cf.\ \bk{Remark 9.2}), namely almost surely: there
is an event $\Omega_c\in\F$ with $P(\Omega_c)=1$ such that $W_0(\omega)=0$ and
$t\mapsto W_t(\omega)$ is continuous on $[0,\infty)$ for every $\omega\in\Omega_c$. We denote by
$\mathscr{Z}_\omega\defeq\{0\le t<\infty;\ W_t(\omega)=0\}$ the set of times at which the path is at
the origin, and by $\Q_+\defeq\Q\cap(0,\infty)$ the set of positive rationals.

\emph{What is to be proved.} We show that there is an event $\Omega^*\in\F$ with $P(\Omega^*)=1$ such
that for every $\omega\in\Omega^*$ and every $0\le T<\infty$ the path $W_\cdot(\omega)$ takes a strictly
positive value, a strictly negative value and the value zero on $(T,\infty)$. Consequently, for
$\omega\in\Omega^*$ the set $\mathscr{Z}_\omega$ is unbounded, and there are times
\begin{equation}\label{eq:2-9-05:a}
  t_1<u_1<t_2<u_2<\cdots,\qquad t_k\to\infty,\qquad W_{t_k}(\omega)=0,\quad W_{u_k}(\omega)\neq 0
  \quad(k\ge 1):
\end{equation}
the path leaves the origin and comes back to it infinitely many times, at arbitrarily large times.

\emph{Step 1: the behaviour near $t=0$, read off the rationals.} For a real-valued process
$X=\{X_t;\ 0\le t<\infty\}$ on $(\Omega,\F,P)$ we set
\[
  E(X)\defeq\bigcap_{n=1}^\infty\ \bigcup_{\substack{q,q'\in\Q_+\\ q\le 1/n,\ q'\le 1/n}}
  \{X_q>0,\ X_{q'}<0\}\in\F .
\]
We claim that $P[E(W)]=1$. According to \bk{Problem 7.18}, for each $n\ge 1$ there is an event of
probability one on which $W$ changes sign infinitely many times in $[0,1/n]$; in particular, there is an event $F_n\subseteq\Omega_c$ with $P(F_n)=1$ such that for
every $\omega\in F_n$ there are $s,s'\in[0,1/n]$ with $W_s(\omega)>0$ and $W_{s'}(\omega)<0$, and
$s,s'\in(0,1/n]$ because $W_0(\omega)=0$. Fix $\omega\in F_n$. By the continuity of the path, $W_\cdot(\omega)$ is strictly positive
on a neighbourhood of $s$ and strictly negative on a neighbourhood of $s'$, and each of these
neighbourhoods contains rationals in $(0,1/n]$. Hence
$F_n\subseteq\bigcup_{q,q'\in\Q_+,\,q\vee q'\le 1/n}\{W_q>0,\ W_{q'}<0\}$, so
$\bigcap_{n\ge1}F_n\subseteq E(W)$ and $P[E(W)]=1$.

\emph{Step 2: time-inversion.} Let $Y$ be the process of \bk{(9.4)}: $Y_0=0$ and $Y_t=tW_{1/t}$ for
$0<t<\infty$. We claim that $P[E(Y)]=1$. By \bk{Lemma 9.4 (ii)}, $Y$ is a standard Brownian motion;
what we need of this is only that the families $(Y_q)_{q\in\Q_+}$ and $(W_q)_{q\in\Q_+}$ have the
same finite-dimensional distributions, and this we verify directly. By \bk{Remark 9.2}, $W$ is a
zero-mean Gaussian process with covariance $E(W_sW_t)=s\wedge t$. For $0<t_1<\dots<t_k<\infty$ the
vector $(Y_{t_1},\dots,Y_{t_k})=(t_1W_{1/t_1},\dots,t_kW_{1/t_k})$ is the image under a linear map of
the zero-mean, jointly normal vector $(W_{1/t_k},\dots,W_{1/t_1})$, hence is itself zero-mean and
jointly normal, with
\[
  E(Y_{t_i}Y_{t_j})=t_it_j\Bigl(\frac1{t_i}\wedge\frac1{t_j}\Bigr)=t_i\wedge t_j
  =E(W_{t_i}W_{t_j}),\qquad 1\le i,j\le k .
\]
A normal distribution on $\R^k$ is determined by its mean vector and covariance matrix (this is the
observation following \bk{Definition 9.1}), so $(Y_{t_1},\dots,Y_{t_k})$ and
$(W_{t_1},\dots,W_{t_k})$ have the same distribution. Now let $\R^{\Q_+}$ carry the product $\sigma$-field
$\B(\R)^{\otimes\Q_+}$, and consider the measurable mappings
$\omega\mapsto(W_q(\omega))_{q\in\Q_+}$ and $\omega\mapsto(Y_q(\omega))_{q\in\Q_+}$ from $\Omega$ into
this space. The measures they induce agree on the finite-dimensional cylinder sets, which form a
$\pi$-system generating the product $\sigma$-field; by Dynkin's $\pi$-$\lambda$ theorem the two
induced measures coincide. The set
\[
  \Gamma\defeq\bigcap_{n=1}^\infty\ \bigcup_{\substack{q,q'\in\Q_+\\ q\le 1/n,\ q'\le 1/n}}
  \bigl\{x\in\R^{\Q_+};\ x_q>0,\ x_{q'}<0\bigr\}
\]
belongs to the product $\sigma$-field, $E(W)=\{(W_q)_{q\in\Q_+}\in\Gamma\}$ and
$E(Y)=\{(Y_q)_{q\in\Q_+}\in\Gamma\}$. Therefore $P[E(Y)]=P[E(W)]=1$, by Step 1.

\emph{Step 3: conclusion.} Put $\Omega^*\defeq E(Y)\cap\Omega_c$, so that $\Omega^*\in\F$ and
$P(\Omega^*)=1$ by Step 2. Fix $\omega\in\Omega^*$ and $0\le T<\infty$, and choose an integer $n\ge 1$
with $n>T$. By the definition of $E(Y)$ there are rationals $q,q'\in(0,1/n]$ with
$Y_q(\omega)>0$ and $Y_{q'}(\omega)<0$, that is,
\[
  q\,W_{1/q}(\omega)>0,\qquad q'\,W_{1/q'}(\omega)<0 .
\]
The numbers $a\defeq 1/q$ and $b\defeq 1/q'$ satisfy $a,b\ge n>T$, $a\neq b$, and
$W_a(\omega)>0>W_b(\omega)$. Since $\omega\in\Omega_c$, the path $W_\cdot(\omega)$ is continuous on the closed
interval with endpoints $a$ and $b$, so by the intermediate value theorem there is a point $t$ strictly between $a$
and $b$ with $W_t(\omega)=0$; in particular $t>T$. Thus on $(T,\infty)$ the path takes a positive
value, a negative value and the value zero, for every $0\le T<\infty$; in particular $\mathscr{Z}_\omega$ is
unbounded.

The times in \eqref{eq:2-9-05:a} are obtained recursively: let $t_1>0$ be a zero of
$W_\cdot(\omega)$, and, $t_k$ being chosen, let $u_k>t_k+1$ be a time with $W_{u_k}(\omega)>0$ and
$t_{k+1}>u_k$ a zero of $W_\cdot(\omega)$; both exist by what was just proved, and
$t_{k+1}>t_k+1$ gives $t_k\to\infty$. Since $P(\Omega^*)=1$, the proof is complete.

\rmk{We have shown that the set $\{\omega;\ \mathscr{Z}_\omega \text{ is unbounded}\}$ contains the
event $\Omega^*$ of probability one. Its trace on $\Omega_c$ is itself an event: for
$\omega\in\Omega_c$, the continuity of the path and the compactness of $[n,m]$ show that
$W_\cdot(\omega)$ has a zero in $[n,m]$ if and only if $\inf_{q\in\Q\cap[n,m]}\abs{W_q(\omega)}=0$,
whence
\[
  \{\omega;\ \mathscr{Z}_\omega \text{ is unbounded}\}\cap\Omega_c
  =\Omega_c\cap\bigcap_{n\ge1}\,\bigcup_{m>n}\Bigl\{\inf_{q\in\Q\cap[n,m]}\abs{W_q}=0\Bigr\}\in\F ,
\]
and this event, which contains $\Omega^*$, has probability one.}

\rmk{The phrase ``infinitely often'' has to be understood as ``at arbitrarily large times'', which
is how the result is used in the proof of \bk{Theorem 9.6 (ii)} (the zero set is unbounded). Read
literally as ``$\mathscr{Z}_\omega$ is an infinite set'', the assertion would already be contained in
\bk{Problem 7.18}: a continuous path which starts at the origin and changes sign infinitely many times
in $[0,\eps]$ has infinitely many zeros in $[0,\eps]$. The time-inversion $t\mapsto 1/t$ of
\bk{Lemma 9.4 (ii)} is what turns this accumulation of zeros at $t=0$ into an accumulation at
$t=\infty$. We have applied \bk{Problem 7.18} to $W$ and transferred its conclusion to $Y$ through the
finite-dimensional distributions, rather than applying it to $Y$ directly; in this way the argument
uses only the Gaussian computation behind \bk{Lemma 9.4 (ii)} and not the almost sure continuity of
$Y$ at $t=0$, which rests on \bk{Problem 9.3}.}
\end{kex}

\begin{kex}{Exercise}{9.17}{}
Throughout, $W=\{W_s,\F_s;\ 0\le s<\infty\}$ is the standard, one-dimensional Brownian motion of
Section~2.9; by \bk{Definition 1.1} every path $s\mapsto W_s(\omega)$ is continuous, so that the Dini
derivates of \bk{Definition 9.16} are defined for every $\omega\in\Omega$. We fix $t\in[0,\infty)$.

\emph{Step 1: reduction to the behaviour at infinity of a sequence.} Let
\[
  B_h\defeq W_{t+h}-W_t,\quad 0\le h<\infty,
  \qquad\qquad
  Y_s\defeq s\,B_{1/s}=s\,\bigl(W_{t+1/s}-W_t\bigr),\quad 0<s<\infty .
\]
Since $(W_{t+h}-W_t)/h=Y_{1/h}$ for $h>0$, and $h\downarrow0$ if and only if $s=1/h\to\infty$,
\bk{(9.12)} and \bk{(9.13)} give, for every $\omega\in\Omega$,
\begin{equation}\label{eq:2-9-17:a}
  D^{+}W_t(\omega)=\varlimsup_{s\to\infty}Y_s(\omega)=\inf_{r>0}\,\sup_{s\ge r}Y_s(\omega),
  \qquad
  D_{+}W_t(\omega)=\varliminf_{s\to\infty}Y_s(\omega)=\sup_{r>0}\,\inf_{s\ge r}Y_s(\omega).
\end{equation}
For every $\omega$ the function $s\mapsto Y_s(\omega)$ is continuous on $(0,\infty)$. Therefore
$\sup_{s\ge r}Y_s=\sup_{s\in\Q,\,s\ge r}Y_s$ and $\inf_{s\ge r}Y_s=\inf_{s\in\Q,\,s\ge r}Y_s$ are
random variables with values in $(-\infty,\infty]$ and $[-\infty,\infty)$, respectively; both are
monotone in $r$, so in \eqref{eq:2-9-17:a} the infimum and the supremum over $r>0$ may be taken over the
positive integers $r$. Consequently $D^{+}W_t$ and $D_{+}W_t$ are extended real-valued random
variables, and the set in \bk{(9.14)} belongs to $\F$. Moreover, \eqref{eq:2-9-17:a} shows that
\begin{equation}\label{eq:2-9-17:b}
  D^{+}W_t\ \ge\ \varlimsup_{n\to\infty}Y_n,\qquad
  D_{+}W_t\ \le\ \varliminf_{n\to\infty}Y_n
  \qquad\text{($n$ running through the positive integers).}
\end{equation}
It suffices then to prove that
\begin{equation}\label{eq:2-9-17:c}
  P\Bigl[\,\varlimsup_{n\to\infty}Y_n=\infty\Bigr]=1
  \qquad\text{and}\qquad
  P\Bigl[\,\varliminf_{n\to\infty}Y_n=-\infty\Bigr]=1 ,
\end{equation}
because the intersection of two events of probability one has probability one.

\emph{Step 2: the sequence $\{Y_n\}_{n\ge1}$ is a random walk with standard normal steps.}
The process $B=\{B_h,\F_{t+h};\ 0\le h<\infty\}$ is a standard Brownian motion in the sense of
\bk{Definition 1.1}: it is continuous and adapted to the filtration $\{\F_{t+h}\}_{h\ge0}$, we have
$B_0=0$, and for $0\le g<h$ the increment $B_h-B_g=W_{t+h}-W_{t+g}$ is independent of $\F_{t+g}$ and
normally distributed with mean zero and variance $h-g$. By \bk{Remark 9.2}, $B$ is a zero-mean
Gaussian process with $E(B_gB_h)=g\wedge h$. Fix $n\ge1$. The vector $(Y_1,\dots,Y_n)$ is the image of
the Gaussian vector $(B_1,B_{1/2},\dots,B_{1/n})$ under a linear map, hence it is Gaussian, with mean
zero and
\[
  E(Y_jY_k)=jk\,E\bigl(B_{1/j}B_{1/k}\bigr)=jk\Bigl(\frac1j\wedge\frac1k\Bigr)=j\wedge k,
  \qquad 1\le j,k\le n .
\]
(This is the time-inversion of \bk{Lemma 9.4}(ii) applied to $B$, restricted to integer times.)
Set $Y_0\defeq0$, so that the formula $E(Y_jY_k)=j\wedge k$ holds for $0\le j,k\le n$, and
$X_k\defeq Y_k-Y_{k-1}$ for $k\ge1$. The vector $(X_1,\dots,X_n)$ is again Gaussian with mean zero,
and for $1\le j\le k\le n$
\begin{align*}
  E(X_jX_k)&=j\wedge k-j\wedge(k-1)-(j-1)\wedge k+(j-1)\wedge(k-1)\\
  &=\begin{cases}
    j-j-(j-1)+(j-1)=0, & j<k,\\
    k-(k-1)-(k-1)+(k-1)=1, & j=k.
  \end{cases}
\end{align*}
A Gaussian vector with identity covariance matrix has independent standard normal components. As
$n$ is arbitrary, $X_1,X_2,\dots$ is a sequence of independent random variables, each with the
standard normal distribution, and
\begin{equation}\label{eq:2-9-17:d}
  Y_n=X_1+\dots+X_n,\qquad n\ge1 .
\end{equation}

\emph{Step 3: proof of \eqref{eq:2-9-17:c}.} Fix a real number $c>0$ and consider the events
\[
  A_n\defeq\bigl\{Y_n>c\sqrt n\,\bigr\},\quad n\ge1,
  \qquad\qquad
  A\defeq\bigcap_{m\ge1}\,\bigcup_{n\ge m}A_n=\{A_n\text{ occurs for infinitely many }n\}.
\]
By \eqref{eq:2-9-17:d}, $Y_n/\sqrt n$ has the standard normal distribution, so
$P(A_n)=p\defeq\int_c^\infty(2\pi)^{-1/2}e^{-x^2/2}\,dx>0$ for every $n$. The events
$\bigcup_{n\ge m}A_n$ decrease to $A$ as $m\uparrow\infty$, whence
\begin{equation}\label{eq:2-9-17:e}
  P(A)=\lim_{m\to\infty}P\Bigl(\,\bigcup_{n\ge m}A_n\Bigr)\ \ge\ \varlimsup_{m\to\infty}P(A_m)=p>0 .
\end{equation}
Now let
\[
  \wt A\defeq\Bigl\{\varlimsup_{n\to\infty}\frac{Y_n}{\sqrt n}\ge c\Bigr\}\ \supseteq\ A .
\]
For every $m\ge1$ we have $Y_m/\sqrt n\to0$ as $n\to\infty$, everywhere on $\Omega$, and thus,
by \eqref{eq:2-9-17:d},
\[
  \varlimsup_{n\to\infty}\frac{Y_n}{\sqrt n}
  =\varlimsup_{n\to\infty}\frac{X_{m+1}+\dots+X_n}{\sqrt n}
\]
is measurable with respect to $\sigma(X_{m+1},X_{m+2},\dots)$. Hence $\wt A$ belongs to the tail
$\sigma$-field $\bigcap_{m\ge1}\sigma(X_{m+1},X_{m+2},\dots)$ of the independent sequence
$\{X_k\}_{k\ge1}$, and Kolmogorov's zero-one law gives $P(\wt A)\in\{0,1\}$.
Since $P(\wt A)\ge P(A)>0$ by \eqref{eq:2-9-17:e}, we conclude that $P(\wt A)=1$. On $\wt A$ there are
infinitely many integers $n$ with $Y_n\ge\frac c2\sqrt n$, and therefore
$\varlimsup_{n\to\infty}Y_n=\infty$ on $\wt A$. This proves the first assertion of
\eqref{eq:2-9-17:c}.

The sequence $\{-X_k\}_{k\ge1}$ also consists of independent standard normal random variables, and its
partial sums are $-Y_n$. The argument just given used only this property of $\{X_k\}$ together with
\eqref{eq:2-9-17:d}; applied to $\{-X_k\}$ it yields $\varlimsup_{n\to\infty}(-Y_n)=\infty$, that is,
$\varliminf_{n\to\infty}Y_n=-\infty$, with probability one. This is the second assertion of
\eqref{eq:2-9-17:c}.

\emph{Conclusion.} By \eqref{eq:2-9-17:b} and \eqref{eq:2-9-17:c}, the event
$\{D^{+}W_t=\infty\text{ and }D_{+}W_t=-\infty\}$, which belongs to $\F$ by Step~1, contains an event of
probability one. This is \bk{(9.14)}.

\rmk{The exceptional null set depends on $t$, and the two equalities in \bk{(9.14)} fail at
certain random times: by \bk{Remark 9.19}, $D^{+}W_t(\omega)\le0$ at every point $t$ of local maximum
of $W_\cdot(\omega)$. This is the phenomenon described in \bk{Remark 9.15}.}
\end{kex}

\begin{kex}{Exercise}{9.21}{}
Let $W=\{W_t,\F_t;\ 0\le t<\infty\}$ be a standard, one-dimensional Brownian motion on $(\Omega,\F,P)$. We say that a function $f:[0,\infty)\to\R$ is \emph{H\"older-continuous with exponent $\gamma$ from the right at $t\ge0$} if there exist $K<\infty$ and $\delta>0$ such that $\abs{f(t+h)-f(t)}\le Kh^\gamma$ for all $h\in[0,\delta]$. This is the weakest reasonable form of the property: it is implied by two-sided H\"older continuity at $t$, and if $f$ is H\"older-continuous with exponent $\gamma$ on some interval $[a,b]$, $a<b$ (in particular, if $f$ is locally H\"older-continuous with exponent $\gamma$ in the sense of \bk{(2.8)} of \bk{Theorem 2.8}, in which case any sufficiently short interval $[a,b]$ will do), then it has the property at $t=a$. We prove:

\begin{claim}
There is an event $C^*\in\F$ with $P(C^*)=0$ such that, for every $\omega\notin C^*$, every $\gamma>\frac12$ and every $t\in[0,\infty)$, the path $W_\cdot(\omega)$ is not H\"older-continuous with exponent $\gamma$ from the right at $t$.
\end{claim}

\emph{Step 1: reduction.} For $\gamma>\frac12$ and integers $m,j,k\ge1$, let
\[
A^{m,\gamma}_{jk}\defeq\bigl\{\omega\in\Omega;\ \abs{W_{t+h}(\omega)-W_t(\omega)}\le jh^\gamma\ \text{for some } t\in[0,m] \text{ and all } h\in[0,1/k]\bigr\};
\]
for $m=1$ this is the set of \bk{(9.19)}. If $W_\cdot(\omega)$ is H\"older-continuous with exponent $\gamma$ from the right at some $t\ge0$, with constants $K,\delta$, then $\omega\in A^{m,\gamma}_{jk}$ for any integers $m\ge t$, $j\ge K$, $k\ge1/\delta$. Moreover, if $\frac12<\gamma'\le\gamma$, then $h^\gamma\le h^{\gamma'}$ for $0\le h\le 1/k\le1$, so that $A^{m,\gamma}_{jk}\subseteq A^{m,\gamma'}_{jk}$. Putting $\gamma_r\defeq\frac12+\frac1r$, and choosing for a given $\gamma>\frac12$ an integer $r$ with $\gamma_r\le\gamma$, we see that the set of $\omega$ for which the path is H\"older-continuous from the right with some exponent $\gamma>\frac12$ at some point $t\ge0$ is contained in
\[
\bigcup_{r=1}^\infty\bigcup_{m=1}^\infty\bigcup_{j=1}^\infty\bigcup_{k=1}^\infty A^{m,\gamma_r}_{jk}.
\]
This is a countable union. Hence it suffices to find, for fixed $\gamma>\frac12$ and fixed integers $m,j,k\ge1$, an event $C\in\F$ with $P(C)=0$ and $A^{m,\gamma}_{jk}\subseteq C$; the union $C^*$ of these events over $\gamma\in\{\gamma_r\}_{r\ge1}$ and $m,j,k\ge1$ then has the properties required in the Claim.

\emph{Step 2: the number of increments.} Fix $\gamma>\frac12$ and $m,j,k\ge1$, and choose an integer $N\ge1$ so large that
\begin{equation}\label{eq:2-9-21:N}
N\bigl(\gamma-\tfrac12\bigr)>1 .
\end{equation}
Set $a\defeq 2j(N+1)^\gamma$.

Let $\omega\in A^{m,\gamma}_{jk}$: there is a number $t\in[0,m]$ with $\abs{W_{t+h}(\omega)-W_t(\omega)}\le jh^\gamma$ for every $0\le h\le1/k$. Take an integer $n\ge(N+1)k$. There exists an integer $i$, $1\le i\le mn$, such that $(i-1)/n\le t\le i/n$. For $\nu=0,1,\dots,N$ we then have
\[
0\le\frac{i+\nu}{n}-t\le\frac{\nu+1}{n}\le\frac{N+1}{n}\le\frac1k ,
\]
so the choice $h=\frac{i+\nu}{n}-t$ is admissible and gives $\abs{W_{(i+\nu)/n}(\omega)-W_t(\omega)}\le j\bigl(\frac{N+1}{n}\bigr)^\gamma$. Consequently, for $\nu=1,\dots,N$,
\[
\abs{W_{(i+\nu)/n}(\omega)-W_{(i+\nu-1)/n}(\omega)}
\le\abs{W_{(i+\nu)/n}(\omega)-W_t(\omega)}+\abs{W_{(i+\nu-1)/n}(\omega)-W_t(\omega)}
\le\frac{a}{n^\gamma}.
\]
Therefore, with
\[
C^{(n)}_i\defeq\bigcap_{\nu=1}^N\Bigl\{\omega\in\Omega;\ \abs{W_{(i+\nu)/n}(\omega)-W_{(i+\nu-1)/n}(\omega)}\le\frac{a}{n^\gamma}\Bigr\}\in\F ,
\]
we have shown
\begin{equation}\label{eq:2-9-21:incl}
A^{m,\gamma}_{jk}\subseteq\bigcup_{i=1}^{mn}C^{(n)}_i\qquad\text{for every integer } n\ge(N+1)k .
\end{equation}

\emph{Step 3: the probability estimate.} Fix $n$ and $i$, and put
\[
Z_\nu\defeq\sqrt n\,\bigl(W_{(i+\nu)/n}-W_{(i+\nu-1)/n}\bigr),\qquad \nu=1,\dots,N .
\]
By \bk{Definition 1.1}, each $Z_\nu$ is standard normal and independent of $\F_{(i+\nu-1)/n}$, whereas $Z_1,\dots,Z_{\nu-1}$ are $\F_{(i+\nu-1)/n}$-measurable; by induction on $\nu$ the random variables $Z_1,\dots,Z_N$ are independent. For $\eps>0$ and a standard normal random variable $Z$,
\[
P[\abs{Z}\le\eps]=\int_{-\eps}^{\eps}\frac{1}{\sqrt{2\pi}}\,e^{-x^2/2}\,dx\le\frac{2\eps}{\sqrt{2\pi}}\le\eps .
\]
Since $C^{(n)}_i=\bigcap_{\nu=1}^N\{\abs{Z_\nu}\le a\,n^{\frac12-\gamma}\}$, we obtain the analogue of \bk{(9.17)}:
\begin{equation}\label{eq:2-9-21:bound}
P\bigl(C^{(n)}_i\bigr)=\prod_{\nu=1}^N P\bigl[\abs{Z_\nu}\le a\,n^{\frac12-\gamma}\bigr]\le a^N n^{-N(\gamma-\frac12)};\qquad i=1,\dots,mn .
\end{equation}

\emph{Step 4: conclusion.} Define, by analogy with \bk{(9.18)},
\[
C\defeq\bigcap_{n=(N+1)k}^\infty\ \bigcup_{i=1}^{mn}C^{(n)}_i\in\F .
\]
By \eqref{eq:2-9-21:incl} we have $A^{m,\gamma}_{jk}\subseteq C$, and by \eqref{eq:2-9-21:bound} and subadditivity
\[
P(C)\le\inf_{n\ge(N+1)k}P\Bigl(\bigcup_{i=1}^{mn}C^{(n)}_i\Bigr)\le\inf_{n\ge(N+1)k}m\,a^N\,n^{1-N(\gamma-\frac12)}=0 ,
\]
because the exponent $1-N(\gamma-\frac12)$ is negative by \eqref{eq:2-9-21:N}. In view of Step 1, the Claim is proved. \qed

\rmk{The only real modification of the proof of \bk{Theorem 9.18} is the choice \eqref{eq:2-9-21:N} of the number $N$ of consecutive increments. Each increment contributes a factor of order $n^{-(\gamma-1/2)}$ in \eqref{eq:2-9-21:bound}, and these have to beat the number $mn$ of events in the union. For $\gamma=1$ the three increments of the book suffice ($N=3$ gives $P(C^{(n)}_i)\le a^3n^{-3/2}$, as in \bk{(9.17)}, hence the bound $m\,a^3n^{-1/2}$ for the union over $i$), and so they do for every $\gamma>\frac56$; but as $\gamma\downarrow\frac12$ one needs $N>1/(\gamma-\frac12)\to\infty$. For $\gamma\le\frac12$ no choice of $N$ makes the bound tend to zero; for $\gamma<\frac12$ this is as it must be, since by \bk{Remark 2.12} almost every path is locally H\"older-continuous with every exponent $\gamma\in(0,\frac12)$. The argument gives no information in the boundary case $\gamma=\frac12$. Finally, since differentiability from the right at $t$, or merely finiteness of $D^+W_t(\omega)$ and $D_+W_t(\omega)$, implies H\"older continuity from the right at $t$ with exponent $1$, the Claim with $\gamma=1$ contains \bk{Theorem 9.18}.}
\end{kex}

\begin{kex}{Problem}{9.22}{}
\bk{Section 2.10} of the book gives only the hint ``use integration by parts''; we carry out the whole computation here.
We fix $x>0$ and write
\[
  I\defeq\int_x^\infty e^{-u^2/2}\,du ,
\]
which is finite and strictly positive, the integrand being continuous, positive and bounded by
$e^{-u/2}$ for $u\ge 1$.

\emph{The upper bound.} For $u\ge x$ we have $1\le u/x$, and $u\mapsto -e^{-u^2/2}$ is a primitive of
$u\mapsto u\,e^{-u^2/2}$ which tends to $0$ as $u\to\infty$. Hence
\[
  I\le\int_x^\infty\frac{u}{x}\,e^{-u^2/2}\,du
   =\frac1x\lim_{b\to\infty}\Bigl[-e^{-u^2/2}\Bigr]_{u=x}^{u=b}
   =\frac1x\,e^{-x^2/2}.
\]

\emph{The lower bound.} For $b>x$ we integrate by parts on $[x,b]$, with the continuously differentiable
functions $u\mapsto 1/u$ and $u\mapsto -e^{-u^2/2}$, whose derivatives are $-1/u^2$ and $u\,e^{-u^2/2}$:
\[
  \int_x^b e^{-u^2/2}\,du
  =\int_x^b\frac1u\cdot u\,e^{-u^2/2}\,du
  =\Bigl[-\frac1u\,e^{-u^2/2}\Bigr]_{u=x}^{u=b}-\int_x^b\frac{1}{u^2}\,e^{-u^2/2}\,du .
\]
As $b\to\infty$ the boundary term at $b$ tends to $0$, and the two integrals increase to the corresponding
integrals over $[x,\infty)$ (monotone convergence; both limits are finite, since
$u^{-2}e^{-u^2/2}\le x^{-2}e^{-u^2/2}$ on $[x,\infty)$). Therefore
\begin{equation}\label{eq:2-9-22:a}
  I=\frac1x\,e^{-x^2/2}-\int_x^\infty\frac{1}{u^2}\,e^{-u^2/2}\,du .
\end{equation}
Since $1/u^2\le 1/x^2$ for $u\ge x$, the last integral is at most $I/x^2$, and \eqref{eq:2-9-22:a} gives
\[
  \frac1x\,e^{-x^2/2}\le I+\frac{1}{x^2}\,I=\frac{1+x^2}{x^2}\,I .
\]
Multiplying by $x^2/(1+x^2)>0$ we obtain
\[
  \frac{x}{1+x^2}\,e^{-x^2/2}\le I ,
\]
which, together with the upper bound, is \bk{(9.20)}.

\rmk{Both inequalities in \bk{(9.20)} are in fact strict, because the comparisons $1\le u/x$ and
$1/u^2\le 1/x^2$ are strict for $u>x$ and the integrands are continuous and positive. The upper bound also
follows from \eqref{eq:2-9-22:a} by discarding the nonnegative integral on the right-hand side. The ratio of
the two bounds is $x^2/(1+x^2)\to 1$ as $x\to\infty$, so that $I\sim x^{-1}e^{-x^2/2}$ as
$x\to\infty$. The proofs of \bk{Theorem 9.23} and \bk{Theorem 9.25}, however, need the exact inequalities
for fixed $x$ and not merely this equivalence: both use the lower bound, in the form
$I\ge e^{-x^2/2}/(x+1/x)$, and the proof of \bk{Theorem 9.25} also uses the upper bound, to estimate the
last summand of \bk{(9.28)}.}
\end{kex}


%% ========================================================================
\solchapter{3}{Stochastic Integration}

\begin{center}
\begin{minipage}{0.86\textwidth}
\small
The 53 problems and exercises of Chapter 3 for which the book prints no solution: 2.12, 2.13, 2.25, 2.28, 2.30, 3.2, 3.10, 3.12, 3.17, 3.18, 3.20, 3.23, 3.24, 3.25, 3.31, 3.32, 3.33, 3.34, 3.35, 3.36, 3.37, 3.38, 3.39, 4.4, 4.16, 4.17, 4.22, 4.23, 4.24, 4.25, 5.6, 5.7, 5.8, 5.9, 5.10, 5.11, 5.18, 5.19, 5.20, 5.21, 6.6, 6.7, 6.18, 6.21, 6.24, 6.25, 6.30, 7.9, 7.10, 7.11, 7.12, 7.13, 7.14.  For 5 more, 3.7, 3.29, 4.5, 6.20, 7.7, the book prints only a hint or solves only some of the parts; what it leaves out is supplied here.  The book's own ``Solutions to Selected Problems'' fully covers 2.5, 2.18, 2.27, 2.29, 3.14, 3.15, 4.7, 4.11, 4.12, 6.12, 6.13, 6.19, 6.29, 7.3, 7.6, which are therefore omitted.
\end{minipage}
\end{center}

\bigskip
\hrule
\bigskip

\solsection{3.2}{Construction of the Stochastic Integral}

\begin{kex}{Problem}{2.12}{}
We show that the stopped process $\{W_{t\wedge T}\}$ converges to $W_T$ in $L^2(\Omega,\F,P)$ as $t\to\infty$; this gives \bk{(2.25)} and, with it, both identities. No assumption on the filtration $\{\F_t\}$ is needed beyond those of \bk{Definition 2.1.1}.

\emph{Step 1: two martingales.} By \bk{Definition 2.1.1}, $W$ is continuous and adapted, $W_0=0$ a.s., and for $0\le s<t$ the increment $W_t-W_s$ is independent of $\F_s$ and normal with mean zero and variance $t-s$. In particular $W_t$ is normal with variance $t$, so $E(W_t^2)=t<\infty$, and for $0\le s<t$ we have, a.s.,
\begin{align*}
  E[W_t-W_s\mid\F_s]&=E(W_t-W_s)=0,\\
  E[W_t^2-W_s^2\mid\F_s]
  &=E[(W_t-W_s)^2\mid\F_s]+2W_s\,E[W_t-W_s\mid\F_s]=t-s ;
\end{align*}
in the second line we used that $W_s$ is $\F_s$-measurable and that $W_s(W_t-W_s)$ is integrable, by the Cauchy--Schwarz inequality. Hence $\{W_t,\F_t;\ 0\le t<\infty\}$ and $\{W_t^2-t,\F_t;\ 0\le t<\infty\}$ are continuous martingales, both with expectation zero at $t=0$.

\emph{Step 2: the stopped processes.} Set $M_t\defeq W_{t\wedge T}$ for $0\le t<\infty$. Applying \bk{Problem 1.3.24}(i) (with $S=0$) to each of the four continuous martingales $\pm W$ and $\pm(W_t^2-t)$, each regarded as a right-continuous submartingale, we find that
\[
  \{M_t,\F_t;\ 0\le t<\infty\}\quad\text{and}\quad
  \{M_t^2-(t\wedge T),\F_t;\ 0\le t<\infty\}
\]
are martingales. Taking expectations and comparing with time $0$, we obtain for every $0\le t<\infty$
\begin{equation}\label{eq:3-2-12:a}
  E(M_t)=0,\qquad E(M_t^2)=E(t\wedge T)\le ET<\infty .
\end{equation}
(Here $M_t^2$ is integrable because $M_t^2-(t\wedge T)$ is, and $0\le t\wedge T\le t$.)

\emph{Step 3: $\{M_t\}$ is Cauchy in $L^2$.} Let $0\le s<t$. The random variable $M_s$ is $\F_s$-measurable and $M_sM_t$ is integrable by the Cauchy--Schwarz inequality and \eqref{eq:3-2-12:a}; therefore
$E(M_sM_t)=E\bigl(M_s\,E[M_t\mid\F_s]\bigr)=E(M_s^2)$, and with \eqref{eq:3-2-12:a},
\begin{equation}\label{eq:3-2-12:b}
  E\bigl[(M_t-M_s)^2\bigr]=E(M_t^2)-2E(M_sM_t)+E(M_s^2)=E(M_t^2)-E(M_s^2)=E(t\wedge T)-E(s\wedge T).
\end{equation}
As $u\to\infty$, $u\wedge T$ increases to $T$, so $E(u\wedge T)\uparrow ET<\infty$ by monotone convergence. The right-hand side of \eqref{eq:3-2-12:b} is at most $ET-E(s\wedge T)$, which tends to zero as $s\to\infty$. Thus
\[
  \lim_{s\to\infty}\ \sup_{t\ge s}\,E\bigl[(M_t-M_s)^2\bigr]=0,
\]
and by the completeness of $L^2(\Omega,\F,P)$ there is a square-integrable random variable $Y$ with $E[(M_t-Y)^2]\to0$ as $t\to\infty$.

\emph{Step 4: identification of the limit.} Since $ET<\infty$, the event $\Omega_0\defeq\{T<\infty\}\in\F$ has probability one. For $\omega\in\Omega_0$ and every $t\ge T(\omega)$ we have $M_t(\omega)=W_{T(\omega)}(\omega)$, so
\[
  W_T=\lim_{n\to\infty}M_n\quad\text{on }\Omega_0 ;
\]
we set $W_T\defeq0$ on the null event $\{T=\infty\}$, on which $W_T$ is not defined. This exhibits $W_T$ as an $\F$-measurable random variable. On the other hand, $M_n\to Y$ in $L^2$, so a subsequence $\{M_{n_k}\}$ converges to $Y$ almost surely. Consequently $Y=W_T$ a.s., that is,
\begin{equation}\label{eq:3-2-12:c}
  E(W_T^2)<\infty\qquad\text{and}\qquad \lim_{t\to\infty}E\bigl[(W_{t\wedge T}-W_T)^2\bigr]=0 .
\end{equation}

\emph{Step 5: conclusion.} Convergence in $L^2$ implies convergence of the first and of the second moments: by the Cauchy--Schwarz inequality and by the triangle inequality in $L^2$,
\begin{align*}
  \abs[\big]{E(W_{t\wedge T})-E(W_T)}
  &\le\bigl(E\bigl[(W_{t\wedge T}-W_T)^2\bigr]\bigr)^{1/2},\\
  \abs[\Big]{\bigl(E(W_{t\wedge T}^2)\bigr)^{1/2}-\bigl(E(W_T^2)\bigr)^{1/2}}
  &\le\bigl(E\bigl[(W_{t\wedge T}-W_T)^2\bigr]\bigr)^{1/2},
\end{align*}
and both right-hand sides tend to zero as $t\to\infty$ by \eqref{eq:3-2-12:c}. This proves \bk{(2.25)}. Combining \bk{(2.25)} with \eqref{eq:3-2-12:a} and with $E(t\wedge T)\uparrow ET$, we conclude
\[
  E(W_T)=\lim_{t\to\infty}E(W_{t\wedge T})=0,\qquad
  E(W_T^2)=\lim_{t\to\infty}E(W_{t\wedge T}^2)=\lim_{t\to\infty}E(t\wedge T)=ET .
\]

\rmk{The hypothesis $ET<\infty$ cannot be weakened to $P[T<\infty]=1$: for the passage time $T_b$ of \bk{(2.6.1)} to a level $b\ne0$ we have $P[T_b<\infty]=1$ (\bk{Remark 2.6.20}) and $W_{T_b}=b$ a.s., so $E(W_{T_b})=b\neq0$; compare \bk{Exercise 2.13}. The argument above is the special case $X\equiv1$, $M=W$ of the construction of \bk{Subsection 3.2.B}: by \bk{(2.24)}, $W_{t\wedge T}=\int_0^t\ind_{\{s\le T\}}\,dW_s$, and \eqref{eq:3-2-12:b} is the isometry \bk{(2.22)} for this integrand.}
\end{kex}

\begin{kex}{Exercise}{2.13}{}
Here $W=\{W_t,\F_t;\,0\le t<\infty\}$ is a standard, one-dimensional Brownian motion, and
\[
  T_b(\omega)=\inf\{t\ge 0;\ W_t(\omega)=b\},\qquad \inf\emptyset\defeq\infty .
\]
Since $W$ has continuous paths and $\{b\}$ is closed, $T_b$ is a stopping time of $\{\F_t\}$ by \bk{Problem 1.2.7}, so $ET_b\in[0,\infty]$ is well defined and \bk{Problem 2.12} is applicable to $T=T_b$ as soon as $ET_b<\infty$.

Fix $b\neq 0$ and suppose, contrary to the claim, that $ET_b<\infty$. Then $P[T_b<\infty]=1$. For every $\omega$ with $T_b(\omega)<\infty$, the set $\{t\ge 0;\ W_t(\omega)=b\}$ is nonempty and, being the preimage of the closed set $\{b\}$ under the continuous map $t\mapsto W_t(\omega)$, it is closed in $[0,\infty)$; hence it contains its infimum, that is,
\[
  W_{T_b(\omega)}(\omega)=b .
\]
Thus $W_{T_b}=b$ holds a.s., and therefore $E(W_{T_b})=b\neq 0$. On the other hand, the first Wald identity of \bk{Problem 2.12}, applied to the stopping time $T_b$ with $ET_b<\infty$, gives $E(W_{T_b})=0$. This contradiction shows that $ET_b=\infty$.

\rmk{Only the first Wald identity is needed; the second one would give $ET_b=E(W_{T_b}^2)=b^2$, which is not by itself contradictory. For $b=0$ we have $T_0=0$ because $W_0=0$ a.s., so the restriction $b\neq0$ is necessary. The result is consistent with \bk{(2.6.2)}--\bk{(2.6.3)}: $T_b$ is a.s.\ finite, but its density $\abs{b}(2\pi t^3)^{-1/2}e^{-b^2/2t}$ decays like $t^{-3/2}$ as $t\to\infty$, so that $\int_0^\infty t\,P[T_b\in dt]=\infty$.}
\end{kex}

\begin{kex}{Problem}{2.25}{}
Throughout, $\{\F_t\}$ satisfies the usual conditions (the standing assumption of \bk{Subsection 3.2.D}), and we fix $\alpha,\beta\in\R$ and set $L\defeq\alpha M+\beta N$. The plan is to show that $X\in\cP^*(L)$, so that the left-hand side is defined, and then to identify $\alpha I^M(X)+\beta I^N(X)$ as $I^L(X)$ by means of the characterization in \bk{Proposition 2.24}. For a process $Z$ and a stopping time $T$ we write $Z^T\defeq\{Z_{t\wedge T};\,0\le t<\infty\}$.

\emph{Step 1: three facts on localization.} Let $Z,Z'\in\cM^{c,\loc}$, with localizing sequences $\{S_k\}$, $\{S_k'\}$ as in \bk{Definition 1.5.15}.

\textbf{(a)} \emph{For every stopping time $T$, $Z^T\in\cM^{c,\loc}$.} Indeed, $(Z^T)^{S_k}=(Z^{S_k})^T$ is the continuous martingale $Z^{S_k}$ stopped at $T$, hence a martingale by \bk{Problem 1.3.24}\,(i) (applied to $Z^{S_k}$ and to $-Z^{S_k}$), and $Z^T_0=Z_0=0$ a.s.

\textbf{(b)} \emph{$aZ+bZ'\in\cM^{c,\loc}$ for all $a,b\in\R$.} The stopping times $S_k\wedge S_k'$ are nondecreasing in $k$ and tend to infinity a.s.; by \bk{Problem 1.3.24}\,(i) again, $Z^{S_k\wedge S_k'}=(Z^{S_k})^{S_k'}$ and $(Z')^{S_k\wedge S_k'}$ are martingales, and so is their linear combination $(aZ+bZ')^{S_k\wedge S_k'}$.

\textbf{(c)} \emph{If there is a constant $c$ with $P[\,\abs{Z_t}\le c\ \text{for all } t\ge0\,]=1$, then $Z\in\cM_2^c$.} For $0\le s<t$ we have $E[Z_{t\wedge S_k}\,|\,\F_s]=Z_{s\wedge S_k}$ a.s.; as $k\to\infty$, $Z_{t\wedge S_k}\to Z_t$ and $Z_{s\wedge S_k}\to Z_s$ a.s., with $\abs{Z_{t\wedge S_k}}\le c$, so the dominated convergence theorem for conditional expectations yields $E[Z_t\,|\,\F_s]=Z_s$ a.s. Thus $Z$ is a bounded, continuous martingale with $Z_0=0$ a.s.

By (b), $L\in\cM^{c,\loc}$. We also recall that for $Z,Z'\in\cM_2^c$ the cross-variation of \bk{Definition 1.5.5} is an adapted, continuous process of bounded variation, null at $t=0$, for which $ZZ'-\qv{Z,Z'}$ is a martingale; by the uniqueness assertion of \bk{Problem 1.5.17} it therefore agrees with the cross-variation of $Z$ and $Z'$ regarded as elements of $\cM^{c,\loc}$ (\bk{Definition 1.5.18}).

\emph{Step 2: bilinearity of the cross-variation on $\cM^{c,\loc}$.} Let $K\in\cM^{c,\loc}$. We claim that
\begin{equation}\label{eq:3-2-25:bil}
  \qv{L,K}_t=\alpha\qv{M,K}_t+\beta\qv{N,K}_t;\qquad 0\le t<\infty,\ \text{a.s. } P .
\end{equation}
The process $A\defeq\alpha\qv{M,K}+\beta\qv{N,K}$ is adapted, continuous, of bounded variation on bounded intervals, and $A_0=0$ a.s. Moreover,
\[
  LK-A=\alpha\bigl(MK-\qv{M,K}\bigr)+\beta\bigl(NK-\qv{N,K}\bigr)
\]
belongs to $\cM^{c,\loc}$ by \bk{Problem 1.5.17} and (b). The uniqueness assertion of \bk{Problem 1.5.17} gives $\qv{L,K}=A$ up to indistinguishability, which is \eqref{eq:3-2-25:bil}. Since $\qv{K,L}=\qv{L,K}$ (again by uniqueness), the cross-variation is linear in its second argument as well.

\emph{Step 3: $X\in\cP^*(L)$.} Applying Step 2 twice, to $L=\alpha M+\beta N$ and to $L'\defeq\alpha M-\beta N$, we obtain a.s., for all $t\ge0$,
\[
  \qv{L}_t=\alpha^2\qv{M}_t+2\alpha\beta\qv{M,N}_t+\beta^2\qv{N}_t,\qquad
  \qv{L'}_t=\alpha^2\qv{M}_t-2\alpha\beta\qv{M,N}_t+\beta^2\qv{N}_t,
\]
and therefore $\qv{L}+\qv{L'}=2\alpha^2\qv{M}+2\beta^2\qv{N}$. Let $\Omega_0$ be the event of probability one on which this identity holds for all $t$ and the four processes $\qv{L},\qv{L'},\qv{M},\qv{N}$ are continuous and nondecreasing. For $\omega\in\Omega_0$, denote by $\mu_L,\mu_{L'},\mu_M,\mu_N$ the measures on $\B([0,\infty))$ with distribution functions $\qv{L}(\omega),\qv{L'}(\omega),\qv{M}(\omega),\qv{N}(\omega)$. The measures $\mu_L+\mu_{L'}$ and $2\alpha^2\mu_M+2\beta^2\mu_N$ are finite on bounded sets and agree on all intervals $(s,t]$, hence coincide; in particular $\mu_L\le 2\alpha^2\mu_M+2\beta^2\mu_N$, and so
\[
  \int_0^T X_t^2\,d\qv{L}_t\le 2\alpha^2\int_0^T X_t^2\,d\qv{M}_t+2\beta^2\int_0^T X_t^2\,d\qv{N}_t;
  \qquad 0\le T<\infty,
\]
on $\Omega_0$. The right-hand side is a.s.\ finite for every $T$, because $X\in\cP^*(M)\cap\cP^*(N)$. As $X$ is progressively measurable, $X\in\cP^*(L)$ and $I^L(X)\in\cM^{c,\loc}$ is defined (\bk{Definition 2.23}).

\emph{Step 4: the integrals in \bk{(2.30)}.} For $Z,K\in\cM^{c,\loc}$ and a.e.\ $\omega$, the function $t\mapsto\qv{Z,K}_t(\omega)$ is continuous and of bounded variation on bounded intervals; $\int_0^t X_u\,d\qv{Z,K}_u$ is the Lebesgue--Stieltjes integral with respect to the associated signed measure on $[0,t]$. That it converges absolutely for $X\in\cP^*(Z)$ is the content of the following extension of \bk{Proposition 2.14}, which is implicit in the statement of \bk{Proposition 2.24}.

\begin{lem}
Let $Z,K\in\cM^{c,\loc}$, let $X$ be a progressively measurable process, and let $\check\xi_s$ denote the total variation of $\xi\defeq\qv{Z,K}$ on $[0,s]$. Then, a.s.,
\begin{equation}\label{eq:3-2-25:kw}
  \int_0^t\abs{X_s}\,d\check\xi_s\le\Bigl(\int_0^t X_s^2\,d\qv{Z}_s\Bigr)^{1/2}\bigl(\qv{K}_t\bigr)^{1/2};
  \qquad 0\le t<\infty .
\end{equation}
\end{lem}

\begin{proof}
For $n\ge1$ let $T_n\defeq\inf\{t\ge0;\ \abs{Z_t}\vee\abs{K_t}\ge n\}$. This is the first time the continuous process $(Z,K)$ enters a closed subset of $\R^2$, so it is a stopping time (\bk{Problem 1.2.7}); the sequence $\{T_n\}$ is nondecreasing, and $T_n\uparrow\infty$ because continuous paths are bounded on bounded intervals. The processes $Z^{(n)}\defeq Z^{T_n}$ and $K^{(n)}\defeq K^{T_n}$ are in $\cM^{c,\loc}$ by (a), and are bounded by $n$ on the event $\{Z_0=K_0=0\}$ of probability one; by (c), they belong to $\cM_2^c$. Since $(ZK-\qv{Z,K})^{T_n}=Z^{(n)}K^{(n)}-\qv{Z,K}^{T_n}$ is in $\cM^{c,\loc}$ by (a), the uniqueness in \bk{Problem 1.5.17} gives, a.s.,
\[
  \qv{Z^{(n)},K^{(n)}}_t=\qv{Z,K}_{t\wedge T_n},\qquad
  \qv{Z^{(n)}}_t=\qv{Z}_{t\wedge T_n},\qquad
  \qv{K^{(n)}}_t=\qv{K}_{t\wedge T_n};\qquad 0\le t<\infty
\]
(the last two by taking $K=Z$, respectively $Z=K$). In particular, the total variation of $\qv{Z^{(n)},K^{(n)}}$ on $[0,s]$ is $\check\xi_{s\wedge T_n}$.

For $m\ge1$ the process $X^{(m)}\defeq\abs{X}\wedge m$ is progressively measurable and, for every $T<\infty$,
\[
  E\int_0^T\bigl(X^{(m)}_s\bigr)^2\,d\qv{Z^{(n)}}_s\le m^2E\qv{Z^{(n)}}_T=m^2E\bigl[(Z^{(n)}_T)^2\bigr]\le m^2n^2,
\]
so that $X^{(m)}\in\cL^*(Z^{(n)})$; likewise $Y\equiv1$ is in $\cL^*(K^{(n)})$. \bk{Proposition 2.14}, applied to $Z^{(n)},K^{(n)}\in\cM_2^c$ with the integrands $X^{(m)}$ and $Y$, yields an event of probability one, which we may take to be the same for all $m,n\ge1$, on which
\[
  \int_0^{t}X^{(m)}_s\,d\check\xi_{s\wedge T_n}\le\Bigl(\int_0^{t}\bigl(X^{(m)}_s\bigr)^2\,d\qv{Z}_{s\wedge T_n}\Bigr)^{1/2}\bigl(\qv{K}_{t\wedge T_n}\bigr)^{1/2};
  \qquad 0\le t<\infty .
\]
For $t\le T_n$ the stopping may be removed from the three integrators. Letting first $m\to\infty$ (monotone convergence on both sides) and then $n\to\infty$, we obtain \eqref{eq:3-2-25:kw} for all $t\ge0$.
\end{proof}

\emph{Step 5: identification.} Put $\Phi\defeq\alpha I^M(X)+\beta I^N(X)$; by (b), $\Phi\in\cM^{c,\loc}$. Fix $K\in\cM_2^c$. By Step 2 (with $I^M(X)$, $I^N(X)$ in place of $M$, $N$) and by \bk{Proposition 2.24} applied to the pairs $(M,X)$ and $(N,X)$, we have a.s., for all $t\ge0$,
\begin{equation}\label{eq:3-2-25:phi}
  \qv{\Phi,K}_t=\alpha\qv{I^M(X),K}_t+\beta\qv{I^N(X),K}_t
  =\alpha\int_0^tX_u\,d\qv{M,K}_u+\beta\int_0^tX_u\,d\qv{N,K}_u .
\end{equation}
Let $\Omega_1$ be the event of probability one on which \eqref{eq:3-2-25:bil} holds for all $t$, the processes $\qv{M,K}$, $\qv{N,K}$, $\qv{L,K}$ are continuous and of bounded variation on bounded intervals, \eqref{eq:3-2-25:kw} holds for the pairs $(Z,K)=(M,K)$ and $(N,K)$, and $\int_0^tX_u^2\,d\qv{M}_u+\int_0^tX_u^2\,d\qv{N}_u<\infty$ for all $t$. Fix $\omega\in\Omega_1$ and $t\ge0$, and let $\nu_M,\nu_N,\nu_L$ be the finite signed measures on $\B([0,t])$ determined by $\qv{M,K}(\omega)$, $\qv{N,K}(\omega)$, $\qv{L,K}(\omega)$. By \eqref{eq:3-2-25:bil}, the finite signed measures $\nu_L$ and $\alpha\nu_M+\beta\nu_N$ agree on every interval $(s,s']\subseteq[0,t]$ and assign mass zero to $\{0\}$; these sets form a $\pi$-system generating $\B([0,t])$, and the class of Borel sets on which two finite signed measures agree is a $\lambda$-system when they agree on $[0,t]$, so $\nu_L=\alpha\nu_M+\beta\nu_N$ by Dynkin's $\pi$-$\lambda$ theorem. Consequently the total variation measures satisfy $\abs{\nu_L}\le\abs{\alpha}\abs{\nu_M}+\abs{\beta}\abs{\nu_N}$. By \eqref{eq:3-2-25:kw}, $X(\omega)$ is integrable on $[0,t]$ with respect to $\abs{\nu_M}$ and $\abs{\nu_N}$, hence also with respect to $\abs{\nu_L}$, and
\[
  \int_0^tX_u\,d\qv{L,K}_u=\alpha\int_0^tX_u\,d\qv{M,K}_u+\beta\int_0^tX_u\,d\qv{N,K}_u .
\]
Combining this with \eqref{eq:3-2-25:phi}, we conclude that
\[
  \qv{\Phi,K}_t=\int_0^tX_u\,d\qv{L,K}_u;\qquad 0\le t<\infty,\ \text{a.s. } P,
\]
for every $K\in\cM_2^c$. Thus $\Phi\in\cM^{c,\loc}$ satisfies \bk{(2.30)} with $L$ in place of $M$, for every $K\in\cM_2^c$. Since $L\in\cM^{c,\loc}$ and $X\in\cP^*(L)$ (Step 3), the uniqueness assertion of \bk{Proposition 2.24} gives $\Phi=I^L(X)$, that is,
\[
  I^{\alpha M+\beta N}(X)=\alpha I^M(X)+\beta I^N(X),
\]
the equality being understood, as always for stochastic integrals, up to indistinguishability.

\rmk{The hypothesis $X\in\cP^*(M)\cap\cP^*(N)$ is what makes the right-hand side meaningful; Step 3 shows that it also suffices for the left-hand side. The converse fails: with $N=-M$ and $\alpha=\beta=1$ every progressively measurable $X$ is in $\cP^*(\alpha M+\beta N)=\cP^*(0)$, whatever $\cP^*(M)$ may be.}
\end{kex}

\begin{kex}{Problem}{2.28}{}
Throughout, $W=\{W_t,\F_t;\,0\le t<\infty\}$ is a standard Brownian motion with respect to a filtration satisfying the usual conditions (the standing assumption of \bk{Section 3.2}); thus $W_0=0$ a.s.\ and, for $0\le s<t$, the increment $W_t-W_s$ is independent of $\F_s$ and normal with mean zero and variance $t-s$. Here $\qv{W}_t=t$, so the integral $I(X)=\int_0^{\cdot}X_u\,dW_u$ is defined for every $X\in\cP$ (\bk{Definition 2.23} together with \bk{Remark 2.11}), it is a continuous, adapted process, and $\int_s^tX_u\,dW_u$ means $I_t(X)-I_s(X)$. By \bk{(2.31)} there is an event $\Omega^*$ with $P(\Omega^*)=1$ on which $\int_0^TX_u^2\,du<\infty$ for every $T<\infty$; on $\Omega^*$ all the quantities in \bk{(2.35)} are finite and
\begin{equation}\label{eq:3-2-28:add}
  \zeta_t(X)=\zeta_s(X)+\zeta_t^s(X),\qquad 0\le s<t<\infty ,
\end{equation}
while on the null set $\Omega\setminus\Omega^*$ we set $\zeta_t(X)\defeq 0$ for all $t$. We treat simple integrands first, and then pass to the limit with the help of \bk{Problem 2.27}.

\begin{lem}
Let $\G$ be a sub-$\sigma$-field of $\F$, let $\xi$ be a real-valued, $\G$-measurable random variable, and let $Z$ be a normal random variable with mean zero and variance $h>0$ which is independent of $\G$. Then
\[
  E\Bigl[\exp\Bigl(\xi Z-\tfrac12\xi^2h\Bigr)\Bigm|\G\Bigr]=1\qquad\text{a.s.}
\]
\end{lem}

\begin{proof}
The function $f(x,z)\defeq\exp(xz-\frac12x^2h)$ is nonnegative and continuous on $\R^2$, and for every $x\in\R$
\[
  \int_{\R}f(x,z)\,\frac{e^{-z^2/2h}}{\sqrt{2\pi h}}\,dz
  =\int_{\R}\frac{1}{\sqrt{2\pi h}}\exp\Bigl(-\frac{(z-xh)^2}{2h}\Bigr)dz=1 .
\]
Fix $A\in\G$. The random vector $(\ind_A,\xi)$ is $\G$-measurable, hence independent of $Z$, so the law of $(\ind_A,\xi,Z)$ on $\R^3$ is the product of the law $\mu$ of $(\ind_A,\xi)$ and the normal law of $Z$. By the Fubini--Tonelli theorem for nonnegative integrands,
\[
  E\bigl[\ind_Af(\xi,Z)\bigr]
  =\int_{\R^2}a\Bigl(\int_{\R}f(x,z)\,\frac{e^{-z^2/2h}}{\sqrt{2\pi h}}\,dz\Bigr)\mu(da,dx)
  =\int_{\R^2}a\,\mu(da,dx)=P(A).
\]
Since $A\in\G$ is arbitrary, the nonnegative random variable $f(\xi,Z)$ is integrable and its conditional expectation given $\G$ is $1$.
\end{proof}

\emph{Step 1: $X\in\cL_0$.} Let $X$ be simple, with $\{t_i\}_{i=0}^\infty$, $\{\xi_i\}_{i=0}^\infty$ and $C$ as in \bk{Definition 2.3}. The process $X$ is adapted and left-continuous, hence progressively measurable, and $\abs{X}\le C$; thus $X\in\cL^*(W)\subseteq\cP^*$, we may take $\Omega^*=\Omega$, and \bk{Definition 2.23} applies to $X$ directly. We check that the integral $I(X)$ of that definition coincides with the integral $I^W(X)$ constructed earlier for the integrator $W\in\cM_2^c$ and integrands in $\cL^*(W)$. Since $W\in\cM_2^c$, we may choose $S_n\equiv\infty$ in the construction leading to \bk{(2.34)}, which does not depend on the choice of $\{S_n\}$. Then \bk{(2.32)}, \bk{(2.33)} read
\[
  T_n=R_n=n\wedge\inf\Bigl\{t\ge0;\ \int_0^tX_u^2\,du\ge n\Bigr\},\qquad
  W^{(n)}_t=W_{t\wedge T_n},\qquad X^{(n)}_t=X_t\ind_{\{T_n\ge t\}},
\]
and $T_n\ge n\wedge(n/C^2)$ because $\int_0^tX_u^2\,du\le C^2t$; in particular $T_n\to\infty$ everywhere. For all $t\ge0$ we have $W^{(n)}_{t\wedge T_n}=W_{t\wedge T_n}$ and $X^{(n)}_{t\wedge T_n}=X_{t\wedge T_n}\ind_{\{T_n\ge t\wedge T_n\}}=X_{t\wedge T_n}$, with $W,W^{(n)}\in\cM_2^c$, $X\in\cL^*(W)$ and $X^{(n)}\in\cL^*(W^{(n)})$. \bk{Corollary 2.21}, applied with the stopping time $T_n$, therefore gives, almost surely, $I^{W^{(n)}}_{t\wedge T_n}(X^{(n)})=I^W_{t\wedge T_n}(X)$ for all $t\ge0$, that is, $I^{W^{(n)}}_t(X^{(n)})=I^W_t(X)$ for $0\le t\le T_n$. Discarding the union over $n$ of the exceptional null sets and using \bk{(2.34)} together with $T_n\to\infty$, we conclude that almost surely $I_t(X)=I^W_t(X)$ for all $t\ge0$. Consequently $I(X)$ is given by \bk{(2.5)}:
\[
  I_t(X)=\sum_{i=0}^\infty\xi_i\,(W_{t\wedge t_{i+1}}-W_{t\wedge t_i}),\qquad 0\le t<\infty .
\]
Suppose that $t_i\le u<v\le t_{i+1}$ for some $i\ge0$. In the difference $I_v(X)-I_u(X)$ the terms with index $i'<i$ cancel (both equal $\xi_{i'}(W_{t_{i'+1}}-W_{t_{i'}})$), the terms with index $i'>i$ vanish, and the term with index $i$ contributes $\xi_i(W_v-W_u)$. Furthermore $X_r=\xi_i$ for $r\in(u,v]\subseteq(t_i,t_{i+1}]$. Hence
\begin{equation}\label{eq:3-2-28:piece}
  \zeta_v^u(X)=\xi_i\,(W_v-W_u)-\tfrac12\,\xi_i^2\,(v-u)\qquad\text{whenever } t_i\le u<v\le t_{i+1}.
\end{equation}
Here $\xi_i$ is $\F_{t_i}$-measurable, hence $\F_u$-measurable, and $W_v-W_u$ is independent of $\F_u$ and normal with mean zero and variance $v-u$. The Lemma, with $\G=\F_u$, yields
\begin{equation}\label{eq:3-2-28:one}
  E\bigl[\exp\bigl(\zeta_v^u(X)\bigr)\bigm|\F_u\bigr]=1\quad\text{a.s.},\qquad\text{whenever } t_i\le u<v\le t_{i+1}.
\end{equation}

Now fix $0\le s<t<\infty$. Because $t_i\uparrow\infty$ strictly, only finitely many of the $t_i$ lie in $(s,t)$; adjoining $s$ and $t$ to them we obtain points $s=u_0<u_1<\dots<u_k=t$ such that every interval $[u_{j-1},u_j]$ is contained in some $[t_i,t_{i+1}]$. Put $Y_j\defeq\exp(\zeta_{u_j}^{u_{j-1}}(X))$, $1\le j\le k$. Each $Y_j$ is nonnegative, $\F_{u_j}$-measurable (because $I(X)$ is adapted and, by \eqref{eq:3-2-28:piece}, so is the Lebesgue integral), and satisfies $E[Y_j\mid\F_{u_{j-1}}]=1$ a.s.\ by \eqref{eq:3-2-28:one}. We claim that
\begin{equation}\label{eq:3-2-28:prod}
  E\bigl[Y_1Y_2\cdots Y_j\bigm|\F_s\bigr]=1\quad\text{a.s.},\qquad 1\le j\le k .
\end{equation}
For $j=1$ this is \eqref{eq:3-2-28:one}. If it holds for $j-1$, then, the conditional expectation of nonnegative random variables being well defined and the nonnegative product $Y_1\cdots Y_{j-1}$ being $\F_{u_{j-1}}$-measurable, the tower property gives
\[
  E\bigl[Y_1\cdots Y_j\bigm|\F_s\bigr]
  =E\Bigl[Y_1\cdots Y_{j-1}\,E\bigl[Y_j\bigm|\F_{u_{j-1}}\bigr]\Bigm|\F_s\Bigr]
  =E\bigl[Y_1\cdots Y_{j-1}\bigm|\F_s\bigr]=1\quad\text{a.s.}
\]
By \eqref{eq:3-2-28:add}, $\exp(\zeta_t(X))=\exp(\zeta_s(X))\,Y_1\cdots Y_k$, and $\exp(\zeta_s(X))$ is nonnegative, finite and $\F_s$-measurable; therefore \eqref{eq:3-2-28:prod} with $j=k$ gives
\begin{equation}\label{eq:3-2-28:mart}
  E\bigl[\exp\bigl(\zeta_t(X)\bigr)\bigm|\F_s\bigr]=\exp\bigl(\zeta_s(X)\bigr)\quad\text{a.s.},\qquad 0\le s<t<\infty .
\end{equation}
Taking $s=0$ and recalling that $\zeta_0(X)=0$, we get $E\exp(\zeta_t(X))=1<\infty$ for every $t$. Thus $\exp(\zeta(X))$ is adapted and integrable and satisfies \eqref{eq:3-2-28:mart}: \emph{it is a martingale when $X\in\cL_0$}.

\emph{Step 2: $X\in\cP$.} By the last assertion of \bk{Problem 2.27} there is a sequence $\{X^{(n)}\}_{n=1}^\infty\subseteq\cL_0$ such that, for every $T>0$, almost surely
\begin{equation}\label{eq:3-2-28:conv}
  \lim_{n\to\infty}\int_0^T\abs{X^{(n)}_u-X_u}^2\,du=0
  \qquad\text{and}\qquad
  \lim_{n\to\infty}\sup_{0\le t\le T}\abs{I_t(X^{(n)})-I_t(X)}=0 .
\end{equation}
Intersecting the corresponding events over $T\in\N$ and with $\Omega^*$, we obtain an event $\Omega^{**}\subseteq\Omega^*$ of probability one on which \eqref{eq:3-2-28:conv} holds for all $T>0$ simultaneously. Fix $\omega\in\Omega^{**}$ and $t\ge0$. The triangle inequality in $L^2([0,t],du)$ gives
\[
  \Bigl|\Bigl(\int_0^t\bigl(X^{(n)}_u\bigr)^2du\Bigr)^{1/2}-\Bigl(\int_0^tX_u^2\,du\Bigr)^{1/2}\Bigr|
  \le\Bigl(\int_0^t\abs{X^{(n)}_u-X_u}^2\,du\Bigr)^{1/2}\xrightarrow[n\to\infty]{}0 ,
\]
where $\int_0^tX_u^2\,du<\infty$; hence $\int_0^t(X^{(n)}_u)^2\,du\to\int_0^tX_u^2\,du$. Together with the second relation in \eqref{eq:3-2-28:conv} this shows
\begin{equation}\label{eq:3-2-28:lim}
  \lim_{n\to\infty}\zeta_t\bigl(X^{(n)}\bigr)=\zeta_t(X)\quad\text{on }\Omega^{**},\qquad\text{for every } t\ge0 .
\end{equation}

\emph{Adaptedness.} By Step 1 each $\zeta_t(X^{(n)})$ is $\F_t$-measurable. By \eqref{eq:3-2-28:lim}, $\zeta_t(X)$ is the a.s.\ limit of these random variables, and $\F_t$ contains all $P$-null sets of $\F$; hence $\zeta_t(X)$, and with it $\exp(\zeta_t(X))$, is $\F_t$-measurable.

\emph{Supermartingale inequality.} Fix $0\le s<t<\infty$. The random variables $\exp(\zeta_t(X^{(n)}))$ are nonnegative and converge a.s.\ to $\exp(\zeta_t(X))$, so by Fatou's lemma for conditional expectations, by \eqref{eq:3-2-28:mart} applied to each $X^{(n)}\in\cL_0$, and by \eqref{eq:3-2-28:lim} at time $s$,
\[
  E\bigl[\exp\bigl(\zeta_t(X)\bigr)\bigm|\F_s\bigr]
  \le\liminf_{n\to\infty}E\bigl[\exp\bigl(\zeta_t(X^{(n)})\bigr)\bigm|\F_s\bigr]
  =\liminf_{n\to\infty}\exp\bigl(\zeta_s(X^{(n)})\bigr)
  =\exp\bigl(\zeta_s(X)\bigr)\quad\text{a.s.}
\]

\emph{Integrability.} With $s=0$ the last display and $\zeta_0(X)=0$ give $E\exp(\zeta_t(X))\le1<\infty$ for every $t\ge0$.

Consequently $\{\exp(\zeta_t(X)),\F_t;\ 0\le t<\infty\}$ is a (nonnegative) supermartingale for every $X\in\cP$, and by Step 1 it is a martingale if $X\in\cL_0$.

\rmk{The argument shows $E\exp(\zeta_t(X))\le1$, and a supermartingale with constant expectation is a martingale; so for $X\in\cP$ the process $\exp(\zeta(X))$ is a martingale if and only if $E\exp(\zeta_t(X))=1$ for all $t\ge0$. Only convergence at fixed times was used in Step 2; the continuity of the paths of $\zeta(X)$ on $\Omega^*$ plays no role in the proof.}
\end{kex}

\begin{kex}{Exercise}{2.30}{}
Throughout, $s\ge 0$ is fixed, $\{\F_t\}$ satisfies the usual conditions, $Z$ is real-valued (if $Z$ is only a.s.\ finite we change it on a null set, which belongs to $\F_0$), and we write $J\defeq I^M(X)\in\cM^{c,\loc}$, so that $\int_s^t X_u\,dM_u = J_t-J_s$ for $t\ge s$.

\emph{Meaning of the left-hand side.} The process $\{ZX_u;\ 0\le u<\infty\}$ need not be adapted, because $Z$ is not $\F_u$-measurable for $u<s$; only its restriction to $(s,\infty)$ matters, and we set
\[
  Y_u(\omega)\defeq Z(\omega)\,\ind_{(s,\infty)}(u)\,X_u(\omega),\qquad
  \int_s^t ZX_u\,dM_u\defeq \int_0^t Y_u\,dM_u = I^M_t(Y),\quad t\ge s .
\]
We check that $Y\in\cP^*$. Let $H_u(\omega)\defeq Z(\omega)\ind_{(s,\infty)}(u)$ and fix $t\ge 0$. If $t\le s$, then $H$ vanishes on $[0,t]\times\Omega$; if $t>s$, then for every $B\in\B(\R)$ with $0\notin B$,
\[
  \{(u,\omega)\in[0,t]\times\Omega;\ H_u(\omega)\in B\}=(s,t]\times\{Z\in B\}\in\B([0,t])\otimes\F_t,
\]
because $\{Z\in B\}\in\F_s\subseteq\F_t$, and $\{H=0\}$ is the complement in $[0,t]\times\Omega$ of $\{H\in\R\setminus\{0\}\}$. Thus $H$ is progressively measurable, and so is the product $Y=HX$. Moreover, for every $T\in[0,\infty)$,
\[
  \int_0^T Y_u^2\,d\qv{M}_u = Z^2\int_{(s,T]}X_u^2\,d\qv{M}_u\le Z^2\int_0^TX_u^2\,d\qv{M}_u<\infty\quad\text{a.s.},
\]
by \bk{(2.31)} for $X$ and the finiteness of $Z$. Hence $Y\in\cP^*$ and $I^M(Y)$ is defined by \bk{Definition 2.23}. We shall prove that
\begin{equation}\label{eq:3-2-30:a}
  I^M_t(Y)=Z\,(J_t-J_{t\wedge s});\qquad 0\le t<\infty,\ \text{a.s.},
\end{equation}
which for $t\ge s$ is the assertion of the exercise.

\begin{lem}
Let $K\in\cM^{c,\loc}$ and let $\zeta$ be a real-valued, $\F_s$-measurable random variable. Then the process $\Psi_t\defeq\zeta\,(K_t-K_{t\wedge s})$, $0\le t<\infty$, belongs to $\cM^{c,\loc}$.
\end{lem}

\begin{proof}
$\Psi$ is continuous, $\Psi_0=0$, and $\Psi$ is adapted: $\Psi_t=0$ for $t\le s$, while for $t>s$ the variables $\zeta$, $K_t$, $K_s$ are $\F_t$-measurable. By \bk{Definition 1.5.15} there is a nondecreasing sequence $\{\rho_n\}$ of stopping times with $\rho_n\to\infty$ a.s., such that $K^{(n)}_t\defeq K_{t\wedge\rho_n}$ is a martingale for every $n$. Set
\[
  \tau_n\defeq\begin{cases} s &\text{on }\{\abs{\zeta}>n\},\\ \infty&\text{on }\{\abs{\zeta}\le n\},\end{cases}
  \qquad \theta_n\defeq\rho_n\wedge\tau_n .
\]
Each $\tau_n$ is a stopping time, since $\{\tau_n\le t\}$ is empty for $t<s$ and equals $\{\abs{\zeta}>n\}\in\F_s\subseteq\F_t$ for $t\ge s$. The sets $\{\abs{\zeta}>n\}$ decrease in $n$, so $\{\tau_n\}$ is nondecreasing, and $\tau_n(\omega)=\infty$ as soon as $n\ge\abs{\zeta(\omega)}$. Therefore $\{\theta_n\}$ is a nondecreasing sequence of stopping times with $\theta_n\to\infty$ a.s.

Let $\zeta_n\defeq\zeta\ind_{\{\abs{\zeta}\le n\}}$, a bounded $\F_s$-measurable random variable. We claim that
\begin{equation}\label{eq:3-2-30:b}
  \Psi_{t\wedge\theta_n}=\zeta_n\bigl(K^{(n)}_t-K^{(n)}_{t\wedge s}\bigr);\qquad 0\le t<\infty .
\end{equation}
Indeed, $\Psi_{t\wedge\theta_n}=\zeta\,(K_{t\wedge\theta_n}-K_{t\wedge\theta_n\wedge s})$. On $\{\abs{\zeta}>n\}$ we have $\theta_n\le s$, so $t\wedge\theta_n\wedge s=t\wedge\theta_n$ and the left side of \eqref{eq:3-2-30:b} is zero, as is the right side. On $\{\abs{\zeta}\le n\}$ we have $\theta_n=\rho_n$ and $\zeta=\zeta_n$, and \eqref{eq:3-2-30:b} holds again.

It remains to see that the right-hand side $\Psi^{(n)}_t$ of \eqref{eq:3-2-30:b} is a martingale. It is adapted (as $\Psi$ was), and integrable because $\abs{\zeta_n}\le n$ and $K^{(n)}_t$, $K^{(n)}_{t\wedge s}$ are integrable. Let $0\le r<t$. If $t\le s$, then $\Psi^{(n)}_t=\Psi^{(n)}_r=0$. If $s\le r$, then $\zeta_n$ and $K^{(n)}_s$ are $\F_r$-measurable, and
\[
  E\bigl[\Psi^{(n)}_t\,\big|\,\F_r\bigr]=\zeta_n\bigl(E[K^{(n)}_t\,|\,\F_r]-K^{(n)}_s\bigr)=\zeta_n\bigl(K^{(n)}_r-K^{(n)}_s\bigr)=\Psi^{(n)}_r\quad\text{a.s.}
\]
If $r<s<t$, the case just treated (with $r$ replaced by $s$) gives $E[\Psi^{(n)}_t\,|\,\F_s]=\Psi^{(n)}_s=0$, whence $E[\Psi^{(n)}_t\,|\,\F_r]=E\bigl[E[\Psi^{(n)}_t\,|\,\F_s]\,\big|\,\F_r\bigr]=0=\Psi^{(n)}_r$ a.s. Thus $\{\Psi_{t\wedge\theta_n},\F_t;\ 0\le t<\infty\}$ is a martingale for every $n$, and $\Psi\in\cM^{c,\loc}$.
\end{proof}

\emph{Proof of \eqref{eq:3-2-30:a}.} Let $\Phi_t\defeq Z\,(J_t-J_{t\wedge s})$. By the Lemma (with $K=J$, $\zeta=Z$), $\Phi\in\cM^{c,\loc}$. According to \bk{Proposition 2.24}, $I^M(Y)$ is the unique element of $\cM^{c,\loc}$ satisfying \bk{(2.30)} with integrand $Y$ for every $N\in\cM_2^c$; hence it suffices to show that, for every $N\in\cM_2^c$,
\begin{equation}\label{eq:3-2-30:c}
  \qv{\Phi,N}_t=\int_0^tY_u\,d\qv{M,N}_u;\qquad 0\le t<\infty,\ \text{a.s.}
\end{equation}
Fix $N\in\cM_2^c\subseteq\cM^{c,\loc}$. \bk{Proposition 2.24}, applied to $X$, gives
\[
  A_t\defeq\qv{J,N}_t=\int_0^tX_u\,d\qv{M,N}_u;\qquad 0\le t<\infty,\ \text{a.s.},
\]
where the Lebesgue--Stieltjes integral on the right is a.s.\ well defined: denoting by $\check\xi_t$ the total variation of $\xi\defeq\qv{M,N}$ on $[0,t]$, we have
\begin{equation}\label{eq:3-2-30:d}
  \int_0^t\abs{X_u}\,d\check\xi_u<\infty;\qquad 0\le t<\infty,\ \text{a.s.}
\end{equation}
To verify \eqref{eq:3-2-30:d}, let $T_n$, $M^{(n)}$, $X^{(n)}$ be as in \bk{(2.32)}, \bk{(2.33)}, so that $M^{(n)}\in\cM_2^c$, $X^{(n)}\in\cL^*(M^{(n)})$ and $T_n\to\infty$ a.s., and put $N^{(n)}_t\defeq N_{t\wedge T_n}$. The process $N^{(n)}$ is a continuous martingale by \bk{Problem 1.3.24(i)}, applied to $N$ and to $-N$, and it is square-integrable because the optional sampling theorem (\bk{Problem 1.3.23(i)}) for the continuous submartingale $N^2$ and the bounded stopping times $t\wedge T_n\le t$ gives $E\bigl[N^2_{t\wedge T_n}\bigr]\le E\bigl[N_t^2\bigr]<\infty$; thus $N^{(n)}\in\cM_2^c$, and the constant process $1$ belongs to $\cL^*(N^{(n)})$ since $E\qv{N^{(n)}}_T=E\bigl[(N^{(n)}_T)^2\bigr]<\infty$. Next, $\Lambda\defeq MN-\xi$ is in $\cM^{c,\loc}$ by \bk{Problem 1.5.17}; if $\{\sigma_k\}$ is a localizing sequence for $\Lambda$, then $\{\Lambda_{t\wedge T_n\wedge\sigma_k},\F_t\}$ is, for every $k$, the martingale $\{\Lambda_{t\wedge\sigma_k},\F_t\}$ stopped at $T_n$, hence a martingale by \bk{Problem 1.3.24(i)}. Consequently $M^{(n)}N^{(n)}-\xi_{\cdot\wedge T_n}=\Lambda_{\cdot\wedge T_n}\in\cM^{c,\loc}$, and since $\xi_{\cdot\wedge T_n}$ is adapted, continuous, of bounded variation and null at $t=0$, the uniqueness assertion of \bk{Problem 1.5.17} yields $\qv{M^{(n)},N^{(n)}}=\xi_{\cdot\wedge T_n}$ a.s. The total variation of the latter process on $[0,t]$ is $\check\xi_{t\wedge T_n}$. Therefore \bk{Proposition 2.14}, applied to $M^{(n)},N^{(n)}\in\cM_2^c$ and the integrands $X^{(n)}$, $1$, gives a.s., for all $t\ge0$,
\[
  \int_0^{t\wedge T_n}\abs{X_u}\,d\check\xi_u
  =\int_0^t\abs{X^{(n)}_u}\,d\check\xi_{u\wedge T_n}
  \le\Bigl(\int_0^t\bigl(X^{(n)}_u\bigr)^2\,d\qv{M^{(n)}}_u\Bigr)^{1/2}\bigl(\qv{N^{(n)}}_t\bigr)^{1/2}<\infty ,
\]
the right-hand side being a.s.\ finite for all $t$ because both factors are nondecreasing in $t$ and have finite expectation at every fixed time. As a.s.\ $T_n>t$ for all sufficiently large $n$, \eqref{eq:3-2-30:d} follows. Since $\abs{Y_u}\le\abs{Z}\,\abs{X_u}$, the integral in \eqref{eq:3-2-30:c} is a.s.\ well defined too, and for a.e.\ $\omega$ and all $t\ge 0$,
\[
  \int_0^tY_u\,d\qv{M,N}_u=Z\int_{(s,t]}X_u\,d\qv{M,N}_u=Z\,(A_t-A_{t\wedge s})\eqqcolon C_t
\]
(the interval $(s,t]$ being empty for $t\le s$). The process $C$ is continuous with $C_0=0$; it is adapted, since $C_t=0$ for $t\le s$ and $Z$, $A_t$, $A_s$ are $\F_t$-measurable for $t>s$; and its total variation on $[0,t]$ is at most $\abs{Z}$ times that of $A$, hence finite. By the uniqueness assertion of \bk{Problem 1.5.17}, \eqref{eq:3-2-30:c} will follow once we show that $\Phi N-C\in\cM^{c,\loc}$.

By \bk{Problem 1.5.17}, $L\defeq JN-A$ belongs to $\cM^{c,\loc}$. For every $t\ge 0$,
\begin{align*}
  \Phi_tN_t-C_t
  &=Z\bigl[(J_tN_t-A_t)-(J_{t\wedge s}N_{t\wedge s}-A_{t\wedge s})-J_{t\wedge s}(N_t-N_{t\wedge s})\bigr]\\
  &=Z\,(L_t-L_{t\wedge s})-ZJ_s\,(N_t-N_{t\wedge s}),
\end{align*}
where we used that $N_t-N_{t\wedge s}=0$ for $t\le s$ and $J_{t\wedge s}=J_s$ for $t\ge s$. The first term is in $\cM^{c,\loc}$ by the Lemma with $K=L$, $\zeta=Z$; the second by the Lemma with $K=N$ and $\zeta=ZJ_s$, which is real-valued and $\F_s$-measurable. If $\{\theta'_n\}$ and $\{\theta''_n\}$ are localizing sequences for these two processes, then $\{\theta'_n\wedge\theta''_n\}$ localizes both, because a continuous martingale stopped at a stopping time is again a martingale (\bk{Problem 1.3.24(i)}, applied to the martingale and to its negative); hence the difference $\Phi N-C$ is in $\cM^{c,\loc}$. This proves \eqref{eq:3-2-30:c}, and \bk{Proposition 2.24} yields $\Phi=I^M(Y)$ up to indistinguishability, which is \eqref{eq:3-2-30:a}. For $s\le t<\infty$ it reads
\[
  \int_s^tZX_u\,dM_u=Z\,(J_t-J_s)=Z\int_s^tX_u\,dM_u\quad\text{a.s.},
\]
the exceptional null set being the same for all $t\ge s$.

\rmk{The definition of the left-hand side adopted above agrees with the convention $\int_s^t\wt Y_u\,dM_u=I_t(\wt Y)-I_s(\wt Y)$ whenever the latter makes sense: if $\wt Y\in\cP^*$ coincides with $ZX$ on $(s,\infty)\times\Omega$ (for instance $\wt Y=ZX$ itself, when $Z$ happens to be $\F_0$-measurable), then $Y=\wt Y-\wt Y\ind_{[0,s]}$, and \bk{(2.17)}, \bk{(2.24)} with $T\equiv s$ (both valid for integrands in $\cP^*$ and $M\in\cM^{c,\loc}$, by localization) give $I_t(Y)=I_t(\wt Y)-I_{t\wedge s}(\wt Y)$. No boundedness or integrability of $Z$ is needed: this is why the argument is carried out in $\cM^{c,\loc}$, with the stopping times $\tau_n$ cutting off the event $\{\abs{Z}>n\}$ at time $s$.}
\end{kex}

\solsection{3.3}{The Change-of-Variable Formula}

\begin{kex}{Problem}{3.2}{}
Throughout, $\{\F_t\}$ satisfies the usual conditions; this is the standing assumption of Section 3.3, made in the sentence immediately preceding \bk{Definition 3.1}, and it is used in Step 2 below. Let $\Omega_0$ be an event of probability one on which \bk{(3.1)} and the second decomposition hold for all $t\ge 0$, and on which $M_0=\wt M_0=0$ (recall that members of $\cM^{c,\loc}$ vanish at $t=0$ a.s., \bk{Definition 1.5.15}). We set
\[
  N_t \defeq M_t-\wt M_t,\qquad 0\le t<\infty ,
\]
a continuous, adapted process with
\begin{equation}\label{eq:3-3-02:a}
  N_t=\wt B_t-B_t=\wt B_t-A^+_t+A^-_t,\quad 0\le t<\infty,\qquad N_0=0,\qquad\text{on }\Omega_0 .
\end{equation}
It suffices to show that $P[N_t=0;\ \forall\, 0\le t<\infty]=1$: on the intersection of this event with $\Omega_0$ we have $M_t=\wt M_t$ and, by \eqref{eq:3-3-02:a}, $B_t=\wt B_t$ for all $t\ge 0$.

\emph{Step 1: the total variation of $N$.} For $t\ge 0$, $m\ge 1$ and $\omega\in\Omega$ let
\[
  D^m_t \defeq \sum_{k=1}^{2^m}\abs[\big]{N_{tk2^{-m}}-N_{t(k-1)2^{-m}}},
  \qquad V_t\defeq \sup_{m\ge 1} D^m_t\in[0,\infty].
\]
We claim that, for every $\omega$, $V_t(\omega)$ is the total variation of the continuous function $s\mapsto N_s(\omega)$ on $[0,t]$, i.e.\ the supremum of $\sum_{j=1}^p\abs{N_{s_j}-N_{s_{j-1}}}$ over all partitions $0=s_0<s_1<\dots<s_p=t$. Each $D^m_t$ is such a sum, so $V_t$ does not exceed the total variation. Conversely, let $t>0$ and let $0=s_0<\dots<s_p=t$ be given; put $s_j^m\defeq t2^{-m}\lfloor 2^m s_j/t\rfloor$. Then $0=s^m_0\le s^m_1\le\dots\le s^m_p=t$ are points of the grid $\{tk2^{-m}\}_{k=0}^{2^m}$, so the triangle inequality gives $\sum_{j=1}^p\abs{N_{s^m_j}-N_{s^m_{j-1}}}\le D^m_t\le V_t$; since $s^m_j\to s_j$ as $m\to\infty$ and $N(\omega)$ is continuous, $\sum_{j=1}^p\abs{N_{s_j}-N_{s_{j-1}}}\le V_t$. This proves the claim (the case $t=0$ being trivial, with $V_0=0$). We record three consequences.
\begin{enumerate}[label=\textup{(\alph*)}]
\item $V_t$ is $\F_t$-measurable for every $t$, being a countable supremum of $\F_t$-measurable random variables.
\item For every $\omega$, the function $t\mapsto V_t(\omega)$ is nondecreasing (it is a total variation) and lower semicontinuous (it is a supremum of the continuous functions $t\mapsto D^m_t(\omega)$); hence it is left-continuous on $(0,\infty)$: $V_s\le V_t$ for $s<t$ and $\liminf_{s\uparrow t}V_s\ge V_t$.
\item On $\Omega_0$ we have $V_t\le \check{\wt B}_t+A^+_t+A^-_t<\infty$ for every $t\ge 0$, where $\check{\wt B}_t$ is the total variation of $\wt B$ on $[0,t]$; this follows from \eqref{eq:3-3-02:a}, the hypothesis on $\wt B$, and the monotonicity of $A^\pm$.
\end{enumerate}

\emph{Step 2: localization.} For $n\ge1$ let $S_n\defeq\inf\{t\ge 0:\ V_t>n\}$ $(\inf\emptyset=\infty)$. Since $V$ is nondecreasing, $\{S_n<t\}=\bigcup_{q\in\Q,\,0\le q<t}\{V_q>n\}\in\F_t$ by (a), so $S_n$ is an optional time, and therefore a stopping time because $\{\F_t\}$ is right-continuous (\bk{Proposition 1.2.3}). We have
\begin{equation}\label{eq:3-3-02:b}
  V_{t\wedge S_n}\le n\qquad\text{for all } t\ge0 \text{ and all }\omega .
\end{equation}
Indeed, $V_s\le n$ for $s<S_n$ by the definition of $S_n$; and if $0<S_n<\infty$, then $V_{S_n}=\lim_{s\uparrow S_n}V_s\le n$ by (b), whereas $V_{S_n}=V_0=0$ if $S_n=0$. Moreover, $V_t\le n$ implies $S_n\ge t$, so by (c)
\begin{equation}\label{eq:3-3-02:c}
  \lim_{n\to\infty}S_n=\infty\qquad\text{on }\Omega_0 .
\end{equation}
Let $\{T_n\}$ and $\{\wt T_n\}$ be nondecreasing sequences of stopping times, tending to infinity a.s., such that $\{M_{t\wedge T_n},\F_t\}$ and $\{\wt M_{t\wedge \wt T_n},\F_t\}$ are martingales for every $n$ (\bk{Definition 1.5.15}), and put $\tau_n\defeq S_n\wedge T_n\wedge \wt T_n$, a stopping time. By \bk{Problem 1.3.24}\,(i), applied to the continuous martingales $\pm M_{\cdot\wedge T_n}$ and $\pm\wt M_{\cdot\wedge\wt T_n}$ and to the stopping time $\tau_n$, the processes $M_{t\wedge\tau_n}=M_{(t\wedge\tau_n)\wedge T_n}$ and $\wt M_{t\wedge\tau_n}$ are martingales, and hence so is
\[
  K_t\defeq N_{t\wedge\tau_n}=M_{t\wedge\tau_n}-\wt M_{t\wedge\tau_n},\qquad 0\le t<\infty
\]
(we suppress the dependence on $n$). Fix $t>0$. For any partition $0=s_0<\dots<s_p=t$, the points $s_j\wedge\tau_n$ form a nondecreasing sequence in $[0,t\wedge\tau_n]$, so by Step 1 and \eqref{eq:3-3-02:b},
\begin{equation}\label{eq:3-3-02:d}
  \sum_{j=1}^p\abs{K_{s_j}-K_{s_{j-1}}}
  =\sum_{j=1}^p\abs{N_{s_j\wedge\tau_n}-N_{s_{j-1}\wedge\tau_n}}
  \le V_{t\wedge\tau_n}\le V_{t\wedge S_n}\le n .
\end{equation}
In particular $\abs{K_s-K_0}\le n$ for all $s\ge0$, and since $K_0=N_0=0$ a.s., the continuous martingale $K$ satisfies $\abs{K_s}\le n$ a.s.\ for every $s\ge 0$.

\emph{Step 3: $K_t=0$ a.s.} Let $s_j\defeq tj2^{-m}$, $0\le j\le 2^m$. For $i<j$, the bounded random variable $K_{s_i}-K_{s_{i-1}}$ is $\F_{s_{j-1}}$-measurable and $E[K_{s_j}-K_{s_{j-1}}\mid\F_{s_{j-1}}]=0$, so the increments of $K$ over the partition are orthogonal in $L^2$. Using $K_0=0$ a.s.\ and \eqref{eq:3-3-02:d}, we obtain
\[
  E\bigl(K_t^2\bigr)=E\Bigl[\Bigl(\sum_{j=1}^{2^m}(K_{s_j}-K_{s_{j-1}})\Bigr)^{2}\Bigr]
  =E\sum_{j=1}^{2^m}(K_{s_j}-K_{s_{j-1}})^2
  \le n\,E\Bigl[\max_{1\le j\le 2^m}\abs{K_{s_j}-K_{s_{j-1}}}\Bigr].
\]
The maximum on the right is bounded by $n$ (by \eqref{eq:3-3-02:d}) and tends to zero as $m\to\infty$ for every $\omega$, because the path $s\mapsto K_s(\omega)$ is uniformly continuous on $[0,t]$. By the bounded convergence theorem the right-hand side tends to zero, whereas the left-hand side does not depend on $m$. Hence $E(K_t^2)=0$, that is,
\[
  N_{t\wedge\tau_n}=0\quad\text{a.s., for every } t\ge0 \text{ and } n\ge 1 .
\]

\emph{Step 4: conclusion.} By \eqref{eq:3-3-02:c} and the choice of $\{T_n\}$, $\{\wt T_n\}$, we have $\tau_n\to\infty$ a.s., so for fixed $t$ and a.e.\ $\omega$ we have $t\wedge\tau_n(\omega)=t$ for all large $n$; thus $N_t=\lim_{n\to\infty}N_{t\wedge\tau_n}=0$ a.s. Consequently $P[N_q=0;\ \forall\, q\in\Q\cap[0,\infty)]=1$, and by the continuity of the paths of $N$, $P[N_t=0;\ \forall\,0\le t<\infty]=1$. As observed at the beginning, this gives $M_t=\wt M_t$ and $B_t=\wt B_t$ for all $0\le t<\infty$, $P$-a.s.

\rmk{The heart of the argument is that a continuous local martingale which starts at zero and has paths of finite variation on bounded intervals is indistinguishable from zero. This is also the content of \bk{Exercise 1.5.21} (which rests on \bk{Theorem 1.5.8} and \bk{Problems 1.5.11, 1.5.12}); the direct computation of Step 3 makes the proof independent of the theory of quadratic variation.}
\end{kex}

\begin{kex}{Problem}{3.7}{}
The book's Section~3.8 prints only the Taylor expansion that replaces the one in Step~2 of the proof
of \bk{Theorem 3.3}; the complete proof (localization, the convergence of every sum, and the passage
to indistinguishability) is given here.

\emph{The statement and its meaning.} The filtration $\{\F_t\}$ satisfies the usual conditions.
We are given $M^{(1)},\dots,M^{(d)}\in\cM^{c,\loc}$; processes $B^{(1)},\dots,B^{(d)}$ as in
\bk{Definition 3.1}, i.e.\ $B^{(i)}=A^{(i),+}-A^{(i),-}$ with $A^{(i),\pm}$ continuous, nondecreasing,
adapted, null at $0$, the total variation of $B^{(i)}$ on $[0,t]$ being
$\check B^{(i)}_t=A^{(i),+}_t+A^{(i),-}_t$; an $\F_0$-measurable random vector $X_0$; and
$X_t=X_0+M_t+B_t$. The function $f\colon[0,\infty)\times\R^d\to\R$ is of class $C^{1,2}$: the
derivatives $\partial f/\partial t$ (one-sided at $t=0$), $\partial f/\partial x_i$ and
$\partial^2 f/\partial x_i\partial x_j$ exist and are continuous on $[0,\infty)\times\R^d$. We write
\[
  f_t\defeq\frac{\partial f}{\partial t},\qquad f_i\defeq\frac{\partial f}{\partial x_i},\qquad
  f_{ij}\defeq\frac{\partial^2 f}{\partial x_i\partial x_j},
\]
and we have to prove that, a.s.,
\begin{equation}\label{eq:3-3-07:ito}
\begin{split}
  f(t,X_t)={}& f(0,X_0)+\int_0^t f_t(s,X_s)\,ds+\sum_{i=1}^d\int_0^t f_i(s,X_s)\,dB^{(i)}_s
  +\sum_{i=1}^d\int_0^t f_i(s,X_s)\,dM^{(i)}_s\\
  &+\frac12\sum_{i=1}^d\sum_{j=1}^d\int_0^t f_{ij}(s,X_s)\,d\qv{M^{(i)},M^{(j)}}_s,
  \qquad 0\le t<\infty .
\end{split}
\end{equation}
The exceptional null sets of the hypotheses belong to $\F_0$ (usual conditions); setting $M$, $B$,
$A^{(i),\pm}$, $\qv{M^{(i)}}$ and $\qv{M^{(i)},M^{(j)}}$ equal to zero on them changes neither
adaptedness nor \eqref{eq:3-3-07:ito}, so we assume that for \emph{every} $\omega$ these processes
are continuous and start at zero, that $\qv{M^{(i)}}$ is nondecreasing and that
$\qv{M^{(i)},M^{(j)}}$ is of bounded variation on bounded intervals. As in \bk{Remark 3.4}, all
terms of \eqref{eq:3-3-07:ito} make sense: for a continuous $g$ on $[0,\infty)\times\R^d$ the process
$g(s,X_s)$ is continuous and adapted, hence progressively measurable, and its paths are bounded on
every $[0,t]$; so it belongs to $\cP^*(M^{(i)})$ and $\int_0^t g(s,X_s)\,dM^{(i)}_s$ is a continuous
local martingale (\bk{Definition 2.23}), whereas the integrals with respect to $ds$, $dB^{(i)}_s$ and
$d\qv{M^{(i)},M^{(j)}}_s$ are Lebesgue--Stieltjes integrals of bounded Borel functions with respect
to continuous functions of bounded variation and are therefore continuous functions of $t$.

For a stopping time $T$ and a process $Z$ we write $Z^T_t\defeq Z_{t\wedge T}$.

\emph{Step 1. Four facts about stopping.} Let $T$ be a stopping time of $\{\F_t\}$.

\textbf{(a)} \emph{If $N\in\cM^{c,\loc}$, then $N^T\in\cM^{c,\loc}$.} Indeed, if $\{S_k\}$ is a
sequence as in \bk{Definition 1.5.15} for $N$, then $(N^T)^{S_k}=(N^{S_k})^T$ is a martingale by
the optional sampling theorem (\bk{Problem 1.3.24}), and $N^T_0=0$.

\textbf{(b)} \emph{If $M,N\in\cM^{c,\loc}$, then $\qv{M^T,N^T}=\qv{M,N}^T$; in particular
$\qv{M^T}=\qv{M}^T$.} Indeed, $\qv{M,N}^T$ is adapted, continuous, of bounded variation and null at
$0$, and $M^TN^T-\qv{M,N}^T=(MN-\qv{M,N})^T$ belongs to $\cM^{c,\loc}$ by (a); the claim follows from
the uniqueness assertion of \bk{Problem 1.5.17}.

\textbf{(c)} \emph{Let $M\in\cM^{c,\loc}$ be such that $M^T\in\cM_2^c$, let $Y\in\cP^*(M)$ and
$Z\in\cL^*(M^T)$, and suppose that for $P$-a.e.\ $\omega$ we have $Y_s(\omega)=Z_s(\omega)$ for all
$s$ with $0<s\le T(\omega)$, $s<\infty$. Then for $P$-a.e.\ $\omega$,}
\begin{equation}\label{eq:3-3-07:loc}
  \int_0^t Y_s\,dM_s=\int_0^t Z_s\,dM^T_s\qquad\text{for all } 0\le t\le T(\omega),\ t<\infty,
\end{equation}
\emph{where the left-hand side is the integral of \bk{Definition 2.23} and the right-hand side the
integral of \bk{Definition 2.9}.}
To see this, recall the construction preceding \bk{Definition 2.23}: there are stopping times
$T_m'\uparrow\infty$ a.s.\ (denoted $T_n$ in \bk{(2.32)}) such that $M^{(m)}\defeq M^{T_m'}\in\cM_2^c$,
$Y^{(m)}_s\defeq Y_s\ind_{\{T_m'\ge s\}}$ belongs to $\cL^*(M^{(m)})$, and, a.s.,
\begin{equation}\label{eq:3-3-07:def}
  I^M_t(Y)=I^{M^{(m)}}_t(Y^{(m)})\qquad\text{for } 0\le t\le T_m' .
\end{equation}
Put $Z'_s\defeq Z_s\ind_{(0,\infty)}(s)+Y_0\ind_{\{0\}}(s)$. This process is progressively
measurable ($Y_0$ is $\F_0$-measurable), and $\int_0^u (Z'_s-Z_s)^2\,d\qv{M^T}_s=0$ for all $u$
because the continuous function $\qv{M^T}$ does not charge $\{0\}$. Hence $Z'\in\cL^*(M^T)$, $Z'$ and
$Z$ are equivalent in $\cL^*(M^T)$, and $I^{M^T}(Z')=I^{M^T}(Z)$ by \bk{(2.15)} and \bk{(2.17)}
(\bk{Proposition 2.10}). With $\tau_m\defeq T\wedge T_m'$ we have, for all $t$,
$M^{(m)}_{t\wedge\tau_m}=M_{t\wedge\tau_m}=M^T_{t\wedge\tau_m}$, and a.s.\
$Y^{(m)}_{t\wedge\tau_m}=Y_{t\wedge\tau_m}=Z'_{t\wedge\tau_m}$: if $t\wedge\tau_m=0$ both sides are
$Y_0$, and otherwise $0<t\wedge\tau_m\le T$. \bk{Corollary 2.21} now gives, a.s.,
$I^{M^{(m)}}_{t\wedge\tau_m}(Y^{(m)})=I^{M^T}_{t\wedge\tau_m}(Z')=I^{M^T}_{t\wedge\tau_m}(Z)$ for all
$t$. Together with \eqref{eq:3-3-07:def} this is \eqref{eq:3-3-07:loc} for
$0\le t\le T\wedge T_m'$, and we let $m\to\infty$.

\textbf{(d)} \emph{Consistency of the two definitions.} Let $M\in\cM_2^c$ and $Y\in\cL^*(M)$. Then
$Y\in\cP^*(M)$: since $E\int_0^u Y_s^2\,d\qv{M}_s<\infty$ for every integer $u$, a.s.\
$\int_0^u Y_s^2\,d\qv{M}_s<\infty$ for all $u<\infty$. Taking $T\equiv\infty$ (so that $M^T=M$) and
$Z=Y$ in (c), we see that the integral of $Y$ with respect to $M$ in the sense of
\bk{Definition 2.23} is indistinguishable from the one in the sense of \bk{Definition 2.9}. For
integrators in $\cM_2^c$ and integrands in $\cL^*$ we may therefore use either definition, and the
properties \bk{(2.5)}, \bk{(2.14)}, \bk{(2.17)} of the integral of \bk{Definition 2.9} are
available.

\emph{Step 2. Reduction to the bounded case.} We shall prove in Steps 3--6 the following.

\begin{claim}
Suppose that, for some constant $K<\infty$, we have a.s.
\begin{equation}\label{eq:3-3-07:bdd}
  \abs{X^{(i)}_0}\le K,\qquad \abs{M^{(i)}_t}\le K,\qquad \check B^{(i)}_t\le K,\qquad
  \qv{M^{(i)}}_t\le K;\qquad 0\le t<\infty,\ 1\le i\le d,
\end{equation}
and that each $M^{(i)}$ is a martingale. Then \eqref{eq:3-3-07:ito} holds a.s.
\end{claim}

Under \eqref{eq:3-3-07:bdd} each $M^{(i)}$ belongs to $\cM_2^c$ and each integrand $f_i(s,X_s)$ to
$\cL^*(M^{(i)})$ (this is verified in Steps~3 and~4), so by Step~1(d) it is immaterial whether the
stochastic integrals in the Claim are read in the sense of \bk{Definition 2.23} or of
\bk{Definition 2.9}; in Steps~3--6 we work with \bk{Definition 2.9}.

Granting the Claim, we treat the general case. For $n\ge1$ let
\begin{gather*}
  \zeta_t\defeq\max_{1\le i\le d}\max\bigl\{\abs{M^{(i)}_t},\,\check B^{(i)}_t,\,\qv{M^{(i)}}_t\bigr\},
  \qquad S_n\defeq\inf\{t\ge0;\ \zeta_t\ge n\},\\
  \Omega_n\defeq\Bigl\{\max_{1\le i\le d}\abs{X^{(i)}_0}<n\Bigr\},\qquad
  T_n\defeq S_n\ind_{\Omega_n}
\end{gather*}
(with $\inf\emptyset=\infty$ and $\infty\cdot 0=0$). The process $\zeta$ is continuous and adapted
with $\zeta_0=0$, so $S_n$ is a stopping time (\bk{Problem 1.2.7}), and so is $T_n$, because
$\{T_n\le t\}=\Omega_n^c\cup\{S_n\le t\}\in\F_t$. Since the paths of $\zeta$ are bounded on bounded
intervals, $S_n\uparrow\infty$; and $\omega\in\Omega_n$ for all large $n$. Define
\[
  X^{[n]}_0\defeq X_0\ind_{\Omega_n},\qquad M^{[n]}\defeq M^{T_n},\qquad B^{[n]}\defeq B^{T_n},\qquad
  X^{[n]}_t\defeq X^{[n]}_0+M^{[n]}_t+B^{[n]}_t
\]
(componentwise). By continuity $\zeta_{t\wedge T_n}\le n$ for all $t$. Hence the components of
$M^{[n]}$ are bounded by $n$; they belong to $\cM^{c,\loc}$ by Step~1(a), and a bounded local
martingale is of class $DL$, hence a martingale (\bk{Problem 1.5.19}(i)). The components of $B^{[n]}$
are adapted (\bk{Proposition 1.2.18}), continuous, with decomposition
$(A^{(i),+})^{T_n}-(A^{(i),-})^{T_n}$ as in \bk{Definition 3.1} and total variation
$\check B^{(i)}_{t\wedge T_n}\le n$; by Step~1(b),
$\qv{M^{[n],(i)}}_t=\qv{M^{(i)}}_{t\wedge T_n}\le n$; and $X^{[n]}_0$ is $\F_0$-measurable with
components bounded by $n$. Thus \eqref{eq:3-3-07:bdd} holds with $K=n$, and the Claim yields a null
set off which, for all $0\le t<\infty$,
\begin{equation}\label{eq:3-3-07:iton}
\begin{split}
  f(t,X^{[n]}_t)={}& f(0,X^{[n]}_0)+\int_0^t f_t(s,X^{[n]}_s)\,ds
  +\sum_{i=1}^d\int_0^t f_i(s,X^{[n]}_s)\,dB^{[n],(i)}_s
  +\sum_{i=1}^d\int_0^t f_i(s,X^{[n]}_s)\,dM^{[n],(i)}_s\\
  &+\frac12\sum_{i=1}^d\sum_{j=1}^d\int_0^t f_{ij}(s,X^{[n]}_s)\,d\qv{M^{[n],(i)},M^{[n],(j)}}_s .
\end{split}
\end{equation}
Let $\omega\in\Omega_n$ and $0\le t\le T_n(\omega)$, $t<\infty$. Then $X^{[n]}_s(\omega)=X_s(\omega)$,
$B^{[n]}_s(\omega)=B_s(\omega)$ for $0\le s\le t$, and
$\qv{M^{[n],(i)},M^{[n],(j)}}_s=\qv{M^{(i)},M^{(j)}}_s$ for $0\le s\le t$ by Step~1(b) (a.s.,
simultaneously for all $s$, both sides being continuous). So the left-hand side, the first three
terms on the right and the last term of \eqref{eq:3-3-07:iton} coincide with the corresponding terms of
\eqref{eq:3-3-07:ito}. For the stochastic integrals we apply Step~1(c) with $M=M^{(i)}$, $T=T_n$,
$Y_s=f_i(s,X_s)$ and $Z_s=f_i(s,X^{[n]}_s)$: here $Z$ is continuous, adapted and bounded on each
$[0,u]\times\Omega$ (the components of $X^{[n]}$ are bounded by $3n$), so $Z\in\cL^*(M^{T_n})$; and
$Y_s=Z_s$ for $0<s\le T_n(\omega)$ for every $\omega$, since this is the above remark when
$\omega\in\Omega_n$ and there is no such $s$ when $\omega\notin\Omega_n$. By
\eqref{eq:3-3-07:loc}, whose right-hand side is the fourth term of \eqref{eq:3-3-07:iton} read in
the sense of \bk{Definition 2.9}, the fourth terms coincide as well, off a null set. Consequently, off a null set
$\Lambda_n$, \eqref{eq:3-3-07:ito} holds for all $\omega\in\Omega_n$ and all finite $t\le T_n(\omega)$.
For $\omega\notin\bigcup_n\Lambda_n$ and $t\in[0,\infty)$ we choose $n$ with $\omega\in\Omega_n$ and
$T_n(\omega)=S_n(\omega)\ge t$, and obtain \eqref{eq:3-3-07:ito}. This proves the theorem, once the
Claim is established.

\emph{Step 3. Proof of the Claim: Taylor expansion.} We assume \eqref{eq:3-3-07:bdd} and, modifying
the processes on a null set of $\F_0$ as before, that it holds for every $\omega$. Then
$M^{(i)}\in\cM_2^c$, and $X_s\in Q\defeq[-3K,3K]^d$ for all $s$. On $\cM_2^c$ the cross-variation of
\bk{Problem 1.5.17} coincides with that of \bk{Definition 1.5.5}, by the uniqueness in
\bk{Problem 1.5.17}. Since both sides of \eqref{eq:3-3-07:ito} are continuous in $t$ and agree at
$t=0$, it suffices (\bk{Problem 1.1.5}) to prove \eqref{eq:3-3-07:ito} a.s.\ for each fixed $t>0$.
Fix $t>0$. The functions $f_t,f_i,f_{ij}$ are bounded on the compact set $[0,t]\times Q$, say by
$C$, and uniformly continuous there; so the modulus
\[
  \rho(\delta)\defeq\max_{g}\,\sup\bigl\{\abs{g(s,x)-g(s',x')};\ s,s'\in[0,t],\ x,x'\in Q,\
  \abs{s-s'}\le\delta,\ \abs{x-x'}\le\delta\bigr\},
\]
the maximum being over $g\in\{f_t,f_i,f_{ij};\,1\le i,j\le d\}$, is nondecreasing in $\delta$, bounded
by $2C$, and $\rho(\delta)\to0$ as $\delta\downarrow0$. For a continuous process $Z$ (possibly
vector-valued) let
\[
  m(Z;\delta)\defeq\sup\{\abs{Z_r-Z_s};\ r,s\in[0,t],\ \abs{r-s}\le\delta\},
\]
which is a random variable (the supremum may be restricted to rational $r,s$) and tends to $0$ as
$\delta\downarrow0$ for every $\omega$. For a partition $\Pi=\{t_0,\dots,t_m\}$,
$0=t_0<t_1<\dots<t_m=t$, with mesh $\norm{\Pi}=\max_k(t_k-t_{k-1})$, we write
$\Delta_kZ\defeq Z_{t_k}-Z_{t_{k-1}}$ and
\[
  \kappa(\Pi)\defeq\rho\bigl(\max\{\norm{\Pi},\,m(X;\norm{\Pi})\}\bigr),\qquad
  0\le\kappa(\Pi)\le 2C .
\]
Then, for $g\in\{f_t,f_i,f_{ij}\}$, $1\le k\le m$ and $s\in[t_{k-1},t_k]$,
\begin{equation}\label{eq:3-3-07:kappa}
  \abs{g(s,X_s)-g(t_{k-1},X_{t_{k-1}})}\le\kappa(\Pi),\qquad
  \abs{g(s,X_{t_k})-g(s,X_s)}\le\kappa(\Pi).
\end{equation}

Let $s\in[0,t]$ and $x,y\in Q$. The function $\varphi(\theta)\defeq f(s,x+\theta(y-x))$ is of class
$C^2$ on $[0,1]$, and Taylor's formula with integral remainder,
$\varphi(1)=\varphi(0)+\varphi'(0)+\int_0^1(1-\theta)\varphi''(\theta)\,d\theta$, gives
\begin{gather}
  f(s,y)-f(s,x)=\sum_{i=1}^d f_i(s,x)(y_i-x_i)
  +\frac12\sum_{i=1}^d\sum_{j=1}^d f_{ij}(s,x)(y_i-x_i)(y_j-x_j)+r(s,x,y),\label{eq:3-3-07:taylor}\\
  r(s,x,y)\defeq\sum_{i=1}^d\sum_{j=1}^d(y_i-x_i)(y_j-x_j)\int_0^1(1-\theta)
  \bigl[f_{ij}(s,x+\theta(y-x))-f_{ij}(s,x)\bigr]\,d\theta .\notag
\end{gather}
Since $Q$ is convex, $x+\theta(y-x)\in Q$, and therefore
\begin{equation}\label{eq:3-3-07:rem}
  \abs{r(s,x,y)}\le\frac12\,\rho(\abs{y-x})\Bigl(\sum_{i=1}^d\abs{y_i-x_i}\Bigr)^2
  \le\frac d2\,\rho(\abs{y-x})\,\abs{y-x}^2 .
\end{equation}
The function $r$ is continuous, being defined by \eqref{eq:3-3-07:taylor} as a difference of
continuous functions. Now, as in the book's Section~3.8, we split
\[
  f(t_k,X_{t_k})-f(t_{k-1},X_{t_{k-1}})
  =\bigl[f(t_k,X_{t_k})-f(t_{k-1},X_{t_k})\bigr]+\bigl[f(t_{k-1},X_{t_k})-f(t_{k-1},X_{t_{k-1}})\bigr],
\]
write the first bracket as $\int_{t_{k-1}}^{t_k}f_t(s,X_{t_k})\,ds$ (fundamental theorem of
calculus), apply \eqref{eq:3-3-07:taylor} to the second, substitute
$\Delta_kX^{(i)}=\Delta_kM^{(i)}+\Delta_kB^{(i)}$ in the first-order terms, and sum over $k$:
\begin{equation}\label{eq:3-3-07:sum}
  f(t,X_t)-f(0,X_0)=J_0(\Pi)+\sum_{i=1}^d\bigl[J_1^{i}(\Pi)+J_2^{i}(\Pi)\bigr]
  +\frac12\sum_{i=1}^d\sum_{j=1}^dJ_3^{ij}(\Pi)+R(\Pi),
\end{equation}
where
\begin{align*}
  J_0(\Pi)&\defeq\sum_{k=1}^m\int_{t_{k-1}}^{t_k}f_t(s,X_{t_k})\,ds, &
  J_1^{i}(\Pi)&\defeq\sum_{k=1}^m f_i(t_{k-1},X_{t_{k-1}})\,\Delta_kB^{(i)},\\
  J_2^{i}(\Pi)&\defeq\sum_{k=1}^m f_i(t_{k-1},X_{t_{k-1}})\,\Delta_kM^{(i)}, &
  J_3^{ij}(\Pi)&\defeq\sum_{k=1}^m f_{ij}(t_{k-1},X_{t_{k-1}})\,\Delta_kX^{(i)}\Delta_kX^{(j)},
\end{align*}
and $R(\Pi)\defeq\sum_{k=1}^m r(t_{k-1},X_{t_{k-1}},X_{t_k})$. All these are random variables. We fix
a sequence $\{\Pi^{(n)}\}_{n=1}^\infty$ of partitions of $[0,t]$ with $\norm{\Pi^{(n)}}\to0$. Then
$\kappa(\Pi^{(n)})\to0$ for every $\omega$, and by bounded convergence
\begin{equation}\label{eq:3-3-07:kap0}
  \lim_{n\to\infty}E\bigl[\kappa(\Pi^{(n)})^2\bigr]=0 .
\end{equation}
We show that each term on the right of \eqref{eq:3-3-07:sum}, evaluated along $\Pi^{(n)}$, converges
in probability to the corresponding term of \eqref{eq:3-3-07:ito}, and that $R(\Pi^{(n)})\to0$ in
probability. Since the left-hand side of \eqref{eq:3-3-07:sum} does not depend on $n$, this proves
\eqref{eq:3-3-07:ito} a.s.\ for our fixed $t$.

\emph{Step 4. The terms $J_0$, $J_1$, $J_2$ and $R$.} By \eqref{eq:3-3-07:kappa},
\[
  \Bigl\lvert J_0(\Pi)-\int_0^tf_t(s,X_s)\,ds\Bigr\rvert
  \le\sum_{k=1}^m\int_{t_{k-1}}^{t_k}\abs{f_t(s,X_{t_k})-f_t(s,X_s)}\,ds\le t\,\kappa(\Pi),
\]
and, since $\Delta_kB^{(i)}=\int_{(t_{k-1},t_k]}dB^{(i)}_s$,
\[
  \Bigl\lvert J_1^{i}(\Pi)-\int_0^tf_i(s,X_s)\,dB^{(i)}_s\Bigr\rvert
  =\Bigl\lvert\sum_{k=1}^m\int_{(t_{k-1},t_k]}\bigl[f_i(t_{k-1},X_{t_{k-1}})-f_i(s,X_s)\bigr]\,dB^{(i)}_s
  \Bigr\rvert\le\kappa(\Pi)\,\check B^{(i)}_t\le K\kappa(\Pi).
\]
Both bounds tend to zero along $\Pi^{(n)}$, for every $\omega$.

For $J_2^i$, let $Y_s\defeq f_i(s,X_s)$, $0\le s<\infty$. This process is continuous and adapted,
and bounded on $[0,u]\times\Omega$ for every $u$ (as $f_i$ is bounded on $[0,u]\times Q$); with
$\qv{M^{(i)}}\le K$ this shows $Y\in\cL^*(M^{(i)})$. The process
\[
  Y^\Pi_s\defeq f_i(0,X_0)\ind_{\{0\}}(s)+\sum_{k=1}^mf_i(t_{k-1},X_{t_{k-1}})\ind_{(t_{k-1},t_k]}(s),
  \qquad 0\le s<\infty,
\]
is simple (extend the partition by $t_{m+l}\defeq t+l$, with coefficients $0$; the coefficients
are bounded by $C$ and $f_i(t_{k-1},X_{t_{k-1}})$ is $\F_{t_{k-1}}$-measurable), and by \bk{(2.5)},
$I^{M^{(i)}}_t(Y^\Pi)=J_2^i(\Pi)$. By \bk{(2.17)}, \bk{(2.14)} (valid on $\cL^*$ by
\bk{Proposition 2.10}) and \eqref{eq:3-3-07:kappa},
\[
  E\Bigl[J_2^i(\Pi)-\int_0^tf_i(s,X_s)\,dM^{(i)}_s\Bigr]^2
  =E\int_0^t\abs{Y^\Pi_s-Y_s}^2\,d\qv{M^{(i)}}_s\le K\,E\bigl[\kappa(\Pi)^2\bigr],
\]
which tends to zero along $\Pi^{(n)}$ by \eqref{eq:3-3-07:kap0}.

For the remainder, write $V^{i}(\Pi)\defeq\sum_{k=1}^m(\Delta_kM^{(i)})^2$. Since
$\sum_k(\Delta_kB^{(i)})^2\le(\sum_k\abs{\Delta_kB^{(i)}})^2\le K^2$, we have
$\sum_k\abs{\Delta_kX}^2\le2\sum_{i=1}^d\bigl[V^i(\Pi)+K^2\bigr]$; and
$\rho(\abs{\Delta_kX})\le\kappa(\Pi)$ because $\abs{\Delta_kX}\le m(X;\norm{\Pi})$. Thus
\eqref{eq:3-3-07:rem}, the Cauchy--Schwarz inequality and \bk{Lemma 1.5.9}
($E[V^i(\Pi)]^2\le6K^4$) give
\[
  E\abs{R(\Pi)}\le d\sum_{i=1}^dE\bigl[\kappa(\Pi)\bigl(V^i(\Pi)+K^2\bigr)\bigr]
  \le d^2\bigl(\sqrt6+1\bigr)K^2\sqrt{E\bigl[\kappa(\Pi)^2\bigr]},
\]
which tends to zero along $\Pi^{(n)}$ by \eqref{eq:3-3-07:kap0}.

\emph{Step 5. A quadratic variation lemma.} Fix $i,j$ and put $h_s\defeq f_{ij}(s,X_s)$, so that
$\abs{h_s}\le C$ for $0\le s\le t$. \emph{Let $N\in\cM_2^c$ satisfy $\abs{N_s}\le L$ and
$\qv{N}_s\le L$ for all $s\in[0,t]$, a.s. Then}
\begin{equation}\label{eq:3-3-07:qv}
  \sum_{k}h_{t_{k-1}}(\Delta_kN)^2\ \longrightarrow\ \int_0^th_s\,d\qv{N}_s
  \qquad\text{in } L^1(\Omega,\F,P),\ \text{along } \Pi^{(n)} .
\end{equation}
Put $D_k\defeq(\Delta_kN)^2-\Delta_k\qv{N}$; these random variables are bounded, $D_k$ is
$\F_{t_k}$-measurable, and $E[D_k\mid\F_{t_{k-1}}]=0$ because $N$ and $N^2-\qv{N}$ are martingales
(the discussion preceding \bk{Lemma 1.5.9}). For $k<l$ the variable $h_{t_{k-1}}D_kh_{t_{l-1}}$ is
bounded and $\F_{t_{l-1}}$-measurable, so $E[h_{t_{k-1}}D_kh_{t_{l-1}}D_l]=0$, and consequently
\begin{align*}
  E\Bigl[\sum_{k=1}^mh_{t_{k-1}}D_k\Bigr]^2=\sum_{k=1}^mE\bigl[h_{t_{k-1}}^2D_k^2\bigr]
  &\le2C^2E\Bigl[\sum_{k=1}^m(\Delta_kN)^4+\sum_{k=1}^m(\Delta_k\qv{N})^2\Bigr]\\
  &\le2C^2\Bigl(E\,V^{(4)}_t(\Pi)+L\,E\bigl[m(\qv{N};\norm{\Pi})\bigr]\Bigr),
\end{align*}
where $V^{(4)}_t(\Pi)=\sum_k(\Delta_kN)^4$. Along $\Pi^{(n)}$ the first expectation tends to zero by
\bk{Lemma 1.5.10}, and the second by bounded convergence, since $\qv{N}$ is continuous and
$m(\qv{N};\cdot)\le L$. Moreover, by \eqref{eq:3-3-07:kappa} and the monotonicity of $\qv{N}$,
\[
  \Bigl\lvert\sum_{k=1}^mh_{t_{k-1}}\Delta_k\qv{N}-\int_0^th_s\,d\qv{N}_s\Bigr\rvert
  =\Bigl\lvert\sum_{k=1}^m\int_{(t_{k-1},t_k]}\bigl[h_{t_{k-1}}-h_s\bigr]\,d\qv{N}_s\Bigr\rvert
  \le L\,\kappa(\Pi),
\]
which tends to zero along $\Pi^{(n)}$ a.s.\ and boundedly, hence in $L^1$. Adding the two
statements we obtain \eqref{eq:3-3-07:qv}.

\emph{Step 6. The terms $J_3^{ij}$.} Expanding
$\Delta_kX^{(i)}\Delta_kX^{(j)}$ with $\Delta_kX=\Delta_kM+\Delta_kB$, we get
\begin{gather*}
  J_3^{ij}(\Pi)=\sum_{k=1}^mh_{t_{k-1}}\Delta_kM^{(i)}\Delta_kM^{(j)}+J_4^{ij}(\Pi),\\
  J_4^{ij}(\Pi)\defeq\sum_{k=1}^mh_{t_{k-1}}\bigl[\Delta_kM^{(i)}\Delta_kB^{(j)}
  +\Delta_kB^{(i)}\Delta_kM^{(j)}+\Delta_kB^{(i)}\Delta_kB^{(j)}\bigr].
\end{gather*}
Each of the three sums in $J_4^{ij}$ contains an increment of some $B^{(l)}$, whose absolute values
add up to at most $\check B^{(l)}_t\le K$; bounding the other factor by its maximum over $k$ we find
\[
  \abs{J_4^{ij}(\Pi)}\le CK\bigl[m(M^{(i)};\norm{\Pi})+m(M^{(j)};\norm{\Pi})
  +m(B^{(i)};\norm{\Pi})\bigr],
\]
which tends to zero along $\Pi^{(n)}$ for every $\omega$, by the continuity of $M$ and $B$.

For the first sum we polarize. The martingales $N^{\pm}\defeq M^{(i)}\pm M^{(j)}$ belong to
$\cM_2^c$ and are bounded by $2K$. By \bk{Problem 1.5.7}(i),(ii),
$\qv{N^\pm}=\qv{M^{(i)}}\pm2\qv{M^{(i)},M^{(j)}}+\qv{M^{(j)}}$, so that
$\qv{N^+}+\qv{N^-}=2\qv{M^{(i)}}+2\qv{M^{(j)}}\le4K$, and since both summands are nonnegative,
$\qv{N^\pm}\le4K$. Step~5 applies to $N^+$ and $N^-$ with $L=4K$. Since
$\Delta_kM^{(i)}\Delta_kM^{(j)}=\frac14[(\Delta_kN^+)^2-(\Delta_kN^-)^2]$ and
$\qv{M^{(i)},M^{(j)}}=\frac14[\qv{N^+}-\qv{N^-}]$ (\bk{Definition 1.5.5}), we conclude from
\eqref{eq:3-3-07:qv} that, along $\Pi^{(n)}$,
in $L^1(\Omega,\F,P)$,
\[
  \sum_{k}h_{t_{k-1}}\Delta_kM^{(i)}\Delta_kM^{(j)}\ \longrightarrow\
  \frac14\Bigl[\int_0^th_s\,d\qv{N^+}_s-\int_0^th_s\,d\qv{N^-}_s\Bigr]
  =\int_0^tf_{ij}(s,X_s)\,d\qv{M^{(i)},M^{(j)}}_s .
\]
Hence $J_3^{ij}(\Pi^{(n)})\to\int_0^tf_{ij}(s,X_s)\,d\qv{M^{(i)},M^{(j)}}_s$ in probability.

Steps 4 and 6 supply all the convergences announced at the end of Step~3. This proves the Claim, and
with it \bk{Theorem 3.6}.

\rmk{The mixed terms $M$--$B$ and $B$--$B$ of the second-order sum vanish in the limit
($J_4^{ij}\to0$), which is why only the cross-variations $\qv{M^{(i)},M^{(j)}}$ appear in
\bk{(3.5)}. Using the integral form of the Taylor remainder instead of an intermediate point
$\eta_k$, as here, avoids the measurable selection of $\eta_k$ needed in Step~2 of the proof of
\bk{Theorem 3.3}.}
\end{kex}

\begin{kex}{Problem}{3.10}{}
We recall the setting of \bk{Example 3.9}: $W=\{W_t,\F_t;\ 0\le t<\infty\}$ is a standard Brownian motion, $\{\F_t\}$ satisfies the usual conditions, $X\in\cP$ (\bk{Definition 2.22} with $M=W$, $\qv{W}_t=t$), that is, $X$ is measurable, adapted and $\int_0^T X_s^2\,ds<\infty$ a.s.\ for every $T$, and
\[
  Z_t=e^{\zeta_t},\qquad \zeta_t=N_t+B_t,\qquad
  N_t\defeq\int_0^t X_s\,dW_s,\qquad B_t\defeq-\frac12\int_0^t X_s^2\,ds .
\]
In integral form, the assertion to be proved is that $Y_tX_t$ is an admissible integrand for $W$, that $\int_0^t\abs{Y_s}X_s^2\,ds<\infty$, and that, a.s.,
\begin{equation}\label{eq:3-3-10:a}
  Y_t=1+\int_0^t Y_sX_s^2\,ds-\int_0^t Y_sX_s\,dW_s;\qquad 0\le t<\infty .
\end{equation}
As the book does throughout \bk{Example 3.9}, we use the results of \bk{Section 3.2} for the integrator $W$ with integrands in the class $\cP$ rather than $\cP^*$; this is legitimate by \bk{Remark 2.11} and the paragraph following \bk{Definition 2.22}, because $t\mapsto\qv{W}_t=t$ is absolutely continuous.

\emph{Step 1. The processes involved.}
The stochastic integral $N$ belongs to $\cM^{c,\loc}$ and, by \bk{(2.19)} (a sample path property, which remains valid for integrands in $\cP$ by localization), $\qv{N}_t=\int_0^tX_s^2\,ds$. Since $X$ is measurable, each $\int_0^nX_s^2\,ds$ is an $\F$-measurable, $[0,\infty]$-valued random variable (Fubini--Tonelli). Let $\Omega_0\in\F$ be the event on which $\int_0^nX_s^2\,ds<\infty$ for every $n\in\N$; then $P(\Omega_0)=1$, and $\Omega_0\in\F_0$ because $\F_0$ contains the $P$-null sets. On $\Omega_0$ the function $t\mapsto B_t$ is continuous, nonincreasing and vanishes at $t=0$; we set $B\equiv0$ on $\Omega\setminus\Omega_0$, that is,
\[
  B_t=-\frac12\,\ind_{\Omega_0}\int_0^tX_s^2\,ds;\qquad0\le t<\infty .
\]
We check that $B$ is adapted. \bk{Remark 2.11} establishes the $\F_t$-measurability of $\int_0^tX_s^2\,ds$ only for $X\in\cL$, whereas our $X$ need not satisfy $E\int_0^TX_s^2\,ds<\infty$; we therefore truncate. For $k\in\N$ the process $X^{(k)}\defeq X\ind_{\{\abs{X}\le k\}}$ is measurable, adapted and bounded, so by \bk{Remark 2.11} (the case of Brownian motion and a bounded integrand, applied on $[0,T]$ for each $T<\infty$) the random variable $\int_0^t(X^{(k)}_s)^2\,ds$ is $\F_t$-measurable for every $t\ge0$. Because $(X^{(k)}_s)^2\uparrow X_s^2$ as $k\to\infty$, the monotone convergence theorem gives, for every $\omega\in\Omega$,
\[
  \int_0^tX_s^2\,ds=\lim_{k\to\infty}\int_0^t\bigl(X^{(k)}_s\bigr)^2\,ds ,
\]
so the left-hand side is an $\F_t$-measurable, $[0,\infty]$-valued random variable; as $\Omega_0\in\F_0\subseteq\F_t$, it follows that $B_t$ is $\F_t$-measurable. Thus $B$ is a continuous, adapted process of bounded variation with $B_0=0$, and $\zeta=N+B$ is a continuous semimartingale in the sense of \bk{Definition 3.1}, with $\zeta_0=0$. Consequently, for every $\omega$ the path $t\mapsto\zeta_t(\omega)$ is bounded on each interval $[0,T]$, whence $\inf_{0\le t\le T}\zeta_t>-\infty$ and $\inf_{0\le t\le T}Z_t=\exp(\inf_{0\le t\le T}\zeta_t)>0$; this is the remark made in the statement. Thus
\[
  Y_t=\frac1{Z_t}=e^{-\zeta_t};\qquad 0\le t<\infty
\]
is a well-defined, strictly positive, continuous and adapted process with $Y_0=1$, and
\begin{equation}\label{eq:3-3-10:b}
  Y_T^*\defeq\sup_{0\le t\le T}Y_t<\infty\quad\text{for every }T<\infty\text{ and every }\omega\in\Omega .
\end{equation}

\emph{Step 2. Associativity of the stochastic integral.}
\begin{lem}
Let $H$ be a continuous, adapted process. Then $H\in\cP^*(N)$, $HX\in\cP$, and a.s.
\[
  \int_0^tH_s\,dN_s=\int_0^tH_sX_s\,dW_s;\qquad 0\le t<\infty .
\]
\end{lem}
\begin{proof}
Being continuous and adapted, $H$ is progressively measurable (\bk{Proposition 1.1.13}), and $HX$ is measurable and adapted. With $H_T^*\defeq\sup_{0\le s\le T}\abs{H_s}<\infty$ we have on $\Omega_0$
\[
  \int_0^TH_s^2\,d\qv{N}_s=\int_0^TH_s^2X_s^2\,ds\le(H_T^*)^2\int_0^TX_s^2\,ds<\infty ,
\]
so that $H\in\cP^*(N)$ and $HX\in\cP$; both integrals $\Phi\defeq I^N(H)$ and $I^W(HX)$ are therefore defined and belong to $\cM^{c,\loc}$.

Let $\wt N\in\cM_2^c$ be arbitrary. By \bk{Proposition 2.24}, applied to the integral $N=I^W(X)$, we have a.s.\ $\qv{N,\wt N}_t=\int_0^tX_u\,d\qv{W,\wt N}_u$ for all $t\ge0$; that is, for almost every $\omega$ the signed measure $d\qv{N,\wt N}$ on $[0,\infty)$ has density $X$ with respect to $d\qv{W,\wt N}$. Hence, for every Borel function $h$ which is bounded on bounded intervals, $\int_0^th(u)\,d\qv{N,\wt N}_u=\int_0^th(u)X_u\,d\qv{W,\wt N}_u$. Using this with $h=H_\cdot(\omega)$ and \bk{Proposition 2.24} for the integral $\Phi=I^N(H)$, we obtain a.s.
\[
  \qv{\Phi,\wt N}_t=\int_0^tH_u\,d\qv{N,\wt N}_u=\int_0^tH_uX_u\,d\qv{W,\wt N}_u;\qquad0\le t<\infty .
\]
Since this holds for every $\wt N\in\cM_2^c$, the uniqueness assertion of \bk{Proposition 2.24} (for the integrator $W$ and the integrand $HX$) yields $\Phi=I^W(HX)$.
\end{proof}

\emph{Step 3. It\^o's rule.}
We apply \bk{Theorem 3.3} to the function $f(x)=e^{-x}$, which is of class $C^2$ on $\R$ with $f'=-f$ and $f''=f$, and to the continuous semimartingale $\zeta=N+B$, $\zeta_0=0$. Since $f(\zeta_t)=Y_t$, \bk{(3.3)} gives, a.s.,
\begin{equation}\label{eq:3-3-10:c}
  Y_t=1-\int_0^tY_s\,dN_s-\int_0^tY_s\,dB_s+\frac12\int_0^tY_s\,d\qv{N}_s;\qquad0\le t<\infty .
\end{equation}
We identify the three integrals. By the Lemma of Step~2 with $H=Y$, the process $YX$ belongs to $\cP$ and
\[
  \int_0^tY_s\,dN_s=\int_0^tY_sX_s\,dW_s .
\]
On $\Omega_0$ the functions $t\mapsto B_t$ and $t\mapsto\qv{N}_t$ are absolutely continuous with densities $-\frac12X_t^2$ and $X_t^2$, respectively, and by \eqref{eq:3-3-10:b} we have $\int_0^tY_sX_s^2\,ds\le Y_t^*\int_0^tX_s^2\,ds<\infty$; therefore the Lebesgue--Stieltjes integrals in \eqref{eq:3-3-10:c} are
\[
  \int_0^tY_s\,dB_s=-\frac12\int_0^tY_sX_s^2\,ds,\qquad
  \int_0^tY_s\,d\qv{N}_s=\int_0^tY_sX_s^2\,ds .
\]
Substituting these three identities into \eqref{eq:3-3-10:c}, we find that a.s.
\[
  Y_t=1-\int_0^tY_sX_s\,dW_s+\Bigl(\frac12+\frac12\Bigr)\int_0^tY_sX_s^2\,ds
     =1+\int_0^tY_sX_s^2\,ds-\int_0^tY_sX_s\,dW_s;\qquad0\le t<\infty ,
\]
which is \eqref{eq:3-3-10:a}, i.e.\ $dY_t=Y_tX_t^2\,dt-Y_tX_t\,dW_t$, $Y_0=1$.

\rmk{In the differential notation of \bk{Example 3.9} the computation reads: $d\zeta_t=X_t\,dW_t-\frac12X_t^2\,dt$, $(d\zeta_t)^2=X_t^2\,dt$, and
$dY_t=-Y_t\,d\zeta_t+\frac12Y_t(d\zeta_t)^2=-Y_tX_t\,dW_t+\frac12Y_tX_t^2\,dt+\frac12Y_tX_t^2\,dt$.
Thus $Z$ has no drift, by \bk{(3.6)}, whereas its reciprocal $Y$ has the nondecreasing drift $\int_0^tY_sX_s^2\,ds$.}
\end{kex}

\begin{kex}{Problem}{3.12}{}
\emph{Meaning of the integrals.} As in \bk{Remark 3.5}, integration with respect to a continuous
semimartingale is integration with respect to its two parts:
\begin{equation}\label{eq:3-3-12:def}
  \int_0^t X_s\,dY_s \defeq \int_0^t X_s\,dN_s + \int_0^t X_s\,dC_s,
  \qquad
  \int_0^t Y_s\,dX_s \defeq \int_0^t Y_s\,dM_s + \int_0^t Y_s\,dB_s .
\end{equation}
All four integrals are well defined. Indeed, $X$ and $Y$ are continuous and adapted, hence
progressively measurable, and for every $\omega$ outside a null set and every $t$ the paths
$s\mapsto X_s(\omega)$, $s\mapsto Y_s(\omega)$ are bounded on $[0,t]$, so that
\[
  \int_0^t X_s^2\,d\qv{N}_s \le \max_{0\le s\le t} X_s^2\cdot\qv{N}_t<\infty,
  \qquad
  \int_0^t Y_s^2\,d\qv{M}_s \le \max_{0\le s\le t} Y_s^2\cdot\qv{M}_t<\infty .
\]
Thus the integrals $\int_0^t X_s\,dN_s$ and $\int_0^t Y_s\,dM_s$ exist as stochastic integrals with
respect to continuous local martingales (Section~3.2, exactly as in \bk{Remark 3.4}) and are
continuous local martingales, while $\int_0^t X_s\,dC_s$ and $\int_0^t Y_s\,dB_s$ exist pathwise as
Lebesgue--Stieltjes integrals of continuous functions with respect to continuous functions of
bounded variation on $[0,t]$.

\emph{Application of It\^o's rule.} We apply \bk{Theorem 3.6} with $d=2$ to the vector of local
martingales $(M^{(1)},M^{(2)})\defeq(M,N)$, the vector of adapted processes of bounded variation
$(B^{(1)},B^{(2)})\defeq(B,C)$, which vanishes at $t=0$, and the initial vector $(X_0,Y_0)$, which is
$\F_0$-measurable because $X$ and $Y$ are adapted; the resulting two-dimensional process is
$(X_t,Y_t)$ of \bk{(3.7)}. For the function we take
\[
  f(t,x,y)\defeq xy ;\qquad (t,x,y)\in[0,\infty)\times\R^2,
\]
which is a polynomial and therefore of class $C^{1,2}$, with
\[
  \frac{\partial f}{\partial t}=0,\quad
  \frac{\partial f}{\partial x}=y,\quad
  \frac{\partial f}{\partial y}=x,\quad
  \frac{\partial^2 f}{\partial x^2}=\frac{\partial^2 f}{\partial y^2}=0,\quad
  \frac{\partial^2 f}{\partial x\,\partial y}=\frac{\partial^2 f}{\partial y\,\partial x}=1 .
\]
Formula \bk{(3.5)} then reads: a.s., for all $0\le t<\infty$,
\begin{equation}\label{eq:3-3-12:ito}
  X_tY_t = X_0Y_0 + \int_0^t Y_s\,dB_s + \int_0^t X_s\,dC_s
  + \int_0^t Y_s\,dM_s + \int_0^t X_s\,dN_s
  + \frac12\Bigl(\int_0^t d\qv{M,N}_s + \int_0^t d\qv{N,M}_s\Bigr).
\end{equation}
Since $MN=NM$, the uniqueness assertion of \bk{Problem 1.5.17} gives $\qv{N,M}=\qv{M,N}$ up to
indistinguishability, and $\qv{M,N}_0=0$ a.s., so the last term of \eqref{eq:3-3-12:ito} equals
$\qv{M,N}_t$. Grouping the remaining integrals according to \eqref{eq:3-3-12:def}, we obtain, a.s.,
\[
  X_tY_t = X_0Y_0 + \int_0^t Y_s\,dX_s + \int_0^t X_s\,dY_s + \qv{M,N}_t ;
  \qquad 0\le t<\infty,
\]
which is \bk{(3.8)}. Since both sides of \bk{(3.8)} are continuous processes, the exceptional null
set does not depend on $t$.

\rmk{The formula does not depend on the choice of the decompositions in \bk{(3.7)}: by
\bk{Problem 3.2} the parts $M,B$ and $N,C$ are determined by $X$ and $Y$ up to
indistinguishability. In particular, for $X=Y$ we recover
$X_t^2=X_0^2+2\int_0^tX_s\,dX_s+\qv{M}_t$, in agreement with \bk{(3.3)} for $f(x)=x^2$ and with
\bk{Example 3.8}.}
\end{kex}

\begin{kex}{Exercise}{3.17}{}
Let $\{\F_t\}$ be the filtration of $W$ in \bk{Definition 2.5.1}. No filtration is prescribed for the processes built from $M=(M^{(1)},M^{(2)},M^{(3)})$, so we say that a continuous process $Y$ with values in $\R^d$ is a $d$-dimensional Brownian motion if $\{Y_t,\F^Y_t;\ t\ge0\}$ is one in the sense of \bk{Definition 2.5.1}, where $\F^Y_t\defeq\sigma(Y_u;\ 0\le u\le t)$; in Step~3 we shall see that $M$ is not a Brownian motion relative to any filtration whatsoever. We write $\eps_i\defeq\sgn(W^{(i)}_1)$, so that $X=\eps_1\eps_2\eps_3$, and
\[
  N\defeq\bigcup_{i=1}^3\{W^{(i)}_1=0\},\qquad P(N)=0 .
\]
On $\Omega\setminus N$ each $\eps_i$ takes the values $\pm1$, with $\eps_i=1$ exactly when $W^{(i)}_1>0$, whichever value is assigned to $\sgn(0)$; in every case $\abs{X}\le1$. The paths of $M$ are continuous and $M_0=0$ a.s., because $W_0=0$ a.s.

\emph{Step 1: two preliminary facts.}
\textbf{(a)} \emph{The processes $W^{(1)},W^{(2)},W^{(3)}$ are independent, and each has the finite-dimensional distributions of a standard one-dimensional Brownian motion.} Indeed, let $0=t_0<t_1<\dots<t_n$. By \bk{Definition 2.5.1}(ii), $W_{t_j}-W_{t_{j-1}}$ is independent of $\F_{t_{j-1}}$, hence of $(W_{t_1}-W_{t_0},\dots,W_{t_{j-1}}-W_{t_{j-2}})$, and has the law $N(0,(t_j-t_{j-1})I_3)$; by induction on $j$ the $3n$ real random variables $W^{(i)}_{t_j}-W^{(i)}_{t_{j-1}}$ are independent, with $W^{(i)}_{t_j}-W^{(i)}_{t_{j-1}}\sim N(0,t_j-t_{j-1})$. Since $W_0=0$ a.s., the three vectors $(W^{(i)}_{t_1},\dots,W^{(i)}_{t_n})$, $i=1,2,3$, are independent, and each has the law of $(B_{t_1},\dots,B_{t_n})$ for a standard Brownian motion $B$. The events $\{(W^{(i)}_{t_1},\dots,W^{(i)}_{t_n})\in\Gamma\}$, $\Gamma\in\B(\R^n)$, form for each $i$ a $\pi$-system generating $\sigma(W^{(i)}_t;\ t\ge0)$, so Dynkin's $\pi$-$\lambda$ theorem shows that these three $\sigma$-fields are independent.

\textbf{(b)} \emph{If $Y$ is a continuous process with values in $\R^2$ which has the same finite-dimensional distributions as $B\defeq(W^{(1)},W^{(3)})$, then $Y$ is a two-dimensional Brownian motion starting at the origin.} First, $P[Y_0=0]=P[B_0=0]=1$. Let $0\le s_1<\dots<s_n\le s<t$. The vector $B_t-B_s$ consists of two components of $W_t-W_s$, so it has the law $N(0,(t-s)I_2)$ and is independent of $\F_s$, hence of $(B_{s_1},\dots,B_{s_n})$. The vector $(Y_{s_1},\dots,Y_{s_n},Y_t-Y_s)$ has the same law as $(B_{s_1},\dots,B_{s_n},B_t-B_s)$; therefore $Y_t-Y_s\sim N(0,(t-s)I_2)$ is independent of $(Y_{s_1},\dots,Y_{s_n})$. The events $\{(Y_{s_1},\dots,Y_{s_n})\in\Gamma\}$ form a $\pi$-system generating $\F^Y_s$, and by Dynkin's $\pi$-$\lambda$ theorem $Y_t-Y_s$ is independent of $\F^Y_s$. Thus $\{Y_t,\F^Y_t\}$ satisfies \bk{Definition 2.5.1}.

\emph{Step 2: the pairs.}
The pair $(M^{(1)},M^{(2)})=(W^{(1)},W^{(2)})$ is adapted to $\{\F_t\}$, and for $s<t$ its increment over $[s,t]$ consists of two components of $W_t-W_s$; it is thus independent of $\F_s\supseteq\F^{(M^{(1)},M^{(2)})}_s$ with law $N(0,(t-s)I_2)$. Hence $(M^{(1)},M^{(2)})$ is a two-dimensional Brownian motion.

Consider $(M^{(1)},M^{(3)})=(W^{(1)},XW^{(3)})$, and put $\cH\defeq\sigma(W^{(1)}_t,W^{(3)}_t;\ t\ge0)$. By Step 1(a), $\eps_2$ is independent of $\cH$ and $P[\eps_2=1]=P[\eps_2=-1]=\tfrac12$. For $A\in\cH$ the events $A\cap\{\eps_1\eps_3=\pm1\}$ belong to $\cH$, and up to subsets of $N$,
\[
  A\cap\{X=1\}=\bigl(A\cap\{\eps_1\eps_3=1\}\cap\{\eps_2=1\}\bigr)\cup\bigl(A\cap\{\eps_1\eps_3=-1\}\cap\{\eps_2=-1\}\bigr),
\]
so that
\begin{equation}\label{eq:3-3-17:a}
  P[A\cap\{X=1\}]=\tfrac12P[A\cap\{\eps_1\eps_3=1\}]+\tfrac12P[A\cap\{\eps_1\eps_3=-1\}]=\tfrac12P(A),
\end{equation}
and consequently $P[A\cap\{X=-1\}]=\tfrac12P(A)$ as well. Now fix $0\le t_1<\dots<t_n$ and a bounded Borel function $f\colon\R^n\times\R^n\to\R$, and write $U\defeq(W^{(1)}_{t_1},\dots,W^{(1)}_{t_n})$, $V\defeq(W^{(3)}_{t_1},\dots,W^{(3)}_{t_n})$. Since $f(U,V)$ and $f(U,-V)$ are bounded and $\cH$-measurable, \eqref{eq:3-3-17:a} extends from indicators to these random variables and gives
\[
  E f(U,XV)=E\bigl[f(U,V)\ind_{\{X=1\}}\bigr]+E\bigl[f(U,-V)\ind_{\{X=-1\}}\bigr]
  =\tfrac12E f(U,V)+\tfrac12E f(U,-V).
\]
By Step 1(a), $U$ and $V$ are independent and $V$ is a centred Gaussian vector, so $-V$ has the same law as $V$ and is independent of $U$; hence $(U,-V)$ and $(U,V)$ have the same law, and $Ef(U,XV)=Ef(U,V)$. This says that $(M^{(1)},M^{(3)})$ has the same finite-dimensional distributions as $(W^{(1)},W^{(3)})$; being continuous, it is a two-dimensional Brownian motion by Step 1(b).

The definition of $M$ and the hypotheses on $W$ are symmetric in the indices $1$ and $2$, so the same argument, with $\eps_1$ in the role of $\eps_2$ and $(W^{(2)},W^{(3)})$ in the role of $(W^{(1)},W^{(3)})$, shows that $(M^{(2)},M^{(3)})$ is a two-dimensional Brownian motion.

\emph{Step 3: the triple is not a Brownian motion.}
Suppose that $\{M_t,\G_t\}$ were a three-dimensional Brownian motion for some filtration $\{\G_t\}$. Since $M_0=0$ a.s., the vector $M_1$ would coincide a.s.\ with the increment $M_1-M_0$ and thus have the law $N(0,I_3)$; in particular,
\[
  P\bigl[M^{(1)}_1>0,\ M^{(2)}_1>0,\ M^{(3)}_1<0\bigr]=\tfrac18 .
\]
But on $\Omega\setminus N$ we have $\eps_3W^{(3)}_1=\abs{W^{(3)}_1}$, whence
\begin{equation}\label{eq:3-3-17:b}
  M^{(3)}_1=\eps_1\eps_2\,\abs{W^{(3)}_1}\qquad\text{on }\Omega\setminus N,
\end{equation}
and on $\{M^{(1)}_1>0,\ M^{(2)}_1>0\}\setminus N$ this gives $\eps_1=\eps_2=1$ and $M^{(3)}_1>0$. The probability displayed above is therefore $0$, a contradiction. So $M$ is not a three-dimensional Brownian motion, relative to any filtration.

\emph{Step 4: why \bk{Theorem 3.16} is not contradicted.}
\bk{Theorem 3.16} requires \emph{one} filtration $\{\F_t\}$ to which the whole vector is adapted and relative to which \emph{all} components are continuous local martingales; the cross-variations in \bk{(3.11)} are then those of \bk{Problem 1.5.17} for this filtration. What Step 2 provides is weaker. For each pair $k\neq j$, let $\F^{(k,j)}_t\defeq\sigma(M^{(k)}_u,M^{(j)}_u;\ u\le t)$. For $s<t$, the increments $\Delta^k\defeq M^{(k)}_t-M^{(k)}_s$ and $\Delta^j\defeq M^{(j)}_t-M^{(j)}_s$ are independent of $\F^{(k,j)}_s$, independent of each other, centred, with variance $t-s$; hence $M^{(k)}$ and $M^{(j)}$ are square-integrable martingales of $\{\F^{(k,j)}_t\}$, and from
\[
  M^{(k)}_tM^{(j)}_t-M^{(k)}_sM^{(j)}_s=\Delta^k\Delta^j+M^{(k)}_s\Delta^j+M^{(j)}_s\Delta^k,\qquad
  (M^{(k)}_t)^2-(M^{(k)}_s)^2=(\Delta^k)^2+2M^{(k)}_s\Delta^k,
\]
the processes $M^{(k)}M^{(j)}$ and $(M^{(k)}_t)^2-t$ are martingales of $\{\F^{(k,j)}_t\}$. So the relations \bk{(3.11)} hold for each pair, but each time relative to a \emph{different} filtration $\{\F^{(k,j)}_t\}$, none of which carries the third component. We now show that no single filtration does the job; more precisely:

\begin{claim}
Let $\{\G_t\}$ be any filtration to which $M$ is adapted. Then $M^{(3)}$ is not a local martingale of $\{\G_t\}$. In particular, the hypotheses of \bk{Theorem 3.16} fail for $M$ and every filtration.
\end{claim}

\begin{proof}
Let $\F^M_t\defeq\sigma(M_u;\ u\le t)\subseteq\G_t$, and suppose that $M^{(3)}$ is a local martingale of $\{\G_t\}$, with stopping times $T_n\uparrow\infty$ a.s.\ as in \bk{Definition 1.5.15}. For $0\le s<t$ we have $E[M^{(3)}_{t\wedge T_n}\mid\G_s]=M^{(3)}_{s\wedge T_n}$ a.s. Moreover $\abs{M^{(3)}_{t\wedge T_n}}\le Z_t\defeq\sup_{u\le t}\abs{W^{(3)}_u}$ because $\abs{X}\le1$, and $E Z_t^2\le4E(W^{(3)}_t)^2=4t$ by Doob's maximal inequality (\bk{Theorem 1.3.8}(iv)) applied to the continuous, nonnegative submartingale $\abs{W^{(3)}}$ of $\{\F_t\}$ (the process $W^{(3)}$ is a martingale of $\{\F_t\}$, its increments being centred and independent of the past). Letting $n\to\infty$, the dominated convergence theorem for conditional expectations gives $E[M^{(3)}_t\mid\G_s]=M^{(3)}_s$ a.s. Conditioning further on $\F^M_s\subseteq\G_s$, and using that $M^{(3)}_s$ is $\F^M_s$-measurable, we obtain
\begin{equation}\label{eq:3-3-17:c}
  E\bigl[(M^{(3)}_t-M^{(3)}_s)\,\zeta\bigr]=0\qquad\text{for every bounded, $\F^M_s$-measurable }\zeta .
\end{equation}
We contradict \eqref{eq:3-3-17:c} with $0<s<1$, $t=1$ and $\zeta\defeq\sgn(M^{(1)}_s)\sgn(M^{(2)}_s)=\sgn(W^{(1)}_s)\sgn(W^{(2)}_s)$.

Let $B$ stand for any of $W^{(1)},W^{(2)},W^{(3)}$, and set
\[
  h\defeq E\abs{B_1}=\sqrt{2/\pi},\qquad c\defeq E\bigl[\sgn(B_1)\sgn(B_s)\bigr];
\]
these numbers do not depend on the component, by Step 1(a). The Gaussian pair $(B_s-sB_1,\,B_1)$ has covariance $s-s=0$, so $B_s-sB_1$ is independent of $B_1$ and centred, and
\[
  E\bigl[\sgn(B_1)B_s\bigr]=s\,E\bigl[\sgn(B_1)B_1\bigr]+E\bigl[\sgn(B_1)\bigr]E[B_s-sB_1]=s\,h .
\]
Next, $B_1-B_s\sim N(0,1-s)$ is independent of $B_s$, so $E[\sgn(B_1)\mid B_s]=g(B_s)$ a.s., where for $x\in\R$ and a standard normal variable $\xi$,
\[
  g(x)\defeq P\bigl[x+\sqrt{1-s}\,\xi>0\bigr]-P\bigl[x+\sqrt{1-s}\,\xi<0\bigr]
  =\sgn(x)\,P\bigl[\abs{\xi}<\abs{x}/\sqrt{1-s}\,\bigr]\qquad(x\neq0).
\]
Thus $g(x)\sgn(x)>0$ for $x\neq0$, and since $P[B_s=0]=0$,
\[
  c=E\bigl[g(B_s)\sgn(B_s)\bigr]>0 .
\]
Now we compute, using \eqref{eq:3-3-17:b}, $X=\eps_1\eps_2\eps_3$ a.s., and the independence of the three components:
\begin{align*}
  E\bigl[M^{(3)}_1\zeta\bigr]
  &=E\bigl[\eps_1\sgn(W^{(1)}_s)\bigr]\,E\bigl[\eps_2\sgn(W^{(2)}_s)\bigr]\,E\abs{W^{(3)}_1}=c^2h,\\
  E\bigl[M^{(3)}_s\zeta\bigr]
  &=E\bigl[\eps_1\sgn(W^{(1)}_s)\bigr]\,E\bigl[\eps_2\sgn(W^{(2)}_s)\bigr]\,E\bigl[\eps_3W^{(3)}_s\bigr]=c^2sh .
\end{align*}
Hence $E[(M^{(3)}_1-M^{(3)}_s)\zeta]=c^2h(1-s)>0$, which contradicts \eqref{eq:3-3-17:c}.
\end{proof}

In summary: $M$ is continuous, each component is a Brownian motion and each pair of components satisfies \bk{(3.11)} in its own filtration, but there is no filtration relative to which $M$ is adapted and all three components are local martingales. The reason is visible in \eqref{eq:3-3-17:b}: the sign of $M^{(3)}_1$ is the product of the signs of $M^{(1)}_1$ and $M^{(2)}_1$, so that the paths of $M^{(1)}$ and $M^{(2)}$ up to a time $s<1$ carry information about the future increment $M^{(3)}_1-M^{(3)}_s$. Nor does the filtration $\{\F_t\}$ of $W$ help: the random variable $X$, which is defined through $W_1$, enters $M^{(3)}_t$ already for $0<t<1$, so either $M^{(3)}$ fails to be adapted to $\{\F_t\}$, or it is adapted and then, by the Claim, it is not a local martingale of $\{\F_t\}$. Since the hypotheses of \bk{Theorem 3.16} are not satisfied, the fact that $M$ is not a Brownian motion does not contradict the theorem; it shows rather that the requirement of a common filtration in that theorem cannot be dropped.
\end{kex}

\begin{kex}{Problem}{3.18}{}
Write $Q=(q_{kj})_{1\le k,j\le d}$, so that
\[
  \wt{W}^{(k)}_t=\sum_{j=1}^d q_{kj}W^{(j)}_t;\qquad 0\le t<\infty,\ 1\le k\le d .
\]
We show that $\wt{W}=\{\wt{W}_t,\F_t;\ 0\le t<\infty\}$, with the \emph{same} filtration $\{\F_t\}$, is a $d$-dimensional Brownian motion starting at the origin in the sense of \bk{Definition 2.5.1}, by verifying the hypotheses of P.~L\'evy's \bk{Theorem 3.16}.

\emph{Step 1: $\wt{W}$ is continuous, adapted, and starts at the origin.} For every $\omega$ the path $t\mapsto \wt{W}_t(\omega)$ is the image of the continuous path $t\mapsto W_t(\omega)$ under the linear, hence continuous, map $x\mapsto Qx$ of $\R^d$; and $\wt{W}_t$ is $\F_t$-measurable, being a linear combination of the $\F_t$-measurable random variables $W^{(1)}_t,\dots,W^{(d)}_t$. Moreover $\{W_0=0\}\subseteq\{\wt{W}_0=0\}$, so $P[\wt{W}_0=0]=1$.

\emph{Step 2: the components are continuous local martingales.} Throughout Section 3.3 the filtration $\{\F_t\}$ is assumed to satisfy the usual conditions (this standing assumption is made at the beginning of the section and is in force in \bk{Theorem 3.16}); consequently the definitions and results of Section 1.5, which are stated under the same assumption, are available. Put
\[
  M^{(j)}_t\defeq W^{(j)}_t-W^{(j)}_0,\qquad
  \wt{M}^{(k)}_t\defeq \wt{W}^{(k)}_t-\wt{W}^{(k)}_0=\sum_{j=1}^d q_{kj}M^{(j)}_t .
\]
By \bk{Problem 2.5.5}, each $\{M^{(j)}_t,\F_t;\ 0\le t<\infty\}$ is a continuous, square-integrable martingale with $M^{(j)}_0=0$, that is, $M^{(j)}\in\cM_2^c$, and
\begin{equation}\label{eq:3-3-18:a}
  \qv{M^{(j)},M^{(m)}}_t=\delta_{jm}\,t;\qquad 0\le t<\infty,\ 1\le j,m\le d .
\end{equation}
The space $\cM_2^c$ is a vector space, so $\wt{M}^{(k)}\in\cM_2^c\subseteq\cM^{c,\loc}$ for every $1\le k\le d$ (a martingale is a local martingale, with the localizing sequence $T_n=n$).

\emph{Step 3: the cross-variations.} The cross-variation of \bk{Definition 1.5.5} is a bilinear form on $\cM_2$ (\bk{Problem 1.5.7}). Hence, by \eqref{eq:3-3-18:a} and the orthogonality relation $QQ^T=I_d$, we have a.s.\ for all $0\le t<\infty$ and $1\le k,l\le d$:
\[
  \qv{\wt{M}^{(k)},\wt{M}^{(l)}}_t
  =\sum_{j=1}^d\sum_{m=1}^d q_{kj}q_{lm}\qv{M^{(j)},M^{(m)}}_t
  =t\sum_{j=1}^d q_{kj}q_{lj}
  =t\,(QQ^T)_{kl}=\delta_{kl}\,t .
\]
The cross-variation in \bk{(3.11)} is the one defined for elements of $\cM^{c,\loc}$ in \bk{Problem 1.5.17} and \bk{Definition 1.5.18}, and for $X,Y\in\cM_2^c$ it coincides with that of \bk{Definition 1.5.5}. Indeed, $\qv{X,Y}=\frac14\bigl[\qv{X+Y}-\qv{X-Y}\bigr]$ is adapted, vanishes at $t=0$, and is the difference of two nondecreasing processes which are continuous because $X\pm Y\in\cM_2^c$ (see the discussion preceding \bk{Definition 1.5.3}); thus it is continuous and of bounded variation on each $[0,t]$. Moreover $XY-\qv{X,Y}$ is a continuous martingale vanishing at $t=0$, hence an element of $\cM^{c,\loc}$. By the uniqueness assertion of \bk{Problem 1.5.17}, $\qv{X,Y}$ is the cross-variation of $X$ and $Y$ regarded as elements of $\cM^{c,\loc}$. The display above is therefore condition \bk{(3.11)} for the process $\wt{W}$.

\emph{Step 4: conclusion.} By Steps 1--3, $\wt{W}$ is a continuous, adapted process in $\R^d$ whose components, after subtraction of their initial values, are in $\cM^{c,\loc}$ relative to $\{\F_t\}$ and satisfy \bk{(3.11)}. \bk{Theorem 3.16} now shows that $\{\wt{W}_t,\F_t;\ 0\le t<\infty\}$ is a $d$-dimensional Brownian motion; its initial distribution is the unit mass at the origin by Step 1, so it starts at the origin.
\rmk{The assertion can also be checked directly from \bk{Definition 2.5.1}, without \bk{Theorem 3.16}. Let $0\le s<t$. The increment $\wt{W}_t-\wt{W}_s=Q(W_t-W_s)$ is a Borel function of $W_t-W_s$, and is therefore independent of $\F_s$. Since $W_t-W_s$ has characteristic function $u\mapsto e^{-\frac12\norm{u}^2(t-s)}$, we obtain for $u\in\R^d$, using $(u,Qx)=(Q^Tu,x)$ and $\norm{Q^Tu}^2=(u,QQ^Tu)=\norm{u}^2$,
\[
  E\bigl[e^{i(u,\wt{W}_t-\wt{W}_s)}\bigr]
  =E\bigl[e^{i(Q^Tu,\,W_t-W_s)}\bigr]
  =e^{-\frac12\norm{Q^Tu}^2(t-s)}
  =e^{-\frac12\norm{u}^2(t-s)} ,
\]
so $\wt{W}_t-\wt{W}_s$ is normal with mean zero and covariance matrix $(t-s)I_d$, by the uniqueness theorem for characteristic functions. Together with Step 1 this is all that \bk{Definition 2.5.1} requires. This direct argument uses no property of the filtration, so the assertion holds whether or not $\{\F_t\}$ satisfies the usual conditions.}

\rmk{Neither argument uses $W_0=0$ beyond Step 1. If $W$ is a $d$-dimensional Brownian motion with initial distribution $\mu$, then Steps 2--4 apply verbatim and show that $\wt{W}=QW$ is a $d$-dimensional Brownian motion with initial distribution $\mu\circ Q^{-1}$; in particular, if $W$ starts at $x\in\R^d$, then $\wt{W}$ starts at $Qx$. This is the form in which the result is used for \bk{(3.15)}. The origin is singled out in the statement because it is the only starting point $x$ with $Qx=x$ for every orthogonal $Q$, i.e.\ the only one for which $\wt{W}$ and $W$ have the same law for every $Q$.}
\end{kex}

\begin{kex}{Problem}{3.20}{}
Throughout, $W=\{W_t,\F_t;\,0\le t<\infty\}$, $(\Omega,\F)$, $\{P^x\}_{x\in\R^d}$ is the $d$-dimensional Brownian family of \bk{Definition 3.19}, $R_t=\norm{W_t}$, and $\wh P^r=P^{\phi(r)}$ for $r\ge0$, where
\[
  \phi:[0,\infty)\to\R^d,\qquad \phi(r)\defeq(r,0,\dots,0).
\]

\emph{The statement to be proved.} \bk{Definition 2.6.3}, modified for the state space $[0,\infty)$, requires that the process $R=\{R_t,\F_t;\,0\le t<\infty\}$ be progressively measurable and that
\begin{enumerate}[label=\textup{(\alph*)}]
\item for each $F\in\F$, the mapping $r\mapsto\wh P^r(F)$ from $[0,\infty)$ into $[0,1]$ is universally measurable;
\item $\wh P^r[R_0=r]=1$ for every $r\ge0$;
\item for $r\ge0$, $t\ge0$, $\Gamma\in\B([0,\infty))$ and any optional time $S$ of $\{\F_t\}$,
\[
  \wh P^r[R_{S+t}\in\Gamma\mid\F_{S+}]=\wh P^r[R_{S+t}\in\Gamma\mid R_S],\qquad \wh P^r\text{-a.s. on }\{S<\infty\};
\]
\item for $r\ge0$, $t\ge0$, $\Gamma\in\B([0,\infty))$ and any optional time $S$ of $\{\F_t\}$,
\[
  \wh P^r[R_{S+t}\in\Gamma\mid R_S=\rho]=\wh P^{\rho}[R_t\in\Gamma],\qquad \wh P^rR_S^{-1}\text{-a.e. }\rho\in[0,\infty).
\]
\end{enumerate}
Here, as in \bk{Remark 2.6.4}, $\{R_{S+t}\in\Gamma\}\defeq\{S<\infty,\ R_{S+t}\in\Gamma\}$, the random variable $R_S$ is defined on $\{S<\infty\}$, $\sigma(R_S)$ is the $\sigma$-field of \bk{Problem 1.1.17}, consisting of the sets $\{S<\infty,\ R_S\in A\}$ and $\{S<\infty,\ R_S\in A\}\cup\{S=\infty\}$, $A\in\B([0,\infty))$, and $\wh P^rR_S^{-1}(A)\defeq\wh P^r[S<\infty,\ R_S\in A]$ is a measure on $\B([0,\infty))$ of total mass $\wh P^r[S<\infty]$. The left-hand side of (d) denotes any Borel-measurable function $g:[0,\infty)\to\R$ for which $g(R_S)=\wh P^r[R_{S+t}\in\Gamma\mid R_S]$ holds $\wh P^r$-a.s.\ on $\{S<\infty\}$; two such functions agree $\wh P^rR_S^{-1}$-a.e., so (d) asserts that $\rho\mapsto\wh P^\rho[R_t\in\Gamma]$ is, up to a $\wh P^rR_S^{-1}$-null set, such a function.

\emph{Step 1: progressive measurability, \textup{(a)} and \textup{(b)}.} By \bk{Theorem 2.6.15}, $W$ with $\{P^x\}_{x\in\R^d}$ is a strong Markov family; in particular $W$ is progressively measurable with respect to $\{\F_t\}$. For each $t\ge0$ the map $(s,\omega)\mapsto R_s(\omega)$ on $[0,t]\times\Omega$ is the composition of the $\B([0,t])\otimes\F_t$-measurable map $(s,\omega)\mapsto W_s(\omega)$ with the continuous map $\norm{\cdot}:\R^d\to[0,\infty)$, and is therefore $\B([0,t])\otimes\F_t$-measurable. Thus $R$ is progressively measurable.

For (a), fix $F\in\F$ and set $f(x)\defeq P^x(F)$, $x\in\R^d$; this function is universally measurable by \bk{Definition 2.5.8}(i), and $\wh P^r(F)=f(\phi(r))$. Let $B\in\B(\R)$, so that $A\defeq f^{-1}(B)\in\mathscr{U}(\R^d)$, and let $\mu$ be a finite measure on $\B([0,\infty))$. Since $\phi$ is continuous, hence Borel-measurable, $\nu\defeq\mu\phi^{-1}$ is a finite measure on $\B(\R^d)$. As $A$ belongs to the completion of $\B(\R^d)$ with respect to $\nu$ (\bk{Definition 2.5.6}), there are $B_1,B_2\in\B(\R^d)$ with $B_1\subseteq A\subseteq B_2$ and $\nu(B_2\setminus B_1)=0$. Then $\phi^{-1}(B_1)\subseteq\phi^{-1}(A)\subseteq\phi^{-1}(B_2)$, the two outer sets are in $\B([0,\infty))$, and $\mu\bigl(\phi^{-1}(B_2)\setminus\phi^{-1}(B_1)\bigr)=\nu(B_2\setminus B_1)=0$. Hence $(f\circ\phi)^{-1}(B)=\phi^{-1}(A)$ lies in the completion of $\B([0,\infty))$ with respect to $\mu$; $\mu$ being arbitrary, $\phi^{-1}(A)\in\mathscr{U}([0,\infty))$, and (a) follows.

For (b): $\wh P^r[R_0=r]\ge P^{\phi(r)}[W_0=\phi(r)]=1$, because $\norm{\phi(r)}=r$.

\emph{Step 2: the transition function.} Fix $t\ge0$ and $\Gamma\in\B([0,\infty))$, and put
\[
  \wt\Gamma\defeq\{y\in\R^d;\ \norm{y}\in\Gamma\}\in\B(\R^d),\qquad
  h(x)\defeq P^x[W_t\in\wt\Gamma]=P^x[R_t\in\Gamma],\quad x\in\R^d,
\]
and $\wh h(\rho)\defeq h(\phi(\rho))=\wh P^\rho[R_t\in\Gamma]$, $\rho\ge0$. We claim that
\begin{equation}\label{eq:3-3-20:rad}
  h(x)=\wh h(\norm{x}),\quad x\in\R^d,\qquad\text{and}\qquad \wh h:[0,\infty)\to[0,1]\ \text{is Borel-measurable.}
\end{equation}
Under $P^x$ the process $W$ is a $d$-dimensional Brownian motion starting at $x$, so by \bk{Definition 2.5.1} the law $P^xW_t^{-1}$ is normal with mean $x$ and covariance matrix $tI_d$ (the unit mass at $x$ if $t=0$). Given $x\in\R^d$, choose an orthogonal matrix $Q$ with $Qx=\phi(\norm{x})$. Under $P^x$ the vector $QW_t$ is normal with mean $Qx$ and covariance $tQI_dQ^T=tI_d$, i.e., it has the law of $W_t$ under $P^{Qx}$; since $\norm{QW_t}=\norm{W_t}$,
\[
  h(x)=P^x[\norm{QW_t}\in\Gamma]=P^{Qx}[\norm{W_t}\in\Gamma]=h(\phi(\norm{x}))=\wh h(\norm{x}).
\]
(This is \bk{(3.15)}, read at the single time $t$.) As for measurability: if $t=0$, then $\wh h=\ind_\Gamma$ by (b); if $t>0$, then
\[
  \wh h(\rho)=\int_{\R^d}\ind_{\wt\Gamma}(y)\,(2\pi t)^{-d/2}\exp\Bigl(-\frac{\norm{y-\phi(\rho)}^2}{2t}\Bigr)\,dy,
\]
and the integrand is a nonnegative Borel-measurable function of $(\rho,y)\in[0,\infty)\times\R^d$, so $\wh h$ is Borel-measurable by the Fubini--Tonelli theorem. This proves \eqref{eq:3-3-20:rad}.

\emph{Step 3: the strong Markov property of $W$, transferred to $R$.} Fix $r\ge0$ and an optional time $S$ of $\{\F_t\}$, and keep $t$, $\Gamma$ as in Step~2. In the notation \bk{(2.5.7)} for the Brownian family, $h=U_t\ind_{\wt\Gamma}$. By \bk{Theorem 2.6.15} and \bk{Proposition 2.6.7}, condition (e) of the latter holds for $W$; with $x=\phi(r)$ and the set $\wt\Gamma$ it reads
\[
  \wh P^r[W_{S+t}\in\wt\Gamma\mid\F_{S+}]=h(W_S),\qquad\wh P^r\text{-a.s. on }\{S<\infty\}.
\]
Now $\{W_{S+t}\in\wt\Gamma\}=\{S<\infty,\ R_{S+t}\in\Gamma\}=\{R_{S+t}\in\Gamma\}$, and $h(W_S)=\wh h(R_S)$ on $\{S<\infty\}$ by \eqref{eq:3-3-20:rad}. Moreover $\{R_{S+t}\in\Gamma\}\subseteq\{S<\infty\}$ and $\{S=\infty\}\in\F_{S+}$, so that
\[
  E^{\phi(r)}\bigl[\ind_{\{S=\infty\}}\wh P^r[R_{S+t}\in\Gamma\mid\F_{S+}]\bigr]=\wh P^r[S=\infty,\ R_{S+t}\in\Gamma]=0,
\]
and the nonnegative random variable $\wh P^r[R_{S+t}\in\Gamma\mid\F_{S+}]$ vanishes $\wh P^r$-a.s.\ on $\{S=\infty\}$. With
\[
  Z\defeq\wh h(R_S)\,\ind_{\{S<\infty\}}
\]
(defined as $0$ on $\{S=\infty\}$) we have therefore shown that
\begin{equation}\label{eq:3-3-20:e}
  \wh P^r[R_{S+t}\in\Gamma\mid\F_{S+}]=Z,\qquad\wh P^r\text{-a.s.}
\end{equation}

\emph{Step 4: \textup{(c)} and \textup{(d)}.} We first check that $\sigma(R_S)\subseteq\F_{S+}$ and that $Z$ is $\sigma(R_S)$-measurable. The optional time $S$ is a stopping time of $\{\F_{t+}\}$, and $R$ is progressively measurable with respect to $\{\F_{t+}\}$; by \bk{Proposition 1.2.18} (cf.\ \bk{Remark 2.6.5}), $\{S<\infty,\ R_S\in A\}\in\F_{S+}$ for every $A\in\B([0,\infty))$, and also $\{S=\infty\}\in\F_{S+}$. Hence every set of $\sigma(R_S)$ is in $\F_{S+}$. Next, for $B\in\B(\R)$, the event $\{Z\in B\}$ equals $\{S<\infty,\ R_S\in\wh h^{-1}(B)\}$ if $0\notin B$, and equals $\{S<\infty,\ R_S\in\wh h^{-1}(B)\}\cup\{S=\infty\}$ if $0\in B$; since $\wh h^{-1}(B)\in\B([0,\infty))$ by \eqref{eq:3-3-20:rad}, in both cases $\{Z\in B\}\in\sigma(R_S)$.

By the tower property and \eqref{eq:3-3-20:e}, $\wh P^r$-a.s.,
\[
  \wh P^r[R_{S+t}\in\Gamma\mid R_S]
  =E^{\phi(r)}\bigl[\wh P^r[R_{S+t}\in\Gamma\mid\F_{S+}]\bigm|\sigma(R_S)\bigr]
  =E^{\phi(r)}[Z\mid\sigma(R_S)]=Z .
\]
Comparing with \eqref{eq:3-3-20:e}, we obtain
\[
  \wh P^r[R_{S+t}\in\Gamma\mid\F_{S+}]=Z=\wh P^r[R_{S+t}\in\Gamma\mid R_S],\qquad\wh P^r\text{-a.s.},
\]
which is (c) (even without the restriction to $\{S<\infty\}$). Furthermore, $\wh h$ is a Borel-measurable function with
\[
  \wh h(R_S)=Z=\wh P^r[R_{S+t}\in\Gamma\mid R_S],\qquad\wh P^r\text{-a.s. on }\{S<\infty\},
\]
so $\wh h$ is a version of $\rho\mapsto\wh P^r[R_{S+t}\in\Gamma\mid R_S=\rho]$; any other version $g$ satisfies $g(R_S)=\wh h(R_S)$ a.s.\ on $\{S<\infty\}$, i.e., $\wh P^rR_S^{-1}\{\rho;\ g(\rho)\ne\wh h(\rho)\}=0$. Since $\wh h(\rho)=\wh P^\rho[R_t\in\Gamma]$, this is (d). Conditions (a)--(d) hold for every $d\ge2$, and the Bessel family with dimension $d$ is a strong Markov family.

\rmk{The only property of the function $\norm{\cdot}$ used above, besides its continuity, is \eqref{eq:3-3-20:rad}: the transition probability $P^x[\norm{W_t}\in\Gamma]$ depends on $x$ only through $\norm{x}$. This is where rotational invariance (\bk{Problem 3.18}, here needed only for the normal law of $W_t$) enters. The argument does not use $d\ge2$: for $d=1$ it shows that $\abs{W}$ is a strong Markov family.}
\end{kex}

\begin{kex}{Problem}{3.23}{}
We use the notation of \bk{Definition 3.19} and, as in the proofs of \bk{Propositions 3.21} and \bk{3.22}, write $P$ for $\hat P^r$; thus $R_t=\norm{W_t}$, where $W$ is a $d$-dimensional Brownian motion with $P[W_0=(r,0,\dots,0)]=1$, and the filtration satisfies the usual conditions, as is assumed throughout \bk{Section 3.3}. The paths of $R$ are continuous and nonnegative, so $m=\inf_{t\in\Q,\,t\ge0}R_t$ is a random variable with values in $[0,r]$, a.s. We set
\[
  f(x)\defeq
  \begin{cases}
    \log x, & d=2,\\
    x^{2-d}, & d\ge 3,
  \end{cases}
  \qquad 0<x<\infty .
\]
In both cases $f$ is strictly monotone, of class $C^\infty$ on $(0,\infty)$, and
\begin{equation}\label{eq:3-3-23:harm}
  \frac{d-1}{2x}\,f'(x)+\frac12 f''(x)=0;\qquad 0<x<\infty .
\end{equation}
Indeed, for $d=2$ the left-hand side is $\frac{1}{2x^2}-\frac{1}{2x^2}$, and for $d\ge3$ we have $f'(x)=(2-d)x^{1-d}$, $f''(x)=(2-d)(1-d)x^{-d}$, so the left-hand side is $\frac12(2-d)x^{-d}\bigl[(d-1)+(1-d)\bigr]=0$.

For $0<a<r<b<\infty$ we introduce the hitting times
\[
  T_a\defeq\inf\{t\ge0;\ R_t=a\},\qquad S_b\defeq\inf\{t\ge0;\ R_t=b\},
\]
with $\inf\emptyset=\infty$; they are stopping times because $R$ is continuous and adapted and $\{a\}$, $\{b\}$ are closed (\bk{Problem 1.2.7}).

\emph{Step 1. For $0<a<r<b<\infty$ we have $P[S_b<\infty]=1$ and}
\begin{equation}\label{eq:3-3-23:exit}
  P[T_a<S_b]=\frac{f(b)-f(r)}{f(b)-f(a)} .
\end{equation}
Since $R_t\ge\abs{W^{(1)}_t}$ and $P$-almost every path of the one-dimensional Brownian motion $W^{(1)}$ is unbounded (\bk{Theorem 2.9.23}), almost every path of $R$ is unbounded; as $R_0=r<b$ a.s., the intermediate value theorem gives $P[S_b<\infty]=1$.

By \bk{Proposition 3.21}, $R$ is a continuous semimartingale with decomposition $R_t=r+B_t+A_t$, where $B$ is a standard, one-dimensional Brownian motion relative to $\{\F_t\}$ and $A_t=\int_0^t\frac{d-1}{2R_s}\,ds=R_t-r-B_t$ is a continuous, adapted, nondecreasing process (finite for all $t$, a.s.). Fix a function $g\colon\R\to\R$ of class $C^2$ which agrees with $f$ on $[a/2,\infty)$. It\^o's rule (\bk{Theorem 3.3}), together with $\qv{B}_t=t$, yields a.s.
\[
  g(R_t)=g(r)+\int_0^t g'(R_s)\,dB_s
   +\int_0^t\Bigl[\frac{d-1}{2R_s}\,g'(R_s)+\frac12 g''(R_s)\Bigr]ds;\qquad 0\le t<\infty .
\]
For a positive integer $n$ let $\tau_n\defeq T_a\wedge S_b\wedge n$, a bounded stopping time. For $0\le s\le\tau_n$ we have $a\le R_s\le b$ (continuity of $R$ and $a<R_0=r<b$), so on this stochastic interval $g$, $g'$, $g''$ may be replaced by $f$, $f'$, $f''$ evaluated at $R_s$, and by \eqref{eq:3-3-23:harm} the Lebesgue integral vanishes up to time $\tau_n$. Evaluating at $t=\tau_n$ we obtain
\begin{equation}\label{eq:3-3-23:ito}
  f(R_{\tau_n})=f(r)+\int_0^{\tau_n}g'(R_s)\,dB_s,\qquad\text{a.s.}
\end{equation}
The process $X_s\defeq g'(R_s)\ind_{\{s\le\tau_n\}}$ is adapted and left-continuous, hence progressively measurable, and it is bounded by $C\defeq\max_{a\le x\le b}\abs{f'(x)}$. Consequently $X\in\cL^*$ for the Brownian motion $B$, the integral $I(X)=\int_0^{\cdot}X_s\,dB_s$ is a martingale in $\cM_2^c$ with $I_0(X)=0$, and by \bk{(2.24)} (valid for integrands in $\cP^*$, as noted before \bk{Proposition 2.24}) and $\tau_n\le n$,
\[
  \int_0^{\tau_n}g'(R_s)\,dB_s=I_n(X)\quad\text{a.s.},\qquad\text{so that}\qquad
  E\int_0^{\tau_n}g'(R_s)\,dB_s=0 .
\]
Taking expectations in \eqref{eq:3-3-23:ito} gives $E f(R_{\tau_n})=f(r)$ for every $n\ge1$. As $n\to\infty$ we have $\tau_n\to T_a\wedge S_b$, which is finite a.s.; hence $R_{\tau_n}\to R_{T_a\wedge S_b}$ a.s., and $\abs{f(R_{\tau_n})}\le\max_{a\le x\le b}\abs{f(x)}$. The bounded convergence theorem gives $E f(R_{T_a\wedge S_b})=f(r)$. On $\{T_a\wedge S_b<\infty\}$ the times $T_a$ and $S_b$ are different (at a common finite value $t$ we would have $a=R_t=b$), and $R_{T_a}=a$ on $\{T_a<\infty\}$, $R_{S_b}=b$ on $\{S_b<\infty\}$, because the level sets of a continuous path are closed. Therefore, with $p\defeq P[T_a<S_b]$, we have $P[S_b<T_a]=1-p$ and
\[
  f(r)=f(a)\,p+f(b)\,(1-p),
\]
which is \eqref{eq:3-3-23:exit}, since $f(a)\ne f(b)$.

\emph{Step 2. For $0<a<r$,}
\begin{equation}\label{eq:3-3-23:hit}
  P[T_a<\infty]=\lim_{k\to\infty}\frac{f(k)-f(r)}{f(k)-f(a)}
  =\begin{cases} 1, & d=2,\\[2pt] (a/r)^{d-2}, & d\ge3.\end{cases}
\end{equation}
Let $k$ range over the integers larger than $r$. If $S_k(\omega)\le t$, then $\max_{0\le s\le t}R_s(\omega)\ge k$; since every path of $R$ is bounded on $[0,t]$, this fails for large $k$, and so $S_k\to\infty$ everywhere. Moreover $S_k\le S_{k+1}$ a.s.\ (on $\{R_0=r\}$ the path must pass through level $k$ before reaching $k+1$). Hence, up to null sets, the events $\{T_a<S_k\}$ increase with $k$, each of them is contained in $\{T_a<\infty\}$, and every $\omega$ with $T_a(\omega)<\infty$ belongs to $\{T_a<S_k\}$ for all large $k$. Thus $P[T_a<\infty]=\lim_{k\to\infty}P[T_a<S_k]$, and \eqref{eq:3-3-23:exit} gives the first equality in \eqref{eq:3-3-23:hit}. For $d=2$,
\[
  \frac{\log k-\log r}{\log k-\log a}\longrightarrow 1,
\]
and for $d\ge3$, since $k^{2-d}\to0$,
\[
  \frac{k^{2-d}-r^{2-d}}{k^{2-d}-a^{2-d}}\longrightarrow\frac{r^{2-d}}{a^{2-d}}=\Bigl(\frac ar\Bigr)^{d-2}.
\]

\emph{Step 3. Conclusion.} Let $0<a<r$. On $\{T_a<\infty\}$ we have $m\le R_{T_a}=a$; conversely, if $m(\omega)<a$ and $R_0(\omega)=r$, then $R_t(\omega)<a<R_0(\omega)$ for some $t$, and the intermediate value theorem gives $T_a(\omega)<\infty$. Hence
\begin{equation}\label{eq:3-3-23:sandwich}
  P[m<a]\le P[T_a<\infty]\le P[m\le a];\qquad 0<a<r .
\end{equation}

\textbf{(i)} Let $d=2$. By \eqref{eq:3-3-23:hit} and \eqref{eq:3-3-23:sandwich}, $P[m\le a]=1$ for every $a\in(0,r)$. Taking $a=1/j$ for the integers $j>1/r$ and intersecting these countably many events, we obtain $P[m\le0]=1$; since $m\ge0$, this means $m=0$ a.s.

\textbf{(ii)} Let $d\ge3$. For $0<c<r$, \eqref{eq:3-3-23:sandwich} and \eqref{eq:3-3-23:hit} give $P[m\le c]\ge(c/r)^{d-2}$, and also, for every $a\in(c,r)$,
\[
  P[m\le c]\le P[m<a]\le\Bigl(\frac ar\Bigr)^{d-2};
\]
letting $a\downarrow c$ we conclude $P[m\le c]=(c/r)^{d-2}$. For $c=0$ the same upper bound, $P[m\le0]\le(a/r)^{d-2}$ for all $a\in(0,r)$, yields $P[m\le0]=0=(0/r)^{d-2}$ because $d-2\ge1$. For $c=r$ we have $m\le R_0=r$ a.s., so $P[m\le r]=1$. This proves
\[
  P[m\le c]=\Bigl(\frac cr\Bigr)^{d-2};\qquad 0\le c\le r .
\]
In other words, $m/r$ has the density $(d-2)u^{d-3}$ on $(0,1)$, which is the beta distribution with parameters $(d-2,1)$.

\rmk{For $d=2$, \bk{Proposition 3.22} shows that the infimum $m=0$ is a.s.\ not attained. The function $f$ is, up to an affine change, the unique solution of \eqref{eq:3-3-23:harm}; equivalently, $x\mapsto f(\norm{x})$ is the radial harmonic function on $\R^d\setminus\{0\}$, which is why $f(R_{t\wedge T_a\wedge S_b})$ is a bounded martingale.}
\end{kex}

\begin{kex}{Problem}{3.24}{}
We use the notation of \bk{Definition 3.19}: $W=\{W_t,\F_t;\ 0\le t<\infty\}$, $\{P^x\}_{x\in\R^d}$ is a $d$-dimensional Brownian family on $(\Omega,\F)$, $R_t=\norm{W_t}$, and $\wh{P}^y\defeq P^{(y,0,\dots,0)}$ for $y\ge 0$; the measure $P$ of the statement is $\wh{P}^r$. Every path of $R$ is continuous. Consequently
\[
  \liminf_{t\to\infty}R_t=\lim_{n\to\infty}\ \inf_{q\in\Q,\ q\ge n}R_q
\]
is a random variable with values in $[0,\infty]$, and the assertion is that it equals $\infty$, a.s.\ $P$.

\emph{Step 1: the passage times of high levels are finite.} For every integer $k>r$ let
\[
  S_k\defeq\inf\{t\ge 0;\ R_t=k\}.
\]
This is the hitting time of the closed set $\{k\}$ by the continuous, adapted process $R$, hence a stopping time of $\{\F_t\}$ (\bk{Problem 1.2.7}), and in particular an optional time. Under $P$, the process $W^{(1)}_t-W^{(1)}_0$, $0\le t<\infty$, is a standard, one-dimensional Brownian motion, so \bk{Theorem 2.9.23}\,(iii) gives $\limsup_{t\to\infty}W^{(1)}_t=\infty$, a.s.\ $P$. Since $R_t\ge\abs{W^{(1)}_t}$, we have $\sup_{0\le t<\infty}R_t=\infty$, a.s.\ $P$. On the intersection of this event with $\{R_0=r\}$, which has probability one, the continuous path $t\mapsto R_t$ starts below every integer $k>r$ and exceeds it at some time, and thus takes the value $k$ by the intermediate value theorem. Therefore
\begin{equation}\label{eq:3-3-24:a}
  P[S_k<\infty \text{ for every integer } k>r]=1 .
\end{equation}
On $\{S_k<\infty\}$ we have $R_{S_k}=k$, because the path is continuous and the infimum defining $S_k$ is attained.

\emph{Step 2: the process after $S_k$.} Fix a real number $c>0$ and let
\[
  G_c\defeq\Bigl\{\omega\in[0,\infty)^{[0,\infty)};\ \inf_{q\in\Q\cap[0,\infty)}\omega(q)\le c\Bigr\}
  =\bigcap_{j=1}^\infty\ \bigcup_{q\in\Q\cap[0,\infty)}\Bigl\{\omega;\ \omega(q)<c+\frac1j\Bigr\}.
\]
This set belongs to the product $\sigma$-field $\B\bigl([0,\infty)^{[0,\infty)}\bigr)$, since it is obtained by countable operations from sets which depend on one coordinate only. For a continuous function $\omega$ the infimum over the nonnegative rationals coincides with the infimum over $[0,\infty)$; hence, for every $y>0$, the continuity of the paths of $R$ and \bk{Problem 3.23}\,(ii) (which requires $d\ge 3$ and a starting point $y>0$) yield
\begin{equation}\label{eq:3-3-24:b}
  \wh{P}^y[R_\cdot\in G_c]=\wh{P}^y\Bigl[\inf_{0\le t<\infty}R_t\le c\Bigr]=\Bigl(\frac{c}{y}\Bigr)^{d-2};\qquad 0<c\le y .
\end{equation}

By \bk{Problem 3.20}, the Bessel family $R$, $\{\wh{P}^y\}_{y\ge 0}$ is a strong Markov family with state space $[0,\infty)$. Accordingly, it satisfies conditions (c$'$) and (d$'$) of \bk{Proposition 2.6.6}; the proof of that proposition, which is the proof of \bk{Proposition 2.5.15} with the constant $s$ replaced by an optional time, makes no use of the particular state space. Let $k$ be an integer with $k>\max\{r,c\}$, and put
\[
  A_k\defeq\{S_k<\infty,\ R_{S_k+\cdot}\in G_c\}
  =\Bigl\{S_k<\infty,\ \inf_{0\le t<\infty}R_{S_k+t}\le c\Bigr\},
\]
in accordance with the convention of \bk{Remark 2.6.4}; the second expression is again justified by the continuity of the paths. By \eqref{eq:3-3-24:a} and the last observation of Step 1, the measure $P R_{S_k}^{-1}$ is the unit mass at the point $k$. Condition (d$'$), with $x=r$, $S=S_k$ and $F=G_c$, holds for $P R_{S_k}^{-1}$-almost every $y$, and therefore holds at $y=k$:
\[
  P[R_{S_k+\cdot}\in G_c\mid R_{S_k}=k]=\wh{P}^k[R_\cdot\in G_c].
\]
Integrating the function $y\mapsto P[R_{S_k+\cdot}\in G_c\mid R_{S_k}=y]$ with respect to $P R_{S_k}^{-1}$ and using \eqref{eq:3-3-24:b} with $y=k$, we obtain
\begin{equation}\label{eq:3-3-24:c}
  P(A_k)=\wh{P}^k[R_\cdot\in G_c]=\Bigl(\frac{c}{k}\Bigr)^{d-2};\qquad k>\max\{r,c\}.
\end{equation}

\emph{Step 3: conclusion.} Let $\Omega_0$ be the event in \eqref{eq:3-3-24:a}, and let $\omega\in\Omega_0$ satisfy $\liminf_{t\to\infty}R_t(\omega)<c$. For every integer $k>\max\{r,c\}$ there is then a time $t\ge S_k(\omega)$ with $R_t(\omega)<c$, so that $\inf_{0\le u<\infty}R_{S_k(\omega)+u}(\omega)\le c$, i.e., $\omega\in A_k$. From \eqref{eq:3-3-24:c} and $d-2\ge 1$ we conclude that
\[
  P\Bigl[\liminf_{t\to\infty}R_t<c\Bigr]\le P(A_k)=\Bigl(\frac{c}{k}\Bigr)^{d-2}\longrightarrow 0\qquad\text{as } k\to\infty ,
\]
and thus $P[\liminf_{t\to\infty}R_t<c]=0$. Taking the union over $c=1,2,\dots$, we arrive at
\[
  P\Bigl[\liminf_{t\to\infty}R_t<\infty\Bigr]=0,\qquad\text{i.e.,}\qquad P\Bigl[\lim_{t\to\infty}R_t=\infty\Bigr]=1 .
\]
\qed

\rmk{The starting point $r=0$, excluded in \bk{Problem 3.23}, causes no difficulty here, because \bk{Problem 3.23}\,(ii) is applied only to the process after the time $S_k$, which starts afresh at $k>0$. The hypothesis $d\ge 3$ cannot be relaxed: for $d=2$ and $r>0$, \bk{Problem 3.23}\,(i) gives $\inf_{0\le t<\infty}R_t=0$ a.s., whereas $R_t>0$ for all $t$ a.s.\ by \bk{Proposition 3.22}; hence $\liminf_{t\to\infty}R_t=0$, a.s.}
\end{kex}

\begin{kex}{Exercise}{3.25}{}
Throughout, $\{\F_t\}$ satisfies the usual conditions, as is assumed in Sections~3.2 and~3.3 of the
book. We write
\[
  A\defeq E\int_0^T\abs{X_t}^{2m}\,dt<\infty .
\]

\emph{Step 1. The martingale $M$.} Hypothesis \bk{(3.24)} involves $X$ only on
$[0,T]$ and says nothing about $X_t$ for $t>T$; in particular, $X$ itself need not belong to any of
the classes of integrands of Section~3.2. We set
\[
  \wt X_t(\omega)\defeq X_t(\omega)\ind_{[0,T]}(t);\qquad 0\le t<\infty,\ \omega\in\Omega,
\]
which is again a measurable, adapted process, and we understand the integral in \bk{(3.25)} as
$\int_0^T\wt X_t\,dW_t$, the value at time $T$ of the stochastic integral of $\wt X$. This is
consistent with the theory of Section~3.2: if $X$ happens to belong to the class $\cP$ of
\bk{Definition 2.22} (the largest class of integrands for $W$ considered in the book, containing
$\cL$), then $\int_0^T X_t\,dW_t=\int_0^T\wt X_t\,dW_t$ a.s. Indeed, this is \bk{(2.24)} of
\bk{Proposition 2.10}, applied with the constant stopping time $T$ and with $t=T$; \bk{(2.24)} is
valid for integrands in $\cL$ by \bk{Remark 2.11}, because $\qv{W}_t=t$ is absolutely continuous,
and the discussion following \bk{Definition 2.23} extends it by localization to integrands in
$\cP^*$; the same localization applies to integrands in $\cP$, for which the integral with respect
to $W$ is defined, again because $\qv{W}_t=t$ is absolutely continuous, as indicated in the paragraph
following \bk{Definition 2.22} (with the aid of \bk{Remark 2.11}).
Since the right-hand side of \bk{(3.25)} and the quantity $A$ are the same for $X$ and $\wt X$, we
may and do write $X$ for $\wt X$ from now on, i.e., we assume $X_t=0$ for $t>T$. H\"older's inequality on
$[0,T]\times\Omega$ with respect to the product of Lebesgue measure and $P$, of total mass $T$, and
with the exponents $m$ and $m/(m-1)$ (nothing is needed if $m=1$), gives
\begin{equation}\label{eq:3-3-25:L2}
  E\int_0^\infty X_t^2\,dt=E\int_0^T X_t^2\,dt\le T^{(m-1)/m}A^{1/m}<\infty .
\end{equation}
Hence $X$ belongs to the class $\cL=\cL(W)$ of Section~3.2, and \bk{Remark 2.11} shows that the
stochastic integral
\[
  M_t\defeq\int_0^t X_s\,dW_s;\qquad 0\le t<\infty,
\]
is defined, belongs to $\cM_2^c$, and has all the properties listed in \bk{Proposition 2.10}; in
particular, by \bk{(2.19)},
\begin{equation}\label{eq:3-3-25:qv}
  \qv{M}_t=\int_0^t X_s^2\,ds;\qquad 0\le t<\infty,\quad\as
\end{equation}
The left-hand side of \bk{(3.25)} is $E\abs{M_T}^{2m}$. By \eqref{eq:3-3-25:L2},
$\int_0^T X_s^2\,ds<\infty$ a.s. After discarding a $P$-null event, which belongs to $\F_0$, we may
and do assume that every path of $M$ is continuous and starts at $M_0=0$.

\emph{Step 2. It\^o's rule.} The function $f(x)\defeq\abs{x}^{2m}$ is of class $C^2$ on $\R$ because
$2m\ge2$, with
\[
  f'(x)=2m\abs{x}^{2m-1}\sgn(x),\qquad f''(x)=2m(2m-1)\abs{x}^{2m-2}
\]
(for $x\ne0$ this is elementary; at the origin $f(h)/h\to0$ and $f'(h)/h=2m\abs{h}^{2m-2}$ tends to
$0$ if $m>1$ and equals $2$ if $m=1$, so the formulas hold at $x=0$ as well, with the convention
$\abs{0}^0=1$, and $f''$ is continuous). The process $M\in\cM_2^c\subseteq\cM^{c,\loc}$ is a
continuous semimartingale (\bk{Definition 3.1}) whose decomposition \bk{(3.1)} has initial value
$M_0=0$, local martingale part $M$ itself, and $B\equiv0$ (that is, $A^+\equiv A^-\equiv0$ in
\bk{(3.2)}). By \bk{Theorem 3.3} and \bk{Remark 3.4}, and since $f(M_0)=f(0)=0$, the process
\[
  N_t\defeq\int_0^t f'(M_s)\,dM_s;\qquad 0\le t<\infty,
\]
is a well-defined element of $\cM^{c,\loc}$, and a.s.
\[
  \abs{M_t}^{2m}=N_t+\frac12\int_0^t f''(M_s)\,d\qv{M}_s;\qquad 0\le t<\infty .
\]
On the event where \eqref{eq:3-3-25:qv} holds for all $t$, the measure $d\qv{M}_s$ on $[0,\infty)$
has density $X_s^2$ with respect to Lebesgue measure, so that a.s.
\begin{equation}\label{eq:3-3-25:ito}
  \abs{M_t}^{2m}=N_t+m(2m-1)\int_0^t \abs{M_s}^{2m-2}X_s^2\,ds;\qquad 0\le t<\infty .
\end{equation}

\emph{Step 3. Localization.} For $n\ge1$ let
\[
  \tau_n\defeq\inf\{t\ge0;\ \abs{M_t}\ge n\},
\]
a stopping time of $\{\F_t\}$ by \bk{Problem 1.2.7}, since $M$ is continuous and adapted and
$\{x;\,\abs{x}\ge n\}$ is closed. Because $M_0=0$ and $M$ is continuous, $\abs{M_s}\le n$ for
$0\le s\le\tau_n$, $s<\infty$. We set
\[
  g_n(t)\defeq E\abs{M_{t\wedge\tau_n}}^{2m}\le n^{2m};\qquad 0\le t\le T,
\]
and claim that
\begin{equation}\label{eq:3-3-25:gn}
  g_n(t)=m(2m-1)\,E\int_0^{t\wedge\tau_n}\abs{M_s}^{2m-2}X_s^2\,ds;\qquad 0\le t\le T .
\end{equation}
(The integrand $(s,\omega)\mapsto\ind_{\{s\le\tau_n(\omega)\}}\abs{M_s(\omega)}^{2m-2}X_s^2(\omega)$
is nonnegative and $\B([0,T])\otimes\F$-measurable, so the right-hand side makes sense by Tonelli's
theorem.) In view of \eqref{eq:3-3-25:ito} evaluated at the time $t\wedge\tau_n$, the claim amounts
to $E N_{t\wedge\tau_n}=0$. To prove this, let $\{S_k\}_{k=1}^\infty$ be a nondecreasing sequence of
stopping times with $S_k\uparrow\infty$ a.s.\ such that $\{N_{t\wedge S_k},\F_t;\ 0\le t<\infty\}$
is a martingale for every $k$ (\bk{Definition 1.5.15}). By \bk{Problem 1.3.24}(i) applied to this
continuous martingale and the stopping time $\tau_n$, the process
$\{N_{t\wedge S_k\wedge\tau_n},\F_t;\ 0\le t<\infty\}$ is a martingale, whence
\begin{equation}\label{eq:3-3-25:zero}
  E N_{t\wedge S_k\wedge\tau_n}=E N_0=0;\qquad 0\le t\le T,\ k\ge1 .
\end{equation}
On the other hand, \eqref{eq:3-3-25:ito} at the time $t\wedge S_k\wedge\tau_n\le T\wedge\tau_n$ and
the bound $\abs{M_s}\le n$ for $s\le\tau_n$ give, a.s.,
\[
  \abs{N_{t\wedge S_k\wedge\tau_n}}
  \le n^{2m}+m(2m-1)\,n^{2m-2}\int_0^T X_s^2\,ds ;\qquad 0\le t\le T,\ k\ge1,
\]
and the right-hand side is integrable by \eqref{eq:3-3-25:L2} and does not depend on $k$. Since
$N_{t\wedge S_k\wedge\tau_n}\to N_{t\wedge\tau_n}$ a.s.\ as $k\to\infty$, dominated convergence in
\eqref{eq:3-3-25:zero} yields $E N_{t\wedge\tau_n}=0$, and \eqref{eq:3-3-25:gn} follows by taking
expectations in \eqref{eq:3-3-25:ito} at the time $t\wedge\tau_n$ (all three terms are integrable).

The right-hand side of \eqref{eq:3-3-25:gn} is nondecreasing in $t$. Therefore
\begin{equation}\label{eq:3-3-25:mono}
  g_n(s)\le g_n(T)\le n^{2m}<\infty;\qquad 0\le s\le T .
\end{equation}

\emph{Step 4. The estimate for the stopped process.} We show that
\begin{equation}\label{eq:3-3-25:stopped}
  g_n(T)=E\abs{M_{T\wedge\tau_n}}^{2m}\le\bigl(m(2m-1)\bigr)^m\,T^{m-1}A;\qquad n\ge1 .
\end{equation}
If $m=1$, then \eqref{eq:3-3-25:gn} with $t=T$ reads $g_n(T)=E\int_0^{T\wedge\tau_n}X_s^2\,ds\le A$,
which is \eqref{eq:3-3-25:stopped}. Let now $m>1$. We apply H\"older's inequality on
$[0,T]\times\Omega$, with respect to the product of Lebesgue measure and $P$ and with the conjugate
exponents $m/(m-1)$ and $m$, to the functions $\ind_{\{s\le\tau_n\}}\abs{M_s}^{2m-2}$ and $X_s^2$:
\begin{equation}\label{eq:3-3-25:holder}
  E\int_0^{T\wedge\tau_n}\abs{M_s}^{2m-2}X_s^2\,ds
  \le\Bigl(E\int_0^{T}\ind_{\{s\le\tau_n\}}\abs{M_s}^{2m}\,ds\Bigr)^{(m-1)/m}
     \Bigl(E\int_0^T\abs{X_s}^{2m}\,ds\Bigr)^{1/m}.
\end{equation}
For $s\le\tau_n$ we have $M_s=M_{s\wedge\tau_n}$; hence, by Tonelli's theorem and
\eqref{eq:3-3-25:mono},
\[
  E\int_0^{T}\ind_{\{s\le\tau_n\}}\abs{M_s}^{2m}\,ds
  \le\int_0^T E\abs{M_{s\wedge\tau_n}}^{2m}\,ds=\int_0^T g_n(s)\,ds\le T\,g_n(T).
\]
Substituting this and \eqref{eq:3-3-25:holder} into \eqref{eq:3-3-25:gn} with $t=T$, we obtain
\[
  g_n(T)\le m(2m-1)\,T^{(m-1)/m}\,g_n(T)^{(m-1)/m}A^{1/m}.
\]
If $g_n(T)=0$, \eqref{eq:3-3-25:stopped} is trivial. Otherwise $0<g_n(T)<\infty$ by
\eqref{eq:3-3-25:mono}, and we may divide by $g_n(T)^{(m-1)/m}$ to get
$g_n(T)^{1/m}\le m(2m-1)\,T^{(m-1)/m}A^{1/m}$; raising this to the power $m$ gives
\eqref{eq:3-3-25:stopped}.

\emph{Step 5. Removing the localization.} Every path of $M$ is bounded on $[0,T]$, so for each
$\omega$ we have $\tau_n(\omega)>T$ for all sufficiently large $n$, and therefore
$M_{T\wedge\tau_n}\to M_T$ as $n\to\infty$. Fatou's lemma and \eqref{eq:3-3-25:stopped} yield
\[
  E\Bigl|\int_0^T X_t\,dW_t\Bigr|^{2m}=E\abs{M_T}^{2m}
  \le\liminf_{n\to\infty}E\abs{M_{T\wedge\tau_n}}^{2m}
  \le\bigl(m(2m-1)\bigr)^m\,T^{m-1}\,E\int_0^T\abs{X_t}^{2m}\,dt,
\]
which is \bk{(3.25)}.

\rmk{For $m=1$ the constant is $1$ and \bk{(3.25)} holds with equality, by \bk{(2.14)}. For $m>1$
the constant is not sharp: with $X\equiv1$ and $m=2$ the left-hand side is $E(W_T^4)=3T^2$, whereas
the right-hand side is $36\,T^2$. The submartingale property of $\abs{M_t}^{2m}$ mentioned in the
hint was replaced in Step~3 by the monotonicity of $g_n$, read off directly from
\eqref{eq:3-3-25:gn}.}
\end{kex}

\begin{kex}{Problem}{3.29}{}
The right-hand inequality of \bk{(3.38)} is proved in the book's Section 3.8, with $\Lambda_m=K_m d^{m+1}$, where $K_m$ is the constant of \bk{(3.37)}; the left-hand inequality is left there to the reader (``a similar proof''), and we prove it here, with
\[
  \lambda_m\defeq k_m\,d^{-(m+1)},
\]
where $k_m>0$ is the constant on the left-hand side of \bk{(3.37)}.

\emph{Conventions.} The stopping time $T$ may take the value $+\infty$. As in \bk{Theorem 3.28}, for a continuous process $X$ with values in $\R$ or $\R^d$ the random variable $\norm{X}^*_T$ (written $X^*_T$ if $X$ is real-valued, in accordance with \bk{(3.26)}) is $\sup\{\norm{X_s};\ 0\le s\le T,\ s<\infty\}$, and for a nondecreasing process $B$ we put $B_T=B_\infty\defeq\lim_{t\to\infty}B_t$ on $\{T=\infty\}$. These random variables take values in $[0,\infty]$; all the expectations below are expectations of nonnegative random variables, and the inequalities between them are understood in $[0,\infty]$. In this sense \bk{(3.37)} holds for every $N\in\cM^{c,\loc}$ and every stopping time $T$.

\emph{Step 1: an elementary inequality.} For $m>0$ and $x_1,\dots,x_d\in[0,\infty]$ we have
\begin{equation}\label{eq:3-3-29:a}
  (x_1+\dots+x_d)^m\le d^m\,(x_1^m+\dots+x_d^m).
\end{equation}
Indeed, $x_1+\dots+x_d\le d\max_{1\le i\le d}x_i$, and since $x\mapsto x^m$ is nondecreasing on $[0,\infty]$,
\[
  (x_1+\dots+x_d)^m\le d^m\bigl(\max_{1\le i\le d}x_i\bigr)^m
  =d^m\max_{1\le i\le d}x_i^m\le d^m\,(x_1^m+\dots+x_d^m).
\]

\emph{Step 2: comparison of the maximal functions.} For every $1\le i\le d$ and $0\le s<\infty$ we have $\abs{M^{(i)}_s}\le\norm{M_s}$. Taking the supremum over $\{s;\ 0\le s\le T(\omega),\ s<\infty\}$ for each $\omega$ gives
\begin{equation}\label{eq:3-3-29:b}
  \bigl(M^{(i)}\bigr)^*_T\le\norm{M}^*_T;\qquad 1\le i\le d .
\end{equation}

\emph{Step 3: the lower bound.} Applying \eqref{eq:3-3-29:a} pointwise with $x_i=\qv{M^{(i)}}_T\ge 0$ and taking expectations, then using the left-hand side of \bk{(3.37)} for each $M^{(i)}\in\cM^{c,\loc}$ and the stopping time $T$, and finally \eqref{eq:3-3-29:b}, we obtain
\begin{align*}
  E(A_T^m)
  &\le d^m\sum_{i=1}^d E\bigl(\qv{M^{(i)}}_T^m\bigr)
   \le \frac{d^m}{k_m}\sum_{i=1}^d E\Bigl[\bigl((M^{(i)})^*_T\bigr)^{2m}\Bigr]\\
  &\le \frac{d^m}{k_m}\sum_{i=1}^d E\bigl[(\norm{M}^*_T)^{2m}\bigr]
   =\frac{d^{m+1}}{k_m}\,E\bigl[(\norm{M}^*_T)^{2m}\bigr].
\end{align*}
This is $\lambda_m E(A_T^m)\le E(\norm{M}^*_T)^{2m}$ with $\lambda_m=k_m d^{-(m+1)}$, and it holds also when one or both sides are infinite. Together with the upper bound proved in the book, \bk{(3.38)} is established.
\rmk{The constants $\lambda_m=k_m d^{-(m+1)}$ and $\Lambda_m=K_m d^{m+1}$ obtained in this way depend on $m$ and on the dimension $d$, but on nothing else: not on the probability space, the filtration, the vector $M$ or the stopping time $T$. This is the sense in which they are ``universal''.}
\end{kex}

\begin{kex}{Exercise}{3.31}{}
Throughout, as everywhere in \bk{Section 3.3}, the filtration $\{\F_t\}$ satisfies the usual conditions; $M\in\cM^{c,\loc}$ is a continuous local martingale with $M_0=0$ a.s.\ (\bk{Definition 1.5.15}), and $\qv{M}$ is its quadratic variation, a continuous, nondecreasing, adapted process with $\qv{M}_0=0$ (\bk{Problem 1.5.17}). The process $\qv{M}$, and likewise every stochastic integral $I^M(X)$ of \bk{Definition 2.23}, is determined only up to indistinguishability, and its defining properties need hold only off a $P$-null set. We adopt the following convention: whenever such a process $Z$ has continuous paths (nondecreasing ones, in the case of $\qv{M}$) off an event $N\in\F$ with $P(N)=0$, we replace it by $Z\ind_{N^c}$. Since $\{\F_t\}$ satisfies the usual conditions, $N\in\F_0$, so $Z\ind_{N^c}$ is again adapted; it is indistinguishable from $Z$, and \emph{all} its paths are continuous (and nondecreasing, for $\qv{M}$). In the same way we take $M$ itself with all paths continuous. None of this affects any statement below, each of which is an assertion about indistinguishable processes.

\emph{Step 1. The polynomials $H_n$.}
For fixed $x,y\in\R$ the function $\alpha\mapsto\exp(\alpha x-\frac12\alpha^2y)$ is entire, so its Taylor series at $\alpha=0$, whose coefficients are $H_n(x,y)/n!$ by the very definition of $H_n$, converges to it for every $\alpha$:
\begin{equation}\label{eq:3-3-31:a}
  \exp\Bigl(\alpha x-\frac12\alpha^2y\Bigr)=\sum_{n=0}^\infty\frac{\alpha^n}{n!}\,H_n(x,y);\qquad \alpha,x,y\in\R .
\end{equation}
On the other hand, multiplying the absolutely convergent series $e^{\alpha x}=\sum_{j\ge0}\alpha^jx^j/j!$ and $e^{-\alpha^2y/2}=\sum_{k\ge0}\alpha^{2k}(-y/2)^k/k!$ and collecting the terms with $j+2k=n$, we obtain another power series in $\alpha$ for the left-hand side of \eqref{eq:3-3-31:a}. By uniqueness of the coefficients of a power series,
\begin{equation}\label{eq:3-3-31:b}
  H_n(x,y)=n!\sum_{0\le k\le n/2}\frac{x^{n-2k}}{(n-2k)!}\,\frac{(-y/2)^k}{k!};\qquad n=0,1,2,\dots,
\end{equation}
which shows that $H_n$ is a polynomial in $(x,y)$ (and reproduces $H_0,\dots,H_4$ as listed in the statement). In particular $H_n$ is of class $C^2$ on $\R^2$, and $H_n(0,0)=0$ for $n\ge1$, because every term of \eqref{eq:3-3-31:b} contains the factor $x^{n-2k}$ or the factor $y^k$ with a positive exponent.

We verify \bk{(3.40)} and \bk{(3.41)} from \eqref{eq:3-3-31:b}. Differentiating in $x$ annihilates the term with $n-2k=0$ (if any), and for $n\ge1$
\[
  \frac{\partial}{\partial x}H_n(x,y)=n\cdot(n-1)!\sum_{0\le k\le (n-1)/2}\frac{x^{n-1-2k}}{(n-1-2k)!}\,\frac{(-y/2)^k}{k!}=nH_{n-1}(x,y),
\]
which is \bk{(3.40)}. Applying it twice, $\frac{\partial^2}{\partial x^2}H_n=n(n-1)H_{n-2}$ for $n\ge2$, while $\frac{\partial^2}{\partial x^2}H_n=0$ for $n=0,1$. Differentiating \eqref{eq:3-3-31:b} in $y$ annihilates the term $k=0$, and with $j=k-1$ we find for $n\ge2$
\[
  \frac{\partial}{\partial y}H_n(x,y)=-\frac12\,n!\sum_{0\le j\le (n-2)/2}\frac{x^{n-2-2j}}{(n-2-2j)!}\,\frac{(-y/2)^{j}}{j!}=-\frac12\,n(n-1)H_{n-2}(x,y),
\]
while $\frac{\partial}{\partial y}H_n=0$ for $n=0,1$. In all cases $\frac{\partial}{\partial y}H_n+\frac12\frac{\partial^2}{\partial x^2}H_n=0$, which is \bk{(3.41)}.

\emph{Step 2. Proof of} (ii). For every $\omega\in\Omega$, $t\ge0$ and $\alpha\in\R$, the identity \eqref{eq:3-3-31:a} with $x=M_t(\omega)$, $y=\qv{M}_t(\omega)$ is the expansion (ii); the series converges absolutely. No property of $M$ is used here.

\emph{Step 3. Proof of} (i). The iterated integral is understood recursively: $I^{(0)}_t\defeq1$ and
\[
  I^{(n)}_t\defeq\int_0^tI^{(n-1)}_s\,dM_s;\qquad 0\le t<\infty,\ n\ge1,
\]
so that $I^{(n)}_t$ is the left-hand side of (i). We prove by induction on $n\ge0$ the assertion
\begin{quote}
$(A_n)$: $I^{(n)}$ is a well-defined continuous, adapted process, and a.s.\ $I^{(n)}_t=\frac1{n!}H_n(M_t,\qv{M}_t)$ for all $0\le t<\infty$.
\end{quote}
$(A_0)$ is trivial since $H_0\equiv1$. Let $n\ge1$ and assume $(A_{n-1})$.

By $(A_{n-1})$ and our convention on versions, $I^{(n-1)}$ is adapted and every one of its paths is continuous; hence it is progressively measurable (\bk{Proposition 1.1.13}), and each of its paths is bounded on every interval $[0,T]$. As every path of $\qv{M}$ is finite and nondecreasing, $\int_0^T(I^{(n-1)}_s)^2\,d\qv{M}_s\le\sup_{0\le s\le T}(I^{(n-1)}_s)^2\cdot\qv{M}_T<\infty$ for every $\omega$ and every $T<\infty$. Thus $I^{(n-1)}\in\cP^*$ (\bk{Definition 2.22}), and by \bk{Definition 2.23} the integral $I^{(n)}=I^M(I^{(n-1)})$ is well defined and belongs to $\cM^{c,\loc}$; in particular it is adapted and, in the version fixed by our convention, all its paths are continuous. The same reasoning shows that the process $Y_t\defeq H_{n-1}(M_t,\qv{M}_t)$, which is adapted and has only continuous paths because $H_{n-1}$ is continuous and $M$, $\qv{M}$ are, belongs to $\cP^*$. By $(A_{n-1})$ the processes $Y$ and $(n-1)!\,I^{(n-1)}$ are indistinguishable, so they define the same element of $\cP^*$, and consequently, by the linearity \bk{(2.17)} of the integral in its integrand, which remains valid for $M\in\cM^{c,\loc}$ and integrands in $\cP^*$ (see the paragraph following \bk{Definition 2.23}),
\begin{equation}\label{eq:3-3-31:c}
  \int_0^tH_{n-1}(M_s,\qv{M}_s)\,dM_s=(n-1)!\,I^{(n)}_t;\qquad0\le t<\infty,\quad\text{a.s.}
\end{equation}

We now apply It\^o's rule in the form of \bk{Theorem 3.6} with $d=2$, with the vector of local martingales $(M^{(1)},M^{(2)})\defeq(M,0)$ in $\cM^{c,\loc}$, the vector $(B^{(1)},B^{(2)})\defeq(0,\qv{M})$ of continuous, adapted processes of bounded variation with $B_0=0$, the initial vector $X_0\defeq(0,0)$, and the function $f(t,x,y)\defeq H_n(x,y)$, which does not depend on $t$ and is a polynomial, hence of class $C^{1,2}$ on $[0,\infty)\times\R^2$. Then $X_t=X_0+M_t+B_t=(M_t,\qv{M}_t)$, and $\qv{M^{(i)},M^{(j)}}\equiv0$ unless $i=j=1$, in which case it is $\qv{M}$. Since $\partial f/\partial t=0$, $B^{(1)}\equiv0$, $M^{(2)}\equiv0$ and $f(0,X_0)=H_n(0,0)=0$ (Step~1; here $n\ge1$), formula \bk{(3.5)} reads, a.s.,
\begin{align*}
  H_n(M_t,\qv{M}_t)
  &=\int_0^t\frac{\partial H_n}{\partial x}(M_s,\qv{M}_s)\,dM_s
   +\int_0^t\Bigl[\frac{\partial H_n}{\partial y}+\frac12\frac{\partial^2H_n}{\partial x^2}\Bigr](M_s,\qv{M}_s)\,d\qv{M}_s\\
  &=n\int_0^tH_{n-1}(M_s,\qv{M}_s)\,dM_s;\qquad0\le t<\infty,
\end{align*}
where the Lebesgue--Stieltjes integral vanishes by \bk{(3.41)} and the integrand of the stochastic integral was rewritten by \bk{(3.40)}. Combining this with \eqref{eq:3-3-31:c}, we obtain a.s.
\[
  H_n(M_t,\qv{M}_t)=n\,(n-1)!\,I^{(n)}_t=n!\,I^{(n)}_t;\qquad0\le t<\infty,
\]
which is $(A_n)$. This completes the induction and proves (i): the two sides are indistinguishable processes, for every $n\ge1$ (and trivially for $n=0$, with the empty integral read as $1$).

\rmk{The hypothesis $M_0=0$ contained in $M\in\cM^{c,\loc}$ is needed for (i): for a continuous local martingale with $M_0\ne0$ the left-hand side of (i) vanishes at $t=0$, whereas the right-hand side equals $M_0^n/n!$. Together, (i) and (ii) exhibit the exponential local martingale $Z_t=\exp(\alpha M_t-\frac12\alpha^2\qv{M}_t)$ as the sum $\sum_{n\ge0}\alpha^nI^{(n)}_t$ of its iterated integrals, in agreement with the equation $Z_t=1+\alpha\int_0^tZ_s\,dM_s$ which It\^o's rule gives for $Z$ (compare \bk{(3.6)}).}

\rmk{The relation with the Hermite polynomials stated in parentheses is checked as follows. Since $\alpha x-\frac12\alpha^2=\frac12x^2-\frac12(x-\alpha)^2$ and $\frac{\partial^n}{\partial\alpha^n}g(x-\alpha)=(-1)^ng^{(n)}(x-\alpha)$ for smooth $g$, we have
\[
  H_n(x,1)=e^{x^2/2}\,\frac{\partial^n}{\partial\alpha^n}e^{-(x-\alpha)^2/2}\Big|_{\alpha=0}=(-1)^ne^{x^2/2}\frac{d^n}{dx^n}e^{-x^2/2}=\sqrt{n!}\,h_n(x).
\]
Moreover, \eqref{eq:3-3-31:b} gives $H_n(cx,c^2y)=c^nH_n(x,y)$ for $c>0$; with $c=\sqrt y$, $y>0$, this yields $H_n(x,y)=y^{n/2}H_n(x/\sqrt y,1)=\sqrt{n!}\,y^{n/2}h_n(x/\sqrt y)$. The formula thus holds for $y>0$, the only case in which its right-hand side is defined.}
\end{kex}

\begin{kex}{Exercise}{3.32}{}
We shall show that \bk{(3.42)} holds, a.s.\ for all $0\le t<\infty$ simultaneously, with
\begin{equation}\label{eq:3-3-32:b}
  b(x)=\sigma(x)\Bigl[\rho+\tfrac12\,\sigma'(x)-c\int_0^x\frac{dy}{\sigma(y)}\Bigr];\qquad x\in\R ,
\end{equation}
which is continuous because $\sigma$ is of class $C^1$ and $1/\sigma$ is continuous.

\emph{Setting.} As throughout Section 3.3, $W=\{W_t,\F_t;\ 0\le t<\infty\}$ is a standard, one-dimensional Brownian motion on a probability space $(\Omega,\F,P)$ whose filtration $\{\F_t\}$ satisfies the usual conditions, and $\xi_0$ is an $\F_0$-measurable random variable (it has to be, for $\xi$ to be adapted). We write
\[
  h(x)\defeq\int_0^x\frac{dy}{\sigma(y)},\qquad
  M_t\defeq\int_0^t e^{cs}\,dW_s,\qquad
  B_t\defeq\rho\int_0^t e^{cs}\,ds ,
\]
so that $f(t,x)=e^{ct}h(x)$ and $\xi_t=\xi_0+M_t+B_t$. The deterministic, continuous integrand $e^{cs}$ belongs to $\cL^*(W)$, so $M\in\cM_2^c$ and, as recalled at the beginning of the proof of \bk{Corollary 2.20}, $\qv{M}_t=\int_0^t e^{2cs}\,ds$. The process $B$ is deterministic, continuous and of bounded variation on bounded intervals, with $B_0=0$. Thus $\xi$ is a continuous semimartingale in the sense of \bk{Definition 3.1}.

\emph{Step 1: the function $g$.} Since $1/\sigma$ is continuous and strictly positive, $h$ is of class $C^1$ with $h'=1/\sigma>0$; since $\sigma$ is of class $C^1$, $h$ is in fact of class $C^2$. The hypothesis that $1/\sigma$ is not integrable at $\pm\infty$ means $\int_0^{\infty}dy/\sigma(y)=\int_{-\infty}^0dy/\sigma(y)=\infty$, i.e.\ $h(x)\to\pm\infty$ as $x\to\pm\infty$. Hence $h$ is a strictly increasing bijection of $\R$ onto $\R$. Let $k\defeq h^{-1}\colon\R\to\R$. By the inverse function theorem $k$ is differentiable with
\begin{equation}\label{eq:3-3-32:k}
  k'(y)=\frac{1}{h'(k(y))}=\sigma(k(y)),\qquad\text{and therefore}\qquad
  k''(y)=\sigma'(k(y))\,k'(y)=\sigma'(k(y))\,\sigma(k(y)) ;
\end{equation}
the first formula shows that $k'$ is of class $C^1$, so $k$ is of class $C^2$. For fixed $t$, the equation $f(t,x)=y$ reads $h(x)=e^{-ct}y$, so that $f(t,\cdot)$ is a bijection of $\R$ onto $\R$ with inverse
\[
  g(t,y)=k\bigl(e^{-ct}y\bigr);\qquad 0\le t<\infty,\ y\in\R .
\]
This formula defines a function of class $C^2$ on $\R\times\R$; in particular $g$ is of class $C^{1,2}$ on $[0,\infty)\times\R$, and by \eqref{eq:3-3-32:k}, with the abbreviation $x=g(t,y)$ (so that $e^{-ct}y=h(x)$),
\begin{equation}\label{eq:3-3-32:d}
  \frac{\partial g}{\partial t}(t,y)=-c\,e^{-ct}y\,k'(e^{-ct}y)=-c\,h(x)\sigma(x),\quad
  \frac{\partial g}{\partial y}(t,y)=e^{-ct}\sigma(x),\quad
  \frac{\partial^2 g}{\partial y^2}(t,y)=e^{-2ct}\sigma'(x)\sigma(x).
\end{equation}

\emph{Step 2: It\^o's rule.} The process $X_t=g(t,\xi_t)$ is continuous and adapted, and $X_0=g(0,\xi_0)=k(\xi_0)$. We apply \bk{Theorem 3.6} with $d=1$ to the semimartingale $\xi=\xi_0+M+B$ and to the function $g$: a.s., for all $0\le t<\infty$,
\[
  X_t=X_0+\int_0^t\frac{\partial g}{\partial t}(s,\xi_s)\,ds
  +\int_0^t\frac{\partial g}{\partial y}(s,\xi_s)\,dB_s
  +\int_0^t\frac{\partial g}{\partial y}(s,\xi_s)\,dM_s
  +\frac12\int_0^t\frac{\partial^2 g}{\partial y^2}(s,\xi_s)\,d\qv{M}_s .
\]
The first, second and fourth integrals are, for each $\omega$, Lebesgue--Stieltjes integrals of continuous functions of $s$; since $dB_s=\rho e^{cs}\,ds$ and $d\qv{M}_s=e^{2cs}\,ds$, the formulas \eqref{eq:3-3-32:d}, evaluated at $y=\xi_s$, $x=X_s$, give
\begin{align*}
  \int_0^t\frac{\partial g}{\partial t}(s,\xi_s)\,ds&=-c\int_0^t h(X_s)\sigma(X_s)\,ds,\qquad
  \int_0^t\frac{\partial g}{\partial y}(s,\xi_s)\,dB_s=\rho\int_0^t\sigma(X_s)\,ds,\\
  \frac12\int_0^t\frac{\partial^2 g}{\partial y^2}(s,\xi_s)\,d\qv{M}_s&=\frac12\int_0^t\sigma'(X_s)\sigma(X_s)\,ds .
\end{align*}
Their sum is $\int_0^t b(X_s)\,ds$, with $b$ as in \eqref{eq:3-3-32:b}. It remains to identify the stochastic integral: with $Y_s\defeq\frac{\partial g}{\partial y}(s,\xi_s)=e^{-cs}\sigma(X_s)$, we must show that
\begin{equation}\label{eq:3-3-32:assoc}
  \int_0^t Y_s\,dM_s=\int_0^t Y_s e^{cs}\,dW_s=\int_0^t\sigma(X_s)\,dW_s;\qquad 0\le t<\infty,\ \text{a.s.}
\end{equation}
Once this is done, \bk{(3.42)} follows.

\emph{Step 3: proof of \eqref{eq:3-3-32:assoc}.} This is the extension of \bk{Corollary 2.20} to integrands which are only locally square-integrable (the extension invoked in \bk{Example 3.9}); we derive it from the definition of the integral in \bk{Definition 2.23}. The process $Y$ is continuous and adapted, hence progressively measurable, and
$\int_0^t Y_s^2\,d\qv{M}_s=\int_0^t Y_s^2e^{2cs}\,ds=\int_0^t\sigma^2(X_s)\,ds<\infty$ for every $t$ and every $\omega$, by continuity. Thus $Y\in\cP^*(M)$ and $\{Y_se^{cs}\}\in\cP^*(W)$, and the stopping times
\[
  T_n\defeq n\wedge\inf\Bigl\{0\le t<\infty;\ \int_0^t\sigma^2(X_s)\,ds\ge n\Bigr\};\qquad n\ge1,
\]
which increase to infinity, are the stopping times $R_n$ of the construction preceding \bk{(2.32)}, simultaneously for the pair $(Y,M)$ and for the pair $(Ye^{c\cdot},W)$. Since $M$ and $W$ belong to $\cM_2^c$, we may take $S_n\equiv\infty$ in that construction (which does not depend on the choice of the localizing sequence $\{S_n\}$), so that \bk{(2.32)} gives the above $T_n$ in both cases. Set $Y^{(n)}_s\defeq Y_s\ind_{\{s\le T_n\}}$; this process is progressively measurable and
\[
  E\int_0^\infty\bigl(Y^{(n)}_s\bigr)^2\,d\qv{M}_s=E\int_0^{T_n}\sigma^2(X_s)\,ds\le n ,
\]
so $Y^{(n)}\in\cL^*(M)$ and $\{Y^{(n)}_se^{cs}\}\in\cL^*(W)$. By \bk{(2.34)}, on $\{t\le T_n\}$ the integral $\int_0^tY_s\,dM_s$ is the integral of $Y^{(n)}$ with respect to the stopped martingale $M_{\cdot\wedge T_n}$, and by \bk{Corollary 2.21} (applied to $M$, $M_{\cdot\wedge T_n}\in\cM_2^c$, the common integrand $Y^{(n)}$ and the stopping time $T_n$) the latter agrees on $\{t\le T_n\}$ with $I^M_t(Y^{(n)})$. In the same way, $\int_0^tY_se^{cs}\,dW_s=I^W_t\bigl(Y^{(n)}e^{c\cdot}\bigr)$ on $\{t\le T_n\}$. Now \bk{Corollary 2.20}, applied to $W\in\cM_2^c$, the integrand $e^{c\cdot}\in\cL^*(W)$, $M=I^W(e^{c\cdot})$ and $Y^{(n)}\in\cL^*(M)$, yields $I^M(Y^{(n)})=I^W\bigl(Y^{(n)}e^{c\cdot}\bigr)$. Consequently, a.s.,
\[
  \int_0^t Y_s\,dM_s=\int_0^t Y_se^{cs}\,dW_s\qquad\text{for all } 0\le t\le T_n ,
\]
and letting $n\to\infty$ we obtain the first equality in \eqref{eq:3-3-32:assoc}; the second is the identity $Y_se^{cs}=\sigma(X_s)$. This completes the proof of \bk{(3.42)} with the drift \eqref{eq:3-3-32:b}.
\rmk{Two special cases confirm \eqref{eq:3-3-32:b}. If $\sigma\equiv1$ then $h(x)=x$, $X_t=e^{-ct}\xi_t$ and $b(x)=\rho-cx$: $X$ is the Ornstein--Uhlenbeck process with drift, $dX_t=(\rho-cX_t)\,dt+dW_t$, as one checks directly by integration by parts. If $c=\rho=0$ then $X_t=k(\xi_0+W_t)$ and $b=\frac12\sigma\sigma'$, the familiar correction term in $dX_t=\sigma(X_t)\,dW_t+\frac12\sigma(X_t)\sigma'(X_t)\,dt$ for a function of Brownian motion solving $k'=\sigma(k)$. In general, $h(X_t)=e^{-ct}\xi_t$ is an Ornstein--Uhlenbeck process with drift, and $X$ is its image under $k=h^{-1}$.}
\end{kex}

\begin{kex}{Exercise}{3.33}{}
\emph{Setting.} As throughout the book's Section 3.3, $W=\{W_t,\F_t;\ 0\le t<\infty\}$ is a Brownian motion on $(\Omega,\F,P)$ with respect to a filtration $\{\F_t\}$ satisfying the usual conditions (if only $W$ is given, we take for $\{\F_t\}$ the augmentation of $\{\F^W_t\}$ under $P$: each of its $\sigma$-fields contains all $P$-null sets by \bk{Definition 2.7.2}, it is right-continuous by \bk{Proposition 2.7.7}, which applies because $W$ is a strong Markov process by \bk{Theorem 2.6.15}, and $W$ is still a Brownian motion with respect to it by \bk{Theorem 2.7.9}); in particular $W\in\cM^{c,\loc}$ and $\qv{W}_t=t$. Discarding a $P$-null event, which belongs to $\F_0$, we may assume that every path of $W$ is continuous. We put
\[
  c\defeq\delta\mu-\tfrac12\delta^2,
  \qquad
  Z_t\defeq\exp\bigl[\delta W^{(\mu)}_t-\tfrac12\delta^2t\bigr]=\exp[\delta W_t+ct],
  \qquad
  A_t\defeq\int_0^t\frac{ds}{Z_s};\qquad 0\le t<\infty .
\]

\emph{Step 1: the processes $A$ and $X$.} For every $\omega$ the function $s\mapsto 1/Z_s(\omega)=\exp[-\delta W_s(\omega)-cs]$ is continuous and strictly positive on $[0,\infty)$. Hence $A_t(\omega)$ is well defined as a Riemann integral, $A_0=0$, and $t\mapsto A_t(\omega)$ is nondecreasing and continuously differentiable, with derivative $1/Z_t(\omega)$. Moreover, $A_t$ is the pointwise limit, as $n\to\infty$, of the Riemann sums $\sum_{k=1}^{n}\frac tn\,Z_{kt/n}^{-1}$, each of which is $\F_t$-measurable; thus $A$ is adapted to $\{\F_t\}$. So $A$ is a continuous, adapted process of bounded variation with $A_0=0$.

For fixed $t$ and $\omega$, the integrand in the definition of $X_t$ equals $Z_t(\omega)/Z_s(\omega)$, a continuous function of $s\in[0,t]$, and the factor $Z_t(\omega)$ does not depend on $s$. Therefore $X_t$ is well defined and
\begin{equation}\label{eq:3-3-33:a}
  X_t=Z_tA_t;\qquad 0\le t<\infty .
\end{equation}
In particular $X$ is a continuous, adapted (hence progressively measurable, by \bk{Proposition 1.1.13}) process, and $\int_0^tX_s^2\,ds<\infty$ for every $t$ and every $\omega$, because continuous paths are bounded on $[0,t]$. Thus $X\in\cP$, and the stochastic integral $\int_0^tX_s\,dW_s$ in the equation to be proved is well defined as an element of $\cM^{c,\loc}$; the Lebesgue integral $\int_0^t(1+\delta\mu X_s)\,ds$ is well defined for the same reason.

\emph{Step 2: It\^o's rule.} We apply \bk{Theorem 3.6} with $d=2$ to the vector semimartingale
\[
  (W_t,A_t)=(0,0)+(M^{(1)}_t,M^{(2)}_t)+(B^{(1)}_t,B^{(2)}_t),
  \qquad M^{(1)}=W,\quad M^{(2)}\equiv0,\quad B^{(1)}\equiv0,\quad B^{(2)}=A,
\]
whose components satisfy the hypotheses of that theorem by Step 1, and to the function
\[
  f(t,x,a)\defeq a\,e^{\delta x+ct};\qquad (t,x,a)\in[0,\infty)\times\R^2 ,
\]
which is of class $C^\infty$, hence of class $C^{1,2}$. Its relevant derivatives are
\[
  \frac{\partial f}{\partial t}=c\,f,\qquad
  \frac{\partial f}{\partial x}=\delta f,\qquad
  \frac{\partial f}{\partial a}=e^{\delta x+ct},\qquad
  \frac{\partial^2 f}{\partial x^2}=\delta^2 f .
\]
By \eqref{eq:3-3-33:a} we have $f(t,W_t,A_t)=Z_tA_t=X_t$ and $f(0,W_0,A_0)=0$. Since $M^{(2)}\equiv0$, the only cross-variation which does not vanish identically is $\qv{M^{(1)},M^{(1)}}_s=\qv{W}_s=s$, and the only nonzero terms in \bk{(3.5)} are those with $\partial f/\partial t$, with $(\partial f/\partial a)\,dB^{(2)}$, with $(\partial f/\partial x)\,dM^{(1)}$ and with $\partial^2f/\partial x^2$. Consequently \bk{(3.5)} reads: a.s.,
\begin{equation}\label{eq:3-3-33:b}
  X_t=c\int_0^tX_s\,ds+\int_0^tZ_s\,dA_s+\delta\int_0^tX_s\,dW_s+\frac{\delta^2}{2}\int_0^tX_s\,ds;
  \qquad 0\le t<\infty .
\end{equation}

\emph{Step 3: conclusion.} For every $\omega$ the function $s\mapsto A_s(\omega)$ is continuously differentiable with derivative $1/Z_s(\omega)$, so the Lebesgue--Stieltjes measure $dA_s(\omega)$ on $[0,t]$ has density $1/Z_s(\omega)$ with respect to Lebesgue measure, and
\[
  \int_0^tZ_s\,dA_s=\int_0^tZ_s\cdot\frac1{Z_s}\,ds=t .
\]
Furthermore $c+\tfrac12\delta^2=\delta\mu$. Substituting these two facts into \eqref{eq:3-3-33:b}, we obtain, a.s.,
\[
  X_t=t+\delta\mu\int_0^tX_s\,ds+\delta\int_0^tX_s\,dW_s
     =\int_0^t(1+\delta\mu X_s)\,ds+\delta\int_0^tX_s\,dW_s;\qquad 0\le t<\infty ,
\]
which is the Shiryaev--Roberts equation; the exceptional null event does not depend on $t$, i.e., the two sides are indistinguishable.

\rmk{The same computation can also be organised via the integration-by-parts formula: by \bk{Theorem 3.6} the process $Z$ satisfies $Z_t=1+\delta\mu\int_0^tZ_s\,ds+\delta\int_0^tZ_s\,dW_s$ (compare \bk{Example 3.9} when $\mu=0$), and formula \bk{(3.8)} of \bk{Problem 3.12}, applied to $Z$ and to the process $A$ of bounded variation, for which the cross-variation term vanishes and $A_0=0$, gives $Z_tA_t=\int_0^tZ_s\,dA_s+\int_0^tA_s\,dZ_s$. We preferred the two-dimensional It\^o rule because it does not require rewriting $\int_0^tA_s\,dZ_s$ as an integral with respect to $W$. As a check, taking expectations formally in the equation gives $m'(t)=1+\delta\mu\,m(t)$, $m(0)=0$ for $m(t)=E X_t$, i.e.\ $m(t)=(e^{\delta\mu t}-1)/(\delta\mu)$ (and $m(t)=t$ if $\delta\mu=0$); this agrees with the direct computation from the definition of $X$, since $E\exp[\delta(W_t-W_s)-\tfrac12\delta^2(t-s)]=1$ gives $E X_t=\int_0^te^{\delta\mu(t-s)}\,ds$ by Tonelli's theorem.}
\end{kex}

\begin{kex}{Exercise}{3.34}{}
Let $W=\{W_t,\F_t;\ 0\le t<\infty\}$ be the given Brownian motion, on a filtration satisfying the
usual conditions, and write
\[
  I^{\beta}_t \defeq \int_0^t e^{\beta s}\,dW_s,\qquad
  Y^{\beta}_t \defeq e^{-\beta t}I^{\beta}_t ;\qquad 0\le t<\infty,\ \beta>0 .
\]
The integrand $s\mapsto e^{\beta s}$ is nonrandom and continuous, hence progressively measurable,
and $\int_0^t e^{2\beta s}\,d\qv{W}_s=\int_0^t e^{2\beta s}\,ds<\infty$ for every $t$; thus
$I^{\beta}$ is a well-defined continuous martingale for each fixed $\beta$.

\emph{Step 1: a pathwise representation.} Fix $\beta>0$. We apply the integration by parts formula
\bk{(3.8)} of \bk{Problem 3.12} to the continuous semimartingales
\[
  X_t\defeq e^{\beta t}=1+B_t,\quad B_t\defeq\int_0^t\beta e^{\beta s}\,ds ,
  \qquad\text{and}\qquad Y_t\defeq W_t ,
\]
that is, with $X_0=1$, $M\equiv0$, $Y_0=0$, $N=W$, $C\equiv0$ in \bk{(3.7)}; here $B$ is nonrandom
(hence adapted), continuous, increasing, with $B_0=0$. Since $\qv{M,N}=\qv{0,W}\equiv0$ and
$X_0Y_0=0$, formula \bk{(3.8)} gives, a.s.,
\[
  \int_0^t e^{\beta s}\,dW_s=e^{\beta t}W_t-\int_0^t W_s\,dB_s
  =e^{\beta t}W_t-\beta\int_0^t W_s e^{\beta s}\,ds ;\qquad 0\le t<\infty,
\]
where the last integral is a pathwise Lebesgue--Stieltjes integral with respect to the absolutely
continuous function $B$, i.e.\ a Lebesgue integral with respect to $\beta e^{\beta s}\,ds$.
Multiplying by $e^{-\beta t}$ we obtain: for every $\beta>0$ there is an event $\Omega_\beta$ with
$P(\Omega_\beta)=1$ on which
\begin{equation}\label{eq:3-3-34:rep}
  Y^{\beta}_t=\Phi_\beta(W)(t)\quad\text{for all }0\le t<\infty,\qquad\text{where}\quad
  \Phi_\beta(w)(t)\defeq w(t)-\beta\int_0^t e^{-\beta(t-s)}w(s)\,ds .
\end{equation}

\emph{Step 2: a deterministic estimate.} Let $w:[0,T]\to\R$ be continuous with $w(0)=0$, and put
\[
  K\defeq\max_{0\le s\le T}\abs{w(s)},\qquad
  m(\delta)\defeq\max\bigl\{\abs{w(t)-w(s)};\ s,t\in[0,T],\ \abs{t-s}\le\delta\bigr\},\quad\delta>0 .
\]
We claim that for all $\beta>0$ and $\delta>0$,
\begin{equation}\label{eq:3-3-34:det}
  \sup_{0\le t\le T}\abs{\Phi_\beta(w)(t)}\le 2m(\delta)+3Ke^{-\beta\delta}.
\end{equation}
Indeed, $\beta\int_0^te^{-\beta(t-s)}\,ds=1-e^{-\beta t}$, so that for $0\le t\le T$
\[
  \Phi_\beta(w)(t)=e^{-\beta t}w(t)+\beta\int_0^t e^{-\beta(t-s)}\bigl(w(t)-w(s)\bigr)\,ds .
\]
For the first term: if $t\le\delta$, then $\abs{w(t)}=\abs{w(t)-w(0)}\le m(\delta)$ and
$e^{-\beta t}\le1$; if $t>\delta$, then $e^{-\beta t}\abs{w(t)}\le Ke^{-\beta\delta}$. In both cases
$e^{-\beta t}\abs{w(t)}\le m(\delta)+Ke^{-\beta\delta}$. For the second term we split the domain of
integration into $\{s\in[0,t];\ t-s\le\delta\}$, on which $\abs{w(t)-w(s)}\le m(\delta)$, and
$\{s\in[0,t];\ t-s>\delta\}$, on which $\abs{w(t)-w(s)}\le2K$ (the latter set is empty if
$t\le\delta$):
\[
  \beta\int_0^t e^{-\beta(t-s)}\abs{w(t)-w(s)}\,ds
  \le m(\delta)\,\beta\int_0^{t}e^{-\beta(t-s)}\,ds
     +2K\,\beta\int_0^{(t-\delta)^+}e^{-\beta(t-s)}\,ds
  \le m(\delta)+2Ke^{-\beta\delta},
\]
because $\beta\int_0^{t-\delta}e^{-\beta(t-s)}\,ds=e^{-\beta\delta}-e^{-\beta t}\le e^{-\beta\delta}$
when $t>\delta$. Adding the two bounds gives \eqref{eq:3-3-34:det}.

Letting $\beta\to\infty$ in \eqref{eq:3-3-34:det} with $\delta$ fixed yields
$\limsup_{\beta\to\infty}\sup_{0\le t\le T}\abs{\Phi_\beta(w)(t)}\le2m(\delta)$ for every
$\delta>0$; and $m(\delta)\downarrow0$ as $\delta\downarrow0$ because $w$ is uniformly continuous on
the compact interval $[0,T]$. Consequently
\begin{equation}\label{eq:3-3-34:lim}
  \lim_{\beta\to\infty}\ \sup_{0\le t\le T}\abs{\Phi_\beta(w)(t)}=0
  \qquad\text{for every continuous } w:[0,T]\to\R \text{ with } w(0)=0,
\end{equation}
the limit being taken along all real $\beta$.

\emph{Step 3: conclusion.} Let $\Omega^*$ be the event, of probability one, on which
$t\mapsto W_t(\omega)$ is continuous and $W_0(\omega)=0$. By \eqref{eq:3-3-34:lim} applied to
$w=W_\cdot(\omega)\big|_{[0,T]}$,
\begin{equation}\label{eq:3-3-34:pw}
  \lim_{\beta\to\infty}\ \sup_{0\le t\le T}\abs{\Phi_\beta(W(\omega))(t)}=0
  \qquad\text{for every }\omega\in\Omega^* .
\end{equation}
For each $\beta$ the stochastic integral $I^\beta$ is defined only up to indistinguishability, so
the statement about the whole family $\{Y^\beta\}_{\beta>0}$ has to be read as follows.

\textbf{(a)} For every $\beta>0$, the process $\{e^{\beta t}\Phi_\beta(W)(t);\ 0\le t<\infty\}$ is,
by \eqref{eq:3-3-34:rep}, a version of $I^\beta$ (indistinguishable from any other), and it is
defined simultaneously for all $\beta>0$ and all $\omega\in\Omega^*$. With this choice of the
integrals, $Y^\beta_t(\omega)=\Phi_\beta(W(\omega))(t)$ for all $\beta$, $t$ and
$\omega\in\Omega^*$, and \eqref{eq:3-3-34:pw} is exactly the assertion
\[
  \lim_{\beta\to\infty}\ \sup_{0\le t\le T}
  \biggl|e^{-\beta t}\int_0^te^{\beta s}\,dW_s\biggr|=0,\quad\text{a.s.}
\]

\textbf{(b)} For arbitrary versions of the integrals $I^\beta$, and any sequence
$\beta_n\to\infty$, the event $\Omega^*\cap\bigcap_{n}\Omega_{\beta_n}$ has probability one, and on
it $\sup_{0\le t\le T}\abs{Y^{\beta_n}_t}=\sup_{0\le t\le T}\abs{\Phi_{\beta_n}(W)(t)}\to0$ as
$n\to\infty$, by \eqref{eq:3-3-34:rep} and \eqref{eq:3-3-34:pw}. \qed

\rmk{Some choice of versions as in \textbf{(a)} is unavoidable if the limit is taken along all real
$\beta$: altering $I^{\beta}$ on a $P$-null set $N_\beta$ for each $\beta$ yields integrals that
are still legitimate individually, but when the sets $N_\beta$ have a non-null union (as can happen
on a probability space carrying uncountably many disjoint null sets, e.g.\ one supporting an
independent uniform random variable $U$, with $N_\beta=\{U=\beta-\lfloor\beta\rfloor\}$ and
$I^\beta\equiv e^{\beta t}$ on $N_\beta$) the limit along the continuum fails on that union. The
representation \eqref{eq:3-3-34:rep} singles out the natural version, jointly continuous in
$(\beta,t)$.}

\rmk{For fixed $\beta$, $Y^\beta$ is the Ornstein--Uhlenbeck process solving
$dY_t=-\beta Y_t\,dt+dW_t$, $Y_0=0$ (cf.\ \bk{Example 3.11}), with
$E(Y^\beta_t)^2=e^{-2\beta t}\int_0^te^{2\beta s}\,ds=(1-e^{-2\beta t})/(2\beta)\le1/(2\beta)$; this
gives convergence to zero in $L^2$ for each fixed $t$, in agreement with the uniform almost sure
statement proved above.}
\end{kex}

\begin{kex}{Exercise}{3.35}{}
Here $W=\{W_t,\F_t;\ 0\le t<\infty\}$ is a standard, one-dimensional Brownian motion, $T$ is a stopping time of $\{\F_t\}$, and, in accordance with the standing assumption made at the beginning of \bk{Section 3}, the filtration $\{\F_t\}$ satisfies the usual conditions. Under the hypothesis $E\sqrt{T}<\infty$ the expectation $ET$ may be infinite (see the remark at the end), so the second identity is an equality in $[0,\infty]$; we prove that $W_T$ is integrable with $E(W_T)=0$, and that $E(W_T^2)=ET$, both sides being finite or both infinite.

\emph{Step 1: preliminaries.} Since $E\sqrt{T}<\infty$, the event $\{T=\infty\}$ is null; on it we set $W_T\defeq0$. The continuous, adapted processes $W$ and $\{W_t^2-t\}$ are martingales: both are integrable, and for $0\le s<t<\infty$ the increment $W_t-W_s$ is independent of $\F_s$ and normally distributed with mean zero and variance $t-s$ (\bk{Definition 2.1.1}), so that, a.s.,
\begin{align*}
  E[W_t\mid\F_s]&=W_s+E(W_t-W_s)=W_s,\\
  E[W_t^2-t\mid\F_s]&=W_s^2+2W_s\,E(W_t-W_s)+E\bigl[(W_t-W_s)^2\bigr]-t=W_s^2-s .
\end{align*}
Therefore \bk{Problem 1.3.24}(i), applied to $\pm W$ and to $\pm(W_t^2-t)$, shows that
\[
  M_t\defeq W_{t\wedge T}\quad\text{and}\quad M_t^2-(t\wedge T);\qquad 0\le t<\infty,
\]
are martingales with respect to $\{\F_t\}$. Both vanish a.s.\ at $t=0$, and therefore
\begin{equation}\label{eq:3-3-35:a}
  E(M_t)=0,\qquad E(M_t^2)=E(t\wedge T);\qquad 0\le t<\infty .
\end{equation}
For every $\omega\in\{T<\infty\}$ we have $M_t(\omega)=W_{T(\omega)}(\omega)$ as soon as $t\ge T(\omega)$; hence
\begin{equation}\label{eq:3-3-35:b}
  \lim_{t\to\infty}M_t=W_T\quad\text{a.s.,}
\end{equation}
and $W_T$ is $\F$-measurable; indeed, $W_T=\lim_{n\to\infty}M_n\ind_{\{T<\infty\}}$ holds everywhere, each $M_n$ is $\F_n$-measurable, and $\{T<\infty\}\in\F$.

\emph{Step 2: the maximal function is integrable.} The process $W$ belongs to $\cM^{c,\loc}$, and $\qv{W}_t=t$: indeed, $t\mapsto t$ is continuous, adapted, nondecreasing, null at $t=0$, and $W_t^2-t$ is a martingale, so the claim follows from the uniqueness assertion of \bk{Problem 1.5.17}. We apply the right-hand inequality of \bk{(3.37)} in \bk{Theorem 3.28} to $M=W$, with $m=\tfrac12$ and the stopping time $T$. With the notation \bk{(3.26)}, i.e.\ $W_T^*=\max_{0\le s\le T}\abs{W_s}$ on $\{T<\infty\}$, we obtain
\begin{equation}\label{eq:3-3-35:c}
  E(W_T^*)\le K_{1/2}\,E\bigl(\qv{W}_T^{1/2}\bigr)=K_{1/2}\,E\sqrt{T}<\infty .
\end{equation}
Moreover, since $t\wedge T\le T$,
\begin{equation}\label{eq:3-3-35:d}
  \abs{M_t}=\abs{W_{t\wedge T}}\le W_T^*\quad\text{and}\quad\abs{W_T}\le W_T^*\qquad\text{a.s., for every }0\le t<\infty .
\end{equation}

\emph{Step 3: the first identity.} By \eqref{eq:3-3-35:c} and \eqref{eq:3-3-35:d}, $W_T$ is integrable and the family $\{M_t\}_{t\ge0}$ is dominated by the integrable random variable $W_T^*$. From \eqref{eq:3-3-35:b}, the dominated convergence theorem and \eqref{eq:3-3-35:a},
\[
  E(W_T)=\lim_{n\to\infty}E(M_n)=0 .
\]

\emph{Step 4: $E(W_T^2)\le ET$.} By \eqref{eq:3-3-35:b}, Fatou's lemma, \eqref{eq:3-3-35:a} and the monotone convergence theorem (as $n\wedge T\uparrow T$),
\[
  E(W_T^2)=E\Bigl(\lim_{n\to\infty}M_n^2\Bigr)\le\liminf_{n\to\infty}E(M_n^2)
  =\lim_{n\to\infty}E(n\wedge T)=ET .
\]

\emph{Step 5: $ET\le E(W_T^2)$.} If $E(W_T^2)=\infty$ there is nothing to prove, so let $E(W_T^2)<\infty$. Fix $0\le t<\infty$ and $A\in\F_t$. For every integer $n\ge t$ the martingale property of $M$ gives $E(M_t\ind_A)=E(M_n\ind_A)$; letting $n\to\infty$ and using \eqref{eq:3-3-35:b}, \eqref{eq:3-3-35:c}, \eqref{eq:3-3-35:d} and dominated convergence, we get $E(M_t\ind_A)=E(W_T\ind_A)$. As $M_t$ is $\F_t$-measurable and $A\in\F_t$ is arbitrary,
\[
  M_t=E[W_T\mid\F_t]\quad\text{a.s.}
\]
The conditional Jensen inequality (for the convex function $x\mapsto x^2$ and the square-integrable random variable $W_T$) yields $M_t^2\le E[W_T^2\mid\F_t]$ a.s., and taking expectations, with \eqref{eq:3-3-35:a},
\[
  E(t\wedge T)=E(M_t^2)\le E(W_T^2);\qquad 0\le t<\infty .
\]
Letting $t\to\infty$ and using monotone convergence, we arrive at $ET\le E(W_T^2)$.

Steps 4 and 5 together give $E(W_T^2)=ET$ in $[0,\infty]$; in particular $E(W_T^2)<\infty$ if and only if $ET<\infty$, in which case we recover \bk{Problem 2.12}.

\rmk{(a) The condition $E\sqrt{T}<\infty$ does not imply $ET<\infty$. Let $B$ be a standard Brownian motion and $\xi\ge0$ a random variable independent of $B$ with $E\sqrt{\xi}<\infty=E\xi$, on a complete probability space with collection of null events $\cN$. Put $\G_t\defeq\sigma\bigl(\sigma(\xi)\cup\F_t^B\cup\cN\bigr)$ and $\F_t\defeq\bigcap_{\eps>0}\G_{t+\eps}$; this filtration satisfies the usual conditions. For $s<t$ and $0<\eps<t-s$ the increment $B_t-B_{s+\eps}$ is independent of $\G_{s+\eps}\supseteq\F_s$, and letting $\eps\downarrow0$ we see from the continuity of $B$ that $B_t-B_s$ is independent of $\F_s$; thus $B$ is a Brownian motion with respect to $\{\F_t\}$. Now $T=\xi$ is a stopping time of $\{\F_t\}$, because $\{\xi\le t\}\in\sigma(\xi)\subseteq\F_t$, and by the above $E(B_\xi)=0$, $E(B_\xi^2)=E\xi=\infty$. (b) For the first identity the exponent $\tfrac12$ cannot be lowered: for the passage time $T_b$ to a level $b\neq0$ we have $E(W_{T_b})=b\ne0$, although the density \bk{(2.6.3)} of $T_b$ is of order $t^{-3/2}$ as $t\to\infty$, so that $E(T_b^{\alpha})<\infty$ for every $0<\alpha<\tfrac12$ (and $E\sqrt{T_b}=\infty$, as it must be).}
\end{kex}

\begin{kex}{Exercise}{3.36}{}
We use the notation of \bk{Definition 3.19} and write $P$ for $\wh{P}^0=P^{(0,\dots,0)}$: under $P$, the process $W=\{W_t,\F_t;\ 0\le t<\infty\}$ is a $d$-dimensional Brownian motion starting at the origin, and $R_t=\norm{W_t}$. As in the proof of \bk{Proposition 3.21}, stochastic integrals with respect to the components of $W$ are taken relative to $\{\F_t\}$, which is assumed to satisfy the usual conditions under $P$ (the standing assumption of the book's Chapter~3). Every path of $R$ is continuous, and by \bk{Proposition 3.22} the event
\[
  \Omega_0\defeq\{R_t>0 \text{ for every } 0<t<\infty\}
\]
has probability one. On $\Omega_0$ the path $t\mapsto M_t=R_t^{2-d}$ is finite and continuous on $[1,\infty)$; on the null set $\Omega\setminus\Omega_0$ we set $M\equiv 0$, which does not affect adaptedness because $\F_1$ contains the $P$-null sets. The time set being $[1,\infty)$, the words ``martingale'' and ``local martingale'' refer to the filtration $\{\F_t;\ 1\le t<\infty\}$; in other words, \bk{Definition 1.5.15} is applied to the process $\{M_{1+u},\F_{1+u};\ 0\le u<\infty\}$. Thus (i) asks for a nondecreasing sequence $\{S_k\}_{k=1}^\infty$ of stopping times of $\{\F_t\}$ with values in $[1,\infty]$ such that $P[\lim_{k\to\infty}S_k=\infty]=1$ and $\{M_{t\wedge S_k},\F_t;\ 1\le t<\infty\}$ is a martingale for every $k\ge 1$.

We prove (ii) first, because the integrability of $M_1$ is needed in (i).

\textbf{(ii)} Fix $t\ge 1$. Under $P$ the random vector $Z\defeq W_t/\sqrt t$ has independent, standard normal components, so $\norm{Z}^2$ has the chi-square distribution with $d$ degrees of freedom, whose density is $y^{\frac d2-1}e^{-y/2}/\bigl(2^{d/2}\Gamma(\tfrac d2)\bigr)$ on $(0,\infty)$. Hence, for every real number $q$,
\begin{equation}\label{eq:3-3-36:a}
  E\norm{Z}^{-q}=\frac{1}{2^{d/2}\Gamma(\frac d2)}\int_0^\infty y^{\frac{d-q}{2}-1}e^{-y/2}\,dy
  =\begin{cases}
     \dfrac{\Gamma\bigl(\frac{d-q}{2}\bigr)}{2^{q/2}\,\Gamma\bigl(\frac d2\bigr)}; & q<d,\\[2.2ex]
     \infty; & q\ge d,
   \end{cases}
\end{equation}
where for $q<d$ we substituted $y=2u$ and used the definition of the gamma function, and for $q\ge d$ the integrand is at least $e^{-1/2}y^{-1}$ on $(0,1)$. Let now $0<p<d/(d-2)$ and put $q\defeq p(d-2)\in(0,d)$. Since $M_t^p=R_t^{-q}=t^{-q/2}\norm{Z}^{-q}$, a.s., \eqref{eq:3-3-36:a} gives
\begin{equation}\label{eq:3-3-36:b}
  E(M_t^p)=\frac{\Gamma\bigl(\frac{d-p(d-2)}{2}\bigr)}{(2t)^{p(d-2)/2}\,\Gamma\bigl(\frac d2\bigr)};
  \qquad 1\le t<\infty,\quad 0<p<\frac{d}{d-2}.
\end{equation}
The right-hand side is decreasing in $t$, and therefore
\[
  \sup_{1\le t<\infty}E(M_t^p)=E(M_1^p)=\frac{\Gamma\bigl(\frac{d-p(d-2)}{2}\bigr)}{2^{p(d-2)/2}\,\Gamma\bigl(\frac d2\bigr)}<\infty .
\]
For the uniform integrability, we note that $d/(d-2)>1$ and choose $p\in(1,d/(d-2))$. Since $M_t\ge 0$, we have $M_t\ind_{\{M_t>\lambda\}}\le\lambda^{1-p}M_t^p$ for every $\lambda>0$, whence
\[
  \sup_{1\le t<\infty}E\bigl[M_t\ind_{\{M_t>\lambda\}}\bigr]\le\lambda^{1-p}E(M_1^p)\longrightarrow 0\quad\text{as }\lambda\to\infty .
\]
Thus the family $\{M_t;\ 1\le t<\infty\}$ is uniformly integrable.

\rmk{The range of $p$ is sharp: by \eqref{eq:3-3-36:a}, $E(M_t^p)=\infty$ for every $t$ if $p\ge d/(d-2)$. In particular $M_t$ is not square-integrable when $d\ge 4$.}

\textbf{(i)} \emph{Step 1: the function.} Let $f(x)\defeq\norm{x}^{2-d}$ for $x\in\R^d\setminus\{0\}$. This function is of class $C^\infty$ on $\R^d\setminus\{0\}$, with
\[
  \frac{\partial f}{\partial x_i}(x)=-(d-2)\frac{x_i}{\norm{x}^{d}},\qquad
  \frac{\partial^2 f}{\partial x_i^2}(x)=-(d-2)\Bigl(\frac{1}{\norm{x}^{d}}-\frac{d\,x_i^2}{\norm{x}^{d+2}}\Bigr);\qquad 1\le i\le d,
\]
so that
\begin{equation}\label{eq:3-3-36:c}
  \sum_{i=1}^d\frac{\partial^2 f}{\partial x_i^2}(x)=-(d-2)\Bigl(\frac{d}{\norm{x}^d}-\frac{d\norm{x}^2}{\norm{x}^{d+2}}\Bigr)=0,
  \qquad
  \norm{\nabla f(x)}=\frac{d-2}{\norm{x}^{d-1}};\qquad x\ne 0 .
\end{equation}
Since $f$ is not defined at the origin, \bk{Theorem 3.6} cannot be applied to it directly, and we modify it near the origin. Fix an integer $k\ge 1$, put $a\defeq(2k)^{-2}$ and $\psi(y)\defeq y^{-(d-2)/2}$ for $y>0$, and let
\[
  \varphi_k(y)\defeq
  \begin{cases}
    \psi(y); & y\ge a,\\
    \psi(a)+\psi'(a)(y-a)+\tfrac12\psi''(a)(y-a)^2; & y<a .
  \end{cases}
\]
Then $\varphi_k\colon\R\to\R$ is of class $C^2$, and consequently $f_k(x)\defeq\varphi_k(\norm{x}^2)$ is of class $C^2$ on $\R^d$, being the composition of $\varphi_k$ with a polynomial. Moreover $f_k=f$ on the open set $G_k\defeq\{x\in\R^d;\ \norm{x}>1/(2k)\}$, so the first and second partial derivatives of $f_k$ coincide with those of $f$ on $G_k$. Finally, $\nabla f_k$ is bounded on $\R^d$: it is continuous on the compact ball $\{\norm{x}\le 1/(2k)\}$, and on $G_k$ we have $\norm{\nabla f_k(x)}\le(d-2)(2k)^{d-1}$ by \eqref{eq:3-3-36:c}.

\emph{Step 2: It\^o's rule.} Under $P$ the components $W^{(1)},\dots,W^{(d)}$ are independent standard Brownian motions relative to $\{\F_t\}$; they belong to $\cM^{c,\loc}$ (indeed to $\cM_2^c$) and $\qv{W^{(i)},W^{(j)}}_t=\delta_{ij}t$. \bk{Theorem 3.6}, applied with $X=W$, $B\equiv 0$ and the function $f_k$ (which does not depend on $t$), gives, a.s.~$P$,
\begin{equation}\label{eq:3-3-36:d}
  f_k(W_t)=f_k(0)+\sum_{i=1}^d\int_0^t\frac{\partial f_k}{\partial x_i}(W_s)\,dW^{(i)}_s
  +\frac12\int_0^t\sum_{i=1}^d\frac{\partial^2 f_k}{\partial x_i^2}(W_s)\,ds;\qquad 0\le t<\infty .
\end{equation}
The integrands $X^{(i)}_s\defeq\frac{\partial f_k}{\partial x_i}(W_s)$ are continuous, adapted and bounded; hence they are progressively measurable and satisfy $E\int_0^T(X^{(i)}_s)^2\,ds<\infty$ for every $T$, i.e.\ $X^{(i)}\in\cL^*(W^{(i)})$, and the stochastic integrals $I^{(i)}_t\defeq\int_0^tX^{(i)}_s\,dW^{(i)}_s$ belong to $\cM_2^c$.

\emph{Step 3: the stopping times.} Let
\[
  S_k\defeq\inf\{t\ge 1;\ R_t\le 1/k\}\qquad(\inf\emptyset=\infty).
\]
We have $\{S_k\le t\}=\emptyset$ for $t<1$, and for $t\ge 1$, by the continuity of the paths of $R$,
\[
  \{S_k\le t\}=\Bigl\{\min_{1\le s\le t}R_s\le\frac1k\Bigr\}
  =\bigcap_{j=1}^\infty\ \bigcup_{s\in(\Q\cap[1,t])\cup\{t\}}\Bigl\{R_s<\frac1k+\frac1j\Bigr\}\in\F_t .
\]
Thus $S_k$ is a stopping time of $\{\F_t\}$ with $S_k\ge 1$, and clearly $S_k\le S_{k+1}$. On $\Omega_0$, for every $T\ge 1$ the continuous, strictly positive function $R$ has a strictly positive minimum $m_T$ on $[1,T]$, so that $S_k>T$ as soon as $1/k<m_T$. Therefore $\lim_{k\to\infty}S_k=\infty$ on $\Omega_0$, i.e.\ a.s.~$P$.

\emph{Step 4: the stopped process.} Fix $k$ and define
\[
  Y^{(i)}_s\defeq\ind_{\{1<s\le S_k\}}\,X^{(i)}_s,\qquad
  N_t\defeq\sum_{i=1}^d\int_0^tY^{(i)}_s\,dW^{(i)}_s;\qquad 0\le t<\infty .
\]
The process $s\mapsto\ind_{\{1<s\le S_k\}}$ is adapted and left-continuous, hence progressively measurable; so $Y^{(i)}$ is bounded and progressively measurable, and $N\in\cM_2^c$. Moreover $N_1=0$, a.s.: by \bk{(2.24)} with the constant stopping time $1$, we have $\int_0^1Y^{(i)}_s\,dW^{(i)}_s=\int_0^1\ind_{\{s\le 1\}}Y^{(i)}_s\,dW^{(i)}_s$, and the integrand on the right is identically zero. Applying \bk{(2.24)} with the stopping times $S_k$ and $1$, and using linearity and $S_k\ge 1$, we obtain for every $t\ge 1$
\begin{equation}\label{eq:3-3-36:e}
  I^{(i)}_{t\wedge S_k}-I^{(i)}_{1}
  =\int_0^t\bigl(\ind_{\{s\le S_k\}}-\ind_{\{s\le 1\}}\bigr)X^{(i)}_s\,dW^{(i)}_s
  =\int_0^tY^{(i)}_s\,dW^{(i)}_s,\quad\text{a.s.}
\end{equation}
We claim that
\begin{equation}\label{eq:3-3-36:f}
  M_{t\wedge S_k}=M_1+N_t;\qquad 1\le t<\infty,\quad\text{a.s.~}P .
\end{equation}
By continuity of both sides in $t$ it suffices to prove this for each fixed $t\ge 1$. Let $\Omega_1\subseteq\Omega_0$ be an event of probability one on which \eqref{eq:3-3-36:d} holds for all times and \eqref{eq:3-3-36:e} holds for the given $t$ and all $i$, and fix $\omega\in\Omega_1$. By the definition of $S_k$, we have $R_s(\omega)>1/k$ for $1\le s<S_k(\omega)$, and hence, by continuity, $R_s(\omega)\ge 1/k$ and $W_s(\omega)\in G_k$ for every $s$ in the interval $J(\omega)\defeq\{s;\ 1\le s\le S_k(\omega),\ s<\infty\}$, provided $S_k(\omega)>1$.

\emph{Case $S_k(\omega)>1$.} The times $1$ and $t\wedge S_k(\omega)$ belong to $J(\omega)$, so $f_k(W_1)=f(W_1)=M_1$ and $f_k(W_{t\wedge S_k})=M_{t\wedge S_k}$ at $\omega$. Evaluating \eqref{eq:3-3-36:d} at the times $t\wedge S_k(\omega)$ and $1$ and subtracting, we find
\[
  M_{t\wedge S_k}-M_1=\sum_{i=1}^d\bigl(I^{(i)}_{t\wedge S_k}-I^{(i)}_1\bigr)
  +\frac12\int_1^{t\wedge S_k}\sum_{i=1}^d\frac{\partial^2f_k}{\partial x_i^2}(W_s)\,ds .
\]
The last integral vanishes, because for $s\in J(\omega)$ the point $W_s(\omega)$ lies in $G_k$, where the second derivatives of $f_k$ are those of $f$ and \eqref{eq:3-3-36:c} applies. By \eqref{eq:3-3-36:e} the sum equals $N_t$, which proves \eqref{eq:3-3-36:f} at $\omega$.

\emph{Case $S_k(\omega)=1$.} Here $M_{t\wedge S_k}=M_1$, and \eqref{eq:3-3-36:e} gives $\int_0^tY^{(i)}_s\,dW^{(i)}_s=I^{(i)}_1-I^{(i)}_1=0$ at $\omega$, so $N_t(\omega)=0$ and \eqref{eq:3-3-36:f} holds again.

\emph{Step 5: conclusion.} By \eqref{eq:3-3-36:b} with $p=1$, the random variable $M_1$ is integrable; it is $\F_1$-measurable. Since $N$ is a martingale, \eqref{eq:3-3-36:f} shows that $M_{t\wedge S_k}$ is integrable and $\F_t$-measurable for every $t\ge 1$, and that for $1\le s\le t$
\[
  E\bigl[M_{t\wedge S_k}\,\big|\,\F_s\bigr]=M_1+E[N_t\,|\,\F_s]=M_1+N_s=M_{s\wedge S_k},\quad\text{a.s.~}P .
\]
Hence $\{M_{t\wedge S_k},\F_t;\ 1\le t<\infty\}$ is a martingale for every $k\ge 1$. Together with Step~3, this shows that $\{M_t,\F_t;\ 1\le t<\infty\}$ is a continuous local martingale.

\textbf{(iii)} Taking $p=1$ in \eqref{eq:3-3-36:b}, we obtain
\[
  E(M_t)=\frac{1}{(2t)^{(d-2)/2}\,\Gamma\bigl(\frac d2\bigr)};\qquad 1\le t<\infty,
\]
which is finite, strictly positive and, because $d-2>0$, strictly decreasing in $t$. A martingale has constant expectation; therefore $\{M_t,\F_t;\ 1\le t<\infty\}$ is not a martingale, although by (i) and (ii) it is a uniformly integrable local martingale. The same argument shows that $M$ is not a martingale on any interval $[1,T]$ with $T>1$, nor with respect to any other filtration.

\rmk{This is consistent with \bk{Problem 1.5.19}: being a nonnegative local martingale, $M$ is a supermartingale, so $E(M_t)$ must be nonincreasing; and since $M$ is not a martingale, it cannot be of class $DL$, even though it is uniformly integrable. For $d=3$ the formula reads $E(1/R_t)=\sqrt{2/(\pi t)}$.}
\end{kex}

\begin{kex}{Exercise}{3.37}{}
We use the notation of \bk{Definition 3.19} and write $P$ for $\hat P^0$: thus $R_t=\norm{W_t}$, where $W=\{W_t=(W^{(1)}_t,W^{(2)}_t),\F_t;\ 0\le t<\infty\}$ is a two-dimensional Brownian motion with $P[W_0=0]=1$, and $\{\F_t\}$ satisfies the usual conditions, as is assumed throughout \bk{Section 3.3}. By \bk{Proposition 3.22} the event $\Omega_0\defeq\{R_t>0;\ \forall\,0<t<\infty\}$ has probability one; it belongs to $\F_0$, because $\F_0$ contains the $P$-null sets. We set $X_t\defeq-\log R_t$ on $\Omega_0$ and $X_t\defeq0$ off $\Omega_0$, for $1\le t<\infty$. Then $X=\{X_t,\F_t;\ 1\le t<\infty\}$ is a continuous, adapted process, and ``local martingale'' is understood as in \bk{Definition 1.5.15}, with the time-parameter set $[1,\infty)$ in place of $[0,\infty)$.

\emph{Step 1. The law of $R_t$, and the moments of $e^{X_t}$.} Fix $t>0$. Under $P$ the coordinates $W^{(1)}_t$, $W^{(2)}_t$ are independent and normal with mean zero and variance $t$, so that, in polar coordinates,
\[
  P[R_t^2>x]=\int_0^{2\pi}\!\!\int_{\sqrt x}^\infty\frac{1}{2\pi t}\,e^{-\rho^2/2t}\rho\,d\rho\,d\theta=e^{-x/2t};\qquad 0\le x<\infty .
\]
In other words, $\xi_t\defeq R_t^2/2t$ has the exponential distribution with parameter one, for every $t>0$. Consequently, for every real $\alpha$,
\begin{equation}\label{eq:3-3-37:mom}
  E e^{\alpha X_t}=E\bigl[R_t^{-\alpha}\bigr]=(2t)^{-\alpha/2}\,E\bigl[\xi_t^{-\alpha/2}\bigr]
  =(2t)^{-\alpha/2}\int_0^\infty x^{-\alpha/2}e^{-x}\,dx ,
\end{equation}
and the last integral is finite if and only if $-\alpha/2>-1$, i.e.\ $\alpha<2$, in which case it equals $\Gamma\bigl(1-\frac\alpha2\bigr)$. Thus
\[
  E e^{\alpha X_t}=(2t)^{-\alpha/2}\,\Gamma\Bigl(1-\frac\alpha2\Bigr)<\infty;\qquad -\infty<\alpha<2,\ t\ge1,
\]
which is the asserted integrability (and $E e^{\alpha X_t}=\infty$ for $\alpha\ge2$). As a check, $\alpha=0$ gives $1$ and $\alpha=-2$ gives $E R_t^2=2t$. Since $\abs{x}\le e^{x}+e^{-x}$, we obtain in particular
\begin{equation}\label{eq:3-3-37:int}
  E\abs{X_t}\le E e^{X_t}+E e^{-X_t}<\infty;\qquad 1\le t<\infty .
\end{equation}

\emph{Step 2. $X$ is not a martingale.} For $t\ge1$ we have, a.s., $X_t=-\frac12\log R_t^2=-\frac12\log(2t)-\frac12\log\xi_t$, where $\xi_t$ has the same (exponential) distribution as $\xi_1$, and $\log\xi_1$ is integrable by \eqref{eq:3-3-37:int}. Therefore
\[
  E X_t=-\tfrac12\log(2t)-\tfrac12E\log\xi_1=E X_1-\tfrac12\log t;\qquad 1\le t<\infty .
\]
The function $t\mapsto E X_t$ is thus strictly decreasing on $[1,\infty)$, whereas a martingale has constant expectation. Hence $X$ is not a martingale (with respect to any filtration).

\emph{Step 3. $X$ is a local martingale.} We shift the time origin to $t=1$. Let
\[
  \G_u\defeq\F_{1+u},\qquad \rho_u\defeq R_{1+u},\qquad \beta_u\defeq B_{1+u}-B_1;\qquad 0\le u<\infty,
\]
where $B=\{B_t,\F_t;\ 0\le t<\infty\}$ is the standard, one-dimensional Brownian motion of \bk{Proposition 3.21}. The filtration $\{\G_u\}$ satisfies the usual conditions. The process $\beta$ is continuous and $\{\G_u\}$-adapted, $\beta_0=0$, and for $0\le u<v$ the increment $\beta_v-\beta_u=B_{1+v}-B_{1+u}$ is independent of $\F_{1+u}=\G_u$ and normal with mean zero and variance $v-u$; thus $\beta=\{\beta_u,\G_u;\ 0\le u<\infty\}$ is a standard Brownian motion. Subtracting \bk{(3.16)} (with $d=2$, $r=0$) at time $1$ from the same equation at time $1+u$, we obtain a.s.
\begin{equation}\label{eq:3-3-37:semi}
  \rho_u=\rho_0+\beta_u+A_u,\qquad A_u\defeq\int_0^u\frac{ds}{2\rho_s};\qquad 0\le u<\infty .
\end{equation}
Here $A$ is continuous and nondecreasing with $A_0=0$, finite because the integral in \bk{(3.16)} is, and $\{\G_u\}$-adapted because $A_u=\rho_u-\rho_0-\beta_u$ a.s.\ and $\G_0$ contains the null sets. So $\rho$ is a continuous semimartingale relative to $\{\G_u\}$ in the sense of \bk{Definition 3.1}, with local martingale part $\beta$, $\qv{\beta}_u=u$.

For each integer $n\ge1$ define
\[
  S_n\defeq\inf\{u\ge0;\ \rho_u\le 1/n\},\qquad T_n\defeq 1+S_n ,
\]
with $\inf\emptyset=\infty$. Since $\rho$ is continuous and adapted and $[0,1/n]$ is closed, $S_n$ is a stopping time of $\{\G_u\}$ (\bk{Problem 1.2.7}); hence $\{T_n\le t\}=\{S_n\le t-1\}\in\G_{t-1}=\F_t$ for $t\ge1$, and $\{T_n\le t\}=\emptyset$ for $t<1$, so $T_n$ is a stopping time of $\{\F_t\}$. Clearly $S_n\le S_{n+1}$. On $\Omega_0$ the path $u\mapsto\rho_u$ is continuous and strictly positive on $[0,\infty)$, so for every $K<\infty$ it is bounded below on $[0,K]$ by a positive number, which gives $S_n>K$ for all large $n$. Therefore
\begin{equation}\label{eq:3-3-37:Tn}
  T_n\le T_{n+1}\ \ (n\ge1),\qquad P\bigl[\lim_{n\to\infty}T_n=\infty\bigr]=1 .
\end{equation}
We also note that, by continuity,
\begin{equation}\label{eq:3-3-37:lower}
  \rho_s\ge 1/n\quad\text{whenever } 0<s\le S_n
\end{equation}
(if $S_n>0$, then $\rho_s>1/n$ for $0\le s<S_n$ and $\rho_{S_n}=1/n$ if $S_n<\infty$; if $S_n=0$ there is no such $s$).

Fix $n$, and let $g\colon\R\to\R$ be a function of class $C^2$ with $g(x)=-\log x$ for $x\ge 1/n$ (for instance, $g$ equal on $(-\infty,1/n)$ to the second-order Taylor polynomial of $-\log$ at $1/n$). It\^o's rule (\bk{Theorem 3.3}), applied to $g$ and the semimartingale \eqref{eq:3-3-37:semi}, gives a.s.
\begin{equation}\label{eq:3-3-37:ito}
  g(\rho_u)=g(\rho_0)+\int_0^u g'(\rho_s)\,d\beta_s
  +\int_0^u\Bigl[\frac{g'(\rho_s)}{2\rho_s}+\frac12\,g''(\rho_s)\Bigr]ds;\qquad 0\le u<\infty .
\end{equation}
For $x\ge1/n$ we have $g'(x)=-1/x$ and $g''(x)=1/x^2$, so $\frac{1}{2x}g'(x)+\frac12g''(x)=0$; by \eqref{eq:3-3-37:lower}, the Lebesgue integral in \eqref{eq:3-3-37:ito} vanishes when $u$ is replaced by $u\wedge S_n$. We now evaluate \eqref{eq:3-3-37:ito} at $u\wedge S_n$ and claim that, a.s.,
\begin{equation}\label{eq:3-3-37:stopped}
  X_{1+(u\wedge S_n)}=X_1+\int_0^{u\wedge S_n}g'(\rho_s)\,d\beta_s;\qquad 0\le u<\infty .
\end{equation}
Indeed, on $\Omega_0\cap\{S_n>0\}$ we have $\rho_0>1/n$ and $\rho_{u\wedge S_n}\ge1/n$, so $g(\rho_0)=-\log R_1=X_1$ and $g(\rho_{u\wedge S_n})=X_{1+(u\wedge S_n)}$, and \eqref{eq:3-3-37:stopped} is \eqref{eq:3-3-37:ito} at $u\wedge S_n$. On $\{S_n=0\}$ both sides of \eqref{eq:3-3-37:stopped} are equal to $X_1$, the stochastic integral being a continuous process which vanishes at time zero.

The integrand $\{g'(\rho_s);\ 0\le s<\infty\}$ is continuous and adapted, hence in $\cP^*$ for the Brownian motion $\beta$. By \bk{(2.24)} (a sample-path property, valid for integrands in $\cP^*$, as noted before \bk{Proposition 2.24}),
\[
  \int_0^{u\wedge S_n}g'(\rho_s)\,d\beta_s=\int_0^u g'(\rho_s)\ind_{\{s\le S_n\}}\,d\beta_s
  =\int_0^u H_s\,d\beta_s,\qquad H_s\defeq-\frac{1}{\rho_s}\,\ind_{\{0<s\le S_n\}} ,
\]
a.s.\ for all $u$; the second equality holds because the two integrands agree for every $s>0$, by \eqref{eq:3-3-37:lower}, and thus differ at most on the Lebesgue-null set $\{s=0\}$. The process $H$ is $\{\G_u\}$-adapted and left-continuous (the indicator of $(0,S_n]$ is adapted and left-continuous, and off the null set $\Omega\setminus\Omega_0$ we may write $H_s=-\rho_s^{-1}\ind_{\{0<s\le S_n\}}$ with $\rho^{-1}$ continuous), hence progressively measurable, and $\abs{H_s}\le n$ by \eqref{eq:3-3-37:lower}. Therefore $E\int_0^u H_s^2\,ds\le n^2u<\infty$ for every $u$, i.e.\ $H\in\cL^*$, and $N\defeq\{\int_0^uH_s\,d\beta_s,\G_u;\ 0\le u<\infty\}$ is a martingale in $\cM_2^c$.

By \eqref{eq:3-3-37:stopped}, $X_{1+(u\wedge S_n)}=X_1+N_u$ a.s.\ for all $u\ge0$. Since $X_1$ is $\G_0$-measurable and integrable by \eqref{eq:3-3-37:int}, the process $\{X_1+N_u,\G_u;\ 0\le u<\infty\}$ is a martingale. Writing $t=1+u$, so that $t\wedge T_n=1+(u\wedge S_n)$ and $\G_u=\F_t$, we conclude that
\[
  \{X_{t\wedge T_n},\F_t;\ 1\le t<\infty\}\quad\text{is a (continuous) martingale for every } n\ge1 .
\]
Together with \eqref{eq:3-3-37:Tn}, this shows that $X=\{X_t,\F_t;\ 1\le t<\infty\}$ is a continuous local martingale.

\rmk{By \eqref{eq:3-3-37:stopped} and \eqref{eq:3-3-37:Tn}, $X_t=X_1-\int_1^tR_s^{-1}\,dB_s$ for $t\ge1$. The exercise shows that integrability of every $X_t$, and even finiteness of the exponential moments $Ee^{\alpha X_t}$, $\alpha<2$, does not make a continuous local martingale a martingale (cf.\ \bk{Remark 1.5.16}). Step 1 used nothing about $t\ge1$ beyond $t>0$; the restriction to $[1,\infty)$ is needed only because $X_0=-\log R_0=+\infty$.}
\end{kex}

\begin{kex}{Exercise}{3.38}{}
All processes are considered on a probability space $(\Omega,\F,P)$ with a filtration $\{\F_t\}$ satisfying the usual conditions (the standing assumption of \bk{Section 3.3}), and ``local martingale'' is understood in the sense of \bk{Definition 1.5.15}, which does not require the process to vanish at $t=0$. Every path of $X$ and of $A$ is continuous, every path of $A$ is nondecreasing, and we write $\Omega_0\defeq\{X_0=A_0=0\}$, an event of probability one. Both processes are adapted. In (i), $\log Z^{(1)}_t=X_t-\frac12A_t$ and $\log Z^{(-1)}_t=-X_t-\frac12A_t$ are $\F_t$-measurable, and hence so are $A_t$ and $X_t$; in (ii), $X$ is adapted by assumption and $A_t=2\bigl(X_t-\log Z^{(1)}_t\bigr)$. We prove (ii) first and then reduce (i) to it. The equality $\qv{X}=A$ is meant up to indistinguishability, as in \bk{Problem 1.5.17}.

\textbf{(ii)} The continuous local martingale $X$ satisfies $X_0=0$ a.s., so $X\in\cM^{c,\loc}$ and $\qv{X}$ is defined (\bk{Problem 1.5.17}); we choose a version of $\qv{X}$ all of whose paths are continuous and nondecreasing. The process
\[
  Y_t\defeq X_t-\tfrac12A_t,\qquad 0\le t<\infty,
\]
is a continuous semimartingale in the sense of \bk{Definition 3.1}, with $Y_0=0$ a.s., local martingale part $M=X$ and $B=A^+-A^-$, where $A^+\equiv0$ and $A^-=\frac12A$. We write $Z\defeq Z^{(1)}=e^{Y}$. It\^o's rule (\bk{Theorem 3.3}) with $f(y)=e^y$ gives, a.s.,
\begin{equation}\label{eq:3-3-38:a}
  Z_t=1+\int_0^tZ_s\,dX_s-\frac12\int_0^tZ_s\,dA_s+\frac12\int_0^tZ_s\,d\qv{X}_s,\qquad 0\le t<\infty .
\end{equation}
Here $Z$ is continuous and adapted, so the stochastic integral $I_t\defeq\int_0^tZ_s\,dX_s$ is defined and belongs to $\cM^{c,\loc}$ (\bk{Definition 2.23}), and the last two integrals are Lebesgue--Stieltjes integrals, formed path by path. We set
\[
  D_t\defeq\frac12\int_0^tZ_s\,d\qv{X}_s-\frac12\int_0^tZ_s\,dA_s,\qquad
  N_t\defeq Z_t-1-I_t,\qquad 0\le t<\infty,
\]
so that \eqref{eq:3-3-38:a} reads
\begin{equation}\label{eq:3-3-38:b}
  N_t=D_t,\qquad 0\le t<\infty,\quad\text{a.s.}
\end{equation}

\emph{The process $N$ is in $\cM^{c,\loc}$.} By hypothesis there are stopping times $S_k\uparrow\infty$ a.s.\ such that $\{Z_{t\wedge S_k},\F_t\}$ is a martingale for every $k$; then so is $\{Z_{t\wedge S_k}-1,\F_t\}$, and since $Z_0-1=0$ on $\Omega_0$ we have $Z-1\in\cM^{c,\loc}$. Let $\{R_k\}$ be a localizing sequence for $I$. By \bk{Problem 1.3.24}\,(i), applied to the continuous martingales $\pm(Z_{\cdot\wedge S_k}-1)$ and $\pm I_{\cdot\wedge R_k}$ and to the stopping time $S_k\wedge R_k$, the processes $Z_{t\wedge S_k\wedge R_k}-1$ and $I_{t\wedge S_k\wedge R_k}$ are martingales, hence so is $N_{t\wedge S_k\wedge R_k}$; moreover $S_k\wedge R_k\uparrow\infty$ a.s.\ and $N$ is continuous with $N_0=0$ a.s. Thus $N\in\cM^{c,\loc}$.

\emph{The process $D$ is continuous, adapted, and the difference of two nondecreasing processes.} For every $\omega$ and every $t$, the path $s\mapsto Z_s(\omega)$ is continuous and strictly positive, so $t\mapsto\int_0^tZ_s\,d\qv{X}_s$ and $t\mapsto\int_0^tZ_s\,dA_s$ are finite, nondecreasing, and continuous (the integrators are continuous). They are adapted, because each is the limit as $m\to\infty$ of the Riemann--Stieltjes sums $\sum_{j=1}^{m}Z_{t(j-1)/m}\bigl(G_{tj/m}-G_{t(j-1)/m}\bigr)$, $G=\qv{X}$ or $G=A$, which are $\F_t$-measurable.

\emph{$D$ vanishes identically.} By \eqref{eq:3-3-38:b} the continuous semimartingale $N$ has the two decompositions
\[
  N_t=N_t+0=0+D_t,\qquad 0\le t<\infty,\quad\text{a.s.},
\]
the first of the form \bk{(3.1)} with $M=N\in\cM^{c,\loc}$ and $B\equiv0$, the second with $\wt M\equiv0\in\cM^{c,\loc}$ and $\wt B=D$, a continuous, adapted process of finite total variation on every bounded interval. By \bk{Problem 3.2}, $B=\wt B$; that is, there is an event $\Omega_1\subseteq\Omega_0$ of probability one, on which also $\qv{X}_0=0$, such that
\begin{equation}\label{eq:3-3-38:c}
  \int_0^tZ_s\,d\qv{X}_s=\int_0^tZ_s\,dA_s,\qquad 0\le t<\infty,\quad\text{on }\Omega_1 .
\end{equation}

\emph{Conclusion.} Fix $\omega\in\Omega_1$ and $t>0$, and consider on $\bigl((0,t],\B((0,t])\bigr)$ the finite measures
\[
  \nu^+(C)\defeq\int_CZ_s(\omega)\,d\qv{X}_s(\omega),\qquad \nu^-(C)\defeq\int_CZ_s(\omega)\,dA_s(\omega).
\]
By \eqref{eq:3-3-38:c}, $\nu^+((a,b])=\nu^-((a,b])$ for all $0\le a<b\le t$. The intervals $(a,b]$, together with the empty set, form a $\pi$-system which contains $(0,t]$ and generates $\B((0,t])$, so $\nu^+=\nu^-$ by Dynkin's $\pi$-$\lambda$ theorem. The function $s\mapsto 1/Z_s(\omega)$ is continuous, hence bounded, on $[0,t]$, and therefore
\[
  \qv{X}_t(\omega)=\int_{(0,t]}\frac{1}{Z_s(\omega)}\,\nu^+(ds)=\int_{(0,t]}\frac{1}{Z_s(\omega)}\,\nu^-(ds)=A_t(\omega),
\]
where we have used $\qv{X}_0(\omega)=A_0(\omega)=0$ and the fact that $\nu^+$ and $\nu^-$ have the density $Z(\omega)$ with respect to the measures induced on $(0,t]$ by $\qv{X}(\omega)$ and $A(\omega)$. Hence $\qv{X}_t=A_t$ for all $0\le t<\infty$ on $\Omega_1$, which is the assertion.

\textbf{(i)} Since $Z^{(1)}$ is a local martingale by hypothesis, part (ii) yields $\qv{X}=A$ as soon as we know that $X\in\cM^{c,\loc}$. This is what we now prove.

\emph{Step 1: localization.} For $n\ge1$ let
\[
  T_n\defeq\inf\{t\ge0:\ \abs{X_t}\ge n\ \text{or}\ A_t\ge n\}\qquad(\inf\emptyset=\infty),
\]
the first time the continuous, adapted, $\R^2$-valued process $(X,A)$ enters the closed set $\{(x,a):\abs{x}\ge n\text{ or }a\ge n\}$. By \bk{Problem 1.2.7}, $T_n$ is a stopping time. The sequence $\{T_n\}$ is nondecreasing, and $T_n(\omega)\to\infty$ for every $\omega$, because the paths of $X$ and $A$ are bounded on every bounded interval. On $\Omega_0$ we have $T_n>0$, $\abs{X_s}<n$ and $0\le A_s<n$ for $s<T_n$, and therefore, by continuity,
\begin{equation}\label{eq:3-3-38:d}
  \abs{X_{t\wedge T_n}}\le n,\qquad 0\le A_{t\wedge T_n}\le n,\qquad 0\le t<\infty,\quad\text{on }\Omega_0 .
\end{equation}

\emph{Step 2: the stopped exponentials are martingales.} Fix $n\ge1$ and $\theta\in\R$. The random variable $Z^{(\theta)}_{t\wedge T_n}$ is $\F_t$-measurable, since $Z^{(\theta)}$ is continuous and adapted, hence progressively measurable (\bk{Propositions 1.1.13} and \bk{1.2.18}), and by \eqref{eq:3-3-38:d}
\begin{equation}\label{eq:3-3-38:e}
  0<Z^{(\theta)}_{t\wedge T_n}\le e^{\abs{\theta}n},\qquad 0\le t<\infty,\quad\text{on }\Omega_0 .
\end{equation}
Let $\{S_k\}_{k=1}^\infty$ be a nondecreasing sequence of stopping times with $S_k\to\infty$ a.s.\ such that $\{Z^{(\theta)}_{t\wedge S_k},\F_t\}$ is a martingale for every $k$. By \bk{Problem 1.3.24}\,(i), $\{Z^{(\theta)}_{t\wedge S_k\wedge T_n},\F_t\}$ is a martingale, so for $0\le s<t$,
\[
  E\bigl[Z^{(\theta)}_{t\wedge S_k\wedge T_n}\bigm|\F_s\bigr]=Z^{(\theta)}_{s\wedge S_k\wedge T_n}\qquad\text{a.s.}
\]
As $k\to\infty$ we have $Z^{(\theta)}_{u\wedge S_k\wedge T_n}\to Z^{(\theta)}_{u\wedge T_n}$ a.s.\ for $u=s,t$, and all these random variables are bounded by $e^{\abs{\theta}n}$ on $\Omega_0$, by \eqref{eq:3-3-38:e} applied at the times $u\wedge S_k$. The bounded convergence theorem for conditional expectations gives
\begin{equation}\label{eq:3-3-38:f}
  E\bigl[Z^{(\theta)}_{t\wedge T_n}\bigm|\F_s\bigr]=Z^{(\theta)}_{s\wedge T_n}\quad\text{a.s.},\qquad 0\le s<t<\infty,\ \theta\in\R .
\end{equation}

\emph{Step 3: differentiation at $\theta=0$.} Fix $n\ge1$ and $0\le s<t$. For $u\in\{s,t\}$ and $\omega\in\Omega_0$ consider the smooth function
\[
  g_u(\theta)\defeq Z^{(\theta)}_{u\wedge T_n}=\exp\bigl(\theta X_{u\wedge T_n}-\tfrac12\theta^2A_{u\wedge T_n}\bigr),\qquad
  g_u'(\theta)=\bigl(X_{u\wedge T_n}-\theta A_{u\wedge T_n}\bigr)g_u(\theta).
\]
We have $g_u(0)=1$, $g_u'(0)=X_{u\wedge T_n}$, and $\abs{g_u'(\theta)}\le2ne^{n}$ for $\abs{\theta}\le1$, by \eqref{eq:3-3-38:d} and \eqref{eq:3-3-38:e}. Consequently, by the mean value theorem,
\begin{equation}\label{eq:3-3-38:g}
  \abs[\big]{m\bigl(g_u(1/m)-1\bigr)}\le2ne^{n}\quad(m\ge1),\qquad \lim_{m\to\infty}m\bigl(g_u(1/m)-1\bigr)=X_{u\wedge T_n},\qquad\text{on }\Omega_0 .
\end{equation}
Now \eqref{eq:3-3-38:f} with $\theta=1/m$ gives
\[
  E\bigl[m\bigl(g_t(1/m)-1\bigr)\bigm|\F_s\bigr]=m\bigl(g_s(1/m)-1\bigr)\qquad\text{a.s.},\quad m\ge1 .
\]
Letting $m\to\infty$ and using \eqref{eq:3-3-38:g} and the bounded convergence theorem for conditional expectations on the left-hand side, we obtain
\[
  E\bigl[X_{t\wedge T_n}\bigm|\F_s\bigr]=X_{s\wedge T_n}\qquad\text{a.s.}
\]
Since $X_{t\wedge T_n}$ is $\F_t$-measurable and bounded by $n$ a.s.\ (\eqref{eq:3-3-38:d}), the process $\{X_{t\wedge T_n},\F_t;\ 0\le t<\infty\}$ is a martingale, for every $n\ge1$. As $T_n\uparrow\infty$, $X$ is continuous and $X_0=0$ a.s., we conclude that $X\in\cM^{c,\loc}$ (\bk{Definition 1.5.15}).

\emph{Step 4.} Now $X$ and $Z^{(1)}$ are local martingales, and part (ii) gives $\qv{X}=A$.

\rmk{Only the values $\theta=\pm1$ (for the adaptedness of $X$ and $A$), $\theta=1$ (for Step 4) and a sequence $\theta_m\to0$, $\theta_m\neq0$ (for Step 3) were used in (i). One can also bypass part (ii): differentiating \eqref{eq:3-3-38:f} twice at $\theta=0$, with the same domination, shows that $X^2_{t\wedge T_n}-A_{t\wedge T_n}$ is a martingale for every $n$, so that $X^2-A\in\cM^{c,\loc}$ and $\qv{X}=A$ by the uniqueness assertion of \bk{Problem 1.5.17}.}
\end{kex}

\begin{kex}{Exercise}{3.39}{}
\emph{Interpretation and a correction.} We write $X_t=X_0+M_t+B_t$ for the decomposition \bk{(3.1)}, and, as in \bk{Remark 3.5}, $\int_0^t f(X_s)\,dX_s\defeq\int_0^t f(X_s)\,dM_s+\int_0^t f(X_s)\,dB_s$. In keeping with \bk{Definition 3.1}, a ``process of bounded variation'' is understood to be a process $B^{(n)}$ whose paths $s\mapsto B^{(n)}_s(\omega)$ are, for $P$-a.e.\ $\omega$, \emph{continuous} and of finite total variation on every bounded interval; the integral on the left-hand side is the Lebesgue--Stieltjes integral, computed path by path. Moreover, we assume that the convergence $B^{(n)}_t\to X_t$, a.s., holds for every $t\ge 0$, i.e., also at $t=0$:
\begin{equation}\label{eq:3-3-39:hyp}
  P\Bigl[\lim_{n\to\infty}B^{(n)}_t=X_t\Bigr]=1\quad\text{for every } 0\le t<\infty .
\end{equation}
The remark at the end shows that neither the convergence at $t=0$ nor the continuity of the paths of $B^{(n)}$ can be dispensed with. We prove the assertion under \eqref{eq:3-3-39:hyp}.

Let $F(x)\defeq\int_0^x f(y)\,dy$, $x\in\R$. Since $f$ is of class $C^1$, the function $F$ is of class $C^2$, with $F'=f$ and $F''=f'$.

\begin{lem}
Let $b\colon[0,t]\to\R$ be continuous and of finite total variation $V$ on $[0,t]$. Then
\begin{equation}\label{eq:3-3-39:chain}
  \int_0^t f(b(s))\,db(s)=F(b(t))-F(b(0)).
\end{equation}
\end{lem}
\begin{proof}
Let $\mu$ be the signed measure on $((0,t],\B((0,t]))$ with $\mu((u,v])=b(v)-b(u)$ for $0\le u\le v\le t$, and $\abs{\mu}$ its total variation measure, so that $\abs{\mu}((0,t])=V$; the left-hand side of \eqref{eq:3-3-39:chain} is $\int_{(0,t]}f(b(s))\,\mu(ds)$, which is well defined because $f\circ b$ is continuous and bounded on $[0,t]$. Let $K\defeq\max_{0\le s\le t}\abs{b(s)}<\infty$ and, for $\delta>0$,
\[
  \rho(\delta)\defeq\sup\{\abs{f(x)-f(y)};\ x,y\in[-K,K],\ \abs{x-y}\le\delta\},
\]
so that $\rho(\delta)\to0$ as $\delta\downarrow0$ by the uniform continuity of $f$ on $[-K,K]$. Fix $\eps>0$. By the uniform continuity of $b$ on $[0,t]$ there is a partition $0=s_0<s_1<\dots<s_k=t$ with $\abs{b(u)-b(v)}\le\eps$ whenever $u,v\in[s_{i-1},s_i]$, $1\le i\le k$. By the mean value theorem there is, for each $i$, a point $\xi_i$ in the closed interval with endpoints $b(s_{i-1})$, $b(s_i)$ such that
\[
  F(b(s_i))-F(b(s_{i-1}))=f(\xi_i)\,\bigl(b(s_i)-b(s_{i-1})\bigr)=\int_{(s_{i-1},s_i]}f(\xi_i)\,\mu(ds).
\]
For $s\in(s_{i-1},s_i]$ we have $\xi_i,b(s)\in[-K,K]$ and $\abs{\xi_i-b(s)}\le\eps$, because both numbers lie in an interval of length at most $\eps$, namely the one spanned by the values of $b$ on $[s_{i-1},s_i]$. Summing over $i$, we obtain
\[
  \Bigl|F(b(t))-F(b(0))-\int_{(0,t]}f(b(s))\,\mu(ds)\Bigr|
  \le\sum_{i=1}^k\int_{(s_{i-1},s_i]}\abs{f(\xi_i)-f(b(s))}\,\abs{\mu}(ds)\le\rho(\eps)\,V .
\]
Letting $\eps\downarrow0$ yields \eqref{eq:3-3-39:chain}.
\end{proof}

\emph{Proof of the assertion.} Fix $t>0$. For each $n\ge1$ let $N_n$ be a $P$-null event off which the path $s\mapsto B^{(n)}_s(\omega)$ is continuous and of finite variation on $[0,t]$. Off $N_n$ the Lemma, applied to $b=B^{(n)}_\cdot(\omega)$, gives
\begin{equation}\label{eq:3-3-39:a}
  \int_0^t f(B^{(n)}_s)\,dB^{(n)}_s=F(B^{(n)}_t)-F(B^{(n)}_0).
\end{equation}
By \eqref{eq:3-3-39:hyp}, applied at the two time points $t$ and $0$, there is a $P$-null event $N'$ off which $B^{(n)}_t\to X_t$ and $B^{(n)}_0\to X_0$ as $n\to\infty$; since $F$ is continuous, off the null event $N'\cup\bigcup_{n\ge1}N_n$ the right-hand side of \eqref{eq:3-3-39:a} converges, and
\begin{equation}\label{eq:3-3-39:b}
  \lim_{n\to\infty}\int_0^t f(B^{(n)}_s)\,dB^{(n)}_s=F(X_t)-F(X_0)\quad\text{a.s. } P.
\end{equation}
On the other hand, $F$ is of class $C^2$ and $X$ is a continuous semimartingale with decomposition \bk{(3.1)}, so It\^o's rule (\bk{Theorem 3.3}) applies to $F(X)$ and gives, a.s.\ $P$,
\begin{equation}\label{eq:3-3-39:c}
  F(X_t)-F(X_0)=\int_0^t f(X_s)\,dM_s+\int_0^t f(X_s)\,dB_s+\frac12\int_0^t f'(X_s)\,d\qv{M}_s .
\end{equation}
The first two terms on the right-hand side of \eqref{eq:3-3-39:c} add up to $\int_0^t f(X_s)\,dX_s$. Combining \eqref{eq:3-3-39:b} and \eqref{eq:3-3-39:c} we conclude that
\[
  \lim_{n \to \infty} \int_0^t f(B_s^{(n)})\,dB_s^{(n)} = \int_0^t f(X_s)\,dX_s + \frac{1}{2} \int_0^t f'(X_s)\,d\qv{M}_s
  \quad\text{a.s. } P,
\]
for every fixed $t>0$.
\rmk{\textbf{(a)} As printed, the hypothesis requires $B^{(n)}_t\to X_t$ only for $t>0$, and this is not sufficient. Take $X\equiv0$ (so $M\equiv0$, $B\equiv0$), the deterministic continuous, nonincreasing processes $B^{(n)}_s\defeq(1-ns)^+$, and $f(x)=x$. Then $B^{(n)}_t=0$ for $n\ge1/t$, so $B^{(n)}_t\to X_t$ for every $t>0$, but $B^{(n)}_0=1\not\to X_0$; by \eqref{eq:3-3-39:chain}, $\int_0^t B^{(n)}_s\,dB^{(n)}_s=\tfrac12\bigl((B^{(n)}_t)^2-(B^{(n)}_0)^2\bigr)=-\tfrac12$ for $n\ge 1/t$, whereas the right-hand side of the assertion is $0$. Thus convergence at $t=0$ has to be added to the hypotheses (it holds, for instance, when $B^{(n)}_0=X_0$ for all $n$).
\textbf{(b)} Continuity of the paths of $B^{(n)}$ cannot be dropped either. Take again $X\equiv0$ and $f(x)=x$, and the right-continuous step functions $B^{(n)}_s\defeq\ind_{[1/n,\,2/n)}(s)$, which are of bounded variation, vanish at $s=0$, and satisfy $B^{(n)}_t=0=X_t$ for $n\ge2/t$. The Lebesgue--Stieltjes measure of $B^{(n)}$ is $\delta_{1/n}-\delta_{2/n}$, so that $\int_{(0,t]}B^{(n)}_s\,dB^{(n)}_s=B^{(n)}_{1/n}-B^{(n)}_{2/n}=1$ for $n\ge2/t$, which does not converge to $0$.
\textbf{(c)} The processes $B^{(n)}$ need not be adapted: only their paths and the a.s.\ convergence at the two time points $0$ and $t$ enter the argument. The correction term $\frac12\int_0^t f'(X_s)\,d\qv{M}_s$ shows that It\^o integrals are not limits of ordinary integrals along smooth approximations of the integrator; when $f$ is of class $C^2$, so that $f(X)$ is a continuous semimartingale with local martingale part $\int_0^\cdot f'(X_s)\,dM_s$, the limit is the Fisk--Stratonovich integral $\int_0^t f(X_s)\circ dX_s$ of \bk{Definition 3.13}.}
\end{kex}

\solsection{3.4}{Representations of Continuous Martingales in Terms of Brownian Motion}

\begin{kex}{Exercise}{4.4}{}
\emph{Conventions.} The cross-variations $A^{(i,j)}$ are those of \bk{Problem 1.5.17}, which is stated for elements of $\cM^{c,\loc}$ and under the standing assumption of \bk{Section 1.5} that $\{\F_t\}$ satisfies the usual conditions. Accordingly we take $M^{(i)}\in\cM^{c,\loc}$, $1\le i\le d$ (so $M^{(i)}_0=0$ a.s., as in \bk{Theorem 4.2}), with respect to a filtration $\{\F_t\}$ satisfying the usual conditions; this is needed in all three parts merely to give the statement a meaning. By \bk{Problem 1.5.17}, $A^{(i,j)}$ is adapted and, for $P$-a.e.\ $\omega$, the path $t\mapsto A^{(i,j)}_t(\omega)$ is continuous, of bounded variation on every interval $[0,t]$, and vanishes at $t=0$; it is determined only up to indistinguishability. The definition of $T_s(\omega)$ requires these path properties for the $\omega$ in question, so we fix versions for which they hold for \emph{every} $\omega\in\Omega$: we replace $A^{(i,j)}$ by zero on a $P$-null event of $\F$ off which they hold; this event belongs to $\F_0$, so adaptedness is preserved. For the same reason we may and do assume that every path of every $M^{(i)}$ is continuous.

\emph{Step 0: the functions $A(\omega)$ and $T(\omega)$.} We fix $\omega$ and suppress it. Let $f:[0,\infty)\to\R$ be continuous and of bounded variation on bounded intervals, and let $\check f_t$ be its total variation on $[0,t]$. We need two facts from real analysis.

\textbf{(a)} \emph{$\check f_t=\sup_{m\ge1}\sum_{k=1}^{2^m}\abs{f(tk2^{-m})-f(t(k-1)2^{-m})}$.} The right-hand side is a supremum of sums over particular partitions of $[0,t]$, hence is at most $\check f_t$. Conversely, given $0=s_0<\dots<s_p=t$ (with $t>0$), let $s^m_l\defeq t2^{-m}\lfloor 2^ms_l/t\rfloor$; these are points of the $m$-th grid with $0=s_0^m\le\dots\le s^m_p=t$, so by the triangle inequality $\sum_{l}\abs{f(s^m_l)-f(s^m_{l-1})}$ is at most the $m$-th dyadic sum, and letting $m\to\infty$ (continuity of $f$) we obtain $\sum_l\abs{f(s_l)-f(s_{l-1})}\le$ the right-hand side.

\textbf{(b)} \emph{$t\mapsto\check f_t$ is nondecreasing and continuous, with $\check f_0=0$ and $\abs{f(t)-f(u)}\le\check f_t-\check f_u$ for $0\le u\le t$.} Write $V(a,b)$ for the total variation of $f$ on $[a,b]$; it is additive over adjacent intervals, so $\check f_t-\check f_u=V(u,t)\ge\abs{f(t)-f(u)}$, which gives monotonicity and the inequality. For continuity, let $0\le a<b$ and $\eps>0$, and choose a partition $a=s_0<s_1<\dots<s_p=b$ with $\sum_{l=1}^p\abs{f(s_l)-f(s_{l-1})}>V(a,b)-\eps$. Inserting a further point can only increase the sum, so by the continuity of $f$ at $a$ we may assume $\abs{f(s_1)-f(a)}<\eps$. Then $V(a,b)-\eps<\eps+V(s_1,b)$, i.e.\ $V(a,s_1)<2\eps$; hence $\lim_{u\downarrow a}V(a,u)=0$, which is right-continuity of $\check f$ at $a$. Symmetrically (inserting a point near $b$), $\lim_{u\uparrow b}V(u,b)=0$, which is left-continuity at $b$.

Applying (b) to each $A^{(i,j)}(\omega)$ we see that, for every $\omega$, the function
\[
  \varphi_t\defeq A_t+t,\qquad 0\le t<\infty,
\]
is continuous with $\varphi_0=0$, and
\begin{equation}\label{eq:3-4-04:a}
  \varphi_t-\varphi_u\ \ge\ t-u\qquad\text{for } 0\le u\le t .
\end{equation}
Thus $\varphi$ is strictly increasing and $\varphi_t\ge t\to\infty$, so $\varphi$ is a bijection of $[0,\infty)$ onto itself. Its inverse $s\mapsto T_s$ is therefore well defined, strictly increasing and continuous, and satisfies
\begin{equation}\label{eq:3-4-04:b}
  T_0=0,\qquad T_s\le s,\qquad 0\le T_s-T_r\le s-r \quad (0\le r\le s),\qquad T_s\le t\iff s\le A_t+t ;
\end{equation}
the inequalities follow from \eqref{eq:3-4-04:a} with $t=T_s$, $u=T_r$, and the equivalence from the monotonicity of $\varphi$ and $\varphi_{T_s}=s$.

\textbf{(i)} By (a), for each $t$ the random variable $\check A^{(i,j)}_t$ is a countable supremum of finite sums of absolute values of $\F_t$-measurable random variables, hence is $\F_t$-measurable; so is $A_t$. By \eqref{eq:3-4-04:b},
\[
  \{T_s\le t\}=\{A_t+t\ge s\}\in\F_t,\qquad 0\le t<\infty ,
\]
so $T_s$ is a stopping time of $\{\F_t\}$, for every $s\ge0$.

\textbf{(ii)} Each $\G_s=\F_{T_s}$ is a sub-$\sigma$-field of $\F$ (\bk{Problem 1.2.13}), and for $r\le s$ we have $T_r\le T_s$ on $\Omega$, whence $\G_r\subseteq\G_s$ by \bk{Lemma 1.2.15}; thus $\{\G_s\}$ is a filtration. Since $T_0\equiv0$, \bk{Problem 1.2.13} gives $\G_0=\F_0$, which contains all the $P$-null events of $\F$. It remains to prove right-continuity. Let $s\ge0$ and $B\in\bigcap_{u>s}\G_u$; then $B\in\F$ and $B\cap\{T_{s+1/n}\le t\}\in\F_t$ for all $n\ge1$, $t\ge0$. Because $u\mapsto T_u$ is strictly increasing and continuous,
\[
  \{T_s<t\}=\bigcup_{n=1}^\infty\{T_{s+1/n}\le t\},\qquad 0\le t<\infty
\]
(if $T_s<t$ then $T_{s+1/n}<t$ for large $n$ by continuity; if $T_{s+1/n}\le t$ then $T_s<T_{s+1/n}\le t$). Hence $B\cap\{T_s<t\}\in\F_t$ for every $t\ge0$, and therefore, for every $m\ge1$,
\[
  B\cap\{T_s\le t\}=\bigcap_{k=m}^\infty\bigl(B\cap\{T_s<t+\tfrac1k\}\bigr)\in\F_{t+1/m}.
\]
So $B\cap\{T_s\le t\}\in\F_{t+}=\F_t$ by the right-continuity of $\{\F_t\}$, i.e.\ $B\in\F_{T_s}=\G_s$. Thus $\bigcap_{u>s}\G_u=\G_s$, and $\{\G_s\}$ satisfies the usual conditions.

\textbf{(iii)} The class $\cM^{c,\loc}$ for the time-changed processes is of course understood with respect to $\{\G_s\}$, which satisfies the usual conditions by (ii). Everything rests on the following lemma.

\begin{lem}
Let $L=\{L_t,\F_t;\ 0\le t<\infty\}\in\cM^{c,\loc}$ have every path continuous. Then $\{L_{T_s},\G_s;\ 0\le s<\infty\}\in\cM^{c,\loc}$.
\end{lem}
\begin{proof}
Put $\wt L_s\defeq L_{T_s}$. Its paths are continuous, as compositions of continuous functions, and $\wt L_0=L_0=0$ a.s. The process $L$ is continuous and adapted, hence progressively measurable (\bk{Proposition 1.1.13}), and $T_s$ is a finite stopping time, so $\wt L_s$ is $\F_{T_s}$-measurable by \bk{Proposition 1.2.18}; thus $\wt L$ is adapted to $\{\G_s\}$.

By \bk{Definition 1.5.15} there is a nondecreasing sequence $\{S_n\}$ of stopping times of $\{\F_t\}$ with $P[\lim_n S_n=\infty]=1$ such that $X^{(n)}\defeq\{L_{t\wedge S_n},\F_t;\ 0\le t<\infty\}$ is a martingale for each $n$. We set
\[
  \sigma_n\defeq\begin{cases} A_{S_n}+S_n=\varphi_{S_n}, & S_n<\infty,\\ \infty, & S_n=\infty.\end{cases}
\]
Since $\varphi$ is strictly increasing with inverse $T$, we have for every $s\ge 0$
\begin{equation}\label{eq:3-4-04:c}
  \{\sigma_n\le s\}=\{S_n\le T_s\},\qquad T_{s\wedge\sigma_n}=T_s\wedge S_n .
\end{equation}
(For the second identity: if $\sigma_n\ge s$ then $S_n\ge T_s$ and both sides equal $T_s$; if $\sigma_n<s$ then $S_n<\infty$ and $T_{\sigma_n}=S_n<T_s$.) By \bk{Lemma 1.2.16}, $\{S_n\le T_s\}\in\F_{T_s}=\G_s$, so each $\sigma_n$ is a stopping time of $\{\G_s\}$. The sequence $\{\sigma_n\}$ is nondecreasing because $\{S_n\}$ is and $\varphi$ is increasing, and $\sigma_n\ge S_n$ by \eqref{eq:3-4-04:a}, so $P[\lim_n\sigma_n=\infty]=1$.

Fix $n$. By \eqref{eq:3-4-04:c}, $\wt L_{s\wedge\sigma_n}=L_{T_s\wedge S_n}=X^{(n)}_{T_s}$ for all $s\ge0$. The process $X^{(n)}$ is a continuous martingale of $\{\F_t\}$, hence progressively measurable, so $X^{(n)}_{T_s}$ is $\G_s$-measurable by \bk{Proposition 1.2.18}. For $0\le r\le s$, $T_r\le T_s$ are stopping times of $\{\F_t\}$ bounded by the constant $s$ (see \eqref{eq:3-4-04:b}). The optional sampling theorem for bounded stopping times, \bk{Problem 1.3.23}\,(i) (cf.\ \bk{Theorem 1.3.22}), applied to the right-continuous submartingales $X^{(n)}$ and $-X^{(n)}$, shows that $X^{(n)}_{T_s}$ is integrable and
\[
  E\bigl[\wt L_{s\wedge\sigma_n}\,\big|\,\G_r\bigr]=E\bigl[X^{(n)}_{T_s}\,\big|\,\F_{T_r}\bigr]=X^{(n)}_{T_r}=\wt L_{r\wedge\sigma_n}\qquad\text{a.s. } P .
\]
Hence $\{\wt L_{s\wedge\sigma_n},\G_s;\ 0\le s<\infty\}$ is a martingale for every $n$, and $\wt L\in\cM^{c,\loc}$ by \bk{Definition 1.5.15}.
\end{proof}

Applying the Lemma to $L=M^{(i)}$ we obtain $N^{(i)}\in\cM^{c,\loc}$ for $1\le i\le d$.

Now fix $1\le i,j\le d$ and define $C_s\defeq A^{(i,j)}_{T_s}$, $0\le s<\infty$. For $0\le r\le s$ and every $\omega$, we obtain from (b), from the fact that every $\check A^{(k,l)}$ is nondecreasing, and from $A_{T_u}=u-T_u$,
\begin{equation}\label{eq:3-4-04:d}
  \abs{C_s-C_r}\le\check A^{(i,j)}_{T_s}-\check A^{(i,j)}_{T_r}\le A_{T_s}-A_{T_r}=(s-r)-(T_s-T_r)\le s-r .
\end{equation}
Thus every path of $C$ is Lipschitz continuous with constant $1$; in particular it is continuous, of bounded variation on bounded intervals, and absolutely continuous. Moreover $C_0=A^{(i,j)}_0=0$, and $C$ is adapted to $\{\G_s\}$ by \bk{Proposition 1.1.13} and \bk{Proposition 1.2.18}, because $A^{(i,j)}$ is continuous and adapted to $\{\F_t\}$.

By \bk{Problem 1.5.17}, $L\defeq M^{(i)}M^{(j)}-A^{(i,j)}$ belongs to $\cM^{c,\loc}$ (relative to $\{\F_t\}$), and all its paths are continuous. The Lemma shows that
\[
  L_{T_s}=N^{(i)}_sN^{(j)}_s-C_s,\qquad 0\le s<\infty,
\]
is in $\cM^{c,\loc}$ relative to $\{\G_s\}$. So $C$ is an adapted, continuous process of bounded variation with $C_0=0$ such that $N^{(i)}N^{(j)}-C\in\cM^{c,\loc}$, and the uniqueness assertion of \bk{Problem 1.5.17} (applied on $(\Omega,\F,P)$ with the filtration $\{\G_s\}$) yields
\[
  P\bigl[\qv{N^{(i)},N^{(j)}}_s=A^{(i,j)}_{T_s};\ \forall\,0\le s<\infty\bigr]=1 .
\]
By \eqref{eq:3-4-04:d}, for $P$-a.e.\ $\omega$ the function $s\mapsto\qv{N^{(i)},N^{(j)}}_s(\omega)$ is Lipschitz continuous with constant $1$, hence absolutely continuous.

\rmk{The vector $N=(N^{(1)},\dots,N^{(d)})$, with the filtration $\{\G_s\}$, therefore satisfies the hypotheses of \bk{Theorem 4.2}. We have in fact proved more than absolute continuity: the densities $\frac{d}{ds}\qv{N^{(i)},N^{(j)}}_s$, which exist for Lebesgue-a.e.\ $s$, are bounded by $1$ in absolute value. The original processes are recovered from $N$ by the inverse time-change, $M^{(i)}_t=N^{(i)}_{A_t+t}$.}
\end{kex}

\begin{kex}{Problem}{4.5}{}
The book's Section 3.8 proves the right-continuity in (i), part (ii), part (iv) and, in (vi), the identity \bk{(4.14)} for indicators of Borel subsets of $[a,b]$ (its equation \bk{(8.4)}, obtained from the case of intervals $[t_1,t_2)$ by a Dynkin-system argument); it declares (iii) and (v) immediate and the passage from indicators to general $G$ in (vi) straightforward, and it does not treat the finiteness of $T$ on $[0,S)$ or the second assertion of (i). We supply these missing arguments here, and add a caveat about \bk{(4.13)} when $A$ attains the value $S$. In (vi), Steps~1 and~2 redo the indicator case with the details the book omits (the domain of $G\circ T$, the endpoint $b$ and the point $s=A(b)$), and Step~3 is the part the book leaves out. Throughout, $S=A(\infty)=\lim_{t\to\infty}A(t)=\sup_{t\ge0}A(t)$, and for $0\le s<S$ we write
\[
  U_s\defeq\{t\ge0;\ A(t)>s\},\qquad\text{so that } T(s)=\inf U_s .
\]
Since $A$ is nondecreasing, $U_s$ is an \emph{up-set}: $t\in U_s$ and $t'\ge t$ imply $t'\in U_s$. Since $A$ is continuous, $U_s$ is open in $[0,\infty)$, and $0\notin U_s$ because $A(0)=0\le s$. Hence
\begin{equation}\label{eq:3-4-05:a}
  U_s=(T(s),\infty),\qquad\text{i.e.}\qquad A(t)>s\iff t>T(s)\qquad(0\le s<S,\ 0\le t<\infty),
\end{equation}
as soon as we know that $U_s\ne\emptyset$, which is the first point of (i).

\textbf{(i)} Let $0\le s<S$. Because $s<S=\sup_{t\ge0}A(t)$, there is a $t$ with $A(t)>s$; thus $U_s\neq\emptyset$ and $T(s)=\inf U_s\in[0,\infty)$. If $0\le s_1\le s_2<S$, then $U_{s_2}\subseteq U_{s_1}$, whence $T(s_1)\le T(s_2)$: the function $T$ is nondecreasing on $[0,S)$ (and on $[0,\infty)$, since $T=\infty$ on $[S,\infty)$). The right-continuity on $[0,S)$ is proved in the book's Section 3.8.

Suppose now that $A(t)<S$ for every $t\ge0$, and let $K\in[0,\infty)$ be given. Then $A(K)<S$, and for every $s\in[A(K),S)$ and every $t\le K$ we have $A(t)\le A(K)\le s$, so that $U_s\subseteq(K,\infty)$ and $T(s)\ge K$. As $K$ is arbitrary and $T$ is nondecreasing, $\lim_{s\uparrow S}T(s)=\infty$.

\rmk{The hypothesis of the second assertion cannot be dropped: if $A(t_0)=S$ for some $t_0<\infty$ (which forces $S<\infty$), then $U_s\supseteq[t_0,\infty)$ for every $s<S$, hence $T(s)\le t_0$ on $[0,S)$ and $\lim_{s\uparrow S}T(s)\le t_0<\infty$, although $T(S)=\infty$.}

\textbf{(ii)} This is proved in the book's Section 3.8. (It also follows from \eqref{eq:3-4-05:a}: for $s<S$, $A(t)>s$ for $t>T(s)$ and $A(t)\le s$ for $t\le T(s)$, so the continuity of $A$ at $T(s)$ gives $A(T(s))=s$.)

\textbf{(iii)} Fix $t\in[0,\infty)$ and put $R\defeq\sup\{\tau\ge t\colon A(\tau)=A(t)\}\in[t,\infty]$; the set contains $\tau=t$.

\emph{Case $A(t)=S$.} Then $T(A(t))=T(S)=\infty$ by definition. On the other hand $A(t)\le A(\tau)\le S=A(t)$ for every $\tau\ge t$, so the set defining $R$ is $[t,\infty)$ and $R=\infty$.

\emph{Case $A(t)<S$.} We apply \eqref{eq:3-4-05:a} with $s=A(t)$: for $\tau\ge t$ we have $A(\tau)\ge A(t)$, and therefore
\[
  A(\tau)=A(t)\iff A(\tau)\le A(t)\iff\tau\notin U_{A(t)}\iff\tau\le T(A(t)).
\]
In particular (take $\tau=t$) $t\le T(A(t))$, and the set defining $R$ is the interval $[t,T(A(t))]$, whose supremum is $T(A(t))$. Thus $R=T(A(t))$ in both cases.

\textbf{(iv)} The continuity of $\varphi\circ T$ on $[0,S)$ is proved in the book's Section 3.8, and so is \bk{(4.13)} for those $t$ with $T(A(t))<\infty$, that is, with $A(t)<S$: by (ii) $A(T(A(t)))=A(t)$, by (iii) $T(A(t))\ge t$, and the property of $\varphi$ gives $\varphi(T(A(t)))=\varphi(t)$.

\rmk{If $A(t)=S$ for some finite $t$, then $T(A(t))=\infty$ and the left-hand side of \bk{(4.13)} is not defined by the statement as printed. In this case put $t_0\defeq\min\{t\ge0;\ A(t)=S\}$ (the minimum exists by continuity of $A$). Then $A\equiv S$ on $[t_0,\infty)$, the property of $\varphi$ gives $\varphi(t)=\varphi(t_0)$ for all $t\ge t_0$, and so $\varphi(\infty)\defeq\lim_{t\to\infty}\varphi(t)=\varphi(t_0)$ exists. With this convention \bk{(4.13)} holds for every $t\in[0,\infty)$: for $t\ge t_0$ both sides equal $\varphi(t_0)$. This is the convention under which \bk{(4.13)} is used in \bk{Problem 4.7}, where $\varphi=M_\cdot(\omega)$ and $M_\infty=\lim_{t\to\infty}M_t$; in \bk{Theorem 4.6} one has $S=\infty$ and the question does not arise.}

\textbf{(v)} Let $0\le t,s<\infty$.

If $s<A(t)$, then $s<S$ and $t\in U_s=(T(s),\infty)$ by \eqref{eq:3-4-05:a}, so $T(s)<t$. Conversely, if $T(s)<t$, then $T(s)<\infty$, hence $s<S$ by the definition of $T$, and $t\in(T(s),\infty)=U_s$, i.e.\ $A(t)>s$. This proves $s<A(t)\Leftrightarrow T(s)<t$.

If $T(s)\le t$, then again $s<S$, and for every $u>t$ we have $u>T(s)$, hence $A(u)>s$ by \eqref{eq:3-4-05:a}. Letting $u\downarrow t$ and using the continuity of $A$ we obtain $A(t)\ge s$.

\rmk{The second implication cannot be reversed: if $A$ is constant on an interval $[t,t']$ with $t<t'$ and $A(t)<S$, then for $s=A(t)$ we have $s\le A(t)$ but $T(s)\ge t'>t$ by \textup{(iii)}.}

\textbf{(vi)} Here ``measurable'' means Borel measurable, $\int_a^bG(t)\,dA(t)$ is the integral of $G$ over $[a,b]$ with respect to the Lebesgue--Stieltjes measure $\mu_A$ of $A$ (the unique measure on $\B([0,\infty))$ with $\mu_A([0,v])=A(v)$ for $0\le v<\infty$; equivalently, $\mu_A(\{0\})=A(0)=0$ and $\mu_A((u,v])=A(v)-A(u)$ for $0\le u\le v<\infty$), and the right-hand side is a Lebesgue integral. The measure $\mu_A$ has no atoms: $\mu_A(\{0\})=0$ by definition, and for $u>0$ the continuity of $A$ gives $\mu_A(\{u\})=\lim_{n\to\infty}\mu_A((u-\tfrac1n,u])=\lim_{n\to\infty}\bigl(A(u)-A(u-\tfrac1n)\bigr)=0$. Consequently $\mu_A([u,v))=\mu_A([u,v])=A(v)-A(u)$ for all $0\le u\le v<\infty$. If $a=b$, or more generally if $A(a)=A(b)$, both sides of \bk{(4.14)} vanish, because $\mu_A([a,b])=A(b)-A(a)=0$. We assume henceforth $A(a)<A(b)$ and put $J\defeq[A(a),A(b))$.

\emph{Step 1. The integrand on the right-hand side.} For $s\in J$ we have $s<A(b)\le S$, so $T(s)<b$ by (v); and $s\ge A(a)$, i.e.\ not $s<A(a)$, so $T(s)\ge a$ by (v). Hence $T$ maps $J$ into $[a,b)$, and $G(T(s))$ is defined for every $s\in J$. (At the single point $s=A(b)$ it need not be: $T(A(b))\ge b$ by (iii), with strict inequality possible, and $T(A(b))=\infty$ if $A(b)=S$. A single point is Lebesgue-null, so the right-hand side of \bk{(4.14)} is understood as $\int_JG(T(s))\,ds$.) Being nondecreasing, $T\colon J\to[a,b)$ is Borel measurable, and thus so is $G\circ T$ on $J$.

\emph{Step 2. $\mu_A$ restricted to $[a,b]$ is the image of Lebesgue measure on $J$ under $T$.} Define two finite measures on $\B([a,b])$ by
\[
  \mu(C)\defeq\mu_A(C),\qquad\nu(C)\defeq\meas\{s\in J;\ T(s)\in C\},
\]
where $\meas$ denotes Lebesgue measure. For $a\le t_1<t_2\le b$ and $s\in J$, part (v) gives $T(s)<t_2\Leftrightarrow s<A(t_2)$ and $T(s)\ge t_1\Leftrightarrow s\ge A(t_1)$; since $A(a)\le A(t_1)\le A(t_2)\le A(b)$, this yields
\[
  \nu([t_1,t_2))=\meas[A(t_1),A(t_2))=A(t_2)-A(t_1)=\mu([t_1,t_2)).
\]
The class of intervals $[t_1,t_2)$ with $a\le t_1<t_2\le b$, together with the empty set, is closed under finite intersections, contains $[a,b)$, and generates $\B([a,b))$. Let $\mathscr{D}$ be the collection of all $C\in\B([a,b))$ with $\mu(C)=\nu(C)$. It contains $[a,b)$; if $C,C'\in\mathscr{D}$ and $C'\subseteq C$, then $\mu(C\setminus C')=\mu(C)-\mu(C')=\nu(C)-\nu(C')=\nu(C\setminus C')$, the subtraction being legitimate because $\mu$ and $\nu$ are finite; and if $C_1\subseteq C_2\subseteq\cdots$ are in $\mathscr{D}$, then $\mu(\bigcup_nC_n)=\lim_n\mu(C_n)=\lim_n\nu(C_n)=\nu(\bigcup_nC_n)$ by continuity from below. Thus $\mathscr{D}$ is a Dynkin system of subsets of $[a,b)$ (\bk{Definition 2.1.2}) which contains the above class, and the \bk{Dynkin System Theorem 2.1.3} shows that $\mathscr{D}=\B([a,b))$; that is, $\mu=\nu$ on $\B([a,b))$. Moreover $\mu(\{b\})=0$ because $\mu_A$ has no atoms, and $\nu(\{b\})=0$ because $T(J)\subseteq[a,b)$ by Step~1. Therefore
\begin{equation}\label{eq:3-4-05:b}
  \int_a^b\ind_C(t)\,dA(t)=\mu(C)=\nu(C)=\int_J\ind_C(T(s))\,ds;\qquad C\in\B([a,b]).
\end{equation}

\emph{Step 3. General $G$.} By linearity, \eqref{eq:3-4-05:b} gives \bk{(4.14)} for every nonnegative simple Borel function $G=\sum_{i=1}^nc_i\ind_{C_i}$ on $[a,b]$. If $G\colon[a,b]\to[0,\infty]$ is Borel measurable, we choose nonnegative simple Borel functions $G_n\uparrow G$ pointwise on $[a,b]$; then $G_n\circ T\uparrow G\circ T$ pointwise on $J$, and the monotone convergence theorem, applied on both sides, yields
\[
  \int_a^bG(t)\,dA(t)=\lim_{n\to\infty}\int_a^bG_n(t)\,dA(t)=\lim_{n\to\infty}\int_JG_n(T(s))\,ds=\int_JG(T(s))\,ds ,
\]
both sides possibly being $+\infty$. Finally, if $G$ is bounded, real-valued and Borel measurable, say $\abs{G}\le c$, we apply the preceding identity to $G^+$ and $G^-$; the four integrals are bounded by $c\,(A(b)-A(a))<\infty$, so we may subtract, and \bk{(4.14)} follows for $G=G^+-G^-$.
\end{kex}

\begin{kex}{Problem}{4.16}{}
\emph{Setting and conventions.} $W=\{W_t,\F_t;\ 0\le t<\infty\}$ is a $d$-dimensional Brownian motion and $\{\F_t\}$ is the augmentation under $P$ of $\{\F^W_t\}$ (\bk{Definition 2.7.2}); we write $\cN$ for the collection of $P$-null sets used in this augmentation, so that $\cN\subseteq\F_0$. The filtration $\{\F_t\}$ is right-continuous by \bk{Proposition 2.7.7} (Brownian motion is a strong Markov process, \bk{Theorem 2.6.15}), hence it satisfies the usual conditions, and $W$ is a Brownian motion with respect to it (\bk{Theorem 2.7.9}); in particular $W^{(j)}\in\cM_2^c$ and $\qv{W^{(j)}}_t=t$. As everywhere in the book, ``$M$ is continuous'' is to be understood as ``$P$-almost every path of $M$ is continuous'', which is also the form of the conclusion of \bk{Theorem 4.15}; the condition displayed in the statement is likewise understood to hold a.s., and \bk{(4.39)} to hold a.s.\ for all $t$ simultaneously. Two facts will be used repeatedly.

\textbf{(F1)} \emph{If $X$ is an adapted RCLL process, the event $C_X\defeq\{\omega;\ t\mapsto X_t(\omega)\text{ is continuous on }[0,\infty)\}$ belongs to $\F_\infty$.} Indeed, an RCLL function is continuous on $[0,T]$ if and only if it is uniformly continuous on $\Q\cap[0,T]$ (necessity is the uniform continuity of a continuous function on a compact interval; sufficiency follows because, by right-continuity, the modulus of continuity on the rationals of $[0,T]$ is inherited by all pairs of points of $[0,T)$). Hence
\[
  C_X=\bigcap_{T\in\N}\ \bigcap_{k\ge1}\ \bigcup_{m\ge1}\
      \bigcap_{\substack{r,s\in\Q\cap[0,T]\\ \abs{r-s}<1/m}}\bigl\{\abs{X_r-X_s}\le 1/k\bigr\}\in\F_\infty .
\]

\textbf{(F2)} \emph{Every $F\in\F_\infty$ with $P(F)=0$ belongs to $\cN\subseteq\F_0$.} Indeed, by \bk{Problems 2.7.3} and \bk{2.7.5}, $\F_\infty=\sigma(\F^W_\infty\cup\cN)$ and there is $G\in\F^W_\infty$ with $F\triangle G\in\cN$; then $P(G)=0$, so $G\in\cN$ and $F\subseteq G\cup(F\triangle G)\in\cN$.

\medskip
\textbf{Part 1: an RCLL martingale $M$ is continuous.} It suffices to show, for each $T\in\N$, that a.e.\ path of $M$ is continuous on $[0,T]$. Fix such a $T$.

\emph{Step 1: bounded approximations of $M_T$.} Since $E\abs{M_T}<\infty$, dominated convergence gives $E\bigl[\abs{M_T}\ind_{\{\abs{M_T}>k\}}\bigr]\to0$ as $k\to\infty$. For each $n\ge1$ we choose $k_n$ with $E\bigl[\abs{M_T}\ind_{\{\abs{M_T}>k_n\}}\bigr]\le3^{-n}$ and set $\xi_n\defeq M_T\ind_{\{\abs{M_T}\le k_n\}}$ (this is the random variable called $M^{(n)}_T$ in the hint). It is $\F_T$-measurable, bounded by $k_n$, and
\begin{equation}\label{eq:3-4-16:a}
  E\abs{\xi_n-M_T}\le3^{-n}.
\end{equation}

\emph{Step 2: the martingales $M^{(n)}$ are continuous.} For every $t\in[0,\infty)$ let $X^{(n)}_t$ be a version of $E[\xi_n\mid\F_{t\wedge T}]$. Then $\{X^{(n)}_t,\F_t;\ 0\le t<\infty\}$ is a martingale (tower property), $t\mapsto EX^{(n)}_t$ is constant, and the filtration satisfies the usual conditions; by \bk{Theorem 1.3.13}, $X^{(n)}$ has a modification $M^{(n)}=\{M^{(n)}_t,\F_t;\ 0\le t<\infty\}$, every path of which is RCLL, and which is a martingale. Because $\abs{M^{(n)}_t}\le k_n$ a.s.\ for every $t$, this martingale is square-integrable, and
\begin{equation}\label{eq:3-4-16:b}
  M^{(n)}_T=X^{(n)}_T=E[\xi_n\mid\F_T]=\xi_n\quad\text{a.s.}
\end{equation}
The process $\wt M^{(n)}_t\defeq M^{(n)}_t-M^{(n)}_0$ is an RCLL, square-integrable martingale of $\{\F_t\}$ (the random variable $M^{(n)}_0$ is bounded and $\F_0$-measurable) with $\wt M^{(n)}_0=0$. \bk{Theorem 4.15} applies to it and shows that $\wt M^{(n)}$, and therefore $M^{(n)}$, is a.s.\ continuous.

\emph{Step 3: uniform convergence.} The process $\{M_t-M^{(n)}_t,\F_t;\ 0\le t<\infty\}$ is a martingale, so by \bk{Proposition 1.3.6} (with $\varphi(x)=\abs{x}$) the process $\abs{M-M^{(n)}}$ is a submartingale; all its paths are right-continuous. \bk{Theorem 1.3.8(i)}, applied on the interval $[\sigma,\tau]=[0,T]$ with $\lambda=2^{-n}$, together with \eqref{eq:3-4-16:b} and \eqref{eq:3-4-16:a}, yields
\[
  P\Bigl[\sup_{0\le t\le T}\abs[\big]{M_t-M^{(n)}_t}\ge2^{-n}\Bigr]
  \le2^{n}\,E\abs[\big]{M_T-M^{(n)}_T}\le\Bigl(\frac23\Bigr)^{n}.
\]
(The supremum is a random variable: by right-continuity it equals the supremum over $(\Q\cap[0,T])\cup\{T\}$.) The right-hand side is summable in $n$, so by the Borel--Cantelli lemma there is an event $\Omega_T$ with $P(\Omega_T)=1$ such that for every $\omega\in\Omega_T$ we have $\sup_{0\le t\le T}\abs{M_t(\omega)-M^{(n)}_t(\omega)}<2^{-n}$ for all sufficiently large $n$. Removing from $\Omega_T$ the countable union of the null events on which some $M^{(n)}$ fails to be continuous, we obtain an event of probability one on which $t\mapsto M_t(\omega)$ is, on $[0,T]$, the uniform limit of the continuous functions $t\mapsto M^{(n)}_t(\omega)$, hence continuous on $[0,T]$. Intersecting over $T\in\N$ we conclude that $P$-a.e.\ path of $M$ is continuous on $[0,\infty)$.

\rmk{The hypothesis $M_0=0$ was not used in Part~1: every RCLL martingale of the augmented Brownian filtration is a.s.\ continuous.}

\medskip
\textbf{Part 2: an RCLL local martingale $M$ with $M_0=0$ a.s.}
By \bk{Definition 1.5.15} there is a nondecreasing sequence $\{T_k\}_{k\ge1}$ of stopping times of $\{\F_t\}$ with $P[\lim_kT_k=\infty]=1$ such that $\{M_{t\wedge T_k},\F_t;\ 0\le t<\infty\}$ is a martingale for every $k$.

\rmk{The assumption $M_0=0$ of the first sentence of the problem is kept in force here; it is necessary for \bk{(4.39)}, whose right-hand side vanishes at $t=0$. It is not needed for the continuity of $M$, as Step~4 shows.}

\emph{Step 4: continuity.} For each $k$, the paths of $t\mapsto M_{t\wedge T_k}$ are RCLL, so by Part~1 (and the remark following it) there is an event $\Omega_k$ of probability one on which $t\mapsto M_{t\wedge T_k}$ is continuous. Let $\omega\in\bigcap_k\Omega_k\cap\{\lim_kT_k=\infty\}$ and $t\ge0$; choosing $k$ with $T_k(\omega)>t$, we see that $M_\cdot(\omega)$ coincides with the continuous function $M_{\cdot\wedge T_k(\omega)}(\omega)$ on the neighbourhood $[0,T_k(\omega))$ of $t$ in $[0,\infty)$, and is therefore continuous at $t$. Thus $M$ is a.s.\ continuous.

\emph{Step 5: a continuous version and a bounded localization.} Let $\Omega^*\defeq C_M\cap\{M_0=0\}$. By (F1), $\Omega\setminus\Omega^*$ is a $P$-null event in $\F_\infty$, hence it belongs to $\F_0$ by (F2). Consequently $\wt M_t\defeq M_t\ind_{\Omega^*}$ defines an \emph{adapted} process, all of whose paths are continuous, with $\wt M_0\equiv0$, and indistinguishable from $M$. For $n\ge1$ put
\[
  S_n\defeq\inf\{t\ge0;\ \abs{\wt M_t}\ge n\},\qquad S_0\defeq0 .
\]
By \bk{Problem 1.2.7}, each $S_n$ is a stopping time of $\{\F_t\}$. The sequence is nondecreasing, $\abs{\wt M_{t\wedge S_n}}\le n$ for all $t$ and $\omega$, and $\abs{\wt M_{S_n}}=n$ on $\{S_n<\infty\}$ by continuity; since a continuous path is bounded on every compact interval, $\lim_nS_n(\omega)=\infty$ for \emph{every} $\omega$.

We claim that $M^{(n)}\defeq\{\wt M_{t\wedge S_n},\F_t;\ 0\le t<\infty\}$ is a martingale. (The symbol $M^{(n)}$ is used from now on with this new meaning.) Fix $n$. For each $k$, \bk{Problem 1.3.24(i)}, applied to the right-continuous martingale $\{M_{t\wedge T_k},\F_t\}$ and to its negative with the stopping time $S_n$, shows that $\{M_{t\wedge T_k\wedge S_n},\F_t;\ 0\le t<\infty\}$ is a martingale; hence so is $\{\wt M_{t\wedge T_k\wedge S_n},\F_t\}$, the two processes being indistinguishable. For $0\le s<t$ and $A\in\F_s$ we thus have
\[
  E\bigl[\ind_A\wt M_{t\wedge T_k\wedge S_n}\bigr]=E\bigl[\ind_A\wt M_{s\wedge T_k\wedge S_n}\bigr];\qquad k\ge1 .
\]
The integrands are bounded by $n$ and, as $k\to\infty$, converge a.s.\ to $\ind_A\wt M_{t\wedge S_n}$ and $\ind_A\wt M_{s\wedge S_n}$, respectively (on $\{\lim_kT_k=\infty\}$ they are eventually equal to these). The bounded convergence theorem gives $E[\ind_A M^{(n)}_t]=E[\ind_AM^{(n)}_s]$, which proves the claim. Thus $M^{(n)}$ is a continuous martingale, bounded by $n$, with $M^{(n)}_0=0$; in particular it is a square-integrable martingale with RCLL paths.

\emph{Step 6: representation of $M^{(n)}$ and consistency.} By \bk{Theorem 4.15}, for each $n\ge1$ there are progressively measurable processes $Y^{(n,1)},\dots,Y^{(n,d)}$ such that
\begin{equation}\label{eq:3-4-16:c}
  E\int_0^T\bigl(Y^{(n,j)}_t\bigr)^2\,dt<\infty\quad(1\le j\le d,\ 0<T<\infty),
  \qquad
  M^{(n)}_t=\sum_{j=1}^d\int_0^tY^{(n,j)}_s\,dW^{(j)}_s\quad(0\le t<\infty),
\end{equation}
the second relation holding a.s. The first relation says that $Y^{(n,j)}\in\cL^*(W^{(j)})$.

Let $1\le n\le m$. Since $S_n\le S_m$, we have $M^{(m)}_{t\wedge S_n}=M^{(n)}_t$ for all $t$. The process $\ind_{\{t\le S_n\}}$ is adapted and left-continuous, hence progressively measurable (\bk{Proposition 1.1.13}), so $Y^{(m,j)}_t\ind_{\{t\le S_n\}}$ is progressively measurable and satisfies \bk{(4.38)}. By \bk{(2.24)} of \bk{Proposition 2.10} and \eqref{eq:3-4-16:c} for $m$, for each $t$ we have a.s.
\[
  M^{(n)}_t=M^{(m)}_{t\wedge S_n}=\sum_{j=1}^d\int_0^{t\wedge S_n}Y^{(m,j)}_s\,dW^{(j)}_s
  =\sum_{j=1}^d\int_0^tY^{(m,j)}_s\ind_{\{s\le S_n\}}\,dW^{(j)}_s ,
\]
and then a.s.\ for all $t$ simultaneously, both sides being continuous in $t$. Thus the processes $Y^{(m,j)}_t\ind_{\{t\le S_n\}}$, $1\le j\le d$, also satisfy \bk{(4.38)} and \bk{(4.39)} for the martingale $M^{(n)}$, and the uniqueness assertion of \bk{Theorem 4.15} yields
\begin{equation}\label{eq:3-4-16:d}
  \sum_{j=1}^d\int_0^\infty\abs[\big]{Y^{(n,j)}_t-Y^{(m,j)}_t\ind_{\{t\le S_n\}}}^2\,dt=0\quad\text{a.s.};\qquad 1\le n\le m .
\end{equation}
Let $\Omega^{**}$ be the event of probability one on which \eqref{eq:3-4-16:d} holds for all pairs $1\le n\le m$.

\emph{Step 7: the integrands $Y^{(j)}$.} For $1\le j\le d$ define
\[
  Y^{(j)}_t\defeq\sum_{n=1}^\infty Y^{(n,j)}_t\,\ind_{\{S_{n-1}<t\le S_n\}};\qquad 0\le t<\infty .
\]
For each $(t,\omega)$ at most one term of the series is different from zero. Each process $\ind_{\{S_{n-1}<t\le S_n\}}=\ind_{\{S_{n-1}<t\}}-\ind_{\{S_n<t\}}$ is adapted and left-continuous, hence progressively measurable; therefore the partial sums, and their pointwise limit $Y^{(j)}$, are progressively measurable.

We claim that for every $\omega\in\Omega^{**}$, every $n\ge1$ and every $j$,
\begin{equation}\label{eq:3-4-16:e}
  Y^{(j)}_t(\omega)\,\ind_{\{t\le S_n(\omega)\}}=Y^{(n,j)}_t(\omega)\quad\text{for Lebesgue-a.e.\ } t\in[0,\infty).
\end{equation}
Indeed, fix $\omega\in\Omega^{**}$ and $n$. By \eqref{eq:3-4-16:d} with $m=n$, $Y^{(n,j)}_t=Y^{(n,j)}_t\ind_{\{t\le S_n\}}=0$ for a.e.\ $t>S_n$, which is \eqref{eq:3-4-16:e} on $(S_n,\infty)$. The set $(0,S_n]\cap[0,\infty)$ is the disjoint union of the sets $(S_{k-1},S_k]\cap[0,\infty)$, $1\le k\le n$; on the $k$-th of them $Y^{(j)}_t=Y^{(k,j)}_t$ by definition, and \eqref{eq:3-4-16:d} for the pair $(k,n)$ gives $Y^{(k,j)}_t=Y^{(n,j)}_t\ind_{\{t\le S_k\}}=Y^{(n,j)}_t$ for a.e.\ $t$ in this set. Since $\{0\}$ is Lebesgue-null, \eqref{eq:3-4-16:e} follows.

From \eqref{eq:3-4-16:e} and \eqref{eq:3-4-16:c} we obtain, for every $n\ge1$ and $T\in\N$,
\[
  \int_0^{T\wedge S_n}\bigl(Y^{(j)}_t\bigr)^2\,dt=\int_0^T\bigl(Y^{(n,j)}_t\bigr)^2\,dt<\infty\quad\text{a.s.,}
\]
the finiteness because the right-hand side has finite expectation. Hence there is an event of probability one on which the left-hand side is finite for all $n\ge1$, $T\in\N$ and $1\le j\le d$. Since $S_n(\omega)\uparrow\infty$ for every $\omega$, for given $\omega$ and $T$ there is $n$ with $S_n(\omega)>T$, and we conclude that, a.s.,
\[
  \int_0^T\bigl(Y^{(j)}_t\bigr)^2\,dt<\infty;\qquad 1\le j\le d,\quad 0\le T<\infty .
\]
Thus $Y^{(j)}\in\cP^*(W^{(j)})$ (\bk{Definition 2.22}), and the stochastic integral $\int_0^tY^{(j)}_s\,dW^{(j)}_s$ is defined as an element of $\cM^{c,\loc}$ by \bk{Definition 2.23}.

\emph{Step 8: the representation \bk{(4.39)}.} Fix $n\ge1$ and $j$, and write $I(X)\defeq\int_0^\cdot X_s\,dW^{(j)}_s$. We show that, a.s.,
\begin{equation}\label{eq:3-4-16:f}
  I_{t\wedge S_n}\bigl(Y^{(j)}\bigr)=I_t\bigl(Y^{(n,j)}\bigr);\qquad 0\le t<\infty .
\end{equation}
Let us recall how the left-hand side is defined. Since $W^{(j)}$ is itself in $\cM_2^c$, in the construction preceding \bk{Definition 2.23} we may take the localizing times of $W^{(j)}$ to be identically $+\infty$ (the definition does not depend on this choice), so that the stopping times of \bk{(2.32)} are
\[
  \rho_m\defeq m\wedge\inf\Bigl\{0\le t<\infty;\ \int_0^t\bigl(Y^{(j)}_s\bigr)^2\,ds\ge m\Bigr\};\qquad m\ge1,
\]
which are nondecreasing in $m$ with $\lim_m\rho_m=\infty$ a.s., and the integrator and integrand of \bk{(2.33)} are
\[
  W^{(j,m)}_t\defeq W^{(j)}_{t\wedge\rho_m},\qquad Z^{(m)}_t\defeq Y^{(j)}_t\ind_{\{t\le\rho_m\}};\qquad 0\le t<\infty .
\]
The process $Z^{(m)}$ is progressively measurable (as in Step~6), and $W^{(j,m)}\in\cM_2^c$ with $\qv{W^{(j,m)}}_t=t\wedge\rho_m$: indeed, by \bk{Problem 1.3.24(i)} both $W^{(j)}_{t\wedge\rho_m}$ and $(W^{(j)}_{t\wedge\rho_m})^2-t\wedge\rho_m$ are continuous martingales, and the quadratic variation is unique (\bk{Definition 1.5.3}). By the definition of $\rho_m$,
\[
  \int_0^\infty\bigl(Z^{(m)}_s\bigr)^2\,d\qv{W^{(j,m)}}_s=\int_0^\infty\bigl(Z^{(m)}_s\bigr)^2\,ds
  =\int_0^{\rho_m}\bigl(Y^{(j)}_s\bigr)^2\,ds\le m,
\]
so $Z^{(m)}$ belongs both to $\cL^*(W^{(j,m)})$ and to $\cL^*(W^{(j)})$. By \bk{(2.34)}, a.s.,
\[
  I_u\bigl(Y^{(j)}\bigr)=I^{W^{(j,m)}}_u\bigl(Z^{(m)}\bigr)\quad\text{for all } m\ge1 \text{ and all } 0\le u\le\rho_m ,
\]
where the right-hand side is the integral of $Z^{(m)}$ with respect to the \emph{stopped} martingale $W^{(j,m)}$. We pass to the integrator $W^{(j)}$ by \bk{Corollary 2.21}, applied with $M=W^{(j)}$, $\wt M=W^{(j,m)}$, the same integrand $X=\wt X=Z^{(m)}$, and the stopping time $T=\rho_m$; its hypotheses hold because $W^{(j)}_{t\wedge\rho_m}=W^{(j,m)}_{t\wedge\rho_m}$ for all $t$ and $\omega$. It gives, a.s., $I^{W^{(j,m)}}_{t\wedge\rho_m}(Z^{(m)})=I_{t\wedge\rho_m}(Z^{(m)})$ for all $0\le t<\infty$. Consequently there is an event of probability one on which
\[
  I_u\bigl(Y^{(j)}\bigr)=I_u\bigl(Y^{(j)}\ind_{\{\cdot\le\rho_m\}}\bigr)\quad\text{for all } m\ge1 \text{ and all } 0\le u\le\rho_m .
\]
Fix $t\ge0$ and $m\ge1$. On the event $\{t\le\rho_m\}$ we may take $u=t\wedge S_n$ in this relation and obtain, a.s.\ on $\{t\le\rho_m\}$,
\begin{align*}
  I_{t\wedge S_n}\bigl(Y^{(j)}\bigr)
  &=I_{t\wedge S_n}\bigl(Y^{(j)}\ind_{\{\cdot\le\rho_m\}}\bigr)
   =I_t\bigl(Y^{(j)}\ind_{\{\cdot\le\rho_m\}}\ind_{\{\cdot\le S_n\}}\bigr)\\
  &=I_t\bigl(Y^{(n,j)}\ind_{\{\cdot\le\rho_m\}}\bigr)
   =I_{t\wedge\rho_m}\bigl(Y^{(n,j)}\bigr)=I_t\bigl(Y^{(n,j)}\bigr).
\end{align*}
Here all integrands except the first belong to $\cL^*(W^{(j)})$; the second and the fourth equalities are \bk{(2.24)} of \bk{Proposition 2.10}, applied with the stopping times $S_n$ and $\rho_m$, respectively; and the third holds because, by \eqref{eq:3-4-16:e}, for a.e.\ $\omega$ the two progressively measurable integrands agree for Lebesgue-a.e.\ $s$, so that by linearity and the isometry \bk{(2.14)} the difference of the two integrals has second moment zero. Letting $m\to\infty$ and using $P[\bigcup_m\{t\le\rho_m\}]=1$, we obtain \eqref{eq:3-4-16:f} a.s.\ for each fixed $t$, and then a.s.\ for all $t$ simultaneously, because both sides are continuous in $t$.

Summing \eqref{eq:3-4-16:f} over $j$ and using \eqref{eq:3-4-16:c}, we get, a.s.\ for all $t\ge0$ and all $n\ge1$,
\[
  \sum_{j=1}^d\int_0^{t\wedge S_n}Y^{(j)}_s\,dW^{(j)}_s=\sum_{j=1}^d\int_0^tY^{(n,j)}_s\,dW^{(j)}_s=M^{(n)}_t=\wt M_{t\wedge S_n}.
\]
On the event of probability one where this holds, given $t\ge0$ we choose $n$ with $S_n>t$ and obtain
\[
  \wt M_t=\sum_{j=1}^d\int_0^tY^{(j)}_s\,dW^{(j)}_s;\qquad 0\le t<\infty .
\]
Since $M$ and $\wt M$ are indistinguishable, this is \bk{(4.39)} for $M$, with progressively measurable integrands $Y^{(j)}$ satisfying $\int_0^T(Y^{(j)}_t)^2\,dt<\infty$ a.s.\ for all $T$ and $j$.
\end{kex}

\begin{kex}{Problem}{4.17}{}
Throughout, $W$ is a $d$-dimensional Brownian motion and $\{\F_t\}$ is the augmentation under $P$ of $\{\F^W_t\}$ (\bk{Definition 2.7.2}); we assume that $W$ starts at a fixed point $x\in\R^d$, i.e.\ $P[W_0=x]=1$ (the second remark at the end shows that this is needed). We apply \bk{Theorem 4.15} to the martingale generated by $\xi$.

\emph{Step 1. The filtration.}
The process $\{W_t,\F^W_t;\ 0\le t<\infty\}$ is a $d$-dimensional Brownian motion, hence strongly Markov (\bk{Theorem 2.6.15} and the paragraph following \bk{Remark 2.7.10}), so $\{\F_t\}$ is right-continuous by \bk{Proposition 2.7.7}; moreover, each $\F_t$ contains the collection $\cN$ of $P$-null sets by construction. Hence $\{\F_t\}$ satisfies the usual conditions. Furthermore, $\F_0$ is $P$-trivial: by \bk{Problem 2.7.3}, every $F\in\F_0=\sigma(\sigma(W_0)\cup\cN)$ differs by a null set from some $G=\{W_0\in\Gamma\}$, $\Gamma\in\B(\R^d)$, and $P(G)=\ind_\Gamma(x)\in\{0,1\}$. Consequently
\begin{equation}\label{eq:3-4-17:a}
  E(\xi\,|\,\F_0)=E(\xi)\quad\text{a.s.}
\end{equation}

\emph{Step 2. The martingale.}
For each $t\ge0$ let $X_t$ be a version of $E(\xi\,|\,\F_t)$. Then $\{X_t,\F_t;\ 0\le t<\infty\}$ is a martingale, and $t\mapsto EX_t=E\xi$ is constant. By \bk{Theorem 1.3.13}, whose hypothesis on the filtration was verified in Step 1, $X$ has a modification $N=\{N_t,\F_t;\ 0\le t<\infty\}$ which is adapted, has RCLL paths, and is a submartingale; since $EN_t=E\xi$ does not depend on $t$, the nonnegative random variables $E(N_t\,|\,\F_s)-N_s$, $s\le t$, have expectation zero, so $N$ is a martingale. We set
\[
  M_t\defeq N_t-N_0;\qquad 0\le t<\infty .
\]
Because $N_0$ is $\F_0$-measurable and integrable, $M=\{M_t,\F_t;\ 0\le t<\infty\}$ is a martingale with RCLL paths and $M_0=0$. It is square-integrable: by Jensen's inequality for conditional expectations, $N_t^2\le E(\xi^2\,|\,\F_t)$ a.s., whence $EN_t^2\le E\xi^2<\infty$ and $EM_t^2\le 2EN_t^2+2EN_0^2\le 4E\xi^2<\infty$ for every $t\ge0$.

\emph{Step 3. The representation.}
\bk{Theorem 4.15} provides progressively measurable processes $Y^{(j)}=\{Y^{(j)}_t,\F_t;\ 0\le t<\infty\}$, $1\le j\le d$, which satisfy \bk{(4.38)} for every finite horizon and \bk{(4.39)}; in particular, a.s.,
\[
  M_T=\sum_{j=1}^d\int_0^TY^{(j)}_s\,dW^{(j)}_s .
\]
Since $\xi$ is $\F_T$-measurable, $N_T=E(\xi\,|\,\F_T)=\xi$ a.s., and $N_0=E(\xi\,|\,\F_0)=E(\xi)$ a.s.\ by \eqref{eq:3-4-17:a}. Therefore
\[
  \xi=N_0+M_T=E(\xi)+\sum_{j=1}^d\int_0^TY^{(j)}_s\,dW^{(j)}_s\qquad\text{a.s. }P,
\]
which is \bk{(4.47)}.

\rmk{Only the values of $Y^{(j)}$ on $[0,T]$ enter \bk{(4.47)}. In fact $Y^{(j)}=0$ for Lebesgue-a.e.\ $t>T$, a.s.: for $t\ge T$ we have $M_t=M_T$ a.s., so by \bk{(4.39)} and the isometry of the stochastic integral, $\sum_{j=1}^dE\int_T^t(Y^{(j)}_s)^2\,ds=E(M_t-M_T)^2=0$.}

\rmk{\bk{Definition 2.5.1} allows a $d$-dimensional Brownian motion to have an arbitrary initial distribution $\mu$, and the proof of \bk{Theorem 4.15} does not use that $\mu$ is a point mass, because there $M_0=0$ is assumed. For \bk{(4.47)}, however, this matters. Steps~2 and~3 use Step~1 only through \eqref{eq:3-4-17:a}, so in general they prove
\[
  \xi=E(\xi\,|\,\F_0)+\sum_{j=1}^d\int_0^TY^{(j)}_s\,dW^{(j)}_s\qquad\text{a.s. }P
\]
(the integrals being taken with respect to the martingales $W^{(j)}-W^{(j)}_0$), and \bk{(4.47)} as printed holds precisely when $E(\xi\,|\,\F_0)$ is a.s.\ constant. It fails otherwise: if $\mu$ is not a point mass, choose $\Gamma\in\B(\R^d)$ with $0<\mu(\Gamma)<1$ and let $\xi=\ind_\Gamma(W_0)$, which is bounded and $\F_0$-measurable. If \bk{(4.47)} held with processes satisfying \bk{(4.38)}, the right-hand stochastic integrals would be the values at time $T$ of square-integrable martingales vanishing at $t=0$, and taking conditional expectations given $\F_0$ would give $\xi=E(\xi)=\mu(\Gamma)$ a.s., which is false. Thus the statement tacitly assumes that $W$ starts at a fixed point (for instance, that $W$ is standard), as we did.}
\end{kex}

\begin{kex}{Exercise}{4.22}{}
Parts (i) and (ii) rest on the following consequence of \bk{Problem 1.5.24}.

\begin{lem}
Let $N=\{N_t,\F_t;\ 0\le t<\infty\}$ be in $\cM^{c,\loc}$ with $E\qv{N}_\infty<\infty$. Then $N_\infty\defeq\lim_{t\to\infty}N_t$ exists a.s. If $N_\infty=0$ a.s., then $\qv{N}_\infty=0$ a.s.
\end{lem}

\begin{proof}
By \bk{Problem 1.5.24}\,(i), the limit $N_\infty$ exists a.s.\ and $E N_\infty^2=E\qv{N}_\infty$. If $N_\infty=0$ a.s., the nonnegative random variable $\qv{N}_\infty$ has expectation zero, hence vanishes a.s.
\end{proof}

We recall from the beginning of \bk{Subsection 3.4.E} that for an integrand $Z$ with $E\int_0^\infty Z_t^2\,d\qv{M}_t<\infty$ the symbol $\int_0^\infty Z_t\,dM_t$ denotes the a.s.\ limit, as $t\to\infty$, of $I^M_t(Z)=\int_0^tZ_s\,dM_s$; this limit exists by the Lemma, applied to $N=I^M(Z)$, for which $\qv{N}_t=\int_0^tZ_s^2\,d\qv{M}_s$.

\medskip
\textbf{(i)} Here $W$ is a $d$-dimensional Brownian motion and $\{\F_t\}$ is the augmentation of $\{\F^W_t\}$, which satisfies the usual conditions. By \bk{Problem 2.5.5}, the processes $M^{(j)}\defeq W^{(j)}-W^{(j)}_0$, $1\le j\le d$, with respect to which the stochastic integrals are taken, belong to $\cM_2^c$ and
\begin{equation}\label{eq:3-4-22:a}
  \qv{M^{(j)},M^{(k)}}_t=\delta_{jk}\,t;\qquad 0\le t<\infty,\ 1\le j,k\le d .
\end{equation}
Let $Z^{(j)}\defeq Y^{(j)}-\wt Y^{(j)}$. This process is progressively measurable, and $(Z^{(j)}_t)^2\le2(Y^{(j)}_t)^2+2(\wt Y^{(j)}_t)^2$ shows $E\int_0^\infty(Z^{(j)}_t)^2\,dt<\infty$; thus $Z^{(j)}\in\cL^*_\infty(W)\subseteq\cL^*(M^{(j)})$. We set
\[
  N_t\defeq\sum_{j=1}^d\int_0^tZ^{(j)}_s\,dW^{(j)}_s;\qquad 0\le t<\infty .
\]
Then $N\in\cM_2^c$, and by the bilinearity of the cross-variation, the formula \bk{(2.26)} (\bk{Proposition 2.17}) and \eqref{eq:3-4-22:a},
\[
  \qv{N}_t=\sum_{j=1}^d\sum_{k=1}^d\int_0^tZ^{(j)}_sZ^{(k)}_s\,d\qv{M^{(j)},M^{(k)}}_s
          =\int_0^t\sum_{j=1}^d\bigl(Z^{(j)}_s\bigr)^2\,ds;\qquad 0\le t<\infty,\ \text{a.s.}
\]
In particular, $E\qv{N}_\infty=\sum_{j=1}^dE\int_0^\infty(Z^{(j)}_s)^2\,ds<\infty$ by monotone convergence. By the linearity \bk{(2.17)} of the integral, a.s.\ for all $t$,
\[
  N_t=\sum_{j=1}^d\int_0^tY^{(j)}_s\,dW^{(j)}_s-\sum_{j=1}^d\int_0^t\wt Y^{(j)}_s\,dW^{(j)}_s ,
\]
and letting $t\to\infty$ we obtain from the two representations of $\xi$ that
\[
  N_\infty=\sum_{j=1}^d\int_0^\infty Y^{(j)}_s\,dW^{(j)}_s-\sum_{j=1}^d\int_0^\infty\wt Y^{(j)}_s\,dW^{(j)}_s
          =\bigl(\xi-E(\xi)\bigr)-\bigl(\xi-E(\xi)\bigr)=0\qquad\text{a.s.}
\]
The Lemma gives $\qv{N}_\infty=\int_0^\infty\sum_{j=1}^d\abs{Y^{(j)}_t-\wt Y^{(j)}_t}^2\,dt=0$ a.s.

\medskip
\textbf{(ii)} Now $M\in\cM^{c,\loc}$, and $\cL^*_\infty(M)$ is the class of $\{\F_t\}$-progressively measurable processes $Z$ with $E\int_0^\infty Z_t^2\,d\qv{M}_t<\infty$, as in \bk{Proposition 4.19}. Let $Z\defeq X-\wt X$. As in (i), $Z$ is progressively measurable and $Z_t^2\le2X_t^2+2\wt X_t^2$, so that $E\int_0^\infty Z_t^2\,d\qv{M}_t<\infty$; in particular $\int_0^TZ_t^2\,d\qv{M}_t<\infty$ a.s.\ for every $T$, i.e., $X,\wt X,Z\in\cP^*(M)$ (\bk{Definition 2.22}). Hence $N\defeq I^M(Z)$ is defined and belongs to $\cM^{c,\loc}$ (\bk{Definition 2.23}), and the sample-path properties \bk{(2.17)} and \bk{(2.19)}, which remain valid for integrands in $\cP^*$ (the paragraph following \bk{Definition 2.23}), give, a.s.\ for all $t$,
\[
  N_t=\int_0^tX_s\,dM_s-\int_0^t\wt X_s\,dM_s,\qquad \qv{N}_t=\int_0^tZ_s^2\,d\qv{M}_s .
\]
Thus $E\qv{N}_\infty<\infty$, and, letting $t\to\infty$,
\[
  N_\infty=\int_0^\infty X_s\,dM_s-\int_0^\infty\wt X_s\,dM_s=\bigl(\xi-E(\xi)\bigr)-\bigl(\xi-E(\xi)\bigr)=0\qquad\text{a.s.}
\]
The Lemma yields $\qv{N}_\infty=\int_0^\infty\abs{X_t-\wt X_t}^2\,d\qv{M}_t=0$ a.s.

\medskip
\textbf{(iii)} Let $W=\{W_t,\F_t;\ 0\le t<\infty\}$ be a standard, one-dimensional Brownian motion, where $\{\F_t\}$ satisfies the usual conditions (as in \bk{Theorem 4.20}). The idea is to integrate the constant $1$ up to time $\tfrac12$, and then to integrate $(1-t)^{-1}$, which makes the integral oscillate so strongly as $t\uparrow1$ that it returns to zero before time $1$; at that moment we switch the integrand off.

\emph{Step 1. An auxiliary martingale on $[0,1)$.} For integers $n\ge3$ let $t_n\defeq1-\frac1n$ and
\[
  g_n(s)\defeq\frac{1}{1-s}\,\ind_{(1/2,\,t_n]}(s);\qquad 0\le s<\infty .
\]
Each $g_n$ is a Borel function bounded by $n$, vanishing off $(\tfrac12,t_n]$; regarded as a (nonrandom) process it is progressively measurable and belongs to $\cL^*(W)$. Let $R^{(n)}\defeq I(g_n)\in\cM_2^c$, and choose versions all of whose paths are continuous. For $m\ge n$ we have $g_n=g_m\ind_{[0,t_n]}$, so \bk{(2.24)} (with the constant stopping time $t_n$) gives $R^{(n)}_t=R^{(m)}_{t\wedge t_n}$ a.s.\ for every $t$; furthermore $E\bigl(R^{(n)}_{1/2}\bigr)^2=\int_0^{1/2}g_n^2(s)\,ds=0$ by \bk{(2.16)}, hence $R^{(n)}_t=E\bigl[R^{(n)}_{1/2}\,\big|\,\F_t\bigr]=0$ a.s.\ for $t\le\tfrac12$. By the continuity of paths there is an event $\Omega_1$ with $P(\Omega_1)=1$ on which
\begin{equation}\label{eq:3-4-22:b}
  R^{(n)}_t=R^{(m)}_{t\wedge t_n}\quad\text{and}\quad R^{(n)}_s=0;\qquad 0\le t<\infty,\ 0\le s\le\tfrac12,\ m\ge n\ge3 .
\end{equation}
Since $\F_0$ contains all $P$-null sets, $\Omega_1\in\F_0$. We define
\[
  N_t\defeq R^{(n)}_t\,\ind_{\Omega_1}\quad\text{for }0\le t\le t_n,\ n\ge3 ;
\]
by \eqref{eq:3-4-22:b} this does not depend on the choice of $n$ with $t_n\ge t$, and it defines a process $\{N_t;\ 0\le t<1\}$ which is adapted to $\{\F_t\}$, has continuous paths, and vanishes on $[0,\tfrac12]$. Moreover, by \eqref{eq:3-4-22:b} with $m=n$,
\begin{equation}\label{eq:3-4-22:c}
  R^{(n)}_t=N_{t\wedge t_n}\quad\text{on }\Omega_1;\qquad 0\le t<\infty,\ n\ge3 .
\end{equation}

\emph{Step 2. $N$ is a time-changed Brownian motion.} Let
\[
  t(u)\defeq1-\frac{1}{u+2},\qquad B_u\defeq N_{t(u)},\qquad \cH_u\defeq\F_{t(u)};\qquad 0\le u<\infty .
\]
The function $t(\cdot)$ is continuous and strictly increasing from $t(0)=\tfrac12$ to $1$, so $\{\cH_u\}$ is a filtration satisfying the usual conditions, and $B$ is a continuous, $\{\cH_u\}$-adapted process with $B_0=0$. Let $0\le u_1<u_2$ and choose $n\ge3$ with $t_n\ge t(u_2)$. Then $B_{u_i}=R^{(n)}_{t(u_i)}$ a.s., and since $R^{(n)}\in\cM_2^c$, the martingale property and \bk{(2.16)} give, a.s.,
\begin{align*}
  E\bigl[B_{u_2}-B_{u_1}\,\big|\,\cH_{u_1}\bigr]
    &=E\bigl[R^{(n)}_{t(u_2)}-R^{(n)}_{t(u_1)}\,\big|\,\F_{t(u_1)}\bigr]=0,\\
  E\bigl[(B_{u_2}-B_{u_1})^2\,\big|\,\cH_{u_1}\bigr]
    &=\int_{t(u_1)}^{t(u_2)}\frac{ds}{(1-s)^2}=\frac{1}{1-t(u_2)}-\frac{1}{1-t(u_1)}=u_2-u_1 .
\end{align*}
Consequently $\{B_u,\cH_u;\ 0\le u<\infty\}$ is a continuous, square-integrable martingale, and $E[B_{u_2}^2-B_{u_1}^2\,|\,\cH_{u_1}]=E[(B_{u_2}-B_{u_1})^2\,|\,\cH_{u_1}]=u_2-u_1$ shows that $\{B_u^2-u\}$ is a martingale as well, i.e., $\qv{B}_u=u$. By P.~L\'evy's \bk{Theorem 3.16}, $B$ is a standard Brownian motion. The law of the iterated logarithm (\bk{Theorem 2.9.23}\,(iii), (iv)) provides an event $\Omega_2\subseteq\Omega_1$ with $P(\Omega_2)=1$ such that
\begin{equation}\label{eq:3-4-22:d}
  \varlimsup_{u\to\infty}B_u=+\infty,\qquad \varliminf_{u\to\infty}B_u=-\infty\qquad\text{on }\Omega_2 .
\end{equation}

\emph{Step 3. The stopping time.} Let
\[
  \tau\defeq\inf\bigl\{t\in[\tfrac12,1);\ W_{1/2}+N_t=0\bigr\},\qquad\text{with }\inf\emptyset\defeq1 ,
\]
so that $\tfrac12\le\tau\le1$ everywhere. For $\omega\in\Omega_2$, the continuous function $u\mapsto W_{1/2}(\omega)+B_u(\omega)$ takes both positive and negative values by \eqref{eq:3-4-22:d}, hence vanishes at some $u\ge0$ by the intermediate value theorem; therefore
\begin{equation}\label{eq:3-4-22:e}
  \tau<1\quad\text{and}\quad W_{1/2}+N_\tau=0\qquad\text{on }\Omega_2 ,
\end{equation}
the second assertion because the zero set of the continuous function $t\mapsto W_{1/2}+N_t$ on $[\tfrac12,1)$ is closed in $[\tfrac12,1)$. Furthermore, $\tau$ is a stopping time of $\{\F_t\}$: we have $\{\tau\le t\}=\emptyset$ for $t<\tfrac12$, $\{\tau\le t\}=\Omega$ for $t\ge1$, and, for $\tfrac12\le t<1$, by continuity and compactness,
\[
  \{\tau\le t\}=\Bigl\{\min_{1/2\le s\le t}\abs{W_{1/2}+N_s}=0\Bigr\}
     =\bigcap_{k=1}^\infty\ \bigcup_{q\in\Q\cap[1/2,\,t]}\bigl\{\abs{W_{1/2}+N_q}<\tfrac1k\bigr\}\in\F_t .
\]

\emph{Step 4. The process $Y$.} We define
\[
  Y_t\defeq\ind_{[0,1/2]}(t)+\frac{1}{1-t}\,\ind_{(1/2,\,1)}(t)\,\ind_{\{t\le\tau\}};\qquad 0\le t<\infty .
\]
The process $\{\ind_{\{t\le\tau\}};\ 0\le t<\infty\}$ is adapted, because $\tau$ is a stopping time, and has left-continuous paths, so it is progressively measurable by \bk{Proposition 1.1.13}; multiplying by a Borel function of $t$ and adding one preserves this property, so $Y$ is progressively measurable. On $\Omega_2$ we have $\tau<1$ and
\[
  \int_0^1Y_t^2\,dt=\frac12+\int_{1/2}^{\tau}\frac{dt}{(1-t)^2}=\frac{1}{1-\tau}-\frac32\in\bigl[\tfrac12,\infty\bigr) .
\]
Thus $0<\int_0^1Y_t^2\,dt<\infty$ a.s. Since $Y$ vanishes on $[1,\infty)$, we also have $\int_0^TY_t^2\,dt<\infty$ a.s.\ for every $T$, i.e., $Y\in\cP^*$ and the continuous local martingale $I(Y)=\{\int_0^tY_s\,dW_s;\ 0\le t<\infty\}$ is defined.

\emph{Step 5. $\int_0^1Y_t\,dW_t=0$ a.s.} Fix $n\ge3$ and let $G_n\defeq\ind_{[0,1/2]}+g_n$, a bounded Borel function. Since $\ind_{[0,1/2]}=\ind_{\{0\}}+\ind_{(0,1/2]}$ is a simple process, whose integral is $W_{t\wedge1/2}-W_0=W_{t\wedge1/2}$, the linearity \bk{(2.17)} gives
\begin{equation}\label{eq:3-4-22:f}
  I_t(G_n)=W_{t\wedge1/2}+R^{(n)}_t;\qquad 0\le t<\infty,\ \text{a.s.}
\end{equation}
Next, we claim that
\begin{equation}\label{eq:3-4-22:g}
  Y_t\,\ind_{\{t\le t_n\}}=G_n(t)\,\ind_{\{t\le\tau\}};\qquad 0\le t<\infty .
\end{equation}
Indeed, for $t\le\tfrac12$ both sides equal $1$ because $\tau\ge\tfrac12$; for $\tfrac12<t\le t_n$ both sides equal $(1-t)^{-1}\ind_{\{t\le\tau\}}$; for $t>t_n$ both sides vanish. We now apply \bk{(2.24)}, which holds for integrands in $\cP^*$ (the paragraph following \bk{Definition 2.23}), twice: to $Y$ with the constant stopping time $t_n$, and to $G_n$ with the stopping time $\tau$. Together with \eqref{eq:3-4-22:g}, \eqref{eq:3-4-22:f}, $\tfrac12\le\tau\le1$ and \eqref{eq:3-4-22:c}, this yields, a.s.,
\[
  I_{t_n}(Y)=I_1\bigl(Y\ind_{[0,t_n]}\bigr)=I_1\bigl(G_n\ind_{[0,\tau]}\bigr)=I_{1\wedge\tau}(G_n)
     =W_{1/2}+R^{(n)}_\tau=W_{1/2}+N_{\tau\wedge t_n} .
\]
Let $\Omega_3\subseteq\Omega_2$ be an event of probability one on which this holds for every $n\ge3$ and $t\mapsto I_t(Y)$ is continuous. For $\omega\in\Omega_3$ we have $\tau(\omega)<1$, hence $\tau(\omega)\wedge t_n=\tau(\omega)$ for all sufficiently large $n$, and for these $n$ \eqref{eq:3-4-22:e} gives $I_{t_n}(Y)(\omega)=W_{1/2}(\omega)+N_\tau(\omega)=0$. Letting $n\to\infty$ and using the continuity of $I(Y)$ at $t=1$, we conclude
\[
  \int_0^1Y_t\,dW_t=I_1(Y)=\lim_{n\to\infty}I_{t_n}(Y)=0\qquad\text{a.s.}
\]

\emph{Non-uniqueness in \bk{Theorem 4.20}.} Take $T=1$ there, and let $\xi$ and a representing process $Y^0$ be as in that theorem: $\int_0^1(Y^0_t)^2\,dt<\infty$ and $\xi=\int_0^1Y^0_t\,dW_t$ a.s. Then $Y^0+Y$ is progressively measurable, $\int_0^1(Y^0_t+Y_t)^2\,dt<\infty$ a.s., and $\int_0^1(Y^0_t+Y_t)\,dW_t=\xi$ a.s.\ by linearity, although $\int_0^1\abs{(Y^0_t+Y_t)-Y^0_t}^2\,dt=\int_0^1Y_t^2\,dt>0$ a.s. Hence the representing integrand is not unique, not even up to sets of $dt\times dP$-measure zero.

\rmk{There is no conflict with (i): necessarily $E\int_0^1Y_t^2\,dt=\infty$, for otherwise $I(Y)$ would be a continuous local martingale with $E\qv{I(Y)}_\infty=E\int_0^1Y_t^2\,dt<\infty$ and $\lim_{t\to\infty}I_t(Y)=I_1(Y)=0$ a.s.\ (recall that $Y$ vanishes on $[1,\infty)$), and the Lemma would force $\int_0^1Y_t^2\,dt=0$ a.s. This is why no such example exists within the class $\cL^*_\infty(W)$ of \bk{Proposition 4.18}.}
\end{kex}

\begin{kex}{Exercise}{4.23}{}
The assertion is false. We exhibit two continuous martingales $M$, $N$ on the same filtered probability space, both starting at zero, whose quadratic variation processes coincide, but for which the random variables $M_t$ and $N_t$ fail to have the same distribution for some $t\ge0$; a fortiori the laws of $M$ and $N$, i.e.\ the measures they induce on $\bigl(C[0,\infty),\B(C[0,\infty))\bigr)$, are different.

\emph{Step 1. The example.}
Let $W=\{W_t,\F_t^W;\ 0\le t<\infty\}$ be a standard, one-dimensional Brownian motion (\bk{Definition 2.1.1}) on a probability space $(\Omega,\F_\infty^W,P)$; thus every path of $W$ is continuous, and $W_0=0$ a.s. Let $\mathscr{N}$ be the collection of subsets of $P$-null sets of $\F_\infty^W$, and let
\[
  \F_t\defeq\sigma\bigl(\F_t^W\cup\mathscr{N}\bigr),\quad 0\le t<\infty,\qquad \F\defeq\sigma\bigl(\F_\infty^W\cup\mathscr{N}\bigr),
\]
be the augmented filtration of \bk{Definition 2.7.2}, with $P$ extended to $\F$ as in \bk{Problem 2.7.3}, so that $(\Omega,\F,P)$ is complete. Every $P$-negligible event of $\F$ belongs to $\mathscr{N}\subseteq\F_0$, and $\{\F_t\}$ is right-continuous by \bk{Proposition 2.7.7}, because $W$ is a strong Markov process (\bk{Theorem 2.6.15}). Hence $\{\F_t\}$ satisfies the usual conditions (\bk{Definition 1.2.25}), which are in force throughout Section~1.5 and are thus part of the definition of $\cM^{c,\loc}$ and of $\qv{M}$; and $\{W_t,\F_t;\ 0\le t<\infty\}$ is still a standard Brownian motion (\bk{Theorem 2.7.9}). From now on we work with this filtration. We set
\[
  T\defeq\inf\{t\ge0;\ W_t=1\},\qquad
  M_t\defeq W_{t\wedge T},\qquad N_t\defeq-M_t;\qquad 0\le t<\infty,
\]
with $\inf\emptyset=\infty$. Since $\{1\}$ is closed and $W$ has continuous paths, $T$ is a stopping time of $\{\F_t\}$ by \bk{Problem 1.2.7}. By \bk{Problem 1.3.24}\,(i), applied to the continuous martingales $W$ and $-W$, the process $M$ is a martingale, and so is $N$. Both are continuous and adapted, and $M_0=N_0=0$ a.s., so $M$ and $N$ belong to $\cM^{c,\loc}$ (\bk{Remark 1.5.16}).

\emph{Step 2. The quadratic variations coincide.}
The process $\{W_t^2-t,\F_t;\ 0\le t<\infty\}$ is a continuous martingale: it is integrable and adapted, and for $0\le s<t$ the increment $W_t-W_s$ is independent of $\F_s$ with mean zero and variance $t-s$, so that
\[
  E\bigl[W_t^2-W_s^2\,\big|\,\F_s\bigr]
  =E\bigl[(W_t-W_s)^2\,\big|\,\F_s\bigr]+2W_s\,E\bigl[W_t-W_s\,\big|\,\F_s\bigr]=t-s\quad\text{a.s.}
\]
A second application of \bk{Problem 1.3.24}\,(i) shows that
\[
  \bigl\{M_t^2-(t\wedge T),\ \F_t;\ 0\le t<\infty\bigr\}
\]
is a continuous martingale; it vanishes a.s.\ at $t=0$, and thus belongs to $\cM^{c,\loc}$. The process $A_t\defeq t\wedge T$ is continuous, nondecreasing, null at $t=0$, and adapted because $T$ is a stopping time. By the uniqueness assertion of \bk{Problem 1.5.17}, $\qv{M}_t=t\wedge T$ for all $t\ge0$, a.s. Since $N^2=M^2$, the same argument yields $\qv{N}_t=t\wedge T$ for all $t\ge0$, a.s. Hence $\qv{M}$ and $\qv{N}$ are indistinguishable; indeed, the same process $A$ serves as the quadratic variation of both.

\emph{Step 3. The laws differ.}
Suppose, to the contrary, that $M_t$ and $N_t$ have the same distribution for every $t\ge0$. On the event $\{W_0=0\}$, which has probability one, the intermediate value theorem and the definition of $T$ give $W_s<1$ for $0\le s<T$, and the continuity of $W$ gives $W_T=1$ on $\{T<\infty\}$; consequently, on $\{W_0=0\}$,
\begin{equation}\label{eq:3-4-23:a}
  M_t\le1\quad\text{for all }t\ge0,\qquad\text{and}\qquad M_t=1\quad\text{for }t\ge T\text{ if }T<\infty.
\end{equation}
Therefore $P[M_t\le1]=P[N_t\ge-1]=1$, and by our assumption $P[M_t\ge-1]=P[N_t\ge-1]=1$. Together with \eqref{eq:3-4-23:a}, this gives
\begin{equation}\label{eq:3-4-23:b}
  \abs{M_t}\le1\quad\text{a.s., for every }t\ge0.
\end{equation}
The martingale of Step~2 has expectation zero, so \eqref{eq:3-4-23:b} implies $E(t\wedge T)=E(M_t^2)\le1$ for every $t\ge0$. Letting $t\to\infty$, we obtain $E(T)\le1$ from the monotone convergence theorem; in particular, $P[T<\infty]=1$. Now \eqref{eq:3-4-23:a} shows that $\lim_{n\to\infty}M_n=1$ a.s., and because of \eqref{eq:3-4-23:b} the bounded convergence theorem yields
\[
  1=\lim_{n\to\infty}E(M_n)=E(M_0)=0,
\]
where we have used the martingale property of $M$. This contradiction shows that there exists a $t>0$ for which $M_t$ and $N_t=-M_t$ have different distributions. The processes $M$ and $N$ of Step~1 thus satisfy all the hypotheses of the assertion, but do not have the same law.

\rmk{The requirement $M_0=N_0=0$ a.s.\ built into the definition of $\cM^{c,\loc}$ (\bk{Definition 1.5.15}) rules out the trivial counterexample $N=M+c$, which is why a random quadratic variation was needed above. When the common quadratic variation is the deterministic function $\qv{M}_t=\qv{N}_t=t$, the assertion is true: by P.~L\'evy's \bk{Theorem 3.16}, both $M$ and $N$ are then standard Brownian motions, and so both induce Wiener measure. In general, however, the Brownian motion $B$ in the time-change representation $M_t=B_{\qv{M}_t}$ (\bk{Theorem 4.6} when $\qv{M}_\infty=\infty$ a.s.; \bk{Problem 4.7}, on a possibly extended probability space, when $P[\qv{M}_\infty<\infty]>0$, as is the case in our example, where $\qv{M}_\infty=T$ and $P[T<\infty]\ge P[W_1\ge1]>0$) depends on $M$, and the joint law of the pair $(B,\qv{M})$ is not determined by $\qv{M}$ alone. In the example, $M$ and $N$ are obtained by running the Brownian motions $W$ and $-W$, respectively, with the same clock $t\wedge T$; this clock is the first passage time of $W$ to the level $1$, but that of $-W$ to the level $-1$.}
\end{kex}

\begin{kex}{Exercise}{4.24}{}
\emph{Setting.} As everywhere in this chapter, $X$ and $Y$ are continuous semimartingales (\bk{Definition 3.1}) on a probability space $(\Omega,\F,P)$ with a filtration $\{\F_t\}$ which satisfies the usual conditions; $W$ and $V$ are standard, one-dimensional Brownian motions with respect to $\{\F_t\}$ (no assumption is made on their joint law), and $X_0$, $Y_0$ are $\F_0$-measurable. We read the hypotheses as follows: for every $t\ge0$ we have $\mu_t\le m(t)$ and $\abs{\sigma_t}\le\rho(t)$ a.s., and
\begin{equation}\label{eq:3-4-24:hyp}
  \int_0^t\abs{\mu_s}\,ds<\infty\ \ \text{a.s.},\qquad \int_0^t\abs{m(s)}\,ds<\infty,\qquad R(t)\defeq\int_0^t\rho^2(s)\,ds<\infty;\qquad 0\le t<\infty
\end{equation}
(see the first remark at the end). We set $a(t)\defeq\int_0^tm(s)\,ds$; the functions $a$ and $R$ are continuous, and $R$ is nondecreasing with $R(0)=0$. The set $\Lambda\defeq\{(s,\omega);\ \mu_s(\omega)>m(s)\text{ or }\abs{\sigma_s(\omega)}>\rho(s)\}$ belongs to $\B([0,\infty))\otimes\F$ and each of its $s$-sections is $P$-null; by the Fubini--Tonelli theorem, for $P$-a.e.\ $\omega$ the section $\{s;\ (s,\omega)\in\Lambda\}$ has Lebesgue measure zero. Consequently
\begin{equation}\label{eq:3-4-24:ae}
  \mu_s\le m(s)\quad\text{and}\quad\abs{\sigma_s}\le\rho(s)\quad\text{for Lebesgue-a.e. }s\ge0,\ \text{a.s.}
\end{equation}
In particular $\int_0^t\sigma_s^2\,ds\le R(t)<\infty$ a.s., so $\sigma\in\cP^*(W)$, and the deterministic function $\rho$ belongs to $\cP^*(V)$; both stochastic integrals in the statement are defined, in the sense of \bk{Definition 2.23}.

\emph{Step 1. The law of a local martingale with deterministic quadratic variation.}

\begin{lem}[1]
Let $M\in\cM^{c,\loc}$ on a filtered probability space $(\Omega',\cH,Q)$, $\{\cH_t\}$ satisfying the usual conditions, and suppose that $\qv{M}_t=R(t)$; $0\le t<\infty$ holds a.s., where $R\colon[0,\infty)\to[0,\infty)$ is a deterministic, continuous, nondecreasing function with $R(0)=0$. Then, for every $t\ge0$ and every $\cH_0$-measurable random variable $\xi$, the random variables $\xi$ and $M_t$ are independent, and $M_t$ is normal with mean zero and variance $R(t)$ \textup{(}equal to zero a.s.\ if $R(t)=0$\textup{)}.
\end{lem}

\begin{proof}
Fix $u\in\R$ and $t\ge0$. We apply It\^o's rule in the form of \bk{Theorem 3.6} with $d=2$ to the vector of continuous semimartingales $(M_s,R(s))$ -- its first component is a local martingale, its second a continuous, adapted (indeed deterministic) process of bounded variation -- and to the functions $g(x_1,x_2)=\cos(ux_1)e^{u^2x_2/2}$ and $h(x_1,x_2)=\sin(ux_1)e^{u^2x_2/2}$, which are of class $C^2(\R^2)$. Since $\frac{\partial g}{\partial x_2}=\frac{u^2}{2}g$, $\frac{\partial^2 g}{\partial x_1^2}=-u^2g$ and $d\qv{M}_s=dR(s)$, the terms involving $dR$ cancel, and we obtain a.s.
\[
  U_s\defeq g(M_s,R(s))=1+\int_0^s\frac{\partial g}{\partial x_1}(M_r,R(r))\,dM_r;\qquad 0\le s<\infty.
\]
Thus $U$ is a continuous local martingale, and so is $U'_s\defeq h(M_s,R(s))$, by the same computation (with $U_0'=0$). Let $\{T_n\}$ be a sequence of stopping times with $T_n\uparrow\infty$ a.s., such that the stopped processes $U_{\cdot\wedge T_n}$ are martingales. Then $E[U_{t\wedge T_n}\,|\,\cH_0]=U_0=1$ a.s., and since $\abs{U_s}\le e^{u^2R(t)/2}$ for $0\le s\le t$, the bounded convergence theorem for conditional expectations gives $E[U_t\,|\,\cH_0]=1$ a.s.; similarly $E[U'_t\,|\,\cH_0]=0$ a.s. Because $U_t+iU_t'=\exp\{iuM_t+u^2R(t)/2\}$, this means
\[
  E\bigl[e^{iuM_t}\,\big|\,\cH_0\bigr]=e^{-u^2R(t)/2}\quad\text{a.s.}
\]
For $v\in\R$ we multiply by the bounded, $\cH_0$-measurable random variable $e^{iv\xi}$ and take expectations:
\[
  E\bigl[e^{iv\xi+iuM_t}\bigr]=E\bigl[e^{iv\xi}\bigr]\,e^{-u^2R(t)/2};\qquad (v,u)\in\R^2.
\]
The right-hand side is the characteristic function of the product of the law of $\xi$ and the normal law with mean zero and variance $R(t)$. By the uniqueness theorem for characteristic functions on $\R^2$, this product measure is the law of $(\xi,M_t)$.
\end{proof}

Consequently, in the situation of the Lemma, for every $b\in\R$ the law of $\xi+b+M_t$ is the convolution of the law of $\xi$ with the normal law of mean $b$ and variance $R(t)$; \emph{it depends only on the law of $\xi$ and on the numbers $b$, $R(t)$.}

\emph{Step 2. The extension.}
Let $\wh W=\{\wh W_t,\wh\F_t;\ 0\le t<\infty\}$ be a standard Brownian motion on a probability space $(\wh\Omega,\wh\F,\wh P)$, with all its paths continuous and with a filtration that satisfies the usual conditions, and let $(\wt\Omega,\wt\F,\wt P)$, $\{\wt\F_t\}$ be the extension constructed in \bk{Remark 4.1}: $\wt\Omega=\Omega\times\wh\Omega$, $\wt P=P\times\wh P$ on $\G=\F\otimes\wh\F$, $\G_t=\F_t\otimes\wh\F_t$, $\wt\F_t=\bigcap_{s>t}\sigma(\G_s\cup\cN)$, where $\cN$ is the collection of $\wt P$-null sets of $\G$, and $\wt\F=\sigma(\G\cup\cN)$. For a random variable or a process $Z$ on $\Omega$ (respectively, on $\wh\Omega$) we denote by $\wt Z$ its \emph{lift} $\wt Z(\omega,\wh\omega)=Z(\omega)$ (respectively, $Z(\wh\omega)$). The lift of a random variable has the same law under $\wt P$ as the random variable has under $P$ (respectively, $\wh P$); the lift of an $\{\F_t\}$-progressively measurable process is $\{\wt\F_t\}$-progressively measurable, because $\G_t\subseteq\wt\F_t$; the lift of a stopping time of $\{\F_t\}$ or of $\{\wh\F_t\}$ is a stopping time of $\{\wt\F_t\}$, for the same reason; and a relation which holds $P$-a.s.\ on $\Omega$ holds $\wt P$-a.s.\ for the lifts.

\begin{lem}[2]
Let $M\in\cM^{c,\loc}$ on $(\Omega,\F,P)$, $\{\F_t\}$ and $N\in\cM^{c,\loc}$ on $(\wh\Omega,\wh\F,\wh P)$, $\{\wh\F_t\}$. Then the processes $\wt M$, $\wt N$ and $\wt M\wt N$ belong to $\cM^{c,\loc}$ on $(\wt\Omega,\wt\F,\wt P)$, $\{\wt\F_t\}$. Moreover,
\[
  \qv{\wt M}=\wt{\qv{M}}\qquad\text{and}\qquad\qv{\wt M,\wt N}\equiv0,\qquad\wt P\text{-a.s.}
\]
\end{lem}

\begin{proof}
Modifying $M$ and $N$ on null sets (which belong to $\F_0$, $\wh\F_0$), we may assume that all their paths are continuous and start at zero.

\emph{(a)} Suppose first that $M$ and $N$ are bounded, continuous martingales on their respective spaces (not necessarily null at $t=0$), and set $L\defeq\wt M\wt N$. This process is bounded, continuous, and $L_t$ is $\G_t$-measurable. For $0\le s<t$, $A\in\F_s$ and $B\in\wh\F_s$, the Fubini--Tonelli theorem and the martingale properties of $M$ and $N$ give
\[
  \wt E\bigl[L_t\ind_{A\times B}\bigr]=E[M_t\ind_A]\cdot\wh E[N_t\ind_B]
  =E[M_s\ind_A]\cdot\wh E[N_s\ind_B]=\wt E\bigl[L_s\ind_{A\times B}\bigr].
\]
The collection of sets $C\in\wt\F$ with $\wt E[(L_t-L_s)\ind_C]=0$ is a Dynkin system, because $L_t-L_s$ is integrable; it contains the $\pi$-system of rectangles $A\times B$ as above, which includes $\wt\Omega$ and generates $\G_s$. By Dynkin's $\pi$-$\lambda$ theorem it contains $\G_s$, and then also $\sigma(\G_s\cup\cN)$, since every set in the latter $\sigma$-field differs from a set in $\G_s$ by a $\wt P$-null set. Now let $C\in\wt\F_s$. For every $r\in(s,t)$ we have $C\in\sigma(\G_r\cup\cN)$, thus $\wt E[L_t\ind_C]=\wt E[L_r\ind_C]$; letting $r\downarrow s$ and using the continuity and boundedness of $L$, we get $\wt E[L_t\ind_C]=\wt E[L_s\ind_C]$. Hence $L$ is a bounded, continuous martingale with respect to $\{\wt\F_t\}$. With $N\equiv1$, this shows that $\wt M$ is such a martingale, and by symmetry so is $\wt N$.

\emph{(b)} In the general case, let $S_n\defeq\inf\{t\ge0;\ \abs{M_t}\ge n\}$ and $T_n\defeq\inf\{t\ge0;\ \abs{N_t}\ge n\}$. These are stopping times (\bk{Problem 1.2.7}), nondecreasing in $n$, with $\lim_nS_n=\lim_nT_n=\infty$ because the paths are bounded on bounded intervals. The stopped processes $M^{(n)}_t\defeq M_{t\wedge S_n}$ and $N^{(n)}_t\defeq N_{t\wedge T_n}$ are continuous local martingales bounded by $n$, hence martingales (stop them along a localizing sequence and pass to the limit in the martingale identity, by bounded convergence). By (a), their lifts and the product $L^{(n)}$ of their lifts are bounded, continuous $\{\wt\F_t\}$-martingales. The lifts $\wt S_n$, $\wt T_n$ are stopping times of $\{\wt\F_t\}$ which increase to infinity, and so does $\tau_n\defeq\wt S_n\wedge\wt T_n$. We have
\[
  \wt M_{t\wedge\wt S_n}=\wt{M^{(n)}}_t,\qquad \wt N_{t\wedge\wt T_n}=\wt{N^{(n)}}_t,\qquad
  (\wt M\wt N)_{t\wedge\tau_n}=L^{(n)}_{t\wedge\tau_n},
\]
and the last process is a martingale by the optional sampling theorem (\bk{Problem 1.3.24}\,(i)). Therefore $\wt M$, $\wt N$ and $\wt M\wt N$ are continuous local martingales of $\{\wt\F_t\}$, null at $t=0$.

\emph{(c)} Since $\wt M\wt N-0\in\cM^{c,\loc}$, the uniqueness assertion of \bk{Problem 1.5.17} yields $\qv{\wt M,\wt N}\equiv0$. Finally, $M^2-\qv{M}$ belongs to $\cM^{c,\loc}$ on $\Omega$, so by (b) its lift $\wt M^2-\wt{\qv{M}}$ belongs to $\cM^{c,\loc}$ on $\wt\Omega$; the process $\wt{\qv{M}}$ is adapted to $\{\wt\F_t\}$, a.s.\ continuous, nondecreasing and null at $t=0$, so the same uniqueness assertion gives $\qv{\wt M}=\wt{\qv{M}}$.
\end{proof}

As in \bk{Remark 4.1}, we now drop the tildes on lifted objects. By that Remark, $\wh W$ is a standard Brownian motion with respect to $\{\wt\F_t\}$ on the extended space (this also follows from Lemma~2 and P.~L\'evy's \bk{Theorem 3.16}, because $\qv{\wh W}_t=t$).

\emph{Step 3. The processes of the hint.}
On $\Omega$ we consider the progressively measurable processes
\[
  I_t\defeq\int_0^t\sigma_s\,dW_s\in\cM^{c,\loc},\qquad
  \theta_t\defeq\sqrt{\bigl(\rho^2(t)-\sigma_t^2\bigr)^+},
\]
for which $\qv{I}_t=\int_0^t\sigma_s^2\,ds$ and $0\le\theta\le\rho$. In view of \eqref{eq:3-4-24:ae}, $\theta_s=\sqrt{\rho^2(s)-\sigma_s^2}$ and
\begin{equation}\label{eq:3-4-24:sum}
  \sigma_s^2+\theta_s^2=\rho^2(s)\quad\text{for Lebesgue-a.e. }s\ge0,\ \text{a.s.}
\end{equation}
On the extended space, the (lifted) process $\theta$ is progressively measurable with respect to $\{\wt\F_t\}$ and satisfies $\int_0^t\theta_s^2\,ds\le R(t)<\infty$; thus $\theta\in\cP^*(\wh W)$ and the stochastic integral
\[
  J_t\defeq\int_0^t\theta_s\,d\wh W_s;\qquad 0\le t<\infty
\]
is defined, belongs to $\cM^{c,\loc}$ with respect to $\{\wt\F_t\}$, and has $\qv{J}_t=\int_0^t\theta_s^2\,ds$. By Lemma~2 (with $M=I$, $N=\wh W$), the lifted process $I$ belongs to $\cM^{c,\loc}$ with respect to $\{\wt\F_t\}$, with $\qv{I}_t=\int_0^t\sigma_s^2\,ds$ and $\qv{\wh W,I}\equiv0$. The characterization of the stochastic integral in \bk{Proposition 2.24} (formula \bk{(2.30)} with $\Phi=J$ and $N=I$) now gives
\[
  \qv{J,I}_t=\int_0^t\theta_s\,d\qv{\wh W,I}_s=0;\qquad 0\le t<\infty,\ \text{a.s.}
\]
For $i=1,2$ we set, on the extended space,
\[
  M^{(i)}_t\defeq I_t+(-1)^iJ_t,\qquad Y^{(i)}_t\defeq Y_0+a(t)+M^{(i)}_t;\qquad 0\le t<\infty .
\]
These are the processes of the hint (with the lift of $\int_0^\cdot\sigma_s\,dW_s$ as their third term). Each $M^{(i)}$ is in $\cM^{c,\loc}$ with respect to $\{\wt\F_t\}$ and, by the bilinearity of the cross-variation and \eqref{eq:3-4-24:sum},
\[
  \qv{M^{(i)}}_t=\qv{I}_t+\qv{J}_t+2(-1)^i\qv{I,J}_t=\int_0^t\bigl(\sigma_s^2+\theta_s^2\bigr)\,ds=R(t);\qquad 0\le t<\infty,\ \text{a.s.}
\]
The random variable $Y_0$ (lifted) is measurable with respect to $\F_0\otimes\{\emptyset,\wh\Omega\}\subseteq\wt\F_0$. Lemma~1 and the observation following it, applied on the extended space with $\xi=Y_0$, $b=a(t)$, show that the law of $Y^{(i)}_t$ under $\wt P$ is the convolution of the law of $Y_0$ with the normal law of mean $a(t)$ and variance $R(t)$. On the original space, $N_t\defeq\int_0^t\rho(s)\,dV_s$ is in $\cM^{c,\loc}$ with respect to $\{\F_t\}$ with $\qv{N}_t=R(t)$, $Y_0$ is $\F_0$-measurable, and $Y_t=Y_0+a(t)+N_t$; by Lemma~1 again, the law of $Y_t$ under $P$ is the same convolution. Therefore, for every $t\ge0$,
\begin{equation}\label{eq:3-4-24:law}
  Y^{(1)}_t,\ Y^{(2)}_t\ \text{(under $\wt P$)}\quad\text{and}\quad Y_t\ \text{(under $P$)}\quad\text{have the same law.}
\end{equation}

\emph{Step 4. The comparison.}
Fix $t\ge0$. On the extended space we have, $\wt P$-a.s.,
\begin{equation}\label{eq:3-4-24:avg}
  \frac{Y^{(1)}_t+Y^{(2)}_t}{2}=Y_0+a(t)+I_t=X_t+(Y_0-X_0)+\int_0^t\bigl(m(s)-\mu_s\bigr)\,ds\ \ge\ X_t,
\end{equation}
where $X$ denotes the lift of the given process: the first equality holds because the terms $\pm J_t$ cancel, the second by the definition of $X$ and \eqref{eq:3-4-24:hyp}, and the inequality because $X_0\le Y_0$ a.s.\ and the integrand $m(s)-\mu_s$ is nonnegative for Lebesgue-a.e.\ $s$, a.s., by \eqref{eq:3-4-24:ae}.

\emph{The tail probabilities.} Let $c\in\R$. By \eqref{eq:3-4-24:avg}, $X_t\le\max\bigl(Y^{(1)}_t,Y^{(2)}_t\bigr)$ a.s., so that, up to a $\wt P$-null set, $\{X_t\ge c\}\subseteq\{Y^{(1)}_t\ge c\}\cup\{Y^{(2)}_t\ge c\}$. Since the lift of $X_t$ has under $\wt P$ the law of $X_t$ under $P$, \eqref{eq:3-4-24:law} yields
\[
  P(X_t\ge c)=\wt P(X_t\ge c)\le\wt P\bigl(Y^{(1)}_t\ge c\bigr)+\wt P\bigl(Y^{(2)}_t\ge c\bigr)=2P(Y_t\ge c).
\]

\emph{The expectations.} A convex function $f\colon\R\to\R$ is continuous, hence Borel-measurable. Since $f$ is nondecreasing and convex, \eqref{eq:3-4-24:avg} gives
\begin{equation}\label{eq:3-4-24:f}
  f(X_t)\le f\Bigl(\tfrac12Y^{(1)}_t+\tfrac12Y^{(2)}_t\Bigr)\le\tfrac12f\bigl(Y^{(1)}_t\bigr)+\tfrac12f\bigl(Y^{(2)}_t\bigr),\qquad\wt P\text{-a.s.}
\end{equation}
The function $f^+=\max(f,0)$ is also nondecreasing and convex, so \eqref{eq:3-4-24:f} holds for it as well, and taking expectations of nonnegative random variables we get from \eqref{eq:3-4-24:law}
\[
  Ef^+(X_t)\le\tfrac12\wt Ef^+\bigl(Y^{(1)}_t\bigr)+\tfrac12\wt Ef^+\bigl(Y^{(2)}_t\bigr)=Ef^+(Y_t).
\]
Suppose that $Ef^+(Y_t)<\infty$. Then $f(X_t)$, $f(Y^{(1)}_t)$ and $f(Y^{(2)}_t)$ have integrable positive parts; their expectations are well defined in $[-\infty,\infty)$, and the expectation is monotone and additive on such random variables. Taking expectations in \eqref{eq:3-4-24:f} and using \eqref{eq:3-4-24:law} once more, we obtain
\[
  Ef(X_t)\le\tfrac12\wt Ef\bigl(Y^{(1)}_t\bigr)+\tfrac12\wt Ef\bigl(Y^{(2)}_t\bigr)=Ef(Y_t).
\]
If, on the other hand, $Ef^+(Y_t)=\infty$ and $Ef(Y_t)$ is defined, then $Ef(Y_t)=\infty$ and there is nothing to prove. Thus $Ef(X_t)\le Ef(Y_t)$ holds whenever the right-hand side is well defined (and the left-hand side is then well defined, except possibly when $Ef(Y_t)=\infty$). This is the case, for instance, if $f$ is bounded from below, or if $Y_0$ is integrable: indeed $f(y)\ge f(0)+\kappa y$ for some $\kappa\ge0$ and all $y\in\R$, so $f^-(Y_t)\le\abs{f(0)}+\kappa\abs{Y_t}$, and $Y_t-Y_0$ is normal.

\rmk{\emph{On the integrability hypothesis.} Since $m(s)\ge\mu_s\ge-\abs{\mu_s}$, the integrand $\abs{\mu_s}+m(s)+\rho^2(s)$ of the statement is nonnegative, but the finiteness of its integral does not imply that of the integrals of its three terms: with $\mu_s=m(s)=-1/s$ and $\sigma=\rho\equiv0$ the integrand vanishes identically, whereas $\int_0^tm(s)\,ds=-\infty$ and neither $X$ nor $Y$ is defined. The condition has to be understood as \eqref{eq:3-4-24:hyp}, which is exactly what is needed for the integrals defining $X$ and $Y$ to make sense. Note that, given $\int_0^t\abs{\mu_s}\,ds<\infty$ a.s.\ and $\mu\le m$, we have $m^-\le\mu^-$, so that the requirement on $m$ in \eqref{eq:3-4-24:hyp} amounts to $\int_0^tm^+(s)\,ds<\infty$.}

\rmk{\emph{On the setting.} The Brownian motion $V$ and the filtration with respect to which $Y$ is a semimartingale entered only through \eqref{eq:3-4-24:law}, i.e., through the law of $Y_t$. Hence the argument applies verbatim if $V$ is a Brownian motion with respect to another filtration $\{\cH_t\}$ on $(\Omega,\F,P)$ which satisfies the usual conditions, as long as $Y_0$ is measurable with respect to both $\F_0$ and $\cH_0$: the first is needed for $Y^{(1)}$, $Y^{(2)}$ to be adapted to $\{\wt\F_t\}$ and for Lemma~1 to apply to them, the second for Lemma~1 to apply to $Y$.}
\end{kex}

\begin{kex}{Exercise}{4.25}{}
As everywhere in this section, $M$ is adapted to a filtration $\{\F_t\}$ which satisfies the usual conditions. The process $N\defeq M-m$ belongs to $\cM^{c,\loc}$, and $\qv{M}$ stands for $\qv{N}$; we write $S\defeq\qv{M}_\infty=\lim_{t\to\infty}\qv{M}_t\le\infty$. We shall prove more than is asked, namely that both limits exist and that
\begin{align}
  x\,P\Bigl[\sup_{t\ge0}M_t\ge x\Bigr]&=m\quad\text{for every }x\ge m,\label{eq:3-4-25:sup}\\
  y\,P\bigl[\sqrt{\qv{M}_\infty}\ge y\bigr]&=y\sqrt{\frac{2}{\pi}}\int_0^{m/y}e^{-z^2/2}\,dz\quad\text{for every }y>0.\label{eq:3-4-25:qv}
\end{align}

\emph{Step 1. Time-change.}
By \bk{Theorem 4.6} (if $S=\infty$ a.s.) or \bk{Problem 4.7} (if $P[S<\infty]>0$), applied to $N$, there is a standard, one-dimensional Brownian motion $B=\{B_s;\ 0\le s<\infty\}$ such that, a.s.,
\begin{equation}\label{eq:3-4-25:tc}
  M_t=m+B_{\qv{M}_t};\qquad 0\le t<\infty.
\end{equation}
In the second case $B$ is defined on an extension of the original probability space (\bk{Remark 4.1}), to which $M$ and $\qv{M}$ are carried over without change of their joint law. Since \eqref{eq:3-4-25:sup} and \eqref{eq:3-4-25:qv} involve only this law, it suffices to prove them on the extension, which we denote again by $(\Omega,\F,P)$. For $b\in\R$ let
\[
  T_b\defeq\inf\{s\ge0;\ B_s=b\},\qquad T'\defeq\inf\{s\ge0;\ B_s<-m\}
\]
(with $\inf\emptyset=\infty$). We denote by $\Omega_0$ the event on which $B_0=0$, $s\mapsto B_s$ is continuous, $M$ and $\qv{M}$ are continuous with $\qv{M}_0=0$ and $\qv{M}$ nondecreasing, \eqref{eq:3-4-25:tc} holds, $M_t\ge0$ for all $t\ge0$, and $\lim_{t\to\infty}M_t=0$; then $P(\Omega_0)=1$.

\emph{Step 2. The passage times.}
According to \bk{(2.6.2)}, which is established for $b>0$ in the proof of \bk{Proposition 2.6.19} and holds for $b<0$ as well because $-B$ is also a standard Brownian motion, we have for $b\ne0$, $t>0$:
\begin{equation}\label{eq:3-4-25:pass}
  P[T_b<t]=\sqrt{\frac{2}{\pi}}\int_{\abs{b}/\sqrt t}^\infty e^{-z^2/2}\,dz,
  \qquad\text{and consequently}\qquad P[T_b<\infty]=1
\end{equation}
(let $t\to\infty$; cf.\ \bk{Remark 2.6.20}). We claim that
\begin{equation}\label{eq:3-4-25:Tprime}
  T'=T_{-m}<\infty\quad\text{a.s.}
\end{equation}
We argue on $\Omega_0$, and set $b_n\defeq-m-\frac1n$ for $n\ge1$. If $B_s<-m$ for some $s\ge0$, then, since $B_0=0>-m$, the intermediate value theorem provides $u\in[0,s)$ with $B_u=-m$; hence $T_{-m}\le T'$. If $T_{b_n}<\infty$, then $B_{T_{b_n}}=b_n<-m$ by continuity, so $T'\le T_{b_n}$. Conversely, if $B_s<-m$, then $B_s\le b_n$ for all sufficiently large $n$, and the intermediate value theorem gives $T_{b_n}\le s$ for these $n$. It follows that, for every $t>0$, the event $\{T'<t\}$ is the union of the events $\{T_{b_n}<t\}$, $n\ge1$; these increase with $n$, because $T_{b_{n+1}}\le T_{b_n}$ (again by the intermediate value theorem, as $0>b_{n+1}>b_n$). From \eqref{eq:3-4-25:pass} and the continuity of the integral in its lower limit,
\[
  P[T'<t]=\lim_{n\to\infty}P[T_{b_n}<t]
  =\lim_{n\to\infty}\sqrt{\frac{2}{\pi}}\int_{(m+\frac1n)/\sqrt t}^\infty e^{-z^2/2}\,dz
  =P[T_{-m}<t].
\]
Because $\{T'<t\}\subseteq\{T_{-m}<t\}$ up to a null set, this gives $P[T_{-m}<t\le T']=0$ for every $t>0$. Taking the union over the rational $t>0$ we obtain $P[T_{-m}<T']=0$, and \eqref{eq:3-4-25:Tprime} follows from $T_{-m}\le T'$ and \eqref{eq:3-4-25:pass}.

\emph{Step 3. $\qv{M}_\infty=T_{-m}$, a.s.}
Let $\omega\in\Omega_0$. The function $t\mapsto\qv{M}_t(\omega)$ is continuous and nondecreasing, starts at zero and increases to $S(\omega)$; by the intermediate value theorem its range $R(\omega)\defeq\{\qv{M}_t(\omega);\ 0\le t<\infty\}$ satisfies
\begin{equation}\label{eq:3-4-25:range}
  [0,S(\omega))\subseteq R(\omega)\subseteq[0,S(\omega)]\cap[0,\infty).
\end{equation}
For every $s\in[0,S(\omega))$ there is thus a $t\ge0$ with $s=\qv{M}_t(\omega)$, and \eqref{eq:3-4-25:tc} gives $m+B_s(\omega)=M_t(\omega)\ge0$. Therefore $B_s(\omega)\ge-m$ for all $0\le s<S(\omega)$, which means
\[
  S\le T'\quad\text{on }\Omega_0.
\]
In conjunction with \eqref{eq:3-4-25:Tprime}, this shows $S\le T_{-m}<\infty$ a.s. On the other hand, for $\omega\in\Omega_0\cap\{S<\infty\}$ we have $\qv{M}_t(\omega)\to S(\omega)$ as $t\to\infty$, and the continuity of $B_\cdot(\omega)$ at $S(\omega)$ together with \eqref{eq:3-4-25:tc} yields
\[
  m+B_{S(\omega)}(\omega)=\lim_{t\to\infty}\bigl(m+B_{\qv{M}_t(\omega)}(\omega)\bigr)=\lim_{t\to\infty}M_t(\omega)=0,
\]
so that $T_{-m}\le S$ on $\Omega_0\cap\{S<\infty\}$. We conclude that
\begin{equation}\label{eq:3-4-25:S}
  \qv{M}_\infty=S=T_{-m}<\infty\quad\text{a.s.}
\end{equation}

\emph{Step 4. Proof of \eqref{eq:3-4-25:qv} and computation of $\sigma$.}
Let $y>0$. From \eqref{eq:3-4-25:S}, \eqref{eq:3-4-25:pass} with $b=-m$, $t=y^2$, and $\sqrt{2/\pi}\int_0^\infty e^{-z^2/2}\,dz=1$,
\[
  P\bigl[\sqrt{\qv{M}_\infty}\ge y\bigr]=1-P[T_{-m}<y^2]
  =1-\sqrt{\frac{2}{\pi}}\int_{m/y}^\infty e^{-z^2/2}\,dz
  =\sqrt{\frac{2}{\pi}}\int_0^{m/y}e^{-z^2/2}\,dz,
\]
which is \eqref{eq:3-4-25:qv}. Since $e^{-m^2/2y^2}\le e^{-z^2/2}\le1$ for $0\le z\le m/y$, we obtain
\[
  \sqrt{\frac{2}{\pi}}\;m\,e^{-m^2/2y^2}\le y\,P\bigl[\sqrt{\qv{M}_\infty}\ge y\bigr]\le\sqrt{\frac{2}{\pi}}\;m,
\]
and letting $y\to\infty$ we see that the limit $\sigma$ exists and equals $\sqrt{2/\pi}\;m$; that is, $\sqrt{\pi/2}\,\sigma=m$.

\emph{Step 5. Proof of \eqref{eq:3-4-25:sup} and computation of $l$.}
For $x=m$ the identity \eqref{eq:3-4-25:sup} is obvious, because $M_0=m$. Let $x>m$, and let $\omega\in\Omega_0$ be such that $S(\omega)=T_{-m}(\omega)<\infty$; by \eqref{eq:3-4-25:S}, almost every $\omega$ has this property. From \eqref{eq:3-4-25:tc}, \eqref{eq:3-4-25:range} and the continuity of $B_\cdot(\omega)$ on the compact interval $[0,S(\omega)]$,
\[
  \sup_{t\ge0}M_t(\omega)=m+\sup_{s\in R(\omega)}B_s(\omega)=m+\max_{0\le s\le T_{-m}(\omega)}B_s(\omega).
\]
If this quantity is at least $x$, then $B_s(\omega)\ge x-m>0$ for some $s\in[0,T_{-m}(\omega)]$, and the intermediate value theorem yields $T_{x-m}(\omega)\le s\le T_{-m}(\omega)$; equality is impossible here, since $B_{T_{-m}(\omega)}(\omega)=-m\ne x-m$. Conversely, if $T_{x-m}(\omega)<T_{-m}(\omega)$, then the maximum is at least $B_{T_{x-m}(\omega)}(\omega)=x-m$. Consequently,
\begin{equation}\label{eq:3-4-25:exit}
  P\Bigl[\sup_{t\ge0}M_t\ge x\Bigr]=P[T_{x-m}<T_{-m}].
\end{equation}
Now $W\defeq m+B$ is a Brownian motion starting at $m$; its passage times to the levels $x$ and $0$ are $T_{x-m}$ and $T_{-m}$, respectively, and the event $\{T_{x-m}<T_{-m}\}$ is determined by the path of $W$, whose law on $\bigl(C[0,\infty),\B(C[0,\infty))\bigr)$ is that of a Brownian motion under the measure $P^m$ of a one-dimensional Brownian family. \bk{Exercise 2.8.13} (with $a=x$ and starting point $m\in[0,x]$) therefore shows that the right-hand side of \eqref{eq:3-4-25:exit} equals $m/x$, and \eqref{eq:3-4-25:sup} is proved. In particular the limit $l$ exists, and $l=m$.

Steps 4 and 5 give $l=\sqrt{\pi/2}\,\sigma=m$.

\rmk{The identities \eqref{eq:3-4-25:sup} and \eqref{eq:3-4-25:S} say that $m/\sup_{t\ge0}M_t$ is uniformly distributed on $(0,1]$, and that $\qv{M}_\infty$ has the law \bk{(2.6.3)} of the first passage time of Brownian motion to the level $m$; in particular $E\bigl[\sup_{t\ge0}M_t\bigr]=E\sqrt{\qv{M}_\infty}=\infty$, since both tails are of the exact order $1/x$. The script \texttt{code/c3\_4\_25.py} checks \eqref{eq:3-4-25:sup} and \eqref{eq:3-4-25:qv} by simulation for $M_t=m\exp(W_t-\frac t2)$.}
\end{kex}

\solsection{3.5}{The Girsanov Theorem}

\begin{kex}{Problem}{5.6}{}
We fix $T\in[0,\infty)$ and $i\in\{1,\dots,d\}$, and write $Z=Z(X)$, $\wt{E}_T$ for the expectation
under $\wt{P}_T$, and $\qv{\cdot,\cdot}^{\sim}$ for cross-variations computed under $\wt{P}_T$.

\emph{Step 0. The two integrals are well defined.}
The random variable $Z_T$ of \bk{(5.2)} is strictly positive, so \bk{(5.4)} gives
$P(A)=\wt{E}_T[\ind_A Z_T^{-1}]$ for $A\in\F_T$: the measures $P$ and $\wt{P}_T$ have the same null
sets in $\F_T$. Consequently ``a.s.\ $P$'' and ``a.s.\ $\wt{P}_T$'' mean the same thing for events of
$\F_T$, and the filtration $\{\F_t\}_{0\le t\le T}$ satisfies the usual conditions on
$(\Omega,\F_T,\wt{P}_T)$ as well. All processes below are considered on $[0,T]$ only; the results of
\bk{Section 3.2} apply to them once they are stopped at time $T$.

The integrand $Y$ is measurable and adapted, but not necessarily progressively measurable. By
\bk{Proposition 1.1.12} it has a progressively measurable modification $\wh{Y}$; by Tonelli's
theorem, $E\int_0^T \ind_{\{Y_t\ne \wh{Y}_t\}}\,dt=\int_0^T P[Y_t\neq \wh{Y}_t]\,dt=0$, so that
\begin{equation}\label{eq:3-5-06:mod}
  \operatorname{meas}\{t\in[0,T];\ Y_t\neq \wh{Y}_t\}=0,\qquad \text{a.s. } P \text{ and } \wt{P}_T
\end{equation}
(the event in \eqref{eq:3-5-06:mod} belongs to $\F$ and has $P$-probability one, hence belongs to
$\F_0\subseteq\F_T$ by the usual conditions). Thus $Y$ and $\wh{Y}$ are equivalent, in the sense of
\bk{Definition 2.22}, as integrands for any continuous local martingale whose quadratic variation
is $t$ on $[0,T]$, under $P$ as well as under $\wt{P}_T$; and
$\int_0^T \wh{Y}_t^2\,dt=\int_0^T Y_t^2\,dt<\infty$ a.s.\ $P$ and $\wt{P}_T$. Now $W^{(i)}$ is a
Brownian motion under $P$, and $\wt{W}^{(i)}$ of \bk{(5.6)} is a Brownian motion on $[0,T]$ under
$\wt{P}_T$ by \bk{Theorem 5.1}; both have quadratic variation $t$ under the respective measure. As
explained after \bk{Definition 2.22} (see also \bk{Remark 2.11}), the integral of
$Y\in\mathscr{P}$ with respect to either of them is the integral of the progressively measurable
representative $\wh{Y}\in\mathscr{P}^*$ of its equivalence class (the construction of the
integral, by localization and approximation by simple processes in the metric of
\bk{Definition 2.1}, depends on the integrand only through its equivalence class), and
\bk{Proposition 2.24} holds for it. Finally, \eqref{eq:3-5-06:mod} shows that
$\int_0^t Y_sX_s^{(i)}\,ds=\int_0^t \wh{Y}_sX_s^{(i)}\,ds$, $0\le t\le T$, almost surely, and this
integral is well defined and finite because, by the Cauchy--Schwarz inequality and \bk{(5.1)},
\[
  \int_0^T \abs{Y_sX^{(i)}_s}\,ds\le
  \Bigl(\int_0^T Y_s^2\,ds\Bigr)^{1/2}\Bigl(\int_0^T (X^{(i)}_s)^2\,ds\Bigr)^{1/2}<\infty
  \qquad\text{a.s.}
\]
Hence nothing in the asserted identity changes if $Y$ is replaced by $\wh{Y}$, and \emph{we assume
from now on that $Y$ is progressively measurable.}

\emph{Step 1. The right-hand side is a local martingale under $\wt{P}_T$.}
Let
\[
  M_t\defeq\int_0^t Y_s\,dW^{(i)}_s,\qquad 0\le t\le T,
\]
the It\^o integral under $P$; it belongs to $\cM^{c,\loc}_T$. Recall that
$\qv{W^{(i)},W^{(j)}}_t=\delta_{ij}t$. By \bk{Proposition 2.24} (under $P$), applied with
$N=W^{(j)}$ and then with $N=M$, we have a.s.\ $P$, for $0\le t\le T$ and $1\le j\le d$,
\begin{equation}\label{eq:3-5-06:cv}
  \qv{M,W^{(j)}}_t=\int_0^t Y_s\,d\qv{W^{(i)},W^{(j)}}_s=\delta_{ij}\int_0^t Y_s\,ds,
  \qquad
  \qv{M}_t=\int_0^t Y_s\,d\qv{W^{(i)},M}_s=\int_0^t Y_s^2\,ds .
\end{equation}
The first relation says that the function $t\mapsto\qv{M,W^{(j)}}_t$ is absolutely continuous
with density $\delta_{ij}Y_t$, so
$\int_0^t X^{(j)}_s\,d\qv{M,W^{(j)}}_s=\delta_{ij}\int_0^t X^{(i)}_sY_s\,ds$, and the process
associated with $M$ in \bk{(5.8)} is
\[
  \wt{M}_t\defeq M_t-\sum_{j=1}^d\int_0^t X^{(j)}_s\,d\qv{M,W^{(j)}}_s
  =\int_0^t Y_s\,dW^{(i)}_s-\int_0^t Y_sX^{(i)}_s\,ds,\qquad 0\le t\le T,
\]
which is exactly the right-hand side of the identity to be proved. Likewise, by \bk{(5.6)}, the
process associated in \bk{(5.8)} with $N\defeq W^{(i)}\in\cM^{c,\loc}_T$ is
$\wt{N}=\wt{W}^{(i)}$. \bk{Proposition 5.4} now yields
$\wt{M}\in\wt{\cM}^{c,\loc}_T$ together with
\begin{equation}\label{eq:3-5-06:tcv}
  \qv{\wt{M}}^{\sim}_t=\qv{M}_t=\int_0^t Y_s^2\,ds,
  \qquad
  \qv{\wt{M},\wt{W}^{(i)}}^{\sim}_t=\qv{M,W^{(i)}}_t=\int_0^t Y_s\,ds;
  \qquad 0\le t\le T,
\end{equation}
a.s.\ $P$ and $\wt{P}_T$, where we have used \eqref{eq:3-5-06:cv}.

\emph{Step 2. Identification with the It\^o integral under $\wt{P}_T$.}
Let
\[
  \Phi_t\defeq\int_0^t Y_s\,d\wt{W}^{(i)}_s,\qquad 0\le t\le T,
\]
be the It\^o integral under $\wt{P}_T$; it belongs to $\wt{\cM}^{c,\loc}_T$. We apply
\bk{Proposition 2.24} on $(\Omega,\F_T,\wt{P}_T)$ to the local martingale $\wt{W}^{(i)}$, whose
quadratic variation under $\wt{P}_T$ is $\qv{\wt{W}^{(i)}}^{\sim}_t=t$. Relation \bk{(2.30)} with
$N=\wt{W}^{(i)}$ gives
\[
  \qv{\Phi,\wt{W}^{(i)}}^{\sim}_t=\int_0^t Y_s\,d\qv{\wt{W}^{(i)}}^{\sim}_s=\int_0^t Y_s\,ds;
  \qquad 0\le t\le T,\quad\text{a.s. }\wt{P}_T,
\]
so that, by the symmetry of the cross-variation, $t\mapsto\qv{\wt{W}^{(i)},\Phi}^{\sim}_t$ is
absolutely continuous with density $Y_t$; by \eqref{eq:3-5-06:tcv}, the same is true of
$t\mapsto\qv{\wt{W}^{(i)},\wt{M}}^{\sim}_t$. Relation \bk{(2.30)} with $N=\wt{M}$, and then with
$N=\Phi$, now gives
\[
  \qv{\Phi,\wt{M}}^{\sim}_t=\int_0^t Y_s\,d\qv{\wt{W}^{(i)},\wt{M}}^{\sim}_s=\int_0^t Y_s^2\,ds,
  \qquad
  \qv{\Phi}^{\sim}_t=\int_0^t Y_s\,d\qv{\wt{W}^{(i)},\Phi}^{\sim}_s=\int_0^t Y_s^2\,ds,
\]
for $0\le t\le T$, a.s.\ $\wt{P}_T$. By the bilinearity of the cross-variation and
\eqref{eq:3-5-06:tcv},
\[
  \qv{\Phi-\wt{M}}^{\sim}_t=\qv{\Phi}^{\sim}_t-2\qv{\Phi,\wt{M}}^{\sim}_t+\qv{\wt{M}}^{\sim}_t=0;
  \qquad 0\le t\le T,\quad\text{a.s. }\wt{P}_T .
\]
Thus $D\defeq\Phi-\wt{M}\in\wt{\cM}^{c,\loc}_T$ has $D_0=0$ and quadratic variation zero. If
$\{\tau_n\}$ is a sequence of stopping times with $\tau_n\uparrow T$, eventually equal to $T$,
$\wt{P}_T$-a.s., such that each $D_{\cdot\wedge\tau_n}$ is a bounded martingale under $\wt{P}_T$
(for instance $\tau_n=T\wedge\inf\{t;\ \abs{D_t}\ge n\}$), then
$D^2_{\cdot\wedge\tau_n}-\qv{D}^{\sim}_{\cdot\wedge\tau_n}=D^2_{\cdot\wedge\tau_n}$ is a martingale,
whence $\wt{E}_T[D^2_{t\wedge\tau_n}]=0$ for all $n$ and $t\in[0,T]$. Letting $n\to\infty$ we get
$D_t=0$ a.s.\ $\wt{P}_T$ for each $t\in[0,T]$, and by the continuity of $D$,
\[
  \wt{P}_T\bigl[\Phi_t=\wt{M}_t,\ 0\le t\le T\bigr]=1 .
\]
The event in question is the countable intersection over rational $t\in[0,T]$ (and $t=T$) of the
events $\{\Phi_t=\wt{M}_t\}\in\F_T$, intersected with the event that both paths are continuous;
it lies in $\F_T$, and therefore has $P$-probability one as well, by Step~0. In other words,
\[
  \int_0^t Y_s\,d\wt{W}_s^{(i)} = \int_0^t Y_s\,dW_s^{(i)} - \int_0^t Y_s X_s^{(i)}\,ds;
  \qquad 0 \le t \le T, \quad \text{a.s. } P \text{ and } \wt{P}_T .
\]

\rmk{The argument used only \bk{Proposition 5.4}, with the pair $(M,W^{(i)})$, and the
characterization \bk{(2.30)} of the stochastic integral under each of the two measures; the
converse \bk{Proposition 5.5} was not needed. The result says that a continuous semimartingale
integral $\int Y\,d\wt{W}^{(i)}$ in the sense of $P$ (an It\^o integral plus a Lebesgue integral)
coincides with the It\^o integral computed under $\wt{P}_T$, so the two interpretations discussed
before the problem agree.}
\end{kex}

\begin{kex}{Problem}{5.7}{}
We work in the setting of Subsection~C: $\Omega=C[0,\infty)$, $W$ is the coordinate process, $P$ is Wiener measure on
$\F^W_\infty=\B(C[0,\infty))$, and $Z_t\defeq\exp(\mu W_t-\tfrac12\mu^2t)$. For $\mu\neq0$, $P^{(\mu)}$ is the measure of
\bk{Corollary 5.2} (with $X\equiv\mu$), characterized by
\begin{equation}\label{eq:3-5-07:pmu}
  P^{(\mu)}(A)=E[\ind_A Z_t];\qquad A\in\F^W_t,\ 0\le t<\infty,
\end{equation}
and under it $\wt W_t\defeq W_t-\mu t$ is a standard Brownian motion; for $\mu=0$ we have $Z\equiv1$ and we set
$P^{(0)}\defeq P$, so that \eqref{eq:3-5-07:pmu} holds as well. The process $Z$ is a continuous martingale of the filtration
$\{\F^W_t\}$ under $P$: it is adapted, and for $s<t$ the increment $W_t-W_s$ is independent of $\F^W_s$ and normal with
variance $t-s$, so $E[Z_t\,|\,\F^W_s]=Z_s\,E[e^{\mu(W_t-W_s)-\mu^2(t-s)/2}]=Z_s$. Being continuous and adapted, $Z$ is
progressively measurable (\bk{Proposition 1.1.13}), so $Z_T$ is defined and $\F^W_T$-measurable on $\{T<\infty\}\in\F^W_T$
(\bk{Proposition 1.2.18}). Since $P[T<\infty]=1$, the random variable in \bk{(5.14)} is defined $P$-a.s.\ and is
nonnegative, so that its expectation, which is $E[\ind_{\{T<\infty\}}Z_T]$, makes sense in $[0,\infty]$.

\emph{Step 1: the identity $E[Z_T]=P^{(\mu)}[T<\infty]$.}
Fix $t\in[0,\infty)$. By \bk{Lemma 1.2.16}, applied to the stopping times $T$ and $S\equiv t$ (for which $\F^W_S=\F^W_t$ by
\bk{Problem 1.2.13}), we have $\{T\le t\}\in\F^W_T\cap\F^W_t=\F^W_{T\wedge t}$. The stopping times $T\wedge t\le t$ are
bounded, so the optional sampling theorem in the form of \bk{Problem 1.3.23}\,(i), applied to the continuous martingales
$Z$ and $-Z$, gives $E[Z_t\,|\,\F^W_{T\wedge t}]=Z_{T\wedge t}$ a.s.\ $P$. Therefore, using \eqref{eq:3-5-07:pmu} for the
event $\{T\le t\}\in\F^W_t$ and the fact that $Z_{T\wedge t}=Z_T$ on $\{T\le t\}$,
\begin{equation}\label{eq:3-5-07:t}
  P^{(\mu)}[T\le t]=E[\ind_{\{T\le t\}}Z_t]=E[\ind_{\{T\le t\}}Z_{T\wedge t}]=E[\ind_{\{T\le t\}}Z_T].
\end{equation}
Now let $t\uparrow\infty$ through the integers. The events $\{T\le t\}$ increase to $\{T<\infty\}\in\F^W_\infty$, so the
left-hand side of \eqref{eq:3-5-07:t} increases to $P^{(\mu)}[T<\infty]$ by continuity from below of the measure
$P^{(\mu)}$, and the right-hand side increases to $E[\ind_{\{T<\infty\}}Z_T]$ by monotone convergence
($Z_T\ge0$). Because $P[T<\infty]=1$, we conclude that
\begin{equation}\label{eq:3-5-07:key}
  E\bigl[\exp(\mu W_T-\tfrac12\mu^2T)\bigr]=E[\ind_{\{T<\infty\}}Z_T]=P^{(\mu)}[T<\infty].
\end{equation}

\emph{Step 2: the equivalence.} By \eqref{eq:3-5-07:key}, the left-hand side of \bk{(5.14)} is always a number in $[0,1]$,
and it equals $1$ if and only if $P^{(\mu)}[T<\infty]=1$. Thus \bk{(5.14)} and \bk{(5.15)} are equivalent.

\emph{Step 3: the stopping time $S_b$.} Let $\mu b<0$; in particular $\mu\ne0$ and $b\neq0$. The process $\wt W_t=W_t-\mu t$
is continuous and adapted to $\{\F^W_t\}$, and $S_b$ is its first hitting time of the closed set $\{b\}$; hence $S_b$ is a
stopping time of $\{\F^W_t\}$ (\bk{Problem 1.2.7}). For $\nu\in\R$ let $\varphi_\nu:\Omega\to\Omega$ be the mapping
$\varphi_\nu(\omega)(t)\defeq\omega(t)-\nu t$. It is $\F^W_\infty/\F^W_\infty$-measurable, because each coordinate
$W_t\circ\varphi_\nu=W_t-\nu t$ is, and it is a bijection with inverse $\varphi_{-\nu}$. With
$T_b\defeq\inf\{t\ge0;\ W_t=b\}$, the passage time of $W$ to the level $b$, we have
\begin{equation}\label{eq:3-5-07:Sb}
  S_b=T_b\circ\varphi_\mu .
\end{equation}
We shall use the following fact: \emph{if $Q$ is a probability measure on $(\Omega,\F^W_\infty)$ under which the continuous
process $W\circ\varphi_\nu=\{W_t-\nu t\}$ is a standard Brownian motion, then $Q\circ\varphi_\nu^{-1}=P$.} Indeed, the
finite-dimensional distributions of the coordinate process under $Q\circ\varphi_\nu^{-1}$ are those of $W\circ\varphi_\nu$
under $Q$, i.e.\ those of Brownian motion, and a probability measure on $\B(C[0,\infty))$ is determined by its values on
the finite-dimensional cylinder sets, which form a $\pi$-system generating this $\sigma$-field (Dynkin's
$\pi$-$\lambda$ theorem).

\textbf{(a)} \emph{$P^{(\mu)}[S_b<\infty]=1$.} Under $Q=P^{(\mu)}$ the process $W\circ\varphi_\mu=\wt W$ is a standard
Brownian motion, so $P^{(\mu)}\circ\varphi_\mu^{-1}=P$, and \eqref{eq:3-5-07:Sb} gives
\[
  P^{(\mu)}[S_b<\infty]=P^{(\mu)}[T_b\circ\varphi_\mu<\infty]=P[T_b<\infty]=1 .
\]
The last equality follows from \bk{(5.10)}: as $\alpha\downarrow0$ we have
$e^{-\alpha T_b}\uparrow\ind_{\{T_b<\infty\}}$, so by monotone convergence
$P[T_b<\infty]=\lim_{\alpha\downarrow0}e^{-\abs{b}\sqrt{2\alpha}}=1$. Thus \bk{(5.15)} holds for $T=S_b$.

\textbf{(b)} \emph{$P[S_b<\infty]=1$.} Under $Q=P^{(-\mu)}$ the process $W\circ\varphi_{-\mu}=\{W_t+\mu t\}$ is a standard
Brownian motion, so $P^{(-\mu)}\circ\varphi_{-\mu}^{-1}=P$; applying $\varphi_\mu^{-1}=\varphi_{-\mu}$ to events, this
says $P\circ\varphi_\mu^{-1}=P^{(-\mu)}$. Hence, by \eqref{eq:3-5-07:Sb} and by \bk{(5.13)} with $-\mu$ in place of $\mu$,
\[
  P[S_b<\infty]=P^{(-\mu)}[T_b<\infty]=\exp\bigl[-\mu b-\abs{\mu b}\bigr]=1,
\]
because $-\mu b>0$. So $S_b$ satisfies the standing hypothesis of the problem.

By (a), (b) and Step 2, the Wald identity \bk{(5.14)} holds for $T=S_b$ whenever $\mu b<0$. Since
$W_{S_b}=b+\mu S_b$ on $\{S_b<\infty\}$ by continuity of the paths, it reads
\[
  E\bigl[\exp\bigl(\tfrac12\mu^2S_b\bigr)\bigr]=e^{-\mu b}.
\]

\rmk{Part (a) did not use the sign of $\mu b$: for every $b\in\R$ and $\mu\ne0$ we have $P^{(\mu)}[S_b<\infty]=1$ (trivially
so for $b=0$, where $S_0=0$). The assumption $\mu b<0$ is what guarantees the hypothesis $P[S_b<\infty]=1$. If $\mu b>0$,
then (b) gives $P[S_b<\infty]=e^{-2\mu b}<1$; identity \eqref{eq:3-5-07:t} and its limit
$E[\ind_{\{T<\infty\}}Z_T]=P^{(\mu)}[T<\infty]$ did not use $P[T<\infty]=1$, so in that case one still has
$E[\ind_{\{S_b<\infty\}}\exp(\mu W_{S_b}-\tfrac12\mu^2S_b)]=1$.}
\end{kex}

\begin{kex}{Problem}{5.8}{}
We set $h(t;b,\mu)\defeq 0$ for $t\le 0$, so that the claim reads
\begin{equation}\label{eq:3-5-08:claim}
  h(t;b_1+b_2,\mu)=\int_0^t h(s;b_1,\mu)\,h(t-s;b_2,\mu)\,ds,\qquad t>0 .
\end{equation}
Since $h(t;b,\mu)=h(t;-b,-\mu)$, it suffices to treat $b_1>0$, $b_2>0$, which we assume from now on.

\emph{Step 1: the case $\mu=0$.} Let $W=\{W_t,\F_t;\,0\le t<\infty\}$ be a standard, one-dimensional
Brownian motion on $(\Omega,\F,P)$ and, for $b\neq 0$, let $T_b=\inf\{t\ge 0;\ W_t=b\}$ be the
passage time to the level $b$; it is the hitting time of the closed set $\{b\}$ by a continuous
adapted process, hence a stopping time of $\{\F_t\}$ (\bk{Problem 1.2.7}).
By \bk{(5.9)}, $T_b$ has the density $h(\,\cdot\,;b,0)$ on $(0,\infty)$, and letting
$\alpha\downarrow 0$ in \bk{(5.10)} we obtain from the monotone convergence theorem
$P[T_b<\infty]=\lim_{\alpha\downarrow 0}E\bigl[e^{-\alpha T_b}\ind_{\{T_b<\infty\}}\bigr]=1$.

Let $S\defeq T_{b_1}$ and $\Omega_0\defeq\{S<\infty,\ W_0=0\}$, so that $P(\Omega_0)=1$. Define
$B_t\defeq W_{S+t}-W_S$ on $\Omega_0$ and $B_t\defeq 0$ on $\Omega\setminus\Omega_0$, for
$0\le t<\infty$. By \bk{Theorem 2.6.16}, applied to the a.s.\ finite optional time $S$ of
$\{\F_t\}$, the process $B$ is a standard Brownian motion independent of $\F_{S+}$. The random
variable $S$ is $\F_{S+}$-measurable, because
$\{S\le s\}\cap\{S\le t\}=\{S\le s\wedge t\}\in\F_{s\wedge t}\subseteq\F_{t+}$ for all $s,t\ge 0$; therefore $S$ is independent of
$\F^B_\infty=\sigma(B_t;\ 0\le t<\infty)$. Put
\[
  \tau\defeq\inf\{t\ge 0;\ B_t=b_2\}.
\]
Because $B$ is continuous and $B_0=0<b_2$, we have
$\{\tau\le t\}=\bigcap_{n\ge 1}\bigcup_{q\in\Q\cap[0,t]}\{B_q>b_2-1/n\}\in\F^B_\infty$
for every $t\ge 0$, so $\tau$ is an $\F^B_\infty$-measurable random variable and is
independent of $S$. Let $\varphi(w)\defeq\inf\{t\ge 0;\ w(t)\ge b_2\}$ for $w\in C[0,\infty)$; the same formula, with the
coordinates $w(q)$ in place of $B_q$, shows that $\varphi\colon C[0,\infty)\to[0,\infty]$ is
$\B(C[0,\infty))$-measurable. By the intermediate value theorem, $\varphi(w)=\inf\{t\ge0;\ w(t)=b_2\}$
whenever $w(0)<b_2$; hence $\tau=\varphi(B)$ everywhere and $T_{b_2}=\varphi(W)$ on $\{W_0=0\}$, i.e.\ a.s.
Since $B$ and $W$ induce the same (Wiener) measure on
$C[0,\infty)$, the random variable $\tau$ has the same distribution as $T_{b_2}$, namely the
density $h(\,\cdot\,;b_2,0)$.

We claim that
\begin{equation}\label{eq:3-5-08:sum}
  T_{b_1+b_2}=S+\tau\qquad\text{on }\Omega_0 .
\end{equation}
Fix $\omega\in\Omega_0$. By continuity of the path and the definition of $S$ we have
$W_S(\omega)=b_1$ and $W_t(\omega)<b_1<b_1+b_2$ for $0\le t<S(\omega)$ (recall $W_0(\omega)=0<b_1$, and use the
intermediate value theorem). Hence
$W_t(\omega)\ne b_1+b_2$ for $0\le t\le S(\omega)$, whereas for $t\ge S(\omega)$,
\[
  W_t(\omega)=b_1+b_2\iff B_{t-S(\omega)}(\omega)=b_2 .
\]
Taking the infimum over such $t$ gives $T_{b_1+b_2}(\omega)=S(\omega)+\tau(\omega)$
(both sides being $+\infty$ if $B_\cdot(\omega)$ never reaches $b_2$). This proves
\eqref{eq:3-5-08:sum}.

Now $S$ and $\tau$ are independent with densities $f\defeq h(\,\cdot\,;b_1,0)$ and
$g\defeq h(\,\cdot\,;b_2,0)$, so by \eqref{eq:3-5-08:sum} and the Fubini--Tonelli theorem
(substituting $r=s+u$ in the inner integral), for every $t>0$,
\[
  P[T_{b_1+b_2}\le t]=P[S+\tau\le t]
  =\int_0^\infty\!\!\int_0^\infty \ind_{\{s+u\le t\}}f(s)g(u)\,du\,ds
  =\int_0^t\Bigl(\int_0^r f(s)g(r-s)\,ds\Bigr)dr .
\]
On the other hand, by \bk{(5.9)}, $P[T_{b_1+b_2}\le t]=\int_0^t h(r;b_1+b_2,0)\,dr$. Thus the two
nonnegative functions $r\mapsto h(r;b_1+b_2,0)$ and $r\mapsto (f*g)(r)=\int_0^r f(s)g(r-s)\,ds$
have the same integral over every interval $(0,t]$, and consequently they agree for Lebesgue-almost
every $r>0$. Both are continuous on $(0,\infty)$: this is evident for the first; for the second, note
that $f$ and $g$, extended by zero to $(-\infty,0]$, are bounded and continuous on $\R$
(because $t^{-3/2}e^{-b^2/2t}\to 0$ as $t\downarrow 0$ and as $t\to\infty$) and integrable, so
$r\mapsto\int_\R f(s)g(r-s)\,ds$ is continuous by dominated convergence, with the dominating
function $\bigl(\sup_{\R} g\bigr)f$. Two continuous functions which agree almost everywhere on $(0,\infty)$
agree everywhere, and we have proved
\begin{equation}\label{eq:3-5-08:zero}
  h(t;b_1+b_2,0)=\int_0^t h(s;b_1,0)\,h(t-s;b_2,0)\,ds,\qquad t>0 .
\end{equation}

\emph{Step 2: general $\mu\in\R$.} Expanding the square $(b-\mu t)^2=b^2-2\mu bt+\mu^2t^2$, we see
that
\[
  h(t;b,\mu)=\exp\Bigl(\mu b-\tfrac12\mu^2 t\Bigr)\,h(t;b,0),\qquad t>0,\ b\ne 0 ;
\]
this is precisely the way in which \bk{(5.12)} is obtained from \bk{(5.11)} and \bk{(5.9)}. Therefore,
for $t>0$ and $0<s<t$,
\[
  h(s;b_1,\mu)\,h(t-s;b_2,\mu)
  =\exp\Bigl(\mu(b_1+b_2)-\tfrac12\mu^2t\Bigr)\,h(s;b_1,0)\,h(t-s;b_2,0),
\]
the exponential factor being independent of $s$. Integrating over $s\in(0,t)$ and using
\eqref{eq:3-5-08:zero}, we obtain
\[
  \int_0^t h(s;b_1,\mu)\,h(t-s;b_2,\mu)\,ds
  =\exp\Bigl(\mu(b_1+b_2)-\tfrac12\mu^2t\Bigr)\,h(t;b_1+b_2,0)=h(t;b_1+b_2,\mu),
\]
which is \eqref{eq:3-5-08:claim}.

\rmk{The hypothesis $b_1b_2>0$ cannot be dropped: if $b_1$ and $b_2$ have opposite signs, the level
$b_1$ need not be crossed before $b_1+b_2$ is reached, and \eqref{eq:3-5-08:sum} fails. Indeed, for
$\mu=0$ the Laplace transforms \bk{(5.10)} of the two sides of \eqref{eq:3-5-08:claim} would be
$e^{-\abs{b_1+b_2}\sqrt{2\alpha}}$ and $e^{-(\abs{b_1}+\abs{b_2})\sqrt{2\alpha}}$, which differ when
$b_1b_2<0$. For $b_1b_2>0$ and $\mu b_1<0$ the identity \eqref{eq:3-5-08:claim} is one between
defective densities; integrating it over $(0,\infty)$ gives
$e^{2\mu(b_1+b_2)}=e^{2\mu b_1}e^{2\mu b_2}$, in agreement with \bk{(5.13)}.}
\end{kex}

\begin{kex}{Exercise}{5.9}{}
We work in the setting of Subsection~C: $W$ is the coordinate mapping process on $\Omega=C[0,\infty)$, $P$ is Wiener measure on $(\Omega,\F^W_\infty)$, and $P^{(\mu)}$ is the measure on $\F^W_\infty$ of \bk{Corollary 5.2} with $P^{(\mu)}(A)=E[\ind_A Z_t]$ for $A\in\F^W_t$, where $Z_t=\exp(\mu W_t-\tfrac12\mu^2t)$. For $a\neq 0$ we write $T_a=\inf\{t\ge 0;\ W_t=a\}$ for the passage time to the level $a$.

\emph{Step 1: measurability and the starting point.} Every path of $W$ is continuous, so $W_*=\inf_{t\in\Q,\,t>0}W_t$ is an $\F^W_\infty$-measurable random variable with values in $[-\infty,\infty)$, and $\{T_a<\infty\}=\bigcup_{n\ge 1}\{T_a\le n\}\in\F^W_\infty$, because $T_a$ is a stopping time of $\{\F^W_t\}$ (\bk{Problem 1.2.7}: $T_a$ is the hitting time of the closed set $\{a\}$ by a continuous adapted process). The event $\Omega_0\defeq\{W_0=0\}$ belongs to $\F^W_0$, and $Z_0=1$, so $P^{(\mu)}(\Omega_0)=E[\ind_{\Omega_0}Z_0]=P(\Omega_0)=1$.

\emph{Step 2: comparison with passage times.} Fix $b>0$. We claim that
\begin{equation}\label{eq:3-5-09:a}
  \Omega_0\cap\{W_*<-b\}\;\subseteq\;\Omega_0\cap\{T_{-b}<\infty\}\;\subseteq\;\{W_*\le -b\}.
\end{equation}
For the first inclusion, let $\omega\in\Omega_0$ with $W_*(\omega)<-b$. There is $t>0$ with $W_t(\omega)<-b<0=W_0(\omega)$, and by the intermediate value theorem the continuous function $s\mapsto W_s(\omega)$ takes the value $-b$ at some $s\in(0,t)$; hence $T_{-b}(\omega)\le s<\infty$. For the second, let $\omega\in\Omega_0$ with $T_{-b}(\omega)<\infty$. The set $\{t\ge 0;\ W_t(\omega)=-b\}$ is closed and nonempty, so it contains its infimum $s\defeq T_{-b}(\omega)$, i.e.\ $W_s(\omega)=-b$; moreover $s>0$ because $W_0(\omega)=0\neq -b$. Therefore $W_*(\omega)\le W_s(\omega)=-b$.

\emph{Step 3: the tail of $-W_*$.} Since $\mu>0$, formula \bk{(5.13)} applied to the level $-b\neq 0$ gives
\[
  P^{(\mu)}[T_{-b}<\infty]=\exp\bigl[\mu(-b)-\abs{\mu(-b)}\bigr]=e^{-2\mu b},\qquad b>0 .
\]
Because $P^{(\mu)}(\Omega_0)=1$, \eqref{eq:3-5-09:a} yields
\begin{equation}\label{eq:3-5-09:b}
  P^{(\mu)}[W_*<-b]\;\le\; e^{-2\mu b}\;\le\; P^{(\mu)}[W_*\le -b],\qquad b>0 .
\end{equation}
Let $b>0$ and $0<b'<b$. Then $\{W_*\le -b\}\subseteq\{W_*<-b'\}$, so the first inequality in \eqref{eq:3-5-09:b}, used at $b'$, gives $P^{(\mu)}[W_*\le -b]\le e^{-2\mu b'}$; letting $b'\uparrow b$ and combining with the second inequality in \eqref{eq:3-5-09:b} we obtain
\begin{equation}\label{eq:3-5-09:c}
  P^{(\mu)}[-W_*\ge b]=e^{-2\mu b},\qquad b>0 .
\end{equation}

\emph{Step 4: conclusion.} The events $\{-W_*\ge n\}$ decrease to $\{W_*=-\infty\}$ as $n\to\infty$, so \eqref{eq:3-5-09:c} gives $P^{(\mu)}[W_*=-\infty]=\lim_{n\to\infty}e^{-2\mu n}=0$: the random variable $-W_*$ is $P^{(\mu)}$-a.s.\ finite. The events $\{-W_*\ge 1/n\}$ increase to $\{-W_*>0\}$, so $P^{(\mu)}[-W_*>0]=\lim_{n\to\infty}e^{-2\mu/n}=1$. Finally, for $b>0$ the events $\{-W_*\ge b+1/n\}$ increase to $\{-W_*>b\}$, whence $P^{(\mu)}[-W_*>b]=e^{-2\mu b}$ as well. Consequently
\[
  P^{(\mu)}[-W_*\le b]=1-e^{-2\mu b}=\int_0^b 2\mu e^{-2\mu x}\,dx,\qquad b>0,
  \qquad\text{and}\qquad P^{(\mu)}[-W_*\le 0]=0 .
\]
Thus the distribution function of $-W_*$ under $P^{(\mu)}$ is that of the exponential law with parameter $2\mu$, which is the assertion $P^{(\mu)}[-W_*\in db]=2\mu e^{-2\mu b}\,db$, $b>0$.

\rmk{The infimum over $t>0$ coincides on $\Omega_0$ with the infimum over $t\ge 0$, since $W_0=\lim_{t\downarrow 0}W_t\ge W_*$. As a consistency check, integrating the density \bk{(5.12)} for the level $-b$ over $t\in(0,\infty)$ must give the total mass $e^{-2\mu b}$ of \bk{(5.13)}; indeed, by \bk{(5.9)}, \bk{(5.10)} with $\alpha=\mu^2/2$,
$\int_0^\infty \frac{b}{\sqrt{2\pi t^3}}\exp\bigl[-\frac{(b+\mu t)^2}{2t}\bigr]dt
 =e^{-\mu b}E\bigl[e^{-\mu^2T_{-b}/2}\bigr]=e^{-\mu b}e^{-\mu b}$.}
\end{kex}

\begin{kex}{Exercise}{5.10}{}
Here $b\neq 0$ and $\mu\in\R$ are fixed, $T_b=\inf\{t\ge 0;\ W_t=b\}$ is the passage time of
Subsection~C, $P^{(\mu)}$ is the probability measure on $\F^W_\infty$ introduced before
\bk{(5.11)}, and $E^{(\mu)}$ denotes expectation with respect to $P^{(\mu)}$. Since
$P^{(\mu)}[T_b=\infty]$ may be positive (see \bk{(5.13)}), the random variable
$e^{-\alpha T_b}$ is understood to be $0$ on $\{T_b=\infty\}$, as is natural for $\alpha>0$:
\begin{equation}\label{eq:3-5-10:conv}
  E^{(\mu)} e^{-\alpha T_b}=E^{(\mu)}\bigl[e^{-\alpha T_b}\ind_{\{T_b<\infty\}}\bigr].
\end{equation}
For $\mu=0$ we have $P^{(0)}=P$ on $\F^W_\infty$ and the assertion is \bk{(5.10)}, so we
assume $\mu\neq0$. Because $W_0=0\neq b$ a.s.~$P$, we have $P[T_b=0]=0$, and then
$P^{(\mu)}[T_b=0]=E[\ind_{\{T_b=0\}}Z_0]=0$ as well, since $\{T_b=0\}\in\F^W_0$.

\emph{Step 1: the law of $T_b$ under $P^{(\mu)}$.} On the Borel subsets $B$ of $(0,\infty)$
we consider the two measures
\[
  \nu(B)\defeq P^{(\mu)}[T_b\in B],\qquad
  \lambda(B)\defeq\int_B \exp\Bigl(\mu b-\tfrac12\mu^2 s\Bigr)\,P[T_b\in ds].
\]
Both are finite: $\nu(B)\le 1$ and $\lambda(B)\le e^{\mu b}$. By \bk{(5.11)} and
$P^{(\mu)}[T_b=0]=P[T_b=0]=0$,
\[
  \nu\bigl((0,t]\bigr)=P^{(\mu)}[T_b\le t]=\lambda\bigl((0,t]\bigr),\qquad 0<t<\infty,
\]
and letting $t\uparrow\infty$ we obtain $\nu\bigl((0,\infty)\bigr)=\lambda\bigl((0,\infty)\bigr)$.
The intervals $(0,t]$, $0<t<\infty$, form a class closed under finite intersections which
generates $\B\bigl((0,\infty)\bigr)$, so Dynkin's $\pi$-$\lambda$ theorem gives $\nu=\lambda$.
Consequently, for every Borel function $f:(0,\infty)\to[0,\infty)$,
\begin{equation}\label{eq:3-5-10:law}
\begin{split}
  E^{(\mu)}\bigl[f(T_b)\ind_{\{0<T_b<\infty\}}\bigr]
  &=\int_{(0,\infty)}f\,d\nu
  =\int_{(0,\infty)}f(s)\exp\Bigl(\mu b-\tfrac12\mu^2 s\Bigr)\,P[T_b\in ds]\\
  &=e^{\mu b}E\Bigl[f(T_b)e^{-\frac12\mu^2T_b}\ind_{\{0<T_b<\infty\}}\Bigr];
\end{split}
\end{equation}
the first and the last equality are the change-of-variable formula for the image of
$P^{(\mu)}$, respectively of $P$, under $T_b$, restricted to $(0,\infty)$.

\emph{Step 2: the Laplace transform.} We take $f(s)=e^{-\alpha s}$ in
\eqref{eq:3-5-10:law}. On the left-hand side $P^{(\mu)}[T_b=0]=0$ and
\eqref{eq:3-5-10:conv} show that we obtain $E^{(\mu)}e^{-\alpha T_b}$. On the right-hand
side $P[0<T_b<\infty]=1$, because $P[T_b=0]=0$ and $P[T_b<\infty]=1$ (the density in
\bk{(5.9)} integrates to one; equivalently, let $\alpha\downarrow0$ in \bk{(5.10)}).
Therefore
\[
  E^{(\mu)}e^{-\alpha T_b}=e^{\mu b}\,E\exp\Bigl[-\Bigl(\alpha+\tfrac12\mu^2\Bigr)T_b\Bigr].
\]
Since $\alpha+\frac12\mu^2>0$, we may apply \bk{(5.10)} with $\alpha+\frac12\mu^2$ in place of
$\alpha$:
\[
  E^{(\mu)}e^{-\alpha T_b}
  =e^{\mu b}\exp\Bigl(-\abs{b}\sqrt{2\alpha+\mu^2}\Bigr)
  =\exp\Bigl(\mu b-\abs{b}\sqrt{\mu^2+2\alpha}\Bigr),\qquad \alpha>0.
\]

\rmk{As a check, letting $\alpha\downarrow0$ and using monotone convergence in
\eqref{eq:3-5-10:conv}, we recover $P^{(\mu)}[T_b<\infty]=\exp(\mu b-\abs{\mu b})$, which is
\bk{(5.13)}. The formula also says that the (possibly defective) density
$h(\,\cdot\,;b,\mu)$ of \bk{(5.12)} has Laplace transform
$\int_0^\infty e^{-\alpha t}h(t;b,\mu)\,dt=\exp\bigl(\mu b-\abs{b}\sqrt{\mu^2+2\alpha}\bigr)$;
the exponent is additive in $b$ on each half-line, in agreement with the convolution identity
of \bk{Problem 5.8}.}
\end{kex}

\begin{kex}{Exercise}{5.11}{}
We work in the setting of Subsection~C: $W$ is the coordinate mapping process on
$\Omega=C[0,\infty)$, $P$ is Wiener measure on $\F^W_\infty$, and for $\mu\in\R$ the measure
$P^{(\mu)}$ is the unique probability measure on $\F^W_\infty$ with
$P^{(\mu)}(A)=E[\ind_A\exp(\mu W_t-\frac12\mu^2t)]$ for $A\in\F^W_t$, $0\le t<\infty$; under
$P^{(\mu)}$ the process $\{W_t-\mu t,\ \F^W_t;\ 0\le t<\infty\}$ is a Brownian motion
(\bk{Corollary 5.2}). We write
\[
  Z_t\defeq\exp\bigl(\nu W_t-\tfrac12\nu^2t\bigr),\qquad R\defeq R_c=\inf\{t\ge0;\ Z_t=c\},
  \qquad b\defeq\frac{\log c}{\nu}>0,
\]
with $\inf\emptyset=\infty$. The process $Z$ is continuous and adapted to $\{\F^W_t\}$, and
$\{c\}$ is closed, so $R$ is indeed a stopping time of $\{\F^W_t\}$ (\bk{Problem 1.2.7}).
Every path of $Z$ is continuous; hence $Z_R=c$ on $\{R<\infty\}$, and on the event
$\{W_0=0\}=\{Z_0=1\}$ the intermediate value theorem gives $Z_s<c$ for $0\le s<R$. Thus
\begin{equation}\label{eq:3-5-11:bd}
  0<Z_{t\wedge R}\le c,\quad 0\le t<\infty,\qquad\text{on }\{W_0=0\}.
\end{equation}

\textbf{(i) The probability $P[R_c<\infty]$.} Under $P$, the process
$\{Z_t,\F^W_t;\ 0\le t<\infty\}$ is a continuous martingale: for $0\le s<t$ the increment
$W_t-W_s$ is independent of $\F^W_s$ and normal with mean zero and variance $t-s$, so that
$E[Z_t\,|\,\F^W_s]=Z_s\,e^{-\frac12\nu^2(t-s)}E[e^{\nu(W_t-W_s)}]=Z_s$ a.s. For fixed $t$ the
stopping time $t\wedge R$ is bounded by $t$, and the optional sampling theorem for bounded
times (\bk{Problem 1.3.23}\,(i), with $S=0$ and $T=t\wedge R$) yields
\begin{equation}\label{eq:3-5-11:os}
  E[Z_{t\wedge R}]=E[Z_0]=1,\qquad 0\le t<\infty .
\end{equation}
We now let $t\to\infty$. On $\{R<\infty\}$ we have $Z_{t\wedge R}=Z_R=c$ for all $t\ge R$.
On $\{R=\infty\}$ we have $Z_{t\wedge R}=Z_t$, and by the strong law of large numbers
\bk{(2.9.2)} of \bk{Problem 2.9.3},
\[
  \nu W_t-\tfrac12\nu^2t=t\Bigl(\nu\,\frac{W_t}{t}-\frac{\nu^2}{2}\Bigr)\longrightarrow-\infty,
  \quad\text{i.e.}\quad Z_t\longrightarrow0,\qquad\text{a.s. }P,
\]
because $\nu\neq0$. Consequently
$\lim_{t\to\infty}Z_{t\wedge R}=c\,\ind_{\{R<\infty\}}$ a.s.~$P$. Since $P[W_0=0]=1$, the
bound \eqref{eq:3-5-11:bd} holds a.s.~$P$, and the bounded convergence theorem applied to
\eqref{eq:3-5-11:os} gives $1=c\,P[R<\infty]$, that is,
\[
  P[R_c<\infty]=\frac1c .
\]

\textbf{(ii) The expectation $E^{(\nu)}R_c$.} Let
$\wt W_t\defeq W_t-\nu t$, so that $\{\wt W_t,\F^W_t;\ 0\le t<\infty\}$ is a Brownian motion
under $P^{(\nu)}$; in particular $P^{(\nu)}[W_0=0]=P^{(\nu)}[\wt W_0=0]=1$. Since
$\nu W_t-\frac12\nu^2t=\nu\bigl(\wt W_t+\frac12\nu t\bigr)$, we have
\begin{equation}\label{eq:3-5-11:Y}
  R=\inf\{t\ge0;\ Y_t=b\},\qquad\text{where}\quad Y_t\defeq\wt W_t+\tfrac12\nu t
  =\tfrac1\nu\log Z_t ,
\end{equation}
and \eqref{eq:3-5-11:bd} reads
\begin{equation}\label{eq:3-5-11:Ybd}
  Y_{t\wedge R}\le b,\quad 0\le t<\infty,\qquad\text{a.s. }P^{(\nu)};
  \qquad Y_R=b\ \text{ on }\{R<\infty\}.
\end{equation}

\emph{Step 1: $E^{(\nu)}R\le 2b/\nu$.} Under $P^{(\nu)}$ the Brownian motion $\wt W$ is a
continuous martingale with respect to $\{\F^W_t\}$, so \bk{Problem 1.3.23}\,(i) gives
$E^{(\nu)}[\wt W_{t\wedge R}]=E^{(\nu)}[\wt W_0]=0$. From the definition of $Y$ and
\eqref{eq:3-5-11:Ybd},
\[
  \frac\nu2\,E^{(\nu)}[t\wedge R]=E^{(\nu)}[Y_{t\wedge R}]-E^{(\nu)}[\wt W_{t\wedge R}]
  =E^{(\nu)}[Y_{t\wedge R}]\le b,\qquad 0\le t<\infty
\]
(all three random variables are integrable, the first because it is bounded by $t$). Letting
$t\uparrow\infty$, the monotone convergence theorem yields
\begin{equation}\label{eq:3-5-11:ub}
  E^{(\nu)}R\le\frac{2b}{\nu}<\infty;\qquad\text{in particular}\quad P^{(\nu)}[R<\infty]=1 .
\end{equation}

\emph{Step 2: equality.} Under $P^{(\nu)}$ the process
$\{\wt W_t^2-t,\F^W_t;\ 0\le t<\infty\}$ is also a continuous martingale (for $s<t$ write
$\wt W_t^2=\wt W_s^2+2\wt W_s(\wt W_t-\wt W_s)+(\wt W_t-\wt W_s)^2$ and use the independence of
$\wt W_t-\wt W_s$ from $\F^W_s$), so \bk{Problem 1.3.23}\,(i) and \eqref{eq:3-5-11:ub} give
\[
  E^{(\nu)}\bigl[\wt W_{t\wedge R}^2\bigr]=E^{(\nu)}[t\wedge R]\le\frac{2b}{\nu},
  \qquad 0\le t<\infty .
\]
The family $\{\wt W_{t\wedge R}\}_{0\le t<\infty}$ is therefore bounded in $L^2(P^{(\nu)})$,
hence uniformly integrable. By \eqref{eq:3-5-11:ub}, \eqref{eq:3-5-11:Ybd} and the
continuity of the paths,
\[
  \lim_{t\to\infty}\wt W_{t\wedge R}=\wt W_R=Y_R-\tfrac12\nu R=b-\tfrac12\nu R,
  \qquad\text{a.s. }P^{(\nu)} .
\]
Almost sure convergence together with uniform integrability implies convergence of the
expectations, and we obtain
\[
  0=\lim_{t\to\infty}E^{(\nu)}[\wt W_{t\wedge R}]=E^{(\nu)}\bigl[b-\tfrac12\nu R\bigr]
  =b-\frac\nu2\,E^{(\nu)}R .
\]
Hence
\[
  E^{(\nu)}R_c=\frac{2b}{\nu}=\frac{2\log c}{\nu^2}.
\]

\rmk{Both results agree with the formulas of Subsection~C. By \eqref{eq:3-5-11:Y}, under
$P^{(\nu)}$ the time $R_c$ is the passage time to the level $b=(\log c)/\nu$ of the Brownian
motion $Y$ with drift $\nu/2$, so it has under $P^{(\nu)}$ the law which $T_b$ has under
$P^{(\nu/2)}$; likewise, under $P$ it has the law which $T_b$ has under $P^{(-\nu/2)}$,
because $Y_t=W_t-\frac12\nu t$. Now \bk{(5.13)} with $\mu=-\nu/2$ gives
$P^{(-\nu/2)}[T_b<\infty]=e^{-\nu b}=1/c$, and the density \bk{(5.12)} with $\mu=\nu/2$ has
mean $b/\mu=2b/\nu$: indeed, differentiating the Laplace transform
$\exp\bigl(\mu b-b\sqrt{\mu^2+2\alpha}\bigr)$ of \bk{Exercise 5.10} at $\alpha=0+$ gives
$-b/\mu$. Finally, since $\{R\le t\}\in\F^W_t\cap\F^W_R$ and $Z_{t\wedge R}=c$ on this event,
the argument of \bk{(5.11)} gives $P^{(\nu)}[R\le t]=c\,P[R\le t]$ for every $t$; letting
$t\to\infty$ we see that the two assertions $P[R<\infty]=1/c$ and $P^{(\nu)}[R<\infty]=1$ are
equivalent.}
\end{kex}

\begin{kex}{Exercise}{5.18}{}
\emph{Setting.} We work under the standing assumptions of Subsection~A:
$W=\{W_t,\F_t;\,0\le t<\infty\}$ is a one-dimensional Brownian motion with $P[W_0=0]=1$,
the filtration $\{\F_t\}$ satisfies the usual conditions, and the process of the exercise is
the restriction of $W$ to $[0,1]$. We write $Y_t\defeq t+W_t^2$, a continuous, adapted
process, and we read the definition of $T$ as
\[
  T=1\wedge\inf\{t\ge0;\ Y_t=1\},
\]
which agrees with the printed one whenever the set $\{0\le t\le1;\ Y_t=1\}$ is not empty and
equals $1$ otherwise. By \bk{Problem 1.2.7} (the hitting time of the closed set $\{1\}$ by a
continuous, adapted process) and \bk{Lemma 1.2.9}, $T$ is a stopping time of $\{\F_t\}$ with
$0\le T\le1$. We extend $X$ to $[0,\infty)$ by $X_t\defeq0$ for $t\ge1$, so that
\[
  X_t=-\frac{2}{(1-t)^2}\,W_t\,\ind_{\{t\le T\}}\,\ind_{[0,1)}(t),\qquad 0\le t<\infty .
\]
The process $\ind_{\{t\le T\}}$ is adapted, because $\{t\le T\}=\{T<t\}^c\in\F_t$, and
left-continuous; $W$ is adapted and continuous; hence both are progressively measurable
(\bk{Proposition 1.1.13}), and so is $X$, their product with a Borel function of $t$. For
a progressively measurable process $H$ with $P[\int_0^t H_u^2\,du<\infty]=1$ for every
$t$, we write $I_t(H)=\int_0^t H_u\,dW_u$ for its stochastic integral, a continuous local
martingale.

We shall use the event
\[
  \Omega_0\defeq\{W_0=0\}\cap\{W_1\neq0\},\qquad P(\Omega_0)=1,
\]
which has full probability because $W_1$ is normally distributed with variance $1$.

\medskip
\textbf{(i)} Fix $\omega\in\Omega_0$. The continuous function $t\mapsto Y_t(\omega)$
satisfies $Y_0(\omega)=0<1$ and $Y_1(\omega)=1+W_1^2(\omega)>1$, so by the intermediate
value theorem the closed set $\{t\in[0,1];\ Y_t(\omega)=1\}$ is not empty and contains
neither $0$ nor $1$. Its smallest element is $T(\omega)$; therefore
\begin{equation}\label{eq:3-5-18:T}
  0<T<1,\qquad T+W_T^2=1,\qquad t+W_t^2\le 1\ \text{ for } 0\le t\le T,
  \qquad\text{on }\Omega_0 .
\end{equation}
For the last assertion: if $Y_t(\omega)>1$ for some $t\le T(\omega)$, the intermediate value
theorem would give $u\in(0,t)$ with $Y_u(\omega)=1$, contradicting the minimality of
$T(\omega)$. In particular $P[T<1]=1$. From \eqref{eq:3-5-18:T} we have $W_t^2\le1-t$ for
$t\le T$, and thus on $\Omega_0$
\begin{equation}\label{eq:3-5-18:int}
  \int_0^1X_t^2\,dt=\int_0^T\frac{4W_t^2}{(1-t)^4}\,dt\le\int_0^T\frac{4}{(1-t)^3}\,dt
  =\frac{2}{(1-T)^2}-2<\infty .
\end{equation}
Since $X_t=0$ for $t\ge1$, the process $X$ satisfies \bk{(5.1)}; hence $I(X)$ and the
process $Z(X)$ of \bk{(5.2)} are well defined on $[0,\infty)$.

\medskip
\textbf{(ii)} The function $(t,x)\mapsto x^2/(1-t)^2$ is singular at $t=1$, so we apply
It\^o's rule up to a fixed time $s\in(0,1)$ and then let $s\uparrow1$.

\emph{Step 1: It\^o's rule.} Fix $s\in(0,1)$ and let $A_t\defeq(1-t\wedge s)^{-1}$,
$0\le t<\infty$. This is a deterministic, continuous, nondecreasing function with $A_0=1$ and
$dA_t=(1-t)^{-2}\ind_{[0,s]}(t)\,dt$. We apply \bk{Theorem 3.6} with $d=2$ to the continuous
semimartingale $(A_t,W_t)$---local martingale part $(0,W_t)$, bounded variation part
$(A_t-1,0)$, initial value $(1,W_0)$---and to the function $f(a,x)=a^2x^2$, which is of
class $C^2$ and does not depend on the time variable. Since $\qv{W}_t=t$ and $W_0=0$ a.s.,
we obtain, a.s.,
\[
  A_t^2W_t^2=\int_0^t2A_uW_u^2\,dA_u+\int_0^t2A_u^2W_u\,dW_u+\int_0^tA_u^2\,du,
  \qquad 0\le t<\infty .
\]
With $H^s_u\defeq 2W_u(1-u\wedge s)^{-2}$, a continuous adapted process, this reads for
$0\le t\le s$:
\begin{equation}\label{eq:3-5-18:ito}
  \frac{W_t^2}{(1-t)^2}=2\int_0^t\frac{W_u^2}{(1-u)^3}\,du+I_t(H^s)+\frac{1}{1-t}-1 ,
\end{equation}
because $\int_0^t(1-u)^{-2}\,du=(1-t)^{-1}-1$. The exceptional null set does not depend on
$t$, so \eqref{eq:3-5-18:ito} may be evaluated at the random time $t=T\wedge s\le s$.

\emph{Step 2: the stochastic integral at time $T\wedge s$.} By \bk{(2.24)}, which remains
valid for integrals with respect to continuous local martingales (see the paragraph preceding
\bk{Proposition 2.24}), applied to the stopping time $T\wedge s\le1$ and to $t=1$,
\[
  I_{T\wedge s}(H^s)=I_{1\wedge(T\wedge s)}(H^s)
  =I_1\bigl(H^s\ind_{\{\cdot\,\le T\wedge s\}}\bigr)\quad\text{a.s.}
\]
For $u\le s$ we have $u<1$ and $u\wedge s=u$, so that
$H^s_u\ind_{\{u\le T\wedge s\}}=2W_u(1-u)^{-2}\ind_{\{u\le T\}}\ind_{\{u\le s\}}
=-X_u\ind_{\{u\le s\}}$ for every $u\ge0$. By linearity of the integral and \bk{(2.24)} once
more, now with the constant stopping time $s$,
\[
  I_{T\wedge s}(H^s)=-I_1\bigl(X\ind_{\{\cdot\,\le s\}}\bigr)=-I_{1\wedge s}(X)=-I_s(X)
  \quad\text{a.s.}
\]
Consequently \eqref{eq:3-5-18:ito} at $t=T\wedge s$ becomes
\begin{equation}\label{eq:3-5-18:stopped}
  \frac{W_{T\wedge s}^2}{(1-T\wedge s)^2}
  =2\int_0^{T\wedge s}\frac{W_u^2}{(1-u)^3}\,du-I_s(X)+\frac{1}{1-T\wedge s}-1
  \quad\text{a.s.}
\end{equation}

\emph{Step 3: $s\uparrow1$.} Let $s_k=1-1/k$, $k\ge2$. There is an event
$\Omega_1\subseteq\Omega_0$ of probability one on which \eqref{eq:3-5-18:stopped} holds
for every $s=s_k$ and $t\mapsto I_t(X)$ is continuous. Fix $\omega\in\Omega_1$. Since
$T(\omega)<1$, we have $T\wedge s_k=T$ for all sufficiently large $k$, and for these $k$ the
left-hand side of \eqref{eq:3-5-18:stopped} is $W_T^2/(1-T)^2=1/(1-T)$ by
\eqref{eq:3-5-18:T}. The terms $1/(1-T)$ cancel, and we are left with
\[
  I_{s_k}(X)=-1+2\int_0^T\frac{W_u^2}{(1-u)^3}\,du\qquad\text{for all large }k .
\]
Letting $k\to\infty$ and using the continuity of $I(X)$ at $t=1$, we conclude that
\begin{equation}\label{eq:3-5-18:I}
  \int_0^1X_t\,dW_t=-1+2\int_0^T\frac{W_t^2}{(1-t)^3}\,dt\quad\text{a.s.},
\end{equation}
where the integral on the right is finite on $\Omega_0$ because $T<1$ there. Subtracting
$\frac12\int_0^1X_t^2\,dt=2\int_0^TW_t^2(1-t)^{-4}\,dt$, which is finite by
\eqref{eq:3-5-18:int}, we obtain
\[
  \int_0^1X_t\,dW_t-\frac12\int_0^1X_t^2\,dt
  =-1-2\int_0^T\left[\frac{1}{(1-t)^4}-\frac{1}{(1-t)^3}\right]W_t^2\,dt\le-1
  \quad\text{a.s.},
\]
the inequality holding because $(1-t)^{-4}\ge(1-t)^{-3}$ for $0\le t<1$ and $T<1$.

\medskip
\textbf{(iii)} \emph{$Z(X)$ is not a martingale.} The process $Z(X)$ of \bk{(5.2)} is a
continuous local martingale with $Z_0(X)=1$ and a supermartingale (see the discussion
preceding \bk{(5.17)}). By part (ii),
\[
  Z_1(X)=\exp\Bigl[\int_0^1X_t\,dW_t-\frac12\int_0^1X_t^2\,dt\Bigr]\le e^{-1}
  \quad\text{a.s.},
\]
so that $EZ_1(X)\le e^{-1}<1=EZ_0(X)$. A martingale has constant expectation; therefore
$\{Z_t(X),\F_t;\,0\le t\le1\}$ is not a martingale.

\emph{The stopped processes are martingales.} Fix $n\ge1$; then $0\le\sigma_n<1$ and
$(1-\sigma_n)^{-2}=n$. Let $X^{(n)}_t\defeq X_t\ind_{\{t\le\sigma_n\}}$, $0\le t<\infty$,
a progressively measurable process. On $\Omega_0$ we have $W_t^2\le1-t$ for $t\le T$ by
\eqref{eq:3-5-18:T}, and thus
\[
  \frac12\int_0^\infty\bigl(X^{(n)}_t\bigr)^2\,dt
  =2\int_0^{T\wedge\sigma_n}\frac{W_t^2}{(1-t)^4}\,dt
  \le2\int_0^{\sigma_n}\frac{dt}{(1-t)^3}=\frac{1}{(1-\sigma_n)^2}-1=n-1 .
\]
Hence $X^{(n)}$ satisfies \bk{(5.1)} and
$E\exp\bigl(\frac12\int_0^{t}(X^{(n)}_u)^2\,du\bigr)\le e^{n-1}<\infty$ for every
$0\le t<\infty$, which is the Novikov condition \bk{(5.20)}. By \bk{Corollary 5.13},
$\{Z_t(X^{(n)}),\F_t;\,0\le t<\infty\}$ is a martingale. On the other hand, \bk{(2.24)} with
the constant stopping time $\sigma_n$ gives, for each $t$,
$I_{t\wedge\sigma_n}(X)=I_t(X^{(n)})$ a.s., and trivially
$\int_0^{t\wedge\sigma_n}X_u^2\,du=\int_0^t(X^{(n)}_u)^2\,du$; therefore
\[
  Z_{t\wedge\sigma_n}(X)=Z_t\bigl(X^{(n)}\bigr)\quad\text{a.s., for each }0\le t<\infty .
\]
The random variable $Z_{t\wedge\sigma_n}(X)$ is $\F_t$-measurable, and a modification of a
martingale by adapted random variables is again a martingale (integrability and the
conditional expectations are unaffected by changes on null sets). Thus
$\{Z_{t\wedge\sigma_n}(X),\F_t;\,0\le t<\infty\}$ is a martingale, and so is its restriction
to $0\le t\le1$.

\rmk{Since $Z_t(X)=Z_{t\wedge\sigma_n}(X)$ for $t\le\sigma_n$ and $\sigma_n\uparrow1$, the
above shows $EZ_t(X)=1$ for every $t<1$: the process $\{Z_t(X);\,0\le t<1\}$ is a
martingale, and the martingale property is lost exactly at the endpoint $t=1$, where
$EZ_1(X)\le e^{-1}$. In particular the family $\{Z_{\sigma_n}(X)\}_{n\ge1}$, which converges
a.s.\ to $Z_1(X)$, is not uniformly integrable, and by \bk{Corollary 5.13} the Novikov
condition must fail on $[0,1]$: $E\exp\bigl(\frac12\int_0^1X_t^2\,dt\bigr)=\infty$.}
\end{kex}

\begin{kex}{Exercise}{5.19}{}
\emph{Setting.} As throughout \bk{Section 5}, $W=(W^{(1)},\dots,W^{(d)})$ is $d$-dimensional, $d\ge1$. Since $W$ is adapted to $\{\F_t\}$ and $W_t-W_s$ is independent of $\F_s\supseteq\F^W_s$ for $0\le s<t$, $W$ is also a Brownian motion relative to its own filtration $\{\F^W_t\}$. Following \bk{Definition 2.7.2}, we write
\[
  \cN\defeq\{F\subseteq\Omega;\ \exists\,G\in\F^W_\infty\text{ with }F\subseteq G,\ P(G)=0\}
\]
for the collection of null sets used in the augmentation, so that $\F_t=\sigma(\F^W_t\cup\cN)$ for $0\le t<\infty$; in particular $\cN\subseteq\F_0\subseteq\F$, and every member of $\cN$ has $P$-measure zero. ``Mutually absolutely continuous with respect to $P$'' refers to the restriction of $P$ to $\F_T$. We record three facts.

\textbf{(a)} \emph{Every $F\in\F_0$ has $P(F)\in\{0,1\}$.} Indeed, by \bk{Problem 2.7.3} (with $\G=\F^W_0$) there is $G\in\F^W_0=\sigma(W_0)$ with $F\triangle G\in\cN$, so $P(F)=P(G)$; and $G=\{W_0\in\Gamma\}$ for some $\Gamma\in\B(\R^d)$, so $P(G)=\ind_\Gamma(0)$ because $P[W_0=0]=1$.

\textbf{(b)} \emph{Every $F\in\F_\infty$ with $P(F)=0$ belongs to $\cN\subseteq\F_0$.} Indeed, $\F_\infty=\sigma(\F^W_\infty\cup\cN)$ by \bk{Problem 2.7.5}, so by \bk{Problem 2.7.3} (with $\G=\F^W_\infty$) there is $G\in\F^W_\infty$ with $F\triangle G\in\cN$; then $P(G)=0$, i.e.\ $G\in\cN$, and $F\subseteq G\cup(F\triangle G)$ is in $\cN$.

\textbf{(c)} \emph{On the probability space $(\Omega,\F_\infty,P|_{\F_\infty})$ the filtration $\{\F_t\}$ satisfies the usual conditions of \bk{Definition 1.2.25}.} Indeed, $W$ is a Brownian motion relative to $\{\F^W_t\}$, hence a strong Markov process (\bk{Theorem 2.6.15}), so its augmented filtration $\{\F_t\}$ is right-continuous by \bk{Proposition 2.7.7}; and by (b) every $P$-negligible event of $\F_\infty$ lies in $\F_0$.

The exercise allows an arbitrary $\sigma$-field $\F\supseteq\F_\infty$, which may contain $P$-null events outside $\cN$; relative to such an $\F$ the second requirement of \bk{Definition 1.2.25} fails, although it is a standing assumption of \bk{Section 5}. However, everything that enters the statement lives on $\F_\infty$: $W$ is adapted to $\{\F_t\}$, each $\wt P_T$ is a measure on $\F_T\subseteq\F_\infty$, and the assertions to be proved (measurability and adaptedness of $X$, \bk{(5.1)}, the martingale property of $Z(X)$ relative to $\{\F_t\}$, and \bk{(5.4)}) involve only events of $\F_\infty$ and their $P$-probabilities. \emph{We therefore work throughout on $(\Omega,\F_\infty,P|_{\F_\infty})$}, still writing $P$ and $E$; on this space $W$ is a Brownian motion relative to $\{\F_t\}$, (c) holds, and the stochastic integrals below, in particular the one in \bk{(5.2)}, are those constructed on it. (If $\{\F_t\}$ does satisfy the usual conditions relative to $\F$, then every $P$-null event of $\F$ is in $\F_0\subseteq\F_\infty$, and passing to $\F_\infty$ changes nothing.) A process which is measurable with respect to $\B([0,\infty))\otimes\F_\infty$ is a fortiori measurable with respect to $\B([0,\infty))\otimes\F$. Finally, the components $W^{(j)}$ belong to $\cM_2^c$ and $\qv{W^{(i)},W^{(j)}}_t=\delta_{ij}t$ (\bk{Problem 2.5.5}).

\medskip
\emph{Step 1: the density martingale.} For each $t\ge0$ let $\xi_t\ge0$ be an $\F_t$-measurable version of the Radon--Nikod\'ym derivative of $\wt P_t$ with respect to $P|_{\F_t}$, which exists because $\wt P_t\ll P$ on $\F_t$:
\begin{equation}\label{eq:3-5-19:a}
  \wt P_t(A)=E[\ind_A\xi_t];\qquad A\in\F_t .
\end{equation}
Then $E\xi_t=1$, and $P[\xi_t>0]=1$: the event $\{\xi_t=0\}\in\F_t$ has $\wt P_t$-measure $E[\ind_{\{\xi_t=0\}}\xi_t]=0$, hence $P$-measure zero because $P\ll\wt P_t$ on $\F_t$. For $0\le t\le T$ and $A\in\F_t\subseteq\F_T$, \eqref{eq:3-5-19:a} and \bk{(5.5)} give $E[\ind_A\xi_T]=\wt P_T(A)=\wt P_t(A)=E[\ind_A\xi_t]$; thus $\{\xi_t,\F_t;\ 0\le t<\infty\}$ is a martingale. Since $t\mapsto E\xi_t$ is constant, hence right-continuous, and the usual conditions hold by (c), \bk{Theorem 1.3.13} provides a modification $Z'=\{Z'_t,\F_t;\ 0\le t<\infty\}$ of $\xi$ which is adapted, has RCLL paths, and is a martingale. For each fixed $t$ we have $Z'_t=\xi_t$ a.s.; consequently $Z'_t$ may replace $\xi_t$ in \eqref{eq:3-5-19:a}, and $P[Z'_t>0]=1$. Moreover, $Z'_0$ is $\F_0$-measurable, so by (a) it is a.s.\ equal to a constant, and this constant is $EZ'_0=1$.

\emph{Step 2: representation.} The process $M'_t\defeq Z'_t-Z'_0$ is an RCLL martingale of $\{\F_t\}$ with $M'_0=0$ (the constant process $Z'_0$ is a martingale, $Z'_0$ being integrable and $\F_0$-measurable), in particular an RCLL local martingale (\bk{Remark 1.5.16}). By \bk{Problem 4.16}, $M'$ is a.s.\ continuous, and there are progressively measurable processes $Y^{(j)}=\{Y^{(j)}_t,\F_t;\ 0\le t<\infty\}$, $1\le j\le d$, such that, a.s.,
\begin{equation}\label{eq:3-5-19:b}
  \int_0^T\bigl(Y^{(j)}_t\bigr)^2\,dt<\infty\quad(1\le j\le d,\ 0\le T<\infty),\qquad
  M'_t=\sum_{j=1}^d\int_0^tY^{(j)}_s\,dW^{(j)}_s\quad(0\le t<\infty).
\end{equation}

\emph{Step 3: a version with good paths.} Let $C$ be the set of $\omega$ for which $t\mapsto Z'_t(\omega)$ is continuous on $[0,\infty)$, $Z'_0(\omega)=1$, and $Z'_q(\omega)\ge0$ for all rational $q\ge0$. An RCLL function is continuous on $[0,n]$ if and only if it is uniformly continuous on $\Q\cap[0,n]$ (by right-continuity, a modulus of continuity on the rationals of $[0,n]$ is inherited by all pairs of points of $[0,n)$, and then the left limit at $n$ equals the value at $n$, the latter being a rational point). Therefore
\[
  C=\{Z'_0=1\}\cap\bigcap_{q\in\Q\cap[0,\infty)}\{Z'_q\ge0\}\cap\bigcap_{n\ge1}\,\bigcap_{k\ge1}\,\bigcup_{m\ge1}
     \bigcap_{\substack{r,s\in\Q\cap[0,n]\\ \abs{r-s}<1/m}}\bigl\{\abs{Z'_r-Z'_s}\le 1/k\bigr\}\in\F_\infty,
\]
and $P(C)=1$ by Steps 1 and 2. By (b), $\Omega\setminus C\in\F_0$. Hence
\[
  Z_t\defeq Z'_t\ind_C+\ind_{\Omega\setminus C};\qquad 0\le t<\infty,
\]
is an adapted process, indistinguishable from $Z'$, with $Z_0\equiv1$ and such that \emph{every} path is continuous and nonnegative. It is a martingale, $P[Z_T>0]=1$ for each $T$, and by \eqref{eq:3-5-19:a}
\begin{equation}\label{eq:3-5-19:c}
  \wt P_T(A)=E[\ind_AZ_T];\qquad A\in\F_T,\quad 0\le T<\infty .
\end{equation}
Finally $M\defeq Z-1$ is a continuous martingale with $M_0=0$, so $M\in\cM^{c,\loc}$, and it is indistinguishable from $M'$; by \eqref{eq:3-5-19:b}, $Y^{(j)}\in\cP^*$ (\bk{Definition 2.22}, with integrator $W^{(j)}$) and
\begin{equation}\label{eq:3-5-19:d}
  M_t=Z_t-1=\sum_{j=1}^dI^{W^{(j)}}_t\bigl(Y^{(j)}\bigr);\qquad 0\le t<\infty,\quad\text{a.s.}
\end{equation}

\emph{Step 4: $Z$ never vanishes.} Let $\tau\defeq\inf\{t\ge0;\ Z_t=0\}$. Since $\{0\}$ is closed and the paths of $Z$ are continuous, $\tau$ is a stopping time of $\{\F_t\}$ (\bk{Problem 1.2.7}), and $Z_\tau=0$ on $\{\tau<\infty\}$. Fix $T\in[0,\infty)$. The process $\{Z_{t\wedge T},\F_t;\ 0\le t<\infty\}$ is a continuous martingale with last element $Z_T$. The optional sampling theorem (\bk{Theorem 1.3.22}, applied to this martingale and to its negative, with the stopping times $\tau\wedge T\le\infty$) gives
\[
  E\bigl[Z_T\mid\F_{\tau\wedge T}\bigr]=Z_{\tau\wedge T}\quad\text{a.s.}
\]
By \bk{Lemma 1.2.16}, $\{\tau\le T\}\in\F_\tau\cap\F_T=\F_{\tau\wedge T}$, and on this event $Z_{\tau\wedge T}=Z_\tau=0$. Therefore
\[
  E\bigl[Z_T\ind_{\{\tau\le T\}}\bigr]=E\bigl[Z_{\tau\wedge T}\ind_{\{\tau\le T\}}\bigr]=0 ,
\]
and since $Z_T>0$ a.s., $P[\tau\le T]=0$. Letting $T\to\infty$ through the integers we obtain $P(\Omega_+)=1$, where
\[
  \Omega_+\defeq\{\tau=\infty\}=\{Z_t>0\ \text{for all } t\ge0\}.
\]

\emph{Step 5: the process $X$ and condition \bk{(5.1)}.} Define, for $1\le j\le d$,
\[
  X^{(j)}_t(\omega)\defeq\frac{Y^{(j)}_t(\omega)}{Z_t(\omega)}\,\ind_{\{Z_t(\omega)>0\}}
  \quad\bigl(\text{understood as }0\text{ where }Z_t(\omega)=0\bigr);\qquad 0\le t<\infty .
\]
The process $Z$ is adapted with continuous paths, hence progressively measurable (\bk{Proposition 1.1.13}); $Y^{(j)}$ is progressively measurable; and $(y,z)\mapsto(y/z)\ind_{\{z>0\}}$ is a Borel function on $\R^2$. Thus $X^{(j)}$ is progressively measurable, in particular measurable and adapted. For $\omega\in\Omega_+$ and $T<\infty$ the number $m_T(\omega)\defeq\min_{0\le t\le T}Z_t(\omega)$ is strictly positive, so that
\[
  \int_0^T\bigl(X^{(j)}_t\bigr)^2\,dt\le\frac1{m_T^2}\int_0^T\bigl(Y^{(j)}_t\bigr)^2\,dt ,
\]
which is finite a.s.\ by \eqref{eq:3-5-19:b}. This is \bk{(5.1)}, and $X^{(j)}\in\cP^*$. Moreover,
\begin{equation}\label{eq:3-5-19:e}
  Y^{(j)}_t(\omega)=Z_t(\omega)X^{(j)}_t(\omega);\qquad 0\le t<\infty,\quad 1\le j\le d,\quad\omega\in\Omega_+ .
\end{equation}

\emph{Step 6: $Z=Z(X)$.} Let
\[
  N\defeq\sum_{j=1}^dI^{W^{(j)}}\bigl(X^{(j)}\bigr)\in\cM^{c,\loc},\qquad A\defeq\qv{N},\qquad
  U_t\defeq\exp\bigl[-N_t+\tfrac12A_t\bigr].
\]
The formula \bk{(2.26)} for the cross-variation of two stochastic integrals remains valid for integrators in $\cM^{c,\loc}$ and integrands in $\cP^*$ (see the paragraph preceding \bk{Proposition 2.24}). Together with the bilinearity of the cross-variation, $\qv{W^{(i)},W^{(j)}}_t=\delta_{ij}t$ and \eqref{eq:3-5-19:d}, it yields, a.s.\ for all $t\ge0$,
\begin{equation}\label{eq:3-5-19:f}
  \qv{M}_t=\int_0^t\norm{Y_s}^2\,ds,\qquad
  \qv{M,N}_t=\int_0^t\sum_{j=1}^dY^{(j)}_sX^{(j)}_s\,ds,\qquad
  A_t=\qv{N}_t=\int_0^t\norm{X_s}^2\,ds .
\end{equation}
In particular, by \bk{(5.2)}, the process $Z(X)$ is indistinguishable from $\exp[N-\frac12A]=1/U$.

We apply It\^o's rule (\bk{Theorem 3.6}) to the function $f(z,x,a)\defeq z\,e^{-x+a/2}$, which is of class $C^2$ on $\R^3$, and to the three-dimensional process $(Z,N,A)=(1,0,0)+(M,N,0)+(0,0,A)$: here $M,N\in\cM^{c,\loc}$, and $A$ is adapted, continuous and nondecreasing with $A_0=0$ (\bk{Problem 1.5.17}). We have
\[
  f_z=e^{-x+a/2},\quad f_x=-f,\quad f_a=\tfrac12f,\quad f_{zz}=0,\quad f_{zx}=-e^{-x+a/2},\quad f_{xx}=f ,
\]
and only $M$ and $N$ contribute to the second-order terms. Since $f(Z_t,N_t,A_t)=Z_tU_t$ and $f(Z_0,N_0,A_0)=1$, we get, a.s.\ for all $t\ge0$,
\begin{equation}\label{eq:3-5-19:g}
  Z_tU_t=1+L_t+\frac12\int_0^tZ_sU_s\,dA_s-\int_0^tU_s\,d\qv{M,N}_s+\frac12\int_0^tZ_sU_s\,d\qv{N}_s ,
\end{equation}
where
\[
  L_t\defeq\int_0^tU_s\,dM_s-\int_0^tZ_sU_s\,dN_s ;
\]
the integrands $U$ and $ZU$ are adapted and continuous, hence in $\cP^*$, and $L\in\cM^{c,\loc}$. By \eqref{eq:3-5-19:f} and \eqref{eq:3-5-19:e}, on an event of probability one the sum of the three Lebesgue--Stieltjes integrals in \eqref{eq:3-5-19:g} equals
\[
  \int_0^tU_s\Bigl(Z_s\norm{X_s}^2-\sum_{j=1}^dY^{(j)}_sX^{(j)}_s\Bigr)ds
  =\int_0^tU_s\sum_{j=1}^dX^{(j)}_s\bigl(Z_sX^{(j)}_s-Y^{(j)}_s\bigr)\,ds=0 .
\]
Hence $Z_tU_t=1+L_t$ for all $t\ge0$, a.s. Next, by \bk{(2.26)} again, bilinearity, \eqref{eq:3-5-19:f} and \eqref{eq:3-5-19:e},
\begin{align*}
  \qv{L}_t&=\int_0^tU_s^2\,d\qv{M}_s-2\int_0^tZ_sU_s^2\,d\qv{M,N}_s+\int_0^tZ_s^2U_s^2\,d\qv{N}_s\\
          &=\int_0^tU_s^2\sum_{j=1}^d\bigl(Y^{(j)}_s-Z_sX^{(j)}_s\bigr)^2\,ds=0;\qquad 0\le t<\infty,\quad\text{a.s.}
\end{align*}
Thus $L^2=L^2-\qv{L}$ is a nonnegative, continuous local martingale (\bk{Problem 1.5.17}), hence a supermartingale by \bk{Problem 1.5.19(ii)}, and $E[L_t^2]\le E[L_0^2]=0$ for every $t$. By continuity, $L_t=0$ for all $t\ge0$, a.s. We have proved that $Z_tU_t=1$, that is,
\[
  Z_t=\frac1{U_t}=\exp\Bigl[\sum_{j=1}^d\int_0^tX^{(j)}_s\,dW^{(j)}_s-\frac12\int_0^t\norm{X_s}^2\,ds\Bigr]=Z_t(X);\qquad 0\le t<\infty,\quad\text{a.s.}
\]

\emph{Conclusion.} The process $X=(X^{(1)},\dots,X^{(d)})$ is measurable, adapted and satisfies \bk{(5.1)} (Step 5). The process $Z(X)$ of \bk{(5.2)} is adapted and indistinguishable from the martingale $Z$, so it is itself a martingale; and for each $0\le T<\infty$ we have $Z_T(X)=Z_T$ a.s., so that \eqref{eq:3-5-19:c} reads
\[
  \wt P_T(A)=E\bigl[\ind_AZ_T(X)\bigr];\qquad A\in\F_T ,
\]
which is \bk{(5.4)}.

\rmk{The hypothesis $P\ll\wt P_T$ was used only to obtain $\xi_T>0$ a.s.\ (Step~1), and it cannot be dropped, since $Z_T(X)>0$ a.s.\ for every $X$ satisfying \bk{(5.1)}. The hypothesis that $\{\F_t\}$ is the augmented Brownian filtration enters through \bk{Problem 4.16} and through fact (a): on a larger filtration the density process need neither be continuous nor start at $1$. By \bk{Theorem 5.1}, the process $W_t-\int_0^tX_s\,ds$, $0\le t\le T$, is then a Brownian motion under $\wt P_T$: on the Brownian filtration every consistent family of equivalent changes of measure is of Girsanov type.}
\end{kex}

\begin{kex}{Exercise}{5.20}{}
As throughout the book's Section~3.5, the filtration $\{\F_t\}$ satisfies the usual
conditions, and $T\in[0,\infty)$ is fixed. The process $\wt M$ is indexed by $[0,T]$, so the
class in the conclusion is the class $\wt{\cM}^{c,\loc}_T$ introduced before
\bk{Proposition 5.4}: we have to show that $\{\wt M_t,\F_t;\ 0\le t\le T\}$ is a continuous
local martingale on $(\Omega,\F_T,\wt P_T)$ with $\wt P_T[\wt M_0=0]=1$. Since $Z_0=1$ a.s.,
the process $Z$ is not itself in $\cM^{c,\loc}$; in the statement $\qv{Z,M}$ stands for
$\qv{Z-1,M}$, where $Z-1\in\cM^{c,\loc}$ by Step~1 below.

\emph{Step 1: the equation for $Z$ and the two expressions for $\wt M$.}
The process $X\defeq L-\frac12\qv{L}$ is a continuous semimartingale with $X_0=0$ a.s., local
martingale part $L$ and bounded variation part $-\frac12\qv{L}$. It\^o's rule
(\bk{Theorem 3.3}) with $f(x)=e^x$ gives, a.s.\ for all $t\ge0$,
\begin{equation}\label{eq:3-5-20:Z}
  Z_t=1+\int_0^tZ_s\,dL_s-\frac12\int_0^tZ_s\,d\qv{L}_s+\frac12\int_0^tZ_s\,d\qv{L}_s
     =1+\Lambda_t,\qquad \Lambda_t\defeq\int_0^tZ_s\,dL_s .
\end{equation}
Here $Z$ is continuous and adapted, hence progressively measurable, and
$\int_0^tZ_s^2\,d\qv{L}_s<\infty$ for every $t$ by the continuity of its paths; thus
$Z\in\cP^*(L)$ and $\Lambda\in\cM^{c,\loc}$ (\bk{Definition 2.23}). In particular $Z$ is a
continuous local martingale with $Z_0=1$ a.s., so $EZ_T=EZ_0=1$ by the martingale hypothesis
and $\wt P_T$ is indeed a probability measure on $(\Omega,\F_T)$, absolutely continuous with
respect to $P$. By \bk{Proposition 2.24} (the identity \bk{(2.30)} with $N=M$) we have, a.s.,
\begin{equation}\label{eq:3-5-20:cross}
  \qv{Z,M}_t\defeq\qv{\Lambda,M}_t=\int_0^tZ_s\,d\qv{L,M}_s;\qquad 0\le t<\infty .
\end{equation}
Fix $\omega$ outside the exceptional null set of \eqref{eq:3-5-20:cross} and $t\ge0$. The
function $s\mapsto\qv{L,M}_s(\omega)$ is continuous and of bounded variation on $[0,t]$, and
\eqref{eq:3-5-20:cross} says that the signed measure $d\qv{Z,M}_s(\omega)$ on $[0,t]$ has
density $Z_s(\omega)$ with respect to $d\qv{L,M}_s(\omega)$. Since $s\mapsto 1/Z_s(\omega)$ is
continuous, hence bounded, on $[0,t]$, we obtain
\[
  \int_0^t\frac{1}{Z_s}\,d\qv{Z,M}_s=\int_0^t\frac{1}{Z_s}\,Z_s\,d\qv{L,M}_s=\qv{L,M}_t ,
\]
which is the equality of the two expressions for $\wt M_t$.

The process $C\defeq-\qv{L,M}=-\frac14\bigl[\qv{L+M}-\qv{L-M}\bigr]$ is adapted, a.s.\
continuous, of bounded variation on bounded intervals, and $C_0=0$ a.s. Hence
$\wt M=M+C$ is adapted, a.s.\ continuous, with $\wt M_0=0$ a.s. The exceptional set is
$P$-null, hence belongs to $\F_0$ and is also $\wt P_T$-null; redefining $\wt M$ to be
identically zero on it changes neither the adaptedness of $\wt M$ nor any of the assertions
to be proved. We therefore assume that \emph{every} path of $\wt M$ is continuous and starts
at~$0$. In particular $\wt P_T[\wt M_0=0]=1$, and $\wt M$ is progressively measurable.

\emph{Step 2: $Z\wt M$ is a local martingale under $P$.}
We apply the integration-by-parts formula \bk{(3.8)} of \bk{Problem 3.12} to the continuous
semimartingales $Z=1+\Lambda$ (with zero bounded variation part) and $\wt M=M+C$:
\[
  Z_t\wt M_t=\int_0^tZ_s\,dM_s+\int_0^tZ_s\,dC_s+\int_0^t\wt M_s\,d\Lambda_s+\qv{\Lambda,M}_t .
\]
By \eqref{eq:3-5-20:cross}, $\int_0^tZ_s\,dC_s=-\int_0^tZ_s\,d\qv{L,M}_s=-\qv{\Lambda,M}_t$,
so the second and fourth terms cancel and
\begin{equation}\label{eq:3-5-20:N}
  N_t\defeq Z_t\wt M_t=\int_0^tZ_s\,dM_s+\int_0^t\wt M_s\,d\Lambda_s;\qquad 0\le t<\infty,
  \quad\text{a.s.~}P .
\end{equation}
The integrands $Z$ and $\wt M$ are continuous and adapted, so they belong to $\cP^*(M)$ and
$\cP^*(\Lambda)$, respectively, and both integrals in \eqref{eq:3-5-20:N} are in
$\cM^{c,\loc}$. Thus $N\in\cM^{c,\loc}$ under $P$.

\emph{Step 3: localization.} For $n\ge1$ let
\[
  \tau_n\defeq\inf\{t\ge0;\ \abs{\wt M_t}\ge n\}\qquad(\inf\emptyset=\infty).
\]
Because $\wt M$ is adapted with continuous paths and $(-\infty,-n]\cup[n,\infty)$ is closed,
$\tau_n$ is a stopping time of $\{\F_t\}$ (\bk{Problem 1.2.7}). The sequence
$\{\tau_n\}$ is nondecreasing, and since every path of $\wt M$ is bounded on every bounded
interval, $\lim_{n\to\infty}\tau_n(\omega)=\infty$ for \emph{every} $\omega\in\Omega$; in
particular $\tau_n(\omega)>T$ for all sufficiently large $n$. Moreover, as $\wt M_0=0$,
\begin{equation}\label{eq:3-5-20:bdd}
  \abs{\wt M_{t\wedge\tau_n}}\le n;\qquad 0\le t<\infty,\ n\ge1 .
\end{equation}
We fix $n$ and write $\tau=\tau_n$.

\begin{claim}
$\{N_{t\wedge\tau},\F_t;\ 0\le t\le T\}$ is a martingale under $P$.
\end{claim}
\begin{proof}
Let $\{\rho_k\}_{k\ge1}$ be a nondecreasing sequence of stopping times with
$\rho_k\uparrow\infty$ a.s.\ such that $\{N_{t\wedge\rho_k},\F_t;\ 0\le t<\infty\}$ is a
martingale for every $k$ (\bk{Definition 1.5.15}). By \bk{Problem 1.3.24}(i), so is
$\{N_{t\wedge\rho_k\wedge\tau},\F_t;\ 0\le t<\infty\}$; hence for $0\le s\le t\le T$ and
$A\in\F_s$,
\begin{equation}\label{eq:3-5-20:k}
  E\bigl[\ind_AN_{t\wedge\tau\wedge\rho_k}\bigr]=E\bigl[\ind_AN_{s\wedge\tau\wedge\rho_k}\bigr];
  \qquad k\ge1 .
\end{equation}
For $u\in[0,T]$ put $\sigma_k\defeq u\wedge\tau\wedge\rho_k$, a stopping time bounded by
$T$. From \eqref{eq:3-5-20:bdd} and the optional sampling theorem for the right-continuous
martingale $Z$ and bounded stopping times (\bk{Problem 1.3.23}(i)),
\[
  \abs{N_{\sigma_k}}=Z_{\sigma_k}\abs{\wt M_{\sigma_k}}\le nZ_{\sigma_k}
  =n\,E\bigl[Z_T\,\big|\,\F_{\sigma_k}\bigr]\quad\text{a.s.}
\]
The conditional expectations of the integrable random variable $Z_T$ with respect to
arbitrary sub-$\sigma$-fields form a uniformly integrable family, so
$\{N_{u\wedge\tau\wedge\rho_k}\}_{k\ge1}$ is uniformly integrable. As
$N_{u\wedge\tau\wedge\rho_k}\to N_{u\wedge\tau}$ a.s.\ when $k\to\infty$, the limit is
integrable and the convergence takes place in $L^1(P)$. Letting $k\to\infty$ in
\eqref{eq:3-5-20:k} with $u=t$ and $u=s$ we get
$E[\ind_AN_{t\wedge\tau}]=E[\ind_AN_{s\wedge\tau}]$, and $N_{t\wedge\tau}$ is
$\F_t$-measurable.
\end{proof}

\emph{Step 4: $\{\wt M_{t\wedge\tau},\F_t;\ 0\le t\le T\}$ is a martingale under $\wt P_T$.}
This process is adapted, continuous and bounded by \eqref{eq:3-5-20:bdd}, hence
$\wt P_T$-integrable. We first record the following form of the Bayes rule of
\bk{Lemma 5.3}: if $\sigma\le T$ is a stopping time and $Y$ is a bounded,
$\F_\sigma$-measurable random variable, then
\begin{equation}\label{eq:3-5-20:bayes}
  \wt E_T[Y]=E[YZ_T]=E\bigl[Y\,E(Z_T\,|\,\F_\sigma)\bigr]=E[YZ_\sigma],
\end{equation}
where the last equality is again \bk{Problem 1.3.23}(i).

Let $0\le s\le t\le T$ and $A\in\F_s$, and set $B\defeq A\cap\{\tau>s\}$. On $\{\tau\le s\}$
we have $\wt M_{t\wedge\tau}=\wt M_\tau=\wt M_{s\wedge\tau}$, so
\begin{equation}\label{eq:3-5-20:B}
  \wt E_T\bigl[\ind_A(\wt M_{t\wedge\tau}-\wt M_{s\wedge\tau})\bigr]
  =\wt E_T\bigl[\ind_B\wt M_{t\wedge\tau}\bigr]-\wt E_T\bigl[\ind_B\wt M_{s\wedge\tau}\bigr].
\end{equation}
The event $B$ belongs to $\F_s$, because $\{\tau>s\}\in\F_s$, and it belongs to
$\F_{s\wedge\tau}$: for $u<s$ we have $B\cap\{s\wedge\tau\le u\}=\emptyset$, and for
$u\ge s$ we have $B\cap\{s\wedge\tau\le u\}=B\in\F_s\subseteq\F_u$. Consequently
$B\in\F_{s\wedge\tau}\subseteq\F_{t\wedge\tau}$ (\bk{Lemma 1.2.15}). Since $\wt M$ is
progressively measurable, $\wt M_{u\wedge\tau}$ is $\F_{u\wedge\tau}$-measurable for
$u\in\{s,t\}$ (\bk{Proposition 1.2.18}), and it is bounded. Thus
\eqref{eq:3-5-20:bayes} applies with $\sigma=u\wedge\tau$ and
$Y=\ind_B\wt M_{u\wedge\tau}$, and yields
\[
  \wt E_T\bigl[\ind_B\wt M_{u\wedge\tau}\bigr]
  =E\bigl[\ind_B\wt M_{u\wedge\tau}Z_{u\wedge\tau}\bigr]=E\bigl[\ind_BN_{u\wedge\tau}\bigr];
  \qquad u\in\{s,t\}.
\]
By the Claim and $B\in\F_s$, $E[\ind_BN_{t\wedge\tau}]=E[\ind_BN_{s\wedge\tau}]$, so the
right-hand side of \eqref{eq:3-5-20:B} vanishes:
\[
  \wt E_T\bigl[\ind_A\wt M_{t\wedge\tau}\bigr]=\wt E_T\bigl[\ind_A\wt M_{s\wedge\tau}\bigr];
  \qquad A\in\F_s,\ 0\le s\le t\le T .
\]
This is the martingale property of $\{\wt M_{t\wedge\tau_n},\F_t;\ 0\le t\le T\}$ under
$\wt P_T$.

\emph{Conclusion.} The process $\{\wt M_t,\F_t;\ 0\le t\le T\}$ is adapted and continuous,
$\wt P_T[\wt M_0=0]=1$, and $\{\tau_n\}$ is a nondecreasing sequence of stopping times of
$\{\F_t\}$ with $\lim_{n\to\infty}\tau_n=\infty$ everywhere (in particular
$\wt P_T$-a.s.) such that each stopped process
$\{\wt M_{t\wedge\tau_n},\F_t;\ 0\le t\le T\}$ is a (bounded) martingale on
$(\Omega,\F_T,\wt P_T)$. Hence $\wt M\in\wt{\cM}^{c,\loc}_T$.

\rmk{(a) As printed, the conclusion reads ``$\wt M\in\wt{\cM}^{c,\loc}$'', a class which the
book does not define; since $\wt M$ is considered on $[0,T]$ and $\wt P_T$ lives on $\F_T$,
the intended class is $\wt{\cM}^{c,\loc}_T$, and this is what we have proved.
(b) \emph{Relation to \bk{Proposition 5.4}.} In the setting of \bk{Theorem 5.1}, suppose
in addition that the processes $X^{(i)}$ are progressively measurable, so that
$X^{(i)}\in\cP^*(W^{(i)})$ by \bk{(5.1)} and \bk{Proposition 2.24} is applicable, and take
$L_t=\sum_{i=1}^d\int_0^tX^{(i)}_s\,dW^{(i)}_s$. By \bk{(2.30)} and
$\qv{W^{(i)},W^{(j)}}_t=\delta_{ij}t$ we have $\qv{L}_t=\int_0^t\norm{X_s}^2\,ds$, so
$Z=Z(X)$ of \bk{(5.2)}, and $\qv{L,M}_t=\sum_{i=1}^d\int_0^tX^{(i)}_s\,d\qv{M,W^{(i)}}_s$
for $M\in\cM^{c,\loc}$, so that $\wt M$ is the process \bk{(5.8)}. \bk{Proposition 5.4}
takes $M$ in $\cM^{c,\loc}_T$, i.e.\ defined on $[0,T]$ only; such an $M$ is covered by
passing to $M^T_t\defeq M_{t\wedge T}$, $0\le t<\infty$, which agrees with $M$ on $[0,T]$.
Indeed, if $\{\tau_n\}$ are stopping times with $\tau_n\uparrow\infty$ a.s.\ for which
$\{M_{t\wedge\tau_n},\F_t;\ 0\le t\le T\}$ are martingales, then
$\{M_{t\wedge T\wedge\tau_n},\F_t;\ 0\le t<\infty\}$ is a martingale: for $s\le t$ we have
$E[M_{t\wedge T\wedge\tau_n}\,|\,\F_s]=M_{s\wedge T\wedge\tau_n}$ a.s.\ by the martingale
property on $[0,T]$ when $s<T$, and trivially when $s\ge T$, because then both random
variables equal the $\F_T$-measurable $M_{T\wedge\tau_n}$. Thus $M^T\in\cM^{c,\loc}$, and
the first assertion of \bk{Proposition 5.4} is, for progressively measurable $X$, a special
case of the exercise. For $X$ merely measurable and adapted, as allowed in
\bk{Theorem 5.1}, the integrals $\int X^{(i)}\,dW^{(i)}$ are those of the extension of the
stochastic integral with respect to Brownian motion to such integrands, for which
\bk{Proposition 2.24} is not stated; we make no claim in that generality.}
\end{kex}

\begin{kex}{Exercise}{5.21}{}
The exercise cannot be solved as printed. If ``the setting of \bk{Corollary 5.2}'' is meant
to include the hypothesis of that corollary that $Z(X)$ is a martingale, there is nothing to
show. If it means only the canonical space, Wiener measure and a progressively measurable
$X$, the assertion is false: Step 3 gives a counterexample. What the hypothesis ``for
\emph{every} $\omega$'' does yield is proved in Steps 4 and 5: a necessary and sufficient
condition for $Z(X)$ to be a martingale, and the fact that $Z(X)$, once it is a martingale,
is \emph{uniformly integrable}.

\rmk{In Step 3 we exhibit, for $d=1$, a left-continuous, $\{\F^W_t\}$-adapted process $X$ on
the canonical space with $\int_0^\infty X_t^2(\omega)\,dt<\infty$ for every $\omega\in\Omega$
and $EZ_1(X)=2\bigl(1-\Phi(1)\bigr)<1$, where $\Phi$ is the standard normal distribution
function; so the hypotheses do not imply the martingale property. The Remark following
\bk{Corollary 5.2} singles out the uniform integrability of the martingale $Z(X)$ as the
property that decides whether $\wt P$ is comparable with $P$ on $\F^W_\infty$, and the
hypothesis of the exercise is imposed on the whole half-line; we therefore regard the
Proposition of Step 5 as the corrected statement.}

\emph{Conventions.} $\Omega=C[0,\infty)^d$, $W$ is the coordinate process, $P$ is Wiener
measure on $\F^W_\infty=\B(C[0,\infty)^d)$, and $X$ is progressively measurable with respect
to $\{\F^W_t\}$. We denote by $\{\F_t\}$ the augmentation of $\{\F^W_t\}$ under $P$; it
satisfies the usual conditions and $W$ is a Brownian motion relative to it
(\bk{Proposition 2.7.7}, \bk{Theorem 2.7.9}). Stochastic integrals with respect to $W$,
and hence $Z(X)$, are constructed on $(\Omega,\F_\infty,P)$, $\{\F_t\}$, as at the beginning of
Section 3.5. By the discussion preceding \bk{Proposition 5.12}, $Z(X)$ is a nonnegative,
continuous $\{\F_t\}$-supermartingale with $Z_0(X)=1$, and it is a martingale if and only if
\bk{(5.17)} holds, i.e.\ $EZ_T(X)=1$ for every $0\le T<\infty$. Since $\F_s=\sigma(\F^W_s\cup\cN)$, where $\cN$ is the
collection of $P$-null sets, \bk{Problem 2.7.3} shows that every $F\in\F_s$ differs from
some $G\in\F^W_s$ by a $P$-null set. Two consequences follow. First, for an integrable $Y$
the random variable $E[Y|\F^W_s]$ satisfies $E[\ind_FY]=E[\ind_GY]=E[\ind_GE[Y|\F^W_s]]
=E[\ind_FE[Y|\F^W_s]]$ for all $F\in\F_s$, so $E[Y|\F_s]=E[Y|\F^W_s]$ a.s. Second,
approximating by simple functions, every $\F_s$-measurable random variable coincides a.s.\
with an $\F^W_s$-measurable one; in particular, for each $t$ there is an
$\F^W_t$-measurable $Z'_t$ with $Z'_t=Z_t(X)$ a.s. The process $Z_t(X)$ itself is in general
only $\{\F_t\}$-adapted, and, as the book does tacitly, we understand ``$Z(X)$ is an
$\{\F^W_t\}$-martingale'' to mean that such a process $Z'$ (for one, equivalently for every,
choice of the $Z'_t$) is an $\{\F^W_t\}$-martingale. By the two consequences just stated,
$E[Z'_t|\F^W_s]=E[Z_t(X)|\F_s]$ a.s.\ for $s\le t$, so this holds if and only if $Z(X)$ is an
$\{\F_t\}$-martingale, i.e.\ if and only if \bk{(5.17)} holds. This is also the sense in
which the hypothesis ``$Z(X)$ is a martingale'' of \bk{Corollary 5.2} is to be read.

\emph{Step 1: truncation.} Let $A_t(\omega)\defeq\int_0^t\norm{X_s(\omega)}^2\,ds$,
$0\le t\le\infty$. By hypothesis $A_\infty(\omega)<\infty$ for every $\omega$, so for every
$\omega$ the function $t\mapsto A_t(\omega)$ is finite, nondecreasing and continuous on
$[0,\infty)$; in particular \bk{(5.1)} holds. Since the restriction of $X$ to
$[0,t]\times\Omega$ is $\B([0,t])\otimes\F^W_t$-measurable, Tonelli's theorem shows that
$A_t$ is $\F^W_t$-measurable. For $n\ge1$ put
\[
  T_n\defeq\inf\{t\ge0;\ A_t\ge n\}\quad(\inf\emptyset=\infty),\qquad
  X^n_t\defeq X_t\,\ind_{\{t\le T_n\}} .
\]
Because $A$ is continuous and nondecreasing, $\{T_n\le t\}=\{A_t\ge n\}\in\F^W_t$, so $T_n$ is
a stopping time of $\{\F^W_t\}$, we have $T_n\le T_{n+1}$, and
\begin{equation}\label{eq:3-5-21:exh}
  \{T_n>T\}=\{A_T<n\}\in\F^W_T,\qquad \{T_n>T\}\uparrow\Omega\quad\text{as }n\to\infty,
  \qquad 0\le T<\infty,
\end{equation}
the convergence holding because $A_T(\omega)<\infty$ for \emph{every} $\omega$. The process
$\ind_{\{t\le T_n\}}$ is adapted to $\{\F^W_t\}$ and has left-continuous paths, hence is
progressively measurable (\bk{Proposition 1.1.13}), and so is $X^n$. Moreover
$\int_0^\infty\norm{X^n_t}^2\,dt=A_{T_n}\le n$ for every $\omega$ (with
$A_{T_n}=A_\infty\le n$ on $\{T_n=\infty\}$). Thus $X^n$ satisfies \bk{(5.1)} and the Novikov
condition \bk{(5.20)}, and \bk{Corollary 5.13} shows that $Z(X^n)$ is a martingale. By
\bk{Corollary 5.2} there is a unique probability measure $\wt P^n$ on $\F^W_\infty$ with
\begin{equation}\label{eq:3-5-21:Pn}
  \wt P^n(A)=E\bigl[\ind_AZ_T(X^n)\bigr];\qquad A\in\F^W_T,\ 0\le T<\infty,
\end{equation}
and under $\wt P^n$ the process $\wt W^n_t\defeq W_t-\int_0^tX^n_s\,ds$, $0\le t<\infty$, is a
$d$-dimensional Brownian motion relative to $\{\F^W_t\}$.

\emph{Step 2: the basic formula.} We claim that, for all $n\ge1$ and $0\le T<\infty$,
\begin{equation}\label{eq:3-5-21:basic}
  E\bigl[\ind_AZ_T(X)\bigr]=\wt P^n(A)\quad\text{for }A\in\F^W_T,\ A\subseteq\{T_n>T\},
  \qquad\text{and}\qquad
  EZ_T(X)=\lim_{n\to\infty}\wt P^n[T_n>T].
\end{equation}
Indeed, the stopping property \bk{(2.24)} of the stochastic integral, which remains valid for
integrands satisfying only \bk{(5.1)} (see the remarks preceding \bk{Proposition 2.24}), gives
$\int_0^{T\wedge T_n}X^{(i)}_s\,dW^{(i)}_s=\int_0^TX^{n,(i)}_s\,dW^{(i)}_s$ a.s., and
$\int_0^{T\wedge T_n}\norm{X_s}^2ds=\int_0^T\norm{X^n_s}^2ds$ holds identically. Hence
\begin{equation}\label{eq:3-5-21:stop}
  Z_{T\wedge T_n}(X)=Z_T(X^n)\quad\text{a.s.},\qquad 0\le T<\infty,
\end{equation}
and in particular $Z_T(X)=Z_T(X^n)$ a.s.\ on $\{T_n>T\}$. Together with
\eqref{eq:3-5-21:Pn} this proves the first assertion of \eqref{eq:3-5-21:basic}. Taking
$A=\{T_n>T\}$, which is admissible by \eqref{eq:3-5-21:exh}, we get
$E[\ind_{\{T_n>T\}}Z_T(X)]=\wt P^n[T_n>T]$, and the second assertion follows from
\eqref{eq:3-5-21:exh} and the monotone convergence theorem.

\emph{Step 3: a counterexample.} Let $d=1$ and, for $\omega\in\Omega=C[0,\infty)$, put
\[
  f_\omega(t)\defeq\omega(t)+\log(1-t)+1,\quad 0\le t<1;\qquad
  \tau(\omega)\defeq\inf\{t\in[0,1);\ f_\omega(t)\le0\}.
\]
Since $\omega$ is bounded on $[0,1]$, we have $f_\omega(t)\to-\infty$ as $t\uparrow1$; hence
$\tau(\omega)<1$ for \emph{every} $\omega\in\Omega$, and $f_\omega(\tau(\omega))\le0$ by
continuity. For $t<1$ we have
$\{\tau\le t\}=\{\min_{0\le u\le t}f_\cdot(u)\le0\}\in\F^W_t$, and $\{\tau\le t\}=\Omega$ for
$t\ge1$; thus $\tau$ is a stopping time of $\{\F^W_t\}$. Define
\[
  X_t(\omega)\defeq\frac{1}{1-t}\,\ind_{\{t\le\tau(\omega)\}},\qquad 0\le t<\infty
\]
(note that $t\le\tau(\omega)$ forces $t<1$). This process is adapted to $\{\F^W_t\}$ and
left-continuous, hence progressively measurable (\bk{Proposition 1.1.13}), and for every
$\omega\in\Omega$
\[
  A_t(\omega)=\int_0^{t\wedge\tau(\omega)}\frac{du}{(1-u)^2}
  =\frac{t\wedge\tau(\omega)}{1-t\wedge\tau(\omega)},\qquad
  \int_0^\infty X_t^2(\omega)\,dt=\frac{\tau(\omega)}{1-\tau(\omega)}<\infty .
\]
So $X$ satisfies all the requirements of the exercise. With $s_n\defeq n/(n+1)$ we have
$A_t\ge n$ if and only if $t\wedge\tau\ge s_n$, whence, in the notation of Step 1,
\[
  T_n=s_n\ \text{ on }\{\tau\ge s_n\},\qquad T_n=\infty\ \text{ on }\{\tau<s_n\},\qquad
  X^n_t=\frac{1}{1-t}\,\ind_{\{t\le\tau\wedge s_n\}},
\]
so that $\{T_n>1\}=\{\tau<s_n\}$ and
$\wt W^n_t=W_t+\log\bigl(1-t\wedge\tau\wedge s_n\bigr)$; in particular
$\wt W^n_t=f_\cdot(t)-1$ for $0\le t\le\tau\wedge s_n$. Let
$R\defeq\inf\{t\ge0;\ \wt W^n_t=-1\}$ and $\Omega_0\defeq\{W_0=0\}\in\F^W_0$. Since
$Z_0(X^n)=1$, \eqref{eq:3-5-21:Pn} with $T=0$ gives
$\wt P^n(\Omega_0)=P(\Omega_0)=1$. Let $\omega\in\Omega_0$, so that $f_\omega(0)=1>0$. If
$\tau(\omega)<s_n$, then $\tau(\omega)>0$ by continuity of $f_\omega$, $f_\omega>0$ on $[0,\tau(\omega))$ by
the definition of $\tau$, and $f_\omega(\tau(\omega))\le0$, so that $f_\omega(\tau(\omega))=0$; thus $\wt W^n_\tau(\omega)=-1$ and
$R(\omega)\le\tau(\omega)<s_n$. If $\tau(\omega)\ge s_n$, then for every $t<s_n$ we have
$t<\tau(\omega)$, hence $f_\omega(t)>0$ and $\wt W^n_t(\omega)=f_\omega(t)-1>-1$, so
$R(\omega)\ge s_n$. Therefore $\{\tau<s_n\}\cap\Omega_0=\{R<s_n\}\cap\Omega_0$. (Off
$\Omega_0$ the identity may fail, e.g.\ if $\omega(t)<-1$ for all $t$; this set is
$\wt P^n$-null.) Under $\wt P^n$ the process $\wt W^n$ is a Brownian motion relative to
$\{\F^W_t\}$ with $\wt W^n_0=W_0=0$ $\wt P^n$-a.s., so $-\wt W^n$ is a standard
one-dimensional Brownian motion on $(\Omega,\F^W_\infty,\wt P^n)$, and $R$ is its passage
time to the level $1$. Hence \bk{(2.6.2)} (proved in \bk{Proposition 2.6.19}) gives
\[
  \wt P^n[T_n>1]=\wt P^n[\tau<s_n]=\wt P^n[R<s_n]=2\Bigl(1-\Phi\bigl(s_n^{-1/2}\bigr)\Bigr).
\]
As $s_n\uparrow1$, the second formula in \eqref{eq:3-5-21:basic} yields
\[
  EZ_1(X)=\lim_{n\to\infty}2\Bigl(1-\Phi\bigl(s_n^{-1/2}\bigr)\Bigr)=2\bigl(1-\Phi(1)\bigr)
  \approx0.317<1 ,
\]
so \bk{(5.17)} fails and $Z(X)$ is not a martingale. (Under $\wt P^n$ the coordinate process
is $\wt W^n_t-\log(1-t)$ up to time $\tau\wedge s_n$; the mass $2\Phi(1)-1$ that is lost
corresponds to Brownian paths $\wt W^n$ that do not reach $-1$ before time $1$, along which
the coordinate process would have to tend to $+\infty$ as $t\uparrow1$, and no such path
exists in $C[0,\infty)$.)

\emph{Step 4: a criterion.} We return to the general setting of the exercise and to the
notation of Steps 1 and 2.

\begin{lem}
Suppose that $X$ is $\{\F^W_t\}$-progressively measurable and
$\int_0^\infty\norm{X_t(\omega)}^2dt<\infty$ for every $\omega\in\Omega$. Then the following
are equivalent:
\begin{enumerate}[label=\textup{(\alph*)}]
\item $Z(X)$ is an $\{\F^W_t\}$-martingale;
\item $\lim_{n\to\infty}\wt P^n[T_n\le T]=0$ for every $0\le T<\infty$;
\item there exists a probability measure $Q$ on $(\Omega,\F^W_\infty)$ such that
$Q(A)=\wt P^n(A)$ for all $n\ge1$, $0\le T<\infty$ and all $A\in\F^W_T$ with
$A\subseteq\{T_n>T\}$.
\end{enumerate}
\end{lem}

\begin{proof}
(a)$\Rightarrow$(c). If $Z(X)$ is a martingale, \bk{Corollary 5.2} provides a probability
measure $\wt P$ on $\F^W_\infty$ with $\wt P(A)=E[\ind_AZ_T(X)]$ for $A\in\F^W_T$. For
$A\in\F^W_T$ with $A\subseteq\{T_n>T\}$, the first assertion of \eqref{eq:3-5-21:basic}
gives $\wt P(A)=\wt P^n(A)$; so $Q=\wt P$ has the required property.

(c)$\Rightarrow$(b). By \eqref{eq:3-5-21:exh} the set $\{T_n>T\}$ belongs to $\F^W_T$, so
$\wt P^n[T_n>T]=Q[T_n>T]$. Again by \eqref{eq:3-5-21:exh}, these sets increase to $\Omega$
as $n\to\infty$, and the $\sigma$-additivity of $Q$ gives
$\wt P^n[T_n>T]=Q[T_n>T]\uparrow Q(\Omega)=1$.

(b)$\Rightarrow$(a). By the second assertion of \eqref{eq:3-5-21:basic},
$EZ_T(X)=\lim_n\bigl(1-\wt P^n[T_n\le T]\bigr)=1$ for every $T$, which is \bk{(5.17)}; hence
the supermartingale $Z(X)$ is a martingale.
\end{proof}

By the first assertion of \eqref{eq:3-5-21:basic}, the set function
$A\mapsto E[\ind_AZ_T(X)]$ agrees with every $\wt P^k$, $k\ge n$, on the sets
$A\in\F^W_T$, $A\subseteq\{T_n>T\}$; the Lemma says that $Z(X)$ is a martingale exactly when
this consistent family comes from a $\sigma$-additive probability measure on $\F^W_\infty$.
In the example of Step 3 it does not: the sets $\{T_n\le1\}=\{\tau\ge s_n\}$ decrease to the
empty set, while $\wt P^n[T_n\le1]=2\Phi\bigl(s_n^{-1/2}\bigr)-1\to2\Phi(1)-1>0$.

\emph{Step 5: the corrected statement.}

\begin{claim}[Proposition]
In the setting of \bk{Corollary 5.2} (so that $Z(X)$ is a martingale and $\wt P$ is the
measure of \bk{(5.7)}), suppose that $\int_0^\infty\norm{X_t(\omega)}^2dt<\infty$ for every
$\omega\in\Omega$. Then $Z(X)$ is a uniformly integrable martingale: the limit
$Z_\infty(X)\defeq\lim_{t\to\infty}Z_t(X)$ exists a.s., $Z_t(X)=E[Z_\infty(X)|\F^W_t]$ a.s.\
for every $0\le t<\infty$, and $EZ_\infty(X)=1$. Moreover $0<Z_\infty(X)<\infty$ a.s., and
$\wt P(A)=E[\ind_AZ_\infty(X)]$ for every $A\in\F^W_\infty$; in particular $P$ and $\wt P$ are
mutually absolutely continuous on $\F^W_\infty$.
\end{claim}

\begin{proof}
\emph{The truncated processes on $[0,\infty]$.} Fix $n\ge1$. The nonnegative, continuous
supermartingales $Z(X^n)$ and $Z(-X^n)$ converge a.s., as $t\to\infty$, to finite limits
$Z_\infty(X^n)$ and $Z_\infty(-X^n)$ (\bk{Problem 1.3.16}). Since
$Z_t(X^n)Z_t(-X^n)=\exp\bigl[-\int_0^t\norm{X^n_s}^2ds\bigr]\ge e^{-n}$ for all $t$, we have
$Z_\infty(X^n)Z_\infty(-X^n)\ge e^{-n}$ and therefore
\begin{equation}\label{eq:3-5-21:pos}
  0<Z_\infty(X^n)<\infty\quad\text{a.s.}
\end{equation}
Furthermore $Z_t(X^n)^2=Z_t(2X^n)\exp\bigl[\int_0^t\norm{X^n_s}^2ds\bigr]\le e^nZ_t(2X^n)$,
and $EZ_t(2X^n)\le1$ because $Z(2X^n)$ is a supermartingale with initial value $1$; thus
$\sup_tEZ_t(X^n)^2\le e^n$, the family $\{Z_t(X^n)\}_{t\ge0}$ is uniformly integrable, and
$Z_t(X^n)\to Z_\infty(X^n)$ in $L^1$. Letting $T\to\infty$ in \eqref{eq:3-5-21:Pn} we obtain
$\wt P^n(A)=E[\ind_AZ_\infty(X^n)]$ for every $A$ in the algebra $\bigcup_{T}\F^W_T$. Both
sides are probability measures on $\F^W_\infty=\sigma\bigl(\bigcup_T\F^W_T\bigr)$, so Dynkin's
$\pi$-$\lambda$ theorem gives
\begin{equation}\label{eq:3-5-21:dens}
  \wt P^n(A)=E\bigl[\ind_AZ_\infty(X^n)\bigr];\qquad A\in\F^W_\infty .
\end{equation}

\emph{The limit of $Z(X)$.} The nonnegative, continuous supermartingale $Z(X)$ converges
a.s.\ to a random variable $Z_\infty(X)$, and $E[Z_\infty(X)|\F_t]\le Z_t(X)$ a.s.\ for every
$t$ (\bk{Problem 1.3.16}); in particular $EZ_\infty(X)\le1$. Let
$\Omega_n\defeq\{T_n=\infty\}=\bigcap_{k\ge1}\{T_n>k\}\in\F^W_\infty$. Applying
\eqref{eq:3-5-21:stop} at integer times we see that, a.s.\ on $\Omega_n$,
$Z_k(X)=Z_k(X^n)$ for all integers $k\ge1$, and therefore
\begin{equation}\label{eq:3-5-21:lim}
  Z_\infty(X)=Z_\infty(X^n)\quad\text{a.s.\ on }\Omega_n .
\end{equation}
Since $\Omega_n\supseteq\{A_\infty<n\}$ and $T_n\le T_{n+1}$, we have
$\Omega_n\uparrow\Omega$ as $n\to\infty$; it is here that we use $A_\infty(\omega)<\infty$ for
\emph{every} $\omega$.

\emph{The expectation of $Z_\infty(X)$.} By \eqref{eq:3-5-21:lim}, \eqref{eq:3-5-21:dens},
the continuity from above of the measure $\wt P^n$, the first assertion of
\eqref{eq:3-5-21:basic} with $A=\{T_n>k\}$ and $T=k$, the defining property \bk{(5.7)} of
$\wt P$, and the continuity from above of $\wt P$,
\begin{align*}
  EZ_\infty(X)\ \ge\ E\bigl[\ind_{\Omega_n}Z_\infty(X^n)\bigr]
  &=\wt P^n(\Omega_n)=\lim_{k\to\infty}\wt P^n[T_n>k]
   =\lim_{k\to\infty}E\bigl[\ind_{\{T_n>k\}}Z_k(X)\bigr]\\
  &=\lim_{k\to\infty}\wt P[T_n>k]=\wt P(\Omega_n).
\end{align*}
Letting $n\to\infty$ and using $\Omega_n\uparrow\Omega$ together with the $\sigma$-additivity of
$\wt P$ on $\F^W_\infty$, we get $EZ_\infty(X)\ge\wt P(\Omega)=1$. Hence $EZ_\infty(X)=1$.

\emph{Conclusion.} For each $t$, the random variable $Z_t(X)-E[Z_\infty(X)|\F_t]$ is
nonnegative a.s.\ and has expectation $EZ_t(X)-EZ_\infty(X)=1-1=0$; hence
$Z_t(X)=E[Z_\infty(X)|\F_t]=E[Z_\infty(X)|\F^W_t]$ a.s., and the family
$\{Z_t(X)\}_{t\ge0}$ is uniformly integrable (\bk{Problem 1.3.20}). Consequently
$\wt P(A)=E[\ind_AZ_T(X)]=E[\ind_AZ_\infty(X)]$ for $A\in\F^W_T$, $0\le T<\infty$, and as in
the derivation of \eqref{eq:3-5-21:dens} this extends to all $A\in\F^W_\infty$. Finally, by
\eqref{eq:3-5-21:lim} and \eqref{eq:3-5-21:pos}, $0<Z_\infty(X)<\infty$ a.s.\ on each
$\Omega_n$, hence a.s.\ on $\bigcup_n\Omega_n=\Omega$. Therefore, for $A\in\F^W_\infty$,
$\wt P(A)=E[\ind_AZ_\infty(X)]=0$ if and only if $P(A)=0$.
\end{proof}

\rmk{The pathwise form of the hypothesis cannot be replaced by
$P\bigl[\int_0^\infty\norm{X_t}^2dt<\infty\bigr]=1$ in this argument, because the event
$\{A_\infty<\infty\}=\bigcup_n\Omega_n$ has to carry full $\wt P$-measure, and $\wt P$ is not
known in advance to be absolutely continuous with respect to $P$ on $\F^W_\infty$. Without
any integrability on the half-line the conclusion fails: for $d=1$ and $X\equiv\mu\ne0$ the
martingale $Z(X)$ is not uniformly integrable, as noted in the Remark following
\bk{Corollary 5.2}.}
\end{kex}

\solsection{3.6}{Local Time and a Generalized It\^o Rule for Brownian Motion}

\begin{kex}{Exercise}{6.6}{}
Let $L$ be a Brownian local time and $\Omega^*$ the event of \bk{Definition 6.3}. We fix
$\omega\in\Omega^*$ and suppress it from the notation; thus $(s,y)\mapsto L_s(y)$ is a continuous,
real-valued function on $[0,\infty)\times\R$, and the occupation time \bk{(6.5)} satisfies \bk{(6.6)}:
\begin{equation}\label{eq:3-6-06:occ}
  \Gamma_s(B)=\int_B 2L_s(y)\,dy;\qquad 0\le s<\infty,\ B\in\B(\R).
\end{equation}
We prove the following statement, which is stronger than the one required.

\begin{claim}
There is no $t\in[0,\infty)$ for which there exist constants $C\in(0,\infty)$ and $\delta>0$ with
\begin{equation}\label{eq:3-6-06:lip}
  \abs{W_{t+h}-W_t}\le Ch;\qquad 0\le h\le\delta .
\end{equation}
\end{claim}

\emph{The claim implies the assertion.} Suppose that $s\mapsto W_s$ is differentiable at some
$t\in(0,\infty)$, with derivative $D\in\R$. By the definition of the derivative there is a
$\delta>0$ such that $\abs{W_{t+h}-W_t-Dh}\le h$ for $0<h\le\delta$, and therefore
\eqref{eq:3-6-06:lip} holds with $C=\abs{D}+1$ (for $h=0$ it is trivial). This contradicts the claim.

\begin{proof}[Proof of the claim]
Suppose that \eqref{eq:3-6-06:lip} holds for some $t\ge 0$, $C>0$, $\delta>0$; set $x\defeq W_t$ and,
for $0<h\le\delta$, let $B_h\defeq[x-Ch,\,x+Ch]$, a compact interval of length $2Ch$. For every
$s\in[t,t+h]$ we have $0\le s-t\le h\le\delta$, so \eqref{eq:3-6-06:lip} gives
$\abs{W_s-x}\le C(s-t)\le Ch$, i.e.\ $W_s\in B_h$. Hence, by \bk{(6.5)},
\begin{equation}\label{eq:3-6-06:h}
  \Gamma_{t+h}(B_h)-\Gamma_t(B_h)=\int_t^{t+h}\ind_{B_h}(W_s)\,ds=h .
\end{equation}
On the other hand, the continuous functions $L_{t+h}(\cdot)$ and $L_t(\cdot)$ are bounded on the
compact set $B_h$, so both integrals in \eqref{eq:3-6-06:occ} with $B=B_h$ and $s\in\{t,t+h\}$ are
finite and may be subtracted:
\begin{equation}\label{eq:3-6-06:diff}
  \Gamma_{t+h}(B_h)-\Gamma_t(B_h)=2\int_{B_h}\bigl(L_{t+h}(y)-L_t(y)\bigr)\,dy
  \le 4Ch\cdot\max_{y\in B_h}\abs{L_{t+h}(y)-L_t(y)} .
\end{equation}
Comparing \eqref{eq:3-6-06:h} with \eqref{eq:3-6-06:diff} and dividing by $4Ch>0$, we obtain
\begin{equation}\label{eq:3-6-06:lower}
  \frac{1}{4C}\le\max_{y\in B_h}\abs{L_{t+h}(y)-L_t(y)};\qquad 0<h\le\delta .
\end{equation}
We now let $h\downarrow 0$. Given $\eps>0$, the continuity of $(s,y)\mapsto L_s(y)$ at the point
$(t,x)$ provides an $\eta>0$ such that $\abs{L_s(y)-L_t(x)}<\eps$ whenever $s\ge 0$,
$\abs{s-t}\le\eta$ and $\abs{y-x}\le\eta$. If $0<h\le\min\{\delta,\eta,\eta/C\}$ and $y\in B_h$, then
both $(t+h,y)$ and $(t,y)$ satisfy these two inequalities, whence
\[
  \abs{L_{t+h}(y)-L_t(y)}\le\abs{L_{t+h}(y)-L_t(x)}+\abs{L_t(x)-L_t(y)}<2\eps .
\]
Thus the right-hand side of \eqref{eq:3-6-06:lower} tends to zero as $h\downarrow 0$, while the
left-hand side is a strictly positive constant. This contradiction proves the claim.
\end{proof}

\rmk{The argument uses only the behaviour of the path to the right of $t$, and only the joint
continuity of $L$ together with \bk{(6.6)}. The claim says that for every $\omega\in\Omega^*$ and
\emph{every} $t\in[0,\infty)$,
\[
  \limsup_{h\downarrow 0}\frac{\abs{W_{t+h}(\omega)-W_t(\omega)}}{h}=\infty ;
\]
indeed, if this upper limit were finite, then $\abs{W_{t+h}-W_t}\le Ch$ would hold for some $C$ and
all sufficiently small $h>0$. In particular the path has a finite right-hand derivative at no point
of $[0,\infty)$. The symmetric argument, with the interval $[t-h,t]$ for $0<h\le\min\{\delta,t\}$ and
the differences $\Gamma_t(B_h)-\Gamma_{t-h}(B_h)$, $L_t(y)-L_{t-h}(y)$, gives the same conclusion for
the left-hand difference quotients at every $t\in(0,\infty)$. Since $P^z(\Omega^*)=1$ for every
$z\in\R$, this recovers the almost sure nowhere differentiability of the Brownian path
(cf.\ \bk{Theorem 2.9.18}), once the existence of local time has been established.}
\end{kex}

\begin{kex}{Problem}{6.7}{}
Both assertions concern one sample point and one time at a time. Throughout we fix $\omega\in\Omega$
and $t\in[0,\infty)$, and we assume that the function $x\mapsto L_t(x,\omega)$, which takes values
in $[0,\infty)$, is Borel-measurable; without this neither \bk{(6.6)} nor \bk{(6.7)} has a meaning,
and it holds for every $\omega\in\Omega^*$ of \bk{Definition 6.3} because of the continuity of
$(t,x)\mapsto L_t(x,\omega)$. The path $s\mapsto W_s(\omega)$ is continuous, hence Borel-measurable,
so $s\mapsto f(W_s(\omega))$ is a nonnegative Borel-measurable function on $[0,t]$ for every
Borel-measurable $f\colon\R\to[0,\infty)$, and both sides of \bk{(6.7)} are well-defined elements of
$[0,\infty]$. We introduce the two set functions
\begin{equation}\label{eq:3-6-07:meas}
  \mu(B)\defeq \Gamma_t(B,\omega)=\int_0^t \ind_B(W_s(\omega))\,ds,
  \qquad
  \nu(B)\defeq \int_B 2L_t(x,\omega)\,dx;\qquad B\in\B(\R).
\end{equation}
Both are measures on $(\R,\B(\R))$: countable additivity follows from the monotone convergence
theorem, applied to the partial sums of $\ind_{\bigcup_k B_k}=\sum_k\ind_{B_k}$ for pairwise
disjoint Borel sets $B_1,B_2,\dots$; and $\mu(\R)=t<\infty$. For every Borel-measurable
$f\colon\R\to[0,\infty)$ we have
\begin{equation}\label{eq:3-6-07:int}
  \int_\R f\,d\mu=\int_0^t f(W_s(\omega))\,ds,
  \qquad
  \int_\R f\,d\nu=2\int_{-\infty}^\infty f(x)L_t(x,\omega)\,dx .
\end{equation}
Indeed, for $f=\ind_B$ these identities are the definitions \eqref{eq:3-6-07:meas}; they extend by
linearity to nonnegative simple functions, and then to an arbitrary Borel-measurable $f\ge 0$ by
choosing simple functions $0\le f_n\uparrow f$ (for instance
$f_n=\min\{n,2^{-n}\lfloor 2^n f\rfloor\}$) and applying the monotone convergence theorem to each
of the four integrals. Thus, for the fixed pair $(t,\omega)$,
\begin{equation}\label{eq:3-6-07:reform}
  \text{\bk{(6.6)} reads } \mu(B)=\nu(B)\ \ \forall\, B\in\B(\R),
  \qquad
  \text{\bk{(6.7)} for $f$ reads } \int_\R f\,d\mu=\int_\R f\,d\nu .
\end{equation}

\textbf{(i)} If \bk{(6.7)} holds for every Borel-measurable $f\colon\R\to[0,\infty)$, then taking
$f=\ind_B$ with $B\in\B(\R)$ gives $\mu(B)=\nu(B)$, which is \bk{(6.6)}. Conversely, if
\bk{(6.6)} holds, then $\mu$ and $\nu$ are the same measure on $\B(\R)$, so
$\int_\R f\,d\mu=\int_\R f\,d\nu$ for every Borel-measurable $f\colon\R\to[0,\infty)$, and by
\eqref{eq:3-6-07:int} this is \bk{(6.7)}. Since $t$ and $\omega$ were arbitrary, the equivalence
holds for every $\omega$ separately, simultaneously for all $t\in[0,\infty)$, and therefore also
for any set of $\omega$'s (such as $\Omega^*$).

\textbf{(ii)} Suppose that \bk{(6.7)} holds for every $h\in\cH$; by \eqref{eq:3-6-07:reform},
\begin{equation}\label{eq:3-6-07:hyp}
  \int_\R h\,d\mu=\int_\R h\,d\nu;\qquad h\in\cH .
\end{equation}
In view of part (i) and \eqref{eq:3-6-07:reform}, it suffices to prove that $\mu=\nu$ on $\B(\R)$.

\emph{Step 1. $\mu$ and $\nu$ agree on bounded open intervals.} Let $-\infty<a<b<\infty$. We choose
rational numbers $p_1<p_2<a$ and $b<p_3<p_4$, and denote by $h_0\in\cH$ the function \bk{(6.8)}
with $(q_1,q_2,q_3,q_4)=(p_1,p_2,p_3,p_4)$; then $h_0=1$ on $[a,b]$. For every integer
$n>2/(b-a)$ we choose rational numbers
\[
  a<q_1^{(n)}<q_2^{(n)}<a+\tfrac1n<b-\tfrac1n<q_3^{(n)}<q_4^{(n)}<b ,
\]
and denote by $h_n\in\cH$ the corresponding function \bk{(6.8)}. It vanishes outside
$[q_1^{(n)},q_4^{(n)}]\subset(a,b)$, takes values in $[0,1]$, and equals $1$ on
$[a+\frac1n,b-\frac1n]$. Consequently
\[
  0\le h_n\le \ind_{(a,b)}\le h_0\quad\text{for all such } n,
  \qquad
  \lim_{n\to\infty}h_n(x)=\ind_{(a,b)}(x)\quad\text{for every } x\in\R :
\]
for $x\notin(a,b)$ all the $h_n(x)$ are zero, and for $x\in(a,b)$ we have $h_n(x)=1$ as soon as
$\frac1n\le\min\{x-a,b-x\}$. The dominating function $h_0$ is integrable with respect to both
measures, because \eqref{eq:3-6-07:hyp} and $0\le h_0\le 1$ give
\[
  \int_\R h_0\,d\nu=\int_\R h_0\,d\mu\le\mu(\R)=t<\infty .
\]
The dominated convergence theorem, applied under $\mu$ and under $\nu$, together with
\eqref{eq:3-6-07:hyp} for $h=h_n$, now yields
\begin{equation}\label{eq:3-6-07:int-ab}
  \mu\bigl((a,b)\bigr)=\lim_{n\to\infty}\int_\R h_n\,d\mu
  =\lim_{n\to\infty}\int_\R h_n\,d\nu=\nu\bigl((a,b)\bigr)<\infty .
\end{equation}

\emph{Step 2. $\mu=\nu$ on $\B(\R)$.} The class
$\cC\defeq\{(a,b);\ -\infty<a<b<\infty\}\cup\{\emptyset\}$ is closed under finite intersections and
generates $\B(\R)$, since every open subset of $\R$ is a countable union of bounded open intervals.
For each integer $n\ge1$ the set functions $\mu_n(B)\defeq\mu(B\cap(-n,n))$ and
$\nu_n(B)\defeq\nu(B\cap(-n,n))$ are measures on $\B(\R)$; by \eqref{eq:3-6-07:int-ab} they are
finite, with the same total mass $\mu((-n,n))=\nu((-n,n))$, and they agree on $\cC$, because
$C\cap(-n,n)\in\cC$ for $C\in\cC$. The class
$\{B\in\B(\R);\ \mu_n(B)=\nu_n(B)\}$ therefore contains $\R$ and the $\pi$-system $\cC$, and it is
closed under proper differences and increasing countable unions (finiteness of $\mu_n,\nu_n$ is
used for the differences); by Dynkin's $\pi$-$\lambda$ theorem it is all of $\B(\R)$. Letting
$n\to\infty$ and using the continuity from below of the measures $\mu$ and $\nu$ along
$B\cap(-n,n)\uparrow B$, we obtain
\[
  \mu(B)=\lim_{n\to\infty}\mu_n(B)=\lim_{n\to\infty}\nu_n(B)=\nu(B);\qquad B\in\B(\R).
\]

\emph{Step 3. Conclusion.} By Step 2 and \eqref{eq:3-6-07:reform}, \bk{(6.6)} holds for the pair
$(t,\omega)$; since $t\in[0,\infty)$ was arbitrary, \bk{(6.6)} holds for this $\omega$, and part
(i) shows that \bk{(6.7)} holds for this $\omega$ and every Borel-measurable function
$f\colon\R\to[0,\infty)$.

\rmk{The argument is carried out for each $\omega$ separately, and the class $\cH$ is countable,
being indexed by quadruples of rational numbers. Hence, if for each $h\in\cH$ the identity
\bk{(6.7)} is known to hold, for all $t\in[0,\infty)$, only outside an event $N_h$ which is null
under every $P^z$, then for every $\omega$ outside the single null event $\bigcup_{h\in\cH}N_h$ for
which $x\mapsto L_t(x,\omega)$ is Borel-measurable for every $t$, the identities \bk{(6.6)} and \bk{(6.7)} hold simultaneously for
all Borel sets $B$, all Borel-measurable $f\ge0$ and all $t\in[0,\infty)$. This is the form in
which the result is used in the construction of Brownian local time.}
\end{kex}

\begin{kex}{Problem}{6.18}{}
We fix $(a,z)\in\R^2$ and work under $P^z$. We shall express $L(a)$ through the running maximum of
a Brownian motion by means of the Skorohod equation, exactly as in the proof of
\bk{Theorem 6.17} but for a general level and a general starting point, and then show that this
running maximum increases to infinity.

\emph{Step 1: the Skorohod representation of $L(a)$.}
Let $B_t\defeq B_t(a)=-\int_0^t\sgn(W_s-a)\,dW_s$, $0\le t<\infty$, be the process of
\bk{Problem 6.13}. We collect four facts, each valid for $P^z$-a.e.\ $\omega\in\Omega$.
\begin{enumerate}[label=\textup{(\alph*)}]
\item The path $t\mapsto B_t(\omega)$ is continuous and $B_0(\omega)=0$, because
$B=\{B_t,\F_t;\ 0\le t<\infty\}$ is a Brownian motion under $P^z$ with $P^z[B_0=0]=1$
\textup{(}\bk{Problem 6.13}\,(i)\textup{)}.
\item The Tanaka formula \bk{(6.25)} holds for all $t$ simultaneously:
$\abs{W_t(\omega)-a}=\abs{z-a}-B_t(\omega)+2L_t(a,\omega)$, $0\le t<\infty$
\textup{(}\bk{Proposition 6.8}, relation \bk{(6.13)}; the hypothesis of that proposition is
satisfied, since Brownian local time exists by \bk{Theorem 6.11}\textup{)}.
\item The path $t\mapsto L_t(a,\omega)$ is continuous and nondecreasing, with $L_0(a,\omega)=0$,
because $L(a)$ is a continuous, increasing process in the sense of \bk{Definition 1.4.4}
\textup{(}\bk{Problem 6.13}\,(ii)\textup{)}.
\item $\int_0^\infty \ind_{\R\setminus\{a\}}(W_t(\omega))\,dL_t(a,\omega)=0$; this is \bk{(6.26)}.
\end{enumerate}
Let $\Omega_1\in\F$ be an event with $P^z(\Omega_1)=1$ on which (a)--(d) hold, and fix
$\omega\in\Omega_1$. We apply \bk{Lemma 6.14} with the number $\abs{z-a}\ge 0$ in place of $z$, with
the continuous function $y(t)\defeq -B_t(\omega)$, which satisfies $y(0)=0$ by (a), and with the
candidate $k(t)\defeq 2L_t(a,\omega)$. By (c), $k$ is continuous and nondecreasing with $k(0)=0$,
which is condition (ii) of the lemma. By (b),
\begin{equation*}
x(t)\defeq \abs{z-a}+y(t)+k(t)=\abs{W_t(\omega)-a}\ \ge 0;\qquad 0\le t<\infty,
\end{equation*}
which is condition (i). Finally $\{s\ge 0;\ x(s)>0\}=\{s\ge 0;\ W_s(\omega)\ne a\}$, so that (d) gives
$\int_0^\infty\ind_{\{x(s)>0\}}\,dk(s)=2\int_0^\infty\ind_{\R\setminus\{a\}}(W_s(\omega))\,dL_s(a,\omega)=0$,
which is condition (iii). The uniqueness assertion of \bk{Lemma 6.14} and the formula \bk{(6.29)}
now yield
\begin{equation}\label{eq:3-6-18:skor}
2L_t(a,\omega)=\max\Bigl[0,\ \max_{0\le s\le t}B_s(\omega)-\abs{z-a}\Bigr];
\qquad 0\le t<\infty,\ \ \omega\in\Omega_1 .
\end{equation}
(For $a=z=0$ this is the first identity in \bk{(6.34)}.)

\emph{Step 2: the supremum of a Brownian path is infinite.}
Let $S\defeq\sup_{0\le t<\infty}B_t$. On the event of (a) we have, by continuity,
$S=\sup_{t\in\Q\cap[0,\infty)}B_t$ and $S\ge B_0=0$; thus $S$ agrees $P^z$-a.s.\ with a
$[0,\infty]$-valued random variable which is a fixed Borel function of the sequence
$(B_t)_{t\in\Q\cap[0,\infty)}$, and its law is determined by the finite-dimensional distributions
of $B$. We claim that
\begin{equation}\label{eq:3-6-18:sup}
P^z[S=\infty]=1 .
\end{equation}

\emph{Scaling.} Fix $c>0$ and let $\wt B_t\defeq cB_{t/c^2}$, $0\le t<\infty$. This process starts
at zero, and for $0\le t_0<t_1<\dots<t_n$ its increments $\wt B_{t_j}-\wt B_{t_{j-1}}
=c\,(B_{t_j/c^2}-B_{t_{j-1}/c^2})$ are independent and normally distributed with mean zero and
variance $c^2(t_j-t_{j-1})/c^2=t_j-t_{j-1}$. Hence $\wt B$ has the same finite-dimensional
distributions as $B$, and consequently $\sup_{t\in\Q\cap[0,\infty)}\wt B_t$ has the same law as $S$.
But $t\mapsto t/c^2$ maps $[0,\infty)$ onto itself, so by continuity of the paths
$\sup_{t\in\Q\cap[0,\infty)}\wt B_t=c\,S$, a.s.\ $P^z$. Therefore $cS$ and $S$ have the same law, and
\begin{equation*}
P^z[S>x]=P^z[cS>x]=P^z[S>x/c];\qquad x>0,\ c>0 .
\end{equation*}
Thus $p\defeq P^z[S>x]$ does not depend on $x\in(0,\infty)$. Letting $x\downarrow 0$ and
$x\uparrow\infty$ along sequences, we obtain $P^z[S>0]=p=P^z[S=\infty]$, and so
\begin{equation}\label{eq:3-6-18:dich}
P^z[0<S<\infty]=0,\qquad\text{i.e.}\qquad S\in\{0,\infty\}\quad\text{a.s.\ }P^z .
\end{equation}

\emph{Exclusion of $S=0$.} Let $B'_t\defeq B_{1+t}-B_1$, $0\le t<\infty$, and
$S'\defeq\sup_{0\le t<\infty}B'_t$. Since $B$ is a Brownian motion relative to $\{\F_t\}$, the
increments of $B'$ are independent, normal with mean zero and variance equal to the length of the
time interval, and independent of $\F_1$; hence $B'$ has the same finite-dimensional distributions
as $B$ and the family $(B'_t)_{t\ge 0}$ is independent of $B_1$. Consequently $S'$ has the same law
as $S$, it satisfies \eqref{eq:3-6-18:dich}, and it is independent of $B_1$. Now
$S\ge\sup_{t\ge 1}B_t=B_1+S'$, so that
\begin{equation*}
\{S=0\}\subseteq\{B_1\le 0,\ S'\le -B_1\}\subseteq\{B_1\le 0,\ S'<\infty\},
\end{equation*}
and the last event differs from $\{B_1\le 0,\ S'=0\}$ by a $P^z$-null set, thanks to
\eqref{eq:3-6-18:dich} for $S'$. With $q\defeq P^z[S=0]=P^z[S'=0]$, independence gives
\begin{equation*}
q\le P^z[B_1\le 0]\,P^z[S'=0]=\tfrac12\,q ,
\end{equation*}
whence $q=0$. Together with \eqref{eq:3-6-18:dich} this proves \eqref{eq:3-6-18:sup}.

\emph{Step 3: conclusion.}
Let $\Omega_2\defeq\Omega_1\cap\{\sup_{t\in\Q\cap[0,\infty)}B_t=\infty\}\in\F$; by Step~2,
$P^z(\Omega_2)=1$. For $\omega\in\Omega_2$ the function $t\mapsto\max_{0\le s\le t}B_s(\omega)$ is
nondecreasing and its limit as $t\to\infty$ is $\sup_{0\le s<\infty}B_s(\omega)=\infty$; then
\eqref{eq:3-6-18:skor} shows that $2L_t(a,\omega)\ge\max_{0\le s\le t}B_s(\omega)-\abs{z-a}
\to\infty$ as $t\to\infty$. Hence
\begin{equation*}
\Omega_2\subseteq\Bigl\{\omega\in\Omega;\ \lim_{t\to\infty}L_t(a,\omega)=\infty\Bigr\}
\qquad\text{and}\qquad P^z(\Omega_2)=1,
\end{equation*}
which is the assertion.

\rmk{The set $\{\lim_{t\to\infty}L_t(a)=\infty\}$ involves uncountably many of the random variables
$L_t(a)$, and \bk{Definition 6.3} restricts the paths of $L$ only on $\Omega^*$; on $\Omega^*$ the
path $t\mapsto L_t(a,\omega)$ is continuous, so the set coincides there with the event
$\bigcap_{K\in\N}\bigcup_{n\in\N}\bigcap_{t\in\Q,\,t\ge n}\{L_t(a)>K\}$. In any case we have shown
that it contains an event of $P^z$-probability one, which is how the assertion is to be read.
The formula \eqref{eq:3-6-18:skor} also shows that, for $\omega\in\Omega_1$, we have
$L_t(a,\omega)=0$ if and only if $\max_{0\le s\le t}B_s(\omega)\le\abs{z-a}$, that is, if and only if
$t\le\inf\{s\ge 0;\ B_s(\omega)>\abs{z-a}\}$.}
\end{kex}

\begin{kex}{Problem}{6.20}{}
Parts (iii) and (iv) are solved in the book's Section 3.8; we prove here parts (i) and (ii), whose
conclusions that solution uses, and we add a remark on two points which it leaves tacit.

\textbf{(i)} Fix $x\in\R$. The set $\{\varphi(y);\ y>x\}$ is bounded from below by $\varphi(x)$ and
the set $\{\varphi(y);\ y<x\}$ is bounded from above by $\varphi(x)$, because $\varphi$ is
nondecreasing. Monotonicity also shows that the one-sided limits exist and are given by
\begin{equation}\label{eq:3-6-20:infsup}
  \varphi_+(x)=\inf_{y>x}\varphi(y)\in[\varphi(x),\infty),\qquad
  \varphi_-(x)=\sup_{y<x}\varphi(y)\in(-\infty,\varphi(x)].
\end{equation}
Indeed, with $m\defeq\inf_{y>x}\varphi(y)$ and $\eps>0$ there is $y_0>x$ with $\varphi(y_0)<m+\eps$,
and then $m\le\varphi(y)\le\varphi(y_0)<m+\eps$ for all $y\in(x,y_0)$; the argument for the left limit
is symmetric. In particular $\varphi_\pm$ are real-valued and \bk{(6.43)} holds.

Next, let $x<y$ and choose any $z\in(x,y)$. By \eqref{eq:3-6-20:infsup},
$\varphi_+(x)\le\varphi(z)\le\varphi_-(y)$, and together with \bk{(6.43)} this gives
\begin{equation}\label{eq:3-6-20:chain}
  \varphi_-(x)\le\varphi(x)\le\varphi_+(x)\le\varphi_-(y)\le\varphi(y)\le\varphi_+(y);\qquad x<y.
\end{equation}
In particular, $\varphi_+$ and $\varphi_-$ are nondecreasing.

\emph{Right-continuity of $\varphi_+$.} Fix $x\in\R$ and $\eps>0$. By \eqref{eq:3-6-20:infsup}
there is $y_0>x$ with $\varphi(y_0)<\varphi_+(x)+\eps$. For every $z\in(x,y_0)$ we have
$\varphi_+(x)\le\varphi_+(z)$ by monotonicity and $\varphi_+(z)\le\varphi(y_0)$ by
\eqref{eq:3-6-20:infsup}, so that
$\varphi_+(x)\le\varphi_+(z)<\varphi_+(x)+\eps$. Hence $\lim_{z\downarrow x}\varphi_+(z)=\varphi_+(x)$.

\emph{Left-continuity of $\varphi_-$.} Fix $x\in\R$ and $\eps>0$, and choose $y_0<x$ with
$\varphi(y_0)>\varphi_-(x)-\eps$. For every $z\in(y_0,x)$ we have
$\varphi(y_0)\le\varphi_-(z)$ by \eqref{eq:3-6-20:infsup} and $\varphi_-(z)\le\varphi_-(x)$ by
monotonicity, so that $\varphi_-(x)-\eps<\varphi_-(z)\le\varphi_-(x)$. Hence
$\lim_{z\uparrow x}\varphi_-(z)=\varphi_-(x)$.

\textbf{(ii)} We compute the remaining one-sided limits of $\varphi_\pm$. For $y<x$,
\eqref{eq:3-6-20:chain} (with the roles of $x$ and $y$ interchanged) gives
$\varphi(y)\le\varphi_+(y)\le\varphi_-(x)$; and for $y>x$ it gives
$\varphi_+(x)\le\varphi_-(y)\le\varphi(y)$. Since $\varphi(y)\to\varphi_-(x)$ as $y\uparrow x$ and
$\varphi(y)\to\varphi_+(x)$ as $y\downarrow x$, we obtain
\begin{equation}\label{eq:3-6-20:cross}
  \lim_{y\uparrow x}\varphi_+(y)=\varphi_-(x),\qquad
  \lim_{y\downarrow x}\varphi_-(y)=\varphi_+(x);\qquad x\in\R.
\end{equation}
Set $C\defeq\{x\in\R;\ \varphi_-(x)=\varphi_+(x)\}$ and $N\defeq\R\setminus C$. By (i) the function
$\varphi_+$ is right-continuous, so it is continuous at $x$ if and only if its left limit at $x$
equals $\varphi_+(x)$, which by \eqref{eq:3-6-20:cross} means $x\in C$. Likewise $\varphi_-$ is
left-continuous, so it is continuous at $x$ if and only if its right limit at $x$, namely
$\varphi_+(x)$, equals $\varphi_-(x)$, i.e., if and only if $x\in C$. Thus $\varphi_+$ and $\varphi_-$
have the same set $C$ of continuity points, and for $x\in C$ the inequalities \bk{(6.43)} force
$\varphi_-(x)=\varphi(x)=\varphi_+(x)$. (The set $C$ is also the set of continuity points of
$\varphi$ itself: $\varphi$ is continuous at $x$ if and only if both one-sided limits
$\varphi_\pm(x)$ equal $\varphi(x)$, and by \bk{(6.43)} this happens exactly when
$\varphi_-(x)=\varphi_+(x)$.)

It remains to show that $N$ is countable. For $x\in N$ the open interval
$I_x\defeq(\varphi_-(x),\varphi_+(x))$ is nonempty, and for $x<y$ in $N$ the intervals $I_x$ and
$I_y$ are disjoint because $\varphi_+(x)\le\varphi_-(y)$ by \eqref{eq:3-6-20:chain}. Choosing a
rational number $q_x\in I_x$ for each $x\in N$ we obtain an injective mapping $x\mapsto q_x$ of $N$
into $\Q$; hence $N$ is countable, and $\varphi_\pm(x)=\varphi(x)$ for every $x\in\R\setminus N$.

\rmk{Two points are tacit in the book's solution of (iii). First, $\Phi$ is well defined: a
nondecreasing function is Borel-measurable (the sets $\{\varphi\le c\}$ are intervals) and bounded
on every bounded interval, hence Lebesgue-integrable there; for $x<0$ the integral $\int_0^x$ means
$-\int_x^0$, so that $\Phi(y)-\Phi(x)=\int_x^y\varphi(u)\,du$ for all $x<y$. Second, that solution
passes from the equalities $\varphi_\pm(x)=D^\pm\Phi(x)$ at the points $x$ of $C$ to all of $\R$ ``by
continuity''; this requires that every $y\in\R$ be the limit of a strictly increasing and of a
strictly decreasing sequence of points of $C$, which holds because the complement $N$ of $C$ is
countable and therefore contains no interval. Given such sequences $x_n\uparrow y$ and
$x_n'\downarrow y$ in $C$, the left-continuity of $\varphi_-$ from (i) and of $D^-\Phi$ from
\bk{Problem 6.19} yields $\varphi_-(y)=\lim_n\varphi_-(x_n)=\lim_n D^-\Phi(x_n)=D^-\Phi(y)$, and the
right-continuity of $\varphi_+$ and $D^+\Phi$ yields $\varphi_+(y)=D^+\Phi(y)$ in the same way.}
\end{kex}

\begin{kex}{Problem}{6.21}{}
By \bk{Problem 6.19}, the one-sided derivatives $D^\pm f(x)$ of \bk{(6.40)} exist and are finite for every $x\in\R$, and they satisfy \bk{(6.41)}:
\begin{equation}\label{eq:3-6-21:a}
D^+f(x)\le D^-f(y)\le D^+f(y);\qquad x<y .
\end{equation}
Since every $y\in\R$ has some $x<y$ to its left, the second inequality in \eqref{eq:3-6-21:a} gives
\begin{equation}\label{eq:3-6-21:b}
D^-f(y)\le D^+f(y);\qquad y\in\R .
\end{equation}

\emph{Step 1: the exceptional set.} We define
\[
N\defeq\{x\in\R;\ D^-f(x)<D^+f(x)\},
\qquad I_x\defeq\bigl(D^-f(x),\,D^+f(x)\bigr)\quad\text{for } x\in N .
\]
Each $I_x$, $x\in N$, is a nonempty open interval. If $x<y$ are two points of $N$, then $D^+f(x)\le D^-f(y)$ by \eqref{eq:3-6-21:a}, so every point of $I_x$ is strictly smaller than every point of $I_y$; in particular $I_x\cap I_y=\emptyset$. Choosing a rational number $q_x\in I_x$ for each $x\in N$, we obtain a mapping $x\mapsto q_x$ from $N$ into $\Q$ which is injective, because the intervals are pairwise disjoint. Hence $N$ is countable.

\emph{Step 2: differentiability off $N$.} Let $x\in\R\setminus N$. By \eqref{eq:3-6-21:b} and the definition of $N$ we have $D^-f(x)=D^+f(x)$, i.e., the two one-sided limits in \bk{(6.40)},
\[
\lim_{h\downarrow0}\frac{f(x+h)-f(x)}{h}\qquad\text{and}\qquad\lim_{h\uparrow0}\frac{f(x+h)-f(x)}{h},
\]
exist, are finite and coincide. This is precisely the statement that the two-sided limit of the difference quotient exists and is finite, that is, $f$ is differentiable at $x$ with $f'(x)=D^+f(x)=D^-f(x)$. This proves \bk{(6.45)}. (Conversely, $f$ is not differentiable at any point of $N$, so $N$ is exactly the set of points of nondifferentiability of $f$.)

\emph{Step 3: the integral representation.} Let $\varphi\colon\R\to\R$ be \emph{any} function with
\begin{equation}\label{eq:3-6-21:c}
D^-f(x)\le\varphi(x)\le D^+f(x);\qquad x\in\R .
\end{equation}
Such a function is nondecreasing: for $x<y$, \eqref{eq:3-6-21:c} and the first inequality in \eqref{eq:3-6-21:a} give
\[
\varphi(x)\le D^+f(x)\le D^-f(y)\le\varphi(y).
\]
Thus $\varphi$ satisfies the hypothesis of \bk{Problem 6.20}, the function $\Phi(x)=\int_0^x\varphi(u)\,du$ is well defined (a monotone function is bounded and Riemann integrable on every compact interval; as usual $\int_0^x=-\int_x^0$ for $x<0$), and \eqref{eq:3-6-21:c} is condition \bk{(6.44)}. Part (iv) of \bk{Problem 6.20} then yields
\begin{equation}\label{eq:3-6-21:d}
f(x)-f(0)=\int_0^x\varphi(u)\,du;\qquad x\in\R .
\end{equation}
In view of \eqref{eq:3-6-21:b}, both $\varphi=D^-f$ and $\varphi=D^+f$ satisfy \eqref{eq:3-6-21:c}, and \eqref{eq:3-6-21:d} becomes
\[
f(x)-f(0)=\int_0^xD^-f(u)\,du=\int_0^xD^+f(u)\,du;\qquad x\in\R .
\]

It remains to interpret and to prove the equality involving $f'$, which by Step~2 is defined on $\R\setminus N$ only. The function $f'$ agrees on $\R\setminus N$ with the nondecreasing, hence Borel-measurable and locally bounded, function $D^-f$, and the countable set $N$ is a Borel set of Lebesgue measure zero. Therefore $f'$ is defined Lebesgue-almost everywhere, it is measurable and locally bounded there, and its Lebesgue integral over the interval with endpoints $0$ and $x$ equals that of $D^-f$; since the Riemann and the Lebesgue integral of the bounded monotone function $D^-f$ over a compact interval coincide, we conclude
\[
\int_0^xf'(u)\,du=\int_0^xD^-f(u)\,du=f(x)-f(0);\qquad x\in\R ,
\]
and \bk{(6.46)} is proved.

\rmk{If one prefers to regard $\int_0^xf'(u)\,du$ as a Riemann integral of an everywhere-defined function, one may extend $f'$ to $N$ by assigning to each $x\in N$ an arbitrary value in $[D^-f(x),D^+f(x)]$. The extension satisfies \eqref{eq:3-6-21:c} on all of $\R$ (with equalities off $N$, by \bk{(6.45)}), so it is nondecreasing and \eqref{eq:3-6-21:d} applies to it directly: the value of the integral does not depend on the choice of the extension.}
\end{kex}

\begin{kex}{Problem}{6.24}{}
We work in the setting of the book's Section 3.6: $\{W_t,\F_t\}$, $(\Omega,\F)$, $\{P^z\}_{z\in\R}$ is the
one-dimensional Brownian family on the canonical space, the filtration satisfies the usual conditions,
$L$ is the Brownian local time of \bk{Definition 6.3} (it exists by \bk{Theorem 6.11}) and $\Omega^*$
is the event of that definition. Under $P^z$ the process $W-z$ is a standard Brownian motion, and
stochastic integrals with respect to $W$ are integrals with respect to this continuous martingale,
whose quadratic variation is $\qv{W}_t=t$.

\emph{Notation and conventions.} The functions $f'$ and $f''$ are given only on $\R\setminus D$, so the
integrands in \bk{(6.53)} are not defined at the times $s$ with $W_s\in D$. We set
\begin{equation}\label{eq:3-6-24:conv}
  \varphi(x)\defeq\begin{cases} f'(x); & x\in\R\setminus D,\\ f'(a_k-); & x=a_k,\end{cases}
  \qquad
  g(x)\defeq\begin{cases} f''(x); & x\in\R\setminus D,\\ 0; & x\in D,\end{cases}
  \qquad
  \delta_k\defeq f'(a_k+)-f'(a_k-),
\end{equation}
we prove \bk{(6.53)} with $f'(W_s)$, $f''(W_s)$ read as $\varphi(W_s)$, $g(W_s)$, and we show in Step~5
that any other choice of the values on $D$ leads to indistinguishable processes. We also put
$a_0\defeq-\infty$, $a_{n+1}\defeq\infty$, so that $\R\setminus D$ is the disjoint union of the open
intervals $I_k\defeq(a_k,a_{k+1})$, $0\le k\le n$, and we fix the point $c\defeq a_1-1\in I_0$.

The functions $\varphi$ and $g$ are continuous at every point of $\R\setminus D$ and have finite
one-sided limits at every point of $D$. Hence they are Borel-measurable and bounded on every compact
interval $[\alpha,\beta]$: each point $x\in[\alpha,\beta]$ has a neighbourhood on which they are
bounded (by continuity if $x\notin D$, by the existence of finite one-sided limits if $x\in D$), and
finitely many such neighbourhoods cover $[\alpha,\beta]$. Being bounded with finitely many
discontinuities on each compact interval, they are Riemann integrable there; as usual
$\int_c^x\defeq-\int_x^c$ for $x<c$.

\emph{Step 1. $f(x)=f(c)+\int_c^x\varphi(u)\,du$ for every $x\in\R$.}
The function $F(x)\defeq f(c)+\int_c^x\varphi(u)\,du$ is continuous on $\R$, because $\varphi$ is
locally bounded, and at every point $x\in\R\setminus D$ it is differentiable with
$F'(x)=\varphi(x)=f'(x)$, because $\varphi$ is continuous at $x$. Thus $f-F$ is continuous on $\R$ and
has derivative zero on each open interval $I_k$; by the mean value theorem it is constant on each
$I_k$, hence, by continuity, on each closure $\ol{I_k}$. As $\ol{I_{k-1}}\cap\ol{I_k}=\{a_k\}$ for
$1\le k\le n$, these constants coincide, and $f-F$ is constant on $\R$. Since $(f-F)(c)=0$, we
conclude $f=F$.

\emph{Step 2. For every $x\in\R$,}
\begin{equation}\label{eq:3-6-24:phi}
  \varphi(x)=\varphi(c)+\int_c^x g(u)\,du+\sum_{k=1}^n\delta_k\,\ind_{(a_k,\infty)}(x).
\end{equation}
Let $G(x)\defeq\int_c^x g(u)\,du$ and $J(x)\defeq\sum_{k=1}^n\delta_k\ind_{(a_k,\infty)}(x)$, and consider
$\psi\defeq\varphi-G-J$. The function $G$ is continuous on $\R$, and differentiable with $G'(x)=g(x)=f''(x)$
at every $x\in\R\setminus D$; the function $J$ is constant on each $I_k$, where it equals
$\sum_{j\le k}\delta_j$. Therefore $\psi$ is differentiable on each $I_k$ with
$\psi'=f''-f''=0$, and there are constants $\gamma_0,\dots,\gamma_n$ with $\psi=\gamma_k$ on $I_k$.
Fix $k\in\{1,\dots,n\}$. Letting $x\uparrow a_k$ in $I_{k-1}$ and $x\downarrow a_k$ in $I_k$, we obtain
from the continuity of $G$
\[
  \gamma_{k-1}=f'(a_k-)-G(a_k)-\sum_{j\le k-1}\delta_j,
  \qquad
  \gamma_k=f'(a_k+)-G(a_k)-\sum_{j\le k}\delta_j=\gamma_{k-1},
\]
by the definition of $\delta_k$; moreover
$\psi(a_k)=\varphi(a_k)-G(a_k)-J(a_k)=f'(a_k-)-G(a_k)-\sum_{j\le k-1}\delta_j=\gamma_{k-1}$, because
$\ind_{(a_j,\infty)}(a_k)=1$ exactly for $j<k$. Hence $\psi$ is constant on $\R$, equal to
$\psi(c)=\varphi(c)$ (note that $G(c)=0$ and $J(c)=0$, as $c<a_1$). This is \eqref{eq:3-6-24:phi}.

\emph{Step 3. $f$ is the difference of two convex functions.}
With $r^+=\max\{r,0\}$ and $r^-=\max\{-r,0\}$ we define, for $x\in\R$,
\begin{align*}
  \varphi_1(x)&\defeq\varphi(c)+\int_c^x g^+(u)\,du+\sum_{k=1}^n\delta_k^+\,\ind_{(a_k,\infty)}(x),
  &
  \varphi_2(x)&\defeq\int_c^x g^-(u)\,du+\sum_{k=1}^n\delta_k^-\,\ind_{(a_k,\infty)}(x),\\
  f_1(x)&\defeq f(c)+\int_c^x\varphi_1(u)\,du,
  &
  f_2(x)&\defeq\int_c^x\varphi_2(u)\,du .
\end{align*}
The functions $g^\pm$ are nonnegative, Borel-measurable and locally bounded, so $\varphi_1$,
$\varphi_2$ are real-valued, nondecreasing and left-continuous (each $\ind_{(a_k,\infty)}$ is
left-continuous and the integrals are continuous in $x$). By \eqref{eq:3-6-24:phi},
$\varphi=\varphi_1-\varphi_2$, and by Step~1,
\begin{equation}\label{eq:3-6-24:diff}
  f=f_1-f_2 .
\end{equation}
For $i=1,2$ the function $x\mapsto\int_0^x\varphi_i(u)\,du$ is convex, with left-derivative
$(\varphi_i)_-=\varphi_i$, by part~(iii) of \bk{Problem 6.20} and the left-continuity of $\varphi_i$; the
function $f_i$ differs from it by an additive constant. Consequently
\begin{equation}\label{eq:3-6-24:Dminus}
  f_1,\ f_2\ \text{are convex, and}\quad D^-f_i(x)=\varphi_i(x);\qquad x\in\R,\ i=1,2,
\end{equation}
and \eqref{eq:3-6-24:diff} exhibits $f$ as the difference of two convex functions.

\emph{Step 4. The second derivative measures.}
Let $\mu_i$ be the second derivative measure of $f_i$ introduced in \bk{(6.47)}. By
\eqref{eq:3-6-24:Dminus} and the definition of $\varphi_i$, for $-\infty<a<b<\infty$,
\[
  \mu_1([a,b))=\varphi_1(b)-\varphi_1(a)=\int_a^b g^+(u)\,du+\sum_{k=1}^n\delta_k^+\,\ind_{[a,b)}(a_k),
\]
since $\ind_{(a_k,\infty)}(b)-\ind_{(a_k,\infty)}(a)=\ind_{[a,b)}(a_k)$; similarly for $\mu_2$, with
$g^-$ and $\delta_k^-$. The set function
$\nu_1(B)\defeq\int_B g^+(u)\,du+\sum_{k=1}^n\delta_k^+\ind_B(a_k)$, $B\in\B(\R)$, is a measure which
is finite on bounded sets and agrees with $\mu_1$ on the class of intervals $[a,b)$. This class is
closed under finite intersections, generates $\B(\R)$, and contains the sets $[-m,m)\uparrow\R$ of
finite measure; by Dynkin's $\pi$-$\lambda$ theorem (applied to the restrictions to $[-m,m)$, then
letting $m\to\infty$) we have $\mu_1=\nu_1$ on $\B(\R)$, and likewise for $\mu_2$. Hence, for every
Borel-measurable $h\colon\R\to[0,\infty)$,
\begin{equation}\label{eq:3-6-24:mu}
\begin{aligned}
  \int_{-\infty}^\infty h(x)\,\mu_1(dx)&=\int_{-\infty}^\infty h(x)g^+(x)\,dx+\sum_{k=1}^n\delta_k^+h(a_k),\\
  \int_{-\infty}^\infty h(x)\,\mu_2(dx)&=\int_{-\infty}^\infty h(x)g^-(x)\,dx+\sum_{k=1}^n\delta_k^-h(a_k).
\end{aligned}
\end{equation}

\emph{Step 5. The processes in \bk{(6.53)} do not depend on the values of $f'$, $f''$ on $D$.}
Fix $z\in\R$. For $\omega\in\Omega^*$ and $0\le t<\infty$, the occupation density formula \bk{(6.6)}
with $B=D$ gives
\begin{equation}\label{eq:3-6-24:null}
  \int_0^t\ind_D(W_s(\omega))\,ds=\int_D 2L_t(x,\omega)\,dx=0,
\end{equation}
because $D$ has Lebesgue measure zero. Let $\wt\varphi$, $\wt g\colon\R\to\R$ be any functions which
agree with $f'$, $f''$, respectively, on $\R\setminus D$. They are Borel-measurable, and
$\wt g-g$, $\wt\varphi-\varphi$ vanish off the finite set $D$, so they are bounded. By
\eqref{eq:3-6-24:null}, for every $\omega\in\Omega^*$,
\[
  \int_0^t\wt g(W_s)\,ds-\int_0^t g(W_s)\,ds=\int_0^t(\wt g-g)(W_s)\,\ind_D(W_s)\,ds=0;\qquad 0\le t<\infty,
\]
both integrals being well defined and finite, since $g$ and $\wt g$ are bounded on the compact range of
$s\mapsto W_s(\omega)$, $0\le s\le t$.

For the stochastic integrals, recall that under $P^z$ the integrator is $M\defeq W-z\in\cM_2^c$, with
$\qv{M}_t=t$ and $\qv{M,N}=\qv{W,N}$ for $N\in\cM_2^c$. The processes $\varphi(W)$ and $\wt\varphi(W)$
are progressively measurable (Borel functions of the continuous, adapted process $W$), and for each
$\omega$ they are bounded on $[0,t]$, because $\varphi$ and $\wt\varphi$ are bounded on the compact range
of the path on $[0,t]$; hence both belong to $\cP^*$, and their integrals are defined by
\bk{Definition 2.23}. Put $X\defeq(\wt\varphi-\varphi)(W)$. This process is progressively measurable and
bounded, say $X_s^2\le C$, so $E^z\int_0^TX_s^2\,ds\le CT<\infty$ and $X\in\cL^*(M)$; moreover
$X_s^2\le C\,\ind_D(W_s)$, so that, by \eqref{eq:3-6-24:null},
$\int_0^tX_s^2\,ds=0$ for all $t\ge0$ on $\Omega^*$. Now let $N\in\cM_2^c$. The process $Y\equiv1$ belongs to
$\cL^*(N)$, since $E^z\qv{N}_T=E^zN_T^2<\infty$. The Kunita--Watanabe inequality
(\bk{Proposition 2.14}, applied to $M,N\in\cM_2^c$, $X\in\cL^*(M)$, $Y\in\cL^*(N)$) gives, with
$\check\xi$ the total variation process of $\qv{W,N}$ and
$\Delta_t\defeq\int_0^t\wt\varphi(W_s)\,d\qv{W,N}_s-\int_0^t\varphi(W_s)\,d\qv{W,N}_s
=\int_0^tX_s\,d\qv{W,N}_s$, for $P^z$-a.e.\ $\omega\in\Omega^*$,
\[
  \abs{\Delta_t}\le\int_0^t\abs{X_s}\,d\check\xi_s
  \le\Bigl(\int_0^tX_s^2\,ds\Bigr)^{1/2}\qv{N}_t^{1/2}=0;\qquad 0\le t<\infty .
\]
(The two Lebesgue--Stieltjes integrals exist because the integrands are bounded on $[0,t]$ and
$\qv{W,N}$ has finite variation on $[0,t]$.) By \bk{Proposition 2.24}, the local martingale
$\Phi\defeq\int_0^{\cdot}\varphi(W_s)\,dW_s\in\cM^{c,\loc}$ satisfies
$\qv{\Phi,N}_t=\int_0^t\varphi(W_s)\,d\qv{W,N}_s$, $0\le t<\infty$, a.s.\ $P^z$, for every
$N\in\cM_2^c$; by the preceding display, it therefore satisfies
$\qv{\Phi,N}_t=\int_0^t\wt\varphi(W_s)\,d\qv{W,N}_s$, $0\le t<\infty$, a.s.\ $P^z$, for every
$N\in\cM_2^c$. The uniqueness assertion of \bk{Proposition 2.24}, applied to the integrand
$\wt\varphi(W)\in\cP^*$, yields that $\Phi$ and $\int_0^\cdot\wt\varphi(W_s)\,dW_s$ are indistinguishable:
\[
  \int_0^t\wt\varphi(W_s)\,dW_s=\int_0^t\varphi(W_s)\,dW_s;\qquad 0\le t<\infty,\ \text{a.s. }P^z .
\]
Thus it suffices to prove \bk{(6.53)} for the particular choice \eqref{eq:3-6-24:conv}.

\emph{Step 6. Proof of \bk{(6.53)}.}
Fix $z\in\R$. The functions $f_1$, $f_2$ are convex, with left-derivatives $\varphi_1$, $\varphi_2$ and
second derivative measures $\mu_1$, $\mu_2$; \bk{Theorem 6.22} gives an event $\Omega_z\subseteq\Omega^*$
with $P^z(\Omega_z)=1$ on which, for $i=1,2$,
\begin{equation}\label{eq:3-6-24:ito}
  f_i(W_t)=f_i(z)+\int_0^t\varphi_i(W_s)\,dW_s+\int_{-\infty}^\infty L_t(x)\,\mu_i(dx);\qquad 0\le t<\infty .
\end{equation}
Fix $\omega\in\Omega_z$ and $t\ge0$. The function $x\mapsto L_t(x,\omega)$ is continuous and nonnegative, so
\eqref{eq:3-6-24:mu} applies to it; and the formula \bk{(6.7)}, which holds on $\Omega^*$ for every
nonnegative Borel-measurable function (\bk{Problem 6.7}), applied to $g^+$ and to $g^-$, gives
\[
  \int_{-\infty}^\infty L_t(x)g^{\pm}(x)\,dx=\frac12\int_0^tg^{\pm}(W_s)\,ds<\infty,
\]
the finiteness because $g^\pm$ are bounded on the compact range of the path on $[0,t]$. Hence
\begin{equation}\label{eq:3-6-24:lt}
\begin{aligned}
  \int_{-\infty}^\infty L_t(x)\,\mu_1(dx)&=\frac12\int_0^tg^+(W_s)\,ds+\sum_{k=1}^n\delta_k^+L_t(a_k),\\
  \int_{-\infty}^\infty L_t(x)\,\mu_2(dx)&=\frac12\int_0^tg^-(W_s)\,ds+\sum_{k=1}^n\delta_k^-L_t(a_k),
\end{aligned}
\end{equation}
and all the terms are finite. We now subtract the identity \eqref{eq:3-6-24:ito} for $i=2$ from that
for $i=1$. On the left we obtain $f(W_t)$ and the constant term is $f(z)$, by \eqref{eq:3-6-24:diff}.
The processes $\varphi_1(W)$, $\varphi_2(W)$ belong to $\cP^*$ (they are progressively measurable, and
bounded on $[0,t]$ for each path, because $\varphi_i$ is monotone, hence locally bounded), so by the
linearity \bk{(2.17)} of the stochastic integral, which remains valid for integrands in $\cP^*$ (see
the discussion preceding \bk{Proposition 2.24}), and by $\varphi_1-\varphi_2=\varphi$,
\[
  \int_0^t\varphi_1(W_s)\,dW_s-\int_0^t\varphi_2(W_s)\,dW_s=\int_0^t\varphi(W_s)\,dW_s;
  \qquad 0\le t<\infty,\ \text{a.s. }P^z .
\]
Finally, by \eqref{eq:3-6-24:lt}, $g^+-g^-=g$ and $\delta_k^+-\delta_k^-=\delta_k$,
\[
  \int_{-\infty}^\infty L_t(x)\,\mu_1(dx)-\int_{-\infty}^\infty L_t(x)\,\mu_2(dx)
  =\frac12\int_0^tg(W_s)\,ds+\sum_{k=1}^n\delta_kL_t(a_k).
\]
Altogether, a.s.\ $P^z$,
\[
  f(W_t)=f(z)+\int_0^t\varphi(W_s)\,dW_s+\frac12\int_0^tg(W_s)\,ds
  +\sum_{k=1}^nL_t(a_k)\,[f'(a_k+)-f'(a_k-)];\qquad 0\le t<\infty,
\]
which, in view of the conventions \eqref{eq:3-6-24:conv} and of Step~5, is \bk{(6.53)}.

\rmk{The signed measure $\mu=\mu_1-\mu_2$, that is,
$\mu(dx)=g(x)\,dx+\sum_{k}\delta_k\,\eps_{a_k}(dx)$ with $\eps_a$ the unit mass at $a$, is the one
attached to $f=f_1-f_2$ in \bk{Corollary 6.23}; Step~6 is the proof of that corollary in the present
case, with $\mu$ computed explicitly. As a check of the constants, $f(x)=\abs{x-a}$ has $D=\{a\}$,
$f'=\pm1$ on the two sides of $a$, $f''=0$ and $f'(a+)-f'(a-)=2$, and \bk{(6.53)} reduces to the
Tanaka formula \bk{(6.13)} with the left-continuous $\sgn$ of \bk{(6.14)}; that the value of the
integrand at $a$ is immaterial is \bk{Remark 6.9}, which Step~5 extends.}
\end{kex}

\begin{kex}{Exercise}{6.25}{}
Throughout, $L=\{L_t(x);\,0\le t<\infty,\ x\in\R\}$ is the Brownian local time of \bk{Definition 6.3}, which exists by \bk{Theorem 6.11} and is the random field appearing in \bk{(6.49)}. We fix $a\in\R$ and $z\in\R$, and we apply \bk{Theorem 6.22} to the three convex functions
\[
f_1(x)\defeq (x-a)^+,\qquad f_2(x)\defeq (x-a)^-=(a-x)^+,\qquad f_3(x)\defeq \abs{x-a}=f_1(x)+f_2(x);
\]
$f_1$ and $f_2$ are convex as maxima of two affine functions, and $f_3$ as their sum. For each of them we have to identify the left-hand derivative $D^-f_i$ of \bk{(6.40)} and the second derivative measure $\mu_i$ of \bk{(6.47)}.

\emph{Step 1: the left-hand derivatives.} For $x>a$ the function $f_1$ coincides with $y\mapsto y-a$ on a neighbourhood of $x$, and for $x<a$ with the zero function, so that $D^-f_1(x)=1$ for $x>a$ and $D^-f_1(x)=0$ for $x<a$. At $x=a$ we have $f_1(a+h)-f_1(a)=0$ for every $h<0$, whence $D^-f_1(a)=0$. Similarly, $f_2$ coincides with $y\mapsto a-y$ near every $x<a$ and with zero near every $x>a$, and $h^{-1}[f_2(a+h)-f_2(a)]=h^{-1}(-h)=-1$ for $h<0$. Therefore, for every $x\in\R$,
\begin{equation}\label{eq:3-6-25:a}
\begin{gathered}
D^-f_1(x)=\ind_{(a,\infty)}(x),\qquad D^-f_2(x)=-\ind_{(-\infty,a]}(x),\\
D^-f_3(x)=\ind_{(a,\infty)}(x)-\ind_{(-\infty,a]}(x)=\sgn(x-a),
\end{gathered}
\end{equation}
where the third identity follows from the first two by addition, and $\sgn$ is the left-continuous function of \bk{(6.14)}: $\sgn(y)=1$ for $y>0$ and $\sgn(y)=-1$ for $y\le 0$.

\emph{Step 2: the second derivative measures.} Let $\delta_a$ be the unit point mass at $a$. For $-\infty<c<b<\infty$, \eqref{eq:3-6-25:a} gives
\begin{align*}
D^-f_1(b)-D^-f_1(c)&=\ind_{(a,\infty)}(b)-\ind_{(a,\infty)}(c)=\ind_{\{c\le a<b\}}=\delta_a([c,b)),\\
D^-f_2(b)-D^-f_2(c)&=\ind_{(-\infty,a]}(c)-\ind_{(-\infty,a]}(b)=\ind_{\{c\le a<b\}}=\delta_a([c,b)),
\end{align*}
and by addition $D^-f_3(b)-D^-f_3(c)=2\delta_a([c,b))$. The intervals $[c,b)$, together with the empty set, form a $\pi$-system which generates $\B(\R)$ and contains the sets $[-n,n)$, $n\ge 1$, which increase to $\R$; two measures on $\B(\R)$ which agree on this $\pi$-system and are finite on each $[-n,n)$ coincide, by Dynkin's $\pi$-$\lambda$ theorem. Hence the measures defined by \bk{(6.47)} are
\begin{equation}\label{eq:3-6-25:b}
\mu_1=\delta_a,\qquad \mu_2=\delta_a,\qquad \mu_3=2\delta_a,
\end{equation}
and consequently, for every $t\ge 0$ and every $\omega$,
\begin{equation}\label{eq:3-6-25:c}
\int_{-\infty}^\infty L_t(x)\,\mu_1(dx)=\int_{-\infty}^\infty L_t(x)\,\mu_2(dx)=L_t(a),\qquad
\int_{-\infty}^\infty L_t(x)\,\mu_3(dx)=2L_t(a).
\end{equation}

\emph{Step 3: the Tanaka formulas.} The integrands in \eqref{eq:3-6-25:a} are bounded Borel functions of $x$, so the processes $D^-f_i(W_s)$ are bounded and progressively measurable, and their stochastic integrals with respect to $W$ are well defined. Substituting \eqref{eq:3-6-25:a} and \eqref{eq:3-6-25:c} into \bk{(6.49)}, written for $f_1$, $f_2$ and $f_3$ in turn, we obtain that each of the following three identities holds for all $0\le t<\infty$, a.s.\ $P^z$:
\begin{align*}
(W_t-a)^+&=(z-a)^+ +\int_0^t \ind_{(a,\infty)}(W_s)\,dW_s+L_t(a),\\
(W_t-a)^-&=(z-a)^- -\int_0^t \ind_{(-\infty,a]}(W_s)\,dW_s+L_t(a),\\
\abs{W_t-a}&=\abs{z-a}+\int_0^t \sgn(W_s-a)\,dW_s+2L_t(a).
\end{align*}
Solving each of them for the local time gives
\begin{align*}
L_t(a)&=(W_t-a)^+-(z-a)^+-\int_0^t \ind_{(a,\infty)}(W_s)\,dW_s,\\
L_t(a)&=(W_t-a)^--(z-a)^-+\int_0^t \ind_{(-\infty,a]}(W_s)\,dW_s,\\
2L_t(a)&=\abs{W_t-a}-\abs{z-a}-\int_0^t \sgn(W_s-a)\,dW_s
\end{align*}
for $0\le t<\infty$, a.s.\ $P^z$, and these are precisely \bk{(6.11)}, \bk{(6.12)} and \bk{(6.13)}. Since $z\in\R$ and $a\in\R$ were arbitrary, the three Tanaka formulas are established.

\rmk{The left-continuous convention \bk{(6.14)} for $\sgn$ is exactly the one produced by the left-hand derivative of $x\mapsto\abs{x-a}$ in \eqref{eq:3-6-25:a}; likewise the open half-line in \bk{(6.11)} and the closed one in \bk{(6.12)} are forced by $D^-f_1(a)=0$ and $D^-f_2(a)=-1$. The derivation is a consistency check rather than a new proof of \bk{Proposition 6.8}: the proof of \bk{Theorem 6.22} rests on the identity \bk{(6.7)} and on the continuity of $x\mapsto L_t(x)$, and the local time with these properties was constructed in \bk{Theorem 6.11} from the right-hand side of \bk{(6.11)}.}
\end{kex}

\begin{kex}{Problem}{6.30}{}
We work in the setting of the book's Section 3.6: $\Omega=C[0,\infty)$, $W$ is the coordinate
process, $L$ is the Brownian local time of \bk{Theorem 6.11}, and $\Omega^*\in\F$ is the event of
\bk{Definition 6.3}, with $P^z(\Omega^*)=1$ for every $z\in\R$. For each $\omega\in\Omega^*$ the
function $(t,a)\mapsto L_t(a,\omega)$ is continuous, and by \bk{Problem 6.7}~(i) the identity
\begin{equation}\label{eq:3-6-30:occ}
  \int_0^t g(W_s(\omega))\,ds=2\int_{-\infty}^\infty g(a)L_t(a,\omega)\,da;\qquad 0\le t<\infty,
\end{equation}
holds for every Borel-measurable $g\colon\R\to[0,\infty)$ (both sides possibly infinite). Moreover,
for $\omega\in\Omega^*$ and every $a\in\R$ the function $t\mapsto L_t(a,\omega)$ is nondecreasing:
if $s<t$, then \bk{(6.6)} gives $\int_B\bigl(L_t(a,\omega)-L_s(a,\omega)\bigr)\,da
=\tfrac12\meas\{s<u\le t;\ W_u(\omega)\in B\}\ge0$ for every bounded $B\in\B(\R)$, so that
$L_t(\cdot,\omega)\ge L_s(\cdot,\omega)$ Lebesgue-almost everywhere, hence everywhere by continuity.

\textbf{Proof of \bk{(6.55)}.} Fix $x\in\R$. The mapping $(s,\omega)\mapsto f(W_s(\omega))=f(\omega(s))$
is $\B([0,\infty))\otimes\B(C[0,\infty))$-measurable, being the composition of the Borel function
$f$ with the continuous mapping $(s,\omega)\mapsto\omega(s)$; by Tonelli's theorem
$\int_0^\infty f(W_s)\,ds$ is a random variable with values in $[0,\infty]$, so the event in
\bk{(6.55)} belongs to $\F$.

Define $\wt L_t(a,\omega)\defeq L_t(a,\omega)\ind_{\Omega^*}(\omega)$ and
\begin{equation*}
  L_\infty(a,\omega)\defeq\sup_{k\in\N}\wt L_k(a,\omega)=\lim_{t\to\infty}\wt L_t(a,\omega)\in[0,\infty];
  \qquad a\in\R,\ \omega\in\Omega.
\end{equation*}
For every $k$, the function $\wt L_k(a,\omega)$ is continuous in $a$ for fixed $\omega$ and
$\F$-measurable in $\omega$ for fixed $a$; hence it is the pointwise limit, as $m\to\infty$, of the
$\B(\R)\otimes\F$-measurable functions
$(a,\omega)\mapsto\wt L_k\bigl(2^{-m}\lfloor 2^m a\rfloor,\omega\bigr)$, and is itself
$\B(\R)\otimes\F$-measurable. Consequently $L_\infty$ is $\B(\R)\otimes\F$-measurable. By
\bk{Problem 6.18} we have $P^x[L_\infty(a)=\infty]=1$ for every $a\in\R$, and Tonelli's theorem
yields
\begin{equation*}
  E^x\Bigl[\meas\{a\in\R;\ L_\infty(a,\cdot)<\infty\}\Bigr]
  =\int_{-\infty}^\infty P^x\bigl[L_\infty(a)<\infty\bigr]\,da=0.
\end{equation*}
Thus there is an event $\Omega^{**}\subseteq\Omega^*$ with $P^x(\Omega^{**})=1$ such that, for every
$\omega\in\Omega^{**}$, the set $N_\omega\defeq\{a;\ L_\infty(a,\omega)<\infty\}$ has Lebesgue
measure zero. Fix $\omega\in\Omega^{**}$. Letting $t=k\to\infty$ in \eqref{eq:3-6-30:occ} with $g=f$
and using the monotone convergence theorem on both sides (recall that $L_k(a,\omega)\uparrow
L_\infty(a,\omega)$), we obtain
\begin{equation*}
  \int_0^\infty f(W_s(\omega))\,ds=2\int_{-\infty}^\infty f(a)L_\infty(a,\omega)\,da
  \ \ge\ 2\int_{\{f>0\}\setminus N_\omega}f(a)\cdot\infty\,da=\infty,
\end{equation*}
because the integrand equals $+\infty$ on the set $\{f>0\}\setminus N_\omega$, which has positive
Lebesgue measure. This proves \bk{(6.55)}.

\rmk{For the second assertion the statement tacitly requires $\|f\|_1<\infty$. Indeed, a function
$f$ with compact support is locally integrable if and only if $\|f\|_1<\infty$; if
$\|f\|_1=\infty$, then \bk{Proposition 6.27} shows that
$P^0\bigl[\int_0^t f(W_s)\,ds<\infty;\ \forall\,0\le t<\infty\bigr]=0$, so the $X^{(n)}$ are not
real-valued continuous processes and the limit $X$ is not defined either. We assume
$0<\|f\|_1<\infty$ from now on. (The positivity of $\|f\|_1$ is the hypothesis
$\meas\{f>0\}>0$; it is not needed for the convergence.)}

\textbf{Proof of \bk{(6.56)}.} Let $K>0$ be such that $f=0$ off $[-K,K]$, and set
$g_n(y)\defeq\sqrt n\,f(\sqrt n\,y)$, so that $g_n$ is Borel-measurable and nonnegative.

\emph{Step 1: the processes as functionals of the path.} For $n\ge1$ let
\begin{equation*}
  A_n\defeq\Bigl\{w\in C[0,\infty);\ \int_0^m g_n(w(u))\,du<\infty,\ \forall\,m\in\N\Bigr\},
\end{equation*}
and define $\Phi_n\colon C[0,\infty)\to C[0,\infty)$ by
\begin{equation*}
  \Phi_n(w)(t)\defeq\ind_{A_n}(w)\int_0^t g_n(w(u))\,du;\qquad 0\le t<\infty.
\end{equation*}
For $w\in A_n$ the function $u\mapsto g_n(w(u))$ is integrable on every bounded interval, so
$\Phi_n(w)$ is indeed continuous. As in the proof of \bk{(6.55)}, Tonelli's theorem shows that
$w\mapsto\int_0^t g_n(w(u))\,du\in[0,\infty]$ is $\B(C[0,\infty))$-measurable for every $t$; hence
$A_n\in\B(C[0,\infty))$ and each coordinate $w\mapsto\Phi_n(w)(t)$ is $\B(C[0,\infty))$-measurable.
Since $\B(C[0,\infty))$ is generated by the finite-dimensional cylinder sets
(\bk{Problem 2.4.2}), $\Phi_n$ is $\B(C[0,\infty))/\B(C[0,\infty))$-measurable.

Let $S_n\colon\Omega\to\Omega$ be the scaling map $(S_n\omega)(u)\defeq n^{-1/2}\omega(nu)$; it is
continuous for the metric \bk{(2.4.1)}, hence Borel-measurable. The change of variable $s=nu$ gives,
for every $\omega\in\Omega$ and $t\ge0$,
\begin{equation}\label{eq:3-6-30:scal}
  \frac1{\sqrt n}\int_0^{nt}f(\omega(s))\,ds=\sqrt n\int_0^t f\bigl(\sqrt n\,(S_n\omega)(u)\bigr)\,du
  =\int_0^t g_n\bigl((S_n\omega)(u)\bigr)\,du .
\end{equation}
Next, for $\omega\in\Omega^*$ we apply \eqref{eq:3-6-30:occ} to $g_n$ and substitute $a=\sqrt n\,b$:
\begin{equation}\label{eq:3-6-30:Y}
\begin{split}
  \int_0^t g_n(\omega(u))\,du&=2\int_{-\infty}^\infty\sqrt n\,f(\sqrt n\,b)L_t(b,\omega)\,db
  =2\int_{-K}^{K}f(a)\,L_t\Bigl(\frac a{\sqrt n},\omega\Bigr)\,da\\
  &\le2\|f\|_1\max_{\abs b\le K}L_t(b,\omega)<\infty .
\end{split}
\end{equation}
Hence $\Omega^*\subseteq A_n$, and therefore $P^0(A_n)=1$.

Under $P^0$ the coordinate process $W$ is a standard Brownian motion, and so is $S_nW=
\{n^{-1/2}W_{nu};\ 0\le u<\infty\}$ by \bk{Lemma 2.9.4}~(i). The two continuous processes thus have
the same finite-dimensional distributions, and since the cylinder sets form a $\pi$-system which
generates $\B(C[0,\infty))$, Dynkin's $\pi$-$\lambda$ theorem gives
\begin{equation}\label{eq:3-6-30:law}
  P^0\bigl[S_n\in\Gamma\bigr]=P^0(\Gamma);\qquad\Gamma\in\B(C[0,\infty)).
\end{equation}
In particular $P^0[S_n\in A_n]=P^0(A_n)=1$, and by \eqref{eq:3-6-30:scal} we have on the event
$\{S_n\in A_n\}$ that $X^{(n)}_t=\Phi_n(S_n)(t)$ for all $t\ge0$. So $X^{(n)}$ is, off a
$P^0$-null event (on which we set it equal to zero), a continuous process, and
\begin{equation*}
  X^{(n)}=\Phi_n\circ S_n,\qquad P^0\text{-a.s.}
\end{equation*}
We also introduce the continuous processes $Y^{(n)}\defeq\Phi_n(W)$, i.e.,
$Y^{(n)}(\omega)=\Phi_n(\omega)$, and $X_t\defeq2\|f\|_1L_t(0)\ind_{\Omega^*}$.

\emph{Step 2: $X^{(n)}$ and $Y^{(n)}$ have the same law under $P^0$.} For
$\Gamma\in\B(C[0,\infty))$ we have $\Phi_n^{-1}(\Gamma)\in\B(C[0,\infty))$ by Step~1, and
\eqref{eq:3-6-30:law} yields
\begin{equation*}
  P^0\bigl[X^{(n)}\in\Gamma\bigr]=P^0\bigl[S_n\in\Phi_n^{-1}(\Gamma)\bigr]
  =P^0\bigl(\Phi_n^{-1}(\Gamma)\bigr)=P^0\bigl[Y^{(n)}\in\Gamma\bigr].
\end{equation*}

\emph{Step 3: $Y^{(n)}\to X$ almost surely in $C[0,\infty)$.} Fix $\omega\in\Omega^*$ and $T>0$.
By \eqref{eq:3-6-30:Y} and $\Omega^*\subseteq A_n$,
$Y^{(n)}_t(\omega)=2\int_{-K}^Kf(a)L_t(a/\sqrt n,\omega)\,da$, and therefore
\begin{equation*}
  \max_{0\le t\le T}\abs[\big]{Y^{(n)}_t(\omega)-X_t(\omega)}
  \le2\int_{-K}^Kf(a)\max_{0\le t\le T}\abs[\Big]{L_t\Bigl(\frac a{\sqrt n},\omega\Bigr)-L_t(0,\omega)}\,da
  \le2\|f\|_1\,\delta_n(\omega),
\end{equation*}
where
$\delta_n(\omega)\defeq\max\bigl\{\abs{L_t(b,\omega)-L_t(0,\omega)};\ 0\le t\le T,\ \abs b\le
K/\sqrt n\bigr\}$. The function $(t,b)\mapsto L_t(b,\omega)$ is continuous, hence uniformly
continuous on the compact set $[0,T]\times[-K,K]$, so $\delta_n(\omega)\to0$ as $n\to\infty$. Since
$T$ is arbitrary, $Y^{(n)}(\omega)\to X(\omega)$ uniformly on bounded intervals, that is,
$\rho\bigl(Y^{(n)}(\omega),X(\omega)\bigr)\to0$ for the metric $\rho$ of \bk{(2.4.1)}, for every
$\omega\in\Omega^*$.

\emph{Step 4: conclusion.} Let $F\colon C[0,\infty)\to\R$ be bounded and continuous. By Step~3,
$F(Y^{(n)})\to F(X)$, $P^0$-a.s., and the bounded convergence theorem together with Step~2 gives
\begin{equation*}
  \lim_{n\to\infty}E^0F\bigl(X^{(n)}\bigr)=\lim_{n\to\infty}E^0F\bigl(Y^{(n)}\bigr)=E^0F(X).
\end{equation*}
According to \bk{Definition 2.4.4}, this is the convergence \bk{(6.56)} under $P^0$.

\rmk{As a check, take $t=1$. By \bk{Theorem 6.17}, $2L_1(0)$ has under $P^0$ the law of
$\max_{0\le s\le1}W_s$, which by the reflection principle is that of $\abs{W_1}$; so \bk{(6.56)}
predicts that $n^{-1/2}\int_0^nf(W_s)\,ds$ converges in distribution to $\|f\|_1\abs{W_1}$. The
means are consistent with this: $E^0\bigl[n^{-1/2}\int_0^nf(W_s)\,ds\bigr]
=n^{-1/2}\int_0^n\!\int f(y)(2\pi s)^{-1/2}e^{-y^2/2s}\,dy\,ds\to\|f\|_1\sqrt{2/\pi}
=\|f\|_1E^0\abs{W_1}$, since $\int_0^n(2\pi s)^{-1/2}\,ds=\sqrt{2n/\pi}$ and the exponential factor
tends to $1$ boundedly for $\abs y\le K$ after the substitution $s=nu$.}
\end{kex}

\solsection{3.7}{Local Time for Continuous Semimartingales}

\begin{kex}{Problem}{7.7}{}
The book's Section 3.8 only says that the solution is a slight modification of that of
\bk{Problem 6.12}, with \bk{Lemma 7.5} providing the continuity in $a$ of the integrand on the
left-hand side; we give the complete argument here.

\emph{Two points of interpretation.} First, for each $a\in\R$ the process
$I_t(a)\defeq\int_0^t\ind_{(a,\infty)}(X_s)\,dM_s$, $0\le t<\infty$, is defined only up to
indistinguishability (the integrand is bounded and progressively measurable, hence in $\cP^*$),
so the left-hand side has no meaning until a version of the random field $I$ has been chosen.
As in the solution of \bk{Problem 6.12}, we take the continuous modification $\hat I$ of
\bk{Lemma 7.5}: there is an event $\Omega^*$ with $P(\Omega^*)=1$ such that
$(t,a)\mapsto\hat I_t(a,\omega)$ is continuous on $[0,\infty)\times\R$ for every
$\omega\in\Omega^*$, and $\hat I_t(a)=I_t(a)$ a.s.\ for every $(t,a)$; each $\hat I_t(a)$ is
$\F_t$-measurable. Since $\{\F_t\}$ satisfies the usual conditions, $\Omega^*\in\F_0$, and after
replacing $\hat I$ by $\hat I\,\ind_{\Omega^*}$ we may assume that $\hat I\equiv 0$ off $\Omega^*$.
Secondly, the hypotheses as printed do not guarantee that either side is finite: a
$\sigma$-finite measure need not be finite on compact sets (see the Remark at the end). We
therefore assume in addition that
\begin{equation}\label{eq:3-7-07:fin}
  \int_{-\infty}^{\infty}h(a)\,\mu(da)<\infty ,
\end{equation}
which holds in particular whenever $\mu$ is finite on compact sets, as is the second derivative
measure of a convex function to which the problem is applied in the proof of \bk{Theorem 7.1}.

\begin{claim}
Assume \eqref{eq:3-7-07:fin}, and set
$F(x)\defeq\int_{-\infty}^{\infty}h(a)\ind_{(a,\infty)}(x)\,\mu(da)$, $x\in\R$. Then $F$ is a
bounded Borel function, the integral $\int_{-\infty}^\infty h(a)\hat I_t(a)\,\mu(da)$ is
well defined and finite for all $t\ge0$ and $\omega\in\Omega$, and
\begin{equation}\label{eq:3-7-07:claim}
  \int_{-\infty}^{\infty}h(a)\,\hat I_t(a)\,\mu(da)=\int_0^tF(X_s)\,dM_s;\qquad 0\le t<\infty,
  \quad\text{a.s. }P.
\end{equation}
\end{claim}

\begin{proof}
\emph{Step 1: notation.} By \eqref{eq:3-7-07:fin}, $\nu(A)\defeq\int_Ah(a)\,\mu(da)$,
$A\in\B(\R)$, is a finite measure on $(\R,\B(\R))$, and integrals of nonnegative or bounded Borel
functions against $h(a)\,\mu(da)$ are integrals against $\nu$. Choose real numbers $c<d$ with
$h=0$ off $[c,d]$; then $\nu(\R\setminus[c,d])=0$. We have $F(x)=\nu\bigl((-\infty,x)\bigr)$, so
$F$ is nondecreasing, hence Borel measurable, and $0\le F\le\nu(\R)$. The process $F(X)$ is
therefore bounded and, being a Borel function of a continuous adapted process, progressively
measurable; thus $F(X)\in\cP^*$ and the right-hand side of \eqref{eq:3-7-07:claim} is a
well-defined element of $\cM^{c,\loc}$.

For $n\ge1$ let $a^{(n)}_k\defeq c+k2^{-n}(d-c)$, $k=0,1,\dots,2^n$, and
$A^{(n)}_0\defeq\{c\}$, $A^{(n)}_k\defeq(a^{(n)}_{k-1},a^{(n)}_k]$ for $1\le k\le 2^n$; these
sets partition $[c,d]$. Define $\varphi_n\colon[c,d]\to[c,d]$ by $\varphi_n(a)\defeq a^{(n)}_k$
for $a\in A^{(n)}_k$. Then
\begin{equation}\label{eq:3-7-07:phi}
  a\le\varphi_n(a)\le a+2^{-n}(d-c);\qquad a\in[c,d],\ n\ge1 .
\end{equation}

\emph{Step 2: the left-hand side.} Define
\begin{equation*}
  L_t\defeq\int_{[c,d]}\hat I_t(a)\,\nu(da),\qquad
  L^{(n)}_t\defeq\int_{[c,d]}\hat I_t\bigl(\varphi_n(a)\bigr)\,\nu(da)
  =\sum_{k=0}^{2^n}\nu\bigl(A^{(n)}_k\bigr)\,\hat I_t\bigl(a^{(n)}_k\bigr);\qquad 0\le t<\infty .
\end{equation*}
For $\omega\in\Omega^*$ and $T<\infty$, the function $(t,a)\mapsto\hat I_t(a,\omega)$ is continuous,
hence bounded and uniformly continuous, on the compact set $[0,T]\times[c,d]$. Consequently
$L_t(\omega)$ is well defined and finite (and equal to the left-hand side of
\eqref{eq:3-7-07:claim}, because $\nu$ does not charge $\R\setminus[c,d]$); $t\mapsto L_t(\omega)$
is continuous on $[0,T]$ by the bounded convergence theorem, $\nu$ being finite; and by
\eqref{eq:3-7-07:phi},
\begin{equation*}
  \sup_{0\le t\le T}\abs[\big]{L^{(n)}_t(\omega)-L_t(\omega)}
  \le\nu(\R)\sup_{\substack{0\le t\le T,\ a,b\in[c,d]\\ \abs{a-b}\le2^{-n}(d-c)}}
  \abs[\big]{\hat I_t(a,\omega)-\hat I_t(b,\omega)}\xrightarrow[n\to\infty]{}0 .
\end{equation*}
For $\omega\notin\Omega^*$ everything vanishes. Thus $L^{(n)}_t\to L_t$ everywhere on $\Omega$; as
each $L^{(n)}_t$ is a finite linear combination of $\F_t$-measurable random variables, $L_t$ is
$\F_t$-measurable, and $L$ is a continuous adapted process.

\emph{Step 3: the approximating sums as stochastic integrals.} Put
\begin{equation*}
  F_n(x)\defeq\sum_{k=0}^{2^n}\nu\bigl(A^{(n)}_k\bigr)\,\ind_{(a^{(n)}_k,\infty)}(x)
  =\int_{[c,d]}\ind_{(\varphi_n(a),\infty)}(x)\,\nu(da);\qquad x\in\R .
\end{equation*}
This is a Borel function with $0\le F_n\le\nu(\R)$, so $F_n(X)\in\cP^*$. Fix $t\ge0$. Since
$\hat I_t(a^{(n)}_k)=I_t(a^{(n)}_k)$ a.s.\ for each of the finitely many $k$, the linearity of
the stochastic integral in the integrand (property \bk{(2.17)}, valid for
$M\in\cM^{c,\loc}$ and integrands in $\cP^*$ by the discussion following \bk{Definition 2.23})
gives
\begin{equation}\label{eq:3-7-07:sum}
  L^{(n)}_t=\sum_{k=0}^{2^n}\nu\bigl(A^{(n)}_k\bigr)\int_0^t\ind_{(a^{(n)}_k,\infty)}(X_s)\,dM_s
  =\int_0^tF_n(X_s)\,dM_s,\qquad\text{a.s. }P .
\end{equation}

\emph{Step 4: convergence of the integrands.} We show that $F_n(x)\to F(x)$ for \emph{every}
$x\in\R$. Fix $x\in\R$ and $a\in[c,d]$. If $a<x$, then by \eqref{eq:3-7-07:phi} we have
$\varphi_n(a)<x$ for all $n$ with $2^{-n}(d-c)<x-a$, so
$\ind_{(\varphi_n(a),\infty)}(x)=1=\ind_{(a,\infty)}(x)$ for all large $n$. If $a\ge x$, then
$\varphi_n(a)\ge a\ge x$ and $\ind_{(\varphi_n(a),\infty)}(x)=0=\ind_{(a,\infty)}(x)$ for all $n$.
Hence $\ind_{(\varphi_n(a),\infty)}(x)\to\ind_{(a,\infty)}(x)$ for every $a\in[c,d]$, and the
bounded convergence theorem for the finite measure $\nu$ yields
\begin{equation*}
  \lim_{n\to\infty}F_n(x)=\int_{[c,d]}\ind_{(a,\infty)}(x)\,\nu(da)=F(x);\qquad x\in\R .
\end{equation*}
(It is for this step that $\varphi_n$ was chosen to approximate from the right: the convergence
holds at every $x$, including the atoms of $\nu$, although it need not be uniform in $x$.)

\emph{Step 5: passage to the limit.} Fix $t\ge0$. For every $\omega$ outside the null set on which
$\qv{M}_t=\infty$ or $s\mapsto\qv{M}_s$ fails to be continuous and nondecreasing, the functions
$s\mapsto\abs{F_n(X_s(\omega))-F(X_s(\omega))}^2$ are bounded by $\nu(\R)^2$ and converge to
zero pointwise on $[0,t]$ by Step 4, so the bounded convergence theorem for the finite measure
$d\qv{M}_s(\omega)$ on $[0,t]$ gives
\begin{equation*}
  \lim_{n\to\infty}\int_0^t\abs[\big]{F_n(X_s)-F(X_s)}^2\,d\qv{M}_s=0,\qquad\text{a.s. }P,
\end{equation*}
and in particular in probability. \bk{Proposition 2.26} (with the stopping time $T\equiv t$)
now shows that $\int_0^tF_n(X_s)\,dM_s\to\int_0^tF(X_s)\,dM_s$ in probability as $n\to\infty$.
On the other hand, by \eqref{eq:3-7-07:sum} and Step 2, $\int_0^tF_n(X_s)\,dM_s=L^{(n)}_t\to L_t$
a.s. Limits in probability being a.s.\ unique, $L_t=\int_0^tF(X_s)\,dM_s$ a.s., for every fixed
$t\ge0$. Both sides are continuous processes ($L$ by Step 2, the stochastic integral by
\bk{Definition 2.23}); they agree a.s.\ simultaneously at all rational $t$, hence at all
$t\in[0,\infty)$. This is \eqref{eq:3-7-07:claim}.
\end{proof}

Since $F(X_s)=\int_{-\infty}^\infty h(a)\ind_{(a,\infty)}(X_s)\,\mu(da)$, the identity
\eqref{eq:3-7-07:claim} is the assertion of the problem, with the null set not depending on $t$.

\rmk{\emph{(a) The choice of the version on the left-hand side is immaterial.} Let $t\ge0$ and let
$J(a,\omega)$ be any $\B(\R)\otimes\F$-measurable function with $J(a,\cdot)=I_t(a)$ a.s.\ for every
$a\in\R$. The function $(a,\omega)\mapsto\hat I_t(a,\omega)$ is continuous in $a$ and
$\F$-measurable in $\omega$, hence $\B(\R)\otimes\F$-measurable, and
$P[J(a,\cdot)\ne\hat I_t(a)]=0$ for every $a$. By Tonelli's theorem the set
$\{(a,\omega);\ J(a,\omega)\ne\hat I_t(a,\omega)\}$ is $\nu\otimes P$-null, so for $P$-a.e.\
$\omega$ we have $J(a,\omega)=\hat I_t(a,\omega)$ for $\nu$-a.e.\ $a$, and then
$\int h(a)J(a,\omega)\,\mu(da)=\int h(a)\hat I_t(a,\omega)\,\mu(da)$.

\emph{(b) The proof does not use the continuity of $h$.} It applies verbatim to every Borel
function $h\ge0$ which vanishes off a compact interval and satisfies \eqref{eq:3-7-07:fin}; the
nonnegativity is used, to make $\nu$ a finite (positive) measure. When $\mu$ is finite on compact
sets, the identity then extends by linearity, writing $h=h^+-h^-$, to every bounded Borel $h$ with
compact support. This covers the interchange
\bk{(7.27)} in the proof of \bk{Theorem 7.1}, where $\mu$ is a finite measure carried by
$[-K,K)$ and no function $h$ appears: one takes any continuous $h\ge0$ with compact support and
$h=1$ on $[-K,K]$ (or simply $h=\ind_{[-K,K]}$).

\emph{(c) Condition \eqref{eq:3-7-07:fin} cannot be omitted.} Let $\mu$ be the counting measure on
$\Q$, which is $\sigma$-finite, and $h(a)=(1-\abs{a})^+$. Then
$\int h(a)\ind_{(a,\infty)}(x)\,\mu(da)=\sum_{q\in\Q,\,q<x}(1-\abs{q})^+=\infty$ for every
$x>-1$: all terms are nonnegative, and with $y\defeq x\wedge 0$, the interval
$\bigl(\tfrac{y-1}{2},y\bigr)\subseteq(-1,x)$ contains infinitely many rationals $q$, for each of
which $(1-\abs{q})^+=1+q\ge\tfrac{1+y}{2}>0$.
With $X=M=W$ a standard Brownian motion, the integrand on the right-hand side of
the problem is thus $+\infty$ on the set $\{s;\ W_s>-1\}$, which has positive Lebesgue measure,
and the right-hand side is not defined.}
\end{kex}

\begin{kex}{Exercise}{7.9}{}
We take for $X$ the reflected Brownian motion. Let $\{W_t,\F_t;\ 0\le t<\infty\}$, $(\Omega,\F)$,
$\{P^z\}_{z\in\R}$ be the one-dimensional Brownian family of \bk{Section 3.6}, whose filtration
$\{\F_t\}$ satisfies the usual conditions, let $L_t(a)$ be the Brownian local time
(\bk{Theorem 6.11}), and work under $P\defeq P^0$. We set
\begin{equation*}
X_t\defeq\abs{W_t},\qquad M_t\defeq\int_0^t\sgn(W_s)\,dW_s,\qquad V_t\defeq 2L_t(0);
\qquad 0\le t<\infty,
\end{equation*}
with $\sgn$ as in \bk{(6.14)}.

\emph{Step 1: $X$ is a continuous semimartingale with decomposition \bk{(7.1)}.}
The integrand $\sgn(W_s)$ is progressively measurable and bounded by $1$, so $M$ is a continuous,
square-integrable martingale; in particular $M\in\cM^{c,\loc}$ and $M_0=0$. By
\bk{Problem 6.13}\,(ii), $L(0)$ is a continuous, increasing process in the sense of
\bk{Definition 1.4.4}: it is adapted to $\{\F_t\}$, $L_0(0)=0$, and $t\mapsto L_t(0)$ is
nondecreasing and continuous, $P^0$-a.s. Hence $V$ is a continuous, nondecreasing, adapted
process with $V_0=0$ a.s., and it is trivially the difference of two such processes. Finally, the
Tanaka formula \bk{(6.13)} with $a=z=0$ reads, $P^0$-a.s.,
\begin{equation*}
X_t=\abs{W_t}=\int_0^t\sgn(W_s)\,dW_s+2L_t(0)=X_0+M_t+V_t;\qquad 0\le t<\infty,
\end{equation*}
which is \bk{(7.1)} with $X_0=0$.

\emph{Step 2: computation of $J_t(a)$.}
By \bk{(6.26)} with $a=0$ there is an event $\Omega_1$ with $P^0(\Omega_1)=1$ such that, for every
$\omega\in\Omega_1$, the path $t\mapsto V_t(\omega)$ is continuous and nondecreasing with
$V_0(\omega)=0$, and the Lebesgue--Stieltjes measure $dV_s(\omega)=2\,dL_s(0,\omega)$ on
$[0,\infty)$ assigns measure zero to the set
$\{s\ge 0;\ W_s(\omega)\ne 0\}=\{s\ge 0;\ X_s(\omega)>0\}$. We fix $\omega\in\Omega_1$ and
$t\in[0,\infty)$, and distinguish two cases.

If $a\ge 0$, then $\{s\in[0,t];\ X_s(\omega)>a\}\subseteq\{s\ge 0;\ X_s(\omega)>0\}$ is a
$dV(\omega)$-null set, and therefore
\begin{equation*}
J_t(a,\omega)=\int_0^t\ind_{(a,\infty)}(X_s(\omega))\,dV_s(\omega)=0.
\end{equation*}

If $a<0$, then $X_s(\omega)\ge 0>a$ for every $s$, the integrand
$\ind_{(a,\infty)}(X_s(\omega))$ is identically equal to $1$, and
\begin{equation*}
J_t(a,\omega)=\int_0^t dV_s(\omega)=V_t(\omega)=2L_t(0,\omega).
\end{equation*}
Thus, for every $\omega\in\Omega_1$,
\begin{equation}\label{eq:3-7-09:J}
J_t(a,\omega)=2L_t(0,\omega)\,\ind_{(-\infty,0)}(a);\qquad (t,a)\in[0,\infty)\times\R.
\end{equation}

\emph{Step 3: the discontinuity.}
According to \bk{Problem 6.13}\,(iii), there is an event $\Omega_2$ with $P^0(\Omega_2)=1$ on
which $L_t(0,\omega)>0$ for all $t>0$. For $\omega\in\Omega_1\cap\Omega_2$ and every $t>0$,
\eqref{eq:3-7-09:J} gives
\begin{equation*}
\lim_{b\uparrow 0}J_t(b,\omega)=2L_t(0,\omega)>0=J_t(0,\omega),
\end{equation*}
so that $a\mapsto J_t(a,\omega)$ has a jump of size $2L_t(0,\omega)$ at $a=0$. Consequently, for
$P^0$-a.e.\ $\omega$ the function $a\mapsto J_t(a,\omega)$ is discontinuous at the origin
simultaneously for all $t>0$; in particular, $J$ is not jointly continuous in $(t,a)$.

\rmk{The example is consistent with \bk{Problem 7.6}: the function \eqref{eq:3-7-09:J} is
continuous in $t$, right-continuous in $a$, and its left limit at $a=0$ is
$2L_t(0)=\int_0^t\ind_{[0,\infty)}(X_s)\,dV_s$, as \bk{(7.20)} requires. The jump equals
$\int_0^t\ind_{\{0\}}(X_s)\,dV_s$: a discontinuity of $J_t(\cdot)$ at the level $a$ occurs exactly
when the bounded-variation part $V$ moves while $X$ sits at $a$. For a fixed level $a$ this a.s.\ does not happen when the paths of $V$ are
absolutely continuous with respect to those of $\qv{M}$, because
$\int_0^\infty\ind_{\{a\}}(X_s)\,d\qv{M}_s=0$ a.s.\ by \bk{Exercise 7.10}; in the example,
$\qv{M}_t=t$ and $V=2L(0)$ is singular with respect to Lebesgue measure.}
\end{kex}

\begin{kex}{Exercise}{7.10}{}
Fix $a\in\R$ and let $\Lambda=\{\Lambda_t(b,\omega);\ (t,b)\in[0,\infty)\times\R,\ \omega\in\Omega\}$ be the semimartingale local time for $X$ provided by \bk{Theorem 7.1}. The function $k\defeq\ind_{\{a\}}:\R\to[0,\infty)$ is Borel-measurable, so \bk{Theorem 7.1}\,(iii) applies to it: for every $n\in\N$ there is an event $\Omega_n$ with $P(\Omega_n)=1$ such that the identity \bk{(7.3)} with $t=n$ holds for every $\omega\in\Omega_n$, that is,
\begin{equation}\label{eq:3-7-10:a}
\int_0^{n}\ind_{\{a\}}\bigl(X_s(\omega)\bigr)\,d\qv{M}_s(\omega)
 =2\int_{-\infty}^{\infty}\ind_{\{a\}}(b)\,\Lambda_n(b,\omega)\,db ;\qquad \omega\in\Omega_n .
\end{equation}
Both sides are meaningful. On the left, $s\mapsto X_s(\omega)$ is continuous, hence $s\mapsto\ind_{\{a\}}(X_s(\omega))$ is a nonnegative Borel function (the indicator of a closed subset of $[0,\infty)$), and it is integrated against the measure on $[0,\infty)$ induced by the continuous, nondecreasing function $s\mapsto\qv{M}_s(\omega)$. On the right, $(t,b,\omega)\mapsto\Lambda_t(b,\omega)$ is nonnegative and jointly measurable by \bk{Theorem 7.1}\,(i), so its section $b\mapsto\Lambda_n(b,\omega)$ is a nonnegative Borel function for every $\omega$.

The integrand on the right-hand side of \eqref{eq:3-7-10:a} vanishes off the singleton $\{a\}$, which has Lebesgue measure zero; therefore that integral is zero for \emph{every} $\omega\in\Omega$. Consequently
\begin{equation*}
\int_0^{n}\ind_{\{a\}}\bigl(X_s(\omega)\bigr)\,d\qv{M}_s(\omega)=0
\qquad\text{for all } n\in\N \text{ and all } \omega\in\Omega^*\defeq\bigcap_{n=1}^{\infty}\Omega_n ,
\end{equation*}
and $P(\Omega^*)=1$. Finally, for each $\omega$ the functions $s\mapsto \ind_{[0,n]}(s)\,\ind_{\{a\}}(X_s(\omega))$ increase to $s\mapsto\ind_{\{a\}}(X_s(\omega))$ as $n\to\infty$, so the monotone convergence theorem gives, for every $\omega\in\Omega^*$,
\begin{equation*}
\int_0^{\infty}\ind_{\{a\}}\bigl(X_s(\omega)\bigr)\,d\qv{M}_s(\omega)
 =\lim_{n\to\infty}\int_0^{n}\ind_{\{a\}}\bigl(X_s(\omega)\bigr)\,d\qv{M}_s(\omega)=0 .
\end{equation*}
This is the assertion.

\rmk{There is no circularity: the proof of \bk{Theorem 7.1}\,(iii) in the book (via \bk{Problem 7.7} and the construction of $\Lambda$) precedes this exercise and does not use it. The same argument shows more generally that $\int_0^\infty\ind_N(X_s)\,d\qv{M}_s=0$ a.s.\ for every Borel set $N\subset\R$ of Lebesgue measure zero.}
\end{kex}

\begin{kex}{Exercise}{7.11}{}
Let $X=\{X_t,\F_t;\ 0\le t<\infty\}$ be a continuous, adapted process whose paths are of bounded variation on every interval $[0,t]$, on a filtered probability space satisfying the usual conditions, as in \bk{(7.1)}. We show in Step 1 that $X$ is a continuous semimartingale whose local martingale part vanishes in every decomposition \bk{(7.1)}. Then we let $\Lambda$ be any semimartingale local time for $X$ in the sense of \bk{Theorem 7.1}, and show that there is an event $\Omega^*$, depending on $\Lambda$, with $P(\Omega^*)=1$ and
\begin{equation}\label{eq:3-7-11:a}
\Lambda_t(a,\omega)=0\qquad\text{for all }(t,a)\in[0,\infty)\times\R\text{ and all }\omega\in\Omega^*.
\end{equation}

\emph{Step 1: $X$ is a continuous semimartingale with $M\equiv0$.} Fix $\omega$ and $t>0$, and let $\check V_t$ be the total variation of $s\mapsto X_s(\omega)$ on $[0,t]$, i.e.\ the supremum of the sums $\sum_{j=1}^m\abs{X_{t_j}-X_{t_{j-1}}}$ over all partitions $0=t_0<t_1<\dots<t_m=t$; put $\check V_0\defeq0$. It is finite and nondecreasing in $t$. We claim that
\begin{equation}\label{eq:3-7-11:b}
\check V_t=\lim_{n\to\infty}S_n,\qquad S_n\defeq\sum_{k=1}^{2^n}\abs[\big]{X_{kt2^{-n}}-X_{(k-1)t2^{-n}}}.
\end{equation}
Indeed, by the triangle inequality the sum over a partition does not decrease when points are added to the partition; hence $S_n\le S_{n+1}\le\check V_t$, and $\lim_n S_n$ exists and is at most $\check V_t$. Conversely, let $0=t_0<\dots<t_m=t$ be a partition and $\eps>0$. For each $n$ let $s^{(n)}_0\defeq0$, $s^{(n)}_m\defeq t$ and, for $0<j<m$, let $s^{(n)}_j$ be the largest point of the form $kt2^{-n}$ not exceeding $t_j$. Then $0=s^{(n)}_0\le s^{(n)}_1\le\dots\le s^{(n)}_m=t$ and $s^{(n)}_j\to t_j$ as $n\to\infty$, so by the continuity of the path we can choose $n$ with $\abs{X_{s^{(n)}_j}-X_{t_j}}<\eps/(2m)$ for all $j$. Since the points $s^{(n)}_j$ belong to the $n$-th dyadic partition of $[0,t]$, the triangle inequality gives
\[
\sum_{j=1}^m\abs{X_{t_j}-X_{t_{j-1}}}\le\sum_{j=1}^m\abs[\big]{X_{s^{(n)}_j}-X_{s^{(n)}_{j-1}}}+\eps\le S_n+\eps\le\lim_{n\to\infty}S_n+\eps .
\]
Taking the supremum over partitions and letting $\eps\downarrow0$ we obtain \eqref{eq:3-7-11:b}. Each $S_n$ is $\F_t$-measurable, and therefore so is $\check V_t$.

The function $t\mapsto\check V_t$ is continuous. To see this, write $\check V[s,u]$ for the total variation of the path on $[s,u]$, so that $\check V_u-\check V_s=\check V[s,u]$ for $s<u$. Fix $s<t$ and $\eps>0$, and choose a partition $s=t_0<t_1<\dots<t_m=t$ whose sum exceeds $\check V[s,t]-\eps$; by the continuity of the path there is $u\in(s,t_1)$ with $\abs{X_{u'}-X_s}<\eps$ for all $u'\in(s,u]$. For such $u'$ the points $u'<t_1<\dots<t_m$ form a partition of $[u',t]$, whence
\[
\check V[u',t]\ge\abs{X_{t_1}-X_{u'}}+\sum_{j=2}^m\abs{X_{t_j}-X_{t_{j-1}}}\ge\sum_{j=1}^m\abs{X_{t_j}-X_{t_{j-1}}}-\eps>\check V[s,t]-2\eps ,
\]
and so $\check V_{u'}-\check V_s=\check V[s,t]-\check V[u',t]<2\eps$: this is right-continuity at $s$. Left-continuity at $t>0$ follows in the same way, with the roles of the endpoints interchanged.

Consequently $V^{\pm}_t\defeq\frac12\bigl(\check V_t\pm(X_t-X_0)\bigr)$ are continuous, adapted processes with $V^\pm_0=0$; they are nondecreasing because $\abs{X_u-X_s}\le\check V_u-\check V_s$ for $s<u$. Thus
\[
X_t=X_0+M_t+V_t,\qquad M\equiv0,\quad V\defeq V^+-V^-=X-X_0,
\]
is a decomposition of the form \bk{(7.1)}, with $M\equiv 0\in\cM^{c,\loc}$ and $\qv{M}\equiv0$.

The semimartingale local time of \bk{Theorem 7.1} refers to a decomposition \bk{(7.1)}, through $\qv{M}$ in \bk{(7.3)}; so we must make sure that no other decomposition has a nontrivial martingale part. Let $X=X_0+M'+V'$ be any decomposition of the form \bk{(7.1)}. The decomposition found above is of the form \bk{(3.1)}, \bk{(3.2)} with $M\equiv0$ and $B=X-X_0$ (the minimality of \bk{(3.2)} is merely a convention for writing $B$ as a difference), and $V'$ is continuous, adapted and of finite total variation on bounded intervals. By \bk{Problem 3.2}, $M'_t=0$ and $V'_t=X_t-X_0$ for all $0\le t<\infty$, a.s. Then $(M')^2$ is indistinguishable from zero, so the zero process is a continuous, nondecreasing, adapted process $A$ with $A_0=0$ for which $(M')^2-A$ is a continuous local martingale, and by the uniqueness of the quadratic variation (\bk{Problem 1.5.17}) we conclude that, almost surely, $\qv{M'}_t=0$ for all $0\le t<\infty$. From now on $M$ denotes the local martingale part of whichever decomposition \bk{(7.1)} the local time $\Lambda$ refers to; we have just shown that there is an event $\Omega_0$ with $P(\Omega_0)=1$ on which $\qv{M}_t=0$ for all $t\ge0$.

\emph{Step 2: $\Lambda_t(\cdot)$ vanishes Lebesgue-almost everywhere.} Let $\Lambda$ be a semimartingale local time for $X$. We apply property (iii) of \bk{Theorem 7.1} with $k\equiv1$: \bk{(7.3)} holds for all $0\le t<\infty$ outside a $P$-null event, and intersecting the complement of that event with $\Omega_0$ we obtain an event $\Omega_1$ with $P(\Omega_1)=1$ such that
\[
\int_{-\infty}^{\infty}\Lambda_t(a,\omega)\,da=\frac12\int_0^t d\qv{M}_s(\omega)=\frac12\qv{M}_t(\omega)=0\qquad\text{for all }0\le t<\infty,\ \omega\in\Omega_1 .
\]
By property (i), $a\mapsto\Lambda_t(a,\omega)$ is Borel-measurable for every fixed $(t,\omega)$, and it is nonnegative. A nonnegative measurable function with zero integral vanishes almost everywhere; therefore, for every $\omega\in\Omega_1$ and every $t\ge0$, the set
\[
Z_t(\omega)\defeq\{a\in\R:\ \Lambda_t(a,\omega)\ne0\}
\]
has Lebesgue measure zero.

\emph{Step 3: $\Lambda_t(\cdot)$ vanishes everywhere.} By property (iv) of \bk{Theorem 7.1}, there is an event $\Omega_2$ with $P(\Omega_2)=1$ such that for $\omega\in\Omega_2$ and all $(t,a)$ we have $\lim_{\tau\to t,\,b\downarrow a}\Lambda_\tau(b,\omega)=\Lambda_t(a,\omega)$; in particular (take $\tau=t$) the function $b\mapsto\Lambda_t(b,\omega)$ is right-continuous on $\R$. Set $\Omega^*\defeq\Omega_1\cap\Omega_2$, and fix $\omega\in\Omega^*$, $t\ge0$ and $a\in\R$. For every $n\ge1$ the interval $(a,a+\frac1n)$ has positive Lebesgue measure, so it is not contained in the null set $Z_t(\omega)$; choose $b_n\in(a,a+\frac1n)\setminus Z_t(\omega)$. Then $b_n\downarrow a$ and $\Lambda_t(b_n,\omega)=0$ for all $n$, whence by right-continuity
\[
\Lambda_t(a,\omega)=\lim_{n\to\infty}\Lambda_t(b_n,\omega)=0 .
\]
This proves \eqref{eq:3-7-11:a}: every semimartingale local time $\Lambda$ of $X$ is identically zero on $[0,\infty)\times\R$ outside a $P$-null event (which depends on $\Lambda$).

\rmk{The result is consistent with the Tanaka--Meyer formula \bk{(7.9)}, which for $M\equiv0$ and $\Lambda\equiv0$ reduces to $\abs{X_t-a}=\abs{X_0-a}+\int_0^t\sgn(X_s-a)\,dX_s$, the pathwise change-of-variable formula for the function $x\mapsto\abs{x-a}$ along a continuous path of bounded variation. Conversely, for a continuous local martingale $M$ which is not constant, \bk{(7.3)} with $k\equiv 1$ gives $\int_{-\infty}^\infty\Lambda_t(a)\,da=\frac12\qv{M}_t$, which is positive for large $t$ with positive probability; so the vanishing of the local time is a feature of the absence of a martingale part.}
\end{kex}

\begin{kex}{Exercise}{7.12}{}
Let $\Lambda=\{\Lambda_t(a,\omega);\ (t,a)\in[0,\infty)\times\R,\ \omega\in\Omega\}$ be a semimartingale local time for $X$, i.e., a nonnegative random field with the properties (i)--(v) of \bk{Theorem 7.1}; we shall use only the joint measurability in (i) and the properties (ii), (iii) and (iv). We have to show that, for $P$-a.e.\ $\omega$, $\Lambda_t(0,\omega)=0$ for all $t\in[0,\infty)$.

\emph{Step 1: the occupation-time identity.} Define $\kappa\colon\R\to[0,\infty)$ by
\begin{equation*}
\kappa(a)\defeq\frac{1}{k(a)}\ \text{ for } a>0,\qquad \kappa(a)\defeq 0\ \text{ for } a\le 0 .
\end{equation*}
Since $k$ is Borel-measurable with values in $(0,\infty)$, the function $\kappa$ is Borel-measurable, nonnegative and real-valued, and $\kappa(X_s)=\ind_{\{X_s>0\}}/k(X_s)$, where the right-hand side is understood, as in the statement, to be zero on $\{X_s\le 0\}$. By \bk{Theorem 7.1}\,(iii), applied to $\kappa$, there is an event $\Omega_1$ with $P(\Omega_1)=1$ such that for every $\omega\in\Omega_1$ and every $n\in\N$,
\begin{equation}\label{eq:3-7-12:a}
\int_0^n \frac{\ind_{\{X_s(\omega)>0\}}}{k(X_s(\omega))}\,d\qv{M}_s(\omega)
 = 2\int_{(0,\infty)}\frac{\Lambda_n(a,\omega)}{k(a)}\,da ,
\end{equation}
both sides being integrals of nonnegative Borel functions with values in $[0,\infty]$ (the section $a\mapsto\Lambda_n(a,\omega)$ is Borel-measurable, being a section of the jointly measurable mapping $(t,a,\omega)\mapsto\Lambda_t(a,\omega)$ of \bk{Theorem 7.1}\,(i)). By hypothesis, for each $n\in\N$ the left-hand side of \eqref{eq:3-7-12:a} is finite outside a $P$-null event; discarding these countably many null events, we obtain an event $\Omega_2\subseteq\Omega_1$ with $P(\Omega_2)=1$ such that
\begin{equation}\label{eq:3-7-12:b}
\int_{(0,\infty)}\frac{\Lambda_n(a,\omega)}{k(a)}\,da<\infty\qquad\text{for all } n\in\N,\ \omega\in\Omega_2 .
\end{equation}

\emph{Step 2: right-continuity at the origin.} By \bk{Theorem 7.1}\,(iv), there is an event $\Omega_3$ with $P(\Omega_3)=1$ such that for every $\omega\in\Omega_3$ and every $(t,a)\in[0,\infty)\times\R$ we have $\lim_{\tau\to t,\,b\downarrow a}\Lambda_\tau(b,\omega)=\Lambda_t(a,\omega)$; taking $\tau=t=n$ and $a=0$, we obtain in particular
\begin{equation}\label{eq:3-7-12:c}
\lim_{b\downarrow 0}\Lambda_n(b,\omega)=\Lambda_n(0,\omega)\qquad\text{for all } n\in\N,\ \omega\in\Omega_3 .
\end{equation}

\emph{Step 3: conclusion.} Let $\omega\in\Omega_2\cap\Omega_3$ and $n\in\N$, and suppose that $\Lambda_n(0,\omega)>0$. By \eqref{eq:3-7-12:c} there exists $\eps=\eps(n,\omega)>0$ such that $\Lambda_n(b,\omega)\ge\frac12\Lambda_n(0,\omega)$ for all $b\in(0,\eps)$. Since the integrand in \eqref{eq:3-7-12:b} is nonnegative, it follows that
\begin{equation*}
\int_{(0,\infty)}\frac{\Lambda_n(a,\omega)}{k(a)}\,da
 \ \ge\ \int_{(0,\eps)}\frac{\Lambda_n(a,\omega)}{k(a)}\,da
 \ \ge\ \frac{\Lambda_n(0,\omega)}{2}\int_{(0,\eps)}\frac{du}{k(u)}=\infty ,
\end{equation*}
by the assumption on $k$; this contradicts \eqref{eq:3-7-12:b}. As $\Lambda$ is nonnegative, we conclude that
\begin{equation*}
\Lambda_n(0,\omega)=0\qquad\text{for all } n\in\N,\ \omega\in\Omega_2\cap\Omega_3 .
\end{equation*}
Finally, by \bk{Theorem 7.1}\,(ii) the mapping $t\mapsto\Lambda_t(0,\omega)$ is nondecreasing with $\Lambda_0(0,\omega)=0$, for every $\omega$. Hence, for $\omega\in\Omega_2\cap\Omega_3$ and $t\in[0,\infty)$, choosing an integer $n\ge t$ we get $0\le\Lambda_t(0,\omega)\le\Lambda_n(0,\omega)=0$. Since $P(\Omega_2\cap\Omega_3)=1$, this shows that
\begin{equation*}
P\bigl[\Lambda_t(0)=0,\ \forall\, 0\le t<\infty\bigr]=1 .
\end{equation*}

\rmk{The argument uses only the right-hand limit in \bk{Theorem 7.1}\,(iv), which is why the hypothesis involves only the set $\{X_s>0\}$ and the conclusion concerns $\Lambda(0)$ and not the left limit $\Lambda(0-)$. If the hypotheses are imposed instead on $\{X_s<0\}$, i.e.\ if $\int_0^t \ind_{\{X_s<0\}}\,d\qv{M}_s/k(-X_s)<\infty$ a.s.\ for every $t$, the same proof (with $\kappa(a)=\ind_{(-\infty,0)}(a)/k(-a)$ and the left limits of \bk{Theorem 7.1}\,(iv)) gives $\Lambda_t(0-)=0$ for all $t$, almost surely.}
\end{kex}

\begin{kex}{Exercise}{7.13}{}
\emph{Setting and conventions.} A continuous local martingale is understood in the sense of
\bk{Definition 1.5.15}, with respect to a filtration $\{\F_t\}$ which satisfies the usual conditions
(as required in \bk{(7.1)}); we do not assume $M_0=0$. If $\{T_n\}$ is a sequence of stopping times
as in that definition, then $M_0=M_{0\wedge T_1}$ is integrable and $\F_0$-measurable, and
$N_{t\wedge T_n}=M_{t\wedge T_n}-M_0$ is a martingale for every $n$; thus
\begin{equation*}
M_t=M_0+N_t,\qquad N\in\cM^{c,\loc},
\end{equation*}
is a continuous semimartingale of the form \bk{(7.1)} with $X_0=M_0$ and $V\equiv0$. We set
$\qv{M}\defeq\qv{N}$ and $\int_0^tH_s\,dM_s\defeq\int_0^tH_s\,dN_s$; this is the reading under
which \bk{(7.28)} of \bk{Remark 7.8}, which allows $M_0\neq0$, is the formula \bk{(7.9)} of
\bk{Theorem 7.1} for the semimartingale $X=M_0+N$, and $\Lambda=\Lambda^M$ is the local time of
this semimartingale. We write $f_t,f_x,f_y,f_{xx}$ for the partial derivatives of $f$, and
$A_t\defeq\qv{M}_t$.

We shall use the following elementary fact: \emph{if $\xi$ is an integrable, $\F_0$-measurable random
variable and $K\in\cM^{c,\loc}$, then $\xi+K$ is a continuous local martingale.} Indeed, if
$\{T_n\}$ is as in \bk{Definition 1.5.15} for $K$, then $\xi+K_{t\wedge T_n}$ is a martingale for
every $n$. (Integrability of $\xi$ is also necessary, since the value at $t=0$ of a martingale is
integrable.)

We shall prove that
\begin{align}
f(\qv{M}_t,\abs{M_t},L_t)&=f(0,\abs{M_0},0)+\int_0^t f_x(\qv{M}_s,\abs{M_s},L_s)\,d\beta_s,
\label{eq:3-7-13:a}\\
f(\qv{M}_t,S_t-M_t,S_t)&=f(0,0,M_0)-\int_0^t f_x(\qv{M}_s,S_s-M_s,S_s)\,dM_s
\label{eq:3-7-13:b}
\end{align}
hold for all $0\le t<\infty$, a.s.\ $P$, where $\beta\in\cM^{c,\loc}$ is the process defined in
\eqref{eq:3-7-13:beta} below and $\qv{\beta}=\qv{M}$. In both formulas the integrand $H$ is a
continuous, adapted process, hence progressively measurable with
$\int_0^t H_s^2\,d\qv{M}_s\le\max_{0\le s\le t}H_s^2\cdot\qv{M}_t<\infty$ for every $t$; so
$H\in\cP^*$ for the integrator in question and the stochastic integral belongs to $\cM^{c,\loc}$.
By the fact just stated, \emph{the two processes of the exercise are continuous local martingales
as soon as their initial values $f(0,\abs{M_0},0)$ and $f(0,0,M_0)$ are integrable}; this is the
case, in particular, when $M_0$ is a.s.\ constant, e.g.\ when $M\in\cM^{c,\loc}$ as in
\bk{(7.1)}, where both initial values equal $f(0,0,0)$. In general, the two processes minus their
initial values belong to $\cM^{c,\loc}$ (see the first remark at the end).

\emph{Step 1: proof of \eqref{eq:3-7-13:a}.} By \bk{(7.28)} with $a=0$ we have, a.s.,
\begin{equation}\label{eq:3-7-13:beta}
\abs{M_t}=\abs{M_0}+\beta_t+L_t,\qquad \beta_t\defeq\int_0^t\sgn(M_s)\,dM_s;\qquad 0\le t<\infty.
\end{equation}
The integrand $\sgn(M_s)$ is bounded and progressively measurable (a Borel function of a continuous,
adapted process), so $\beta\in\cM^{c,\loc}$, and
since $\sgn^2\equiv1$ under the convention \bk{(6.14)}, property \bk{(2.19)} (valid for integrators
in $\cM^{c,\loc}$, as noted before \bk{Proposition 2.24}) gives
\begin{equation*}
\qv{\beta}_t=\int_0^t\sgn^2(M_s)\,d\qv{M}_s=\qv{M}_t=A_t .
\end{equation*}
By \bk{Theorem 7.1}~(i),~(ii), the process $L=2\Lambda(0)$ is adapted, continuous and nondecreasing
with $L_0=0$; the same is true of $A$. We apply It\^o's rule, \bk{Theorem 3.6}, with $d=3$ to the
vector semimartingale $\Theta_t\defeq(A_t,\abs{M_t},L_t)$, which has the $\F_0$-measurable initial
value $\Theta_0=(0,\abs{M_0},0)$, local martingale parts $(0,\beta,0)$ and bounded variation parts
$(A,L,L)$, and to the function $(s;t,x,y)\mapsto f(t,x,y)$,
which does not depend on the time variable $s$ and is of class $C^{1,2}$ because
$f\in C^2(\R^3)$. The only nonvanishing cross-variation is $\qv{\beta}=A$, so we obtain, a.s.,
\begin{align*}
f(\Theta_t)&=f(\Theta_0)+\int_0^t f_t(\Theta_s)\,dA_s+\int_0^t f_x(\Theta_s)\,dL_s
+\int_0^t f_y(\Theta_s)\,dL_s\\
&\qquad+\int_0^t f_x(\Theta_s)\,d\beta_s+\frac12\int_0^t f_{xx}(\Theta_s)\,dA_s\\
&=f(\Theta_0)+\int_0^t f_x(\Theta_s)\,d\beta_s
+\int_0^t\Bigl(f_t+\tfrac12f_{xx}\Bigr)(\Theta_s)\,dA_s+\int_0^t(f_x+f_y)(\Theta_s)\,dL_s
\end{align*}
for all $0\le t<\infty$. The integral with respect to $dA_s$ vanishes because $f_t+\frac12f_{xx}=0$
on $\R^3$. As for the last integral: by \bk{(7.2)} with $a=0$ there is an event of probability one
on which $\int_0^\infty\ind_{\{M_s\neq0\}}\,dL_s=0$, whereas on the set $\{s;\ M_s=0\}$ the
integrand is $(f_x+f_y)(A_s,0,L_s)=0$ by the second hypothesis on $f$. Hence, on that event,
\begin{equation*}
\abs[\Big]{\int_0^t(f_x+f_y)(\Theta_s)\,dL_s}
\le\int_0^t\abs{(f_x+f_y)(\Theta_s)}\,\ind_{\{M_s\neq0\}}\,dL_s=0;\qquad 0\le t<\infty,
\end{equation*}
and \eqref{eq:3-7-13:a} is proved.

\emph{Step 2: the measure $dS$ is carried by $\{s;\ M_s=S_s\}$.} The process $S$ is continuous and
nondecreasing with $S_0=M_0$, and it is
adapted, since $S_t=\sup\{M_s;\ s\in([0,t]\cap\Q)\cup\{t\}\}$. Fix $\omega$ and consider the set
$G\defeq\{s\ge0;\ S_s>M_s\}$, which is open in $[0,\infty)$ by continuity. Let $s\in G$. There is
$\eps>0$ such that $M_u<S_s$ for all $u\in I\defeq(s-\eps,s+\eps)\cap[0,\infty)$. For
$u\in I$, $u\ge s$, we have $S_u=\max\bigl(S_s,\max_{s\le r\le u}M_r\bigr)=S_s$; for $u\in I$,
$u<s$, we have $S_s=\max\bigl(S_u,\max_{u\le r\le s}M_r\bigr)$ with
$\max_{u\le r\le s}M_r<S_s$, whence again $S_u=S_s$. Thus $S$ is locally constant on $G$, hence
constant on each of the countably many connected components (intervals) of $G$, and the
Lebesgue--Stieltjes measure $dS$ assigns zero mass to each of them. Therefore, for every $\omega$,
\begin{equation}\label{eq:3-7-13:supp}
\int_0^\infty\ind_{\{S_s\neq M_s\}}\,dS_s=0 .
\end{equation}

\emph{Step 3: proof of \eqref{eq:3-7-13:b}.} We write $S_t=M_0+(S_t-M_0)$ and
$S_t-M_t=-N_t+(S_t-M_0)$, where $S-M_0$ is continuous, adapted and nondecreasing, with initial
value zero. We apply \bk{Theorem 3.6} as in Step~1, now to the
vector semimartingale $\Xi_t\defeq(A_t,S_t-M_t,S_t)$, which has the $\F_0$-measurable initial value
$\Xi_0=(0,0,M_0)$, local martingale
parts $(0,-N,0)$ and bounded variation parts $(A,S-M_0,S-M_0)$; the only nonvanishing
cross-variation is $\qv{-N}=A$. We obtain, a.s.,
\begin{equation*}
f(\Xi_t)=f(\Xi_0)-\int_0^t f_x(\Xi_s)\,dM_s
+\int_0^t\Bigl(f_t+\tfrac12f_{xx}\Bigr)(\Xi_s)\,dA_s+\int_0^t(f_x+f_y)(\Xi_s)\,dS_s
\end{equation*}
for all $0\le t<\infty$. The $dA_s$-integral vanishes as before. In the last integral the integrand
equals $(f_x+f_y)(A_s,0,S_s)=0$ on $\{s;\ S_s=M_s\}$, and the complement of this set is $dS$-null
by \eqref{eq:3-7-13:supp}; so this integral vanishes as well, and \eqref{eq:3-7-13:b} follows.

\textbf{(i)} The function $f(t,x,y)\defeq x^2-t$ is of class $C^2(\R^3)$, with
$f_t+\frac12f_{xx}=-1+1=0$ on $\R^3$ and $f_x(t,0,y)+f_y(t,0,y)=0+0=0$. Here $f(0,0,M_0)=0$ and
$f_x=2x$, so \eqref{eq:3-7-13:b} reads
\begin{equation*}
(S_t-M_t)^2-\qv{M}_t=-2\int_0^t(S_s-M_s)\,dM_s ;\qquad 0\le t<\infty,
\end{equation*}
and the process $(S_t-M_t)^2-\qv{M}_t$ belongs to $\cM^{c,\loc}$, whatever $M_0$ may be.

\textbf{(ii)} Let $g\in C^1(\R)$. We shall show that, a.s., for all $0\le t<\infty$,
\begin{align}
g(L_t)-\abs{M_t}\,g'(L_t)&=g(0)-\abs{M_0}\,g'(0)-\int_0^t g'(L_s)\,d\beta_s,
\label{eq:3-7-13:c}\\
g(S_t)-(S_t-M_t)\,g'(S_t)&=g(M_0)+\int_0^t g'(S_s)\,dM_s .\label{eq:3-7-13:d}
\end{align}
The integrands are continuous and adapted, so the stochastic integrals belong to $\cM^{c,\loc}$.
The initial value $g(0)-\abs{M_0}\,g'(0)$ is integrable because $M_0$ is; hence the first process
in (ii) is always a continuous local martingale, and the second one is a continuous local martingale
provided $g(M_0)$ is integrable, in particular if $M_0$ is a.s.\ constant.

\emph{The case $g\in C^3(\R)$.} Then $f(t,x,y)\defeq g(y)-x\,g'(y)$ is of class $C^2(\R^3)$, with
$f_t=0$, $f_x=-g'(y)$, $f_{xx}=0$ and $f_y=g'(y)-x\,g''(y)$. Hence $f_t+\frac12f_{xx}=0$ on $\R^3$
and $f_x(t,0,y)+f_y(t,0,y)=-g'(y)+g'(y)=0$; since $f(0,\abs{M_0},0)=g(0)-\abs{M_0}\,g'(0)$,
$f(0,0,M_0)=g(M_0)$ and $f_x=-g'(y)$, the identities
\eqref{eq:3-7-13:a}, \eqref{eq:3-7-13:b} become exactly \eqref{eq:3-7-13:c}, \eqref{eq:3-7-13:d}.

\emph{The general case.} If $g$ is merely of class $C^1$, the function $g(y)-x\,g'(y)$ need not be
of class $C^2$ and the first part of the exercise does not apply directly; we argue by approximation.
By the Weierstrass approximation theorem there is, for every $n\ge1$, a polynomial $p_n$ with
$\sup_{\abs{y}\le n}\abs{p_n(y)-g'(y)}\le1/n$. The polynomial
$g_n(y)\defeq g(0)+\int_0^yp_n(u)\,du$ satisfies $g_n(0)=g(0)$, $g_n'=p_n$ and
\begin{equation*}
\abs{g_n(y)-g(y)}=\abs[\Big]{\int_0^y\bigl(p_n(u)-g'(u)\bigr)\,du}\le\frac{\abs{y}}{n}\le1;
\qquad\abs{y}\le n.
\end{equation*}
Consequently $g_n\to g$ and $g_n'\to g'$ uniformly on every compact subset of $\R$. Fix
$0\le T<\infty$. For every $\omega$, the path $s\mapsto L_s$ maps $[0,T]$ into the compact interval
$[0,L_T]$, so
\begin{equation*}
\int_0^T\abs{g_n'(L_s)-g'(L_s)}^2\,d\qv{\beta}_s
\le\sup_{0\le y\le L_T}\abs{g_n'(y)-g'(y)}^2\cdot\qv{M}_T\xrightarrow[n\to\infty]{}0,
\end{equation*}
and the same holds with $(S,M)$ in place of $(L,\beta)$ and the interval $[M_0,S_T]$ in place of
$[0,L_T]$. By \bk{Proposition 2.26} (with the constant stopping time $T$),
\begin{equation*}
\sup_{0\le t\le T}\abs[\Big]{\int_0^tg_n'(L_s)\,d\beta_s-\int_0^tg'(L_s)\,d\beta_s}
\xrightarrow[n\to\infty]{}0,\qquad
\sup_{0\le t\le T}\abs[\Big]{\int_0^tg_n'(S_s)\,dM_s-\int_0^tg'(S_s)\,dM_s}
\xrightarrow[n\to\infty]{}0
\end{equation*}
in probability, and therefore almost surely along a common subsequence $\{n_k\}$. On the other
hand, for every $\omega$ and every $t$ we have
$g_n(L_t)-\abs{M_t}g_n'(L_t)\to g(L_t)-\abs{M_t}g'(L_t)$,
$g_n(S_t)-(S_t-M_t)g_n'(S_t)\to g(S_t)-(S_t-M_t)g'(S_t)$, $g_n'(0)\to g'(0)$ and
$g_n(M_0)\to g(M_0)$. Each $g_n$ is of class $C^3$, so
\eqref{eq:3-7-13:c} and \eqref{eq:3-7-13:d} hold for $g_n$, for all $n$ simultaneously, outside
a null event; letting $n=n_k\to\infty$ we obtain \eqref{eq:3-7-13:c} and \eqref{eq:3-7-13:d} for $g$
on $[0,T]$, almost surely. Taking $T=1,2,\dots$ completes the proof.

\rmk{If $M_0$ is not a.s.\ constant, the initial values $f(0,\abs{M_0},0)$, $f(0,0,M_0)$ and
$g(M_0)$ need not be integrable. Indeed, let $W$ be a standard Brownian motion, $\xi$ an integrable
random variable independent of $W$, and $\{\F_t\}$ the augmentation of the filtration generated by
$\xi$ and $W$; then $W$ is a martingale of this filtration, which satisfies the usual conditions, and
$M\defeq\xi+W$ is a continuous martingale with $M_0=\xi$. For the function $f(t,x,y)=x^2-t$ of
part~(i) we have $f(0,\abs{M_0},0)=\xi^2$, which is not integrable if $E\,\xi^2=\infty$; for
$g(y)=e^{y^2}$ and $f(t,x,y)=g(y)-x\,g'(y)$, which satisfies both hypotheses by part~(ii), we have
$f(0,0,M_0)=g(M_0)=e^{\xi^2}$, which is not integrable if $E\,e^{\xi^2}=\infty$. The corresponding processes are then not local martingales in the literal
sense of \bk{Definition 1.5.15}, because a martingale has an integrable initial value. What
\eqref{eq:3-7-13:a}--\eqref{eq:3-7-13:d} show in complete generality is that each of these processes
is the sum of its $\F_0$-measurable initial value and an element of $\cM^{c,\loc}$. Note also that
only the values of $f$ on $[0,\infty)\times[0,\infty)\times\R$ enter, because $\qv{M}$, $\abs{M}$
and $S-M$ are nonnegative.}

\rmk{As a check, take $M_0=0$. With $g(y)=y$, \eqref{eq:3-7-13:c} is the
Tanaka formula \eqref{eq:3-7-13:beta} and \eqref{eq:3-7-13:d} reduces to $M_t=M_t$; with
$g(y)=y^2/2$, \eqref{eq:3-7-13:d} states that $M_tS_t-\frac12S_t^2=\int_0^tS_s\,dM_s$, which agrees
with integration by parts, $d(M_tS_t)=S_t\,dM_t+M_t\,dS_t$ and $M_t\,dS_t=S_t\,dS_t$ by
\eqref{eq:3-7-13:supp}.}
\end{kex}

\begin{kex}{Exercise}{7.14}{}
\emph{Conventions.} As in \bk{(7.1)}, the filtration $\{\F_t\}$ satisfies the usual conditions,
$X_t=M_t+V_t$ with $M\in\cM^{c,\loc}$ and $V=A^+-A^-$, where $A^\pm$ are continuous, nondecreasing,
adapted processes with $A^\pm_0=0$; we take this to be the minimal decomposition of $V$, as in
\bk{(3.2)}, so that $\check V\defeq A^++A^-$ is the total variation of $V$. Identities between
processes, such as the one in (iii), are understood to hold for all $0\le t<\infty$, almost surely,
as in \bk{Definition 3.1}. Following \bk{Lemma 6.14}~(iii), condition (i) means that the paths of $V$
do not vary on the set $\{t\ge0;\ X_t\neq0\}=\{t\ge0;\ X_t>0\}$, i.e.,
\begin{equation}\label{eq:3-7-14:flat}
\int_0^\infty\ind_{\{X_s\neq0\}}\,d\check V_s=0,\quad\text{a.s. } P.
\end{equation}
The integral in (ii) is, by definition, $J_t\defeq I_t+B_t$, where
\begin{equation*}
I_t\defeq\int_0^t\ind_{\{X_s\neq0\}}\,dM_s,\qquad
B_t\defeq\int_0^t\ind_{\{X_s\neq0\}}\,dV_s
=\int_0^t\ind_{\{X_s\neq0\}}\,dA^+_s-\int_0^t\ind_{\{X_s\neq0\}}\,dA^-_s .
\end{equation*}
The process $\ind_{\{X_s\neq0\}}$ is bounded and progressively measurable, because $X$ is continuous
and adapted. Hence $I\in\cM^{c,\loc}$, and each of the two Lebesgue--Stieltjes integrals which make up
$B$ is a continuous, nondecreasing, adapted process starting at zero; thus $J$ is a continuous
semimartingale with decomposition \bk{(3.1)}. We observe that
\begin{equation}\label{eq:3-7-14:equiv}
\text{\eqref{eq:3-7-14:flat} holds}\quad\Longleftrightarrow\quad
B_t=0\ \text{ for all } 0\le t<\infty,\ \text{a.s.}
\end{equation}
Indeed, for fixed $\omega$ and $T<\infty$, the signed measure
$\nu(ds)=\ind_{\{X_s\neq0\}}\,dV_s$ on $[0,T]$ has total variation measure
$\ind_{\{X_s\neq0\}}\,d\check V_s$. If the latter vanishes, then $B_t=\nu([0,t])=0$ for $t\le T$;
conversely, if $\nu([0,t])=0$ for all $0\le t\le T$, then $\nu$ vanishes on the $\pi$-system of
intervals $[0,t]$, which generates the Borel sets of $[0,T]$, so $\nu=0$ and its total variation
is zero; we then let $T\to\infty$ through the integers.

We shall use repeatedly the uniqueness of the decomposition of a continuous semimartingale,
\bk{Problem 3.2}: if $M',M''\in\cM^{c,\loc}$ and $B',B''$ are continuous, adapted processes of finite
total variation on bounded intervals, starting at zero, each being the difference of two continuous,
nondecreasing, adapted processes, and if $M'_t+B'_t=M''_t+B''_t$ for all $t$, a.s., then
$M'=M''$ and $B'=B''$ for all $t$, a.s.

\textbf{(i) $\Rightarrow$ (ii).} By \eqref{eq:3-7-14:equiv}, $B\equiv0$ a.s., so $J$ is
indistinguishable from $I\in\cM^{c,\loc}$. Since $\F_0$ contains the $P$-null events, $J$ is adapted,
and it is itself a continuous local martingale (reduced by the same stopping times as $I$).

\textbf{(ii) $\Rightarrow$ (i).} If $J\in\cM^{c,\loc}$, then $J$ admits the two decompositions
$J=I+B$ and $J=J+0$ of the type just described. By \bk{Problem 3.2}, $B_t=0$ for all $t$, a.s., and
\eqref{eq:3-7-14:flat} follows from \eqref{eq:3-7-14:equiv}.

\textbf{(i) $\Rightarrow$ (iii).} \emph{Step 1: $V$ is nondecreasing.} We apply the Tanaka--Meyer
formula \bk{(7.7)} with $a=0$. Because $X\ge0$ and $X_0=0$, it reads
\begin{equation*}
X_t=X_t^+=\int_0^t\ind_{(0,\infty)}(X_s)\,dM_s+\int_0^t\ind_{(0,\infty)}(X_s)\,dV_s+C_t^+(0);
\qquad 0\le t<\infty,\ \text{a.s.},
\end{equation*}
where $C^+(0)$ is an adapted, continuous, nondecreasing process (see the text following
\bk{(7.9)}); evaluating at $t=0$ shows $C_0^+(0)=0$ a.s. The first integral is $I_t$ and the second
is $B_t$, which vanishes identically under (i), by \eqref{eq:3-7-14:equiv}. Therefore
\begin{equation*}
M_t+V_t=X_t=I_t+C_t^+(0);\qquad 0\le t<\infty,\ \text{a.s.},
\end{equation*}
and \bk{Problem 3.2} yields $V_t=C_t^+(0)$ for all $t$, a.s. Consequently there is an event
$\Omega^*$ with $P(\Omega^*)=1$ such that, for every $\omega\in\Omega^*$: the path $V_\cdot(\omega)$
is continuous and nondecreasing with $V_0(\omega)=0$; $M_0(\omega)=0$; $X_t(\omega)=M_t(\omega)
+V_t(\omega)\ge0$ for all $t$; and $\int_0^\infty\ind_{\{X_s(\omega)>0\}}\,dV_s(\omega)=0$ (this
is \eqref{eq:3-7-14:flat}, since now $\check V=V$).

\emph{Step 2: the Skorohod equation.} Fix $\omega\in\Omega^*$ and put $z=0$, $y(t)\defeq M_t(\omega)$,
$k(t)\defeq V_t(\omega)$. Then $y$ is continuous with $y(0)=0$, and $k$ is a continuous function
which satisfies the three conditions of \bk{Lemma 6.14}: $x(t)\defeq z+y(t)+k(t)=X_t(\omega)\ge0$;
$k(0)=0$ and $k$ is nondecreasing; $\int_0^\infty\ind_{\{x(s)>0\}}\,dk(s)=0$. By the uniqueness
assertion of that lemma, $k$ is given by \bk{(6.29)}:
\begin{equation*}
V_t(\omega)=\max\Bigl[0,\max_{0\le s\le t}\bigl(-M_s(\omega)\bigr)\Bigr]
=\max_{0\le s\le t}\bigl(-M_s(\omega)\bigr);\qquad 0\le t<\infty,
\end{equation*}
the last equality because $-M_0(\omega)=0$. Setting $N\defeq-M\in\cM^{c,\loc}$, we conclude that
\begin{equation*}
X_t=M_t+V_t=\max_{0\le s\le t}N_s-N_t;\qquad 0\le t<\infty,
\end{equation*}
on $\Omega^*$, which is (iii).

\textbf{(iii) $\Rightarrow$ (i).} Let $N\in\cM^{c,\loc}$ be as in (iii) and set
$S_t\defeq\max_{0\le s\le t}N_s$. The process $S$ is continuous and nondecreasing, it is adapted
because $S_t=\sup\{N_s;\ s\in([0,t]\cap\Q)\cup\{t\}\}$, and $S_0=N_0=0$ a.s.\ (elements of
$\cM^{c,\loc}$ start at zero, by \bk{Definition 1.5.15}); replacing $S$ by $S-S_0$, which
changes it only on a null event and preserves adaptedness since $\F_0$ contains the null events, we
may take $S_0\equiv0$. Now
\begin{equation*}
M_t+V_t=X_t=-N_t+S_t;\qquad 0\le t<\infty,\ \text{a.s.},
\end{equation*}
with $-N\in\cM^{c,\loc}$, so \bk{Problem 3.2} gives $V_t=S_t$ (and $M_t=-N_t$) for all $t$, a.s.
In particular $V$ is a.s.\ nondecreasing and $\check V=V=S$.

Fix $\omega$ outside the null event on which $N_0\neq0$ or one of the identities $V=S$, $X=S-N$
fails, and apply \bk{Lemma 6.14} with $z=0$ and the continuous function $y(t)\defeq-N_t(\omega)$,
$y(0)=0$. The function \bk{(6.29)} is
\begin{equation*}
k(t)=\max\Bigl[0,\max_{0\le s\le t}N_s(\omega)\Bigr]=S_t(\omega),
\end{equation*}
because $N_0(\omega)=0$, and the corresponding $x(t)=y(t)+k(t)=S_t(\omega)-N_t(\omega)=X_t(\omega)$.
The lemma asserts that this $k$ satisfies its condition (iii), namely
\begin{equation*}
\int_0^\infty\ind_{\{X_s(\omega)>0\}}\,dS_s(\omega)=0 .
\end{equation*}
Since $\check V=S$ a.s., this is \eqref{eq:3-7-14:flat}, and (i) is proved.

Together with the equivalence of (i) and (ii), this shows that the three conditions are equivalent.

\rmk{The proof gives more: under any of the three conditions the process $V$ is nondecreasing,
$V_t=C_t^+(0)=\max_{0\le s\le t}(-M_s)$, and the local martingale in (iii) is necessarily $N=-M$,
by the uniqueness of the decomposition \bk{(7.1)}. Moreover $M=I$, that is,
$\int_0^t\ind_{\{X_s=0\}}\,dM_s=0$ for all $t$, a.s.}
\end{kex}


%% ========================================================================
\solchapter{4}{Brownian Motion and Partial Differential Equations}

\begin{center}
\begin{minipage}{0.86\textwidth}
\small
The 15 problems and exercises of Chapter 4 for which the book prints no solution: 2.4, 2.8, 2.24, 2.26, 2.27, 2.28, 3.8, 3.11, 3.13, 3.14, 3.15, 4.6, 4.7, 4.13, 4.14.  For 1 more, 3.1, the book prints only a hint or solves only some of the parts; what it leaves out is supplied here.  The book's own ``Solutions to Selected Problems'' fully covers 2.16, 2.22, 2.25, 3.2, which are therefore omitted.
\end{minipage}
\end{center}

\bigskip
\hrule
\bigskip

\solsection{4.2}{Harmonic Functions and the Dirichlet Problem}

\begin{kex}{Exercise}{2.4}{}
Throughout, $D\subseteq\R^d$ is open (the standing assumption of the section) and nonempty.

\emph{Step 1: $\partial D\neq\emptyset$, and $\ol{D}$ is compact.}
The set $\ol{D}$ is closed and bounded, hence compact. If $\partial D=\ol{D}\setminus D$ were empty, then $D=\ol{D}$ would be a nonempty subset of $\R^d$ which is both open and closed, so $D=\R^d$ by the connectedness of $\R^d$; this contradicts the boundedness of $D$.

\emph{Step 2: the maximum is attained on $\partial D$.}
The continuous function $u$ attains its maximum over the compact set $\ol{D}$: there is $x_0\in\ol{D}$ with
\[
u(x_0)=M\defeq\max_{x\in\ol{D}}u(x)<\infty .
\]
If $x_0\in\partial D$, there is nothing more to prove. Otherwise $x_0\in D$, and then
$M=u(x_0)\le\sup_{x\in D}u(x)\le M$, so that $u$ achieves its supremum over $D$ at the point $x_0$ of $D$. Since $D$ is open and connected and $u$ is harmonic in $D$, \bk{Corollary 2.3} shows that $u\equiv M$ on $D$. Now take any $y\in\partial D$, which is possible by Step~1. Because $y\in\ol{D}$, there is a sequence $\{y_n\}_{n=1}^\infty\subseteq D$ with $y_n\to y$, and the continuity of $u$ on $\ol{D}$ gives $u(y)=\lim_{n\to\infty}u(y_n)=M$. Thus the maximum $M$ is attained at $y\in\partial D$ in this case as well. In both cases
\begin{equation}\label{eq:4-2-04:a}
\max_{x\in\ol{D}}u(x)=\max_{x\in\partial D}u(x).
\end{equation}

\emph{Step 3: uniqueness.}
Let $v$ be harmonic in $D$, continuous on $\ol{D}$, with $v=u$ on $\partial D$. The function $w\defeq u-v$ is continuous on $\ol{D}$, of class $C^2$ in $D$ with $\Delta w=\Delta u-\Delta v=0$ there, and $w=0$ on $\partial D$; the same is true of $-w$. Applying \eqref{eq:4-2-04:a} to $w$ and to $-w$, we obtain
\[
\max_{x\in\ol{D}}w(x)=\max_{x\in\partial D}w(x)=0,
\qquad
\max_{x\in\ol{D}}\bigl(-w(x)\bigr)=\max_{x\in\partial D}\bigl(-w(x)\bigr)=0 ,
\]
that is, $w\le 0$ and $w\ge 0$ on $\ol{D}$. Hence $w\equiv 0$, and $v=u$ on $\ol{D}$, in particular on $D$.

\rmk{Applying \eqref{eq:4-2-04:a} to $-u$ shows that $u$ attains its minimum over $\ol{D}$ on $\partial D$, too. Moreover, connectedness of $D$ is not needed for either assertion. Indeed, let $D$ be bounded, open and nonempty, and let $G$ be a connected component of $D$; it is open (balls are connected), bounded and connected. We have $\partial G\subseteq\partial D$: a point $y\in\partial G$ lies in $\ol{G}\subseteq\ol{D}$, and if $y$ were in $D$, some ball $y+B_r\subseteq D$ would be a connected subset of $D$ meeting $G$, hence contained in $G$ by the maximality of the component, so that $y\in G$, which is impossible for a boundary point of the open set $G$. Since $u$ is continuous on $\ol{G}\subseteq\ol{D}$ and harmonic in $G$, \eqref{eq:4-2-04:a} for $G$ gives $\sup_G u\le\max_{\partial G}u\le\max_{\partial D}u$; taking the union over the components yields $\sup_D u\le\max_{\partial D}u$, and then $\max_{\ol{D}}u=\max_{\partial D}u$ because $\ol{D}=D\cup\partial D$. Step~3 goes through unchanged.}
\end{kex}

\begin{kex}{Exercise}{2.8}{}
Here $d=2$, $\partial D=\{(x_1,0);\ x_1\in\R\}$, $\ol D=\{(x_1,x_2);\ x_2\ge 0\}$ and $f\equiv 0$ on $\partial D$. We first identify the function given by \bk{(2.12)}, and then exhibit solutions of $(D,f)$ different from it.

\emph{Step 1: the function \bk{(2.12)} is identically zero.}
Fix $x=(x_1,x_2)\in D$ and write $W=(W^{(1)},W^{(2)})$. Under $P^x$ the process $W^{(2)}-x_2$ is a standard one-dimensional Brownian motion, so by the law of the iterated logarithm (\bk{Theorem 2.9.23}\,(iv)) we have $\liminf_{t\to\infty}W^{(2)}_t=-\infty$, a.s.\ $P^x$. Since $W^{(2)}_0=x_2>0$ and $t\mapsto W^{(2)}_t$ is continuous, the intermediate value theorem shows that $W^{(2)}$ vanishes at some finite time, and therefore
\begin{equation}\label{eq:4-2-08:a}
  P^x[\tau_D<\infty]=1;\qquad x\in D,
\end{equation}
where $\tau_D=\inf\{t\ge 0;\ W_t\in D^c\}=\inf\{t\ge 0;\ W^{(2)}_t\le 0\}$ is the exit time \bk{(2.1)}. On $\{\tau_D<\infty\}$ the continuity of the paths gives $W^{(2)}_{\tau_D}=0$, that is, $W_{\tau_D}\in\partial D$ and $f(W_{\tau_D})=0$. Hence the integrability condition \bk{(2.13)} holds trivially, and $E^xf(W_{\tau_D})=0$ for every $x\in D$. For $x\in\partial D$ we have $\tau_D=0$ and $W_{\tau_D}=x$, $P^x$-a.s., so that $E^xf(W_{\tau_D})=f(x)=0$ as well. Consequently, the function of \bk{(2.12)} is
\[
  u_0(x)\defeq E^xf(W_{\tau_D})=0;\qquad x\in\ol D .
\]
It is, of course, a (bounded) solution of $(D,f)$.

\emph{Step 2: an unbounded solution.}
Let
\[
  u(x_1,x_2)\defeq x_2;\qquad (x_1,x_2)\in\ol D .
\]
This function is continuous on $\ol D$, it is the restriction to $D$ of a polynomial of degree one, so it is of class $C^2(D)$ with $\Delta u=0$ in $D$, and $u(x_1,0)=0=f(x_1,0)$ for every $x_1\in\R$. Thus $u$ satisfies \bk{(2.10)}, \bk{(2.11)}: it is a solution to the Dirichlet problem $(D,f)$. It is unbounded, because $u(0,n)=n$ for every $n\ge 1$, and it is not given by \bk{(2.12)}, because $u(x)=x_2>0=u_0(x)$ at \emph{every} point $x\in D$.

The same reasoning applies to $c\,x_2$ for any real $c\neq 0$, and to many other functions, for instance
\[
  x_1x_2,\qquad 3x_1^2x_2-x_2^3,\qquad e^{x_1}\sin x_2,\qquad \sin x_1\,\sinh x_2 ,
\]
each of which is continuous on $\ol D$, harmonic in $D$ (as one checks by differentiating twice), vanishes on $\{x_2=0\}$, is unbounded on $D$ and does not vanish identically in $D$. So $(D,f)$ has infinitely many unbounded solutions, none of which is of the form \bk{(2.12)}.

\rmk{There is no conflict with \bk{Proposition 2.7}. Its hypotheses on $(D,f)$ are satisfied here: $f$ is bounded and \bk{(2.14)} is \eqref{eq:4-2-08:a}; the proposition then asserts that $u_0\equiv 0$ is the only \emph{bounded} solution, and the example shows that the boundedness of the solution cannot be dropped from its statement. For $u(x)=x_2$ one sees exactly where the proof breaks down: $\{W^{(2)}_t,\F_t\}$ is a continuous martingale under $P^x$, hence so is the stopped process $\{u(W_{t\wedge\tau_D})=W^{(2)}_{t\wedge\tau_D},\F_t\}$ by \bk{Problem 1.3.24}, and thus $E^xu(W_{t\wedge\tau_D})=x_2$ for every $t\ge 0$; but $u(W_{t\wedge\tau_D})\to u(W_{\tau_D})=0$, a.s.\ $P^x$, as $t\to\infty$. The family $\{u(W_{t\wedge\tau_D})\}_{t\ge 0}$ is nonnegative and bounded in $L^1(P^x)$, but not uniformly integrable, so the passage to the limit under the expectation, which the proof of \bk{Proposition 2.7} justifies by the bounded convergence theorem, is not permitted.}
\end{kex}

\begin{kex}{Exercise}{2.24}{}
Throughout, $r>0$ is fixed, $\tau\defeq\tau_{B_r}$, and
\[
K(x,y)\defeq\frac{r^{d-2}\bigl(r^2-\norm{x}^2\bigr)}{\norm{x-y}^d};\qquad x\in\R^d,\ y\in\partial B_r,\ x\neq y .
\]
The assertion \bk{(2.26)} means that $P^x[W_\tau\in\Gamma]=\int_\Gamma K(x,y)\,\mu_r(dy)$ for every
$x\in B_r$ and $\Gamma\in\B(\partial B_r)$. It is the probabilistic reading of \bk{Theorem 2.23}; since
\bk{(2.25)} is misprinted and its derivation in the book is only sketched (see the Remark at the end), we do
not use that theorem, but verify directly that the Poisson integral with kernel $K$ solves the Dirichlet
problem in $B_r$, and then appeal to \bk{Proposition 2.7}.

\emph{Step 1: for each $y\in\partial B_r$, the function $K(\cdot,y)$ is harmonic in $\R^d\setminus\{y\}$,
and $K>0$ on $B_r\times\partial B_r$.} Positivity is evident. For $p\in\R$ and $w\in\R^d\setminus\{0\}$ we
have $\frac{\partial}{\partial w_i}\norm{w}^p=p\,w_i\norm{w}^{p-2}$ and
$\frac{\partial^2}{\partial w_i^2}\norm{w}^p=p\norm{w}^{p-2}+p(p-2)w_i^2\norm{w}^{p-4}$, whence
\[
\Delta\norm{w}^p=p\,(p+d-2)\norm{w}^{p-2};\qquad w\neq0 .
\]
In particular $\Delta\norm{w}^{2-d}=0$ and $\Delta\norm{w}^{-d}=2d\norm{w}^{-d-2}$, and therefore, for
$1\le i\le d$,
\[
\Delta\bigl(w_i\norm{w}^{-d}\bigr)=w_i\,\Delta\norm{w}^{-d}+2\,\frac{\partial}{\partial w_i}\norm{w}^{-d}
=2d\,w_i\norm{w}^{-d-2}-2d\,w_i\norm{w}^{-d-2}=0 .
\]
Now fix $y\in\partial B_r$ and put $w=x-y$. Since $\norm{y}=r$, we have
$\norm{x}^2=\norm{w}^2+2(w,y)+r^2$, so that
\[
K(x,y)=r^{d-2}\,\frac{-\norm{w}^2-2(w,y)}{\norm{w}^d}
=-r^{d-2}\Bigl(\norm{w}^{2-d}+2\sum_{i=1}^d y_i\,w_i\norm{w}^{-d}\Bigr),
\]
which is a linear combination of functions of $w=x-y$ that are of class $C^\infty$ and harmonic in
$\{w\neq0\}$. The Laplacian being invariant under translations, $K(\cdot,y)$ is harmonic in
$\R^d\setminus\{y\}$.

\emph{Step 2: $\int_{\partial B_r}K(x,y)\,\mu_r(dy)=1$ for every $x\in B_r$.} For $x=0$ the integrand is
identically $r^{d-2}r^2/r^d=1$, and $\mu_r$ is a probability measure. Let $x\in B_r\setminus\{0\}$, and write
$x=\rho\theta$ with $0<\rho<r$, $\norm{\theta}=1$. On the open ball $G\defeq\{z\in\R^d;\ \norm{z}<r^2/\rho\}$
we define
\[
g(z)\defeq K\Bigl(\frac{\rho}{r}\,z,\;r\theta\Bigr)
=\frac{r^{d-2}\bigl(r^2-\rho^2\norm{z}^2/r^2\bigr)}{\norm{(\rho/r)z-r\theta}^d}.
\]
For $z\in G$ we have $\norm{(\rho/r)z}<r$, so $(\rho/r)z\neq r\theta$; by Step 1 and because
$\Delta\bigl[h(cz)\bigr]=c^2(\Delta h)(cz)$ for a constant $c$, the function $g$ is harmonic in $G$. If
$\norm{y}=r$, then $\norm{(\rho/r)y}=\rho=\norm{x}$ and
\[
\Bigl\lVert\frac{\rho}{r}\,y-r\theta\Bigr\rVert^2=\rho^2-2\rho\,(y,\theta)+r^2=\norm{\rho\theta-y}^2
=\norm{x-y}^2 ,
\]
so that $g(y)=K(x,y)$ for every $y\in\partial B_r$. Since $\ol{B}_r\subset G$, the mean-value property of the
harmonic function $g$ (\bk{Proposition 2.2}, with $a=0$) yields
\[
\int_{\partial B_r}K(x,y)\,\mu_r(dy)=\int_{\partial B_r}g(y)\,\mu_r(dy)=g(0)=\frac{r^{d-2}\,r^2}{r^d}=1 .
\]

\emph{Step 3: the Poisson integral solves the Dirichlet problem.} Let $f:\partial B_r\to\R$ be continuous,
hence bounded, with $\norm{f}_\infty\defeq\max_{\partial B_r}\abs{f}$, and define $u:\ol{B}_r\to\R$ by
\[
u(x)\defeq\int_{\partial B_r}f(y)K(x,y)\,\mu_r(dy);\quad x\in B_r,\qquad\qquad u(x)\defeq f(x);\quad x\in\partial B_r .
\]
We claim that $u$ is a (bounded) solution of the Dirichlet problem $(B_r,f)$. By Steps 1 and 2,
$\abs{u}\le\norm{f}_\infty$ on $\ol{B}_r$, and \bk{(2.11)} holds by definition.

\emph{$u$ is harmonic in $B_r$.} Fix $0<\rho<r$. On the compact set $\ol{B}_\rho\times\partial B_r$ we have
$\norm{x-y}\ge r-\rho>0$, so $K$ and all its partial derivatives in $x$ of first and second order are
continuous there, hence bounded. By the mean-value theorem the difference quotients in $x$ of
$K(\cdot,y)$ and of its first-order partial derivatives are then bounded on $B_\rho$, uniformly in
$y\in\partial B_r$, and the dominated convergence theorem (with respect to the finite measure
$\abs{f}\,d\mu_r$) shows that $u$ has in $B_\rho$ partial derivatives of first and second order, obtained by
differentiating under the integral sign; these derivatives are continuous in $B_\rho$, again by dominated
convergence. Thus $u\in C^2(B_\rho)$ and, by Step 1,
$\Delta u(x)=\int_{\partial B_r}f(y)\,\Delta_xK(x,y)\,\mu_r(dy)=0$ for $x\in B_\rho$. As $\rho<r$ is
arbitrary, $u$ is of class $C^2$ and harmonic in $B_r$; in particular it is continuous in $B_r$.

\emph{$u$ is continuous at every $y_0\in\partial B_r$.} Let $\eps>0$, and choose $\delta>0$ such that
$\abs{f(y)-f(y_0)}<\eps$ for all $y\in\partial B_r$ with $\norm{y-y_0}<\delta$. Let $z\in\ol{B}_r$ with
$\norm{z-y_0}<\delta/2$. If $z\in\partial B_r$, then $\abs{u(z)-u(y_0)}=\abs{f(z)-f(y_0)}<\eps$. If
$z\in B_r$, then by Step 2 and the positivity of $K$,
\[
\abs{u(z)-u(y_0)}=\Bigl\lvert\int_{\partial B_r}\bigl(f(y)-f(y_0)\bigr)K(z,y)\,\mu_r(dy)\Bigr\rvert
\le\eps+2\norm{f}_\infty\int_{\{y\in\partial B_r;\ \norm{y-y_0}\ge\delta\}}K(z,y)\,\mu_r(dy).
\]
For $y$ in the last domain of integration, $\norm{z-y}\ge\norm{y-y_0}-\norm{z-y_0}\ge\delta/2$, while
\[
r^2-\norm{z}^2=(r+\norm{z})(\norm{y_0}-\norm{z})\le2r\norm{y_0-z};
\]
hence $K(z,y)\le2r^{d-1}(2/\delta)^d\norm{z-y_0}$ there, and
\[
\abs{u(z)-u(y_0)}\le\eps+4\norm{f}_\infty r^{d-1}\Bigl(\frac{2}{\delta}\Bigr)^d\norm{z-y_0}.
\]
In either case $\abs{u(z)-u(y_0)}<2\eps$ as soon as $z\in\ol{B}_r$ is sufficiently close to $y_0$. Thus $u$
is continuous on $\ol{B}_r$, and the claim is proved.

\emph{Step 4: identification of the exit distribution.} Fix $x\in B_r$. The set $B_r$ is bounded, so
$P^x[\tau<\infty]=1$ by \bk{(2.2)}; that is, \bk{(2.14)} holds for $D=B_r$. On $\{\tau<\infty\}$ the
continuity of the paths gives $W_\tau\in\partial B_r$: indeed $W_t\in B_r$ for $t<\tau$, so
$W_\tau\in\ol{B}_r$, and $W_\tau\in B_r^c$ because this set is closed and $\tau$ is the infimum of the times
at which $W$ visits it. Moreover $W_\tau\ind_{\{\tau<\infty\}}$ is a random vector, since $\tau$ is a
stopping time (\bk{Problem 1.2.7}) and $W$ is continuous and adapted, hence progressively measurable.
Consequently
\[
\nu_x(\Gamma)\defeq P^x[W_\tau\in\Gamma],\qquad
\lambda_x(\Gamma)\defeq\int_\Gamma K(x,y)\,\mu_r(dy);\qquad\Gamma\in\B(\partial B_r),
\]
are both probability measures on $\bigl(\partial B_r,\B(\partial B_r)\bigr)$, the second one by Step 2 and
the positivity of $K$.

Let $f:\partial B_r\to\R$ be continuous. By Step 3 the function $u$ constructed there is a bounded solution
of $(B_r,f)$; since $f$ is bounded and \bk{(2.14)} holds, \bk{Proposition 2.7} shows that $u$ has the
representation \bk{(2.12)}:
\begin{equation}\label{eq:4-2-24:a}
\int_{\partial B_r}f\,d\nu_x=E^xf(W_\tau)=u(x)=\int_{\partial B_r}f(y)K(x,y)\,\mu_r(dy)
=\int_{\partial B_r}f\,d\lambda_x .
\end{equation}
Now let $F\subseteq\partial B_r$ be closed and nonempty. The functions
$f_n(y)\defeq\bigl(1-n\,\mathrm{dist}(y,F)\bigr)^+$, $n\ge1$, are continuous on $\partial B_r$, take values
in $[0,1]$, and decrease to $\ind_F$ as $n\to\infty$. Applying \eqref{eq:4-2-24:a} to $f_n$ and letting
$n\to\infty$, we obtain $\nu_x(F)=\lambda_x(F)$ from the bounded convergence theorem. The closed subsets of
$\partial B_r$ form a class which is closed under finite intersections and generates $\B(\partial B_r)$, and
the two probability measures $\nu_x$, $\lambda_x$ agree on it; by Dynkin's $\pi$-$\lambda$ theorem they
agree on $\B(\partial B_r)$. This is \bk{(2.26)}:
\[
P^x[W_{\tau_{B_r}}\in\Gamma]=\int_\Gamma\frac{r^{d-2}\bigl(r^2-\norm{x}^2\bigr)}{\norm{x-y}^d}\,\mu_r(dy);
\qquad x\in B_r,\ \Gamma\in\B(\partial B_r).
\]

\rmk{(a) Formulas \bk{(2.24)} and \bk{(2.25)} are misprinted in the book: the factors $(r-\norm{x-c}^2)$
and $(r-\norm{x}^2)$ should read $(r^2-\norm{x-c}^2)$ and $(r^2-\norm{x}^2)$, in agreement with
\bk{(2.23)} and \bk{(2.26)}. Indeed, for $f\equiv1$ the unique solution of $(B_r,f)$ is $u\equiv1$
(\bk{Proposition 2.7}), whereas the right-hand side of \bk{(2.25)} as printed equals
$r^{d-2}\cdot r\cdot r^{-d}=1/r$ at $x=0$. Steps 1--3 above, together with \bk{Proposition 2.7}, prove the
corrected \bk{Theorem 2.23} without the Kelvin transformation and the change of variable \bk{(2.21)};
conversely, once the corrected \bk{(2.25)} is granted, Step 4 alone gives \bk{(2.26)}.

(b) Two checks. For $x=0$, \bk{(2.26)} reduces to the definition \bk{(2.5)} of $\mu_r$. The argument is
valid for every $d\ge1$; for $d=1$, where $\mu_r$ assigns mass $\frac12$ to each of the points $\pm r$,
\bk{(2.26)} reads $P^x[W_\tau=r]=\frac{r^2-x^2}{2r(r-x)}=\frac{r+x}{2r}$ and
$P^x[W_\tau=-r]=\frac{r-x}{2r}$, the familiar exit probabilities of one-dimensional Brownian motion from
$(-r,r)$.}
\end{kex}

\begin{kex}{Exercise}{2.26}{}
Throughout, $D\subseteq\R^d$ is open, $f:\partial D\to\R$ is bounded and continuous (the hypotheses of \bk{Proposition 2.7} without \bk{(2.14)}), $M\defeq\sup_{\partial D}\abs{f}$, and $\{W_t,\F_t\}$, $(\Omega,\F)$, $\{P^x\}_{x\in\R^d}$ is the Brownian family of the section. The functions $u$ and $v$ are considered on $\ol D$.

\emph{Notation on path space.} Let $C\defeq C[0,\infty)^d$ be the space of continuous functions $\omega:[0,\infty)\to\R^d$, with the metric
\[
\rho(\omega_1,\omega_2)\defeq\sum_{n=1}^\infty\frac1{2^n}\max_{0\le t\le n}\bigl(\norm{\omega_1(t)-\omega_2(t)}\wedge1\bigr)
\]
(this is \bk{(2.4.1)} with $\norm{\cdot}$ in place of $\abs{\cdot}$) and its Borel $\sigma$-field $\B(C)$. Convergence in $\rho$ is uniform convergence on compact intervals, that is, convergence of each of the $d$ components in the space $C[0,\infty)$ of \bk{Problem 2.4.1}; so $C$ is homeomorphic to the $d$-fold product of that separable space, and is itself separable. We claim that $\B(C)$ coincides with the $\sigma$-field $\G$ generated by the coordinate mappings $\omega\mapsto\omega(t)$, $t\ge0$ (for $d=1$ this is \bk{Problem 2.4.2}). The coordinate mappings are continuous, hence $\G\subseteq\B(C)$. Conversely, for fixed $\omega_2\in C$ the maximum over $[0,n]$ in the definition of $\rho(\omega,\omega_2)$ equals, by the continuity of the paths, the supremum over the rational $t\in[0,n]$; so $\omega\mapsto\rho(\omega,\omega_2)$ is $\G$-measurable and every open ball belongs to $\G$. As $C$ is separable, every open set is a countable union of open balls, and $\B(C)\subseteq\G$. In particular $W:\Omega\to C$ is $\F/\B(C)$-measurable, because each $W_t$ is $\F$-measurable. Let $P_*$ be Wiener measure on $(C,\B(C))$, the law of a standard $d$-dimensional Brownian motion. Under $P^x$ the process $W-x$ is a standard Brownian motion, so
\begin{equation}\label{eq:4-2-26:law}
E^x\Phi(W)=\int_C\Phi(x+\omega)\,P_*(d\omega),\qquad x\in\R^d,
\end{equation}
for every bounded Borel function $\Phi:C\to\R$. For $\omega\in C$ and $s\ge0$ write $\theta_s\omega\defeq\omega(s+\cdot)\in C$ and
\[
\tau(\omega)\defeq\inf\{t\ge0;\ \omega(t)\in D^c\}\in[0,\infty],
\]
so that $\tau_D=\tau(W)$. Since $D^c$ is closed and $\omega$ is continuous, the infimum is attained when it is finite, and for each $t\ge 0$ the set $\{\tau\le t\}=\{\omega;\ \omega(s)\in D^c\text{ for some }s\in[0,t]\}$ is closed in $C$: if $\omega_k\to\omega$ uniformly on $[0,t]$ and $\omega_k(s_k)\in D^c$ with $s_k\in[0,t]$, then along a subsequence $s_k\to s\in[0,t]$ and $\omega_k(s_k)\to\omega(s)\in D^c$. Hence $\tau$ is Borel measurable. The mapping $\omega\mapsto\omega(\tau(\omega))$ on $\{\tau<\infty\}$ is Borel as well, being the composition of $\omega\mapsto(\tau(\omega),\omega)$ with the continuous evaluation $(t,\omega)\mapsto\omega(t)$ on $[0,\infty)\times C$.

\begin{lem}
Let $\Phi:C\to\R$ be bounded and Borel measurable, and suppose that
\begin{equation}\label{eq:4-2-26:inv}
\Phi(\omega)=\Phi(\theta_s\omega)\quad\text{whenever } s\ge 0 \text{ and } \omega(t)\in D \text{ for all } t\in[0,s].
\end{equation}
Then $h(x)\defeq E^x\Phi(W)$, $x\in\R^d$, is bounded and Borel measurable, and its restriction to $D$ is of class $C^\infty$ and harmonic.
\end{lem}

\begin{proof}
Clearly $\abs{h}\le\sup\abs{\Phi}$. The mapping $(x,\omega)\mapsto x+\omega$ from $\R^d\times C$ to $C$ is continuous, so $(x,\omega)\mapsto\Phi(x+\omega)$ is measurable with respect to $\B(\R^d\times C)=\B(\R^d)\otimes\B(C)$ (both spaces are separable metric spaces). By \eqref{eq:4-2-26:law} and the measurability assertion of the Fubini--Tonelli theorem, $h$ is Borel measurable.

Fix $a\in D$ and $0<r<\infty$ with $a+\ol B_r\subseteq D$, and let $\sigma\defeq\tau_{a+B_r}=\inf\{t\ge0;\ W_t\notin a+B_r\}$. This is a stopping time of $\{\F_t\}$ (\bk{Problem 1.2.7}), and $P^a[\sigma<\infty]=1$ by \bk{(2.2)}. We work on $(\Omega,\F,P^a)$ and set $W'_t\defeq W_{\sigma+t}-W_\sigma$ on $\{\sigma<\infty\}$, and $W'\defeq 0$ on the $P^a$-null set $\{\sigma=\infty\}$. By \bk{Theorem 2.6.16}, $W'$ is a standard $d$-dimensional Brownian motion independent of $\F_{\sigma+}$, and the random vector $W_\sigma\ind_{\{\sigma<\infty\}}$ is $\F_\sigma$-measurable (\bk{Proposition 1.2.18}). Consequently the pair $(W_\sigma,W')$, with values in $\R^d\times C$, has the product law $\nu\otimes P_*$ under $P^a$, where $\nu(dy)\defeq P^a[W_\sigma\in dy]$.

On $\{W_0=a,\ \sigma<\infty\}$ we have $W_t\in a+B_r\subseteq D$ for $0\le t<\sigma$ and $W_\sigma\in a+\partial B_r\subseteq D$, so \eqref{eq:4-2-26:inv} with $s=\sigma$ gives $\Phi(W)=\Phi(\theta_\sigma W)=\Phi(W_\sigma+W')$, $P^a$-a.s. By Fubini's theorem and \eqref{eq:4-2-26:law},
\begin{equation}\label{eq:4-2-26:sm}
h(a)=E^a\Phi(W_\sigma+W')=\int_{\R^d}\Bigl(\int_C\Phi(y+\omega)\,P_*(d\omega)\Bigr)\nu(dy)=E^a h(W_\sigma).
\end{equation}
Finally, with $\rho_r(\omega)\defeq\inf\{t\ge0;\ \norm{\omega(t)}\ge r\}$ we have $\sigma=\rho_r(W-a)$ and $W_\sigma-a=(W-a)(\rho_r(W-a))$, and $W-a$ has law $P_*$ under $P^a$; likewise $W_{\tau_{B_r}}=W(\rho_r(W))$, and $W$ has law $P_*$ under $P^0$. Hence $W_\sigma-a$ under $P^a$ has the same law as $W_{\tau_{B_r}}$ under $P^0$, which is $\mu_r$ by \bk{(2.5)}, and \eqref{eq:4-2-26:sm} becomes
\[
h(a)=\int_{\partial B_r}h(a+x)\,\mu_r(dx).
\]
Thus the bounded Borel function $h|_D$ has the mean-value property of \bk{Definition 2.1}, and \bk{Proposition 2.5} shows that it is of class $C^\infty$ and harmonic. (The proof of \bk{Proposition 2.5} integrates the function against bounded densities over balls; for a bounded Borel function these integrals are well defined.)
\end{proof}

\emph{Step 1: $u$ and $v$ are harmonic in $D$.}
Let $\tilde f:\R^d\to\R$ be equal to $f$ on $\partial D$ and to $0$ elsewhere; it is Borel measurable, because $\partial D$ is closed and $f$ is continuous on it. Put
\[
\Phi_1(\omega)\defeq\ind_{\{\tau(\omega)=\infty\}},\qquad
\Phi_2(\omega)\defeq\tilde f\bigl(\omega(\tau(\omega))\bigr)\ind_{\{\tau(\omega)<\infty\}},\qquad\omega\in C.
\]
Both are bounded and Borel measurable by the remarks above. They satisfy \eqref{eq:4-2-26:inv}: if $\omega(t)\in D$ for all $t\in[0,s]$, then for $t\ge s$ we have $\omega(t)\in D^c$ if and only if $(\theta_s\omega)(t-s)\in D^c$, whence $\tau(\omega)=s+\tau(\theta_s\omega)$ (both sides possibly infinite) and, when these are finite, $\omega(\tau(\omega))=(\theta_s\omega)(\tau(\theta_s\omega))$.

For $x\in\ol D$ we have $W_{\tau_D}\in\partial D$ on $\{W_0=x,\ \tau_D<\infty\}$: indeed $W_{\tau_D}\in D^c$ because $D^c$ is closed, and $W_{\tau_D}\in\ol D$ because either $\tau_D=0$ and $W_0=x\in\ol D$, or $\tau_D>0$ and $W_{\tau_D}$ is the limit of the points $W_t\in D$, $t\uparrow\tau_D$. Therefore
\[
v(x)=E^x\Phi_1(W),\qquad u(x)=E^x\Phi_2(W),\qquad x\in\ol D,
\]
and the Lemma shows that $u$ and $v$ are (of class $C^\infty$ and) harmonic in $D$. Moreover $0\le v\le1$ and $\abs{u}\le M$.

\emph{Step 2: $v(x)\to0$ as $x\to a$, $x\in D$, at a regular point $a\in\partial D$.}
If $a\in\partial D$ is regular, then
\begin{equation}\label{eq:4-2-26:reg}
\lim_{\substack{x\to a\\ x\in D}}P^x[\tau_D>\eps]=0\qquad\text{for every }\eps>0 .
\end{equation}
For $d\ge2$ this is the implication (ii)$\Rightarrow$(iii) of \bk{Theorem 2.12}. The proof of this implication given in the book (upper semicontinuity of $x\mapsto P^x[\sigma_D>\eps]$, the inequality $\tau_D\le\sigma_D$, and $P^a[\sigma_D>\eps]=0$) makes no use of the assumption $d\ge2$, which enters only in (i)$\Rightarrow$(ii) through \bk{Proposition 3.3.22}, nor of \bk{(2.14)}; so \eqref{eq:4-2-26:reg} holds for $d=1$ too. Since $\{\tau_D=\infty\}\subseteq\{\tau_D>\eps\}$, we obtain
\[
0\le\varlimsup_{\substack{x\to a\\ x\in D}}v(x)\le\lim_{\substack{x\to a\\ x\in D}}P^x[\tau_D>\eps]=0 .
\]

\emph{Step 3: $u(x)\to f(a)$ as $x\to a$, $x\in D$, at a regular point $a\in\partial D$.}
This is \bk{(2.15)} for the present $u$; since $\tau_D$ may now be infinite with positive probability, we repeat the argument for (iii)$\Rightarrow$(i) in \bk{Theorem 2.12} with the indicator in place. For $\eps>0$ and $r>0$ let $q(\eps,r)\defeq P_*[\max_{0\le t\le\eps}\norm{\omega(t)}<r]$; this equals $P^x[\max_{0\le t\le\eps}\norm{W_t-W_0}<r]$ for every $x\in\R^d$, and $q(\eps,r)\uparrow1$ as $\eps\downarrow0$, because the paths are continuous and vanish at $t=0$. Set $A\defeq\{\tau_D<\infty,\ \norm{W_{\tau_D}-W_0}<r\}$. Since
$\{\max_{0\le t\le\eps}\norm{W_t-W_0}<r\}\cap\{\tau_D\le\eps\}\subseteq A$, we have for $x\in D$
\[
P^x[A^c]\le1-q(\eps,r)+P^x[\tau_D>\eps],
\]
and letting first $x\to a$ in $D$, using \eqref{eq:4-2-26:reg}, and then $\eps\downarrow0$, we get
\begin{equation}\label{eq:4-2-26:A}
\lim_{\substack{x\to a\\ x\in D}}P^x\bigl[\tau_D<\infty,\ \norm{W_{\tau_D}-x}<r\bigr]=1\qquad\text{for every }r>0 .
\end{equation}
Now let $\eta>0$ and choose $r>0$ such that $\abs{f(y)-f(a)}\le\eta$ for all $y\in\partial D$ with $\norm{y-a}<2r$. Let $x\in D$ with $\norm{x-a}<r$. On $A\cap\{W_0=x\}$ we have $W_{\tau_D}\in\partial D$ and $\norm{W_{\tau_D}-a}<2r$, so $\abs{f(W_{\tau_D})\ind_{\{\tau_D<\infty\}}-f(a)}\le\eta$; on $A^c\cap\{W_0=x\}$ this quantity is at most $2M$. Hence
\[
\abs{u(x)-f(a)}\le E^x\bigl|f(W_{\tau_D})\ind_{\{\tau_D<\infty\}}-f(a)\bigr|\le\eta+2M\,P^x[A^c],
\]
and \eqref{eq:4-2-26:A} yields $\varlimsup_{x\to a,\,x\in D}\abs{u(x)-f(a)}\le\eta$. As $\eta>0$ is arbitrary, $\lim_{x\to a,\,x\in D}u(x)=f(a)$.

\emph{Step 4: $u+\lambda v$ solves $(D,f)$ when every boundary point is regular.}
Fix $\lambda\in\R$ and let $w\defeq u+\lambda v$ on $\ol D$. Then $\abs{w}\le M+\abs{\lambda}$, so $w$ is bounded, and by Step 1 it is of class $C^2$ (indeed $C^\infty$) in $D$ with $\Delta w=\Delta u+\lambda\Delta v=0$ there; this is \bk{(2.10)}. For $x\in\partial D$ we have $P^x[W_0=x]=1$ and $x\in D^c$, so $P^x[\tau_D=0]=1$; therefore $v(x)=0$, $u(x)=f(x)$, and $w=f$ on $\partial D$, which is \bk{(2.11)}. It remains to check that $w$ is continuous on $\ol D$. In the open set $D$ this follows from its smoothness. Let $a\in\partial D$ and $\eta>0$. By Steps 2 and 3 (the point $a$ is regular) there is $\delta_1>0$ with $\abs{w(x)-f(a)}\le\abs{u(x)-f(a)}+\abs{\lambda}v(x)\le\eta$ for $x\in D$, $\norm{x-a}<\delta_1$, and by the continuity of $f$ there is $\delta_2>0$ with $\abs{w(x)-f(a)}=\abs{f(x)-f(a)}\le\eta$ for $x\in\partial D$, $\norm{x-a}<\delta_2$. Since $w(a)=f(a)$, $w$ is continuous at $a$ relative to $\ol D=D\cup\partial D$. Thus $w=u+\lambda v$ is a bounded solution to the Dirichlet problem $(D,f)$.

\rmk{If every point of $\partial D$ is regular and \bk{(2.14)} fails, then $v(a)>0$ for some $a\in D$, and the solutions $u+\lambda v$, $\lambda\in\R$, of Step 4 are pairwise distinct: the bounded solution to $(D,f)$ is then not unique, and the conclusion of \bk{Proposition 2.7} is false without \bk{(2.14)}. If \bk{(2.14)} holds, then $v\equiv0$ on $D$ and $u$ is the function \bk{(2.12)}. The regularity assumption cannot be dropped in Step 4. Take $d=3$ and $D=\R^3\setminus\{0\}$, so that $\partial D=\{0\}$. By \bk{Proposition 3.3.22}, applied to the Bessel process $\norm{W}$, we have $P^x[W_t\ne0\text{ for all }t>0]=1$ for every $x\in\R^3$; hence the origin is an irregular boundary point, \bk{(2.14)} fails, $v=1$ and $u=0$ on $D$, while $v(0)=0$ and $u(0)=f(0)$. Thus $u+\lambda v$ equals $\lambda$ on $D$ and $f(0)$ at the origin, and it is continuous on $\ol D=\R^3$ only for $\lambda=f(0)$.}
\end{kex}

\begin{kex}{Exercise}{2.27}{}
Throughout, $D\subset\R^d$ is open, bounded and nonempty (if $D=\emptyset$ then $\partial D=\emptyset$ and there is nothing to prove). We fix $a\in\partial D$ and construct a barrier at $a$ in the sense of \bk{Definition 2.13}: a continuous function $v\colon\ol{D}\to\R$, harmonic in $D$, positive on $\ol{D}\setminus\{a\}$ and with $v(a)=0$.

\emph{The case $d=1$.} Here no probability is needed: we take $v(x)\defeq\abs{x-a}$, $x\in\ol{D}$. It is continuous on $\ol{D}$, vanishes at $a$ and is positive elsewhere. Since $a\notin D$ and $D$ is open, for every $x\in D$ there is an interval $(x-\eps,x+\eps)\subset D$, which does not contain $a$; on it $v$ coincides with one of the affine functions $y\mapsto y-a$, $y\mapsto a-y$. Hence $v$ is of class $C^2$ in $D$ with $v''=0$ there, i.e.\ $v$ is harmonic in $D$, and $v$ is a barrier at $a$.

\emph{The case $d\ge 2$.} Let
\[
  f(y)\defeq\norm{y-a},\quad y\in\partial D,
  \qquad\qquad
  v(x)\defeq E^x f(W_{\tau_D}),\quad x\in\ol{D},
\]
with $\tau_D$ as in \bk{(2.1)}. This is \bk{(2.12)} for the boundary function $f$, and it is meaningful: for $x\in D$ we have $P^x[\tau_D<\infty]=1$ by \bk{(2.2)}, and then $W_{\tau_D}\in\partial D$ by continuity of the paths (the path is in $D$ on $[0,\tau_D)$, and $W_{\tau_D}\in D^c$ because $D^c$ is closed); for $x\in\partial D$ we have $\tau_D=0$, $P^x$-a.s., so that
\begin{equation}\label{eq:4-2-27:bdry}
  v(x)=f(x)=\norm{x-a},\qquad x\in\partial D.
\end{equation}
The function $f$ is continuous on $\partial D$, and bounded because $\partial D$ is bounded; in particular $0\le v\le\sup_{\partial D}f<\infty$ on $\ol{D}$. We verify the three properties of a barrier.

\emph{Step 1: $v$ is harmonic in $D$.} Since $f$ is bounded, condition \bk{(2.13)} holds, and \bk{Proposition 2.6} asserts that $v$ is harmonic in $D$. In particular $v$ is continuous in the open set $D$.

\emph{Step 2: $v$ is continuous on $\ol{D}$.} It remains to prove continuity at an arbitrary point $b\in\partial D$, relative to $\ol{D}$. Let $\eps>0$. By hypothesis $b$ is regular for $D$, and $f$ is bounded, measurable and continuous at $b$; since $d\ge2$, the implication (ii)$\Rightarrow$(i) of \bk{Theorem 2.12} applies and gives \bk{(2.15)} at the point $b$:
\[
  \lim_{\substack{x\to b\\ x\in D}}v(x)=f(b).
\]
Thus there is $\delta_1>0$ such that $\abs{v(x)-f(b)}<\eps$ for all $x\in D$ with $\norm{x-b}<\delta_1$. On the other hand, by \eqref{eq:4-2-27:bdry} and the triangle inequality, $\abs{v(x)-f(b)}=\bigl|\norm{x-a}-\norm{b-a}\bigr|\le\norm{x-b}<\eps$ for all $x\in\partial D$ with $\norm{x-b}<\eps$. Since $v(b)=f(b)$, we obtain $\abs{v(x)-v(b)}<\eps$ for all $x\in\ol{D}=D\cup\partial D$ with $\norm{x-b}<\delta_1\wedge\eps$. Hence $v$ is continuous on $\ol{D}$.

\emph{Step 3: $v(a)=0$ and $v>0$ on $\ol{D}\setminus\{a\}$.} By \eqref{eq:4-2-27:bdry}, $v(a)=0$ and $v(x)=\norm{x-a}>0$ for $x\in\partial D\setminus\{a\}$. We have seen that $v\ge0$ on $\ol{D}$, so it remains to exclude that $v(x_0)=0$ for some $x_0\in D$. Suppose this happens, let $G$ be the connected component of $D$ containing $x_0$ (an open, bounded, nonempty set), and let $Z\defeq\{x\in G;\ v(x)=0\}$. This set is nonempty, and closed in $G$ because $v$ is continuous. It is also open: if $x\in Z$, choose $r>0$ with $x+\ol{B}_r\subset G$; by Step~1 and \bk{Proposition 2.2}, $v$ has the mean-value property in $D$, so for every $0<\rho\le r$
\[
  0=v(x)=\int_{\partial B_\rho}v(x+y)\,\mu_\rho(dy).
\]
The integrand is continuous and nonnegative, and $\mu_\rho$ is the normalized surface measure on $\partial B_\rho$ (see the discussion following \bk{(2.5)}), which charges every nonempty relatively open subset of the sphere; hence $v(x+y)=0$ for every $y\in\partial B_\rho$. As $\rho\in(0,r]$ is arbitrary, $v=0$ on $x+\ol{B}_r$, that is, $x+B_r\subset Z$. Since $G$ is connected, $Z=G$: the function $v$ vanishes identically on $G$, and by Step~2 it vanishes on $\ol{G}$ as well.

Next, $\partial G\subset\partial D$. Indeed, $\partial G\subset\ol{G}\subset\ol{D}$; and if some $y\in\partial G$ belonged to $D$, the component $G'$ of $D$ containing $y$ would be an open neighbourhood of $y$, hence would meet $G$, so that $G'=G$ and $y\in G$, which is impossible for a boundary point of the open set $G$. Consequently, by \eqref{eq:4-2-27:bdry}, $0=v(y)=\norm{y-a}$ for every $y\in\partial G$; that is, $\partial G\subset\{a\}$ and $\ol{G}\subset G\cup\{a\}$. Since $a\notin D\supset G$, this yields
\[
  \R^d\setminus\{a\}=G\,\cup\,\bigl(\R^d\setminus\ol{G}\bigr),
\]
a union of two disjoint open sets. Both are nonempty: $G\ni x_0$, and $\R^d\setminus\ol{G}\neq\emptyset$ because $\ol{G}$ is bounded. This contradicts the connectedness of $\R^d\setminus\{a\}$, valid for $d\ge2$. Therefore $v>0$ in $D$, and altogether $v>0$ on $\ol{D}\setminus\{a\}$.

Steps 1--3 show that $v$ is a barrier at $a$. As $a\in\partial D$ was arbitrary, every boundary point of $D$ has a barrier.

\rmk{Regularity of the single point $a$ would give only the continuity of $x\mapsto E^x\norm{W_{\tau_D}-a}$ at $a$; Step~2 uses the regularity of \emph{every} boundary point, in accordance with the hypothesis. The case $d=1$ was treated separately because \bk{Theorem 2.12} is stated for $d\ge2$ and because Step~3 uses the connectedness of $\R^d\setminus\{a\}$. Together with \bk{Proposition 2.15}, the exercise shows that for bounded $D$ all boundary points are regular if and only if all of them have barriers.}
\end{kex}

\begin{kex}{Exercise}{2.28}{}
We identify $\mathbb{C}$ with $\R^2$ through $x+iy\leftrightarrow(x,y)$, so that $\abs{z}$ is the
Euclidean norm, and we write $\ol{X}=X^{(1)}+iX^{(2)}$ for the complex-valued process associated
with an $\R^2$-valued process $X=(X^{(1)},X^{(2)})$. Two preliminaries are used throughout.

\emph{Preliminary A (the filtration).} A two-dimensional Brownian motion
$W=\{W_t,\F_t\}$ on $(\Omega,\F,P)$ in the sense of \bk{Definition 2.5.1} remains one when $\{\F_t\}$
is replaced by the filtration $\wt\F_t\defeq\sigma(\F_{t+}\cup\cN)$, where $\cN$ is the collection of
$P$-null sets of the completion of $\F$; this filtration satisfies the usual conditions.
Indeed, for $0\le s<t$, $A\in\F_{s+}$, a bounded continuous $g\colon\R^2\to\R$ and
$0<\eps<t-s$ we have $A\in\F_{s+\eps}$, hence
$E[g(W_t-W_{s+\eps})\ind_A]=E[g(W_t-W_{s+\eps})]\,P(A)$; letting $\eps\downarrow0$ we obtain, by
continuity of the paths and bounded convergence, $E[g(W_t-W_s)\ind_A]=E[g(W_t-W_s)]\,P(A)$, so that
$W_t-W_s$ is independent of $\F_{s+}$; and every set of $\wt\F_s$ differs from a set of $\F_{s+}$ by
a null set. Since $\F_t\subseteq\wt\F_t$, a statement proved for $\{\wt\F_t\}$ concerning the
\emph{paths} of $W$ is a statement about the original process. We therefore assume from now on, as
the book does throughout Chapter~3, that $\{\F_t\}$ satisfies the usual conditions (if it does not,
the stopping times and the filtration $\{\G_s\}$ below are those of $\{\wt\F_t\}$).

\emph{Preliminary B (the law of a Brownian motion).} For $x\in\R^2$ let $Q^x$ be the measure
\bk{(2.5.1)} on $\bigl(C[0,\infty)^2,\B(C[0,\infty)^2)\bigr)$, i.e.\ $Q^x(F)=Q^0(F-x)$ with $Q^0$
the two-dimensional Wiener measure, and $Q^\mu(F)=\int_{\R^2}Q^x(F)\,\mu(dx)$ as in \bk{(2.5.2)}
(the integrand is Borel-measurable by \bk{Problem 2.5.2}). \emph{If $X=\{X_t,\mathscr{H}_t\}$ is a
two-dimensional Brownian motion with initial distribution $\mu$, on any probability space and
with respect to any filtration, then}
\begin{equation}\label{eq:4-2-28:law}
  P[X_\cdot\in F]=Q^\mu(F)=\int_{\R^2}Q^x(F)\,\mu(dx);\qquad F\in\B(C[0,\infty)^2).
\end{equation}
Indeed, by \bk{Definition 2.5.1}, $X-X_0$ is a continuous process with the finite-dimensional
distributions of the coordinate process under $Q^0$; these determine a probability measure on
$\B(C[0,\infty)^2)$ (the finite-dimensional cylinder sets form a $\pi$-system generating this
$\sigma$-field), so $X-X_0$ has law $Q^0$. By \bk{Problem 2.5.5}, $X-X_0$ is independent of
$\mathscr{H}_0$, hence of $X_0$, and Fubini's theorem gives
$P[X_\cdot\in F]=\int Q^0(F-x)\,\mu(dx)=Q^\mu(F)$. The same statement holds in dimension one, with
the measures $P^x$ of the canonical one-dimensional Brownian family.

\medskip
\textbf{(i)} Let $f$ be entire and nonconstant, write $f=u+iv$ with $u,v\colon\R^2\to\R$, and set
\begin{equation}\label{eq:4-2-28:A}
  A_t\defeq\int_0^t\abs{f'(\ol{W}_r)}^2\,dr,\qquad
  T(s)\defeq\inf\{t\ge0;\ A_t>s\};\qquad 0\le t,s<\infty .
\end{equation}
The process $A$ is continuous, nondecreasing and adapted, with $A_0=0$. The ``appropriate condition''
is
\begin{equation}\label{eq:4-2-28:C}
  P\bigl[A_\infty=\infty\bigr]=1,\qquad A_\infty\defeq\lim_{t\to\infty}A_t ,
\end{equation}
and we prove the following.

\begin{claim}
Assume \eqref{eq:4-2-28:C}. Then each $T(s)$ is an a.s.\ finite stopping time of $\{\F_t\}$, and with
$\G_s\defeq\F_{T(s)}$ the process
\[
  Z=\bigl\{Z_s\defeq f\bigl(\ol{W}_{T(s)}\bigr),\ \G_s;\ 0\le s<\infty\bigr\}
\]
is a complex-valued Brownian motion, with initial distribution the law of $f(\ol{W}_0)$; in
particular $\ol{B}\defeq Z-Z_0$ is a complex-valued Brownian motion starting at the origin. Moreover,
almost surely,
\begin{equation}\label{eq:4-2-28:rep}
  f(\ol{W}_t)=Z_{A_t}=f(\ol{W}_0)+\ol{B}_{A_t};\qquad 0\le t<\infty .
\end{equation}
\end{claim}

\begin{proof}
\emph{Step 1: two orthogonal local martingales.} The functions $u,v$ are of class $C^\infty$ and
satisfy the Cauchy--Riemann equations $u_x=v_y$, $u_y=-v_x$; consequently $\Delta u=\Delta v=0$
and $\abs{f'}^2=\abs{u_x+iv_x}^2=u_x^2+u_y^2=v_x^2+v_y^2$. By \bk{Problem 2.5.5} the processes
$W^{(j)}-W^{(j)}_0$ are continuous martingales with
$\qv{W^{(j)}-W^{(j)}_0,W^{(k)}-W^{(k)}_0}_t=\delta_{jk}t$, and $W_0$ is $\F_0$-measurable. It\^o's
rule (\bk{Theorem 3.3.6}, with no bounded variation part) applied to $u$ and to $v$ gives, a.s.\
for all $t$,
\begin{align}
  M^{(1)}_t&\defeq u(W_t)-u(W_0)
     =\int_0^tu_x(W_r)\,dW^{(1)}_r+\int_0^tu_y(W_r)\,dW^{(2)}_r,\label{eq:4-2-28:M1}\\
  M^{(2)}_t&\defeq v(W_t)-v(W_0)
     =-\int_0^tu_y(W_r)\,dW^{(1)}_r+\int_0^tu_x(W_r)\,dW^{(2)}_r,\label{eq:4-2-28:M2}
\end{align}
the terms $\tfrac12\int_0^t\Delta u(W_r)\,dr$ and $\tfrac12\int_0^t\Delta v(W_r)\,dr$ being zero, and
$v_x,v_y$ having been replaced by $-u_y,u_x$. The integrands are continuous and adapted, hence
progressively measurable with paths bounded on bounded intervals, so $M^{(1)},M^{(2)}\in\cM^{c,\loc}$.
Applying \bk{(3.2.30)} twice (\bk{Proposition 3.2.24}) and using bilinearity of the
cross-variation, we find
\begin{equation}\label{eq:4-2-28:qv}
  \qv{M^{(1)}}_t=\qv{M^{(2)}}_t=\int_0^t\bigl(u_x^2+u_y^2\bigr)(W_r)\,dr=A_t,\qquad
  \qv{M^{(1)},M^{(2)}}_t=\int_0^t\bigl(-u_xu_y+u_yu_x\bigr)(W_r)\,dr=0 .
\end{equation}

\emph{Step 2: F.~B.~Knight's theorem.} By \eqref{eq:4-2-28:qv} and \eqref{eq:4-2-28:C} the pair
$M=(M^{(1)},M^{(2)})$ satisfies the hypotheses of \bk{Theorem 3.4.13} with $d=2$, and the two
time-changes $T_1,T_2$ of that theorem both coincide with $T$ of \eqref{eq:4-2-28:A}. Each $T(s)$ is a
stopping time, finite on $\Omega_0\defeq\{A_\infty=\infty\}$; on the null set
$\Omega\setminus\Omega_0\in\F_0$ we put $B\defeq0$. \bk{Theorem 3.4.13} asserts that
\[
  B^{(1)}_s\defeq M^{(1)}_{T(s)},\qquad B^{(2)}_s\defeq M^{(2)}_{T(s)};\qquad 0\le s<\infty,
\]
are \emph{independent} standard, one-dimensional Brownian motions. It follows that
$B=(B^{(1)},B^{(2)})$ is a two-dimensional Brownian motion starting at the origin relative to its own
filtration $\{\F^B_s\}$: for $s<t$, $\Gamma_1,\Gamma_2\in\B(\R)$ and $C_j\in\F^{B^{(j)}}_s$, the
independence of $\F^{B^{(1)}}_\infty$ and $\F^{B^{(2)}}_\infty$ and the Brownian property of each
$B^{(j)}$ give
\[
  P\bigl[B^{(1)}_t-B^{(1)}_s\in\Gamma_1,\,B^{(2)}_t-B^{(2)}_s\in\Gamma_2,\,C_1\cap C_2\bigr]
  =\prod_{j=1}^2P\bigl[B^{(j)}_t-B^{(j)}_s\in\Gamma_j\bigr]\,P[C_j],
\]
and the sets $C_1\cap C_2$ form a $\pi$-system generating $\F^B_s$, the rectangles
$\Gamma_1\times\Gamma_2$ one generating $\B(\R^2)$; by Dynkin's $\pi$-$\lambda$ theorem $B_t-B_s$ is
independent of $\F^B_s$ and has the law $N(0,(t-s)I_2)$. Thus
\begin{equation}\label{eq:4-2-28:B}
  \ol{B}_s=f\bigl(\ol{W}_{T(s)}\bigr)-f(\ol{W}_0),\quad0\le s<\infty,
  \ \text{ is a complex-valued Brownian motion starting at } 0,
\end{equation}
relative to $\{\F^B_s\}$. Finally, \bk{(3.4.17)} of \bk{Theorem 3.4.6}, applied to $M^{(1)}$ and to
$M^{(2)}$, yields $M^{(j)}_t=B^{(j)}_{A_t}$ for all $t\ge0$, a.s., which is
\eqref{eq:4-2-28:rep}.

\emph{Step 3: the filtration $\{\G_s\}$ and the initial condition.} It remains to see that $B$ is a
Brownian motion relative to the larger filtration $\G_s=\F_{T(s)}$, which satisfies the usual
conditions by \bk{Theorem 3.4.6}; then $B$ is independent of $\G_0\supseteq\F_0$ (\bk{Problem 2.5.5}),
$Z_0=f(\ol{W}_0)$ is $\G_0$-measurable, and $Z=Z_0+\ol{B}$ satisfies \bk{Definition 2.5.1} relative
to $\{\G_s\}$. By \bk{Theorem 3.4.6}, each $\{B^{(j)}_s,\G_s\}$ is a standard Brownian motion, so
$B^{(j)}$ is a continuous $\{\G_s\}$-martingale with $\qv{B^{(j)}}_s=s$. By P.~L\'evy's
\bk{Theorem 3.3.16} it suffices to show that $B^{(1)}B^{(2)}$ is a $\{\G_s\}$-martingale, for then
$\qv{B^{(1)},B^{(2)}}\equiv0$. Fix $n\ge1$ and let $\wt M^{(j)}_t\defeq M^{(j)}_{t\wedge T(n)}$.
Then $\qv{\wt M^{(j)}}_\infty=A_{T(n)}=n$ a.s.\ (\bk{Problem 3.4.5}\,(ii)), so by
\bk{Problem 1.5.24} $\wt M^{(j)}$ is a martingale bounded in $L^2$ with
$\sup_tE(\wt M^{(j)}_t)^2\le n$, and Doob's maximal inequality gives
$E\bigl[\sup_{t}(\wt M^{(j)}_t)^2\bigr]\le4n$. Since $\qv{M^{(1)},M^{(2)}}=0$, the product
$N_t\defeq\wt M^{(1)}_t\wt M^{(2)}_t$ is a continuous local martingale; it is dominated by the
integrable random variable $\tfrac12\sup_t(\wt M^{(1)}_t)^2+\tfrac12\sup_t(\wt M^{(2)}_t)^2$, hence
(dominated convergence along a localizing sequence) it is a uniformly integrable martingale. The
optional sampling \bk{Theorem 1.3.22} gives, for $0\le s_1<s_2\le n$,
\[
  E\bigl[B^{(1)}_{s_2}B^{(2)}_{s_2}\,\big|\,\G_{s_1}\bigr]
  =E\bigl[N_{T(s_2)}\,\big|\,\F_{T(s_1)}\bigr]=N_{T(s_1)}=B^{(1)}_{s_1}B^{(2)}_{s_1}\quad\text{a.s.},
\]
and $n$ is arbitrary.
\end{proof}

\begin{rem}
Step~3 shows that, because the two components of $f(\ol{W})$ are run by the \emph{same} clock $A$,
the independence of $B^{(1)}$ and $B^{(2)}$ can also be obtained from L\'evy's theorem. Parts (ii) and
(iii) below use only \eqref{eq:4-2-28:B} and \eqref{eq:4-2-28:rep}, i.e.\ Steps 1 and 2.
\end{rem}

\begin{rem}
Condition \eqref{eq:4-2-28:C} is in fact automatically satisfied when $f$ is entire and nonconstant.
Let $E\defeq\{A_\infty<\infty\}$ and $\tau_k\defeq\inf\{t\ge0;\ A_t\ge k\}$. The stopped local
martingale $M^{(j)}_{\cdot\wedge\tau_k}$ has $\qv{M^{(j)}}_{\tau_k}\le k$, so it converges a.s.\ as
$t\to\infty$ (\bk{Problem 1.5.24}\,(i)); since $E=\bigcup_k\{\tau_k=\infty\}$, the limit
$\lim_{t\to\infty}f(\ol{W}_t)$ exists in $\mathbb{C}$ a.s.\ on $E$. On the other hand, choose
$z_1,z_2$ with $f(z_1)\ne f(z_2)$, put $\delta=\abs{f(z_1)-f(z_2)}/3$ and take $\rho>0$ with
$\abs{f(z)-f(z_k)}<\delta$ for $\abs{z-z_k}<\rho$, $k=1,2$. The set
$G_z\defeq\{w;\ \abs{w(t)-z}<\rho\text{ for some }t\ge0\}$ is open in $C[0,\infty)^2$, and
$Q^x(G_z)=1$ for all $x,z\in\R^2$: this is trivial for $x=z$, and for $x\ne z$ we have
$Q^x(G_z)=Q^{x-z}(G_0)\ge Q^{x-z}[\inf_t\abs{w(t)}=0]=1$ by \bk{Problem 3.3.23}\,(i), \bk{(3.3.15)}
and \eqref{eq:4-2-28:law}. For every integer $n$, $\{W_{n+t},\F_{n+t};\ t\ge0\}$ is a Brownian
motion, so \eqref{eq:4-2-28:law} yields $P[W_{n+\cdot}\in G_{z_k}]=1$. Hence, a.s., there are
arbitrarily large times $t,t'$ with $\abs{f(\ol{W}_t)-f(z_1)}<\delta$ and
$\abs{f(\ol{W}_{t'})-f(z_2)}<\delta$, whence $\abs{f(\ol{W}_t)-f(\ol{W}_{t'})}>\delta$, and
$f(\ol{W}_t)$ does not converge. Therefore $P(E)=0$. (\bk{Problem 3.3.23} is obtained in the book by
the method of \bk{Proposition 3.3.22}; this remark is not used in what follows.)
\end{rem}

\medskip
\textbf{(ii)} We apply (i) to $f(z)=\xi e^z$, which is entire and nonconstant, with $f'=f$ and
$\abs{f'(x+iy)}^2=\abs{\xi}^2e^{2x}$. Thus
\[
  A_t=\abs{\xi}^2\int_0^te^{2W^{(1)}_r}\,dr;\qquad 0\le t<\infty,
\]
and we must verify \eqref{eq:4-2-28:C}. The set
$H\defeq\{w\in C[0,\infty);\ \int_0^\infty e^{2w(r)}\,dr=\infty\}$ is Borel, because
$w\mapsto\int_0^ne^{2w(r)}\,dr$ is continuous on $C[0,\infty)$ for each $n$. The function
$y\mapsto e^{2y}$ is Borel-measurable, nonnegative and strictly positive on a set of positive
Lebesgue measure (on all of $\R$), so \bk{(3.6.55)} of \bk{Problem 3.6.30}, applied to the canonical
one-dimensional Brownian family, gives $P^x(H)=1$ for every $x\in\R$. Now $W^{(1)}$ is a
one-dimensional Brownian motion with initial distribution $\nu\defeq P(W^{(1)}_0)^{-1}$, so by the
one-dimensional version of \eqref{eq:4-2-28:law},
\[
  P[A_\infty=\infty]=P\bigl[W^{(1)}_\cdot\in H\bigr]=\int_{\R}P^x(H)\,\nu(dx)=1 .
\]
The Claim of part (i) therefore applies: with $T(s)=\inf\{t\ge0;\ A_t>s\}$,
\begin{equation}\label{eq:4-2-28:ii}
  M_{T(s)}=\xi e^{\ol{W}_0}+\ol{B}_s\quad(0\le s<\infty),\qquad
  M_t=\xi e^{\ol{W}_0}+\ol{B}_{A_t}\quad(0\le t<\infty),\qquad\text{a.s.},
\end{equation}
where $\ol{B}$ is a complex-valued Brownian motion starting at the origin; and
$\{M_{T(s)},\F_{T(s)};\ 0\le s<\infty\}$ is a complex-valued Brownian motion with initial
distribution the law of $\xi e^{\ol{W}_0}$. Here the integrand of $A$ is strictly positive, so on
$\{A_\infty=\infty\}$ the function $t\mapsto A_t$ is a continuous, strictly increasing bijection of
$[0,\infty)$ onto itself and $T$ is its inverse: $M$ and the Brownian motion $M_{T(\cdot)}$ trace the
same curve in $\mathbb{C}$.

\medskip
\textbf{(iii)} We have to show that a Bessel process $R$ with dimension $d\ge2$ starting at $r\ge0$
satisfies $P[R_t>0;\ \forall\,0<t<\infty]=1$ (\bk{(3.3.19)}).

\emph{Reduction.} By \bk{Definition 3.3.19}, $R_t=\norm{W_t}$, where $W=\{W_t,\F_t\}$ is, under
the measure in question, a $d$-dimensional Brownian motion starting at $(r,0,\dots,0)$. By
\bk{Definition 2.5.1}, $V\defeq(W^{(1)},W^{(2)})$ is then a two-dimensional Brownian motion relative
to $\{\F_t\}$ starting at $x_r\defeq(r,0)$, and $R_t\ge\norm{V_t}$. The set
\[
  F_+\defeq\{w\in C[0,\infty)^2;\ w(t)\ne0,\ \forall\,0<t<\infty\}
\]
is Borel, its complement being the union over $n\ge1$ of the closed sets
$\{w;\ \min_{1/n\le t\le n}\abs{w(t)}=0\}$; similarly
$F\defeq\{w;\ w(t)\ne0,\ \forall\,0\le t<\infty\}$ is Borel. By \eqref{eq:4-2-28:law},
\[
  P[R_t>0;\ \forall\,0<t<\infty]\ \ge\ P[V_\cdot\in F_+]=Q^{x_r}(F_+),
\]
so it suffices to prove that $Q^x(F_+)=1$ for every $x\in\R^2$.

\emph{The case $x\ne0$.} Let $W$ be a two-dimensional Brownian motion \emph{starting at the origin}
on some probability space $(\Omega,\F,P)$, with a filtration satisfying the usual conditions
(Preliminary~A), and let $\xi\defeq x^{(1)}+ix^{(2)}\in\mathbb{C}\setminus\{0\}$. By
\eqref{eq:4-2-28:ii} of part (ii), with $\ol{W}_0=0$, there is a complex-valued Brownian motion
$\ol{B}$ starting at the origin such that, on the event $\{A_\infty=\infty\}$ of probability one,
\[
  \xi+\ol{B}_s=M_{T(s)}=\xi\exp\bigl(\ol{W}_{T(s)}\bigr)\ne0\qquad\text{for every }0\le s<\infty,
\]
because $T(s)<\infty$ and the exponential function has no zero. But $x+B$ is a two-dimensional
Brownian motion starting at $x$ (relative to $\{\F^B_s\}$, by \eqref{eq:4-2-28:B}), so
\eqref{eq:4-2-28:law} gives
\[
  Q^x(F_+)\ \ge\ Q^x(F)=P\bigl[x+B_\cdot\in F\bigr]=1 .
\]

\emph{The case $x=0$.} Let $w$ denote the coordinate process on $C[0,\infty)^2$ and fix $\eps>0$.
Under $Q^0$, the process $\{w(\eps+t),\F^w_{\eps+t};\ 0\le t<\infty\}$ is a two-dimensional
Brownian motion whose initial distribution $\mu_\eps$ is normal with mean zero and covariance
$\eps I_2$; in particular $\mu_\eps(\{0\})=0$. By \eqref{eq:4-2-28:law} and the case already
treated,
\[
  Q^0\bigl[w(t)\ne0;\ \forall\,\eps\le t<\infty\bigr]=Q^0\bigl[w(\eps+\cdot)\in F\bigr]
  =\int_{\R^2\setminus\{0\}}Q^y(F)\,\mu_\eps(dy)=1 .
\]
Taking $\eps=1/n$ and intersecting over $n\ge1$ we obtain $Q^0(F_+)=1$. This completes the new
proof of \bk{Proposition 3.3.22}: the argument of the book, based on It\^o's rule for $\log R_t$ and
on the hitting probabilities \bk{(3.3.21)}--\bk{(3.3.23)}, has been replaced by the observation that
a planar Brownian motion issued from $\xi\ne0$ is a time-change of $\xi e^{\ol{W}}$, a process which
visibly never vanishes.
\end{kex}

\solsection{4.3}{The One-Dimensional Heat Equation}

\begin{kex}{Problem}{3.1}{}
The book's Section 4.5 asserts without proof a pointwise bound on the derivatives of $p$, deduces from it the
domination \bk{(5.3)}, and proves \bk{(3.5)} for $n=0$, $m=1$ only, leaving the general case to the reader. We
prove the bound on the derivatives of $p$ and then \bk{(3.5)} for all $n$, $m$; the argument is self-contained.

Throughout, $S\defeq(0,\frac{1}{2a})\times\R$, and for integers $n,m\ge 0$ we write
\[
  p_{n,m}(t;x,y)\defeq\frac{\partial^{n+m}}{\partial t^n\,\partial x^m}\,p(t;x,y),\qquad t>0,\ x,y\in\R .
\]
The function $(t,x,y)\mapsto p(t;x,y)=(2\pi t)^{-1/2}e^{-(x-y)^2/2t}$ is infinitely differentiable on
$(0,\infty)\times\R^2$, so its mixed partial derivatives do not depend on the order of the differentiations,
and $\frac{\partial}{\partial x}p_{n,m}=p_{n,m+1}$, $\frac{\partial}{\partial t}p_{n,m}=p_{n+1,m}$.

\emph{Step 1: the form of the derivatives of $p$.} We claim that for all $n,m\ge 0$ there is a polynomial
$R_{n,m}(s,z)$ in two variables, of degree at most $2n+m$ in $z$, such that
\begin{equation}\label{eq:4-3-01:a}
  p_{n,m}(t;x,y)=p(t;x,y)\,R_{n,m}\Bigl(\frac1t,\;x-y\Bigr),\qquad t>0,\ x,y\in\R .
\end{equation}
Indeed, $R_{0,0}\equiv 1$. Since
$\frac{\partial}{\partial x}p=-\frac{x-y}{t}\,p$ and
$\frac{\partial}{\partial t}p=\bigl(-\frac{1}{2t}+\frac{(x-y)^2}{2t^2}\bigr)p$, differentiation of
\eqref{eq:4-3-01:a} with respect to $x$, respectively $t$ (note $\frac{\partial}{\partial t}\frac1t=-\frac1{t^2}$),
shows that \eqref{eq:4-3-01:a} holds for $(n,m+1)$, respectively $(n+1,m)$, with
\[
  R_{n,m+1}=-sz\,R_{n,m}+\frac{\partial R_{n,m}}{\partial z},\qquad
  R_{n+1,m}=\Bigl(-\frac s2+\frac{s^2z^2}{2}\Bigr)R_{n,m}-s^2\,\frac{\partial R_{n,m}}{\partial s}.
\]
These are again polynomials; the first recursion raises the degree in $z$ by at most $1$ and the second by at
most $2$. Induction (first on $m$ with $n=0$, then on $n$) proves the claim.

\emph{Step 2: domination.} Fix $\beta>0$, $\eps>0$ and $0<t_0<t_1\defeq\frac{1}{2(a+\eps)}$, and set
$B\defeq\{(t,x);\ t_0<t<t_1,\ \abs{x}<\beta\}$. We show that for all $n,m\ge 0$ there is a constant
$D(n,m)<\infty$, depending also on $a,\beta,\eps,t_0$, such that
\begin{equation}\label{eq:4-3-01:b}
  \abs{f(y)\,p_{n,m}(t;x,y)}\le g_{n,m}(y)\defeq D(n,m)\,\abs{f(y)}\,e^{-ay^2},\qquad (t,x)\in B,\ y\in\R .
\end{equation}
By \bk{(3.3)}, $g_{n,m}$ is integrable over $\R$. Let $(t,x)\in B$, $y\in\R$, and $N\defeq 2n+m$. Because
$(x-y)^2\ge 0$ and $0<t<t_1$, we have $\frac{(x-y)^2}{2t}\ge\frac{(x-y)^2}{2t_1}=(a+\eps)(x-y)^2$; moreover
$(x-y)^2\ge(\abs{y}-\abs{x})^2\ge y^2-2\abs{x}\abs{y}\ge y^2-2\beta\abs{y}$. Therefore
\[
  \frac{(x-y)^2}{2t}\;\ge\;(a+\eps)y^2-2(a+\eps)\beta\abs{y}
  \;\ge\;\Bigl(a+\frac\eps2\Bigr)y^2-\frac{2(a+\eps)^2\beta^2}{\eps},
\]
where the last step uses $2(a+\eps)\beta\abs{y}\le\frac\eps2y^2+\frac{2}{\eps}(a+\eps)^2\beta^2$. Hence, with
$c_1\defeq(2\pi t_0)^{-1/2}\exp\{2(a+\eps)^2\beta^2/\eps\}$,
\[
  p(t;x,y)\le c_1\,e^{-(a+\eps/2)y^2}.
\]
Next write $R_{n,m}(s,z)=\sum_{j=0}^{J}\sum_{k=0}^{N}c_{jk}s^jz^k$. Since $0<\frac1t<\frac1{t_0}$,
$\abs{z}^k\le 1+\abs{z}^N$ for $0\le k\le N$, and
$\abs{x-y}^N\le(\beta+\abs{y})^N\le 2^N(\beta^N+\abs{y}^N)$, we obtain
\[
  \Bigl|R_{n,m}\Bigl(\frac1t,x-y\Bigr)\Bigr|\le\sum_{j,k}\abs{c_{jk}}\,t_0^{-j}\,\bigl(1+\abs{x-y}^N\bigr)
  \le c_2\,(1+\abs{y}^N),\qquad c_2\defeq 2^N(1+\beta^N)\sum_{j,k}\abs{c_{jk}}\,t_0^{-j}.
\]
Finally $c_3\defeq\sup_{y\in\R}(1+\abs{y}^N)e^{-\eps y^2/2}<\infty$. Combining these three estimates with
\eqref{eq:4-3-01:a} gives \eqref{eq:4-3-01:b} with $D(n,m)=c_1c_2c_3$.

\emph{Step 3: the integrals are well defined and continuous.} Let $(t^*,x^*)\in S$. Since $t^*<\frac1{2a}$, we
can choose $\eps>0$ with $t^*<t_1=\frac{1}{2(a+\eps)}$, then $t_0\in(0,t^*)$ and $\beta>\abs{x^*}$; thus every
point of $S$ lies in an open rectangle $B$ of the kind considered in Step~2. Consequently, by
\eqref{eq:4-3-01:b}, for every $(t,x)\in S$ the Borel function $y\mapsto f(y)p_{n,m}(t;x,y)$ is integrable, and
\[
  u_{n,m}(t,x)\defeq\int_{-\infty}^{\infty}f(y)\,p_{n,m}(t;x,y)\,dy,\qquad (t,x)\in S,
\]
is a well-defined real number. For $n=m=0$ this shows that $E^x\abs{f(W_t)}=\int\abs{f(y)}p(t;x,y)\,dy<\infty$,
so that $u=u_{0,0}$ of \bk{(3.4)} is defined on $S$. Moreover each $u_{n,m}$ is continuous on $S$: if
$(t_k,x_k)\to(t^*,x^*)$, then $(t_k,x_k)\in B$ for large $k$, the integrands
$f(y)p_{n,m}(t_k;x_k,y)$ converge to $f(y)p_{n,m}(t^*;x^*,y)$ for every $y$ by the continuity of $p_{n,m}$, and
they are dominated by $g_{n,m}$; the dominated convergence theorem gives $u_{n,m}(t_k,x_k)\to u_{n,m}(t^*,x^*)$.

\emph{Step 4: differentiation under the integral sign.} We use the following standard fact.

\begin{lem}
Let $I\subseteq\R$ be an open interval and $G:I\times\R\to\R$ a function such that $G(\theta,\cdot)$ is Borel
measurable and integrable for every $\theta\in I$, $G(\cdot,y)$ is differentiable on $I$ for every $y\in\R$, and
$\abs{\frac{\partial}{\partial\theta}G(\theta,y)}\le g(y)$ for all $\theta\in I$, $y\in\R$, with $g$ integrable.
Then $\Phi(\theta)\defeq\int_{-\infty}^\infty G(\theta,y)\,dy$ is differentiable on $I$ and
$\Phi'(\theta)=\int_{-\infty}^\infty\frac{\partial}{\partial\theta}G(\theta,y)\,dy$.
\end{lem}
\begin{proof}
Fix $\theta\in I$ and a sequence $h_k\to0$, $h_k\ne0$, $\theta+h_k\in I$. The functions
$q_k(y)\defeq h_k^{-1}[G(\theta+h_k,y)-G(\theta,y)]$ are Borel measurable and converge pointwise to
$\frac{\partial}{\partial\theta}G(\theta,y)$, which is therefore Borel measurable in $y$. By the mean-value
theorem, $q_k(y)=\frac{\partial}{\partial\theta}G(\vartheta,y)$ for some $\vartheta=\vartheta_k(y)$ between
$\theta$ and $\theta+h_k$, hence in $I$, so $\abs{q_k}\le g$. The dominated convergence theorem yields
$h_k^{-1}[\Phi(\theta+h_k)-\Phi(\theta)]=\int q_k(y)\,dy\to\int\frac{\partial}{\partial\theta}G(\theta,y)\,dy$.
\end{proof}

We claim that for all $n,m\ge0$ the partial derivatives of $u_{n,m}$ exist at every point of $S$ and
\begin{equation}\label{eq:4-3-01:c}
  \frac{\partial}{\partial x}u_{n,m}=u_{n,m+1},\qquad \frac{\partial}{\partial t}u_{n,m}=u_{n+1,m}
  \qquad\text{on } S .
\end{equation}
Let $(t^*,x^*)\in S$ and choose $B$ as in Step~3 with $(t^*,x^*)\in B$. For the first identity we apply the
Lemma with $I=(-\beta,\beta)$, $\theta=x$ and $G(x,y)=f(y)p_{n,m}(t^*;x,y)$: for each $x\in I$ this is Borel
measurable and integrable in $y$ by Step~3, it is differentiable in $x$ with derivative
$f(y)p_{n,m+1}(t^*;x,y)$, and this derivative is bounded in absolute value by the integrable function
$g_{n,m+1}(y)$ by \eqref{eq:4-3-01:b}, because $(t^*,x)\in B$ for $x\in I$. Hence
$\frac{\partial}{\partial x}u_{n,m}(t^*,x^*)=u_{n,m+1}(t^*,x^*)$. For the second identity we apply the Lemma with
$I=(t_0,t_1)$, $\theta=t$ and $G(t,y)=f(y)p_{n,m}(t;x^*,y)$, whose derivative in $t$ is $f(y)p_{n+1,m}(t;x^*,y)$,
bounded by $g_{n+1,m}(y)$ since $(t,x^*)\in B$ for $t\in I$. This proves \eqref{eq:4-3-01:c}.

\emph{Step 5: conclusion.} By \eqref{eq:4-3-01:c} and induction on the number of differentiations, every
iterated partial derivative of $u=u_{0,0}$, taken in any order, exists on $S$ and equals $u_{n,m}$, where $n$ is
the number of differentiations with respect to $t$ and $m$ the number with respect to $x$. In particular,
differentiating $m$ times in $x$ and then $n$ times in $t$,
\[
  \frac{\partial^{n+m}}{\partial t^n\,\partial x^m}\,u(t,x)=u_{n,m}(t,x)
  =\int_{-\infty}^{\infty}f(y)\,\frac{\partial^{n+m}}{\partial t^n\,\partial x^m}\,p(t;x,y)\,dy,
  \qquad 0<t<\frac1{2a},\ x\in\R,
\]
which is \bk{(3.5)}. All these derivatives are continuous on $S$ by Step~3, so $u$ is of class $C^\infty$
on $S$.

\rmk{This also gives the remaining assertion made before the problem in the book: by \bk{(3.5)} and \bk{(3.2)},
$\frac{\partial u}{\partial t}=\int f(y)\frac{\partial p}{\partial t}\,dy
=\frac12\int f(y)\frac{\partial^2p}{\partial x^2}\,dy=\frac12\frac{\partial^2u}{\partial x^2}$ on $S$, i.e.\ $u$
satisfies the heat equation \bk{(3.1)}. As a check of Step~1, the recursions give $R_{0,1}=-sz$,
$R_{0,2}=s^2z^2-s$ and $R_{1,0}=\frac12(s^2z^2-s)=\frac12R_{0,2}$, in agreement with \bk{(3.2)}.}
\end{kex}

\begin{kex}{Exercise}{3.8}{}
We prove (ii) first; part (i) is then the special case $f\equiv 0$. Instead of establishing the uniform integrability suggested in the hint directly, we identify the measure in the representation \bk{(3.22)}; the uniform integrability is obtained as a by-product in the closing remark.

Throughout, $u\ge 0$ is of class $C^{1,2}$ on $(0,T)\times\R$, satisfies \bk{(3.1)} there, and
\begin{equation}\label{eq:4-3-08:lim}
\lim_{\substack{t \downarrow 0 \\ y \to x}} u(t,y) = f(x);\quad x \in \R,
\end{equation}
for some function $f\colon\R\to\R$. Being a limit of nonnegative numbers, $f(x)$ is nonnegative for every $x$.

\emph{Step 1: the representation.} The function $u$ satisfies condition (ii)$'$ of \bk{Corollary 3.7}, hence also condition (i)$'$: there is a nondecreasing function $F\colon\R\to\R$ with
\begin{equation}\label{eq:4-3-08:rep}
u(t,x)=\int_{-\infty}^{\infty}p(t;x,y)\,\mu(dy);\quad 0<t<T,\ x\in\R,
\end{equation}
where $\mu$ denotes the Lebesgue--Stieltjes measure $dF$ on $(\R,\B(\R))$. This is a nonnegative measure with $\mu([a,b])<\infty$ for all real $a<b$, because $F$ is real-valued. (If $F$ is not right-continuous, $dF$ is understood as the measure induced by its right-continuous modification $y\mapsto F(y+)$; the integrand in \eqref{eq:4-3-08:rep} being continuous in $y$, the values of $F$ at its jumps do not affect the integral.) We fix once and for all a number $t_0\in(0,T)$ and record that
\begin{equation}\label{eq:4-3-08:fin}
\int_{-\infty}^{\infty}p(t_0;0,y)\,\mu(dy)=u(t_0,0)<\infty .
\end{equation}

\emph{Step 2: the convergence in \eqref{eq:4-3-08:lim} is uniform on compact sets, and $f$ is continuous.} Let $K>0$ and $\eps>0$. By \eqref{eq:4-3-08:lim}, every $x\in[-K,K]$ has a $\delta_x>0$ such that
\begin{equation}\label{eq:4-3-08:loc}
\abs{u(t,y)-f(x)}\le\frac{\eps}{2}\quad\text{whenever } 0<t<\delta_x,\ \abs{y-x}<\delta_x .
\end{equation}
For $\abs{y-x}<\delta_x$ we let $t\downarrow0$ in \eqref{eq:4-3-08:loc}; since $u(t,y)\to f(y)$ by \eqref{eq:4-3-08:lim} applied at the point $y$, this gives $\abs{f(y)-f(x)}\le\eps/2$. In particular $f$ is continuous at $x$ (so the continuity of $f$ assumed in (ii) is in fact a consequence of \eqref{eq:4-3-08:lim}), and together with \eqref{eq:4-3-08:loc},
\[
\abs{u(t,y)-f(y)}\le\eps\quad\text{whenever } 0<t<\delta_x,\ \abs{y-x}<\delta_x .
\]
The intervals $(x-\delta_x,x+\delta_x)$, $x\in[-K,K]$, cover the compact set $[-K,K]$; we select finitely many of them, centred at $x_1,\dots,x_m$, which still cover it, and set $\delta\defeq\min_{1\le i\le m}\delta_{x_i}>0$. Then
\begin{equation}\label{eq:4-3-08:unif}
\sup_{\abs{y}\le K}\abs{u(t,y)-f(y)}\le\eps;\quad 0<t<\delta .
\end{equation}

\emph{Step 3: testing against a continuous function of compact support.} Let $\varphi\colon\R\to[0,\infty)$ be continuous and vanish outside $[-K,K]$, and put $\norm{\varphi}\defeq\max_{x}\varphi(x)$. From \eqref{eq:4-3-08:unif} and the boundedness of the continuous function $f$ on $[-K,K]$,
\begin{equation}\label{eq:4-3-08:left}
\lim_{t\downarrow0}\int_{-\infty}^{\infty}\varphi(x)u(t,x)\,dx=\int_{-\infty}^{\infty}\varphi(x)f(x)\,dx .
\end{equation}
On the other hand, \eqref{eq:4-3-08:rep}, Tonelli's theorem (all integrands are nonnegative and jointly measurable) and the symmetry $p(t;x,y)=p(t;y,x)$ give, for $0<t<T$,
\begin{equation}\label{eq:4-3-08:ton}
\int_{-\infty}^{\infty}\varphi(x)u(t,x)\,dx=\int_{-\infty}^{\infty}\Phi_t(y)\,\mu(dy),
\qquad
\Phi_t(y)\defeq\int_{-\infty}^{\infty}\varphi(x)p(t;y,x)\,dx=E^0\varphi(y+W_t).
\end{equation}
We let $t\downarrow0$ on the right-hand side by dominated convergence.

\emph{Pointwise limit.} Since $\varphi$ is bounded and continuous and $W_t\to W_0=0$, $P^0$-a.s., the bounded convergence theorem gives $\lim_{t\downarrow0}\Phi_t(y)=\varphi(y)$ for every $y\in\R$.

\emph{Domination.} We claim that
\begin{equation}\label{eq:4-3-08:dom}
0\le\Phi_t(y)\le C\,p(t_0;0,y);\quad y\in\R,\ 0<t\le\frac{t_0}{4},
\end{equation}
with $C\defeq\norm{\varphi}\max\bigl\{1/p(t_0;0,2K),\,2\sqrt{2\pi t_0}\bigr\}$. If $\abs{y}\le2K$, then $\Phi_t(y)\le\norm{\varphi}$ and $p(t_0;0,y)\ge p(t_0;0,2K)$, so \eqref{eq:4-3-08:dom} holds. If $\abs{y}>2K$, then $\varphi(y+W_t)\ne0$ forces $\abs{y+W_t}\le K$, hence $\abs{W_t}\ge\abs{y}-K\ge\abs{y}/2$. For a standard normal random variable $Z$ and $a\ge0$, Chebyshev's inequality yields $P[Z\ge a]\le e^{-a^2}E e^{aZ}=e^{-a^2/2}$, so that $P^0[\abs{W_t}\ge b]\le 2e^{-b^2/2t}$ for $b\ge0$. Therefore, for $\abs{y}>2K$ and $0<t\le t_0/4$,
\[
\Phi_t(y)\le\norm{\varphi}\,P^0\Bigl[\abs{W_t}\ge\frac{\abs{y}}{2}\Bigr]
\le2\norm{\varphi}\,e^{-y^2/8t}\le2\norm{\varphi}\,e^{-y^2/2t_0}
=2\norm{\varphi}\sqrt{2\pi t_0}\;p(t_0;0,y),
\]
which is \eqref{eq:4-3-08:dom}.

By \eqref{eq:4-3-08:fin} the dominating function in \eqref{eq:4-3-08:dom} is $\mu$-integrable, and dominated convergence in \eqref{eq:4-3-08:ton} gives
\[
\lim_{t\downarrow0}\int_{-\infty}^{\infty}\varphi(x)u(t,x)\,dx=\int_{-\infty}^{\infty}\varphi(y)\,\mu(dy).
\]
Comparing with \eqref{eq:4-3-08:left} we conclude that
\begin{equation}\label{eq:4-3-08:test}
\int_{-\infty}^{\infty}\varphi(y)\,\mu(dy)=\int_{-\infty}^{\infty}\varphi(y)f(y)\,dy
\end{equation}
for every nonnegative continuous function $\varphi$ with compact support.

\emph{Step 4: identification of $\mu$.} Let $\nu(dy)\defeq f(y)\,dy$; because $f$ is nonnegative and continuous, $\nu$ is a measure on $\B(\R)$ which is finite on compact sets. Given real numbers $a<b$, the continuous functions $\varphi_n(y)\defeq\min\{1,\,n(y-a)^+,\,n(b-y)^+\}$ have compact support $[a,b]$ and increase to $\ind_{(a,b)}$ as $n\to\infty$; by \eqref{eq:4-3-08:test} and monotone convergence, $\mu((a,b))=\nu((a,b))<\infty$. The bounded open intervals, together with the empty set, form a $\pi$-system which generates $\B(\R)$ and contains the sets $(-n,n)\uparrow\R$, on which both measures are finite. By Dynkin's $\pi$-$\lambda$ theorem, $\mu=\nu$ on $\B(\R)$.

\textbf{(ii)} Substituting $\mu(dy)=f(y)\,dy$ in \eqref{eq:4-3-08:rep}, we obtain
\[
u(t,x)=\int_{-\infty}^{\infty}p(t;x,y)f(y)\,dy;\quad 0<t<T,\ x\in\R .
\]

\textbf{(i)} Here $f\equiv0$, which is continuous, and (ii) gives $u(t,x)=0$ for all $(t,x)\in(0,T)\times\R$.

\rmk{\emph{On the hint.} Fix $0<t<T$, $x\in\R$ and a Brownian family as in \bk{Corollary 3.7}. By condition (iii)$'$ of that corollary, $M_s\defeq u(t-s,W_s)$, $0\le s<t$, is a nonnegative martingale under $P^x$, and by \eqref{eq:4-3-08:lim} and the continuity of the Brownian paths, $\lim_{s\uparrow t}M_s=f(W_t)$ holds $P^x$-a.s. Granted the uniform integrability of $\{M_s\}_{0\le s<t}$, one gets at once $u(t,x)=E^xM_0=\lim_{s\uparrow t}E^xM_s=E^xf(W_t)$, which is the assertion of (ii) (and of (i) when $f\equiv0$). Conversely, the result proved above yields this uniform integrability: we have $E^xf(W_t)=u(t,x)=E^xM_s<\infty$ for every $s\in[0,t)$, so Scheff\'e's lemma applied to the nonnegative variables $M_s\to f(W_t)$ (along any sequence $s_n\uparrow t$) shows $M_s\to f(W_t)$ in $L^1(P^x)$; letting $r\uparrow t$ in $M_s=E^x[M_r\,|\,\F_s]$ gives $M_s=E^x[f(W_t)\,|\,\F_s]$, and a family of conditional expectations of one integrable random variable is uniformly integrable. We also note that, unlike \bk{Theorem 3.3}, no growth condition of the type \bk{(3.8)} is imposed here; nonnegativity takes its place.}
\end{kex}

\begin{kex}{Problem}{3.11}{}
Throughout, $\{W_t,\F_t;\,0\le t<\infty\}$, $(\Omega,\F)$, $\{P^x\}_{x\in\R}$ is the Brownian family of the text and $Z_t=v(t,W_t)$ for $t>0$. We fix $b>0$, $s>0$, write $A(t)\defeq A(t,b)$, and set
\[
D\defeq\{\omega\in\Omega;\ W_t(\omega)\ge A(t)\ \text{for some } t\ge s\}.
\]
The left-hand side of (3.36) is understood in the usual way: it is $g(a)$, where $g\colon\R\to[0,1]$ is any Borel-measurable function with $P^0[D\mid\sigma(W_s)]=g(W_s)$, a.s.\ $P^0$. Such a $g$ is determined only up to a set which is null for the distribution of $W_s$ under $P^0$; this is the reason for the qualification ``Lebesgue-almost every $a$''.

\emph{Step 0: the representation \textup{(3.31)}.}
The text asserts without proof that the hypotheses force \bk{(3.31)}; since everything rests on it, we verify this. The function $v$ is positive, of class $C^{1,2}$ on $(0,\infty)\times\R$ and satisfies \bk{(3.12)} there, which is condition (ii) of \bk{Theorem 3.6} with $T=\infty$; by \bk{Proposition 3.9} there is a nondecreasing $F\colon\R\to\R$ such that \bk{(3.25)} holds. Let $\mu_F$ be the Lebesgue--Stieltjes measure on $\B(\R)$ induced by $F$. For $t>0$, $x<0$ and $y\le0$ we have $yx\ge0$, whence
\[
v(t,x)\ge\int_{(-\infty,0]}\exp\Bigl(yx-\tfrac12y^2t\Bigr)\,\mu_F(dy)\ge\int_{(-\infty,0]}e^{-y^2t/2}\,\mu_F(dy)\ge0 .
\]
Letting $x\to-\infty$ and using \bk{(3.30)}, we see that the last integral vanishes; its integrand is strictly positive, so $\mu_F((-\infty,0])=0$. Consequently $v(t,0)=\int_{(0,\infty)}e^{-y^2t/2}\,\mu_F(dy)$, which increases to $\mu_F((0,\infty))$ as $t\downarrow0$ by monotone convergence, and \bk{(3.27)} gives $\mu_F((0,\infty))=1$. Thus $\mu_F$ is a probability measure carried by $(0,\infty)$; replacing $F$ by $y\mapsto\mu_F((-\infty,y])$, which induces the same measure, we obtain a probability distribution function with $F(0+)=F(0)=0$, $F(\infty)=1$ and
\begin{equation}\label{eq:4-3-11:e}
v(t,x)=\int_{(0,\infty)}\exp\Bigl(yx-\tfrac12y^2t\Bigr)\,\mu_F(dy);\qquad 0<t<\infty,\ x\in\R,
\end{equation}
which is \bk{(3.31)}; the symbol $\int_{0+}^\infty\cdots\,dF(y)$ in \bk{(3.31)} and in (3.37) means integration over $(0,\infty)$ with respect to $\mu_F$.

\emph{Step 1: the event $D$ in terms of $Z$.}
Fix $t>0$. The function $v(t,\cdot)$ is continuous, and it is strictly increasing: if $x<x'$, then $e^{yx}<e^{yx'}$ for every $y>0$, and integrating this strict inequality between finite integrals in \eqref{eq:4-3-11:e} against the probability measure $\mu_F$ on $(0,\infty)$ gives $v(t,x)<v(t,x')$. Therefore $A(t)$ is the unique solution of $v(t,A(t))=b$ (see \bk{(3.32)}). Hence, for $t>0$ and $x\in\R$,
\begin{equation}\label{eq:4-3-11:a}
x\ge A(t)\iff v(t,x)\ge b,\qquad x<A(t)\iff v(t,x)<b,
\end{equation}
and therefore $D=\{Z_t\ge b\ \text{for some } t\ge s\}$. For every $\omega$ the path $t\mapsto Z_t(\omega)$ is continuous on $(0,\infty)$, because $v$ and $t\mapsto W_t(\omega)$ are continuous. On a compact interval $[s,n]$, $n\in\N$, $n>s$, such a path attains its maximum, and this maximum is the supremum over the countable dense set $Q_n\defeq\Q\cap[s,n]$. Consequently
\[
D=\bigcup_{n>s}\Bigl\{\max_{s\le t\le n}Z_t\ge b\Bigr\}=\bigcup_{n>s}\Bigl\{\sup_{t\in Q_n}Z_t\ge b\Bigr\}\in\F .
\]
Let $D'\defeq\{\sup_{t>s}Z_t\ge b\}=\{\sup_{t\in\Q,\,t>s}Z_t\ge b\}\in\F$ and $\Omega_0\defeq\{\lim_{t\to\infty}Z_t=0\}$; by \bk{Lemma 3.10}, $P^0(\Omega_0)=1$. By continuity of the paths, $\sup_{t>s}Z_t=\sup_{t\ge s}Z_t$, so that $D\subseteq D'$. Conversely, if $\omega\in D'\cap\Omega_0$, there is an integer $n>s$ with $Z_t(\omega)\le b/2$ for all $t\ge n$. Then $\sup_{t\ge n}Z_t(\omega)\le b/2<b\le\sup_{t\ge s}Z_t(\omega)$, and since $\sup_{t\ge s}Z_t(\omega)$ is the larger of $\max_{s\le t\le n}Z_t(\omega)$ and $\sup_{t\ge n}Z_t(\omega)$, it equals the former. Hence $\max_{s\le t\le n}Z_t(\omega)\ge b$ and $\omega\in D$. Thus
\begin{equation}\label{eq:4-3-11:b}
D\subseteq D',\qquad D'\setminus D\subseteq\Omega\setminus\Omega_0,\qquad P^0(D'\setminus D)=0 .
\end{equation}

\emph{Step 2: conditioning on $\F_s$.}
Consider the process $Y_u\defeq Z_{s+u}=v(s+u,W_{s+u})$ with the filtration $\G_u\defeq\F_{s+u}$, $0\le u<\infty$. Every path of $Y$ is continuous on $[0,\infty)$ and positive; $\{Y_u,\G_u;\,0\le u<\infty\}$ is a martingale on $(\Omega,\F,P^0)$, because $\{Z_t,\F_t;\,0<t<\infty\}$ is one (as shown in the text preceding \bk{Lemma 3.10}); and $Y_\infty=\lim_{u\to\infty}Y_u=0$, a.s.\ $P^0$, by \bk{Lemma 3.10}. These are precisely the hypotheses of \bk{Problem 1.3.28}. (Working with $Y$ rather than with the extended process $Z$ of \bk{Lemma 3.10} avoids the time-origin, where the continuity of $Z$ holds only almost surely.) Part (i) of that problem, at time $0$ for $Y$, yields $P^0[\sup_{u>0}Y_u\ge b\mid\G_0]=Y_0/b$ a.s.\ $P^0$ on $\{Y_0<b\}$, that is, $P^0[D'\mid\F_s]=Z_s/b$, a.s.\ $P^0$ on $\{Z_s<b\}$. In view of \eqref{eq:4-3-11:b}, $\ind_D=\ind_{D'}$ a.s.\ $P^0$, and so
\begin{equation}\label{eq:4-3-11:c}
\ind_{\{Z_s<b\}}\,P^0[D\mid\F_s]=\ind_{\{Z_s<b\}}\,\frac{Z_s}{b},\qquad\text{a.s. } P^0 .
\end{equation}

\emph{Step 3: proof of \textup{(3.36)}.}
By \eqref{eq:4-3-11:a}, $\{Z_s<b\}=\{W_s<A(s)\}\in\sigma(W_s)\subseteq\F_s$, and $Z_s=v(s,W_s)$ is $\sigma(W_s)$-measurable. Let $g$ be as in the opening paragraph. Using the tower property and then \eqref{eq:4-3-11:c}, whose right-hand side is $\sigma(W_s)$-measurable and bounded by $1$, we obtain, a.s.\ $P^0$,
\[
\ind_{\{W_s<A(s)\}}\,g(W_s)
=E^0\bigl[\ind_{\{Z_s<b\}}P^0[D\mid\F_s]\bigm|\sigma(W_s)\bigr]
=\ind_{\{W_s<A(s)\}}\,\frac{v(s,W_s)}{b}.
\]
Hence the Borel set $N\defeq\{a\in\R;\ a<A(s),\ g(a)\neq v(s,a)/b\}$ satisfies $P^0[W_s\in N]=0$. Under $P^0$ the random variable $W_s$ has the density $x\mapsto(2\pi s)^{-1/2}e^{-x^2/2s}$, which is strictly positive on $\R$ because $s>0$; therefore $N$ has Lebesgue measure zero. Since $\{a;\ v(s,a)<b\}=(-\infty,A(s))$ by \eqref{eq:4-3-11:a}, this is the assertion (3.36): $g(a)=v(s,a)/b$ for Lebesgue-almost every $a$ with $v(s,a)<b$.

\emph{Step 4: proof of \textup{(3.37)}.}
On $\{Z_s\ge b\}$ we have $W_s\ge A(s)$, so $\{Z_s\ge b\}\subseteq D$. Taking expectations in \eqref{eq:4-3-11:c} and using \eqref{eq:4-3-11:a} once more, we get
\begin{equation}\label{eq:4-3-11:d}
P^0(D)=P^0[Z_s\ge b]+P^0\bigl(D\cap\{Z_s<b\}\bigr)
=P^0[W_s\ge A(s)]+\frac1b\,E^0\bigl[v(s,W_s)\ind_{\{W_s<A(s)\}}\bigr]
\end{equation}
(which is also what \bk{Problem 1.3.28}(ii) asserts for $Y$). Since $W_s/\sqrt s$ is standard normal under $P^0$, the first term is
\[
P^0[W_s\ge A(s)]=1-\Phi\biggl(\frac{A(s)}{\sqrt s}\biggr).
\]
For the second term we insert the representation \eqref{eq:4-3-11:e}. The integrand below is nonnegative and continuous in $(x,y)$, so Tonelli's theorem permits the interchange of the integrations:
\begin{align*}
E^0\bigl[v(s,W_s)\ind_{\{W_s<A(s)\}}\bigr]
&=\int_{-\infty}^{A(s)}\int_{0+}^{\infty}\exp\Bigl(yx-\tfrac12y^2s\Bigr)\,dF(y)\;\frac{1}{\sqrt{2\pi s}}\,e^{-x^2/2s}\,dx\\
&=\int_{0+}^{\infty}\int_{-\infty}^{A(s)}\frac{1}{\sqrt{2\pi s}}\exp\biggl(-\frac{(x-ys)^2}{2s}\biggr)\,dx\,dF(y),
\end{align*}
where we have used $yx-\tfrac12y^2s-\tfrac{x^2}{2s}=-\tfrac{(x-ys)^2}{2s}$. The substitution $z=(x-ys)/\sqrt s$ turns the inner integral into
\[
\frac{1}{\sqrt{2\pi}}\int_{-\infty}^{(A(s)-ys)/\sqrt s}e^{-z^2/2}\,dz=\Phi\biggl(\frac{A(s)}{\sqrt s}-y\sqrt s\biggr).
\]
Substituting these two expressions into \eqref{eq:4-3-11:d} gives
\[
P^0[W_t\ge A(t,b),\text{ for some } t\ge s]
=1-\Phi\biggl(\frac{A(s,b)}{\sqrt s}\biggr)+\frac1b\int_{0+}^{\infty}\Phi\biggl(\frac{A(s,b)}{\sqrt s}-y\sqrt s\biggr)\,dF(y),
\]
which is (3.37).

\rmk{The qualification ``Lebesgue-almost every $a$ with $v(s,a)<b$'' pertains to (3.36) only; (3.37) contains no $a$ and holds for every $s>0$ and $b>0$. As a check, take $F=\ind_{[\mu,\infty)}$ with $\mu>0$, so that $v(t,x)=\exp(\mu x-\mu^2t/2)$ and $A(t,b)=\beta t+\gamma$ with $\beta=\mu/2$, $\gamma=\mu^{-1}\log b$, $b=e^{2\beta\gamma}$. Then (3.36) reads $e^{\mu a-\mu^2s/2}/b=e^{-2\beta(\gamma-a+\beta s)}$, and (3.37) becomes $1-\Phi(\gamma/\sqrt s+\beta\sqrt s)+e^{-2\beta\gamma}\Phi(\gamma/\sqrt s-\beta\sqrt s)$, because $A(s,b)/\sqrt s-\mu\sqrt s=\gamma/\sqrt s-\beta\sqrt s$; these are \bk{(3.38)} and \bk{(3.39)}. The first of them agrees with the formula $P[B_u-\beta u=c\text{ for some }u\ge0]=e^{-2\beta c}$, $c>0$, for a standard Brownian motion $B$ (this is \bk{(3.5.13)} with drift $-\beta$ and level $c$), applied to $B_u=W_{s+u}-W_s$, which is independent of $W_s$, and $c=\gamma+\beta s-a$.}
\end{kex}

\begin{kex}{Exercise}{3.13}{}
Fix $\beta>0$, $\gamma>0$, and put
\[
  A \defeq \{W_t\ge \beta t+\gamma \text{ for some } t\ge 0\},\qquad
  A_s \defeq \{W_t\ge \beta t+\gamma \text{ for some } t\ge s\},\quad s>0 .
\]
Since the paths of $W$ are continuous, the supremum of $W_t-\beta t$ over a compact interval is
attained and equals the supremum over the rationals in that interval, so that
\[
  A_s=\bigcup_{n\ge s}\Bigl\{\sup_{t\in[s,n]\cap\Q}(W_t-\beta t)\ge\gamma\Bigr\},
  \qquad
  A=\bigcup_{n\ge 1}\Bigl\{\sup_{t\in[0,n]\cap\Q}(W_t-\beta t)\ge\gamma\Bigr\}
\]
are events. We give two derivations: the first from the formula \bk{(3.39)} of the present
section, the second from \bk{Section 3.5.C}, to which the statement alludes.

\emph{First proof.}
The events $A_s$ increase as $s\downarrow 0$, and their union is
$\{W_t\ge\beta t+\gamma \text{ for some } t>0\}$. This set differs from $A$ only by a subset of
$\{W_0\ge\gamma\}$, which is $P^0$-null because $P^0[W_0=0]=1$ and $\gamma>0$. Hence, by
continuity from below of $P^0$ along any sequence $s_k\downarrow 0$, and by \bk{(3.39)} (valid
for every $s>0$, $\beta>0$, $\gamma\in\R$),
\[
  P^0(A)=\lim_{s\downarrow0}P^0(A_s)
  =\lim_{s\downarrow0}\Bigl[\,1-\Phi\Bigl(\frac{\gamma}{\sqrt s}+\beta\sqrt s\Bigr)
     +e^{-2\beta\gamma}\,\Phi\Bigl(\frac{\gamma}{\sqrt s}-\beta\sqrt s\Bigr)\Bigr].
\]
Because $\gamma>0$, both arguments $\gamma/\sqrt s\pm\beta\sqrt s$ tend to $+\infty$ as
$s\downarrow0$, and $\Phi(x)\to1$ as $x\to\infty$; the limit is therefore
$1-1+e^{-2\beta\gamma}=e^{-2\beta\gamma}$.

\emph{Second proof.}
Let $X_t\defeq W_t-\beta t$, $t\ge0$, and for a continuous function
$x\colon[0,\infty)\to\R$ let $T_\gamma(x)\defeq\inf\{t\ge0;\ x(t)=\gamma\}$ (with
$\inf\emptyset=\infty$). The set
\[
  G\defeq\{x\in C[0,\infty);\ T_\gamma(x)<\infty\}
\]
is a Borel subset of $C[0,\infty)$: for $x(0)<\gamma$ the intermediate value theorem shows that
$x$ takes the value $\gamma$ if and only if $\max_{0\le t\le n}x(t)\ge\gamma$ for some $n$, so
that $G\cap\{x(0)<\gamma\}=\{x(0)<\gamma\}\cap\bigcup_{n}\{\sup_{t\in[0,n]\cap\Q}x(t)\ge\gamma\}$,
while $G\cap\{x(0)>\gamma\}$ is described symmetrically with infima and
$G\supseteq\{x(0)=\gamma\}$. By the same intermediate value argument, on $\{W_0=0\}$ we have
\begin{equation}\label{eq:4-3-13:a}
  A=\{X_t\ge\gamma \text{ for some } t\ge0\}=\{X_\cdot\in G\}.
\end{equation}

We now recall the setting of \bk{Section 3.5.C}. There $W$ is a standard Brownian motion under
a measure $P$, and for $\mu\ne0$ the measure $P^{(\mu)}$ on $\F^W_\infty$ is such that
$\wt W_t=W_t-\mu t$ is a standard Brownian motion under $P^{(\mu)}$. Consequently the law of
$W=\wt W+\mu\,\cdot$ under $P^{(\mu)}$, as a measure on
$\bigl(C[0,\infty),\B(C[0,\infty))\bigr)$, is the image of Wiener measure under the continuous
map $x\mapsto x+\mu\,\cdot\,$. Take $\mu=-\beta$. Under $P^0$ the process $W$ is a standard
Brownian motion, so the law of $X=W-\beta\,\cdot$ under $P^0$ is the image of Wiener measure
under the same map. Thus $X$ under $P^0$ and $W$ under $P^{(-\beta)}$ have the same law on
$C[0,\infty)$, and \eqref{eq:4-3-13:a} together with \bk{(3.5.13)}, applied with $\mu=-\beta$
and $b=\gamma$, gives
\[
  P^0(A)=P^0[X_\cdot\in G]=P^{(-\beta)}[W_\cdot\in G]=P^{(-\beta)}[T_\gamma<\infty]
  =\exp\bigl[-\beta\gamma-\abs{\beta\gamma}\bigr]=e^{-2\beta\gamma}.
\]

\rmk{The link with \bk{Exercise 3.5.9} is the following. Since $-W$ is again a standard
Brownian motion under $P^0$, the second proof shows equally that, for every $\gamma>0$,
\[
  e^{-2\beta\gamma}=P^0[W_t+\beta t\le-\gamma \text{ for some } t\ge0]
  =P^{(\beta)}[W_t\le-\gamma \text{ for some } t>0],
\]
the process $W_t+\beta t$ under $P^0$ having the law of $W$ under $P^{(\beta)}$, and the time
$t=0$ being immaterial because $W_0=0>-\gamma$ almost surely. With
$W_*=\inf_{t>0}W_t$ we have the inclusions
$\{-W_*>\gamma\}\subseteq\{W_t\le-\gamma \text{ for some } t>0\}\subseteq\{-W_*\ge\gamma\}$,
whence $P^{(\beta)}[-W_*>\gamma]\le e^{-2\beta\gamma}$ and
$P^{(\beta)}[-W_*>\gamma]\ge P^{(\beta)}[-W_*\ge\gamma']\ge e^{-2\beta\gamma'}$ for every
$\gamma'>\gamma$. Letting $\gamma'\downarrow\gamma$ we obtain
$P^{(\beta)}[-W_*>\gamma]=e^{-2\beta\gamma}$, $\gamma>0$: under $P^{(\beta)}$ the random
variable $-W_*$ is exponentially distributed with parameter $2\beta$, which is the assertion of
that exercise. As a check, the formula gives probability $1$ in the limits $\beta\downarrow0$
(Brownian motion without drift reaches every level) and $\gamma\downarrow0$.}
\end{kex}

\begin{kex}{Exercise}{3.14}{}
We proceed as in the treatment \bk{(3.43)}, \bk{(3.44)} of the absorbed Brownian motion, but extend $f$ to an \emph{even} function instead of an odd one. Throughout, $S \defeq (0,1/2a)\times\R$ and $p(t;x,y) = (2\pi t)^{-1/2}e^{-(x-y)^2/2t}$.

\emph{Step 1. Reduction to \bk{(3.4)}.} Define $g\colon\R\to\R$ by $g(y) \defeq f(\abs{y})$ for $y\neq 0$ and $g(0)\defeq 0$. Then $g$ is Borel-measurable, and
\begin{equation*}
\int_{-\infty}^{\infty} e^{-ay^2}\abs{g(y)}\,dy = 2\int_0^\infty e^{-ay^2}\abs{f(y)}\,dy < \infty
\end{equation*}
by \bk{(3.42)}, so $g$ satisfies \bk{(3.3)} with the same constant $a$. According to \bk{(3.4)} and \bk{Problem 3.1}, the function
\begin{equation}\label{eq:4-3-14:v}
v(t,x) \defeq E^x g(W_t) = \int_{-\infty}^{\infty} g(y)\,p(t;x,y)\,dy;\quad (t,x)\in S,
\end{equation}
is well defined (the integral converges absolutely; this is also contained in \eqref{eq:4-3-14:dom} below), possesses the partial derivatives of all orders on $S$, these are given by \bk{(3.5)}, and $v$ satisfies the heat equation \bk{(3.1)} on $S$. Under $P^x$ the random variable $W_t$, $t>0$, has the density $p(t;x,\cdot)$; in particular $P^x[W_t = 0]=0$, hence $f(\abs{W_t}) = g(W_t)$ holds $P^x$-a.s., the expectation defining $u$ exists, and
\begin{equation}\label{eq:4-3-14:uv}
u(t,x) = v(t,x);\quad 0<t<\frac{1}{2a},\ x>0.
\end{equation}

\emph{Step 2. A domination, and the continuity of the derivatives of $v$.} \bk{Problem 3.1} asserts the existence of the derivatives of $v$; for the class $C^{1,2}$ and for the Neumann condition we need their continuity on $S$, which we now verify. Let $q$ denote any one of the four functions
\begin{equation*}
p,\qquad \frac{\partial p}{\partial x} = \frac{y-x}{t}\,p,\qquad
\frac{\partial^2 p}{\partial x^2} = \Bigl(\frac{(y-x)^2}{t^2}-\frac{1}{t}\Bigr)p,\qquad
\frac{\partial p}{\partial t} = \frac{1}{2}\,\frac{\partial^2 p}{\partial x^2},
\end{equation*}
the last identity being \bk{(3.2)}. Fix $0<t_0<t_1<1/2a$ and $R>0$. Since $1/(2t_1) > a$, we can choose $\delta\in(0,1)$ so that
\begin{equation*}
c \defeq \frac{1}{2}\Bigl(\frac{1-\delta}{2t_1} - a\Bigr) > 0 .
\end{equation*}
For $\abs{x}\le R$ and $y\in\R$ the inequality $2xy \le \delta y^2 + x^2/\delta$ gives
$(x-y)^2 \ge (1-\delta)y^2 - (\delta^{-1}-1)R^2$, and therefore, for $t_0\le t\le t_1$,
\begin{equation*}
p(t;x,y) \le \frac{1}{\sqrt{2\pi t_0}}\exp\Bigl\{-\frac{(x-y)^2}{2t_1}\Bigr\}
\le \frac{1}{\sqrt{2\pi t_0}}\exp\Bigl\{\frac{(\delta^{-1}-1)R^2}{2t_1}\Bigr\}\,e^{-(a+2c)y^2}.
\end{equation*}
The factors multiplying $p$ in the formulas for $q$ are bounded in absolute value by $1 + (\abs{y}+R)/t_0 + (\abs{y}+R)^2/t_0^2 + 1/t_0$, and the product of this polynomial in $\abs{y}$ with $e^{-2cy^2}$ is bounded on $\R$. Consequently there is a constant $C$, depending only on $t_0,t_1,R,a$, with
\begin{equation}\label{eq:4-3-14:dom}
\abs{q(t;x,y)} \le C e^{-ay^2};\quad t_0\le t\le t_1,\ \abs{x}\le R,\ y\in\R .
\end{equation}
Now for each fixed $y$ the function $(t,x)\mapsto q(t;x,y)$ is continuous on $S$, and by \eqref{eq:4-3-14:dom} the integrands $g(y)q(t;x,y)$ are dominated, uniformly in $(t,x)\in[t_0,t_1]\times[-R,R]$, by the integrable function $C e^{-ay^2}\abs{g(y)}$. The dominated convergence theorem shows that $(t,x)\mapsto\int_{-\infty}^\infty g(y)q(t;x,y)\,dy$ is continuous on $[t_0,t_1]\times[-R,R]$, hence on $S$, because $t_0$, $t_1$, $R$ are arbitrary. In view of \bk{(3.5)}, these four integrals are $v$, $\partial v/\partial x$, $\partial^2 v/\partial x^2$ and $\partial v/\partial t$, respectively. Thus $v$ is of class $C^{1,2}$ on $S$.

\emph{Step 3. $u$ is of class $C^{1,2}$ and solves \bk{(3.1)}.} The set $(0,1/2a)\times(0,\infty)$ is an open subset of $S$ on which $u=v$ by \eqref{eq:4-3-14:uv}. Hence the partial derivatives of $u$ exist there and coincide with those of $v$; by Step~2 the function $u$ is of class $C^{1,2}$ on $(0,1/2a)\times(0,\infty)$, and by Step~1 it satisfies $\partial u/\partial t = \tfrac12\,\partial^2u/\partial x^2$ there.

\emph{Step 4. The initial condition \bk{(3.6)}.} Let $x>0$ be a point of continuity of $f$. Since $g=f$ on $(0,\infty)$, a neighbourhood of $x$, the function $g$ is continuous at $x$, and \bk{Problem 3.2}, applied to $g$, yields
$\lim_{t\downarrow 0,\,y\to x} v(t,y) = g(x) = f(x)$. For $y$ sufficiently close to $x$ we have $y>0$ and thus $v(t,y)=u(t,y)$ by \eqref{eq:4-3-14:uv}; therefore
\begin{equation*}
\lim_{\substack{t\downarrow 0\\ y\to x}} u(t,y) = f(x).
\end{equation*}

\emph{Step 5. The Neumann condition.} For $(t,x)\in S$ the substitution $y\mapsto -y$ in \eqref{eq:4-3-14:v}, together with $g(-y)=g(y)$ and $p(t;-x,-y)=p(t;x,y)$, gives
\begin{equation*}
v(t,-x) = \int_{-\infty}^{\infty} g(y)\,p(t;-x,y)\,dy = \int_{-\infty}^{\infty} g(-y)\,p(t;-x,-y)\,dy = v(t,x).
\end{equation*}
Thus, for each $t\in(0,1/2a)$, the function $x\mapsto v(t,x)$ is even and differentiable on $\R$; differentiating the identity $v(t,-x)=v(t,x)$ at $x=0$ we obtain $-\frac{\partial v}{\partial x}(t,0) = \frac{\partial v}{\partial x}(t,0)$, that is,
\begin{equation*}
\frac{\partial v}{\partial x}(t,0) = 0;\quad 0<t<\frac{1}{2a}.
\end{equation*}
Finally, $\partial v/\partial x$ is continuous on $S$ by Step~2, and it agrees with $\partial u/\partial x$ on $(0,1/2a)\times(0,\infty)$. Hence, for every $t\in(0,1/2a)$,
\begin{equation*}
\lim_{\substack{s \to t \\ x \downarrow 0}} \frac{\partial}{\partial x}\,u(s,x)
= \lim_{\substack{s \to t \\ x \downarrow 0}} \frac{\partial v}{\partial x}(s,x)
= \frac{\partial v}{\partial x}(t,0) = 0 .
\end{equation*}

\rmk{The function $u$ is the temperature in a semi-infinite rod with initial temperature $f$ whose end $x=0$ is insulated; the underlying process $\abs{W}$ is Brownian motion reflected at the origin, whose transition density on $(0,\infty)$ is $p_+(t;x,y) \defeq p(t;x,y)+p(t;x,-y)$, so that $u(t,x)=\int_0^\infty f(y)\,p_+(t;x,y)\,dy$, in complete analogy with \bk{(3.44)}.}
\end{kex}

\begin{kex}{Exercise}{3.15}{}
Throughout, $D\defeq(0,\infty)\times(0,b)$, $\norm{\varphi}\defeq\sup\abs{\varphi}$ for a bounded function $\varphi$, $p(t;x,y)=(2\pi t)^{-1/2}e^{-(x-y)^2/2t}$, and $\{W_t,\F_t\}$, $\{P^x\}_{x\in\R}$ is a one-dimensional Brownian family; $T_y$ is the passage time to $y$. The boundary conditions are understood as in the statement: the function equal to $u$ on $D$, to $f$ on $\{0\}\times(0,b)$, to $g$ on $(0,\infty)\times\{0\}$ and to $k$ on $(0,\infty)\times\{b\}$ is continuous on the union $\ol{D}{}'$ of these four sets. This holds if and only if
\begin{equation}\label{eq:4-3-15:bc}
\lim_{\substack{(s,y)\to(0,x_0)\\ (s,y)\in D}}u(s,y)=f(x_0),\qquad
\lim_{\substack{(s,y)\to(t_0,0)\\ (s,y)\in D}}u(s,y)=g(t_0),\qquad
\lim_{\substack{(s,y)\to(t_0,b)\\ (s,y)\in D}}u(s,y)=k(t_0)
\end{equation}
for all $0<x_0<b$, $t_0>0$. Indeed, necessity is obvious; conversely, if \eqref{eq:4-3-15:bc} holds and $\varphi$ denotes the boundary function, then for a boundary point $\pi$ and $\eps>0$ there is an open neighbourhood $U$ of $\pi$ with $\abs{u-\varphi(\pi)}\le\eps$ on $U\cap D$, and every boundary point $\pi'\in U$ satisfies $\abs{\varphi(\pi')-\varphi(\pi)}\le\eps$ because $\varphi(\pi')$ is the limit of $u$ along points of $U\cap D$.

We show in Part~I that \bk{(3.46)} defines a bounded solution, and in Part~II that every bounded solution is given by \bk{(3.46)}.

\textbf{Part I: the function \bk{(3.46)} is a bounded solution.}
The three events appearing in \bk{(3.46)} are disjoint, so the random variable under the expectation is bounded by $K\defeq\max\{\norm f,\norm g,\norm k\}$ (on $\{t<T_0\wedge T_b\}$ we have $W_t\in(0,b)$, $P^x$-a.s.\ for $0<x<b$, by continuity of the paths; on the second event $0<t-T_0$, on the third $0<t-T_b$). Hence $u$ is well defined, $\abs u\le K$, and $u=u_f+v_g+w_k$ on $D$, where
\begin{gather*}
u_f(t,x)\defeq E^x\bigl[f(W_t)\ind_{\{t<T_0\wedge T_b\}}\bigr],\qquad
v_g(t,x)\defeq E^x\bigl[g(t-T_0)\ind_{\{T_0<t\wedge T_b\}}\bigr],\\
w_k(t,x)\defeq E^x\bigl[k(t-T_b)\ind_{\{T_b<t\wedge T_0\}}\bigr].
\end{gather*}
We shall use two elementary facts about passage times. For $0<x<b$, under $P^x$ the process $W-x$ is a standard Brownian motion and so is $x-W$; therefore, by \bk{(2.6.2)},
\begin{equation}\label{eq:4-3-15:pass}
P^x[T_0<s]=P^0[T_x<s]=2P^0[W_s>x]=1-P^0[\abs{W_s}\le x];\quad s>0,
\end{equation}
and by \bk{Exercise 2.8.13},
\begin{equation}\label{eq:4-3-15:exit}
P^x[T_0<T_b]=\frac{b-x}{b},\qquad P^x[T_b<T_0]=\frac{x}{b}.
\end{equation}

\emph{Step 1. The term $u_f$.} By \bk{Proposition 2.8.10} (with $a=b$) and the definition \bk{(2.8.9)} of $p_-$,
\begin{equation*}
u_f(t,x)=\int_0^b f(y)\sum_{n=-\infty}^{\infty}\bigl[p(t;x,y+2nb)-p(t;x,-y-2nb)\bigr]\,dy;\quad (t,x)\in D.
\end{equation*}
The intervals $I_n\defeq(2nb,2nb+b)$ and $J_n\defeq(-2nb-b,-2nb)$, $n\in\Z$, are pairwise disjoint and their union is $\R\setminus b\Z$. Consequently
\begin{equation*}
\sum_{n=-\infty}^\infty\int_0^b\abs{f(y)}\bigl[p(t;x,y+2nb)+p(t;x,-y-2nb)\bigr]\,dy
\le\norm f\int_{-\infty}^\infty p(t;x,z)\,dz=\norm f<\infty,
\end{equation*}
and Fubini's theorem allows us to integrate term by term. The substitutions $z=y+2nb$ and $z=-y-2nb$ give
\begin{equation}\label{eq:4-3-15:F}
u_f(t,x)=\int_{-\infty}^\infty F(z)\,p(t;x,z)\,dz=E^xF(W_t),
\end{equation}
where $F\colon\R\to\R$ is defined by $F(z)\defeq f(z-2nb)$ for $z\in I_n$, $F(z)\defeq-f(-z-2nb)$ for $z\in J_n$, and $F\defeq0$ on $b\Z$. The function $F$ is Borel-measurable, $\abs F\le\norm f$, $F=f$ on $(0,b)=I_0$, and
\begin{equation}\label{eq:4-3-15:odd}
F(-z)=-F(z),\qquad F(z+2b)=F(z),\qquad F(b-z)=-F(b+z);\quad z\in\R.
\end{equation}
(If $z\in I_n$ then $-z\in J_n$ and $F(-z)=-f(z-2nb)=-F(z)$; this gives the first identity. If $z\in I_n$ then $z+2b\in I_{n+1}$, and if $z\in J_n$ then $z+2b\in J_{n-1}$; in both cases the definition gives $F(z+2b)=F(z)$. The third identity follows from the first two: $F(b-z)=-F(z-b)=-F(z+b)$.)

We use \eqref{eq:4-3-15:F} to \emph{define} $u_f$ on all of $(0,\infty)\times\R$. Being bounded, $F$ satisfies \bk{(3.3)} for every $a>0$; by \bk{(3.4)} and \bk{Problem 3.1}, $u_f$ has partial derivatives of all orders on $(0,\infty)\times\R$, given by \bk{(3.5)}, and satisfies the heat equation \bk{(3.1)} there. These derivatives are continuous: let $q$ be one of $p$, $\partial p/\partial x=\frac{y-x}{t}p$, $\partial^2p/\partial x^2=\bigl(\frac{(y-x)^2}{t^2}-\frac1t\bigr)p$, $\partial p/\partial t=\frac12\partial^2p/\partial x^2$; for $0<t_0\le t\le t_1$, $\abs x\le R$, the inequality $(x-y)^2\ge\frac12y^2-R^2$ yields
\begin{equation*}
\abs{q(t;x,y)}\le\Bigl(1+\frac{\abs y+R}{t_0}+\frac{(\abs y+R)^2}{t_0^2}+\frac1{t_0}\Bigr)\frac{e^{R^2/2t_1}}{\sqrt{2\pi t_0}}\,e^{-y^2/4t_1},
\end{equation*}
an integrable function of $y$, and since $F$ is bounded, the dominated convergence theorem shows that $(t,x)\mapsto\int F(y)q(t;x,y)\,dy$ is continuous on $[t_0,t_1]\times[-R,R]$. Thus $u_f$ is of class $C^{1,2}$ on $(0,\infty)\times\R$, in particular on $D$, and solves \bk{(3.1)}.

Since $F$ is continuous at every $x_0\in(0,b)$, \bk{Problem 3.2} gives $\lim_{t\downarrow0,\,y\to x_0}u_f(t,y)=F(x_0)=f(x_0)$. Moreover, for $t>0$, the first and third identities in \eqref{eq:4-3-15:odd} and the symmetry of $p(t;0,\cdot)$ give
\begin{equation*}
u_f(t,0)=\int_{-\infty}^\infty F(z)p(t;0,z)\,dz=0,\qquad u_f(t,b)=\int_{-\infty}^\infty F(b+z)p(t;0,z)\,dz=0,
\end{equation*}
the integrands being odd and integrable. As $u_f$ is continuous on $(0,\infty)\times\R$, we conclude
\begin{equation}\label{eq:4-3-15:uf}
\lim_{(s,y)\to(0,x_0)}u_f(s,y)=f(x_0),\qquad\lim_{(s,y)\to(t_0,0)}u_f(s,y)=0,\qquad\lim_{(s,y)\to(t_0,b)}u_f(s,y)=0.
\end{equation}

\emph{Step 2. The kernel $H$.} Let $h(s,z)=(2\pi)^{-1/2}zs^{-3/2}e^{-z^2/2s}$; $s>0$, $z\in\R$, be the function of \bk{(3.10)}; it is infinitely differentiable and solves \bk{(3.1)}. Formula \bk{(2.8.25)} (with $a=b$) reads
\begin{equation}\label{eq:4-3-15:dens}
P^x[T_0\in ds,\ T_0<T_b]=H(s,x)\,ds,\qquad H(s,x)\defeq\sum_{n=-\infty}^\infty h(s,x+2nb);\quad s>0,\ 0<x<b.
\end{equation}
We set $H(s,x)\defeq0$ for $s\le0$, $0<x<b$, and claim:
\begin{equation}\label{eq:4-3-15:H}
H\ \text{is of class }C^{1,2}\text{ on }\R\times(0,b)\quad\text{and}\quad\frac{\partial H}{\partial s}=\frac12\frac{\partial^2H}{\partial x^2}\ \text{ there.}
\end{equation}
Let $\mathscr{D}$ denote any of the operators $\mathrm{id}$, $\partial/\partial z$, $\partial^2/\partial z^2$, $\partial/\partial s$. We compute
\begin{equation*}
\frac{\partial h}{\partial z}=\frac{s^{-3/2}-z^2s^{-5/2}}{\sqrt{2\pi}}\,e^{-z^2/2s},\qquad
\frac{\partial h}{\partial s}=\frac12\frac{\partial^2h}{\partial z^2}=\frac{z^3s^{-7/2}-3zs^{-5/2}}{2\sqrt{2\pi}}\,e^{-z^2/2s}.
\end{equation*}
Fix $T>0$ and $0<\delta<b/2$. Since $s^{-\beta}\le T^{7/2-\beta}s^{-7/2}$ for $0<s\le T$, $\beta\le7/2$, there is a constant $C$ depending only on $T$ with $\abs{\mathscr{D}h(s,z)}\le C(1+\abs z^3)s^{-7/2}e^{-z^2/2s}$ for $0<s\le T$, $z\in\R$. If in addition $\abs z\ge\delta$, then $e^{-z^2/2s}=e^{-z^2/4s}e^{-z^2/4s}\le e^{-\delta^2/4s}e^{-z^2/4T}$, and so
\begin{equation}\label{eq:4-3-15:hbd}
\abs{\mathscr{D}h(s,z)}\le C\,\psi(s)\,(1+\abs z^3)e^{-z^2/4T};\quad 0<s\le T,\ \abs z\ge\delta,\qquad\psi(s)\defeq s^{-7/2}e^{-\delta^2/4s}.
\end{equation}
The function $\psi$ is bounded on $(0,\infty)$ and $\lim_{s\downarrow0}\psi(s)=0$. For $x\in[\delta,b-\delta]$ and $z_n\defeq x+2nb$ we have $\abs{z_n}\ge\delta$ for all $n\in\Z$, and $\abs n b\le(2\abs n-1)b\le\abs{z_n}\le(2\abs n+1)b$ for $n\ne0$. Therefore \eqref{eq:4-3-15:hbd} yields
\begin{equation}\label{eq:4-3-15:Hbd}
\sum_{n=-\infty}^\infty\abs{\mathscr{D}h(s,x+2nb)}\le C\,A\,\psi(s);\quad 0<s\le T,\ \delta\le x\le b-\delta,
\end{equation}
with $A\defeq\sum_{n\in\Z}\bigl(1+(2\abs n+1)^3b^3\bigr)e^{-n^2b^2/4T}<\infty$, and the $n$-th term is bounded on $(0,T]\times[\delta,b-\delta]$ by $C\norm\psi$ times the $n$-th term of the convergent series $A$. By the Weierstrass test the series $\sum_n\mathscr{D}h(s,x+2nb)$ converge uniformly on $(0,T]\times[\delta,b-\delta]$ for each of the four choices of $\mathscr{D}$. The theorem on termwise differentiation of series now shows that $H$ is of class $C^{1,2}$ on $(0,T)\times(\delta,b-\delta)$, with $\mathscr{D}H=\sum_n\mathscr{D}h(\cdot,\cdot+2nb)$; since each term solves \bk{(3.1)}, so does $H$. As $T$, $\delta$ are arbitrary, \eqref{eq:4-3-15:H} holds on $(0,\infty)\times(0,b)$; it holds trivially on $(-\infty,0)\times(0,b)$. It remains to consider the points $(0,x)$. By \eqref{eq:4-3-15:Hbd},
\begin{equation}\label{eq:4-3-15:H0}
\sup_{\delta\le x\le b-\delta}\abs{\mathscr{D}H(s,x)}\le CA\psi(s)\longrightarrow0\quad\text{as }s\downarrow0 .
\end{equation}
With $\mathscr{D}=\mathrm{id}$ this shows that $H$ is continuous on $\R\times(0,b)$. Since $H(0,\cdot)\equiv0$, the derivatives $\partial H/\partial x$ and $\partial^2H/\partial x^2$ exist and vanish at the points $(0,x)$. For $s>0$ the mean value theorem gives $H(s,x)/s=\frac{\partial H}{\partial s}(\sigma,x)$ for some $\sigma\in(0,s)$, which tends to zero as $s\downarrow0$ by \eqref{eq:4-3-15:H0}; the difference quotient from the left is zero; hence $\frac{\partial H}{\partial s}(0,x)=0$. Thus $\partial H/\partial s$, $\partial H/\partial x$, $\partial^2H/\partial x^2$ exist on $\R\times(0,b)$, vanish for $s\le0$, and are continuous by \eqref{eq:4-3-15:H0}; the heat equation holds trivially for $s\le0$. This proves \eqref{eq:4-3-15:H}.

\emph{Step 3. The term $v_g$ is of class $C^{1,2}$ and solves \bk{(3.1)}.} Let $(t,x)\in D$. Under $P^x$ we have $T_0>0$, and $\{T_0<t\wedge T_b\}=\{T_0<t\}\cap\{T_0<T_b\}$. By \eqref{eq:4-3-15:dens}, applied to the bounded Borel function $s\mapsto g(t-s)\ind_{(0,t)}(s)$ ($H(\cdot,x)$ is bounded, by \eqref{eq:4-3-15:Hbd}),
\begin{equation}\label{eq:4-3-15:abel}
v_g(t,x)=\int_0^tg(t-s)H(s,x)\,ds=\int_0^tg(r)H(t-r,x)\,dr=\int_0^Tg(r)H(t-r,x)\,dr
\end{equation}
for every $T>t$, because $H(t-r,x)=0$ when $r\ge t$. Fix $T>0$ and $0<\delta<b/2$. By \eqref{eq:4-3-15:H} the functions $H$, $\partial H/\partial s$, $\partial H/\partial x$, $\partial^2H/\partial x^2$ are continuous, hence bounded by some constant $C'$, on the compact set $[-T,T]\times[\delta,b-\delta]$, which contains all the points $(t-r,x)$ with $0<t<T$, $0<r<T$, $\delta\le x\le b-\delta$. For such $(t,x)$ the integrand in the last integral of \eqref{eq:4-3-15:abel} is continuously differentiable in $t$ and twice continuously differentiable in $x$, with derivatives bounded by $\norm gC'$, which is integrable over $(0,T)$. The theorem on differentiation under the integral sign (applied once in $t$, and twice in $x$) and the dominated convergence theorem show that $v_g$ is of class $C^{1,2}$ on $(0,T)\times(\delta,b-\delta)$, with
\begin{equation*}
\frac{\partial v_g}{\partial t}(t,x)=\int_0^Tg(r)\frac{\partial H}{\partial s}(t-r,x)\,dr
=\frac12\int_0^Tg(r)\frac{\partial^2H}{\partial x^2}(t-r,x)\,dr=\frac12\frac{\partial^2v_g}{\partial x^2}(t,x).
\end{equation*}
Since $T$ and $\delta$ are arbitrary, $v_g$ is of class $C^{1,2}$ on $D$ and solves \bk{(3.1)} there.

\emph{Step 4. Boundary behaviour of $v_g$.} We claim that for $0<x_0<b$, $t_0>0$:
\begin{equation}\label{eq:4-3-15:vg}
\lim_{\substack{(s,x)\to(0,x_0)\\(s,x)\in D}}v_g(s,x)=0,\qquad\lim_{\substack{(s,x)\to(t_0,b)\\(s,x)\in D}}v_g(s,x)=0,\qquad\lim_{\substack{(s,x)\to(t_0,0)\\(s,x)\in D}}v_g(s,x)=g(t_0).
\end{equation}
For the first limit, \eqref{eq:4-3-15:pass} and Chebyshev's inequality give, for $x\ge x_0/2$,
$\abs{v_g(s,x)}\le\norm gP^x[T_0<s]=2\norm gP^0[W_s>x]\le2\norm gP^0[W_s>x_0/2]\le8\norm gs/x_0^2\to0$ as $s\downarrow0$.
For the second, \eqref{eq:4-3-15:exit} gives $\abs{v_g(s,x)}\le\norm gP^x[T_0<T_b]=\norm g(b-x)/b\to0$ as $x\uparrow b$.
For the third, let $\eps>0$ and choose $\eta\in(0,t_0/2)$ such that $\abs{g(r)-g(t_0)}\le\eps$ whenever $\abs{r-t_0}\le2\eta$. Let $(s,x)\in D$ with $\abs{s-t_0}<\eta$, and put $A\defeq\{T_0<\eta\}\cap\{T_0<T_b\}$. Since $s>t_0-\eta>\eta$, we have $A\subseteq\{T_0<s\wedge T_b\}$, and on $A$ we have $\abs{(s-T_0)-t_0}<2\eta$. Writing
\begin{equation*}
v_g(s,x)-g(t_0)=E^x\bigl[(g(s-T_0)-g(t_0))\ind_{\{T_0<s\wedge T_b\}}\bigr]-g(t_0)P^x[T_0\ge s\wedge T_b]
\end{equation*}
and splitting the expectation over $A$ and $\{T_0<s\wedge T_b\}\setminus A$, we obtain from \eqref{eq:4-3-15:pass}, \eqref{eq:4-3-15:exit}
\begin{equation*}
\abs{v_g(s,x)-g(t_0)}\le\eps+3\norm gP^x(A^c)\le\eps+3\norm g\bigl(P^x[T_0\ge\eta]+P^x[T_b\le T_0]\bigr)
\le\eps+3\norm g\Bigl(\frac{2x}{\sqrt{2\pi\eta}}+\frac xb\Bigr);
\end{equation*}
here $P^x[T_0\ge\eta]=P^0[\abs{W_\eta}\le x]\le2x/\sqrt{2\pi\eta}$ and $P^x[T_b\le T_0]=1-P^x[T_0<T_b]=x/b$. Letting $x\downarrow0$ and then $\eps\downarrow0$ gives the third limit in \eqref{eq:4-3-15:vg}.

\emph{Step 5. The term $w_k$, and conclusion of Part I.} Fix $(t,x)\in D$. Under $P^x$ the process $\wt W\defeq b-W$ is a Brownian motion starting at $b-x$, so its law on $C[0,\infty)$ is the law of $W$ under $P^{b-x}$; its passage times to $0$ and $b$ are $T_b$ and $T_0$, respectively. As $k(t-T_b)\ind_{\{T_b<t\wedge T_0\}}$ is the same Borel-measurable functional of the path $\wt W$ as $k(t-T_0)\ind_{\{T_0<t\wedge T_b\}}$ is of $W$, we get
\begin{equation*}
w_k(t,x)=v_k(t,b-x);\quad(t,x)\in D,
\end{equation*}
where $v_k$ is defined like $v_g$ with $k$ in place of $g$ (this is \bk{(2.8.26)} in integrated form). Steps 3 and 4 apply to $v_k$; since $x\mapsto b-x$ maps $(0,b)$ onto itself and the heat equation is invariant under this reflection, $w_k$ is of class $C^{1,2}$ on $D$, solves \bk{(3.1)}, and by \eqref{eq:4-3-15:vg}
\begin{equation}\label{eq:4-3-15:wk}
\lim_{(s,x)\to(0,x_0)}w_k(s,x)=0,\qquad\lim_{(s,x)\to(t_0,0)}w_k(s,x)=0,\qquad\lim_{(s,x)\to(t_0,b)}w_k(s,x)=k(t_0),
\end{equation}
the limits being taken in $D$. Adding up, $u=u_f+v_g+w_k$ is bounded, of class $C^{1,2}$ on $D$, solves the heat equation there, and by \eqref{eq:4-3-15:uf}, \eqref{eq:4-3-15:vg}, \eqref{eq:4-3-15:wk} satisfies \eqref{eq:4-3-15:bc}. Thus \bk{(3.46)} is a bounded solution of the problem.

\textbf{Part II: uniqueness.} Let $v$ be any bounded solution, and fix $(t,x)\in D$. Take an integer $n_0$ with $1/n_0<\min\{x,b-x,t\}$, and for $n\ge n_0$ define
\begin{equation*}
S_n\defeq\inf\{s\ge0;\ W_s\notin(1/n,b-1/n)\},\qquad\tau_n\defeq S_n\wedge(t-1/n).
\end{equation*}
$S_n$ is the hitting time of a closed set by a continuous adapted process, hence a stopping time (\bk{Problem 1.2.7}), and $\tau_n$ is a stopping time bounded by $t$.

\emph{Step 1. $v(t,x)=E^xv(t-\tau_n,W_{\tau_n})$.} Choose $C^\infty$ functions $\chi_1,\chi_2\colon\R\to[0,1]$ with $\chi_1=1$ on $(-\infty,t-\frac1n]$, $\chi_1=0$ on $[t-\frac1{2n},\infty)$, $\chi_2=1$ on $[\frac1n,b-\frac1n]$, $\chi_2=0$ outside $(\frac1{2n},b-\frac1{2n})$, and set $\varphi(s,y)\defeq\chi_1(s)\chi_2(y)v(t-s,y)$ for $s<t$, $0<y<b$, and $\varphi(s,y)\defeq0$ otherwise. Then $\varphi$ is of class $C^{1,2}$ on $\R\times\R$, $\partial\varphi/\partial y$ is bounded on $[0,\infty)\times\R$ (it is continuous and vanishes outside the compact set $[0,t-\frac1{2n}]\times[\frac1{2n},b-\frac1{2n}]$), and
\begin{equation}\label{eq:4-3-15:phi}
\varphi(s,y)=v(t-s,y),\qquad\frac{\partial\varphi}{\partial s}(s,y)+\frac12\frac{\partial^2\varphi}{\partial y^2}(s,y)
=\Bigl(-\frac{\partial v}{\partial t}+\frac12\frac{\partial^2v}{\partial x^2}\Bigr)(t-s,y)=0
\end{equation}
for $(s,y)\in(-\infty,t-\frac1n]\times[\frac1n,b-\frac1n]$. Under $P^x$, $W=x+(W-x)$ with $W-x$ a standard Brownian motion, so It\^o's rule (\bk{Theorem 3.3.6}) gives, $P^x$-a.s.,
\begin{equation*}
\varphi(\tau_n,W_{\tau_n})=\varphi(0,x)+\int_0^{\tau_n}\Bigl(\frac{\partial\varphi}{\partial s}+\frac12\frac{\partial^2\varphi}{\partial y^2}\Bigr)(s,W_s)\,ds+\int_0^{\tau_n}\frac{\partial\varphi}{\partial y}(s,W_s)\,dW_s .
\end{equation*}
For $0\le s\le\tau_n$ we have $s\le t-\frac1n$ and $W_s\in[\frac1n,b-\frac1n]$, so by \eqref{eq:4-3-15:phi} the Lebesgue integral vanishes, and $\varphi(\tau_n,W_{\tau_n})=v(t-\tau_n,W_{\tau_n})$, $\varphi(0,x)=v(t,x)$. The stochastic integral $\int_0^\cdot\frac{\partial\varphi}{\partial y}(s,W_s)\,dW_s$ has a bounded, continuous, adapted integrand, hence is a continuous martingale in $\cM_2^c$, indexed by $[0,\infty)$ and starting at zero. Since $\tau_n\le t$ is a bounded stopping time, the optional sampling theorem in the form of \bk{Problem 1.3.23}(i), applied to this martingale and to its negative with the stopping times $0\le\tau_n$, shows that its value at $\tau_n$ has expectation zero. Taking expectations, we obtain
\begin{equation}\label{eq:4-3-15:mart}
v(t,x)=E^xv(t-\tau_n,W_{\tau_n});\quad n\ge n_0.
\end{equation}

\emph{Step 2. Passage to the limit.} The sequence $\{S_n\}$ is nondecreasing and $S_n\le T_0\wedge T_b$; let $S\defeq\lim_nS_n$. If $S<\infty$ then, since $W_{S_n}\in\{\frac1n,b-\frac1n\}$, continuity gives $W_S\in\{0,b\}$, so $S\ge T_0\wedge T_b$. Hence $S=T_0\wedge T_b$. We examine three events.

On $\{t<T_0\wedge T_b\}$ we have $S>t$, so $S_n>t$ and $\tau_n=t-\frac1n$ for all large $n$, and $(t-\tau_n,W_{\tau_n})=(\frac1n,W_{t-1/n})$ is a sequence of points of $D$ converging to $(0,W_t)$ with $W_t\in(0,b)$. By the first relation in \eqref{eq:4-3-15:bc} for $v$, $v(t-\tau_n,W_{\tau_n})\to f(W_t)$.

On $\{T_0<t\wedge T_b\}$ we have $S_n\uparrow T_0<t$, so $\tau_n=S_n$ for all large $n$, and $(t-\tau_n,W_{\tau_n})=(t-S_n,W_{S_n})\in D$ converges to $(t-T_0,0)$ with $t-T_0>0$. By the second relation in \eqref{eq:4-3-15:bc}, $v(t-\tau_n,W_{\tau_n})\to g(t-T_0)$.

On $\{T_b<t\wedge T_0\}$ we obtain in the same way $v(t-\tau_n,W_{\tau_n})\to k(t-T_b)$.

These three events are disjoint. On the complement of their union we have $T_0\wedge T_b\le t$; moreover $T_0\wedge T_b<t$ is impossible there, for then $T_0\ne T_b$ (the path cannot be at $0$ and at $b$ simultaneously), and one of the last two events would occur. Thus the complement is contained in $\{T_0\wedge T_b=t\}\subseteq\{W_t\in\{0,b\}\}$, which is $P^x$-null. Consequently, $P^x$-a.s.,
\begin{equation*}
\lim_{n\to\infty}v(t-\tau_n,W_{\tau_n})=f(W_t)\ind_{\{t<T_0\wedge T_b\}}+g(t-T_0)\ind_{\{T_0<t\wedge T_b\}}+k(t-T_b)\ind_{\{T_b<t\wedge T_0\}}.
\end{equation*}
Since $v$ is bounded, the bounded convergence theorem applied to \eqref{eq:4-3-15:mart} shows that $v(t,x)$ equals the right-hand side of \bk{(3.46)}. As $(t,x)\in D$ was arbitrary, every bounded solution coincides with the function \bk{(3.46)}, which is a bounded solution by Part~I.

\rmk{No compatibility between $f$, $g$, $k$ at the corners $(0,0)$, $(0,b)$ is needed, and no limits of $f$ at $0$, $b$ or of $g$, $k$ at $0$ are assumed: the corners are excluded from the boundary conditions, and a Brownian path started in $(0,b)$ almost surely does not lead the space-time point $(t-s,W_s)$ to a corner. Boundedness cannot be dropped from the uniqueness assertion. By Step~2 of Part~I, the function $H$ of \eqref{eq:4-3-15:dens} is of class $C^{1,2}$ on $D$, solves \bk{(3.1)} there, and tends to zero at every point of $\{0\}\times(0,b)$ by \eqref{eq:4-3-15:H0}. For $0<t_1\le s\le t_2$ and $0\le x\le b$ we have $\abs{h(s,x+2nb)}\le(2\pi)^{-1/2}t_1^{-3/2}(2\abs n+1)b\,e^{-(2\abs n-1)^2b^2/2t_2}$, $n\ne0$, so the series defining $H$ converges uniformly on $[t_1,t_2]\times[0,b]$ and $H$ extends continuously to $(0,\infty)\times[0,b]$; since $h(s,\cdot)$ is odd, the terms with indices $n$ and $-n$ cancel at $x=0$, those with indices $n$ and $-n-1$ cancel at $x=b$, and thus $H(s,0)=H(s,b)=0$. Hence $H$ solves the problem with $f=g=k=0$; but $H\not\equiv0$, because $\int_0^\infty H(s,x)\,ds=P^x[T_0<T_b]=(b-x)/b>0$. By Part~II, $H$ is unbounded on $D$ (it blows up near the corner $(0,0)$, like the function $h$ of \bk{(3.10)}).}
\end{kex}

\solsection{4.4}{The Formulas of Feynman and Kac}

\begin{kex}{Exercise}{4.6}{}
We read the growth hypothesis as in \bk{Corollary 4.5}: for each finite $T>0$ there are constants $C=C_T\in\R$ and $a=a_T$ with $0<a<1/(2Td)$ such that $-V(t,x)\le C+a\norm{x}^2$ for all $(t,x)\in[0,T]\times\R^d$. We linearize the equation by the substitution
\[
u(t,x)\defeq e^{-V(t,x)};\qquad (t,x)\in[0,\infty)\times\R^d,
\]
and then apply \bk{Corollary 4.5}.

\emph{Step 1: $u$ solves the linear Cauchy problem \bk{(4.1)}, \bk{(4.2)} with $g\equiv 0$, $f\equiv 1$.}
Since $V$ is real-valued and continuous on $[0,\infty)\times\R^d$ and of class $C^{1,2}$ on $(0,\infty)\times\R^d$, and the exponential function is of class $C^\infty$ on $\R$, the function $u$ is continuous and strictly positive on $[0,\infty)\times\R^d$ and of class $C^{1,2}$ on $(0,\infty)\times\R^d$. By the chain rule we have, on $(0,\infty)\times\R^d$ and for $1\le i\le d$,
\[
\frac{\partial u}{\partial t}=-u\,\frac{\partial V}{\partial t},\qquad
\frac{\partial u}{\partial x_i}=-u\,\frac{\partial V}{\partial x_i},\qquad
\frac{\partial^2 u}{\partial x_i^2}=u\Bigl(\frac{\partial V}{\partial x_i}\Bigr)^{2}-u\,\frac{\partial^2 V}{\partial x_i^2},
\]
and summing the last identity over $i$ gives $\Delta u=u\bigl(\norm{\nabla V}^2-\Delta V\bigr)$. Therefore, by \bk{(4.10)},
\[
\frac{\partial u}{\partial t}+ku
=-u\Bigl(\frac{\partial V}{\partial t}-k\Bigr)
=-u\Bigl(\frac12\Delta V-\frac12\norm{\nabla V}^2\Bigr)
=\frac12\Delta u\qquad\text{on }(0,\infty)\times\R^d ,
\]
which is \bk{(4.1)} with $g\equiv 0$. By \bk{(4.11)}, $u(0,x)=e^{0}=1$ for every $x\in\R^d$, which is \bk{(4.2)} with $f\equiv 1$.

\emph{Step 2: the growth condition.}
Fix $T>0$ and let $C$, $a$ be the constants associated with $T$. For $(t,x)\in[0,T]\times\R^d$,
\[
0<u(t,x)=e^{-V(t,x)}\le e^{C+a\norm{x}^2},
\]
so that, with $K\defeq e^{C}>0$ and $g\equiv0$,
\[
\max_{0\le t\le T}\abs{u(t,x)}+\max_{0\le t\le T}\abs{g(t,x)}\le Ke^{a\norm{x}^2};\qquad\forall\,x\in\R^d,
\]
where $0<a<1/(2Td)$. Note that only a lower bound on $V$ is needed, because only an upper bound on the positive function $u$ is.

\emph{Step 3: conclusion.}
The functions $f\equiv1$, $g\equiv0$ and $k\colon\R^d\to[0,\infty)$ are continuous, $u\colon[0,\infty)\times\R^d\to\R$ is continuous, of class $C^{1,2}$ on $(0,\infty)\times\R^d$, satisfies \bk{(4.1)} and \bk{(4.2)} by Step~1, and satisfies the growth condition of \bk{Corollary 4.5} for each finite $T>0$ by Step~2. Hence \bk{Corollary 4.5} applies, and the representation \bk{(4.9)} reads
\begin{equation}\label{eq:4-4-06:a}
e^{-V(t,x)}=u(t,x)=E^x\left[\exp\left\{-\int_0^t k(W_s)\,ds\right\}\right];\qquad 0\le t<\infty,\ x\in\R^d .
\end{equation}
Both sides of \eqref{eq:4-4-06:a} are strictly positive (the left-hand side because $V$ is real-valued; the right-hand side, independently, because for every $\omega$ the function $s\mapsto k(W_s(\omega))$ is continuous and nonnegative, so that $0\le\int_0^t k(W_s)\,ds<\infty$ and the integrand of the expectation lies in $(0,1]$). Taking logarithms in \eqref{eq:4-4-06:a} we obtain
\[
V(t,x)=-\log E^x\left[\exp\left\{-\int_0^t k(W_s)\,ds\right\}\right];\qquad 0\le t<\infty,\ x\in\R^d,
\]
which is \bk{(4.12)}.

\rmk{Since the expectation in \eqref{eq:4-4-06:a} lies in $(0,1]$, every solution $V$ as in the statement is in fact nonnegative; thus, a posteriori, the growth condition holds with $C=0$ and any $a>0$. In particular, \bk{(4.10)}, \bk{(4.11)} has at most one solution in the class described in the statement.}
\end{kex}

\begin{kex}{Exercise}{4.7}{}
We follow the method of \bk{Problem 2.25}: we stop the Brownian path at its exit times from an increasing sequence of balls, on which all the functions involved are bounded, and apply It\^o's rule; the only new feature is the discount factor
\begin{equation*}
X_t\defeq\exp\{-A_t\},\qquad A_t\defeq\alpha t+\int_0^t k(W_s)\,ds;\quad 0\le t<\infty .
\end{equation*}
Since $k\ge 0$ and $k(W_\cdot)$ is continuous, $A$ is a continuous, nondecreasing process with $A_0=0$, adapted to $\{\F^W_t\}$, and
\begin{equation}\label{eq:4-4-07:X}
0<X_t\le e^{-\alpha t}\le 1;\quad 0\le t<\infty .
\end{equation}
The process $(t,\omega)\mapsto f(W_t(\omega))X_t(\omega)$ is continuous and adapted, hence jointly measurable; by \bk{(4.15)} and Fubini--Tonelli, the random variable $\int_0^\infty\abs{f(W_t)}X_t\,dt$ is $P^x$-integrable, in particular finite a.s.\ $P^x$, and $z(x)$ in \bk{(4.14)} is a well-defined real number for every $x\in\R^d$.

We fix $x\in\R^d$ and work under $P^x$, under which $W$ is a $d$-dimensional Brownian motion with $P^x[W_0=x]=1$.

\emph{The filtration.} All the processes and random times appearing below are adapted to, respectively stopping times of, the filtration $\{\F^W_t\}$ generated by $W$. In order to use the stochastic calculus of Chapter~3, which presupposes the usual conditions, we regard them as adapted to the augmentation $\{\F^x_t\}$ of $\{\F^W_t\}$ under $P^x$: this filtration satisfies the usual conditions (\bk{Theorem 2.6.15} and \bk{Proposition 2.7.7}), and $W$ is a $d$-dimensional Brownian motion relative to it (\bk{Theorem 2.7.9}). The identities between expectations obtained below involve only $\F^W_\infty$-measurable random variables, and so do not depend on this choice.

\emph{Step 1: It\^o's rule.} The function $F(w,a)\defeq\psi(w)e^{-a}$ is of class $C^2$ on $\R^d\times\R$. We apply \bk{Theorem 3.3.6} to $F$ and to the $(d+1)$-dimensional semimartingale $(W,A)$, whose local martingale part is $(W^{(1)}-W_0^{(1)},\dots,W^{(d)}-W_0^{(d)},0)$, with $\qv{W^{(i)},W^{(j)}}_t=\delta_{ij}t$, and whose bounded variation part is $(0,\dots,0,A)$. Because $dA_s=(\alpha+k(W_s))\,ds$, we obtain, a.s.\ $P^x$, for all $0\le t<\infty$:
\begin{align*}
\psi(W_t)X_t
&=\psi(x)-\int_0^t\psi(W_s)X_s\bigl(\alpha+k(W_s)\bigr)ds
 +\sum_{i=1}^d\int_0^t X_s\frac{\partial\psi}{\partial x_i}(W_s)\,dW^{(i)}_s
 +\frac12\int_0^t X_s\,\Delta\psi(W_s)\,ds .
\end{align*}
By \bk{(4.13)} we have $\tfrac12\Delta\psi-(\alpha+k)\psi=-f$ on $\R^d$, and therefore, a.s.\ $P^x$,
\begin{equation}\label{eq:4-4-07:ito}
\psi(W_t)X_t+\int_0^t f(W_s)X_s\,ds=\psi(x)+M_t;\qquad
M_t\defeq\sum_{i=1}^d\int_0^t X_s\frac{\partial\psi}{\partial x_i}(W_s)\,dW^{(i)}_s,
\end{equation}
for all $0\le t<\infty$. Each integrand $X_s\,\frac{\partial\psi}{\partial x_i}(W_s)$ is continuous and adapted, hence progressively measurable with $\int_0^t X_s^2\bigl(\frac{\partial\psi}{\partial x_i}(W_s)\bigr)^2ds<\infty$ for every $t$ and every $\omega$; that is, it belongs to $\cP^*$, and $M\in\cM^{c,\loc}$. (For a general $\psi$ of class $C^2$ these integrands need not belong to $\cL^*$, which is why we localize.)

\emph{Step 2: localization.} For every integer $n>\norm{x}$ let
\begin{equation*}
\tau_n\defeq\inf\{t\ge 0;\ \norm{W_t}\ge n\}.
\end{equation*}
This is the hitting time of a closed set by a continuous process, hence a stopping time of $\{\F^W_t\}$ (\bk{Problem 1.2.7}). Since $\norm{W_0}=\norm{x}<n$ a.s.\ $P^x$, we have $\norm{W_s}\le n$ for $0\le s\le\tau_n$. Every path of $W$ is continuous, hence bounded on every bounded time-interval, so that
\begin{equation}\label{eq:4-4-07:tau}
\tau_n\uparrow\infty\quad\text{as } n\to\infty .
\end{equation}
Let $C_n\defeq\max_{\norm{y}\le n}\norm{\nabla\psi(y)}<\infty$. By the sample-path property \bk{(3.2.24)}, which remains valid for integrands in $\cP^*$ (see the paragraph preceding \bk{Proposition 3.2.24}), we have a.s.\ $P^x$, for all $0\le t<\infty$, $M_{t\wedge\tau_n}=\sum_{i=1}^d\int_0^t Y^{(i)}_s\,dW^{(i)}_s$ with
\begin{equation*}
Y^{(i)}_s\defeq\ind_{\{s\le\tau_n\}}X_s\frac{\partial\psi}{\partial x_i}(W_s),\qquad \abs{Y^{(i)}_s}\le C_n\quad\text{a.s. } P^x,
\end{equation*}
thanks to \eqref{eq:4-4-07:X}. The processes $Y^{(i)}$ are progressively measurable and bounded, so $E^x\int_0^t(Y^{(i)}_s)^2\,ds<\infty$ for every $t$; thus $\{M_{t\wedge\tau_n};\,0\le t<\infty\}$ is a square-integrable martingale starting at zero, and $E^xM_{t\wedge\tau_n}=0$. Evaluating \eqref{eq:4-4-07:ito} at $t\wedge\tau_n$ (the two sides are continuous in $t$, so the identity holds a.s.\ simultaneously for all $t$) we get
\begin{equation*}
\psi(W_{t\wedge\tau_n})X_{t\wedge\tau_n}+\int_0^{t\wedge\tau_n} f(W_s)X_s\,ds=\psi(x)+M_{t\wedge\tau_n}\quad\text{a.s. } P^x .
\end{equation*}
The first term on the left is bounded by $\max_{\norm{y}\le n}\abs{\psi(y)}$, and the second is dominated by the $P^x$-integrable random variable $\int_0^\infty\abs{f(W_s)}X_s\,ds$. Taking expectations with $t=n$ and writing $\sigma_n\defeq n\wedge\tau_n$, we arrive at
\begin{equation}\label{eq:4-4-07:key}
\psi(x)=E^x\bigl[\psi(W_{\sigma_n})X_{\sigma_n}\bigr]+E^x\int_0^{\sigma_n}f(W_s)X_s\,ds;\qquad n>\norm{x}.
\end{equation}

\emph{Step 3: the limit of the second term.} By \eqref{eq:4-4-07:tau}, $\sigma_n\uparrow\infty$, so on the event $\{\int_0^\infty\abs{f(W_s)}X_s\,ds<\infty\}$ of full $P^x$-measure we have $\int_0^{\sigma_n}f(W_s)X_s\,ds\to\int_0^\infty f(W_s)X_s\,ds$ (dominated convergence for the Lebesgue integral in $s$). Moreover, $\abs[\big]{\int_0^{\sigma_n}f(W_s)X_s\,ds}\le\int_0^\infty\abs{f(W_s)}X_s\,ds$, which is $P^x$-integrable by \bk{(4.15)}. The dominated convergence theorem gives
\begin{equation}\label{eq:4-4-07:lim}
\lim_{n\to\infty}E^x\int_0^{\sigma_n}f(W_s)X_s\,ds=E^x\int_0^\infty f(W_s)X_s\,ds=z(x).
\end{equation}

\emph{Step 4: $\psi$ bounded.} Let $K\defeq\sup_{y\in\R^d}\abs{\psi(y)}<\infty$. By \eqref{eq:4-4-07:X},
\begin{equation*}
\abs[\big]{\psi(W_{\sigma_n})X_{\sigma_n}}\le Ke^{-\alpha\sigma_n}\le K,
\end{equation*}
and $e^{-\alpha\sigma_n}\to 0$ as $n\to\infty$, because $\alpha>0$ and $\sigma_n\to\infty$. By the bounded convergence theorem, $E^x[\psi(W_{\sigma_n})X_{\sigma_n}]\to 0$. Letting $n\to\infty$ in \eqref{eq:4-4-07:key} and using \eqref{eq:4-4-07:lim}, we obtain $\psi(x)=z(x)$.

\emph{Step 5: $\psi$ nonnegative.} Now $\psi(W_{\sigma_n})X_{\sigma_n}\ge 0$, so \eqref{eq:4-4-07:key} yields
\begin{equation*}
\psi(x)\ge E^x\int_0^{\sigma_n}f(W_s)X_s\,ds;\qquad n>\norm{x},
\end{equation*}
and letting $n\to\infty$ we conclude from \eqref{eq:4-4-07:lim} that $\psi(x)\ge z(x)$.

Since $x\in\R^d$ was arbitrary, both assertions are proved.

\rmk{The inequality in the second assertion can be strict, and boundedness cannot be dropped in the first: for $d=1$, $k\equiv 0$, $f\equiv 0$ we have $z\equiv 0$, whereas $\psi(x)=\cosh\bigl(x\sqrt{2\alpha}\bigr)$ is a nonnegative (unbounded) solution of class $C^2$ to \bk{(4.13)}. In the terms of \bk{Problem 2.25}: letting $t\to\infty$ in the stopped identity of Step 2 would give the analogue of \bk{(2.27)} for the ball $\{y\in\R^d;\ \norm{y}<n\}$, with boundary datum $\psi$ on the sphere $\{\norm{y}=n\}$; \eqref{eq:4-4-07:key} is its version at the bounded time $\sigma_n$, in which the term $E^x[\psi(W_{\sigma_n})X_{\sigma_n}]$ vanishes in the limit if $\psi$ is bounded and has a sign if $\psi\ge 0$.}
\end{kex}

\begin{kex}{Exercise}{4.13}{}
Here $\{W_s,\F_s;\,s\ge0\}$, $(\Omega,\F)$, $\{P^x\}_{x\in\R^d}$ is a $d$-dimensional Brownian family (\bk{Definition 2.5.8}), and the dot in the display denotes multiplication. We fix $x\in\R^d$, $t>0$, $0<a<1/(2t)$, write $u\cdot v$ for the inner product of $u,v\in\R^d$, and set
\[
B_s\defeq W_s-W_0,\qquad M\defeq\sup_{0\le s\le t}\norm{B_s},\qquad M^*\defeq\sup_{0\le s\le t}\norm{W_s}.
\]
Every path of $W$ is continuous (\bk{Definition 2.5.8}\,(iii) and \bk{Definition 2.5.1}), so both suprema are finite, they are attained, and they may be computed over the countable set $[0,t]\cap\Q$; in particular $M$ and $M^*$ are $\F$-measurable random variables with values in $[0,\infty)$, and $\sup_{0\le s\le t}\norm{W_s}^2=(M^*)^2$.

\emph{Step 1: reduction to one-dimensional Brownian motions.}
Let $\theta\in\R^d$ with $\norm{\theta}=1$, and put $B^\theta_s\defeq\theta\cdot B_s$, $0\le s<\infty$. This process is continuous, adapted to $\{\F_s\}$, and $B^\theta_0=0$. By \bk{Definition 2.5.8}\,(iii), $W$ is under $P^x$ a $d$-dimensional Brownian motion starting at $x$; hence, by \bk{Definition 2.5.1}\,(ii), under $P^x$ and for $0\le r<s$, the increment $W_s-W_r=B_s-B_r$ is independent of $\F_r$ and normally distributed with mean zero and covariance matrix $(s-r)I_d$; hence $B^\theta_s-B^\theta_r=\theta\cdot(B_s-B_r)$ is independent of $\F_r$ and normally distributed with mean zero and variance $(s-r)\norm{\theta}^2=s-r$. Thus $\{B^\theta_s,\F_s;\,0\le s<\infty\}$ is, under $P^x$, a standard one-dimensional Brownian motion (\bk{Definition 2.1.1}), and the estimate \bk{(2.8.3)$'$} of \bk{Problem 2.8.2}, which concerns only the law induced on $(C[0,\infty),\B(C[0,\infty)))$ by a standard one-dimensional Brownian motion (this law is Wiener measure, whatever the underlying probability space and filtration, because it is determined by the finite-dimensional distributions), gives
\begin{equation}\label{eq:4-4-13:a}
P^x\Bigl[\max_{0\le s\le t}\abs{B^\theta_s}\ge c\Bigr]\le\sqrt{\frac{t}{2\pi}}\;\frac{4}{c}\,e^{-c^2/2t};\qquad c>0 .
\end{equation}

\emph{Step 2: a tail estimate for $M$.}
Fix $\eps\in(0,1)$. The unit sphere $S\defeq\{u\in\R^d;\ \norm{u}=1\}$ is compact, so it is covered by finitely many open balls of radius $\eps$ centred at points of $S$: there is a finite set $\Theta\subset S$, with $N\defeq\abs{\Theta}$ elements, such that every $u\in S$ satisfies $\norm{u-\theta}<\eps$ for some $\theta\in\Theta$. For such a pair we have $\theta\cdot u=1-\frac12\norm{u-\theta}^2\ge1-\frac12\eps^2\ge1-\eps$. Applying this to $u=v/\norm{v}$, we obtain
\[
(1-\eps)\norm{v}\le\max_{\theta\in\Theta}\,\theta\cdot v;\qquad v\in\R^d
\]
(for $v=0$ this is trivial). Consequently $(1-\eps)M\le\max_{\theta\in\Theta}\max_{0\le s\le t}\abs{B^\theta_s}$, and \eqref{eq:4-4-13:a} with $c=(1-\eps)b$ yields
\begin{equation}\label{eq:4-4-13:b}
P^x[M\ge b]\le\sum_{\theta\in\Theta}P^x\Bigl[\max_{0\le s\le t}\abs{B^\theta_s}\ge(1-\eps)b\Bigr]
\le\frac{K}{b}\,\exp\Bigl\{-\frac{(1-\eps)^2b^2}{2t}\Bigr\};\qquad b>0,
\end{equation}
where $K\defeq 4N(1-\eps)^{-1}\sqrt{t/2\pi}$ depends only on $d$, $t$ and $\eps$.

\emph{Step 3: exponential moments of $M^2$.}
Let $0<\alpha<(1-\eps)^2/(2t)$. For every $m\ge0$ we have $e^{\alpha m^2}=1+\int_0^\infty2\alpha b\,e^{\alpha b^2}\ind_{\{b\le m\}}\,db$. Taking $m=M$, integrating with respect to $P^x$ and using Tonelli's theorem (the integrand is nonnegative and jointly measurable in $(b,\omega)$), then \eqref{eq:4-4-13:b}, we find
\begin{equation}\label{eq:4-4-13:c}
E^x\bigl[e^{\alpha M^2}\bigr]=1+\int_0^\infty2\alpha b\,e^{\alpha b^2}P^x[M\ge b]\,db
\le1+2\alpha K\int_0^\infty\exp\Bigl\{-\Bigl(\frac{(1-\eps)^2}{2t}-\alpha\Bigr)b^2\Bigr\}\,db<\infty,
\end{equation}
because the coefficient of $b^2$ in the last exponent is strictly negative.

\emph{Step 4: conclusion.}
Since $2at<1$, we can choose $\eps\in(0,1)$ so small that $2at<(1-\eps)^2$, and then $\delta>0$ so small that $\alpha\defeq a(1+\delta)<(1-\eps)^2/(2t)$. By \bk{Definition 2.5.8}\,(ii), $P^x[W_0=x]=1$; on the event $\{W_0=x\}$ we have $\norm{W_s}\le\norm{x}+\norm{B_s}$ for all $s$, hence $M^*\le\norm{x}+M$, and the elementary inequality $2pq\le\delta p^2+\delta^{-1}q^2$ gives
\[
(M^*)^2\le(\norm{x}+M)^2\le(1+\delta)M^2+(1+\delta^{-1})\norm{x}^2;\qquad\text{a.s. }P^x .
\]
Therefore, by \eqref{eq:4-4-13:c},
\[
E^x\Bigl[\exp\Bigl\{a\sup_{0\le s\le t}\norm{W_s}^2\Bigr\}\Bigr]
\le e^{a(1+\delta^{-1})\norm{x}^2}\,E^x\bigl[e^{\alpha M^2}\bigr]<\infty .
\]

\rmk{The range $0<a<1/(2t)$ cannot be enlarged: if $a\ge1/(2t)$, then the first coordinate $W^{(1)}_t$ of $W_t$ is, under $P^0$, normal with mean zero and variance $t$, and
\[
E^0\Bigl[\exp\Bigl\{a\sup_{0\le s\le t}\norm{W_s}^2\Bigr\}\Bigr]\ge E^0\bigl[e^{a(W^{(1)}_t)^2}\bigr]
=\frac{1}{\sqrt{2\pi t}}\int_{-\infty}^\infty e^{(a-\frac{1}{2t})y^2}\,dy=\infty .
\]
The finite set $\Theta$ of directions in Step~2 is one way of passing from dimension one to dimension $d$ without loss in the exponent, but not the only one. Let $B^{(1)},\dots,B^{(d)}$ be the coordinates of $B$, so that $M^2\le\sum_{i=1}^d\sup_{0\le s\le t}\abs{B^{(i)}_s}^2$. For $0=s_0<s_1<\dots<s_n$ the increments $B_{s_k}-B_{s_{k-1}}$, $1\le k\le n$, are under $P^x$ independent with laws $N(0,(s_k-s_{k-1})I_d)$ (\bk{Definition 2.5.1}\,(ii)), so the $nd$ random variables $B^{(i)}_{s_k}-B^{(i)}_{s_{k-1}}$ are independent; hence the $\sigma$-fields $\sigma(B^{(i)}_s;\,0\le s<\infty)$, $1\le i\le d$, generated by the coordinate processes are independent (their finite-dimensional cylinder sets form $\pi$-systems; Dynkin's $\pi$-$\lambda$ theorem), and so are the random variables $\sup_{0\le s\le t}\abs{B^{(i)}_s}$. Therefore, for $0<\alpha<1/(2t)$,
\[
E^x\bigl[e^{\alpha M^2}\bigr]\le\prod_{i=1}^dE^x\Bigl[\exp\Bigl\{\alpha\sup_{0\le s\le t}\abs{B^{(i)}_s}^2\Bigr\}\Bigr]<\infty ,
\]
each factor being finite by \eqref{eq:4-4-13:a} (with $\theta$ the $i$-th unit vector) and the computation of Step~3; Step~4 then applies unchanged. With H\"older's inequality in place of independence, the same coordinate bound would only reach $0<a<1/(2td)$.}
\end{kex}

\begin{kex}{Exercise}{4.14}{}
The statements to be proved are the following.

\emph{\bk{Theorem 4.2}$'$. Let $v$ be as in \bk{Definition 4.1} and assume that \bk{(4.6)} holds for some constants $K>0$ and $0<a<1/(2T)$. Then $v$ admits the representation \bk{(4.7)}; in particular, such a solution is unique.}

\emph{\bk{Corollary 4.5}$'$. Let $f$, $k$, $g$, $u$ be as in \bk{Corollary 4.5}, except that for each finite $T>0$ the constants $K>0$ and $a$ of the growth condition are only required to satisfy $0<a<1/(2T)$. Then $u$ admits the representation \bk{(4.9)}.}

In the book's proof of \bk{Theorem 4.2} the restriction $a<1/(2Td)$ is used only to show that the boundary term on $\{S_n\le r\}$ vanishes as $n\to\infty$, by an estimate on the tail of $P^x[S_n\le T]$. We replace that estimate by a domination: throughout, $T>0$, $K$ and $a\in(0,1/(2T))$ are fixed, $x\in\R^d$, and
\begin{equation}\label{eq:4-4-14:Y}
Y\defeq\exp\Bigl\{a\max_{0\le s\le T}\norm{W_s}^2\Bigr\},\qquad E^xY<\infty .
\end{equation}
Here $Y$ is a random variable because, by continuity of the paths, the maximum is a supremum over the rational $s\in[0,T]$, and the integrability is \bk{Exercise 4.13} with $t=T$. We also write $A_\theta\defeq\int_0^\theta k(W_s)\,ds$; since $k\ge0$ is continuous, $A$ is a continuous, nondecreasing, adapted process with $A_0=0$ and $0<e^{-A_\theta}\le1$. By \bk{(4.6)}, for all $s\in[0,T]$ and $\theta\in[0,T]$,
\begin{equation}\label{eq:4-4-14:dom}
\abs{v(s,W_\theta)}\le Ke^{a\norm{W_\theta}^2}\le KY,\qquad \abs{g(s,W_\theta)}\le KY .
\end{equation}

\textbf{Proof of \bk{Theorem 4.2}$'$.} For $t=T$, \bk{(4.7)} reads $v(T,x)=E^xf(W_0)=f(x)$, which is \bk{(4.5)}. We treat first $0<t<T$ and then $t=0$.

\emph{Step 1: It\^o's rule.} Fix $0<t<T$ and $0<r<r'<T-t$. Let $m\colon[0,\infty)\to[0,r')$ be given by $m(\theta)=\theta$ for $0\le\theta\le r$ and $m(\theta)=r+(r'-r)\tanh\bigl((\theta-r)/(r'-r)\bigr)$ for $\theta>r$; it is of class $C^1$, because both one-sided derivatives at $r$ equal $1$. The function
\[
F(\theta,y,b)\defeq v\bigl(t+m(\theta),y\bigr)e^{-b};\qquad (\theta,y,b)\in[0,\infty)\times\R^d\times\R,
\]
is of class $C^{1,2}$, since $t+m(\theta)$ takes values in $[t,t+r')\subset(0,T)$, the interior of the strip on which $v$ is of class $C^{1,2}$ (see the convention preceding \bk{Theorem 3.3}). Under $P^x$ the process $W-W_0$ is a $d$-dimensional Brownian motion, so its components are in $\cM^{c,\loc}$ with $\qv{W^{(i)},W^{(j)}}_\theta=\delta_{ij}\theta$, and $A$ is an adapted process of bounded variation with $A_0=0$. (Stochastic integrals under $P^x$ are taken with respect to the filtration $\{\F^x_\theta\}$, where $\F^x_\theta$ is the $\sigma$-field generated by $\F_{\theta+}$ and the $P^x$-null sets of the $P^x$-completion of $\F$. This filtration satisfies the usual conditions, and $W-W_0$ is a Brownian motion relative to it: for $0\le s<u$ and $0<\eps<u-s$ the increment $W_u-W_{s+\eps}$ is independent of $\F_{s+\eps}\supseteq\F_{s+}$ and converges to $W_u-W_s$ as $\eps\downarrow0$, so $W_u-W_s$ is independent of $\F_{s+}$; adjoining null sets does not affect independence. Stopping times of $\{\F_\theta\}$ are stopping times of $\{\F^x_\theta\}$.) We apply \bk{Theorem 3.3.6} to the $(d+1)$-dimensional semimartingale $(W,A)$, whose last component has zero martingale part, and to $F$. For $0\le\theta\le r$ we have $m(\theta)=\theta$, $m'(\theta)=1$, and therefore, a.s.\ $P^x$, for all $0\le\theta\le r$,
\begin{align*}
v(t+\theta,W_\theta)e^{-A_\theta}
&=v(t,W_0)+\int_0^\theta e^{-A_s}\Bigl[\frac{\partial v}{\partial t}+\frac12\Delta v-kv\Bigr](t+s,W_s)\,ds\\
&\quad+\sum_{i=1}^d\int_0^\theta e^{-A_s}\frac{\partial v}{\partial x_i}(t+s,W_s)\,dW^{(i)}_s .
\end{align*}
By \bk{(4.4)} the bracket equals $-g(t+s,W_s)$, and $W_0=x$ a.s.\ $P^x$. Hence, a.s.\ $P^x$,
\begin{equation}\label{eq:4-4-14:ito}
v(t+\theta,W_\theta)e^{-A_\theta}+\int_0^\theta g(t+s,W_s)e^{-A_s}\,ds=v(t,x)+N_\theta;\qquad 0\le\theta\le r,
\end{equation}
where $N_\theta\defeq\sum_{i=1}^d\int_0^\theta e^{-A_s}\frac{\partial v}{\partial x_i}\bigl(t+m(s),W_s\bigr)\,dW^{(i)}_s$, $0\le\theta<\infty$, is a continuous local martingale.

\emph{Step 2: localization.} For every integer $n>\norm{x}$ let $S_n\defeq\inf\{\theta\ge0;\ \norm{W_\theta}\ge n\}$, a stopping time as the hitting time of a closed set by a continuous adapted process (\bk{Problem 1.2.7}). Since $\norm{W_0}=\norm{x}<n$ a.s.\ $P^x$, we have $\norm{W_s}\le n$ for $0\le s\le S_n$, a.s.\ $P^x$; and since every path of $W$ is bounded on bounded intervals,
\begin{equation}\label{eq:4-4-14:S}
\text{for every }\omega\text{ there is }n_0(\omega)\text{ with } S_n(\omega)>T\text{ for all }n\ge n_0(\omega).
\end{equation}
Let $C_n$ be the maximum of $\norm{\nabla v}$ over the compact set $[t,t+r]\times\{\norm{y}\le n\}$. By \bk{(3.2.24)} (valid for local martingale integrators, as noted before \bk{Proposition 3.2.24}), $N_{\theta\wedge r\wedge S_n}=\sum_{i=1}^d\int_0^\theta Z^{(i)}_s\,dW^{(i)}_s$ with
\[
Z^{(i)}_s\defeq\ind_{\{s\le r\wedge S_n\}}e^{-A_s}\frac{\partial v}{\partial x_i}(t+s,W_s),\qquad \abs{Z^{(i)}_s}\le C_n\quad\text{a.s. }P^x .
\]
The processes $Z^{(i)}$ are bounded and progressively measurable, so $\{N_{\theta\wedge r\wedge S_n};\ 0\le\theta<\infty\}$ is a square-integrable martingale starting at zero; in particular $E^xN_{r\wedge S_n}=0$. Evaluating \eqref{eq:4-4-14:ito} at $\theta=r\wedge S_n$ and taking expectations (both terms on the left are integrable by \eqref{eq:4-4-14:dom}, \eqref{eq:4-4-14:Y}) we obtain
\begin{equation}\label{eq:4-4-14:n}
v(t,x)=E^x\Bigl[v\bigl(t+r\wedge S_n,W_{r\wedge S_n}\bigr)e^{-A_{r\wedge S_n}}
+\int_0^{r\wedge S_n}g(t+s,W_s)e^{-A_s}\,ds\Bigr].
\end{equation}

\emph{Step 3: $n\to\infty$, then $r\uparrow T-t$.} By \eqref{eq:4-4-14:S}, for every $\omega$ the random variable inside the expectation in \eqref{eq:4-4-14:n} equals, for all $n\ge n_0(\omega)$,
\[
\Phi_r\defeq v(t+r,W_r)e^{-A_r}+\int_0^{r}g(t+s,W_s)e^{-A_s}\,ds .
\]
Because $t+r\wedge S_n\le T$ and $r\wedge S_n\le T$, the bounds \eqref{eq:4-4-14:dom} show that it is dominated in absolute value by $(1+T)KY$, for every $n$. This is the point where \bk{Exercise 4.13} replaces the tail estimate of the book. By \eqref{eq:4-4-14:Y} and the dominated convergence theorem, $v(t,x)=E^x\Phi_r$ for every $0<r<T-t$. Now let $r\uparrow T-t$. The function $v$ is continuous on $[0,T]\times\R^d$ and the paths of $W$ and $A$ are continuous, so $v(t+r,W_r)e^{-A_r}\to v(T,W_{T-t})e^{-A_{T-t}}=f(W_{T-t})e^{-A_{T-t}}$ by \bk{(4.5)}; and $\int_0^rg(t+s,W_s)e^{-A_s}\,ds\to\int_0^{T-t}g(t+s,W_s)e^{-A_s}\,ds$, the integrand being bounded by $KY<\infty$. Again $\abs{\Phi_r}\le(1+T)KY$, and dominated convergence gives
\begin{equation}\label{eq:4-4-14:FK}
v(t,x)=E^x\Bigl[f(W_{T-t})e^{-A_{T-t}}+\int_0^{T-t}g(t+\theta,W_\theta)e^{-A_\theta}\,d\theta\Bigr];\qquad 0<t<T,
\end{equation}
which is \bk{(4.7)} for $0<t<T$. (The integrand $(\theta,\omega)\mapsto g(t+\theta,W_\theta(\omega))e^{-A_\theta(\omega)}$ is continuous in $\theta$ and measurable in $\omega$, hence jointly measurable, and it is dominated by $KY$; so the expectation in \eqref{eq:4-4-14:FK} is well defined and finite.)

\emph{Step 4: $t=0$.} We let $t\downarrow0$ in \eqref{eq:4-4-14:FK}. The left-hand side converges to $v(0,x)$ by continuity of $v$. Write the random variable on the right-hand side as
\[
\Psi_t\defeq f(W_{T-t})e^{-A_{T-t}}+\int_0^{T}\ind_{[0,T-t)}(\theta)\,g(t+\theta,W_\theta)e^{-A_\theta}\,d\theta .
\]
For fixed $\omega$, the first term converges to $f(W_T)e^{-A_T}$ as $t\downarrow0$, by continuity of $f=v(T,\cdot)$ and of the paths. For each $\theta\in[0,T)$ the integrand converges to $g(\theta,W_\theta)e^{-A_\theta}$, by continuity of $g$ on $[0,T]\times\R^d$, and it is bounded by $KY(\omega)$; by bounded convergence on $[0,T]$ the integral converges to $\int_0^Tg(\theta,W_\theta)e^{-A_\theta}\,d\theta$. Finally $\abs{\Psi_t}\le(1+T)KY$ by \eqref{eq:4-4-14:dom} (recall $f=v(T,\cdot)$), so that dominated convergence yields
\[
v(0,x)=E^x\Bigl[f(W_T)e^{-A_T}+\int_0^{T}g(\theta,W_\theta)e^{-A_\theta}\,d\theta\Bigr],
\]
which is \bk{(4.7)} for $t=0$. This completes the proof of \bk{(4.7)} on $[0,T]\times\R^d$. The right-hand side of \bk{(4.7)} depends only on $f$, $g$, $k$; hence two solutions of the Cauchy problem \bk{(4.4)}, \bk{(4.5)} which satisfy \bk{(4.6)}, each with its own constants $K>0$ and $a\in(0,1/(2T))$, coincide.

\textbf{Proof of \bk{Corollary 4.5}$'$.} For $t=0$, \bk{(4.9)} reads $u(0,x)=E^xf(W_0)=f(x)$, which is \bk{(4.2)}. Fix $T>0$, let $K>0$ and $0<a<1/(2T)$ be the constants of the hypothesis for this $T$, and set
\[
v(t,x)\defeq u(T-t,x),\qquad \wt g(t,x)\defeq g(T-t,x);\qquad (t,x)\in[0,T]\times\R^d .
\]
Then $v$ and $\wt g$ are continuous on $[0,T]\times\R^d$, and $v$ is of class $C^{1,2}$ on $[0,T)\times\R^d$ because $u$ is of class $C^{1,2}$ on $(0,\infty)\times\R^d$. For $(t,x)\in[0,T)\times\R^d$ we get from \bk{(4.1)} at the point $(T-t,x)\in(0,\infty)\times\R^d$:
\[
-\frac{\partial v}{\partial t}(t,x)+k(x)v(t,x)=\Bigl[\frac{\partial u}{\partial t}+ku\Bigr](T-t,x)=\Bigl[\frac12\Delta u+g\Bigr](T-t,x)=\frac12\Delta v(t,x)+\wt g(t,x),
\]
and $v(T,x)=u(0,x)=f(x)$ by \bk{(4.2)}. Thus $v$ solves the Cauchy problem of \bk{Definition 4.1} with potential $k$, Lagrangian $\wt g$ and terminal condition $f$. Moreover,
\[
\max_{0\le t\le T}\abs{v(t,x)}+\max_{0\le t\le T}\abs{\wt g(t,x)}
=\max_{0\le t\le T}\abs{u(t,x)}+\max_{0\le t\le T}\abs{g(t,x)}\le Ke^{a\norm{x}^2};\qquad x\in\R^d,
\]
which is \bk{(4.6)} with $0<a<1/(2T)$. By \bk{Theorem 4.2}$'$, formula \bk{(4.7)} holds for $v$, $\wt g$; at $t=0$ it reads
\[
u(T,x)=v(0,x)=E^x\Bigl[f(W_T)\exp\Bigl\{-\int_0^Tk(W_s)\,ds\Bigr\}
+\int_0^Tg(T-\theta,W_\theta)\exp\Bigl\{-\int_0^\theta k(W_s)\,ds\Bigr\}d\theta\Bigr].
\]
Since $T>0$ and $x\in\R^d$ are arbitrary, this is \bk{(4.9)}.

\rmk{(a) The same proof gives \bk{Remark 4.3} with $0<a<1/(2T)$ in \bk{(4.8)}, for $0<t\le T$: if $g\ge0$ and only $v$ is assumed to satisfy the growth condition, then in \eqref{eq:4-4-14:n} the term containing $v$ is still dominated by $KY$, while the integral $\int_0^{r\wedge S_n}g(t+s,W_s)e^{-A_s}\,ds$ is bounded (because $\norm{W_s}\le n$ for $s\le S_n$), nonnegative, and nondecreasing both in $n$ and in $r$; in Step~3 the monotone convergence theorem takes the place of dominated convergence for this term. Step~4 is not available in this situation. For $t=0$ one fixes $0<\eps<r<T$ and argues as in Steps~1--3 with the process $\{v(\eps+\theta,W_{\eps+\theta})e^{-A_{\eps+\theta}};\ 0\le\theta\le r-\eps\}$, the filtration $\{\F^x_{\eps+\theta}\}$ and the Brownian motion $\{W_{\eps+\theta}-W_\eps\}$, stopping at the first time $\theta$ with $\norm{W_{\eps+\theta}}\ge n$ (on $\{\norm{W_\eps}\ge n\}$ this time is zero and the stochastic integrand vanishes; otherwise it is bounded as in Step~2). This gives
$E^x[v(\eps,W_\eps)e^{-A_\eps}]=E^x[f(W_T)e^{-A_T}+\int_\eps^Tg(\theta,W_\theta)e^{-A_\theta}\,d\theta]$. As $\eps\downarrow0$, the left-hand side converges to $v(0,x)$ by dominated convergence (the integrand is bounded by $KY$), and the right-hand side converges to that of \bk{(4.7)} with $t=0$ by monotone convergence. (b) The right-hand side of \bk{(4.7)} need not be finite when $a\ge1/(2T)$: for $f(y)=e^{a\norm{y}^2}$ we have $E^xf(W_T)=\infty$.}
\end{kex}


%% ========================================================================
\solchapter{5}{Stochastic Differential Equations}

\begin{center}
\begin{minipage}{0.86\textwidth}
\small
The 47 problems and exercises of Chapter 5 for which the book prints no solution: 1.2, 2.17, 2.19, 2.20, 2.26, 2.27, 2.28, 3.12, 3.13, 4.3, 4.4, 4.33, 4.34, 4.35, 5.12, 5.14, 5.28, 5.31, 5.34, 5.35, 5.36, 5.37, 5.38, 5.39, 5.40, 5.41, 5.42, 6.1, 6.6, 6.13, 6.14, 6.16, 6.17, 6.18, 6.20, 6.21, 6.22, 7.7, 7.10, 7.11, 7.12, 7.13, 8.10, 8.13, 8.15, 8.20, 8.23.  For 1 more, 6.11, the book prints only a hint or solves only some of the parts; what it leaves out is supplied here.  The book's own ``Solutions to Selected Problems'' fully covers 2.7, 2.10, 2.11, 2.12, 3.15, 3.21, 3.22, 4.7, 4.13, 4.25, 5.3, 5.27, 6.2, 6.15, 7.3, 8.1, 8.17, 8.19, which are therefore omitted.
\end{minipage}
\end{center}

\bigskip
\hrule
\bigskip

\solsection{5.1}{Introduction}

\begin{kex}{Problem}{1.2}{}
We fix $x\in\R^d$ and work on $(\Omega,\F,P^x)$, where the filtration $\{\F_t\}$ satisfies the usual conditions (the standing assumption of \bk{Section 3.2}, under which the stochastic integrals in \bk{(1.9)} are defined), $W$ is an $r$-dimensional Brownian motion relative to $\{\F_t\}$, and $X$ is continuous, adapted and satisfies \bk{(1.9)}; the matrix $a$ is given by \bk{(1.10)}. Replacing $W$ by $W-W_0$, which changes neither the hypotheses ($W_0$ is $\F_0$-measurable, so $W-W_0$ is continuous and adapted; its increments are those of $W$, so it satisfies \bk{Definition 2.5.1}\,(ii), and it satisfies \bk{Definition 2.5.1}\,(i) with $\mu$ the unit mass at $0$; thus it is an $r$-dimensional Brownian motion starting at $0$) nor the integrals in \bk{(1.9)}, we may assume $W_0=0$. Let $K<\infty$ be a common bound for all the $\abs{b_i}$ and $\abs{\sigma_{ij}}$, so that
\begin{equation}\label{eq:5-1-02:abound}
  \abs{a_{ik}(y)}\le rK^2;\qquad y\in\R^d,\ 1\le i,k\le d ,
\end{equation}
and note that the $a_{ik}$ are continuous. We write, for $1\le i\le d$ and $0\le t<\infty$,
\[
  B^{(i)}_t\defeq\int_0^t b_i(X_s)\,ds,\qquad
  M^{(i)}_t\defeq\sum_{j=1}^r\int_0^t\sigma_{ij}(X_s)\,dW^{(j)}_s,
  \qquad\text{so that}\quad X^{(i)}_t=x_i+B^{(i)}_t+M^{(i)}_t .
\]
In particular $X_0=x$, $P^x$-a.s.

\emph{Step 1: the martingales $M^{(i)}$.} By \bk{Problem 2.5.5}, each $W^{(j)}$ is in $\cM_2^c$ and $\qv{W^{(j)},W^{(l)}}_t=t\,\delta_{jl}$. For every bounded, continuous $g\colon\R^d\to\R$ the process $g(X)$ is continuous and adapted, hence progressively measurable (\bk{Proposition 1.1.13}), and bounded; thus $E^x\int_0^T g^2(X_s)\,ds<\infty$ for every $T$, i.e.\ $g(X)\in\cL^*(W^{(j)})$. With $g=\sigma_{ij}$ this shows, by \bk{Proposition 3.2.10}, that each $M^{(i)}$ is in $\cM_2^c$. By the bilinearity of the cross-variation (\bk{Problem 1.5.7}) and \bk{Proposition 3.2.17},
\begin{equation}\label{eq:5-1-02:qv}
  \qv{M^{(i)},M^{(k)}}_t
  =\sum_{j=1}^r\sum_{l=1}^r\int_0^t\sigma_{ij}(X_s)\sigma_{kl}(X_s)\,d\qv{W^{(j)},W^{(l)}}_s
  =\int_0^t a_{ik}(X_s)\,ds;\qquad 0\le t<\infty,
\end{equation}
$P^x$-a.s. Since $M^{(i)}M^{(k)}-\qv{M^{(i)},M^{(k)}}$ is a martingale vanishing at $t=0$ (\bk{Definition 1.5.5}) and $M^{(i)}$ is a martingale with $M^{(i)}_0=0$, we have for every $t\ge0$
\begin{equation}\label{eq:5-1-02:mom}
  E^x M^{(i)}_t=0,\qquad
  E^x\bigl[M^{(i)}_tM^{(k)}_t\bigr]=E^x\int_0^t a_{ik}(X_s)\,ds,\qquad
  E^x\bigl[(M^{(i)}_t)^2\bigr]\le rK^2t ,
\end{equation}
the last by \eqref{eq:5-1-02:abound}.

\emph{Step 2: a limit.} We claim that for every bounded, continuous $g\colon\R^d\to\R$,
\begin{equation}\label{eq:5-1-02:lim}
  \lim_{t\downarrow0}\frac1t\,E^x\int_0^t g(X_s)\,ds=g(x).
\end{equation}
Indeed, for every $\omega$ the function $s\mapsto g(X_s(\omega))$ is continuous, so the Riemann integral $\int_0^t g(X_s)\,ds$ is a limit of Riemann sums, hence a random variable, and by the fundamental theorem of calculus
$\frac1t\int_0^t g(X_s(\omega))\,ds\to g(X_0(\omega))$ as $t\downarrow0$. These random variables are bounded in absolute value by $\sup\abs{g}$, and $g(X_0)=g(x)$, $P^x$-a.s.; the bounded convergence theorem (applied along an arbitrary sequence $t_n\downarrow0$) gives \eqref{eq:5-1-02:lim}.

\emph{Step 3: proof of \bk{(1.3)}.} Since $\abs{B^{(i)}_t}\le Kt$, all terms of $X^{(i)}_t-x_i=B^{(i)}_t+M^{(i)}_t$ are integrable, and \eqref{eq:5-1-02:mom} gives
$E^x[X^{(i)}_t-x_i]=E^x\int_0^t b_i(X_s)\,ds$. By \eqref{eq:5-1-02:lim} with $g=b_i$ this is $t\,b_i(x)+o(t)$ as $t\downarrow0$.

\emph{Step 4: proof of \bk{(1.4)}.} We expand
\[
  (X^{(i)}_t-x_i)(X^{(k)}_t-x_k)
  =B^{(i)}_tB^{(k)}_t+B^{(i)}_tM^{(k)}_t+M^{(i)}_tB^{(k)}_t+M^{(i)}_tM^{(k)}_t .
\]
All four terms are integrable, because the $B^{(i)}_t$ are bounded and the $M^{(i)}_t$ are square-integrable. We have $\abs{B^{(i)}_tB^{(k)}_t}\le K^2t^2$ and, by the Cauchy--Schwarz inequality and \eqref{eq:5-1-02:mom},
\[
  E^x\abs[\big]{B^{(i)}_tM^{(k)}_t}\le Kt\,\Bigl(E^x\bigl[(M^{(k)}_t)^2\bigr]\Bigr)^{1/2}\le \sqrt r\,K^2\,t^{3/2},
\]
and likewise for $M^{(i)}_tB^{(k)}_t$. Hence, by the second relation in \eqref{eq:5-1-02:mom},
\[
  \abs[\Big]{E^x\bigl[(X^{(i)}_t-x_i)(X^{(k)}_t-x_k)\bigr]-E^x\int_0^t a_{ik}(X_s)\,ds}
  \le K^2t^2+2\sqrt r\,K^2t^{3/2}=o(t)\quad\text{as }t\downarrow0,
\]
while $E^x\int_0^t a_{ik}(X_s)\,ds=t\,a_{ik}(x)+o(t)$ by \eqref{eq:5-1-02:lim} with $g=a_{ik}$. This is \bk{(1.4)}.

\emph{Step 5: proof of \bk{(1.1)}.} Let $f\in C^2(\R^d)$ be bounded with bounded first- and second-order derivatives. The process $X=X_0+M+B$ is of the form required in \bk{Theorem 3.3.6}: the $M^{(i)}$ are in $\cM_2^c\subseteq\cM^{c,\loc}$, and the $B^{(i)}$ are continuous, adapted (as limits of Riemann sums of $\F_t$-measurable variables), of bounded variation on bounded intervals, with $B^{(i)}_0=0$. It\^o's rule and \eqref{eq:5-1-02:qv} yield, $P^x$-a.s.,
\begin{align*}
  f(X_t)-f(X_0)
  &=\sum_{i=1}^d\int_0^t\frac{\partial f}{\partial x_i}(X_s)\,b_i(X_s)\,ds
   +\frac12\sum_{i=1}^d\sum_{k=1}^d\int_0^t\frac{\partial^2f}{\partial x_i\partial x_k}(X_s)\,a_{ik}(X_s)\,ds\\
  &\qquad+\sum_{i=1}^dN^{(i)}_t\\
  &=\int_0^t(\cA f)(X_s)\,ds+\sum_{i=1}^dN^{(i)}_t;\qquad 0\le t<\infty,
\end{align*}
where $\cA f$ is given by \bk{(1.2)} and $N^{(i)}_t\defeq\int_0^t\frac{\partial f}{\partial x_i}(X_s)\,dM^{(i)}_s$. The integrand $\frac{\partial f}{\partial x_i}(X)$ is bounded, continuous and adapted, hence progressively measurable, and by \eqref{eq:5-1-02:qv}, \eqref{eq:5-1-02:abound}
\[
  E^x\int_0^T\Bigl(\frac{\partial f}{\partial x_i}(X_s)\Bigr)^2d\qv{M^{(i)}}_s
  \le rK^2\,T\,\sup_{y\in\R^d}\abs[\Big]{\frac{\partial f}{\partial x_i}(y)}^2<\infty;\qquad 0\le T<\infty .
\]
Thus $\frac{\partial f}{\partial x_i}(X)\in\cL^*(M^{(i)})$, and $N^{(i)}$ is a martingale in $\cM_2^c$ by \bk{Proposition 3.2.10}; in particular $E^xN^{(i)}_t=0$. The function $\cA f$ is bounded and continuous, because $b_i$, $a_{ik}$ and the derivatives of $f$ of orders one and two are. Taking expectations in the display above (all terms are bounded or integrable) and recalling $X_0=x$, $P^x$-a.s., we obtain
\[
  E^xf(X_t)-f(x)=E^x\int_0^t(\cA f)(X_s)\,ds;\qquad 0\le t<\infty ,
\]
and \eqref{eq:5-1-02:lim} with $g=\cA f$ gives
$\lim_{t\downarrow0}\frac1t\bigl[E^xf(X_t)-f(x)\bigr]=(\cA f)(x)$, which is \bk{(1.1)}.

\rmk{The book obtains \bk{(1.3)}, \bk{(1.4)} heuristically by inserting the unbounded test functions $f_i$, $f_{ik}$ in \bk{(1.1)}. They are not covered by the class of $f$ in the statement, which is why Steps 3 and 4 argue directly from \bk{(1.9)}.}
\end{kex}

\solsection{5.2}{Strong Solutions}

\begin{kex}{Exercise}{2.17}{}
Throughout, $x^{1/3}$ is the real cube root of $x\in\R$ and $x^{2/3}\defeq (x^{1/3})^2=\abs{x}^{2/3}$, so that
$b(x)=3x^{1/3}$ and $\sigma(x)=3x^{2/3}$ are continuous on $\R$ and $(w^3)^{1/3}=w$, $(w^3)^{2/3}=w^2$ for
$w\in\R$. The setting is that of \bk{Definition 2.1} with $d=r=1$ and $\xi\equiv 0$: $W$ is a standard
Brownian motion on $(\Omega,\F,P)$ with continuous paths, and $\{\F_t\}$ is the augmentation \bk{(2.3)} of
$\{\F^W_t\}$, which satisfies the usual conditions and for which $W$ is still a Brownian motion. We use
the convention $\inf\emptyset=\infty$; thus $\beta_\infty=\infty$ and $X^{(\infty)}\equiv 0$.

\emph{Step 1: $\beta_\theta$ is a stopping time, and a formula for $X^{(\theta)}$.}
Fix $\theta\in[0,\infty]$ and write $\beta=\beta_\theta$. For $t<\theta$ we have $\{\beta\le t\}=\emptyset$.
For $\theta\le t<\infty$, the continuous function $s\mapsto\abs{W_s(\omega)}$ attains its minimum on the
compact interval $[\theta,t]$, and this minimum equals the infimum over the countable dense set
$D_t\defeq(\Q\cap[\theta,t])\cup\{t\}$; consequently
\[
\{\beta\le t\}=\{\exists\, s\in[\theta,t]:\ W_s=0\}=\Bigl\{\inf_{s\in D_t}\abs{W_s}=0\Bigr\}\in\F^W_t\subseteq\F_t .
\]
Hence $\beta$ is a stopping time of $\{\F_t\}$. Moreover, on $\{\beta<\infty\}$ the set
$\{s\ge\theta;\ W_s=0\}$ is nonempty and closed, so it contains its infimum:
\begin{equation}\label{eq:5-2-17:a}
W_\beta=0\quad\text{on }\{\beta<\infty\}.
\end{equation}
It follows that, for \emph{every} $\omega\in\Omega$ and $0\le t<\infty$,
\begin{equation}\label{eq:5-2-17:b}
X^{(\theta)}_t=W_t^3-W_{t\wedge\beta}^3 .
\end{equation}
Indeed, if $t<\beta$ both sides vanish, and if $t\ge\beta$ then $\beta<\infty$, $W_{t\wedge\beta}=W_\beta=0$
by \eqref{eq:5-2-17:a}, and both sides equal $W_t^3$.

\emph{Step 2: conditions \textup{(i)--(iii)} of \bk{Definition 2.1}.}
By \eqref{eq:5-2-17:b}, every path of $X^{(\theta)}$ is continuous, and $X^{(\theta)}_0=0=\xi$. The process
$W$ is continuous and adapted, hence progressively measurable (\bk{Proposition 1.1.13}), so the stopped
process $\{W_{t\wedge\beta}\}$ is adapted to $\{\F_t\}$ (\bk{Proposition 1.2.18}); by \eqref{eq:5-2-17:b},
$X^{(\theta)}$ is adapted to $\{\F_t\}$. Finally, $s\mapsto \abs{b(X^{(\theta)}_s)}+\sigma^2(X^{(\theta)}_s)$
is continuous for every $\omega$, hence bounded on each $[0,t]$, and (iii) holds.

\emph{Step 3: the equation.}
From the definition of $X^{(\theta)}$ we have, for all $\omega$ and $s\ge 0$,
$(X^{(\theta)}_s)^{1/3}=W_s\ind_{\{\beta\le s\}}$ and $(X^{(\theta)}_s)^{2/3}=W_s^2\ind_{\{\beta\le s\}}$;
since $W_s=0$ when $s=\beta$ by \eqref{eq:5-2-17:a}, these can be rewritten as
\begin{equation}\label{eq:5-2-17:c}
(X^{(\theta)}_s)^{1/3}=W_s\,\ind_{\{\beta< s\}},\qquad
(X^{(\theta)}_s)^{2/3}=W_s^2-W_s^2\,\ind_{\{s\le\beta\}} .
\end{equation}
The process $W^2$ is continuous and adapted, hence progressively measurable, and
$E\int_0^T W_s^4\,ds=\int_0^T 3s^2\,ds=T^3<\infty$ for every $T<\infty$; thus $W^2\in\cL^*$ (relative to
$M=W$), and $M_t\defeq\int_0^t W_s^2\,dW_s$ is a well-defined element of $\cM_2^c$. The same is true of
$W^2\ind_{\{\cdot\le\beta\}}$ (a product of $W^2$ with a left-continuous adapted process, dominated by
$W^2$), and by \bk{(3.2.24)} of \bk{Proposition 3.2.10},
\[
\int_0^t W_s^2\,\ind_{\{s\le\beta\}}\,dW_s=M_{t\wedge\beta},\qquad 0\le t<\infty,\ \as
\]
(both sides are continuous in $t$, so the exceptional null set can be chosen independently of $t$).
By \eqref{eq:5-2-17:c} and the linearity of the stochastic integral we obtain, a.s.,
\begin{equation}\label{eq:5-2-17:d}
3\int_0^t (X^{(\theta)}_s)^{1/3}\,ds+3\int_0^t (X^{(\theta)}_s)^{2/3}\,dW_s
=3\int_{t\wedge\beta}^t W_s\,ds+3\bigl(M_t-M_{t\wedge\beta}\bigr),\qquad 0\le t<\infty .
\end{equation}
On the other hand, It\^o's rule (\bk{Theorem 3.3.3}) applied to $f(x)=x^3$ and the continuous
semimartingale $W$, for which $\qv{W}_t=t$, gives
\[
W_t^3=3M_t+3\int_0^t W_s\,ds,\qquad 0\le t<\infty,\ \as
\]
On the event of probability one where this identity holds for all $t$, we may evaluate it at $t$ and at
$t\wedge\beta(\omega)$ and subtract:
\[
W_t^3-W_{t\wedge\beta}^3=3\bigl(M_t-M_{t\wedge\beta}\bigr)+3\int_{t\wedge\beta}^t W_s\,ds,\qquad 0\le t<\infty .
\]
Comparing with \eqref{eq:5-2-17:d} and \eqref{eq:5-2-17:b}, we conclude that, almost surely,
\[
X^{(\theta)}_t=3\int_0^t (X^{(\theta)}_s)^{1/3}\,ds+3\int_0^t (X^{(\theta)}_s)^{2/3}\,dW_s,\qquad 0\le t<\infty .
\]
This is condition (iv) of \bk{Definition 2.1}; together with Step 2, $X^{(\theta)}$ is a strong solution
of the equation, with respect to $W$ and the initial condition $\xi=0$, for every $\theta\in[0,\infty]$.

\emph{Step 4: the solutions are pairwise distinguishable.}
Let $0\le\theta_1<\theta_2\le\infty$, and put $t\defeq\theta_2$ if $\theta_2<\infty$, $t\defeq\theta_1+1$ if
$\theta_2=\infty$; then $\theta_1<t<\infty$. We claim that
\begin{equation}\label{eq:5-2-17:e}
P\bigl[X^{(\theta_2)}_t=0\bigr]=1,\qquad P\bigl[X^{(\theta_1)}_t\neq 0\bigr]>0 .
\end{equation}
For the first assertion: if $\theta_2=\infty$ then $X^{(\theta_2)}\equiv 0$; if $\theta_2<\infty$ then
$t=\theta_2>0$, so $P[W_{\theta_2}\ne0]=1$, and on $\{W_{\theta_2}\neq0\}$ we have
$\beta_{\theta_2}>\theta_2=t$ by \eqref{eq:5-2-17:a}, whence $X^{(\theta_2)}_t=0$.
For the second assertion: if $\theta_1=0$ then $\beta_{\theta_1}=0$, so $X^{(\theta_1)}_t=W_t^3\neq0$ with
probability one, because $t>0$. If $\theta_1>0$, consider
$A\defeq\{0<W_{\theta_1}<1\}\cap\{W_t-W_{\theta_1}<-1\}$; the two random variables involved are independent
and nondegenerate normal, so $P(A)>0$. On $A$ we have $W_{\theta_1}>0>W_t$, so by the intermediate value
theorem the path has a zero in $(\theta_1,t)$; thus $\beta_{\theta_1}<t$ and
$X^{(\theta_1)}_t=W_t^3\neq0$ on $A$. This proves \eqref{eq:5-2-17:e}.

By \eqref{eq:5-2-17:e}, $P[X^{(\theta_1)}_t\neq X^{(\theta_2)}_t]>0$, so $X^{(\theta_1)}$ and
$X^{(\theta_2)}$ are not modifications of one another, let alone indistinguishable. Therefore
$\{X^{(\theta)};\ 0\le\theta\le\infty\}$ is an uncountable family of pairwise different strong solutions,
on the same space, with the same Brownian motion and the same initial condition; strong uniqueness
fails.

\emph{Step 5: the coefficients.}
For $a,b\ge0$ and $0<p\le1$ we have $(a+b)^p\le a^p+b^p$ (for $a+b>0$, the numbers $u=a/(a+b)$ and
$v=b/(a+b)$ lie in $[0,1]$, so $u^p+v^p\ge u+v=1$). Applying this with $p=\tfrac23$ to
$\abs{x}\le\abs{y}+\bigl|\abs{x}-\abs{y}\bigr|$ and using symmetry in $(x,y)$, we get
\[
\abs{\sigma(x)-\sigma(y)}=3\,\bigl|\abs{x}^{2/3}-\abs{y}^{2/3}\bigr|
\le 3\,\bigl|\abs{x}-\abs{y}\bigr|^{2/3}\le 3\abs{x-y}^{2/3}=h(\abs{x-y}),\qquad x,y\in\R,
\]
with $h(u)\defeq3u^{2/3}$. This $h\colon[0,\infty)\to[0,\infty)$ is strictly increasing, $h(0)=0$, and
\[
\int_{(0,\eps)}h^{-2}(u)\,du=\frac19\int_0^\eps u^{-4/3}\,du=\infty\qquad\text{for every }\eps>0 .
\]
So $\sigma$ satisfies \bk{(2.24)} with a function $h$ as required in \bk{Proposition 2.13}, including
\bk{(2.25)} (compare \bk{Example 2.14} with $\alpha=\tfrac23\ge\tfrac12$).

Suppose, on the other hand, that $b$ satisfied the condition \bk{(2.23)$'$} of \bk{Remark 2.16} for some
function $\kappa\colon[0,\infty)\to[0,\infty)$ with $\int_{(0,\eps)}du/\kappa(u)=\infty$ for every $\eps>0$.
Taking $x=u>0$ and $y=0$ in \bk{(2.23)$'$} gives $3u^{1/3}\le\kappa(u)$, whence
\[
\int_{(0,\eps)}\frac{du}{\kappa(u)}\le\frac13\int_0^\eps u^{-1/3}\,du=\frac{\eps^{2/3}}{2}<\infty ,
\]
a contradiction. Hence $b(x)=3x^{1/3}$ satisfies neither \bk{(2.23)$'$} for any admissible $\kappa$ nor,
a fortiori, the Lipschitz condition \bk{(2.23)} (which is the case $\kappa(u)=Ku$). The failure of strong
uniqueness established in Step 4 is therefore consistent with \bk{Proposition 2.13} and
\bk{Remark 2.16}: it is caused by the drift, not by the dispersion coefficient.
\end{kex}

\begin{kex}{Exercise}{2.19}{}
We keep the setting and the conditions (i)--(iii) of \bk{Proposition 2.18}, assume (iv$'$), and write
$\Delta_t\defeq X^{(1)}_t-X^{(2)}_t$. The function interpolating between $b_1$ and $b_2$ will be
denoted by $\beta_m$ rather than $b_m$, to avoid a clash with $b_1,b_2$. We cannot simply apply
\bk{Proposition 2.18} to a pair of equations with drifts $b_1,\beta_m$ and then $\beta_m,b_2$, since no
solution of the equation with drift $\beta_m$ is at our disposal; instead we repeat the proof of
\bk{Proposition 2.18} after localization, $\beta_m$ providing the estimate for which condition (v) was used.

\emph{Step 1: construction of $\beta_m$.} Fix an integer $m\ge 1$ and let $Q_m\defeq[0,m]\times[-m,m]$.
The function $f\defeq\tfrac12(b_1+b_2)$ is continuous by (i), and so is $b_2-b_1$, which is strictly
positive by (iv$'$); since $Q_m$ is compact,
\[
\delta\defeq\min_{Q_m}\tfrac12(b_2-b_1)>0,\qquad M\defeq\max_{Q_m}\abs{f}<\infty ,
\]
and $f$ is uniformly continuous on $Q_m$: there is $\eta>0$ such that $\abs{f(t,x)-f(t,y)}\le\delta$
whenever $(t,x),(t,y)\in Q_m$ and $\abs{x-y}\le\eta$. Put $L_m\defeq 2M/\eta$ and
\[
\beta_m(t,x)\defeq\inf_{\abs{y}\le m}\bigl[f(t,y)+L_m\abs{x-y}\bigr];\qquad 0\le t\le m,\ \abs{x}\le m .
\]
The infimum is finite, being bounded below by $-M$. We claim that, for $0\le t\le m$ and
$\abs{x},\abs{y}\le m$,
\begin{equation}\label{eq:5-2-19:a}
b_1(t,x)\le\beta_m(t,x)\le b_2(t,x),\qquad \abs{\beta_m(t,x)-\beta_m(t,y)}\le L_m\abs{x-y} .
\end{equation}
Indeed, for fixed $t$ and $y$ the function $x\mapsto f(t,y)+L_m\abs{x-y}$ is Lipschitz with constant
$L_m$, and a finite infimum of such functions has the same property; this is the second assertion.
Taking $y=x$ in the infimum gives $\beta_m(t,x)\le f(t,x)\le b_2(t,x)$. For the lower bound let
$\abs{y}\le m$. If $\abs{x-y}\le\eta$, then $f(t,y)+L_m\abs{x-y}\ge f(t,y)\ge f(t,x)-\delta$; if
$\abs{x-y}>\eta$, then $f(t,y)+L_m\abs{x-y}\ge -M+L_m\eta=M\ge f(t,x)$. Hence
$\beta_m(t,x)\ge f(t,x)-\delta\ge f(t,x)-\tfrac12\bigl(b_2(t,x)-b_1(t,x)\bigr)=b_1(t,x)$.

From \eqref{eq:5-2-19:a} we obtain the only property of $\beta_m$ that will be used:
\begin{equation}\label{eq:5-2-19:b}
b_1(s,x_1)-b_2(s,x_2)\le\beta_m(s,x_1)-\beta_m(s,x_2)\le L_m(x_1-x_2);\qquad
0\le s\le m,\ -m\le x_2<x_1\le m .
\end{equation}

\emph{Step 2: localization.} The process $\max\bigl(\abs{X^{(1)}},\abs{X^{(2)}}\bigr)$ is continuous
and adapted, so by \bk{Problem 1.2.7}
\[
T_m\defeq m\wedge\inf\bigl\{t\ge0;\ \abs{X^{(1)}_t}\ge m\ \text{or}\ \abs{X^{(2)}_t}\ge m\bigr\}
\]
is a stopping time of $\{\F_t\}$, and $T_m\uparrow\infty$ as $m\to\infty$, because continuous paths are
bounded on bounded intervals. By continuity of the paths,
\begin{equation}\label{eq:5-2-19:c}
\abs{X^{(1)}_s}\le m,\quad \abs{X^{(2)}_s}\le m,\quad s\le m\qquad\text{for } 0\le s\le T_m,\ \text{on } \{T_m>0\},
\end{equation}
whereas on $\{T_m=0\}$ we have $\Delta_{t\wedge T_m}=\Delta_0\le0$ a.s.\ by (iii). In particular
\begin{equation}\label{eq:5-2-19:d}
0\le\Delta^+_{t\wedge T_m}\le 2m\quad\text{a.s.};\qquad 0\le t<\infty .
\end{equation}
Let $C_m\defeq\max_{Q_m}\bigl(\abs{\sigma}+\abs{b_1}+\abs{b_2}\bigr)$, finite by (i).

\emph{Step 3: the It\^o rule.} Let $\{a_n\}$, $\rho_n$ and $\psi_n$ be as in the proof of
\bk{Proposition 2.13}, constructed from the function $h$ of condition (ii), and let
$\varphi_n\defeq\psi_n\ind_{(0,\infty)}$ as in the proof of \bk{Proposition 2.18}. Since $\rho_n$ is
continuous with support in $(a_n,a_{n-1})$ and $\int\rho_n(u)\,du=1$, we have
$\varphi_n(x)=\int_0^{x^+}\!\int_0^y\rho_n(u)\,du\,dy$, so that $\varphi_n$ vanishes on $(-\infty,a_n]$,
is of class $C^2(\R)$, and
\begin{equation}\label{eq:5-2-19:e}
\begin{gathered}
\varphi_n'=0\ \text{on } (-\infty,0],\qquad 0\le\varphi_n'\le1,\qquad
\varphi_n''(x)=\rho_n(x)\ind_{(0,\infty)}(x),\\
(x-a_{n-1})^+\le\varphi_n(x)\le x^+ ;
\end{gathered}
\end{equation}
the last inequalities hold because $\varphi_n'(y)=1$ for $y\ge a_{n-1}$. From \bk{(2.31)},
\[
\Delta_t=\Delta_0+\int_0^t\bigl[b_1(s,X^{(1)}_s)-b_2(s,X^{(2)}_s)\bigr]ds
+\int_0^t\bigl[\sigma(s,X^{(1)}_s)-\sigma(s,X^{(2)}_s)\bigr]dW_s;\qquad 0\le t<\infty,
\]
a.s., which exhibits $\Delta$ as a continuous semimartingale (the integrands are continuous, adapted
processes). The It\^o rule (\bk{Theorem 3.3.3}), evaluated at the time $t\wedge T_m$, gives a.s.
\begin{align}
\varphi_n(\Delta_{t\wedge T_m})={}&\varphi_n(\Delta_0)
+\int_0^{t\wedge T_m}\varphi_n'(\Delta_s)\bigl[b_1(s,X^{(1)}_s)-b_2(s,X^{(2)}_s)\bigr]ds \notag\\
&+\frac12\int_0^{t\wedge T_m}\varphi_n''(\Delta_s)\bigl[\sigma(s,X^{(1)}_s)-\sigma(s,X^{(2)}_s)\bigr]^2ds \notag\\
&+\int_0^{t}\ind_{\{s\le T_m\}}\varphi_n'(\Delta_s)\bigl[\sigma(s,X^{(1)}_s)-\sigma(s,X^{(2)}_s)\bigr]dW_s .
\label{eq:5-2-19:f}
\end{align}
We examine the four terms on the right.

\textbf{(a)} $\varphi_n(\Delta_0)=0$ a.s., by (iii) and \eqref{eq:5-2-19:e}.

\textbf{(b)} Let $0<s\le T_m$ (so that \eqref{eq:5-2-19:c} is in force). If $\Delta_s\le0$, the first
integrand vanishes, by \eqref{eq:5-2-19:e}. If $\Delta_s>0$, then \eqref{eq:5-2-19:b} with
$x_1=X^{(1)}_s$, $x_2=X^{(2)}_s$, together with $0\le\varphi_n'\le1$, shows that the integrand is at most
$L_m\Delta_s$ when $b_1(s,X^{(1)}_s)-b_2(s,X^{(2)}_s)\ge0$, and it is nonpositive otherwise. In all
cases the integrand is bounded above by $L_m\Delta_s^+$, and in absolute value by $2C_m$. Therefore
\[
\int_0^{t\wedge T_m}\varphi_n'(\Delta_s)\bigl[b_1(s,X^{(1)}_s)-b_2(s,X^{(2)}_s)\bigr]ds
\le L_m\int_0^{t\wedge T_m}\Delta_s^+\,ds\le L_m\int_0^t\Delta^+_{s\wedge T_m}\,ds .
\]

\textbf{(c)} By (ii) and \eqref{eq:5-2-19:e}, and because $0\le\rho_n(x)\le 2/\bigl(nh^2(x)\bigr)$ for $x>0$,
\[
0\le\varphi_n''(\Delta_s)\bigl[\sigma(s,X^{(1)}_s)-\sigma(s,X^{(2)}_s)\bigr]^2
\le\rho_n(\Delta_s)\,h^2(\Delta_s)\,\ind_{\{\Delta_s>0\}}\le\frac2n ,
\]
so the third term of \eqref{eq:5-2-19:f} lies in $[0,t/n]$.

\textbf{(d)} The integrand $\ind_{\{s\le T_m\}}\varphi_n'(\Delta_s)\bigl[\sigma(s,X^{(1)}_s)-\sigma(s,X^{(2)}_s)\bigr]$
is progressively measurable and, by \eqref{eq:5-2-19:c}, bounded in absolute value by $2C_m$ for every
$s>0$. Hence it belongs to $\cL^*$ for $W$, the stochastic integral in \eqref{eq:5-2-19:f} is a
square-integrable martingale starting at zero, and its expectation is zero.

The left-hand side of \eqref{eq:5-2-19:f} is bounded by $2m$, by \eqref{eq:5-2-19:d} and
\eqref{eq:5-2-19:e}, and the terms in (b), (c) are bounded; we may thus take expectations in
\eqref{eq:5-2-19:f} and obtain, using the Fubini--Tonelli theorem for the nonnegative, jointly
measurable integrand $\Delta^+_{s\wedge T_m}$,
\begin{equation}\label{eq:5-2-19:g}
E\varphi_n(\Delta_{t\wedge T_m})\le L_m\int_0^tE\bigl(\Delta^+_{s\wedge T_m}\bigr)ds+\frac tn;\qquad
0\le t<\infty,\ n\ge1 .
\end{equation}

\emph{Step 4: conclusion.} Let $g(t)\defeq E\bigl(\Delta^+_{t\wedge T_m}\bigr)$. By
\eqref{eq:5-2-19:d} and the bounded convergence theorem, $g$ is a continuous function with values in
$[0,2m]$. By \eqref{eq:5-2-19:e}, $\varphi_n(x)\to x^+$ as $n\to\infty$ for every $x\in\R$ (recall
$a_n\downarrow0$), with $0\le\varphi_n(x)\le x^+$; a further application of bounded convergence in
\eqref{eq:5-2-19:g} yields
\[
0\le g(t)\le L_m\int_0^tg(s)\,ds;\qquad 0\le t<\infty .
\]
The Gronwall inequality (\bk{Problem 2.7}, with $\alpha\equiv0$ and $\beta=L_m$) gives $g\equiv0$, that is,
$X^{(1)}_{t\wedge T_m}\le X^{(2)}_{t\wedge T_m}$ a.s., for every $t\ge0$ and every $m\ge1$. Letting
$m\to\infty$ and using $T_m\uparrow\infty$, we get $P[X^{(1)}_t\le X^{(2)}_t]=1$ for every $t\ge0$.
Consequently, the event $\{X^{(1)}_t\le X^{(2)}_t,\ \forall\,t\in\Q\cap[0,\infty)\}$ has probability
one, and by the continuity of the paths of $X^{(1)}$ and $X^{(2)}$ it coincides with the event in
\bk{(2.32)}. This proves \bk{(2.32)}.

\rmk{The strictness in \textup{(iv$'$)} was used only in Step 1, to make room ($\delta>0$) for a
function which is Lipschitz in $x$ on $Q_m$ between $b_1$ and $b_2$. Only this Lipschitz property in the
space variable, as in \bk{(2.23)}, is needed; the function $\beta_m$ constructed above is in fact also
continuous in $(t,x)$, but this plays no role, since $\beta_m$ enters only through the pointwise
inequality \eqref{eq:5-2-19:b}.}
\end{kex}

\begin{kex}{Exercise}{2.20}{}
We work in the setting of \bk{Definition 2.1} with $d=r=1$: $W$ is a Brownian motion on $(\Omega,\F,P)$, $\xi$ is a real random variable independent of $\F^W_\infty$ (no moment condition is imposed on $\xi$), and $\{\F_t\}$ is the augmented filtration \bk{(2.3)}, which satisfies the usual conditions and with respect to which $W$ is a Brownian motion. We have to show that there is a strong solution of \bk{(2.4)$''$} relative to $W$ with initial condition $\xi$, and that strong uniqueness holds in the sense of \bk{Definition 2.3}.

\emph{Step 1. The change of scale.} Since $\sigma$ is positive and of class $C^2$, the function $1/\sigma$ is positive and of class $C^2$, so $f(x)=\int_0^x du/\sigma(u)$ is of class $C^3$ and strictly increasing, with
\[
f'=\frac1\sigma,\qquad f''=-\frac{\sigma'}{\sigma^2}.
\]
The assumption that $1/\sigma$ is integrable neither at $+\infty$ nor at $-\infty$ says precisely that $f(x)\to\pm\infty$ as $x\to\pm\infty$; hence $f$ is a bijection of $\R$ onto $\R$. Let $g\defeq f^{-1}\colon\R\to\R$. As $f'>0$, the inverse function theorem shows that $g$ is differentiable with $g'=1/f'(g)=\sigma(g)$; the right-hand side is of class $C^1$, so $g$ is of class $C^2$ and
\begin{equation}\label{eq:5-2-20:g}
g'=\sigma\circ g,\qquad g''=(\sigma'\circ g)\,g'=(\sigma\sigma')\circ g .
\end{equation}
We define
\[
\beta(y)\defeq\Bigl(\frac{b}{\sigma}-\frac12\,\sigma'\Bigr)\bigl(g(y)\bigr);\qquad y\in\R .
\]
The function $b/\sigma-\tfrac12\sigma'$ is of class $C^1$, hence so is $\beta$, and by \eqref{eq:5-2-20:g}
\[
\beta'(y)=\Bigl(\frac{b'\sigma-b\sigma'}{\sigma^2}-\frac12\,\sigma''\Bigr)\bigl(g(y)\bigr)\,\sigma\bigl(g(y)\bigr)
=\Bigl(b'-\frac12\,\sigma\sigma''-\frac{b\sigma'}{\sigma}\Bigr)\bigl(g(y)\bigr).
\]
By hypothesis there is a constant $K<\infty$ with $\abs{\beta'}\le K$ on $\R$, and the mean value theorem gives
\begin{equation}\label{eq:5-2-20:lip}
\abs{\beta(y)-\beta(z)}\le K\abs{y-z};\qquad y,z\in\R .
\end{equation}

\emph{Step 2. The equation with unit dispersion.} Let $\eta\defeq f(\xi)$, an $\F_0$-measurable random variable. We construct a continuous, $\{\F_t\}$-adapted process $Y$ such that, for \emph{every} $\omega\in\Omega$,
\begin{equation}\label{eq:5-2-20:Y}
Y_t=\eta+\int_0^t\beta(Y_s)\,ds+W_t;\qquad 0\le t<\infty .
\end{equation}
This is the path-by-path Picard--Lindel\"of argument indicated in the book after \bk{(2.34)}; we cannot simply quote \bk{Theorem 2.9}, because $\eta$ need not be square-integrable. Set $Y^{(0)}_t\defeq\eta$ and, for $k\ge0$,
\[
Y^{(k+1)}_t\defeq\eta+\int_0^t\beta\bigl(Y^{(k)}_s\bigr)\,ds+W_t;\qquad 0\le t<\infty .
\]
By induction on $k$, each $Y^{(k)}$ is a continuous, $\{\F_t\}$-adapted process: if $Y^{(k)}$ is, then $s\mapsto\beta(Y^{(k)}_s(\omega))$ is continuous for every $\omega$, so the integral $\int_0^t\beta(Y^{(k)}_s)\,ds$ is well defined, continuous in $t$, and is the limit as $n\to\infty$ of the Riemann sums $\frac tn\sum_{j=1}^n\beta(Y^{(k)}_{jt/n})$, which are $\F_t$-measurable; and $W$ is continuous and adapted.

Fix $\omega\in\Omega$ and $T\in(0,\infty)$, and let $D^{(k)}_t\defeq\max_{0\le s\le t}\abs{Y^{(k)}_s-Y^{(k-1)}_s}$ for $k\ge1$, $0\le t\le T$. Then $D^{(1)}_T=\max_{0\le s\le T}\abs{\beta(\eta)s+W_s}<\infty$, and for $k\ge1$, by \eqref{eq:5-2-20:lip},
\[
\abs[\big]{Y^{(k+1)}_s-Y^{(k)}_s}\le\int_0^s\abs[\big]{\beta\bigl(Y^{(k)}_u\bigr)-\beta\bigl(Y^{(k-1)}_u\bigr)}\,du\le K\int_0^sD^{(k)}_u\,du,
\]
whence $D^{(k+1)}_t\le K\int_0^tD^{(k)}_u\,du$ for $0\le t\le T$, the right-hand side being nondecreasing in $s$. Since $D^{(1)}_t\le D^{(1)}_T$, induction yields
\[
D^{(k)}_t\le D^{(1)}_T\,\frac{(Kt)^{k-1}}{(k-1)!};\qquad 0\le t\le T,\ k\ge1 .
\]
Therefore $\sum_{k\ge1}D^{(k)}_T\le D^{(1)}_Te^{KT}<\infty$, and the functions $Y^{(k)}_\cdot(\omega)$ converge uniformly on $[0,T]$. As $T$ is arbitrary, $Y_t(\omega)\defeq\lim_{k\to\infty}Y^{(k)}_t(\omega)$ exists for all $t\ge0$ and $\omega\in\Omega$, the convergence being uniform on compact intervals; thus $Y$ has continuous paths, and $Y_t$ is $\F_t$-measurable as a pointwise limit of $\F_t$-measurable random variables. By \eqref{eq:5-2-20:lip}, $\beta(Y^{(k)}_s)\to\beta(Y_s)$ uniformly on $[0,T]$, so we may let $k\to\infty$ in the recursion and obtain \eqref{eq:5-2-20:Y}.

\emph{Step 3. Existence.} Put $X_t\defeq g(Y_t)$. This process is continuous and $\{\F_t\}$-adapted, and $X_0=g(f(\xi))=\xi$. By \eqref{eq:5-2-20:Y}, $Y$ is a continuous semimartingale: $Y=Y_0+B+W$ with $Y_0=\eta$ $\F_0$-measurable, $B_t=\int_0^t\beta(Y_s)\,ds$ continuous, adapted, of bounded variation on compact intervals, $B_0=0$, and $W\in\cM^{c,\loc}$ with $\qv{W}_t=t$. Since $g$ is of class $C^2$, It\^o's rule (\bk{Theorem 3.3.3}) and \eqref{eq:5-2-20:g} give, almost surely, for all $0\le t<\infty$,
\begin{align*}
X_t&=g(\eta)+\int_0^tg'(Y_s)\beta(Y_s)\,ds+\int_0^tg'(Y_s)\,dW_s+\frac12\int_0^tg''(Y_s)\,ds\\
&=\xi+\int_0^t\Bigl[\sigma(X_s)\Bigl(\frac{b(X_s)}{\sigma(X_s)}-\frac12\,\sigma'(X_s)\Bigr)+\frac12\,\sigma(X_s)\sigma'(X_s)\Bigr]ds+\int_0^t\sigma(X_s)\,dW_s\\
&=\xi+\int_0^tb(X_s)\,ds+\int_0^t\sigma(X_s)\,dW_s ,
\end{align*}
where we used $g'(Y_s)=\sigma(X_s)$, $g''(Y_s)=\sigma(X_s)\sigma'(X_s)$ and $\beta(Y_s)=(b/\sigma-\tfrac12\sigma')(X_s)$. This is \bk{(2.4)$''$}. Finally, $s\mapsto\abs{b(X_s)}+\sigma^2(X_s)$ is continuous, hence bounded on $[0,t]$ for every $\omega$, so condition (iii) of \bk{Definition 2.1} holds. Thus $X$ is a strong solution of \bk{(2.4)$''$} relative to $W$ with initial condition $\xi$.

\emph{Step 4. Uniqueness.} The functions $b$ and $\sigma$ are continuously differentiable, so for every $n\ge1$ the mean value theorem gives
\[
\abs{b(x)-b(y)}+\abs{\sigma(x)-\sigma(y)}\le K_n\abs{x-y};\qquad\abs{x}\le n,\ \abs{y}\le n,
\]
with $K_n\defeq\max_{\abs{u}\le n}\abs{b'(u)}+\max_{\abs{u}\le n}\abs{\sigma'(u)}<\infty$. This is condition \bk{(2.8)} for the time-independent coefficients $b(t,x)=b(x)$, $\sigma(t,x)=\sigma(x)$, and \bk{Theorem 2.5} shows that strong uniqueness holds for \bk{(2.4)$''$}: any two strong solutions relative to the same $W$ and $\xi$ are indistinguishable. Together with Step 3, the equation has a unique strong solution.

\rmk{The computation of the hint is the converse of Step 3 and exhibits the ``change of scale'' announced before the exercise. If $X$ is any strong solution of \bk{(2.4)$''$}, then $X$ is a continuous semimartingale with $\qv{X}_t=\int_0^t\sigma^2(X_s)\,ds$, and It\^o's rule applied to the $C^2$-function $f$ gives
\[
f(X_t)=f(\xi)+\int_0^t\Bigl[\frac{b(X_s)}{\sigma(X_s)}-\frac12\,\sigma'(X_s)\Bigr]ds+\int_0^t\frac{\sigma(X_s)}{\sigma(X_s)}\,dW_s
=f(\xi)+\int_0^t\beta\bigl(f(X_s)\bigr)\,ds+W_t ,
\]
so that $Y=f(X)$ solves the equation \eqref{eq:5-2-20:Y} of the form \bk{(2.34)}, whose drift $\beta$ satisfies the global Lipschitz condition \eqref{eq:5-2-20:lip}. This yields uniqueness a second time: if $X$ and $\wt X$ are two strong solutions, then $\Delta_t\defeq f(X_t)-f(\wt X_t)$ satisfies $\abs{\Delta_t}\le K\int_0^t\abs{\Delta_s}\,ds$ almost surely for all $t$, hence $\Delta\equiv0$ by the Gronwall inequality (\bk{Problem 2.7}), and $X=\wt X$ because $f$ is injective. The hypothesis on $1/\sigma$ is what makes $f$ onto, so that $g$ is defined on all of $\R$ and $g(Y)$ makes sense for every value of $Y$; the boundedness of $b'-\tfrac12\sigma\sigma''-b\sigma'/\sigma$ is exactly the boundedness of $\beta'$.}
\end{kex}

\begin{kex}{Exercise}{2.26}{}
\emph{Setting.} Let $\beta\defeq\sup_{t,x}\abs{b(t,x)}<\infty$ and let $K$ be the constant of \bk{(2.23)}. The system is an equation of the form \bk{(2.1)} with $d=2$, $r=1$, state $z=(x,y)$, initial vector $\zeta\defeq(\xi,\eta)$, and coefficients
\[
B_\alpha(t,z)\defeq\bigl(y,\ \alpha b(t,x)-\alpha y\bigr),\qquad
\Sigma_\alpha(t,z)\defeq(0,\alpha)^{\mathsf T}.
\]
The vector $\zeta$ is independent of $\F^W_\infty$, and $\{\F_t\}$ denotes the filtration \bk{(2.3)} built from $\zeta$ and $W$; a strong solution is understood in the sense of \bk{Definition 2.1} with this filtration. By \bk{Definition 2.1.1}, $W$ is a continuous process, i.e.\ every sample path of $W$ is continuous (this is the book's usage of the term, cf.\ \bk{Proposition 1.1.13} and \bk{Remark 1.1.14}), whereas $W_0=0$ holds only almost surely. The event
\[
\Omega_0\defeq\bigl\{\abs{\xi}+\abs{\eta}<\infty,\ W_0=0\bigr\}
\]
belongs to $\sigma(\zeta,W_0)\subseteq\F_0$ and has probability one; all the path-by-path statements below are made for $\omega\in\Omega_0$, and on $\Omega\setminus\Omega_0$ we set every process constructed below equal to zero, which affects neither its adaptedness (as $\Omega_0\in\F_0$) nor any almost sure statement.

For $z=(x,y)$, $z'=(x',y')$ and $t\ge 0$,
\begin{equation}\label{eq:5-2-26:lip}
\norm{B_\alpha(t,z)-B_\alpha(t,z')}\le \abs{y-y'}+\alpha K\abs{x-x'}+\alpha\abs{y-y'}\le L_\alpha\norm{z-z'},
\qquad L_\alpha\defeq 1+\alpha+\alpha K ,
\end{equation}
and $\Sigma_\alpha$ is constant; thus the global Lipschitz condition \bk{(2.12)} holds.

\textbf{(i)} \emph{Uniqueness.} By \eqref{eq:5-2-26:lip} the coefficients satisfy \bk{(2.8)} with $K_n=L_\alpha$ for every $n$, so strong uniqueness holds by \bk{Theorem 2.5}: any two strong solutions with initial condition $\zeta$ are indistinguishable.

\emph{Existence.} Since $\norm{B_\alpha(t,z)}^2+\norm{\Sigma_\alpha}^2\le y^2+\alpha^2(\beta+\abs{y})^2+\alpha^2\le (1+2\alpha^2)(1+\beta^2)(1+\norm{z}^2)$, the linear growth condition \bk{(2.13)} holds as well, and existence of a strong solution follows from \bk{Theorem 2.9} together with \bk{Problem 2.12}, which removes the moment assumption \bk{(2.14)}. For part (ii) it is useful to have solutions which satisfy the equations on one and the same event for all $\alpha$ simultaneously, so we also carry out the path-by-path construction which the book describes after \bk{(2.34)}; it is available because the dispersion is constant. Fix $\alpha$ and define, for every $\omega\in\Omega_0$,
\[
Z^{(0)}_t\defeq\zeta,\qquad
Z^{(k+1)}_t\defeq \zeta+\int_0^t B_\alpha\bigl(s,Z^{(k)}_s\bigr)\,ds+(0,\alpha W_t);\qquad 0\le t<\infty,\ k\ge 0 .
\]
By induction, every $Z^{(k)}$ has continuous paths and is adapted to $\{\F_t\}$: if $Z^{(k)}$ has these properties, then $s\mapsto B_\alpha(s,Z^{(k)}_s)$ is continuous ($b$ being continuous), so the integral over $[0,t]$ is the limit of Riemann sums, each of which is $\F_t$-measurable, and $\zeta$ is $\F_0$-measurable. Fix $T<\infty$ and $\omega\in\Omega_0$, and set $D^{(k)}_t\defeq\max_{0\le s\le t}\norm{Z^{(k)}_s-Z^{(k-1)}_s}$. Then $D^{(1)}_T<\infty$ by continuity, and \eqref{eq:5-2-26:lip} gives $D^{(k+1)}_t\le L_\alpha\int_0^tD^{(k)}_s\,ds$ for $k\ge1$, whence by induction
\[
D^{(k)}_t\le D^{(1)}_T\,\frac{(L_\alpha t)^{k-1}}{(k-1)!};\qquad 0\le t\le T,\ k\ge 1 .
\]
The right-hand sides are summable in $k$, so $\{Z^{(k)}_\cdot(\omega)\}_{k}$ is Cauchy in the supremum norm on $[0,T]$ for every $T$, and converges uniformly on compact intervals to a continuous function $Z_\cdot(\omega)$. The process $Z=(X^{(\alpha)},Y^{(\alpha)})$ is adapted, being on $\Omega_0\in\F_0$ a pointwise limit of adapted processes and zero elsewhere, and $Z_0=\zeta$, because $Z^{(k)}_0=\zeta+(0,\alpha W_0)=\zeta$ for every $k$ on $\Omega_0$. By \eqref{eq:5-2-26:lip}, $B_\alpha(s,Z^{(k)}_s)\to B_\alpha(s,Z_s)$ uniformly in $s\in[0,T]$, so we may let $k\to\infty$ in the recursion and obtain, for every $\omega\in\Omega_0$ and every $0\le t<\infty$,
\begin{equation}\label{eq:5-2-26:sys}
X^{(\alpha)}_t=\xi+\int_0^tY^{(\alpha)}_s\,ds,\qquad
Y^{(\alpha)}_t=\eta+\alpha\int_0^t\bigl[b(s,X^{(\alpha)}_s)-Y^{(\alpha)}_s\bigr]ds+\alpha W_t .
\end{equation}
Since $\int_0^t\alpha\,dW_s=\alpha(W_t-W_0)=\alpha W_t$ a.s., this is the integral form \bk{(2.4)} of the system, so (iv) of \bk{Definition 2.1} holds; (i) and (ii) have been checked, and (iii) holds because $s\mapsto B_\alpha(s,Z_s)$ is continuous, hence integrable on $[0,t]$, and $\Sigma_\alpha$ is constant. Thus $(X^{(\alpha)},Y^{(\alpha)})$ is a strong solution, and by the uniqueness part it is the only one up to indistinguishability.

\emph{The limiting equation.} Equation \bk{(2.34)} has drift $b$ and unit dispersion; these satisfy \bk{(2.12)}, \bk{(2.13)} because $b$ is bounded and satisfies \bk{(2.23)}. Hence, by \bk{Theorem 2.9}, \bk{Problem 2.12} and \bk{Theorem 2.5} (or by the path-by-path construction above), \bk{(2.34)} has a strong solution $X$ with initial condition $\xi$ relative to $W$, unique up to indistinguishability. It is adapted to the filtration \bk{(2.3)} generated by $\xi$ and $W$, which is contained in $\{\F_t\}$; in particular $X$ does not involve $\eta$.

\textbf{(ii)} For $T<\infty$ and $\delta>0$ we write
\[
M_T\defeq\max_{0\le t\le T}\abs{W_t},\qquad
m_T(\delta)\defeq\sup\bigl\{\abs{W_t-W_s}:\ s,t\in[0,T],\ \abs{t-s}\le\delta\bigr\} ;
\]
$M_T<\infty$ and $\lim_{\delta\downarrow0}m_T(\delta)=0$ for every $\omega$, by the continuity, hence uniform continuity, of every Brownian path on $[0,T]$.

\begin{lem}
Let $\alpha>0$, let $(X^{(\alpha)},Y^{(\alpha)})$ be a strong solution of the system and $X$ a strong solution of \bk{(2.34)}, and let $\Omega_\alpha$ be the event on which \eqref{eq:5-2-26:sys} and \bk{(2.34)} hold for all $0\le t<\infty$, $\abs{\xi}+\abs{\eta}<\infty$ and $W_0=0$; then $P(\Omega_\alpha)=1$. For every $\omega\in\Omega_\alpha$, $T<\infty$ and $\delta>0$,
\begin{equation}\label{eq:5-2-26:bound}
\sup_{0\le t\le T}\abs[\big]{X^{(\alpha)}_t-X_t}\le e^{KT}\Bigl[\frac{2\abs{\eta}+\beta}{\alpha}+2m_T(\delta)+3M_Te^{-\alpha\delta}\Bigr].
\end{equation}
\end{lem}

\begin{proof}
That $P(\Omega_\alpha)=1$ follows from \bk{Definition 2.1}(ii),(iv) (for both equations; each event in question has probability one), again with $\int_0^t\alpha\,dW_s=\alpha W_t$ a.s., together with $P(\Omega_0)=1$. We fix $\omega\in\Omega_\alpha$ and argue with deterministic continuous functions, writing $X^{(\alpha)}=X^\alpha$, $Y^{(\alpha)}=Y^\alpha$.

\emph{Step 1: a formula for $Y^\alpha$.} Put $V_t\defeq Y^\alpha_t-\alpha W_t$. By the second equation in \eqref{eq:5-2-26:sys}, $V_t=\eta+\alpha\int_0^t[b(s,X^\alpha_s)-Y^\alpha_s]\,ds$ with a continuous integrand; hence $V$ is continuously differentiable, $V_0=\eta$ and
\[
\frac{d}{dt}V_t=\alpha b(t,X^\alpha_t)-\alpha V_t-\alpha^2W_t,\qquad\text{so}\qquad
\frac{d}{dt}\bigl(e^{\alpha t}V_t\bigr)=\alpha e^{\alpha t}\bigl[b(t,X^\alpha_t)-\alpha W_t\bigr].
\]
Integrating and dividing by $\alpha e^{\alpha t}$, we obtain
\begin{equation}\label{eq:5-2-26:Y}
\frac{Y^\alpha_t}{\alpha}=\frac{e^{-\alpha t}\eta}{\alpha}+\int_0^te^{-\alpha(t-s)}b(s,X^\alpha_s)\,ds+R^\alpha_t,
\qquad R^\alpha_t\defeq W_t-\alpha\int_0^te^{-\alpha(t-s)}W_s\,ds .
\end{equation}

\emph{Step 2: estimate of $R^\alpha$.} Since $\alpha\int_0^te^{-\alpha(t-s)}\,ds=1-e^{-\alpha t}$, we have
\[
R^\alpha_t=e^{-\alpha t}W_t+\alpha\int_0^te^{-\alpha(t-s)}(W_t-W_s)\,ds .
\]
Let $0\le t\le T$. If $t\le\delta$, then $e^{-\alpha t}\abs{W_t}\le\abs{W_t-W_0}\le m_T(\delta)$, since $W_0=0$ on $\Omega_\alpha$; if $t>\delta$, then $e^{-\alpha t}\abs{W_t}\le M_Te^{-\alpha\delta}$. In the integral, $\abs{W_t-W_s}\le m_T(\delta)$ for $(t-\delta)^+\le s\le t$ and $\abs{W_t-W_s}\le 2M_T$ for $0\le s<(t-\delta)^+$, while
\[
\alpha\int_{(t-\delta)^+}^te^{-\alpha(t-s)}\,ds\le 1,\qquad
\alpha\int_0^{(t-\delta)^+}e^{-\alpha(t-s)}\,ds\le e^{-\alpha\delta}.
\]
Therefore $\sup_{0\le t\le T}\abs{R^\alpha_t}\le 2m_T(\delta)+3M_Te^{-\alpha\delta}$. The second term on the right of \eqref{eq:5-2-26:Y} is at most $\beta\int_0^te^{-\alpha(t-s)}\,ds\le\beta/\alpha$ in absolute value, so that
\begin{equation}\label{eq:5-2-26:Ybd}
\rho\defeq\sup_{0\le t\le T}\frac{\abs{Y^\alpha_t-\eta}}{\alpha}\le\frac{2\abs{\eta}+\beta}{\alpha}+2m_T(\delta)+3M_Te^{-\alpha\delta}.
\end{equation}

\emph{Step 3: comparison with $X$.} Solving the second equation of \eqref{eq:5-2-26:sys} for $\int_0^tY^\alpha_s\,ds$ and substituting into the first one, we get
\[
X^\alpha_t=\xi+\int_0^tb(s,X^\alpha_s)\,ds+W_t-\frac{Y^\alpha_t-\eta}{\alpha};\qquad 0\le t<\infty,
\]
which is \bk{(2.34)} up to the last term. Subtracting \bk{(2.34)} and using \bk{(2.23)}, the continuous function $g(t)\defeq\abs{X^\alpha_t-X_t}$ satisfies
\[
0\le g(t)\le\rho+K\int_0^tg(s)\,ds;\qquad 0\le t\le T .
\]
The Gronwall inequality (\bk{Problem 2.7}, with the constant function $\rho$ in the role of $\alpha(t)$ and $K$ in the role of $\beta$) yields $g(t)\le\rho+K\int_0^t\rho e^{K(t-s)}\,ds=\rho e^{Kt}\le\rho e^{KT}$ for $0\le t\le T$, and \eqref{eq:5-2-26:bound} follows from \eqref{eq:5-2-26:Ybd}.
\end{proof}

\emph{Conclusion.} Let $(X^{(\alpha)},Y^{(\alpha)})$, $\alpha>0$, be the solutions constructed in (i), which satisfy \eqref{eq:5-2-26:sys} for every $\omega\in\Omega_0$. Then $\Omega_\alpha$ contains, for every $\alpha$, the event $\Omega^*$ of those $\omega\in\Omega_0$ for which $X$ satisfies \bk{(2.34)} for all $t$ (recall that $\Omega_0$ requires $\abs{\xi}+\abs{\eta}<\infty$ and $W_0=0$), and $P(\Omega^*)=1$. For $\omega\in\Omega^*$, $T<\infty$ and $\delta>0$ we may let $\alpha\to\infty$ in \eqref{eq:5-2-26:bound}:
\[
\limsup_{\alpha\to\infty}\ \sup_{0\le t\le T}\abs[\big]{X^{(\alpha)}_t-X_t}\le 2e^{KT}m_T(\delta),
\]
because $\eta(\omega)$ and $M_T(\omega)$ are finite. Letting $\delta\downarrow0$ we conclude that
\[
\lim_{\alpha\to\infty}\ \sup_{0\le t\le T}\abs[\big]{X^{(\alpha)}_t-X_t}=0\qquad\text{for every }\omega\in\Omega^*\text{ and every }T<\infty,
\]
which is the assertion (with an exceptional event that does not depend on $T$).

\rmk{For each $\alpha$ the strong solution is determined only up to indistinguishability, and the limit in (ii) runs over the uncountable set $(0,\infty)$; with arbitrarily chosen versions the exceptional null events $\Omega\setminus\Omega_\alpha$ could have a union which is not null. Statement (ii) is therefore to be read for the versions constructed in (i), as above. For arbitrary versions the lemma still gives the result along every sequence $\alpha_n\to\infty$: on $\bigcap_n\Omega_{\alpha_n}$, an event of probability one, \eqref{eq:5-2-26:bound} holds for all $n$, and the same passage to the limit applies.}

\rmk{The initial velocity $\eta$ enters \eqref{eq:5-2-26:bound} only through the term $2\abs{\eta}/\alpha$ and disappears in the limit, in accordance with the fact that $X$ does not depend on $\eta$. No moment condition on $\xi$, $\eta$ is needed, since the whole argument for (ii) is carried out path by path.}
\end{kex}

\begin{kex}{Exercise}{2.27}{}
No initial condition is printed; we work in the setting of \bk{Definition 2.1}: $W$ is a
one-dimensional Brownian motion, $\xi$ is a real random variable independent of $\F^W_\infty$,
$\{\F_t\}$ is the augmented filtration \bk{(2.3)}, which satisfies the usual conditions and for
which $W$ is a Brownian motion, and we impose $X_0=\xi$. With
\begin{equation*}
\sigma(x)\defeq\sqrt{1+x^2},\qquad b(x)\defeq\sqrt{1+x^2}+\tfrac12 x,\qquad x\in\R,
\end{equation*}
the equation is $dX_t=b(X_t)\,dt+\sigma(X_t)\,dW_t$. We show that
\begin{equation}\label{eq:5-2-27:sol}
X_t=\sinh\bigl(\sinh^{-1}(\xi)+t+W_t\bigr)
=\xi\cosh(t+W_t)+\sqrt{1+\xi^2}\,\sinh(t+W_t),\qquad 0\le t<\infty,
\end{equation}
is a strong solution, and that every solution with $X_0=\xi$ is indistinguishable from it. Here
$\sinh^{-1}(x)=\log\bigl(x+\sqrt{1+x^2}\bigr)$ is the inverse of the increasing bijection
$\sinh\colon\R\to\R$; the second expression in \eqref{eq:5-2-27:sol} follows from the addition
formula for $\sinh$ and $\cosh(\sinh^{-1}\xi)=\sqrt{1+\xi^2}$.

\emph{Step 1: the process \eqref{eq:5-2-27:sol} is a strong solution.}
Let $Z_t\defeq\sinh^{-1}(\xi)+t+W_t$. This is a continuous semimartingale with respect to
$\{\F_t\}$ in the sense of \bk{Definition 3.3.1}: $Z_0=\sinh^{-1}(\xi)$ is $\F_0$-measurable, its
local martingale part is $W$, with $\qv{W}_t=t$, and its part of bounded variation is $B_t=t$.
The function $\sinh$ is of class $C^2$ with $\sinh'=\cosh$ and $\sinh''=\sinh$, so It\^o's rule
(\bk{Theorem 3.3.3}) gives for $X_t=\sinh(Z_t)$, almost surely,
\begin{equation*}
X_t=\xi+\int_0^t\cosh(Z_s)\,dW_s+\int_0^t\cosh(Z_s)\,ds+\frac12\int_0^t\sinh(Z_s)\,ds,
\qquad 0\le t<\infty.
\end{equation*}
Since $\cosh>0$ and $\cosh^2=1+\sinh^2$, we have $\cosh(Z_s)=\sqrt{1+X_s^2}=\sigma(X_s)$, and
$\sinh(Z_s)=X_s$; therefore
\begin{equation*}
X_t=\xi+\int_0^t b(X_s)\,ds+\int_0^t\sigma(X_s)\,dW_s,\qquad 0\le t<\infty,
\end{equation*}
which is \bk{(2.4)}. The process $X$ has continuous paths, is adapted to $\{\F_t\}$ because
$\xi$ and $W_t$ are $\F_t$-measurable, and satisfies $X_0=\sinh(\sinh^{-1}\xi)=\xi$. Finally,
$s\mapsto\abs{b(X_s)}+\sigma^2(X_s)$ is continuous, hence bounded on every $[0,t]$, so condition
(iii) of \bk{Definition 2.1} holds. Thus $X$ is a strong solution. No moment condition on $\xi$
is needed.

\emph{Step 2: uniqueness.}
Let $X$ be any continuous process, adapted to a filtration $\{\F_t\}$ which satisfies the usual
conditions and for which $W$ is a Brownian motion, such that $X_0=\xi$ a.s.\ and \bk{(2.4)}
holds; this covers the strong solutions of \bk{Definition 2.1}, and also the solutions on an
arbitrary filtered space of \bk{Section 5.3}. Then $X$ is a continuous semimartingale with
\begin{equation*}
M_t=\int_0^t\sigma(X_s)\,dW_s\in\cM^{c,\loc},\qquad
\qv{M}_t=\int_0^t\bigl(1+X_s^2\bigr)\,ds,\qquad B_t=\int_0^t b(X_s)\,ds,
\end{equation*}
the integrand $\sigma(X)$ being continuous and adapted, hence in $\cP^*(W)$. The function
$f\defeq\sinh^{-1}$ is of class $C^2(\R)$, with
\begin{equation*}
f'(x)=\frac{1}{\sqrt{1+x^2}}=\frac{1}{\sigma(x)},\qquad
f''(x)=-\frac{x}{(1+x^2)^{3/2}} .
\end{equation*}
By \bk{Theorem 3.3.3}, almost surely, for all $0\le t<\infty$,
\begin{equation}\label{eq:5-2-27:ito}
f(X_t)=f(\xi)+\int_0^t\frac{dM_s}{\sigma(X_s)}
+\int_0^t\frac{b(X_s)}{\sigma(X_s)}\,ds
-\frac12\int_0^t\frac{X_s}{(1+X_s^2)^{3/2}}\,\bigl(1+X_s^2\bigr)\,ds .
\end{equation}
The two Lebesgue integrals combine to
\begin{equation*}
\int_0^t\Bigl(1+\frac{X_s}{2\sqrt{1+X_s^2}}-\frac{X_s}{2\sqrt{1+X_s^2}}\Bigr)ds=t .
\end{equation*}
For the stochastic integral, put $\Phi_t\defeq\int_0^t\sigma(X_s)^{-1}\,dM_s$, which is well
defined in $\cM^{c,\loc}$ because $1/\sigma(X)$ is continuous, adapted and bounded by $1$. For
every $N\in\cM^{c,\loc}$, two applications of \bk{Proposition 3.2.24} (the identity
\bk{(3.2.30)}, first for the integral $\Phi$ with respect to $M$, then for the integral $M$
with respect to $W$) give
\begin{equation*}
\qv{\Phi,N}_t=\int_0^t\frac{1}{\sigma(X_s)}\,d\qv{M,N}_s
=\int_0^t\frac{1}{\sigma(X_s)}\,\sigma(X_s)\,d\qv{W,N}_s=\qv{W,N}_t ,
\end{equation*}
so that $\qv{\Phi-W}=\qv{\Phi-W,\Phi}-\qv{\Phi-W,W}=0$, and hence $\Phi=W$ up to
indistinguishability (a continuous local martingale which starts at zero and has zero quadratic
variation vanishes identically, \bk{Problem 1.5.12} applied after localization). Substituting
into \eqref{eq:5-2-27:ito}, we obtain almost surely
\begin{equation*}
\sinh^{-1}(X_t)=\sinh^{-1}(\xi)+t+W_t,\qquad 0\le t<\infty,
\end{equation*}
and applying $\sinh$ to both sides yields \eqref{eq:5-2-27:sol}. Hence \eqref{eq:5-2-27:sol} is
the unique solution; in particular, for a nonrandom initial condition $X_0=x$,
\begin{equation*}
X_t=x\cosh(t+W_t)+\sqrt{1+x^2}\,\sinh(t+W_t).
\end{equation*}

\rmk{The substitution is the one suggested by the surrounding text. Since
$\sigma\sigma'(x)=x$, the drift is $b=\sigma+\tfrac12\sigma\sigma'$, so the equation is the
Fisk--Stratonovich equation $dX_t=\sigma(X_t)\,dt+\sigma(X_t)\circ dW_t$ of \bk{(2.36)}, and
$\int_0^x du/\sigma(u)=\sinh^{-1}(x)$ is the change of scale of \bk{Exercise 2.20}, which
turns it into $d\sinh^{-1}(X_t)=dt+dW_t$. In the notation of the proof of
\bk{Proposition 2.21} (with $\sqrt{1+x^2}$ in the roles of both $b$ and $\sigma$ of
\bk{(2.36)}), the solution of \bk{(2.37)} is $u(x,y)=\sinh(x+\sinh^{-1}y)$, since
$\partial u/\partial x=\cosh(x+\sinh^{-1}y)=\sqrt{1+u^2}$ and $u(0,y)=y$. Then
$1/\rho(x,y)=\partial u/\partial y=\cosh(x+\sinh^{-1}y)/\sqrt{1+y^2}$, so the function of
\bk{(2.40)} is $f(x,y)=\rho(x,y)\sqrt{1+u(x,y)^2}=\sqrt{1+y^2}$, and the ordinary differential
equation \bk{(2.42)} reads $\dot Y_t=\sqrt{1+Y_t^2}$, $Y_0=\xi$. Its solution is
$Y_t=\sinh(\sinh^{-1}\xi+t)$, and $u(W_t,Y_t)=\sinh(\sinh^{-1}\xi+t+W_t)$ is
\eqref{eq:5-2-27:sol}. Uniqueness also follows from \bk{Theorem 2.5}, because $b$ and $\sigma$
have bounded derivatives; the argument of Step~2 does not need it.

As a check, for $X_0=x$ the formula gives, with $E\,e^{\pm W_t}=e^{t/2}$,
$E X_t=\tfrac12\bigl(x+\sqrt{1+x^2}\bigr)e^{3t/2}+\tfrac12\bigl(x-\sqrt{1+x^2}\bigr)e^{-t/2}$,
whose derivative at $t=0$ is $\sqrt{1+x^2}+\tfrac12x=b(x)$, as it must be.}
\end{kex}

\begin{kex}{Exercise}{2.28}{}
Throughout, $W=\{W_t=(W^{(1)}_t,\ldots,W^{(d)}_t),\F_t;\ 0\le t<\infty\}$ is a $d$-dimensional Brownian motion on a probability space $(\Omega,\F,P)$, with an arbitrary initial distribution in part~(i), and $\{\F_t\}$ satisfies the usual conditions, as is assumed for the stochastic calculus of Chapter~3. By analogy with \bk{(2.4)$'$} and \bk{(2.36)}, equation \bk{(2.36)$''$} means that, almost surely,
\begin{equation}\label{eq:5-2-28:fs}
X^{(i)}_t = X^{(i)}_0 + \int_0^t b_i(s,X_s)\,ds + \sum_{j=1}^d \int_0^t \sigma_{ij}(s,X_s)\circ dW^{(j)}_s;
\qquad 0\le t<\infty,\ 1\le i\le d,
\end{equation}
where the last integrals are Fisk--Stratonovich integrals in the sense of \bk{Definition 3.3.13}; part of the assertion is that each process $\sigma_{ij}(\cdot,X_\cdot)$ is a continuous semimartingale, so that these integrals are defined.

\textbf{(i)} Fix $1\le i,j\le d$ and set
\begin{equation*}
g(t,y)\defeq \frac{\partial u_i}{\partial y_j}(t,y)=\sigma_{ij}\bigl(t,u(t,y)\bigr);\qquad (t,y)\in[0,\infty)\times\R^d .
\end{equation*}

\emph{Step 1: $g$ is of class $C^{1,2}$.} The function $\sigma_{ij}$ has continuous first-order partial derivatives in all of its $d+1$ arguments, hence is differentiable on $[0,\infty)\times\R^d$ (with one-sided derivatives in $t$ at $t=0$), and the same is true of each $u_k$. The chain rule and the hypotheses $\partial u_k/\partial t=b_k(t,u)$, $\partial u_k/\partial y_l=\sigma_{kl}(t,u)$ therefore give
\begin{align*}
\frac{\partial g}{\partial t}(t,y)&=\Bigl[\frac{\partial\sigma_{ij}}{\partial t}+\sum_{k=1}^d\frac{\partial\sigma_{ij}}{\partial x_k}\,b_k\Bigr]\bigl(t,u(t,y)\bigr),
&
\frac{\partial g}{\partial y_l}(t,y)&=\Bigl[\sum_{k=1}^d\frac{\partial\sigma_{ij}}{\partial x_k}\,\sigma_{kl}\Bigr]\bigl(t,u(t,y)\bigr).
\end{align*}
Both are continuous, because $b_k$, $\sigma_{kl}$, the first-order derivatives of $\sigma_{ij}$ and $u$ are. For fixed $t$, the functions $x\mapsto(\partial\sigma_{ij}/\partial x_k)(t,x)$ and $x\mapsto\sigma_{kl}(t,x)$ are of class $C^1$ and so is $y\mapsto u(t,y)$; hence $\partial g/\partial y_l$ is differentiable in $y$, with
\begin{equation*}
\frac{\partial^2 g}{\partial y_m\partial y_l}(t,y)=\Bigl[\sum_{k=1}^d\sum_{p=1}^d\Bigl(\frac{\partial^2\sigma_{ij}}{\partial x_p\partial x_k}\,\sigma_{kl}+\frac{\partial\sigma_{ij}}{\partial x_k}\,\frac{\partial\sigma_{kl}}{\partial x_p}\Bigr)\sigma_{pm}\Bigr]\bigl(t,u(t,y)\bigr),
\end{equation*}
which is again continuous on $[0,\infty)\times\R^d$. Thus $g\in C^{1,2}([0,\infty)\times\R^d)$.

\emph{Step 2: the integrand is a semimartingale, and its Fisk--Stratonovich integral.} We have $W_t=W_0+(W_t-W_0)$, where $W_0$ is $\F_0$-measurable and the components of $W-W_0$ are continuous martingales with $\qv{W^{(k)}-W^{(k)}_0,W^{(l)}-W^{(l)}_0}_t=\delta_{kl}t$ (\bk{Problem 2.5.5}). It\^o's rule (\bk{Theorem 3.3.6}), applied to $g$ and this vector semimartingale, yields a.s.
\begin{equation*}
\sigma_{ij}(t,X_t)=g(t,W_t)=g(0,W_0)+N_t+C_t;\qquad 0\le t<\infty,
\end{equation*}
with
\begin{equation*}
N_t\defeq\sum_{l=1}^d\int_0^t\frac{\partial g}{\partial y_l}(s,W_s)\,dW^{(l)}_s,\qquad
C_t\defeq\int_0^t\Bigl[\frac{\partial g}{\partial s}+\frac12\sum_{l=1}^d\frac{\partial^2 g}{\partial y_l^2}\Bigr](s,W_s)\,ds .
\end{equation*}
Here $N\in\cM^{c,\loc}$, and $C$ is continuous, adapted, of bounded variation on bounded intervals, with $C_0=0$. Hence $\sigma_{ij}(\cdot,X_\cdot)$ is a continuous semimartingale, whose local martingale part $N$ satisfies, by \bk{Proposition 3.2.24} (see \bk{(3.2.30)}),
\begin{equation*}
\qv{W^{(j)}-W^{(j)}_0,N}_t=\sum_{l=1}^d\int_0^t\frac{\partial g}{\partial y_l}(s,W_s)\,d\qv{W^{(l)},W^{(j)}}_s=\int_0^t\frac{\partial g}{\partial y_j}(s,W_s)\,ds .
\end{equation*}
The integrator $W^{(j)}$ has local martingale part $W^{(j)}-W^{(j)}_0$ and no bounded variation part. \bk{Definition 3.3.13} and $\partial g/\partial y_j=\partial^2u_i/\partial y_j^2$ now give
\begin{equation}\label{eq:5-2-28:conv}
\int_0^t\sigma_{ij}(s,X_s)\circ dW^{(j)}_s=\int_0^t\frac{\partial u_i}{\partial y_j}(s,W_s)\,dW^{(j)}_s+\frac12\int_0^t\frac{\partial^2u_i}{\partial y_j^2}(s,W_s)\,ds;\qquad 0\le t<\infty .
\end{equation}

\emph{Step 3: It\^o's rule for $u_i$.} Applying \bk{Theorem 3.3.6} to $u_i\in C^{1,2}$ and the same semimartingale $W$, we obtain a.s., for all $0\le t<\infty$,
\begin{equation*}
X^{(i)}_t=u_i(0,W_0)+\int_0^t\frac{\partial u_i}{\partial s}(s,W_s)\,ds+\sum_{j=1}^d\int_0^t\frac{\partial u_i}{\partial y_j}(s,W_s)\,dW^{(j)}_s+\frac12\sum_{j=1}^d\int_0^t\frac{\partial^2u_i}{\partial y_j^2}(s,W_s)\,ds .
\end{equation*}
In the first integral, $(\partial u_i/\partial s)(s,W_s)=b_i(s,u(s,W_s))=b_i(s,X_s)$; and by \eqref{eq:5-2-28:conv}, summed over $j$, the last two sums add up to $\sum_{j=1}^d\int_0^t\sigma_{ij}(s,X_s)\circ dW^{(j)}_s$. Since $u_i(0,W_0)=X^{(i)}_0$, this is \eqref{eq:5-2-28:fs}. The process $X$ is continuous and $\{\F_t\}$-adapted, being a continuous function of $(t,W_t)$; thus $X$ solves \bk{(2.36)$''$}, with initial condition $X_0=u(0,W_0)$.

\begin{rem}
Equation \bk{(2.36)$''$} is equivalent to the It\^o equation
\begin{equation}\label{eq:5-2-28:ito}
dX^{(i)}_t=\Bigl[b_i+\frac12\sum_{j=1}^d\sum_{k=1}^d\frac{\partial\sigma_{ij}}{\partial x_k}\,\sigma_{kj}\Bigr](t,X_t)\,dt+\sum_{j=1}^d\sigma_{ij}(t,X_t)\,dW^{(j)}_t;\qquad 1\le i\le d,
\end{equation}
for every continuous, adapted, $\R^d$-valued process $X$; this is the multidimensional, time-dependent counterpart of the equivalence of \bk{(2.36)} and \bk{(2.36)$'$} shown in \bk{Remark 2.22}. (Here only the continuity of $b$ and $\sigma\in C^{1,2}$ are used, not the existence of $u$.) Indeed, suppose that $X$ solves either equation. Then $X$ is a continuous semimartingale and, by uniqueness of the decomposition, the local martingale part of $X^{(k)}$ is $M^{(k)}\defeq\sum_{l}\int_0^\cdot\sigma_{kl}(s,X_s)\,dW^{(l)}_s$: for \eqref{eq:5-2-28:ito} this is immediate, and for \eqref{eq:5-2-28:fs} it follows from \bk{Definition 3.3.13}, because each Fisk--Stratonovich integral $\int_0^\cdot\sigma_{kl}(s,X_s)\circ dW^{(l)}_s$ is the It\^o integral $\int_0^\cdot\sigma_{kl}(s,X_s)\,dW^{(l)}_s$ (a continuous local martingale, the integrand being continuous and adapted) plus a continuous adapted process of bounded variation, one half of a cross-variation. By \bk{Proposition 3.2.24}, $\qv{M^{(k)},W^{(j)}}_t=\int_0^t\sigma_{kj}(s,X_s)\,ds$, where here and below the cross-variation with $W^{(j)}$ means the one with its local martingale part $W^{(j)}-W^{(j)}_0$; the integrands $\sigma_{kl}(\cdot,X_\cdot)$ and $(\partial\sigma_{ij}/\partial x_k)(\cdot,X_\cdot)$ used here are continuous and adapted, hence locally bounded and progressively measurable, so all the stochastic integrals are defined. It\^o's rule (\bk{Theorem 3.3.6}) applied to $\sigma_{ij}\in C^{1,2}$ and $X$ shows that $\sigma_{ij}(\cdot,X_\cdot)$ is a continuous semimartingale (so that the Fisk--Stratonovich integrals in \eqref{eq:5-2-28:fs} are defined) with local martingale part $\sum_k\int_0^\cdot(\partial\sigma_{ij}/\partial x_k)(s,X_s)\,dM^{(k)}_s$, whose cross-variation with $W^{(j)}-W^{(j)}_0$ is, by \bk{Proposition 3.2.24} again,
\begin{equation*}
\sum_{k=1}^d\int_0^t\frac{\partial\sigma_{ij}}{\partial x_k}(s,X_s)\,d\qv{M^{(k)},W^{(j)}}_s=\int_0^t\Bigl[\sum_{k=1}^d\frac{\partial\sigma_{ij}}{\partial x_k}\,\sigma_{kj}\Bigr](s,X_s)\,ds .
\end{equation*}
Hence, by \bk{Definition 3.3.13}, a.s.\ for all $t$,
\begin{equation*}
\int_0^t\sigma_{ij}(s,X_s)\circ dW^{(j)}_s=\int_0^t\sigma_{ij}(s,X_s)\,dW^{(j)}_s+\frac12\int_0^t\Bigl[\sum_{k=1}^d\frac{\partial\sigma_{ij}}{\partial x_k}\,\sigma_{kj}\Bigr](s,X_s)\,ds ,
\end{equation*}
and summing over $j$ shows that the integral forms of \eqref{eq:5-2-28:fs} and \eqref{eq:5-2-28:ito} coincide. In particular, the process $X=u(t,W_t)$ of part~(i) also solves \eqref{eq:5-2-28:ito}; for it, the last display is \eqref{eq:5-2-28:conv}, since $\partial^2u_i/\partial y_j^2=\sum_k(\partial\sigma_{ij}/\partial x_k)(t,u)\,\sigma_{kj}(t,u)$ by Step~1.
\end{rem}

\textbf{(ii)} We place ourselves in the setting of \bk{Subsection 5.2.A}: $W$ is a standard, one-dimensional Brownian motion ($W_0=0$ a.s.), $\xi$ is a real random variable independent of $\F^W_\infty$, $\{\F_t\}$ is the augmented filtration \bk{(2.3)}, which satisfies the usual conditions and with respect to which $W$ is a Brownian motion, and $a$ is a real constant. The coefficients are
\begin{equation*}
b(t,x)=\frac{2x}{1+t}-a(1+t)^2,\qquad\sigma(t,x)=a(1+t)^2 ,
\end{equation*}
and no moment condition is imposed on $\xi$. We claim that the unique strong solution with initial condition $\xi$ is
\begin{equation}\label{eq:5-2-28:sol}
X_t=(1+t)^2\bigl[\xi+a(W_t-t)\bigr];\qquad 0\le t<\infty .
\end{equation}

\emph{Step 1: the function $u$.} The coefficient $b$ is continuous and $\sigma$ is of class $C^{1,2}$. The requirement $\partial u/\partial y=\sigma(t,u)=a(1+t)^2$ forces $u(t,y)=a(1+t)^2y+c(t)$, and then $\partial u/\partial t=b(t,u)$ reads
\begin{equation*}
2a(1+t)y+c'(t)=2a(1+t)y+\frac{2c(t)}{1+t}-a(1+t)^2 ,
\quad\text{i.e.,}\quad
\frac{d}{dt}\Bigl[\frac{c(t)}{(1+t)^2}\Bigr]=-a .
\end{equation*}
With $c(0)=0$ this gives $c(t)=-at(1+t)^2$ and
\begin{equation*}
u(t,y)=a(1+t)^2(y-t) .
\end{equation*}
Indeed, $u$ is of class $C^{1,2}$ (it is a polynomial), $\partial u/\partial y=a(1+t)^2=\sigma(t,u)$, and
\begin{equation*}
\frac{\partial u}{\partial t}(t,y)=2a(1+t)(y-t)-a(1+t)^2=\frac{2u(t,y)}{1+t}-a(1+t)^2=b\bigl(t,u(t,y)\bigr).
\end{equation*}

\emph{Step 2: a solution starting at the origin.} By part~(i) with $d=1$, the process $Y_t\defeq u(t,W_t)=a(1+t)^2(W_t-t)$ satisfies $Y_0=0$ and, a.s.,
\begin{equation*}
Y_t=\int_0^tb(s,Y_s)\,ds+\int_0^ta(1+s)^2\circ dW_s;\qquad 0\le t<\infty .
\end{equation*}
The integrand $a(1+s)^2$ is a deterministic function of class $C^1$, that is, a continuous semimartingale whose local martingale part vanishes. The cross-variation term in \bk{(3.3.9)} is therefore zero, the Fisk--Stratonovich integral coincides with the It\^o integral $\int_0^ta(1+s)^2\,dW_s$, and $Y$ satisfies the It\^o equation of the statement.

\emph{Step 3: arbitrary initial condition.} The process $Z_t\defeq\xi(1+t)^2$ satisfies, for every $\omega$,
\begin{equation*}
Z_t=\xi+\int_0^t2\xi(1+s)\,ds=\xi+\int_0^t\frac{2Z_s}{1+s}\,ds .
\end{equation*}
Let $X\defeq Z+Y$, which is the process \eqref{eq:5-2-28:sol}. Since $b(s,z+y)=2z/(1+s)+b(s,y)$, adding the last two displays gives, a.s.,
\begin{equation*}
X_t=\xi+\int_0^tb(s,X_s)\,ds+\int_0^ta(1+s)^2\,dW_s;\qquad 0\le t<\infty ,
\end{equation*}
which is \bk{(2.4)} for our equation. The other requirements of \bk{Definition 2.1} hold as well: $X$ has continuous paths; $X_t$ is a Borel function of $(\xi,W_t)$, hence measurable with respect to $\G_t=\sigma(\xi)\vee\F^W_t\subseteq\F_t$; $X_0=\xi$; and for every $\omega$ and $t$ the function $s\mapsto\abs{b(s,X_s(\omega))}+\sigma^2(s,X_s(\omega))$ is continuous, hence integrable, on $[0,t]$. Thus \eqref{eq:5-2-28:sol} is a strong solution.

\emph{Step 4: uniqueness.} For all $t\ge0$ and $x,y\in\R$,
\begin{equation*}
\abs{b(t,x)-b(t,y)}+\abs{\sigma(t,x)-\sigma(t,y)}=\frac{2\abs{x-y}}{1+t}\le2\abs{x-y} ,
\end{equation*}
so \bk{(2.8)} holds with $K_n=2$ for every $n$, and strong uniqueness holds by \bk{Theorem 2.5}: every strong solution with initial condition $\xi$ is indistinguishable from \eqref{eq:5-2-28:sol}.

\rmk{As a check, \eqref{eq:5-2-28:sol} can be verified without part~(i): by the integration by parts formula \bk{(3.3.8)}, in which the cross-variation vanishes because $(1+t)^2$ is of bounded variation, $d\bigl[(1+t)^2(\xi+aW_t-at)\bigr]=2(1+t)(\xi+aW_t-at)\,dt+(1+t)^2(a\,dW_t-a\,dt)$, and the first term is $\frac{2}{1+t}X_t\,dt$. Note also that $\sigma(t,x)=a(1+t)^2$ is unbounded in $t$, so the linear growth condition \bk{(2.13)} fails when $a\ne0$ and the existence result \bk{Theorem 2.9} does not apply directly; existence is supplied here by the explicit formula.}
\end{kex}

\solsection{5.3}{Weak Solutions}

\begin{kex}{Exercise}{3.12}{}
The statement uses the letter $X$ twice: for the solution of \bk{(3.12)} under $Q$, and for the coordinate process of a Brownian family under $\{P^x\}$. These coincide for the solution constructed in the proof of \bk{Proposition 3.6}, where $Q=Q^0$ is obtained from $P^0$ by a change of measure on the space of the Brownian family; we prove \bk{(3.13)} for this solution in Steps~1--2. In Step~3 we show that the left-hand side of \bk{(3.13)} is the same for \emph{every} weak solution of \bk{(3.12)}, so that \bk{(3.13)} holds for all of them, the right-hand side being computed on the Brownian family. Here $\sgn$ is the function of \bk{Example 3.5}, $\sgn(x)=1$ for $x>0$ and $\sgn(x)=-1$ for $x\le 0$, which is also the function \bk{(3.6.14)}.

Throughout, $\{X_s,\F_s;\,0\le s<\infty\}$, $(\Omega,\F)$, $\{P^x\}_{x\in\R}$ is the one-dimensional Brownian family of \bk{Section 3.6} (on the canonical space, with a filtration satisfying the usual conditions), and $L$ is its local time, which exists by \bk{Theorem 3.6.11}. We fix $t\ge 0$ and $\Gamma\in\B(\R)$. For $t=0$ both sides of \bk{(3.13)} equal $\ind_\Gamma(0)$, because $Q[X_0=0]=1$, $P^0[X_0=0]=1$ and $L_0(0)=0$, $P^0$-a.s.\ (by \bk{(3.6.13)} at time $0$); so let $t>0$.

\emph{Step 1: the solution of Proposition 3.6.} The drift $b(s,x)\defeq-\sgn(x)$ is Borel-measurable and bounded by $1$, so it satisfies \bk{(3.3)} with $K=1$. We run the proof of \bk{Proposition 3.6} with $d=1$, with time horizon $T=t$ and with $\mu=\delta_0$. The process
\begin{equation}\label{eq:5-3-12:Z}
Z_s \defeq \exp\Bigl\{-\int_0^s \sgn(X_u)\,dX_u-\frac12\int_0^s\sgn^2(X_u)\,du\Bigr\}
=\exp\Bigl\{-\int_0^s \sgn(X_u)\,dX_u-\frac s2\Bigr\},\qquad 0\le s\le t,
\end{equation}
where the stochastic integral is taken under $P^0$ with respect to the Brownian motion $X$, is a martingale under $P^0$ (this also follows directly from Novikov's condition, \bk{Corollary 3.5.13}, since $\int_0^t\sgn^2(X_u)\,du=t$). We define the probability measure $Q^0$ on $\F_t$ by
\begin{equation}\label{eq:5-3-12:Q}
Q^0(A)\defeq E^0[\ind_A Z_t],\qquad A\in\F_t .
\end{equation}
By the Girsanov \bk{Theorem 3.5.1}, the process $W_s\defeq X_s+\int_0^s\sgn(X_u)\,du$, $0\le s\le t$, is an $\{\F_s\}$-Brownian motion under $Q^0$, and $Q^0[X_0=0]=1$ because $P^0[X_0=0]=1$ and $Q^0\ll P^0$. Hence $(X,W)$, $(\Omega,\F_t,Q^0)$, $\{\F_s\}_{0\le s\le t}$ is a weak solution of \bk{(3.12)} on $[0,t]$: this is the solution provided by \bk{Proposition 3.6}. Since $\{X_t\in\Gamma\}\in\F_t$, \eqref{eq:5-3-12:Q} gives
\begin{equation}\label{eq:5-3-12:a}
Q^0[X_t\in\Gamma]=E^0\bigl[\ind_{\{X_t\in\Gamma\}}Z_t\bigr],
\end{equation}
which is \bk{(3.11)} for $n=1$, $t_1=T=t$ and $\mu=\delta_0$.

\emph{Step 2: the Tanaka formula.} The Tanaka formula \bk{(3.6.13)} with $a=0$ and $z=0$ reads
\[
2L_t(0)=\abs{X_t}-\int_0^t\sgn(X_s)\,dX_s,\qquad P^0\text{-a.s.},
\]
with the same function $\sgn$ and the same $P^0$-stochastic integral as in \eqref{eq:5-3-12:Z}. Substituting $\int_0^t\sgn(X_s)\,dX_s=\abs{X_t}-2L_t(0)$ into \eqref{eq:5-3-12:Z} we obtain
\[
Z_t=\exp\Bigl\{-\abs{X_t}+2L_t(0)-\frac t2\Bigr\},\qquad P^0\text{-a.s.},
\]
and \eqref{eq:5-3-12:a} becomes
\begin{equation}\label{eq:5-3-12:b}
Q^0[X_t\in\Gamma]=e^{-t/2}E^0\bigl[\ind_{\{X_t\in\Gamma\}}\exp\bigl(-\abs{X_t}+2L_t(0)\bigr)\bigr].
\end{equation}
This is \bk{(3.13)} for the solution of Step~1.

\emph{Step 3: an arbitrary weak solution.} Let now $(\wh X,\wh W)$, $(\wh\Omega,\wh\F,Q)$, $\{\wh\F_s\}$ be any weak solution of \bk{(3.12)}, defined for $0\le s\le T$ with some $T\ge t$ or for $0\le s<\infty$. Restricted to the time interval $[0,t]$ it is a weak solution of equation \bk{(3.2)} with $b(s,x)=-\sgn(x)$ and time horizon $t$, with initial distribution $\delta_0$; so is the solution of Step~1. Condition \bk{(3.6)} holds for both of them, because $\int_0^t\abs{b(s,x_s)}^2\,ds=t<\infty$ for every path $x$. By \bk{Proposition 3.10} (this is the content of \bk{Corollary 3.11}), the pairs $(\wh X_s,\wh W_s)_{0\le s\le t}$ under $Q$ and $(X_s,W_s)_{0\le s\le t}$ under $Q^0$ have the same law; in particular
\[
Q[\wh X_t\in\Gamma]=Q^0[X_t\in\Gamma],
\]
and \eqref{eq:5-3-12:b} yields
\[
Q[\wh X_t\in\Gamma]=e^{-t/2}E^0\bigl[\ind_{\{X_t\in\Gamma\}}\exp\bigl(-\abs{X_t}+2L_t(0)\bigr)\bigr];\qquad\Gamma\in\B(\R),
\]
which is \bk{(3.13)}. \qed

\rmk{For $\Gamma=\R$ the right-hand side of \bk{(3.13)} is $E^0[Z_t]=1$, as it must be. Moreover, by \bk{Theorem 3.6.17} the pair $(\abs{X_t},2L_t(0))$ has under $P^0$ the law of $(M_t-B_t,M_t)$ for a standard Brownian motion $B$ with running maximum $M$, whose density is given by \bk{(2.8.2)}; since $-X$ is again a Brownian motion under $P^0$, with the same local time at the origin as $X$ (by \bk{(3.6.6)} and the continuity of $x\mapsto L_t(x)$), the pairs $(X_t,L_t(0))$ and $(-X_t,L_t(0))$ have the same law, and we obtain the density
\[
Q[X_t\in dx]=\frac{e^{-t/2-\abs{x}}}{\sqrt{2\pi t^3}}\int_0^\infty (b+\abs{x})\exp\Bigl\{b-\frac{(b+\abs{x})^2}{2t}\Bigr\}\,db\,dx ,
\]
which we have checked numerically against a direct simulation of \bk{(3.12)}.}
\end{kex}

\begin{kex}{Problem}{3.13}{}
We create the drift by a change of measure, as in the proof of \bk{Proposition 3.6} but starting from a solution of the driftless equation. The Girsanov theory of \bk{Section 3.5} is set up for a Brownian motion on $[0,\infty)$, while the given Brownian motion lives only on $[0,T]$; we therefore first enlarge the probability space.

\emph{Setting and conventions.} For a process $Y$ we write $Y^T_t\defeq Y_{t\wedge T}$, $t\ge0$. If $\{\F_t\}_{0\le t\le T}$ satisfies the usual conditions, so does $\{\F_{t\wedge T}\}_{t\ge0}$. A continuous adapted process $Y$ on $[0,T]$ with $Y_0=0$ is a continuous local martingale on $[0,T]$ (a member of the class $\cM^{c,\loc}_T$ of p.~193) if $Y^T\in\cM^{c,\loc}$ relative to $\{\F_{t\wedge T}\}_{t\ge0}$; cross-variations on $[0,T]$ are those of the frozen processes, and stochastic integrals over $[0,T]$ are $\int_0^th\,dY\defeq I^{Y^T}_t(h)$, $0\le t\le T$. If $Y,Y'\in\cM^{c,\loc}$, then $Y^T\in\cM^{c,\loc}$ (\bk{Problem 1.3.24}\,(i) along a localizing sequence) and $Y^TY'^T-\qv{Y,Y'}^T=(YY'-\qv{Y,Y'})^T\in\cM^{c,\loc}$, so that $\qv{Y^T,Y'^T}=\qv{Y,Y'}^T$ by the uniqueness in \bk{Problem 1.5.17}.

A weak solution of \bk{(3.14)} for $0\le t\le T$ with initial distribution $\mu$ thus consists of a probability space $(\Omega,\F,P)$ with a filtration $\{\F_t\}_{0\le t\le T}$ satisfying the usual conditions, a continuous adapted $\R^d$-valued process $X=\{X_t;\,0\le t\le T\}$ with $P[X_0\in\Gamma]=\mu(\Gamma)$, $\Gamma\in\B(\R^d)$, and a $d$-dimensional Brownian motion $W=\{W_t,\F_t;\,0\le t\le T\}$ (here $r=d$), such that, $P$-a.s., $\int_0^T\sigma_{ij}^2(s,X_s)\,ds<\infty$ and $X^{(i)}_t=X^{(i)}_0+\sum_j\int_0^t\sigma_{ij}(s,X_s)\,dW^{(j)}_s$ for $0\le t\le T$, $1\le i,j\le d$. Here $W_0=0$ a.s.: this is implicit in \bk{Definition 3.1}, because the stochastic integral of \bk{Definition 3.2.23} is defined only for integrators in $\cM^{c,\loc}$, whose members vanish at time $0$ (\bk{Definition 1.5.15}). (If one prefers to read $\int\sigma\,dW$ as $\int\sigma\,d(W-W_0)$, one replaces $W$ by $W-W_0$, which is again a Brownian motion relative to $\{\F_t\}$, and nothing changes.) The coefficients $b_i,\sigma_{ij}$ are Borel-measurable, as always for \bk{(2.1)}.

Fix such a weak solution. Put $V^{(j)}\defeq(W^{(j)})^T$ and $g_{ij}(s)\defeq\sigma_{ij}(s,X_{s\wedge T})$, $s\ge0$. Since $W$ has independent Gaussian increments relative to $\{\F_t\}$, each $V^{(j)}$ and each $V^{(j)}V^{(k)}-\delta_{jk}(t\wedge T)$ is a continuous martingale relative to $\{\F_{t\wedge T}\}$; thus $V^{(j)}\in\cM^{c,\loc}$ and $\qv{V^{(j)},V^{(k)}}_t=\delta_{jk}(t\wedge T)$. The process $X^T$ is continuous and adapted to $\{\F_{t\wedge T}\}$, hence progressively measurable (\bk{Proposition 1.1.13}), and so is each $g_{ij}$; moreover $\int_0^tg_{ij}^2\,d\qv{V^{(j)}}=\int_0^{t\wedge T}\sigma^2_{ij}(s,X_s)\,ds<\infty$ a.s., so $g_{ij}\in\cP^*(V^{(j)})$ (\bk{(3.2.31)} and \bk{Definition 3.2.22}) and $I^{V^{(j)}}(g_{ij})$ is defined by \bk{Definition 3.2.23}. Let $I^{(i)}\defeq\sum_jI^{V^{(j)}}(g_{ij})$ and $M^{(i)}\defeq(I^{(i)})^T\in\cM^{c,\loc}$. The equation says that $X^{(i)}_t=X^{(i)}_0+M^{(i)}_t$, $0\le t\le T$, a.s. By \bk{(3.2.26)} (valid for integrators in $\cM^{c,\loc}$, p.~147), \bk{Proposition 3.2.24}, $\qv{V^{(j)},V^{(k)}}_t=\delta_{jk}(t\wedge T)$ and the stopping rule above (note $V^{(k)}=(V^{(k)})^T$),
\begin{equation}\label{eq:5-3-13:a}
\qv{M^{(i)}}_t=\sum_{j=1}^d\int_0^{t\wedge T}g_{ij}^2(s)\,ds,\qquad
\qv{M^{(i)},V^{(k)}}_t=\int_0^{t\wedge T}g_{ik}(s)\,ds;\qquad t\ge0 .
\end{equation}
Finally let $K<\infty$, $\lambda>0$ be constants with $\norm{b(t,x)}\le K$ and with all eigenvalues of $\sigma(t,x)\sigma^{tr}(t,x)$ at least $\lambda$, for all $(t,x)$.

\emph{Step 1: identification of a stochastic integral.}
\begin{lem}[identification]
Let $\{\cH_t\}$ satisfy the usual conditions on a probability space, let $\beta^{(1)},\dots,\beta^{(d)}\in\cM^{c,\loc}$ with $\qv{\beta^{(j)},\beta^{(k)}}_t=\delta_{jk}(t\wedge T)$, and let $h_1,\dots,h_d$ be progressively measurable with $\int_0^Th_j^2(s)\,ds<\infty$ a.s. If $N\in\cM^{c,\loc}$ satisfies $\qv{N}_t=\sum_j\int_0^{t\wedge T}h_j^2\,ds$ and $\qv{N,\beta^{(j)}}_t=\int_0^{t\wedge T}h_j\,ds$ for all $j$ and $t\ge0$, then $N=\sum_jI^{\beta^{(j)}}(h_j)$ up to indistinguishability.
\end{lem}
\begin{proof}
Since $\int_0^th_j^2\,d\qv{\beta^{(j)}}=\int_0^{t\wedge T}h_j^2\,ds<\infty$ a.s., $h_j\in\cP^*(\beta^{(j)})$ (\bk{(3.2.31)}, \bk{Definition 3.2.22}), so \bk{Definition 3.2.23} applies and $I\defeq\sum_jI^{\beta^{(j)}}(h_j)\in\cM^{c,\loc}$. By \bk{(3.2.26)}, $\qv{I}_t=\sum_{j,k}\int_0^th_jh_k\,d\qv{\beta^{(j)},\beta^{(k)}}=\sum_j\int_0^{t\wedge T}h_j^2\,ds$; by \bk{Proposition 3.2.24} (relation \bk{(3.2.30)} for every element of $\cM^{c,\loc}$), $\qv{I,N}_t=\sum_j\int_0^th_j\,d\qv{\beta^{(j)},N}=\sum_j\int_0^{t\wedge T}h_j^2\,ds$. Cross-variation is bilinear (uniqueness in \bk{Problem 1.5.17}), so $D\defeq N-I\in\cM^{c,\loc}$ has $\qv{D}=\qv{N}-2\qv{N,I}+\qv{I}=0$; thus $D^2\in\cM^{c,\loc}$. With $\sigma_n\defeq\inf\{t\ge0;\abs{D_t}\ge n\}$, a stopping time by \bk{Problem 1.2.7}, the local martingale $D^2_{\cdot\wedge\sigma_n}$ is bounded by $n^2$ a.s., hence a martingale (dominated convergence along a localizing sequence), and $E D^2_{t\wedge\sigma_n}=ED_0^2=0$. Letting $n\to\infty$ ($\sigma_n\uparrow\infty$ by continuity) gives $D_t=0$ a.s.\ for each $t$, and continuity gives $D\equiv0$ a.s.
\end{proof}

\emph{Step 2: enlarging the probability space.} Let $(C,\cC,P^0)$ be $C[0,\infty)^d$ with its Borel $\sigma$-field and $d$-dimensional Wiener measure, $B_t(\omega')\defeq\omega'(t)$ and $\F^B_t\defeq\sigma(B_u;\,u\le t)$; under $P^0$, $\{B_t,\F^B_t\}$ is a $d$-dimensional Brownian motion starting at the origin (p.~72). Let $\wh\Omega\defeq\Omega\times C$, let $(\wh\Omega,\wh\F,\wh P)$ be the completion of $(\Omega\times C,\F\otimes\cC,P\otimes P^0)$, let $\wh\cN$ be the class of $\wh P$-null sets, and $\pi(\omega,\omega')\defeq\omega$. For a random variable $\xi$ on $\Omega$ put $\wh\xi\defeq\xi\circ\pi$, and similarly for processes; since $\wh P\circ\pi^{-1}=P$, $\wh\xi$ has under $\wh P$ the law of $\xi$ under $P$, and $P$-a.s.\ properties of $\omega$ hold $\wh P$-a.s. Writing $u^+\defeq\max(u,0)$, define for $t\ge0$
\[
\G_t\defeq\sigma(\mathscr{R}_t),\quad
\mathscr{R}_t\defeq\bigl\{A_1\times A_2;\ A_1\in\F_{t\wedge T},\ A_2\in\F^B_{(t-T)^+}\bigr\},\qquad
\wh\F_t\defeq\bigcap_{\eps>0}\sigma(\G_{t+\eps}\cup\wh\cN).
\]
$\{\wh\F_t\}$ is a right-continuous filtration with $\wh\F_0\supseteq\wh\cN$; since $\wh P$ is complete, $\{\wh\F_t\}$ satisfies the usual conditions, and so does $\{\wh\F_{t\wedge T}\}_{t\ge0}$. Each $\mathscr{R}_t$ is a $\pi$-system containing $\wh\Omega$. If $\xi$ is $\F_{t\wedge T}$-measurable, then $\wh\xi$ is measurable with respect to $\G_{t\wedge T}\subseteq\wh\F_{t\wedge T}\subseteq\wh\F_t$.

\begin{lem}[transfer]
Let $0\le s<t$, let $\{\Phi_r;\,s\le r\le t\}$ be complex random variables on $\wh\Omega$ with $\abs{\Phi_r}\le c_0$ a.s.\ and $r\mapsto\Phi_r$ right-continuous at $s$, and let $c:[s,t]\to\mathbb{C}$ be right-continuous at $s$. If $\wh E[\ind_A\Phi_r]=\wh P(A)c(r)$ for all $r\in(s,t)$ and $A\in\mathscr{R}_r$, then $\wh E[\ind_A\Phi_s]=\wh P(A)c(s)$ for all $A\in\wh\F_s$.
\end{lem}
\begin{proof}
Fix $r\in(s,t)$. The class of $A\in\G_r$ with $\wh E[\ind_A\Phi_r]=\wh P(A)c(r)$ contains $\mathscr{R}_r$ and is closed under proper differences (linearity) and increasing unions (bounded convergence); by Dynkin's $\pi$-$\lambda$ theorem it is $\G_r$. The sets $A\in\wh\F$ that differ from some $A'\in\G_r$ by a $\wh P$-null set form a $\sigma$-field containing $\G_r\cup\wh\cN$, hence containing $\sigma(\G_r\cup\wh\cN)\supseteq\wh\F_s$. So for $A\in\wh\F_s$ we get $\wh E[\ind_A\Phi_r]=\wh E[\ind_{A'}\Phi_r]=\wh P(A')c(r)=\wh P(A)c(r)$. Let $r\downarrow s$ and use bounded convergence.
\end{proof}

\textbf{(a) A Brownian motion on $[0,\infty)$.} Define $\wh W_t(\omega,\omega')\defeq W_{t\wedge T}(\omega)+B_{(t-T)^+}(\omega')-B_0(\omega')$, $t\ge0$. It is continuous and $\G_t$-adapted, $\wh W_0=\wh{W_0}=0$ a.s., and $\wh W_{t\wedge T}=\wh{W_{t\wedge T}}$, i.e.\ $(\wh W^{(j)})^T=\wh V^{(j)}$. Let $0\le r<t$, $v\in\R^d$, $A=A_1\times A_2\in\mathscr{R}_r$. The increment $\wh W_t-\wh W_r$ is the sum of $W_{t\wedge T}-W_{r\wedge T}$, a function of $\omega$ which under $P$ is independent of $\F_{r\wedge T}$ and normal with covariance $(t\wedge T-r\wedge T)I_d$, and of $B_{(t-T)^+}-B_{(r-T)^+}$, a function of $\omega'$ which under $P^0$ is independent of $\F^B_{(r-T)^+}$ and normal with covariance $((t-T)^+-(r-T)^+)I_d$. By Fubini's theorem,
\[
\wh E\bigl[\ind_Ae^{i\langle v,\wh W_t-\wh W_r\rangle}\bigr]
=P(A_1)\,P^0(A_2)\,e^{-\norm{v}^2[(t\wedge T-r\wedge T)+((t-T)^+-(r-T)^+)]/2}
=\wh P(A)\,e^{-\norm{v}^2(t-r)/2},
\]
because $u\wedge T+(u-T)^+=u$. The transfer lemma, with $\Phi_r=e^{i\langle v,\wh W_t-\wh W_r\rangle}$ and $c(r)=e^{-\norm{v}^2(t-r)/2}$, yields $\wh E[\ind_Ae^{i\langle v,\wh W_t-\wh W_s\rangle}]=\wh P(A)e^{-\norm{v}^2(t-s)/2}$ for all $s<t$, $A\in\wh\F_s$, $v\in\R^d$. If $\wh P(A)>0$, the probability measure $\wh P(\,\cdot\cap A)/\wh P(A)$ thus gives $\wh W_t-\wh W_s$ the characteristic function of the normal law $\nu_{t-s}\defeq\mathcal N(0,(t-s)I_d)$; hence $\wh P(A\cap\{\wh W_t-\wh W_s\in\Gamma\})=\wh P(A)\,\nu_{t-s}(\Gamma)$ for all $A\in\wh\F_s$, $\Gamma\in\B(\R^d)$. So $\wh W$ is a $d$-dimensional Brownian motion relative to $\{\wh\F_t\}$ with $\wh P[\wh W_0=0]=1$ (\bk{Definition 2.5.1}), and the standing assumptions of \bk{Section 3.5} hold for $(\wh\Omega,\wh\F,\wh P)$, $\{\wh\F_t\}$, $\wh W$.

\textbf{(b) Lifting local martingales.} \emph{If $Y\in\cM^{c,\loc}$ relative to $\{\F_{t\wedge T}\}$ on $(\Omega,\F,P)$, then $\wh Y\in\cM^{c,\loc}$ relative to $\{\wh\F_{t\wedge T}\}$ on $(\wh\Omega,\wh\F,\wh P)$.} Indeed, let $T_n$ localize $Y$ and put $S_n\defeq T_n\wedge\inf\{t\ge0;\abs{Y_t}\ge n\}$: these are nondecreasing stopping times (\bk{Problem 1.2.7}) with $S_n\uparrow\infty$ a.s., and $Y^n\defeq Y_{\cdot\wedge S_n}$ is a martingale (\bk{Problem 1.3.24}\,(i)) with $\abs{Y^n}\le n$ a.s. The $\wh S_n$ are stopping times of $\{\wh\F_{t\wedge T}\}$, since $\{\wh S_n\le t\}=\{S_n\le t\}\times C$ with $\{S_n\le t\}\in\F_{t\wedge T}$; moreover $\wh S_n\uparrow\infty$ $\wh P$-a.s., $\wh Y$ is continuous and adapted with $\wh Y_0=0$ a.s., and $\wh Y_{\cdot\wedge\wh S_n}=\wh{Y^n}$. For $0\le r<t$ and $A=A_1\times A_2\in\mathscr{R}_r$, Fubini's theorem and the martingale property of $Y^n$ give $\wh E[\ind_A(\wh{Y^n_t}-\wh{Y^n_r})]=P^0(A_2)\,E[\ind_{A_1}(Y^n_t-Y^n_r)]=0$. The transfer lemma ($\Phi_r=\wh{Y^n_t}-\wh{Y^n_r}$, $c\equiv0$) gives $\wh E[\ind_A(\wh{Y^n_t}-\wh{Y^n_s})]=0$ for $s<t$ and $A\in\wh\F_s\supseteq\wh\F_{s\wedge T}$, so $\wh{Y^n}$ is a martingale relative to $\{\wh\F_{t\wedge T}\}$. Consequently, if $Y,Y'\in\cM^{c,\loc}$ relative to $\{\F_{t\wedge T}\}$, then $\wh Y\wh Y'-\wh{\qv{Y,Y'}}$, the lift of $YY'-\qv{Y,Y'}$, is in $\cM^{c,\loc}$, while $\wh{\qv{Y,Y'}}$ is continuous, adapted, of bounded variation and null at $0$; by \bk{Problem 1.5.17}, $\qv{\wh Y,\wh Y'}=\wh{\qv{Y,Y'}}$.

\textbf{(c) The lifted solution.} Put $\wh X_t\defeq\wh{X_{t\wedge T}}$, a continuous process adapted to $\{\wh\F_{t\wedge T}\}$, and $\wh g_{ij}(s)\defeq\sigma_{ij}(s,\wh X_s)$, the lift of $g_{ij}$, which is progressively measurable relative to $\{\wh\F_{t\wedge T}\}$. Then $\wh P[\wh X_0\in\Gamma]=\mu(\Gamma)$, and $\wh P$-a.s.\ $\int_0^T\wh g_{ij}^2\,ds<\infty$ and $\wh X^{(i)}_t=\wh X^{(i)}_0+\wh M^{(i)}_t$ for $0\le t\le T$. By (b) and \eqref{eq:5-3-13:a}, $\wh M^{(i)},\wh V^{(j)}\in\cM^{c,\loc}$ relative to $\{\wh\F_{t\wedge T}\}$, $\qv{\wh V^{(j)},\wh V^{(k)}}_t=\delta_{jk}(t\wedge T)$, and
\begin{equation}\label{eq:5-3-13:b}
\qv{\wh M^{(i)}}_t=\sum_{j=1}^d\int_0^{t\wedge T}\wh g_{ij}^2(s)\,ds,\qquad
\qv{\wh M^{(i)},\wh V^{(k)}}_t=\int_0^{t\wedge T}\wh g_{ik}(s)\,ds;\qquad t\ge0 .
\end{equation}
The identification lemma ($\cH_t=\wh\F_{t\wedge T}$, $\beta=\wh V$, $h_j=\wh g_{ij}$, $N=\wh M^{(i)}$) gives $\wh M^{(i)}=\sum_jI^{\wh V^{(j)}}(\wh g_{ij})$. As $\wh V^{(j)}=(\wh W^{(j)})^T$, this says that $(\wh X,\wh W)$, $(\wh\Omega,\wh\F,\wh P)$, $\{\wh\F_t\}_{0\le t\le T}$ is a weak solution of \bk{(3.14)} for $0\le t\le T$ with initial distribution $\mu$, whose Brownian motion $\wh W$ is defined on $[0,\infty)$ relative to $\{\wh\F_t\}_{t\ge0}$. From now on we work with this solution.

\emph{Step 3: the integrand.} Put $\theta(t,x)\defeq\sigma^{-1}(t,x)b(t,x)$. By Cramer's rule the entries of $\sigma^{-1}$ are polynomials in the entries of $\sigma$ divided by $\det\sigma\ne0$, so $\theta$ is Borel-measurable. The symmetric matrix $(\sigma\sigma^{tr})^{-1}$ has its eigenvalues in $(0,1/\lambda]$, whence
\begin{equation}\label{eq:5-3-13:c}
\norm{\theta(t,x)}^2=b^{tr}(\sigma^{-1})^{tr}\sigma^{-1}b
=b^{tr}(\sigma\sigma^{tr})^{-1}b\le\frac{\norm{b(t,x)}^2}{\lambda}\le\frac{K^2}{\lambda}.
\end{equation}
Set $\wh\theta_t\defeq\theta(t,\wh X_t)\ind_{[0,T]}(t)$, $t\ge0$. As $\wh X$ is continuous and adapted, $\wh\theta$ is progressively measurable relative to $\{\wh\F_t\}$, and
\begin{equation}\label{eq:5-3-13:d}
\sum_{k=1}^d\wh g_{ik}(t)\,\wh\theta^{(k)}_t=b_i(t,\wh X_t);\qquad 0\le t\le T,\ 1\le i\le d.
\end{equation}

\emph{Step 4: change of measure.} By \eqref{eq:5-3-13:c}, $\wh\theta$ satisfies \bk{(3.5.1)}, and $\wh E\exp\bigl(\frac12\int_0^{t}\norm{\wh\theta_s}^2ds\bigr)\le e^{K^2T/(2\lambda)}<\infty$ for every $t\ge0$, which is \bk{(3.5.20)}. By Novikov's condition (\bk{Corollary 3.5.13}),
\[
\wh Z_t\defeq\exp\Bigl[\sum_{k=1}^d\int_0^t\wh\theta^{(k)}_s\,d\wh W^{(k)}_s-\frac12\int_0^t\norm{\wh\theta_s}^2ds\Bigr];\qquad t\ge0,
\]
is a martingale, with $\wh Z_0=1$. Define $\wt P(A)\defeq\wh E[\ind_A\wh Z_T]$ for $A\in\wh\F_T$, as in \bk{(3.5.4)}. By the Girsanov \bk{Theorem 3.5.1}, the process
\[
\wt W^{(k)}_t\defeq\wh W^{(k)}_t-\int_0^t\wh\theta^{(k)}_s\,ds;\qquad 1\le k\le d,\ 0\le t\le T,
\]
is a $d$-dimensional Brownian motion $\{\wt W_t,\wh\F_t;\,0\le t\le T\}$ on $(\wh\Omega,\wh\F_T,\wt P)$, with $\wt W_0=0$ a.s.

Because $\wh Z_T>0$, $\wh P$ and $\wt P$ have the same null sets in $\wh\F_T$. Hence $\wh\F_0$ contains all $\wt P$-null sets of $\wh\F_T$, and $\{\wh\F_t\}_{0\le t\le T}$ satisfies the usual conditions on $(\wh\Omega,\wh\F_T,\wt P)$; every property described by an event of $\wh\F_T$ that holds $\wh P$-a.s.\ holds $\wt P$-a.s. In particular $\wh X$ is continuous and adapted, and $\int_0^T\{\abs{b_i(s,\wh X_s)}+\wh g_{ij}^2(s)\}\,ds\le KT+\int_0^T\wh g_{ij}^2(s)\,ds<\infty$, $\wt P$-a.s. Since $\{\wh X_0\in\Gamma\}\in\wh\F_0$ and $\wh E[\wh Z_T\,|\,\wh\F_0]=\wh Z_0=1$,
\[
\wt P[\wh X_0\in\Gamma]=\wh E\bigl[\ind_{\{\wh X_0\in\Gamma\}}\wh E[\wh Z_T\,|\,\wh\F_0]\bigr]=\wh P[\wh X_0\in\Gamma]=\mu(\Gamma);\qquad\Gamma\in\B(\R^d).
\]

\emph{Step 5: the equation.} Fix $i$. The restrictions to $[0,T]$ of $\wh M^{(i)}$ and of $\wh W^{(j)}$ belong to $\cM^{c,\loc}_T$, their frozen versions $\wh M^{(i)}$ and $\wh V^{(j)}$ being in $\cM^{c,\loc}$ relative to $\{\wh\F_{t\wedge T}\}$ by Step 2. We apply \bk{Proposition 3.5.4}, whose hypotheses hold by Steps 2(a) and 4. In the notation \bk{(3.5.8)}, by \eqref{eq:5-3-13:b} and \eqref{eq:5-3-13:d},
\[
\wt M^{(i)}_t=\wh M^{(i)}_t-\sum_{k=1}^d\int_0^t\wh\theta^{(k)}_s\,\wh g_{ik}(s)\,ds=\wh M^{(i)}_t-\int_0^tb_i(s,\wh X_s)\,ds;\qquad 0\le t\le T,
\]
while the process associated with $\wh W^{(j)}$ is $\wh W^{(j)}_t-\int_0^t\wh\theta^{(j)}_s\,ds=\wt W^{(j)}_t$. The proposition asserts that $\wt M^{(i)}\in\wt{\cM}^{c,\loc}_T$ and that cross-variations on $[0,T]$ are unchanged; by \eqref{eq:5-3-13:b}, under $\wt P$,
\[
\qv{\wt M^{(i)}}_t=\sum_{j=1}^d\int_0^{t}\wh g_{ij}^2(s)\,ds,\qquad
\qv{\wt M^{(i)},\wt W^{(j)}}_t=\int_0^{t}\wh g_{ij}(s)\,ds;\qquad 0\le t\le T .
\]
Now apply the identification lemma on $(\wh\Omega,\wh\F_T,\wt P)$ with $\cH_t=\wh\F_{t\wedge T}$, $\beta=\wt W^T$ (a martingale with $\qv{\wt W^{(j),T},\wt W^{(k),T}}_t=\delta_{jk}(t\wedge T)$, since $\wt W$ is a Brownian motion on $[0,T]$), $h_j=\wh g_{ij}$ (progressively measurable, with $\int_0^T\wh g_{ij}^2\,ds<\infty$ $\wt P$-a.s.) and $N=(\wt M^{(i)})^T$, whose quadratic variation and cross-variations with the $\beta^{(j)}$ are the frozen versions of those just displayed. It yields $\wt M^{(i)}_t=\sum_j\int_0^t\wh g_{ij}(s)\,d\wt W^{(j)}_s$, $0\le t\le T$, $\wt P$-a.s. Combining this with $\wh X^{(i)}_t=\wh X^{(i)}_0+\wh M^{(i)}_t$ (valid $\wt P$-a.s.), we obtain, $\wt P$-a.s.,
\[
\wh X_t=\wh X_0+\int_0^tb(s,\wh X_s)\,ds+\int_0^t\sigma(s,\wh X_s)\,d\wt W_s;\qquad 0\le t\le T.
\]
Hence $(\wh X,\wt W)$, $(\wh\Omega,\wh\F_T,\wt P)$, $\{\wh\F_t\}_{0\le t\le T}$ is a weak solution of \bk{(2.1)} for $0\le t\le T$ with initial distribution $\mu$.

\rmk{If the given solution of \bk{(3.14)} is defined for all $t\ge0$, Step 2 is unnecessary: one takes $\theta(t,X_t)\ind_{[0,T]}(t)$ on the original space. The nonsingularity of $\sigma$ and the eigenvalue bound were used only to produce a bounded progressively measurable $\wh\theta$ with \eqref{eq:5-3-13:d}; boundedness could be replaced by any condition guaranteeing that $\wh Z$ is a martingale.}
\end{kex}

\solsection{5.4}{The Martingale Problem of Stroock and Varadhan}

\begin{kex}{Problem}{4.3}{}
Throughout, $D_if$ and $D_{ik}f$ denote $\partial f/\partial x_i$ and $\partial^2 f/\partial x_i\partial x_k$. In \bk{(4.2)$'$} the operator $\cA_s'$ acts on $f(s,\cdot)\in C^2(\R^d)$ and the result is evaluated at the whole path $X$; that is, the integrand is
\begin{equation}\label{eq:5-4-03:integrand}
  \frac{\partial f}{\partial s}(s,X_s)
  +\frac12\sum_{i=1}^d\sum_{k=1}^d a_{ik}(s,X)\,D_{ik}f(s,X_s)
  +\sum_{i=1}^d b_i(s,X)\,D_if(s,X_s).
\end{equation}
By \bk{Definition 3.14} and \bk{Definition 3.1}, the filtration $\{\F_t\}$ satisfies the usual conditions, $X$ is continuous and adapted, $W$ is an $r$-dimensional Brownian motion relative to $\{\F_t\}$, and conditions (iii), (iv) of \bk{Definition 3.14} hold. By \bk{Problem 2.5.5} the processes $W^{(j)}-W^{(j)}_0$ are in $\cM_2^c$ with
\begin{equation}\label{eq:5-4-03:w}
  \qv{W^{(j)}-W^{(j)}_0,\,W^{(l)}-W^{(l)}_0}_t=t\,\delta_{jl};\qquad 1\le j,l\le r,
\end{equation}
and the stochastic integrals $dW^{(j)}$ are integrals with respect to these martingales.

\emph{Step 1: measurability and integrability.} Let $\G_t=\sigma(y(u);\,0\le u\le t)$ be the $\sigma$-fields on $C[0,\infty)^d$ of \bk{Definition 3.5.15}. For fixed $t$ the mapping $\omega\mapsto X_\cdot(\omega)$ from $\Omega$ into $C[0,\infty)^d$ is $\F_t/\G_t$-measurable, because $\G_t$ is generated by the coordinate mappings $y\mapsto y(u)$, $0\le u\le t$, and each $X_u$ with $u\le t$ is $\F_t$-measurable. Consequently $(s,\omega)\mapsto(s,X_\cdot(\omega))$ is measurable from $([0,t]\times\Omega,\B([0,t])\otimes\F_t)$ into $([0,t]\times C[0,\infty)^d,\B([0,t])\otimes\G_t)$ (it suffices to look at preimages of measurable rectangles). Composing it with the restriction to $[0,t]\times C[0,\infty)^d$ of a progressively measurable functional $\mu$, which is $\B([0,t])\otimes\G_t$-measurable, we see that $(s,\omega)\mapsto\mu(s,X_\cdot(\omega))$ is a progressively measurable process. Thus $b_i(s,X)$, $\sigma_{ij}(s,X)$ and, by \bk{(4.6)}, $a_{ik}(s,X)$ are progressively measurable. The processes $\frac{\partial f}{\partial s}(s,X_s)$, $D_if(s,X_s)$, $D_{ik}f(s,X_s)$ are continuous and adapted, hence progressively measurable (\bk{Proposition 1.1.13}), and so the process \eqref{eq:5-4-03:integrand} is progressively measurable as well.

Let $\Omega^*$ be the event on which $\int_0^n\{\abs{b_i(s,X)}+\sigma_{ij}^2(s,X)\}\,ds<\infty$ for all $n\ge1$ and all $i,j$; by condition (iii), $P(\Omega^*)=1$, and $\Omega^*\in\F_0$ because $\F_0$ contains the $P$-null sets. From \bk{(4.6)},
\[
  \abs{a_{ik}(s,y)}\le\frac12\sum_{j=1}^r\bigl\{\sigma_{ij}^2(s,y)+\sigma_{kj}^2(s,y)\bigr\},
\]
so that on $\Omega^*$ the functions $s\mapsto b_i(s,X)$ and $s\mapsto a_{ik}(s,X)$ are Lebesgue integrable on every $[0,t]$; since the three continuous processes above are bounded on $[0,t]$ for each $\omega$, the function \eqref{eq:5-4-03:integrand} is integrable on $[0,t]$ for every $\omega\in\Omega^*$ and $t\ge 0$. Hence $M^f$ is well defined by \bk{(4.2)$'$} on $\Omega^*$ (we set $M^f\equiv0$ off $\Omega^*$), it is continuous, and it is adapted by Fubini's theorem and the progressive measurability of \eqref{eq:5-4-03:integrand}. For the same reasons
\[
  B^{(i)}_t\defeq\int_0^t b_i(s,X)\,ds\ \ \text{on }\Omega^*,\qquad B^{(i)}\defeq0\ \ \text{off }\Omega^*;\qquad 1\le i\le d,
\]
are continuous, adapted processes of bounded variation with $B^{(i)}_0=0$. Because $\sigma_{ij}(\cdot,X)$ is progressively measurable and $\int_0^t\sigma_{ij}^2(s,X)\,ds<\infty$ a.s.\ for every $t$, it belongs to the class $\cP^*$ of \bk{Definition 3.2.22} relative to $W^{(j)}-W^{(j)}_0$, and
\[
  M^{(i)}_t\defeq\sum_{j=1}^r\int_0^t\sigma_{ij}(s,X)\,dW^{(j)}_s;\qquad 0\le t<\infty,\ 1\le i\le d,
\]
is in $\cM^{c,\loc}$ (\bk{Definition 3.2.23}; a finite sum of elements of $\cM^{c,\loc}$ is in $\cM^{c,\loc}$, as one sees by taking the minimum of the localizing sequences). Condition (iv) of \bk{Definition 3.14} states that, almost surely,
\begin{equation}\label{eq:5-4-03:semi}
  X^{(i)}_t=X^{(i)}_0+M^{(i)}_t+B^{(i)}_t;\qquad 0\le t<\infty,\ 1\le i\le d.
\end{equation}

\emph{Step 2: the cross-variations of the $M^{(i)}$.} The cross-variation of \bk{Problem 1.5.17} is bilinear on $\cM^{c,\loc}$: if $Y,Y',Z\in\cM^{c,\loc}$ and $\alpha,\beta\in\R$, then $(\alpha Y+\beta Y')Z-(\alpha\qv{Y,Z}+\beta\qv{Y',Z})$ is in $\cM^{c,\loc}$, and $\alpha\qv{Y,Z}+\beta\qv{Y',Z}$ is continuous, adapted, of bounded variation and vanishes at $t=0$, so it is $\qv{\alpha Y+\beta Y',Z}$ by the uniqueness assertion of that problem. Moreover, the cross-variation formula \bk{(3.2.26)},
\begin{equation}\label{eq:5-4-03:cv}
  \qv{I^M(Y),I^N(Z)}_t=\int_0^tY_uZ_u\,d\qv{M,N}_u;\qquad 0\le t<\infty,\ \text{a.s.},
\end{equation}
is valid for $M,N\in\cM^{c,\loc}$, $Y\in\cP^*(M)$, $Z\in\cP^*(N)$ (see the paragraph preceding \bk{Proposition 3.2.24}). From bilinearity, \eqref{eq:5-4-03:cv}, \eqref{eq:5-4-03:w} and \bk{(4.6)} we obtain, almost surely,
\begin{equation}\label{eq:5-4-03:qvM}
  \qv{M^{(i)},M^{(k)}}_t=\sum_{j=1}^r\sum_{l=1}^r\int_0^t\sigma_{ij}(s,X)\sigma_{kl}(s,X)\,\delta_{jl}\,ds
  =\int_0^ta_{ik}(s,X)\,ds;\qquad 0\le t<\infty.
\end{equation}

\emph{Step 3: $M^f\in\cM^{c,\loc}$.} In view of \eqref{eq:5-4-03:semi}, $X$ is, up to indistinguishability, of the form required in \bk{Theorem 3.3.6}: $X_0$ is $\F_0$-measurable, $M^{(i)}\in\cM^{c,\loc}$, and $B^{(i)}$ is continuous, adapted, of bounded variation, with $B^{(i)}_0=0$. It\^o's rule \bk{(3.3.5)} gives, almost surely, for all $0\le t<\infty$,
\begin{align*}
  f(t,X_t)={}&f(0,X_0)+\int_0^t\frac{\partial f}{\partial s}(s,X_s)\,ds
  +\sum_{i=1}^d\int_0^tD_if(s,X_s)\,dB^{(i)}_s
  +\sum_{i=1}^d\int_0^tD_if(s,X_s)\,dM^{(i)}_s\\
  &+\frac12\sum_{i=1}^d\sum_{k=1}^d\int_0^tD_{ik}f(s,X_s)\,d\qv{M^{(i)},M^{(k)}}_s .
\end{align*}
On $\Omega^*$ the function $t\mapsto B^{(i)}_t$ is absolutely continuous with density $b_i(\cdot,X)$, and by \eqref{eq:5-4-03:qvM} almost every path of $\qv{M^{(i)},M^{(k)}}$ is absolutely continuous with density $a_{ik}(\cdot,X)$; hence the Lebesgue--Stieltjes integrals above are equal to $\int_0^tD_if(s,X_s)\,b_i(s,X)\,ds$ and $\int_0^tD_{ik}f(s,X_s)\,a_{ik}(s,X)\,ds$, respectively. Comparing with \eqref{eq:5-4-03:integrand} and \bk{(4.2)$'$} we arrive at
\begin{equation}\label{eq:5-4-03:rep}
  M^f_t=\sum_{i=1}^d\int_0^tD_if(s,X_s)\,dM^{(i)}_s;\qquad 0\le t<\infty,\ \text{a.s.}
\end{equation}
The integrand $D_if(s,X_s)$ is progressively measurable and, being continuous in $s$, satisfies
\[
  \int_0^t\bigl(D_if(s,X_s)\bigr)^2\,d\qv{M^{(i)}}_s\le\max_{0\le s\le t}\bigl(D_if(s,X_s)\bigr)^2\cdot\qv{M^{(i)}}_t<\infty,
\]
so it belongs to $\cP^*(M^{(i)})$, and each stochastic integral in \eqref{eq:5-4-03:rep} is in $\cM^{c,\loc}$ by \bk{Definition 3.2.23}. Therefore the right-hand side of \eqref{eq:5-4-03:rep} is in $\cM^{c,\loc}$. The process $M^f$ is continuous, adapted and indistinguishable from it, so that stopping $M^f$ at the localizing times of the right-hand side yields martingales, and $M^f_0=0$. Thus $M^f\in\cM^{c,\loc}$, and likewise $M^g\in\cM^{c,\loc}$.

\emph{Step 4: the cross-variation \bk{(4.3)$'$}.} Applying \eqref{eq:5-4-03:rep} to $f$ and to $g$, and then bilinearity, \eqref{eq:5-4-03:cv} and \eqref{eq:5-4-03:qvM}, we get, almost surely, for all $0\le t<\infty$,
\begin{align*}
  \qv{M^f,M^g}_t
  &=\sum_{i=1}^d\sum_{k=1}^d\int_0^tD_if(s,X_s)\,D_kg(s,X_s)\,d\qv{M^{(i)},M^{(k)}}_s\\
  &=\sum_{i=1}^d\sum_{k=1}^d\int_0^ta_{ik}(s,X)\,D_if(s,X_s)\,D_kg(s,X_s)\,ds,
\end{align*}
which is \bk{(4.3)$'$}.

\emph{Step 5: $M^f\in\cM_2^c$ under \bk{(4.7)}.} Suppose that $\abs{D_if}\le C$ on $[0,\infty)\times\R^d$ for $1\le i\le d$ and that \bk{(4.7)} holds, where $\norm{\sigma}^2=\sum_{i,j}\sigma_{ij}^2$ as in \bk{(2.9)}. For $\xi\in\R^d$, \bk{(4.6)} and the Cauchy--Schwarz inequality give
\[
  0\le\sum_{i=1}^d\sum_{k=1}^da_{ik}(s,y)\xi_i\xi_k=\sum_{j=1}^r\Bigl(\sum_{i=1}^d\sigma_{ij}(s,y)\xi_i\Bigr)^2
  \le\norm{\sigma(s,y)}^2\,\norm{\xi}^2 .
\]
Taking $g=f$ in \bk{(4.3)$'$} and $\xi_i=D_if(s,X_s)$, so that $\norm{\xi}^2\le dC^2$, we obtain for $0\le t\le T<\infty$
\begin{equation}\label{eq:5-4-03:qvbound}
  \qv{M^f}_t\le dC^2K_T^2\,t\le dC^2K_T^2\,T\qquad\text{a.s.}
\end{equation}
Let $\{T_n\}_{n=1}^\infty$ be a nondecreasing sequence of stopping times with $T_n\to\infty$ a.s.\ such that $\{M^f_{t\wedge T_n},\F_t;\ 0\le t<\infty\}$ is a martingale for every $n$. Fix $0\le s<t\le T$. By \bk{Problem 1.5.19}(iii), applied to $M^f\in\cM^{c,\loc}$ and to the bounded stopping times $t\wedge T_n$ and $t$, and by \eqref{eq:5-4-03:qvbound},
\[
  E\bigl[(M^f_{t\wedge T_n})^2\bigr]\le E\qv{M^f}_{t\wedge T_n}\le dC^2K_T^2\,T,\qquad
  E\bigl[(M^f_t)^2\bigr]\le E\qv{M^f}_t\le dC^2K_T^2\,T .
\]
Hence the sequence $\{M^f_{t\wedge T_n}\}_{n\ge1}$ is bounded in $L^2$, therefore uniformly integrable, and it converges a.s.\ to $M^f_t$; so the convergence takes place in $L^1$, and the same holds with $s$ in place of $t$. Letting $n\to\infty$ in $E[M^f_{t\wedge T_n}\ind_A]=E[M^f_{s\wedge T_n}\ind_A]$, $A\in\F_s$, we find $E[M^f_t\ind_A]=E[M^f_s\ind_A]$. Thus $M^f$ is a continuous martingale with $M^f_0=0$ and $E(M^f_t)^2<\infty$ for every $t$; that is, $M^f\in\cM_2^c$.

\rmk{(a) Only the spatial derivatives $D_if$ enter Step 5, and it suffices that they be bounded on $[0,T]\times\R^d$ for each $T<\infty$. The hypothesis of the last assertion of \bk{Proposition 4.2} (compact support of $f$, with the $\sigma_{ij}$ bounded on it) has no counterpart here, because $\sigma(t,y)$ depends on the whole path of $y$ up to time $t$ and not only on $y(t)$; the global bound \bk{(4.7)} replaces it.}

\rmk{(b) In the sense of \bk{Remark 4.1}, $f\in C^{1,2}([0,\infty)\times\R^d)$ means that $f$ is continuous on $[0,\infty)\times\R^d$ and that its derivatives exist on $(0,\infty)\times\R^d$ and extend continuously to $t=0$; in Step 3 these extensions are used at $s=0$. If one wishes to apply \bk{Theorem 3.3.6} only to functions of class $C^{1,2}$ on an open set containing $[0,\infty)\times\R^d$, one applies it to $f_\eps(t,x)\defeq f(t+\eps,x)$, $\eps>0$, and lets $\eps\downarrow0$ along a sequence. For fixed $\omega\in\Omega^*$ and $t$, the points $(s+\eps,X_s)$ with $0\le s\le t$, $0\le\eps\le1$ lie in a compact subset of $[0,\infty)\times\R^d$, on which $f$ and the extended derivatives are uniformly continuous; so $f_\eps$ and its derivatives, evaluated at $(s,X_s)$, converge to those of $f$ uniformly in $s\in[0,t]$. The integrals with respect to $ds$ then converge by dominated convergence (the integrands are dominated by constant multiples of $1+\abs{b_i(s,X)}+\abs{a_{ik}(s,X)}$), and the stochastic integrals converge in probability by \bk{Proposition 3.2.26}, because $\int_0^t\abs{D_if_\eps(s,X_s)-D_if(s,X_s)}^2\,d\qv{M^{(i)}}_s\to0$ a.s. This gives the formula of Step 3 for $f$, a.s.\ for each fixed $t$, and then for all $t$ simultaneously by continuity.}
\end{kex}

\begin{kex}{Problem}{4.4}{}
We write $W_t=(W^{(1)}_t,\dots,W^{(d)}_t)$ and, for $f\in C^2(\R^d)$,
\begin{equation*}
M^f_t\defeq f(W_t)-f(W_0)-\frac12\int_0^t\Delta f(W_s)\,ds,\qquad
M^{(i)}_t\defeq W^{(i)}_t-W^{(i)}_0;\qquad 0\le t<\infty,\ 1\le i\le d .
\end{equation*}
``A $d$-dimensional Brownian motion'' is understood in the sense of \bk{Definition 2.5.1}, with the distribution of $W_0$ as initial distribution; thus $W$ is a $d$-dimensional Brownian motion precisely when, for $0\le s<t$, the increment $W_t-W_s$ is independent of $\F_s$ and normal with mean zero and covariance matrix $(t-s)I_d$. As everywhere in the stochastic calculus of Chapters 3 and 5 (see the beginning of \bk{Section 3.3}), the filtration $\{\F_t\}$ is assumed to satisfy the usual conditions; the remark at the end shows that the assertion remains true without them.

Every $M^f$ is continuous, adapted and null at $t=0$: the path $s\mapsto \Delta f(W_s)$ is continuous, so $\int_0^t\Delta f(W_s)\,ds$ is continuous in $t$ and is the limit of Riemann sums which are $\F_t$-measurable.

\emph{Necessity.} Let $W$ be a $d$-dimensional Brownian motion. Fix $0\le s<t$ and $1\le i,k\le d$. The random variable $M^{(i)}_t=W^{(i)}_t-W^{(i)}_0$ is normal with mean zero and variance $t$, hence square-integrable, and $E[M^{(i)}_t-M^{(i)}_s\mid\F_s]=E[W^{(i)}_t-W^{(i)}_s]=0$ by the independence of $W_t-W_s$ and $\F_s$. Thus $M^{(i)}$ is a continuous, square-integrable martingale with $M^{(i)}_0=0$. Moreover,
\begin{equation*}
M^{(i)}_tM^{(k)}_t-M^{(i)}_sM^{(k)}_s
=\bigl(M^{(i)}_t-M^{(i)}_s\bigr)\bigl(M^{(k)}_t-M^{(k)}_s\bigr)
+M^{(i)}_s\bigl(M^{(k)}_t-M^{(k)}_s\bigr)+M^{(k)}_s\bigl(M^{(i)}_t-M^{(i)}_s\bigr),
\end{equation*}
where all products are integrable by the Cauchy--Schwarz inequality. Taking conditional expectations given $\F_s$, and using that $M_s$ is $\F_s$-measurable while $M_t-M_s=W_t-W_s$ is independent of $\F_s$ with mean zero and covariance matrix $(t-s)I_d$, we obtain
\begin{equation*}
E\bigl[M^{(i)}_tM^{(k)}_t-M^{(i)}_sM^{(k)}_s\bigm|\F_s\bigr]=\delta_{ik}(t-s)+0+0 .
\end{equation*}
Hence $\{M^{(i)}_tM^{(k)}_t-\delta_{ik}t,\F_t\}$ is a continuous martingale, and the uniqueness assertion of \bk{Problem 1.5.17} gives
\begin{equation}\label{eq:5-4-04:a}
\qv{M^{(i)},M^{(k)}}_t=\delta_{ik}\,t;\qquad 0\le t<\infty,\ \text{a.s.}
\end{equation}
Now let $f\in C^2(\R^d)$. The process $W_t=W_0+M_t$ is a vector of continuous semimartingales with $\F_0$-measurable $W_0$, local martingale part $M=(M^{(1)},\dots,M^{(d)})$ and bounded variation part zero, so It\^o's rule (\bk{Theorem 3.3.6}, applied to the function $(t,x)\mapsto f(x)$) and \eqref{eq:5-4-04:a} yield, a.s.,
\begin{equation*}
f(W_t)=f(W_0)+\sum_{i=1}^d\int_0^t\frac{\partial f}{\partial x_i}(W_s)\,dM^{(i)}_s
+\frac12\sum_{i=1}^d\int_0^t\frac{\partial^2 f}{\partial x_i^2}(W_s)\,ds;\qquad 0\le t<\infty,
\end{equation*}
that is, $M^f$ is indistinguishable from $I\defeq\sum_{i=1}^d\int_0^{\cdot}\frac{\partial f}{\partial x_i}(W_s)\,dM^{(i)}_s$. Each integrand is continuous and adapted, hence progressively measurable with paths bounded on bounded intervals, so it is an admissible integrand for $M^{(i)}$ (\bk{Remark 3.3.4}) and $I\in\cM^{c,\loc}$. If $\{T_n\}$ is a localizing sequence for $I$, then for every $n$ and $t$ we have $M^f_{t\wedge T_n}=I_{t\wedge T_n}$ a.s., and $M^f_{t\wedge T_n}$ is $\F_t$-measurable; consequently $\{M^f_{t\wedge T_n},\F_t\}$ is a martingale along with $\{I_{t\wedge T_n},\F_t\}$. Therefore $M^f\in\cM^{c,\loc}$.

\emph{Sufficiency.} Suppose $M^f\in\cM^{c,\loc}$ for every $f\in C^2(\R^d)$. With $f(x)=x_i$ we have $\Delta f=0$ and $M^f=M^{(i)}$, so
\begin{equation}\label{eq:5-4-04:b}
M^{(i)}\in\cM^{c,\loc};\qquad 1\le i\le d .
\end{equation}
With $f(x)=x_ix_k$ we have $\Delta f=2\delta_{ik}$, so
\begin{equation}\label{eq:5-4-04:c}
N^{(i,k)}_t\defeq W^{(i)}_tW^{(k)}_t-W^{(i)}_0W^{(k)}_0-\delta_{ik}t;\quad 0\le t<\infty,\qquad\text{is in }\cM^{c,\loc}.
\end{equation}
Expanding $M^{(i)}_tM^{(k)}_t=(W^{(i)}_t-W^{(i)}_0)(W^{(k)}_t-W^{(k)}_0)$ we find the identity
\begin{equation}\label{eq:5-4-04:d}
M^{(i)}_tM^{(k)}_t-\delta_{ik}t=N^{(i,k)}_t-W^{(i)}_0M^{(k)}_t-W^{(k)}_0M^{(i)}_t;\qquad 0\le t<\infty .
\end{equation}

\begin{claim}
If $\xi$ is a real-valued, $\F_0$-measurable random variable and $L\in\cM^{c,\loc}$, then $\xi L\in\cM^{c,\loc}$.
\end{claim}
\begin{proof}
Let $\{T_n\}$ be a localizing sequence for $L$ and set $\rho_n\defeq0$ on $\{\abs{\xi}>n\}$, $\rho_n\defeq\infty$ on $\{\abs{\xi}\le n\}$. Since $\{\rho_n\le t\}=\{\abs{\xi}>n\}\in\F_0\subseteq\F_t$, each $\rho_n$ is a stopping time; the sequence $\{\rho_n\}$ is nondecreasing and $\rho_n=\infty$ as soon as $n\ge\abs{\xi}$. Thus $R_n\defeq T_n\wedge\rho_n$ is a nondecreasing sequence of stopping times with $\lim_nR_n=\infty$ a.s. Because $L_0=0$ a.s., we have, a.s. for each $t$,
\begin{equation*}
\xi L_{t\wedge R_n}=\eta_n L_{t\wedge T_n},\qquad \eta_n\defeq\xi\,\ind_{\{\abs{\xi}\le n\}} .
\end{equation*}
The random variable $\eta_n$ is bounded and $\F_0$-measurable, and $\{L_{t\wedge T_n},\F_t\}$ is a martingale; hence $\eta_nL_{t\wedge T_n}$ is integrable, $\F_t$-measurable, and $E[\eta_nL_{t\wedge T_n}\mid\F_s]=\eta_nL_{s\wedge T_n}$ for $s<t$. So $\{\xi L_{t\wedge R_n},\F_t\}$ is a martingale for every $n$; as $\xi L$ is continuous, adapted and null at $t=0$ a.s., $\xi L\in\cM^{c,\loc}$.
\end{proof}

By \eqref{eq:5-4-04:b}, \eqref{eq:5-4-04:c} and the claim (with $\xi=W^{(i)}_0$, $W^{(k)}_0$), each of the three processes on the right-hand side of \eqref{eq:5-4-04:d} belongs to $\cM^{c,\loc}$. This class is a vector space: if $\{T_n\}$, $\{T_n'\}$ localize $L$, $L'$, then $\{T_n\wedge T_n'\}$ localizes both, because a continuous martingale stopped at a stopping time is again a martingale (\bk{Problem 1.3.24}(i)). Therefore $M^{(i)}M^{(k)}-\delta_{ik}t\in\cM^{c,\loc}$. Since $t\mapsto\delta_{ik}t$ is continuous, adapted, of bounded variation and null at $t=0$, the uniqueness part of \bk{Problem 1.5.17} shows that \eqref{eq:5-4-04:a} holds. Together with \eqref{eq:5-4-04:b}, this is exactly the hypothesis of P.~L\'evy's theorem (\bk{Theorem 3.3.16}) for the continuous, adapted process $W$, and we conclude that $\{W_t,\F_t;\ 0\le t<\infty\}$ is a $d$-dimensional Brownian motion.

\rmk{\emph{Filtrations which do not satisfy the usual conditions.} Let $\{\F_t\}$ be arbitrary, and let $\G_t$ be the $\sigma$-field generated by $\F_{t+}$ and the $P$-negligible sets ($P$ being extended to the completion of $\F$); then $\{\G_t\}$ satisfies the usual conditions, $\F_t\subseteq\G_t$, and the solution above applies to $\{W_t,\G_t\}$. It therefore suffices to show that each of the two properties in the statement holds relative to $\{\F_t\}$ if and only if it holds relative to $\{\G_t\}$; negligible sets affect neither independence nor conditional expectations, so only the passage from $\F_s$ to $\F_{s+}$ matters.

\emph{Brownian property.} If $W_t-W_s$ is independent of $\G_s$, it is independent of $\F_s\subseteq\G_s$. Conversely, let $W$ be a Brownian motion relative to $\{\F_t\}$, $s<t$, $A\in\F_{s+}$ and $g:\R^d\to\R$ bounded and continuous. For $1/n<t-s$ we have $A\in\F_{s+1/n}$, so $E[g(W_t-W_{s+1/n})\ind_A]=E[g(W_t-W_{s+1/n})]\,P(A)$; letting $n\to\infty$ (continuity of paths, bounded convergence) gives $E[g(W_t-W_s)\ind_A]=E[g(W_t-W_s)]\,P(A)$. Hence $W_t-W_s$ is independent of $\F_{s+}$, and of $\G_s$.

\emph{Local martingale property.} Suppose $M^f$ is a local martingale relative to $\{\F_t\}$, with localizing sequence $\{T_n\}$, and put $N_t=M^f_{t\wedge T_n}$. For $s<t$ and $A\in\F_{s+}$ we have $E[N_t\ind_A]=E[N_{s+1/m}\ind_A]$ for $1/m<t-s$; as $m\to\infty$, $N_{s+1/m}\to N_s$ and the family $\{N_u=E[N_t\mid\F_u];\,u\le t\}$ is uniformly integrable, so $E[N_t\ind_A]=E[N_s\ind_A]$. Thus $N$ is a martingale relative to $\{\G_t\}$; the $T_n$ are stopping times of $\{\G_t\}$, and $M^f$ is a local martingale relative to $\{\G_t\}$. Conversely, suppose $M^f$ is a local martingale relative to $\{\G_t\}$, localized by $\{T_n\}$. The times $S_n\defeq\inf\{t\ge0;\ \norm{W_t}\ge n\}$ are stopping times of $\{\F_t\}$ (\bk{Problem 1.2.7}) which increase to infinity, because the paths of $W$ are bounded on bounded intervals. On $\{S_n=0\}$ we have $M^f_{t\wedge S_n}=0$, and on $\{S_n>0\}$ we have $\norm{W_u}\le n$ for $0\le u\le S_n$; consequently
\begin{equation*}
\abs[\big]{M^f_{t\wedge S_n}}\le 2\max_{\norm{x}\le n}\abs{f(x)}+\frac t2\max_{\norm{x}\le n}\abs{\Delta f(x)};\qquad 0\le t<\infty .
\end{equation*}
For $s<t$, the martingale $\{M^f_{u\wedge T_m},\G_u\}$ stopped at $S_n$ is again a martingale (\bk{Problem 1.3.24}(i)), so $E[M^f_{t\wedge S_n\wedge T_m}\mid\G_s]=M^f_{s\wedge S_n\wedge T_m}$; letting $m\to\infty$ with the help of the above bound and the bounded convergence theorem for conditional expectations, we get $E[M^f_{t\wedge S_n}\mid\G_s]=M^f_{s\wedge S_n}$. Since $M^f_{\cdot\wedge S_n}$ is adapted to $\{\F_t\}$ and $\F_s\subseteq\G_s$, conditioning on $\F_s$ shows that $\{M^f_{t\wedge S_n},\F_t\}$ is a martingale for every $n$; that is, $M^f$ is a local martingale relative to $\{\F_t\}$.}
\end{kex}

\begin{kex}{Exercise}{4.33}{}
Throughout, $a\defeq\sigma\sigma^{T}$, so that $a_{ik}=\sum_{j=1}^r\sigma_{ij}\sigma_{kj}$ and $b_i$ are Borel measurable and bounded on compact sets; since the derivatives of $f$ up to order two are continuous, the function $\cA f$ of \bk{(1.2)} is Borel measurable and bounded on compact subsets of $\R^d$. The process $X$ is adapted to $\{\F_t\}$ (this is the meaning of the notation $X=\{X_t,\F_t;\ 0\le t<\infty\}$) and continuous, hence progressively measurable (\bk{Proposition 1.1.13}). We write $Y_t\defeq f(X_t)$ and $\gamma_t\defeq\cA f(X_t)$.

\emph{Step 1. The three processes are well defined, continuous and adapted, and vanish at $t=0$.} Let $h\colon\R^d\to\R$ be Borel measurable and bounded on compact sets, and let $\varphi\colon[0,\infty)\to\R$ be continuous. For every $\omega$ and $t$, the path $s\mapsto X_s(\omega)$ maps $[0,t]$ into a compact set, so $s\mapsto\varphi(s)h(X_s(\omega))$ is a bounded Borel function on $[0,t]$; thus $J_t(\omega)\defeq\int_0^t\varphi(s)h(X_s(\omega))\,ds$ is a finite number, and $t\mapsto J_t(\omega)$ is locally absolutely continuous. For fixed $t$, the map $(s,\omega)\mapsto\varphi(s)h(X_s(\omega))$ is $\B([0,t])\otimes\F_t$-measurable on $[0,t]\times\Omega$ by progressive measurability. For $n\in\N$ the set $\Omega_n\defeq\{\sup_{0\le s\le t}\norm{X_s}\le n\}$ belongs to $\F_t$ (the supremum may be taken over rational $s$ and $s=t$), on $[0,t]\times\Omega_n$ this map is bounded, so $J_t\ind_{\Omega_n}$ is $\F_t$-measurable by Fubini's theorem; since $\Omega_n\uparrow\Omega$, $J_t$ is $\F_t$-measurable. Applying this with $h=\cA f$, $\varphi\equiv1$, and with $h=\alpha f-\cA f$, $\varphi(s)=e^{-\alpha s}$, we see that
\[
A_t\defeq\int_0^t\gamma_s\,ds,\qquad M_t=Y_t-Y_0-A_t,\qquad \Lambda_t
\]
are continuous, $\{\F_t\}$-adapted processes with $M_0=\Lambda_0=0$. If $f$ is bounded away from zero on compact sets, i.e.\ $\inf_K\abs{f}>0$ for every compact $K\subset\R^d$, then $f$ has no zeros and $g\defeq\cA f/f$ is Borel measurable and bounded on compact sets; consequently
\[
D_t\defeq\exp\Bigl\{-\int_0^t g(X_s)\,ds\Bigr\},\qquad N_t=Y_tD_t-Y_0
\]
are continuous and adapted, $D$ is strictly positive, and $N_0=0$.

\emph{Step 2. A pathwise inversion formula.} Let $\cC$ denote the class of continuous, $\{\F_t\}$-adapted processes $C$ which are nowhere zero and of the form $C_t=C_0+\int_0^t c_s\,ds$, where, for every $\omega$, $s\mapsto c_s(\omega)$ is a Borel function bounded on every bounded interval. Each $C\in\cC$ has paths of bounded variation on bounded intervals. For $C\in\cC$ and a continuous adapted process $Z$ we set
\[
I(C,Z)_t\defeq C_tZ_t-\int_0^t Z_s\,dC_s=C_tZ_t-\int_0^t Z_sc_s\,ds;\qquad 0\le t<\infty .
\]
This is again a continuous adapted process: the integral is, for every $\omega$, the limit of the Riemann--Stieltjes sums $\sum_{k=0}^{n-1} Z_{tk/n}\,(C_{t(k+1)/n}-C_{tk/n})$, each of which is $\F_t$-measurable.

\begin{lem}
If $C\in\cC$, then $1/C\in\cC$, and for every continuous adapted process $Z$ we have
\[
I\bigl(1/C,\,I(C,Z)\bigr)_t=Z_t\qquad\text{for all }t\ge0\text{ and all }\omega\in\Omega .
\]
\end{lem}

\begin{proof}
Fix $\omega$. The function $t\mapsto C_t$ is locally absolutely continuous and, being continuous without zeros, is bounded away from zero on every bounded interval. Hence $E\defeq1/C$ is locally absolutely continuous with $E_t=E_0+\int_0^t e_s\,ds$, where $e_s\defeq-c_s/C_s^2$ is Borel and bounded on bounded intervals. As $E$ is also continuous, adapted and nowhere zero, $1/C\in\cC$.

Put $H_t\defeq\int_0^tZ_sc_s\,ds$ and $V\defeq I(C,Z)=CZ-H$. The functions $E$ and $H$ are locally absolutely continuous and $H_0=0$, so the product rule for absolutely continuous functions gives
\[
E_tH_t=\int_0^t\bigl(E_sZ_sc_s+H_se_s\bigr)\,ds .
\]
Therefore, since $E_tC_t=1$,
\begin{align*}
I(E,V)_t&=E_tV_t-\int_0^tV_se_s\,ds
=Z_t-E_tH_t-\int_0^t\bigl(C_sZ_s-H_s\bigr)e_s\,ds\\
&=Z_t-\int_0^tZ_s\bigl(E_sc_s+C_se_s\bigr)\,ds=Z_t,
\end{align*}
because $E_sc_s+C_se_s=c_s/C_s-c_s/C_s=0$.
\end{proof}

\emph{Step 3. The transformation $I(C,\cdot)$ preserves $\cM^{c,\loc}$.} Let $C\in\cC$ and $Z\in\cM^{c,\loc}$. We apply the integration-by-parts formula \bk{(3.3.8)} of \bk{Problem 3.3.12} to the continuous semimartingales $Z=0+Z+0$ and $C=C_0+0+(C-C_0)$, whose local martingale parts are $Z$ and $0$ (note that $C_0$ is $\F_0$-measurable and $C-C_0$ is continuous, adapted, of bounded variation and null at $0$). The cross-variation term vanishes and we obtain, a.s.\ for all $t\ge0$,
\[
I(C,Z)_t=C_tZ_t-\int_0^tZ_s\,dC_s=\int_0^tC_s\,dZ_s .
\]
The integrand $C$ is continuous and adapted, hence progressively measurable with $\int_0^tC_s^2\,d\qv{Z}_s<\infty$ for all $t$; so the right-hand side is a stochastic integral in the sense of \bk{Definition 3.2.23} and belongs to $\cM^{c,\loc}$. The process $I(C,Z)$ is continuous, adapted and indistinguishable from this element of $\cM^{c,\loc}$, hence is itself in $\cM^{c,\loc}$: it vanishes at $0$, and stopped at the same localizing sequence it is, at each time, a.s.\ equal to a martingale. This is the statement recalled in the hint.

Combining this with the Lemma, we have for every $C\in\cC$ and every continuous adapted process $Z$:
\begin{equation}\label{eq:5-4-33:equiv}
Z\in\cM^{c,\loc}\iff I(C,Z)\in\cM^{c,\loc}.
\end{equation}
Indeed, ``$\Rightarrow$'' was just proved; and if $V\defeq I(C,Z)\in\cM^{c,\loc}$, then $Z=I(1/C,V)\in\cM^{c,\loc}$ by the same implication applied to $1/C\in\cC$ and $V$.

\emph{Step 4. $\Lambda=I(C,M)$ with $C_t=e^{-\alpha t}$.} This $C$ is in $\cC$, with $c_s=-\alpha e^{-\alpha s}$. For every $\omega$ and $t$, substituting $M_s=Y_s-Y_0-A_s$ we find
\begin{align*}
I(C,M)_t&=e^{-\alpha t}M_t+\alpha\int_0^te^{-\alpha s}M_s\,ds\\
&=e^{-\alpha t}Y_t+\alpha\int_0^te^{-\alpha s}Y_s\,ds
-Y_0\Bigl(e^{-\alpha t}+\alpha\int_0^te^{-\alpha s}\,ds\Bigr)\\
&\qquad-\Bigl(e^{-\alpha t}A_t+\alpha\int_0^te^{-\alpha s}A_s\,ds\Bigr).
\end{align*}
Here $e^{-\alpha t}+\alpha\int_0^te^{-\alpha s}\,ds=1$, and the product rule for the absolutely continuous functions $t\mapsto e^{-\alpha t}$ and $t\mapsto A_t$ gives $e^{-\alpha t}A_t=\int_0^t\bigl(e^{-\alpha s}\gamma_s-\alpha e^{-\alpha s}A_s\bigr)\,ds$, so that the last bracket equals $\int_0^te^{-\alpha s}\gamma_s\,ds$. Hence
\[
I(C,M)_t=e^{-\alpha t}f(X_t)-f(X_0)+\int_0^te^{-\alpha s}\bigl(\alpha f(X_s)-\cA f(X_s)\bigr)\,ds=\Lambda_t ,
\]
and \eqref{eq:5-4-33:equiv} with $Z=M$ yields $M\in\cM^{c,\loc}\iff\Lambda\in\cM^{c,\loc}$.

\emph{Step 5. $N=I(D,M)$.} Assume now that $f$ is bounded away from zero on compact sets. By Step 1, $D$ is continuous, adapted and strictly positive; for every $\omega$ it is the exponential of a locally absolutely continuous function, hence $D_t=1+\int_0^tc_s\,ds$ with $c_s\defeq-g(X_s)D_s$, which is Borel in $s$ and bounded on bounded intervals. So $D\in\cC$. Note that $\gamma_s=g(X_s)Y_s$. For every $\omega$ and $t$, the product rule for the absolutely continuous functions $A$ and $D$ gives
\[
A_tD_t=\int_0^t\bigl(D_s\gamma_s-A_sg(X_s)D_s\bigr)\,ds=\int_0^tg(X_s)D_s\bigl(Y_s-A_s\bigr)\,ds ,
\]
and $Y_0D_t=Y_0-\int_0^tg(X_s)D_sY_0\,ds$. Therefore, writing $Y_t=Y_0+M_t+A_t$,
\begin{align*}
N_t=Y_tD_t-Y_0&=M_tD_t+\int_0^tg(X_s)D_s\bigl(Y_s-A_s-Y_0\bigr)\,ds\\
&=D_tM_t+\int_0^tM_s\,g(X_s)D_s\,ds=D_tM_t-\int_0^tM_s\,dD_s=I(D,M)_t .
\end{align*}
By \eqref{eq:5-4-33:equiv} with $C=D$ and $Z=M$, we conclude $M\in\cM^{c,\loc}\iff N\in\cM^{c,\loc}$. Together with Step 4, the three conditions $M\in\cM^{c,\loc}$, $\Lambda\in\cM^{c,\loc}$, $N\in\cM^{c,\loc}$ are equivalent.

\rmk{Steps 3--5 give the representations $\Lambda_t=\int_0^te^{-\alpha s}\,dM_s$ and $N_t=\int_0^tD_s\,dM_s$ when $M\in\cM^{c,\loc}$, and conversely $M_t=\int_0^te^{\alpha s}\,d\Lambda_s$, resp.\ $M_t=\int_0^tD_s^{-1}\,dN_s$, when $\Lambda$, resp.\ $N$, is in $\cM^{c,\loc}$. No integrability of $f(X_0)$ is needed, since all three processes start at zero.}
\end{kex}

\begin{kex}{Exercise}{4.34}{}
\emph{Conventions.} The filtration $\{\F_t\}$ of a weak solution satisfies the usual conditions (\bk{Definition 3.1}\,(i), which \bk{Definition 3.14} retains); in particular every $P$-null set of $\F$ belongs to $\F_0$. ``Locally integrable'' means that $P(\Omega_0)=1$, where $\Omega_0\defeq\{\int_0^n\abs{k_u}\,du<\infty\text{ for all }n\in\N\}$. As in \bk{(4.2)$'$}, the integrand of $\Lambda$ is evaluated at $(s,X_s)$, the coefficients in $\cA_s'f$ at $(s,X)$; thus, with
\begin{equation}\label{eq:5-4-34:g}
g_s\defeq\frac{\partial f}{\partial s}(s,X_s)
+\frac12\sum_{i=1}^d\sum_{l=1}^d a_{il}(s,X)\frac{\partial^2 f}{\partial x_i\partial x_l}(s,X_s)
+\sum_{i=1}^d b_i(s,X)\frac{\partial f}{\partial x_i}(s,X_s),
\end{equation}
$Y_t\defeq f(t,X_t)$ and $C_t\defeq\exp\{-\int_0^tk_u\,du\}$, the process in question is
\begin{equation}\label{eq:5-4-34:L}
\Lambda_t=Y_tC_t-Y_0-\int_0^t\bigl(g_s-k_sY_s\bigr)C_s\,ds .
\end{equation}
We read condition \bk{(A.1)$'$} of \bk{Problem 4.13} as: for every $0<T<\infty$ there is a constant $K_T$ with $\norm{b(t,y)}\le K_T$ and $\abs{\sigma_{ij}(t,y)}\le K_T$ for all $i,j$, $0\le t\le T$ and $y\in C[0,\infty)^d$; this is equivalent, up to the value of the constant, to a bound on $\norm{b(t,y)}+\norm{\sigma(t,y)}$.

\emph{Step 1. The processes involved are well defined, continuous and adapted.} For fixed $t$, the map $(u,\omega)\mapsto k_u(\omega)$ on $[0,t]\times\Omega$ is $\B([0,t])\otimes\F_t$-measurable, so by Tonelli's theorem $\int_0^tk_u^{+}\,du$ and $\int_0^tk_u^{-}\,du$ are $\F_t$-measurable with values in $[0,\infty]$. Hence $\Omega_0\in\F$, and since $P(\Omega_0)=1$ we have $\Omega_0\in\F_0$. We set
\[
\kappa_t\defeq\ind_{\Omega_0}\int_0^tk_u\,du,\qquad C_t\defeq e^{-\kappa_t};\qquad 0\le t<\infty,
\]
which fixes the definition of $C$ on the null set $\Omega\setminus\Omega_0$. The process $\kappa$ is adapted and, for every $\omega$, $t\mapsto\kappa_t(\omega)$ is locally absolutely continuous. Therefore $C$ is adapted, continuous and strictly positive; on each bounded interval $\kappa_\cdot(\omega)$ takes values in a compact set, on which the exponential function is Lipschitz, so $C_\cdot(\omega)$ is locally absolutely continuous with
\begin{equation}\label{eq:5-4-34:C}
C_t=1-\int_0^tk_sC_s\,ds;\qquad 0\le t<\infty,\ \omega\in\Omega_0 .
\end{equation}
In particular, $C$ has paths of bounded variation on bounded intervals.

The process $Y$ is continuous and adapted. By \bk{Problem 4.3}, the process
\[
M_t\defeq M^f_t=Y_t-Y_0-\int_0^tg_s\,ds;\qquad 0\le t<\infty
\]
of \bk{(4.2)$'$} belongs to $\cM^{c,\loc}$; here $\int_0^t\abs{g_s}\,ds<\infty$ for all $t$, a.s., because the derivatives of $f$ are continuous, hence bounded along $\{(s,X_s);\,0\le s\le t\}$, while $\abs{a_{il}}\le\frac12\sum_{j=1}^r(\sigma_{ij}^2+\sigma_{lj}^2)$ and condition (iii) of \bk{Definition 3.14} holds (under \bk{(A.1)$'$}, $b(s,X)$ and $a(s,X)$ are even bounded on $[0,t]$). Consequently $A_t\defeq\int_0^tg_s\,ds=Y_t-Y_0-M_t$ is, after modification on a null set, a continuous adapted process of bounded variation with $A_0=0$, and
\begin{equation}\label{eq:5-4-34:Y}
Y_t=Y_0+M_t+A_t;\qquad 0\le t<\infty
\end{equation}
is a continuous semimartingale with $\F_0$-measurable initial value $Y_0=f(0,X_0)$. Finally, on the event of probability one where $\int_0^t(\abs{g_s}+\abs{k_s})\,ds<\infty$ for all $t$, the integrand $(g_s-k_sY_s)C_s$ in \eqref{eq:5-4-34:L} is integrable on bounded intervals, because $Y$ and $C$ are continuous. So $\Lambda$ is defined by \eqref{eq:5-4-34:L} outside a null set, on which we set $\Lambda\equiv0$; it is continuous, and $\Lambda_0=0$.

\emph{Step 2. $\Lambda$ is a stochastic integral.} We apply the integration-by-parts formula \bk{(3.3.8)} of \bk{Problem 3.3.12} to the continuous semimartingales $Y$ of \eqref{eq:5-4-34:Y} and $C=1+0+(C-1)$, whose local martingale part is zero. The cross-variation term vanishes, and we obtain, a.s.\ for all $0\le t<\infty$,
\[
Y_tC_t-Y_0=\int_0^tC_s\,dY_s+\int_0^tY_s\,dC_s
=\int_0^tC_s\,dM_s+\int_0^tC_sg_s\,ds-\int_0^tY_sk_sC_s\,ds .
\]
We have used \eqref{eq:5-4-34:C}: for $\omega\in\Omega_0$ the Lebesgue--Stieltjes measure $dC_s$ has density $-k_sC_s$ with respect to Lebesgue measure, and likewise $dA_s=g_s\,ds$. The integrand $C$ of the stochastic integral is continuous and adapted, hence progressively measurable with $\int_0^tC_s^2\,d\qv{M}_s<\infty$ for all $t$, so $C\in\cP^*(M)$ and $\int_0^\cdot C_s\,dM_s\in\cM^{c,\loc}$ (\bk{Definition 3.2.23}). Comparing with \eqref{eq:5-4-34:L}, we have shown
\begin{equation}\label{eq:5-4-34:rep}
\Lambda_t=\int_0^tC_s\,dM_s;\qquad 0\le t<\infty,\quad\text{a.s.}
\end{equation}
Thus $\Lambda$ is indistinguishable from an element of $\cM^{c,\loc}$. Since $\F_0$ contains the $P$-null sets, $\Lambda$ is adapted, and for a localizing sequence $\{T_n\}$ of the right-hand side of \eqref{eq:5-4-34:rep} each $\{\Lambda_{t\wedge T_n},\F_t\}$ is a martingale. Hence $\Lambda\in\cM^{c,\loc}$.

\emph{Step 3. The quadratic variation of $\Lambda$.} By \bk{Proposition 3.2.24}, relation \bk{(3.2.30)} holds for the integral \eqref{eq:5-4-34:rep} and every $N\in\cM^{c,\loc}$: $\qv{\Lambda,N}_t=\int_0^tC_s\,d\qv{M,N}_s$. With $N=M$ this gives $d\qv{\Lambda,M}_s=C_s\,d\qv{M}_s$, and then with $N=\Lambda$,
\[
\qv{\Lambda}_t=\int_0^tC_s\,d\qv{M,\Lambda}_s=\int_0^tC_s^2\,d\qv{M}_s;\qquad 0\le t<\infty,\quad\text{a.s.}
\]
On the other hand, \bk{(4.3)$'$} with $g=f$ and \bk{(4.6)} yield
\begin{equation}\label{eq:5-4-34:qv}
\qv{M}_t=\int_0^t\sum_{j=1}^r\Bigl(\sum_{i=1}^d\sigma_{ij}(s,X)\frac{\partial f}{\partial x_i}(s,X_s)\Bigr)^2ds
\le\int_0^t\norm{\sigma(s,X)}^2\,\norm{\nabla f(s,X_s)}^2\,ds ,
\end{equation}
by the Cauchy--Schwarz inequality, where $\norm{\sigma}^2=\sum_{i,j}\sigma_{ij}^2$ and $\nabla f$ is the gradient of $f$ in the space variable.

\emph{Step 4. The martingale property.} Assume now that $\norm{\nabla f}\le G$ on $[0,\infty)\times\R^d$ and $k_t(\omega)\ge-\beta$ for all $(t,\omega)$, for some constants $G,\beta\in[0,\infty)$. Then $\kappa_t\ge-\beta t$ and $0<C_t\le e^{\beta t}$. Fix $0<T<\infty$. By \bk{(A.1)$'$}, $\norm{\sigma(s,X)}^2\le dr\,K_T^2$ for $0\le s\le T$, so \eqref{eq:5-4-34:qv} and Step 3 give
\begin{equation}\label{eq:5-4-34:bd}
\qv{\Lambda}_T\le c_T\defeq e^{2\beta T}\,dr\,K_T^2\,G^2\,T<\infty,\quad\text{a.s.}
\end{equation}
The assertion now follows from the following fact.

\begin{lem}
Let $L\in\cM^{c,\loc}$ and suppose that for every $0<T<\infty$ there is a constant $c_T$ with $\qv{L}_T\le c_T$ a.s. Then $L$ is a martingale, and $E(L_T^2)\le c_T$.
\end{lem}

\begin{proof}
We first note that if $Z\in\cM^{c,\loc}$ satisfies $\sup_{0\le s\le t}\abs{Z_s}\le c$ a.s.\ for some constants $t,c$, then $E[Z_t\mid\F_s]=Z_s$ a.s.\ for $0\le s\le t$. Indeed, if $\{S_m\}$ localizes $Z$, then $E[Z_{t\wedge S_m}\mid\F_s]=Z_{s\wedge S_m}$, and since $S_m\uparrow\infty$ a.s.\ we may let $m\to\infty$, using the bounded convergence theorem for conditional expectations.

Let $T_n\defeq\inf\{t\ge0;\ \abs{L_t}\ge n\}$, a stopping time by \bk{Problem 1.2.7}; $T_n\uparrow\infty$ because the paths of $L$ are continuous, and $\abs{L_{t\wedge T_n}}\le n$ since $L_0=0$. If $\{S_m\}$ localizes $L$, then each $\{L_{t\wedge T_n\wedge S_m},\F_t\}$ is a martingale (\bk{Problem 1.3.24}\,(i)), so $L^{(n)}\defeq L_{\cdot\wedge T_n}\in\cM^{c,\loc}$, and by the preceding paragraph $L^{(n)}$ is a bounded martingale. Next, $L^2-\qv{L}\in\cM^{c,\loc}$ by the definition of $\qv{L}$ (\bk{Problem 1.5.17}), and stopping at $T_n$ as before, $Z^{(n)}\defeq(L^{(n)})^2-\qv{L}_{\cdot\wedge T_n}$ is in $\cM^{c,\loc}$. For $t\le T$ we have $\abs{Z^{(n)}_t}\le n^2+c_T$ a.s., so the first paragraph gives $E(Z^{(n)}_T)=Z^{(n)}_0=0$, that is,
\[
E\bigl(L_{T\wedge T_n}^2\bigr)=E\,\qv{L}_{T\wedge T_n}\le c_T;\qquad n\ge1 .
\]
Since $T$ was arbitrary, for each $t>0$ the sequence $\{L_{t\wedge T_n}\}_{n\ge1}$ is bounded in $L^2$ by $c_t$, thus uniformly integrable, and it converges a.s.\ to $L_t$; so $L_{t\wedge T_n}\to L_t$ in $L^1$, and $E(L_t^2)\le c_t$ by Fatou's lemma. For $0\le s<t$ and $A\in\F_s$ we have $E[L_{t\wedge T_n}\ind_A]=E[L_{s\wedge T_n}\ind_A]$, and letting $n\to\infty$ we obtain $E[L_t\ind_A]=E[L_s\ind_A]$. Therefore $L$ is a martingale.
\end{proof}

By \eqref{eq:5-4-34:bd} the Lemma applies to $L=\Lambda$: the process $\{\Lambda_t,\F_t;\ 0\le t<\infty\}$ is a martingale, and in fact $\Lambda\in\cM_2^c$ with $E(\Lambda_T^2)\le c_T$.

\rmk{(a) Condition \bk{(A.1)$'$} is not needed for the first assertion, which rests only on \bk{Problem 4.3}; and of \bk{(A.1)$'$} only the bound \bk{(4.7)} on $\sigma$ enters the second. (b) For the second assertion we have used only the boundedness of the first-order space derivatives of $f$ and the lower bound on $k$; the boundedness of $f$, $\partial f/\partial t$ and of the second derivatives is not needed for the martingale property of $\Lambda$, although under these additional hypotheses and \bk{(A.1)$'$} each of the three terms of \eqref{eq:5-4-34:L} is separately bounded on $[0,T]\times\Omega_0$: then $\abs{Y_tC_t}\le e^{\beta t}\sup\abs{f}$, the integrand is dominated on $[0,T]$ by a constant multiple of $(1+\abs{k_s})C_s$, and $\int_0^t\abs{k_s}C_s\,ds\le1+2\beta\int_0^te^{\beta s}\,ds$ by \eqref{eq:5-4-34:C} and $\abs{k_s}\le k_s+2\beta$. (c) For $k_t\equiv\alpha$ and time-homogeneous coefficients, \eqref{eq:5-4-34:rep} reads $\Lambda_t=\int_0^te^{-\alpha s}\,dM_s$; this is the identity behind the hint of \bk{Exercise 4.33}, obtained from \bk{Problem 3.3.12} with $C_t=e^{-\alpha t}$.}
\end{kex}

\begin{kex}{Exercise}{4.35}{}
As in \bk{Definition 4.15}, $(\Omega,\B)=(C[0,\infty)^d,\B(C[0,\infty)^d))$, $y$ is the coordinate process, $\B_t=\sigma(y(s);\,0\le s\le t)$, the coefficients $b_i$, $\sigma_{ij}$ are Borel measurable, and $\cA$ is the operator \bk{(1.2)} with $a=\sigma\sigma^T$. We fix $x\in\R^d$ and $t>0$ (for $t=0$ the assertion is $E^xf(y(0))=f(x)$, which is \bk{(4.22)}), write $P=P^x$, $E=E^x$, and put $h\defeq\cA f+\lambda f$. Since $f\in C^2(\R^d)$ and $b$, $a$ are Borel measurable and bounded on compact sets, $h$ is Borel measurable and bounded on compact sets, and $h\le c$ by assumption.

\emph{Step 1. Localization.} For every integer $k\ge1$ let
\[
S_k\defeq\inf\{s\ge0:\ \norm{y(s)}\ge k\}\qquad(\inf\emptyset=\infty).
\]
By continuity of the paths, $\{S_k\le s\}=\{\max_{0\le u\le s}\norm{y(u)}\ge k\}\in\B_s$ for every $s\ge0$ (the maximum may be computed over the rationals in $[0,s]$ and $u=s$), so $S_k$ is a stopping time of $\{\B_s\}$; moreover
\begin{equation}\label{eq:5-4-35:ball}
\norm{y(u)}\le k\quad\text{for } 0\le u\le S_k,\ u<\infty,\qquad\text{on }\{\norm{y(0)}\le k\},
\end{equation}
and, since every path is bounded on $[0,t]$, we have $S_k(\omega)>t$ for all $k>\max_{0\le u\le t}\norm{\omega(u)}$. Let $\chi_k\in C^\infty(\R^d)$ have compact support and be equal to $1$ on $\{z:\norm{z}\le k+1\}$, and set $g\defeq g_k\defeq f\chi_k\in C^2_0(\R^d)$. As $g=f$ on the open ball of radius $k+1$, we have
\begin{equation}\label{eq:5-4-35:agree}
g(z)=f(z),\qquad \cA g(z)+\lambda g(z)=h(z)\qquad\text{for }\norm{z}\le k.
\end{equation}
The function $\cA g$ is Borel measurable and bounded, because it vanishes off the compact support of $g$, on which $b$, $a$ and the derivatives of $g$ are bounded. Consequently
\[
M_s\defeq g(y(s))-g(y(0))-\int_0^s\cA g(y(u))\,du;\qquad 0\le s<\infty,
\]
is a continuous, $\{\B_s\}$-adapted process (adaptedness of the integral follows from Fubini's theorem, the integrand being bounded and $\B([0,s])\otimes\B_s$-measurable on $[0,s]\times\Omega$) with
\begin{equation}\label{eq:5-4-35:bd}
\abs{M_s}\le 2\sup\abs{g}+s\,\sup\abs{\cA g};\qquad 0\le s<\infty,
\end{equation}
and \bk{(4.21)} says precisely that $\{M_s,\B_s;\ 0\le s<\infty\}$ is a martingale under $P$.

\emph{Step 2. The discounted martingale.} We claim that
\[
L_s\defeq e^{\lambda s}M_s-\lambda\int_0^se^{\lambda u}M_u\,du;\qquad 0\le s<\infty,
\]
is a continuous $\{\B_s\}$-martingale under $P$. It is continuous, adapted (Fubini's theorem again, $(u,\omega)\mapsto M_u(\omega)$ being continuous in $u$ and $\B_s$-measurable in $\omega$ for $u\le s$) and bounded on bounded time-intervals by \eqref{eq:5-4-35:bd}. Let $0\le s<r<\infty$ and $F\in\B_s$. By \eqref{eq:5-4-35:bd} we may use Fubini's theorem, and then the martingale property $E[\ind_FM_u]=E[\ind_FM_s]$ for $u\ge s$:
\begin{align*}
E\bigl[\ind_F(L_r-L_s)\bigr]
&=e^{\lambda r}E[\ind_FM_r]-e^{\lambda s}E[\ind_FM_s]-\lambda\int_s^re^{\lambda u}E[\ind_FM_u]\,du\\
&=E[\ind_FM_s]\Bigl(e^{\lambda r}-e^{\lambda s}-\lambda\int_s^re^{\lambda u}\,du\Bigr)=0 .
\end{align*}
This proves the claim. Next we compute $L$ pathwise. Write $G_s\defeq g(y(s))$ and $A_s\defeq\int_0^s\cA g(y(u))\,du$, so that $M=G-G_0-A$. The function $s\mapsto A_s$ is Lipschitz continuous with $A_0=0$, hence the product rule for absolutely continuous functions gives $e^{\lambda s}A_s=\int_0^se^{\lambda u}\bigl(\cA g(y(u))+\lambda A_u\bigr)\,du$; also $e^{\lambda s}-\lambda\int_0^se^{\lambda u}\,du=1$. Therefore, for every $\omega\in\Omega$ and $s\ge0$,
\begin{align}
L_s&=e^{\lambda s}G_s-\lambda\int_0^se^{\lambda u}G_u\,du-G_0
-\Bigl(e^{\lambda s}A_s-\lambda\int_0^se^{\lambda u}A_u\,du\Bigr)\notag\\
&=e^{\lambda s}g(y(s))-g(y(0))-\int_0^se^{\lambda u}\bigl(\cA g+\lambda g\bigr)(y(u))\,du .
\label{eq:5-4-35:L}
\end{align}

\emph{Step 3. Optional sampling.} Let $k>\norm{x}$ and $T\defeq t\wedge S_k$, a stopping time of $\{\B_s\}$ bounded by $t$. The optional sampling theorem for bounded optional times, \bk{Problem 1.3.23}\,(i) (which requires only right-continuity of the paths, not the usual conditions on the filtration), applied to the continuous martingales $L$ and $-L$ and the times $0\le T$, gives $E[L_T]=E[L_0]=0$. By \bk{(4.22)}, $y(0)=x$ holds $P$-a.s.; on this event \eqref{eq:5-4-35:ball} gives $\norm{y(u)}\le k$ for $0\le u\le T$, and then \eqref{eq:5-4-35:L}, \eqref{eq:5-4-35:agree} yield
\[
L_T=e^{\lambda T}f(y(T))-f(x)-\int_0^Te^{\lambda u}h(y(u))\,du ,\qquad P\text{-a.s.}
\]
The three terms on the right are bounded random variables ($f$ and $h$ are bounded on the ball of radius $k$, and $T\le t$). Taking expectations and using $h\le c$, $c\ge0$ and $T\le t$, we obtain
\begin{equation}\label{eq:5-4-35:k}
E\bigl[e^{\lambda(t\wedge S_k)}f\bigl(y(t\wedge S_k)\bigr)\bigr]
=f(x)+E\int_0^{t\wedge S_k}e^{\lambda u}h(y(u))\,du
\le f(x)+\frac{c}{\lambda}\bigl(e^{\lambda t}-1\bigr)
\end{equation}
for every integer $k>\norm{x}$.

\emph{Step 4. Passage to the limit.} For every $\omega$ we have $t\wedge S_k(\omega)=t$ for all sufficiently large $k$ (Step 1), so the nonnegative random variables $e^{\lambda(t\wedge S_k)}f(y(t\wedge S_k))$ converge pointwise to $e^{\lambda t}f(y(t))$ as $k\to\infty$. Fatou's lemma and \eqref{eq:5-4-35:k} give
\[
e^{\lambda t}E^xf(y(t))\le\liminf_{k\to\infty}E^x\bigl[e^{\lambda(t\wedge S_k)}f\bigl(y(t\wedge S_k)\bigr)\bigr]\le f(x)+\frac{c}{\lambda}\bigl(e^{\lambda t}-1\bigr),
\]
and multiplying by $e^{-\lambda t}$ we arrive at
\[
E^xf(y(t))\le f(x)e^{-\lambda t}+\frac{c}{\lambda}\bigl(1-e^{-\lambda t}\bigr).
\]
In particular $f(y(t))$ is $P^x$-integrable. Since $x\in\R^d$ and $t>0$ were arbitrary, the proof is complete.

\rmk{The process $L$ of Step 2 is the process $\Lambda$ of \bk{Exercise 4.33} with $\alpha=-\lambda$ (and $g$ in place of $f$); we have verified its martingale property directly because the filtration $\{\B_s\}$ does not satisfy the usual conditions assumed there. The nonnegativity of $f$ enters only through Fatou's lemma in Step 4, and the bound cannot be improved: for $d=1$, $b(x)=-x$, $\sigma\equiv1$ (the Ornstein--Uhlenbeck process) and $f(x)=x^2$ we have $\cA f+2f\equiv1$, and indeed $E^x y(t)^2=x^2e^{-2t}+\tfrac12(1-e^{-2t})$.}
\end{kex}

\solsection{5.5}{A Study of the One-Dimensional Case}

\begin{kex}{Problem}{5.12}{}
Throughout, \bk{(ND)} and \bk{(LI)} are in force, and integrals $\int_u^v$ with $v<u$ are understood as $-\int_v^u$.

\emph{Step 1: local integrability.} We set $f\defeq 2b/\sigma^2$, a Borel-measurable, real-valued function on $\R$ by \bk{(ND)}. By \bk{(LI)}, every $x\in\R$ has an open neighbourhood $(x-\eps_x,x+\eps_x)$ on which $\abs{f}$ is integrable. A compact interval $[u,v]$ is covered by finitely many of these neighbourhoods, so $\int_u^v\abs{f(y)}\,dy<\infty$: the function $f$ is locally integrable. Consequently, for $u,v\in\R$,
\[
  F(u,v)\defeq\int_u^v f(\zeta)\,d\zeta
\]
is a well-defined real number, it satisfies
\begin{equation}\label{eq:5-5-12:add}
  F(u,w)=F(u,v)+F(v,w);\qquad u,v,w\in\R
\end{equation}
(additivity of the Lebesgue integral over adjacent intervals, valid for every ordering of $u,v,w$ thanks to the sign convention), and for fixed $u$ the function $v\mapsto F(u,v)$ is continuous on $\R$ (dominated convergence, $f$ being integrable on compact intervals).

\emph{Step 2: the derivative of $p_c$.} For $c\in\R$ let $g_c(\xi)\defeq\exp\{-F(c,\xi)\}$, $\xi\in\R$. By Step 1 this is a continuous, strictly positive, finite function, and \bk{(5.42)} reads $p_c(x)=\int_c^x g_c(\xi)\,d\xi$. The integrand being continuous, the fundamental theorem of calculus shows that $p_c$ is continuously differentiable with
\begin{equation}\label{eq:5-5-12:der}
  p_c'(x)=g_c(x)=\exp\{-F(c,x)\}\in(0,\infty);\qquad x\in\R .
\end{equation}

\emph{Step 3: proof of \bk{(5.45)}.} Let $a,c\in\R$. By \eqref{eq:5-5-12:add} with $(u,v,w)=(a,c,\xi)$ and by \eqref{eq:5-5-12:der},
\begin{equation}\label{eq:5-5-12:mult}
  g_a(\xi)=\exp\{-F(a,c)\}\exp\{-F(c,\xi)\}=p_a'(c)\,g_c(\xi);\qquad \xi\in\R .
\end{equation}
Hence, for every $x\in\R$, splitting the integral of the continuous function $g_a$ at the point $c$ and then using \eqref{eq:5-5-12:mult},
\[
  p_a(x)=\int_a^x g_a(\xi)\,d\xi
  =\int_a^c g_a(\xi)\,d\xi+\int_c^x g_a(\xi)\,d\xi
  =p_a(c)+p_a'(c)\int_c^x g_c(\xi)\,d\xi
  =p_a(c)+p_a'(c)\,p_c(x),
\]
which is \bk{(5.45)}.

\emph{Step 4: the values at $\pm\infty$.} By \eqref{eq:5-5-12:der} each $p_c$ is strictly increasing, so the limits
$p_c(\infty)=\lim_{x\to\infty}p_c(x)\in(0,\infty]$ and $p_c(-\infty)=\lim_{x\to-\infty}p_c(x)\in[-\infty,0)$ exist; these are the values of the continuous extension of $p_c$ to $[-\infty,\infty]$ introduced after \bk{(5.44)}. In \bk{(5.45)} the number $p_a(c)$ is finite and $p_a'(c)$ is finite and strictly positive. Letting $x\to\pm\infty$ in \bk{(5.45)} we obtain, with the usual arithmetic in $[-\infty,\infty]$,
\[
  p_a(\pm\infty)=p_a(c)+p_a'(c)\,p_c(\pm\infty).
\]
Therefore $p_a(\infty)<\infty$ if and only if $p_c(\infty)<\infty$, and $p_a(-\infty)>-\infty$ if and only if $p_c(-\infty)>-\infty$. Since $a,c\in\R$ were arbitrary, the finiteness or nonfiniteness of $p_c(\infty)$, and likewise of $p_c(-\infty)$, does not depend on $c$.
\end{kex}

\begin{kex}{Exercise}{5.14}{}
We take the Borel-measurable coefficients
\[
  b(x)\defeq \ind_{\{0\}}(x),\qquad \sigma(x)\defeq \ind_{\R\setminus\{0\}}(x);\qquad x\in\R .
\]

\emph{Step 1: \bk{(ND)} fails, \bk{(LI)} holds, and $p$, $q$, $\tilde\sigma$ are well defined.}
Since $\sigma^2(0)=0$, condition \bk{(ND)} fails, but only at the single point $x=0$. For every $y\neq 0$ we have $\abs{b(y)}/\sigma^2(y)=0$; thus the function $\abs{b}/\sigma^2$ is defined and equal to zero off the Lebesgue-null set $\{0\}$, and
\[
  \int_{x-\eps}^{x+\eps}\frac{\abs{b(y)}\,dy}{\sigma^2(y)}=0<\infty\qquad\text{for all }x\in\R,\ \eps>0,
\]
whatever value in $[0,\infty]$ is assigned to the integrand at $y=0$. Hence \bk{(LI)} holds. For the same reason the inner integral in \bk{(5.42)} vanishes for all $c,\xi\in\R$, and the scale function is
\[
  p(x)=\int_c^x 1\,d\xi=x-c;\qquad x\in\R .
\]
It has the continuous, strictly positive derivative $p'\equiv 1$, it maps $\R$ onto $\bigl(p(-\infty),p(\infty)\bigr)=\R$, and its inverse is $q(y)=y+c$. The function of \bk{(5.48)} is therefore
\[
  \tilde\sigma(y)=p'(q(y))\,\sigma(q(y))=\ind_{\R\setminus\{-c\}}(y);\qquad y\in\R ,
\]
and \bk{(5.46)} is the equation
\begin{equation}\label{eq:5-5-14:Y}
  Y_t=Y_0+\int_0^t \ind_{\{Y_s\neq -c\}}\,dW_s;\qquad 0\le t<\infty .
\end{equation}

\emph{Step 2: a solution of \bk{(5.46)}, \bk{(5.47)}.}
Let $W=\{W_t,\F_t;\,0\le t<\infty\}$ be a standard Brownian motion on some probability space $(\Omega,\F,P)$, with a filtration $\{\F_t\}$ satisfying the usual conditions, and set
\[
  Y_t\defeq -c;\qquad 0\le t<\infty .
\]
This process is continuous and adapted, and $Y_0=-c$ satisfies \bk{(5.47)} because $p(-\infty)=-\infty<-c<\infty=p(\infty)$. The integrand $\tilde\sigma(Y_s)=\ind_{\{Y_s\neq -c\}}$ is identically zero, so the stochastic integral in \eqref{eq:5-5-14:Y} is defined and equal to zero, and \eqref{eq:5-5-14:Y} reads $-c=-c$. Thus $Y$ is a (strong, hence also weak) solution of \bk{(5.46)}, \bk{(5.47)}, and it does not explode.

\emph{Step 3: $Y$ is not of the form $p(X)$ for a solution $X$ of \bk{(5.1)}.}
Suppose that, on some filtered probability space carrying a Brownian motion $W$, there were a solution $X$ of
\begin{equation}\label{eq:5-5-14:X}
  X_t=X_0+\int_0^t \ind_{\{X_s=0\}}\,ds+\int_0^t \ind_{\{X_s\neq 0\}}\,dW_s ,
\end{equation}
possibly only up to an explosion time $S$, with $p(X_t)=-c$ for all $0\le t<\infty$, almost surely. Because $p$ is injective on $[-\infty,\infty]$ with $p(0)=-c$, this forces $X_t=0$ for all $0\le t<\infty$, almost surely. Then $X$ does not explode: with $S_n=\inf\{t\ge 0;\ \abs{X_t}\ge n\}$ as in \bk{(5.2)}, we have $S_n=\infty$ a.s.\ for every $n\ge 1$, so that \eqref{eq:5-5-14:X} would have to hold for all $0\le t<\infty$. But for $X\equiv 0$ the stochastic integral in \eqref{eq:5-5-14:X} has integrand zero and vanishes, the Lebesgue integral equals $t$, and \eqref{eq:5-5-14:X} becomes $0=t$, which is false for every $t>0$. Hence no solution $X$ of \bk{(5.1)} satisfies $Y=p(X)$, although $Y$ solves \bk{(5.46)}, \bk{(5.47)}.

\rmk{The example shows which step of the proof of \bk{Proposition 5.13} needs \bk{(ND)}. In the converse part one applies the generalized It\^o rule to $X=q(Y)$ and uses \bk{(5.44)}, \bk{(5.48)} to identify $\tfrac12 q''(Y_s)\,\tilde\sigma^2(Y_s)$ with $b(X_s)$. Here $q''\equiv 0$, so that this term is zero, whereas $b(X_s)=1$ on $\{X_s=0\}$: the identification fails exactly on the set $\{s;\ \sigma(X_s)=0\}$, which for the solution $Y\equiv -c$ is the whole half-line. The values of $b$ on the zero set of $\sigma$ are invisible to \bk{(LI)}, to $p$ and to $\tilde\sigma$ when this set is Lebesgue-null, but not to the equation \bk{(5.1)}.}
\end{kex}

\begin{kex}{Problem}{5.28}{}
Throughout, $I=(\ell,r)$ with $-\infty\le\ell<r\le\infty$, the Borel-measurable functions $\sigma,b\colon I\to\R$ satisfy \bk{(ND)$'$} and \bk{(LI)$'$}, and an integral $\int_u^w$ with $w<u$ is understood as $-\int_w^u$.

\emph{Step 1: local integrability.} By \bk{(LI)$'$} every $x\in I$ has a neighbourhood $(x-\eps_x,x+\eps_x)$, which we may take inside $I$, on which $1/\sigma^2$ and $\abs{b}/\sigma^2$ are integrable. A compact interval $[u,w]\subset I$ is covered by finitely many such neighbourhoods; hence $1/\sigma^2$ and $f\defeq 2b/\sigma^2$ are integrable on every compact subinterval of $I$. Consequently
\[
  F(u,w)\defeq\int_u^w f(\zeta)\,d\zeta;\qquad u,w\in I,
\]
is a real number, $F(u,w)=F(u,t)+F(t,w)$ for all $u,t,w\in I$ (in any order, thanks to the sign convention), and $w\mapsto F(u,w)$ is continuous on $I$ by dominated convergence.

\emph{Step 2: the scale functions.} For $c\in I$, \bk{(5.42)} reads $p_c(x)=\int_c^x\exp\{-F(c,\xi)\}\,d\xi$, $x\in I$. The integrand is continuous, so $p_c$ is continuously differentiable on $I$ with
\begin{equation}\label{eq:5-5-28:pder}
  p_c'(x)=\exp\{-F(c,x)\}\in(0,\infty),\qquad p_c(c)=0,\quad p_c'(c)=1 .
\end{equation}
For $a,c\in I$ the additivity $F(a,x)=F(a,c)+F(c,x)$ gives
\begin{equation}\label{eq:5-5-28:mult}
  p_a'(x)=p_a'(c)\,p_c'(x);\qquad x\in I,
\end{equation}
and integrating \eqref{eq:5-5-28:mult} from $c$ to $x$ we recover \bk{(5.45)} on $I$:
\begin{equation}\label{eq:5-5-28:p}
  p_a(x)=p_a(c)+p_a'(c)\,p_c(x);\qquad x\in I .
\end{equation}

\emph{Step 3: the functions $v_c$.} Fix $c\in I$ and set
\[
  h_c(z)\defeq\frac{2}{p_c'(z)\sigma^2(z)},\qquad
  H_c(u,y)\defeq\int_u^y h_c(z)\,dz;\qquad z,u,y\in I .
\]
The function $1/p_c'$ is continuous and positive on $I$, hence bounded on compact subintervals, so by Step 1 the function $h_c$ is integrable on compact subintervals of $I$; it is strictly positive everywhere, because $0<\sigma^2<\infty$ on $I$. Therefore $H_c$ is real-valued and additive, $H_c(u,y)=H_c(u,t)+H_c(t,y)$, the function $y\mapsto H_c(u,y)$ is continuous and strictly increasing on $I$, and $H_c(u,u)=0$. In this notation \bk{(5.75)} is
\[
  v_c(x)=\int_c^x p_c'(y)\,H_c(c,y)\,dy;\qquad x\in I,
\]
the integral of a continuous function. Thus $v_c$ is well defined, real-valued and continuously differentiable on $I$, with
\begin{equation}\label{eq:5-5-28:vder}
  v_c(c)=0,\qquad v_c'(x)=p_c'(x)\,H_c(c,x);\qquad x\in I .
\end{equation}
Since $p_c'>0$ and $H_c(c,x)$ has the sign of $x-c$, the function $v_c$ is nondecreasing on $[c,r)$ and nonincreasing on $(\ell,c]$; in particular $v_c\ge 0$, and the limits $v_c(r-)\defeq\lim_{x\uparrow r}v_c(x)$ and $v_c(\ell+)\defeq\lim_{x\downarrow\ell}v_c(x)$ exist in $[0,\infty]$.

\emph{Step 4: proof of \bk{(5.76)}.} Let $a,c\in I$. By \eqref{eq:5-5-28:mult}, for $y,z\in I$,
\[
  \frac{p_a'(y)}{p_a'(z)}=\frac{p_c'(y)}{p_c'(z)},
  \qquad\text{that is,}\qquad p_a'(y)\,h_a(z)=p_c'(y)\,h_c(z).
\]
Integrating in $z$ from $a$ to $y$ and using \eqref{eq:5-5-28:vder} with $a$ in place of $c$, we obtain for every $y\in I$
\begin{equation}\label{eq:5-5-28:vader}
  v_a'(y)=p_a'(y)\,H_a(a,y)=p_c'(y)\,H_c(a,y)=p_c'(y)\,H_c(a,c)+p_c'(y)\,H_c(c,y).
\end{equation}
At $y=c$ this gives, because $p_c'(c)=1$ and $H_c(c,c)=0$,
\begin{equation}\label{eq:5-5-28:vac}
  v_a'(c)=H_c(a,c).
\end{equation}
Now $v_a$ is continuously differentiable, so by \eqref{eq:5-5-28:vader}, \eqref{eq:5-5-28:vac}, \eqref{eq:5-5-28:pder} and \eqref{eq:5-5-28:vder}, for every $x\in I$,
\[
  v_a(x)-v_a(c)=\int_c^x v_a'(y)\,dy
  =v_a'(c)\int_c^x p_c'(y)\,dy+\int_c^x p_c'(y)\,H_c(c,y)\,dy
  =v_a'(c)\,p_c(x)+v_c(x),
\]
which is \bk{(5.76)}.

\emph{Step 5: the endpoints.} We treat the right endpoint. Each $p_c$ is strictly increasing, so $p_c(r-)\defeq\lim_{x\uparrow r}p_c(x)\in(0,\infty]$ exists. We first note that, for every $c\in I$,
\begin{equation}\label{eq:5-5-28:impl}
  p_c(r-)=\infty\ \Longrightarrow\ v_c(r-)=\infty
\end{equation}
(this is \bk{(5.73)} of \bk{Problem 5.27}; the argument is one line). Indeed, fix $c'\in(c,r)$; then $H_c(c,y)\ge H_c(c,c')>0$ for $y\ge c'$, and since the integrand in \eqref{eq:5-5-28:vder} is nonnegative on $[c,r)$,
\[
  v_c(x)\ge\int_{c'}^x p_c'(y)\,H_c(c,y)\,dy\ge H_c(c,c')\,\bigl(p_c(x)-p_c(c')\bigr);\qquad c'\le x<r,
\]
and letting $x\uparrow r$ yields \eqref{eq:5-5-28:impl}.

Let now $a,c\in I$. Since $p_a(c)\in\R$ and $p_a'(c)\in(0,\infty)$, \eqref{eq:5-5-28:p} shows that $p_a(r-)=\infty$ if and only if $p_c(r-)=\infty$.

\emph{Case 1: $p_c(r-)=\infty$.} Then also $p_a(r-)=\infty$, and \eqref{eq:5-5-28:impl}, applied to $c$ and to $a$, gives $v_c(r-)=v_a(r-)=\infty$.

\emph{Case 2: $p_c(r-)<\infty$.} The numbers $v_a(c)$, $v_a'(c)$ and $p_c(r-)$ are finite, so we may let $x\uparrow r$ in \bk{(5.76)} to obtain
\[
  v_a(r-)=v_a(c)+v_a'(c)\,p_c(r-)+v_c(r-)
\]
in $(-\infty,\infty]$; hence $v_a(r-)<\infty$ if and only if $v_c(r-)<\infty$.

In both cases $v_a(r-)$ and $v_c(r-)$ are simultaneously finite or infinite. The left endpoint is handled in the same way. The limit $p_c(\ell+)\defeq\lim_{x\downarrow\ell}p_c(x)\in[-\infty,0)$ exists, and \eqref{eq:5-5-28:p} shows that $p_a(\ell+)=-\infty$ if and only if $p_c(\ell+)=-\infty$. For $c'\in(\ell,c)$ and $\ell<x\le c'$ we have $-H_c(c,y)=H_c(y,c)\ge H_c(c',c)>0$ for $y\le c'$, whence, the integrand below being nonnegative on $(\ell,c]$,
\[
  v_c(x)=\int_x^c p_c'(y)\,H_c(y,c)\,dy\ge H_c(c',c)\,\bigl(p_c(c')-p_c(x)\bigr),
\]
so that $p_c(\ell+)=-\infty$ implies $v_c(\ell+)=\infty$ (this is \bk{(5.74)}). If $p_c(\ell+)=-\infty$, then $v_c(\ell+)=v_a(\ell+)=\infty$; if $p_c(\ell+)>-\infty$, we let $x\downarrow\ell$ in \bk{(5.76)}, all of whose terms other than $v_a(x)$ and $v_c(x)$ have finite limits, and get $v_a(\ell+)=v_a(c)+v_a'(c)\,p_c(\ell+)+v_c(\ell+)$. Therefore $v_a(\ell+)<\infty$ if and only if $v_c(\ell+)<\infty$. Since $a,c\in I$ were arbitrary, the finiteness or nonfiniteness of $v_c(r-)$ and of $v_c(\ell+)$ does not depend on $c$.

\rmk{The case distinction in Step~5 is needed: by \eqref{eq:5-5-28:vac}, $v_a'(c)=H_c(a,c)$ is negative when $a>c$, so if $p_c(r-)=\infty$ the right-hand side of \bk{(5.76)} is of the form $-\infty+\infty$ as $x\uparrow r$, and one cannot pass to the limit termwise.}
\end{kex}

\begin{kex}{Exercise}{5.31}{}
\emph{Setting.} As in \bk{Definition 5.1}, $(\Omega,\F,P)$ is a probability space with a filtration $\{\F_t\}$ satisfying the usual conditions, $W$ is a standard Brownian motion relative to $\{\F_t\}$, and $X$ is a continuous, adapted, real-valued process for which \bk{(5.79)} holds for all $0\le t<\infty$, almost surely; the stochastic integral is defined because $X$ is continuous and adapted, hence progressively measurable with $\int_0^t X_s^2\,ds<\infty$. The constants $\mu,\nu$ are real, and we assume $\nu\neq0$ (see the Remark at the end). Equation \bk{(5.79)} is \bk{(5.1)} with
\[
  b(y)=\mu y,\qquad \sigma(y)=\nu y ,
\]
and we take $I=(\ell,r)=(0,\infty)$. On $I$ we have $\sigma^2(y)=\nu^2y^2>0$, which is \bk{(ND)$'$}, and $(1+\abs{b(y)})/\sigma^2(y)=(1+\abs{\mu}y)/(\nu^2y^2)$ is continuous on $(0,\infty)$, hence integrable over every compact subinterval, which is \bk{(LI)$'$}.

\emph{Step 1: a weak solution in $I$.} Let $T_0\defeq\inf\{t\ge0;\ X_t=0\}$, a stopping time by \bk{Problem 1.2.7}, and set $\wt X_t\defeq X_{t\wedge T_0}$. Since $X$ is continuous with $X_0=x>0$, we have $X_t>0$ for $t<T_0$; thus $\wt X$ is a continuous, adapted, $[0,\infty)$-valued process with $\wt X_0=x\in I$, and $\wt X_t=0$ for $t\ge T_0$. Let $\ell_n\downarrow0$ and $r_n\uparrow\infty$ be strictly monotone with $0<\ell_n<r_n<\infty$, and let $S_n\defeq\inf\{t\ge0;\ \wt X_t\notin(\ell_n,r_n)\}$ as in \bk{(5.56)}. Because $\wt X_{T_0}=0$ on $\{T_0<\infty\}$, we have $S_n\le T_0$, so that
\begin{equation}\label{eq:5-5-31:agree}
  \wt X_s=X_s\quad\text{for } 0\le s\le S_n,\ s<\infty,
  \qquad\text{and}\qquad \wt X_{t\wedge S_n}=X_{t\wedge S_n};\quad 0\le t<\infty .
\end{equation}
Condition \bk{(5.3)} holds because $s\mapsto\abs{\mu \wt X_s}+\nu^2\wt X_s^2$ is continuous, hence bounded on $[0,t]$. For \bk{(5.4)} we evaluate \bk{(5.79)} at $t\wedge S_n$ and use the sample-path property \bk{(3.2.24)} of the stochastic integral (valid for integrands in $\cP^*$ by the discussion preceding \bk{Proposition 3.2.24}) together with \eqref{eq:5-5-31:agree}:
\begin{align*}
  \wt X_{t\wedge S_n}
  &= x+\mu\int_0^{t}X_s\ind_{\{s\le S_n\}}\,ds+\nu\int_0^{t}X_s\ind_{\{s\le S_n\}}\,dW_s\\
  &= x+\int_0^{t}b(\wt X_s)\ind_{\{s\le S_n\}}\,ds+\int_0^{t}\sigma(\wt X_s)\ind_{\{s\le S_n\}}\,dW_s ,
\end{align*}
almost surely for each $t$, and then for all $t$ simultaneously since both sides are continuous in $t$. Hence $(\wt X,W)$ is a weak solution in $I=(0,\infty)$ of \bk{(5.1)} in the sense of \bk{Definition 5.20}, with nonrandom initial condition $x\in I$. Its exit time \bk{(5.57)} is
\begin{equation}\label{eq:5-5-31:S}
  S=\inf\{t\ge0;\ \wt X_t\notin(0,\infty)\}=T_0 ,
\end{equation}
because $\wt X$ is real-valued and first vanishes at $T_0$.

\emph{Step 2: the scale function.} Put $\gamma\defeq2\mu/\nu^2$ and choose $c=1$ in \bk{(5.42)}. Since $-2\int_1^\xi b(\zeta)\sigma^{-2}(\zeta)\,d\zeta=-\gamma\log\xi$, we obtain $p'(\xi)=\xi^{-\gamma}$ and
\begin{equation}\label{eq:5-5-31:p}
  p(y)=\int_1^y\xi^{-\gamma}\,d\xi=
  \begin{cases}
    \dfrac{y^{1-\gamma}-1}{1-\gamma}, & \gamma\neq1,\\[1.5ex]
    \log y, & \gamma=1,
  \end{cases}
  \qquad 0<y<\infty .
\end{equation}
Consequently
\begin{equation}\label{eq:5-5-31:plim}
  \begin{aligned}
    \gamma<1:&\quad p(0+)=-\tfrac{1}{1-\gamma}>-\infty,\quad p(\infty-)=\infty;\\
    \gamma>1:&\quad p(0+)=-\infty,\quad p(\infty-)=\tfrac{1}{\gamma-1}<\infty;\\
    \gamma=1:&\quad p(0+)=-\infty,\quad p(\infty-)=\infty .
  \end{aligned}
\end{equation}

\emph{Step 3: Feller's test.} By \bk{(5.65)} with $c=1$,
\[
  v(y)=\int_1^y p'(\eta)\int_1^\eta\frac{2\,dz}{p'(z)\sigma^2(z)}\,d\eta
      =\frac{2}{\nu^2}\int_1^y\eta^{-\gamma}\int_1^\eta z^{\gamma-2}\,dz\,d\eta ;\qquad 0<y<\infty .
\]
If $\gamma=1$, the inner integral is $\log\eta$ and $v(y)=\nu^{-2}(\log y)^2$, so $v(0+)=v(\infty-)=\infty$. If $\gamma\neq1$, the inner integral is $(\eta^{\gamma-1}-1)/(\gamma-1)$, whence
\begin{equation}\label{eq:5-5-31:v}
  v(y)=\frac{2}{\nu^2(\gamma-1)}\int_1^y\bigl(\eta^{-1}-\eta^{-\gamma}\bigr)\,d\eta
      =\frac{2}{\nu^2(\gamma-1)}\bigl(\log y-p(y)\bigr) .
\end{equation}
For $\gamma<1$ this reads $v(y)=\frac{2}{\nu^2(1-\gamma)}\bigl(p(y)-\log y\bigr)$: as $y\downarrow0$, $p(y)$ stays bounded and $-\log y\to\infty$; as $y\to\infty$, $p(y)-\log y=\bigl(y^{1-\gamma}-1\bigr)/(1-\gamma)-\log y\to\infty$, a positive power dominating the logarithm. For $\gamma>1$: as $y\to\infty$, $p(y)$ stays bounded and $\log y\to\infty$; as $y\downarrow0$, $\log y-p(y)=\bigl(y^{-(\gamma-1)}-1\bigr)/(\gamma-1)-\log(1/y)\to\infty$, for the same reason. In all three cases
\[
  v(0+)=v(\infty-)=\infty ,
\]
which is \bk{(5.77)}. \bk{Theorem 5.29}, applied to the weak solution $\wt X$ in $I$ of Step~1, gives $P[S=\infty]=1$; by \eqref{eq:5-5-31:S}, $P[T_0=\infty]=1$. Therefore, almost surely, $X_t\in(0,\infty)$ for all $0\le t<\infty$, and $\wt X=X$.

\emph{Step 4: the three cases.} We apply \bk{Proposition 5.22} to $\wt X$, using \eqref{eq:5-5-31:plim}, and recall that almost surely $S=\infty$ and $\wt X=X$, so that $\lim_{t\uparrow S}$ is $\lim_{t\to\infty}$ and the suprema and infima over $[0,S)$ are over $[0,\infty)$.

\textbf{(i)} If $\mu<\nu^2/2$, then $\gamma<1$, $p(\ell+)>-\infty$ and $p(r-)=\infty$: this is case (b) of \bk{Proposition 5.22}, which gives $\lim_{t\to\infty}X_t=\ell=0$ and $\sup_{0\le t<\infty}X_t<r=\infty$, a.s.

\textbf{(ii)} If $\mu>\nu^2/2$, then $\gamma>1$, $p(\ell+)=-\infty$ and $p(r-)<\infty$: this is case (c), which gives $\lim_{t\to\infty}X_t=r=\infty$ and $\inf_{0\le t<\infty}X_t>\ell=0$, a.s.

\textbf{(iii)} If $\mu=\nu^2/2$, then $\gamma=1$, $p(\ell+)=-\infty$ and $p(r-)=\infty$: this is case (a), which gives $\sup_{0\le t<\infty}X_t=r=\infty$ and $\inf_{0\le t<\infty}X_t=\ell=0$, a.s.

\rmk{The hypothesis $\nu\neq0$ is needed for \bk{(ND)$'$} and is implicit in the statement. For $\nu=0$, \bk{(5.79)} is the linear equation $X_t=x+\mu\int_0^tX_s\,ds$ with the unique solution $X_t=xe^{\mu t}$; then $X_t\in(0,\infty)$ and (i), (ii) remain true, but (iii) fails, since $\mu=\nu^2/2=0$ gives $X\equiv x$. The conclusions agree with the explicit formula $X_t=x\exp\{(\mu-\frac12\nu^2)t+\nu W_t\}$ quoted in the statement: the exponent is a Brownian motion with drift $\mu-\frac12\nu^2$ and nonzero dispersion $\nu$, which tends to $-\infty$, tends to $+\infty$, or oscillates between $-\infty$ and $+\infty$ according as the drift is negative, positive or zero.}
\end{kex}

\begin{kex}{Exercise}{5.34}{}
Throughout, $(X,W)$, $(\Omega,\F,P)$, $\{\F_t\}$ is a weak solution in $I=(\ell,r)=(-\infty,\infty)$ of \bk{(5.1)}, in the sense of \bk{Definition 5.20}, with nonrandom initial condition $X_0=x\in\R$, and $S$ is its explosion time \bk{(5.57)}; this is the setting of \bk{Theorem 5.29} and \bk{Proposition 5.32}. In both parts $\sigma^2\equiv 1>0$ and $b$ is continuous, hence locally bounded, so \bk{(ND)$'$} and \bk{(LI)$'$} hold and \bk{Proposition 5.32} applies. We take $c=0$ in \bk{(5.42)} and \bk{(5.65)}; by \bk{Problems 5.12} and \bk{5.28} this choice does not affect the finiteness of $p(\pm\infty)$, $v(\pm\infty)$.

Both parts rest on one elementary estimate. Let $k\in\{3,4\}$ and $y>0$. For $0\le z\le y$ we have
\[
  y^k-z^k=(y-z)\sum_{j=0}^{k-1}y^{k-1-j}z^{j}\;\ge\;(y-z)\,y^{k-1}\ge 0 ,
\]
and therefore, bounding the integrand first by $1$ and then by $e^{-y^{k-1}(y-z)}$,
\begin{equation}\label{eq:5-5-34:a}
  0\le e^{-y^k}\int_0^y e^{z^k}\,dz=\int_0^y e^{-(y^k-z^k)}\,dz
  \le \min\Bigl\{y,\ \int_0^y e^{-y^{k-1}(y-z)}\,dz\Bigr\}
  \le \min\bigl\{y,\ y^{-(k-1)}\bigr\}.
\end{equation}

\textbf{(a)} Here $-2\int_0^\xi b(\zeta)\,d\zeta=-\xi^3$, so \bk{(5.42)} and \bk{(5.65)} give
\[
  p'(x)=e^{-x^3},\qquad p(x)=\int_0^x e^{-\xi^3}\,d\xi,\qquad
  v(x)=\int_0^x e^{-y^3}\int_0^y 2e^{z^3}\,dz\,dy;\qquad x\in\R .
\]
For $x<0$ we have $p(x)=-\int_x^0 e^{-\xi^3}\,d\xi=-\int_0^{\abs{x}}e^{\eta^3}\,d\eta\le -\abs{x}$, whence
\[
  p(\ell+)=p(-\infty)=-\infty .
\]
For $x>0$ the integrand defining $v$ is nonnegative, so $v$ is nondecreasing on $[0,\infty)$, and \eqref{eq:5-5-34:a} with $k=3$ yields
\[
  v(r-)=v(\infty)=\lim_{x\to\infty}v(x)\le 2\int_0^1 y\,dy+2\int_1^\infty y^{-2}\,dy=3<\infty .
\]
Thus condition (ii) of \bk{Proposition 5.32} is satisfied, and $P[S<\infty]=1$.

\rmk{Since $p(-\infty)=-\infty$, \bk{(5.74)} gives $v(-\infty)=\infty$, so condition (i) of \bk{Proposition 5.32} fails in part (a). The proof of sufficiency in \bk{Proposition 5.32} shows moreover that, almost surely, $S=T_r$: the solution explodes to $+\infty$, as the sign of the drift suggests.}

\textbf{(b)} Now $-2\int_0^\xi b(\zeta)\,d\zeta=-\xi^4$, so
\[
  p'(x)=e^{-x^4},\qquad
  v(x)=\int_0^x e^{-y^4}\int_0^y 2e^{z^4}\,dz\,dy;\qquad x\in\R .
\]
The inner function $g(y)\defeq e^{-y^4}\int_0^y 2e^{z^4}\,dz$ is odd, because the substitution $z\mapsto -z$ gives $\int_0^{-y}e^{z^4}\,dz=-\int_0^{y}e^{z^4}\,dz$; consequently $v(x)=\int_0^x g(y)\,dy$ is even. On $[0,\infty)$ we have $g\ge 0$, so $v$ is nondecreasing there, and \eqref{eq:5-5-34:a} with $k=4$ yields
\[
  v(r-)=v(\infty)\le 2\int_0^1 y\,dy+2\int_1^\infty y^{-3}\,dy=2<\infty,
  \qquad v(\ell+)=v(-\infty)=v(\infty)<\infty .
\]
Hence condition (i) of \bk{Proposition 5.32} holds, and the last assertion of that proposition (justified in \bk{Remark 5.33}) gives $ES<\infty$.
\end{kex}

\begin{kex}{Exercise}{5.35}{}
Here $\ell$ denotes the constant in the dispersion coefficient, and the endpoints of $I$ are $0$ and $\infty$. On $I$ we have $\sigma^2(x)=\ell^2x>0$, and $(1+\abs{b(y)})/\sigma^2(y)=(1+k)/(\ell^2y)$ is bounded on every compact subinterval of $I$; thus \bk{(ND)$'$} and \bk{(LI)$'$} hold. Throughout, $(X,W)$, $(\Omega,\F,P)$, $\{\F_t\}$ is a weak solution in $I$ (\bk{Definition 5.20}) of
\begin{equation}\label{eq:5-5-35:sde}
  dX_t=k\,dt+\ell\sqrt{X_t}\,dW_t,\qquad X_0=x\in I,
\end{equation}
with exit time $S=\lim_{n\to\infty}S_n$ as in \bk{(5.56)}, \bk{(5.57)}.

\textbf{(i)} We take $c=1$ in \bk{(5.42)}; by \bk{Problem 5.12} the finiteness of $p(0+)$ and of $p(\infty-)$ does not depend on this choice. Since $2b(\zeta)/\sigma^2(\zeta)=\mu/\zeta$,
\[
  p'(\xi)=\exp\Bigl\{-\mu\int_1^\xi\frac{d\zeta}{\zeta}\Bigr\}=\xi^{-\mu},
  \qquad
  p(x)=\int_1^x\xi^{-\mu}\,d\xi=
  \begin{cases}
    \log x; & \mu=1,\\[0.5ex]
    \dfrac{x^{1-\mu}-1}{1-\mu}; & \mu\neq 1,
  \end{cases}
  \qquad 0<x<\infty.
\]
If $\mu=1$, then $p(0+)=-\infty$ and $p(\infty-)=\infty$: this is case (a). If $\mu<1$, then $x^{1-\mu}\to 0$ as $x\downarrow 0$ and $x^{1-\mu}\to\infty$ as $x\to\infty$, so $p(0+)=-1/(1-\mu)>-\infty$ and $p(\infty-)=\infty$: this is case (b). If $\mu>1$, then $x^{1-\mu}\to\infty$ as $x\downarrow 0$ and $x^{1-\mu}\to0$ as $x\to\infty$, and $1-\mu<0$, so $p(0+)=-\infty$ and $p(\infty-)=1/(\mu-1)<\infty$: this is case (c).

\textbf{(ii)} The speed measure \bk{(5.51)} is $m(dz)=2\,dz/(p'(z)\sigma^2(z))=(2/\ell^2)z^{\mu-1}\,dz$, so the function of \bk{(5.65)} is, because $2/(\ell^2\mu)=1/k$,
\begin{equation}\label{eq:5-5-35:v}
  v(x)=\int_1^x y^{-\mu}\int_1^y\frac{2z^{\mu-1}}{\ell^2}\,dz\,dy
      =\frac1k\int_1^x\bigl(1-y^{-\mu}\bigr)\,dy
      =\frac{x-1-p(x)}{k};\qquad 0<x<\infty .
\end{equation}

\emph{Case $\mu=1$.} We are in case (a), and \bk{Proposition 5.22}(a) gives $P[S=\infty]=1$ directly: the solution does not explode. (It is, moreover, recurrent in $I$, with $\inf_{t}X_t=0$ and $\sup_tX_t=\infty$ a.s.)

\emph{Case $\mu>1$.} By (i), $p(0+)=-\infty$ and $p(\infty-)=1/(\mu-1)$, so \eqref{eq:5-5-35:v} gives
\[
  v(0+)=\frac{-1-p(0+)}{k}=\infty,\qquad
  v(\infty-)=\lim_{x\to\infty}\frac1k\Bigl(x-1-p(x)\Bigr)=\infty .
\]
Thus \bk{(5.77)} holds and \bk{Theorem 5.29} yields $P[S=\infty]=1$: the solution does not explode. (\bk{Proposition 5.22}(c) adds that $\lim_{t\to\infty}X_t=\infty$ and $\inf_{t}X_t>0$, a.s.)

\emph{Case $\mu<1$.} By (i), $p(0+)=-1/(1-\mu)$ and $p(\infty-)=\infty$, so \eqref{eq:5-5-35:v} gives
\[
  v(0+)=\frac1k\Bigl(-1+\frac1{1-\mu}\Bigr)=\frac{\mu}{k(1-\mu)}=\frac{2}{\ell^2(1-\mu)}<\infty .
\]
Condition (iii) of \bk{Proposition 5.32} is satisfied (its left endpoint is $0$, its right endpoint $\infty$), whence $P[S<\infty]=1$. On the other hand we are in case (b) of \bk{Proposition 5.22}, so that $P[\lim_{t\uparrow S}X_t=0]=P[\sup_{0\le t<S}X_t<\infty]=1$. Since $X$ is continuous with values in $[0,\infty]$, we conclude that
\[
  S<\infty\quad\text{and}\quad X_S=\lim_{t\uparrow S}X_t=0,\qquad\text{a.s.};
\]
that is, the solution leaves $I$ in finite time, and it does so through the origin, never through $\infty$.

\smallskip
For (iii) and (iv) we need one computation, valid for all $k,\ell>0$.

\begin{lem}
For every $n\ge1$ we have, almost surely,
\begin{equation}\label{eq:5-5-35:sqrt}
  \sqrt{X_{t\wedge S_n}}=\sqrt{x}+\frac{4k-\ell^2}{8}\int_0^{t\wedge S_n}\frac{ds}{\sqrt{X_s}}+\frac{\ell}{2}\,W_{t\wedge S_n};
  \qquad 0\le t<\infty .
\end{equation}
\end{lem}
\begin{proof}
Let $0<\ell_n<r_n<\infty$ be the numbers of \bk{Definition 5.20}(iii)$'$ and set $\alpha\defeq\ell_n\wedge x$, $\beta\defeq r_n\vee x$. The process $Y_t\defeq X_{t\wedge S_n}$ takes values in $[\alpha,\beta]\subset(0,\infty)$, and by \bk{(5.4)} it is the continuous semimartingale
\[
  Y_t=x+\int_0^tk\,\ind_{\{s\le S_n\}}\,ds+\int_0^t\ell\sqrt{X_s}\,\ind_{\{s\le S_n\}}\,dW_s ,
  \qquad
  \qv{Y}_t=\int_0^t\ell^2X_s\,\ind_{\{s\le S_n\}}\,ds .
\]
Let $g\in C^2(\R)$ coincide with $y\mapsto\sqrt{y}$ on $[\alpha/2,2\beta]$; then $g'(y)=\tfrac12y^{-1/2}$ and $g''(y)=-\tfrac14y^{-3/2}$ for $y\in[\alpha,\beta]$. It\^o's rule (\bk{Theorem 3.3.3}) applied to $g(Y_t)$ gives, since $Y_s=X_s$ for $s\le S_n$,
\[
  \sqrt{Y_t}=\sqrt{x}+\int_0^{t\wedge S_n}\frac{k}{2\sqrt{X_s}}\,ds
  +\int_0^t\frac{\ell}{2}\,\ind_{\{s\le S_n\}}\,dW_s
  -\frac18\int_0^{t\wedge S_n}X_s^{-3/2}\,\ell^2X_s\,ds ,
\]
and this is \eqref{eq:5-5-35:sqrt}, because $\int_0^t\ind_{\{s\le S_n\}}\,dW_s=W_{t\wedge S_n}$.
\end{proof}

\textbf{(iii)} Let $\ell=2$ and $k=d\in\{2,3,\ldots\}$. \emph{The solution $X$ is the square of a Bessel process with dimension $d$ starting at $\sqrt{x}$}; precisely, (A) $R\defeq\sqrt{X}$ satisfies the integral equation \bk{(3.3.16)} of the Bessel process with dimension $d$, driven by $W$; (B) the square of the Bessel process with dimension $d$ starting at $\sqrt x$ is a weak solution of \eqref{eq:5-5-35:sde}; and (C) $X$ has the same law as the square of that Bessel process.

\emph{(A).} Now $\mu=d/2\ge1$, so $P[S=\infty]=1$ by (ii): almost surely $S_n\uparrow\infty$ and $X_t\in(0,\infty)$ for all $t\ge0$. In \eqref{eq:5-5-35:sqrt} the constant is $(4d-4)/8=(d-1)/2$ and $\ell/2=1$; letting $n\to\infty$ we obtain
\begin{equation}\label{eq:5-5-35:bessel}
  R_t=\sqrt{x}+\int_0^t\frac{d-1}{2R_s}\,ds+W_t;\qquad 0\le t<\infty,\ \text{a.s.},
\end{equation}
where the integral is finite because $R$ is continuous and strictly positive. This is \bk{(3.3.16)} with $r=\sqrt{x}$ and $B=W$. Since $\int_0^t\sqrt{X_s}\,\ind_{\{s\le S_n\}}\,dW_s=\int_0^{t\wedge S_n}\sqrt{X_s}\,dW_s$ and the continuous process $\sqrt X$ is integrable with respect to $W$, letting $n\to\infty$ in \bk{(5.4)} itself we also get
\begin{equation}\label{eq:5-5-35:global}
  X_t=x+d\,t+2\int_0^t\sqrt{X_s}\,dW_s;\qquad 0\le t<\infty,\ \text{a.s.}
\end{equation}

\emph{(B).} Let $\wt W=(\wt W^{(1)},\ldots,\wt W^{(d)})$ be a $d$-dimensional Brownian motion with $\wt W_0=(\sqrt x,0,\ldots,0)$, on a filtered probability space satisfying the usual conditions, and let $\wt R_t\defeq\norm{\wt W_t}$ be the Bessel process with dimension $d$ starting at $\sqrt x$ (\bk{Definition 3.3.19}). By \bk{Proposition 3.3.22} and $\wt R_0=\sqrt x>0$, almost surely $\wt R_t>0$ for all $t\ge0$. By \bk{Proposition 3.3.21}, the process $B=\sum_{i=1}^dB^{(i)}$ with $B^{(i)}_t=\int_0^t(\wt W^{(i)}_s/\wt R_s)\,d\wt W^{(i)}_s$ is a standard Brownian motion. It\^o's rule (\bk{Theorem 3.3.6}) for $f(y)=\norm{y}^2$, which is of class $C^2(\R^d)$ with $\partial_if=2y_i$, $\partial_i\partial_jf=2\delta_{ij}$, gives for $\wt X_t\defeq\wt R_t^2$
\[
  \wt X_t=x+2\sum_{i=1}^d\int_0^t\wt W^{(i)}_s\,d\wt W^{(i)}_s+d\,t .
\]
For each $i$ and each $T<\infty$, the integrand $\wt W^{(i)}/\wt R$ is bounded by one, and $E\int_0^T\wt R_s^2\,d\qv{B^{(i)}}_s\le E\int_0^T\wt R_s^2\,ds=xT+dT^2/2<\infty$; hence \bk{Corollary 3.2.20} applies (with the martingale $\wt W^{(i)}-\wt W^{(i)}_0\in\cM_2^c$, the integrand $\wt W^{(i)}/\wt R\in\cL^*(\wt W^{(i)}-\wt W^{(i)}_0)$ and $\wt R\in\cL^*(B^{(i)})$) and yields $\int_0^t\wt R_s\,dB^{(i)}_s=\int_0^t\wt R_s\,(\wt W^{(i)}_s/\wt R_s)\,d\wt W^{(i)}_s=\int_0^t\wt W^{(i)}_s\,d\wt W^{(i)}_s$. Summing over $i$,
\[
  \wt X_t=x+d\,t+2\int_0^t\sqrt{\wt X_s}\,dB_s;\qquad 0\le t<\infty .
\]
Thus $(\wt X,B)$ satisfies \eqref{eq:5-5-35:global}; it is $(0,\infty)$-valued, so it is a weak solution of \eqref{eq:5-5-35:sde} in $I$ with $S=\infty$, in agreement with (ii). (This also agrees with \bk{Example 5.25}: the Bessel process is in case (a) for $d=2$ and in case (c) for $d\ge3$, exactly as $\mu=d/2$ is equal to, or greater than, one.)

\emph{(C).} By \eqref{eq:5-5-35:global} and (B), both $(X,W)$ and $(\wt X,B)$ are weak solutions, in the sense of \bk{Definition 3.1} and with the same initial condition $x$, of the equation
\[
  dZ_t=d\,dt+2\sqrt{\abs{Z_t}}\,dW_t
\]
on the whole real line. (Condition (iii) of \bk{Definition 2.1}, required in \bk{Definition 3.1}, holds because the paths are continuous and real-valued, so that the integrand $d+4\abs{Z_s}$ is bounded on every $[0,t]$; condition (iv) of \bk{Definition 2.1} is \eqref{eq:5-5-35:global} for $(X,W)$, and the last display of (B) for $(\wt X,B)$. For $X$, which is $(0,\infty)$-valued only off the $P$-null event $\{S<\infty\}$, we first redefine $X_t\defeq x$ for all $t$ on this event, which belongs to $\F_0$ by the usual conditions, so that adaptedness and \eqref{eq:5-5-35:global} are not affected.) Its drift is constant, and its dispersion coefficient satisfies $\bigl\lvert2\sqrt{\abs{y}}-2\sqrt{\abs{z}}\bigr\rvert\le2\sqrt{\abs{y-z}}=h(\abs{y-z})$ with $h(u)\defeq2\sqrt u$, a strictly increasing function with $h(0)=0$ and $\int_{(0,\eps)}h^{-2}(u)\,du=\int_0^\eps du/(4u)=\infty$. Hence pathwise uniqueness holds by \bk{Proposition 2.13} and \bk{Remark 3.3}, and uniqueness in the sense of probability law follows from \bk{Proposition 3.20}. Therefore $X$ has the law of $\wt R^{\,2}$, and $\sqrt X$ has the law of the Bessel process with dimension $d$ starting at $\sqrt x$.

\rmk{For general $k,\ell$, the process $Z\defeq4X/\ell^2$ satisfies $dZ_t=2\mu\,dt+2\sqrt{Z_t}\,dW_t$ up to its exit time from $I$. So whenever $2\mu=4k/\ell^2=d$ is an integer $\ge2$, $X$ is $\ell^2/4$ times the square of a $d$-dimensional Bessel process; the trichotomy of (i), (ii) is the dichotomy ``dimension $=2$ / dimension $>2$'' of \bk{Example 5.25}, extended to ``dimension'' $2\mu<2$, where the origin is attained. Part (iv) is the case of dimension $2\mu=1$.}

\textbf{(iv)} Let $\ell^2=4k$, so that $\mu=\tfrac12<1$: by (ii), $S<\infty$ and $X_S=0$ a.s. \bk{Definition 5.20} imposes nothing on $X$ after the exit time: conditions \bk{(5.3)}, \bk{(5.4)} involve $X$ only on the intervals $[0,S_n]$, and on $[S,\infty)$ a solution may be any continuous adapted $[0,\infty]$-valued continuation. In this part we adopt the stipulation
\begin{equation}\label{eq:5-5-35:stip}
  X_t=X_S;\qquad S\le t<\infty,
\end{equation}
which extends to $I=(0,\infty)$ the stipulation that the book makes after \bk{(5.57)} for $I=(-\infty,\infty)$. Without it, everything below remains true with ``$0\le t<\infty$'' replaced by ``$0\le t\le S$''. We claim that, with
\[
  T\defeq\inf\bigl\{t\ge0;\ W_t=-2\sqrt{x}/\ell\bigr\},
\]
we have, almost surely,
\begin{equation}\label{eq:5-5-35:explicit}
  S=T<\infty\qquad\text{and}\qquad X_t=\Bigl(\sqrt{x}+\frac{\ell}{2}\,W_{t\wedge T}\Bigr)^{2};\quad 0\le t<\infty
\end{equation}
for every solution satisfying \eqref{eq:5-5-35:stip} (for $t\ge S$ the formula expresses $X_t=X_S=0$), and that conversely \eqref{eq:5-5-35:explicit} defines a solution, which satisfies \eqref{eq:5-5-35:stip}. Thus the equation has a strong solution, and any two solutions defined on the same space with the same Brownian motion $W$ have the same exit time $T$ and coincide on $[0,T]$ a.s., hence on $[0,\infty)$ under \eqref{eq:5-5-35:stip}: the solution is the square of a Brownian motion with variance parameter $\ell^2/4=k$ started at $\sqrt x$, stopped when it reaches the origin.

\emph{How the formula is found.} We have $\tfrac12\sigma(y)\sigma'(y)=\tfrac12\,\ell\sqrt{y}\cdot\ell/(2\sqrt{y})=\ell^2/4=k$, so \eqref{eq:5-5-35:sde} is of the form \bk{(2.36)$'$} with the function $b$ of that equation identically zero. \bk{Proposition 2.21} then suggests $X_t=u(W_t,Y_t)$, where $\partial u/\partial w=\sigma(u)=\ell\sqrt u$, $u(0,y)=y$, that is, $u(w,y)=(\sqrt y+\ell w/2)^2$ as long as $\sqrt{y}+\ell w/2\ge0$, and where $Y$ solves an ordinary differential equation with zero right-hand side, so $Y_t\equiv x$. The hypotheses of \bk{Proposition 2.21} are not satisfied ($\sigma$ is not of class $C^2(\R)$ with bounded derivatives), so we verify the result directly.

\emph{Every solution is given by \eqref{eq:5-5-35:explicit}.} Since $4k-\ell^2=0$, the lemma gives $\sqrt{X_{t\wedge S_n}}=\sqrt x+\tfrac{\ell}{2}W_{t\wedge S_n}$ for all $t\ge0$ and $n\ge1$, a.s. As $S_n\uparrow S$ and $t<S$ implies $t<S_n$ for large $n$, we obtain on an event $\Omega^*$ of probability one
\begin{equation}\label{eq:5-5-35:before}
  \sqrt{X_t}=\sqrt{x}+\frac{\ell}{2}\,W_t;\qquad 0\le t<S .
\end{equation}
Fix $\omega\in\Omega^*$. For $t<S(\omega)$ we have $X_t(\omega)>0$, so the right-hand side of \eqref{eq:5-5-35:before} is strictly positive and $W_t(\omega)\neq-2\sqrt x/\ell$; hence $T(\omega)\ge S(\omega)$. If $S(\omega)<\infty$, then letting $t\uparrow S(\omega)$ in \eqref{eq:5-5-35:before} and using the continuity of $X$ (in $[0,\infty]$) and of $W$, we find $\sqrt{X_S(\omega)}=\sqrt x+\tfrac\ell2W_S(\omega)<\infty$; since $X_S(\omega)\in\{0,\infty\}$, necessarily $X_S(\omega)=0$ and $W_S(\omega)=-2\sqrt x/\ell$, so $T(\omega)\le S(\omega)$. Thus $S=T$ a.s.\ (if $S(\omega)=\infty$ this follows from $T\ge S$), \eqref{eq:5-5-35:before} gives the formula in \eqref{eq:5-5-35:explicit} for $t<S$; for $t=S=T<\infty$ both sides vanish, since $X_S=0$ and $W_T=-2\sqrt x/\ell$ by continuity of $W$; and for $t>S$ the left-hand side equals $X_S=0$ by \eqref{eq:5-5-35:stip}, while the right-hand side is $(\sqrt x+\tfrac\ell2W_T)^2=0$. Finally $T<\infty$ a.s., because almost every Brownian path is unbounded from below (\bk{Theorem 2.9.23}); this recovers the conclusion of (ii) without \bk{Proposition 5.32}.

\emph{Formula \eqref{eq:5-5-35:explicit} defines a solution.} Conversely, let $W=\{W_t,\F_t\}$ be any standard Brownian motion with a filtration satisfying the usual conditions, define $T$ as above and $X$ by the formula in \eqref{eq:5-5-35:explicit}. Then $X$ is continuous, adapted (even to the augmented filtration of $W$), $[0,\infty)$-valued, $X_0=x$, and $X_t>0$ for $t<T$, $X_t=0$ for $t\ge T$; hence its exit time from $I$ is $S=T$, \eqref{eq:5-5-35:stip} holds, and $S_n\le T$ for the times of \bk{(5.56)}. It\^o's rule applied to the function $w\mapsto(\sqrt x+\ell w/2)^2$ and the continuous martingale $\{W_{t\wedge S_n};\,t\ge0\}$, whose quadratic variation is $t\wedge S_n$, gives
\[
  X_{t\wedge S_n}=x+\int_0^t\ell\Bigl(\sqrt x+\frac\ell2W_s\Bigr)\ind_{\{s\le S_n\}}\,dW_s+\frac{\ell^2}{4}\,(t\wedge S_n)
  =x+\int_0^tk\,\ind_{\{s\le S_n\}}\,ds+\int_0^t\ell\sqrt{X_s}\,\ind_{\{s\le S_n\}}\,dW_s ,
\]
because $\ell^2/4=k$ and, for $s\le S_n\le T$, $\sqrt x+\tfrac\ell2W_s\ge0$ equals $\sqrt{X_s}$. This is \bk{(5.4)}; and \bk{(5.3)} holds since $\abs{b}+\sigma^2=k+\ell^2X_s$ is bounded on $[0,S_n]$. So $(X,W)$ is a (strong) solution in $I$ in the sense of \bk{Definition 5.20}.

\rmk{As a check on \eqref{eq:5-5-35:v}: $v'(x)=(1-x^{-\mu})/k$ and $v''(x)=\mu x^{-\mu-1}/k$ give $\tfrac12\ell^2x\,v''(x)+k\,v'(x)=\tfrac{\ell^2\mu}{2k}x^{-\mu}+1-x^{-\mu}=1$, with $v(1)=v'(1)=0$, as required by \bk{(5.70)} for $u_1=v$. In (iv), $S=T$ is the first passage time of the standard Brownian motion $\{W_t,\F_t\}$ to the level $a\defeq-2\sqrt x/\ell\neq0$, which is a.s.\ finite but has infinite expectation: by \bk{Proposition 2.6.19}, $T$ has the density \bk{(2.6.3)} with $b=a$, $P[T\in dt]=\abs{a}(2\pi t^3)^{-1/2}e^{-a^2/2t}\,dt$ on $(0,\infty)$, and since $e^{-a^2/2t}\ge e^{-1/2}$ for $t\ge a^2$,
\[
  ET\ge\int_{a^2}^\infty t\,\frac{\abs{a}}{\sqrt{2\pi t^3}}\,e^{-1/2}\,dt=\frac{\abs{a}e^{-1/2}}{\sqrt{2\pi}}\int_{a^2}^\infty\frac{dt}{\sqrt t}=\infty .
\]
This is consistent with \bk{Proposition 5.32}, which asserts $ES<\infty$ only under its condition (i), whereas here $v(\infty-)=\infty$.}
\end{kex}

\begin{kex}{Exercise}{5.36}{}
Here $I=(-\infty,\infty)$, $b(x)=(1+x)(1+x^2)$ and $\sigma(x)=1+x^2$. Both functions are continuous and $\sigma^2>0$, so \bk{(ND)$'$} and \bk{(LI)$'$} hold. We let $(X,W)$, $(\Omega,\F,P)$, $\{\F_t\}$ be a weak solution up to an explosion time in the sense of \bk{Definition 5.1} (to which \bk{Definition 5.20} reduces when $I=\R$), with $S_n=\inf\{t\ge0;\ \abs{X_t}\ge n\}$ and $S=\lim_{n\to\infty}S_n$ as in \bk{(5.2)}, \bk{(5.5)}. We have to show $ES<\infty$. This implies $P[S<\infty]=1$, and on $\{S<\infty\}$ the solution does reach $\pm\infty$ at time $S$: for $n>\abs{X_0}$ the continuity of $X$ gives $\abs{X_{S_n}}=n$, whence $\abs{X_S}=\lim_{n\to\infty}\abs{X_{S_n}}=\infty$. In Step~1 we verify the criterion of \bk{Proposition 5.32} when $X_0=x$ is nonrandom; in Step~2 we compute $ES$ directly, for an arbitrary (possibly random) real-valued initial condition $X_0$.

\emph{Step 1: Feller's test.} We take $c=0$ in \bk{(5.42)}. Since
\[
  \int_0^\xi\frac{2b(\zeta)}{\sigma^2(\zeta)}\,d\zeta=\int_0^\xi\frac{2+2\zeta}{1+\zeta^2}\,d\zeta
  =2\arctan\xi+\log(1+\xi^2),
\]
the scale function has the derivative $p'(\xi)=e^{-2\arctan\xi}/(1+\xi^2)$, and the substitution $\theta=\arctan\xi$ yields
\begin{equation}\label{eq:5-5-36:p}
  p(x)=\int_0^{\arctan x}e^{-2\theta}\,d\theta=\tfrac12\bigl(1-e^{-2\arctan x}\bigr);\qquad
  p(\infty)=\tfrac12(1-e^{-\pi}),\quad p(-\infty)=-\tfrac12(e^{\pi}-1).
\end{equation}
The speed measure \bk{(5.51)} is $m(dx)=2e^{2\arctan x}(1+x^2)^{-1}\,dx$, so that $\int_0^y m(dz)=e^{2\arctan y}-1$ for every $y\in\R$ (with the convention $\int_0^y=-\int_y^0$ for $y<0$), and the first expression in \bk{(5.65)} becomes
\begin{equation}\label{eq:5-5-36:v}
  v(x)=\int_0^x p'(y)\bigl(e^{2\arctan y}-1\bigr)\,dy
  =\int_0^x\frac{1-e^{-2\arctan y}}{1+y^2}\,dy=\arctan x-p(x);\qquad x\in\R.
\end{equation}
Consequently
\[
  v(\infty)=\tfrac12\bigl(\pi-1+e^{-\pi}\bigr)<\infty,\qquad v(-\infty)=\tfrac12\bigl(e^{\pi}-1-\pi\bigr)<\infty,
\]
which is condition (i) of \bk{Proposition 5.32}. By that proposition and \bk{Remark 5.33}, for every nonrandom initial condition $X_0=x\in\R$ we have $ES<\infty$.

\emph{Step 2: direct computation.} Let $\beta\defeq\pi/(2\sinh\pi)$ and
\begin{equation}\label{eq:5-5-36:g}
  g(\theta)\defeq\frac{\pi}{2}\coth\pi-\theta-\beta e^{-2\theta};\qquad \theta\in\Bigl[-\frac{\pi}{2},\frac{\pi}{2}\Bigr].
\end{equation}
Then $g'(\theta)=-1+2\beta e^{-2\theta}$ and $g''(\theta)=-4\beta e^{-2\theta}$, so
\begin{equation}\label{eq:5-5-36:ode}
  g'+\tfrac12 g''=-1,\qquad
  g\Bigl(\pm\frac{\pi}{2}\Bigr)=\frac{\pi}{2\sinh\pi}\bigl(\cosh\pi\mp\sinh\pi-e^{\mp\pi}\bigr)=0 .
\end{equation}
Being strictly concave and equal to zero at both endpoints, $g$ is strictly positive on $(-\pi/2,\pi/2)$; we set $K\defeq\max g<\infty$, and note $\abs{g'}\le 1+2\beta e^{\pi}$. The function $h(x)\defeq g(\arctan x)$, $x\in\R$, is of class $C^2$, satisfies $0<h\le K$ and $\lim_{x\to\pm\infty}h(x)=0$, and has the derivatives
\[
  h'(x)=\frac{g'(\arctan x)}{1+x^2},\qquad h''(x)=\frac{g''(\arctan x)-2x\,g'(\arctan x)}{(1+x^2)^2}.
\]
Therefore $\sigma(x)h'(x)=g'(\arctan x)$ and, by \eqref{eq:5-5-36:ode},
\begin{equation}\label{eq:5-5-36:gen}
  b(x)h'(x)+\tfrac12\sigma^2(x)h''(x)=(1+x)g'(\arctan x)+\tfrac12 g''(\arctan x)-x\,g'(\arctan x)=-1;\qquad x\in\R.
\end{equation}

Fix $n\ge1$. The process $Y_t\defeq X_{t\wedge S_n}$ is a.s.\ real-valued, because $\abs{Y_t}\le n\vee\abs{X_0}$ and $\abs{X_0}<\infty$ a.s.; by \bk{(5.3)} and \bk{(5.4)} it is a continuous semimartingale,
\[
  Y_t=X_0+\int_0^t b(X_s)\ind_{\{s\le S_n\}}\,ds+\int_0^t\sigma(X_s)\ind_{\{s\le S_n\}}\,dW_s;\qquad 0\le t<\infty,
\]
with $X_0$ measurable with respect to $\F_0$. It\^o's rule (\bk{Theorem 3.3.3}), together with $Y_s=X_s$ for $s\le S_n$ and with \eqref{eq:5-5-36:gen}, gives almost surely
\begin{equation}\label{eq:5-5-36:ito}
  h(X_{t\wedge S_n})=h(X_0)-(t\wedge S_n)+\int_0^t g'(\arctan X_s)\ind_{\{s\le S_n\}}\,dW_s;\qquad 0\le t<\infty.
\end{equation}
The integrand in \eqref{eq:5-5-36:ito} is left-continuous, adapted and bounded by $1+2\beta e^{\pi}$, so the stochastic integral is a square-integrable martingale starting at zero and has expectation zero. Since $h$ is bounded, we may take expectations in \eqref{eq:5-5-36:ito}; with $t=n$ we obtain
\begin{equation}\label{eq:5-5-36:exp}
  E(n\wedge S_n)=E\,h(X_0)-E\,h(X_{n\wedge S_n})\le E\,h(X_0)\le K;\qquad n\ge1.
\end{equation}
As $n\to\infty$ we have $n\wedge S_n\uparrow S$, and the monotone convergence theorem gives $ES\le K<\infty$. This proves the claim, for every initial condition.

In fact \eqref{eq:5-5-36:exp} yields the value of $ES$. We now know that $S<\infty$ a.s. For almost every $\omega$ and every integer $n>S(\omega)\vee\abs{X_0(\omega)}$ we have $S_n(\omega)\le S(\omega)<n$ and $\abs{X_{S_n}(\omega)}=n$, hence $h(X_{n\wedge S_n}(\omega))\in\{h(n),h(-n)\}$; thus $h(X_{n\wedge S_n})\to0$ a.s., and $E\,h(X_{n\wedge S_n})\to0$ by the bounded convergence theorem. Letting $n\to\infty$ in \eqref{eq:5-5-36:exp} we conclude
\begin{equation}\label{eq:5-5-36:ES}
  ES=E\,g(\arctan X_0)\le K,\qquad\text{and}\qquad
  ES=\frac{\pi}{2}\coth\pi-\arctan x-\frac{\pi\,e^{-2\arctan x}}{2\sinh\pi}\quad\text{if }X_0=x\in\R .
\end{equation}
For instance, $ES=\frac{\pi}{2}\tanh\frac{\pi}{2}\approx1.4407$ when $X_0=0$.

\rmk{The restriction $X_0\in(-\frac{\pi}{2},\frac{\pi}{2})$ in the statement plays no role: the state space of the equation is all of $\R$, and both steps above hold for every real initial condition. The interval $(-\frac{\pi}{2},\frac{\pi}{2})$ is the range of $\arctan X$, and this is presumably what is meant. Indeed, It\^o's rule applied to $\arctan X_{t\wedge S_n}$ as in Step~2, with $b(x)(1+x^2)^{-1}-x\sigma^2(x)(1+x^2)^{-2}=1$, gives $\arctan X_{t\wedge S_n}=\arctan X_0+(t\wedge S_n)+W_{t\wedge S_n}$: up to time $S$ the process $\arctan X_t$ is the Brownian motion with drift $\arctan X_0+t+W_t$, and $S$ is its first exit time from $(-\frac{\pi}{2},\frac{\pi}{2})$. The function $g$ of \eqref{eq:5-5-36:g} is the expected exit time of this process as a function of its starting point, and \eqref{eq:5-5-36:ES} agrees with the formula $ES=M(x)$ of \bk{Remark 5.33}: a numerical quadrature of $\int_{-\infty}^\infty G(x,y)\,m(dy)$ with $p$, $m$ as in Step~1 reproduces $g(\arctan x)$, as does a simulation of the equation.}
\end{kex}

\begin{kex}{Exercise}{5.37}{}
In each case $(X,W)$, $(\Omega,\F,P)$, $\{\F_t\}$ denotes an arbitrary weak solution in $I=(\ell,r)$ of \bk{(5.1)} in the sense of \bk{Definition 5.20}, with nonrandom initial condition $X_0=x\in I$, and $S$ is the exit time \bk{(5.57)}. In all three cases the coefficients are continuous on $I$ and $\sigma^2>0$ there, so $(1+\abs{b})/\sigma^2$ is continuous, hence locally integrable, on $I$: the conditions \bk{(ND)$'$} and \bk{(LI)$'$} hold, and \bk{Proposition 5.22} and \bk{Theorem 5.29} are applicable. We shall show that
\[
  P[S=\infty]=1\qquad\text{in each of the cases (i), (ii), (iii), for all values of the parameters,}
\]
and describe along the way the behaviour of $X$ near the endpoints of $I$. We write $p$ for the scale function \bk{(5.42)}, $m$ for the speed measure \bk{(5.51)} and $v$ for the function \bk{(5.65)}; we recall that $v\ge 0$, and that for $\ell<x<c$ the formula \bk{(5.65)} reads $v(x)=\int_x^c p'(y)\,m((y,c))\,dy$.

\textbf{(i)} We take $c=0$. Since $2b(\zeta)/\sigma^2(\zeta)=-\zeta/(1+\zeta^2)$, we obtain
\[
  p'(\xi)=\exp\Bigl\{\int_0^\xi\frac{\zeta\,d\zeta}{1+\zeta^2}\Bigr\}=\sqrt{1+\xi^2},
  \qquad
  p(x)=\int_0^x\sqrt{1+\xi^2}\,d\xi;\qquad x\in\R .
\]
Because $p'\ge 1$, we have $p(x)\ge x$ for $x\ge0$ and $p(x)\le x$ for $x\le 0$, so $p(-\infty)=-\infty$ and $p(\infty)=\infty$. We are in case (a) of \bk{Proposition 5.22}: \emph{there is no explosion, $P[S=\infty]=1$, and $X$ is recurrent}, with $\sup_{0\le t<\infty}X_t=\infty$ and $\inf_{0\le t<\infty}X_t=-\infty$ a.s. (Equivalently, $v(\pm\infty)=\infty$ by \bk{Problem 5.27}, and \bk{Theorem 5.29} applies; alternatively, $\abs{b(x)}+\abs{\sigma(x)}\le (1+\sqrt2)(1+\abs{x})$, and \bk{Remark 5.19} applies.)

\textbf{(ii)} We take $c=1$ and set $\gamma\defeq 2\mu/\delta^2\in\R$. Here $2b(\zeta)/\sigma^2(\zeta)=\gamma/\zeta$, so for $0<x<\infty$
\begin{equation}\label{eq:5-5-37:pii}
  p'(\xi)=\exp\Bigl\{-\gamma\int_1^\xi\frac{d\zeta}{\zeta}\Bigr\}=\xi^{-\gamma},
  \qquad
  p(x)=
  \begin{cases}
    \dfrac{x^{1-\gamma}-1}{1-\gamma}; & \gamma\ne 1,\\[1.5ex]
    \log x; & \gamma=1,
  \end{cases}
\end{equation}
and $m(dz)=2\,dz/\bigl(p'(z)\delta^2z^2\bigr)=(2/\delta^2)\,z^{\gamma-2}\,dz$. For $\gamma\neq 1$ we have $\int_1^y z^{\gamma-2}\,dz=(y^{\gamma-1}-1)/(\gamma-1)$ for every $y>0$, and \bk{(5.65)} gives
\begin{equation}\label{eq:5-5-37:vii}
  v(x)=\frac{2}{\delta^2(\gamma-1)}\int_1^x\bigl(y^{-1}-y^{-\gamma}\bigr)\,dy
      =\frac{2}{\delta^2(\gamma-1)}\,\bigl[\log x-p(x)\bigr];\qquad 0<x<\infty,
\end{equation}
whereas for $\gamma=1$ we have $\int_1^y z^{-1}\,dz=\log y$ and $v(x)=(2/\delta^2)\int_1^x y^{-1}\log y\,dy=(\log x)^2/\delta^2$. We distinguish three cases.

\emph{Case $\gamma=1$} ($\mu=\delta^2/2$). Then $p(0+)=-\infty$, $p(\infty)=\infty$, and $v(0+)=v(\infty)=\infty$.

\emph{Case $\gamma<1$} ($\mu<\delta^2/2$). Then $p(0+)=-1/(1-\gamma)>-\infty$ and $p(\infty)=\infty$. As $x\downarrow 0$, the bracket in \eqref{eq:5-5-37:vii} tends to $-\infty$ (because $\log x\to-\infty$ and $p(x)$ stays bounded); as $x\uparrow\infty$ it also tends to $-\infty$ (because $x^{1-\gamma}/\log x\to\infty$). The factor in front of the bracket is negative, so $v(0+)=v(\infty)=\infty$.

\emph{Case $\gamma>1$} ($\mu>\delta^2/2$). Then $p(x)=(1-x^{-(\gamma-1)})/(\gamma-1)$, so $p(0+)=-\infty$ and $p(\infty)=1/(\gamma-1)<\infty$. As $x\downarrow0$ the bracket in \eqref{eq:5-5-37:vii} tends to $+\infty$ (because $x^{-(\gamma-1)}/\abs{\log x}\to\infty$), and as $x\uparrow\infty$ it tends to $+\infty$ as well (because $p$ is bounded above). The factor in front is positive, so again $v(0+)=v(\infty)=\infty$.

Thus \bk{(5.77)} holds for all $\mu\in\R$, $\delta>0$, and \bk{Theorem 5.29} gives $P[S=\infty]=1$: \emph{the solution never reaches $0$ or $\infty$ in finite time.} Since $S=\infty$ a.s., \bk{Proposition 5.22} describes the behaviour as $t\to\infty$: for $\mu=\delta^2/2$ (case (a)) $X$ is recurrent in $(0,\infty)$, with $\inf_{t}X_t=0$ and $\sup_tX_t=\infty$; for $\mu<\delta^2/2$ (case (b)) $\lim_{t\to\infty}X_t=0$ and $\sup_tX_t<\infty$; for $\mu>\delta^2/2$ (case (c)) $\lim_{t\to\infty}X_t=\infty$ and $\inf_tX_t>0$, a.s. Both endpoints are thus inaccessible, although one of them is attracting when $\mu\ne\delta^2/2$. This is the content of \bk{Exercise 5.31} (with $\nu=\delta$), and it agrees with the explicit formula $X_t=x\exp\{(\mu-\tfrac12\delta^2)t+\delta W_t\}$, which is established in \eqref{eq:5-5-37:rep} below (with $\kappa=0$).

\textbf{(iii)} We take $c=1$ and set $\alpha\defeq 2/\delta^2>0$, $\beta\defeq 2\mu/\delta\in\R$. Now
\[
  \frac{2b(\zeta)}{\sigma^2(\zeta)}=\frac{\alpha}{\zeta^2}+\frac{\beta}{\zeta},
  \qquad
  \int_1^\xi\frac{2b(\zeta)}{\sigma^2(\zeta)}\,d\zeta=\alpha\Bigl(1-\frac1\xi\Bigr)+\beta\log\xi ,
\]
so that, for $0<\xi<\infty$ and $0<z<\infty$,
\begin{equation}\label{eq:5-5-37:piii}
  p'(\xi)=e^{-\alpha}\,\xi^{-\beta}e^{\alpha/\xi},
  \qquad
  m(dz)=\frac{2\,dz}{p'(z)\,\delta^2z^2}=\frac{2e^{\alpha}}{\delta^2}\,z^{\beta-2}e^{-\alpha/z}\,dz .
\end{equation}

\emph{The endpoint $0$.} Since $e^{\alpha/\xi}\ge (\alpha/\xi)^k/k!$ for every integer $k\ge 0$, choosing $k\ge 1-\beta$ we get $p'(\xi)\ge (e^{-\alpha}\alpha^k/k!)\,\xi^{-\beta-k}\ge (e^{-\alpha}\alpha^k/k!)\,\xi^{-1}$ for $0<\xi\le 1$, and therefore
\[
  p(0+)=-\int_0^1p'(\xi)\,d\xi=-\infty .
\]
Consequently $v(0+)=\infty$; this is \bk{(5.74)}, and it is seen directly: for $0<x<\tfrac12$,
\[
  v(x)=\int_x^1p'(y)\,m((y,1))\,dy\ \ge\ m\bigl((\tfrac12,1)\bigr)\int_x^{1/2}p'(y)\,dy
  \xrightarrow[x\downarrow0]{}\infty .
\]

\emph{The endpoint $\infty$.} For $\xi\ge1$ we have $1\le e^{\alpha/\xi}\le e^{\alpha}$, hence $e^{-\alpha}\xi^{-\beta}\le p'(\xi)\le\xi^{-\beta}$, and thus
\[
  p(\infty)=\infty\iff\beta\le1\iff\mu\le\delta/2 .
\]
If $\beta\le 1$, then $p(0+)=-\infty$, $p(\infty)=\infty$ and we are in case (a) of \bk{Proposition 5.22}: $P[S=\infty]=1$ and $X$ is recurrent in $(0,\infty)$. If $\beta>1$, then $p(\infty)<\infty$ and we must compute $v(\infty)$. From \eqref{eq:5-5-37:piii} and $e^{-\alpha/z}\ge e^{-\alpha}$ for $z\ge 1$,
\[
  m((1,y))\ \ge\ \frac{2}{\delta^2}\int_1^yz^{\beta-2}\,dz=\frac{2}{\delta^2}\cdot\frac{y^{\beta-1}-1}{\beta-1};
  \qquad y\ge1,
\]
and with $p'(y)\ge e^{-\alpha}y^{-\beta}$ we obtain for $x\ge 1$
\[
  v(x)=\int_1^xp'(y)\,m((1,y))\,dy\ \ge\ \frac{2e^{-\alpha}}{\delta^2(\beta-1)}\int_1^x\bigl(y^{-1}-y^{-\beta}\bigr)\,dy
  \ \ge\ \frac{2e^{-\alpha}}{\delta^2(\beta-1)}\Bigl(\log x-\frac{1}{\beta-1}\Bigr),
\]
which tends to infinity as $x\uparrow\infty$. Hence $v(0+)=v(\infty)=\infty$ also when $\beta>1$, and \bk{Theorem 5.29} gives $P[S=\infty]=1$; we are then in case (c) of \bk{Proposition 5.22}, so $\lim_{t\to\infty}X_t=\infty$ and $\inf_{0\le t<\infty}X_t>0$, a.s.

To summarize: \emph{in case (iii) there is no explosion for any $\mu\in\R$, $\delta>0$; the solution is recurrent in $(0,\infty)$ if $\mu\le\delta/2$, and tends to $\infty$ as $t\to\infty$ (without reaching it in finite time, and staying bounded away from $0$) if $\mu>\delta/2$.}

\emph{The explicit solution (the hint).} The conclusions of (ii) and (iii) can be confirmed without Feller's test. Let $\kappa\in\{0,1\}$, $b(x)=\kappa+\delta\mu x$, $\sigma(x)=\delta x$ on $I=(0,\infty)$; the value $\kappa=1$ is case (iii), and $\kappa=0$ is case (ii) with $\delta\mu$ in place of $\mu$. With $W^{(\mu)}_t\defeq W_t+\mu t$ as in \bk{Exercise 3.3.33}, set
\[
  Z_t\defeq\exp\bigl\{\delta W^{(\mu)}_t-\tfrac12\delta^2t\bigr\},\qquad U_t\defeq 1/Z_t=\exp\bigl\{-\delta W_t-(\delta\mu-\tfrac12\delta^2)t\bigr\};\qquad 0\le t<\infty .
\]
We claim that, almost surely,
\begin{equation}\label{eq:5-5-37:rep}
  X_{t\wedge S}=Z_{t\wedge S}\Bigl(x+\kappa\int_0^{t\wedge S}\frac{ds}{Z_s}\Bigr);\qquad 0\le t<\infty .
\end{equation}
Indeed, It\^o's rule (\bk{Theorem 3.3.3}), applied to the exponential function and the semimartingale $-\delta W_t-(\delta\mu-\frac12\delta^2)t$, gives
\[
  U_t=1-\delta\int_0^tU_s\,dW_s+(\delta^2-\delta\mu)\int_0^tU_s\,ds .
\]
Fix $n$ so large that $x\in(\ell_n,r_n)$, with $\ell_n$, $r_n$, $S_n$ as in \bk{Definition 5.20}, and put $Y_t\defeq X_{t\wedge S_n}$. This process takes values in $[\ell_n,r_n]\subset(0,\infty)$ and, by \bk{(5.4)}, is the continuous semimartingale
\[
  Y_t=x+\int_0^t(\kappa+\delta\mu Y_s)\ind_{\{s\le S_n\}}\,ds+\delta\int_0^tY_s\ind_{\{s\le S_n\}}\,dW_s ,
\]
with bounded integrands. The integration-by-parts formula of \bk{Problem 3.3.12} yields
\begin{align*}
  Y_tU_t&=x+\int_0^tY_s\,dU_s+\int_0^tU_s\,dY_s+\qv{Y,U}_t\\
  &=x+\int_0^tU_s\ind_{\{s\le S_n\}}\bigl[-\delta Y_s+\delta Y_s\bigr]\,dW_s
     +\int_0^tU_s\ind_{\{s\le S_n\}}\bigl[(\delta^2-\delta\mu)Y_s+\kappa+\delta\mu Y_s-\delta^2Y_s\bigr]\,ds\\
  &\qquad{}+\int_0^tY_s\ind_{\{s>S_n\}}\,dU_s ,
\end{align*}
where we have used $\qv{Y,U}_t=-\delta^2\int_0^tY_sU_s\ind_{\{s\le S_n\}}\,ds$. Evaluating at $t\wedge S_n$ in place of $t$, the last integral vanishes and we are left with $X_{t\wedge S_n}U_{t\wedge S_n}=x+\kappa\int_0^{t\wedge S_n}U_s\,ds$ for all $t\ge 0$, a.s. Letting $n\to\infty$ and using $S_n\uparrow S$, the continuity of $X$ (as a $[0,\infty]$-valued process) and of $U$, and $0<Z_t<\infty$, we arrive at \eqref{eq:5-5-37:rep}.

Now the right-hand side of \eqref{eq:5-5-37:rep} is a finite, strictly positive number for every $t$, because $x>0$, $\kappa\ge 0$ and $Z$ is continuous and strictly positive. On the other hand, on $\{S<\infty\}$ the continuity of $X$ and the definition \bk{(5.57)} of $S$ force $X_S\in\{0,\infty\}$. Taking $t\ge S$ in \eqref{eq:5-5-37:rep} we conclude $P[S<\infty]=0$, in agreement with what was found above, and
\[
  X_t=Z_t\Bigl(x+\kappa\int_0^t\frac{ds}{Z_s}\Bigr)
     =x\,e^{\delta W^{(\mu)}_t-\frac12\delta^2t}+\kappa\int_0^t\exp\bigl[\delta\{W^{(\mu)}_t-W^{(\mu)}_s\}-\tfrac12\delta^2(t-s)\bigr]\,ds;\qquad 0\le t<\infty .
\]
For $\kappa=1$ this is $xZ_t$ plus the process of \bk{Exercise 3.3.33} (which corresponds to $x=0$); for $\kappa=0$, with $\mu/\delta$ in place of $\mu$, it is the geometric Brownian motion $x\exp\{(\mu-\frac12\delta^2)t+\delta W_t\}$ of (ii). The formula also exhibits the dichotomy found in (iii): if $\mu>\delta/2$, then $X_t\ge xZ_t=x\exp\{\delta W_t+\delta(\mu-\frac12\delta)t\}\to\infty$ a.s., by the strong law of large numbers for Brownian motion (\bk{Problem 2.9.3}).
\end{kex}

\begin{kex}{Exercise}{5.38}{}
\emph{Conventions.} The function $b$ is Borel-measurable, as are all coefficients in \bk{(5.1)}. Covering a compact set $K\subset\R$ by finitely many of the intervals $(x-\eps,x+\eps)$ of the hypothesis, we see that $\int_Kb^2(y)\,dy<\infty$, and then $\int_K\abs{b(y)}\,dy\le\bigl(\meas(K)\int_Kb^2(y)\,dy\bigr)^{1/2}<\infty$ by the Cauchy--Schwarz inequality: both $b^2$ and $b$ are locally integrable. To avoid a clash of notation, the Brownian motion which drives the equation is denoted by $B$, and $W$ is reserved for the Brownian motion of (ii) and (iii); thus a solution is a triple $(X,B)$, $(\Omega,\F,P)$, $\{\F_t\}$ as in \bk{Definition 5.1}, with $\sigma\equiv1$, with $S_n$ and $S$ as in \bk{(5.2)}, \bk{(5.5)}, and with $P[X_0\in G]=\mu(G)$; $G\in\B(\R)$. The Brownian motion $W=\{W_t,\G_t;\,0\le t<\infty\}$ of (ii), (iii) is defined on a space $(\Theta,\G,Q)$, possibly different from $(\Omega,\F,P)$, and $\{\G_t\}$ satisfies the usual conditions, as it must for the stochastic integral $\int_0^tb(W_s)\,dW_s$ (an integral with respect to the standard Brownian motion $W-W_0$) to be defined; expectations with respect to $Q$ are written $E^Q$. We write $C\defeq C[0,\infty)$, $\B_t\defeq\B_t(C[0,\infty))$, and $\varphi_T\colon C\to C$ for the truncation mapping $(\varphi_Tz)(s)=z(s\wedge T)$ of \bk{(3.19)}; a function on $C$ is $\B_T$-measurable if and only if it is Borel-measurable and of the form $f\circ\varphi_T$. We put $\bar b\defeq b$ on $\R$ and $\bar b(\pm\infty)\defeq0$. Finally, on $\{S>T\}$ the path $t\mapsto X_{t\wedge T}$ is real-valued and continuous, and the event on the left-hand side of the formula in (ii) is
\[
  \{X_\cdot\in\Gamma,\ S>T\}\defeq\{S>T\}\cap\{X_{\cdot\wedge T}\in\Gamma\}\in\F_T ;
\]
this is the only possible meaning, because the indicator of $\Gamma\in\B_T$ depends on a path only through its restriction to $[0,T]$, and it is an event of $\F_T$ because $\{S>T\}=\{\max_{0\le s\le T}\abs{X_s}<\infty\}\in\F_T$ and each coordinate $X_{t\wedge T}$ is $\F_T$-measurable.

\textbf{(i)} With $\sigma\equiv1$, the condition \bk{(ND)} holds, $\sigma^{-2}\equiv1$ is locally integrable at every point, and \bk{(LI)} is the local integrability of $b$ just established. The claim is \bk{Theorem 5.15}.

\textbf{(ii)} Let $(X,B)$, $(\Omega,\F,P)$, $\{\F_t\}$ be a weak solution up to the explosion time $S$ with initial distribution $\mu$, and fix $T\in(0,\infty)$. We recall from \bk{(5.6)} and the stipulation following it that, almost surely, $X_t\in\R$ for $t<S$ and $X_t=X_S=\pm\infty$ for $t\ge S$, and that $S_n<S$ on $\{S<\infty\}$ (if $\abs{X_0}<n$ then $\abs{X_{S_n}}=n$ on $\{S_n<\infty\}$, and otherwise $S_n=0<S$). We set
\[
  A_t\defeq\int_0^t\bar b^2(X_s)\,ds\in[0,\infty],\qquad m_t\defeq\max_{0\le s\le t}\abs{X_s}\in[0,\infty];\qquad 0\le t<\infty .
\]
Both processes are nondecreasing and $\{\F_t\}$-adapted ($X$ is continuous and adapted, hence progressively measurable, and Tonelli's theorem applies to $A$); $m$ is continuous in $[0,\infty]$, and $A$ is left-continuous by monotone convergence.

\emph{Step 1: almost surely, $A_t<\infty$ for every $t<S$; consequently $A$ is finite and continuous on $[0,S)$.} Fix $n\ge1$. By \bk{(5.3)}, \bk{(5.4)} there is an event of probability one on which $\abs{X_0}<\infty$ and
\[
  X^{(n)}_t\defeq X_{t\wedge S_n}=X_0+M_t+V_t;\qquad M_t\defeq B_{t\wedge S_n},\quad V_t\defeq\int_0^{t\wedge S_n}b(X_s)\,ds,\qquad 0\le t<\infty .
\]
Here $M\in\cM_2^c$ with $\qv{M}_t=t\wedge S_n$, and $V$ is the difference of the continuous, nondecreasing, adapted processes $\int_0^{t\wedge S_n}b^{\pm}(X_s)\,ds$, which are finite by \bk{(5.3)} (on the exceptional null event, which belongs to $\F_0$, we set $X^{(n)}\equiv0$, $V\equiv 0$). Thus $X^{(n)}$ is a continuous semimartingale of the form \bk{(3.7.1)}, and \bk{Theorem 3.7.1} provides a local time $\Lambda$ for it. Let $j\ge n$ be an integer and $k_j\defeq b^2\ind_{[-j,j]}$. By parts (iii) and (iv) of that theorem, almost surely,
\[
  \int_0^{t\wedge S_n}k_j(X_s)\,ds=\int_0^tk_j(X^{(n)}_s)\,d\qv{M}_s=2\int_{-j}^{j}b^2(a)\Lambda_t(a)\,da
  \le2\sup_{\abs{a}\le j}\Lambda_t(a)\cdot\int_{-j}^jb^2(a)\,da<\infty
\]
for all $t\ge0$, since the function $a\mapsto\Lambda_t(a)$ is RCLL, hence bounded on $[-j,j]$. On $\{\abs{X_0}\le j\}$ we have $\abs{X_s}\le j$ for $0\le s\le S_n$ (if $\abs{X_0}<n$ then $\abs{X_s}\le n\le j$ up to time $S_n$, and otherwise $S_n=0$), so the left-hand side is $A_{t\wedge S_n}$. Letting $j\to\infty$ we obtain $A_{t\wedge S_n}<\infty$ for all $t\ge 0$ and all $n\ge1$, almost surely. If $t<S$, then $t<S_n$ for some $n$, because $S_n\uparrow S$; hence $A_t<\infty$, and on $[0,t]$ the function $A$ is the indefinite integral of an integrable function. This proves Step~1.

\emph{Step 2: localization.} For $k\ge1$ we define
\[
  D_k\defeq\{t\ge0;\ m_t\ge k\ \text{or}\ A_t\ge k\},\qquad\tau_k\defeq T\wedge\inf D_k .
\]
Since $m$ and $A$ are nondecreasing, $\{\tau_k<t\}=\bigcup_{q\in\Q,\,q<t}\{m_q\ge k\text{ or }A_q\ge k\}\in\F_t$ for $0<t\le T$, so $\tau_k$ is an optional time, hence a stopping time, of the right-continuous filtration $\{\F_t\}$. We collect its properties.
\begin{enumerate}[label=\textup{(\alph*)}]
\item $\tau_k\le S_k\wedge T$, because $S_k\in D_k$ on $\{S_k<\infty\}$; in particular, almost surely $X_s\in\R$ for $0\le s\le\tau_k$.
\item $A_{\tau_k}\le k$: this is trivial if $\tau_k=0$, and otherwise $A_t<k$ for $t<\tau_k$ and $A$ is left-continuous.
\item Almost surely, $\tau_k\in D_k$ on $\{\tau_k<T\}$. Indeed, let $t_i\downarrow\tau_k$ with $t_i\in D_k$. If $m_{t_i}\ge k$ for infinitely many $i$, then $m_{\tau_k}\ge k$ by continuity. Otherwise $m_{t_i}<k$ and $A_{t_i}\ge k$ for some $i$ and all larger ones; then $t_i<S_k\le S$, $A$ is continuous on $[0,t_i]$ by Step~1, and $A_{\tau_k}\ge k$.
\item $\{\tau_k=T\}=\{m_t<k\text{ and }A_t<k\text{ for all }t<T\}$ increases with $k$, and
  \begin{equation}\label{eq:5-5-38:union}
    \bigcup_{k=1}^\infty\{\tau_k=T\}=\{S>T\}\quad\text{modulo a $P$-null event.}
  \end{equation}
  Indeed, on $\{S>T\}$ we have $m_T<\infty$ and, by Step~1, $A_T<\infty$ a.s., so $\tau_k=T$ as soon as $k>m_T\vee A_T$; conversely, $\tau_k=T$ implies $\abs{X_T}=\lim_{t\uparrow T}\abs{X_t}\le k$, thus $S>T$.
\end{enumerate}

\emph{Step 3: change of measure.} Fix $k\ge1$ and let $H_s\defeq\bar b(X_s)\ind_{\{s\le\tau_k\}}$, a progressively measurable process with $\int_0^\infty H_s^2\,ds=A_{\tau_k}\le k$ by (b). By Novikov's condition (\bk{Corollary 3.5.13}), the process
\[
  \xi_t\defeq\exp\Bigl\{-\int_0^tH_s\,dB_s-\frac12\int_0^tH_s^2\,ds\Bigr\};\qquad0\le t<\infty,
\]
is a martingale; it is $Z(-H)$ in the notation \bk{(3.5.2)}. Let $\wt P_k(\Lambda)\defeq E[\ind_\Lambda\xi_T]$; $\Lambda\in\F_T$. Because $\xi_T>0$, the measures $P$ and $\wt P_k$ are equivalent on $\F_T$, and $\{\F_t\}_{0\le t\le T}$ satisfies the usual conditions also under $\wt P_k$. According to the Girsanov \bk{Theorem 3.5.1},
\[
  \wt B_t\defeq B_t+\int_0^tH_s\,ds=B_t+\int_0^{t\wedge\tau_k}\bar b(X_s)\,ds;\qquad0\le t\le T,
\]
is a Brownian motion $\{\wt B_t,\F_t;\,0\le t\le T\}$ on $(\Omega,\F_T,\wt P_k)$. We set $Y_t\defeq X_0+\wt B_t$; $0\le t\le T$. Since $\wt B$ is independent of $\F_0$ under $\wt P_k$ and $\wt P_k[X_0\in G]=E[\ind_{\{X_0\in G\}}E(\xi_T\,|\,\F_0)]=\mu(G)$, \emph{the process $\{Y_t,\F_t;\,0\le t\le T\}$ is, under $\wt P_k$, a Brownian motion on $[0,T]$ with initial distribution $\mu$.} Evaluating \bk{(5.4)} (with $n=k$) at $t\wedge\tau_k$ and using (a), we obtain, almost surely,
\begin{equation}\label{eq:5-5-38:XY}
  X_{t\wedge\tau_k}=X_0+\int_0^{t\wedge\tau_k}b(X_s)\,ds+B_{t\wedge\tau_k}=Y_{t\wedge\tau_k};\qquad0\le t\le T .
\end{equation}

\emph{Step 4: functionals on path space.} For $w\in C$ we define, in analogy with Step~2,
\[
  I_t(w)\defeq\int_0^tb^2(w(s))\,ds,\qquad
  \rho_k(w)\defeq T\wedge\inf\Bigl\{t\ge0;\ \max_{0\le s\le t}\abs{w(s)}\ge k\ \text{or}\ I_t(w)\ge k\Bigr\},
\]
and $h(s,w)\defeq b(w(s))\ind_{\{s\le\rho_k(w)\}}$. Exactly as in Step~2, $I_t$ is $\B_t$-measurable, $\{\rho_k<s\}\in\B_s$ for $0<s\le T$, the event $\{\rho_k=T\}=\{\max_{s\le t}\abs{w(s)}<k,\ I_t(w)<k;\ \forall\,t<T\}$ increases with $k$, and $\int_0^\infty h^2(s,w)\,ds=I_{\rho_k(w)}(w)\le k$ for \emph{every} $w\in C$. The process $h$ on $C$ is the product of the progressively measurable process $b(w(s))$ and of the left-continuous, $\{\B_s\}$-adapted process $\ind_{\{s\le\rho_k\}}$; it is therefore a progressively measurable functional in the sense of \bk{Definition 3.5.15}, it vanishes for $s>T$, and $h(s,w)=h(s,\varphi_Tw)$, $\rho_k(w)=\rho_k(\varphi_Tw)$. We claim that
\begin{equation}\label{eq:5-5-38:rho}
  \rho_k(Y_{\cdot\wedge T})=\tau_k,\qquad I_{\rho_k}(Y_{\cdot\wedge T})=A_{\tau_k},\qquad h(s,Y_{\cdot\wedge T})=H_s;\ \ 0\le s<\infty,
\end{equation}
almost surely. By \eqref{eq:5-5-38:XY} and (a), a.s.\ $X_s=Y_s\in\R$ for $s\le\tau_k$, whence $\max_{s\le t}\abs{Y_s}=m_t$ and $I_t(Y_{\cdot\wedge T})=A_t$ for $t\le\tau_k$. Therefore no $t<\tau_k$ satisfies the condition which defines $\rho_k(Y_{\cdot\wedge T})$, so $\rho_k(Y_{\cdot\wedge T})\ge\tau_k$; and on $\{\tau_k<T\}$ the time $t=\tau_k$ does satisfy it, by (c), so $\rho_k(Y_{\cdot\wedge T})\le\tau_k$. This gives the first identity in \eqref{eq:5-5-38:rho}, and the other two follow from it.

Let $P_*$ be the probability measure induced on $(C,\B(C))$ by a Brownian motion with initial distribution $\mu$ (it is determined by $\mu$, because the finite-dimensional distributions are). If $\{U_t,\cH_t;\,0\le t\le T\}$ is a Brownian motion on $[0,T]$ with initial distribution $\mu$, then $U_{\cdot\wedge T}$ induces the measure $P_*\varphi_T^{-1}$, which agrees with $P_*$ on $\B_T$.

\begin{lem}
There exists a $\B_T$-measurable function $F\colon C\to\R$ with the following property: if $U=\{U_t,\cH_t;\,0\le t\le T\}$ is a Brownian motion on $[0,T]$ with initial distribution $\mu$, defined on a probability space $(\Theta',\cH,Q')$ with a filtration satisfying the usual conditions, then
\[
  \int_0^Th(s,U_{\cdot\wedge T})\,dU_s=F(U_{\cdot\wedge T}),\qquad\text{a.s. }Q' .
\]
\end{lem}
\begin{proof}
Let $\norm{g}^2\defeq\int_C\int_0^Tg^2(s,w)\,ds\,P_*(dw)$ for measurable $g\colon[0,T]\times C\to\R$; thus $\norm{h}^2\le k$. The truncations $h\ind_{\{\abs{h}\le N\}}$ are bounded, progressively measurable for $\{\B_t\}$, and converge to $h$ in this norm as $N\to\infty$, by dominated convergence. By \bk{Lemma 3.2.4}, applied on $(C,\B(C),P_*)$ with the filtration $\{\B_t\}$ (for a progressively measurable process, parts (a) and (b) of its proof use neither completeness nor right-continuity of the filtration), each truncation is a limit in the norm $\norm{\cdot}$ of simple processes in the sense of \bk{Definition 3.2.3}, which we restrict to $(0,T]$. Hence there are
\[
  g_n(s,w)=\sum_{i=0}^{r_n-1}g_{n,i}(w)\ind_{(t^n_i,t^n_{i+1}]}(s),\qquad 0=t^n_0<\dots<t^n_{r_n}=T,
\]
with $g_{n,i}$ bounded and $\B_{t^n_i}$-measurable, such that $\norm{g_n-h}\to0$. Put
$F_n(w)\defeq\sum_{i}g_{n,i}(w)\bigl(w(t^n_{i+1})-w(t^n_i)\bigr)$, a $\B_T$-measurable function.

Let $U$ be as in the statement. The mapping $\theta\mapsto U_{\cdot\wedge T}(\theta)$ is measurable from $(\Theta',\cH_t)$ to $(C,\B_t)$ for every $t\le T$, so $g_n(\cdot,U_{\cdot\wedge T})$ is a simple process and $h(\cdot,U_{\cdot\wedge T})$ a progressively measurable process for $\{\cH_t\}$; by the definition \bk{(3.2.5)} of the integral of a simple process, $\int_0^Tg_n(s,U_{\cdot\wedge T})\,dU_s=F_n(U_{\cdot\wedge T})$. Since $U_{\cdot\wedge T}$ induces a measure which agrees with $P_*$ on $\B_T$, the It\^o isometry yields
\begin{equation}\label{eq:5-5-38:iso}
  E^{Q'}\Bigl[\Bigl(F_n(U_{\cdot\wedge T})-\int_0^Th(s,U_{\cdot\wedge T})\,dU_s\Bigr)^2\Bigr]
  =E^{Q'}\int_0^T\bigl(g_n-h\bigr)^2(s,U_{\cdot\wedge T})\,ds=\norm{g_n-h}^2 .
\end{equation}
Applied to one particular $U$ (for instance, to the Brownian motion $W$ of the statement, restricted to $[0,T]$), \eqref{eq:5-5-38:iso} shows that $\{F_n(U_{\cdot\wedge T})\}$ is a Cauchy sequence in $L^2(Q')$, i.e., that $\{F_n\}$ is a Cauchy sequence in $L^2(C,\B_T,P_*)$. Let $F$ be its limit in this space; $F$ does not depend on $U$. For an arbitrary $U$ as in the statement, $F_n(U_{\cdot\wedge T})\to F(U_{\cdot\wedge T})$ in $L^2(Q')$, and \eqref{eq:5-5-38:iso} identifies the limit with $\int_0^Th(s,U_{\cdot\wedge T})\,dU_s$.
\end{proof}

With $F$ as in the Lemma we introduce the nonnegative, $\B_T$-measurable function
\[
  G_k(w)\defeq\exp\Bigl\{F(w)-\frac12I_{\rho_k(w)}(w)\Bigr\}\cdot\ind_{\{\rho_k(w)=T\}};\qquad w\in C .
\]

\emph{Step 5: the left-hand side.} We show that
\begin{equation}\label{eq:5-5-38:left}
  P\bigl[X_{\cdot\wedge T}\in\Gamma,\ \tau_k=T\bigr]=\int_CG_k\ind_\Gamma\,dP_*;\qquad\Gamma\in\B_T .
\end{equation}
By \bk{Problem 3.5.6} (applied with the process $-H$ in the role of the process $X$ of \bk{Theorem 3.5.1}, and with $H$ as integrand), the It\^o integral of $H$ with respect to the $\wt P_k$-Brownian motion $\wt B$ satisfies
\[
  \int_0^TH_s\,d\wt B_s=\int_0^TH_s\,dB_s+\int_0^TH_s^2\,ds=\int_0^TH_s\,dB_s+A_{\tau_k},\qquad\text{a.s. }P\text{ and }\wt P_k .
\]
On the other hand, $H_s=h(s,Y_{\cdot\wedge T})$ by \eqref{eq:5-5-38:rho}, $dY=d\wt B$, and the Lemma applies to $U=Y$ under $\wt P_k$; therefore $\int_0^TH_s\,d\wt B_s=F(Y_{\cdot\wedge T})$ a.s. Using \eqref{eq:5-5-38:rho} once more, we conclude that
\[
  \frac1{\xi_T}=\exp\Bigl\{\int_0^TH_s\,dB_s+\frac12A_{\tau_k}\Bigr\}
  =\exp\Bigl\{F(Y_{\cdot\wedge T})-\frac12I_{\rho_k}(Y_{\cdot\wedge T})\Bigr\},
  \qquad
  \frac{\ind_{\{\tau_k=T\}}}{\xi_T}=G_k(Y_{\cdot\wedge T}),
\]
almost surely. Moreover, on $\{\tau_k=T\}$ we have $X_{\cdot\wedge T}=Y_{\cdot\wedge T}$ by \eqref{eq:5-5-38:XY}, and $G_k(Y_{\cdot\wedge T})=0$ off this event. Consequently,
\[
  P\bigl[X_{\cdot\wedge T}\in\Gamma,\ \tau_k=T\bigr]
  =\wt E_k\Bigl[\frac1{\xi_T}\ind_{\{\tau_k=T\}}\ind_\Gamma(X_{\cdot\wedge T})\Bigr]
  =\wt E_k\bigl[G_k(Y_{\cdot\wedge T})\ind_\Gamma(Y_{\cdot\wedge T})\bigr]=\int_CG_k\ind_\Gamma\,dP_* ,
\]
because $G_k\ind_\Gamma$ is $\B_T$-measurable and, under $\wt P_k$, the process $Y$ is a Brownian motion on $[0,T]$ with initial distribution $\mu$. This is \eqref{eq:5-5-38:left}.

\emph{Step 6: the right-hand side.} Let $W=\{W_t,\G_t;\,0\le t<\infty\}$ be as in the statement. The random variable $W_0$ is $\G_0$-measurable and $\beta\defeq W-W_0$ is a standard Brownian motion independent of $\G_0$; hence $(W_0,\beta)$ induces the measure $\mu\times P^0_*$ on $\R\times C$, where $P^0_*$ is Wiener measure. For every $x\in\R$ the function $b^2(x+\cdot)$ is locally integrable, so the Engelbert--Schmidt zero-one law (\bk{Proposition 3.6.27}, implication (iii)$\Rightarrow$(ii)) and Fubini's theorem give
\begin{equation}\label{eq:5-5-38:ES}
  Q\bigl[I_t(W)<\infty;\ \forall\,0\le t<\infty\bigr]
  =\int_\R P^0_*\Bigl[\int_0^tb^2(x+w(s))\,ds<\infty;\ \forall\,0\le t<\infty\Bigr]\mu(dx)=1
\end{equation}
(the set $\{(x,w);\ \int_0^nb^2(x+w(s))\,ds<\infty,\ \forall\,n\ge1\}$ is product-measurable by Tonelli's theorem). Thus $\{b(W_s)\}$ is a progressively measurable process in the class $\cP^*$ of \bk{Definition 3.2.22}, the stochastic integral $\int_0^tb(W_s)\,dW_s$ and the process $Z$ of (iii) are well defined, and
\begin{equation}\label{eq:5-5-38:Wk}
  \{\rho_k(W)=T\}\uparrow\Theta\quad\text{as }k\to\infty,\ \text{modulo a $Q$-null event,}
\end{equation}
since $\rho_k(W)=T$ whenever $k>\max_{s\le T}\abs{W_s}\vee I_T(W)$. Now $\rho\defeq\rho_k(W)$ is a stopping time of $\{\G_t\}$, and by the property \bk{(3.2.24)} of the stochastic integral (which remains valid for integrands in $\cP^*$) and the Lemma (with $U=W$ on $[0,T]$),
\[
  \int_0^{T\wedge\rho}b(W_s)\,dW_s=\int_0^Tb(W_s)\ind_{\{s\le\rho\}}\,dW_s=\int_0^Th(s,W)\,dW_s=F(W),\qquad\text{a.s. }Q .
\]
On $\{\rho=T\}$ the left-hand side is $\int_0^Tb(W_s)\,dW_s$ and $I_\rho(W)=I_T(W)$; hence $Z_T\ind_{\{\rho_k(W)=T\}}=G_k(W)$ a.s., and
\begin{equation}\label{eq:5-5-38:right}
  E^Q\bigl[Z_T\ind_{\{W_\cdot\in\Gamma\}}\ind_{\{\rho_k(W)=T\}}\bigr]=E^Q\bigl[G_k(W)\ind_\Gamma(W)\bigr]=\int_CG_k\ind_\Gamma\,dP_* .
\end{equation}

\emph{Step 7: conclusion.} From \eqref{eq:5-5-38:left} and \eqref{eq:5-5-38:right},
\[
  P\bigl[X_{\cdot\wedge T}\in\Gamma,\ \tau_k=T\bigr]=E^Q\bigl[Z_T\ind_{\{W_\cdot\in\Gamma\}}\ind_{\{\rho_k(W)=T\}}\bigr];\qquad k\ge1.
\]
We let $k\to\infty$. The left-hand side increases to $P[X_{\cdot\wedge T}\in\Gamma,\ S>T]$ by \eqref{eq:5-5-38:union}, and the right-hand side increases to $E^Q[Z_T\ind_{\{W_\cdot\in\Gamma\}}]$ by \eqref{eq:5-5-38:Wk} and the monotone convergence theorem. This is the generalized Girsanov formula.

\rmk{The Engelbert--Schmidt zero-one law is what makes the right-hand side meaningful, and what makes the localization exhaust the whole space on the Brownian side, \eqref{eq:5-5-38:Wk}. On the side of the solution, the corresponding fact is Step~1, and there the localizing events exhaust only $\{S>T\}$; this is the source of the possible defect of mass in the formula. For bounded $b$ there is no explosion and the formula reduces to \bk{(3.11)}.}

\textbf{(iii)} \emph{The martingale property.} By \eqref{eq:5-5-38:ES}, $\{b(W_s)\}$ is an integrand in $\cP^*$ for the standard Brownian motion $W-W_0$, and by \bk{Problem 3.2.28} the process $\{Z_t,\G_t;\,0\le t<\infty\}$ is a (nonnegative) supermartingale, with $Z_0=1$. The formula of (ii) with $\Gamma=C[0,\infty)$ reads
\begin{equation}\label{eq:5-5-38:mass}
  E^Q[Z_T]=P[S>T];\qquad0<T<\infty .
\end{equation}
If $Z$ is a martingale, then $P[S>T]=E^Q[Z_T]=E^Q[Z_0]=1$ for every $T$, and $P[S=\infty]=\lim_{T\to\infty}P[S>T]=1$. Conversely, if $P[S=\infty]=1$, then $E^Q[Z_t]=1$ for all $t\ge0$; for $0\le s<t$ the random variable $Z_s-E^Q[Z_t\,|\,\G_s]$ is nonnegative a.s.\ by the supermartingale property and has expectation $1-1=0$, so it vanishes a.s., and $Z$ is a martingale. (By (i), or by \eqref{eq:5-5-38:mass} itself, the number $P[S=\infty]$ is the same for all solutions with initial distribution $\mu$.)

\emph{The criterion.} We now show that $P[S=\infty]=1$ if and only if $\varphi(\infty)=\varphi(-\infty)=\infty$. Note that the integrand of $\varphi$ is positive, and that for $x<0$ the two outer integrals are both taken in the negative direction, so $\varphi\ge0$, and $\varphi$ is monotone on each of $(-\infty,0]$ and $[0,\infty)$; the limits $\varphi(\pm\infty)\in[0,\infty]$ exist. With $\sigma\equiv1$ and $I=\R$, the conditions \bk{(ND)$'$}, \bk{(LI)$'$} hold, and the scale function \bk{(5.42)} and the function \bk{(5.65)} with base point $c\in\R$ are
\begin{align*}
  p_c'(y)&=\exp\Bigl\{-2\int_c^yb(u)\,du\Bigr\},\\
  v_c(x)&=\int_c^xp_c'(y)\int_c^y\frac{2\,dz}{p_c'(z)}\,dy=2\int_c^x\int_c^y\exp\Bigl\{-2\int_z^yb(u)\,du\Bigr\}\,dz\,dy ,
\end{align*}
so that $v_0=2\varphi$. For nonrandom initial conditions the assertion is thus exactly Feller's test, \bk{Theorem 5.29}. Since that theorem is stated only for $X_0=x$, we repeat its argument for a general initial distribution, localizing on events of $\F_0$. For $c\in\R$ let $u=u_c$ be the function of \bk{Lemma 5.26}: $u$ has an absolutely continuous derivative, $u(c)=1$, $u$ is strictly increasing on $(c,\infty)$ and strictly decreasing on $(-\infty,c)$, $\frac12u''+bu'=u$ a.e., and $1+v_c\le u\le e^{v_c}$. We extend $u$ to $[-\infty,\infty]$ by its limits $u(\pm\infty)\in(1,\infty]$.

Let $(X,B)$ be a solution as in (ii), and fix $n\ge1$. We apply the generalized It\^o rule of \bk{Problem 3.7.3} to $u$ and the semimartingale $X^{(n)}=X_0+M+V$ of Step~1, taking for $u''$ the function $2(u-bu')$ (the choice of the version is immaterial by \bk{(3.7.3)}); as $dM_s=\ind_{\{s\le S_n\}}dB_s$, $dV_s=b(X_s)\ind_{\{s\le S_n\}}ds$ and $d\qv{M}_s=\ind_{\{s\le S_n\}}ds$, we get
\[
  u(X_{t\wedge S_n})=u(X_0)+\int_0^{t\wedge S_n}u'(X_s)\,dB_s+\int_0^{t\wedge S_n}u(X_s)\,ds;\qquad0\le t<\infty,
\]
and then, integrating by parts with the continuous function $e^{-(t\wedge S_n)}$ of bounded variation,
\begin{equation}\label{eq:5-5-38:K}
  K^{(n)}_t\defeq e^{-(t\wedge S_n)}u(X_{t\wedge S_n})=u(X_0)+\int_0^te^{-s}u'(X_s)\ind_{\{s\le S_n\}}\,dB_s;\qquad0\le t<\infty,
\end{equation}
almost surely. The last two displays are the two displays in the proof of \bk{Theorem 5.29} (the representations of $Z^{(n)}$ and of $M^{(n)}$), with $\sigma\equiv1$; there $\tau_n=n$, and this additional stopping is not needed here. As $n\to\infty$,
\begin{equation}\label{eq:5-5-38:Klim}
  K^{(n)}_t\longrightarrow e^{-(t\wedge S)}u(X_{t\wedge S})\in[0,\infty],
\end{equation}
because $X$ is continuous in $[-\infty,\infty]$ and so is $u$.

\emph{Sufficiency.} Assume $\varphi(\pm\infty)=\infty$ and take $c=0$; then $u(\pm\infty)\ge1+2\varphi(\pm\infty)=\infty$. Fix integers $n\ge j\ge1$, a number $t>0$, and put $G_j\defeq\{\abs{X_0}\le j\}\in\F_0$. On $G_j$ we have $\abs{X_s}\le n$ for $s\le S_n$, so the process $\ind_{G_j}e^{-s}u'(X_s)\ind_{\{s\le S_n\}}$ is progressively measurable and bounded by $\max_{[-n,n]}\abs{u'}$; its stochastic integral is a martingale, and equals $\ind_{G_j}$ times the integral in \eqref{eq:5-5-38:K} (\bk{Exercise 3.2.30}). Taking expectations in \eqref{eq:5-5-38:K} multiplied by $\ind_{G_j}$, and then using \eqref{eq:5-5-38:Klim} and Fatou's lemma, we obtain
\[
  E\bigl[\ind_{G_j}e^{-(t\wedge S)}u(X_{t\wedge S})\bigr]\le\liminf_{n\to\infty}E\bigl[\ind_{G_j}K^{(n)}_t\bigr]=E\bigl[\ind_{G_j}u(X_0)\bigr]\le\max_{[-j,j]}u<\infty .
\]
But on $\{S\le t\}$ we have $X_{t\wedge S}=X_S=\pm\infty$ a.s., by \bk{(5.6)}, and there the integrand on the left equals $\ind_{G_j}e^{-S}\cdot\infty$. Hence $P[G_j\cap\{S\le t\}]=0$; letting $j\to\infty$ and $t\to\infty$, and recalling $\abs{X_0}<\infty$ a.s., we get $P[S<\infty]=0$.

\emph{Necessity.} Assume $\varphi(\infty)<\infty$, i.e., $v_0(\infty)<\infty$. Choose real numbers $c<j$ with $\mu((c,j])>0$. By \bk{Problem 5.28}, $v_c(\infty)<\infty$ as well, so $u=u_c$ satisfies $u(\infty)\le e^{v_c(\infty)}<\infty$, and $1\le u\le u(\infty)$ on $[c,\infty]$. Let $G\defeq\{c<X_0\le j\}\in\F_0$, so that $P(G)>0$, and let $T_c\defeq\inf\{t\ge0;\,X_t=c\}$, a stopping time. For $n\ge j\vee\abs{c}$ we have $X_s\in[c,n]$ on $G$ for $s\le S_n\wedge T_c$; stopping \eqref{eq:5-5-38:K} at $T_c$ (by \bk{(3.2.24)}) and multiplying by $\ind_G$, we see as before that the integrand $\ind_Ge^{-s}u'(X_s)\ind_{\{s\le S_n\wedge T_c\}}$ is bounded, and thus
\[
  E\bigl[\ind_Gu(X_0)\bigr]=E\bigl[\ind_Ge^{-(t\wedge S_n\wedge T_c)}u(X_{t\wedge S_n\wedge T_c})\bigr];\qquad t\ge0 .
\]
The integrand on the right is bounded by $u(\infty)$, and $X_{t\wedge S\wedge T_c}\in[c,\infty]$ on $G$; by \eqref{eq:5-5-38:Klim} (applied at the time $t\wedge T_c$) and bounded convergence,
\[
  E\bigl[\ind_Gu(X_0)\bigr]=E\bigl[\ind_Ge^{-(t\wedge S\wedge T_c)}u(X_{t\wedge S\wedge T_c})\bigr];\qquad t\ge0 .
\]
Suppose now that $P[S=\infty]=1$. Then the right-hand side is $E[\ind_Ge^{-(t\wedge T_c)}u(X_{t\wedge T_c})]$. As $t\to\infty$, the integrand converges to $\ind_Ge^{-T_c}u(c)=\ind_Ge^{-T_c}$ on $\{T_c<\infty\}$, and to zero on $\{T_c=\infty\}$, where it is bounded by $e^{-t}u(\infty)$. By bounded convergence,
\[
  E\bigl[\ind_Gu(X_0)\bigr]=E\bigl[\ind_Ge^{-T_c}\ind_{\{T_c<\infty\}}\bigr]\le P(G) .
\]
This is impossible, because $u(X_0)>u(c)=1$ on $G$ and $P(G)>0$. Therefore $P[S=\infty]<1$. If $\varphi(-\infty)<\infty$, the same argument applies with $G=\{-j\le X_0<c\}$, where $c>-j$ are such that $\mu([-j,c))>0$, and with $u_c$ bounded and strictly decreasing on $[-\infty,c]$.

We have shown that $P[S=\infty]=1$ if and only if $\varphi(\infty)=\varphi(-\infty)=\infty$, for every initial distribution $\mu$; together with the first part of (iii), this completes the proof.
\end{kex}

\begin{kex}{Exercise}{5.39}{}
The setting is that of the paragraph containing \bk{(5.58)}, \bk{(5.59)}: $I=(\ell,r)$, the Borel-measurable coefficients $b,\sigma\colon I\to\R$ satisfy \bk{(ND)$'$} and \bk{(LI)$'$}, $(X,W)$, $(\Omega,\F,P)$, $\{\F_t\}$ is a weak solution in $I$ of \bk{(5.1)} (\bk{Definition 5.20}) with $X_0=x$, and $\ell<a<x<b<r$. We write $T\defeq T_{a,b}=\inf\{t\ge 0;\ X_t\notin(a,b)\}$ and $D\defeq p(b)-p(a)>0$. Of $h$ we shall use only that it is nonnegative, bounded and Borel-measurable on $[a,b]$; put $\norm{h}\defeq\sup_{[a,b]}h$.

\emph{Step 1: the function $H$.} By \bk{(LI)$'$} and the compactness of $[a,b]\subset I$ we have $\int_a^b(1+\abs{b(y)})\sigma^{-2}(y)\,dy<\infty$. The derivative $p'$ of the scale function is continuous and strictly positive, hence bounded away from zero on $[a,b]$; thus the speed measure \bk{(5.51)} satisfies $m([a,b])<\infty$. Since $0\le G_{a,b}\le D$ on $[a,b]^2$, the function
\[
  H(\xi)\defeq\int_a^b G_{a,b}(\xi,y)h(y)\,m(dy);\qquad a\le\xi\le b,
\]
is well defined, nonnegative and bounded, and $H(a)=H(b)=0$ because $G_{a,b}(a,\cdot)=G_{a,b}(b,\cdot)=0$ by \bk{(5.52)}. For $a\le y\le\xi\le b$ the identity
\[
  \bigl(p(y)-p(a)\bigr)\bigl(p(b)-p(\xi)\bigr)
  =\bigl(p(\xi)-p(a)\bigr)\bigl(p(b)-p(y)\bigr)-\bigl(p(\xi)-p(y)\bigr)\bigl(p(b)-p(a)\bigr)
\]
holds (expand both sides), whereas for $\xi\le y$ we have $G_{a,b}(\xi,y)=(p(\xi)-p(a))(p(b)-p(y))/D$. Consequently, exactly as in \bk{(5.55)},
\begin{equation}\label{eq:5-5-39:a}
  H(\xi)=-\int_a^\xi\bigl(p(\xi)-p(y)\bigr)h(y)\,m(dy)+\frac{p(\xi)-p(a)}{D}\,C,
  \qquad C\defeq\int_a^b\bigl(p(b)-p(y)\bigr)h(y)\,m(dy).
\end{equation}
Let $\varphi(z)\defeq\int_a^z h(y)\,m(dy)=\int_a^z 2h(y)\bigl(p'(y)\sigma^2(y)\bigr)^{-1}dy$, an absolutely continuous function on $[a,b]$. Writing $p(\xi)-p(y)=\int_y^\xi p'(z)\,dz$ and using Tonelli's theorem, we obtain $\int_a^\xi(p(\xi)-p(y))h(y)\,m(dy)=\int_a^\xi p'(z)\varphi(z)\,dz$; the integrand being continuous, \eqref{eq:5-5-39:a} shows that $H$ is continuously differentiable on $[a,b]$ (one-sided derivatives at the endpoints) with
\begin{equation}\label{eq:5-5-39:b}
  H'(\xi)=p'(\xi)\Bigl(\frac{C}{D}-\varphi(\xi)\Bigr);\qquad a\le\xi\le b.
\end{equation}
Now $p'=\exp\{-2\int_c^{\,\cdot}b(\zeta)\sigma^{-2}(\zeta)\,d\zeta\}$ is absolutely continuous on $[a,b]$, with derivative $p''=-2b\sigma^{-2}p'$ Lebesgue-almost everywhere (\bk{(5.43)}). Being a product of two absolutely continuous functions on a compact interval, $H'$ is absolutely continuous on $[a,b]$, and almost everywhere in $(a,b)$ its derivative is
\[
  p''(\xi)\Bigl(\frac{C}{D}-\varphi(\xi)\Bigr)-p'(\xi)\,\frac{2h(\xi)}{p'(\xi)\sigma^2(\xi)}
  =-\frac{2}{\sigma^2(\xi)}\bigl(b(\xi)H'(\xi)+h(\xi)\bigr).
\]
We therefore \emph{define}, at every point of $(a,b)$,
\begin{equation}\label{eq:5-5-39:c}
  H''(\xi)\defeq-\frac{2}{\sigma^2(\xi)}\bigl(b(\xi)H'(\xi)+h(\xi)\bigr),\quad\text{so that}\quad
  b(\xi)H'(\xi)+\tfrac12\sigma^2(\xi)H''(\xi)=-h(\xi);\quad a<\xi<b.
\end{equation}
This $H''$ is a Borel-measurable, integrable function on $(a,b)$ with $H'(\xi)=H'(a)+\int_a^\xi H''(z)\,dz$. Finally we extend $H$ to a function $f\colon\R\to\R$ by $f=H$ on $[a,b]$, $f(\xi)=H'(a)(\xi-a)$ for $\xi<a$ and $f(\xi)=H'(b)(\xi-b)$ for $\xi>b$. Then $f$ is continuously differentiable, and $f'$ is absolutely continuous on every compact interval, with density $f''\defeq H''\ind_{(a,b)}$.

\emph{Step 2: the stopped process is a semimartingale.} The process $X$ is continuous and adapted with values in $[\ell,r]\subseteq[-\infty,\infty]$, so it may take the values $\pm\infty$ when $\ell=-\infty$ or $r=\infty$. Extend $\arctan$ to $[-\infty,\infty]$ by $\arctan(\pm\infty)\defeq\pm\pi/2$; this is a strictly increasing homeomorphism of $[-\infty,\infty]$ onto $[-\pi/2,\pi/2]$. Hence $Z_t\defeq\arctan X_t$ defines a real-valued process with continuous paths, adapted to $\{\F_t\}$, and $X_t\notin(a,b)$ if and only if $Z_t\in\Gamma\defeq\R\setminus(\arctan a,\arctan b)$. Thus $T$ is the hitting time of the closed set $\Gamma\subseteq\R$ by $Z$, and \bk{Problem 1.2.7} (with $d=1$) shows that $T$ is a stopping time. Choose $n$ so large that $\ell_n<a$ and $b<r_n$; then $T\le S_n$, with $S_n$ as in \bk{(5.56)}. By \bk{(5.3)},
\begin{equation}\label{eq:5-5-39:d}
  \int_0^{t\wedge T}\bigl\{\abs{b(X_s)}+\sigma^2(X_s)\bigr\}\,ds<\infty;\qquad 0\le t<\infty,\ \text{a.s.},
\end{equation}
and \bk{(5.4)}, evaluated at $t\wedge T$, together with the behaviour of stochastic integrals under stopping, gives
\[
  Y_t\defeq X_{t\wedge T}=x+V_t+M_t,\qquad
  V_t\defeq\int_0^{t\wedge T}b(X_s)\,ds,\qquad
  M_t\defeq\int_0^t\sigma(X_s)\ind_{\{s\le T\}}\,dW_s;\qquad 0\le t<\infty,
\]
almost surely. By \eqref{eq:5-5-39:d}, $M\in\cM^{c,\loc}$ with $\qv{M}_t=\int_0^{t\wedge T}\sigma^2(X_s)\,ds$, and $V$ is the difference of the continuous, nondecreasing, adapted processes $\int_0^{\cdot\wedge T}b^{\pm}(X_s)\,ds$. Thus $Y$ is a continuous semimartingale of the form \bk{(3.7.1)}; it takes values in $[a,b]$, and $Y_s=X_s\in(a,b)$ for $0\le s<T$.

\emph{Step 3: the generalized It\^o rule.} \bk{Problem 3.7.3} applies to $f$ and $Y$. The density of $f'$ appearing there may be taken to be our $f''$: if $g$ is another Borel-measurable density of $f'$, then $k\defeq\abs{f''-g}$ vanishes Lebesgue-almost everywhere, and \bk{(3.7.3)} yields $\int_0^t k(Y_s)\,d\qv{M}_s=0$ for all $t\ge0$, almost surely. Hence, almost surely, for all $0\le t<\infty$,
\begin{align*}
  H(X_{t\wedge T})=f(Y_t)
  &=H(x)+\int_0^t f'(Y_s)\,dM_s+\int_0^{t\wedge T}\Bigl\{b(X_s)H'(X_s)+\tfrac12\sigma^2(X_s)H''(X_s)\Bigr\}\,ds\\
  &=H(x)+\int_0^{t}H'(X_s)\sigma(X_s)\ind_{\{s\le T\}}\,dW_s-\int_0^{t\wedge T}h(X_s)\,ds .
\end{align*}
Here we have used that $X_s\in(a,b)$ for every $s\in[0,t\wedge T]$ with the possible exception of $s=T$, so that $f'(Y_s)=H'(X_s)$, $f''(Y_s)=H''(X_s)$ and \eqref{eq:5-5-39:c} apply for Lebesgue-almost every such $s$; the $ds$-integral of the first line converges absolutely by \eqref{eq:5-5-39:d}, since $\abs{bH'}+\tfrac12\sigma^2\abs{H''}\le 2\abs{b}\sup_{[a,b]}\abs{H'}+\norm{h}$ on $(a,b)$.

\emph{Step 4: expectations.} Let $\tau_n\defeq\inf\{t\ge0;\ \qv{M}_t\ge n\}$, a stopping time; by \eqref{eq:5-5-39:d}, $\tau_n\uparrow\infty$ almost surely as $n\to\infty$. Since $H'$ is bounded on $[a,b]$, say by $K$, the integrand $H'(X_s)\sigma(X_s)\ind_{\{s\le T\wedge\tau_n\}}$ satisfies $E\int_0^\infty\abs{H'(X_s)}^2\sigma^2(X_s)\ind_{\{s\le T\wedge\tau_n\}}\,ds\le K^2n$, so the stochastic integral of Step 3, stopped at $\tau_n$, is a square-integrable martingale starting at zero. Taking expectations in Step 3 at time $t\wedge\tau_n$ we get
\[
  E^x\int_0^{t\wedge\tau_n\wedge T}h(X_s)\,ds=H(x)-E^xH(X_{t\wedge\tau_n\wedge T}).
\]
We let $n\to\infty$: the left-hand side converges by the monotone convergence theorem ($h\ge0$), the right-hand side by the bounded convergence theorem ($H$ is bounded and continuous on $[a,b]$, and $X$ is continuous). Thus
\begin{equation}\label{eq:5-5-39:e}
  E^x\int_0^{t\wedge T}h(X_s)\,ds=H(x)-E^xH(X_{t\wedge T});\qquad 0\le t<\infty.
\end{equation}
For $h\equiv1$ this is \bk{(5.58)}, and the book deduces from it that $E^xT<\infty$; in particular $T<\infty$ almost surely, and then $X_T\in\{a,b\}$ by the continuity of $X$ and $X_0=x\in(a,b)$. Letting $t\to\infty$ in \eqref{eq:5-5-39:e}, we have $H(X_{t\wedge T})\to H(X_T)=0$ almost surely, boundedly, so that $E^xH(X_{t\wedge T})\to0$ by bounded convergence, while the left-hand side increases to $E^x\int_0^Th(X_s)\,ds$ by monotone convergence. Therefore
\[
  E^x\int_0^{T_{a,b}}h(X_t)\,dt=H(x)=\int_a^bG_{a,b}(x,y)h(y)\,m(dy).
\]

\rmk{(a) The relation ``$(a,b)\subsetneqq I$'' of the statement has to be read as $\ell<a<b<r$, as in the text preceding \bk{(5.58)}: the Green's function \bk{(5.52)} is defined only for $[a,b]\subseteq I$, and $p(a)$, $p(b)$ must be finite. (b) Piecewise continuity of $h$ was not used; the formula holds for every bounded, Borel-measurable $h\colon I\to[0,\infty)$, and then, by monotone convergence applied to $h\wedge n$, for every Borel-measurable $h\colon I\to[0,\infty]$, both sides being possibly infinite. (c) For $h\equiv1$ the formula is \bk{(5.59)}. For Brownian motion ($p(y)=y$, $m(dy)=2\,dy$) and $h\equiv 1$ the right-hand side is $\frac{2}{b-a}\bigl\{(b-x)\tfrac12(x-a)^2+(x-a)\tfrac12(b-x)^2\bigr\}=(x-a)(b-x)$, the familiar value of $E^xT_{a,b}$.}
\end{kex}

\begin{kex}{Exercise}{5.40}{}
Throughout, \bk{(ND)$'$} and \bk{(LI)$'$} hold, and $m$ is a finite measure on $I$, absolutely continuous with respect to Lebesgue measure by \bk{(5.51)}; in particular $m$ has no atoms. The statement speaks of ``the'' solution starting at $x$ and of its law $P^x$, and steps (i), (iv) compare the process after a stopping time with a solution started afresh. We therefore begin by making precise the family $\{P^x\}_{x\in I}$ and its strong Markov property; this is the only place where the boundedness of $b$ and $\sigma$ on compact subintervals of $I$ is used.

\emph{Step 0: the laws $P^x$ and the strong Markov property.}
Since $p$ is continuous and strictly increasing on $I$ with $p(\ell+)=-\infty$, $p(r-)=\infty$, it maps $I$ onto $\R$; let $q\colon\R\to I$ be its (continuous, strictly increasing) inverse, and set, as in \bk{(5.48)},
\[
  \wt\sigma(y)\defeq p'(q(y))\,\sigma(q(y));\qquad y\in\R .
\]
This Borel-measurable function vanishes nowhere, by \bk{(ND)$'$} and $p'>0$; it is bounded on compact subsets of $\R$, because $q$ maps a compact subset of $\R$ into a compact subinterval of $I$, on which $p'$ is continuous and $\sigma$ is bounded by assumption; and $\wt\sigma^{-2}$ is locally integrable, since the substitution $y=p(\xi)$ gives, for $\ell<\alpha<\beta<r$,
\[
  \int_{p(\alpha)}^{p(\beta)}\frac{dy}{\wt\sigma^2(y)}=\int_\alpha^\beta\frac{d\xi}{p'(\xi)\sigma^2(\xi)}=\tfrac12\,m([\alpha,\beta])<\infty .
\]
Thus the sets $I(\wt\sigma)$ and $Z(\wt\sigma)$ of Subsection A are both empty, and by \bk{Theorem 5.7} the equation
\begin{equation}\label{eq:5-5-40:Y}
  Y_t=\eta+\int_0^t\wt\sigma(Y_s)\,dW_s;\qquad 0\le t<\infty,
\end{equation}
has, for every $\eta\in\R$, a weak solution, which is unique in the sense of probability law. We denote by $Q^\eta$ the law of this solution on $(\Omega,\B)\defeq(C[0,\infty),\B(C[0,\infty)))$, by $y(\cdot)$ the coordinate process, by $\B_t\defeq\sigma(y(s);\,0\le s\le t)$ its filtration and by $(\theta_s\omega)(t)=\omega(s+t)$ the shifts. Equation \eqref{eq:5-5-40:Y} is of the form \bk{(4.20)} with $d=r=1$, drift zero and dispersion $\wt\sigma$; it is well posed (\bk{Definition 4.14}) and its coefficients are bounded on compact sets, so \bk{Theorem 4.20} provides, for every stopping time $T$ of $\{\B_t\}$, every $F\in\B$ and every $\eta\in\R$,
\begin{equation}\label{eq:5-5-40:SM}
  Q^\eta\bigl[\theta_T^{-1}F\,\big|\,\B_T\bigr](\omega)=Q^{\omega(T(\omega))}[F],\qquad Q^\eta\text{-a.e. }\omega\in\{T<\infty\}.
\end{equation}

On $\Omega$ we define the continuous, $I$-valued, $\{\B_t\}$-adapted process $X_t(\omega)\defeq q(\omega(t))$, and we set
\[
  P^x\defeq Q^{p(x)},\qquad E^x\defeq\text{expectation under }P^x;\qquad x\in I .
\]
We claim that \emph{under $P^x$, the process $X$ has the law of every weak solution in $I$ of \bk{(5.1)} with initial condition $x$, and that such solutions exist.} Indeed, let $(X',W')$ be a weak solution in $I$ with $X'_0=x$ (\bk{Definition 5.20}). By \bk{Proposition 5.22}(a) its exit time is almost surely infinite, so after modification on a null set $X'$ is $I$-valued. The generalized It\^o rule gives \bk{(5.49)} for $Y'\defeq p(X')$, with $S_n$ as in \bk{(5.56)} (this is how \bk{(5.49)} is used in the proof of \bk{Proposition 5.22}); since $S_n\uparrow\infty$ almost surely, condition \bk{(5.3)} in \bk{Definition 5.20} yields, almost surely, $\int_0^t\sigma^2(X'_s)\,ds<\infty$ for every $t\ge0$. As $p'$ is continuous on $I$ and $\{X'_s;\,0\le s\le t\}$ is a compact subset of $I$, it follows that, almost surely,
\[
  \int_0^t\wt\sigma^2(Y'_s)\,ds=\int_0^t\bigl(p'(X'_s)\bigr)^2\sigma^2(X'_s)\,ds\le\max_{0\le s\le t}\bigl(p'(X'_s)\bigr)^2\int_0^t\sigma^2(X'_s)\,ds<\infty;\qquad t\ge0 ,
\]
so that the continuous local martingale $M_t\defeq\int_0^t\wt\sigma(Y'_s)\,dW'_s$ is defined for all $t\ge0$, and the stochastic integral in \bk{(5.49)} equals $M_{t\wedge S_n}$. Since $S_n\uparrow\infty$, for almost every $\omega$ and every $t$ we have $t\wedge S_n=t$ for all large $n$; hence letting $n\to\infty$ in \bk{(5.49)} we find that $Y'$ solves \eqref{eq:5-5-40:Y} with $\eta=p(x)$. Hence $Y'$ has law $Q^{p(x)}$, and $X'=q(Y')$ has the law of $X=q(y(\cdot))$ under $P^x$. Conversely, let $(Y',W')$ be a weak solution of \eqref{eq:5-5-40:Y} with $\eta=p(x)$, and put $X'\defeq q(Y')$, a continuous, adapted, $I$-valued process with $X'_0=x$. Choose strictly monotone sequences $\{\ell_n\}$, $\{r_n\}$ with $\ell<\ell_n<r_n<r$, $\ell_n\downarrow\ell$, $r_n\uparrow r$, and let $S_n$ be defined by \bk{(5.56)} for $X'$. We verify the two requirements of \bk{Definition 5.20}(iii)$'$. As for \bk{(5.3)}, fix $n\ge1$ and $t\ge0$. If $S_n=0$, the integral $\int_0^{t\wedge S_n}\{\abs{b(X'_s)}+\sigma^2(X'_s)\}\,ds$ vanishes. If $S_n>0$, then $X'_s\in(\ell_n,r_n)$ for $0\le s<S_n$ and, by continuity, $X'_s\in[\ell_n,r_n]$ for $0\le s\le t\wedge S_n$; since $\abs{b}$ and $\abs{\sigma}$ are bounded on the compact subinterval $[\ell_n,r_n]$ of $I$ by some constant $C_n$, the integral is at most $(C_n+C_n^2)\,t<\infty$. Thus \bk{(5.3)} holds, even for every $\omega$. As for \bk{(5.4)}, the computation which concludes the proof of \bk{Proposition 5.13} (the generalized It\^o rule for $q$, \bk{(5.44)} and \bk{(5.48)}, stopped at $S_n$) gives, almost surely, for every $n\ge1$,
\[
  X'_{t\wedge S_n}=x+\int_0^{t\wedge S_n}b(X'_s)\,ds+\int_0^{t\wedge S_n}\sigma(X'_s)\,dW'_s;\qquad 0\le t<\infty,
\]
where the Lebesgue integral converges absolutely by what we have just shown; this is \bk{(5.4)}. Hence $X'$ is a weak solution in $I$ of \bk{(5.1)} with $X'_0=x$, and the claim is proved. Consequently:

\begin{itemize}
\item[(A)] \emph{Every identity proved in the book or in \bk{Exercise 5.39} for the expectation of a nonnegative Borel functional of the path of a weak solution in $I$ starting at $x$ holds for $E^x$ and the process $X$ on $\Omega$.} In particular \bk{Exercise 5.39} (which contains \bk{(5.59)} as the case $h\equiv1$) reads: for $\ell<a<x<b<r$, $T_{a,b}\defeq\inf\{t\ge0;\ X_t\notin(a,b)\}$ and every bounded, piecewise continuous $h\colon I\to[0,\infty)$,
\begin{equation}\label{eq:5-5-40:G}
  E^x\int_0^{T_{a,b}}h(X_t)\,dt=\int_a^bG_{a,b}(x,y)h(y)\,m(dy).
\end{equation}
\item[(B)] \emph{Let $T$ be a stopping time of $\{\B_t\}$ and $x,c\in I$ with $P^x[T<\infty,\ X_T=c]=1$, and let $g\colon[0,\infty)\times\Omega\to[0,\infty]$ be $\B([0,\infty))\otimes\B$-measurable. Then $s\mapsto E^c g(s,\cdot)$ is Borel-measurable and}
\begin{equation}\label{eq:5-5-40:B}
  E^x\bigl[g(T,\theta_T\,\cdot)\bigr]=\int_{[0,\infty)}E^c\bigl[g(s,\cdot)\bigr]\,P^x[T\in ds].
\end{equation}
\end{itemize}
To see (B), take first $g(s,\omega)=\ind_B(s)\ind_F(\omega)$ with $B\in\B([0,\infty))$, $F\in\B$. Since $\{T\in B\}\in\B_T$ and $\omega(T(\omega))=p(c)$ for $P^x$-a.e.\ $\omega$, \eqref{eq:5-5-40:SM} yields $E^x[\ind_{\{T\in B\}}\ind_F\circ\theta_T]=P^x[T\in B]\,P^c[F]$, which is \eqref{eq:5-5-40:B}. The measurable rectangles form a $\pi$-system generating $\B([0,\infty))\otimes\B$, and the bounded measurable $g$ for which the measurability assertion and \eqref{eq:5-5-40:B} hold form a vector space containing the constants and closed under bounded monotone limits; by the monotone class theorem it contains all bounded measurable $g$, and nonnegative $g$ follow by monotone convergence. (The map $\omega\mapsto(T(\omega),\theta_{T(\omega)}\omega)$ is measurable on $\{T<\infty\}$ because $(s,\omega)\mapsto\theta_s\omega$ is continuous.)

For $a\in I$ the time $T_a=\inf\{t\ge0;\ X_t=a\}=\inf\{t\ge0;\ y(t)=p(a)\}$ is a stopping time of $\{\B_t\}$ by \bk{Problem 1.2.7}. So is $S_{a;b}$ for $a\neq b$: writing $D_t\defeq(\Q\cap[0,t])\cup\{t\}$, we have
\[
  \{S_{a;b}\le t\}=\bigcap_{n\ge1}\ \bigcup_{\substack{q_1,q_2\in D_t\\ q_1\le q_2}}
  \bigl\{\abs{X_{q_1}-a}<\tfrac1n,\ \abs{X_{q_2}-b}<\tfrac1n\bigr\}\in\B_t .
\]
Indeed, if $S_{a;b}\le t$, then $s_1\defeq T_a<s_2\defeq S_{a;b}\le t$ with $X_{s_1}=a$, $X_{s_2}=b$, and continuity provides $q_1\le q_2$ in $D_t$ as close to $s_1,s_2$ as we wish; conversely, if $\omega$ belongs to the right-hand side, a subsequence of the corresponding pairs $(q_1^{(n)},q_2^{(n)})$ converges to some $s_1\le s_2\le t$ with $X_{s_1}=a$, $X_{s_2}=b$, whence $T_a\le s_1$ and $S_{a;b}\le s_2\le t$. By continuity, $X_{T_a}=a$ on $\{T_a<\infty\}$ and $X_{S_{a;b}}=b$ on $\{S_{a;b}<\infty\}$, and
\begin{equation}\label{eq:5-5-40:S}
  S_{a;b}=T_a+T_b\circ\theta_{T_a}\quad\text{on }\{T_a<\infty\}.
\end{equation}

\emph{Step 1: occupation integrals up to one-sided hitting times.} Let $h\colon I\to[0,\infty)$ be bounded and piecewise continuous. We show that
\begin{align}
  E^x\int_0^{T_a}h(X_t)\,dt&=\int_{(a,r)}\bigl(p(x\wedge y)-p(a)\bigr)h(y)\,m(dy);&&\ell<a<x<r,\label{eq:5-5-40:Ta}\\
  E^x\int_0^{T_b}h(X_t)\,dt&=\int_{(\ell,b)}\bigl(p(b)-p(x\vee y)\bigr)h(y)\,m(dy);&&\ell<x<b<r.\label{eq:5-5-40:Tb}
\end{align}
Fix $a<x$ and a sequence $x<b_n\uparrow r$. For every $\omega$ and $t<\infty$ the set $\{X_s(\omega);\,0\le s\le t\}$ is a compact subset of $I$, so $T_{b_n}(\omega)>t$ for all large $n$; thus $T_{b_n}\uparrow\infty$, and $T_{a,b_n}=T_a\wedge T_{b_n}\uparrow T_a$. By monotone convergence the left-hand side of \eqref{eq:5-5-40:G}, with $b=b_n$, increases to the left-hand side of \eqref{eq:5-5-40:Ta}. On the right-hand side, \bk{(5.52)} gives for $a<y<b_n$
\[
  G_{a,b_n}(x,y)=\bigl(p(x\wedge y)-p(a)\bigr)\,\frac{p(b_n)-p(x\vee y)}{p(b_n)-p(a)},
\]
where the last fraction lies in $[0,1]$ and tends to $1$ as $n\to\infty$, because $p(b_n)\to p(r-)=\infty$. Hence $G_{a,b_n}(x,y)\ind_{(a,b_n)}(y)\to(p(x\wedge y)-p(a))\ind_{(a,r)}(y)$, the convergence being dominated by the constant $p(x)-p(a)$, which is $m$-integrable since $m(I)<\infty$ and $h$ is bounded. Dominated convergence gives \eqref{eq:5-5-40:Ta}. The proof of \eqref{eq:5-5-40:Tb} is symmetric: with $x>a_n\downarrow\ell$ we have $T_{a_n}\uparrow\infty$, $T_{a_n,b}\uparrow T_b$, and
\[
  G_{a_n,b}(x,y)=\bigl(p(b)-p(x\vee y)\bigr)\,\frac{p(x\wedge y)-p(a_n)}{p(b)-p(a_n)}\longrightarrow p(b)-p(x\vee y)\qquad(a_n<y<b),
\]
dominated by $p(b)-p(x)$, because $p(a_n)\to p(\ell+)=-\infty$.

\textbf{(i)} Take $h\equiv1$ in \eqref{eq:5-5-40:Ta}, \eqref{eq:5-5-40:Tb} and split the integrals at $x$ (recall that $m$ has no atoms). For $\ell<a<x$,
\begin{align*}
  E^xT_a&=\int_{(a,x)}\bigl(p(y)-p(a)\bigr)\,m(dy)+\bigl(p(x)-p(a)\bigr)m([x,r))\\
  &=-\int_a^x\bigl(p(x)-p(y)\bigr)\,m(dy)+\bigl(p(x)-p(a)\bigr)\,m((a,r)),
\end{align*}
upon writing $p(y)-p(a)=(p(x)-p(a))-(p(x)-p(y))$. For $x<b<r$,
\begin{align*}
  E^xT_b&=\bigl(p(b)-p(x)\bigr)m((\ell,x])+\int_{(x,b)}\bigl(p(b)-p(y)\bigr)\,m(dy)\\
  &=-\int_x^b\bigl(p(y)-p(x)\bigr)\,m(dy)+\bigl(p(b)-p(x)\bigr)\,m((\ell,b)),
\end{align*}
upon writing $p(b)-p(y)=(p(b)-p(x))-(p(y)-p(x))$. Both expectations are finite: they are bounded by $(p(x)-p(a))m(I)$ and $(p(b)-p(x))m(I)$, respectively.

The third identity is the case $h\equiv1$ of the following one, which we prove for every bounded, piecewise continuous $h\colon I\to[0,\infty)$ and $\ell<x<u<r$:
\begin{equation}\label{eq:5-5-40:cycle}
\begin{split}
  E^x\int_0^{S_{u;x}}h(X_t)\,dt&=E^x\int_0^{T_u}h(X_t)\,dt+E^u\int_0^{T_x}h(X_t)\,dt\\
  &=\bigl(p(u)-p(x)\bigr)\int_Ih(y)\,m(dy).
\end{split}
\end{equation}
We have just seen that $E^xT_u<\infty$, so $P^x[T_u<\infty,\ X_{T_u}=u]=1$. The functional $\Gamma(\omega)\defeq\int_0^{T_x(\omega)}h(X_t(\omega))\,dt$ is $\B$-measurable with values in $[0,\infty]$ (Tonelli's theorem, the integrand $h(X_t(\omega))\ind_{\{t<T_x(\omega)\}}$ being jointly measurable), and by \eqref{eq:5-5-40:S}
\[
  \Gamma\circ\theta_{T_u}=\int_0^{T_x\circ\theta_{T_u}}h(X_{T_u+t})\,dt=\int_{T_u}^{S_{u;x}}h(X_t)\,dt\qquad\text{on }\{T_u<\infty\}.
\]
Applying (B) with $T=T_u$, $c=u$ and $g(s,\omega)=\Gamma(\omega)$ we get $E^x\int_{T_u}^{S_{u;x}}h(X_t)\,dt=E^u\Gamma$, which is the first equality in \eqref{eq:5-5-40:cycle}. For the second we use \eqref{eq:5-5-40:Tb} (with $b=u$) and \eqref{eq:5-5-40:Ta} (with starting point $u$ and $a=x$):
\begin{multline*}
  E^x\int_0^{T_u}h(X_t)\,dt+E^u\int_0^{T_x}h(X_t)\,dt\\
  =\int_I\Bigl\{\bigl(p(u)-p(x\vee y)\bigr)\ind_{(\ell,u)}(y)+\bigl(p(u\wedge y)-p(x)\bigr)\ind_{(x,r)}(y)\Bigr\}h(y)\,m(dy).
\end{multline*}
The expression in braces equals $p(u)-p(x)$ for every $y\in I$: for $y\le x$ it is $(p(u)-p(x))+0$, for $x<y<u$ it is $(p(u)-p(y))+(p(y)-p(x))$, and for $y\ge u$ it is $0+(p(u)-p(x))$. This proves \eqref{eq:5-5-40:cycle}. With $h\equiv1$ we obtain
\[
  E^xS_{u;x}=E^xT_u+E^uT_x=\bigl(p(u)-p(x)\bigr)\,m(I)<\infty .
\]

\textbf{(ii)} For $z\in I$ the function $h=\ind_{(\ell,z)}$ is bounded, nonnegative and piecewise continuous, and \eqref{eq:5-5-40:cycle} gives
\[
  E^x\int_0^{S_{u;x}}\ind_{(\ell,z)}(X_t)\,dt=\bigl(p(u)-p(x)\bigr)\,m((\ell,z)).
\]

\textbf{(iii)} Fix $\ell<x<u<r$, $z\in I\setminus\{x\}$, and write $S\defeq S_{u;x}$. We note first that
\begin{equation}\label{eq:5-5-40:noatom}
  P^x[S=t]=0\quad\text{for every }t\ge0,\qquad P^x[0<S<\infty]=1 .
\end{equation}
Indeed, $S\ge T_u>0$ everywhere on $\{X_0=x\}$, because $X_0=x\ne u$; $S<\infty$ almost surely by (i); and for $t>0$ we have $\{S=t\}\subseteq\{X_t=x\}$, an event of $P^x$-probability zero by \bk{(5.80)}.

Let $t\ge0$ and $t_n\to t$, $t_n\ge0$. On the event $\{X_t\ne z\}$ we have $\ind_{\{X_{t_n}<z\}}\to\ind_{\{X_t<z\}}$ by the continuity of the paths, and on $\{S\ne t\}$ we have $\ind_{\{S>t_n\}}\to\ind_{\{S>t\}}$. Now $P^x[X_t=z]=0$: for $t>0$ this is \bk{(5.80)}, and for $t=0$ it holds because $X_0=x\ne z$ (this is where $z\ne x$ is needed). Together with \eqref{eq:5-5-40:noatom}, $\ind_{\{X_{t_n}<z,\ S>t_n\}}\to\ind_{\{X_t<z,\ S>t\}}$ holds $P^x$-almost surely, and the bounded convergence theorem gives $a(t_n)\to a(t)$. Thus $a$ is continuous on $[0,\infty)$; in particular the maxima over $[n-1,n]$ are attained.

For $n-1\le t\le n$ we have $0\le a(t)\le P^x[S>t]\le P^x[S>n-1]$, hence by (i)
\[
  \sum_{n=1}^\infty\max_{n-1\le t\le n}a(t)\le\sum_{n=0}^\infty P^x[S>n]\le1+\int_0^\infty P^x[S>t]\,dt=1+E^xS<\infty .
\]

\textbf{(iv)} Keep $x,u,z,S$ as in (iii), and let $\nu(B)\defeq P^x[S\in B]$, $B\in\B([0,\infty))$. By \eqref{eq:5-5-40:noatom} and (i), $\nu$ is a probability measure on $[0,\infty)$ without atoms, with $\nu(\{0\})=0$ and mean
\[
  \mu\defeq\int_{[0,\infty)}s\,\nu(ds)=E^xS=\bigl(p(u)-p(x)\bigr)\,m(I)\in(0,\infty).
\]

\emph{The renewal equation.} Fix $t\ge0$ and decompose $F(t)=a(t)+P^x[X_t<z,\ S\le t]$. On $\{S\le t\}$ we have $X_t=X_{t-S}\circ\theta_S$, so with the jointly measurable function $g(s,\omega)\defeq\ind_{[0,t]}(s)\,\ind_{\{q(\omega((t-s)^+))<z\}}$ (the map $(s,\omega)\mapsto\omega((t-s)^+)$ is continuous) we have $\ind_{\{X_t<z,\ S\le t\}}=g(S,\theta_S\,\cdot)$ on $\{S<\infty\}$. Since $P^x[S<\infty,\ X_S=x]=1$, (B) applies with $T=S$, $c=x$; as $E^xg(s,\cdot)=\ind_{[0,t]}(s)F(t-s)$, we obtain that $F$ is Borel-measurable and
\begin{equation}\label{eq:5-5-40:renewal}
  F(t)=a(t)+\int_{[0,t]}F(t-s)\,\nu(ds);\qquad0\le t<\infty .
\end{equation}

\emph{Solution of the renewal equation.} For a bounded Borel function $f\ge0$ and a finite measure $\rho$ on $[0,\infty)$ write $(f*\rho)(t)\defeq\int_{[0,t]}f(t-s)\,\rho(ds)$, and let $\nu^{*k}$ be the $k$-fold convolution of $\nu$ ($\nu^{*0}=\delta_0$), i.e.\ the law of the sum of $k$ independent random variables with law $\nu$; by Tonelli's theorem $(f*\nu^{*k})*\nu=f*\nu^{*(k+1)}$. Iterating \eqref{eq:5-5-40:renewal} we get, for every $n\ge1$,
\[
  F=\sum_{k=0}^{n-1}a*\nu^{*k}+F*\nu^{*n}.
\]
With $\beta\defeq\int e^{-s}\,\nu(ds)<1$ (recall $\nu(\{0\})=0$), Chebyshev's inequality gives $\nu^{*n}([0,t])\le e^t\beta^n$, so $0\le(F*\nu^{*n})(t)\le e^t\beta^n\to0$, and the \emph{renewal measure} $U\defeq\sum_{k=0}^\infty\nu^{*k}$ satisfies $U([0,t])\le e^t/(1-\beta)<\infty$. By monotone convergence,
\[
  F(t)=\sum_{k=0}^\infty(a*\nu^{*k})(t)=\int_{[0,t]}a(t-s)\,U(ds);\qquad0\le t<\infty .
\]

\emph{The renewal theorem.} In the form given by Feller (1971), Chapter XI, Section 1, the renewal theorem states: if $\nu$ is a probability measure on $[0,\infty)$ with $\nu(\{0\})<1$ which is not arithmetic (not concentrated on a set $\{0,\lambda,2\lambda,\dots\}$, $\lambda>0$), with mean $\mu\le\infty$ and renewal measure $U$, and if $a$ is directly Riemann integrable, then $\int_{[0,t]}a(t-s)\,U(ds)\to\mu^{-1}\int_0^\infty a(s)\,ds$ as $t\to\infty$. Our $\nu$ has no atoms, so it assigns measure zero to every countable set and is not arithmetic; $a$ is directly Riemann integrable by (iii). Therefore, using Tonelli's theorem and then (ii),
\begin{align*}
  \lim_{t\to\infty}F(t)&=\frac1\mu\int_0^\infty P^x[X_t<z,\ S>t]\,dt
  =\frac1\mu\,E^x\int_0^{S}\ind_{(\ell,z)}(X_t)\,dt\\
  &=\frac{(p(u)-p(x))\,m((\ell,z))}{(p(u)-p(x))\,m(I)}=\frac{m((\ell,z))}{m(I)} .
\end{align*}
This is \bk{(5.81)} for all $x,z\in I$ with $z\ne x$ (any $u\in(x,r)$ may be used).

\emph{The case $z=x$.} For $z'<x<z''$ in $I$ we have $P^x[X_t<z']\le P^x[X_t<x]\le P^x[X_t<z'']$, so by what has been proved
\[
  \frac{m((\ell,z'))}{m(I)}\le\liminf_{t\to\infty}P^x[X_t<x]\le\limsup_{t\to\infty}P^x[X_t<x]\le\frac{m((\ell,z''))}{m(I)} .
\]
Letting $z'\uparrow x$ and $z''\downarrow x$, both bounds converge to $m((\ell,x))/m(I)$, because $m$ is finite and $m(\{x\})=0$. Hence \bk{(5.81)} holds for all $x,z\in I$.

\rmk{(a) The boundedness of $b$ and $\sigma$ on compact subintervals of $I$ was used only in Step 0: for $\sigma$, through the local boundedness of $\wt\sigma$, which \bk{Theorem 4.20} requires of the coefficients of \eqref{eq:5-5-40:Y}; for both $b$ and $\sigma$, in the verification of \bk{(5.3)} for $X'=q(Y')$, which shows that weak solutions in $I$ exist. (b) As a check, for the Ornstein--Uhlenbeck equation $dX_t=-X_t\,dt+dW_t$ on $I=\R$ we have, with $c=0$ in \bk{(5.42)}, $p'(x)=e^{x^2}$, $p(\pm\infty)=\pm\infty$, $m(dx)=2e^{-x^2}dx$, $m(I)=2\sqrt\pi$; the law of $X_t$ is Gaussian, so \bk{(5.80)} holds, and \bk{(5.81)} asserts that $X_t$ converges in distribution to the normal law with mean $0$ and variance $\tfrac12$, as is seen directly from $X_t=xe^{-t}+\int_0^te^{-(t-s)}\,dW_s$. For $x=0$, $u=1$ the formulas of (i), (ii) give $E^0T_1\approx4.038$, $E^1T_0\approx1.147$, $E^0S_{1;0}=2\sqrt\pi\int_0^{1}e^{s^2}ds\approx5.185$ and, for $z=\tfrac12$, $E^0\int_0^{S_{1;0}}\ind_{(-\infty,z)}(X_t)\,dt\approx3.942$; a simulation of the equation by the Euler scheme reproduces these four values to within its discretization and sampling errors (a few percent, the grid detecting the hitting times slightly late).}
\end{kex}

\begin{kex}{Exercise}{5.41}{}
We recall the setting of \bk{Proposition 2.13}: $d=r=1$, the coefficients $b,\sigma\colon[0,\infty)\times\R\to\R$ are Borel-measurable and satisfy
\[
  \abs{b(t,x)-b(t,y)}\le K\abs{x-y},\qquad \abs{\sigma(t,x)-\sigma(t,y)}\le h(\abs{x-y});
  \qquad 0\le t<\infty,\ x,y\in\R,
\]
where $K>0$ and $h\colon[0,\infty)\to[0,\infty)$ is strictly increasing with $h(0)=0$ and $\int_{(0,\eps)}h^{-2}(u)\,du=\infty$ for every $\eps>0$. In accordance with \bk{Definition 2.3}, let $W$ be a one-dimensional Brownian motion on some $(\Omega,\F,P)$, let $\xi$ be a random variable independent of $W$, let $\{\F_t\}$ be the filtration \bk{(2.3)}, which satisfies the usual conditions, and let $X^{(1)}$, $X^{(2)}$ be two strong solutions of \bk{(2.1)} relative to $W$ with initial condition $\xi$ (\bk{Definition 2.1}). We have to show that $P[X^{(1)}_t=X^{(2)}_t;\ 0\le t<\infty]=1$.

\emph{Step 1: the difference is a continuous semimartingale.} We set
\[
  X\defeq X^{(1)}-X^{(2)},\qquad
  \beta_s\defeq b\bigl(s,X^{(1)}_s\bigr)-b\bigl(s,X^{(2)}_s\bigr),\qquad
  \gamma_s\defeq \sigma\bigl(s,X^{(1)}_s\bigr)-\sigma\bigl(s,X^{(2)}_s\bigr).
\]
The processes $X^{(i)}$ are continuous and $\{\F_t\}$-adapted, hence progressively measurable, and $b$, $\sigma$ are Borel-measurable; therefore $\beta$ and $\gamma$ are progressively measurable. The hypotheses on the coefficients give, for every $\omega\in\Omega$ and $0\le s\le t<\infty$,
\begin{equation}\label{eq:5-5-41:a}
  \abs{\beta_s}\le K\abs{X_s}\le K\max_{0\le u\le t}\abs{X_u}<\infty,\qquad
  \abs{\gamma_s}\le h(\abs{X_s})\le h\Bigl(\max_{0\le u\le t}\abs{X_u}\Bigr)<\infty,
\end{equation}
where we have used the continuity of $X$ and the monotonicity of $h$. Consequently
\[
  V^{\pm}_t\defeq\int_0^t\beta_s^{\pm}\,ds,\qquad V_t\defeq V^+_t-V^-_t=\int_0^t\beta_s\,ds;\qquad 0\le t<\infty,
\]
are well defined and finite for every $\omega$; $V^{\pm}$ are continuous, nondecreasing, null at $t=0$, and adapted (by Fubini's theorem and the progressive measurability of $\beta^\pm$). Likewise $\int_0^t\gamma_s^2\,ds<\infty$ for every $t$ and $\omega$, so the stochastic integral
\[
  M_t\defeq\int_0^t\gamma_s\,dW_s;\qquad 0\le t<\infty,
\]
is defined and belongs to $\cM^{c,\loc}$ (\bk{Section 3.2.D}), with
\begin{equation}\label{eq:5-5-41:b}
  \qv{M}_t=\int_0^t\gamma_s^2\,ds;\qquad 0\le t<\infty,\ \as
\end{equation}
(this is \bk{(3.2.19)}, which remains valid for integrators in $\cM^{c,\loc}$, as noted before \bk{Proposition 3.2.24}). Subtracting the two instances of \bk{(2.4)}, in which $X^{(1)}_0=\xi=X^{(2)}_0$ a.s.\ by \bk{Definition 2.1}(ii) and all four integrals are well defined by \bk{Definition 2.1}(iii), and using the linearity of the Lebesgue and of the stochastic integral, we obtain
\begin{equation}\label{eq:5-5-41:c}
  X_0=0\quad\text{and}\quad X_t=M_t+V_t;\qquad 0\le t<\infty,\ \as
\end{equation}
Thus $X$ is a continuous semimartingale with a decomposition of the form \bk{(3.7.1)}, relative to a filtration which satisfies the usual conditions. Let $\Lambda=\{\Lambda_t(a)\}$ be its semimartingale local time (\bk{Theorem 3.7.1}).

\emph{Step 2: the local time of $X$ at the origin vanishes.} We define $k\colon(0,\infty)\to(0,\infty)$ by $k(u)\defeq h^2(u)$. This function is strictly positive, because $h$ is strictly increasing with $h(0)=0$; it is Borel-measurable, because $h$ is monotone; and $\int_{(0,\eps)}du/k(u)=\infty$ for every $\eps>0$ by \bk{(2.25)}. For every $s$ with $X_s>0$ the second bound in \eqref{eq:5-5-41:a} reads $\gamma_s^2\le h^2(X_s)=k(X_s)$; in conjunction with \eqref{eq:5-5-41:b} this gives, for every $t\in(0,\infty)$,
\[
  \int_0^t\frac{d\qv{M}_s}{k(X_s)}\,\ind_{\{X_s>0\}}
  =\int_0^t\frac{\gamma_s^2}{h^2(X_s)}\,\ind_{\{X_s>0\}}\,ds\le t<\infty,\qquad\as
\]
The hypotheses of \bk{Exercise 3.7.12} are thus satisfied, and we conclude that
\begin{equation}\label{eq:5-5-41:d}
  P\bigl[\Lambda_t(0)=0;\ 0\le t<\infty\bigr]=1 .
\end{equation}

\emph{Step 3: the Tanaka--Meyer formula.} The formula \bk{(3.7.9)} with $a=0$ (that is, \bk{(3.7.4)} for the convex function $f(x)=\abs{x}$, whose left-hand derivative is the function $\sgn$ of \bk{(3.6.14)} and whose second derivative measure is twice the unit mass at the origin), together with \eqref{eq:5-5-41:c} and \eqref{eq:5-5-41:d}, gives
\begin{equation}\label{eq:5-5-41:e}
  \abs{X_t}=N_t+\int_0^t\sgn(X_s)\,\beta_s\,ds;\qquad 0\le t<\infty,\ \as,
  \qquad\text{where } N_t\defeq\int_0^t\sgn(X_s)\,dM_s .
\end{equation}
Here we have written $\int_0^t\sgn(X_s)\,dV_s=\int_0^t\sgn(X_s)\,dV^+_s-\int_0^t\sgn(X_s)\,dV^-_s=\int_0^t\sgn(X_s)\beta_s\,ds$, and $N\in\cM^{c,\loc}$ with $N_0=0$, the integrand $\sgn(X)$ being progressively measurable and bounded. From \eqref{eq:5-5-41:a} and $\abs{\sgn}\equiv1$ we have
\begin{equation}\label{eq:5-5-41:f}
  \Bigl|\int_0^t\sgn(X_s)\,\beta_s\,ds\Bigr|\le K\int_0^t\abs{X_s}\,ds;\qquad 0\le t<\infty .
\end{equation}

\emph{Step 4: localization and the Gronwall inequality.} For $n\ge1$ let $\tau_n\defeq\inf\{t\ge0;\ \abs{X_t}\ge n\}$. This is a stopping time of $\{\F_t\}$, as the first entrance time of a closed set by a continuous adapted process (\bk{Problem 1.2.7}), and $\tau_n\uparrow\infty$ as $n\to\infty$ for every $\omega$, because the paths of $X$ are bounded on compact time-intervals. Since $X_0=0$ a.s., we have $\abs{X_{s\wedge\tau_n}}\le n$ for all $s\ge0$, a.s. We fix $n$ and $t$. From \eqref{eq:5-5-41:e} and \eqref{eq:5-5-41:f},
\begin{equation}\label{eq:5-5-41:g}
  \abs{N_{s\wedge\tau_n}}\le\abs{X_{s\wedge\tau_n}}+K\int_0^{s\wedge\tau_n}\abs{X_u}\,du\le n(1+Kt);\qquad 0\le s\le t,\ \as
\end{equation}
Let $\{\sigma_k\}_{k=1}^\infty$ be a nondecreasing sequence of stopping times with $P[\lim_{k\to\infty}\sigma_k=\infty]=1$ such that $\{N_{s\wedge\sigma_k},\F_s;\ 0\le s<\infty\}$ is a (continuous) martingale for every $k$ (\bk{Definition 1.5.15}). The martingale $\{N_{s\wedge t\wedge\sigma_k},\F_s;\ 0\le s<\infty\}$ is continuous and has the last element $N_{t\wedge\sigma_k}$; the optional sampling theorem (\bk{Theorem 1.3.22}), applied to it and to the stopping time $t\wedge\tau_n$, gives $E[N_{t\wedge\tau_n\wedge\sigma_k}]=E[N_0]=0$. As $k\to\infty$ we have $N_{t\wedge\tau_n\wedge\sigma_k}\to N_{t\wedge\tau_n}$ a.s., and these random variables are bounded by the constant in \eqref{eq:5-5-41:g}; the bounded convergence theorem yields
\[
  E\bigl[N_{t\wedge\tau_n}\bigr]=0 .
\]
We now evaluate \eqref{eq:5-5-41:e} at $t\wedge\tau_n$, take expectations (all terms are bounded), and use \eqref{eq:5-5-41:f}:
\begin{align*}
  E\abs{X_{t\wedge\tau_n}}
  =E\int_0^{t\wedge\tau_n}\sgn(X_s)\,\beta_s\,ds
  &\le K\,E\int_0^{t\wedge\tau_n}\abs{X_{s\wedge\tau_n}}\,ds\\
  &\le K\,E\int_0^{t}\abs{X_{s\wedge\tau_n}}\,ds
  =K\int_0^t E\abs{X_{s\wedge\tau_n}}\,ds,
\end{align*}
the last step by Tonelli's theorem (the map $(s,\omega)\mapsto X_{s\wedge\tau_n(\omega)}(\omega)$ is jointly measurable, being continuous in $s$ and measurable in $\omega$). Hence the function $g(t)\defeq E\abs{X_{t\wedge\tau_n}}$ satisfies
\[
  0\le g(t)\le K\int_0^t g(s)\,ds;\qquad 0\le t<\infty .
\]
Moreover $g$ is continuous, by the continuity of the paths of $X$ and the bounded convergence theorem ($\abs{X_{t\wedge\tau_n}}\le n$ a.s.). The Gronwall inequality (\bk{Problem 2.7}, with $\alpha\equiv0$ and $\beta=K$, on every interval $[0,T]$) gives $g\equiv0$, that is,
\[
  X_{t\wedge\tau_n}=0\quad\as,\qquad\text{for every } t\ge0 \text{ and } n\ge1 .
\]

\emph{Step 5: conclusion.} Letting $n\to\infty$ and recalling that $\tau_n\uparrow\infty$, we obtain $X_t=0$ a.s.\ for every fixed $t\ge0$; thus $P[X_t=0;\ t\in\Q\cap[0,\infty)]=1$, and the continuity of the paths of $X$ yields
\[
  P\bigl[X^{(1)}_t=X^{(2)}_t;\ 0\le t<\infty\bigr]=P\bigl[X_t=0;\ 0\le t<\infty\bigr]=1 .
\]
Since $W$, $\xi$ and the two strong solutions were arbitrary, strong uniqueness holds for equation \bk{(2.1)}, which is the assertion of \bk{Proposition 2.13}.

\rmk{The condition \bk{(2.24)}--\bk{(2.25)} on $\sigma$ enters only in Step~2, where it replaces the approximation of $\abs{x}$ by the functions $\psi_n$ of \bk{(2.26)} in the book's proof: both devices serve to show that the ``second-order term'' in the It\^o expansion of $\abs{X_t}$, here $2\Lambda_t(0)$, vanishes. No use was made of the particular form \bk{(2.3)} of the filtration, beyond the usual conditions; nor was the reduction to the integrability condition \bk{(2.27)} needed, since the stopped local martingale $N_{\cdot\wedge\tau_n}$ is bounded on bounded time-intervals by \eqref{eq:5-5-41:g}.}
\end{kex}

\begin{kex}{Exercise}{5.42}{}
\emph{Setting and conventions.} The maximum $X^{(1)}\vee X^{(2)}$ and the local time of $X^{(1)}-X^{(2)}$ make sense only for real-valued processes, so the state space is one-dimensional, $d=1$; the driving Brownian motion $W=(W^{(1)},\dots,W^{(r)})$ may have any dimension $r\ge 1$, and $b,\ \sigma_j\colon[0,\infty)\times\R\to\R$, $1\le j\le r$, are Borel-measurable. The semimartingale local time $\Lambda$ of \bk{Theorem 3.7.1} is defined for a continuous semimartingale relative to one filtration satisfying the usual conditions. Accordingly, the hypothesis is understood as follows.
\begin{itemize}
\item[\textup{(H)}] Whenever $(X^{(1)},W)$ and $(X^{(2)},W)$ are weak solutions of \bk{(2.1)} (\bk{Definition 3.1}) on the same space $(\Omega,\F,P)$, relative to the same filtration $\{\F_t\}$, with $P[X^{(1)}_0=X^{(2)}_0]=1$, the local time $\Lambda(0)$ at the origin of the continuous semimartingale $X=X^{(1)}-X^{(2)}$ satisfies $P[\Lambda_t(0)=0;\ \forall\, 0\le t<\infty]=1$.
\end{itemize}
In Steps 1--4 we prove that two such solutions are indistinguishable. \bk{Definition 3.2} allows the two solutions to be adapted to different filtrations; Step 5 removes this restriction.

Fix two weak solutions as in (H). For $i=1,2$ we have, almost surely,
\begin{equation}\label{eq:5-5-42:a}
  X^{(i)}_t=X^{(i)}_0+\int_0^t b(s,X^{(i)}_s)\,ds+\sum_{j=1}^r\int_0^t\sigma_j(s,X^{(i)}_s)\,dW^{(j)}_s;\qquad 0\le t<\infty,
\end{equation}
and
\begin{equation}\label{eq:5-5-42:b}
  P\Bigl[\int_0^t\bigl\{\abs{b(s,X^{(i)}_s)}+\sigma_j^2(s,X^{(i)}_s)\bigr\}\,ds<\infty\Bigr]=1;\qquad 0\le t<\infty,\ 1\le j\le r .
\end{equation}
The processes $X^{(i)}$ are continuous and adapted, hence progressively measurable, and so are the processes
\[
  \beta_s\defeq b(s,X^{(1)}_s)-b(s,X^{(2)}_s),\qquad K^{(j)}_s\defeq\sigma_j(s,X^{(1)}_s)-\sigma_j(s,X^{(2)}_s),\qquad H_s\defeq\ind_{(0,\infty)}(X_s),
\]
because $b$, $\sigma_j$ are Borel-measurable.

\emph{Step 1. $X$ is a continuous semimartingale of the form \bk{(3.7.1)}.} By \eqref{eq:5-5-42:b} we have $\int_0^t\abs{\beta_s}\,ds<\infty$ and $\int_0^t (K^{(j)}_s)^2\,ds<\infty$ for all $t$, almost surely. Hence
\[
  M_t\defeq\sum_{j=1}^r N^{(j)}_t,\qquad N^{(j)}_t\defeq\int_0^t K^{(j)}_s\,dW^{(j)}_s,\qquad
  V_t\defeq\int_0^t\beta_s^+\,ds-\int_0^t\beta_s^-\,ds
\]
are well defined, $M\in\cM^{c,\loc}$, and $V$ is the difference of two continuous, nondecreasing processes starting at zero (set equal to zero on the null event where one of the integrals is infinite for some $t$); these are adapted, because $\beta^{\pm}$ are progressively measurable and $\F_0$ contains the $P$-null sets. Subtracting the two equations \eqref{eq:5-5-42:a} we obtain $X_t=X_0+M_t+V_t$, $0\le t<\infty$, almost surely, with $X_0=X^{(1)}_0-X^{(2)}_0=0$ a.s. The filtration satisfies the usual conditions by \bk{Definition 3.1}\,(i), so \bk{Theorem 3.7.1} applies to $X$.

\emph{Step 2. The Tanaka--Meyer formula.} The function $f(x)=x^+$ is convex, with left-hand derivative $D^-f=\ind_{(0,\infty)}$ and with second derivative measure \bk{(3.6.47)} equal to the unit mass at the origin, since $\mu([a,b))=\ind_{(0,\infty)}(b)-\ind_{(0,\infty)}(a)=\ind_{\{a\le 0<b\}}$. Therefore \bk{(3.7.4)} (equivalently, \bk{(3.7.7)} and \bk{(3.7.10)} with $a=0$) gives, almost surely,
\[
  X_t^+=X_0^++\int_0^tH_s\,dM_s+\int_0^tH_s\,dV_s+\Lambda_t(0);\qquad 0\le t<\infty .
\]
Here $X_0^+=0$ a.s., and $\Lambda_\cdot(0)\equiv 0$ a.s.\ by (H). Since $V$ is absolutely continuous with density $\beta$, we have $\int_0^tH_s\,dV_s=\int_0^tH_s\beta_s\,ds$. For the stochastic integral we claim that
\begin{equation}\label{eq:5-5-42:c}
  \int_0^tH_s\,dM_s=\sum_{j=1}^r\int_0^tH_sK^{(j)}_s\,dW^{(j)}_s;\qquad 0\le t<\infty,\ \text{a.s.}
\end{equation}
Indeed, $H$ is bounded, so it is integrable with respect to each of $M$, $N^{(1)},\dots,N^{(r)}$, and \bk{Problem 3.2.25} gives $\int_0^\cdot H\,dM=\sum_j\int_0^\cdot H\,dN^{(j)}$. Fix $j$; the process $HK^{(j)}$ is progressively measurable with $\int_0^t(H_sK^{(j)}_s)^2\,ds\le\int_0^t(K^{(j)}_s)^2\,ds<\infty$ a.s. For every $L\in\cM^{c,\loc}$, three applications of \bk{(3.2.30)}, in the form valid for local martingales (\bk{Proposition 3.2.24}), yield
\[
  \Bigl\langle\int_0^\cdot H\,dN^{(j)},L\Bigr\rangle_t=\int_0^tH_s\,d\qv{N^{(j)},L}_s
  =\int_0^tH_sK^{(j)}_s\,d\qv{W^{(j)},L}_s=\Bigl\langle\int_0^\cdot HK^{(j)}\,dW^{(j)},L\Bigr\rangle_t ,
\]
the middle equality because $\qv{N^{(j)},L}_t=\int_0^tK^{(j)}_s\,d\qv{W^{(j)},L}_s$. The uniqueness assertion of \bk{Proposition 3.2.24} now gives $\int_0^\cdot H\,dN^{(j)}=\int_0^\cdot HK^{(j)}\,dW^{(j)}$, and \eqref{eq:5-5-42:c} follows. We have proved that, almost surely,
\begin{equation}\label{eq:5-5-42:d}
  X_t^+=\int_0^tH_s\beta_s\,ds+\sum_{j=1}^r\int_0^tH_sK^{(j)}_s\,dW^{(j)}_s;\qquad 0\le t<\infty .
\end{equation}

\emph{Step 3. $Y\defeq X^{(1)}\vee X^{(2)}$ solves \bk{(2.1)}.} The process $Y$ is continuous and $\{\F_t\}$-adapted, and $Y=X^{(2)}+X^+$. Moreover, $Y_s=X^{(1)}_s$ on $\{H_s=1\}=\{X^{(1)}_s>X^{(2)}_s\}$ and $Y_s=X^{(2)}_s$ on $\{H_s=0\}$, so that for every Borel function $g\colon[0,\infty)\times\R\to\R$,
\begin{equation}\label{eq:5-5-42:e}
  g(s,Y_s)=g(s,X^{(2)}_s)+H_s\bigl[g(s,X^{(1)}_s)-g(s,X^{(2)}_s)\bigr];\qquad 0\le s<\infty .
\end{equation}
In particular $\abs{b(s,Y_s)}\le\abs{b(s,X^{(1)}_s)}+\abs{b(s,X^{(2)}_s)}$ and $\sigma_j^2(s,Y_s)\le\sigma_j^2(s,X^{(1)}_s)+\sigma_j^2(s,X^{(2)}_s)$, so \eqref{eq:5-5-42:b} holds with $Y$ in place of $X^{(i)}$; this is condition (iii) of \bk{Definition 2.1}. Adding \eqref{eq:5-5-42:a} for $i=2$ to \eqref{eq:5-5-42:d}, and using \eqref{eq:5-5-42:e} with $g=b$ and with $g=\sigma_j$ together with the linearity of the integrals, we obtain almost surely
\begin{align*}
  Y_t&=X^{(2)}_0+\int_0^t\bigl[b(s,X^{(2)}_s)+H_s\beta_s\bigr]\,ds+\sum_{j=1}^r\int_0^t\bigl[\sigma_j(s,X^{(2)}_s)+H_sK^{(j)}_s\bigr]\,dW^{(j)}_s\\
  &=Y_0+\int_0^tb(s,Y_s)\,ds+\sum_{j=1}^r\int_0^t\sigma_j(s,Y_s)\,dW^{(j)}_s;\qquad 0\le t<\infty,
\end{align*}
where we have used $Y_0=X^{(1)}_0\vee X^{(2)}_0=X^{(2)}_0$ a.s. Thus $(Y,W)$, $(\Omega,\F,P)$, $\{\F_t\}$ is a weak solution of \bk{(2.1)}, and $P[Y_0=X^{(1)}_0=X^{(2)}_0]=1$.

\emph{Step 4. Indistinguishability.} The three weak solutions $(Y,W)$, $(X^{(1)},W)$, $(X^{(2)},W)$ have the same initial distribution. By uniqueness in the sense of probability law (\bk{Definition 3.4}), the processes $Y$, $X^{(1)}$, $X^{(2)}$ have the same law; in particular, for each fixed $t\ge0$ and $i=1,2$, the random variables $Y_t$ and $X^{(i)}_t$ have the same distribution, whence
\[
  E\bigl[\arctan Y_t-\arctan X^{(i)}_t\bigr]=0 .
\]
The integrand is bounded, and it is nonnegative because $Y_t\ge X^{(i)}_t$; hence it vanishes almost surely, and since $\arctan$ is strictly increasing, $Y_t=X^{(i)}_t$ a.s. Consequently $P[X^{(1)}_t=X^{(2)}_t]=1$ for every $t\ge 0$, thus $P[X^{(1)}_t=X^{(2)}_t;\ \forall\,t\in\Q\cap[0,\infty)]=1$, and by the continuity of the paths
\[
  P\bigl[X^{(1)}_t=X^{(2)}_t;\ \forall\,0\le t<\infty\bigr]=1 .
\]
This is pathwise uniqueness for pairs of solutions relative to a common filtration.

\emph{Step 5. Solutions relative to different filtrations.} Let now $(X^{(1)},W)$, $(\Omega,\F,P)$, $\{\F^{(1)}_t\}$ and $(X^{(2)},W)$, $(\Omega,\F,P)$, $\{\F^{(2)}_t\}$ be weak solutions with $P[X^{(1)}_0=X^{(2)}_0]=1$, as in \bk{Definition 3.2}; they have a common initial distribution $\mu$. We apply to them the construction of \bk{Section 5.3.D}, with $(\Omega_j,\F_j,\nu_j)=(\Omega,\F,P)$, $j=1,2$, and with the Brownian motion of each of the two solutions (the process denoted there by $W^{(j)}$, not to be confused with the coordinates of $W$ used above) equal to our $W$; apart from this we use the notation of that section, writing $\wt P$ for the measure \bk{(3.23)} on the canonical space. On that space, the processes $(x+y_1,w)$ and $(x+y_2,w)$ are weak solutions of \bk{(2.1)} relative to one and the same filtration, with the same Brownian motion $w$ and the same initial value $x$, as recorded in the proof of \bk{Proposition 3.20}. Steps 1--4 apply to this pair and give \bk{(3.24)}: $\wt P[y_1=y_2]=1$. This is the only way in which pathwise uniqueness enters \bk{Problem 3.22}, which therefore provides a measurable function $k\colon\R\times C[0,\infty)^r\to C[0,\infty)$ satisfying \bk{(3.28)}:
\[
  \wt P\bigl[y_1=y_2=k(x,w)\bigr]=1 .
\]
The set $A\defeq\{(x,w,y)\in\Theta;\ y=k(x,w)\}$ belongs to $\B(\Theta)$, being the inverse image of the diagonal of $C[0,\infty)\times C[0,\infty)$ (a closed subset of a separable metric space, hence a member of the product $\sigma$-field) under the measurable mapping $(x,w,y)\mapsto(k(x,w),y)$. From \bk{(3.21)$'$} we conclude that, for $j=1,2$,
\[
  P\bigl[X^{(j)}_\cdot-X^{(j)}_0=k(X^{(j)}_0,W_\cdot)\bigr]=\wt P\bigl[(x,w,y_j)\in A\bigr]=1 .
\]
Since $X^{(1)}_0=X^{(2)}_0$ a.s., we obtain $X^{(1)}_\cdot=X^{(1)}_0+k(X^{(1)}_0,W_\cdot)=X^{(2)}_0+k(X^{(2)}_0,W_\cdot)=X^{(2)}_\cdot$ almost surely. Hence pathwise uniqueness holds for \bk{(2.1)} in the sense of \bk{Definition 3.2}.

\rmk{The hypothesis on the local time cannot be dropped. In \bk{Example 3.5} (the equation $dX_t=\sgn(X_t)\,dW_t$, $X_0=0$, with $r=1$) uniqueness in the sense of probability law holds, and if $(X,W)$, $(\Omega,\F,P)$, $\{\F_t\}$ is a weak solution then so is $(-X,W)$, $(\Omega,\F,P)$, $\{\F_t\}$; moreover $X$ is a Brownian motion relative to $\{\F_t\}$. The local time at the origin of the difference $Z\defeq X-(-X)=2X$ is not identically zero. Indeed, $Z$ has the decomposition \bk{(3.7.1)} with $Z_0=0$, $M=2X$, $V\equiv0$, and \bk{(3.7.9)} with $a=0$ gives, since $\sgn(2x)=\sgn(x)$,
\[
  2\abs{X_t}=2\int_0^t\sgn(X_s)\,dX_s+2\Lambda_t(0);\qquad 0\le t<\infty,\ \text{a.s.}
\]
The integrand $\sgn(X_s)$ is bounded and progressively measurable, so $\int_0^\cdot\sgn(X_s)\,dX_s$ is a square-integrable martingale with respect to the Brownian motion $X$ (\bk{Proposition 3.2.10}) and has expectation zero. If $\Lambda_\cdot(0)$ vanished identically a.s., taking expectations would give $E\abs{X_t}=0$, which is false for $t>0$ because $X_t$ is normal with mean zero and variance $t$. Hence $P[\Lambda_t(0)>0]>0$ for some (in fact every) $t>0$, and pathwise uniqueness fails, as \bk{Example 3.5} shows.}
\end{kex}

\solsection{5.6}{Linear Equations}

\begin{kex}{Problem}{6.1}{}
Here $X$ is the process of \bk{(6.6)}, the (pathwise unique) solution of \bk{(6.1)}. As in Section 5.2, $\{\F_t\}$ denotes the augmented filtration \bk{(2.3)} generated by $\xi$ and $W$: because $\xi$ and $W$ are independent, $W=\{W_t,\F_t;\,0\le t<\infty\}$ is an $r$-dimensional Brownian motion (with $W_0=0$), the initial vector $X_0=\xi$ is $\F_0$-measurable, and $\{\F_t\}$ satisfies the usual conditions (see the discussion following \bk{(2.3)}). The coefficients $A$, $a$, $\sigma$ are nonrandom, measurable and locally bounded. Integrals and expectations of vector- or matrix-valued quantities are taken entry by entry, so that constant matrices may be moved in and out of them; $\norm{\cdot}$ denotes the Euclidean norm of a vector and of a matrix regarded as a vector of its entries.

\emph{Step 1: the deterministic ingredients.} The fundamental solution $\Phi$ of \bk{(6.3)} is absolutely continuous on every bounded interval, in particular continuous, and $\Phi(t)$ is nonsingular for every $t\ge 0$. Since matrix inversion is continuous on the set of nonsingular matrices (Cramer's rule), $\Phi^{-1}$ is continuous as well, and both $\Phi$ and $\Phi^{-1}$ are bounded on every interval $[0,T]$. Consequently the functions
\[
  b(u)\defeq\Phi^{-1}(u)a(u),\qquad \Psi(u)\defeq\Phi^{-1}(u)\sigma(u);\qquad 0\le u<\infty,
\]
with values in $\R^d$ and in the $(d\times r)$ matrices, respectively, are Borel-measurable and bounded on every $[0,T]$. We set
\[
  \beta(t)\defeq\int_0^t b(u)\,du,\qquad
  Q(t)\defeq\int_0^t\Psi(u)\Psi^T(u)\,du;\qquad 0\le t<\infty,
\]
which are well-defined, nonrandom and absolutely continuous on bounded intervals.

\emph{Step 2: the stochastic integral.} By \bk{Problem 2.5.5}, the components $W^{(1)},\dots,W^{(r)}$ of $W$ belong to $\cM_2^c$ and $\qv{W^{(j)},W^{(l)}}_t=\delta_{jl}\,t$. Each entry $\Psi_{ij}$ is a nonrandom, measurable function, hence, regarded as a process, it is progressively measurable, and $E\int_0^T\Psi_{ij}^2(u)\,du=\int_0^T\Psi_{ij}^2(u)\,du<\infty$ for every $T<\infty$; thus $\Psi_{ij}\in\cL^*(W^{(j)})$. By \bk{Proposition 3.2.10}, the processes
\[
  M^{(i)}_t\defeq\sum_{j=1}^r\int_0^t\Psi_{ij}(u)\,dW^{(j)}_u;\qquad 0\le t<\infty,\ 1\le i\le d,
\]
belong to $\cM_2^c$; they are the components of the vector $M_t=\int_0^t\Phi^{-1}(u)\sigma(u)\,dW_u$ appearing in \bk{(6.6)}. The bilinearity of the cross-variation (\bk{Problem 1.5.7}) and \bk{Proposition 3.2.17} give, for $1\le i,k\le d$,
\[
  \qv{M^{(i)},M^{(k)}}_t
  =\sum_{j=1}^r\sum_{l=1}^r\int_0^t\Psi_{ij}(u)\Psi_{kl}(u)\,d\qv{W^{(j)},W^{(l)}}_u
  =\int_0^t\sum_{j=1}^r\Psi_{ij}(u)\Psi_{kj}(u)\,du=Q_{ik}(t),
\]
a nonrandom quantity. Now $M^{(i)}M^{(k)}-\qv{M^{(i)},M^{(k)}}$ is a martingale vanishing at $t=0$ (\bk{Definition 1.5.5}), so $E[M^{(i)}_sM^{(k)}_s]=Q_{ik}(s)$ for every $s\ge 0$. If $0\le s\le t$, then $M^{(i)}_s$ is $\F_s$-measurable and square-integrable, $M^{(k)}_t$ is square-integrable, and the martingale property yields
\[
  E\bigl[M^{(i)}_sM^{(k)}_t\bigr]=E\bigl[M^{(i)}_s\,E[M^{(k)}_t\,|\,\F_s]\bigr]=E\bigl[M^{(i)}_sM^{(k)}_s\bigr]=Q_{ik}(s).
\]
Interchanging the roles of $(s,i)$ and $(t,k)$ when $t<s$, and noting that $Q$ is symmetric, we conclude that
\begin{equation}\label{eq:5-6-01:a}
  EM_t=0,\qquad E\bigl[M_sM_t^T\bigr]=Q(s\wedge t);\qquad 0\le s,t<\infty .
\end{equation}

\emph{Step 3: proof of \bk{(6.10)}.} In the notation just introduced, \bk{(6.6)} reads
\begin{equation}\label{eq:5-6-01:b}
  X_t=\Phi(t)\bigl[X_0+\beta(t)+M_t\bigr];\qquad 0\le t<\infty .
\end{equation}
Since $E\norm{X_0}^2<\infty$, $\beta(t)$ is nonrandom and $E\norm{M_t}^2=\tr Q(t)<\infty$, every component of $X_t$ is square-integrable; hence $m(t)$, and by the Cauchy--Schwarz inequality also $\rho(s,t)$, are well defined and finite. Taking expectations in \eqref{eq:5-6-01:b} and using $EM_t=0$, we obtain $m(t)=\Phi(t)[m(0)+\beta(t)]$, which is \bk{(6.10)}.

\emph{Step 4: proof of \bk{(6.11)}.} Subtracting \bk{(6.10)} from \eqref{eq:5-6-01:b}, we have
\[
  X_t-m(t)=\Phi(t)\bigl[Y+M_t\bigr],\qquad Y\defeq X_0-m(0),
\]
and therefore, for $0\le s,t<\infty$,
\begin{equation}\label{eq:5-6-01:c}
  \rho(s,t)=\Phi(s)\Bigl(E\bigl[YY^T\bigr]+E\bigl[YM_t^T\bigr]+E\bigl[M_sY^T\bigr]+E\bigl[M_sM_t^T\bigr]\Bigr)\Phi^T(t),
\end{equation}
all four expectations being finite by the Cauchy--Schwarz inequality. The first one is $\rho(0,0)=V(0)$. The vector $Y$ is $\F_0$-measurable and square-integrable, and each $M^{(k)}$ is a square-integrable martingale with $M^{(k)}_0=0$; thus
\[
  E\bigl[Y_iM^{(k)}_t\bigr]=E\bigl[Y_i\,E[M^{(k)}_t\,|\,\F_0]\bigr]=E\bigl[Y_iM^{(k)}_0\bigr]=0;\qquad 1\le i,k\le d,\ t\ge 0,
\]
so the second and third expectations in \eqref{eq:5-6-01:c} vanish. The fourth is $Q(s\wedge t)$ by \eqref{eq:5-6-01:a}. Hence
\[
  \rho(s,t)=\Phi(s)\bigl[V(0)+Q(s\wedge t)\bigr]\Phi^T(t)
  =\Phi(s)\biggl[V(0)+\int_0^{s\wedge t}\Phi^{-1}(u)\sigma(u)\bigl(\Phi^{-1}(u)\sigma(u)\bigr)^T\,du\biggr]\Phi^T(t),
\]
which is \bk{(6.11)}.

\emph{Step 5: the equations \bk{(6.12)} and \bk{(6.13)}.} Because $A$, $a$, $\sigma$ are only measurable, these equations are to be understood in the same sense as \bk{(6.2)} and \bk{(6.3)}: $m$ and $V$ are absolutely continuous on every bounded interval and satisfy the equations at Lebesgue-almost every $t\ge 0$ (equivalently, in integrated form). To verify this, recall three facts about real functions on a bounded interval $[0,T]$: the indefinite integral of a bounded measurable function $g$ is absolutely continuous, with derivative $g(t)$ at almost every $t$ (Lebesgue's differentiation theorem); sums and products of absolutely continuous functions are absolutely continuous; and the product rule holds at every point at which the factors are differentiable. By definition of $\Phi$ there is a Lebesgue-null set $N_1$ off which $\Phi$ is differentiable with $\dot\Phi(t)=A(t)\Phi(t)$, hence also $\frac{d}{dt}\Phi^T(t)=\Phi^T(t)A^T(t)$; and there is a Lebesgue-null set $N_2$ off which $\beta$ and $Q$ are differentiable with $\dot\beta(t)=\Phi^{-1}(t)a(t)$ and $\dot Q(t)=\Phi^{-1}(t)\sigma(t)\sigma^T(t)\bigl(\Phi^{-1}(t)\bigr)^T$.

By \bk{(6.10)}, $m=\Phi\,[m(0)+\beta]$ has entries which are finite sums of products of absolutely continuous functions, so $m$ is absolutely continuous on bounded intervals, and for $t\notin N_1\cup N_2$,
\[
  \dot m(t)=A(t)\Phi(t)\bigl[m(0)+\beta(t)\bigr]+\Phi(t)\Phi^{-1}(t)a(t)=A(t)m(t)+a(t).
\]
Setting $s=t$ in \bk{(6.11)} gives $V=\Phi\,[V(0)+Q]\,\Phi^T$ (this is \bk{(6.14)}), which is likewise absolutely continuous on bounded intervals, and for $t\notin N_1\cup N_2$,
\begin{align*}
  \dot V(t)
  &=A(t)\Phi(t)\bigl[V(0)+Q(t)\bigr]\Phi^T(t)
   +\Phi(t)\,\Phi^{-1}(t)\sigma(t)\sigma^T(t)\bigl(\Phi^{-1}(t)\bigr)^T\Phi^T(t)\\
  &\qquad+\Phi(t)\bigl[V(0)+Q(t)\bigr]\Phi^T(t)A^T(t)\\
  &=A(t)V(t)+\sigma(t)\sigma^T(t)+V(t)A^T(t),
\end{align*}
where we have used $\bigl(\Phi^{-1}(t)\bigr)^T\Phi^T(t)=\bigl(\Phi(t)\Phi^{-1}(t)\bigr)^T=I$. These are \bk{(6.12)} and \bk{(6.13)}, with initial values $m(0)=EX_0$ and $V(0)=E[(X_0-m(0))(X_0-m(0))^T]$.

\rmk{If $A$, $a$ and $\sigma$ are continuous, then $\Phi$, $\beta$ and $Q$ are continuously differentiable, the null sets $N_1$, $N_2$ are empty, and \bk{(6.12)}, \bk{(6.13)} hold at every $t\ge 0$. In general they cannot hold everywhere: for $d=r=1$, $A\equiv 0$, $a\equiv 0$, $\sigma=\ind_{[1,\infty)}$ and $X_0=0$ we get $V(t)=(t-1)^+$, which is not differentiable at $t=1$. The independence of $X_0$ and $W$ assumed in \bk{(6.1)} entered only through the fact that $W$ is a Brownian motion relative to a filtration $\{\F_t\}$ for which $X_0$ is $\F_0$-measurable.}
\end{kex}

\begin{kex}{Problem}{6.6}{}
Throughout, $A$ is a constant $(d\times d)$ matrix all of whose eigenvalues have negative real parts, $\sigma$ is a constant $(d\times r)$ matrix, and $e^{tA}$ is the matrix exponential \bk{(6.17)}; these are the standing assumptions of the paragraph preceding the problem. We write $Q\defeq\sigma\sigma^T$ and $\norm{\cdot}$ for the operator norm on matrices. We use the elementary properties of the exponential series: $e^{sA}e^{tA}=e^{(s+t)A}$ for all $s,t\in\R$ (so that $e^{tA}$ is invertible with inverse $e^{-tA}$), $(e^{tA})^T=e^{tA^T}$, and $\frac{d}{dt}e^{tA}=Ae^{tA}=e^{tA}A$, the series being differentiable term by term because it and its termwise derivative converge locally uniformly in $t$. Integrals of matrix-valued functions are taken entry by entry; consequently, for constant matrices $B$, $C$ of suitable sizes and an integrable matrix function $F$, we have $B\bigl(\int F(s)\,ds\bigr)C=\int BF(s)C\,ds$.

\emph{Step 1: the integral \bk{(6.19)} converges absolutely.} Let $A=PJP^{-1}$ be the (complex) Jordan decomposition of $A$. Then $e^{sA}=Pe^{sJ}P^{-1}$, and each entry of $e^{sJ}$ is either zero or of the form $\frac{s^k}{k!}e^{\mu s}$, where $\mu$ is an eigenvalue of $A$ and $0\le k\le d-1$. Choose $\lambda>0$ with $\operatorname{Re}\mu<-\lambda$ for every eigenvalue $\mu$ of $A$. Since $s^ke^{(\operatorname{Re}\mu+\lambda)s}$ is bounded on $[0,\infty)$, there is a constant $c<\infty$ such that
\begin{equation}\label{eq:5-6-06:a}
  \norm{e^{sA}}\le c\,e^{-\lambda s};\qquad 0\le s<\infty .
\end{equation}
Because $\norm{e^{sA^T}}=\norm{(e^{sA})^T}=\norm{e^{sA}}$, we obtain $\norm{e^{sA}Qe^{sA^T}}\le c^2\norm{Q}e^{-2\lambda s}$, and the continuous function $s\mapsto e^{sA}Qe^{sA^T}$ is (entrywise) absolutely integrable on $[0,\infty)$. Thus
\[
  V=\int_0^\infty e^{sA}Qe^{sA^T}\,ds
\]
is a well-defined $(d\times d)$ matrix.

\emph{Step 2: $V(t)\equiv V$.} Let $V(0)=V$ in \bk{(6.14)$'$} and fix $t\ge 0$. For the first term, we move the constant matrices $e^{tA}$ and $e^{tA^T}$ under the integral sign, use the semigroup property, and substitute $u=t+s$:
\[
  e^{tA}Ve^{tA^T}
  =\int_0^\infty e^{(t+s)A}Qe^{(t+s)A^T}\,ds
  =\int_t^\infty e^{uA}Qe^{uA^T}\,du .
\]
For the second term, in the same way and with the substitution $u=t-s$,
\[
  e^{tA}\biggl[\int_0^t e^{-sA}Qe^{-sA^T}\,ds\biggr]e^{tA^T}
  =\int_0^t e^{(t-s)A}Qe^{(t-s)A^T}\,ds
  =\int_0^t e^{uA}Qe^{uA^T}\,du .
\]
Adding the two expressions, \bk{(6.14)$'$} gives
\[
  V(t)=\int_0^t e^{uA}Qe^{uA^T}\,du+\int_t^\infty e^{uA}Qe^{uA^T}\,du=V=V(0);\qquad 0\le t<\infty .
\]

\emph{Step 3: the equation \bk{(6.20)}.} For an arbitrary initial matrix $V(0)$, the function of \bk{(6.14)$'$} is $V(t)=e^{tA}K(t)e^{tA^T}$, with
\[
  K(t)\defeq V(0)+\int_0^t e^{-sA}Qe^{-sA^T}\,ds .
\]
The integrand is continuous, so $K$ is continuously differentiable with $\dot K(t)=e^{-tA}Qe^{-tA^T}$ by the fundamental theorem of calculus. The product rule and $\frac{d}{dt}e^{tA}=Ae^{tA}$, $\frac{d}{dt}e^{tA^T}=e^{tA^T}A^T$ yield
\[
  \dot V(t)=Ae^{tA}K(t)e^{tA^T}+e^{tA}\bigl(e^{-tA}Qe^{-tA^T}\bigr)e^{tA^T}+e^{tA}K(t)e^{tA^T}A^T
  =AV(t)+V(t)A^T+\sigma\sigma^T,
\]
which is \bk{(6.13)$'$}. With the choice $V(0)=V$ of \bk{(6.19)}, Step 2 shows that $V(t)\equiv V$ is constant, so $\dot V(t)\equiv 0$, and \bk{(6.13)$'$} becomes
\[
  0=AV+VA^T+\sigma\sigma^T,
\]
which is \bk{(6.20)}. Hence $V(t)\equiv V$ is a constant solution of both \bk{(6.13)$'$} and \bk{(6.14)$'$}.

\rmk{The equation \bk{(6.20)} can also be read off from \bk{(6.19)} directly: $\frac{d}{ds}\bigl(e^{sA}Qe^{sA^T}\bigr)=A\,e^{sA}Qe^{sA^T}+e^{sA}Qe^{sA^T}A^T$, and integrating over $[0,T]$ gives $e^{TA}Qe^{TA^T}-Q=A\,V_T+V_TA^T$ with $V_T\defeq\int_0^Te^{sA}Qe^{sA^T}\,ds$. As $T\to\infty$ we have $V_T\to V$ by Step 1 and $e^{TA}Qe^{TA^T}\to 0$ by \eqref{eq:5-6-06:a}, whence $-Q=AV+VA^T$.}
\end{kex}

\begin{kex}{Problem}{6.11}{}
The book's Section 5.9 records that $Z_1,\ldots,Z_n$ are jointly normal and uncorrelated and states their means and variances without computation, leaving the derivation of \bk{(6.28)} from these facts to the reader; we verify the moments and carry out that derivation in full. The assertion \bk{(6.28)} means that the random vector $(X_{t_1},\ldots,X_{t_n})$ has the function on its right-hand side as a density with respect to Lebesgue measure on $\R^n$. We write
\[
  s_i\defeq T-t_i,\qquad \delta_i\defeq t_i-t_{i-1}=s_{i-1}-s_i,\qquad
  v_i\defeq\frac{1}{s_i}-\frac{1}{s_{i-1}}=\frac{\delta_i}{s_is_{i-1}};\qquad i=1,\ldots,n,
\]
so that $T=s_0>s_1>\cdots>s_n>0$ and $\delta_i>0$, $v_i>0$.

\emph{Step 1: the law of $(Z_1,\ldots,Z_n)$.} Let $U_i\defeq X_{t_i}/s_i$; $i=0,1,\ldots,n$, so that $U_0=a/T$ is a constant and $Z_i=U_i-U_{i-1}$. By \bk{Corollary 6.10}, the vector $(X_{t_1},\ldots,X_{t_n})$ is normal, with $E(X_t)=a(1-\frac{t}{T})+b\frac{t}{T}$ by \bk{(6.27)} and $\Cov(X_s,X_t)=s-\frac{st}{T}=\frac{s(T-t)}{T}$ for $0\le s\le t<T$ by \bk{(6.25)}. Consequently $(Z_1,\ldots,Z_n)$, being the image of $(X_{t_1},\ldots,X_{t_n})$ under an affine map, is a normal vector as well. With $q(t)\defeq\frac{1}{T-t}-\frac1T=\frac{t}{T(T-t)}$ we find, for $0\le i\le j\le n$,
\[
  E(U_i)=\frac{a}{T}+\frac{b\,t_i}{T(T-t_i)}=\frac aT+b\,q(t_i),\qquad
  \Cov(U_i,U_j)=\frac{t_i(T-t_j)}{T(T-t_i)(T-t_j)}=q(t_i);
\]
that is, $\Cov(U_i,U_j)=q(t_{i\wedge j})$ for all $0\le i,j\le n$ (both sides vanish if $i\wedge j=0$). Therefore
\begin{align*}
  E(Z_i)&=b\bigl(q(t_i)-q(t_{i-1})\bigr)=b\,v_i,\\
  \Var(Z_i)&=q(t_i)-2q(t_{i-1})+q(t_{i-1})=v_i,\\
  \Cov(Z_i,Z_j)&=q(t_i)-q(t_i)-q(t_{i-1})+q(t_{i-1})=0;\qquad 1\le i<j\le n.
\end{align*}
A normal vector with diagonal covariance matrix has independent components (its characteristic function factorizes), so $Z_1,\ldots,Z_n$ are independent, and $Z_i$ is normal with mean $bv_i$ and variance $v_i>0$. Hence $(Z_1,\ldots,Z_n)$ has the density
\begin{equation}\label{eq:5-6-11:a}
  g(z_1,\ldots,z_n)=\prod_{i=1}^n\frac{1}{\sqrt{2\pi v_i}}\exp\Bigl\{-\frac{(z_i-bv_i)^2}{2v_i}\Bigr\};\qquad (z_1,\ldots,z_n)\in\R^n.
\end{equation}

\emph{Step 2: change of variables.} Let $\varphi:\R^n\to\R^n$ be the affine map $x=(x_1,\ldots,x_n)\mapsto z=(z_1,\ldots,z_n)$ given by
\[
  z_i=\frac{x_i}{s_i}-\frac{x_{i-1}}{s_{i-1}};\qquad i=1,\ldots,n,\quad x_0=a.
\]
It is a bijection, with inverse $x_i=s_i\bigl(\frac aT+z_1+\cdots+z_i\bigr)$, and its Jacobian matrix is lower triangular with diagonal entries $1/s_i$, so that $\det D\varphi=\prod_{i=1}^n s_i^{-1}>0$. Since $(Z_1,\ldots,Z_n)=\varphi(X_{t_1},\ldots,X_{t_n})$, we have for every $\Gamma\in\B(\R^n)$, by the change-of-variables formula for the Lebesgue integral,
\[
  P[(X_{t_1},\ldots,X_{t_n})\in\Gamma]=P[(Z_1,\ldots,Z_n)\in\varphi(\Gamma)]
  =\int_{\varphi(\Gamma)}g(z)\,dz=\int_\Gamma g(\varphi(x))\prod_{i=1}^n\frac{1}{s_i}\,dx .
\]
Thus $(X_{t_1},\ldots,X_{t_n})$ has the density
\begin{equation}\label{eq:5-6-11:b}
  f(x_1,\ldots,x_n)=\prod_{i=1}^n\frac{1}{s_i\sqrt{2\pi v_i}}\cdot
  \exp\Bigl\{-\frac12\sum_{i=1}^n\frac{(z_i-bv_i)^2}{v_i}\Bigr\},\qquad z=\varphi(x),
\end{equation}
and it remains to identify $f$ with the right-hand side of \bk{(6.28)}.

\emph{Step 3: the algebra.} For the constant factor, $s_i^2v_i=\delta_is_i/s_{i-1}$ and the product telescopes:
\begin{equation}\label{eq:5-6-11:c}
  \prod_{i=1}^n\frac{1}{s_i\sqrt{2\pi v_i}}=\prod_{i=1}^n\frac{1}{\sqrt{2\pi\delta_i}}\sqrt{\frac{s_{i-1}}{s_i}}
  =\sqrt{\frac{T}{T-t_n}}\;\prod_{i=1}^n\frac{1}{\sqrt{2\pi(t_i-t_{i-1})}} .
\end{equation}
For the exponent, put $y_i\defeq x_i-b$; $i=0,\ldots,n$. Since $v_i=\frac1{s_i}-\frac1{s_{i-1}}$, we have $z_i-bv_i=\frac{y_i}{s_i}-\frac{y_{i-1}}{s_{i-1}}$, whence
\[
  \frac{(z_i-bv_i)^2}{v_i}=\frac{(s_{i-1}y_i-s_iy_{i-1})^2}{\delta_i\,s_is_{i-1}} .
\]
We claim that
\begin{equation}\label{eq:5-6-11:d}
  \frac{(s_{i-1}y_i-s_iy_{i-1})^2}{\delta_i\,s_is_{i-1}}
  =\frac{(y_i-y_{i-1})^2}{\delta_i}+\frac{y_i^2}{s_i}-\frac{y_{i-1}^2}{s_{i-1}} .
\end{equation}
Indeed, multiplying the right-hand side by $\delta_is_is_{i-1}$ and using $\delta_i=s_{i-1}-s_i$ we obtain
\begin{align*}
  &s_is_{i-1}(y_i-y_{i-1})^2+(s_{i-1}-s_i)s_{i-1}y_i^2-(s_{i-1}-s_i)s_iy_{i-1}^2\\
  &\qquad=s_{i-1}^2y_i^2-2s_is_{i-1}y_iy_{i-1}+s_i^2y_{i-1}^2=(s_{i-1}y_i-s_iy_{i-1})^2 .
\end{align*}
Summing \eqref{eq:5-6-11:d} over $i=1,\ldots,n$, the last two terms telescope, and with $y_i-y_{i-1}=x_i-x_{i-1}$, $y_n=x_n-b$, $y_0=a-b$, $s_n=T-t_n$, $s_0=T$ we get
\begin{equation}\label{eq:5-6-11:e}
  \sum_{i=1}^n\frac{(z_i-bv_i)^2}{v_i}=\sum_{i=1}^n\frac{(x_i-x_{i-1})^2}{t_i-t_{i-1}}+\frac{(x_n-b)^2}{T-t_n}-\frac{(a-b)^2}{T} .
\end{equation}
Substituting \eqref{eq:5-6-11:c} and \eqref{eq:5-6-11:e} into \eqref{eq:5-6-11:b} and recalling \bk{(2.2.6)}, we arrive at
\begin{align*}
  f(x_1,\ldots,x_n)
  &=\prod_{i=1}^n\frac{1}{\sqrt{2\pi(t_i-t_{i-1})}}\exp\Bigl\{-\frac{(x_i-x_{i-1})^2}{2(t_i-t_{i-1})}\Bigr\}\cdot
  \frac{\dfrac{1}{\sqrt{2\pi(T-t_n)}}\exp\Bigl\{-\dfrac{(x_n-b)^2}{2(T-t_n)}\Bigr\}}{\dfrac{1}{\sqrt{2\pi T}}\exp\Bigl\{-\dfrac{(a-b)^2}{2T}\Bigr\}}\\
  &=\prod_{i=1}^np(t_i-t_{i-1};x_{i-1},x_i)\cdot\frac{p(T-t_n;x_n,b)}{p(T;a,b)},
\end{align*}
which is \bk{(6.28)}.

\rmk{For $n=1$ and $t_1=t$, completing the square in $x$ shows that the right-hand side of \bk{(6.28)} is the normal density with mean $a(1-\frac tT)+b\frac tT$ and variance $\frac{t(T-t)}{T}$, in agreement with \bk{(6.27)} and \bk{(6.25)}.}
\end{kex}

\begin{kex}{Problem}{6.13}{}
We fix $a\in\R$, $T>0$ and $0=t_0<t_1<\cdots<t_n<T$, and we write $x=(x_1,\ldots,x_n)$, $x_0=a$, $\xi\defeq(W_{t_1},\ldots,W_{t_n})$. With $p$ the Gaussian kernel \bk{(2.2.6)}, we set, for $x\in\R^n$ and $b\in\R$,
\[
  f(x,b)\defeq\prod_{i=1}^n p(t_i-t_{i-1};x_{i-1},x_i)\cdot p(T-t_n;x_n,b),
  \qquad
  q_b(x)\defeq\frac{f(x,b)}{p(T;a,b)},
\]
so that $q_b(x)$ is the density on the right-hand side of \bk{(6.28)}; it is well defined because $p(T;a,b)>0$. Both $f$ and $(x,b)\mapsto q_b(x)$ are continuous and strictly positive on $\R^{n+1}$.

\emph{The precise statement.} The event $\{W_T=b\}$ has $P^a$-probability zero, so the conditional distribution of $\xi$ given $W_T=b$ is to be understood as a regular conditional distribution of $\xi$ given $W_T$ under $P^a$: a family $\{Q_b\}_{b\in\R}$ of probability measures on $(\R^n,\B(\R^n))$ such that, for every $\Gamma\in\B(\R^n)$, the function $b\mapsto Q_b(\Gamma)$ is Borel-measurable and
\begin{equation}\label{eq:5-6-13:a}
  P^a[\xi\in\Gamma,\ W_T\in\Delta]=E^a\bigl[Q_{W_T}(\Gamma)\,\ind_{\{W_T\in\Delta\}}\bigr];\qquad \Delta\in\B(\R),
\end{equation}
that is, $Q_{W_T}(\Gamma)$ is a version of $P^a[\xi\in\Gamma\mid W_T]$. We shall prove:
\begin{enumerate}
  \item[\textbf{(i)}] the measures $Q_b(\Gamma)\defeq\int_\Gamma q_b(x)\,dx$; $b\in\R$, form such a family;
  \item[\textbf{(ii)}] if $\{\wt{Q}_b\}_{b\in\R}$ is any other such family, then $\wt{Q}_b=Q_b$ for Lebesgue-almost every $b\in\R$.
\end{enumerate}
Together, (i) and (ii) are the assertion of the problem. Since, by \bk{Definition 6.12}, a Brownian bridge from $a$ to $b$ on $[0,T]$ is an a.s.\ continuous process whose finite-dimensional distributions at times $t_1<\cdots<t_n$ in $(0,T)$ are the measures $Q_b$, they also justify the italicized sentence of the statement.

\emph{Step 1: the joint density of $(\xi,W_T)$ under $P^a$.} We put $t_{n+1}\defeq T$. By \bk{Definition 2.5.8}, we have $P^a[W_0=a]=1$ and, under $P^a$, the process $W$ is a Brownian motion; hence the increments $D_i\defeq W_{t_i}-W_{t_{i-1}}$; $1\le i\le n+1$, are independent under $P^a$ (for each $i$, the increment $D_i$ is independent of $\F_{t_{i-1}}$, with respect to which $D_1,\ldots,D_{i-1}$ are measurable), and $D_i$ is normal with mean zero and variance $t_i-t_{i-1}>0$. Therefore the vector $D=(D_1,\ldots,D_{n+1})$ has under $P^a$ the density $\prod_{i=1}^{n+1}p(t_i-t_{i-1};0,d_i)$ with respect to Lebesgue measure on $\R^{n+1}$. Now $P^a$-a.s.
\[
  (W_{t_1},\ldots,W_{t_n},W_T)=\Lambda(D),\qquad
  \Lambda(d_1,\ldots,d_{n+1})\defeq\bigl(a+d_1,\ a+d_1+d_2,\ \ldots,\ a+d_1+\cdots+d_{n+1}\bigr),
\]
and $\Lambda$ is an affine bijection of $\R^{n+1}$ whose linear part is lower triangular with unit diagonal, so it has Jacobian determinant $1$, and its inverse is $(x_1,\ldots,x_{n+1})\mapsto(x_1-x_0,x_2-x_1,\ldots,x_{n+1}-x_n)$ with $x_0=a$. By the change-of-variables formula and $p(t;0,y-z)=p(t;z,y)$, the vector $(\xi,W_T)$ has under $P^a$ the density
\[
  \prod_{i=1}^{n+1}p(t_i-t_{i-1};x_{i-1},x_i)=f(x,b)\qquad(x_{n+1}=b)
\]
with respect to Lebesgue measure on $\R^{n+1}$; in other words,
\begin{equation}\label{eq:5-6-13:b}
  P^a[\xi\in\Gamma,\ W_T\in\Delta]=\int_\Delta\int_\Gamma f(x,b)\,dx\,db;\qquad\Gamma\in\B(\R^n),\ \Delta\in\B(\R),
\end{equation}
the order of integration being immaterial by Tonelli's theorem. Moreover, $W_T-a=W_T-W_0$ is $P^a$-a.s.\ normal with mean zero and variance $T$, so that
\begin{equation}\label{eq:5-6-13:c}
  P^a[W_T\in\Delta]=\int_\Delta p(T;a,b)\,db;\qquad\Delta\in\B(\R).
\end{equation}

\emph{Step 2: $Q_b$ is a probability measure for every $b\in\R$.} We have to show that
\begin{equation}\label{eq:5-6-13:d}
  \int_{\R^n}f(x,b)\,dx=p(T;a,b)\qquad\text{for every }b\in\R .
\end{equation}
Taking $\Gamma=\R^n$ in \eqref{eq:5-6-13:b} and comparing with \eqref{eq:5-6-13:c}, we see that the nonnegative Borel functions $b\mapsto\int_{\R^n}f(x,b)\,dx$ (measurable by Tonelli's theorem) and $b\mapsto p(T;a,b)$ have the same integral over every Borel set $\Delta$; hence they agree for Lebesgue-almost every $b$. The right-hand side of \eqref{eq:5-6-13:d} is continuous in $b$. So is the left-hand side: $b\mapsto f(x,b)$ is continuous for each $x$, and
\begin{equation}\label{eq:5-6-13:e}
  0<f(x,b)\le\frac{1}{\sqrt{2\pi(T-t_n)}}\,g(x),\qquad g(x)\defeq\prod_{i=1}^n p(t_i-t_{i-1};x_{i-1},x_i),
\end{equation}
where $g$ does not depend on $b$ and is integrable over $\R^n$ (it is the density of $\xi$ under $P^a$, by the argument of Step 1 applied to $D_1,\ldots,D_n$); thus the dominated convergence theorem applies. Two continuous functions which agree Lebesgue-almost everywhere agree everywhere, because the complement of a Lebesgue-null set is dense. This proves \eqref{eq:5-6-13:d}, and consequently $Q_b(\R^n)=\int_{\R^n}q_b(x)\,dx=1$ for every $b\in\R$.

\emph{Step 3: proof of (i).} Let $\Gamma\in\B(\R^n)$. The function $(x,b)\mapsto q_b(x)\ind_\Gamma(x)$ is nonnegative and Borel-measurable on $\R^{n+1}$, so $b\mapsto Q_b(\Gamma)=\int_\Gamma q_b(x)\,dx$ is Borel-measurable by Tonelli's theorem. For $\Delta\in\B(\R)$, the bounded Borel function $h(b)\defeq Q_b(\Gamma)\ind_\Delta(b)$ satisfies, by \eqref{eq:5-6-13:c} (which says that $W_T$ has density $p(T;a,\cdot)$ under $P^a$), the definition of $q_b$, and \eqref{eq:5-6-13:b},
\[
  E^a\bigl[Q_{W_T}(\Gamma)\ind_{\{W_T\in\Delta\}}\bigr]
  =\int_\Delta Q_b(\Gamma)\,p(T;a,b)\,db
  =\int_\Delta\int_\Gamma f(x,b)\,dx\,db
  =P^a[\xi\in\Gamma,\ W_T\in\Delta].
\]
This is \eqref{eq:5-6-13:a}. Since $Q_{W_T}(\Gamma)$ is $\sigma(W_T)$-measurable and every event of $\sigma(W_T)$ is of the form $\{W_T\in\Delta\}$ with $\Delta\in\B(\R)$, the random variable $Q_{W_T}(\Gamma)$ is a version of $P^a[\xi\in\Gamma\mid W_T]$, and (i) is proved.

\emph{Step 4: proof of (ii).} Let $\{\wt{Q}_b\}_{b\in\R}$ be another family of probability measures on $(\R^n,\B(\R^n))$ such that $b\mapsto\wt{Q}_b(\Gamma)$ is Borel-measurable and \eqref{eq:5-6-13:a} holds with $\wt{Q}$ in place of $Q$, for every $\Gamma\in\B(\R^n)$. Let $\mathscr{R}$ be the countable collection of all rectangles $\prod_{i=1}^n(\alpha_i,\beta_i]$ with rational $\alpha_i<\beta_i$, together with the empty set; $\mathscr{R}$ is closed under finite intersections and generates $\B(\R^n)$. Fix $\Gamma\in\mathscr{R}$ and let $\Delta_\Gamma^{\pm}\defeq\{b\in\R;\ \pm(\wt{Q}_b(\Gamma)-Q_b(\Gamma))>0\}\in\B(\R)$. Subtracting the two versions of \eqref{eq:5-6-13:a} with $\Delta=\Delta_\Gamma^{\pm}$ and using \eqref{eq:5-6-13:c}, we obtain
\[
  \int_{\Delta_\Gamma^{\pm}}\bigl(\wt{Q}_b(\Gamma)-Q_b(\Gamma)\bigr)\,p(T;a,b)\,db=0 .
\]
The integrand has a strict sign on $\Delta_\Gamma^{\pm}$, because $p(T;a,b)>0$ for every $b$; hence $\Delta_\Gamma^{+}$ and $\Delta_\Gamma^{-}$ have Lebesgue measure zero. The set $N\defeq\bigcup_{\Gamma\in\mathscr{R}}(\Delta_\Gamma^+\cup\Delta_\Gamma^-)$ is a countable union of Lebesgue-null sets, hence Lebesgue-null, and for $b\notin N$ the probability measures $\wt{Q}_b$ and $Q_b$ agree on the $\pi$-system $\mathscr{R}$. By Dynkin's $\pi$-$\lambda$ theorem, $\wt{Q}_b=Q_b$ on $\B(\R^n)$ for every $b\notin N$. This proves (ii) and completes the solution.

\rmk{The qualification ``for Lebesgue-almost every $b$'' cannot be removed from (ii): changing $Q_b$ for $b$ in a Lebesgue-null set affects neither the measurability in $b$ nor \eqref{eq:5-6-13:a}. The family $\{Q_b\}$ of (i) is, however, singled out among all versions by its continuity in $b$, and it is the limit of elementary conditional probabilities for \emph{every} $b\in\R$:
\[
  \lim_{\eps\downarrow0}P^a\bigl[\xi\in\Gamma\bigm|\abs{W_T-b}<\eps\bigr]=Q_b(\Gamma);\qquad\Gamma\in\B(\R^n).
\]
Indeed, by \eqref{eq:5-6-13:b} and \eqref{eq:5-6-13:c} the conditional probability on the left equals $\int_{b-\eps}^{b+\eps}F(c)\,dc\big/\int_{b-\eps}^{b+\eps}p(T;a,c)\,dc$ with $F(c)\defeq\int_\Gamma f(x,c)\,dx$; the function $F$ is continuous by the bound \eqref{eq:5-6-13:e} and dominated convergence, $p(T;a,\cdot)$ is continuous and positive, and so the quotient converges to $F(b)/p(T;a,b)=Q_b(\Gamma)$ as $\eps\downarrow0$.}
\end{kex}

\begin{kex}{Problem}{6.14}{}
Here $T>0$ and $a,b\in\R$ are fixed, and we write $B\defeq B^{a\to b}$, $m(t)\defeq a(1-t/T)+b\,t/T$ as in \bk{(6.27)} and $\rho(s,t)\defeq(s\wedge t)-st/T$ as in \bk{(6.25)}. According to \bk{Definition 6.12}, we have to show that $B$ is almost surely continuous on $[0,T]$ and that, for $0<t_1<\cdots<t_n<T$, the distribution of $(B_{t_1},\ldots,B_{t_n})$ is the one given by \bk{(6.28)}.

\emph{Step 1: continuity.} Every path of $W$ is continuous (\bk{Definition 2.1.1}), and for every $\omega$ the map $t\mapsto B_t(\omega)$ is the sum of $t\mapsto W_t(\omega)$ and of the affine function $t\mapsto a+(b-a-W_T(\omega))\,t/T$. Hence every path of $B$ is continuous on $[0,T]$. \rmk{Since $W_0=0$ a.s., we also have $B_0=a$ a.s.\ and $B_T=b$ surely, as befits a bridge from $a$ to $b$; the formula \bk{(6.28)}, however, involves only times in the open interval $(0,T)$.}

\emph{Step 2: $B$ is Gaussian, with mean function $m$ and covariance function $\rho$.} Let $0<t_1<\cdots<t_n<T$ and put $t_0\defeq0$, $t_{n+1}\defeq T$. By \bk{Definition 2.1.1}, the increments $D_i\defeq W_{t_i}-W_{t_{i-1}}$; $1\le i\le n+1$, are independent (each $D_i$ is independent of $\F_{t_{i-1}}$, with respect to which $D_1,\ldots,D_{i-1}$ are measurable) and normal with mean zero; thus $(D_1,\ldots,D_{n+1})$ is a Gaussian vector. Because $W_0=0$ a.s., we have $W_{t_k}=D_1+\cdots+D_k$ a.s.\ for $1\le k\le n+1$, and therefore
\[
  B_{t_k}=m(t_k)+\sum_{i=1}^{k}D_i-\frac{t_k}{T}\sum_{i=1}^{n+1}D_i\quad\text{a.s.};\qquad 1\le k\le n.
\]
The vector $(B_{t_1},\ldots,B_{t_n})$ is thus almost surely equal to an affine image of a Gaussian vector, so it is itself Gaussian (possibly degenerate, a priori). Its mean vector is $(m(t_1),\ldots,m(t_n))$, because $E(W_t)=0$ for every $t\ge0$. For the covariances, we recall that $E(W_sW_t)=s\wedge t$ for $s,t\ge0$: if $s\le t$, then $E(W_sW_t)=E(W_s^2)+E[W_s(W_t-W_s)]=s$, the increment $W_t-W_s$ having mean zero and being independent of $W_s$. Consequently, for $0\le s,t\le T$,
\begin{align*}
  \Cov(B_s,B_t)
  &=E\Bigl[\Bigl(W_s-\frac{s}{T}W_T\Bigr)\Bigl(W_t-\frac{t}{T}W_T\Bigr)\Bigr]\\
  &=E(W_sW_t)-\frac{t}{T}E(W_sW_T)-\frac{s}{T}E(W_tW_T)+\frac{st}{T^2}E(W_T^2)\\
  &=(s\wedge t)-\frac{st}{T}-\frac{st}{T}+\frac{st}{T}=\rho(s,t).
\end{align*}

\emph{Step 3: identification of the finite-dimensional distributions.} This step is the remark following \bk{Definition 6.12}; we give the argument, which compares $B$ with the process of \bk{(6.26)}. Since \bk{Corollary 6.10} and \bk{Problem 6.11} belong to a setting in which the filtration satisfies the usual conditions, whereas $\{\F_t\}$ is not assumed to satisfy them, we first produce a standard Brownian motion on a filtered space that does. Let $\{\F^W_t\}$ be the filtration generated by $W$. As $\F^W_s\subseteq\F_s$, the increment $W_t-W_s$ is independent of $\F^W_s$ for $0\le s<t$, so $\{W_t,\F^W_t;\ t\ge0\}$ is a one-dimensional Brownian motion on $(\Omega,\F^W_\infty,P)$, with initial distribution $\mu\defeq\delta_0$ because $W_0=0$ a.s. Let $\{\F^\mu_t\}$ be its augmented filtration and $(\Omega,\F^\mu,P^\mu)$ the completed space of \bk{Problem 2.7.3}, on which $P^\mu$ extends $P|_{\F^W_\infty}$. By \bk{Theorem 2.7.9}, $\{W_t,\F^\mu_t;\ t\ge0\}$ is still a standard Brownian motion. The filtration $\{\F^\mu_t\}$ is right-continuous by \bk{Proposition 2.7.7}, which applies because a Brownian motion is a strong Markov process (\bk{Theorem 2.6.15}). Moreover, $\F^\mu_0$ contains every $P^\mu$-negligible event of $\F^\mu$: if $F\in\F^\mu$ and $P^\mu(F)=0$, choose $G\in\F^W_\infty$ with $F\triangle G\in\mathscr{N}^\mu$ (\bk{Problem 2.7.3}); then $P(G)=P^\mu(F)=0$, and $F\triangle G\subseteq H$ for some $H\in\F^W_\infty$ with $P(H)=0$, so that $F\subseteq G\cup H\in\F^W_\infty$ with $P(G\cup H)=0$; that is, $F\in\mathscr{N}^\mu\subseteq\F^\mu_0$. Thus $\{\F^\mu_t\}$ satisfies the usual conditions (\bk{Definition 1.2.25}).

Now let $X$ be the process \bk{(6.26)} built from the Brownian motion $\{W_t,\F^\mu_t\}$ on $(\Omega,\F^\mu,P^\mu)$. By \bk{Corollary 6.10}, $X$ is a Gaussian process under $P^\mu$ with mean function $m$ and covariance function $\rho$; by \bk{Problem 6.11}, the $P^\mu$-distribution of $(X_{t_1},\ldots,X_{t_n})$ is given by \bk{(6.28)} whenever $0<t_1<\cdots<t_n<T$. Only this distribution will be used below; in particular it is irrelevant that $X$ lives on the completed space. Now the distribution of a Gaussian vector in $\R^n$ with mean vector $\gamma$ and covariance matrix $\Sigma$ is determined by $\gamma$ and $\Sigma$, since its characteristic function is $u\mapsto\exp\bigl(i\,u\cdot\gamma-\tfrac12\,u\cdot\Sigma u\bigr)$ and a probability measure on $\R^n$ is determined by its characteristic function. By Step 2, the $P$-distribution of $(B_{t_1},\ldots,B_{t_n})$ and the $P^\mu$-distribution of $(X_{t_1},\ldots,X_{t_n})$ are Gaussian with the same mean vector $(m(t_k))_{1\le k\le n}$ and the same covariance matrix $(\rho(t_j,t_k))_{1\le j,k\le n}$; hence they coincide, and the distribution of $(B_{t_1},\ldots,B_{t_n})$ is the one given by \bk{(6.28)}.

By Steps 1 and 3 and \bk{Definition 6.12}, the process $B^{a\to b}$ is a Brownian bridge from $a$ to $b$ on $[0,T]$.
\end{kex}

\begin{kex}{Exercise}{6.16}{}
\emph{Setting.} On a probability space $(\Omega,\F,P)$ with a filtration $\{\F_t\}$ satisfying the usual conditions, $W$ is a standard Brownian motion relative to $\{\F_t\}$, and $X$ is a continuous, adapted process which satisfies \bk{(5.79)},
\[
  X_t=x+\mu\int_0^tX_s\,ds+\nu\int_0^tX_s\,dW_s;\qquad 0\le t<\infty,
\]
almost surely, where $x>0$ and $\mu,\nu$ are real constants. We shall prove that, for every real $k$ (in particular for $k>1$), the process $Y=X^k$ is well defined, strictly positive, and satisfies
\begin{equation}\label{eq:5-6-16:sde}
  dY_t=\bigl(k\mu+\tfrac12k(k-1)\nu^2\bigr)Y_t\,dt+k\nu\,Y_t\,dW_t,\qquad Y_0=x^k,
\end{equation}
that is, $Y$ is again a geometric Brownian motion, with $x,\mu,\nu$ replaced by $x^k$, $k\mu+\frac12k(k-1)\nu^2$, $k\nu$; and that
\begin{equation}\label{eq:5-6-16:moment}
  E\bigl(X_t^k\bigr)=x^k\exp\Bigl\{\bigl(k\mu+\tfrac12k(k-1)\nu^2\bigr)t\Bigr\};\qquad 0\le t<\infty .
\end{equation}

\emph{Step 1: $X$ is strictly positive.} For a non-integer $k$ the power $X_t^k$ makes sense only if $X_t>0$, and the function $y\mapsto y^k$ is of class $C^2$ only on $(0,\infty)$ in general; so we first show that
\begin{equation}\label{eq:5-6-16:explicit}
  X_t=x\exp\bigl\{(\mu-\tfrac12\nu^2)t+\nu W_t\bigr\};\qquad 0\le t<\infty,\quad\text{a.s.}
\end{equation}
(This is the formula quoted in \bk{Exercise 5.31}, and \bk{(6.34)}; we verify directly that \emph{every} solution of \bk{(5.79)} has this form.) The process $X$ is a continuous semimartingale $X=x+M+B$ with $M_t=\nu\int_0^tX_s\,dW_s\in\cM^{c,\loc}$ and $B_t=\mu\int_0^tX_s\,ds$ adapted, continuous and of bounded variation; moreover $\qv{W}_t=t$ and $\qv{W,M}_t=\nu\int_0^tX_s\,ds$. The function
\[
  f(t,w,y)\defeq y\,e(t,w),\qquad e(t,w)\defeq\exp\bigl\{-(\mu-\tfrac12\nu^2)t-\nu w\bigr\},
\]
is of class $C^{1,2}$ on $[0,\infty)\times\R^2$, with
\[
  \frac{\partial f}{\partial t}=-(\mu-\tfrac12\nu^2)\,y\,e,\quad
  \frac{\partial f}{\partial w}=-\nu\,y\,e,\quad
  \frac{\partial f}{\partial y}=e,\quad
  \frac{\partial^2 f}{\partial w^2}=\nu^2\,y\,e,\quad
  \frac{\partial^2 f}{\partial w\,\partial y}=-\nu\,e,\quad
  \frac{\partial^2 f}{\partial y^2}=0 .
\]
We apply It\^o's rule, \bk{Theorem 3.3.6}, to $f$ and the two-dimensional semimartingale $(W,X)$. Writing $e_s\defeq e(s,W_s)$, we obtain, almost surely for all $0\le t<\infty$,
\begin{align*}
  X_te_t-x
  &=\int_0^t\Bigl[-(\mu-\tfrac12\nu^2)X_se_s+\mu X_se_s+\tfrac12\nu^2X_se_s-\nu e_s\cdot\nu X_s\Bigr]ds
    +\int_0^t\bigl[-\nu X_se_s+e_s\,\nu X_s\bigr]dW_s\\
  &=0 .
\end{align*}
(In the $ds$-integral the four terms come, in this order, from $\partial f/\partial t$, from $dB$, from $\frac12(\partial^2f/\partial w^2)\,d\qv{W}$ and from $(\partial^2f/\partial w\,\partial y)\,d\qv{W,M}$; in the $dW$-integral, from $(\partial f/\partial w)\,dW$ and $(\partial f/\partial y)\,dM$.) Hence $X_t=x/e_t$, which is \eqref{eq:5-6-16:explicit}. In particular, after discarding a $P$-null event, $X_t\in(0,\infty)$ for all $t$, and $Y_t=X_t^k$ is a well-defined, continuous, adapted, strictly positive process for every real $k$.

\emph{Step 2: the equation for $Y$.} By \eqref{eq:5-6-16:explicit}, $Y_t=g(t,W_t)$ with
\[
  g(t,w)\defeq x^k\exp\bigl\{k(\mu-\tfrac12\nu^2)t+k\nu w\bigr\},
\]
a function of class $C^{1,2}$ on $[0,\infty)\times\R$ with $\partial g/\partial t=k(\mu-\frac12\nu^2)g$, $\partial g/\partial w=k\nu g$ and $\partial^2g/\partial w^2=k^2\nu^2g$. \bk{Theorem 3.3.6}, applied to $g$ and $W$, gives almost surely
\[
  Y_t=x^k+\int_0^t\bigl[k(\mu-\tfrac12\nu^2)+\tfrac12k^2\nu^2\bigr]Y_s\,ds+k\nu\int_0^tY_s\,dW_s;\qquad 0\le t<\infty .
\]
Since $k(\mu-\frac12\nu^2)+\frac12k^2\nu^2=k\mu+\frac12k(k-1)\nu^2$, this is \eqref{eq:5-6-16:sde}, i.e.
\begin{equation}\label{eq:5-6-16:int}
  Y_t=x^k+a_k\int_0^tY_s\,ds+k\nu\int_0^tY_s\,dW_s,\qquad a_k\defeq k\mu+\tfrac12k(k-1)\nu^2 .
\end{equation}
(The same equation results from the formal computation $dY=kX^{k-1}\,dX+\frac12k(k-1)X^{k-2}\,d\qv{X}$ with $d\qv{X}_t=\nu^2X_t^2\,dt$, which is It\^o's rule for the function $y\mapsto y^k$; Step~1 is what makes it legitimate, because $X$ never leaves $(0,\infty)$, where this function is $C^2$.)

\emph{Step 3: an a priori bound.} We claim that for every real $k$,
\begin{equation}\label{eq:5-6-16:bound}
  E\bigl(X_t^k\bigr)\le x^k e^{\abs{a_k}t};\qquad 0\le t<\infty .
\end{equation}
Fix $k$, and for integers $n>x^k$ let $T_n\defeq\inf\{t\ge0;\ Y_t\ge n\}$, a stopping time by \bk{Problem 1.2.7}. By continuity, $0<Y_{t\wedge T_n}\le n$ for all $t$, and $T_n\uparrow\infty$ a.s.\ as $n\to\infty$, because the continuous function $s\mapsto Y_s(\omega)$ is bounded on every compact interval. Evaluating \eqref{eq:5-6-16:int} at $t\wedge T_n$ gives
\[
  Y_{t\wedge T_n}=x^k+a_k\int_0^tY_s\ind_{\{s\le T_n\}}\,ds+k\nu\int_0^tY_s\ind_{\{s\le T_n\}}\,dW_s .
\]
The integrand $Y_s\ind_{\{s\le T_n\}}$ is progressively measurable and bounded by $n$, so the stochastic integral is a martingale in $\cM_2^c$ starting at zero and has expectation zero. Taking expectations, using $0\le Y_s\ind_{\{s\le T_n\}}\le Y_{s\wedge T_n}$ and Tonelli's theorem, we find that $g_n(t)\defeq E(Y_{t\wedge T_n})$ satisfies
\[
  0\le g_n(t)\le x^k+\abs{a_k}\int_0^tg_n(s)\,ds;\qquad 0\le t<\infty .
\]
The function $g_n$ is continuous by the bounded convergence theorem. The Gronwall inequality of \bk{Problem 2.7} therefore yields $g_n(t)\le x^k+\abs{a_k}\int_0^tx^ke^{\abs{a_k}(t-s)}\,ds=x^ke^{\abs{a_k}t}$. Since $Y_{t\wedge T_n}\to Y_t$ a.s.\ as $n\to\infty$, Fatou's lemma gives \eqref{eq:5-6-16:bound}.

\emph{Step 4: computation of the moment.} Fix $k$. Applying \eqref{eq:5-6-16:bound} with $2k$ in place of $k$, we have $E(Y_s^2)=E(X_s^{2k})\le x^{2k}e^{\abs{a_{2k}}s}$, and thus, by Tonelli's theorem,
\[
  E\int_0^tY_s^2\,ds=\int_0^tE(Y_s^2)\,ds\le x^{2k}\,t\,e^{\abs{a_{2k}}t}<\infty;\qquad 0\le t<\infty .
\]
Being continuous and adapted, $Y$ is progressively measurable; consequently $Y\in\cL^*$, the stochastic integral $\int_0^{\cdot}Y_s\,dW_s$ is a martingale in $\cM_2^c$, and $E\int_0^tY_s\,dW_s=0$. Moreover $E\int_0^tY_s\,ds=\int_0^tE(Y_s)\,ds\le x^k\,t\,e^{\abs{a_k}t}<\infty$ by \eqref{eq:5-6-16:bound} and Tonelli's theorem. Taking expectations in \eqref{eq:5-6-16:int}, we see that $m(t)\defeq E(Y_t)$ is finite and satisfies
\[
  m(t)=x^k+a_k\int_0^tm(s)\,ds;\qquad 0\le t<\infty .
\]
The function $m$ is Borel measurable (Tonelli) and bounded on compact intervals by \eqref{eq:5-6-16:bound}; the identity then shows first that $m$ is continuous, and next that $m$ is continuously differentiable with $m'(t)=a_km(t)$, $m(0)=x^k$. Hence $\frac{d}{dt}\bigl(e^{-a_kt}m(t)\bigr)=0$, so $m(t)=x^ke^{a_kt}$, which is \eqref{eq:5-6-16:moment}.

\rmk{The restriction $k>1$ plays no role: \eqref{eq:5-6-16:sde} and \eqref{eq:5-6-16:moment} hold for every real $k$. For $k=1$ we recover $E(X_t)=xe^{\mu t}$, and for $k=2$ we get $E(X_t^2)=x^2e^{(2\mu+\nu^2)t}$, so that $\Var(X_t)=x^2e^{2\mu t}\bigl(e^{\nu^2t}-1\bigr)$. As an independent check, \eqref{eq:5-6-16:explicit} and the formula $E\,e^{\lambda W_t}=e^{\lambda^2t/2}$ for the moment generating function of the normal random variable $W_t$ give directly
\[
  E\bigl(X_t^k\bigr)=x^ke^{k(\mu-\frac12\nu^2)t}\,E\,e^{k\nu W_t}=x^k\exp\bigl\{k(\mu-\tfrac12\nu^2)t+\tfrac12k^2\nu^2t\bigr\}
  =x^ke^{a_kt},
\]
in agreement with \eqref{eq:5-6-16:moment}.}
\end{kex}

\begin{kex}{Exercise}{6.17}{}
Throughout, $W=\{W_t=(W^{(1)}_t,\ldots,W^{(d)}_t),\F_t;\,0\le t<\infty\}$ is a $d$-dimensional Brownian motion (\bk{Definition 2.5.1}) with $W_0=0$ a.s., and \bk{(6.26)}, \bk{(6.29)} are read as equations between vectors, i.e., componentwise: the $k$-th component of $X$ (resp.\ of $B^{a\to b}$) is given by \bk{(6.26)} (resp.\ \bk{(6.29)}) with $a_k$, $b_k$, $W^{(k)}$ in place of $a$, $b$, $W$. We write $p_1$ for the one-dimensional kernel \bk{(2.2.6)}, so that
\begin{equation}\label{eq:5-6-17:a}
  p(t;x,y)=\prod_{k=1}^d p_1(t;x_k,y_k);\qquad x=(x_1,\ldots,x_d),\ y=(y_1,\ldots,y_d)\in\R^d,\ t>0 .
\end{equation}
The requirement on the finite-dimensional distributions of a $d$-dimensional bridge $Y$ is that, for $0=t_0<t_1<\cdots<t_n<T$ and $x_0=a$, the random vector $(Y_{t_1},\ldots,Y_{t_n})$ have the density
\begin{equation}\label{eq:5-6-17:b}
  f^{a\to b}_{t_1,\ldots,t_n}(x_1,\ldots,x_n)\defeq\prod_{i=1}^n p(t_i-t_{i-1};x_{i-1},x_i)\cdot\frac{p(T-t_n;x_n,b)}{p(T;a,b)};
  \qquad (x_1,\ldots,x_n)\in(\R^d)^n,
\end{equation}
with respect to Lebesgue measure on $(\R^d)^n$. We let $m(t)\defeq a(1-\frac tT)+b\frac tT\in\R^d$ and $\rho(s,t)\defeq(s\wedge t)-\frac{st}{T}$; $0\le s,t\le T$, as in \bk{(6.27)}, \bk{(6.25)}.

\emph{Step 0: preliminaries.} \textbf{(a)} \emph{The Gaussian structure of $W$.} Let $0=s_0<s_1<\cdots<s_m$. By \bk{Definition 2.5.1}, each increment $D_i\defeq W_{s_i}-W_{s_{i-1}}$ is independent of $\F_{s_{i-1}}$, with respect to which $D_1,\ldots,D_{i-1}$ are measurable, and is normal with mean zero and covariance matrix $(s_i-s_{i-1})I_d$. Hence $D_1,\ldots,D_m$ are independent, and $(D_1,\ldots,D_m)$ is a Gaussian vector in $\R^{dm}$. Since $W_0=0$ a.s., $W_{s_i}=D_1+\cdots+D_i$ a.s., so $(W_{s_1},\ldots,W_{s_m})$ is a.s.\ equal to a linear image of a Gaussian vector and is therefore Gaussian, with mean zero. For $0\le s\le t$ and $1\le j,k\le d$,
\begin{equation}\label{eq:5-6-17:c}
  E\bigl(W^{(j)}_sW^{(k)}_t\bigr)=E\bigl(W^{(j)}_sW^{(k)}_s\bigr)+E\bigl[W^{(j)}_s\bigl(W^{(k)}_t-W^{(k)}_s\bigr)\bigr]=\delta_{jk}\,s=\delta_{jk}\,(s\wedge t),
\end{equation}
because $W_s=W_s-W_0$ a.s.\ has covariance matrix $sI_d$ and $W_t-W_s$ has mean zero and is independent of $W_s$. Thus the family $\{W^{(k)}_t;\,t\ge0,\,1\le k\le d\}$ is centered Gaussian, and components with different indices $k$ are uncorrelated. Since jointly Gaussian vectors which are mutually uncorrelated are independent (their joint characteristic function factorizes), \emph{for any finite set of times $s_1,\ldots,s_m$ the $d$ random vectors $(W^{(k)}_{s_1},\ldots,W^{(k)}_{s_m})$; $k=1,\ldots,d$, are independent.} Moreover, each $W^{(k)}=\{W^{(k)}_t,\F_t;\,0\le t<\infty\}$ is a standard, one-dimensional Brownian motion in the sense of \bk{Definition 2.1.1}: it is continuous and adapted, $W^{(k)}_0=0$ a.s., and $W^{(k)}_t-W^{(k)}_s$ is independent of $\F_s$ and normal with mean zero and variance $t-s$.

\textbf{(b)} \emph{The law of an a.s.\ continuous process.} Let $C[0,T]^d$ be the space of continuous functions $y\colon[0,T]\to\R^d$ with the supremum norm; its Borel $\sigma$-field $\B(C[0,T]^d)$ is generated by the coordinate maps $y\mapsto y(t)$; $0\le t\le T$ (the proof is that of \bk{Problem 2.4.2}). If $Y=\{Y_t;\,0\le t\le T\}$ is an a.s.\ continuous process, i.e., there is an event $\Omega_0$ with $P(\Omega_0)=1$ on which all paths are continuous, we regard $Y$ as the map $\omega\mapsto Y_\cdot(\omega)$ for $\omega\in\Omega_0$, $\omega\mapsto0$ for $\omega\notin\Omega_0$, from $\Omega$ into $C[0,T]^d$. This map is measurable, because its composition with every coordinate map is, and the measure it induces on $\B(C[0,T]^d)$, the \emph{law} of $Y$, does not depend on the choice of $\Omega_0$. For such a $Y$, the expression $EF(Y_\cdot)$ is the integral of $F$ with respect to the law of $Y$.

\begin{lem}
Two almost surely continuous, $\R^d$-valued processes $Y$, $Y'$ on $[0,T]$ $($possibly defined on different probability spaces$)$, for which $(Y_{t_1},\ldots,Y_{t_n})$ and $(Y'_{t_1},\ldots,Y'_{t_n})$ have the same distribution whenever $0<t_1<\cdots<t_n<T$, have the same law on $C[0,T]^d$.
\end{lem}
\begin{proof}
Let $\cC$ be the class of sets $\{y\in C[0,T]^d;\ (y(t_1),\ldots,y(t_n))\in\Gamma\}$ with $n\ge1$, $0<t_1<\cdots<t_n<T$ and $\Gamma\in\B((\R^d)^n)$. It is closed under finite intersections, and the two laws agree on it by hypothesis. The $\sigma$-field $\sigma(\cC)$ makes every coordinate map $y\mapsto y(t)$; $0<t<T$, measurable, hence also the maps $y\mapsto y(0)=\lim_{j\to\infty}y(T/j)$ and $y\mapsto y(T)=\lim_{j\to\infty}y(T-T/j)$; thus $\sigma(\cC)=\B(C[0,T]^d)$. By the Dynkin System Theorem (\bk{Theorem 2.1.3}), the two laws agree on $\B(C[0,T]^d)$.
\end{proof}

Consequently, all $d$-dimensional Brownian bridges from $a$ to $b$ on $[0,T]$ have one and the same law on $C[0,T]^d$, which we denote by $Q^{a\to b}$; part (i) shows that such processes exist. In (ii) and (iii) the symbol $B^{a\to b}$ may thus stand for the process \bk{(6.29)} or for any other bridge from $a$ to $b$: the assertions concern $Q^{a\to b}$ and $Q^{b\to a}$ only.

\textbf{(i)} \emph{The process $B^{a\to b}$ of \bk{(6.29)}.} Every path of $B\defeq B^{a\to b}$ is continuous on $[0,T]$, being the sum of a path of $W$ and of an affine function of $t$. By Step 0(a), $W^{(k)}$ is a standard one-dimensional Brownian motion, so by \bk{Problem 6.14} the component $B^{(k)}$ is a one-dimensional Brownian bridge from $a_k$ to $b_k$ on $[0,T]$: for $0=t_0<t_1<\cdots<t_n<T$, the vector $\xi^k\defeq(B^{(k)}_{t_1},\ldots,B^{(k)}_{t_n})$ has the density
\[
  g_k(x_{1,k},\ldots,x_{n,k})\defeq\prod_{i=1}^n p_1(t_i-t_{i-1};x_{i-1,k},x_{i,k})\cdot\frac{p_1(T-t_n;x_{n,k},b_k)}{p_1(T;a_k,b_k)};\qquad x_{0,k}\defeq a_k,
\]
on $\R^n$. Now $\xi^k$ is a Borel function of $(W^{(k)}_{t_1},\ldots,W^{(k)}_{t_n},W^{(k)}_T)$, so $\xi^1,\ldots,\xi^d$ are independent by Step 0(a). Therefore $(\xi^1,\ldots,\xi^d)$ has the density $\prod_{k=1}^dg_k(x_{1,k},\ldots,x_{n,k})$ on $(\R^n)^d$, and by \eqref{eq:5-6-17:a}, with $x_i\defeq(x_{i,1},\ldots,x_{i,d})$,
\begin{equation}\label{eq:5-6-17:d}
  \prod_{k=1}^dg_k(x_{1,k},\ldots,x_{n,k})=f^{a\to b}_{t_1,\ldots,t_n}(x_1,\ldots,x_n).
\end{equation}
Since $(B_{t_1},\ldots,B_{t_n})$ is obtained from $(\xi^1,\ldots,\xi^d)$ by a permutation of the $nd$ coordinates, which preserves Lebesgue measure, $(B_{t_1},\ldots,B_{t_n})$ has the density \eqref{eq:5-6-17:b}. Hence $B^{a\to b}$ is a $d$-dimensional Brownian bridge from $a$ to $b$ on $[0,T]$.

\emph{The process $X$ of \bk{(6.26)}.} As in the construction of the stochastic integral (\bk{Section 3.2}, where the filtration is assumed to satisfy the usual conditions of \bk{Definition 1.2.25}), we assume that $\{\F_t\}$ satisfies the usual conditions. For $1\le k\le d$ the component
\[
  X^{(k)}_t=a_k\Bigl(1-\frac tT\Bigr)+b_k\frac tT+(T-t)\,M^{(k)}_t;\quad 0\le t<T,\qquad X^{(k)}_T=b_k,
  \qquad\text{where } M^{(k)}_t\defeq\int_0^t\frac{dW^{(k)}_s}{T-s},
\]
is the process \bk{(6.26)} built from the standard one-dimensional Brownian motion $W^{(k)}$. By \bk{Corollary 6.10} it is almost surely continuous on $[0,T]$, and by \bk{Problem 6.11} the vector $\eta^k\defeq(X^{(k)}_{t_1},\ldots,X^{(k)}_{t_n})$ has the density $g_k$ displayed above, for $0=t_0<t_1<\cdots<t_n<T$. The intersection of the $d$ events of probability one on which the components are continuous is an event of probability one on which $X$ is continuous. In view of \eqref{eq:5-6-17:d} and of the argument just given for $B^{a\to b}$, it remains only to show that $\eta^1,\ldots,\eta^d$ are independent, and for this it suffices that the vectors $\mu^k\defeq(M^{(k)}_{t_1},\ldots,M^{(k)}_{t_n})$; $k=1,\ldots,d$, be independent, because $\eta^k$ is an affine function of $\mu^k$.

Put $\tau\defeq t_n<T$ and $h(s)\defeq(T-s)^{-1}\ind_{[0,\tau]}(s)$; $s\ge0$. This nonrandom, bounded, measurable function belongs to $\cL^*(W^{(k)})$, and $M^{(k)}_t=I^{(k)}_t(h)$ a.s.\ for $0\le t\le\tau$, where $I^{(k)}$ denotes stochastic integration with respect to $W^{(k)}$ (by \bk{(3.2.14)} and \bk{(3.2.17)}, $I^{(k)}_t$ depends only on the restriction of the integrand to $[0,t]$). For $N\ge1$ let
\[
  h_N(s)\defeq\frac1T\,\ind_{\{0\}}(s)+\sum_{j=0}^{N-1}\frac{1}{T-j\tau/N}\,\ind_{(j\tau/N,\,(j+1)\tau/N]}(s);\qquad s\ge0 .
\]
This is a simple process in the sense of \bk{Definition 3.2.3}: take the strictly increasing partition points $r_j\defeq j\tau/N$ for $0\le j\le N$ and $r_j\defeq\tau+j-N$ for $j>N$ (so $r_0=0$ and $r_j\to\infty$), and the nonrandom coefficients $c_j\defeq(T-j\tau/N)^{-1}$ for $0\le j<N$, $c_j\defeq0$ for $j\ge N$, which are bounded by $(T-\tau)^{-1}$ and trivially $\F_{r_j}$-measurable; then $h_N(s)=c_0\ind_{\{0\}}(s)+\sum_{j\ge0}c_j\ind_{(r_j,r_{j+1}]}(s)$, the value $h_N(0)=1/T$ being $c_0$, as the definition requires. (The fixed times $t_1<\cdots<t_n=\tau$ keep their meaning throughout.) By \bk{(3.2.5)}
\[
  I^{(k)}_t(h_N)=\sum_{j=0}^{N-1}\frac{1}{T-j\tau/N}\Bigl(W^{(k)}_{t\wedge(j+1)\tau/N}-W^{(k)}_{t\wedge j\tau/N}\Bigr).
\]
Hence $\mu^k_N\defeq(I^{(k)}_{t_1}(h_N),\ldots,I^{(k)}_{t_n}(h_N))$ is a Borel function of the values of $W^{(k)}$ at finitely many times, the same times for every $k$, so that $\mu^1_N,\ldots,\mu^d_N$ are independent by Step 0(a). Because $s\mapsto(T-s)^{-1}$ is uniformly continuous on $[0,\tau]$, and since $h_N(0)=h(0)=1/T$, we have $\eps_N\defeq\sup_{0\le s\le\tau}\abs{h_N(s)-h(s)}\le\sup\{\abs{(T-s)^{-1}-(T-s')^{-1}};\ 0\le s'\le s\le\tau,\ s-s'\le\tau/N\}\to0$ as $N\to\infty$, and $h_N=h=0$ on $(\tau,\infty)$. By the linearity \bk{(3.2.17)} and the isometry \bk{(3.2.14)}, both valid for integrands in $\cL^*$ by \bk{Proposition 3.2.10}, and since $\qv{W^{(k)}}_s=s$,
\[
  E\bigl(I^{(k)}_{t_i}(h_N)-M^{(k)}_{t_i}\bigr)^2=\int_0^{t_i}\bigl(h_N(s)-h(s)\bigr)^2ds\le\tau\,\eps_N^2\xrightarrow[N\to\infty]{}0;\qquad 1\le i\le n,\ 1\le k\le d.
\]
Thus $(\mu^1_N,\ldots,\mu^d_N)\to(\mu^1,\ldots,\mu^d)$ in $L^2$, hence in distribution, and for $u_1,\ldots,u_d\in\R^n$
\[
  E\exp\Bigl\{i\sum_{k=1}^du_k\cdot\mu^k\Bigr\}=\lim_{N\to\infty}E\exp\Bigl\{i\sum_{k=1}^du_k\cdot\mu^k_N\Bigr\}
  =\lim_{N\to\infty}\prod_{k=1}^dE\,e^{iu_k\cdot\mu^k_N}=\prod_{k=1}^dE\,e^{iu_k\cdot\mu^k}.
\]
The joint characteristic function of $(\mu^1,\ldots,\mu^d)$ factorizes, so $\mu^1,\ldots,\mu^d$ are independent. This completes the proof that $X$ is a $d$-dimensional Brownian bridge from $a$ to $b$ on $[0,T]$.

\textbf{(ii)} Let $B^{a\to b}$ and $B^{b\to a}$ be given by \bk{(6.29)}, and set
\[
  \beta_t\defeq W_t-\frac tT\,W_T,\qquad Y_t\defeq B^{b\to a}_{T-t}=b\,\frac tT+a\Bigl(1-\frac tT\Bigr)+\beta_{T-t}=m(t)+\beta_{T-t};\qquad 0\le t\le T .
\]
By Step 0(a), $\{\beta^{(k)}_t;\,0\le t\le T,\,1\le k\le d\}$ is a centered Gaussian family (finitely many of its members form a linear image of a Gaussian vector), and by \eqref{eq:5-6-17:c}, for $0\le s,t\le T$,
\begin{equation}\label{eq:5-6-17:e}
  E\bigl(\beta^{(j)}_s\beta^{(k)}_t\bigr)=\delta_{jk}\Bigl[(s\wedge t)-\frac tT\,s-\frac sT\,t+\frac{st}{T^2}\,T\Bigr]=\delta_{jk}\,\rho(s,t).
\end{equation}
The family $\{\beta^{(k)}_{T-t}\}$ is centered Gaussian as well, with
$E\bigl(\beta^{(j)}_{T-s}\beta^{(k)}_{T-t}\bigr)=\delta_{jk}\,\rho(T-s,T-t)$, and
\[
  \rho(T-s,T-t)=T-(s\vee t)-\frac{(T-s)(T-t)}{T}=s+t-(s\vee t)-\frac{st}{T}=\rho(s,t).
\]
Hence, for $0<t_1<\cdots<t_n<T$, the vectors $(Y_{t_1},\ldots,Y_{t_n})$ and $(B^{a\to b}_{t_1},\ldots,B^{a\to b}_{t_n})$ $=(m(t_i)+\beta_{t_i})_{1\le i\le n}$ are both Gaussian in $\R^{nd}$, with the same mean vector $(m(t_i))_{1\le i\le n}$ and the same covariances $\delta_{jk}\,\rho(t_i,t_l)$. A Gaussian distribution being determined by its mean vector and covariance matrix, the two vectors have the same distribution. Both $Y$ and $B^{a\to b}$ have continuous paths, so by the Lemma of Step 0 they have the same law on $C[0,T]^d$; this is the claim.

For arbitrary bridges the conclusion persists: the time reversal $R\colon C[0,T]^d\to C[0,T]^d$, $(Ry)(t)\defeq y(T-t)$, is an isometry and hence Borel measurable, and what we have just shown reads $Q^{b\to a}R^{-1}=Q^{a\to b}$. If $\wt B$ is any Brownian bridge from $b$ to $a$ and $\wh B$ any Brownian bridge from $a$ to $b$, then the law of $\{\wt B_{T-t};\,0\le t\le T\}=R(\wt B)$ is $Q^{b\to a}R^{-1}=Q^{a\to b}$, the law of $\wh B$.

\textbf{(iii)} Define $\Phi\colon\R^d\times C[0,T]^d\to C[0,T]^d$ by
\[
  \Phi(b,y)(t)\defeq a+y(t)+\frac tT\,(b-a);\qquad 0\le t\le T .
\]
For each $t$ the map $(b,y)\mapsto\Phi(b,y)(t)$ is $\B(\R^d)\otimes\B(C[0,T]^d)$-measurable; since the coordinate maps generate $\B(C[0,T]^d)$, the map $\Phi$ is $\B(\R^d)\otimes\B(C[0,T]^d)/\B(C[0,T]^d)$-measurable, and $F\circ\Phi$ is a bounded, $\B(\R^d)\otimes\B(C[0,T]^d)$-measurable function. With $\beta$ as in (ii), regarded as a random element of $C[0,T]^d$ (all its paths are continuous), we have for every $\omega$ and every $b\in\R^d$
\begin{equation}\label{eq:5-6-17:f}
  \Phi(b,\beta_\cdot)=B^{a\to b}_\cdot\ \text{ of \bk{(6.29)}},\qquad
  \Phi(a+W_T,\beta_\cdot)(t)=a+W_t-\frac tT\,W_T+\frac tT\,W_T=a+W_t .
\end{equation}

\emph{The random element $\beta$ is independent of the random vector $W_T$.} Let $0\le t_1<\cdots<t_n\le T$. The vector $(\beta_{t_1},\ldots,\beta_{t_n},W_T)$ is Gaussian by Step 0(a), and by \eqref{eq:5-6-17:c}
\[
  E\bigl(\beta^{(j)}_tW^{(k)}_T\bigr)=\delta_{jk}\Bigl[(t\wedge T)-\frac tT\,T\Bigr]=0;\qquad 0\le t\le T,\ 1\le j,k\le d .
\]
Thus $(\beta_{t_1},\ldots,\beta_{t_n})$ and $W_T$ are jointly Gaussian and uncorrelated, hence independent. Consequently, for fixed $\Gamma\in\B(\R^d)$, the finite measures $A\mapsto P[\beta\in A,\,W_T\in\Gamma]$ and $A\mapsto P[\beta\in A]\,P[W_T\in\Gamma]$ on $\B(C[0,T]^d)$ agree on the finite-dimensional cylinder sets, which form a class closed under finite intersections that contains $C[0,T]^d$ and generates $\B(C[0,T]^d)$; by the Dynkin System Theorem they agree on $\B(C[0,T]^d)$. This is the asserted independence.

\emph{Proof of \bk{(6.35)}.} Let $Q$ denote the law of $\beta$ on $C[0,T]^d$. Since $W_T=W_T-W_0$ a.s.\ is normal with mean zero and covariance matrix $TI_d$, the vector $a+W_T$ has the density $b\mapsto(2\pi T)^{-d/2}\exp\{-\norm{b-a}^2/2T\}=p(T;a,b)$ on $\R^d$. By the independence just proved, the pair $(a+W_T,\beta)$ induces on $\B(\R^d)\otimes\B(C[0,T]^d)$ the product measure $p(T;a,b)\,db\otimes Q(dy)$. Using \eqref{eq:5-6-17:f} and Fubini's theorem for the bounded, product-measurable function $F\circ\Phi$, we obtain that $b\mapsto\int F(\Phi(b,y))\,Q(dy)=EF(\Phi(b,\beta_\cdot))=EF(B^{a\to b}_\cdot)$ is Borel measurable and that
\begin{align*}
  EF(a+W_\cdot)&=E\bigl[(F\circ\Phi)(a+W_T,\beta_\cdot)\bigr]
  =\int_{\R^d}\biggl(\int_{C[0,T]^d}F(\Phi(b,y))\,Q(dy)\biggr)p(T;a,b)\,db\\
  &=\int_{\R^d}EF(B^{a\to b}_\cdot)\,p(T;a,b)\,db .
\end{align*}
Here $a+W_\cdot$ denotes the path $\{a+W_t;\,0\le t\le T\}$, and $B^{a\to b}$ is the process \bk{(6.29)}; by (i) and Step 0(b), $EF(B^{a\to b}_\cdot)=\int F\,dQ^{a\to b}$, so \bk{(6.35)} holds with any $d$-dimensional Brownian bridge from $a$ to $b$ in the integrand.

\rmk{Formula \bk{(6.35)} says that $\{Q^{a\to b}\}_{b\in\R^d}$ is a regular conditional distribution of the path $a+W_\cdot$ given $a+W_T$. Indeed, for $A\in\B(C[0,T]^d)$ the map $b\mapsto Q^{a\to b}(A)=E\ind_A(B^{a\to b}_\cdot)$ is Borel measurable, as shown above with $F=\ind_A$. For $\Gamma\in\B(\R^d)$ apply \bk{(6.35)} to the bounded, measurable function $F(y)\defeq\ind_A(y)\ind_\Gamma(y(T))$ (the coordinate map $y\mapsto y(T)$ is measurable); since every path of \bk{(6.29)} satisfies $B^{a\to b}_T=b$, we obtain
\[
  P\bigl[a+W_\cdot\in A,\ a+W_T\in\Gamma\bigr]=\int_\Gamma Q^{a\to b}(A)\,p(T;a,b)\,db=E\bigl[Q^{a\to b}(A)\big|_{b=a+W_T}\,\ind_{\{a+W_T\in\Gamma\}}\bigr],
\]
the last equality because $a+W_T$ has the density $p(T;a,\cdot)$. As $\Gamma$ ranges over $\B(\R^d)$ the events $\{a+W_T\in\Gamma\}$ exhaust $\sigma(W_T)$, so $Q^{a\to b}(A)|_{b=a+W_T}$ is a version of $P[a+W_\cdot\in A\,|\,W_T]$.}
\end{kex}

\begin{kex}{Exercise}{6.18}{}
For $t=0$ we have $P(X^{(a)}_0\in A)=\ind_A(a)$ and \bk{(6.38)} is trivial, so we fix $t>0$ throughout. We write
\[
  V\defeq\tfrac12\bigl(\Delta\Phi+\norm{\nabla\Phi}^2\bigr),\qquad
  c\defeq\sup_{x\in\R^d}\abs{V(x)}<\infty ;
\]
$V$ is continuous, and it is bounded because $\nabla\Phi$ and the second partial derivatives of $\Phi$ are bounded. We denote by $C[0,t]^d$ the space of continuous functions $y\colon[0,t]\to\R^d$ with the supremum norm; its Borel $\sigma$-field is generated by the coordinate maps $y\mapsto y(s)$; $0\le s\le t$ (as in \bk{Problem 2.4.2}). On this space we consider the functional
\[
  G(y)\defeq\exp\Bigl\{-\int_0^t V\bigl(y(s)\bigr)\,ds\Bigr\};\qquad y\in C[0,t]^d,
\]
which is continuous (by bounded convergence), hence Borel-measurable, and satisfies $e^{-ct}\le G\le e^{ct}$. Finally, $W=\{W_s;\ 0\le s\le t\}$ denotes a $d$-dimensional Brownian motion with $W_0=0$ a.s., and $p(t;x,y)=(2\pi t)^{-d/2}\exp\{-\norm{x-y}^2/2t\}$.

\emph{Step 1: the identity \bk{(6.39)} for every Borel function $f\colon\R^d\to[0,\infty]$.} We shall show
\begin{equation}\label{eq:5-6-18:a}
  E f\bigl(X^{(a)}_t\bigr)=e^{-\Phi(a)}\,E\Bigl[f(a+W_t)\,e^{\Phi(a+W_t)}\,G(a+W_\cdot)\Bigr];\qquad a\in\R^d ,
\end{equation}
which is \bk{(6.39)}. Let $X=\{X_s,\F_s;\ 0\le s\le t\}$, $(\Omega,\F)$, $\{P^x\}_{x\in\R^d}$ be a $d$-dimensional Brownian family, as in the proof of \bk{Proposition 3.6} with $T=t$ and $b(s,x)=\nabla\Phi(x)$, a bounded continuous function, so that \bk{(3.3)} holds; the filtration $\{\F_s\}$ is chosen to satisfy the usual conditions (\bk{Remark 3.7}), so that stochastic integrals with respect to $X$ and the It\^o rule are available under each $P^a$. With
\[
  Z_s\defeq\exp\Bigl\{\sum_{j=1}^d\int_0^s\frac{\partial\Phi}{\partial x_j}(X_u)\,dX^{(j)}_u-\frac12\int_0^s\norm{\nabla\Phi(X_u)}^2\,du\Bigr\};\qquad 0\le s\le t,
\]
that proof shows that $Z$ is a martingale under each $P^a$ and that, under the probability measure $Q^a$ with $dQ^a/dP^a=Z_t$, the process $X$ together with the $Q^a$-Brownian motion \bk{(3.4)} is a weak solution of \bk{(6.36)} on $[0,t]$ with $Q^a[X_0=a]=1$. On the other hand, the strong solution $X^{(a)}$, restricted to $[0,t]$, with its driving Brownian motion and the augmented filtration of the latter, is also a weak solution of \bk{(6.36)} on $[0,t]$ with initial distribution $\delta_a$. Since the drift $\nabla\Phi$ is bounded, \bk{Corollary 3.11} asserts that uniqueness in the sense of probability law holds; thus $X^{(a)}_t$ has under $P$ the same distribution as $X_t$ under $Q^a$, and (this is \bk{(3.11)} with $n=1$, $t_1=T=t$, and $\delta_a$ as initial distribution)
\begin{equation}\label{eq:5-6-18:b}
  P\bigl[X^{(a)}_t\in\Gamma\bigr]=Q^a[X_t\in\Gamma]=E^a\bigl[\ind_{\{X_t\in\Gamma\}}Z_t\bigr];\qquad\Gamma\in\B(\R^d).
\end{equation}
Under $P^a$, the process $X$ is a $d$-dimensional Brownian motion starting at $a$, and $\Phi$ is of class $C^2$; the It\^o rule (\bk{Theorem 3.3.6}) gives, $P^a$-a.s.,
\[
  \Phi(X_t)-\Phi(a)=\sum_{j=1}^d\int_0^t\frac{\partial\Phi}{\partial x_j}(X_u)\,dX^{(j)}_u+\frac12\int_0^t\Delta\Phi(X_u)\,du .
\]
Solving for the stochastic integral and substituting into the definition of $Z_t$, we obtain
\begin{equation}\label{eq:5-6-18:c}
  Z_t=\exp\Bigl\{\Phi(X_t)-\Phi(a)-\int_0^tV(X_u)\,du\Bigr\}=e^{-\Phi(a)}\,e^{\Phi(X_t)}\,G(X_\cdot),\qquad P^a\text{-a.s.}
\end{equation}
From \eqref{eq:5-6-18:b}, \eqref{eq:5-6-18:c} and the fact that the law of $\{X_s;\ 0\le s\le t\}$ under $P^a$ on $C[0,t]^d$ is that of $\{a+W_s;\ 0\le s\le t\}$, we conclude that \eqref{eq:5-6-18:a} holds for $f=\ind_\Gamma$. Both sides of \eqref{eq:5-6-18:a} are linear in $f$ and continuous along increasing sequences (monotone convergence), so \eqref{eq:5-6-18:a} holds for nonnegative simple functions and then for all Borel $f\colon\R^d\to[0,\infty]$. In particular, with $f\equiv1$,
\begin{equation}\label{eq:5-6-18:d}
  E\Bigl[e^{\Phi(a+W_t)}\,G(a+W_\cdot)\Bigr]=e^{\Phi(a)};\qquad a\in\R^d .
\end{equation}

\emph{Step 2: bridges.} For $a,b\in\R^d$ let $B^{a\to b}$ be the $d$-dimensional Brownian bridge from $a$ to $b$ on $[0,t]$ given by \bk{(6.29)} with $T=t$,
\[
  B^{a\to b}_s=a\Bigl(1-\frac st\Bigr)+b\,\frac st+\Bigl(W_s-\frac st\,W_t\Bigr);\qquad 0\le s\le t
\]
(\bk{Exercise 6.17}\,(i)). For each $s$ the map $(a,b,\omega)\mapsto B^{a\to b}_s(\omega)$ is $\B(\R^d)\otimes\B(\R^d)\otimes\F$-measurable; since the coordinate maps generate the Borel $\sigma$-field of $C[0,t]^d$, the map $(a,b,\omega)\mapsto B^{a\to b}_\cdot(\omega)\in C[0,t]^d$ is measurable as well. Consequently, for every Borel-measurable $F\colon C[0,t]^d\to[0,\infty]$, the function $(a,b)\mapsto EF(B^{a\to b}_\cdot)$ is Borel-measurable on $\R^d\times\R^d$, by Tonelli's theorem. In particular,
\[
  g(a,b)\defeq E\,G\bigl(B^{a\to b}_\cdot\bigr);\qquad a,b\in\R^d,
\]
is a Borel-measurable function with $e^{-ct}\le g\le e^{ct}$. We need two facts.

\textbf{(a)} \emph{$g(a,b)=g(b,a)$ for all $a,b\in\R^d$.} Let $\rho\colon C[0,t]^d\to C[0,t]^d$ be the time reversal $(\rho y)(s)\defeq y(t-s)$, an isometry, hence Borel-measurable. The change of variable $s\mapsto t-s$ in the integral defining $G$ shows $G\circ\rho=G$. By \bk{Exercise 6.17}\,(ii), the processes $B^{a\to b}$ and $\rho(B^{b\to a})$ have the same law on $C[0,t]^d$. Therefore
\[
  g(a,b)=E\,G\bigl(B^{a\to b}_\cdot\bigr)=E\,(G\circ\rho)\bigl(B^{b\to a}_\cdot\bigr)=E\,G\bigl(B^{b\to a}_\cdot\bigr)=g(b,a).
\]

\textbf{(b)} \emph{For every Borel function $h\colon\R^d\to[0,\infty]$ and every $a\in\R^d$,}
\begin{equation}\label{eq:5-6-18:e}
  E\bigl[h(a+W_t)\,G(a+W_\cdot)\bigr]=\int_{\R^d}h(b)\,g(a,b)\,p(t;a,b)\,db .
\end{equation}
Indeed, for $n\ge1$ the functional $F_n(y)\defeq\bigl(h(y(t))\wedge n\bigr)G(y)$ is bounded and Borel-measurable on $C[0,t]^d$, so \bk{Exercise 6.17}\,(iii) yields
\[
  E\,F_n(a+W_\cdot)=\int_{\R^d}E\,F_n\bigl(B^{a\to b}_\cdot\bigr)\,p(t;a,b)\,db
  =\int_{\R^d}\bigl(h(b)\wedge n\bigr)\,g(a,b)\,p(t;a,b)\,db ,
\]
where we have used $B^{a\to b}_t=b$. Letting $n\to\infty$ and using the monotone convergence theorem on both sides, we arrive at \eqref{eq:5-6-18:e}.

\emph{Step 3: proof of \bk{(6.38)}.} Let $A\in\B(\R^d)$. Applying \eqref{eq:5-6-18:a} with $f=\ind_A$ and then \eqref{eq:5-6-18:e} with $h=\ind_Ae^{\Phi}$, we get
\begin{equation}\label{eq:5-6-18:f}
  P\bigl(X^{(a)}_t\in A\bigr)=e^{-\Phi(a)}\int_A e^{\Phi(b)}\,g(a,b)\,p(t;a,b)\,db;\qquad a\in\R^d .
\end{equation}
The right-hand side is a Borel-measurable function of $a$ (Tonelli's theorem), so the integral in \bk{(6.38)} is well defined. Multiplying \eqref{eq:5-6-18:f} by $e^{2\Phi(a)}$, integrating over $a\in\R^d$, interchanging the order of integration (Tonelli's theorem: the integrand is nonnegative and jointly measurable), and using the symmetries $g(a,b)=g(b,a)$ of (a) and $p(t;a,b)=p(t;b,a)$, we find
\begin{align*}
  \int_{\R^d}P\bigl(X^{(a)}_t\in A\bigr)\,\mu(da)
  &=\int_{\R^d}\int_A e^{\Phi(a)}\,e^{\Phi(b)}\,g(a,b)\,p(t;a,b)\,db\,da\\
  &=\int_A e^{\Phi(b)}\Bigl[\int_{\R^d}e^{\Phi(a)}\,g(b,a)\,p(t;b,a)\,da\Bigr]db .
\end{align*}
By \eqref{eq:5-6-18:e}, with the roles of $a$ and $b$ interchanged and $h=e^{\Phi}$, and then by \eqref{eq:5-6-18:d}, the inner integral equals
\[
  E\Bigl[e^{\Phi(b+W_t)}\,G(b+W_\cdot)\Bigr]=e^{\Phi(b)} .
\]
Hence
\[
  \int_{\R^d}P\bigl(X^{(a)}_t\in A\bigr)\,\mu(da)=\int_Ae^{2\Phi(b)}\,db=\mu(A),
\]
both sides being possibly equal to $+\infty$. This is \bk{(6.38)}.

\rmk{The hint suggests establishing \bk{(6.39)} for $f\in C_0^\infty(\R^d)$ and concluding with the determining-class property of \bk{Problem 4.25}. Since \bk{(3.11)} is available for indicators of Borel sets, we have obtained \bk{(6.39)} directly for all nonnegative Borel functions, and \bk{Problem 4.25} is not needed. This also avoids a small difficulty: \bk{Definition 4.24} concerns \emph{finite} measures, whereas $\mu$ need not be finite (for $\Phi\equiv0$ it is Lebesgue measure, which is indeed invariant for Brownian motion).}
\end{kex}

\begin{kex}{Exercise}{6.20}{}
With $Z_t\defeq(X_t,Y_t)^T$, the system is the equation \bk{(6.1)} with $d=2$, $r=1$, $a(t)\equiv 0$ and the constant matrices
\[
  A\defeq\begin{pmatrix}0&1\\-\beta&-\alpha\end{pmatrix},\qquad
  b\defeq\begin{pmatrix}0\\ \sigma\end{pmatrix}
\]
in the roles of $A(t)$ and $\sigma(t)$; that is, $dZ_t=AZ_t\,dt+b\,dW_t$. As in \bk{(6.1)}, $W=\{W_t,\F_t;\,0\le t<\infty\}$ is a Brownian motion on a probability space $(\Omega,\F,P)$ whose filtration satisfies the usual conditions, and the initial vector $Z_0=(X_0,Y_0)^T$ is $\F_0$-measurable, hence independent of $W$.

\textbf{(i)} \emph{Step 1: the fundamental solution.} Set $\Delta\defeq\alpha^2-4\beta$ and define $S\colon\R\to\R$ by
\[
  S(t)\defeq
  \begin{cases}
    \dfrac{\sinh(\gamma t)}{\gamma},\quad \gamma\defeq\tfrac12\sqrt{\Delta}; & \Delta>0,\\[1.2ex]
    t; & \Delta=0,\\[0.6ex]
    \dfrac{\sin(\omega t)}{\omega},\quad \omega\defeq\tfrac12\sqrt{-\Delta}; & \Delta<0.
  \end{cases}
\]
In each case $S''=\frac{\Delta}{4}S$, $S(0)=0$ and $S'(0)=1$. Let
\begin{equation}\label{eq:5-6-20:u}
  u(t)\defeq e^{-\alpha t/2}S(t),\qquad\text{so that}\qquad
  u'(t)=e^{-\alpha t/2}\bigl(S'(t)-\tfrac{\alpha}{2}S(t)\bigr).
\end{equation}
Then $u''=e^{-\alpha t/2}\bigl(S''-\alpha S'+\frac{\alpha^2}{4}S\bigr)$, and therefore
\[
  u''+\alpha u'+\beta u
  =e^{-\alpha t/2}\Bigl(S''-\bigl(\tfrac{\alpha^2}{4}-\beta\bigr)S\Bigr)=0,\qquad u(0)=0,\quad u'(0)=1 .
\]
We claim that
\begin{equation}\label{eq:5-6-20:exp}
  e^{tA}=\Psi(t)\defeq
  \begin{pmatrix}u'(t)+\alpha u(t)&u(t)\\[0.3ex]-\beta u(t)&u'(t)\end{pmatrix};\qquad t\in\R .
\end{equation}
Indeed, $\Psi(0)=I$, and using $u''=-\alpha u'-\beta u$ (whence $u''+\alpha u'=-\beta u$) we find
\[
  \Psi'(t)=\begin{pmatrix}-\beta u&u'\\-\beta u'&u''\end{pmatrix},\qquad
  A\Psi(t)=\begin{pmatrix}-\beta u&u'\\-\beta(u'+\alpha u)+\alpha\beta u&-\beta u-\alpha u'\end{pmatrix}
  =\begin{pmatrix}-\beta u&u'\\-\beta u'&u''\end{pmatrix}.
\]
Thus $\Psi$ solves $\dot\Phi=A\Phi$, $\Phi(0)=I$, and so does $t\mapsto e^{tA}$ by \bk{(6.17)}; the solution of this linear matrix equation being unique (\bk{(6.3)}), \eqref{eq:5-6-20:exp} follows.

\emph{Step 2: the solution.} According to \bk{(6.6)} and \bk{(6.17)}, the process
\begin{equation}\label{eq:5-6-20:Z}
  Z_t=e^{tA}\Bigl[Z_0+\int_0^t e^{-sA}b\,dW_s\Bigr];\qquad 0\le t<\infty
\end{equation}
is a continuous, adapted solution of the system, and by \bk{Theorem 2.5} (the coefficients $z\mapsto Az$ and $b$ are globally Lipschitz) every solution on the same space, with the same Brownian motion and the same initial vector, is indistinguishable from it. The integrand in \eqref{eq:5-6-20:Z} is deterministic and continuous, so the stochastic integral is well defined. For fixed $t$, the entries of the constant matrix $e^{tA}$ may be taken inside the stochastic integral by linearity, and $e^{tA}e^{-sA}b=e^{(t-s)A}b=\sigma\bigl(u(t-s),u'(t-s)\bigr)^T$ by \eqref{eq:5-6-20:exp}. Reading off the two components of \eqref{eq:5-6-20:Z} we obtain, for every $t\ge 0$, almost surely,
\begin{equation}\label{eq:5-6-20:sol}
\begin{aligned}
  X_t&=\bigl(u'(t)+\alpha u(t)\bigr)X_0+u(t)\,Y_0+\sigma\int_0^t u(t-s)\,dW_s,\\
  Y_t&=-\beta u(t)\,X_0+u'(t)\,Y_0+\sigma\int_0^t u'(t-s)\,dW_s,
\end{aligned}
\end{equation}
with $u$ given by \eqref{eq:5-6-20:u}. Explicitly, in the underdamped case $\alpha^2<4\beta$,
\begin{gather*}
  u(t)=e^{-\alpha t/2}\,\frac{\sin\omega t}{\omega},\qquad
  u'(t)=e^{-\alpha t/2}\Bigl(\cos\omega t-\frac{\alpha}{2\omega}\sin\omega t\Bigr),\\
  u'(t)+\alpha u(t)=e^{-\alpha t/2}\Bigl(\cos\omega t+\frac{\alpha}{2\omega}\sin\omega t\Bigr);
\end{gather*}
in the overdamped case $\alpha^2>4\beta$ the same formulas hold with $\sinh\gamma t$, $\cosh\gamma t$, $\gamma$ in place of $\sin\omega t$, $\cos\omega t$, $\omega$; and in the critical case $\alpha^2=4\beta$,
\[
  u(t)=t\,e^{-\alpha t/2},\qquad u'(t)=\bigl(1-\tfrac{\alpha t}{2}\bigr)e^{-\alpha t/2},\qquad
  u'(t)+\alpha u(t)=\bigl(1+\tfrac{\alpha t}{2}\bigr)e^{-\alpha t/2}.
\]

\rmk{The stochastic integrals in \eqref{eq:5-6-20:sol} can be removed. For fixed $t$, It\^o's rule applied to the product $u(t-s)W_s$, $0\le s\le t$, gives $u(0)W_t=-\int_0^tu'(t-s)W_s\,ds+\int_0^tu(t-s)\,dW_s$, and likewise with $u'$ in place of $u$. Since $u(0)=0$ and $u'(0)=1$,
\[
  \int_0^tu(t-s)\,dW_s=\int_0^tu'(t-s)W_s\,ds,\qquad
  \int_0^tu'(t-s)\,dW_s=W_t+\int_0^tu''(t-s)W_s\,ds\qquad\text{a.s.,}
\]
which expresses $(X_t,Y_t)$ as an explicit continuous functional of $(X_0,Y_0)$ and the path of $W$ on $[0,t]$.}

\textbf{(ii)} The eigenvalues of $A$ are the roots $\lambda_\pm=\frac12\bigl(-\alpha\pm\sqrt{\Delta}\bigr)$ of $\lambda^2+\alpha\lambda+\beta=0$. If $\Delta<0$, they have real part $-\alpha/2<0$; if $\Delta\ge 0$, they are real with sum $-\alpha<0$ and product $\beta>0$, hence both negative. Thus all eigenvalues of $A$ have negative real parts, the matrix
\[
  V\defeq\int_0^\infty e^{sA}\,bb^T e^{sA^T}ds
\]
of \bk{(6.19)} is well defined, and by \bk{Problem 6.6} it satisfies \bk{(6.20)}: $AV+VA^T=-bb^T$.

\emph{Computation of $V$.} We first observe that the equation $AD+DA^T=-bb^T$ has at most one solution $D$. If $D$ is the difference of two solutions, then $AD+DA^T=0$, so
$\frac{d}{dt}\bigl(e^{tA}De^{tA^T}\bigr)=e^{tA}(AD+DA^T)e^{tA^T}=0$ and $D=e^{tA}De^{tA^T}$ for all $t\ge 0$; by \eqref{eq:5-6-20:exp} and \eqref{eq:5-6-20:u} every entry of $e^{tA}$ tends to zero as $t\to\infty$ (the functions $S$, $S'$ grow at most like $e^{\gamma t}$ with $\gamma<\alpha/2$ when $\Delta>0$, at most linearly when $\Delta=0$, and are bounded when $\Delta<0$), whence $D=0$. Now let
\begin{equation}\label{eq:5-6-20:V}
  V_0\defeq\frac{\sigma^2}{2\alpha}\begin{pmatrix}1/\beta&0\\0&1\end{pmatrix}.
\end{equation}
Then
\[
  AV_0=\frac{\sigma^2}{2\alpha}\begin{pmatrix}0&1\\-1&-\alpha\end{pmatrix},\qquad
  AV_0+V_0A^T=AV_0+(AV_0)^T=\begin{pmatrix}0&0\\0&-\sigma^2\end{pmatrix}=-bb^T,
\]
and by uniqueness $V=V_0$.

\emph{Stationarity.} Suppose now that $(X_0,Y_0)$ is $\F_0$-measurable and normally distributed with mean zero and covariance matrix \eqref{eq:5-6-20:V}; in other words, $X_0$ and $Y_0$ are independent, independent of $W$, with
\[
  X_0\sim N\Bigl(0,\frac{\sigma^2}{2\alpha\beta}\Bigr),\qquad Y_0\sim N\Bigl(0,\frac{\sigma^2}{2\alpha}\Bigr).
\]
The hypotheses of \bk{Theorem 6.7} are satisfied: the dispersion matrix $b$ is constant, $a(t)\equiv 0$, the eigenvalues of $A$ have negative real parts, and $Z_0$ is bivariate normal with $m(0)=0$ and $E(Z_0Z_0^T)=V$ as in \bk{(6.19)}. Hence $Z=(X,Y)^T$ is a zero-mean Gaussian process (\bk{Problem 6.2}) whose covariance function \bk{(6.8)} is given by \bk{(6.21)}:
\begin{equation}\label{eq:5-6-20:rho}
  \rho(s,t)=E\bigl(Z_sZ_t^T\bigr)=
  \begin{cases}
    e^{(s-t)A}V; & 0\le t\le s<\infty,\\
    Ve^{(t-s)A^T}; & 0\le s\le t<\infty.
  \end{cases}
\end{equation}
In particular, $\rho(s+h,t+h)=\rho(s,t)$ for all $s,t,h\ge 0$. For $0\le t_1<\dots<t_n$ and $h\ge 0$, the random vectors $(Z_{t_1},\dots,Z_{t_n})$ and $(Z_{t_1+h},\dots,Z_{t_n+h})$ are both normal with mean zero and with the same covariance matrix $\bigl(\rho(t_i,t_j)\bigr)_{i,j}$, so they have the same distribution: $(X_t,Y_t)$ is a stationary Gaussian process.

\rmk{This is the only normal initial law which makes $Z$ stationary. Stationarity forces $m(t)\equiv m(0)$ and $V(t)\equiv V(0)$; then \bk{(6.12)} gives $Am(0)=0$, i.e.\ $m(0)=0$ because $\det A=\beta\neq 0$, and \bk{(6.13)$'$} gives $AV(0)+V(0)A^T=-bb^T$, i.e.\ $V(0)=V_0$ by the uniqueness shown above. The law $N(0,V_0)$ has density proportional to $\exp\{-\frac{2\alpha}{\sigma^2}(\frac12y^2+\frac12\beta x^2)\}$, the Gibbs density associated with the energy of the oscillator.}

\textbf{(iii)} We insert \eqref{eq:5-6-20:exp} and \eqref{eq:5-6-20:V} into \eqref{eq:5-6-20:rho}. For $t\ge 0$ and $\tau\ge 0$,
\[
  \rho(t+\tau,t)=e^{\tau A}V
  =\frac{\sigma^2}{2\alpha}
  \begin{pmatrix}\dfrac{u'(\tau)+\alpha u(\tau)}{\beta}&u(\tau)\\[1.6ex]-u(\tau)&u'(\tau)\end{pmatrix},
\]
that is,
\begin{equation}\label{eq:5-6-20:cov}
\begin{aligned}
  E\bigl(X_{t+\tau}X_t\bigr)&=\frac{\sigma^2}{2\alpha\beta}\bigl(u'(\tau)+\alpha u(\tau)\bigr), &
  E\bigl(X_{t+\tau}Y_t\bigr)&=\frac{\sigma^2}{2\alpha}\,u(\tau),\\
  E\bigl(Y_{t+\tau}X_t\bigr)&=-\frac{\sigma^2}{2\alpha}\,u(\tau), &
  E\bigl(Y_{t+\tau}Y_t\bigr)&=\frac{\sigma^2}{2\alpha}\,u'(\tau),
\end{aligned}
\end{equation}
and $\rho(t,t+\tau)=Ve^{\tau A^T}=\rho(t+\tau,t)^T$. With the function $S$ of part (i) and \eqref{eq:5-6-20:u}, for all $s,t\ge 0$ and $\tau\defeq\abs{t-s}$:
\begin{align*}
  \Cov(X_s,X_t)&=\frac{\sigma^2}{2\alpha\beta}\,e^{-\alpha\tau/2}\Bigl(S'(\tau)+\frac{\alpha}{2}S(\tau)\Bigr),\\
  \Cov(Y_s,Y_t)&=\frac{\sigma^2}{2\alpha}\,e^{-\alpha\tau/2}\Bigl(S'(\tau)-\frac{\alpha}{2}S(\tau)\Bigr),\\
  \Cov(X_s,Y_t)&=\frac{\sigma^2}{2\alpha}\,\sgn(s-t)\,e^{-\alpha\tau/2}S(\tau).
\end{align*}
In the underdamped case $\alpha^2<4\beta$, $\omega=\frac12\sqrt{4\beta-\alpha^2}$, these read
\begin{align*}
  \Cov(X_s,X_t)&=\frac{\sigma^2}{2\alpha\beta}\,e^{-\alpha\tau/2}\Bigl(\cos\omega\tau+\frac{\alpha}{2\omega}\sin\omega\tau\Bigr),\\
  \Cov(Y_s,Y_t)&=\frac{\sigma^2}{2\alpha}\,e^{-\alpha\tau/2}\Bigl(\cos\omega\tau-\frac{\alpha}{2\omega}\sin\omega\tau\Bigr),\\
  \Cov(X_s,Y_t)&=\frac{\sigma^2}{2\alpha\omega}\,\sgn(s-t)\,e^{-\alpha\tau/2}\sin\omega\tau;
\end{align*}
in the overdamped case $\alpha^2>4\beta$ one replaces $\cos\omega\tau$, $\sin\omega\tau$, $\omega$ by $\cosh\gamma\tau$, $\sinh\gamma\tau$, $\gamma=\frac12\sqrt{\alpha^2-4\beta}$; and in the critical case $\alpha^2=4\beta$,
\begin{gather*}
  \Cov(X_s,X_t)=\frac{\sigma^2}{2\alpha\beta}\Bigl(1+\frac{\alpha\tau}{2}\Bigr)e^{-\alpha\tau/2},\qquad
  \Cov(Y_s,Y_t)=\frac{\sigma^2}{2\alpha}\Bigl(1-\frac{\alpha\tau}{2}\Bigr)e^{-\alpha\tau/2},\\
  \Cov(X_s,Y_t)=\frac{\sigma^2}{2\alpha}\,(s-t)\,e^{-\alpha\tau/2}.
\end{gather*}

\rmk{Two checks. At $\tau=0$ the matrix in \eqref{eq:5-6-20:cov} reduces to $V$, since $u(0)=0$ and $u'(0)=1$. Moreover, $X_{t+\tau}=X_t+\int_t^{t+\tau}Y_r\,dr$ gives $E(X_{t+\tau}X_t)=E(X_t^2)+\int_0^\tau E(Y_{t+r}X_t)\,dr$ by Fubini's theorem, and indeed $\frac{d}{d\tau}\frac{\sigma^2}{2\alpha\beta}(u'+\alpha u)=\frac{\sigma^2}{2\alpha\beta}(u''+\alpha u')=-\frac{\sigma^2}{2\alpha}u$, in agreement with the entry $E(Y_{t+\tau}X_t)$.}
\end{kex}

\begin{kex}{Exercise}{6.21}{}
Here $d=r=1$, $W$ is a standard, one-dimensional Brownian motion with respect to a filtration $\{\F_t\}$, and $A\colon[0,\infty)\to\R$ is measurable and locally bounded, as stipulated after \bk{(6.1)}; the equation reads $dX_t=A(t)X_t\,dt+\sigma\,dW_t$, $X_0=x$.

\emph{Step 1: the solution.} Since $A$ is measurable and bounded on every interval $[0,t]$, the function
\[
  \Phi(t)\defeq\exp\Bigl\{\int_0^tA(u)\,du\Bigr\};\quad 0\le t<\infty
\]
is well defined, strictly positive and continuous. The function $t\mapsto\int_0^tA(u)\,du$ is absolutely continuous on every compact interval, with derivative $A(t)$ for Lebesgue-almost every $t$, and the exponential function is Lipschitz continuous on bounded sets; hence $\Phi$ is absolutely continuous on every compact interval, with $\dot\Phi(t)=A(t)\Phi(t)$ for almost every $t$, and $\Phi(0)=1$. Thus $\Phi$ is the fundamental solution of \bk{(6.3)}, which is unique. As observed in the text following \bk{(6.6)}, pathwise uniqueness holds for \bk{(6.1)} and its solution is given by \bk{(6.6)}; in the present case,
\begin{equation}\label{eq:5-6-21:a}
  X_t=\Phi(t)\bigl[x+\sigma M_t\bigr],\qquad M_t\defeq\int_0^t\frac{dW_s}{\Phi(s)};\quad 0\le t<\infty.
\end{equation}

\emph{Step 2: the second moment.} The integrand $(s,\omega)\mapsto1/\Phi(s)$ does not depend on $\omega$ and is continuous in $s$; it is therefore progressively measurable, and $\int_0^T\Phi^{-2}(s)\,ds<\infty$ for every $T<\infty$, so that it belongs to the class $\cL^*$ associated with $W$ (for which $\qv{W}_s=s$). By \bk{Proposition 3.2.10}, the process $M$ is a square-integrable martingale which satisfies \bk{(3.2.12)}--\bk{(3.2.14)}:
\[
  EM_t=EM_0=0,\qquad EM_t^2=\int_0^t\frac{ds}{\Phi^2(s)};\quad 0\le t<\infty.
\]
Squaring \eqref{eq:5-6-21:a} and taking expectations, we obtain for every $t\ge0$
\begin{equation}\label{eq:5-6-21:b}
  EX_t^2=\Phi^2(t)\bigl[x^2+2x\sigma\,EM_t+\sigma^2EM_t^2\bigr]
  =x^2\,\Phi^2(t)+\sigma^2\int_0^t\Bigl(\frac{\Phi(t)}{\Phi(s)}\Bigr)^{2}ds.
\end{equation}

\emph{Step 3: the bound.} For $0\le s\le t$, the assumption $A(u)\le-\alpha$; $0\le u<\infty$, gives
\[
  0<\frac{\Phi(t)}{\Phi(s)}=\exp\Bigl\{\int_s^tA(u)\,du\Bigr\}\le e^{-\alpha(t-s)},
\]
and in particular $0<\Phi(t)\le e^{-\alpha t}$. Substituting these bounds into \eqref{eq:5-6-21:b}, we conclude that
\[
  EX_t^2\le x^2e^{-2\alpha t}+\sigma^2\int_0^te^{-2\alpha(t-s)}\,ds
  =x^2e^{-2\alpha t}+\frac{\sigma^2}{2\alpha}\bigl(1-e^{-2\alpha t}\bigr)
  =\frac{\sigma^2}{2\alpha}+\Bigl(x^2-\frac{\sigma^2}{2\alpha}\Bigr)e^{-2\alpha t}
\]
for every $t\ge0$, as claimed.

\rmk{If $A(t)\equiv-\alpha$, then $\Phi(t)=e^{-\alpha t}$ and both inequalities of Step 3 are equalities: the bound is attained by the Ornstein--Uhlenbeck process of \bk{Example 6.8}. In general, the right-hand side increases (if $x^2<\sigma^2/2\alpha$) or decreases (if $x^2>\sigma^2/2\alpha$) to the limit $\sigma^2/2\alpha$ as $t\to\infty$, so that $\sup_{t\ge0}EX_t^2\le x^2\vee(\sigma^2/2\alpha)$.}
\end{kex}

\begin{kex}{Exercise}{6.22}{}
\emph{Conventions.} As in \bk{Section 3.2}, the filtration $\{\F_t\}$ satisfies the usual conditions. For the integrals in \bk{(6.40)} to be defined, the processes $A$, $S^{(p)}$ have to be measurable as well as adapted, and we assume this. By \bk{Proposition 1.1.12} each of their entries has a progressively measurable modification, which we may truncate at the level of the bound of the original entry. A measurable process and a measurable modification of it differ only on a subset of $[0,T]\times\Omega$ whose sections at fixed $t$ are $P$-null; by the Fubini--Tonelli theorem this set is null for the product of Lebesgue measure and $P$, so replacing $A$, $S^{(p)}$ by these modifications does not affect any of the integrals below, up to indistinguishability. We assume therefore that $A$, $S^{(p)}$ are progressively measurable, and we extend them by $A(t)=S^{(p)}(t)=0$ for $t>T$, so that all processes are defined on $[0,\infty)$ and all solutions are constant on $[T,\infty)$. Stochastic integrals with respect to $W^{(p)}$ are, as usual, integrals with respect to the martingale $W^{(p)}-W^{(p)}_0$, which belongs to $\cM_2^c$ by \bk{Problem 2.5.5}; since only the increments of $W$ enter \bk{(6.40)} and \bk{(6.41)}, replacing $W$ by $W-W_0$ (an $r$-dimensional Brownian motion with respect to the same filtration) changes none of the equations, and we assume $W_0=0$. Then $W^{(p)}\in\cM_2^c$ and $\qv{W^{(p)},W^{(q)}}_t=\delta_{pq}t$ by \bk{Problem 2.5.5}. We write $\norm{M}^2\defeq\sum_{i,j}M_{ij}^2$ for a $(d\times d)$ matrix $M$, as in \bk{(2.9)}; the Cauchy--Schwarz inequality gives $\norm{MN}\le\norm{M}\,\norm{N}$. We fix a constant $K$ with $\norm{A(t,\omega)}\le K$, $\norm{S^{(p)}(t,\omega)}\le K$ for all $(t,\omega)$ and $1\le p\le r$, and set
\[
  D(t)\defeq\sum_{p=1}^r\bigl(S^{(p)}(t)\bigr)^2-A(t),\qquad\text{so that}\quad \norm{D(t,\omega)}\le rK^2+K.
\]
By a \emph{solution} of \bk{(6.40)} we mean a continuous, adapted, matrix-valued process $X$ which satisfies \bk{(6.40)} for all $0\le t<\infty$, almost surely; by a solution of \bk{(6.41)} we mean a continuous, adapted, matrix-valued process $Y$ which satisfies, almost surely,
\begin{equation}\label{eq:5-6-22:Y}
  Y(t)=I+\int_0^tY(s)D(s)\,ds-\sum_{p=1}^r\int_0^tY(s)S^{(p)}(s)\,dW^{(p)}_s;\quad 0\le t<\infty.
\end{equation}
A continuous, adapted process is progressively measurable (\bk{Proposition 1.1.13}) and has paths bounded on $[0,T]$; hence the entries of $S^{(p)}X$ and $YS^{(p)}$ are progressively measurable, vanish on $(T,\infty)$ and are pathwise bounded, so that they belong to the class $\cP^*$ of \bk{Definition 3.2.22} for $W^{(p)}$, and all the integrals above are well defined.

\rmk{\bk{Theorems 2.5} and \bk{2.9} deal with nonrandom coefficients $b(t,x)$, $\sigma(t,x)$, whereas here the coefficients $A(t,\omega)x$, $S^{(p)}(t,\omega)x$ depend on $\omega$; the proofs of these theorems carry over, but the theorems cannot be quoted as they stand. For this reason we do not rely on the assertion of existence and uniqueness in the statement: existence is proved in Step 1 and uniqueness is a by-product of Step 2. We shall show: \emph{there exist a solution $X$ of \bk{(6.40)} and a solution $Y$ of \bk{(6.41)}; each of them is unique up to indistinguishability; and almost surely, $X(t)$ is invertible with $X^{-1}(t)=Y(t)$ for every $0\le t<\infty$.} This contains the claim of the exercise.}

\emph{Step 1: existence.}

\begin{lem}
Let $B_0,B_1,\dots,B_r$ be progressively measurable, $(d\times d)$ matrix-valued processes which vanish on $(T,\infty)\times\Omega$ and satisfy $\norm{B_p(t,\omega)}\le L$ for some constant $L$. Then there exists a continuous, adapted, matrix-valued process $Z$ with $E\bigl[\sup_{0\le t<\infty}\norm{Z(t)}^2\bigr]<\infty$, such that almost surely
\begin{equation}\label{eq:5-6-22:Z}
  Z(t)=I+\int_0^tZ(s)B_0(s)\,ds+\sum_{p=1}^r\int_0^tZ(s)B_p(s)\,dW^{(p)}_s;\quad 0\le t<\infty.
\end{equation}
\end{lem}

\begin{proof}
Let $\mathscr{H}$ be the set of continuous, adapted, matrix-valued processes $H$ with $\norm{H}_*^2\defeq E\bigl[\sup_{0\le t<\infty}\norm{H(t)}^2\bigr]<\infty$. For $H\in\mathscr{H}$, the entries of $HB_p$ are progressively measurable and $\norm{H(s)B_p(s)}\le L\norm{H(s)}\ind_{[0,T]}(s)$, whence $E\int_0^\infty\norm{H(s)B_p(s)}^2ds\le L^2T\norm{H}_*^2<\infty$. Consequently, the entries of $HB_p$ belong to $\cL^*$ for $W^{(p)}$, and by \bk{Proposition 3.2.10} the process
\[
  J(H)(t)\defeq\int_0^tH(s)B_0(s)\,ds+\sum_{p=1}^r\int_0^tH(s)B_p(s)\,dW^{(p)}_s;\quad 0\le t<\infty
\]
is well defined, continuous and adapted, its stochastic integrals being matrices of martingales in $\cM_2^c$; it is constant on $[T,\infty)$, and $J$ is linear. We claim that, with $C\defeq(r+1)L^2(T+4r)$,
\begin{equation}\label{eq:5-6-22:est}
  E\Bigl[\sup_{0\le s\le t}\norm{J(H)(s)}^2\Bigr]\le C\int_0^tE\norm{H(s)}^2\,ds;\quad 0\le t\le T.
\end{equation}
Indeed, $J(H)$ is a sum of $r+1$ terms, so $\norm{J(H)(s)}^2$ is at most $r+1$ times the sum of the squared norms of these terms. For the first one, the Cauchy--Schwarz inequality yields, for $0\le s\le t\le T$,
\[
  \norm[\Big]{\int_0^sH(u)B_0(u)\,du}^2\le\Bigl(\int_0^s\norm{H(u)B_0(u)}\,du\Bigr)^2\le TL^2\int_0^t\norm{H(u)}^2\,du.
\]
For the others we apply Doob's maximal inequality (\bk{Theorem 1.3.8}(iv)) and the isometry \bk{(3.2.14)} of \bk{Proposition 3.2.10} to each entry of the matrix $N^{(p)}\defeq\int_0^{\cdot}H(u)B_p(u)\,dW^{(p)}_u$ of martingales in $\cM_2^c$, and sum over the entries:
\[
  E\Bigl[\sup_{0\le s\le t}\norm{N^{(p)}(s)}^2\Bigr]\le\sum_{i,j}4\,E\bigl(N^{(p)}_{ij}(t)\bigr)^2
  =4\,E\int_0^t\norm{H(u)B_p(u)}^2\,du\le4L^2E\int_0^t\norm{H(u)}^2\,du.
\]
Taking expectations, using Tonelli's theorem and adding the bound $TL^2\int_0^tE\norm{H(u)}^2du$ for the first term to the $r$ bounds $4L^2\int_0^tE\norm{H(u)}^2du$ for the others, we obtain \eqref{eq:5-6-22:est}, since $(r+1)(TL^2+4rL^2)=C$. Because $J(H)$ is constant on $[T,\infty)$, \eqref{eq:5-6-22:est} with $t=T$ gives
\begin{equation}\label{eq:5-6-22:cont}
  \norm{J(H)}_*^2\le CT\norm{H}_*^2;\quad H\in\mathscr{H},
\end{equation}
and in particular $J$ maps $\mathscr{H}$ into itself.

Now set $Z^{(0)}\equiv I$ and $Z^{(k+1)}\defeq I+J(Z^{(k)})$; $k\ge0$. All these processes belong to $\mathscr{H}$, and by linearity $Z^{(k+1)}-Z^{(k)}=J\bigl(Z^{(k)}-Z^{(k-1)}\bigr)$ for $k\ge1$, whereas $Z^{(1)}-Z^{(0)}=J(I)$. The functions
\[
  g_k(t)\defeq E\Bigl[\sup_{0\le s\le t}\norm{Z^{(k+1)}(s)-Z^{(k)}(s)}^2\Bigr];\quad 0\le t\le T
\]
are nondecreasing, and \eqref{eq:5-6-22:est} gives $g_0(t)\le C\norm{I}^2t=Cd\,t$ and $g_k(t)\le C\int_0^tg_{k-1}(s)\,ds$ for $k\ge1$; by induction,
\[
  g_k(t)\le d\,\frac{(Ct)^{k+1}}{(k+1)!};\quad 0\le t\le T,\ k\ge0.
\]
Since $Z^{(k+1)}-Z^{(k)}$ is constant on $[T,\infty)$, we have $\norm{Z^{(k+1)}-Z^{(k)}}_*=\sqrt{g_k(T)}$, and $\sum_k\sqrt{g_k(T)}<\infty$ by the ratio test. By Minkowski's inequality and the monotone convergence theorem, the random variable $\sum_{k\ge0}\sup_{t}\norm{Z^{(k+1)}(t)-Z^{(k)}(t)}$ has finite second moment; in particular it is finite on an event $\Omega^*$ with $P(\Omega^*)=1$. On $\Omega^*$ the sequence $Z^{(k)}$ converges uniformly on $[0,\infty)$; we denote its limit by $Z$ on $\Omega^*$ and set $Z\equiv I$ on $\Omega\setminus\Omega^*$. Then $Z$ is continuous, it is adapted because $\Omega^*\in\F_0$ under the usual conditions, and
\[
  \norm{Z-Z^{(k)}}_*\le\sum_{j\ge k}\sqrt{g_j(T)}\xrightarrow[k\to\infty]{}0;
\]
in particular $Z\in\mathscr{H}$. From \eqref{eq:5-6-22:cont} and the linearity of $J$ we get $\norm{J(Z)-J(Z^{(k)})}_*\to0$. Letting $k\to\infty$ in $Z^{(k+1)}=I+J(Z^{(k)})$, where both sides converge in the norm $\norm{\cdot}_*$, we conclude that $\norm{Z-I-J(Z)}_*=0$, which is \eqref{eq:5-6-22:Z}.
\end{proof}

The Lemma with $B_0=D$, $B_p=-S^{(p)}$ and $L=rK^2+K$ provides a solution $Y$ of \bk{(6.41)}, i.e., of \eqref{eq:5-6-22:Y}. The Lemma with $B_0=A^{\top}$, $B_p=(S^{(p)})^{\top}$ and $L=K$ provides a process $Z$ satisfying \eqref{eq:5-6-22:Z}; transposing this identity entry by entry, we see that $X\defeq Z^{\top}$ is a solution of \bk{(6.40)}.

\emph{Step 2: $YX=I$.} Let $X$ be \emph{any} solution of \bk{(6.40)} and $Y$ \emph{any} solution of \bk{(6.41)}. We show that, almost surely,
\begin{equation}\label{eq:5-6-22:YX}
  Y(t)X(t)=I;\quad 0\le t<\infty.
\end{equation}
The entries of $X$ and $Y$ are continuous semimartingales: for $1\le i,j,k\le d$,
\begin{align*}
  X_{kj}&=\delta_{kj}+M_{kj}+B_{kj},&
  M_{kj}(t)&\defeq\sum_{p=1}^r\int_0^t\bigl(S^{(p)}X\bigr)_{kj}(s)\,dW^{(p)}_s,&
  B_{kj}(t)&\defeq\int_0^t(AX)_{kj}(s)\,ds,\\
  Y_{ik}&=\delta_{ik}+N_{ik}+C_{ik},&
  N_{ik}(t)&\defeq-\sum_{p=1}^r\int_0^t\bigl(YS^{(p)}\bigr)_{ik}(s)\,dW^{(p)}_s,&
  C_{ik}(t)&\defeq\int_0^t(YD)_{ik}(s)\,ds,
\end{align*}
where $M_{kj},N_{ik}\in\cM^{c,\loc}$ (\bk{Definition 3.2.23}) and $B_{kj},C_{ik}$ are continuous, adapted and of bounded variation on bounded intervals, being integrals of progressively measurable, pathwise bounded integrands. We need two facts.

\textbf{(a)} \emph{If $\sigma\in\cP^*$, $N\defeq\int_0^{\cdot}\sigma_s\,dW^{(p)}_s$ and $U$ is a continuous, adapted process, then $\int_0^tU_s\,dN_s=\int_0^tU_s\sigma_s\,dW^{(p)}_s$; $0\le t<\infty$, a.s.} Both integrals are defined, since $U$ and $U\sigma$ are progressively measurable and $\int_0^tU_s^2\,d\qv{N}_s=\int_0^tU_s^2\sigma_s^2\,ds<\infty$. For every $\wt N\in\cM_2^c$, \bk{Proposition 3.2.24} applied to $N$ gives $\qv{N,\wt N}_t=\int_0^t\sigma_s\,d\qv{W^{(p)},\wt N}_s$, and then, applied to $\int U\,dN$ and to $\int U\sigma\,dW^{(p)}$,
\[
  \Bigl\langle\int_0^{\cdot}U_s\,dN_s,\wt N\Bigr\rangle_t=\int_0^tU_s\,d\qv{N,\wt N}_s=\int_0^tU_s\sigma_s\,d\qv{W^{(p)},\wt N}_s
  =\Bigl\langle\int_0^{\cdot}U_s\sigma_s\,dW^{(p)}_s,\wt N\Bigr\rangle_t .
\]
The uniqueness assertion of \bk{Proposition 3.2.24} yields (a). (This is \bk{Corollary 3.2.20} for local martingales.)

\textbf{(b)} \emph{If $\sigma,\rho\in\cP^*$, then $\bigl\langle\int_0^{\cdot}\sigma_s\,dW^{(p)}_s,\int_0^{\cdot}\rho_s\,dW^{(q)}_s\bigr\rangle_t=\delta_{pq}\int_0^t\sigma_s\rho_s\,ds$.} This is the cross-variation formula \bk{(3.2.26)}, which remains valid for local martingales and integrands in $\cP^*$ (see the paragraph preceding \bk{Proposition 3.2.24}), together with $\qv{W^{(p)},W^{(q)}}_t=\delta_{pq}t$, which is \bk{Problem 2.5.5} applied to $W^{(p)}-W^{(p)}_0=W^{(p)}$ (recall that $W_0=0$).

We now apply the integration-by-parts formula \bk{(3.3.8)} of \bk{Problem 3.3.12} to the semimartingales $Y_{ik}$ and $X_{kj}$: almost surely, for all $0\le t<\infty$,
\begin{equation}\label{eq:5-6-22:ibp}
  Y_{ik}(t)X_{kj}(t)=\delta_{ik}\delta_{kj}+\int_0^tY_{ik}(s)\,dX_{kj}(s)+\int_0^tX_{kj}(s)\,dY_{ik}(s)+\qv{N_{ik},M_{kj}}_t ,
\end{equation}
where $\int Y_{ik}\,dX_{kj}=\int Y_{ik}\,dM_{kj}+\int Y_{ik}\,dB_{kj}$ and similarly for $\int X_{kj}\,dY_{ik}$. By (a) and the linearity of the stochastic integral in the integrator (\bk{Problem 3.2.25}), and because $B_{kj}$, $C_{ik}$ are absolutely continuous with the indicated densities,
\begin{align*}
  \int_0^tY_{ik}\,dX_{kj}&=\sum_{p=1}^r\int_0^tY_{ik}\bigl(S^{(p)}X\bigr)_{kj}\,dW^{(p)}_s+\int_0^tY_{ik}(AX)_{kj}\,ds,\\
  \int_0^tX_{kj}\,dY_{ik}&=-\sum_{p=1}^r\int_0^t\bigl(YS^{(p)}\bigr)_{ik}X_{kj}\,dW^{(p)}_s+\int_0^t(YD)_{ik}X_{kj}\,ds,
\end{align*}
and by (b) and the bilinearity of the cross-variation,
\[
  \qv{N_{ik},M_{kj}}_t=-\sum_{p=1}^r\int_0^t\bigl(YS^{(p)}\bigr)_{ik}\bigl(S^{(p)}X\bigr)_{kj}\,ds .
\]
We substitute these three expressions into \eqref{eq:5-6-22:ibp} and sum over $k=1,\dots,d$. Since
\[
  \sum_kY_{ik}\bigl(S^{(p)}X\bigr)_{kj}=\bigl(YS^{(p)}X\bigr)_{ij}=\sum_k\bigl(YS^{(p)}\bigr)_{ik}X_{kj},
\]
the stochastic integrals cancel, and we are left with
\begin{align*}
  (YX)_{ij}(t)&=\delta_{ij}+\int_0^t\Bigl(YAX+YDX-\sum_{p=1}^rYS^{(p)}S^{(p)}X\Bigr)_{ij}(s)\,ds\\
  &=\delta_{ij}+\int_0^t\Bigl(Y\Bigl[A+D-\sum_{p=1}^r\bigl(S^{(p)}\bigr)^2\Bigr]X\Bigr)_{ij}(s)\,ds=\delta_{ij},
\end{align*}
by the definition of $D$. This holds for all $1\le i,j\le d$ and all $0\le t<\infty$ outside a single $P$-null set, and \eqref{eq:5-6-22:YX} is proved.

\emph{Step 3: conclusion.} Let $X$ be a solution of \bk{(6.40)}, and let $Y$ be the solution of \bk{(6.41)} constructed in Step 1. By \eqref{eq:5-6-22:YX} there is an event of probability one on which $Y(t)X(t)=I$ for every $0\le t<\infty$. A square matrix with a left inverse is invertible, and the left inverse is the inverse; thus, on this event, $X(t)$ is invertible and $X^{-1}(t)=Y(t)$ for every $t\ge0$. Substituting $Y=X^{-1}$ into \eqref{eq:5-6-22:Y}, we obtain \bk{(6.41)}.

Finally, uniqueness. If $X$ and $X'$ are solutions of \bk{(6.40)}, then with $Y$ as above, \eqref{eq:5-6-22:YX} gives $X(t)=Y^{-1}(t)=X'(t)$ for all $t\ge0$, almost surely. Likewise, if $Y$ and $Y'$ are solutions of \bk{(6.41)} and $X$ is the solution of \bk{(6.40)} constructed in Step 1, then $Y(t)=X^{-1}(t)=Y'(t)$ for all $t\ge0$, almost surely.

\rmk{For $d=r=1$ the solution of \bk{(6.40)} is $X(t)=\exp\bigl\{\int_0^tS(s)\,dW_s+\int_0^t\bigl(A(s)-\tfrac12S^2(s)\bigr)ds\bigr\}$, and \bk{(6.41)} reduces to $d(1/X)=(1/X)\bigl[(S^2-A)\,dt-S\,dW\bigr]$, in agreement with \bk{Problem 3.3.10} (the case $A\equiv0$). For $d\ge2$ the matrices do not commute in general, and the order of the factors in \bk{(6.40)}, \bk{(6.41)} matters: $X$ is multiplied by the coefficients from the left, $X^{-1}$ from the right.}
\end{kex}

\solsection{5.7}{Connections with Partial Differential Equations}

\begin{kex}{Problem}{7.7}{}
The constant $\rho$ in \bk{(7.19)} is understood to be strictly positive, as the bound on $\mu$ requires. We keep all the other hypotheses of \bk{Theorem 7.6}: the continuity and the conditions \bk{(7.10)}, \bk{(7.11)} of $f$, $g$, $k\ge 0$, the standing assumptions \bk{(7.2)}--\bk{(7.4)}, and the fact that $v$ is continuous on $[0,T]\times\R^d$, of class $C^{1,2}([0,T)\times\R^d)$, and solves \bk{(7.12)}, \bk{(7.13)}. We have to show that \bk{(7.15)} holds when \bk{(7.14)} is replaced by \bk{(7.19)}, \bk{(7.20)}.

\emph{Step 1: where \bk{(7.14)} enters the proof of \bk{Theorem 7.6}.} For $t=T$ the representation \bk{(7.15)} is just \bk{(7.13)}, so we fix $(t,x)\in[0,T)\times\R^d$, a weak solution $(X,W)$, $(\Omega,\F,P)$, $\{\F_s\}$ of \bk{(7.1)} with $X_t=x$, and $\tau_n\defeq\inf\{s\ge t;\ \norm{X_s}\ge n\}$. The proof of \bk{Theorem 7.6} establishes the identity \bk{(7.16)} for every integer $n>\norm{x}$ by means of It\^o's rule and of \bk{(7.12)}, \bk{(7.13)} alone, and then lets $n\to\infty$. The first and the third terms on the right-hand side of \bk{(7.16)} are shown there to converge to the two terms on the right-hand side of \bk{(7.15)}, by dominated convergence (using \bk{(7.10)}\,(i), \bk{(7.11)}\,(i), $k\ge 0$ and the moment estimate \bk{(7.17)}, which rests only on the linear growth condition contained in \bk{(7.2)}) or by monotone convergence (under \bk{(7.10)}\,(ii), \bk{(7.11)}\,(ii)); and we have $\tau_n\uparrow\infty$ a.s.\ because the paths of $X$ are continuous, hence bounded on $[t,T]$. None of this involves the growth of $v$. The condition \bk{(7.14)} is used at one place only, namely to prove that the second term of \bk{(7.16)} tends to zero:
\begin{equation}\label{eq:5-7-07:goal}
  \lim_{n\to\infty} E^{t,x}\Bigl[\,\abs{v(\tau_n,X_{\tau_n})}
  \exp\Bigl\{-\int_t^{\tau_n} k(\theta,X_\theta)\,d\theta\Bigr\}\ind_{\{\tau_n\le T\}}\Bigr]=0 .
\end{equation}
It therefore suffices to derive \eqref{eq:5-7-07:goal} from \bk{(7.19)} and \bk{(7.20)}; the rest of the proof of \bk{Theorem 7.6} then applies verbatim and yields \bk{(7.15)}.

\emph{Step 2: reduction to a tail estimate.} Let $n>\norm{x}$. On $\{\tau_n\le T\}$ we have $t\le\tau_n\le T$ and, by the continuity of the paths of $X$ and $\norm{X_t}=\norm{x}<n$, also $\norm{X_{\tau_n}}=n$. Thus \bk{(7.20)} gives $\abs{v(\tau_n,X_{\tau_n})}\le Me^{\mu n^2}$ on this event, and since $k\ge 0$ the expectation in \eqref{eq:5-7-07:goal} is bounded by
\begin{equation}\label{eq:5-7-07:bd}
  Me^{\mu n^2}\,P^{t,x}[\tau_n\le T]
  =Me^{\mu n^2}\,P^{t,x}\Bigl[\max_{t\le\theta\le T}\norm{X_\theta}\ge n\Bigr].
\end{equation}
If $\norm{X_\theta}\ge n$, then $\abs{X^{(i)}_\theta}\ge n/\sqrt d$ for at least one index $i$, because $\norm{X_\theta}^2=\sum_{i=1}^d\abs{X^{(i)}_\theta}^2$. Writing $m_n$ for the integer part of $n/\sqrt d$, we obtain
\begin{equation}\label{eq:5-7-07:union}
  P^{t,x}\Bigl[\max_{t\le\theta\le T}\norm{X_\theta}\ge n\Bigr]
  \le\sum_{i=1}^d P^{t,x}\Bigl[\max_{t\le\theta\le T}\abs{X^{(i)}_\theta}\ge m_n\Bigr].
\end{equation}

\emph{Step 3: application of \bk{Problem 3.4.12}.} Fix $1\le i\le d$ and set, for $0\le u<\infty$,
\[
  Y_u\defeq X^{(i)}_{t+u},\qquad
  C_u\defeq\int_t^{t+u}b_i(\theta,X_\theta)\,d\theta,\qquad
  N_u\defeq\sum_{j=1}^r\int_t^{t+u}\sigma_{ij}(\theta,X_\theta)\,dW^{(j)}_\theta,
  \qquad\G_u\defeq\F_{t+u}.
\]
The filtration $\{\G_u\}$ satisfies the usual conditions, since $\{\F_s\}$ does, and $\{W_{t+u}-W_t,\G_u\}$ is an $r$-dimensional Brownian motion. By \bk{(7.1)}, $Y_u=x_i+N_u+C_u$ for all $u\ge 0$, a.s. The integrands are continuous in $\theta$ and, by \bk{(7.19)}, bounded: $\abs{b_i}\le\rho$ and $\sigma_{ij}^2\le\rho$. Hence $C$ is a continuous, $\{\G_u\}$-adapted process of bounded variation with $C_0=0$, and $N$ is a continuous (square-integrable) martingale of $\{\G_u\}$ with $N_0=0$ and
\[
  \qv{N}_u=\sum_{j=1}^r\int_t^{t+u}\sigma_{ij}^2(\theta,X_\theta)\,d\theta ,
\]
because $\qv{W^{(j)},W^{(l)}}_\theta=\delta_{jl}\,\theta$. Consequently, by \bk{(7.19)},
\[
  \abs{C_u}+\qv{N}_u\le\int_t^{t+u}\Bigl(\abs{b_i(\theta,X_\theta)}
  +\sum_{j=1}^r\sigma_{ij}^2(\theta,X_\theta)\Bigr)d\theta\le\rho u;\qquad 0\le u<\infty,
\]
almost surely. These are exactly the hypotheses of \bk{Problem 3.4.12} for the semimartingale $Y$, with the same constant $\rho$. That problem provides an integer $n_i\ge 1$ (depending on $x_i$, $\rho$ and $T$) such that
\[
  P^{t,x}\Bigl[\max_{0\le u\le T}\abs{Y_u}\ge m\Bigr]\le\exp\Bigl\{-\frac{m^2}{18\rho T}\Bigr\}
  \qquad\text{for every integer } m\ge n_i .
\]
Since $\max_{t\le\theta\le T}\abs{X^{(i)}_\theta}=\max_{0\le u\le T-t}\abs{Y_u}\le\max_{0\le u\le T}\abs{Y_u}$, we conclude from \eqref{eq:5-7-07:union} that
\begin{equation}\label{eq:5-7-07:tail}
  P^{t,x}\Bigl[\max_{t\le\theta\le T}\norm{X_\theta}\ge n\Bigr]
  \le d\exp\Bigl\{-\frac{m_n^2}{18\rho T}\Bigr\}
  \qquad\text{whenever } m_n\ge\max_{1\le i\le d}n_i ,
\end{equation}
and the condition on $m_n$ is satisfied for all sufficiently large $n$, because $m_n\to\infty$.

\emph{Step 4: conclusion.} For $n\ge\sqrt d$ we have $m_n>n/\sqrt d-1\ge 0$, whence
\[
  m_n^2\ge\Bigl(\frac{n}{\sqrt d}-1\Bigr)^2\ge\frac{n^2}{d}-\frac{2n}{\sqrt d}\,.
\]
Combining this with \eqref{eq:5-7-07:bd} and \eqref{eq:5-7-07:tail}, we see that for all sufficiently large $n$ the expectation in \eqref{eq:5-7-07:goal} is at most
\[
  Md\exp\Bigl\{\mu n^2-\frac{m_n^2}{18\rho T}\Bigr\}
  \le Md\exp\Bigl\{-\Bigl(\frac{1}{18\rho Td}-\mu\Bigr)n^2+\frac{n}{9\rho T\sqrt d}\Bigr\}.
\]
The coefficient of $n^2$ in the exponent is strictly negative, precisely because $\mu<1/(18\rho Td)$; thus the right-hand side tends to zero as $n\to\infty$. This proves \eqref{eq:5-7-07:goal}, and with it the representation \bk{(7.15)} on $[0,T]\times\R^d$ under \bk{(7.19)}, \bk{(7.20)} in place of \bk{(7.14)}. In particular, the solution of the Cauchy problem \bk{(7.12)}, \bk{(7.13)} is unique among the functions satisfying \bk{(7.20)}.

\rmk{The factor $d$ in the bound for $\mu$ comes from the passage from the Euclidean norm to the coordinates in Step~2: \bk{Problem 3.4.12} controls each coordinate at the level $n/\sqrt d$. The strictness of the inequality $\mu<1/(18\rho Td)$ is what absorbs both the factor $d$ in \eqref{eq:5-7-07:tail} and the rounding of $n/\sqrt d$ to the integer $m_n$, which is needed since \bk{Problem 3.4.12} is stated for integer levels.}
\end{kex}

\begin{kex}{Exercise}{7.10}{}
\emph{Setting and conventions.} We assume \bk{(7.2)}--\bk{(7.4)}, and that $k\colon[0,T]\times\R^d\to[0,\infty)$ is continuous. We fix $(t,x)\in[0,T)\times\R^d$ and a weak solution $(X,W)$, $(\Omega,\F,P)$, $\{\F_s\}$ of \bk{(7.1)} with $X_t=x$, and we suppress the superscript $(t,x)$. The random variable $Y$ is defined on the same probability space (which we enlarge, if necessary, by taking its product with $\bigl([0,\infty),\B([0,\infty)),e^{-y}\,dy\bigr)$; this does not change the law of $X$), has $P[Y>c]=e^{-c}$ for $c\ge 0$, and is independent of the $\sigma$-field $\sigma(X_s;\ s\ge t)$. Since $k$ is given on $[0,T]\times\R^d$ only, the infimum defining $\rho$ is taken over $s\in[t,T]$, with $\inf\emptyset=\infty$. We set
\begin{equation*}
  K_s\defeq\int_t^s k(\theta,X_\theta)\,d\theta,\qquad Z_s\defeq e^{-K_s};\qquad t\le s\le T .
\end{equation*}
For every $\omega$ the function $\theta\mapsto k(\theta,X_\theta(\omega))$ is continuous and nonnegative on $[t,T]$, so $s\mapsto K_s(\omega)$ is finite, continuous and nondecreasing with $K_t=0$, and $0<Z_s\le 1$. The integrand $(\theta,\omega)\mapsto k(\theta,X_\theta(\omega))$ is $\B([t,T])\otimes\sigma(X_s;\ s\ge t)$-measurable, because $X$ has continuous paths; hence each $K_s$ and each $Z_s$ is $\sigma(X_\theta;\ \theta\ge t)$-measurable by Fubini's theorem.

Concerning \bk{Definition 7.9}, the formula \bk{(7.22)} presupposes that $\xi\mapsto G(t,x;\tau,\xi)$ is Borel-measurable, and we read \bk{(7.23)} in the sense in which the book uses it when it derives \bk{(7.26)} from \bk{Theorem 7.6}: the function $u$ of \bk{(7.22)}, extended by $u(\tau,\cdot)\defeq f$, is continuous on $[0,\tau]\times\R^d$ (see the Remark at the end).

\emph{Step 1: the event $\{\rho>\tau\}$.} Let $t<\tau\le T$. We claim that
\begin{equation}\label{eq:5-7-10:event}
  \{\rho>\tau\}=\{K_\tau<Y\}.
\end{equation}
If $K_\tau\ge Y$, then $\tau$ belongs to the set whose infimum is $\rho$, so $\rho\le\tau$. If $K_\tau<Y$, then $K_s\le K_\tau<Y$ for all $s\in[t,\tau]$ by monotonicity; if $\tau=T$ the set defining $\rho$ is empty and $\rho=\infty>\tau$, while if $\tau<T$ the continuity of $K$ provides $\delta>0$ with $K_s<Y$ on $[t,\tau+\delta]$, so that $\rho\ge\tau+\delta>\tau$. This proves \eqref{eq:5-7-10:event}; in particular $\{\rho>\tau\}$ is an event. As $\Delta\notin\R^d$, the definition of $\wt X$ and \eqref{eq:5-7-10:event} give
\begin{equation*}
  \{\wt X_\tau\in A\}=\{X_\tau\in A,\ \rho>\tau\}=\{X_\tau\in A,\ K_\tau<Y\};\qquad A\in\B(\R^d).
\end{equation*}
The pair $(X_\tau,K_\tau)$ is $\sigma(X_s;\ s\ge t)$-measurable, hence independent of $Y$, and the function $(a,c,y)\mapsto\ind_A(a)\ind_{\{c<y\}}$ is nonnegative and Borel. By Fubini's theorem applied to the product of the law of $(X_\tau,K_\tau)$ and the law of $Y$, and by $P[Y>c]=e^{-c}$ for $c\ge 0$,
\begin{equation}\label{eq:5-7-10:kill}
  P[\wt X_\tau\in A]=E\bigl[\ind_A(X_\tau)\,e^{-K_\tau}\bigr]
  =E^{t,x}\Bigl[\ind_A(X_\tau)\exp\Bigl\{-\int_t^\tau k(\theta,X_\theta)\,d\theta\Bigr\}\Bigr];
  \qquad A\in\B(\R^d).
\end{equation}

\emph{Step 2: continuous test functions.} Let $t<\tau\le T$ and $f\in C_0(\R^d)$, i.e., $f$ is continuous with compact support (\bk{Remark 4.1}). We show that
\begin{equation}\label{eq:5-7-10:test}
  \int_{\R^d}G(t,x;\tau,\xi)f(\xi)\,d\xi=E^{t,x}\bigl[f(X_\tau)Z_\tau\bigr].
\end{equation}
Let $u$ be given by \bk{(7.22)} on $[0,\tau)\times\R^d$ and $u(\tau,\cdot)\defeq f$. We apply \bk{Theorem 7.6} with the time horizon $\tau$ in place of $T$, with $u$ in place of $v$, with the restriction of $k$ to $[0,\tau]\times\R^d$, with the terminal function $f$, and with $g\equiv 0$. Its hypotheses hold: \bk{(7.2)}--\bk{(7.4)} are in force; $k\ge 0$ is continuous; $f$ is continuous and bounded, so \bk{(7.10)}\,(i) holds, and $g\equiv0$ satisfies \bk{(7.11)}\,(i); by \bk{Definition 7.9}, $u$ is continuous on $[0,\tau]\times\R^d$, of class $C^{1,2}([0,\tau)\times\R^d)$, and satisfies \bk{(7.21)}, which is \bk{(7.12)} with $g\equiv0$, as well as the terminal condition $u(\tau,\cdot)=f$; finally $u$ is bounded on $[0,\tau)\times\R^d$ by \bk{Definition 7.9} and on $\{\tau\}\times\R^d$ because $f$ is, so that \bk{(7.14)} holds. The representation \bk{(7.15)} at the point $(t,x)$ reads $u(t,x)=E^{t,x}[f(X_\tau)Z_\tau]$, and this is \eqref{eq:5-7-10:test}.

\emph{Step 3: proof of \bk{(7.30)}.} Fix $t<\tau\le T$ and define two measures on $\B(\R^d)$ by
\begin{equation*}
  \mu(A)\defeq E^{t,x}\bigl[\ind_A(X_\tau)Z_\tau\bigr],\qquad
  \nu(A)\defeq\int_A G(t,x;\tau,\xi)\,d\xi .
\end{equation*}
Both are measures, by monotone convergence and because $G\ge 0$, and $\mu(\R^d)=E^{t,x}[Z_\tau]\le 1$. Approximating a bounded Borel function by simple functions we get $\int h\,d\mu=E^{t,x}[h(X_\tau)Z_\tau]$ and, for $h\ge0$, $\int h\,d\nu=\int_{\R^d}G(t,x;\tau,\xi)h(\xi)\,d\xi$; thus \eqref{eq:5-7-10:test} says
\begin{equation}\label{eq:5-7-10:munu}
  \int_{\R^d}f\,d\nu=\int_{\R^d}f\,d\mu;\qquad f\in C_0(\R^d),\ f\ge0 .
\end{equation}
For a compact set $C\subset\R^d$ and $n\ge1$ let $f_n(\xi)\defeq\bigl(1-n\,\mathrm{dist}(\xi,C)\bigr)^+$. Each $f_n$ is continuous, has compact support and values in $[0,1]$, and $f_n\downarrow\ind_C$ pointwise because $C$ is closed. By \eqref{eq:5-7-10:munu}, $\nu(C)\le\int f_1\,d\nu=\int f_1\,d\mu\le1$, so $f_1$ is integrable with respect to both measures, and dominated convergence in \eqref{eq:5-7-10:munu} with $f=f_n$ yields
\begin{equation*}
  \nu(C)=\mu(C)\quad\text{for every compact } C\subset\R^d .
\end{equation*}
Taking for $C$ the closed ball of radius $m$ centred at the origin and letting $m\to\infty$, we obtain $\nu(\R^d)=\mu(\R^d)\le1$. Thus $\mu$ and $\nu$ are finite measures of the same total mass which agree on the compact sets; these form a class closed under finite intersections which generates $\B(\R^d)$. By Dynkin's $\pi$-$\lambda$ theorem, $\mu=\nu$ on $\B(\R^d)$. In conjunction with \eqref{eq:5-7-10:kill}, this is
\begin{equation*}
  P[\wt X_\tau\in A]=\mu(A)=\int_A G(t,x;\tau,\xi)\,d\xi;\qquad A\in\B(\R^d),\quad 0\le t<\tau\le T,
\end{equation*}
which is \bk{(7.30)}. In particular, $\int_{\R^d}G(t,x;\tau,\xi)\,d\xi=P[\rho>\tau]=E^{t,x}[Z_\tau]\le 1$, and for $k\equiv0$ we have $\rho=\infty$ a.s.\ and recover \bk{(7.26)}.

The equality $\mu=\nu$ extends from indicators to simple functions by linearity, to nonnegative Borel functions by monotone convergence, and to Borel functions $h$ with $E^{t,x}[\abs{h(X_\tau)}Z_\tau]<\infty$ by decomposition into positive and negative parts:
\begin{equation}\label{eq:5-7-10:borel}
  E^{t,x}\bigl[h(X_\tau)Z_\tau\bigr]=\int_{\R^d}G(t,x;\tau,\xi)h(\xi)\,d\xi;\qquad 0\le t<\tau\le T,
\end{equation}
valid for every Borel $h\colon\R^d\to\R$ which is nonnegative (both sides possibly $+\infty$) or satisfies $E^{t,x}[\abs{h(X_\tau)}Z_\tau]<\infty$ (the right-hand integral then converging absolutely).

\emph{Step 4: the form \bk{(7.29)} of the solution \bk{(7.15)}.} Let now $f$, $g$, $v$ be as in \bk{Theorem 7.6}: $f$ and $g$ are continuous and satisfy \bk{(7.10)}, \bk{(7.11)}, and $v$ is a solution of \bk{(7.12)}, \bk{(7.13)} which satisfies \bk{(7.14)}. For $t=T$ the representation \bk{(7.15)} is \bk{(7.13)}, and \bk{(7.29)} is meaningless since $G(T,x;\tau,\xi)$ is not defined; so let $0\le t<T$. By \bk{Theorem 7.6},
\begin{equation}\label{eq:5-7-10:fk}
  v(t,x)=E^{t,x}\Bigl[f(X_T)Z_T+\int_t^Tg(s,X_s)Z_s\,ds\Bigr].
\end{equation}

We treat the two terms separately. If \bk{(7.10)}\,(i) holds, then $\abs{f(X_T)}Z_T\le L(1+\norm{X_T}^{2\lambda})$ is integrable by \bk{(7.17)}; if \bk{(7.10)}\,(ii) holds, then $f(X_T)Z_T\ge0$. In either case \eqref{eq:5-7-10:borel} with $\tau=T$, $h=f$ gives
\begin{equation}\label{eq:5-7-10:fterm}
  E^{t,x}\bigl[f(X_T)Z_T\bigr]=\int_{\R^d}G(t,x;T,\xi)f(\xi)\,d\xi .
\end{equation}
The function $(s,\omega)\mapsto g(s,X_s(\omega))Z_s(\omega)$ is $\B([t,T])\otimes\F$-measurable, since $g$ is continuous and $X$, $Z$ have continuous paths. If \bk{(7.11)}\,(i) holds, then by \bk{(7.17)}
\begin{equation*}
  E^{t,x}\int_t^T\abs{g(s,X_s)}Z_s\,ds\le L(T-t)\Bigl(1+E^{t,x}\Bigl[\max_{t\le\theta\le T}\norm{X_\theta}^{2\lambda}\Bigr]\Bigr)<\infty,
\end{equation*}
and $E^{t,x}[\abs{g(s,X_s)}Z_s]<\infty$ for every $s\in[t,T]$; Fubini's theorem and \eqref{eq:5-7-10:borel} with $\tau=s$, $h=g(s,\cdot)$ give
\begin{equation}\label{eq:5-7-10:gterm}
  E^{t,x}\int_t^Tg(s,X_s)Z_s\,ds=\int_t^TE^{t,x}\bigl[g(s,X_s)Z_s\bigr]\,ds
  =\int_t^T\int_{\R^d}G(t,x;\tau,\xi)g(\tau,\xi)\,d\xi\,d\tau .
\end{equation}
If \bk{(7.11)}\,(ii) holds, the same identities follow from Tonelli's theorem and the nonnegative case of \eqref{eq:5-7-10:borel}. (The single point $\tau=t$, where \eqref{eq:5-7-10:borel} is not available, is of measure zero. The inner integral in \eqref{eq:5-7-10:gterm} is a measurable function of $\tau$, being equal to $E^{t,x}[g(\tau,X_\tau)Z_\tau]$, so the iterated integral makes sense as written.)

Finally, in each of the four combinations of \bk{(7.10)} and \bk{(7.11)}, each of the two random variables $f(X_T)Z_T$ and $\int_t^Tg(s,X_s)Z_s\,ds$ is either integrable or nonnegative, so the expectation of their sum is the sum of their expectations, with values in $(-\infty,\infty]$; as the left-hand side $v(t,x)$ of \eqref{eq:5-7-10:fk} is finite, both expectations are finite. Adding \eqref{eq:5-7-10:fterm} and \eqref{eq:5-7-10:gterm} we conclude from \eqref{eq:5-7-10:fk} that
\begin{equation*}
  v(t,x)=\int_{\R^d}G(t,x;T,\xi)f(\xi)\,d\xi+\int_t^T\int_{\R^d}G(t,x;\tau,\xi)g(\tau,\xi)\,d\xi\,d\tau;
  \qquad 0\le t<T,\ x\in\R^d,
\end{equation*}
which is \bk{(7.29)}, now for the Cauchy problem \bk{(7.12)}, \bk{(7.13)} with the potential term.

\rmk{As printed, \bk{(7.23)} is a limit in $t$ for fixed $x$, which by itself does not make the extension of $u$ by $u(\tau,\cdot)=f$ continuous on $[0,\tau]\times\R^d$. \bk{Theorem 7.6} requires this continuity, so it is needed already for the book's own derivation of \bk{(7.25)}, \bk{(7.26)}, and Step~2 uses it in exactly the same way; this is why we read \bk{(7.23)} as $\lim_{(s,y)\to(\tau,x),\,s<\tau}u(s,y)=f(x)$, a property enjoyed by the fundamental solutions constructed in the references quoted after \bk{Definition 7.9}. The continuity enters the proof of \bk{Theorem 7.6} only to guarantee $u(s,X_s)\to f(X_\tau)$ a.s.\ as $s\uparrow\tau$. Note also that Steps~1--3 use no hypothesis on $f$, $g$, $v$: \bk{(7.30)} holds under \bk{(7.2)}--\bk{(7.4)} for continuous $k\ge0$ as soon as a fundamental solution exists.}
\end{kex}

\begin{kex}{Exercise}{7.11}{}
\emph{Conventions.} As in \bk{Theorem 7.6}, a solution of the Cauchy problem is a function $v$ which is continuous on $[0,T]\times\R^d$, of class $C^{1,2}([0,T)\times\R^d)$, and satisfies the two displayed relations; $f$ and $g$ are continuous. The function $b=(b_1,\dots,b_d)$ is assumed Borel-measurable (no continuity is needed), and we set $\beta\defeq\sup_{[0,T]\times\R^d}\norm{b}<\infty$. Condition \bk{(7.20)} reads $\abs{v(s,y)}\le Me^{\mu\norm{y}^2}$ on $[0,T]\times\R^d$, with $0<\mu<1/(18\rho Td)$ and $\rho$ the constant of \bk{(7.19)}; here the dispersion matrix is the identity, so \bk{(7.19)} is $\abs{b_i}+1\le\rho$ and thus $\rho\ge 1$. All we shall use is
\begin{equation}\label{eq:5-7-11:mu}
  8\mu T<1 .
\end{equation}
For $t=T$ the asserted formula reduces to $v(T,x)=f(x)$, because $P^x[W_0=x]=1$. We therefore fix $(t,x)\in[0,T)\times\R^d$, write $\ell\defeq T-t>0$, and work on the probability space $(\Omega,\F,P^x)$.

\emph{Step 0: the filtration.} By \bk{Definition 2.5.8} a $d$-dimensional Brownian family $\{W_s,\F_s;\ s\ge0\}$ is defined for all times $s\ge0$; the range $0\le t\le T$ in the statement only indicates that the formula involves $W$ on $[0,T-t]$. Hence $W$ and $\{\F_s\}$ are available on $[0,\infty)$, and under $P^x$ the process $W$ is a $d$-dimensional Brownian motion on $[0,\infty)$ starting at $x$. The stochastic integral of Chapter~3 requires the usual conditions, which $\{\F_s\}$ need not satisfy. Let $\G_s$ be the $\sigma$-field generated by $\F_{s+}$ and the $P^x$-null sets (of the $P^x$-completion of $\F$). Then $\{\G_s\}$ satisfies the usual conditions, and $B_s\defeq W_s-x$ is still a standard $d$-dimensional Brownian motion relative to it: for $s<u$, $A\in\F_{s+}$ and $\varphi$ bounded and continuous on $\R^d$ we have $A\in\F_{s+1/k}$ and hence $E^x[\ind_A\varphi(W_u-W_{s+1/k})]=P^x(A)\,E^x[\varphi(W_u-W_{s+1/k})]$ for all large $k$; letting $k\to\infty$ (bounded convergence, continuity of paths) we see that $W_u-W_s$ is independent of $\F_{s+}$, hence of $\G_s$. All stochastic integrals below are taken with respect to $B$ and $\{\G_s\}$, so that $\int(\,\cdot\,,dW_\theta)=\int(\,\cdot\,,dB_\theta)$. The process
\[
  X_\theta\defeq b\bigl((t+\theta)\wedge T,\,W_\theta\bigr);\qquad 0\le\theta<\infty,
\]
is progressively measurable (as the composition of a Borel function with a continuous adapted process) and satisfies $\norm{X_\theta}\le\beta$. With the notation \bk{(3.5.2)} we set
\[
  Z_s\defeq Z_s(X)=\exp\Bigl\{\int_0^s(X_\theta,dB_\theta)-\frac12\int_0^s\norm{X_\theta}^2\,d\theta\Bigr\},
\]
which for $0\le s\le\ell$ is exactly the exponential appearing in the statement. We have to prove
\begin{equation}\label{eq:5-7-11:goal}
  v(t,x)=E^x\Bigl[f(W_\ell)Z_\ell+\int_0^\ell g(t+s,W_s)Z_s\,ds\Bigr].
\end{equation}

\emph{Step 1: two maximal estimates.} Put
\[
  Z^*\defeq\max_{0\le s\le\ell}Z_s,\qquad \Gamma\defeq\max_{0\le s\le\ell}e^{\mu\norm{W_s}^2}.
\]
We claim that $E^x[(Z^*)^2]<\infty$ and $E^x[\Gamma^2]<\infty$, so that $E^x[\Gamma Z^*]<\infty$ by the Cauchy--Schwarz inequality.

Since $X$ and $2X$ are bounded, they satisfy the Novikov condition \bk{(3.5.20)}, and \bk{Corollary 3.5.13} shows that $Z(X)$ and $Z(2X)$ are martingales; in particular $E^x[Z_s(2X)]=1$. Directly from the definition,
\[
  Z_s^2=Z_s(2X)\exp\Bigl\{\int_0^s\norm{X_\theta}^2\,d\theta\Bigr\}\le e^{\beta^2 s}Z_s(2X),
\]
so $E^x[Z_s^2]\le e^{\beta^2 s}$. Doob's maximal inequality (\bk{Theorem 1.3.8}\,(iv) with $p=2$), applied to the nonnegative continuous martingale $Z$, gives
\begin{equation}\label{eq:5-7-11:Z}
  E^x[(Z^*)^2]\le 4E^x[Z_\ell^2]\le 4e^{\beta^2 T}.
\end{equation}

For the second estimate, recall that for a standard $d$-dimensional Brownian motion $B$ and $a<1/(2s)$ we have $E^x[e^{a\norm{B_s}^2}]=(1-2as)^{-d/2}$. By \eqref{eq:5-7-11:mu}, $Y_s\defeq e^{2\mu\norm{B_s}^2}$, $0\le s\le\ell$, satisfies $E^x[Y_s^2]=(1-8\mu s)^{-d/2}<\infty$. Since $E^x[e^{2\mu\norm{B_s}^2}]=\infty$ for $s\ge1/(4\mu)$, we work with the stopped process $B_{s\wedge\ell}$, $0\le s<\infty$. Its components are $\{\G_s\}$-martingales on $[0,\infty)$: for $s<u$ we have $E^x[B^{(i)}_{u\wedge\ell}\,|\,\G_s]=B^{(i)}_{s\wedge\ell}$, trivially if $s\ge\ell$ and by the martingale property of $B^{(i)}$ if $s<\ell$. The function $\varphi(y)\defeq e^{2\mu\norm{y}^2}$ is convex on $\R^d$, and $E^x[\varphi(B_{s\wedge\ell})]=(1-4\mu(s\wedge\ell))^{-d/2}\le(1-4\mu\ell)^{-d/2}<\infty$ for every $s\ge0$ by \eqref{eq:5-7-11:mu}. Therefore \bk{Problem 1.3.7} shows that $\{Y_{s\wedge\ell},\G_s;\ 0\le s<\infty\}$ is a nonnegative submartingale; its paths are continuous, and $E^x[Y_\ell^2]<\infty$. Using $\norm{W_s}^2\le 2\norm{x}^2+2\norm{B_s}^2$ and Doob's maximal inequality (\bk{Theorem 1.3.8}\,(iv) with $p=2$ and $[\sigma,\tau]=[0,\ell]$) applied to $Y_{\cdot\wedge\ell}$,
\begin{equation}\label{eq:5-7-11:G}
  E^x[\Gamma^2]\le e^{4\mu\norm{x}^2}E^x\Bigl[\max_{0\le s\le\ell}Y_s^2\Bigr]
  \le 4e^{4\mu\norm{x}^2}E^x[Y_\ell^2]=4e^{4\mu\norm{x}^2}(1-8\mu\ell)^{-d/2}<\infty .
\end{equation}

\emph{Step 2: It\^o's rule.} Fix $S\in(0,\ell)$ and let
\[
  h(s)\defeq s\ \ (0\le s\le S),\qquad
  h(s)\defeq S+\frac{\ell-S}{2}\Bigl(1-\exp\Bigl\{-\frac{2(s-S)}{\ell-S}\Bigr\}\Bigr)\ \ (s>S).
\]
This is a $C^1$ function from $[0,\infty)$ into $[0,(S+\ell)/2]\subset[0,\ell)$, so $u(s,y)\defeq v(t+h(s),y)$ is of class $C^{1,2}$ on $[0,\infty)\times\R^d$ and $F(s,y,z)\defeq u(s,y)\,z$ is of class $C^{1,2}$ on $[0,\infty)\times\R^{d+1}$. On $[0,S]\times\R^d$ we have $u(s,y)=v(t+s,y)$, and the partial differential equation gives
\begin{equation}\label{eq:5-7-11:pde}
  \frac{\partial u}{\partial s}(s,y)+\frac12\Delta u(s,y)+\bigl(b(t+s,y),\nabla u(s,y)\bigr)=-g(t+s,y);
  \qquad 0\le s\le S,\ y\in\R^d .
\end{equation}
The process $(W^{(1)},\dots,W^{(d)},Z)$ is a vector of continuous semimartingales: $W=x+B$ and, by \bk{(3.5.3)}, $Z_s=1+\sum_{i=1}^d\int_0^sZ_\theta X^{(i)}_\theta\,dB^{(i)}_\theta$. From $\qv{B^{(i)},B^{(j)}}_s=\delta_{ij}s$ (\bk{Problem 2.5.5}) and \bk{Proposition 3.2.24} we get $\qv{B^{(i)},Z}_s=\int_0^sZ_\theta X^{(i)}_\theta\,d\theta$, where we have used that $\qv{\int_0^\cdot Z_\theta X^{(i)}_\theta\,dB^{(i)}_\theta,N}=\int_0^\cdot Z_\theta X^{(i)}_\theta\,d\qv{B^{(i)},N}_\theta$ for every $N\in\cM^{c,\loc}$, and hence $d\qv{Z,N}_\theta=\sum_{i}Z_\theta X^{(i)}_\theta\,d\qv{B^{(i)},N}_\theta$. Moreover, stochastic integrals with respect to $Z$ (that is, with respect to $Z-1\in\cM^{c,\loc}$) are stochastic integrals with respect to the $B^{(i)}$ with the integrand multiplied by $ZX^{(i)}$: this is the local-martingale version of \bk{Corollary 3.2.20}, and it follows from \bk{Proposition 3.2.24} by the same argument. Indeed, let $Y$ be continuous and adapted; then $\int_0^R Y_\theta^2\,d\qv{Z}_\theta=\int_0^RY_\theta^2Z_\theta^2\norm{X_\theta}^2\,d\theta<\infty$ and $\int_0^R(Y_\theta Z_\theta X^{(i)}_\theta)^2\,d\theta<\infty$ for every $R>0$, a.s.\ (so $Y\in\cP^*(Z-1)$ and $YZX^{(i)}\in\cP^*(B^{(i)})$), and $\Psi\defeq\sum_i\int_0^\cdot Y_\theta Z_\theta X^{(i)}_\theta\,dB^{(i)}_\theta\in\cM^{c,\loc}$ satisfies, for every $N\in\cM^{c,\loc}$,
\[
  \qv{\Psi,N}_s=\sum_{i=1}^d\int_0^sY_\theta Z_\theta X^{(i)}_\theta\,d\qv{B^{(i)},N}_\theta=\int_0^sY_\theta\,d\qv{Z,N}_\theta ,
\]
so that $\Psi=\int_0^\cdot Y_\theta\,dZ_\theta$ by the uniqueness assertion of \bk{Proposition 3.2.24}. Since $\partial^2F/\partial z^2=0$, $\partial F/\partial z=u$ and $\partial^2F/\partial y_i\partial z=\partial u/\partial y_i$, \bk{Theorem 3.3.6} yields, a.s.\ for all $s\ge0$,
\begin{align*}
  u(s,W_s)Z_s={}&v(t,x)
  +\int_0^sZ_\theta\Bigl[\frac{\partial u}{\partial s}+\frac12\Delta u+(X_\theta,\nabla u)\Bigr](\theta,W_\theta)\,d\theta\\
  &+\sum_{i=1}^d\int_0^sZ_\theta\Bigl[\frac{\partial u}{\partial y_i}(\theta,W_\theta)
     +u(\theta,W_\theta)X^{(i)}_\theta\Bigr]dB^{(i)}_\theta .
\end{align*}
For $\theta\le S$ we have $X_\theta=b(t+\theta,W_\theta)$, so by \eqref{eq:5-7-11:pde} the first integrand equals $-Z_\theta\,g(t+\theta,W_\theta)$.

Now let $n>\norm{x}$ be an integer, $\tau_n\defeq\inf\{s\ge0;\ \norm{W_s}\ge n\}$ (a stopping time, by the continuity of the paths) and $\sigma\defeq S\wedge\tau_n$. Evaluating at $s=\sigma$,
\begin{equation}\label{eq:5-7-11:ito}
  v(t+\sigma,W_\sigma)Z_\sigma+\int_0^\sigma g(t+\theta,W_\theta)Z_\theta\,d\theta
  =v(t,x)+\sum_{i=1}^d\int_0^S\Phi^{(i)}_\theta\,dB^{(i)}_\theta,
\end{equation}
where $\Phi^{(i)}_\theta\defeq\ind_{\{\theta\le\sigma\}}Z_\theta\bigl[\partial u/\partial y_i(\theta,W_\theta)+u(\theta,W_\theta)X^{(i)}_\theta\bigr]$. Here we have used the identity \bk{(3.2.24)}, $I_{s\wedge\sigma}(Y)=I_s(Y\ind_{\{\cdot\le\sigma\}})$ a.s., taken at $s=S$ (note $S\wedge\sigma=\sigma$) for the integrands $Y=Z\bigl[\partial u/\partial y_i(\cdot,W)+u(\cdot,W)X^{(i)}\bigr]$ with respect to $B^{(i)}$; these integrands are progressively measurable with a.s.\ locally bounded paths, hence belong to $\cP^*(B^{(i)})$, and \bk{(3.2.24)} remains valid for integrands in $\cP^*$ with respect to local martingales, as stated in the remark preceding \bk{Proposition 3.2.24}. On $\{\theta\le\sigma\}$ we have $\theta\le S$ and $\norm{W_\theta}\le n$ ($P^x$-a.s., because $W_0=x$), and $u$, $\nabla u$ are continuous, hence bounded by some constant $K$, on the compact set $[0,S]\times\{\norm{y}\le n\}$. Thus $\abs{\Phi^{(i)}_\theta}\le K(1+\beta)Z_\theta$ and, by Step~1,
\[
  E^x\int_0^S\bigl(\Phi^{(i)}_\theta\bigr)^2d\theta\le K^2(1+\beta)^2\,S\,e^{\beta^2T}<\infty .
\]
Hence the stochastic integrals in \eqref{eq:5-7-11:ito} are square-integrable martingales evaluated at time $S$ and have expectation zero. The same bounds show that both terms on the left of \eqref{eq:5-7-11:ito} are integrable ($g$ is bounded on $[t,t+S]\times\{\norm{y}\le n\}$). Taking expectations,
\begin{equation}\label{eq:5-7-11:pre}
  v(t,x)=E^x\bigl[v(t+\sigma,W_\sigma)Z_\sigma\bigr]+E^x\int_0^\sigma g(t+\theta,W_\theta)Z_\theta\,d\theta ;
  \qquad\sigma=S\wedge\tau_n .
\end{equation}

\emph{Step 3: passage to the limit.} Apply \eqref{eq:5-7-11:pre} with $S=S_n\defeq\ell(1-1/n)$, and write $\sigma_n\defeq S_n\wedge\tau_n$ for integers $n>\max\{\norm{x},1\}$. The sequence $\{\sigma_n\}$ is nondecreasing; since the path $s\mapsto W_s$ is bounded on $[0,\ell]$, we have $\tau_n>\ell$ for all large $n$ (depending on $\omega$), and therefore $\sigma_n\uparrow\ell$, $P^x$-a.s.

\emph{The first term.} By the continuity of $v$ on $[0,T]\times\R^d$ and of the paths of $W$ and $Z$,
\[
  v(t+\sigma_n,W_{\sigma_n})Z_{\sigma_n}\xrightarrow[n\to\infty]{}v(T,W_\ell)Z_\ell=f(W_\ell)Z_\ell,\qquad P^x\text{-a.s.}
\]
By \bk{(7.20)}, $\abs{v(t+\sigma_n,W_{\sigma_n})Z_{\sigma_n}}\le M\,\Gamma Z^*$, which is integrable by Step~1. Dominated convergence gives
\[
  \lim_{n\to\infty}E^x\bigl[v(t+\sigma_n,W_{\sigma_n})Z_{\sigma_n}\bigr]=E^x\bigl[f(W_\ell)Z_\ell\bigr],
\]
and $f(W_\ell)Z_\ell$ is integrable.

\emph{The second term.} The integrand $(\theta,\omega)\mapsto g(t+\theta,W_\theta)Z_\theta$ is continuous in $\theta$ and adapted, hence jointly measurable. If $g$ satisfies \bk{(7.11)}\,(i), let $C\defeq\sup_{r\ge0}(1+r^{2\lambda})e^{-\mu r^2}<\infty$; then
\[
  \abs{g(t+\theta,W_\theta)}Z_\theta\le LC\,\Gamma Z^*;\qquad 0\le\theta\le\ell,
\]
and the right-hand side is integrable with respect to $d\theta\times P^x$ on $[0,\ell]\times\Omega$. Since $\ind_{\{\theta\le\sigma_n\}}\to\ind_{\{\theta\le\ell\}}$ for a.e.\ $\theta$, dominated convergence yields
\begin{equation}\label{eq:5-7-11:g}
  \lim_{n\to\infty}E^x\int_0^{\sigma_n}g(t+\theta,W_\theta)Z_\theta\,d\theta
  =E^x\int_0^{\ell}g(t+\theta,W_\theta)Z_\theta\,d\theta .
\end{equation}
If instead $g$ satisfies \bk{(7.11)}\,(ii), the random variables $\int_0^{\sigma_n}g(t+\theta,W_\theta)Z_\theta\,d\theta$ are nonnegative and nondecreasing in $n$, with limit $\int_0^\ell g(t+\theta,W_\theta)Z_\theta\,d\theta$, and \eqref{eq:5-7-11:g} holds by monotone convergence; the limit is finite, being equal to $v(t,x)-E^x[f(W_\ell)Z_\ell]$ by \eqref{eq:5-7-11:pre} and the previous paragraph.

Letting $n\to\infty$ in \eqref{eq:5-7-11:pre} with $\sigma=\sigma_n$, we obtain \eqref{eq:5-7-11:goal}; in both cases the random variable inside the expectation is integrable, and the expectation may be split into the two terms treated above. This is the asserted representation.

\rmk{(a) The proof uses the bound \bk{(7.20)} only with $\mu<1/(8T)$, cf.\ \eqref{eq:5-7-11:mu}, and it does not use \bk{(7.10)} at all: since $f=v(T,\cdot)$, the terminal function automatically satisfies $\abs{f(y)}\le Me^{\mu\norm{y}^2}$. (b) By the Girsanov theorem (\bk{Theorem 3.5.1}), under the probability measure $\wt P(A)\defeq E^x[\ind_AZ_\ell]$ on $\G_\ell$ the process $W$ solves on $[0,\ell]$ the equation $dW_s=b(t+s,W_s)\,ds+d\wt W_s$ with a Brownian motion $\wt W$, and since $Z$ is a martingale, the right-hand side of \eqref{eq:5-7-11:goal} equals $\wt E\bigl[f(W_\ell)+\int_0^\ell g(t+s,W_s)\,ds\bigr]$. Thus the formula is the representation \bk{(7.15)} (with $k\equiv0$ and the identity as dispersion matrix) for the diffusion with drift $b$, written on the Wiener space; the direct argument above has the advantage of requiring neither the continuity of $b$ nor the conditions \bk{(7.2)}--\bk{(7.4)} of \bk{Theorem 7.6}.}
\end{kex}

\begin{kex}{Exercise}{7.12}{}
\emph{Step 1: the process, its generator and the two equations.} One-dimensional Brownian motion with constant drift $\mu\in\R$ is the solution of \bk{(7.1)} with $d=r=1$, $b(t,x)\equiv\mu$ and $\sigma(t,x)\equiv 1$, namely
\begin{equation}\label{eq:5-7-12:X}
  X^{(t,x)}_s = x+\mu(s-t)+(W_s-W_t);\qquad t\le s<\infty .
\end{equation}
These coefficients are constant, so \bk{(7.2)} holds, and \eqref{eq:5-7-12:X} is the only process satisfying \bk{(7.1)}, so \bk{(7.3)}, \bk{(7.4)} hold as well. Here $a(t,x)=\sigma^2(t,x)\equiv 1$, and the operator $\cA_t$ of \bk{(4.1)} does not depend on $t$:
\[
  \cA f(x)=\tfrac12 f''(x)+\mu f'(x) .
\]
With $k\equiv 0$, the backward Kolmogorov equation \bk{(7.21)} in the backward variables $(t,x)$ reads
\begin{equation}\label{eq:5-7-12:back}
  -\frac{\partial u}{\partial t}(t,x)=\frac12\,\frac{\partial^2 u}{\partial x^2}(t,x)+\mu\,\frac{\partial u}{\partial x}(t,x),
\end{equation}
and the forward Kolmogorov equation \bk{(7.24)} in the forward variables $(\tau,\xi)$ reads, because $a\equiv 1$ and $b\equiv\mu$ are constants,
\begin{equation}\label{eq:5-7-12:forw}
  \frac{\partial \psi}{\partial \tau}(\tau,\xi)=\frac12\,\frac{\partial^2 \psi}{\partial \xi^2}(\tau,\xi)-\mu\,\frac{\partial \psi}{\partial \xi}(\tau,\xi) .
\end{equation}

\emph{Step 2: the transition probability density.} Let $0\le t<\tau<\infty$ and $x\in\R$. In \eqref{eq:5-7-12:X} the increment $W_\tau-W_t$ is normal with mean zero and variance $\tau-t>0$, so $X^{(t,x)}_\tau$ is normal with mean $x+\mu(\tau-t)$ and variance $\tau-t$. Consequently
\begin{equation}\label{eq:5-7-12:G}
  P\bigl[X^{(t,x)}_\tau\in A\bigr]=\int_A G(t,x;\tau,\xi)\,d\xi;\qquad A\in\B(\R),
\end{equation}
where
\[
  G(t,x;\tau,\xi)\defeq\frac{1}{\sqrt{2\pi(\tau-t)}}\exp\Bigl\{-\frac{\bigl(\xi-x-\mu(\tau-t)\bigr)^2}{2(\tau-t)}\Bigr\} .
\]
This is \bk{(7.26)}: $G$ is the transition probability density of the process. In terms of the Gaussian kernel
\[
  p(s,y)\defeq\frac{1}{\sqrt{2\pi s}}\,e^{-y^2/2s};\qquad s>0,\ y\in\R,
\]
we have
\begin{equation}\label{eq:5-7-12:Gp}
  G(t,x;\tau,\xi)=p\bigl(\tau-t,\ \xi-x-\mu(\tau-t)\bigr);\qquad t<\tau,\ x\in\R,\ \xi\in\R .
\end{equation}
The function $p$ is infinitely differentiable on $(0,\infty)\times\R$. From $\log p(s,y)=-\tfrac12\log(2\pi s)-y^2/2s$ we compute
\[
  \frac{\partial p}{\partial y}=-\frac{y}{s}\,p,\qquad
  \frac{\partial^2 p}{\partial y^2}=\Bigl(\frac{y^2}{s^2}-\frac1s\Bigr)p,\qquad
  \frac{\partial p}{\partial s}=\Bigl(\frac{y^2}{2s^2}-\frac{1}{2s}\Bigr)p,
\]
so that $p$ solves the heat equation
\begin{equation}\label{eq:5-7-12:heat}
  \frac{\partial p}{\partial s}=\frac12\,\frac{\partial^2 p}{\partial y^2}\qquad\text{on }(0,\infty)\times\R .
\end{equation}

\emph{Step 3: the backward equation.} Fix $(\tau,\xi)\in(0,\infty)\times\R$ and set $\varphi(t,x)\defeq G(t,x;\tau,\xi)$ for $0\le t<\tau$, $x\in\R$. By \eqref{eq:5-7-12:Gp}, $\varphi$ is the composition of $p$ with the affine map $(t,x)\mapsto(s,y)=(\tau-t,\ \xi-x-\mu(\tau-t))$, for which $\partial s/\partial t=-1$, $\partial y/\partial t=\mu$, $\partial s/\partial x=0$, $\partial y/\partial x=-1$. Hence $\varphi$ is of class $C^{1,2}([0,\tau)\times\R)$ (indeed $C^\infty$), and the chain rule gives, with all derivatives of $p$ evaluated at $(s,y)$,
\[
  \frac{\partial\varphi}{\partial t}=-\frac{\partial p}{\partial s}+\mu\,\frac{\partial p}{\partial y},\qquad
  \frac{\partial\varphi}{\partial x}=-\frac{\partial p}{\partial y},\qquad
  \frac{\partial^2\varphi}{\partial x^2}=\frac{\partial^2 p}{\partial y^2}.
\]
Using \eqref{eq:5-7-12:heat} we obtain
\[
  -\frac{\partial\varphi}{\partial t}
  =\frac{\partial p}{\partial s}-\mu\,\frac{\partial p}{\partial y}
  =\frac12\,\frac{\partial^2 p}{\partial y^2}-\mu\,\frac{\partial p}{\partial y}
  =\frac12\,\frac{\partial^2\varphi}{\partial x^2}+\mu\,\frac{\partial\varphi}{\partial x},
\]
which is \eqref{eq:5-7-12:back}: for fixed $(\tau,\xi)$, the transition density satisfies the backward Kolmogorov equation in the backward variables $(t,x)$.

\emph{Step 4: the forward equation.} Fix $(t,x)\in[0,\infty)\times\R$ and set $\psi(\tau,\xi)\defeq G(t,x;\tau,\xi)$ for $\tau>t$, $\xi\in\R$. Now $\psi$ is the composition of $p$ with the affine map $(\tau,\xi)\mapsto(s,y)=(\tau-t,\ \xi-x-\mu(\tau-t))$, for which $\partial s/\partial\tau=1$, $\partial y/\partial\tau=-\mu$, $\partial s/\partial\xi=0$, $\partial y/\partial\xi=1$. Hence $\psi$ is of class $C^{1,2}((t,\infty)\times\R)$ and
\[
  \frac{\partial\psi}{\partial\tau}=\frac{\partial p}{\partial s}-\mu\,\frac{\partial p}{\partial y},\qquad
  \frac{\partial\psi}{\partial\xi}=\frac{\partial p}{\partial y},\qquad
  \frac{\partial^2\psi}{\partial\xi^2}=\frac{\partial^2 p}{\partial y^2}.
\]
By \eqref{eq:5-7-12:heat},
\[
  \frac{\partial\psi}{\partial\tau}
  =\frac12\,\frac{\partial^2 p}{\partial y^2}-\mu\,\frac{\partial p}{\partial y}
  =\frac12\,\frac{\partial^2\psi}{\partial\xi^2}-\mu\,\frac{\partial\psi}{\partial\xi},
\]
which is \eqref{eq:5-7-12:forw}: for fixed $(t,x)$, the transition density satisfies the forward Kolmogorov equation in the forward variables $(\tau,\xi)$.

\rmk{The function $G$ of \eqref{eq:5-7-12:G} is in fact a fundamental solution of \eqref{eq:5-7-12:back} in the sense of \bk{Definition 7.9}, for every $T>0$. It is nonnegative. Let $\tau\in(0,T]$, let $f\colon\R\to\R$ be continuous with compact support $K$, and let $u(t,x)\defeq\int_{\R}G(t,x;\tau,\xi)f(\xi)\,d\xi$ as in \bk{(7.22)}. Since $G(t,x;\tau,\cdot)$ is a probability density, $\abs{u}\le\max\abs{f}$. The functions $G$, $\partial G/\partial t$, $\partial G/\partial x$, $\partial^2G/\partial x^2$ are continuous in $(t,x,\xi)$ on $[0,\tau)\times\R\times\R$, hence bounded on $C\times K$ for every compact $C\subset[0,\tau)\times\R$; thus we may differentiate under the integral sign over $K$ (dominated convergence), the resulting derivatives of $u$ are continuous, $u$ is of class $C^{1,2}([0,\tau)\times\R)$, and $u$ satisfies \eqref{eq:5-7-12:back} by Step 3. Finally, by \eqref{eq:5-7-12:G}, $u(t,x)=Ef\bigl(x+\mu(\tau-t)+\sqrt{\tau-t}\,Z\bigr)$ with $Z$ standard normal, and this converges to $f(x)$ as $t\uparrow\tau$ by the bounded convergence theorem and the continuity of $f$; this is \bk{(7.23)}.}
\end{kex}

\begin{kex}{Exercise}{7.13}{}
Throughout, the standing assumptions \bk{(7.2)}--\bk{(7.4)} are in force, $a=\sigma\sigma^T$, and $\cA$ is the operator of \bk{(1.2)}. We fix $r>0$ and $x\in\R^d$ with $\norm{x}>r$, and a weak solution $(X,W)$, $(\Omega,\F,P^x)$, $\{\F_t\}$ of \bk{(7.1)} with initial time $0$ and $X_0=x$; thus $\{\F_t\}$ satisfies the usual conditions, $W$ is an $r'$-dimensional Brownian motion ($r'$ being the number of columns of $\sigma$), $X$ is continuous and adapted, and
\begin{equation}\label{eq:5-7-13:sde}
  X_t=x+\int_0^t b(X_s)\,ds+\int_0^t\sigma(X_s)\,dW_s;\qquad 0\le t<\infty,
\end{equation}
almost surely. By \bk{(7.4)} the law of $X$, and hence each of the quantities $P^x(\tau_r<\infty)$ and $E^x\tau_r$, does not depend on the choice of the weak solution. We write $\tau\defeq\tau_r$ and, for $R>\norm{x}$,
\[
  \sigma_R\defeq\inf\{t\ge 0;\ \norm{X_t}\ge R\},\qquad \rho_R\defeq\tau\wedge\sigma_R,\qquad
  K_R\defeq\{y\in\R^d;\ r\le\norm{y}\le R\}.
\]
Both $\tau$ and $\sigma_R$ are hitting times of closed sets by a continuous, adapted process, hence stopping times of $\{\F_t\}$ (\bk{Problem 1.2.7}), and so is $\rho_R$. The following facts are consequences of the continuity of the paths of $X$ and of $r<\norm{x}<R$:
\begin{gather}
  X_s\in K_R\quad\text{for } 0\le s\le\rho_R,\ s<\infty;\label{eq:5-7-13:in}\\
  \norm{X_\tau}=r\ \text{ on }\{\tau<\infty\},\qquad \norm{X_{\sigma_R}}=R\ \text{ on }\{\sigma_R<\infty\};\label{eq:5-7-13:bdry}\\
  \sigma_R\uparrow\infty\quad\text{as } R\uparrow\infty .\label{eq:5-7-13:expl}
\end{gather}
Indeed, $r<\norm{X_s}<R$ for $s<\rho_R$ by the definition of the two times, and the closed sets $\ol B_r$, $\{\norm{y}\ge R\}$ contain $X_\tau$, $X_{\sigma_R}$, respectively, when these times are finite; this gives \eqref{eq:5-7-13:in}, \eqref{eq:5-7-13:bdry}. For \eqref{eq:5-7-13:expl}, $\sigma_R$ is nondecreasing in $R$, and $\sigma_R>t$ as soon as $R>\max_{0\le s\le t}\norm{X_s}$, which is finite. In particular, $\tau\neq\sigma_R$ on $\{\rho_R<\infty\}$.

\emph{Step 1: a supermartingale-type inequality.} Suppose that, for some constant $c\ge 0$, we have $\cA f\le -c$ on $\R^d\setminus\{0\}$. We claim that for every $R>\norm{x}$ and $0\le t<\infty$,
\begin{equation}\label{eq:5-7-13:key}
  E^x f(X_{t\wedge\rho_R})+c\,E^x(t\wedge\rho_R)\le f(x).
\end{equation}
Since $f$ is defined only away from the origin, we first modify it near the origin. Let $\psi\colon\R\to[0,1]$ be a function of class $C^\infty$ with $\psi=0$ on $(-\infty,r^2/9]$ and $\psi=1$ on $[r^2/4,\infty)$, and set $g(y)\defeq\psi(\norm{y}^2)f(y)$ for $y\neq 0$, $g(0)\defeq 0$. Then $g$ vanishes on the ball $\{\norm{y}\le r/3\}$ and is of class $C^2$ on $\R^d\setminus\{0\}$, so $g\in C^2(\R^d)$; and $g=f$ on the open set $\{\norm{y}>r/2\}\supset K_R$, so that $g$, its first and second partial derivatives, and $\cA g$ agree there with those of $f$. By \eqref{eq:5-7-13:sde}, $X^{(i)}=x_i+M^{(i)}+B^{(i)}$ with
\[
  M^{(i)}_t=\sum_{j}\int_0^t\sigma_{ij}(X_s)\,dW^{(j)}_s\in\cM^{c,\loc},\qquad
  B^{(i)}_t=\int_0^t b_i(X_s)\,ds,
\]
and $\qv{M^{(i)},M^{(k)}}_t=\int_0^t a_{ik}(X_s)\,ds$, the integrands being continuous and adapted. It\^o's rule (\bk{Theorem 3.3.6}) applied to $g$ gives, a.s.,
\[
  g(X_t)=g(x)+\int_0^t\cA g(X_s)\,ds+\sum_{i=1}^d\sum_{j}\int_0^t\frac{\partial g}{\partial x_i}(X_s)\,\sigma_{ij}(X_s)\,dW^{(j)}_s;
  \qquad 0\le t<\infty.
\]
We evaluate this at $t\wedge\rho_R$ and use \eqref{eq:5-7-13:in} to replace $g$ by $f$:
\begin{equation}\label{eq:5-7-13:ito}
  f(X_{t\wedge\rho_R})=f(x)+\int_0^{t\wedge\rho_R}\cA f(X_s)\,ds
  +\sum_{i=1}^d\sum_{j}\int_0^{t}\ind_{\{s\le\rho_R\}}\frac{\partial f}{\partial x_i}(X_s)\,\sigma_{ij}(X_s)\,dW^{(j)}_s .
\end{equation}
The functions $(\partial f/\partial x_i)\sigma_{ij}$ are continuous on the compact set $K_R$, hence bounded there; by \eqref{eq:5-7-13:in} the integrands of the stochastic integrals in \eqref{eq:5-7-13:ito} are bounded, progressively measurable processes, so these integrals are square-integrable martingales starting at zero and have expectation zero. Likewise $f$ is bounded on $K_R$, so $f(X_{t\wedge\rho_R})$ is integrable, and $\int_0^{t\wedge\rho_R}\cA f(X_s)\,ds\le -c\,(t\wedge\rho_R)$ by \eqref{eq:5-7-13:in} and the hypothesis. Taking expectations in \eqref{eq:5-7-13:ito} yields \eqref{eq:5-7-13:key}.

\textbf{(i)} Fix $R>\norm{x}$. The ball $B_R$ is an open, bounded domain containing $x$, and by hypothesis it satisfies \bk{(7.9)} for some $1\le\ell\le d$; its exit time $\inf\{t\ge0;\ X_t\notin B_R\}$ is $\sigma_R$. Hence \bk{Lemma 7.4} gives $E^x\sigma_R<\infty$, and in particular
\begin{equation}\label{eq:5-7-13:fin}
  \rho_R\le\sigma_R<\infty\quad\text{a.s. } P^x.
\end{equation}
We use \eqref{eq:5-7-13:key} with $c=0$, which is \bk{(7.31)}: $E^xf(X_{t\wedge\rho_R})\le f(x)$ for all $t\ge0$. As $t\to\infty$ we have $f(X_{t\wedge\rho_R})\to f(X_{\rho_R})$ a.s.\ by \eqref{eq:5-7-13:fin}, and these random variables are bounded by $\max_{K_R}\abs{f}$; the bounded convergence theorem gives
\begin{equation}\label{eq:5-7-13:stop}
  E^xf(X_{\rho_R})\le f(x).
\end{equation}
By \eqref{eq:5-7-13:bdry}, \eqref{eq:5-7-13:fin} and the definition of $F$, we have a.s.\ $f(X_{\rho_R})\ge F(R)$ on $\{\sigma_R<\tau\}$ and $f(X_{\rho_R})\ge F(r)$ on $\{\tau<\sigma_R\}$, and these two events exhaust $\Omega$ up to a null set. Therefore \eqref{eq:5-7-13:stop} implies
\[
  F(R)\,P^x(\sigma_R<\tau)+F(r)\,\bigl(1-P^x(\sigma_R<\tau)\bigr)\le f(x),
\]
and, since $F(R)>F(r)$,
\begin{equation}\label{eq:5-7-13:prob}
  P^x(\sigma_R<\tau)\le\frac{f(x)-F(r)}{F(R)-F(r)} .
\end{equation}
Finally, $\{\tau=\infty\}\subseteq\{\sigma_R<\tau\}\cup\{\sigma_R=\infty\}$, and the last event is null by \eqref{eq:5-7-13:fin}. Thus $P^x(\tau=\infty)$ is bounded by the right-hand side of \eqref{eq:5-7-13:prob} for every $R>\norm{x}$; letting $R\to\infty$ and using $F(R)\to\infty$ we obtain $P^x(\tau_r=\infty)=0$, which is \bk{(7.32)}.

\textbf{(ii)} We take $f(x)\defeq\log\norm{x}$ on $\R^d\setminus\{0\}$. This function is of class $C^\infty$ there, $F(r)=\log r$ is strictly increasing, and $\lim_{r\to\infty}F(r)=\infty$. For $x\neq0$ we compute
\[
  \frac{\partial f}{\partial x_i}(x)=\frac{x_i}{\norm{x}^2},\qquad
  \frac{\partial^2 f}{\partial x_i\partial x_k}(x)=\frac{\delta_{ik}}{\norm{x}^2}-\frac{2x_ix_k}{\norm{x}^4},
\]
so that
\[
  \cA f(x)=\frac{1}{\norm{x}^2}\biggl[\frac12\sum_{i=1}^d a_{ii}(x)-\frac{(x,a(x)x)}{\norm{x}^2}+(x,b(x))\biggr].
\]
Under \bk{(7.33)} the bracket is nonpositive, i.e., \bk{(7.31)} holds, and \bk{(7.32)} follows from part (i).

\textbf{(iii)} Now $\cA f\le-1$ on $\R^d\setminus\{0\}$, and \eqref{eq:5-7-13:key} with $c=1$ reads
\[
  E^x(t\wedge\rho_R)\le f(x)-E^xf(X_{t\wedge\rho_R});\qquad R>\norm{x},\ 0\le t<\infty .
\]
By \eqref{eq:5-7-13:in} we have $\norm{X_{t\wedge\rho_R}}\ge r$, and because $F$ is increasing,
$f(X_{t\wedge\rho_R})\ge F(\norm{X_{t\wedge\rho_R}})\ge F(r)$. Consequently
\[
  E^x(t\wedge\tau\wedge\sigma_R)\le f(x)-F(r);\qquad R>\norm{x},\ 0\le t<\infty,
\]
a bound which does not depend on $t$ or $R$. We let first $R\uparrow\infty$, then $t\uparrow\infty$; by \eqref{eq:5-7-13:expl} the random variables $t\wedge\tau\wedge\sigma_R$ increase to $\tau$, and the monotone convergence theorem gives
\[
  E^x\tau_r\le f(x)-F(r)<\infty ,
\]
which is \bk{(7.34)}.

\rmk{The proof of (iii) uses neither part (i) nor condition \bk{(7.9)}; it uses of $F$ only that it is nondecreasing. The proof of (i) uses \bk{(7.9)} only for the balls $B_R$, and of $F$ only that $F(R)\to\infty$; there \bk{(7.9)} cannot be dropped, since for $b\equiv0$, $\sigma\equiv0$ every $f$ satisfies \bk{(7.31)} while $X\equiv x$ never reaches $\ol B_r$. As a check on (ii): for Brownian motion ($b\equiv 0$, $a$ the identity matrix) condition \bk{(7.33)} reads $d/2\le1$, in agreement with the recurrence of Brownian motion in dimensions $d=1,2$ and its transience for $d\ge3$.}
\end{kex}

\solsection{5.8}{Applications to Economics}

\begin{kex}{Exercise}{8.10}{}
We assume $T>0$ (see the remark at the end), and we fix $x>0$. Recall that $\{\F_t\}$ is the augmentation of the filtration generated by the one-dimensional Brownian motion $W=\{W_t,\F_t;\,0\le t\le T\}$, so that it satisfies the usual conditions. With $r\equiv0$, $b_1\equiv0$, $\sigma\equiv1$, formula \bk{(8.13)} for the wealth process associated with a pair $(\pi,C)$ and the initial endowment $x$ reads
\begin{equation}\label{eq:5-8-10:a}
  X_t=x-\int_0^tC_s\,ds+\int_0^t\pi_s\,dW_s;\qquad 0\le t\le T .
\end{equation}
According to \bk{Definitions 8.2, 8.5} and \bk{8.9}, we have to find a measurable, adapted process $\pi$ with $\int_0^T\pi_t^2\,dt<\infty$ a.s., such that with $C\equiv0$ (which is forced by condition (i) of \bk{Definition 8.9}, and trivially satisfies \bk{(8.12)}) the process $X$ of \eqref{eq:5-8-10:a} satisfies $X_t\ge0$ for all $0\le t\le T$, and $X_T=0$, almost surely. The idea is to invest the amount $(T-s)^{-1/2}$ in the stock at time $s$, an amount which grows so fast as $s\uparrow T$ that the gains process is a Brownian motion run on the whole half-line, and to stop investing as soon as the wealth reaches zero.

Stochastic integrals with respect to $W$ are understood in the sense of Section 3.2, relative to the martingale $\{W_{t\wedge T},\F_{t\wedge T};\,0\le t<\infty\}\in\cM_2^c$, whose quadratic variation is $t\wedge T$; only times $t\le T$ will matter.

\emph{Step 1: a martingale on $[0,T)$.} Let
\[
  h(s)\defeq(T-s)^{-1/2}\ \text{ for }0\le s<T,\qquad h(T)\defeq 0,\qquad
  \theta(u)\defeq T(1-e^{-u})\ \text{ for }0\le u<\infty ,
\]
so that $\theta$ is a continuous, strictly increasing bijection of $[0,\infty)$ onto $[0,T)$ and
\begin{equation}\label{eq:5-8-10:b}
  \int_0^{t}h^2(s)\,ds=\log\frac{T}{T-t}\quad(0\le t<T),\qquad\text{in particular}\quad \int_0^{\theta(u)}h^2(s)\,ds=u\quad(u\ge 0).
\end{equation}
For $n\ge1$ put $t_n\defeq\theta(n)$ and $h_n\defeq h\,\ind_{[0,t_n]}$. Each $h_n$ is a bounded, nonrandom Borel function, hence a progressively measurable process in $\cL^*$, and by \bk{Proposition 3.2.10} the process $I(h_n)=\{\int_0^th_n(s)\,dW_s;\ 0\le t<\infty\}$ belongs to $\cM_2^c$ and has quadratic variation $\qv{I(h_n)}_t=\int_0^{t\wedge T}h_n^2(s)\,ds$ (see \bk{(3.2.19)}). For $m\ge n$, relation \bk{(3.2.24)} with the constant stopping time $t_n$ gives $I_{t\wedge t_n}(h_m)=I_t(h_m\ind_{[0,t_n]})=I_t(h_n)$ a.s.\ for every $t$, and, both sides being continuous in $t$, there is an event $\Omega_0$ with $P(\Omega_0)=1$ on which
\[
  I_t(h_m)=I_t(h_n)\qquad\text{for all } 0\le t\le t_n \text{ and all } m\ge n\ge1 .
\]
We define $M_t\defeq I_t(h_n)$ for $0\le t\le t_n$ on $\Omega_0$, and $M\equiv0$ on $\Omega\setminus\Omega_0$. Since $t_n\uparrow T$, this defines a process $M=\{M_t;\,0\le t<T\}$ with continuous paths and $M_0=0$; it is adapted because $\F_0$ contains all $P$-null sets. (Of course, $M_t=\int_0^t(T-s)^{-1/2}\,dW_s$.)

\emph{Step 2: $M$ is a time-changed Brownian motion.} Set
\[
  B_u\defeq M_{\theta(u)},\qquad \G_u\defeq\F_{\theta(u)};\qquad 0\le u<\infty .
\]
The process $B$ is continuous, adapted to $\{\G_u\}$, and $B_0=0$. The filtration $\{\G_u\}$ satisfies the usual conditions: $\G_0=\F_0$ contains the null sets, and $\bigcap_{v>u}\G_v=\bigcap_{t>\theta(u)}\F_t=\F_{\theta(u)}=\G_u$ because $\theta$ is continuous and strictly increasing. Let $0\le u<v$ and choose an integer $n\ge v$. Then $B_u=I_{\theta(u)}(h_n)$ and $B_v=I_{\theta(v)}(h_n)$ a.s., and since $I(h_n)$ and $I(h_n)^2-\qv{I(h_n)}$ are martingales with respect to $\{\F_t\}$, we obtain from \eqref{eq:5-8-10:b}
\[
  E[B_v\,|\,\G_u]=B_u,\qquad E[B_v^2-v\,|\,\G_u]=B_u^2-u,\qquad \text{a.s.}
\]
(both $B_v$ and $B_v^2$ are integrable because $I(h_n)\in\cM_2^c$). Thus $B$ is a continuous, square-integrable martingale relative to $\{\G_u\}$, and $B_u^2-u$ is a martingale; by the uniqueness of the quadratic variation, $\qv{B}_u=u$. L\'evy's \bk{Theorem 3.3.16} now shows that $\{B_u,\G_u;\,0\le u<\infty\}$ is a standard, one-dimensional Brownian motion. Consequently, by \bk{Remark 2.6.20}, its passage time to the level $-x$,
\[
  \rho\defeq\inf\{u\ge0;\ B_u=-x\},
\]
is almost surely finite.

\emph{Step 3: the strategy.} Define
\[
  \tau\defeq\inf\{t\in[0,T);\ M_t=-x\}\quad(\inf\emptyset\defeq T),\qquad
  \pi_s\defeq h(s)\,\ind_{\{s\le\tau\}}\ \ (0\le s\le T),\qquad C\equiv0 .
\]
Because $\theta$ is an increasing bijection of $[0,\infty)$ onto $[0,T)$, we have $\tau=\theta(\rho)$ on $\{\rho<\infty\}$, and therefore, by Step~2,
\begin{equation}\label{eq:5-8-10:c}
  P[\tau<T]=1 .
\end{equation}
By the continuity of $M$ and $M_0=0>-x$, we have on $\{\tau<T\}$:
\begin{equation}\label{eq:5-8-10:d}
  M_\tau=-x\qquad\text{and}\qquad M_t>-x\ \text{ for } 0\le t<\tau .
\end{equation}
Moreover, $\tau$ is a stopping time of $\{\F_t\}$: for $0\le t<T$ the continuity of $M$ gives
$\{\tau\le t\}=\{\min_{0\le s\le t}M_s\le -x\}=\{\inf_{s\in([0,t]\cap\Q)\cup\{t\}}M_s\le-x\}\in\F_t$, and $\{\tau\le T\}=\Omega$.

The process $\ind_{\{s\le\tau\}}$ is adapted, since $\{s\le\tau\}=\{\tau<s\}^c$ and $\{\tau<s\}=\bigcup_{k\ge1}\{\tau\le s-1/k\}\in\F_s$ (where $\{\tau\le s-1/k\}\defeq\emptyset$ if $s-1/k<0$), and it has left-continuous paths; it is therefore progressively measurable, and so is its product $\pi$ with the nonrandom Borel function $h$. In particular $\pi$ is measurable and adapted. From \eqref{eq:5-8-10:b} and \eqref{eq:5-8-10:c},
\[
  \int_0^T\pi_s^2\,ds=\int_0^{\tau}h^2(s)\,ds=\log\frac{T}{T-\tau}<\infty\qquad\text{a.s.,}
\]
which is \bk{(8.11)}. Hence $\pi$ is a portfolio process and $C$ a consumption process in the sense of \bk{Definition 8.2}, and the stochastic integral
\[
  N_t\defeq\int_0^t\pi_s\,dW_s;\qquad 0\le t\le T,
\]
is defined as in \bk{Definition 3.2.23}; it is a continuous local martingale, and the wealth process \eqref{eq:5-8-10:a} is $X_t=x+N_t$, $0\le t\le T$, with $X_0=x$.

\emph{Step 4: computation of the wealth.} We claim that, almost surely,
\begin{equation}\label{eq:5-8-10:e}
  N_t=M_{t\wedge\tau}\qquad\text{for all } 0\le t<T .
\end{equation}
Fix $n\ge1$. The sample path property \bk{(3.2.24)} holds for the integral of \bk{Definition 3.2.23} as well (see the remarks preceding \bk{Proposition 3.2.24}). Applying it to $\pi$ and the constant stopping time $t_n$, and then to $h_n\in\cL^*$ and the stopping time $\tau$, we get for every $t\ge0$, a.s.,
\[
  N_{t\wedge t_n}=\int_0^t\pi_s\ind_{\{s\le t_n\}}\,dW_s=\int_0^th_n(s)\ind_{\{s\le\tau\}}\,dW_s=I_{t\wedge\tau}(h_n),
\]
where the middle equality holds because the two integrands coincide. For $t\le t_n$ the left-hand side is $N_t$ and the right-hand side is $M_{t\wedge\tau}$, a.s. Both sides being continuous in $t$, the identity $N_t=M_{t\wedge\tau}$ holds a.s.\ simultaneously for all $t\in[0,t_n]$; letting $n$ run through the positive integers, we obtain \eqref{eq:5-8-10:e}.

\emph{Step 5: conclusion.} Let $\Omega_1$ be the event of probability one on which $\tau<T$, \eqref{eq:5-8-10:e} holds, and $t\mapsto N_t$ is continuous on $[0,T]$. On $\Omega_1$ we have, by \eqref{eq:5-8-10:d} and \eqref{eq:5-8-10:e},
\[
  X_t=x+M_{t\wedge\tau}\;
  \begin{cases}
    >0, & 0\le t<\tau,\\
    =0, & \tau\le t<T,
  \end{cases}
  \qquad\text{and}\qquad X_T=\lim_{t\uparrow T}X_t=x+M_\tau=0 .
\]
Thus $X_t\ge0$ for all $0\le t\le T$, a.s., which is \bk{(8.14)}: the pair $(\pi,C)$ is admissible for the initial endowment $x$. Furthermore $C_t=0=g_t$ for $0\le t\le T$ and $X_T=0=f_T$ a.s., which are conditions (i) and (ii) of \bk{Definition 8.9}. Therefore $(\pi,C)$ is a hedging strategy against the contingent claim $(g,f_T)=(0,0)$ with $X_0=x$.

\rmk{(a) The section allows $0\le T<\infty$, but for $T=0$ the claim fails, since then $X_T=X_0=x\ne0=f_T$; the hypothesis $T>0$ is needed. (b) The wealth process just constructed is a nonnegative continuous local martingale with $EX_T=0<x=X_0$, hence a supermartingale but not a martingale; in the present case $\wt P=P$, and this is the situation of strict inequality in \bk{(8.22)}. The strategy ``throws away'' the amount $x$: it shows that a claim may be hedged from several initial endowments, which is why the fair price is defined as the \emph{smallest} of them (here $0$).}
\end{kex}

\begin{kex}{Exercise}{8.13}{}
We work under condition \bk{(8.3)} and the assumptions preceding it, with a fixed contingent claim $(g,f_T)$ as in \bk{Definition 8.7}. As in Chapter~3, where the stochastic integral is constructed under this standing assumption (first paragraph of \bk{Section 3.2}), the filtration $\{\F_t\}$ is taken to satisfy the usual conditions of \bk{Definition 1.2.25}; we use only that every $\F_t$ contains all $P$-negligible events of $\F$. Section~5.8 rests on this convention: the stochastic integrals in \bk{(8.13)} and \bk{(8.20)} have integrands that are merely measurable and adapted, and such integrands are admitted by the argument of \bk{Remark 3.2.11}, which invokes exactly this property. (For the augmentation $\F_t=\sigma(\F^W_t\cup\cN)$ of $\{\F^W_t\}$ prescribed in Section~5.8, this property holds for instance when $\F$ is the $P$-completion of $\F^W_T$, since then every $P$-negligible event of $\F$ is contained in a $P$-null event of $\F^W_T$, hence belongs to $\cN$.)

\begin{lem}
Let $\varphi=\{\varphi_s;\ 0\le s\le T\}$ be a measurable, adapted process which is either nonnegative or bounded. Then, for every $t\in[0,T]$, the random variable $\Phi_t\defeq\int_0^t\varphi_s\,ds$ \textup{(}with values in $[0,\infty]$ in the nonnegative case\textup{)} is $\F_t$-measurable.
\end{lem}

\begin{proof}
Let $I$ be $[0,\infty)$, resp.\ $[-c,c]$ with $c$ a bound for $\abs{\varphi}$. By \bk{Proposition 1.1.12}, $\varphi$ has a progressively measurable modification; composing it with the continuous map $y\mapsto(y\vee0)$, resp.\ $y\mapsto(y\vee(-c))\wedge c$, which is the identity on $I$, we obtain a progressively measurable modification $\tilde\varphi$ of $\varphi$ with values in $I$. The restriction of $(s,\omega)\mapsto\tilde\varphi_s(\omega)$ to $[0,t]\times\Omega$ is $\B([0,t])\otimes\F_t$-measurable, so $\tilde\Phi_t\defeq\int_0^t\tilde\varphi_s\,ds$ is $\F_t$-measurable by the Fubini--Tonelli theorem, while $\Phi_t$ is $\F$-measurable for the same reason. The set $D\defeq\{(s,\omega)\in[0,t]\times\Omega;\ \varphi_s(\omega)\ne\tilde\varphi_s(\omega)\}$ belongs to $\B([0,t])\otimes\F$ and its $s$-sections are $P$-null, so by the Fubini--Tonelli theorem the event $N\defeq\{\omega;\ \meas\{s\in[0,t];\ (s,\omega)\in D\}>0\}$ belongs to $\F$ and $P(N)=0$. Off $N$ we have $\Phi_t=\tilde\Phi_t$. For a Borel set $B$, therefore, $\{\Phi_t\in B\}=(\{\tilde\Phi_t\in B\}\setminus N)\cup(\{\Phi_t\in B\}\cap N)$; the first set lies in $\F_t$, and the second is a $P$-negligible event of $\F$, hence lies in $\F_t$ as well.
\end{proof}

We write
\[
  \beta_t\defeq e^{-\int_0^t r(u)\,du};\qquad 0\le t\le T,
  \qquad\text{so that}\qquad
  Q=\beta_Tf_T+\int_0^T\beta_sg_s\,ds
\]
by \bk{(8.30)}. Since $r$ is measurable, adapted and bounded, say $\abs{r}\le K$, the Lemma shows that $\beta$ is adapted; it is continuous, and $e^{-KT}\le\beta_t\le e^{KT}$. The process $\beta g$ is therefore nonnegative, measurable and adapted, and by the Lemma $\int_0^t\beta_sg_s\,ds$ is $\F_t$-measurable for each $t$; in particular $Q$ is $\F_T$-measurable ($f_T$ being $\F_T$-measurable), so $\wt EQ$ and $\wt E[Q\mid\F_t]$ are defined with $\wt P$ as in \bk{(8.18)}. Since $Z_T>0$, $P$ and $\wt P$ have the same null events in $\F_T$, and since $\F_T$ contains every $P$-negligible event of $\F$, a property holds $P$-a.s.\ if and only if it holds $\wt P$-a.s.; we simply write ``a.s.'' By \bk{Lemma 8.11}, $x\defeq\wt EQ$ is finite. We shall prove the following statement, which contains both assertions of the exercise.

\begin{claim}
Let $(\pi,C)$ and $(\pi',C')$ be two hedging strategies against $(g,f_T)$, both with initial wealth $x=\wt EQ$, and let $X$, $X'$ be the corresponding wealth processes. Then
\begin{enumerate}
\item[\textup{(a)}] $P[X_t=X'_t;\ 0\le t\le T]=1$, and for every $t\in[0,T]$ we have, almost surely,
\begin{equation}\label{eq:5-8-13:val}
  X_t=\wt E\biggl[e^{-\int_t^T r(u)\,du}f_T+\int_t^T e^{-\int_t^s r(u)\,du}g_s\,ds\biggm|\F_t\biggr];
\end{equation}
\item[\textup{(b)}] $C=C'$ and $\pi=\pi'$ hold $\meas\times P$-a.e.\ on $[0,T]\times\Omega$.
\end{enumerate}
\end{claim}

Taking for $(\pi',C')$ the hedging strategy constructed in the proof of \bk{Theorem 8.12}, we see from (b) that every hedging strategy with initial wealth $\wt EQ$ agrees with it $\meas\times P$-a.e., and from (a) that its wealth process is indistinguishable from the continuous process $\xi$ of \bk{(8.31)} and is given by \bk{(8.32)}, which is \eqref{eq:5-8-13:val}. Because $\wt EQ$ is the fair price (\bk{Theorem 8.12}), this is what was to be shown.

\begin{proof}[Proof of the Claim]
\emph{Step 1: the discounted gains process is a $\wt P$-martingale.} Let $(\pi,C)$ be a hedging strategy with initial wealth $x$ and wealth process $X$, and set
\[
  M_t\defeq x+\int_0^t\beta_s\,\pi^T(s)\sigma(s)\,d\wt W_s;\qquad 0\le t\le T .
\]
By \bk{(8.20)},
\begin{equation}\label{eq:5-8-13:M}
  \beta_tX_t+\int_0^t\beta_sC_s\,ds=M_t;\qquad 0\le t\le T,\quad\text{a.s.},
\end{equation}
where the exceptional null event may be taken independent of $t$, because both sides are continuous in $t$ (the wealth process \bk{(8.13)} is continuous, and $\int_0^T\beta_sC_s\,ds<\infty$ a.s.\ by \bk{(8.12)}). As observed in the book after \bk{(8.20)}, $M$ is a continuous local martingale under $\wt P$ (the integrand $\beta\sigma^T\pi$ is measurable, adapted and square-integrable on $[0,T]$, a.s., by \bk{(8.11)} and the boundedness of $\beta$ and $\sigma$; thus its components belong to the class $\cP$ of \bk{Definition 3.2.22} relative to $\wt W^{(j)}$, and integrands from $\cP$ are allowed because $\qv{\wt W^{(j)}}_t=t$ is absolutely continuous, see the paragraph following \bk{Definition 3.2.22} and \bk{Remark 3.2.11}). Moreover, by \eqref{eq:5-8-13:M}, almost surely $M_t\ge0$ for all $t\in[0,T]$, since $X\ge0$ by admissibility (\bk{Definition 8.5}, \bk{(8.14)}) and $C\ge0$ because $C$ is a consumption process (\bk{Definition 8.2}). By \bk{Problem 1.5.19(ii)}, $M$ is a supermartingale under $\wt P$; in particular $\wt EM_t\le M_0=x<\infty$ for $0\le t\le T$.

By \bk{Definition 8.9} there is an event of probability one on which $C_t=g_t$ for all $t\in[0,T]$ and $X_T=f_T$. On the intersection of this event with the event where \eqref{eq:5-8-13:M} holds,
\begin{equation}\label{eq:5-8-13:MT}
  M_T=\beta_Tf_T+\int_0^T\beta_sg_s\,ds=Q,
  \qquad\text{whence}\qquad \wt EM_T=\wt EQ=x=M_0 .
\end{equation}
Fix $t\in[0,T]$. The random variable $M_t-\wt E[M_T\mid\F_t]$ is nonnegative a.s.\ by the supermartingale property, and its $\wt P$-expectation is $\wt EM_t-\wt EM_T=\wt EM_t-x\le0$. Therefore it vanishes a.s., and with \eqref{eq:5-8-13:MT} we obtain
\begin{equation}\label{eq:5-8-13:mart}
  M_t=\wt E[Q\mid\F_t]\quad\text{a.s.};\qquad 0\le t\le T .
\end{equation}

\emph{Step 2: proof of} (a). Let $M'$ be defined as $M$, with $\pi'$ in place of $\pi$; then \eqref{eq:5-8-13:M} holds with $(C',X',M')$ in place of $(C,X,M)$. By \eqref{eq:5-8-13:mart}, applied to both strategies, $M_t=M'_t$ a.s.\ for every $t\in[0,T]$; hence, almost surely, $M_t=M'_t$ for all rational $t\in[0,T]$, and then for all $t\in[0,T]$, since both processes are continuous. On the intersection of this event with the events of probability one on which $C_t=g_t=C'_t$ for all $t\in[0,T]$ and on which \eqref{eq:5-8-13:M} holds for both strategies, we have
\[
  X_t=\frac1{\beta_t}\Bigl(M_t-\int_0^t\beta_sg_s\,ds\Bigr)
     =\frac1{\beta_t}\Bigl(M'_t-\int_0^t\beta_sg_s\,ds\Bigr)=X'_t;\qquad 0\le t\le T,
\]
which is the first assertion of (a). For the second, fix $t\in[0,T]$ and put $Y\defeq\int_0^t\beta_sg_s\,ds$. As shown after the Lemma, $Y$ is $\F_t$-measurable, and $0\le Y\le Q$, so $Y$ is $\wt P$-integrable. Taking out what is known and using \eqref{eq:5-8-13:mart} and \eqref{eq:5-8-13:M}, we obtain, almost surely,
\[
  \wt E[\,Q-Y\mid\F_t]=M_t-Y=\beta_tX_t ,
\]
the last equality on the event, of probability one, where \eqref{eq:5-8-13:M} holds and $C=g$ on $[0,t]$. Since $Q-Y=\beta_Tf_T+\int_t^T\beta_sg_s\,ds$ a.s., we conclude that, almost surely,
\[
  X_t=\frac1{\beta_t}\,\wt E\Bigl[\,\beta_Tf_T+\int_t^T\beta_sg_s\,ds\Bigm|\F_t\Bigr],
\]
and \eqref{eq:5-8-13:val} follows by taking the bounded, $\F_t$-measurable factor $1/\beta_t$ inside the conditional expectation, since $\beta_s/\beta_t=e^{-\int_t^s r(u)\,du}$ for $t\le s\le T$.

\emph{Step 3: proof of} (b). The set $\{(t,\omega)\in[0,T]\times\Omega;\ C_t(\omega)\ne C'_t(\omega)\}$ is product-measurable, and for almost every $\omega$ its $\omega$-section is empty, because $C_t=g_t=C'_t$ for all $t\in[0,T]$, a.s. By Fubini's theorem it is a $\meas\times P$-null set.

Now let $\psi(s)\defeq\beta_s\,\sigma^T(s)\bigl(\pi(s)-\pi'(s)\bigr)$, an $\R^d$-valued, measurable, adapted process with $\int_0^T\norm{\psi(s)}^2ds<\infty$ a.s. By the definitions of $M$ and $M'$, the linearity of the stochastic integral, and Step 2,
\[
  N_t\defeq\sum_{j=1}^d\int_0^t\psi_j(s)\,d\wt W^{(j)}_s=M_t-M'_t=0;\qquad 0\le t\le T,\quad\text{a.s.}
\]
The process $N$ is a continuous local martingale under $\wt P$, relative to $\{\F_t\}$, whose quadratic variation is
\[
  \qv{N}_t=\sum_{j=1}^d\sum_{k=1}^d\int_0^t\psi_j(s)\psi_k(s)\,d\qv{\wt W^{(j)},\wt W^{(k)}}_s
          =\int_0^t\norm{\psi(s)}^2ds;\qquad 0\le t\le T,
\]
by \bk{(3.2.19)} and \bk{(3.2.26)} (sample path properties which remain valid for integrals with respect to local martingales, \bk{Section 3.2.D}; the components $\psi_j$ belong to the class $\cP$ of \bk{Definition 3.2.22} relative to $\wt W^{(j)}$, which suffices because $\qv{\wt W^{(j)}}_t=t$ is absolutely continuous, see the paragraph following \bk{Definition 3.2.22} and \bk{Remark 3.2.11}) and the fact that $\wt W$ is a $d$-dimensional Brownian motion under $\wt P$ (\bk{(8.19)}). On the other hand, $N$ is indistinguishable from the zero process, whose quadratic variation is zero; by the uniqueness assertion of \bk{Problem 1.5.17}, $\qv{N}_T=0$ a.s. Consequently there is an event $\Omega_0$ of probability one such that, for $\omega\in\Omega_0$,
\[
  \int_0^T\beta_s^2(\omega)\,\norm[\big]{\sigma^T(s,\omega)\bigl(\pi(s,\omega)-\pi'(s,\omega)\bigr)}^2ds=0 ,
\]
and such that $\sigma^T(s,\omega)$ is invertible for every $s\in[0,T]$ (\bk{Problem 8.1}). For $\omega\in\Omega_0$, since $\beta_s(\omega)>0$, we have $\sigma^T(s,\omega)(\pi(s,\omega)-\pi'(s,\omega))=0$, and therefore $\pi(s,\omega)=\pi'(s,\omega)$, for Lebesgue-almost every $s\in[0,T]$. The set $\{(s,\omega);\ \pi(s,\omega)\ne\pi'(s,\omega)\}$ is product-measurable and its $\omega$-sections have Lebesgue measure zero for $\omega\in\Omega_0$; by Fubini's theorem, it is a $\meas\times P$-null set.
\end{proof}

\rmk{The uniqueness is tied to the initial wealth $x=\wt EQ$: for $x>\wt EQ$ we have $\wt EM_T=\wt EQ<x=M_0$, so $M$ is not a $\wt P$-martingale and Step~1 breaks down. Conversely, nothing better than $\meas\times P$-a.e.\ uniqueness of $\pi$ can hold, since altering $\pi$ on a $\meas\times P$-null set (keeping it measurable and adapted) changes neither \bk{(8.11)} nor, up to indistinguishability, the wealth process \bk{(8.13)}. The right-hand side of \eqref{eq:5-8-13:val} is defined, for each $t$, only up to a null set; the valuation process is its continuous modification.}
\end{kex}

\begin{kex}{Exercise}{8.15}{}
\emph{Conventions and hypotheses.} We read ``piecewise $C^2$'' as follows: there is a set $D\subset(0,\infty)$ with finitely many points in every bounded interval, such that on the closure of each of the intervals into which $D$ divides $[0,\infty)$ the function $h$ coincides with a function of class $C^2$. Thus $h''$ is defined and continuous on $(0,\infty)\setminus D$, the one-sided limits $h'(a\pm)$ exist at every $a\in D$ and are the one-sided derivatives of $h$ at $a$, and $h'(0)=h'(0+)$ is the right-hand derivative at the origin. Since $h$ is convex, $h'$ is nondecreasing on $[0,\infty)\setminus D$, so that
\[
  h''\ge 0 \ \text{ on } (0,\infty)\setminus D,
  \qquad \theta_a\defeq h'(a+)-h'(a-)\ge 0\quad\text{for } a\in D .
\]
The statement speaks of ``the contingent claim $(0,f_T)$''; according to \bk{Definition 8.7} this includes the integrability condition \bk{(8.29)}, i.e.\ $E[h(P_1(T))^\mu]<\infty$ for some $\mu>1$, which we assume (see the second remark at the end). The initial stock price $p_1$ is nonnegative. We shall prove that, for every $t\in[0,T]$,
\begin{equation}\label{eq:5-8-15:main}
  X_t=\wt{E}\bigl[e^{-r(T-t)}h(P_1(T))\bigm|\F_t\bigr]
  =\int_0^\infty h''(q)\,v_{q,T}(t,P_1(t))\,dq+\sum_{a\in D}\theta_a\,v_{a,T}(t,P_1(t))
  \quad\text{a.s.}
\end{equation}
When $h$ is continuously differentiable, all the $\theta_a$ vanish and \eqref{eq:5-8-15:main} is \bk{(8.37)}; the first remark at the end shows that \bk{(8.37)} is false without this proviso.

\emph{Step 1: the first equality.} The random variable $f_T=h(P_1(T))$ is nonnegative and $\F_T$-measurable, because $h$ is continuous on each of the countably many pieces, hence Borel measurable; with \bk{(8.29)}, the pair $(0,f_T)$ is a contingent claim. \bk{Example 8.14} is a special case of the model of \bk{Subsection 8.A} with $d=1$, and \bk{(8.3)} holds with $\eps=\sigma^2$. By \bk{Theorem 8.12} and \bk{Exercise 8.13}, the valuation process is the process $X$ of \bk{(8.32)}, which for $g\equiv 0$ and $r(\cdot)\equiv r$ reads $X_t=\wt{E}[e^{-r(T-t)}f_T\mid\F_t]$; the conditional expectation is well defined and a.s.\ finite because $\wt{E}f_T<\infty$ by \bk{Lemma 8.11}. This is the first equality in \eqref{eq:5-8-15:main}.

\emph{Step 2: decomposition of $h$ into option payoffs.} We claim that
\begin{equation}\label{eq:5-8-15:dec}
  h(x)=\int_0^\infty h''(q)\,(x-q)^+\,dq+\sum_{a\in D}\theta_a\,(x-a)^+;\qquad 0\le x<\infty .
\end{equation}
Let $y\in(0,\infty)\setminus D$ and let $a_1<\dots<a_m$ be the points of $D$ in $(0,y)$; put $a_0=0$, $a_{m+1}=y$. On each $[a_{k-1},a_k]$ the function $h$ agrees with a $C^2$ function, so the fundamental theorem of calculus gives $h'(a_k-)-h'(a_{k-1}+)=\int_{a_{k-1}}^{a_k}h''(q)\,dq$ for $1\le k\le m+1$ (with $h'(a_{m+1}-)=h'(y)$). Summing over $k$ and using $h'(0+)=0$, we obtain
\begin{equation}\label{eq:5-8-15:hprime}
  h'(y)=\int_0^y h''(q)\,dq+\sum_{a\in D,\ a<y}\theta_a;\qquad y\in(0,\infty)\setminus D .
\end{equation}
Next, $h$ is continuous on $[0,\infty)$ (it is continuous on each closed piece, and adjacent pieces share an endpoint) and continuously differentiable on each closed piece, so the same telescoping argument yields $h(x)=h(0)+\int_0^x h'(y)\,dy=\int_0^x h'(y)\,dy$. Inserting \eqref{eq:5-8-15:hprime} and interchanging the order of integration and summation, which is legitimate by Tonelli's theorem because $h''\ge 0$ and $\theta_a\ge 0$, we find
\[
  h(x)=\int_0^\infty h''(q)\int_0^x\ind_{\{q<y\}}\,dy\,dq
  +\sum_{a\in D}\theta_a\int_0^x\ind_{\{a<y\}}\,dy ,
\]
and $\int_0^x\ind_{\{q<y\}}\,dy=(x-q)^+$. This proves \eqref{eq:5-8-15:dec}.

\emph{Step 3: the stock price under $\wt{P}$.} By \bk{(8.19)} and \bk{Theorem 3.5.1}, $\{\wt{W}_t,\F_t;\ 0\le t\le T\}$ is a Brownian motion on $(\Omega,\F_T,\wt{P})$, and by \bk{Example 8.14}, $dP_1(t)=rP_1(t)\,dt+\sigma P_1(t)\,d\wt{W}_t$, $P_1(0)=p_1$. Let $Y_t\defeq\exp\{-\sigma\wt{W}_t-(r-\tfrac12\sigma^2)t\}$. It\^o's rule gives $dY_t=Y_t\,[(\sigma^2-r)\,dt-\sigma\,d\wt{W}_t]$ and $d\qv{P_1,Y}_t=-\sigma^2P_1(t)Y_t\,dt$, whence, by the integration-by-parts formula,
\[
  d\bigl(P_1(t)Y_t\bigr)=P_1(t)Y_t\bigl[(\sigma^2-r)\,dt-\sigma\,d\wt{W}_t
  +r\,dt+\sigma\,d\wt{W}_t-\sigma^2\,dt\bigr]=0 .
\]
Hence $P_1(t)Y_t=p_1$, i.e.\ $P_1(t)=p_1\exp\{\sigma\wt{W}_t+(r-\tfrac12\sigma^2)t\}$ for $0\le t\le T$, a.s. Fix $t\in[0,T)$, write $\tau\defeq T-t>0$ and
\[
  \xi\defeq\exp\bigl\{\sigma(\wt{W}_T-\wt{W}_t)+(r-\tfrac12\sigma^2)\tau\bigr\},
  \qquad\text{so that}\qquad P_1(T)=P_1(t)\,\xi\quad\text{a.s.}
\]
Under $\wt{P}$, the random variable $\xi$ is independent of $\F_t$ and has the law of $\exp\{\sigma\sqrt{\tau}\,Z+(r-\tfrac12\sigma^2)\tau\}$ with $Z$ standard normal, while $P_1(t)$ is $\F_t$-measurable with values in $[0,\infty)$. Therefore, for every Borel function $\varphi\colon[0,\infty)\to[0,\infty)$, the function $x\mapsto\wt{E}\varphi(x\xi)\in[0,\infty]$ is Borel measurable (Tonelli) and
\begin{equation}\label{eq:5-8-15:markov}
  \wt{E}\bigl[e^{-r\tau}\varphi(P_1(T))\bigm|\F_t\bigr]=U_\varphi(P_1(t))\quad\text{a.s.},
  \qquad U_\varphi(x)\defeq e^{-r\tau}\,\wt{E}\,\varphi(x\xi);\quad x\ge 0 .
\end{equation}
(This is the standard rule for the conditional expectation of a nonnegative function of an $\F_t$-measurable variable and a variable independent of $\F_t$.)

\emph{Step 4: the option payoff.} Let $q>0$ and $\varphi_q(x)\defeq(x-q)^+$. We verify that $U_{\varphi_q}(x)=v_{q,T}(t,x)$ for all $x\ge 0$, where $v_{q,T}$ is the function of \bk{(8.34)}. For $x=0$ both sides vanish ($\rho_\pm(\tau,0)=-\infty$ and $\Phi(-\infty)=0$). Let $x>0$. With $\rho_\pm=\rho_\pm(\tau,x)$ as in \bk{(8.34)} we have $\rho_+-\rho_-=\sigma\sqrt{\tau}$ and
\[
  x\exp\bigl\{\sigma\sqrt{\tau}\,z+(r-\tfrac12\sigma^2)\tau\bigr\}>q
  \iff z>-\rho_- .
\]
Consequently, completing the square in the first integral,
\begin{align*}
  e^{-r\tau}\,\wt{E}(x\xi-q)^+
  &=x\int_{-\rho_-}^\infty e^{\sigma\sqrt{\tau}z-\sigma^2\tau/2}\,\frac{e^{-z^2/2}}{\sqrt{2\pi}}\,dz
    -qe^{-r\tau}\int_{-\rho_-}^\infty\frac{e^{-z^2/2}}{\sqrt{2\pi}}\,dz\\
  &=x\int_{-\rho_--\sigma\sqrt{\tau}}^\infty\frac{e^{-w^2/2}}{\sqrt{2\pi}}\,dw-qe^{-r\tau}\Phi(\rho_-)
   =x\,\Phi(\rho_+)-qe^{-r\tau}\Phi(\rho_-)=v_{q,T}(t,x).
\end{align*}
In particular $(q,x)\mapsto v_{q,T}(t,x)$ is nonnegative, and continuous on $(0,\infty)\times(0,\infty)$.

\emph{Step 5: conclusion.} Let $t\in[0,T)$. We apply \eqref{eq:5-8-15:dec} at the point $x\xi$, multiply by $e^{-r\tau}$ and take $\wt{E}$-expectations; all terms being nonnegative and jointly measurable in $(q,\omega)$, Tonelli's theorem and Step 4 give, for every $x\ge 0$,
\[
  U_h(x)=\int_0^\infty h''(q)\,U_{\varphi_q}(x)\,dq+\sum_{a\in D}\theta_a\,U_{\varphi_a}(x)
  =\int_0^\infty h''(q)\,v_{q,T}(t,x)\,dq+\sum_{a\in D}\theta_a\,v_{a,T}(t,x)
\]
as an identity in $[0,\infty]$. Substituting $x=P_1(t)$ and using \eqref{eq:5-8-15:markov} with $\varphi=h$, we obtain the second equality in \eqref{eq:5-8-15:main}; both sides are a.s.\ finite by Step 1. For $t=T$ we have $v_{q,T}(T,x)=(x-q)^+$ by \bk{(8.34)}, and \eqref{eq:5-8-15:main} is \eqref{eq:5-8-15:dec} evaluated at $x=P_1(T)$, since $X_T=f_T$ a.s. This proves \eqref{eq:5-8-15:main}, and hence \bk{(8.37)} when $h\in C^1([0,\infty))$.

\rmk{\emph{On the hypothesis ``piecewise $C^2$''.} If $h'$ is allowed to jump, \bk{(8.37)} is false as printed. The function $h(x)=(x-1)^+$ is convex, piecewise $C^2$ (indeed piecewise linear) and satisfies $h(0)=h'(0)=0$; here $h''=0$ off $D=\{1\}$, so the right-hand side of \bk{(8.37)} is zero, whereas by \bk{(8.33)}, \bk{(8.36)} the valuation process is $v_{1,T}(t,P_1(t))$, which is strictly positive for $t<T$ when $p_1>0$ (by Step 4, $v_{1,T}(t,x)=e^{-r\tau}\wt{E}(x\xi-1)^+>0$ for $x>0$, because $\wt{P}[x\xi>1]>0$). The same positivity shows that for $p_1>0$ and $t<T$ the two right-hand sides of \bk{(8.37)} and \eqref{eq:5-8-15:main} agree only if every $\theta_a$ is zero. Thus \bk{(8.37)} requires $h$ to be continuously differentiable with a piecewise continuous second derivative, which is presumably what is meant; in general the jumps of $h'$ contribute the sum in \eqref{eq:5-8-15:main}, i.e.\ $h''(q)\,dq$ has to be replaced by the second-derivative measure of $h$.}

\rmk{\emph{On the integrability condition.} Under $P$ we have $P_1(T)=p_1\exp\{\int_0^Tb_1(s)\,ds-\tfrac12\sigma^2T+\sigma W_T\}$ with $b_1$ bounded, so $P_1(T)$ has moments of all orders, and \bk{(8.29)} holds whenever $h(x)\le K(1+x^m)$ for some constants $K$, $m$. It is not implied by the hypotheses on $h$: for $h(x)=e^{x^2}-1$ and $p_1>0$ we have $P_1(T)\ge ce^{\sigma W_T}$ for some constant $c>0$ and $E\exp\{c^2e^{2\sigma W_T}\}=\infty$, so that $f_T$ is not integrable and $(0,f_T)$ is not a contingent claim in the sense of \bk{Definition 8.7}. Steps 2--5 do not use \bk{(8.29)}: the second equality in \eqref{eq:5-8-15:main} always holds as an identity between $[0,\infty]$-valued random variables.}
\end{kex}

\begin{kex}{Exercise}{8.20}{}
We work in the setting of the book's Subsection 5.8.D: constant coefficients \bk{(8.47)}, $\theta=\sigma^{-1}(b-r\mathbf{1})\neq 0$, and $G$, $S$ the functions of the standing assumption \bk{(8.48)}--\bk{(8.52)}. Both cases are treated at once by means of the function
\begin{equation}\label{eq:5-8-20:f}
  f(t)\defeq\int_t^T e^{-k(s-t)}\,ds=
  \begin{cases}
    \dfrac{1-e^{-k(T-t)}}{k}, & k\neq 0,\\[1.2ex]
    T-t, & k=0,
  \end{cases}
  \qquad 0\le t\le T,
\end{equation}
and the exponent
\[
  p\defeq\frac{\delta}{\delta-1}<0 ,\qquad\text{for which}\qquad p-1=\frac{1}{\delta-1},\qquad p(p-1)=\frac{\delta}{(1-\delta)^2} .
\]
The function $f$ is continuously differentiable on $[0,T]$, with
\begin{equation}\label{eq:5-8-20:ode}
  -f'(t)+kf(t)=1;\quad 0\le t\le T,\qquad f(T)=0 ,
\end{equation}
(differentiate either expression in \eqref{eq:5-8-20:f}), and the integral in \eqref{eq:5-8-20:f} shows that $0<f(t)\le f_*\defeq T\max(1,e^{-kT})$ for $0\le t<T$. We have to show that, for $0\le t\le T$, $y>0$ and for $0\le t\le T$, $x\ge0$ (the range on which $V$ is defined by \bk{(8.57)}),
\begin{equation}\label{eq:5-8-20:goal}
  G(t,y)=f(t)\Bigl(\frac{y}{\delta}\Bigr)^{p},\qquad S(t,y)=\delta\,G(t,y),\qquad V(t,x)=e^{-\beta t}f(t)^{1-\delta}x^{\delta} .
\end{equation}

\emph{Step 1: the data.} The function $U(c)=c^\delta$ maps $[0,\infty)$ into $[0,\infty)$, is continuous, strictly increasing and strictly concave, with $U(0)=0$; it is three times continuously differentiable on $(0,\infty)$, with $U'(c)=\delta c^{\delta-1}$, so that $U'(0)=\infty$ and $U'(\infty)=0$. It is thus a utility function in the sense of Subsection 5.8.C. Solving $\delta c^{\delta-1}=y$ gives the inverse of $U'$:
\[
  I(y)=\Bigl(\frac{y}{\delta}\Bigr)^{1/(\delta-1)}=\Bigl(\frac{y}{\delta}\Bigr)^{p-1};\qquad 0<y<\infty ,
\]
and therefore, for $y>0$,
\begin{equation}\label{eq:5-8-20:rhs}
  U(I(y))=\Bigl(\frac{y}{\delta}\Bigr)^{\delta/(\delta-1)}=\Bigl(\frac{y}{\delta}\Bigr)^{p},\qquad
  yI(y)=\delta\,\frac{y}{\delta}\Bigl(\frac{y}{\delta}\Bigr)^{p-1}=\delta\Bigl(\frac{y}{\delta}\Bigr)^{p}=\delta\,U(I(y)) .
\end{equation}

\emph{Step 2: the candidates solve the Cauchy problems.} Define $\Gamma(t,y)\defeq f(t)(y/\delta)^p$ and $\Sigma(t,y)\defeq\delta\,\Gamma(t,y)$ on $[0,T]\times(0,\infty)$. Both are nonnegative, and since $f\in C^1([0,T])$ and $y\mapsto y^p$ is infinitely differentiable on $(0,\infty)$, the derivatives $\Gamma_t,\Gamma_{ty},\Gamma_y,\Gamma_{yy},\Gamma_{yyy}$ exist and are jointly continuous on $[0,T]\times(0,\infty)$; thus $\Gamma$ and $\Sigma$ are of class $C^{1,3}$ in the sense explained after \bk{(8.51)}. With $\varphi(y)\defeq(y/\delta)^p$ we have $y\varphi'(y)=p\,\varphi(y)$ and $y^2\varphi''(y)=p(p-1)\varphi(y)$, so the operator $L$ of the book gives
\[
  L\Gamma(t,y)=\Bigl[-f'(t)+\Bigl(\beta-(\beta-r)p-\tfrac12\norm{\theta}^2p(p-1)\Bigr)f(t)\Bigr]\Bigl(\frac{y}{\delta}\Bigr)^{p} .
\]
Now
\[
  \beta-(\beta-r)p=\beta+\frac{(\beta-r)\delta}{1-\delta}=\frac{\beta-r\delta}{1-\delta},
  \qquad
  \tfrac12\norm{\theta}^2p(p-1)=\frac{1}{1-\delta}\cdot\frac12\norm{\theta}^2\frac{\delta}{1-\delta},
\]
so the coefficient of $f(t)$ is exactly $k$, and by \eqref{eq:5-8-20:ode} and \eqref{eq:5-8-20:rhs},
\[
  L\Gamma(t,y)=\bigl(-f'(t)+kf(t)\bigr)\Bigl(\frac{y}{\delta}\Bigr)^{p}=U(I(y)),\qquad \Gamma(T,y)=0;\qquad 0\le t\le T,\ y>0 .
\]
This is \bk{(8.48)}, \bk{(8.49)}. Since $L$ is linear, $L\Sigma=\delta\,L\Gamma=\delta\,U(I(y))=yI(y)$ by \eqref{eq:5-8-20:rhs}, and $\Sigma(T,y)=0$: this is \bk{(8.50)}, \bk{(8.51)}.

For the growth condition \bk{(8.52)}, set $\lambda\defeq 1-p=1/(1-\delta)>0$ and observe that for every real $q$ with $\abs{q}\le\lambda$ and every $y>0$ we have $y^q\le 1+y^{-\lambda}+y^{\lambda}$ (if $y\ge1$, then $y^q\le y^{\abs q}\le y^\lambda$; if $y<1$, then $y^q\le y^{-\abs q}\le y^{-\lambda}$). Because $\abs{p}<\lambda$ and $\abs{p-1}=\lambda$, we obtain for $y>0$
\begin{align*}
  \max_{0\le t\le T}\Gamma(t,y)&\le f_*\delta^{-p}\bigl(1+y^{-\lambda}+y^{\lambda}\bigr),\\
  \max_{0\le t\le T}\abs{\Gamma_y(t,y)}&=\max_{0\le t\le T}f(t)\,\frac{\abs p}{\delta}\Bigl(\frac{y}{\delta}\Bigr)^{p-1}\le f_*\abs p\,\delta^{-p}\bigl(1+y^{-\lambda}+y^{\lambda}\bigr),
\end{align*}
and the same bounds multiplied by $\delta$ for $\Sigma$ and $\abs{\Sigma_y}$. Hence $\Gamma,\Gamma_y,\Sigma,\Sigma_y$ satisfy \bk{(8.52)} with this $\lambda$ and $M=f_*(1+\abs p)\delta^{-p}$. (The derivatives $\Gamma_y$, $\Sigma_y$ are negative, so for them \bk{(8.52)} is of course understood for the absolute value.) Consequently, for $U(c)=c^\delta$ the standing assumption of Subsection 5.8.D is fulfilled, by the pair $(\Gamma,\Sigma)$.

\emph{Step 3: identification of $G$ and $S$.} Let $(G,S)$ be any pair of functions as in the standing assumption. \bk{Problem 8.19}, applied to $H=G$ with $g(t,y)=U(I(y))$ and to $H=S$ with $g(t,y)=yI(y)$ (both $g$ are continuous and nonnegative on $[0,T]\times(0,\infty)$, and $G$, $S$ are nonnegative, of class $C^{1,2}$ and satisfy \bk{(8.52)}), yields the representations \bk{(8.55)}, \bk{(8.56)}, whose right-hand sides do not involve $G$ or $S$. The same problem applies to $\Gamma$ and $\Sigma$, by Step 2. Therefore, for $0\le t\le T$ and $y>0$,
\[
  G(t,y)=E\int_t^T e^{-\beta(s-t)}U\bigl(I(Y^{(t,y)}_s)\bigr)\,ds=\Gamma(t,y),
\]
and likewise $S(t,y)=\Sigma(t,y)=\delta\,\Gamma(t,y)=\delta\,G(t,y)$. In view of \eqref{eq:5-8-20:f}, these are the first two formulas of \eqref{eq:5-8-20:goal}, i.e.\ the asserted formulas for $G$ and $S$ in both cases $k\ne0$ and $k=0$.

\emph{Step 4: the value function.} Fix $0\le t<T$, so that $f(t)>0$. By \bk{(8.61)} and Step 3,
\[
  \mathscr{X}(t,y)=\frac{S(t,y)}{y}=\frac{\delta f(t)}{y}\Bigl(\frac{y}{\delta}\Bigr)^{p}=f(t)\Bigl(\frac{y}{\delta}\Bigr)^{p-1}=f(t)\Bigl(\frac{y}{\delta}\Bigr)^{1/(\delta-1)};\qquad 0<y<\infty .
\]
As $1/(\delta-1)<0$, this is a continuous, strictly decreasing map of $(0,\infty)$ onto $(0,\infty)$, in agreement with the book's assertion that $\mathscr{X}(t,\cdot)$ is strictly decreasing on $[0,\infty]$ with $\mathscr{X}(t,0)=\infty$, $\mathscr{X}(t,\infty)=0$. For $0<x<\infty$ the equation $f(t)(y/\delta)^{1/(\delta-1)}=x$ has the unique solution $y=\delta\,(x/f(t))^{\delta-1}\in(0,\infty)$, so the inverse function of \bk{(8.62)} is
\[
  \mathscr{Y}(t,x)=\delta\Bigl(\frac{x}{f(t)}\Bigr)^{\delta-1};\qquad 0<x<\infty .
\]
Substituting into \bk{(8.66)} and using Step 3 together with $(\delta-1)p=\delta$, we obtain for $0\le t<T$, $x>0$:
\[
  V(t,x)=e^{-\beta t}G\bigl(t,\mathscr{Y}(t,x)\bigr)
  =e^{-\beta t}f(t)\Bigl(\frac{\mathscr{Y}(t,x)}{\delta}\Bigr)^{p}
  =e^{-\beta t}f(t)\Bigl(\frac{x}{f(t)}\Bigr)^{(\delta-1)p}
  =e^{-\beta t}f(t)^{1-\delta}x^{\delta} .
\]
With \eqref{eq:5-8-20:f} this is
\[
  V(t,x)=e^{-\beta t}\Bigl(\frac{1-e^{-k(T-t)}}{k}\Bigr)^{1-\delta}x^\delta\quad(k\neq0),
  \qquad
  V(t,x)=e^{-\beta t}(T-t)^{1-\delta}x^\delta\quad(k=0),
\]
which is the third formula of \eqref{eq:5-8-20:goal} on the range $0\le t<T$, $x>0$ of \bk{(8.66)}.

\emph{Step 5: the boundary cases $t=T$ and $x=0$.} For every $0\le t\le T$ and $x\ge0$ the pair $(\pi,C)\equiv(0,0)$ is admissible for $(t,x)$, since \bk{(8.58)} then gives $X_s=x+\int_t^s rX_u\,du$, i.e.\ $X_s=xe^{r(s-t)}\ge0$; so the supremum in \bk{(8.57)} is over a nonempty set. If $t=T$, the integral in \bk{(8.57)} is over the degenerate interval $[T,T]$ and vanishes for every admissible pair, so $V(T,x)=0$; and the right-hand side of \eqref{eq:5-8-20:goal} is $e^{-\beta T}f(T)^{1-\delta}x^\delta=0$, because $f(T)=0$ and $1-\delta>0$. Now let $0\le t<T$ and $x=0$, and let $(\pi,C)$ be admissible for $(t,0)$. By \bk{(8.59)},
\[
  E\int_t^T e^{-r(s-t)}Z^t_sC_s\,ds\le0 .
\]
The integrand is nonnegative, since $C_s\ge0$, and $e^{-r(s-t)}Z^t_s>0$ everywhere; a nonnegative product-measurable function with integral at most $0$ vanishes almost everywhere, so $C_s(\omega)=0$ for $(\mathrm{Leb}\otimes P)$-a.e.\ $(s,\omega)\in[t,T]\times\Omega$. As $U(0)=0$, the function $(s,\omega)\mapsto e^{-\beta s}U(C_s(\omega))$ also vanishes a.e., and by Fubini--Tonelli $E\int_t^Te^{-\beta s}U(C_s)\,ds=0$. Hence $V(t,0)=0=e^{-\beta t}f(t)^{1-\delta}\cdot0^\delta$, since $0^\delta=0$. Together with Step 4, this completes the proof of \eqref{eq:5-8-20:goal}.

\rmk{The formulas for $G$ and $S$ can also be read off directly from \bk{(8.55)}, \bk{(8.56)}, which confirms Step 3. Let $Z^t_s$ be as in \bk{(8.54)}. Since $W$ is a $d$-dimensional Brownian motion, the coordinates of $W_s-W_t$ are independent and $N(0,s-t)$-distributed, so $\theta^{T}(W_s-W_t)$ is normal with mean zero and variance $\norm\theta^2(s-t)$; hence $E(Z^t_s)^p=\exp\{\tfrac12\norm\theta^2(p^2-p)(s-t)\}$, and by \bk{(8.53)} and \eqref{eq:5-8-20:rhs},
\[
  e^{-\beta(s-t)}E\,U\bigl(I(Y^{(t,y)}_s)\bigr)
  =\Bigl(\frac y\delta\Bigr)^{p}\exp\Bigl\{-\bigl[\beta-(\beta-r)p-\tfrac12\norm\theta^2p(p-1)\bigr](s-t)\Bigr\}
  =\Bigl(\frac y\delta\Bigr)^{p}e^{-k(s-t)} .
\]
Integrating over $t\le s\le T$ (Fubini--Tonelli, the integrand being nonnegative) gives $G(t,y)=f(t)(y/\delta)^p$. Similarly $Z^t_sI(Y^{(t,y)}_s)=(y/\delta)^{p-1}e^{(p-1)(\beta-r)(s-t)}(Z^t_s)^{p}$, whence $e^{-r(s-t)}E[Z^t_sI(Y^{(t,y)}_s)]=(y/\delta)^{p-1}e^{-k(s-t)}$ and $S(t,y)=y\,f(t)(y/\delta)^{p-1}=\delta f(t)(y/\delta)^p$. The case $k=0$ is also the limit $k\to0$ of the case $k\ne0$, since $(1-e^{-k(T-t)})/k\to T-t$.}
\end{kex}

\begin{kex}{Exercise}{8.23}{}
We are in the setting of the book's Subsection 5.8.D: constant coefficients \bk{(8.47)}, $\theta=\sigma^{-1}(b-r\mathbf{1})\ne0$, $U(c)=c^\delta$ with $0<\delta<1$, and $\{\F_s\}$ the augmentation of the filtration generated by the $d$-dimensional Brownian motion $W$. We fix an initial condition $(t,x)\in[0,T)\times(0,\infty)$; see the remark at the end for $x=0$. We write
\[
  \gamma\defeq\frac{1}{1-\delta}>1,\qquad
  f(s)\defeq\int_s^Te^{-k(u-s)}\,du=
  \begin{cases}\bigl(1-e^{-k(T-s)}\bigr)/k,&k\ne0,\\ T-s,&k=0,\end{cases}
  \qquad 0\le s\le T,
\]
with $k=\gamma\bigl(\beta-r\delta\bigr)-\tfrac12\norm\theta^2\delta\gamma^2$ the constant of \bk{Exercise 8.20}. The function $f$ is continuously differentiable on $[0,T]$, positive on $[0,T)$, and
\begin{equation}\label{eq:5-8-23:f}
  f'(s)=kf(s)-1;\quad 0\le s\le T,\qquad f(T)=0 .
\end{equation}
We have $I(y)=(y/\delta)^{-\gamma}$ for $0<y<\infty$, and \bk{Exercise 8.20} together with \bk{(8.61)}, \bk{(8.62)} gives, for $0\le s<T$,
\begin{equation}\label{eq:5-8-23:XY}
  \mathscr{X}(s,y)=\frac{S(s,y)}{y}=f(s)\,I(y),\qquad
  \mathscr{Y}(s,\xi)=\delta\Bigl(\frac{\xi}{f(s)}\Bigr)^{\delta-1},\qquad
  V(s,\xi)=e^{-\beta s}f(s)^{1-\delta}\xi^\delta<\infty
\end{equation}
for $0<y<\infty$, $0<\xi<\infty$.

\emph{Step 1: the feedback law.} From \eqref{eq:5-8-23:XY} we read off, for $0\le s<T$ and $\xi>0$,
\[
  I\bigl(\mathscr{Y}(s,\xi)\bigr)=\frac{\xi}{f(s)},\qquad
  \mathscr{Y}_x(s,\xi)=-(1-\delta)\frac{\mathscr{Y}(s,\xi)}{\xi},\qquad
  \frac{\mathscr{Y}(s,\xi)}{\mathscr{Y}_x(s,\xi)}=-\gamma\,\xi ,
\]
so that the feedback formulas \bk{(8.81)}, \bk{(8.82)} become
\begin{equation}\label{eq:5-8-23:fb}
  C^*_s=\frac{X^*_s}{f(s)},\qquad
  \pi^*_s=\gamma\,(\sigma\sigma^T)^{-1}(b-r\mathbf{1})\,X^*_s ;\qquad t\le s<T .
\end{equation}
Both are linear in the wealth: one consumes at the rate $1/f(s)$ per unit of wealth, a deterministic rate which increases to infinity as $s\uparrow T$, and one keeps in the stocks the constant proportions $\gamma(\sigma\sigma^T)^{-1}(b-r\mathbf{1})$ of the wealth. Since
\begin{equation}\label{eq:5-8-23:alg}
  (\pi^*_s)^T\sigma=\gamma X^*_s(b-r\mathbf{1})^T(\sigma^T)^{-1}=\gamma X^*_s\theta^T,\qquad
  (b-r\mathbf{1})^T\pi^*_s=\gamma X^*_s\,\theta^T\theta=\gamma\norm\theta^2X^*_s ,
\end{equation}
substitution of \eqref{eq:5-8-23:fb} into the wealth equation \bk{(8.58)} gives the linear equation
\begin{equation}\label{eq:5-8-23:sde}
  dX_s=X_s\Bigl[\Bigl(r+\gamma\norm\theta^2-\frac{1}{f(s)}\Bigr)ds+\gamma\,\theta^T\,dW_s\Bigr];\quad t\le s<T,\qquad X_t=x .
\end{equation}
The identities \bk{(8.81)}, \bk{(8.82)} are derived in the book only for $\meas\times P$-almost every $(s,\omega)$, for an optimal pair whose uniqueness is not discussed, and \bk{(8.82)} presupposes $X^*_s>0$. We therefore do not rely on them: in Steps 2--4 we solve \eqref{eq:5-8-23:sde}, verify directly that the resulting pair is optimal, and show that it is the only optimal pair.

\emph{Step 2: solution of the wealth equation.} Define, for $t\le s\le T$,
\begin{equation}\label{eq:5-8-23:Lam}
  \Lambda_s\defeq\exp\Bigl\{\gamma\,\theta^T(W_s-W_t)+\gamma\bigl(\tfrac12\norm\theta^2-\beta+r\bigr)(s-t)\Bigr\}
  =e^{-\gamma(\beta-r)(s-t)}\bigl(Z^t_s\bigr)^{-\gamma},
\end{equation}
with $Z^t_s$ as in \bk{(8.54)}, and
\begin{equation}\label{eq:5-8-23:X}
  X^*_s\defeq x\,\frac{f(s)}{f(t)}\,\Lambda_s
  =x\,\frac{f(s)}{f(t)}\exp\Bigl\{\frac{\theta^T(W_s-W_t)}{1-\delta}+\frac{\tfrac12\norm\theta^2-\beta+r}{1-\delta}\,(s-t)\Bigr\};\qquad t\le s\le T .
\end{equation}
When $k\ne0$ the factor $f(s)/f(t)$ is $\bigl(1-e^{-k(T-s)}\bigr)/\bigl(1-e^{-k(T-t)}\bigr)$, and it is $(T-s)/(T-t)$ when $k=0$. The process $X^*$ is continuous and $\{\F_s\}$-adapted on $[t,T]$, and strictly positive on $[t,T)$. By It\^o's rule,
\[
  d\Lambda_s=\Lambda_s\Bigl[\gamma\,\theta^T\,dW_s+\Bigl(\gamma\bigl(\tfrac12\norm\theta^2-\beta+r\bigr)+\tfrac12\gamma^2\norm\theta^2\Bigr)ds\Bigr],
\]
and since $f$ is of class $C^1$ on $[0,T]$ and satisfies \eqref{eq:5-8-23:f}, the product rule gives on $[t,T]$
\[
  dX^*_s=-\frac{x}{f(t)}\Lambda_s\,ds
  +X^*_s\Bigl[\Bigl(k+\gamma\bigl(\tfrac12\norm\theta^2-\beta+r\bigr)+\tfrac12\gamma^2\norm\theta^2\Bigr)ds+\gamma\,\theta^T\,dW_s\Bigr].
\]
Using $\gamma(1-\delta)=1$ we compute
\[
  k+\gamma\bigl(\tfrac12\norm\theta^2-\beta+r\bigr)+\tfrac12\gamma^2\norm\theta^2
  =\gamma r(1-\delta)+\tfrac12\gamma^2\norm\theta^2(1-\delta)+\tfrac12\gamma\norm\theta^2
  =r+\gamma\norm\theta^2 .
\]
Consequently, with
\begin{equation}\label{eq:5-8-23:C}
  C^*_s\defeq\frac{x}{f(t)}\,\Lambda_s,\qquad \pi^*_s\defeq\gamma\,(\sigma\sigma^T)^{-1}(b-r\mathbf{1})\,X^*_s;\qquad t\le s\le T,
\end{equation}
we have
\begin{equation}\label{eq:5-8-23:dX}
  dX^*_s=\bigl(rX^*_s-C^*_s+\gamma\norm\theta^2X^*_s\bigr)ds+\gamma X^*_s\,\theta^T\,dW_s;\qquad t\le s\le T,\quad X^*_t=x .
\end{equation}
As $C^*_s=X^*_s/f(s)$ for $t\le s<T$, the process $X^*$ solves \eqref{eq:5-8-23:sde}, and the processes \eqref{eq:5-8-23:C} are related to $X^*$ by \eqref{eq:5-8-23:fb}.

The solution of \eqref{eq:5-8-23:sde} is unique. Indeed, let $X$ be a continuous, adapted process on $[t,T)$ satisfying \eqref{eq:5-8-23:sde}, and write $\mu(s)\defeq r+\gamma\norm\theta^2-1/f(s)$, a continuous function on $[t,T)$. Since $X^*>0$ on $[t,T)$, It\^o's rule gives $d(1/X^*_s)=(1/X^*_s)\bigl[(\gamma^2\norm\theta^2-\mu(s))\,ds-\gamma\,\theta^T\,dW_s\bigr]$ and then
\[
  d\Bigl(\frac{X_s}{X^*_s}\Bigr)=\frac{X_s}{X^*_s}\Bigl[\mu(s)\,ds+\gamma\,\theta^T\,dW_s+\bigl(\gamma^2\norm\theta^2-\mu(s)\bigr)ds-\gamma\,\theta^T\,dW_s-\gamma^2\norm\theta^2\,ds\Bigr]=0 ,
\]
so $X_s/X^*_s=X_t/X^*_t=1$ for $t\le s<T$, a.s. Thus \eqref{eq:5-8-23:X} is the solution of the wealth equation under the feedback law \eqref{eq:5-8-23:fb}.

\emph{Step 3: the pair $(\pi^*,C^*)$ of \eqref{eq:5-8-23:C} is optimal for $(t,x)$, and $X^*$ is its wealth process.} The processes $\pi^*$ and $C^*$ are continuous and adapted on $[t,T]$, hence measurable with $\int_t^T\norm{\pi^*_s}^2\,ds<\infty$ and $\int_t^TC^*_s\,ds<\infty$, and $C^*\ge0$; they are a portfolio and a consumption process. By \eqref{eq:5-8-23:alg}, equation \eqref{eq:5-8-23:dX} is exactly the wealth equation \bk{(8.58)} for the pair $(\pi^*,C^*)$; since for a given pair $(\pi,C)$ the difference $D$ of two continuous solutions of \bk{(8.58)} satisfies $D_s=r\int_t^sD_u\,du$ and therefore vanishes, $X^*$ is the wealth process of $(\pi^*,C^*)$. It is nonnegative, so $(\pi^*,C^*)$ is admissible for $(t,x)$.

Next, $C^*$ is the consumption process \bk{(8.64)}. Indeed $I(cy)=c^{-\gamma}I(y)$ for $c,y>0$, so by \bk{(8.63)}, \eqref{eq:5-8-23:XY} and \eqref{eq:5-8-23:Lam},
\[
  I\bigl(\eta^{(t,x)}_s\bigr)=I\bigl(\mathscr{Y}(t,x)\bigr)\bigl(Z^t_se^{(\beta-r)(s-t)}\bigr)^{-\gamma}=\frac{x}{f(t)}\,\Lambda_s=C^*_s;\qquad t\le s\le T .
\]
By \bk{(8.65)}, $E\int_t^Te^{-\beta s}U(C^*_s)\,ds=V(t,x)$; the admissible pair $(\pi^*,C^*)$ thus attains the supremum in \bk{(8.57)}, i.e., it is optimal.

\emph{Step 4: uniqueness of the optimal pair.} Let $(\pi,C)$ be any pair that is admissible for $(t,x)$ and optimal, and let $X\ge0$ be its wealth process. We show that $C=C^*$ and $\pi=\pi^*$ for $\meas\times P$-a.e.\ $(s,\omega)\in[t,T]\times\Omega$, and that $X$ and $X^*$ are indistinguishable; in particular, the statements of the exercise hold for every optimal pair.

Put $H_s\defeq e^{-r(s-t)}Z^t_s$, so that $dH_s=-H_s(r\,ds+\theta^T\,dW_s)$. Writing \bk{(8.58)} as $dX_s=(rX_s-C_s+\pi_s^T\sigma\theta)\,ds+\pi_s^T\sigma\,dW_s$ (recall $b-r\mathbf{1}=\sigma\theta$), It\^o's product rule yields
\begin{equation}\label{eq:5-8-23:M}
  M_s\defeq H_sX_s+\int_t^sH_uC_u\,du=x+\int_t^sH_u\bigl(\pi_u^T\sigma-X_u\theta^T\bigr)\,dW_u;\qquad t\le s\le T .
\end{equation}
The integrand on the right is measurable, adapted and has a.s.\ square-integrable paths, because $H$ and $X$ are continuous and $\int_t^T\norm{\pi_u}^2du<\infty$; hence $M$ is a continuous local martingale, and it is nonnegative. By \bk{Problem 1.5.19(ii)} it is a supermartingale, so
\begin{equation}\label{eq:5-8-23:budget}
  E\bigl[H_TX_T\bigr]+E\int_t^TH_uC_u\,du=E\,M_T\le x .
\end{equation}
The same applies to $(\pi^*,C^*)$ with $M^*_s\defeq H_sX^*_s+\int_t^sH_uC^*_u\,du$; moreover, Step 3 gives $C^*=I(\eta^{(t,x)})$, and for $y\defeq\mathscr{Y}(t,x)$ the definitions \bk{(8.53)} and \bk{(8.63)} show that $\eta^{(t,x)}_s=y\,e^{(\beta-r)(s-t)}Z^t_s=Y^{(t,y)}_s$ for $t\le s\le T$; hence $C^*=I(Y^{(t,y)})$, and \bk{(8.60)}, \bk{(8.62)} yield
\begin{equation}\label{eq:5-8-23:eq}
  E\int_t^TH_uC^*_u\,du=\mathscr{X}\bigl(t,\mathscr{Y}(t,x)\bigr)=x .
\end{equation}

\textbf{(i)} \emph{Consumption.} Let $\eta_s\defeq\eta^{(t,x)}_s=\mathscr{Y}(t,x)e^{(\beta-r)(s-t)}Z^t_s$, so that $e^{-\beta s}\eta_s=y_0H_s$ with the constant $y_0\defeq e^{-\beta t}\mathscr{Y}(t,x)>0$. For each $y>0$ the function $c\mapsto U(c)-yc$ is strictly concave on $[0,\infty)$ with derivative $\delta c^{\delta-1}-y$, which vanishes only at $c=I(y)$; hence
\begin{equation}\label{eq:5-8-23:conc}
  U(I(y))-yI(y)\ge U(c)-yc\quad\text{for all }c\ge0,\ \text{with equality only for }c=I(y) .
\end{equation}
The four quantities $E\int_t^Te^{-\beta s}U(C^*_s)\,ds=E\int_t^Te^{-\beta s}U(C_s)\,ds=V(t,x)$, $E\int_t^TH_sC^*_s\,ds$ and $E\int_t^TH_sC_s\,ds$ are finite, by \eqref{eq:5-8-23:XY}, \eqref{eq:5-8-23:eq} and \eqref{eq:5-8-23:budget}. We may therefore write
\begin{align*}
  0&=E\int_t^Te^{-\beta s}\bigl[U(C^*_s)-U(C_s)\bigr]ds\\
   &=E\int_t^Te^{-\beta s}\bigl[\bigl(U(C^*_s)-\eta_sC^*_s\bigr)-\bigl(U(C_s)-\eta_sC_s\bigr)\bigr]ds
     +y_0\Bigl(E\int_t^TH_sC^*_s\,ds-E\int_t^TH_sC_s\,ds\Bigr).
\end{align*}
The first integrand is nonnegative by \eqref{eq:5-8-23:conc} (recall $C^*_s=I(\eta_s)$), and the last bracket is nonnegative by \eqref{eq:5-8-23:eq} and \eqref{eq:5-8-23:budget}. Hence both terms vanish. From the first, the integrand vanishes $\meas\times P$-a.e., and the equality case of \eqref{eq:5-8-23:conc} gives $C_s=I(\eta_s)=C^*_s$ for $\meas\times P$-a.e.\ $(s,\omega)$.

\textbf{(ii)} \emph{Wealth.} By (i) and \eqref{eq:5-8-23:eq}, $E\int_t^TH_uC_u\,du=x$, so \eqref{eq:5-8-23:budget} gives $E[H_TX_T]\le0$; as $H_T>0$ and $X_T\ge0$, we get $X_T=0$ a.s.\ and $E\,M_T=x=E\,M_t$. A supermartingale with constant expectation is a martingale; thus for $t\le s\le T$, a.s.,
\[
  H_sX_s+\int_t^sH_uC_u\,du=M_s=E\bigl[M_T\,\big|\,\F_s\bigr]=E\Bigl[\int_t^TH_uC_u\,du\,\Big|\,\F_s\Bigr],
\]
that is, $X_s=H_s^{-1}E\bigl[\int_s^TH_uC_u\,du\,\big|\,\F_s\bigr]$ a.s.
The pair $(\pi^*,C^*)$ is also admissible and optimal (Step 3), so the same formula holds for $X^*$ and $C^*$. Since $\int_s^TH_uC_u\,du=\int_s^TH_uC^*_u\,du$ a.s.\ by (i) and Fubini's theorem, we obtain $X_s=X^*_s$ a.s.\ for every $s\in[t,T]$, and by continuity of paths $X$ and $X^*$ are indistinguishable.

\textbf{(iii)} \emph{Portfolio.} Subtracting the equations \bk{(8.58)} of the two pairs and using (i), (ii), we get
\[
  \int_t^s(\pi_u-\pi^*_u)^T\sigma\,dW_u=-\int_t^s(b-r\mathbf{1})^T(\pi_u-\pi^*_u)\,du;\qquad t\le s\le T,\ \text{a.s.}
\]
The left-hand side is a continuous local martingale and the right-hand side has paths of bounded variation, so the quadratic variation of the left-hand side vanishes: $\int_t^T\norm{\sigma^T(\pi_u-\pi^*_u)}^2du=0$ a.s. Since $\sigma$ is nonsingular, $\pi=\pi^*$ for $\meas\times P$-a.e.\ $(s,\omega)$.

\emph{Step 5: conclusion.} By Steps 3 and 4, for every initial condition $(t,x)\in[0,T)\times(0,\infty)$ the optimal pair is unique up to $\meas\times P$-null sets, its wealth process is (indistinguishable from) the process \eqref{eq:5-8-23:X}, and
\[
  C^*_s=\frac{X^*_s}{f(s)},\qquad \pi^*_s=\frac{1}{1-\delta}\,(\sigma\sigma^T)^{-1}(b-r\mathbf{1})\,X^*_s ;\qquad t\le s<T,
\]
which are linear functions of $X^*_s$ (with $1/f(s)=k/(1-e^{-k(T-s)})$ if $k\ne0$ and $1/f(s)=1/(T-s)$ if $k=0$), in agreement with \bk{(8.81)}, \bk{(8.82)}. Finally, in \eqref{eq:5-8-23:X} the factor $\Lambda_T$ is finite and $f(T)=0$ by \eqref{eq:5-8-23:f}; therefore
\[
  X^*_T=x\,\frac{f(T)}{f(t)}\,\Lambda_T=0\quad\text{a.s.,}\qquad\text{whereas}\qquad X^*_s>0\ \text{ for }t\le s<T .
\]
The investor is never bankrupt before $T$ and exhausts the wealth exactly at the terminal time; this is as it should be, since terminal wealth carries no utility. (The relation $X_T=0$ a.s.\ for an arbitrary optimal pair was also obtained independently in Step 4(ii).)

\rmk{For $x=0$ the wealth process of an admissible pair is a nonnegative process for which \eqref{eq:5-8-23:budget} reads $E[H_TX_T]+E\int_t^TH_uC_u\,du\le0$; hence $C=0$ a.e., $M\equiv0$ in \eqref{eq:5-8-23:M}, $X\equiv0$, and then $\pi=0$ a.e.\ as in Step 4(iii). The relations \eqref{eq:5-8-23:fb} hold trivially in this case. As a check on \eqref{eq:5-8-23:C}: $\theta^T(W_s-W_t)$ is normal with mean zero and variance $\norm\theta^2(s-t)$, whence $E\bigl[H_sC^*_s\bigr]=(x/f(t))e^{-k(s-t)}$ and $E\bigl[e^{-\beta s}(C^*_s)^\delta\bigr]=e^{-\beta t}(x/f(t))^{\delta}e^{-k(s-t)}$; integrating over $[t,T]$ gives $x$, in accordance with \eqref{eq:5-8-23:eq}, and $e^{-\beta t}f(t)^{1-\delta}x^\delta=V(t,x)$, in accordance with \bk{Exercise 8.20}.}
\end{kex}


%% ========================================================================
\solchapter{6}{P.~L\'evy's Theory of Brownian Local Time}

\begin{center}
\begin{minipage}{0.86\textwidth}
\small
The 14 problems and exercises of Chapter 6 for which the book prints no solution: 2.2, 2.11, 2.19, 2.20, 3.4, 3.5, 3.8, 4.4, 4.5, 4.6, 4.10, 4.15, 5.3, 5.4.  For 1 more, 2.5, the book prints only a hint or solves only some of the parts; what it leaves out is supplied here.  The book's own ``Solutions to Selected Problems'' fully covers 2.12, 2.16, 2.24, 3.7, 3.11, 4.11, which are therefore omitted.
\end{minipage}
\end{center}

\bigskip
\hrule
\bigskip

\solsection{6.2}{Alternate Representations of Brownian Local Time}

\begin{kex}{Problem}{2.2}{}
\emph{Step 1: the functional equation.} Since $G$ is strictly positive, $g \defeq \log G\colon I\to\R$ is well defined; it is monotone in the same sense as $G$, and \bk{(2.13)} becomes
\begin{equation}\label{eq:6-2-02:add}
g(b+c)=g(b)+g(c);\qquad b,c\in I .
\end{equation}
Replacing $G$ by $1/G$ (which is again positive, monotone and satisfies \bk{(2.13)}) if necessary, we may assume that $g$ is nondecreasing; if we show $1/G(b)=e^{-\mu b}$, then $G(b)=e^{-\lambda b}$ with $\lambda=-\mu$.

Taking $b=c=0$ in \eqref{eq:6-2-02:add} gives $g(0)=0$, and induction on $n$ gives $g(nb)=n\,g(b)$ for all $n\in\N$ and $b\in I$. Set $a\defeq g(1)$. For a rational number $r=m/n$ with $m\ge 0$ and $n\ge 1$ integers we get $n\,g(r)=g(nr)=g(m)=m\,g(1)$ (trivially so when $m=0$), that is, $g(r)=ar$. If $I=\R$, then \eqref{eq:6-2-02:add} with $c=-b$ gives $g(-b)=-g(b)$, so that $g(r)=ar$ holds for negative rationals as well. Thus
\begin{equation*}
g(r)=ar\qquad\text{for every } r\in I\cap\Q .
\end{equation*}
Now let $b\in I$ be arbitrary, and choose sequences $\{p_n\}$, $\{q_n\}$ in $I\cap\Q$ with $p_n\le b\le q_n$, $p_n\to b$ and $q_n\to b$ (if $I=[0,\infty)$ and $b=0$, take $p_n=0$). Because $g$ is nondecreasing,
\begin{equation*}
a\,p_n=g(p_n)\le g(b)\le g(q_n)=a\,q_n ,
\end{equation*}
and letting $n\to\infty$ we obtain $g(b)=ab$. Hence $G(b)=e^{ab}=e^{-\lambda b}$ for all $b\in I$, with $\lambda\defeq -a=-\log G(1)\in\R$.

\emph{Step 2: the random variable $\tau_1$.} We make precise the objects appearing in \bk{(2.12)}. Let $\Omega_0$ be the event on which $M$ of \bk{(2.2)} is continuous and nondecreasing with $M_0=0$, $\lim_{t\to\infty}M_t=\infty$, on which the first identity of \bk{(2.1)} holds for all $t\ge 0$, and on which $S_0=0$. We have $P(\Omega_0)=1$: $M$ is the running maximum of the Brownian motion $B$, which is unbounded from above almost surely; \bk{(2.1)} holds almost surely; and $S_0=T_0=0$ almost surely because $S$ is a modification of $T$ (\bk{Theorem 2.1}). We fix $\omega\in\Omega_0$ in this paragraph. For $0\le b<c<\infty$ the times $T_b\le S_b$ are finite, and $S_b\le T_c$ because $M_{T_c}=c>b$; moreover, if $M_t>b$ then by the intermediate value theorem $T_c\le t$ for every $c\in(b,M_t)$. Therefore $T_{b+}\defeq\lim_{c\downarrow b}T_c=S_b$, and the jump of $T$ at $b$ has size $S_b-T_b$. Fix $\ell>0$ and let
\begin{equation*}
J\defeq\{x\ge 0;\ S_x-T_x>\ell\}
\end{equation*}
be the set of locations of the jumps of $T$ whose size exceeds $\ell$. If $x_1<\dots<x_k$ are points of $J\cap[0,b]$, then the intervals $(T_{x_i},S_{x_i})$ are pairwise disjoint (since $S_{x_i}\le T_{x_{i+1}}$) and contained in $[0,S_b]$, whence $k\ell<S_b<\infty$. So $J\cap[0,b]$ is finite for every $b$, and $J$ has a smallest element whenever it is nonempty. The ``minimal $b$ for which $T$ has a jump on $[0,b]$ of size exceeding $\ell$'' is thus
\begin{equation*}
\tau_1\defeq\min J\quad(\text{with }\tau_1\defeq\infty\text{ if }J=\emptyset),\qquad\text{and}\qquad \{\tau_1>b\}=\{J\cap[0,b]=\emptyset\};\quad 0\le b<\infty .
\end{equation*}
Since $S_0-T_0=0$ on $\Omega_0$, we have $0\notin J$, and therefore
\begin{equation}\label{eq:6-2-02:pos}
\tau_1>0\quad\text{on }\Omega_0 .
\end{equation}

\emph{$\tau_1$ is a random variable.} We claim that for $\omega\in\Omega_0$ and $0\le b<\infty$,
\begin{equation}\label{eq:6-2-02:meas}
J\cap[0,b]\neq\emptyset
\iff
\exists\,\eps\in\Q\cap(0,\infty)\ \ \forall\,n\ge 1\ \ \exists\,k\in\{0,\dots,n-1\}:\quad
S_{(k+1)b/n}-T_{kb/n}>\ell+\eps .
\end{equation}
Indeed, if $x\in J\cap[0,b]$, choose a rational $\eps>0$ with $S_x-T_x>\ell+\eps$; for each $n$ there is $k\in\{0,\dots,n-1\}$ with $kb/n\le x\le(k+1)b/n$, and then $S_{(k+1)b/n}-T_{kb/n}\ge S_x-T_x>\ell+\eps$ because $S$ and $T$ are nondecreasing. Conversely, suppose the right-hand side holds with indices $k_n$. Along a subsequence, $k_nb/n\to x$ for some $x\in[0,b]$, and then also $(k_n+1)b/n\to x$. For every $\delta>0$ and all large $n$ in the subsequence we have $(x-\delta)^+\le k_nb/n$ and $(k_n+1)b/n\le x+\delta$, so $S_{x+\delta}-T_{(x-\delta)^+}>\ell+\eps$. Letting $\delta\downarrow 0$ and using the right-continuity of $S$ and the left-continuity of $T$ (and $T_{(x-\delta)^+}=T_0$ if $x=0$), we get $S_x-T_x\ge\ell+\eps>\ell$, i.e., $x\in J\cap[0,b]$. Each $S_c$ and $T_c$ is a random variable, so \eqref{eq:6-2-02:meas} shows that $\{\tau_1>b\}\cap\Omega_0\in\F$ for every $b\ge0$; the filtration satisfies the usual conditions, so $\tau_1$ (defined arbitrarily off $\Omega_0$) is a random variable with values in $(0,\infty]$.

\emph{$\tau_1$ is finite with positive probability.} Consider the event
\begin{equation*}
E\defeq\Bigl\{W_1>1,\ \min_{1\le s\le 1+2\ell}(W_s-W_1)>-1\Bigr\}.
\end{equation*}
The process $\{W_1-W_{1+u};\,u\ge0\}$ is a standard Brownian motion independent of $W_1$, and the second condition defining $E$ says that this Brownian motion does not reach the level $1$ during $[0,2\ell]$, which is implied by its not reaching the level $1$ before time $3\ell$. By \bk{(2.6.2)}, therefore,
\begin{equation*}
P(E)\ge P[W_1>1]\cdot\bigl(1-2P[W_{3\ell}>1]\bigr)>0 .
\end{equation*}
On $E\cap\Omega_0$ we have $W_s>0$ for all $s\in[1,1+2\ell]$, hence $B_s<M_s$ on this interval by \bk{(2.1)}. It follows that $M$ is constant on $[1,1+2\ell]$: if we had $M_t>M_1$ for some $t\in(1,1+2\ell]$, the maximum of $B$ over $[0,t]$ would be attained at some $s\in(1,t]$, and then $B_s=M_t=M_s$, a contradiction. With $x\defeq M_1$ we thus have $T_x\le 1$ and $S_x\ge 1+2\ell$, so $S_x-T_x\ge 2\ell>\ell$, $x\in J$ and $\tau_1\le x<\infty$. Consequently
\begin{equation}\label{eq:6-2-02:fin}
P[\tau_1<\infty]\ge P(E)>0 .
\end{equation}

\emph{Step 3: the distribution of $\tau_1$.} Let $G(b)\defeq P[\tau_1>b]$, $0\le b<\infty$. This function is nonincreasing, and \bk{(2.12)} says that
\begin{equation}\label{eq:6-2-02:mult}
G(b+c)=G(b)\,G(c);\qquad 0\le b,c<\infty
\end{equation}
(for those $c$ with $G(c)>0$ this is \bk{(2.12)} multiplied by $P[\tau_1>c]$; if $G(c)=0$ the conditional probability in \bk{(2.12)} is not defined, but \eqref{eq:6-2-02:mult} holds trivially because $0\le G(b+c)\le G(c)=0$). In order to apply Step~1 we must check that $G$ is strictly positive. Suppose that $G(c)=0$ for some $c\ge0$. By \eqref{eq:6-2-02:mult}, $G(c/2)^2=G(c)=0$, and inductively $G(c2^{-n})=0$ for all $n\ge 0$. The events $\{\tau_1>c2^{-n}\}$ increase to $\{\tau_1>0\}$ as $n\to\infty$, so $P[\tau_1>0]=\lim_{n}G(c2^{-n})=0$, which contradicts \eqref{eq:6-2-02:pos}. Hence $G\colon[0,\infty)\to(0,1]$ is monotone and satisfies \bk{(2.13)} with $I=[0,\infty)$, and Step~1 provides $\lambda\in\R$ with
\begin{equation*}
P[\tau_1>b]=e^{-\lambda b};\qquad 0\le b<\infty .
\end{equation*}
Because $e^{-\lambda}=G(1)\le 1$, we have $\lambda\ge 0$. If $\lambda$ were zero, we would have $P[\tau_1>n]=1$ for every $n\in\N$ and thus $P[\tau_1=\infty]=1$, contradicting \eqref{eq:6-2-02:fin}. Therefore $\lambda>0$; in particular $\tau_1<\infty$ almost surely, and $\tau_1$ is exponentially distributed with parameter $\lambda$.

\rmk{The book derives \bk{(2.12)} only heuristically, and Step~3 takes it as given, as the statement of the problem does; what Steps~2 and~3 add is that the two facts needed besides \bk{(2.12)}, namely $P[\tau_1>0]=1$ (for the positivity of $G$) and $P[\tau_1<\infty]>0$ (for $\lambda>0$), hold. The value of $\lambda$ is not determined by this argument.}
\end{kex}

\begin{kex}{Problem}{2.5}{}
The book's Section 6.6 constructs the random measure $\nu$ and sketches the verification of condition (i) of \bk{Definition 2.3} and of $E\nu(C)=\lambda(C)$; it leaves condition (ii) (``straightforward'') to the reader. We recall the construction, add the few measurability facts on which it rests, and prove (i) and (ii) in full.

\emph{The construction.} Let $\lambda$ be $\sigma$-finite with $\lambda(H)=\infty$. There are pairwise disjoint sets $D_1,D_2,\dots\in\cH$ with union $H$ and $\lambda(D_n)<\infty$; since $\sum_n\lambda(D_n)=\infty$, infinitely many of them have positive measure. Adjoining the union of the $\lambda$-null ones (again a $\lambda$-null set in $\cH$) to one of the others, we obtain a partition $\{H_i\}_{i=1}^\infty\subseteq\cH$ of $H$ with $0<\lambda(H_i)<\infty$ for every $i\ge 1$. Put
\[
\lambda_i(C)\defeq\lambda(C\cap H_i),\qquad C\in\cH,
\]
a finite measure on $(H,\cH)$ of total mass $\lambda_i(H)=\lambda(H_i)$, and note that $\sum_{i=1}^\infty\lambda_i=\lambda$. On a suitable probability space $(\Omega,\F,P)$ (a countable product of probability spaces) there are independent random variables $N_i$, $\xi^{(i)}_j$; $i,j\ge 1$, where $N_i$ is Poisson with mean $\lambda(H_i)$ and $\xi^{(i)}_j$ is $H$-valued with distribution $\lambda_i/\lambda_i(H)$. For $i\ge 1$, $\omega\in\Omega$ and $C\in\cH$ we set
\[
\nu_i(C)\defeq\sum_{j=1}^{N_i}\ind_C\bigl(\xi^{(i)}_j\bigr)
=\sum_{j=1}^{\infty}\ind_{\{j\le N_i\}}\ind_C\bigl(\xi^{(i)}_j\bigr),
\qquad
\nu(C)\defeq\sum_{i=1}^\infty\nu_i(C).
\]
For each $\omega$, $\nu_i$ is a finite sum of unit point masses, hence an $\N_0$-valued measure on $(H,\cH)$, and $\nu$, being a countable sum of such measures, is a measure with values in $\N_0\cup\{\infty\}$ (countable additivity follows by interchanging two summations of nonnegative terms). For fixed $C$, the second expression for $\nu_i(C)$ shows that it is a random variable, and so is $\nu(C)$, as the limit of an increasing sequence of random variables. Thus $\nu$ is a mapping of the kind considered in \bk{Definition 2.3}. By the case of finite total mass treated in the proof of \bk{Theorem 2.4}, applied to $\lambda_i$, each $\nu_i$ is a Poisson random measure with intensity $\lambda_i$. Monotone convergence gives $E\nu(C)=\sum_i\lambda_i(C)=\lambda(C)$.

\emph{Condition (i).} Fix $C\in\cH$ and put $S_n\defeq\sum_{i=1}^n\nu_i(C)$, $\Lambda_n\defeq\sum_{i=1}^n\lambda_i(C)$; $n\ge 1$. The random variables $\nu_1(C),\nu_2(C),\dots$ are independent (this is the case $m=1$ of Step 1 below, whose proof does not use condition (i) for $\nu$), and $\nu_i(C)$ is Poisson with mean $\lambda_i(C)$ if $\lambda_i(C)>0$, whereas $\nu_i(C)=0$ a.s.\ if $\lambda_i(C)=0$, because then $E\nu_i(C)=0$ and $\nu_i(C)\ge 0$. A sum of finitely many independent Poisson random variables is Poisson (multiply the generating functions $E s^{X}=e^{-\mu(1-s)}$, $0\le s\le 1$), and a.s.\ vanishing summands change nothing; therefore
\begin{equation}\label{eq:6-2-05:a}
P[S_n\le k]=e^{-\Lambda_n}\sum_{j=0}^k\frac{\Lambda_n^j}{j!};\qquad k\in\N_0,\ n\ge 1,
\end{equation}
which is also correct when $\Lambda_n=0$ (with $0^0=1$). Now $S_n\uparrow\nu(C)$ and $\Lambda_n\uparrow\lambda(C)$ as $n\to\infty$. Since the sequence $\{S_n\}$ is nondecreasing with limit $\nu(C)$, we have $\{S_n\le k\}\downarrow\{\nu(C)\le k\}$, and so $P[\nu(C)\le k]=\lim_{n\to\infty}P[S_n\le k]$ for every $k\in\N_0$. We let $n\to\infty$ in \eqref{eq:6-2-05:a}. If $\lambda(C)<\infty$, the right-hand side, a continuous function of $\Lambda_n\in[0,\infty)$, converges and
\[
P[\nu(C)\le k]=e^{-\lambda(C)}\sum_{j=0}^k\frac{\lambda(C)^j}{j!};\qquad k\in\N_0.
\]
Letting $k\to\infty$ we see that $P[\nu(C)<\infty]=1$, and taking differences in $k$ that $P[\nu(C)=k]=e^{-\lambda(C)}\lambda(C)^k/k!$; thus $\nu(C)$ is Poisson with mean $\lambda(C)$ (and $\nu(C)=0$ a.s.\ if $\lambda(C)=0$). If $\lambda(C)=\infty$, then $\Lambda_n\to\infty$, the right-hand side of \eqref{eq:6-2-05:a} tends to zero for each fixed $k$, and so $P[\nu(C)\le k]=0$ for every $k\in\N_0$; that is, $\nu(C)=\infty$ a.s. This is condition (i) of \bk{Definition 2.3}. \rmk{The book's solution treats the case $\lambda(C)=\infty$ with the second half of the Borel--Cantelli lemma; the limit in \eqref{eq:6-2-05:a} handles both cases at once.}

\emph{Condition (ii).} Let $C_1,\dots,C_m$ be pairwise disjoint sets in $\cH$.

\emph{Step 1: the random variables $\nu_i(C_k)$; $i\ge 1$, $1\le k\le m$, are independent.} It suffices to show this for $i\le n$, $1\le k\le m$, with $n\ge 1$ arbitrary. The random vector $V_i\defeq\bigl(\nu_i(C_1),\dots,\nu_i(C_m)\bigr)$ is measurable with respect to $\F^{(i)}\defeq\sigma\bigl(N_i,\xi^{(i)}_1,\xi^{(i)}_2,\dots\bigr)$. The $\sigma$-fields $\F^{(1)},\F^{(2)},\dots$ are generated by disjoint subfamilies of an independent family of random variables, and are therefore independent; hence $V_1,\dots,V_n$ are independent. Moreover, the components of each $V_i$ are independent, because $\nu_i$ is a Poisson random measure and so satisfies condition (ii) of \bk{Definition 2.3}. Consequently, for arbitrary $B_{ik}\subseteq\N_0$,
\begin{align*}
P\bigl[\nu_i(C_k)\in B_{ik};\ 1\le i\le n,\ 1\le k\le m\bigr]
&=\prod_{i=1}^n P\bigl[\nu_i(C_k)\in B_{ik};\ 1\le k\le m\bigr]\\
&=\prod_{i=1}^n\prod_{k=1}^m P\bigl[\nu_i(C_k)\in B_{ik}\bigr],
\end{align*}
which is the asserted independence.

\emph{Step 2: grouping.} For $1\le k\le m$ let $\G_k\defeq\sigma\bigl(\nu_i(C_k);\ i\ge 1\bigr)$ and let $\cP_k$ be the collection of events of the form $\bigcap_{i=1}^n\{\nu_i(C_k)\in B_i\}$ with $n\ge 1$ and $B_1,\dots,B_n\subseteq\N_0$. Each $\cP_k$ is closed under finite intersections, contains $\Omega$ and generates $\G_k$, and by Step 1 we have
\[
P(A_1\cap\dots\cap A_m)=P(A_1)\cdots P(A_m)\qquad\text{whenever } A_k\in\cP_k,\ 1\le k\le m.
\]
We extend this identity to $A_k\in\G_k$ one index at a time. Fix $A_2\in\cP_2,\dots,A_m\in\cP_m$; the collection of $A_1\in\F$ for which the identity holds contains $\cP_1$ (in particular $\Omega$), is closed under complementation (subtract the identity for $A_1$ from the one for $\Omega$) and under countable disjoint unions (by the countable additivity of $P$), so it is a Dynkin system and contains $\sigma(\cP_1)=\G_1$ by Dynkin's $\pi$-$\lambda$ theorem. Next we fix $A_1\in\G_1$, $A_3\in\cP_3,\dots,A_m\in\cP_m$ and repeat the argument for $A_2$, and so on. After $m$ steps the identity holds for all $A_k\in\G_k$, $1\le k\le m$; that is, $\G_1,\dots,\G_m$ are independent.

\emph{Step 3: conclusion.} For each $k$, the $\N_0\cup\{\infty\}$-valued random variable $\nu(C_k)=\lim_{n\to\infty}\sum_{i=1}^n\nu_i(C_k)$ is $\G_k$-measurable. By Step 2, $\nu(C_1),\dots,\nu(C_m)$ are independent, which is condition (ii) of \bk{Definition 2.3}. Together with condition (i) and $E\nu(C)=\lambda(C)$; $C\in\cH$, this shows that $\nu$ is a Poisson random measure with intensity measure $\lambda$, and completes the proof of \bk{Theorem 2.4}.
\end{kex}

\begin{kex}{Problem}{2.11}{}
Write $D \defeq D[0,\infty)$, and let $\mathscr{R}$ be the collection of all sets displayed in the statement; we call them \emph{rectangles}. By \bk{Definition 2.9}, $\B(D)$ is the smallest $\sigma$-field of subsets of $D$ containing all sets of the form \bk{(2.2.1)}, that is, all sets
\begin{equation}\label{eq:6-2-11:cyl}
\{y \in D;\ (y(s_1),\dots,y(s_m)) \in A\},
\qquad m \ge 1,\quad s_1,\dots,s_m \in [0,\infty),\quad A \in \B(\R^m),
\end{equation}
where the times $s_j$ need be neither ordered nor distinct. We denote by $\cC$ the collection of the sets \eqref{eq:6-2-11:cyl}, so that $\B(D)=\sigma(\cC)$.

\emph{Step 1: $\mathscr{R}$ is closed under pairwise intersection.} Let
\begin{equation*}
R = \{y \in D;\ y(t_i) \in \Gamma_i,\ i=1,\dots,n\},
\qquad
R' = \{y \in D;\ y(t'_j) \in \Gamma'_j,\ j=1,\dots,n'\}
\end{equation*}
be two rectangles, and let $u_1 < u_2 < \dots < u_p$ be the elements of the finite set $\{t_1,\dots,t_n\} \cup \{t'_1,\dots,t'_{n'}\}$, arranged in increasing order. For $k=1,\dots,p$ we set $\Lambda_k \defeq \Gamma_i$ if $u_k = t_i$ for some $i$ and $\Lambda_k \defeq \R$ otherwise, and $\Lambda'_k \defeq \Gamma'_j$ if $u_k = t'_j$ for some $j$ and $\Lambda'_k \defeq \R$ otherwise; these are well defined because the $t_i$ are distinct, and so are the $t'_j$. Since a condition of the form $y(u_k) \in \R$ is vacuous,
\begin{equation*}
R = \{y \in D;\ y(u_k) \in \Lambda_k,\ k=1,\dots,p\},
\qquad
R' = \{y \in D;\ y(u_k) \in \Lambda'_k,\ k=1,\dots,p\},
\end{equation*}
and therefore
\begin{equation*}
R \cap R' = \{y \in D;\ y(u_1) \in \Lambda_1 \cap \Lambda'_1, \dots, y(u_p) \in \Lambda_p \cap \Lambda'_p\},
\end{equation*}
with $0 \le u_1 < \dots < u_p < \infty$ and $\Lambda_k \cap \Lambda'_k \in \B(\R)$. Thus $R \cap R' \in \mathscr{R}$.

\emph{Step 2: $\sigma(\mathscr{R}) = \B(D)$.} Every rectangle is of the form \eqref{eq:6-2-11:cyl}, with $m=n$, $s_i = t_i$ and $A = \Gamma_1 \times \dots \times \Gamma_n \in \B(\R^n)$; hence $\mathscr{R} \subseteq \cC$ and $\sigma(\mathscr{R}) \subseteq \sigma(\cC) = \B(D)$.

For the reverse inclusion it suffices to show $\cC \subseteq \sigma(\mathscr{R})$. Fix $m \ge 1$ and $s_1,\dots,s_m \in [0,\infty)$, define $\pi : D \to \R^m$ by $\pi(y) \defeq (y(s_1),\dots,y(s_m))$, and let
\begin{equation*}
\mathscr{A} \defeq \{A \in \B(\R^m);\ \pi^{-1}(A) \in \sigma(\mathscr{R})\}.
\end{equation*}
Because $\pi^{-1}$ commutes with complements and countable unions, and $\pi^{-1}(\R^m) = D$, the collection $\mathscr{A}$ is a $\sigma$-field of subsets of $\R^m$. For $\Gamma_1,\dots,\Gamma_m \in \B(\R)$ we have
\begin{equation*}
\pi^{-1}(\Gamma_1 \times \dots \times \Gamma_m) = \bigcap_{j=1}^m \{y \in D;\ y(s_j) \in \Gamma_j\},
\end{equation*}
and each of the $m$ sets on the right is a rectangle (with $n=1$), so that this finite intersection belongs to $\sigma(\mathscr{R})$. Hence $\mathscr{A}$ contains all measurable rectangles $\Gamma_1 \times \dots \times \Gamma_m$; these include the products of open intervals, and every open subset of $\R^m$ is a countable union of such products, so they generate $\B(\R^m)$. Consequently $\mathscr{A} = \B(\R^m)$, which says that every set \eqref{eq:6-2-11:cyl} with these $s_1,\dots,s_m$ lies in $\sigma(\mathscr{R})$. As $m$ and $s_1,\dots,s_m$ were arbitrary, $\cC \subseteq \sigma(\mathscr{R})$, and thus $\B(D) = \sigma(\cC) \subseteq \sigma(\mathscr{R})$.

\emph{Step 3: the sets on which $P$ and $\wt{P}$ agree form a Dynkin system.} Let
\begin{equation*}
\cD \defeq \{A \in \B(D);\ P(A) = \wt{P}(A)\}.
\end{equation*}
We verify the three conditions of \bk{Definition 2.1.2}, with $\Omega = D$.
\begin{enumerate}[label=(\roman*)]
\item $D \in \cD$, because $P(D) = 1 = \wt{P}(D)$.
\item If $A, B \in \cD$ and $B \subseteq A$, then $A \setminus B \in \B(D)$ and, both measures being finite,
$P(A \setminus B) = P(A) - P(B) = \wt{P}(A) - \wt{P}(B) = \wt{P}(A \setminus B)$; hence $A \setminus B \in \cD$.
\item If $\{A_k\}_{k=1}^\infty \subseteq \cD$ and $A_1 \subseteq A_2 \subseteq \cdots$, then $\bigcup_{k} A_k \in \B(D)$ and, by the continuity from below of the two measures,
$P\bigl(\bigcup_{k} A_k\bigr) = \lim_{k\to\infty} P(A_k) = \lim_{k\to\infty} \wt{P}(A_k) = \wt{P}\bigl(\bigcup_{k} A_k\bigr)$; hence $\bigcup_k A_k \in \cD$.
\end{enumerate}

\emph{Step 4: conclusion.} By hypothesis $\mathscr{R} \subseteq \cD$. By Step 1 the collection $\mathscr{R}$ is closed under pairwise intersection, and by Step 3 $\cD$ is a Dynkin system containing it; the Dynkin System \bk{Theorem 2.1.3} then yields $\sigma(\mathscr{R}) \subseteq \cD$. By Step 2, $\sigma(\mathscr{R}) = \B(D)$, so that $\cD = \B(D)$: the measures $P$ and $\wt{P}$ agree on $\B(D[0,\infty))$.
\end{kex}

\begin{kex}{Exercise}{2.19}{}
We fix $\delta>0$ and write
\[
\lambda \defeq \int_\delta^\infty \frac{d\ell}{\sqrt{2\pi\ell^3}}
=\frac{1}{\sqrt{2\pi}}\Bigl[-2\ell^{-1/2}\Bigr]_{\ell=\delta}^{\ell=\infty}
=\sqrt{\frac{2}{\pi\delta}},
\]
\[
N_b\defeq N_b^\delta=\nu\bigl((0,b]\times[\delta,\infty)\bigr),
\qquad X_b\defeq X_b(\delta)=N_b-\lambda b ;\qquad 0\le b<\infty,
\]
where $\nu$ is the random measure of \bk{(2.36)}; here $b\sqrt{2/\pi\delta}$ in the statement means $b\sqrt{2/(\pi\delta)}=\lambda b$, as in the two displays and in \bk{(2.37)}. By \bk{Corollary 2.13} and the computation of the L\'evy measure of $S$ preceding \bk{(2.36)}, $\nu$ is a Poisson random measure on $(0,\infty)^2$ (\bk{Definition 2.3}) with intensity measure $dt\,d\ell/\sqrt{2\pi\ell^3}$.

\emph{Step 1: the paths of $N$.} Let $\Omega_0$ be the event on which $S_b<\infty$ for every $b\ge 0$; since $S_b=\inf\{t\ge0;\,B_t>b\}$ by \bk{(2.3)} and the Brownian motion $B$ satisfies $\sup_{t}B_t=\infty$ a.s., we have $P(\Omega_0)=1$. For $\omega\in\Omega_0$ the function $b\mapsto S_b(\omega)$ is real-valued, nondecreasing and right-continuous, being the right-continuous inverse \bk{(2.3)} of the continuous, nondecreasing function $M_\cdot(\omega)$. The jumps of size at least $\delta$ of such a function in $(0,b]$ have total size at most $S_b(\omega)-S_0(\omega)$, so there are at most $S_b(\omega)/\delta$ of them: $N_b(\omega)<\infty$ for all $b$, and the set $D(\omega)$ of the locations of the jumps of size at least $\delta$ is a subset of $(0,\infty)$ with finitely many points in every bounded interval. Since $N_b(\omega)=\#\bigl(D(\omega)\cap(0,b]\bigr)$, the path $b\mapsto N_b(\omega)$ is integer-valued, nondecreasing, right-continuous with left limits, and $N_0(\omega)=\nu_\omega(\emptyset)=0$. Thus the paths of $N$, and those of $X$, are right-continuous on $\Omega_0$. Moreover, $\nu_\omega$ being a measure,
\begin{equation}\label{eq:6-2-19:a}
N_b-N_a=\nu\bigl((a,b]\times[\delta,\infty)\bigr)\quad\text{on }\Omega_0;\qquad 0\le a<b<\infty .
\end{equation}

\rmk{As everywhere in this section of the book (compare \bk{Definition 2.6}, which asks only that almost every path of a subordinator be right-continuous), the null event $\Omega\setminus\Omega_0$, on which $\nu$ is not properly defined by \bk{(2.36)}, is disregarded; ``right-continuous'' is to be understood accordingly.}

\emph{Step 2: $N$ is a Poisson process with intensity $\lambda$ relative to $\{\F^N_b\}$, and $\F^X_b=\F^N_b$.} For $0\le a<b$, the set $C=(a,b]\times[\delta,\infty)$ has intensity $(b-a)\int_\delta^\infty d\ell/\sqrt{2\pi\ell^3}=\lambda(b-a)<\infty$, so by \bk{Definition 2.3}(i) and \eqref{eq:6-2-19:a} the increment $N_b-N_a$ is a Poisson random variable with mean $\lambda(b-a)$. If $0=b_0<b_1<\dots<b_n$, the sets $(b_{k-1},b_k]\times[\delta,\infty)$, $1\le k\le n$, are pairwise disjoint, and \bk{Definition 2.3}(ii) with \eqref{eq:6-2-19:a} shows that the increments $N_{b_k}-N_{b_{k-1}}$, $1\le k\le n$, are independent. Hence $N$ has independent increments, and they are stationary because the law of $N_b-N_a$ depends on $b-a$ only.

Fix $0\le a<b$. For $0\le c_1<\dots<c_m\le a$ the vector $(N_{c_1},\dots,N_{c_m})$ is a function of the increments of $N$ over $(0,c_1],(c_1,c_2],\dots,(c_{m-1},c_m]$, and is therefore independent of $N_b-N_a$ (we may add the interval $(c_m,a]$ to the list). The events $\{(N_{c_1},\dots,N_{c_m})\in\Gamma\}$, with $m\ge1$, $c_i$ as above and $\Gamma\in\B(\R^m)$, form a $\pi$-system which generates $\F^N_a=\sigma(N_c;\,0\le c\le a)$, and the events independent of $N_b-N_a$ form a $\lambda$-system; by Dynkin's $\pi$-$\lambda$ theorem,
\begin{equation}\label{eq:6-2-19:b}
N_b-N_a \text{ is independent of } \F^N_a \text{ and Poisson distributed with mean } \lambda(b-a).
\end{equation}
Together with Step 1 this says that $\{N_b,\F^N_b;\,0\le b<\infty\}$ is a Poisson process with intensity $\lambda$ in the sense of \bk{Definition 1.3.3}. Finally, $X_b=N_b-\lambda b$ and $N_b=X_b+\lambda b$ differ by a constant for each $b$, so $\F^X_b=\F^N_b$ for every $b$; we write $\F_b$ for this $\sigma$-field.

\emph{Step 3: $X$ is a square-integrable martingale with stationary, independent increments.} For $0\le a<b$ we have $X_b-X_a=(N_b-N_a)-\lambda(b-a)$, so by \eqref{eq:6-2-19:b} the increment $X_b-X_a$ is independent of $\F_a$, its law depends only on $b-a$, and
\begin{equation}\label{eq:6-2-19:c}
E(X_b-X_a)=0,\qquad E(X_b-X_a)^2=\Var(N_b-N_a)=\lambda(b-a),
\end{equation}
because a Poisson random variable with mean $\theta$ has variance $\theta$. In particular (take $a=0$, $X_0=0$) $EX_b^2=\lambda b<\infty$, the process $X$ is adapted to $\{\F_b\}$, and
\[
E[X_b\mid\F_a]=X_a+E(X_b-X_a)=X_a\quad\text{a.s.};\qquad 0\le a<b<\infty .
\]
So $\{X_b,\F_b;\,0\le b<\infty\}$ is a right-continuous (Step 1), square-integrable martingale; this is the compensated Poisson process of \bk{Problem 1.3.4}. Its increments over disjoint intervals are those of $N$ shifted by constants, hence independent and stationary.

\emph{Step 4: $X_b^2-\lambda b$ is a martingale.} Let $0\le a<b$. The product $X_a(X_b-X_a)$ is integrable by the Cauchy--Schwarz inequality, and from $X_b^2-X_a^2=(X_b-X_a)^2+2X_a(X_b-X_a)$, the independence of $X_b-X_a$ and $\F_a$, the $\F_a$-measurability of $X_a$ and \eqref{eq:6-2-19:c} we obtain
\[
E\bigl[X_b^2-X_a^2\,\big|\,\F_a\bigr]=E(X_b-X_a)^2+2X_a\,E(X_b-X_a)=\lambda(b-a)\quad\text{a.s.}
\]
Thus $\{X_b^2-\lambda b,\F_b;\,0\le b<\infty\}$ is a martingale, with right-continuous paths on $\Omega_0$.

\emph{Step 5: the quadratic variation.} The quadratic variation of \bk{Definition 1.5.3} is defined for elements of $\cM_2$, relative to a filtration which satisfies the usual conditions (\bk{Definition 1.2.25}); the filtration $\{\F_b\}$ does not, so we first enlarge it, as in \bk{Example 1.5.4}. We work on the completion of $(\Omega,\F,P)$; let $\cN$ be the collection of its $P$-null sets, $\G_b\defeq\sigma(\F_b\cup\cN)$ and $\G_{b+}\defeq\bigcap_{c>b}\G_c$. The filtration $\{\G_{b+}\}$ is right-continuous and $\G_{0+}\supseteq\cN$, so it satisfies the usual conditions, and $X$ is adapted to it.

\begin{claim}
Let $\{Y_b,\F_b;\,0\le b<\infty\}$ be a martingale, almost every path of which is right-continuous. Then $\{Y_b,\G_{b+};\,0\le b<\infty\}$ is a martingale.
\end{claim}
\begin{proof}
Let $0\le a<b$. The class of sets $G$ for which there is $F\in\F_c$ with $G\,\triangle\,F\in\cN$ is a $\sigma$-field containing $\F_c\cup\cN$, hence contains $\G_c$; consequently $E[(Y_b-Y_c)\ind_G]=E[(Y_b-Y_c)\ind_F]=0$ for $G\in\G_c$, $c<b$, that is, $Y_c=E[Y_b\mid\G_c]$ a.s. Now let $G\in\G_{a+}$ and choose $a_n\downarrow a$ with $a<a_n<b$. Then $G\in\G_{a_n}$, so $E[Y_b\ind_G]=E[Y_{a_n}\ind_G]$ for every $n$. The family $\{Y_{a_n}\}=\{E[Y_b\mid\G_{a_n}]\}$ of conditional expectations of one integrable random variable is uniformly integrable, and $Y_{a_n}\to Y_a$ a.s. by right-continuity; hence $E[Y_{a_n}\ind_G]\to E[Y_a\ind_G]$ and $E[Y_b\ind_G]=E[Y_a\ind_G]$. Since $Y_a$ is $\G_{a+}$-measurable and integrable, $E[Y_b\mid\G_{a+}]=Y_a$ a.s.
\end{proof}

Applying the claim to $Y=X$ (Step 3) and to $Y_b=X_b^2-\lambda b$ (Step 4), we see that $X=\{X_b,\G_{b+};\,0\le b<\infty\}$ is a right-continuous, square-integrable martingale with $X_0=0$, i.e.\ $X\in\cM_2$, and that $X^2-A$ is a martingale for $A_b\defeq\lambda b$. The process $A$ is deterministic, hence adapted, with $A_0=0$ and continuous, nondecreasing paths; it is thus an increasing process in the sense of \bk{Definition 1.4.4}, and it is natural because every continuous increasing process is natural (\bk{Remark 1.4.6}(ii)). By the uniqueness assertion of \bk{Definition 1.5.3} (uniqueness in the Doob--Meyer decomposition of $X^2$),
\[
\qv{X(\delta)}_b=\lambda b=b\sqrt{\frac{2}{\pi\delta}};\qquad 0\le b<\infty,
\]
up to indistinguishability. This is \bk{Example 1.5.4} applied to the Poisson process $N$ of Step 2.

\emph{Step 6: the law of large numbers.} This is the strong law for Poisson processes of \bk{Remark 1.3.10}, applied to $N$; for completeness we give the short argument. By Step 2 the random variables $N_k-N_{k-1}$, $k\ge1$, are independent and Poisson distributed with mean $\lambda$, so the strong law of large numbers for independent, identically distributed, integrable random variables gives $N_n/n=\frac1n\sum_{k=1}^n(N_k-N_{k-1})\to\lambda$ a.s.\ as $n\to\infty$ through the integers. On $\Omega_0$ the paths of $N$ are nondecreasing, so for $b\ge1$ and $n\defeq\lfloor b\rfloor$,
\[
\frac{n}{n+1}\cdot\frac{N_n}{n}\;\le\;\frac{N_b}{b}\;\le\;\frac{n+1}{n}\cdot\frac{N_{n+1}}{n+1}.
\]
Both bounds converge to $\lambda$ as $b\to\infty$ on the intersection of $\Omega_0$ with the event $\{N_n/n\to\lambda\}$, which has probability one. Therefore
\[
\lim_{b\to\infty}\frac1b\,N_b^\delta=\lambda=\sqrt{\frac{2}{\pi\delta}}\quad\text{a.s.}
\]
\end{kex}

\begin{kex}{Exercise}{2.20}{}
Throughout, $\eps>0$ is fixed, $\nu$ is the Poisson random measure \bk{(2.36)} on $(0,\infty)^2$, whose intensity measure is $\lambda(db\times d\ell)=db\,\mu(d\ell)$ with $\mu(d\ell)=d\ell/\sqrt{2\pi\ell^3}$, and we write $V_b$ for $V_b(\eps)$, $\mu_\eps(d\ell)\defeq \ind_{(0,\eps)}(\ell)\,\mu(d\ell)$ and
\begin{equation}\label{eq:6-2-20:m}
m_1\defeq\int_{(0,\eps)}\ell\,\mu(d\ell)=\int_0^\eps\frac{d\ell}{\sqrt{2\pi\ell}}=\sqrt{\frac{2\eps}{\pi}},
\qquad
m_2\defeq\int_{(0,\eps)}\ell^2\,\mu(d\ell)=\int_0^\eps\sqrt{\frac{\ell}{2\pi}}\,d\ell=\frac{2}{3}\sqrt{\frac{\eps^3}{2\pi}}.
\end{equation}
Thus $V_b=\textit{\L}_b^\eps-bm_1$, and \bk{(2.44)} reads $\qv{V}_b=bm_2$. The book asserts, after \bk{(2.39)}, that $\textit{\L}^{\eps}$ is a subordinator with the moment generating function \bk{(2.40)} (with $\delta=0$); since the proof is only indicated there, and since we need the moments as well, we derive everything from \bk{Definition 2.3}.

\emph{Step 1: Approximation of the increments.} For $0\le a<b<\infty$ we set
\begin{equation*}
\Lambda(a,b]\defeq\int_{(0,\eps)}\ell\,\nu((a,b]\times d\ell)\in[0,\infty],
\qquad\text{so that}\qquad
\textit{\L}^\eps_b=\Lambda(0,b],\quad \textit{\L}^\eps_b=\textit{\L}^\eps_a+\Lambda(a,b],
\end{equation*}
the last identity because $\nu_\omega$ is a measure for every $\omega$. For $n\ge1$ and $\ell>0$ let $\varphi_n(\ell)\defeq 2^{-n}(\lceil 2^n\ell\rceil-1)$ be the largest element of $2^{-n}\Z$ which is strictly smaller than $\ell$. Since $2^{-n}\Z\subseteq 2^{-(n+1)}\Z$, we have $0\le\varphi_n(\ell)\le\varphi_{n+1}(\ell)<\ell$, and $\ell-\varphi_n(\ell)\le 2^{-n}$; hence $\varphi_n(\ell)\uparrow\ell$. The function $\varphi_n$ equals $k2^{-n}$ on $(k2^{-n},(k+1)2^{-n}]$; with $J^n_k\defeq (k2^{-n},(k+1)2^{-n}]\cap(0,\eps)$ and $K_n\defeq\lceil 2^n\eps\rceil$ (so that $J_k^n=\emptyset$ for $k\ge K_n$), we define the random variables
\begin{equation}\label{eq:6-2-20:Xn}
X_n(a,b]\defeq\int_{(0,\eps)}\varphi_n(\ell)\,\nu((a,b]\times d\ell)=\sum_{k=1}^{K_n}k2^{-n}\,\nu\bigl((a,b]\times J^n_k\bigr).
\end{equation}
By the monotone convergence theorem, applied for each $\omega$ to the measure $\nu_\omega((a,b]\times\cdot\,)$,
\begin{equation}\label{eq:6-2-20:mono}
0\le X_n(a,b]\uparrow\Lambda(a,b]\quad\text{as }n\to\infty,\quad\text{for every }\omega\in\Omega .
\end{equation}
For $k\ge1$ we have $\lambda((a,b]\times J_k^n)=(b-a)\mu(J_k^n)\le(b-a)\mu([2^{-n},\infty))<\infty$. Therefore, by \bk{Definition 2.3}, the random variables $\nu((a,b]\times J^n_k)$; $1\le k\le K_n$, are independent and Poisson with expectations $(b-a)\mu(J_k^n)$ (a Poisson variable with expectation zero being a.s.\ zero). A Poisson random variable has variance equal to its expectation, so from \eqref{eq:6-2-20:Xn}
\begin{equation*}
E X_n(a,b]=(b-a)\sum_{k=1}^{K_n}k2^{-n}\mu(J^n_k)=(b-a)\int\varphi_n\,d\mu_\eps,
\qquad
\Var\bigl(X_n(a,b]\bigr)=(b-a)\int\varphi_n^2\,d\mu_\eps .
\end{equation*}
Letting $n\to\infty$ and using monotone convergence on both sides ($X_n(a,b]\uparrow\Lambda(a,b]$, $X_n(a,b]^2\uparrow\Lambda(a,b]^2$, $\varphi_n\uparrow\ell$, $\varphi_n^2\uparrow\ell^2$) in $EX_n=(b-a)\int\varphi_n\,d\mu_\eps$ and in $E X_n^2=\Var(X_n)+(EX_n)^2$, we obtain from \eqref{eq:6-2-20:m}
\begin{equation}\label{eq:6-2-20:mom}
E\Lambda(a,b]=(b-a)m_1,\qquad E\bigl(\Lambda(a,b]^2\bigr)=(b-a)m_2+(b-a)^2m_1^2<\infty,
\qquad \Var\bigl(\Lambda(a,b]\bigr)=(b-a)m_2 .
\end{equation}
In particular $\Lambda(a,b]<\infty$ a.s. Two further consequences of \eqref{eq:6-2-20:Xn}, \eqref{eq:6-2-20:mono} are the following.

\textbf{(a)} \emph{The law of $\Lambda(a,b]$ depends on $(a,b)$ only through $b-a$.} Indeed, by \eqref{eq:6-2-20:Xn} the law of $X_n(a,b]$ is that of $\sum_{k=1}^{K_n}k2^{-n}\xi_k$ with $\xi_1,\dots,\xi_{K_n}$ independent Poisson variables of expectations $(b-a)\mu(J_k^n)$, so it depends only on $b-a$; and for $f$ bounded and continuous on $\R$, $Ef(\Lambda(a,b])=\lim_{n}Ef(X_n(a,b])$ by \eqref{eq:6-2-20:mono} and bounded convergence.

\textbf{(b)} \emph{For $0=b_0<b_1<\dots<b_r$, the random variables $\Lambda(b_{i-1},b_i]$; $1\le i\le r$, are independent.} Indeed, for fixed $n$ the sets $(b_{i-1},b_i]\times J^n_k$; $1\le i\le r$, $1\le k\le K_n$, are pairwise disjoint, so by \bk{Definition 2.3}(ii) the variables $\nu((b_{i-1},b_i]\times J^n_k)$ are independent, and so are the variables $X_n(b_{i-1},b_i]$; $1\le i\le r$, which are functions of disjoint subfamilies of them. For bounded continuous $f_1,\dots,f_r$ we let $n\to\infty$ in $E\prod_i f_i(X_n(b_{i-1},b_i])=\prod_iEf_i(X_n(b_{i-1},b_i])$, using \eqref{eq:6-2-20:mono} and bounded convergence, and obtain $E\prod_i f_i(\Lambda(b_{i-1},b_i])=\prod_iEf_i(\Lambda(b_{i-1},b_i])$, which is the asserted independence.

\emph{Step 2: The paths.} Let $\Omega_0$ be the event on which $\textit{\L}^\eps_b<\infty$ and $N_b^{1/k}<\infty$ for all positive integers $b$ and $k$. We have $P(\Omega_0)=1$, by \eqref{eq:6-2-20:mom} and because $EN_b^{\delta}=b\sqrt{2/(\pi\delta)}<\infty$ by \bk{(2.37)}. Since $\textit{\L}^\eps_b$ and $N^\delta_b=\nu((0,b]\times[\delta,\infty))$ are nondecreasing in $b$, and $N_b^\delta$ is nonincreasing in $\delta$, on $\Omega_0$ we have $\textit{\L}^\eps_b<\infty$ and $N^\delta_b<\infty$ for all real $b\ge0$ and $\delta>0$. For $\omega\in\Omega_0$ the path $b\mapsto \textit{\L}^\eps_b(\omega)$ vanishes at $b=0$, is nondecreasing, and is right-continuous: if $b_j\downarrow b$, then
\begin{equation*}
\textit{\L}^\eps_{b_j}-\textit{\L}^\eps_b=\int_{(0,\infty)^2}\ell\,\ind_{(b,b_j]\times(0,\eps)}(s,\ell)\,\nu(ds\times d\ell)\longrightarrow0
\end{equation*}
by dominated convergence, the integrands being dominated by the $\nu_\omega$-integrable function $\ell\,\ind_{(0,b_1]\times(0,\eps)}$ and converging pointwise to zero. Replacing $\nu_\omega$ by the zero measure for $\omega$ in the null event $\Omega\setminus\Omega_0\in\F$ yields again a Poisson random measure with intensity $\lambda$, and alters neither the almost sure statements \bk{(2.45)}, \bk{(2.46)} nor any of the distributional statements below. We do so; then \emph{every} path of $\textit{\L}^\eps$, and of $V$, is real-valued and right-continuous, and $V_0=0$.

\emph{Step 3: Stationary, independent increments; the martingale property.} For $0\le a<b$ we have $V_b-V_a=\Lambda(a,b]-(b-a)m_1$. By (a) and (b) of Step~1, $V$ has stationary and independent increments, and by \eqref{eq:6-2-20:mom}
\begin{equation}\label{eq:6-2-20:inc}
E(V_b-V_a)=0,\qquad E(V_b-V_a)^2=(b-a)m_2;\qquad\text{in particular } EV_b^2=bm_2<\infty .
\end{equation}
Fix $0\le a<b$. For $0\le c_1<\dots<c_r=a$, the vector $(V_{c_1},\dots,V_{c_r})$ is a function of the increments of $V$ over $[0,c_1],[c_1,c_2],\dots,[c_{r-1},c_r]$, of which $V_b-V_a$ is independent. The events $\{(V_{c_1},\dots,V_{c_r})\in\Gamma\}$, $\Gamma\in\B(\R^r)$, with $r$ and $c_1<\dots<c_r=a$ arbitrary, form a collection which is closed under finite intersections and generates $\F^V_a$; by Dynkin's $\pi$-$\lambda$ theorem, $V_b-V_a$ is independent of $\F^V_a$. With \eqref{eq:6-2-20:inc} this gives
\begin{align*}
E[V_b-V_a\mid\F^V_a]&=E(V_b-V_a)=0,\\
E[V_b^2-V_a^2\mid\F^V_a]&=E[(V_b-V_a)^2\mid\F^V_a]+2V_aE[V_b-V_a\mid\F^V_a]=(b-a)m_2,
\end{align*}
a.s.\ (the product $V_a(V_b-V_a)$ is integrable by the Cauchy--Schwarz inequality). Hence $\{V_b,\F^V_b\}$ is a right-continuous martingale with $V_0=0$ and $EV_b^2<\infty$, i.e., a square-integrable martingale, and
\begin{equation}\label{eq:6-2-20:V2}
\{V_b^2-bm_2,\ \F^V_b;\ 0\le b<\infty\}\quad\text{is a martingale.}
\end{equation}

\emph{Step 4: The quadratic variation \bk{(2.44)}.} \bk{Definition 1.5.3} presupposes a filtration satisfying the usual conditions, which $\{\F^V_b\}$ need not do. Let $\cN$ be the collection of $P$-null sets of (the completion of) $\F$, and $\wt{\F}_b\defeq\bigcap_{c>b}\sigma(\F^V_c\cup\cN)$; the filtration $\{\wt\F_b\}$ satisfies the usual conditions and contains $\{\F^V_b\}$. We claim that $V_b-V_a$ is independent of $\wt\F_a$ for $0\le a<b$. For $a<c<b$, the increment $V_b-V_c$ is independent of $\F^V_c$, hence of $\sigma(\F^V_c\cup\cN)$ (every set in this $\sigma$-field differs from a set in $\F^V_c$ by a null set), hence of $\wt\F_a$. Thus, for $A\in\wt\F_a$ and $f$ bounded and continuous, $E[f(V_b-V_c)\ind_A]=Ef(V_b-V_c)\,P(A)$; letting $c\downarrow a$ and using the right-continuity of $V$ and bounded convergence, we get $E[f(V_b-V_a)\ind_A]=Ef(V_b-V_a)\,P(A)$, which proves the claim. The computation of Step~3 may now be repeated with $\wt\F_a$ in place of $\F^V_a$: $\{V_b,\wt\F_b\}$ is a right-continuous martingale with $V_0=0$, $EV_b^2<\infty$, i.e., $V\in\cM_2$, and $\{V_b^2-bm_2,\wt\F_b\}$ is a martingale. The deterministic process $A_b\defeq bm_2$ is adapted, increasing in the sense of \bk{Definition 1.4.4}, continuous, and hence natural (\bk{Remark 1.4.6}(ii)). By the uniqueness in \bk{Definition 1.5.3},
\begin{equation*}
\qv{V(\eps)}_b=bm_2=\frac{2b}{3}\sqrt{\frac{\eps^3}{2\pi}};\qquad 0\le b<\infty,
\end{equation*}
which is \bk{(2.44)}. (With respect to $\{\F^V_b\}$ itself, \eqref{eq:6-2-20:V2} is the corresponding statement.)

\emph{Step 5: The moment generating function \bk{(2.47)}.} Let $\alpha\ge0$ and $b>0$. From \eqref{eq:6-2-20:Xn}, the independence of the Poisson variables $\nu((0,b]\times J^n_k)$, and the formula $Ee^{-\theta\xi}=\exp[-\kappa(1-e^{-\theta})]$ for a Poisson variable $\xi$ with expectation $\kappa$,
\begin{equation*}
Ee^{-\alpha X_n(0,b]}=\prod_{k=1}^{K_n}\exp\bigl[-b\mu(J^n_k)(1-e^{-\alpha k2^{-n}})\bigr]
=\exp\Bigl[-b\int(1-e^{-\alpha\varphi_n})\,d\mu_\eps\Bigr].
\end{equation*}
As $n\to\infty$, the left-hand side converges to $Ee^{-\alpha \textit{\L}^\eps_b}$ by \eqref{eq:6-2-20:mono} and bounded convergence, and the integral on the right increases to $\int(1-e^{-\alpha\ell})\,\mu_\eps(d\ell)\le\alpha m_1<\infty$ by monotone convergence. This is \bk{(2.40)} with $\delta=0$; multiplying by $e^{\alpha bm_1}$ and recalling \eqref{eq:6-2-20:m} we obtain \bk{(2.47)}:
\begin{equation*}
Ee^{-\alpha V_b}=\exp\Bigl[\alpha b\sqrt{\frac{2\eps}{\pi}}-b\int_{(0,\eps)}(1-e^{-\alpha\ell})\frac{d\ell}{\sqrt{2\pi\ell^3}}\Bigr]
=\exp\Bigl[b\int_{(0,\eps)}(e^{-\alpha\ell}-1+\alpha\ell)\,\mu(d\ell)\Bigr].
\end{equation*}

\emph{Step 6: The representation \bk{(2.45)}.} Fix $\omega\in\Omega_0$ and $b\ge0$, and let $\mathfrak{m}(A)\defeq\nu_\omega((0,b]\times A)$; $A\in\B((0,\infty))$. This is a measure with values in $\N_0\cup\{\infty\}$ and $\mathfrak{m}([\delta,\infty))=N^\delta_b(\omega)<\infty$ for every $\delta>0$; in particular, its restriction to $(0,\eps)=\bigcup_k[1/k,\eps)$ is $\sigma$-finite. Writing $\ell=\int_{(0,\eps)}\ind_{\{x<\ell\}}\,dx$ for $0<\ell<\eps$, and applying Tonelli's theorem to the nonnegative Borel function $(x,\ell)\mapsto\ind_{\{x<\ell\}}$ on $(0,\eps)^2$, we find
\begin{equation}\label{eq:6-2-20:ton}
\textit{\L}^\eps_b=\int_{(0,\eps)}\ell\,\mathfrak{m}(d\ell)=\int_{(0,\eps)}\int_{(0,\eps)}\ind_{\{x<\ell\}}\,dx\,\mathfrak{m}(d\ell)=\int_{(0,\eps)}\mathfrak{m}\bigl((x,\eps)\bigr)\,dx .
\end{equation}
For $0<x<\eps$ all the quantities in
\begin{equation*}
\mathfrak{m}\bigl((x,\eps)\bigr)=\mathfrak{m}\bigl([x,\infty)\bigr)-\mathfrak{m}\bigl([\eps,\infty)\bigr)-\mathfrak{m}(\{x\})=N^x_b-N^\eps_b-\mathfrak{m}(\{x\})
\end{equation*}
are finite. Moreover, $\mathfrak{m}(\{x\})\ne0$ only for $x$ in a countable set: an atom of the integer-valued measure $\mathfrak{m}$ has mass at least one, so $[1/k,\infty)$ contains at most $N^{1/k}_b<\infty$ atoms, for every $k$. A countable set is Lebesgue-null, and $x\mapsto N^x_b$ is nonincreasing, hence Borel measurable; thus \eqref{eq:6-2-20:ton} becomes
\begin{equation*}
\textit{\L}^\eps_b=\int_{(0,\eps)}\bigl(N^x_b-N^\eps_b\bigr)\,dx .
\end{equation*}
The integrand is nonnegative, the left-hand side is finite and so is $\eps N_b^\eps$; hence $\int_{(0,\eps)}N^x_b\,dx<\infty$ and $\textit{\L}^\eps_b=-\eps N^\eps_b+\int_{(0,\eps)}N^\ell_b\,d\ell$. This holds for every $b\ge0$ and every $\omega\in\Omega_0$, i.e., \bk{(2.45)} holds simultaneously for all $0\le b<\infty$, almost surely.

\emph{Step 7: The law of large numbers \bk{(2.46)}.} By Step~1, the random variables $\Lambda(k-1,k]=\textit{\L}^\eps_k-\textit{\L}^\eps_{k-1}$; $k\ge1$, are independent and identically distributed, with expectation $m_1$. Kolmogorov's strong law of large numbers gives $\lim_{n\to\infty}\textit{\L}^\eps_n/n=m_1$ a.s., along the integers. For real $b\ge1$ let $n=\lfloor b\rfloor$; since $\textit{\L}^\eps$ is nondecreasing,
\begin{equation*}
\frac{n}{n+1}\cdot\frac{\textit{\L}^\eps_n}{n}\le\frac{\textit{\L}^\eps_b}{b}\le\frac{n+1}{n}\cdot\frac{\textit{\L}^\eps_{n+1}}{n+1},
\end{equation*}
and both bounds converge to $m_1$ as $b\to\infty$ on the event where $\textit{\L}^\eps_n/n\to m_1$. Therefore $\lim_{b\to\infty}\textit{\L}^\eps_b/b=m_1=\sqrt{2\eps/\pi}$ a.s., which is \bk{(2.46)}.

\rmk{The moments \eqref{eq:6-2-20:inc} can also be read off from \bk{(2.47)}, as the hint suggests: by $0\le e^{-u}-1+u\le u^2/2$ for $u\ge0$ and dominated convergence, the exponent $b\int_{(0,\eps)}(e^{-\alpha\ell}-1+\alpha\ell)\,\mu(d\ell)$ equals $\tfrac12\alpha^2bm_2+o(\alpha^2)$ as $\alpha\downarrow0$, in agreement with $EV_b=0$ and $EV_b^2=bm_2$. A simulation of the jumps of size in $[10^{-5},\eps)$, with $\eps=\tfrac12$, $b=2$ and $2\cdot10^4$ paths, gave mean $1.130$ and variance $0.191$ for $\textit{\L}^\eps_b$, against $bm_1=1.128$ and $bm_2=0.188$, and reproduced \bk{(2.47)} for $\alpha=0.5,\,2,\,5$ within the sampling error.}
\end{kex}

\solsection{6.3}{Two Independent Reflected Brownian Motions}

\begin{kex}{Problem}{3.4}{}
We fix $t>0$ and write
\[
  g(a,b)\defeq \frac{2(a+b)}{\sqrt{2\pi t^3}}\exp\Bigl\{-\frac{(a+b)^2}{2t}\Bigr\};
  \qquad a\ge 0,\ b\ge 0 .
\]
We denote by $\mu$ the law of $(W_t,2L_t(0))$ and by $\nu$ the law of $(\abs{W_t},2L_t(0))$ under
$P^0$; these are probability measures on $\R\times[0,\infty)$ and on $[0,\infty)^2$, respectively.
The assertions to be proved are that $\nu$ has the density $g$ on $(0,\infty)^2$, and $\mu$ the
density $\tfrac12 g(\abs{a},b)$ on $\R\times(0,\infty)$, with respect to Lebesgue measure; we shall
see that these open sets carry the whole mass, so that the densities describe $\nu$ and $\mu$
completely.

\emph{Step 1: proof of \bk{(3.18)}.} Let $M_t=\max_{0\le s\le t}W_s$ be the running maximum
\bk{(2.8.1)}, and let $\lambda$ be the law of $(W_t,M_t)$ under $P^0$. Since $M_t\ge W_t$ and
$M_t\ge W_0=0$, a.s.\ $P^0$, the measure $\lambda$ is concentrated on the set
$D\defeq\{(x,y)\in\R^2;\ x\le y,\ y\ge 0\}$, and \bk{Proposition 2.8.1} asserts that on $D$ it has
the density
\[
  f(x,y)=\frac{2(2y-x)}{\sqrt{2\pi t^3}}\exp\Bigl\{-\frac{(2y-x)^2}{2t}\Bigr\};
  \qquad (x,y)\in D,
\]
that is, $\lambda(G)=\int_{G\cap D} f(x,y)\,dx\,dy$ for every $G\in\B(\R^2)$.

By \bk{Theorem 3.6.17}, the processes $\{(M_s-W_s,M_s);\,0\le s<\infty\}$ and
$\{(\abs{W_s},2L_s(0));\,0\le s<\infty\}$ have the same law under $P^0$; in particular, their values
at the time $t$ have the same distribution. Thus $\nu$ is the image of $\lambda$ under the linear
bijection $T(x,y)\defeq(y-x,y)$ of $\R^2$, which has determinant of absolute value $1$, inverse
$T^{-1}(a,b)=(b-a,b)$, and maps $D$ onto $[0,\infty)^2$. For $H\in\B([0,\infty)^2)$, the change of
variables $(x,y)=T^{-1}(a,b)$ gives
\[
  \nu(H)=\lambda\bigl(T^{-1}H\bigr)=\int_{T^{-1}H} f(x,y)\,dx\,dy
        =\int_H f(b-a,b)\,da\,db=\int_H g(a,b)\,da\,db,
\]
because $2b-(b-a)=a+b$. Hence $\nu$ is absolutely continuous with respect to Lebesgue measure on
$[0,\infty)^2$, with density $g$. The boundary of the quadrant is a Lebesgue null set, so
$\nu\bigl((0,\infty)^2\bigr)=1$ and $g$, restricted to $(0,\infty)^2$, is the density of $\nu$. This is
\bk{(3.18)}.

\emph{Step 2: symmetry.} We show that $\mu$ is invariant under the reflection
$\rho(a,b)\defeq(-a,b)$. Let $\Omega^*$ be the event of \bk{Definition 3.6.3}; it has
$P^0$-probability one. For $\omega\in\Omega^*$ and $\eps>0$, the occupation-density identity
\bk{(3.6.6)} with $B=[-\eps,\eps]$ gives
\[
  \frac1{2\eps}\,\meas\{0\le s\le t;\ \abs{W_s(\omega)}\le\eps\}
  =\frac1{2\eps}\int_{-\eps}^{\eps}2L_t(x,\omega)\,dx ,
\]
and the right-hand side converges to $2L_t(0,\omega)$ as $\eps\downarrow 0$, by the continuity of
$x\mapsto L_t(x,\omega)$. Consequently,
\begin{equation}\label{eq:6-3-04:phi}
  2L_t(0)=\Phi(W)\quad\text{a.s. }P^0,
\end{equation}
where
\[
  \Phi(w)\defeq\liminf_{n\to\infty}\frac n2\,\meas\Bigl\{0\le s\le t;\ \abs{w(s)}\le\frac1n\Bigr\};
  \qquad w\in C[0,\infty).
\]
The mapping $(s,w)\mapsto w(s)$ is continuous on $[0,t]\times C[0,\infty)$, hence jointly
measurable, and Tonelli's theorem shows that
$w\mapsto\int_0^t\ind_{\{\abs{w(s)}\le 1/n\}}\,ds$ is $\B(C[0,\infty))$-measurable for every $n$;
therefore $\Phi\colon C[0,\infty)\to[0,\infty]$ is measurable. Obviously $\Phi(-w)=\Phi(w)$. Now
$-W$ is, like $W$, a standard Brownian motion under $P^0$, so $W$ and $-W$ induce the same measure
(Wiener measure) on $\bigl(C[0,\infty),\B(C[0,\infty))\bigr)$. Applying to both the measurable
mapping $w\mapsto(w(t),\Phi(w))$, we see that $(W_t,\Phi(W))$ and
$(-W_t,\Phi(-W))=(-W_t,\Phi(W))$ have the same law under $P^0$, and by
\eqref{eq:6-3-04:phi} this means
\begin{equation}\label{eq:6-3-04:sym}
  \mu(G)=\mu(\rho G);\qquad G\in\B(\R\times[0,\infty)).
\end{equation}

\emph{Step 3: proof of \bk{(3.17)}.} Let $Q_+\defeq(0,\infty)^2$ and $Q_-\defeq\rho Q_+
=(-\infty,0)\times(0,\infty)$. For a Borel set $H\subseteq Q_+$ we have
$\{(\abs{W_t},2L_t(0))\in H\}=\{(W_t,2L_t(0))\in H\}\cup\{(W_t,2L_t(0))\in\rho H\}$, a disjoint
union, whence by \eqref{eq:6-3-04:sym} and Step 1
\begin{equation}\label{eq:6-3-04:half}
  \mu(H)=\mu(\rho H)=\frac12\,\nu(H)=\int_H\frac12\,g(a,b)\,da\,db;\qquad H\in\B(Q_+).
\end{equation}
In particular $\mu(Q_+)=\mu(Q_-)=\tfrac12\nu(Q_+)=\tfrac12$, so $\mu$ is concentrated on
$Q_+\cup Q_-$ (in agreement with $P^0[W_t=0]=0$). Let now $G\in\B(\R\times(0,\infty))$.
Then $\mu(G)=\mu(G\cap Q_+)+\mu(G\cap Q_-)$, and $\mu(G\cap Q_-)=\mu(\rho(G\cap Q_-))$ by
\eqref{eq:6-3-04:sym}. Using
\eqref{eq:6-3-04:half} for $G\cap Q_+$ and for $\rho(G\cap Q_-)\subseteq Q_+$, and then the
substitution $a\mapsto -a$ in the second integral, we obtain
\begin{align*}
  \mu(G)&=\mu(G\cap Q_+)+\mu\bigl(\rho(G\cap Q_-)\bigr)\\
  &=\int_{G\cap Q_+}\frac12\,g(a,b)\,da\,db+\int_{G\cap Q_-}\frac12\,g(-a,b)\,da\,db\\
  &=\int_G\frac12\,g(\abs{a},b)\,da\,db,
\end{align*}
the line $\{a=0\}$ being Lebesgue null. Thus $\mu$ has on $\R\times(0,\infty)$ the density
\[
  \frac12\,g(\abs{a},b)=\frac{b+\abs{a}}{\sqrt{2\pi t^3}}\exp\Bigl\{-\frac{(b+\abs{a})^2}{2t}\Bigr\},
\]
which is \bk{(3.17)}.

\rmk{As a check, the marginals come out right. With $c=b+\abs{a}$ and
$\int_c^\infty u\,e^{-u^2/2t}\,du=t\,e^{-c^2/2t}$, integration of \bk{(3.17)} over $b\in(0,\infty)$
gives $(2\pi t)^{-1/2}e^{-a^2/2t}$, the density of $W_t$, and integration of \bk{(3.18)} over
$a\in(0,\infty)$ gives $2(2\pi t)^{-1/2}e^{-b^2/2t}$, $b>0$, the density of $M_t$ \bk{(2.8.3)} and
therefore of $2L_t(0)$ \bk{(3.6.28)}. The same computation, done twice, yields the tail formula
$P^0[W_t>a;\ 2L_t(0)>b]=P^0[W_t>a+b]$ for $a\ge0$, $b\ge0$.}
\end{kex}

\begin{kex}{Exercise}{3.5}{}
Fix $t>0$ and $\Gamma\in\B(\R)$. In \bk{Exercise 5.3.12} the process $X$ is, under $P^0$, a standard
Brownian motion starting at the origin and $L_t(0)$ is its local time at the origin, so the right-hand
side of \bk{(5.3.13)} is
\begin{equation}\label{eq:6-3-05:rhs}
R_t(\Gamma)\defeq e^{-t/2}\,E^0\bigl[\ind_{\{X_t\in\Gamma\}}\exp\bigl(-\abs{X_t}+2L_t(0)\bigr)\bigr]
 = e^{-t/2}\,E^0\bigl[g\bigl(X_t,2L_t(0)\bigr)\bigr]
\end{equation}
with $g(a,b)\defeq \ind_\Gamma(a)\,e^{-\abs{a}+b}$,
an expectation of a nonnegative Borel function of the pair $(X_t,2L_t(0))$; it is thus well defined in
$[0,\infty]$.

\emph{Step 1: the law of $(X_t,2L_t(0))$.} By \bk{Problem 3.4}, formula \bk{(3.17)}, the function
\[
f_t(a,b)\defeq \frac{b+\abs{a}}{\sqrt{2\pi t^3}}\exp\Bigl\{-\frac{(b+\abs{a})^2}{2t}\Bigr\},
\qquad a\in\R,\ b>0,
\]
is the density of the $P^0$-distribution of $(X_t,2L_t(0))$ on $\R\times(0,\infty)$. This distribution
has no mass outside $\R\times(0,\infty)$, since $f_t$ already integrates to one there: with the
substitution $v=b+\abs{a}$,
\[
\int_\R\int_0^\infty f_t(a,b)\,db\,da
 =\int_\R\frac{1}{\sqrt{2\pi t^3}}\int_{\abs{a}}^\infty v\,e^{-v^2/2t}\,dv\,da
 =\int_\R\frac{1}{\sqrt{2\pi t}}\,e^{-a^2/2t}\,da=1 .
\]
Consequently $E^0[h(X_t,2L_t(0))]=\int_\R\int_0^\infty h(a,b)f_t(a,b)\,db\,da$ for every Borel
function $h:\R\times[0,\infty)\to[0,\infty]$, the order of integration being immaterial by Tonelli's
theorem. Applied to $h=g$ in \eqref{eq:6-3-05:rhs}, this gives
\begin{equation}\label{eq:6-3-05:q}
R_t(\Gamma)=\int_\Gamma q_t(a)\,da,\qquad
q_t(a)\defeq e^{-t/2}\,e^{-\abs{a}}\int_0^\infty e^{b}\,
\frac{b+\abs{a}}{\sqrt{2\pi t^3}}\exp\Bigl\{-\frac{(b+\abs{a})^2}{2t}\Bigr\}\,db ,
\end{equation}
and it remains to evaluate the inner integral.

\emph{Step 2: computation of $q_t$.} Fix $a\in\R$ and substitute $v=b+\abs{a}$ in
\eqref{eq:6-3-05:q}; then $e^{-\abs{a}}e^{b}=e^{-2\abs{a}}e^{v}$, and since
\[
v-\frac{v^2}{2t}-\frac{t}{2}=-\frac{(v-t)^2}{2t},
\]
we obtain
\[
q_t(a)=\frac{e^{-2\abs{a}}}{\sqrt{2\pi t^3}}\int_{\abs{a}}^\infty v\,
\exp\Bigl\{-\frac{(v-t)^2}{2t}\Bigr\}\,dv .
\]
We write $v=(v-t)+t$. The first part is integrated explicitly,
\[
\int_{\abs{a}}^\infty (v-t)\exp\Bigl\{-\frac{(v-t)^2}{2t}\Bigr\}\,dv
=\Bigl[-t\exp\Bigl\{-\frac{(v-t)^2}{2t}\Bigr\}\Bigr]_{v=\abs{a}}^{v=\infty}
=t\exp\Bigl\{-\frac{(\abs{a}-t)^2}{2t}\Bigr\},
\]
so that
\[
q_t(a)=\frac{1}{\sqrt{2\pi t}}\biggl[e^{-2\abs{a}}\exp\Bigl\{-\frac{(\abs{a}-t)^2}{2t}\Bigr\}
+e^{-2\abs{a}}\int_{\abs{a}}^\infty\exp\Bigl\{-\frac{(v-t)^2}{2t}\Bigr\}\,dv\biggr].
\]
Finally $-2\abs{a}-\dfrac{(\abs{a}-t)^2}{2t}=-\dfrac{(\abs{a}+t)^2}{2t}$, and the last display is exactly
\bk{(3.19)}. Together with \eqref{eq:6-3-05:q} this proves that the right-hand side of \bk{(5.3.13)}
equals $\int_\Gamma q_t(a)\,da$.

\rmk{As a check, $q_t$ is a probability density, as it must be in view of \bk{(5.3.13)}. Indeed, by
Tonelli's theorem,
\[
\int_\R e^{-2\abs{a}}\int_{\abs{a}}^\infty e^{-(v-t)^2/2t}\,dv\,da
=\int_0^\infty\bigl(1-e^{-2v}\bigr)e^{-(v-t)^2/2t}\,dv
=\int_0^\infty e^{-(v-t)^2/2t}\,dv-\int_0^\infty e^{-(v+t)^2/2t}\,dv ,
\]
while $\int_\R e^{-(\abs{a}+t)^2/2t}\,da=2\int_0^\infty e^{-(v+t)^2/2t}\,dv$; hence
\[
\sqrt{2\pi t}\int_\R q_t(a)\,da=\int_0^\infty e^{-(v-t)^2/2t}\,dv+\int_0^\infty e^{-(v+t)^2/2t}\,dv
=\int_\R e^{-(v-t)^2/2t}\,dv=\sqrt{2\pi t}.
\]
Moreover, for fixed $a\in\R$ we can let $t\to\infty$ in \bk{(3.19)}. The first term satisfies
$0\le(2\pi t)^{-1/2}e^{-(\abs{a}+t)^2/2t}\le(2\pi t)^{-1/2}e^{-t/2}\to0$, and the substitution
$v=t+z\sqrt{t}$ gives
\[
\frac{1}{\sqrt{2\pi t}}\int_{\abs{a}}^\infty e^{-(v-t)^2/2t}\,dv
=\int_{(\abs{a}-t)/\sqrt{t}}^\infty\frac{1}{\sqrt{2\pi}}\,e^{-z^2/2}\,dz\ \longrightarrow\ 1 ,
\]
because the lower limit $(\abs{a}-t)/\sqrt{t}$ tends to $-\infty$. Hence
$\lim_{t\to\infty}q_t(a)=e^{-2\abs{a}}$ for every $a\in\R$, and this limit is again a probability
density on $\R$, since $\int_\R e^{-2\abs{a}}\,da=2\int_0^\infty e^{-2a}\,da=1$.}
\end{kex}

\begin{kex}{Exercise}{3.8}{}
We have to show that, for every $t>0$,
\begin{equation}\label{eq:6-3-08:goal}
P^0[\Gamma_+(t)\le\theta]=\int_0^{\theta/t}\frac{ds}{\pi\sqrt{s(1-s)}}
=\frac{2}{\pi}\arcsin\sqrt{\frac{\theta}{t}};\qquad 0\le\theta\le t,
\end{equation}
where $\Gamma_+(t)=\meas\{0\le s\le t;\ W_s>0\}$ as in \bk{(3.3)}. The second equality in \eqref{eq:6-3-08:goal} is the substitution $s=\sin^2\varphi$, $0\le\varphi\le\arcsin\sqrt{\theta/t}$, under which $ds=2\sin\varphi\cos\varphi\,d\varphi$ and $\sqrt{s(1-s)}=\sin\varphi\cos\varphi$; so only the equality of the two outer members has to be proved. Throughout, $\tau>0$ is fixed.

\emph{Step 1: the law of $L_+(\tau)$.} By \bk{(3.14)} we have $L_+(\tau)=\max_{0\le u\le\tau}B_+(u)$, a.s.\ $P^0$, where $B_+$ of \bk{(3.10)} is a standard, one-dimensional Brownian motion under $P^0$. Hence \bk{(2.8.3)} gives
\begin{equation}\label{eq:6-3-08:lawL}
P^0[L_+(\tau)\in db]=\sqrt{\frac{2}{\pi\tau}}\,e^{-b^2/2\tau}\,db;\qquad b>0.
\end{equation}

\emph{Step 2: the Laplace transform of $\Gamma_+^{-1}(\tau)-\tau$.} The function $\Gamma_+$ is continuous and nondecreasing with $\Gamma_+(0)=0$ and $0\le\Gamma_+(t)-\Gamma_+(s)\le t-s$ for $s\le t$, and $\lim_{t\to\infty}\Gamma_+(t)=\infty$ a.s.\ $P^0$ by \bk{(3.6)}. Consequently the random variable $\Gamma_+^{-1}(\tau)=\inf\{t\ge0;\ \Gamma_+(t)>\tau\}$ of \bk{(3.9)} is a.s.\ finite, and it satisfies $\Gamma_+^{-1}(\tau)\ge\tau$ because $\Gamma_+(t)\le t$. Thus
\[
\xi\defeq\Gamma_+^{-1}(\tau)-\tau
\]
is a random variable with values in $[0,\infty)$, a.s.\ $P^0$. Fix $\alpha>0$. Taking expectations in the formula \bk{(3.22)} of \bk{Theorem 3.6} and using \eqref{eq:6-3-08:lawL}, we obtain
\begin{equation}\label{eq:6-3-08:lap1}
E^0e^{-\alpha\xi}=E^0e^{-\sqrt{2\alpha}\,L_+(\tau)}
=\int_0^\infty e^{-b\sqrt{2\alpha}}\sqrt{\frac{2}{\pi\tau}}\,e^{-b^2/2\tau}\,db.
\end{equation}
According to \bk{(2.8.5)}, \bk{(2.8.6)}, for every $b>0$,
\[
e^{-b\sqrt{2\alpha}}=\int_0^\infty e^{-\alpha s}\,\frac{b}{\sqrt{2\pi s^3}}\,e^{-b^2/2s}\,ds .
\]
We substitute this into \eqref{eq:6-3-08:lap1} and interchange the order of integration, which Tonelli's theorem allows since the integrand is nonnegative and jointly continuous in $(b,s)\in(0,\infty)^2$:
\begin{equation}\label{eq:6-3-08:lap2}
E^0e^{-\alpha\xi}=\int_0^\infty e^{-\alpha s}g(s)\,ds,\qquad
g(s)\defeq\frac{1}{\pi\sqrt{s^3\tau}}\int_0^\infty b\,
\exp\Bigl\{-\frac{b^2}{2}\Bigl(\frac1s+\frac1\tau\Bigr)\Bigr\}\,db .
\end{equation}
Since $\int_0^\infty b\,e^{-cb^2/2}\,db=1/c$ for $c>0$, and here $c=(s+\tau)/s\tau$, we find
\begin{equation}\label{eq:6-3-08:g}
g(s)=\frac{1}{\pi\sqrt{s^3\tau}}\cdot\frac{s\tau}{s+\tau}=\frac{\sqrt{\tau}}{\pi\sqrt{s}\,(s+\tau)};\qquad s>0.
\end{equation}
With the substitution $s=\tau u^2$ we have, for every $x\ge0$,
\begin{equation}\label{eq:6-3-08:G}
\int_0^x g(s)\,ds=\frac{2}{\pi}\int_0^{\sqrt{x/\tau}}\frac{du}{1+u^2}=\frac{2}{\pi}\arctan\sqrt{\frac{x}{\tau}},
\end{equation}
and letting $x\to\infty$ we see that $g$ is a probability density on $(0,\infty)$. Relation \eqref{eq:6-3-08:lap2} holds for every $\alpha>0$ and says that the distribution of $\xi$ and the probability measure $g(s)\,ds$, both carried by $[0,\infty)$, have the same Laplace transform; by the uniqueness theorem for Laplace transforms of finite measures on $[0,\infty)$ they coincide. In particular the distribution of $\xi$ has no atoms, and from \eqref{eq:6-3-08:G}
\begin{equation}\label{eq:6-3-08:lawxi}
P^0\bigl[\Gamma_+^{-1}(\tau)<\tau+x\bigr]=P^0[\xi<x]=\frac{2}{\pi}\arctan\sqrt{\frac{x}{\tau}};\qquad x\ge0.
\end{equation}

\emph{Step 3: inversion.} For $t>0$ we claim that
\begin{equation}\label{eq:6-3-08:inv}
\{\Gamma_+^{-1}(\tau)<t\}=\{\Gamma_+(t)>\tau\}.
\end{equation}
Indeed, if $\Gamma_+^{-1}(\tau)<t$, then by the definition of the infimum there is $s<t$ with $\Gamma_+(s)>\tau$, and $\Gamma_+(t)\ge\Gamma_+(s)>\tau$ by monotonicity. Conversely, if $\Gamma_+(t)>\tau$, then by continuity of $\Gamma_+$ there is $s<t$ with $\Gamma_+(s)>\tau$, whence $\Gamma_+^{-1}(\tau)\le s<t$.

\emph{Step 4: conclusion.} Let $t>0$ and $0<\theta<t$. Applying \eqref{eq:6-3-08:inv} and then \eqref{eq:6-3-08:lawxi}, with $\tau=\theta$ and $x=t-\theta>0$, we get
\[
P^0[\Gamma_+(t)\le\theta]=1-P^0\bigl[\Gamma_+^{-1}(\theta)<t\bigr]
=1-\frac{2}{\pi}\arctan\sqrt{\frac{t-\theta}{\theta}}
=\frac{2}{\pi}\arctan\sqrt{\frac{\theta}{t-\theta}},
\]
the last step by the identity $\arctan y+\arctan(1/y)=\pi/2$, $y>0$. Finally, the angle $\varphi\defeq\arctan\sqrt{\theta/(t-\theta)}\in(0,\pi/2)$ satisfies $\sin^2\varphi=\tan^2\varphi/(1+\tan^2\varphi)=\theta/t$, that is, $\varphi=\arcsin\sqrt{\theta/t}$. This proves \eqref{eq:6-3-08:goal} for $0<\theta<t$. For $\theta=t$ both sides of \eqref{eq:6-3-08:goal} are equal to $1$, because $\Gamma_+(t)\le t$ and $\arcsin 1=\pi/2$. For $\theta=0$ the right-continuity of the distribution function of $\Gamma_+(t)$ yields $P^0[\Gamma_+(t)\le0]=\lim_{\theta\downarrow0}(2/\pi)\arcsin\sqrt{\theta/t}=0$, which is the value of the right-hand side. The arc-sine law \eqref{eq:6-3-08:goal} is established.

\rmk{The argument uses \bk{Theorem 3.6} only through the unconditional consequence \eqref{eq:6-3-08:lap1} of \bk{(3.22)}, and the proof of that theorem (which rests on \bk{Theorem 3.1}, \bk{Problem 2.7.19}, \bk{Problem 3.7} and \bk{(2.8.6)}; the third and fourth identities of \bk{(3.25)}, delegated to \bk{Problem 3.7}, involve only \bk{(3.9)}, \bk{(3.13)} and the monotonicity and continuity of $L$ and $\Gamma_\pm$) makes no use of \bk{Proposition 4.4.11}; the derivation is therefore independent of the Feynman--Kac argument of \bk{Section 4.4}. The density \eqref{eq:6-3-08:g} has a transparent meaning in the light of \bk{(3.23)}: $\xi=T^-_{L_+(\tau)}$ is the passage time of the Brownian motion $B_-$ to the independent random level $L_+(\tau)$, and \eqref{eq:6-3-08:lawxi} can be read as $P^0[\sqrt{x}\,\abs{Z_1}>\sqrt{\tau}\,\abs{Z_2}]$ for independent standard normal random variables $Z_1$, $Z_2$, a probability which the rotational invariance of the law of $(Z_1,Z_2)$ evaluates at once as $(2/\pi)\arctan\sqrt{x/\tau}$.}
\end{kex}

\solsection{6.4}{Elastic Brownian Motion}

\begin{kex}{Problem}{4.4}{}
Throughout, $T_b=\inf\{t\ge 0;\ W_t=b\}$, $\Gamma_+(t)=\meas\{0\le s\le t;\ W_s>0\}=\int_0^t\ind_{(0,\infty)}(W_s)\,ds$ as in \bk{(3.3)}, and $L_t(0)$ is the local time of \bk{(4.1)}, normalised as in \bk{Section 3.6}. We write
\begin{equation*}
\kappa\defeq\sqrt{2\alpha},\qquad \theta\defeq\sqrt{2(\alpha+\beta)},\qquad
\Delta\defeq\frac{\gamma+\kappa}{\theta}\,\sinh(b\theta)+\cosh(b\theta) ,
\end{equation*}
so that $\Delta>1$ is the common denominator in \bk{(4.15)}. For every $x\le b$ we have
\begin{equation}\label{eq:6-4-04:fin}
P^x[T_b<\infty]=1 ,
\end{equation}
because under $P^x$ the process $W-x$ is a standard Brownian motion, which reaches the level $b-x\ge 0$ in finite time almost surely (let $t\to\infty$ in \bk{(2.6.2)}).

\emph{Step 1: the function $f$.} Define $f\colon\R\to\R$ by
\begin{equation*}
f(y)\defeq
\begin{cases}
\dfrac1\Delta\Bigl[\dfrac{\gamma+\kappa}{\theta}\,\sinh(\theta y)+\cosh(\theta y)\Bigr]; & y\ge 0,\\[2.5ex]
\dfrac1\Delta\,e^{\kappa y}; & y<0 .
\end{cases}
\end{equation*}
For $x\le b$, $f(x)$ is the right-hand side of \bk{(4.15)}. The function $f$ is continuous on $\R$ with $f(0)=1/\Delta$, and $f'$, $f''$ exist and are continuous on $\R\setminus\{0\}$, with
\begin{align*}
f'(y)&=\frac{1}{\Delta}\bigl[(\gamma+\kappa)\cosh(\theta y)+\theta\sinh(\theta y)\bigr], & f''(y)&=\theta^2 f(y); & y&>0,\\
f'(y)&=\kappa f(y), & f''(y)&=\kappa^2 f(y); & y&<0.
\end{align*}
The one-sided limits $f'(0\pm)$, $f''(0\pm)$ exist and are finite, and
\begin{equation}\label{eq:6-4-04:jump}
f'(0+)-f'(0-)=\frac{\gamma+\kappa}{\Delta}-\frac{\kappa}{\Delta}=\gamma f(0) .
\end{equation}
Since $\theta^2=2(\alpha+\beta)$ and $\kappa^2=2\alpha$,
\begin{equation}\label{eq:6-4-04:ode}
\tfrac12 f''(y)=\bigl(\alpha+\beta\,\ind_{(0,\infty)}(y)\bigr)f(y);\qquad y\in\R\setminus\{0\}.
\end{equation}
Finally, $f'>0$ on $(-\infty,0)\cup(0,\infty)$, so the continuous function $f$ is strictly increasing on $\R$; as $f>0$ and $f(b)=1$, we obtain
\begin{equation}\label{eq:6-4-04:bd}
0<f(y)\le 1\quad\text{and}\quad 0<f'(y)\le C\defeq\frac{1}{\Delta}\bigl[(\gamma+\kappa)\cosh(\theta b)+\theta\sinh(\theta b)\bigr];
\qquad y\in(-\infty,b]\setminus\{0\}
\end{equation}
(for $y<0$ we have $f'(y)\le\kappa/\Delta\le C$). We set $f'(0)\defeq f'(0-)$ and $f''(0)\defeq 0$; this convention is immaterial, because the set $\{s\ge0;\ W_s=0\}$ has Lebesgue measure zero, a.s.\ $P^x$ (by Fubini--Tonelli, $E^x\int_0^\infty\ind_{\{0\}}(W_s)\,ds=\int_0^\infty P^x[W_s=0]\,ds=0$), and therefore neither the Lebesgue integrals nor the stochastic integrals below depend on it.

\emph{Step 2: a local martingale.} Fix $x\in\R$. By \bk{Problem 3.6.24}, applied with the single exceptional point $a_1=0$, and by \eqref{eq:6-4-04:jump}, the process $U_t\defeq f(W_t)$ satisfies
\begin{equation}\label{eq:6-4-04:ito}
U_t=f(x)+\int_0^t f'(W_s)\,dW_s+\frac12\int_0^t f''(W_s)\,ds+\gamma f(0)\,L_t(0);\qquad 0\le t<\infty,\ \text{a.s. }P^x .
\end{equation}
In particular $U$ is a continuous semimartingale under $P^x$. Let
\begin{equation*}
A_t\defeq\alpha t+\beta\,\Gamma_+(t)+\gamma L_t(0),\qquad Z_t\defeq e^{-A_t},\qquad Y_t\defeq f(W_t)\,Z_t=U_tZ_t;\qquad 0\le t<\infty .
\end{equation*}
The process $A$ is adapted, continuous and nondecreasing with $A_0=0$ (for $L(0)$, adaptedness, continuity and $L_0(0)=0$ are part of \bk{Proposition 3.6.8}, and the monotonicity of $t\mapsto L_t(0)$ is \bk{Problem 3.6.13}(ii)); hence $0<Z\le1$. It\^o's rule (\bk{Theorem 3.3.6}) for the function $(u,a)\mapsto ue^{-a}$ and the pair of semimartingales $(U,A)$, the second of which has no local martingale part, gives
\begin{equation*}
Y_t=f(x)+\int_0^t Z_s\,dU_s-\int_0^t U_sZ_s\,dA_s ,
\end{equation*}
and substituting \eqref{eq:6-4-04:ito} we obtain, a.s.\ $P^x$, for all $0\le t<\infty$,
\begin{align*}
Y_t=f(x)&+\int_0^t Z_sf'(W_s)\,dW_s
+\int_0^t Z_s\Bigl[\tfrac12 f''(W_s)-\bigl(\alpha+\beta\,\ind_{(0,\infty)}(W_s)\bigr)f(W_s)\Bigr]ds\\
&+\gamma\int_0^t Z_s\bigl[f(0)-f(W_s)\bigr]\,dL_s(0) .
\end{align*}
By \eqref{eq:6-4-04:ode} the integrand of the $ds$-integral vanishes whenever $W_s\ne0$, hence for Lebesgue-almost every $s$, and this integral is zero. The integrand of the $dL_s(0)$-integral vanishes on $\{s;\ W_s=0\}$, while the measure $dL_s(0)$ does not charge $\{s;\ W_s\ne0\}$ by \bk{(3.6.26)} of \bk{Problem 3.6.13}(ii); this integral is zero as well. Consequently
\begin{equation}\label{eq:6-4-04:Y}
Y_t=f(x)+\int_0^t Z_sf'(W_s)\,dW_s;\qquad 0\le t<\infty,\ \text{a.s. }P^x .
\end{equation}

\emph{Step 3: stopping at $T_b$.} Let now $x\le b$. For $0\le s\le T_b$ we have $W_s\le b$ by the continuity of the paths, so by \eqref{eq:6-4-04:bd} the integrand $Z_sf'(W_s)\ind_{\{s\le T_b\}}$ is bounded by $C$ (recall that $f'(0)=f'(0-)\le C$). The stopped process $\{Y_{t\wedge T_b}-f(x);\ 0\le t<\infty\}$ is thus a stochastic integral of a bounded progressively measurable integrand with respect to $W$, hence a square-integrable martingale starting at zero, and
\begin{equation*}
E^x\,Y_{t\wedge T_b}=f(x);\qquad 0\le t<\infty .
\end{equation*}
Moreover $0<Y_{t\wedge T_b}=f(W_{t\wedge T_b})Z_{t\wedge T_b}\le1$ by \eqref{eq:6-4-04:bd}, and by \eqref{eq:6-4-04:fin} and the continuity of $Y$,
\begin{equation*}
\lim_{t\to\infty}Y_{t\wedge T_b}=Y_{T_b}=f(b)\,e^{-A_{T_b}}=\exp\bigl(-\alpha T_b-\beta\,\Gamma_+(T_b)-\gamma L_{T_b}(0)\bigr),\qquad\text{a.s. }P^x .
\end{equation*}
The bounded convergence theorem yields
\begin{equation*}
E^x\exp\bigl(-\alpha T_b-\beta\,\Gamma_+(T_b)-\gamma L_{T_b}(0)\bigr)=f(x);\qquad x\le b ,
\end{equation*}
which is \bk{(4.15)}.

\rmk{The argument is that of \bk{Theorem 4.3}, written out for the half-line $(-\infty,b]$ on which the process is stopped at its \emph{upper} end-point. One may instead quote \bk{Theorem 4.3} for the Brownian family $-W$ (whose local time at the origin is again $L(0)$, by \bk{(4.1)}), with the interval $[-b,\infty)$, $n=1$, $a_1=0$, $\gamma_1=\gamma$, $k=\beta\ind_{(-\infty,0)}$ and the bounded function $\tilde v(y)=f(-y)$; relations \eqref{eq:6-4-04:ode}, \eqref{eq:6-4-04:jump} and $f(b)=1$ are precisely \bk{(4.11)}, \bk{(4.12)} and \bk{(4.13)}. In doing so one should read the parenthetical clause of \bk{Theorem 4.3} as ``\bk{(4.14)} should be omitted if $c=\infty$'': the printed ``\bk{(4.13)}'' is a misprint, since it is the condition $\tilde v(c)=0$ that is void for $c=\infty$, whereas $\tilde v(b)=1$ is indispensable. Steps 2 and 3 use only $\alpha>0$, $\beta\ge0$, $\gamma\ge0$; for $\beta=\gamma=0$ the function $f$ reduces to $f(y)=e^{-(b-y)\sqrt{2\alpha}}$, $y\le b$, and \bk{(4.15)} to the familiar formula $E^xe^{-\alpha T_b}=e^{-(b-x)\sqrt{2\alpha}}$ of \bk{(2.8.6)}.}

\emph{Step 4: the limit \bk{(4.16)}.} Take $x=0$ and fix $\gamma>0$; \bk{(4.15)} reads $E^0\exp(-\alpha T_b-\beta\Gamma_+(T_b)-\gamma L_{T_b}(0))=1/\Delta$. Under $P^0$ we have $0\le\Gamma_+(T_b)\le T_b<\infty$ a.s.\ by \eqref{eq:6-4-04:fin}, so as $(\alpha,\beta)\to(0,0)$ in $(0,\infty)^2$ the integrand, which is bounded by $1$, converges a.s.\ to $\exp(-\gamma L_{T_b}(0))$; by bounded convergence the left-hand side tends to $E^0\exp(-\gamma L_{T_b}(0))$. On the right-hand side $\theta\downarrow0$ and $\kappa\downarrow0$, so that $\sinh(b\theta)/\theta\to b$, $\cosh(b\theta)\to1$ and
\begin{equation*}
\Delta=(\gamma+\kappa)\,\frac{\sinh(b\theta)}{\theta}+\cosh(b\theta)\longrightarrow\gamma b+1 .
\end{equation*}
This proves \bk{(4.16)}. Since
\begin{equation*}
\int_0^\infty e^{-\gamma\ell}\,\frac1b\,e^{-\ell/b}\,d\ell=\frac{1/b}{\gamma+1/b}=\frac{1}{1+\gamma b};\qquad\gamma>0 ,
\end{equation*}
the nonnegative random variable $L_{T_b}(0)$ has under $P^0$ the same Laplace transform as the exponential distribution with density $b^{-1}e^{-\ell/b}$, $\ell>0$. By the uniqueness theorem for Laplace transforms of probability measures on $[0,\infty)$, $L_{T_b}(0)$ is exponentially distributed under $P^0$ with parameter $1/b$, and in particular $E^0L_{T_b}(0)=b$.

\emph{Step 5: the limit \bk{(4.17)}.} Fix $\beta>0$ and $0\le x\le b$. Under $P^x$ we have $T_b<\infty$ and $L_{T_b}(0)<\infty$ a.s., so as $(\alpha,\gamma)\to(0,0)$ in $(0,\infty)^2$ the left-hand side of \bk{(4.15)} converges to $E^xe^{-\beta\Gamma_+(T_b)}$, again by bounded convergence. On the right-hand side $\theta\to\sqrt{2\beta}>0$ and $(\gamma+\kappa)/\theta\to0$, so the numerator tends to $\cosh(x\sqrt{2\beta})$ and the denominator to $\cosh(b\sqrt{2\beta})$. This is \bk{(4.17)}, and $x=0$ gives
\begin{equation}\label{eq:6-4-04:cosh}
E^0e^{-\beta\Gamma_+(T_b)}=\frac{1}{\cosh(b\sqrt{2\beta})} ,
\end{equation}
which is \bk{(4.4.23)}.

\emph{Step 6: \bk{(4.4.23)} from \bk{(2.8.29)} and \bk{Theorem 3.1}.} We work under $P^0$. Recall $\Gamma_+^{-1}(\tau)=\inf\{t\ge0;\ \Gamma_+(t)>\tau\}$ from \bk{(3.9)} and $W_+(\tau)=W_{\Gamma_+^{-1}(\tau)}$ from \bk{(3.12)}. The idea is that $\Gamma_+(T_b)$ is the first time the process $W_+$, which runs on the clock $\Gamma_+$, reaches the level $b$.

\emph{(a) A pathwise identity.} Fix $\omega$ with $T_b(\omega)<\infty$ and put $\sigma\defeq\Gamma_+(T_b)$. Since $W_{T_b}=b>0$, the continuous path $W$ is strictly positive on a neighbourhood of $T_b$ in $[0,\infty)$, and on this neighbourhood the nondecreasing function $\Gamma_+$ is strictly increasing. Hence
\begin{equation}\label{eq:6-4-04:mono}
\Gamma_+(t)<\sigma\ \text{ for } 0\le t<T_b,\qquad \Gamma_+(t)>\sigma\ \text{ for } t>T_b .
\end{equation}
The second half of \eqref{eq:6-4-04:mono} and the continuity of $\Gamma_+$ give $\Gamma_+^{-1}(\sigma)=T_b$, whence $W_+(\sigma)=W_{T_b}=b$. If $0\le\tau<\sigma$, then $t\defeq\Gamma_+^{-1}(\tau)\le T_b$ is finite and $\Gamma_+(t)=\tau<\sigma$ by continuity of $\Gamma_+$, so $t<T_b$ by \eqref{eq:6-4-04:mono}, and therefore $W_+(\tau)=W_t<b$ (as $W_0=0<b$, the path stays strictly below $b$ before $T_b$). We have shown that
\begin{equation}\label{eq:6-4-04:hit}
\Gamma_+(T_b)=\inf\{\tau\ge0;\ W_+(\tau)=b\}\qquad\text{on }\{T_b<\infty\},\ \text{hence a.s. }P^0 .
\end{equation}

\emph{(b) Identification of the law.} By \bk{Theorem 3.1}, $W_+$ is a reflected Brownian motion under $P^0$; precisely, by \bk{(3.15)} we have a.s.\ $W_+(\tau)=\max_{0\le u\le\tau}B_+(u)-B_+(\tau)$ for all $\tau\ge0$, where $B_+$ is a standard Brownian motion, and by \bk{Theorem 3.6.17} this continuous process has the same law on $C[0,\infty)$ as $\abs{W}$ under $P^0$. The mapping $H(y)\defeq\inf\{\tau\ge0;\ y(\tau)=b\}$ is measurable on $\{y\in C[0,\infty);\ y(0)=0\}$, since $\{H\le t\}=\{\max_{0\le s\le t}y(s)\ge b\}$ there. Therefore, by \eqref{eq:6-4-04:hit}, $\Gamma_+(T_b)=H(W_+)$ has under $P^0$ the same distribution as
\begin{equation*}
H(\abs{W})=\inf\{t\ge0;\ \abs{W_t}=b\}=T_b\wedge T_{-b} .
\end{equation*}
Finally, under $P^0$ the process $W+b$ is a Brownian motion starting at $b$, so $T_{-b}\wedge T_b$ under $P^0$ has the same law as $T_0\wedge T_{2b}$ under $P^b$. Formula \bk{(2.8.29)} with $a=2b$, $x=b$ and $\alpha=\beta$ now gives
\begin{equation*}
E^0e^{-\beta\Gamma_+(T_b)}=E^0e^{-\beta(T_{-b}\wedge T_b)}=E^be^{-\beta(T_0\wedge T_{2b})}
=\frac{\cosh\bigl((b-\tfrac{2b}{2})\sqrt{2\beta}\bigr)}{\cosh\bigl(\tfrac{2b}{2}\sqrt{2\beta}\bigr)}=\frac{1}{\cosh(b\sqrt{2\beta})} ,
\end{equation*}
in agreement with \eqref{eq:6-4-04:cosh}: \bk{(4.4.23)} follows from \bk{(2.8.29)} and \bk{Theorem 3.1} with no computation. Neither \bk{(2.8.29)} nor \bk{Theorem 3.1} depends on \bk{(4.4.23)} or on the present problem, so the derivation is not circular.
\end{kex}

\begin{kex}{Exercise}{4.5}{}
We work under $P^0$. Under this measure $W$ is a Brownian motion with respect to the given filtration, hence also with respect to the smaller filtration $\{\F^W_t\}$ which it generates: $W$ is adapted to $\{\F^W_t\}$, and for $s<t$ the increment $W_t-W_s$ is independent of the given $\sigma$-field at time $s$, which contains $\F^W_s$. We let $\{\F_t\}$ be the augmentation of $\{\F^W_t\}$ under $P^0$ (\bk{Definition 2.7.2}), with $P^0$ extended to the completion of $\F^W_\infty$ as in \bk{Problem 2.7.3}. By construction every $\F_t$ contains all $P^0$-null sets; the filtration $\{\F_t\}$ is right-continuous by \bk{Proposition 2.7.7}, which applies because the Brownian motion $\{W_t,\F^W_t;\ t\ge0\}$ is a strong Markov process (\bk{Theorem 2.6.15}); and $W$ is still a Brownian motion relative to $\{\F_t\}$ (\bk{Theorem 2.7.9}). Thus $\{\F_t\}$ satisfies the usual conditions, and the processes $M$, $X=M-W$ introduced below, being functionals of $W$, are adapted to it. We give a direct proof by It\^o's rule applied to the pair $(M,\,M-W)$; the connection with the hint is explained in the remark at the end. Fix $\alpha>0$, $\gamma>0$, and write
\begin{equation*}
\lambda\defeq\sqrt{2\alpha},\qquad X_t\defeq M_t-W_t\ \ge 0,\qquad
g(x)\defeq \lambda\cosh(\lambda x)+\gamma\sinh(\lambda x);\quad x\in\R .
\end{equation*}
The function $g$ is of class $C^\infty$ and satisfies
\begin{equation}\label{eq:6-4-05:g}
\tfrac12 g''=\alpha g\ \text{ on }\R,\qquad g(0)=\lambda,\qquad g'(0)=\gamma\lambda=\gamma\, g(0),
\qquad g(b)=\lambda\cosh(\lambda b)+\gamma\sinh(\lambda b)>0 .
\end{equation}
The process $M$ is continuous, adapted and nondecreasing with $M_0=0$, $P^0$-a.s., so $X$ is a continuous, adapted, nonnegative process with $X_0=0$, and $R_b=\inf\{t\ge0;\ X_t=b\}$ (with $\inf\emptyset=\infty$) is a stopping time, being the first time a continuous adapted process enters a closed set (\bk{Problem 1.2.7}). On $\{R_b=\infty\}$ we interpret $\exp(-\alpha R_b-\gamma M_{R_b})$ as $0$; we shall see in Step~3 that $P^0[R_b<\infty]=1$.

\emph{Step 1: the measure $dM$ is carried by $\{t;\ X_t=0\}$.} Fix a path for which $W$ is continuous. The set $G\defeq\{t>0;\ X_t>0\}$ is open, hence a countable union of disjoint open intervals. Let $[u,v]$ be a compact interval contained in $G$. We have $M_v=\max\bigl(M_u,\max_{u\le s\le v}W_s\bigr)$. If $\max_{u\le s\le v}W_s$ exceeded $M_u$, and $s^*\in[u,v]$ were a point where this maximum is attained, then $W_{s^*}\ge\max(M_u,\max_{u\le s\le s^*}W_s)=M_{s^*}$, i.e.\ $X_{s^*}=0$, contradicting $s^*\in G$. Thus $M_v=M_u$. Consequently $M$ is constant on every compact subinterval of every component of $G$, so the Lebesgue--Stieltjes measure $dM$ of the continuous nondecreasing function $M$ assigns measure zero to each component of $G$, and
\begin{equation}\label{eq:6-4-05:flat}
\int_0^\infty \ind_{\{X_s>0\}}\,dM_s=0 .
\end{equation}

\emph{Step 2: a martingale.} The triple $(t,M_t,X_t)$ is a vector of continuous semimartingales: $t$ and $M$ are continuous, adapted and of bounded variation on bounded intervals, and $X=-W+M$ has local martingale part $-W$ with $\qv{-W}_t=t$. We apply It\^o's rule (\bk{Theorem 3.3.6}) to the $C^\infty$ function $F(t,m,x)=e^{-\alpha t-\gamma m}g(x)$ and obtain, for the process
\begin{equation*}
Y_t\defeq e^{-\alpha t-\gamma M_t}\,g(X_t);\quad 0\le t<\infty,
\end{equation*}
with the abbreviation $D_s\defeq e^{-\alpha s-\gamma M_s}\in(0,1]$, the identity
\begin{align*}
Y_t=g(0)&+\int_0^t D_s\bigl[\tfrac12g''(X_s)-\alpha g(X_s)\bigr]ds
-\int_0^t D_s\,g'(X_s)\,dW_s\\
&+\int_0^t D_s\bigl[g'(X_s)-\gamma g(X_s)\bigr]dM_s;\quad 0\le t<\infty,\ P^0\text{-a.s.}
\end{align*}
The $ds$-integral vanishes by \eqref{eq:6-4-05:g}. In the $dM$-integral, the integrand equals $D_s[g'(0)-\gamma g(0)]=0$ on $\{s;\ X_s=0\}$ by \eqref{eq:6-4-05:g}, and the complement $\{s;\ X_s>0\}$ is $dM$-null by \eqref{eq:6-4-05:flat}; so this integral vanishes as well. Therefore
\begin{equation}\label{eq:6-4-05:Y}
Y_{t\wedge R_b}=g(0)-\int_0^{t}\ind_{\{s\le R_b\}}D_s\,g'(X_s)\,dW_s;\quad 0\le t<\infty,\ P^0\text{-a.s.}
\end{equation}
For $s\le R_b$ we have $0\le X_s\le b$ (the continuous process $X$ starts at $0$ and cannot exceed $b$ before reaching it), so the integrand in \eqref{eq:6-4-05:Y} is progressively measurable and bounded in absolute value by $\max_{0\le x\le b}\abs{g'(x)}<\infty$. The stochastic integral in \eqref{eq:6-4-05:Y} is thus a square-integrable martingale starting at zero, and taking expectations we get
\begin{equation}\label{eq:6-4-05:E}
E^0\,Y_{t\wedge R_b}=g(0)=\lambda;\quad 0\le t<\infty.
\end{equation}

\emph{Step 3: proof of \textup{(4.18)}.} Since $g$ is positive and increasing on $[0,\infty)$, we have $0<Y_{t\wedge R_b}\le g(b)$ for all $t$. As $t\to\infty$: on $\{R_b<\infty\}$, $Y_{t\wedge R_b}\to Y_{R_b}=g(b)\exp(-\alpha R_b-\gamma M_{R_b})$, because $X_{R_b}=b$ by continuity; on $\{R_b=\infty\}$, $Y_{t\wedge R_b}=Y_t\le e^{-\alpha t}g(b)\to0$. The bounded convergence theorem and \eqref{eq:6-4-05:E} give
\begin{equation*}
g(b)\,E^0\bigl[\ind_{\{R_b<\infty\}}\exp(-\alpha R_b-\gamma M_{R_b})\bigr]=\lambda ,
\end{equation*}
that is, in view of \eqref{eq:6-4-05:g},
\begin{equation}\label{eq:6-4-05:main}
E^0\bigl[\ind_{\{R_b<\infty\}}\exp(-\alpha R_b-\gamma M_{R_b})\bigr]
=\frac{\sqrt{2\alpha}}{\sqrt{2\alpha}\cosh(b\sqrt{2\alpha})+\gamma\sinh(b\sqrt{2\alpha})};
\quad\alpha>0,\ \gamma>0.
\end{equation}
Now let $\alpha\downarrow0$ in \eqref{eq:6-4-05:main} with $\gamma>0$ fixed. On the left, $\ind_{\{R_b<\infty\}}e^{-\alpha R_b}\uparrow\ind_{\{R_b<\infty\}}$, and the monotone convergence theorem applies. On the right, dividing numerator and denominator by $\lambda=\sqrt{2\alpha}$ and using $\cosh(b\lambda)\to1$, $\sinh(b\lambda)/\lambda\to b$ as $\lambda\downarrow0$, the limit is $1/(1+\gamma b)$. Hence
\begin{equation}\label{eq:6-4-05:lim}
E^0\bigl[\ind_{\{R_b<\infty\}}\exp(-\gamma M_{R_b})\bigr]=\frac{1}{1+\gamma b};\quad\gamma>0.
\end{equation}
Letting $\gamma\downarrow0$ in \eqref{eq:6-4-05:lim} and using monotone convergence once more (recall $M_{R_b}\ge0$), we obtain $P^0[R_b<\infty]=1$. The indicator in \eqref{eq:6-4-05:main} may therefore be dropped, and \eqref{eq:6-4-05:main} is (4.18).

\emph{Step 4: proof of \textup{(4.19)} and the law of $M_{R_b}$.} Since $P^0[R_b<\infty]=1$, the identity \eqref{eq:6-4-05:lim} is (4.19). The right-hand side of (4.19) is the Laplace transform of the exponential distribution with mean $b$:
\begin{equation*}
\int_0^\infty e^{-\gamma m}\,\frac1b\,e^{-m/b}\,dm=\frac{1/b}{\gamma+1/b}=\frac1{1+\gamma b};\quad\gamma>0 .
\end{equation*}
A probability measure on $[0,\infty)$ is determined by its Laplace transform, so $P^0[M_{R_b}\in dm]=b^{-1}e^{-m/b}\,dm$ on $(0,\infty)$; in particular $E^0M_{R_b}=\int_0^\infty (m/b)\,e^{-m/b}\,dm=b$.

\rmk{\emph{Checks, and the role of the hint.} As $\gamma\downarrow0$, (4.18) becomes $E^0e^{-\alpha R_b}=1/\cosh(b\sqrt{2\alpha})$. This agrees with the hint: by \bk{Theorem 3.6.17}, the pair $(M-W,\,M)$ has under $P^0$ the same law as $(\abs{W},\,2L(0))$, and $(R_b,M_{R_b})$ is obtained from the first pair by the same measurable functional of the paths which produces $(\tau_b,\,2L_{\tau_b}(0))$ from the second, where $\tau_b\defeq\inf\{t\ge0;\ \abs{W_t}=b\}$. Thus (4.18) and (4.19) are equivalent to
\begin{equation*}
E^0\exp\bigl(-\alpha\tau_b-2\gamma L_{\tau_b}(0)\bigr)=\frac{\sqrt{2\alpha}}{\sqrt{2\alpha}\cosh(b\sqrt{2\alpha})+\gamma\sinh(b\sqrt{2\alpha})},
\qquad E^0\exp\bigl(-2\gamma L_{\tau_b}(0)\bigr)=\frac1{1+\gamma b},
\end{equation*}
and for $\gamma\downarrow0$ the first of these is the familiar transform $E^0e^{-\alpha\tau_b}=1/\cosh(b\sqrt{2\alpha})$ of the exit time of $W$ from $(-b,b)$. The argument above did not use \bk{Theorem 3.6.17}; the function $g$ plays for $(M-W,M)$ the role that the solution $\tilde v$ of \bk{(4.11)--(4.14)} plays in \bk{Theorem 4.3} for $(W,L(0))$, the condition $g'(0)=\gamma g(0)$ replacing the jump condition \bk{(4.12)}.}
\end{kex}

\begin{kex}{Exercise}{4.6}{}
Throughout, $\Gamma_+(t)=\meas\{0\le s\le t;\ W_s>0\}=\int_0^t\ind_{(0,\infty)}(W_s)\,ds$ as in \bk{(3.3)},
$L_t(0)$ is the local time \bk{(4.1)}, and we abbreviate $\mu\defeq\sqrt{2\alpha}$, $\nu\defeq\sqrt{2(\alpha+\beta)}$ and
\[
g_a(t;b,\tau)\defeq\frac{b}{\pi\tau^{1/2}(t-\tau)^{3/2}}
\exp\Bigl\{-\frac{b^2}{2(t-\tau)}-\frac{(a+b)^2}{2\tau}\Bigr\};
\qquad b>0,\ 0<\tau<t .
\]

\textbf{(i)} We apply \bk{Theorem 4.1} with $n=1$, $a_1=0$, $\gamma_1=\gamma$ and
\[
f=\ind_{[a,\infty)},\qquad k=\beta\,\ind_{(0,\infty)} .
\]
Both functions are piecewise continuous in the sense of \bk{Definition 4.4.8}, $k\ge 0$, $D_f=\{a\}$,
$D_k=\{0\}$, and \bk{(4.6)} holds because $0\le f\le 1$. Since
$\int_0^t k(W_s)\,ds=\beta\Gamma_+(t)$, the function $u$ of \bk{(4.5)} is
\[
u(x)=E^x\int_0^\infty \ind_{[a,\infty)}(W_t)\,e^{-\alpha t-\beta\Gamma_+(t)-\gamma L_t(0)}\,dt ,
\]
and the left-hand side of (4.20) is $u(0)$. We exhibit a bounded function $\tilde u$ with the properties
required in \bk{Theorem 4.1}. Equation \bk{(4.7)} reads
\[
\tfrac12\tilde u''=\alpha\tilde u\ \text{ on }(-\infty,0),\qquad
\tfrac12\tilde u''=(\alpha+\beta)\tilde u\ \text{ on }(0,a),\qquad
\tfrac12\tilde u''=(\alpha+\beta)\tilde u-1\ \text{ on }(a,\infty),
\]
so we set
\[
\tilde u(x)\defeq
\begin{cases}
A\,e^{\mu x}; & x\le 0,\\[2pt]
B\,e^{\nu x}+C\,e^{-\nu x}; & 0\le x\le a,\\[2pt]
\dfrac{2}{\nu^2}+D\,e^{-\nu(x-a)}; & x\ge a,
\end{cases}
\]
with the constants
\[
B\defeq\frac{e^{-\nu a}}{\nu^2},\qquad
A\defeq\frac{2\nu B}{\gamma+\mu+\nu},\qquad
C\defeq A-B,\qquad
D\defeq C e^{-\nu a}-\frac1{\nu^2}.
\]
Each of the three pieces solves the corresponding equation above (recall $\mu^2=2\alpha$,
$\nu^2=2(\alpha+\beta)$, so that the constant $2/\nu^2=1/(\alpha+\beta)$ is a particular solution on
$(a,\infty)$); hence $\tilde u$ is $C^2$ on $\R\setminus\{0,a\}$ and satisfies \bk{(4.7)} there. We check the
matching conditions.
\begin{itemize}
\item \emph{Continuity at $0$:} $B+C=A$ by the definition of $C$.
\item \emph{Continuity at $a$:} $Be^{\nu a}+Ce^{-\nu a}=\nu^{-2}+Ce^{-\nu a}=2\nu^{-2}+D$.
\item \emph{$C^1$ at $a$:} the left derivative is $\nu(Be^{\nu a}-Ce^{-\nu a})=\nu(\nu^{-2}-Ce^{-\nu a})=-\nu D$,
which is the right derivative.
\item \emph{Condition \bk{(4.8)} at $0$:} $\tilde u'(0+)-\tilde u'(0-)=\nu(B-C)-\mu A=\nu(2B-A)-\mu A
=2\nu B-(\mu+\nu)A=\gamma A=\gamma\tilde u(0)$, because $2\nu B=(\gamma+\mu+\nu)A$.
\end{itemize}
Thus $\tilde u$ is continuous on $\R$, $C^1$ on $\R\setminus\{0\}$, $C^2$ on $\R\setminus\{0,a\}$, and satisfies
\bk{(4.7)}, \bk{(4.8)}. It is bounded: $\abs{\tilde u}\le\abs{A}$ on $(-\infty,0]$, $\tilde u$ is continuous
on the compact interval $[0,a]$, and $\abs{\tilde u}\le 2\nu^{-2}+\abs{D}$ on $[a,\infty)$. Since $f$ is
bounded as well, \bk{Theorem 4.1} gives $u=\tilde u$; in particular
\[
u(0)=A=\frac{2e^{-a\nu}}{\nu\,(\gamma+\mu+\nu)},
\]
which is (4.20).

\textbf{(ii)} Fix $a>0$. The assertion (4.21) means that for every $t>0$ and every Borel set
$\Lambda\subseteq[0,\infty)^2$,
\begin{equation}\label{eq:6-4-06:meaning}
P^0\bigl[W_t\ge a;\ (L_t(0),\Gamma_+(t))\in\Lambda\bigr]
=\iint_{\Lambda\cap((0,\infty)\times(0,t))} g_a(t;b,\tau)\,db\,d\tau .
\end{equation}
Fix $\beta>0$, $\gamma>0$ and define, for $t>0$,
\[
\varphi(t)\defeq E^0\bigl[\ind_{[a,\infty)}(W_t)\,e^{-\beta\Gamma_+(t)-\gamma L_t(0)}\bigr],
\qquad
\psi(t)\defeq\int_0^t\!\!\int_0^\infty e^{-\beta\tau-\gamma b}\,g_a(t;b,\tau)\,db\,d\tau .
\]

\emph{Step 1: the two Laplace transforms agree.} The process
$(t,\omega)\mapsto\ind_{[a,\infty)}(W_t)e^{-\alpha t-\beta\Gamma_+(t)-\gamma L_t(0)}$ is nonnegative and
jointly measurable (it is built from continuous adapted processes), so Tonelli's theorem and (4.20) give,
for every $\alpha>0$,
\begin{equation}\label{eq:6-4-06:lt1}
\int_0^\infty e^{-\alpha t}\varphi(t)\,dt=\frac{2e^{-a\nu}}{\nu\,(\gamma+\mu+\nu)} .
\end{equation}
On the other hand, writing
\[
g_a(t;b,\tau)=\frac{2}{\sqrt{2\pi\tau}}\,e^{-(a+b)^2/2\tau}\cdot
\frac{b}{\sqrt{2\pi(t-\tau)^3}}\,e^{-b^2/2(t-\tau)},
\qquad e^{-\alpha t}=e^{-\alpha\tau}e^{-\alpha(t-\tau)},
\]
and using Tonelli's theorem (all integrands are nonnegative) together with the substitution $r=t-\tau$ in the
innermost integral, we obtain
\begin{align*}
\int_0^\infty e^{-\alpha t}\psi(t)\,dt
&=\int_0^\infty e^{-\gamma b}\int_0^\infty e^{-(\alpha+\beta)\tau}\frac{2}{\sqrt{2\pi\tau}}\,e^{-(a+b)^2/2\tau}
\biggl(\int_0^\infty e^{-\alpha r}\frac{b}{\sqrt{2\pi r^3}}\,e^{-b^2/2r}\,dr\biggr)d\tau\,db\\
&=\int_0^\infty e^{-\gamma b}\,e^{-b\mu}\biggl(\int_0^\infty e^{-(\alpha+\beta)\tau}\frac{2}{\sqrt{2\pi\tau}}\,
e^{-(a+b)^2/2\tau}\,d\tau\biggr)db\\
&=\int_0^\infty e^{-\gamma b}\,e^{-b\mu}\,\frac{2}{\nu}\,e^{-(a+b)\nu}\,db
=\frac{2e^{-a\nu}}{\nu\,(\gamma+\mu+\nu)} .
\end{align*}
Here the second equality is the formula displayed in the statement (with $\lambda=\alpha$), and the third is
the formula of \bk{Remark 4.4.10} with $\alpha+\beta$ in place of $\alpha$ and $\xi=a+b$. Comparing with
\eqref{eq:6-4-06:lt1},
\begin{equation}\label{eq:6-4-06:lt2}
\int_0^\infty e^{-\alpha t}\varphi(t)\,dt=\int_0^\infty e^{-\alpha t}\psi(t)\,dt<\infty;\qquad\alpha>0 .
\end{equation}

\emph{Step 2: $\varphi$ and $\psi$ are continuous on $(0,\infty)$.} For $\varphi$: fix $t>0$ and let
$s\to t$. The paths of $\Gamma_+$ and $L_\cdot(0)$ are continuous, and
$\ind_{[a,\infty)}(W_s)\to\ind_{[a,\infty)}(W_t)$ on the event $\{W_t\ne a\}$, which has $P^0$-probability
one. Since the integrand is bounded by $1$, the bounded convergence theorem gives $\varphi(s)\to\varphi(t)$.

For $\psi$: the substitution $\tau=ts$ gives $\psi(t)=\int_0^1\!\int_0^\infty e^{-\beta ts-\gamma b}h_t(b,s)\,db\,ds$
with
\[
h_t(b,s)\defeq t\,g_a(t;b,ts)=\frac{b}{\pi t\,s^{1/2}(1-s)^{3/2}}
\exp\Bigl\{-\frac1{2t}\Bigl[\frac{b^2}{1-s}+\frac{(a+b)^2}{s}\Bigr]\Bigr\} ,
\]
which is continuous in $t>0$ for fixed $(b,s)$. For $0<t_0\le t\le t_1$ we have
\[
e^{-\beta ts-\gamma b}h_t(b,s)\le \frac{t_1}{t_0}\,e^{-\gamma b}h_{t_1}(b,s),
\]
and the right-hand side is integrable on
$(0,\infty)\times(0,1)$: indeed $\sup_{\tau>0}\tau^{-1/2}e^{-c/\tau}=(2ec)^{-1/2}$ for $c>0$, and
$\int_0^\infty \frac{b}{\sqrt{2\pi r^3}}e^{-b^2/2r}\,dr=1$ (let $\lambda\downarrow0$ in the formula of the
statement and use monotone convergence), whence
\[
\int_0^1 h_{t_1}(b,s)\,ds=\int_0^{t_1} g_a(t_1;b,\tau)\,d\tau
\le\frac{2}{\sqrt{2\pi e}\,(a+b)}\int_0^{t_1}\frac{b}{\sqrt{2\pi r^3}}\,e^{-b^2/2r}\,dr
\le\frac{2}{\sqrt{2\pi e}\;a},
\]
and $\int_0^\infty e^{-\gamma b}\,db<\infty$. By dominated convergence $\psi$ is continuous on $[t_0,t_1]$,
hence on $(0,\infty)$.

\emph{Step 3: $\varphi=\psi$.} By \eqref{eq:6-4-06:lt2} with $\alpha=1$, the measures
$e^{-t}\varphi(t)\,dt$ and $e^{-t}\psi(t)\,dt$ on $[0,\infty)$ are finite, and by \eqref{eq:6-4-06:lt2}
with $\alpha+1$ in place of $\alpha$ their Laplace transforms coincide on $(0,\infty)$. By the uniqueness
theorem for Laplace transforms of finite measures on $[0,\infty)$ the two measures are equal, so
$\varphi=\psi$ Lebesgue-almost everywhere; both functions being continuous on $(0,\infty)$, we conclude
$\varphi(t)=\psi(t)$ for every $t>0$.

\emph{Step 4: conclusion.} Fix $t>0$ and consider the two measures on $\B([0,\infty)^2)$
\[
m_1(\Lambda)\defeq P^0\bigl[W_t\ge a;\ (L_t(0),\Gamma_+(t))\in\Lambda\bigr],\qquad
m_2(\Lambda)\defeq\iint_{\Lambda\cap((0,\infty)\times(0,t))}g_a(t;b,\tau)\,db\,d\tau .
\]
Step 3, valid for all $\beta>0$, $\gamma>0$, says precisely that
\[
\iint e^{-\gamma b-\beta\tau}\,m_1(db,d\tau)=\varphi(t)=\psi(t)=\iint e^{-\gamma b-\beta\tau}\,m_2(db,d\tau);
\qquad\beta>0,\ \gamma>0 .
\]
The measure $m_1$ is finite, with total mass at most $1$; letting $\beta\downarrow0$, $\gamma\downarrow0$ in
this identity and using monotone convergence on both sides, we see that
$m_2([0,\infty)^2)=m_1([0,\infty)^2)\le1$, so $m_2$ is finite as well. By the uniqueness theorem for Laplace transforms of finite measures on $[0,\infty)^2$, $m_1=m_2$, which is
\eqref{eq:6-4-06:meaning}, i.e., (4.21).

\textbf{(iii)} Fix $t>0$ and recall
$f(a,b;t,\tau)=\dfrac{b(a+b)}{\pi\tau^{3/2}(t-\tau)^{3/2}}
\exp\Bigl\{-\dfrac{b^2}{2(t-\tau)}-\dfrac{(a+b)^2}{2\tau}\Bigr\}$ from \bk{(3.29)}. We have to show that,
with $U\defeq(0,\infty)\times(0,t)$,
\begin{equation}\label{eq:6-4-06:pos}
P^0\bigl[(W_t,L_t(0),\Gamma_+(t))\in\Theta\bigr]
=\iiint_{\Theta\cap((0,\infty)\times U)}f(a,b;t,\tau)\,da\,db\,d\tau
\end{equation}
for every $\Theta\in\B\bigl((0,\infty)\times[0,\infty)^2\bigr)$, and
\begin{equation}\label{eq:6-4-06:neg}
P^0\bigl[(W_t,L_t(0),\Gamma_+(t))\in\Theta\bigr]
=\iiint_{\Theta\cap((-\infty,0)\times U)}f(-a,b;t,t-\tau)\,da\,db\,d\tau
\end{equation}
for every $\Theta\in\B\bigl((-\infty,0)\times[0,\infty)^2\bigr)$.
(Since $P^0[W_t=0]=0$, these two identities describe the whole law of the triple.)

\emph{The case $a>0$.} For $b>0$, $0<\tau<t$ fixed, $x\mapsto g_x(t;b,\tau)$ is continuously differentiable
on $(0,\infty)$ with
\[
-\frac{\partial}{\partial x}g_x(t;b,\tau)=\frac{x+b}{\tau}\,g_x(t;b,\tau)=f(x,b;t,\tau),
\qquad\lim_{x\to\infty}g_x(t;b,\tau)=0,
\]
so that $g_a(t;b,\tau)=\int_a^\infty f(x,b;t,\tau)\,dx$ for every $a>0$. Substituting this into
\eqref{eq:6-4-06:meaning} and using Tonelli's theorem, we see that \eqref{eq:6-4-06:pos} holds for every set
of the form $\Theta=[a,\infty)\times\Lambda$ with $a>0$ and $\Lambda\in\B([0,\infty)^2)$. These sets form a
$\pi$-system which generates $\B((0,\infty)\times[0,\infty)^2)$ and contains the sequence
$[1/n,\infty)\times[0,\infty)^2$, $n\ge1$, increasing to $(0,\infty)\times[0,\infty)^2$. Both sides of
\eqref{eq:6-4-06:pos} are measures in $\Theta$, and they are finite and equal on each
$[1/n,\infty)\times[0,\infty)^2$. By Dynkin's $\pi$-$\lambda$ theorem, applied to the traces of both measures
on $[1/n,\infty)\times[0,\infty)^2$, they agree on all Borel subsets of that set, and letting $n\to\infty$
yields \eqref{eq:6-4-06:pos} in general.

\emph{The case $a<0$.} Let $\wt W\defeq-W$, which is a standard Brownian motion under $P^0$, and denote by
$\wt L_t(0)$, $\wt\Gamma_+(t)$ its local time at the origin and its occupation time of $(0,\infty)$. From
\bk{(4.1)} we have $\wt L_t(0)=L_t(0)$, a.s.\ $P^0$, because $\abs{\wt W_s}=\abs{W_s}$. Moreover
\[
\wt\Gamma_+(t)=\meas\{0\le s\le t;\ W_s<0\}=t-\Gamma_+(t)-\meas\{0\le s\le t;\ W_s=0\}=t-\Gamma_+(t),
\quad\text{a.s. }P^0,
\]
since $E^0\meas\{0\le s\le t;\ W_s=0\}=\int_0^tP^0[W_s=0]\,ds=0$ by Tonelli's theorem. Now
$(W_t,L_t(0),\Gamma_+(t))$ is a.s.\ equal to $H(W)$ for the Borel-measurable functional
\[
H(w)\defeq\Bigl(w(t),\ \limsup_{n\to\infty}\tfrac n4\meas\{0\le s\le t;\ \abs{w(s)}\le\tfrac1n\},\
\meas\{0\le s\le t;\ w(s)>0\}\Bigr)
\]
on $C[0,\infty)$, again by \bk{(4.1)}; and $W$, $\wt W$ have the same law under $P^0$. Therefore
$(W_t,L_t(0),\Gamma_+(t))$ has the same law as
$H(\wt W)=(\wt W_t,\wt L_t(0),\wt\Gamma_+(t))=(-W_t,L_t(0),t-\Gamma_+(t))$, a.s. Given
$\Theta\in\B((-\infty,0)\times[0,\infty)^2)$, put
$\Theta'\defeq\{(x,b,\sigma)\in(0,\infty)\times[0,\infty)\times\R;\ (-x,b,t-\sigma)\in\Theta\}$. Because
$0\le\Gamma_+(t)\le t$, we get from this equivalence in law and from \eqref{eq:6-4-06:pos}
(applied to $\Theta'\cap((0,\infty)\times[0,\infty)^2)$)
\begin{align*}
P^0\bigl[(W_t,L_t(0),\Gamma_+(t))\in\Theta\bigr]
&=P^0\bigl[(-W_t,L_t(0),t-\Gamma_+(t))\in\Theta\bigr]
=P^0\bigl[(W_t,L_t(0),\Gamma_+(t))\in\Theta'\bigr]\\
&=\iiint_{\Theta'\cap((0,\infty)\times U)}f(x,b;t,\sigma)\,dx\,db\,d\sigma\\
&=\iiint_{\Theta\cap((-\infty,0)\times U)}f(-a,b;t,t-\tau)\,da\,db\,d\tau,
\end{align*}
where the last step is the change of variables $x=-a$, $\sigma=t-\tau$, which maps $(0,\infty)\times U$ onto
$(-\infty,0)\times U$ and has Jacobian of absolute value one. This is \eqref{eq:6-4-06:neg}, and
\eqref{eq:6-4-06:pos}, \eqref{eq:6-4-06:neg} together are \bk{(3.30)}.

\rmk{The derivation of \bk{(3.30)} given here is independent of \bk{Proposition 3.9}: it uses only
\bk{Theorem 4.1} (which rests on the generalized It\^o rule \bk{(3.6.53)}), the two elementary Laplace
transforms quoted in the statement, and \bk{(4.1)}. As a check, integrating \eqref{eq:6-4-06:meaning} over
$\Lambda=[0,\infty)^2$ must give $P^0[W_t\ge a]$; numerically, for $t=1$, $a=0.3$ the double integral of
$g_a$ equals $0.38209$, which is $P^0[W_1\ge0.3]=\int_{0.3}^\infty e^{-x^2/2}\,dx/\sqrt{2\pi}$.}
\end{kex}

\begin{kex}{Exercise}{4.10}{}
We fix $\beta>0$, $b>0$ and work in the setting of Section~6.4: $W=\{W_t,\F_t;\,0\le t<\infty\}$,
$(\Omega,\F)$, $\{P^x\}_{x\in\R}$ is a one-dimensional Brownian family with local time $L$, and
$\{R_t,\wt{\F}_t\}$, $(\wt\Omega,\wt\F)$, $\{\wt P^r\}_{r\ge0}$ is a Bessel family of dimension $2$
as in \bk{Theorem 4.7}; thus $R_t=\sqrt{(B^{(1)}_t)^2+(B^{(2)}_t)^2}$, where
$B=(B^{(1)},B^{(2)})$ is, under $\wt P^0$, a two-dimensional Brownian motion starting at the origin
(\bk{Definition 3.3.19}). We set
\[
  \varphi \defeq E^0\exp\Bigl\{-\beta\int_0^b W_t^2\,dt\Bigr\}\in(0,1],
\]
the integral being a finite random variable because $W$ has continuous paths. The idea is that, by
the Ray--Knight theorem, the occupation time $\Gamma_+(T_b)$ of \bk{(4.17)} is distributed as
$\int_0^b R_t^2\,dt$, the sum of two independent copies of $\int_0^b W_t^2\,dt$. For $n\ge1$ and
$1\le k\le n$ we write $t_k\defeq kb/n$ and introduce the Riemann sums
\[
  S_n\defeq\frac bn\sum_{k=1}^n W_{t_k}^2,\qquad
  \Lambda_n\defeq\frac bn\sum_{k=1}^n 2L_{T_b}(b-t_k),\qquad
  S^{(i)}_n\defeq\frac bn\sum_{k=1}^n \bigl(B^{(i)}_{t_k}\bigr)^2\quad(i=1,2).
\]

\emph{Step 1: $\Gamma_+(T_b)=2\int_0^b L_{T_b}(b-t)\,dt$, a.s.\ $P^0$.} Let $\Omega^*$ be the event
of \bk{Definition 3.6.3} and $\Omega^{**}\defeq\Omega^*\cap\{W_0=0,\ T_b<\infty\}$; then
$P^0(\Omega^{**})=1$ by \bk{Remark 2.6.20}. Fix $\omega\in\Omega^{**}$. Since $W_0(\omega)=0<b$ and
the path is continuous, $W_s(\omega)\le b$ for $0\le s\le T_b(\omega)$, and therefore, with the
notation \bk{(3.6.5)} and with \bk{(3.3)},
\[
  \Gamma_+(T_b)=\meas\{0\le s\le T_b;\ W_s>0\}=\meas\{0\le s\le T_b;\ 0<W_s\le b\}
  =\Gamma_{T_b}\bigl((0,b]\bigr).
\]
The identity \bk{(3.6.6)} holds at $\omega$ for every $t\in[0,\infty)$ and every Borel set; taking
$t=T_b(\omega)$ and the set $(0,b]$ we obtain
\begin{equation}\label{eq:6-4-10:occ}
  \Gamma_+(T_b)=\int_0^b 2L_{T_b}(x)\,dx=\int_0^b 2L_{T_b}(b-t)\,dt
  \qquad\text{on }\Omega^{**}.
\end{equation}
Moreover, $t\mapsto 2L_{T_b(\omega)}(b-t,\omega)$ is continuous on $[0,b]$, because
$(t,x)\mapsto L_t(x,\omega)$ is continuous. Hence its Riemann sums converge to its integral:
\begin{equation}\label{eq:6-4-10:riem1}
  \lim_{n\to\infty}\Lambda_n=\Gamma_+(T_b)\qquad\text{a.s.\ }P^0 .
\end{equation}
(On $\Omega^{**}$ the variable $L_{T_b}(a)$ is the value at the finite random time $T_b$ of the
process $t\mapsto L_t(a)$, which is adapted with continuous paths; so each $\Lambda_n$ is a random
variable on $\Omega^{**}$, and we set it equal to zero off this event.)

\emph{Step 2: use of the Ray--Knight theorem.} By \bk{Theorem 4.7}, the random vector
$\bigl(L_{T_b}(b-t_1),\dots,L_{T_b}(b-t_n)\bigr)$ under $P^0$ has the same distribution as
$\bigl(\tfrac12R^2_{t_1},\dots,\tfrac12R^2_{t_n}\bigr)$ under $\wt P^0$. Applying to both the same
continuous function on $\R^n$ we see that $\Lambda_n$ under $P^0$ has the same distribution as
\[
  \frac bn\sum_{k=1}^n R^2_{t_k}=S^{(1)}_n+S^{(2)}_n\qquad\text{under }\wt P^0,
\]
so that
\begin{equation}\label{eq:6-4-10:rk}
  E^0 e^{-\beta\Lambda_n}=\wt E^0\exp\bigl\{-\beta\bigl(S^{(1)}_n+S^{(2)}_n\bigr)\bigr\};
  \qquad n\ge1.
\end{equation}

\emph{Step 3: the right-hand side of \eqref{eq:6-4-10:rk} equals $\bigl(E^0e^{-\beta S_n}\bigr)^2$.}
With $t_0\defeq0$, property (ii) of \bk{Definition 2.5.1} shows that under $\wt P^0$ the increment
$B_{t_k}-B_{t_{k-1}}$ is independent of $\wt\F_{t_{k-1}}$, hence of
$(B_{t_j}-B_{t_{j-1}})_{j<k}$, and is normal with mean zero and covariance matrix
$(t_k-t_{k-1})I_2$; by induction on $k$, the $2n$ random variables
\[
  B^{(i)}_{t_k}-B^{(i)}_{t_{k-1}};\qquad i=1,2,\quad 1\le k\le n,
\]
are independent, and the two with index $k$ are normal with mean zero and variance $t_k-t_{k-1}$.
Since $B_0=0$ a.s.\ $\wt P^0$, the vectors $(B^{(1)}_{t_k})_{k\le n}$ and $(B^{(2)}_{t_k})_{k\le n}$,
obtained from the two groups of increments by partial summation, are independent, and each has the
distribution of the vector $(W_{t_k})_{k\le n}$ under $P^0$, whose increments are likewise
independent, centred normal with variances $t_k-t_{k-1}$, and for which $W_0=0$ a.s. Consequently
$S^{(1)}_n$ and $S^{(2)}_n$ are independent under $\wt P^0$, and each is distributed as $S_n$
under $P^0$. With \eqref{eq:6-4-10:rk} this gives
\begin{equation}\label{eq:6-4-10:sq}
  E^0 e^{-\beta\Lambda_n}=\bigl(E^0 e^{-\beta S_n}\bigr)^2;\qquad n\ge1.
\end{equation}

\emph{Step 4: passage to the limit.} For every $\omega\in\Omega$ the function $t\mapsto W_t^2(\omega)$
is continuous on $[0,b]$, so $S_n(\omega)\to\int_0^b W_t^2(\omega)\,dt$ as $n\to\infty$. The random
variables $e^{-\beta S_n}$ and $e^{-\beta\Lambda_n}$ take values in $(0,1]$, because $S_n\ge0$ and
$\Lambda_n\ge0$. Letting $n\to\infty$ in \eqref{eq:6-4-10:sq}, the bounded convergence theorem
together with \eqref{eq:6-4-10:riem1} yields
\[
  E^0 e^{-\beta\Gamma_+(T_b)}=\varphi^2 .
\]
On the other hand, formula \bk{(4.17)} of \bk{Problem 4.4}, taken at $x=0$, reads
$E^0e^{-\beta\Gamma_+(T_b)}=1/\cosh(b\sqrt{2\beta})$. Therefore
$\varphi^2=1/\cosh(b\sqrt{2\beta})$, and since $\varphi>0$,
\[
  E^0\exp\Bigl\{-\beta\int_0^b W_t^2\,dt\Bigr\}=\frac{1}{\sqrt{\cosh(b\sqrt{2\beta})}} .
\]

\rmk{The proof of \bk{Theorem 4.7} (\bk{Lemmas 4.8, 4.9}) does not use the present exercise, so
there is no circularity. As a check of the formula, Tonelli's theorem gives
$E^0\int_0^b W_t^2\,dt=\int_0^b t\,dt=b^2/2$, and for the centred Gaussian process $W$,
$E^0\bigl(\int_0^bW_t^2\,dt\bigr)^2=\int_0^b\!\!\int_0^b\bigl(st+2(s\wedge t)^2\bigr)\,ds\,dt
=\tfrac{b^4}{4}+\tfrac{b^4}{3}=\tfrac{7b^4}{12}$; this agrees with the expansion
$\bigl(\cosh(b\sqrt{2\beta})\bigr)^{-1/2}
=\bigl(1+\beta b^2+\tfrac16\beta^2b^4+\cdots\bigr)^{-1/2}
=1-\tfrac{b^2}{2}\beta+\tfrac{7b^4}{24}\beta^2+O(\beta^3)$ as $\beta\downarrow0$.}
\end{kex}

\begin{kex}{Exercise}{4.15}{}
\rmk{The closing brace after ``$u<\infty$'' in the display of (i) is a misprint: as in \bk{(4.43)}, the
conditioning is on the $\sigma$-field $\sigma(W_+(u);\,0\le u<\infty)$.}

\emph{Setting and conventions.} We are in the setting of \bk{Section 3.6}, which \bk{Sections 6.3}
and \bk{6.4} retain: $W$ is the coordinate process on $\Omega=C[0,\infty)$, $\{\F_t\}$ is the
augmented filtration \bk{(2.7.3)}, $\F=\F_\infty$, $\theta_s$ are the shift operators,
$T_x=\inf\{t\ge 0;\ W_t=x\}$, and $L$ is Brownian local time with the event $\Omega^*$ of
\bk{Definition 3.6.3}. Since $P^z(\Omega\setminus\Omega^*)=0$ for every $z$, the set $\Omega^*$ belongs
to $\F_0$, and we may assume $L_t(x,\omega)=0$ for $\omega\notin\Omega^*$; then $(t,x)\mapsto
L_t(x,\omega)$ is continuous for \emph{every} $\omega$, each $L(x)$ is a continuous, adapted (hence
progressively measurable) process, and $L_T(x)\ind_{\{T<\infty\}}$ is $\F$-measurable for every random
time $T$, and $\F_T$-measurable if $T$ is a stopping time. We shall use the following facts, valid
for every $x,z\in\R$ and $P^z$-a.e.\ $\omega$:
\begin{enumerate}[label=(\alph*)]
\item $t\mapsto L_t(x)$ is nondecreasing, and $\int_0^\infty \ind_{\R\setminus\{x\}}(W_t)\,dL_t(x)=0$
(\bk{Problem 3.6.13}\,(ii)); consequently $L(x)$ is constant on every interval $[s,t]$ such that
$W_u\ne x$ for all $u\in[s,t)$;
\item if $S$ is a stopping time of $\{\F_t\}$ with $P^z[S<\infty,\ W_S=c]=1$ for some $c\in\R$, then
$L_{S+t}(y,\omega)=L_S(y,\omega)+L_t(y,\theta_S\omega)$ for all $t\ge 0$ and $y\in\R$ (the additive
functional property of \bk{Remark 3.6.5} at the time $S$; we prove it after the Lemma below);
\item $\varlimsup_{t\to\infty}W_t=\infty$, $\varliminf_{t\to\infty}W_t=-\infty$ (by the law of the
iterated logarithm, \bk{Theorem 2.9.23}), so that every $T_y$ is finite, and
$\lim_{t\to\infty}L_t(x)=\infty$ (\bk{Problem 3.6.18}).
\end{enumerate}
We also need the strong Markov property in the following form.

\begin{lem}
Let $z,c\in\R$ and let $S$ be a stopping time of $\{\F_t\}$ with $P^z[S<\infty,\ W_S=c]=1$. Then
\[
E^z\bigl[Z\cdot(Y\circ\theta_S)\bigr]=E^z[Z]\cdot E^c[Y]
\]
for all bounded random variables $Z$, $Y$ such that $Z$ is $\F_S$-measurable and $Y$ is
$\F$-measurable.
\end{lem}
\begin{proof}
Under $P^z$ the process $\{W_t,\F_t\}$ is a Brownian motion (\bk{Theorem 2.7.15}), so by
\bk{Theorem 2.6.16} the process $\{W_{S+t}-W_S;\ t\ge 0\}$ is a standard Brownian
motion independent of $\F_{S+}\supseteq\F_S$. Since $W_S=c$ a.s., the $C[0,\infty)$-valued random
element $\theta_S\omega=W_{S+\cdot}(\omega)$ is independent of $\F_S$ and its law on
$\B(C[0,\infty))$ is $P^c$. This is the assertion when $Y$ is $\B(C[0,\infty))$-measurable. In
general, $\F$ is contained in the $P^c$-completion of $\B(C[0,\infty))$ (see \bk{(2.7.3)}), so there
are bounded Borel $Y_1\le Y\le Y_2$ with $P^c[Y_1<Y_2]=0$; then $Y_1\circ\theta_S\le Y\circ\theta_S\le
Y_2\circ\theta_S$ and $P^z[Y_1\circ\theta_S<Y_2\circ\theta_S]=P^c[Y_1<Y_2]=0$. Thus $Y\circ\theta_S$
is measurable with respect to the $P^z$-completion of $\F$ (to which $E^z$ extends), and the assertion
for $Y$ follows from that for $Y_1$.
\end{proof}

\emph{Proof of} (b). The Lemma with $Z=1$ and $Y=\ind_{\Omega^*}$ gives
$P^z[\theta_S\omega\in\Omega^*]=P^c(\Omega^*)=1$. Let $\omega$ be such that $S(\omega)<\infty$,
$\omega\in\Omega^*$ and $\theta_S\omega\in\Omega^*$, and fix $t\ge 0$. With the occupation time
$\Gamma_t(B)$ of \bk{(3.6.5)} we have from \bk{(3.6.6)}, for every bounded $B\in\B(\R)$,
\[
\int_B 2L_t(y,\theta_S\omega)\,dy=\Gamma_t(B,\theta_S\omega)
=\Gamma_{S+t}(B,\omega)-\Gamma_S(B,\omega)
=\int_B 2\bigl[L_{S+t}(y,\omega)-L_S(y,\omega)\bigr]\,dy .
\]
The two integrands are continuous functions of $y$ whose integrals over every bounded interval agree;
hence they agree for every $y\in\R$, which is (b).

\textbf{(i)} We work under $P^0$ and write $L_t=L_t(0)$, $\Lambda\defeq L_{T_a}(0)$. Let $\cN$ be the
collection of $P^0$-null sets of $\F$ and
\[
\G\defeq\sigma\bigl(\sigma(W_+(u);\ 0\le u<\infty)\cup\cN\bigr).
\]
Every set of $\G$ differs from a set of $\sigma(W_+(u);\,0\le u<\infty)$ by a null set (the sets with
this property form a $\sigma$-field containing the generators of $\G$). Therefore a version of the
conditional expectation of an integrable random variable given $\sigma(W_+(u);\,0\le u<\infty)$ is
also a version of its conditional expectation given $\G$, and it suffices to prove that
\begin{equation}\label{eq:6-4-15:goal}
\Lambda\ \text{is $\G$-measurable}\qquad\text{and}\qquad
E^0\bigl[e^{-\beta L_{T_a}(b)}\ind_G\bigr]=E^0\bigl[\varphi(\Lambda)\ind_G\bigr]\quad
\text{for all } G\in\G,
\end{equation}
where $\varphi(s)\defeq\exp\{-\beta s/(1+\beta\abs{b})\}$, $s\ge 0$.

\emph{Step 1: $\Lambda$ is $\G$-measurable.} Recall $\Gamma_+(t)=\meas\{0\le s\le t;\ W_s>0\}$,
$\Gamma_+^{-1}(\tau)=\inf\{t\ge 0;\ \Gamma_+(t)>\tau\}$, $W_+(\tau)=W_{\Gamma_+^{-1}(\tau)}$ and
$L_+(\tau)=L(\Gamma_+^{-1}(\tau))$ from \bk{(3.3)}, \bk{(3.9)}, \bk{(3.12)}, \bk{(3.13)}. Let
$S\defeq\Gamma_+(T_a)$; we argue on the event $\{T_a<\infty\}$ of probability one. The function
$\Gamma_+$ is continuous and nondecreasing, so $\Gamma_+(t)\le S$ for $t\le T_a$; and since
$W_{T_a}=a>0$, the path is positive on some interval $[T_a,T_a+\delta]$, so $\Gamma_+(t)>S$ for
$t>T_a$. Hence $\Gamma_+^{-1}(S)=T_a$, and
\begin{equation}\label{eq:6-4-15:S}
W_+(S)=W_{T_a}=a,\qquad L_+(S)=L_{T_a}=\Lambda .
\end{equation}
If $0\le\tau<S=\Gamma_+(T_a)$, then by continuity there is $t<T_a$ with $\Gamma_+(t)>\tau$, so
$\Gamma_+^{-1}(\tau)<T_a$ and $W_+(\tau)=W_{\Gamma_+^{-1}(\tau)}<a$. Together with \eqref{eq:6-4-15:S}
this shows $S=\inf\{\tau\ge 0;\ W_+(\tau)=a\}$ and $\{S>\tau\}=\{W_+(u)<a\ \text{for all }u\in[0,\tau]\}$.
Now $W_+$ and $L_+$ have continuous paths a.s., because a.s.\ \bk{(3.14)} and \bk{(3.15)} hold for
all $\tau$ simultaneously, with the continuous process $B_+$. Thus, up to a null set,
\[
\{S>\tau\}=\Bigl\{\sup\nolimits_{u\in(\Q\cap[0,\tau])\cup\{\tau\}}W_+(u)<a\Bigr\}\in\G,
\qquad 0\le\tau<\infty,
\]
and $S$ is $\G$-measurable. By \bk{Remark 3.2}, $L_+(\tau)$ is $\G$-measurable for every $\tau$. With
$S_n\defeq 2^{-n}(\lfloor 2^nS\rfloor+1)\downarrow S$, the random variables
$L_+(S_n)=\sum_{k\ge 1}L_+(k2^{-n})\ind_{\{S_n=k2^{-n}\}}$ are $\G$-measurable and converge a.s.\ to
$L_+(S)=\Lambda$; as $\G$ contains the null sets, $\Lambda$ is $\G$-measurable.

\emph{Step 2: $L_{T_a}(b)=N_\Lambda$ a.s.} Here $\rho_s=\inf\{t\ge 0;\ L_t>s\}$ and $N_s=L_{\rho_s}(b)$
are as in \bk{Lemma 4.12}. By (a) and (c), a.s.\ $L$ is continuous, nondecreasing and unbounded, so
that $\rho_s<\infty$ for all $s$ and $s\mapsto\rho_s$ is nondecreasing and right-continuous (if
$t>\rho_s$, there is $t'<t$ with $L_{t'}>s$, whence $\rho_{s'}\le t'<t$ for $s\le s'<L_{t'}$);
consequently $s\mapsto N_s$ is right-continuous. Let
\[
R\defeq\inf\{t\ge T_a;\ W_t=0\}=\inf\{t\ge 0;\ W_t=0,\ \max\nolimits_{0\le s\le t}W_s\ge a\},
\]
the hitting time of the closed set $\{0\}\times[a,\infty)$ by a continuous, adapted, two-dimensional
process, hence a stopping time; $R<\infty$ and $W_R=0$ a.s., by (c). On $[T_a,R)$ we have $W>0>b$, so
by (a), a.s.,
\begin{equation}\label{eq:6-4-15:flat}
L_R(0)=L_{T_a}(0)=\Lambda,\qquad L_R(b)=L_{T_a}(b).
\end{equation}
The event $A\defeq\bigcap_{n\ge 1}\{L_{1/n}(0)>0\}\in\F$ has $P^0(A)=1$ by \bk{Problem 3.6.13}\,(iii).
The Lemma (with $z=c=0$, $S=R$, $Z=1$, $Y=\ind_A$) gives $P^0[\theta_R\omega\in A]=P^0(A)=1$, and
then (b) with $S=R$ yields $L_{R+1/n}-L_R=L_{1/n}(0)\circ\theta_R>0$ for all $n\ge 1$, a.s. Since $L$ is
nondecreasing, we conclude from this and \eqref{eq:6-4-15:flat} that a.s.\ $L_t\le\Lambda$ for $t\le R$
and $L_t>\Lambda$ for $t>R$; that is, $\rho_\Lambda=R$, and by \eqref{eq:6-4-15:flat} again
\begin{equation}\label{eq:6-4-15:N}
N_\Lambda=L_R(b)=L_{T_a}(b),\qquad\text{a.s. }P^0 .
\end{equation}

\emph{Step 3: conditioning.} By \bk{Lemma 4.12}, formula \bk{(4.42)}, and $N_0=0$ a.s., we have
$E^0[e^{-\beta N_s}]=\varphi(s)$ for every $s\ge 0$. By \bk{Remark 4.13}, for each $s\ge 0$ the random
variable $N_s$ coincides a.s.\ with a functional of $W_-$, i.e., with a random variable measurable
with respect to $\sigma\bigl(\sigma(W_-(u);\,0\le u<\infty)\cup\cN\bigr)$; by \bk{Theorem 3.1} this
$\sigma$-field is independent of $\sigma(W_+(u);\,0\le u<\infty)$, hence of $\G$. Thus each $N_s$ is
independent of $\G$.

Let $\Lambda_n\defeq 2^{-n}(\lfloor 2^n\Lambda\rfloor+1)$; these random variables are $\G$-measurable
by Step~1, take values in $\{k2^{-n};\ k\ge 1\}$, and decrease to $\Lambda$. For $G\in\G$ the events
$G\cap\{\Lambda_n=k2^{-n}\}$ are in $\G$, so
\begin{align*}
E^0\bigl[e^{-\beta N_{\Lambda_n}}\ind_G\bigr]
&=\sum_{k\ge 1}E^0\bigl[e^{-\beta N_{k2^{-n}}}\ind_{G\cap\{\Lambda_n=k2^{-n}\}}\bigr]\\
&=\sum_{k\ge 1}\varphi(k2^{-n})\,P^0\bigl[G\cap\{\Lambda_n=k2^{-n}\}\bigr]
=E^0\bigl[\varphi(\Lambda_n)\ind_G\bigr].
\end{align*}
As $n\to\infty$, $N_{\Lambda_n}\to N_\Lambda$ a.s.\ by right-continuity of $N$, and
$\varphi(\Lambda_n)\to\varphi(\Lambda)$ by continuity of $\varphi$; all integrands are bounded by $1$.
The bounded convergence theorem and \eqref{eq:6-4-15:N} give
\[
E^0\bigl[e^{-\beta L_{T_a}(b)}\ind_G\bigr]=E^0\bigl[e^{-\beta N_\Lambda}\ind_G\bigr]
=E^0\bigl[\varphi(\Lambda)\ind_G\bigr],\qquad G\in\G,
\]
which is \eqref{eq:6-4-15:goal}. Hence
$E^0[\exp\{-\beta L_{T_a}(b)\}\mid W_+(u);\ 0\le u<\infty]
=\exp\{-\beta L_{T_a}(0)/(1+\beta\abs{b})\}$, a.s.\ $P^0$.

\textbf{(ii)} \emph{Step 1: the case of the starting point $0$.} The random variable
$e^{-\alpha\Lambda}$ is bounded and $\G$-measurable by Step~1 of (i), so (i) gives
\[
E^0\bigl[e^{-\beta L_{T_a}(b)-\alpha L_{T_a}(0)}\bigr]
=E^0\Bigl[e^{-\alpha\Lambda}\,E^0\bigl[e^{-\beta L_{T_a}(b)}\bigm|\G\bigr]\Bigr]
=E^0\bigl[e^{-\gamma L_{T_a}(0)}\bigr],\qquad
\gamma\defeq\alpha+\frac{\beta}{1+\beta\abs{b}}>0 .
\]
By \bk{Problem 4.11}, formula \bk{(4.40)}, the right-hand side equals $(1+\gamma a)^{-1}$, so that
\begin{equation}\label{eq:6-4-15:zero}
E^0\bigl[e^{-\beta L_{T_a}(b)-\alpha L_{T_a}(0)}\bigr]
=\frac{1+\beta\abs{b}}{(1+\alpha a)(1+\beta\abs{b})+\beta a}.
\end{equation}

\emph{Step 2: the starting point $b$.} Under $P^b$ we have a.s.\ $T_0<T_a<\infty$, because $b<0<a$ and
the path is continuous and satisfies (c); hence $T_a=T_0+T_a\circ\theta_{T_0}$. By (a), $L(0)$ is
constant on $[0,T_0]$, i.e.\ $L_{T_0}(0)=0$, and by (b) applied with $S=T_0$, $c=0$,
$t=T_a(\theta_{T_0}\omega)$ and $y\in\{b,0\}$ we obtain, a.s.\ $P^b$,
\[
L_{T_a}(b)=L_{T_0}(b)+\bigl(L_{T_a}(b)\bigr)\circ\theta_{T_0},\qquad
L_{T_a}(0)=\bigl(L_{T_a}(0)\bigr)\circ\theta_{T_0}.
\]
Thus, with the bounded $\F$-measurable random variable
$Y\defeq\exp\{-\beta L_{T_a}(b)-\alpha L_{T_a}(0)\}\ind_{\{T_a<\infty\}}$ and the bounded
$\F_{T_0}$-measurable random variable $Z\defeq\exp\{-\beta L_{T_0}(b)\}\ind_{\{T_0<\infty\}}$, the Lemma
(with $z=b$, $c=0$, $S=T_0$) gives
\begin{equation}\label{eq:6-4-15:split}
E^b\bigl[e^{-\beta L_{T_a}(b)-\alpha L_{T_a}(0)}\bigr]=E^b\bigl[Z\cdot(Y\circ\theta_{T_0})\bigr]
=E^b\bigl[e^{-\beta L_{T_0}(b)}\bigr]\cdot E^0\bigl[e^{-\beta L_{T_a}(b)-\alpha L_{T_a}(0)}\bigr].
\end{equation}
To compute the first factor, let $\omega$ be such that $\omega\in\Omega^*$ and $\omega-b\in\Omega^*$;
by \bk{(3.6.3)} the set of such $\omega$ has $P^b$-probability one. The representation
\bk{(3.6.2)}, valid on $\Omega^*$ (\bk{Remark 3.6.5}), shows that
\[
L_t(b,\omega)=\lim_{\eps\downarrow 0}\frac{1}{4\eps}\meas\{0\le s\le t;\ \abs{\omega(s)-b}\le\eps\}
=L_t(0,\omega-b),\qquad 0\le t<\infty,
\]
and clearly $T_0(\omega)=T_{-b}(\omega-b)$. Since the law of $\omega-b$ under $P^b$ is $P^0$, we get
from \bk{(4.40)}
\[
E^b\bigl[e^{-\beta L_{T_0}(b)}\bigr]=E^0\bigl[e^{-\beta L_{T_{\abs{b}}}(0)}\bigr]
=\frac{1}{1+\beta\abs{b}} .
\]
Substituting this and \eqref{eq:6-4-15:zero} into \eqref{eq:6-4-15:split}, we find
\begin{equation}\label{eq:6-4-15:result}
E^b\bigl[\exp\{-\beta L_{T_a}(b)-\alpha L_{T_a}(0)\}\bigr]
=\frac{1}{(1+\alpha a)(1+\beta\abs{b})+\beta a}.
\end{equation}

\emph{Step 3: identification with $f_2$.} From \bk{(4.22)}, \bk{(4.23)} with $n=1$ we have
$f_1(t_1,\gamma_1)=f_0+t_1\gamma_1f_0=1+\gamma_1t_1$, and with $n=2$,
\[
f_2(t_1,\gamma_1;t_2,\gamma_2)=f_1(t_2-t_1,\gamma_2)+t_1\bigl[\gamma_1f_1(t_2-t_1,\gamma_2)+\gamma_2f_0\bigr]
=(1+\gamma_1t_1)\bigl(1+\gamma_2(t_2-t_1)\bigr)+\gamma_2t_1 .
\]
For $t_1=a$, $\gamma_1=\alpha$, $t_2=a+\abs{b}$, $\gamma_2=\beta$ this is
$f_2(a,\alpha;a+\abs{b},\beta)=(1+\alpha a)(1+\beta\abs{b})+\beta a$, the denominator in
\eqref{eq:6-4-15:result}. This proves (ii).

\rmk{The result agrees with \bk{Lemma 4.9}: translating by $\abs{b}$ as in Step~2, the left-hand side of
\eqref{eq:6-4-15:result} is $E^0[\exp\{-\delta_1L_{T_{b'}}(\xi_1)-\delta_2L_{T_{b'}}(\xi_2)\}]$ with
$b'=a+\abs{b}$, $\xi_1=0$, $\xi_2=\abs{b}$, $\delta_1=\beta$, $\delta_2=\alpha$, and \bk{(4.27)} gives
$1/f_2(b'-\xi_2,\delta_2;\,b'-\xi_1,\delta_1)=1/f_2(a,\alpha;\,a+\abs{b},\beta)$. Moreover, for
$\alpha=0$ formula \eqref{eq:6-4-15:result} reduces to $(1+\beta(a+\abs{b}))^{-1}$, in accordance with
\bk{(4.40)} and the translation invariance used in Step~2.}
\end{kex}

\solsection{6.5}{An Application: Transition Probabilities of Brownian Motion with Two-Valued Drift}

\begin{kex}{Exercise}{5.3}{}
Throughout, $t>0$, and for $y\in\R$ we write
\[
  \varphi_t(y)\defeq\frac{1}{\sqrt{2\pi t}}\,e^{-y^2/2t},\qquad
  \Phi_t(y)\defeq\Phi\Bigl(\frac{y}{\sqrt t}\Bigr)=\int_{-\infty}^{y}\varphi_t(u)\,du,\qquad
  \ol\Phi_t(y)\defeq 1-\Phi_t(y)=\Phi_t(-y),
\]
and we introduce the two functions
\begin{equation}\label{eq:6-5-03:MK}
  M(t,y)\defeq (y^2+t)\,\Phi_t(y)+t\,y\,\varphi_t(y),\qquad
  K(t,y)\defeq t\,\varphi_t(y)-y\,\ol\Phi_t(y)
\end{equation}
on $(0,\infty)\times\R$. We shall use repeatedly the identity, valid for all real $r$,
\begin{equation}\label{eq:6-5-03:shift}
  e^{2r}\varphi_t(r+t)=\varphi_t(r-t),
\end{equation}
which holds because $(r+t)^2-(r-t)^2=4rt$, and the elementary integrals ($c\in\R$)
\begin{equation}\label{eq:6-5-03:int}
  \int_c^\infty\varphi_t(w)\,dw=\ol\Phi_t(c),\qquad
  \int_c^\infty w\,\varphi_t(w)\,dw=t\,\varphi_t(c),\qquad
  \int_c^\infty w^2\varphi_t(w)\,dw=t\,c\,\varphi_t(c)+t\,\ol\Phi_t(c);
\end{equation}
the second follows from $\partial_w\varphi_t(w)=-(w/t)\varphi_t(w)$, and the third from the second by an integration by parts.

\emph{Step 1: the function $g$.} For $t>0$ and \emph{all} $r\in\R$ we set
\begin{equation}\label{eq:6-5-03:g}
  g(t,r)\defeq \tfrac12+M(t,r-t)-\tfrac12\,\Phi_t(r-t)
  -e^{2r}\Bigl[K(t,r+t)+\tfrac12\,\ol\Phi_t(r+t)\Bigr].
\end{equation}
Inserting \eqref{eq:6-5-03:MK} and using \eqref{eq:6-5-03:shift} in the form $e^{2r}\,t\,\varphi_t(r+t)=t\,\varphi_t(r-t)$, we obtain
\begin{equation}\label{eq:6-5-03:gexp}
  g(t,r)=\tfrac12+t\,\varphi_t(r-t)\,(r-t-1)
  +\bigl\{(r-t)^2+t-\tfrac12\bigr\}\Phi_t(r-t)
  +e^{2r}\bigl(r+t-\tfrac12\bigr)\ol\Phi_t(r+t).
\end{equation}
Since $t\,\varphi_t(y)=\sqrt{t/2\pi}\,e^{-y^2/2t}$, the right-hand side of \bk{(5.15)} is exactly $g(t,\abs{x})$. The claim \bk{(5.15)} is thus
\begin{equation}\label{eq:6-5-03:claim}
  v(t,x)=g(t,\abs{x});\qquad t>0,\ x\in\R.
\end{equation}

\emph{Step 2: proof of \eqref{eq:6-5-03:claim}.} By the definition \bk{(5.6)}, $\wt p_t(x,\cdot)$ is the density of $Z_t$ under $\wt P^x$, so
$v(t,x)=\int_{-\infty}^\infty z^2\,\wt p_t(x,z)\,dz$, an integral of a nonnegative function with values in $[0,\infty]$. With $\theta_1=-\theta_0=1$ the symmetry relation \bk{(5.13)} reads $\wt p_t(x,z)=\wt p_t(-x,-z)$, whence $v(t,x)=v(t,-x)$ by the change of variable $z\mapsto -z$. It suffices therefore to prove \eqref{eq:6-5-03:claim} for $x\ge0$.

Let $x\ge 0$ be fixed and put $a\defeq x-t$. For $\theta=1$, $\delta=0$, \bk{Proposition 5.1}(i) and \bk{(5.14)} give, for $r>0$,
\begin{align*}
  \wt p_t(x,r)&=\varphi_t(x-r-t)+e^{-2r}\int_{x+r}^\infty\varphi_t(u-t)\,du
               =\varphi_t(r-a)+e^{-2r}\,\ol\Phi_t(r+a),\\
  \wt p_t(x,-r)&=e^{2x}\varphi_t(x+r+t)+e^{-2r}\int_{x+r}^\infty\varphi_t(u-t)\,du
               =e^{-2r}\varphi_t(r+a)+e^{-2r}\,\ol\Phi_t(r+a),
\end{align*}
where in the second line we have used \eqref{eq:6-5-03:shift} with $r$ replaced by $x+r$: $e^{2x}\varphi_t(x+r+t)=e^{-2r}\varphi_t(x+r-t)$. Consequently
\begin{equation}\label{eq:6-5-03:split}
  v(t,x)=\int_0^\infty r^2\bigl[\wt p_t(x,r)+\wt p_t(x,-r)\bigr]dr
  =\int_0^\infty r^2\varphi_t(r-a)\,dr+\int_0^\infty r^2\,q(r)\,dr,
\end{equation}
with
\[
  q(r)\defeq e^{-2r}\bigl[\varphi_t(r+a)+2\,\ol\Phi_t(r+a)\bigr]
  =-\frac{d}{dr}\Bigl[e^{-2r}\,\ol\Phi_t(r+a)\Bigr].
\]
For the first integral in \eqref{eq:6-5-03:split} we substitute $w=r-a$ and use \eqref{eq:6-5-03:int} with $c=-a$, together with $\ol\Phi_t(-a)=\Phi_t(a)$ and $\varphi_t(-a)=\varphi_t(a)$:
\begin{align}
  \int_0^\infty r^2\varphi_t(r-a)\,dr
  &=\int_{-a}^\infty (w^2+2aw+a^2)\,\varphi_t(w)\,dw \notag\\
  &=-t\,a\,\varphi_t(a)+t\,\Phi_t(a)+2a\,t\,\varphi_t(a)+a^2\Phi_t(a)=M(t,a).
  \label{eq:6-5-03:I1}
\end{align}
For the second, two integrations by parts (all boundary terms at infinity vanish because of the factor $e^{-2r}$, and those at $0$ are displayed) give
\begin{align*}
  \int_0^\infty r^2 q(r)\,dr
  &=\int_0^\infty 2r\,e^{-2r}\,\ol\Phi_t(r+a)\,dr\\
  &=\Bigl[-\bigl(r+\tfrac12\bigr)e^{-2r}\,\ol\Phi_t(r+a)\Bigr]_{r=0}^{\infty}
    -\int_0^\infty\bigl(r+\tfrac12\bigr)e^{-2r}\varphi_t(r+a)\,dr\\
  &=\tfrac12\,\ol\Phi_t(a)-\int_0^\infty\bigl(r+\tfrac12\bigr)e^{-2r}\varphi_t(r+a)\,dr,
\end{align*}
since $\frac{d}{dr}\bigl[-(r+\frac12)e^{-2r}\bigr]=2r\,e^{-2r}$ and $\frac{d}{dr}\ol\Phi_t(r+a)=-\varphi_t(r+a)$. Now $r+a=(x+r)-t$, so \eqref{eq:6-5-03:shift} yields $e^{-2r}\varphi_t(r+a)=e^{2x}\varphi_t(r+x+t)$, and with $b\defeq x+t$, the substitution $u=r+b$ and \eqref{eq:6-5-03:int},
\[
  \int_0^\infty\bigl(r+\tfrac12\bigr)\varphi_t(r+b)\,dr
  =\int_b^\infty\bigl(u-b+\tfrac12\bigr)\varphi_t(u)\,du
  =t\,\varphi_t(b)-b\,\ol\Phi_t(b)+\tfrac12\,\ol\Phi_t(b)
  =K(t,b)+\tfrac12\,\ol\Phi_t(b).
\]
Therefore
\begin{equation}\label{eq:6-5-03:I2}
  \int_0^\infty r^2 q(r)\,dr
  =\tfrac12-\tfrac12\,\Phi_t(x-t)-e^{2x}\Bigl[K(t,x+t)+\tfrac12\,\ol\Phi_t(x+t)\Bigr].
\end{equation}
Adding \eqref{eq:6-5-03:I1} (with $a=x-t$) and \eqref{eq:6-5-03:I2} in \eqref{eq:6-5-03:split}, and comparing with \eqref{eq:6-5-03:g}, we obtain $v(t,x)=g(t,x)$ for $x\ge 0$. This proves \eqref{eq:6-5-03:claim}, i.e.\ \bk{(5.15)}; in particular $v$ is finite.

\emph{Step 3: a differential equation for $g$.} We have $\partial_y\Phi_t=\varphi_t$, $\partial_y\varphi_t=-(y/t)\varphi_t$, and
\[
  \partial_t\Phi_t(y)=-\frac{y}{2t}\,\varphi_t(y),\qquad
  \partial_t\varphi_t(y)=\Bigl(\frac{y^2}{2t^2}-\frac{1}{2t}\Bigr)\varphi_t(y).
\]
Hence $\partial_t\Phi_t=\frac12\partial_{yy}\Phi_t$, and the same holds for $\ol\Phi_t=1-\Phi_t$. For the functions of \eqref{eq:6-5-03:MK} we compute
\begin{align*}
  M_y&=2y\,\Phi_t+(y^2+t)\varphi_t+t\,\varphi_t-y^2\varphi_t=2y\,\Phi_t+2t\,\varphi_t, &
  M_{yy}&=2\Phi_t+2y\,\varphi_t-2y\,\varphi_t=2\Phi_t,\\
  K_y&=-y\,\varphi_t-\ol\Phi_t+y\,\varphi_t=-\ol\Phi_t, &
  K_{yy}&=\varphi_t,
\end{align*}
and
\begin{align*}
  M_t&=\Phi_t-(y^2+t)\frac{y}{2t}\,\varphi_t+y\,\varphi_t
       +t\,y\Bigl(\frac{y^2}{2t^2}-\frac{1}{2t}\Bigr)\varphi_t=\Phi_t=\tfrac12 M_{yy},\\
  K_t&=\varphi_t+t\Bigl(\frac{y^2}{2t^2}-\frac{1}{2t}\Bigr)\varphi_t-y\cdot\frac{y}{2t}\,\varphi_t
      =\tfrac12\,\varphi_t=\tfrac12 K_{yy}.
\end{align*}
Thus the functions
$f\defeq M-\frac12\Phi_t$ and $h\defeq K+\frac12\ol\Phi_t$ are of class $C^\infty$ on $(0,\infty)\times\R$ and satisfy the heat equation $f_t=\frac12f_{yy}$, $h_t=\frac12h_{yy}$. By \eqref{eq:6-5-03:g}, $g(t,r)=\frac12+f(t,r-t)-e^{2r}h(t,r+t)$, so $g$ is $C^\infty$ on $(0,\infty)\times\R$, and with $f$ and its derivatives evaluated at $(t,r-t)$, $h$ and its derivatives at $(t,r+t)$:
\begin{align*}
  g_t&=f_t-f_y-e^{2r}(h_t+h_y), &
  g_r&=f_y-e^{2r}(2h+h_y), &
  g_{rr}&=f_{yy}-e^{2r}(4h+4h_y+h_{yy}),
\end{align*}
whence
\begin{equation}\label{eq:6-5-03:pde}
  \tfrac12 g_{rr}-g_r=\tfrac12 f_{yy}-f_y-e^{2r}\bigl(\tfrac12h_{yy}+h_y\bigr)=g_t
  \qquad\text{on }(0,\infty)\times\R.
\end{equation}
Moreover, $f_y=2y\,\Phi_t+2t\,\varphi_t-\frac12\varphi_t$ and $2h+h_y=2t\,\varphi_t-2y\,\ol\Phi_t+\ol\Phi_t-\ol\Phi_t-\frac12\varphi_t$, so that, by \eqref{eq:6-5-03:shift} once more, the terms containing $\varphi_t$ cancel in $g_r$ and
\begin{equation}\label{eq:6-5-03:gr}
  g_r(t,r)=2(r-t)\,\Phi_t(r-t)+2e^{2r}(r+t)\,\ol\Phi_t(r+t).
\end{equation}
In particular,
\begin{equation}\label{eq:6-5-03:neumann}
  g_r(t,0)=-2t\,\Phi_t(-t)+2t\,\ol\Phi_t(t)=0;\qquad t>0.
\end{equation}

\emph{Step 4: the equation \bk{(5.16)}.} By \eqref{eq:6-5-03:claim}, $v(t,x)=g(t,\abs{x})$ with $g$ smooth on $(0,\infty)\times\R$. Hence $v$ and $v_t(t,x)=g_t(t,\abs{x})$ are continuous on $(0,\infty)\times\R$, and for $x\ne0$,
\[
  v_x(t,x)=(\sgn x)\,g_r(t,\abs{x}),\qquad v_{xx}(t,x)=g_{rr}(t,\abs{x}).
\]
At $x=0$ we have, by \eqref{eq:6-5-03:neumann} and the differentiability of $g$ and $g_r$ in $r$ at $r=0$,
\begin{align*}
  \frac{v(t,k)-v(t,0)}{k}&=(\sgn k)\,\frac{g(t,\abs{k})-g(t,0)}{\abs{k}}\xrightarrow[k\to0]{}0,\\
  \frac{v_x(t,k)-0}{k}&=\frac{g_r(t,\abs{k})-g_r(t,0)}{\abs{k}}\xrightarrow[k\to0]{}g_{rr}(t,0),
\end{align*}
so $v_x(t,0)=0$ and $v_{xx}(t,0)=g_{rr}(t,0)$. The formulas $v_x(t,x)=(\sgn x)g_r(t,\abs{x})$ (with any value of $\sgn0$) and $v_{xx}(t,x)=g_{rr}(t,\abs{x})$ therefore hold for all $x\in\R$, and both functions are continuous on $(0,\infty)\times\R$ (for $v_x$, because $g_r(t,0)=0$). Thus $v$ is of class $C^{1,2}$ on $(0,\infty)\times\R$, and for every $(t,x)$, by \eqref{eq:6-5-03:pde} evaluated at $(t,\abs{x})$,
\begin{align*}
  \tfrac12v_{xx}(t,x)-(\sgn x)\,v_x(t,x)
  &=\tfrac12g_{rr}(t,\abs{x})-(\sgn x)^2g_r(t,\abs{x})\\
  &=\tfrac12g_{rr}(t,\abs{x})-g_r(t,\abs{x})=g_t(t,\abs{x})=v_t(t,x);
\end{align*}
for $x\neq0$ we have used $(\sgn x)^2=1$, and for $x=0$ the fact that $g_r(t,0)=0$, which makes the value assigned to $\sgn 0$ irrelevant. This is \bk{(5.16)}, in the classical sense, at every point of $(0,\infty)\times\R$.

\emph{Step 5: the limits.} Fix $x\in\R$, put $r=\abs{x}\ge0$ and use \eqref{eq:6-5-03:gexp}. We need only the bound $1-\Phi(z)\le e^{-z^2/2}$ for $z\ge0$, which holds since $u^2\ge z^2+(u-z)^2$ for $u\ge z\ge 0$ gives $1-\Phi(z)\le e^{-z^2/2}\,(2\pi)^{-1/2}\int_z^\infty e^{-(u-z)^2/2}du=\frac12e^{-z^2/2}$.

\textbf{(i)} $t\to\infty$. Let $t>r$. From $(r\pm t)^2/2t\ge \mp r+t/2$ we get
\begin{gather*}
  t\,\varphi_t(r-t)\le\sqrt{t}\,e^{r-t/2},\qquad
  \Phi_t(r-t)=\ol\Phi_t(t-r)\le e^{-(t-r)^2/2t}\le e^{r-t/2},\\
  \ol\Phi_t(r+t)\le e^{-(r+t)^2/2t}\le e^{-r-t/2}.
\end{gather*}
Each of the last three terms in \eqref{eq:6-5-03:gexp} is thus bounded in absolute value by a polynomial in $t$ and $\sqrt t$ times $e^{r-t/2}$, and tends to zero. Hence $\lim_{t\to\infty}v(t,x)=\frac12$.

\textbf{(ii)} $t\downarrow0$. The second term in \eqref{eq:6-5-03:gexp} is bounded in absolute value by $\sqrt{t/2\pi}\,(r+t+1)$ and tends to zero. If $r>0$, then $(r-t)/\sqrt t\to\infty$ and $(r+t)/\sqrt t\to\infty$, so $\Phi_t(r-t)\to1$, $\ol\Phi_t(r+t)\to0$ and
\[
  \lim_{t\downarrow0}v(t,x)=\tfrac12+0+\bigl(r^2-\tfrac12\bigr)\cdot1+e^{2r}\bigl(r-\tfrac12\bigr)\cdot0=r^2=x^2.
\]
If $r=0$, then $\Phi_t(-t)=\Phi(-\sqrt t)\to\frac12$ and $\ol\Phi_t(t)=1-\Phi(\sqrt t)\to\frac12$, so
\[
  \lim_{t\downarrow0}v(t,0)=\tfrac12+0+\bigl(-\tfrac12\bigr)\cdot\tfrac12+\bigl(-\tfrac12\bigr)\cdot\tfrac12=0=x^2.
\]
This completes the proof.

\rmk{Two consistency checks. First, by \bk{(5.14)} and the computation at the beginning of Step 2, the density of $\abs{Z_t}$ under $\wt P^x$, $x\ge0$, is $\varphi_t(r-a)+q(r)$ on $(0,\infty)$, and it has total mass $\Phi_t(a)+\ol\Phi_t(a)=1$, as it must. Secondly, the limit $\frac12$ in \textbf{(i)} is the second moment $\int_{-\infty}^\infty z^2e^{-2\abs{z}}dz$ of the probability density $e^{-2\abs{z}}$, and this density is the pointwise limit of $\wt p_t(x,\cdot)$ on $\R\setminus\{0\}$ as $t\to\infty$: for $x\ge0$ and $r>0$, in the two expressions for $\wt p_t(x,\pm r)$ in Step 2 (where $a=x-t$) the terms $\varphi_t(r-a)$ and $e^{-2r}\varphi_t(r+a)$ are at most $(2\pi t)^{-1/2}$ and tend to zero, whereas $\ol\Phi_t(r+a)=1-\Phi\bigl((x+r-t)/\sqrt t\bigr)$ tends to $1$ because $(x+r-t)/\sqrt t\to-\infty$; hence $\wt p_t(x,\pm r)\to e^{-2r}$, and the case $x<0$ follows from $\wt p_t(x,z)=\wt p_t(-x,-z)$. The relation \eqref{eq:6-5-03:neumann}, $v_x(t,0)=0$, is what makes $v$ a classical solution of \bk{(5.16)} across the switching point $x=0$, and \eqref{eq:6-5-03:gr} is the explicit expression of $v_x(t,x)=(\sgn x)\,g_r(t,\abs{x})$ needed in \bk{Exercise 5.4}.}
\end{kex}

\begin{kex}{Exercise}{5.4}{}
Throughout, $\theta_1=-\theta_0=1$ and $\delta=0$, so that $\mathscr{U}$ is the set of Borel-measurable $b\colon[0,\infty)\times\R\to[-1,1]$ and the function of \bk{(5.4)} is $u(x)=1$ for $x\le 0$, $u(x)=-1$ for $x>0$. We write $\varphi=\Phi'$ for the standard normal density and use \bk{Exercise 5.3}: the function $v(t,x)\defeq\wt{E}^xZ_t^2$ ($t>0$, $x\in\R$), where $Z$ solves \bk{(5.5)}, is given by \bk{(5.15)}. Thus $v(t,x)=g(t,\abs{x})$, where for $t>0$, $r\in\R$,
\begin{gather*}
  g(t,r)\defeq\frac12+\sqrt{t}\,(r-t-1)\,\varphi(a)+\Bigl\{(r-t)^2+t-\frac12\Bigr\}\Phi(a)
     +e^{2r}\Bigl(r+t-\frac12\Bigr)\bigl[1-\Phi(c)\bigr],\\
  a=a(t,r)\defeq\frac{r-t}{\sqrt t},\qquad c=c(t,r)\defeq\frac{r+t}{\sqrt t}.
\end{gather*}
The function $g$ is infinitely differentiable on $(0,\infty)\times\R$.

\emph{Step 1. The derivatives of $g$.} We have $c^2-a^2=4r$, hence
\begin{equation}\label{eq:6-5-04:a}
  e^{2r}\varphi(c)=\varphi(a),
\end{equation}
and $\varphi'(y)=-y\varphi(y)$, $a_r=c_r=t^{-1/2}$, $a_t=-c/(2t)$, $c_t=-a/(2t)$, $a\sqrt t=r-t$, $c\sqrt t=r+t$. Differentiating in $r$ and using \eqref{eq:6-5-04:a},
\begin{align*}
  g_r&=\frac{\varphi(a)}{\sqrt t}\Bigl[t-(r-t-1)(r-t)+(r-t)^2+t-\frac12-\Bigl(r+t-\frac12\Bigr)\Bigr]\\
     &\qquad+2(r-t)\Phi(a)+e^{2r}\,(2r+2t)\bigl[1-\Phi(c)\bigr],
\end{align*}
and the bracket vanishes identically. Therefore
\begin{align}
  g_r&=2(r-t)\Phi(a)+2e^{2r}(r+t)\bigl[1-\Phi(c)\bigr],\label{eq:6-5-04:b}\\
  g_{rr}&=2\Phi(a)+\frac{2(r-t)}{\sqrt t}\varphi(a)+2e^{2r}(2r+2t+1)\bigl[1-\Phi(c)\bigr]-\frac{2(r+t)}{\sqrt t}\,e^{2r}\varphi(c)\notag\\
        &=2\Phi(a)-4\sqrt t\,\varphi(a)+2e^{2r}(2r+2t+1)\bigl[1-\Phi(c)\bigr].\label{eq:6-5-04:c}
\end{align}
Differentiating in $t$ and using \eqref{eq:6-5-04:a} again, the coefficient of $\varphi(a)$ in $g_t$ is
\begin{align*}
  &\frac{r-t-1}{2\sqrt t}-\sqrt t+\frac{(r-t-1)\,ac}{2\sqrt t}-\Bigl(a^2t+t-\frac12\Bigr)\frac{c}{2t}+\Bigl(c\sqrt t-\frac12\Bigr)\frac{a}{2t}\\
  &\quad=\Bigl(\frac a2-\frac1{2\sqrt t}\Bigr)-\sqrt t+\Bigl(\frac{a^2c}{2}-\frac{ac}{2\sqrt t}\Bigr)
     +\Bigl(-\frac{a^2c}{2}-\frac c2+\frac{c}{4t}\Bigr)+\Bigl(\frac{ac}{2\sqrt t}-\frac{a}{4t}\Bigr)\\
  &\quad=\frac{a-c}{2}-\sqrt t+\frac{c-a}{4t}-\frac{1}{2\sqrt t}=-2\sqrt t,
\end{align*}
because $c-a=2\sqrt t$. Consequently
\begin{equation}\label{eq:6-5-04:d}
  g_t=-2\sqrt t\,\varphi(a)+\bigl(1-2(r-t)\bigr)\Phi(a)+e^{2r}\bigl[1-\Phi(c)\bigr]=\tfrac12 g_{rr}-g_r
  \qquad\text{on }(0,\infty)\times\R,
\end{equation}
the second equality by comparison with \eqref{eq:6-5-04:b}, \eqref{eq:6-5-04:c}. Moreover, since $a(t,0)=-\sqrt t$, $c(t,0)=\sqrt t$,
\begin{equation}\label{eq:6-5-04:e}
  g_r(t,0)=-2t\,\Phi(-\sqrt t)+2t\bigl[1-\Phi(\sqrt t)\bigr]=0;\qquad t>0.
\end{equation}

\emph{Step 2. $g_r(t,r)>0$ for $t>0$, $r>0$, and a bound.} Let $M(y)\defeq(1-\Phi(y))/\varphi(y)$, $y\in\R$, and $h(y)\defeq yM(y)$. Then $M'(y)=yM(y)-1$ and $h'(y)=(1+y^2)M(y)-y$. For $y>0$ we have $h'(y)>0$: the function $k(y)\defeq1-\Phi(y)-y\varphi(y)/(1+y^2)$ satisfies
\[
  k'(y)=-\varphi(y)-\varphi(y)\,\frac{1-2y^2-y^4}{(1+y^2)^2}=-\frac{2\varphi(y)}{(1+y^2)^2}<0,\qquad \lim_{y\to\infty}k(y)=0,
\]
so $k>0$ on $\R$, which is $h'>0$ on $(0,\infty)$. Also, for $y>0$, $1-\Phi(y)=\int_y^\infty\varphi(z)\,dz<\int_y^\infty (z/y)\varphi(z)\,dz=\varphi(y)/y$, i.e.
\begin{equation}\label{eq:6-5-04:f}
  0<h(y)<1;\qquad y>0,
\end{equation}
whereas $h(y)\le 0$ for $y\le 0$. Now $\Phi(a)=1-\Phi(-a)=M(-a)\varphi(a)$ and, by \eqref{eq:6-5-04:a}, $e^{2r}[1-\Phi(c)]=M(c)\varphi(a)$; thus \eqref{eq:6-5-04:b} reads
\begin{equation}\label{eq:6-5-04:g}
  g_r=2\sqrt t\,\bigl\{a\Phi(a)+c\,e^{2r}[1-\Phi(c)]\bigr\}=2\sqrt t\,\varphi(a)\bigl\{h(c)-h(-a)\bigr\}.
\end{equation}
Let $r>0$. Then $c>0$ and $c-(-a)=2r/\sqrt t>0$. If $-a\le 0$, then $h(-a)\le 0<h(c)$; if $-a>0$, then $h(-a)<h(c)$ because $h$ is strictly increasing on $(0,\infty)$. In both cases \eqref{eq:6-5-04:g} gives $g_r(t,r)>0$. Finally, for $r\ge 0$ and $0<t\le T$, \eqref{eq:6-5-04:b}, the identity $2e^{2r}(r+t)[1-\Phi(c)]=2\sqrt t\,\varphi(a)h(c)$, \eqref{eq:6-5-04:f} and $\varphi\le(2\pi)^{-1/2}$ yield
\begin{equation}\label{eq:6-5-04:h}
  \abs{g_r(t,r)}\le 2(r+t)+2\sqrt t\,\varphi(a)\le 2r+2T+\sqrt T.
\end{equation}

\emph{Step 3. $v$ is of class $C^{1,2}$ on $(0,\infty)\times\R$ and satisfies \bk{(5.17)}.} Since $v(t,x)=g(t,\abs x)$, the derivative $v_t(t,x)=g_t(t,\abs x)$ exists and is continuous on $(0,\infty)\times\R$. For $x\neq0$ we have $v_x(t,x)=\sgn(x)\,g_r(t,\abs x)$ and $v_{xx}(t,x)=g_{rr}(t,\abs x)$. At $x=0$, by \eqref{eq:6-5-04:e},
\begin{align*}
  \frac{v(t,y)-v(t,0)}{y}&=\frac{\abs y}{y}\cdot\frac{g(t,\abs y)-g(t,0)}{\abs y}\xrightarrow[y\to0]{}0,\\
  \frac{v_x(t,y)-0}{y}&=\frac{g_r(t,\abs y)-g_r(t,0)}{\abs y}\xrightarrow[y\to0]{}g_{rr}(t,0),
\end{align*}
so that $v_x(t,0)=0$ and $v_{xx}(t,0)=g_{rr}(t,0)$. Hence $v_{xx}(t,x)=g_{rr}(t,\abs x)$ everywhere, which is continuous, and $v_x$ is continuous on $(0,\infty)\times\R$: this is clear off the line $\{x=0\}$, and at a point $(t_0,0)$ it follows from $\abs{v_x(t,x)}=\abs{g_r(t,\abs x)}\to g_r(t_0,0)=0$ as $(t,x)\to(t_0,0)$. Thus $v\in C^{1,2}((0,\infty)\times\R)$, and by Step 2 and \eqref{eq:6-5-04:e},
\begin{equation}\label{eq:6-5-04:i}
  x\,v_x(t,x)=\abs x\,g_r(t,\abs x)\ge0,\qquad\text{so}\qquad \abs{v_x(t,x)}=g_r(t,\abs x);\qquad t>0,\ x\in\R.
\end{equation}
For every real $p$ we have $\min_{\abs u\le1}(up)=-\abs p$. Therefore, by \eqref{eq:6-5-04:i} and \eqref{eq:6-5-04:d} (with $r=\abs x\ge0$),
\[
  \frac12v_{xx}(t,x)+\min_{\abs u\le1}\bigl(u\,v_x(t,x)\bigr)=\frac12g_{rr}(t,\abs x)-g_r(t,\abs x)=g_t(t,\abs x)=v_t(t,x)
\]
for all $(t,x)\in(0,\infty)\times\R$, which is \bk{(5.17)}. (Equation \eqref{eq:6-5-04:d} is \bk{(5.16)} of \bk{Exercise 5.3}; what \bk{(5.17)} adds to it is the sign information \eqref{eq:6-5-04:i}, by which the minimum in \bk{(5.17)} is attained at $u=u(x)$.)

\emph{Step 4. Behaviour of $v$ for small $t$.} Let $(Z,W)$, $(\Omega,\F,\wt P^y)$, $\{\F_t\}$ be a weak solution of \bk{(5.5)} with $Z_0=y$. Since $\abs u\le1$, we have $\abs{Z_\eps-y}\le\eps+\abs{W_\eps}$, whence $\abs{Z_\eps^2-y^2}\le(\eps+\abs{W_\eps})(2\abs y+\eps+\abs{W_\eps})$ and, because $\wt E^y\abs{W_\eps}\le\sqrt\eps$ and $\wt E^yW_\eps^2=\eps$,
\begin{equation}\label{eq:6-5-04:j}
  \abs{v(\eps,y)-y^2}\le\wt E^y\abs{Z_\eps^2-y^2}\le2\abs y(\eps+\sqrt\eps)+2\eps^2+2\eps;\qquad\eps>0,\ y\in\R.
\end{equation}

\emph{Step 5. Proof of \bk{(5.18)}.} For $t=0$ both sides of \bk{(5.18)} equal $x^2$. Fix $t>0$, $x\in\R$, $b\in\mathscr U$ and a weak solution $(X,W)$, $(\Omega,\F,\wt P^x)$, $\{\F_t\}$ of \bk{(5.1)}; thus $W$ is a standard Brownian motion with respect to $\{\F_t\}$ under $\wt P^x$, and $X$ is continuous, adapted, and
\begin{equation}\label{eq:6-5-04:k}
  \abs{X_s}\le\abs x+s+\abs{W_s}\le\abs x+t+W^*,\qquad 0\le s\le t,\qquad W^*\defeq\max_{0\le s\le t}\abs{W_s},
\end{equation}
where $\wt E^x(W^*)^2\le4t$ by Doob's maximal inequality (\bk{Theorem 1.3.8}(iv), applied to the continuous, nonnegative submartingale $\abs{W}$ with $p=2$).

Let $0<\eps<t$. In order to apply It\^o's rule we choose a continuously differentiable function $\rho\colon[0,\infty)\to[\eps/2,t]$ with $\rho(s)=t-s$ for $0\le s\le t-\eps$ (for instance $\rho(s)=\frac\eps2+\frac\eps2\exp\{-2(s-t+\eps)/\eps\}$ for $s>t-\eps$) and set $f(s,y)\defeq v(\rho(s),y)$. By Step 3, $f$ is of class $C^{1,2}$ on $[0,\infty)\times\R$. \bk{Theorem 3.3.6}, applied to $f$ and the continuous semimartingale $X$ with local martingale part $W$ and bounded-variation part $\int_0^\cdot b(r,X_r)\,dr$, gives, $\wt P^x$-a.s., for $0\le s\le t-\eps$,
\begin{equation}\label{eq:6-5-04:l}
  v(t-s,X_s)=v(t,x)+\int_0^s\Bigl[-v_t+b(r,X_r)\,v_x+\frac12v_{xx}\Bigr](t-r,X_r)\,dr+\int_0^sv_x(t-r,X_r)\,dW_r.
\end{equation}
Since $\abs{b}\le1$, \bk{(5.17)} shows that the integrand of the Lebesgue integral satisfies
\[
  \Bigl[-v_t+b(r,X_r)\,v_x+\frac12v_{xx}\Bigr](t-r,X_r)\ \ge\ \Bigl[-v_t+\frac12v_{xx}+\min_{\abs u\le1}(uv_x)\Bigr](t-r,X_r)=0 .
\]
The process $H_r\defeq v_x(t-r,X_r)\ind_{[0,t-\eps]}(r)$ is progressively measurable (it is adapted, and continuous on $[0,t-\eps]$) and, by \eqref{eq:6-5-04:i}, \eqref{eq:6-5-04:h} with $T=t$, and \eqref{eq:6-5-04:k}, $\abs{H_r}\le2\abs{X_r}+2t+\sqrt t\le2\abs x+4t+\sqrt t+2W^*$; hence $\wt E^x\int_0^tH_r^2\,dr<\infty$, the stochastic integral $\int_0^\cdot H_r\,dW_r$ is a square-integrable martingale, and its expectation at time $t-\eps$ is zero. The left-hand side of \eqref{eq:6-5-04:l} at $s=t-\eps$ is integrable by \eqref{eq:6-5-04:j} and \eqref{eq:6-5-04:k}. Taking expectations in \eqref{eq:6-5-04:l} at $s=t-\eps$ we obtain
\begin{equation}\label{eq:6-5-04:m}
  \wt E^x\,v(\eps,X_{t-\eps})\ \ge\ v(t,x);\qquad0<\eps<t.
\end{equation}
Now we let $\eps\downarrow0$. By \eqref{eq:6-5-04:j} and \eqref{eq:6-5-04:k},
\[
  \wt E^x\bigl|v(\eps,X_{t-\eps})-X_{t-\eps}^2\bigr|\le2(\eps+\sqrt\eps)\bigl(\abs x+t+\wt E^xW^*\bigr)+2\eps^2+2\eps\xrightarrow[\eps\downarrow0]{}0,
\]
while $X_{t-\eps}^2\to X_t^2$ by continuity of the paths, with $X_{t-\eps}^2\le(\abs x+t+W^*)^2$, an integrable bound; by dominated convergence $\wt E^xX_{t-\eps}^2\to\wt E^xX_t^2$. Hence \eqref{eq:6-5-04:m} yields
\[
  \wt E^xX_t^2\ \ge\ v(t,x)=\wt E^xZ_t^2 ,
\]
the last equality being the definition of $v$; here $\wt E^xZ_t^2$ depends only on the law of $Z$, which is uniquely determined by \bk{(5.5)} (uniqueness in the sense of probability law, recalled after \bk{(5.1)}; note that $u\in\mathscr U$, so \bk{(5.5)} is the case $b(t,y)=u(y)$ of \bk{(5.1)}). This is \bk{(5.18)}.

\emph{Step 6. Optimality for \bk{(5.3)}.} Let $\alpha>0$ and $b\in\mathscr U$, with $X$ as in Step 5. The process $X$ is continuous and adapted, hence jointly measurable in $(t,\omega)$, and the integrand $e^{-\alpha t}X_t^2$ is nonnegative; by Tonelli's theorem and \bk{(5.18)},
\[
  J(x;b)=\wt E^x\int_0^\infty e^{-\alpha t}X_t^2\,dt=\int_0^\infty e^{-\alpha t}\,\wt E^xX_t^2\,dt
  \ \ge\ \int_0^\infty e^{-\alpha t}\,\wt E^xZ_t^2\,dt=J(x;b_*)
\]
for every $x\in\R$, where $b_*(t,y)\defeq u(y)$. Since $b_*\in\mathscr U$, the minimum in \bk{(5.3)} is attained at $b_*$, and the corresponding solution $Z$ of \bk{(5.5)} is the optimally controlled process; this holds for every discount rate $\alpha>0$, and indeed $Z$ minimizes the second moment at each fixed time.

\rmk{The formulas \eqref{eq:6-5-04:b}--\eqref{eq:6-5-04:d}, the vanishing \eqref{eq:6-5-04:e} and the positivity of $g_r$ were confirmed against finite differences of \bk{(5.15)} on a grid, and \bk{(5.18)} by simulation for two competing drifts (\texttt{code/c6\_5\_04.py}).}
\end{kex}

\end{document}
